% DOKUMENTENKLASSE
\documentclass[a4paper,11pt]{report}

% PAPIERFORMAT 
\usepackage[a4paper]{geometry}			% verkleinert die margins
\usepackage{fullpage}				% verkleinert den größten Teil der margins um etwa 1,5cm
	\frenchspacing				% kein zusätzlicher Abstand nach Punkt

% SPRACHE
\usepackage[utf8]{inputenc}				% für spezielle Zeichen, Umlaute etc.
\usepackage[T1]{fontenc}
\usepackage[english]{babel}				% [ngerman] setzt die Spracheinstellung auf Neue dt. Rechtschreibung

% SILBENTRENNUNG
 \tolerance=1000 
 \hyphenpenalty=1000

% LINKDARSTELLUNG

\usepackage{hyphenat}	
\usepackage{hyperref, bookmark}

% ZEICHEN
\usepackage{pifont}						% spezielle (Pfeile etc.)
\usepackage{textcomp}						% Gradzeichen
\usepackage{tipa}						% IPA-Darstellung

% FARBEN 
\usepackage[usenames,dvipsnames]{color}	% für die 68 Standardfarben

% SPALTEN
\usepackage{multicol}						% Spalten im Dokument mit \begin{multicols}{2, 3, 4 … 10}
% \usepackage{draftwatermark}				% kreiert Wasserzeichen 'Draft' (kann mit colorpackage kollidieren, ansonst nach geometry einfügen)
% \SetWatermarkScale{5}						% vergrößert Schriftart
% \SetWatermarkLightness{0.9}				% verändert Helligkeit des Textes 0=schwarz, 1=weiß
%\SetWatermarkText{\textsc{Rohfassung}}	% verändert Standarttext

% EINZÜGE
\setlength{\parindent}{0cm}				% lässt Paragraphen NICHT (0cm) einrücken

% DEFINIERTE BEFEHLE						kreiert neue (verkürzte) Befehle
\newcommand{\hi}{\hangindent=0.3cm}		
\newcommand{\B}[1]{\textbf{#1}}
\newcommand{\I}[1]{\textipa{#1}}
\newcommand{\on}[1]{\oldstylenums{#1}}
\newcommand{\LINK}[2]{{\color{blue}\href{#1}{\textsc{#2}}}}
\newcommand{\cp}[1]{{\color{Plum}{#1}}}
\newcommand{\cb}[1]{{\color{blue}{#1}}}
\newcommand{\HT}[2]{\hypertarget{#1}{#2}}
\newcommand{\HL}[2]{\hyperlink{#1}{#2}}
\renewcommand{\thefootnote}{\fnsymbol{footnote}}

% DOKUMENT

\begin{document}
\title{\B{An Annotated Na'vi Dictionary}\footnote{This work is licensed under the Creative Commons Attribution-ShareAlike License, version 3.0. To view a copy of this license, visit http://creativecommons.org/licenses/by-sa/3.0/.} \\
\emph{Aysìkenonghu a Lì'upuk leNa'vi}}  		
\author{Compiled by Stefan G. M\"uller (Plumps)}

\date{\today}

\maketitle
% ===== Blindbefehle ======
% \I{[]}
% \cp{}
% \B{}
% \emph{}
% $\cdot$
% \LINK{}{}
% \HT{}{}						akzeptiert kein ì (= I), ä (= A)
% \HL{}{\B{
%  \ding{96} botany, flora 
%  \ding{214} opposite of 
%  \ding{43} see, compare with entry
%  \ding{246} military, combat, archery
% \\  \on{}. \B{}, .
% (\emph{a. subjective})
% (\emph{b. agentive})
% (\emph{c. patientive})
% (\emph{d. genitive})
% (\emph{e. dative})
% (\emph{f. topic})

% ======= ToC ??? ==========

% Test

% \include{38_testbearbeitung}
 
%  \end{document}

% ======= Introduction ========

% \section*{Introduction}

% \emph{Faylì'u a plltxe san Zola'u nìprrte' sìk zìya'u ye'rìn …}

Since its first inception in December of \on{2009}, the \LINK{https://files.learnnavi.org/dicts/NaviDictionary.pdf}{Taronyu} dictionary, compiled by Taronyu (\emph{alu} Richard Littauer), has been---and still remains---a vital part of the Na'vi community. Through its many iterations it is now taken care of by Markì (\emph{alu} Mark Miller) and still provides the first look at the main body of core vocabulary for the Na'vi language.

My activities as a \emph{karyu}, `teacher' in the community soon showed me that beginners would often take a word from the dictionary and apply it with the little grammatical knowledge that they had to form their first sentences. As helpful as this can be without examples to live by how a word is specifically used in Na'vi this could only end in confusion. 

Venturing into other languages and dictionaries I soon came to the realisation that I was dissatisfied with the information that the dictionary provided. In my opinion it could do so much more. My dictionaries of e.g. English, French, Swedish, and Irish all had at least a few example sentences to show the user how a word was in use, whether there were special cases, proverbs or idioms that this word was a part of. Especially with a conlang that was still in development and where there are no native speakers, I had the strong feeling that in order to provide good Na'vi one should lead with good examples and provide the necessary sources to examplify.

As late as \on{2011} I had also begun to collect all official example sentences by \emph{Karyu Pawl}, Dr. Paul Frommer who by then had started the \LINK{http://naviteri.org/}{Na'viteri} blog. These resources gave me a new spark to attempt another kind of dictionary.

% =====================

% Speaking a language is not only knowing the vocabulary and be familiar with the grammar. It also entails knowing which words go together. Not only proverbs and idioms but also certain parts of speech only go with certain verbs.

% It's not always clear how these go together and a look in a common dictionary will only show you the word. If you stack up the words you will get a hotch\hyp{}potch of Na'vi that more resembles English (or your respective native language) than cohesive Na'vi. This is especially true with prepositions. Anybody who has ventured into learning a new language will know that one word doesn't always translate in your language.

% The best way to learn is\slash know is by example. 

% Over the years since \emph{Avatar} has been released and the Learn Na'vi Community first formed there have been many attempts at a Na'vi dictionary--from \emph{Taronyu} to \emph{The \on{500} Most Commonly Used Na'vi Words}. Most of them only concentrate on the words rather than their interaction with each other. And sometimes the meaning of a word is obscured by the fact that the translation itself can have different meanings.

% This work tries to fill that gap. To lead (and teach) by example. 

% It has been attempted to find an official example of every given word in use and give it in the appropriate entry. `Official' is defined as coming directly from the source, i.e. Dr. Frommer himself (most often nowadays via his blog naviteri.org). This Dictionary also tries to be a point of reference. Therefore a link to every example is given. It may not always be the first mention of a word but sometimes subsequent information can help and be of value to interested parties.





% ======= How-to ========

% \include{02_howto}

% ===== section IPA =====

 \pdfbookmark[1]{On Pronunciation}{IPA}

\subsection*{On Pronunciation}

There are different dialects spoken on Pandora (usually referred to by clan or area where it is spoken). It is needless to say that one is not better than the other or that there is anything like a standard or preferred dialect. 

Since the first contact between humans and Na'vi presumably happened with the Omatikaya Clan, the study of the language also started with the Omatikayan Dialect, which is usually referred to as `Forest Na'vi' (\emph{Lì'fya Na'rìngä}). Its pronunciation and vocabulary is the best known variety of Na'vi (so far) and is therefore predominantly concentrated on in this dictionary. Every pronunciation refers to Forest Na'vi first, possible variations are noted after the IPA or directly in the lexical entry.

\subsection*{The Sounds of Na'vi -- IPA}

The Na'vi alphabet has \on{33} distinct sounds that are written with the following letters:
\begin{center}
\B{'  a  aw  ay  ä  e  ew  ey  f  h  i  ì  k  kx  l  ll m \\ 
n  ng  o  p  px  r  rr  s  t  tx  ts  u  v  w  y  z}
\end{center}

The letters b, c, d, j, and q don't exist, g and x only appear in combination---they don't have a sound value of their own.

The pronunciaton of Na'vi is quite regular, although for an English speaker, it has some unusual sounds and sound combinations that can best be learned by listening to sound samples.

\subsubsection*{Vowels}

\begin{tabular}{ c l c }
    &  & IPA \\
  \B{a} & as in `father' & \I{a} \\
  \B{ä} & as in `cat' & \I{\ae} \\
  \B{e} & as in `when' & \I{E} \\
  \B{i} & as the long /i/ in `feat' or `beat' & \I{i} \\
  \B{ì} & as the short /i/ in `fit' or `bit' & \I{I} \\
  \B{o} & as in `so', `or' and `fort' & \I{o} \\
  \B{u} & as in `food' and `boot' (in open and closed syllables) & \I{u} \\ 
    & or `foot' and `book' (in closed syllables) & \I{U} \\
\end{tabular} \\

Note: The difference between open and closed syllables is whether the syllable ends in a vowel (open), e.g. \B{lu} \I{[lu]} `be' or a consonant (closed), e.g. \B{tsun} \I{[\t{ts}un]} or \I{\t{[ts}Un]} `can, be able to'.

\subsubsection*{Diphthongs}

\begin{tabular}{ r l l l r l l}
    &  & IPA  \\
  \B{aw} & as in `house' & \I{aw} \\
  \B{ay} & as in `my' & \I{aj} \\
  \B{ew} & no exact English equivalent; close to the /eow/ in `Beowulf' & \I{Ew} \\
  \B{ey} & as in `hey' or `say' & \I{Ej} \\
\end{tabular}

\subsubsection*{Pseudovowels}

\begin{tabular}{ c l c }
    &  & IPA \\
  \B{ll} & light l as in `fill' or `people' & \I{\s{l}} \\
  \B{rr} & very strongly trilled at the hard gums behind the front teeth & \I{\s{r}} \\
\end{tabular} \\

Note: In contrast to the vowels, \B{ll} and \B{rr} can never start a syllable. They must be preceded by a consonant, e.g. \B{'rr}.pxom, sä.\B{pll}.txe or consonant cluster, e.g. nì.\B{fkrr}.

\subsubsection*{Consonants}

\begin{tabular}{ r l l }
    &  & IPA \\
  \B{'} & = `glottal stop', the sharp break in `uh-oh' & \I{P} \\
  \B{f} & as in `fill' or `cliff' & \I{f} \\
  \B{h} & never silent, always pronounced & \I{h} \\
  \B{k} & unaspirated as the k in `skin' at the start of a syllable & \I{k} \\
             & unreleased at the end of a syllable & \I{k\textcorner} \\
  \B{kx} & pronounced k while holding breath & \I{k'} \\
  \B{l} & as in `life' or `level' & \I{l} \\
  \B{m} & as in `man' or `from' & \I{m} \\
  \B{n} & as in `now' or `on' & \I{n} \\
  \B{ng} & one sound as in `singer', no audible g sound & \I{N} \\
  \B{p} & unaspirated as the p in `sport' at the start of a syllable & \I{p} \\
             & unreleased at the end of a syllable & \I{p\textcorner} \\
  \B{px} & pronounced p while holding breath & \I{p'} \\ 
  \B{r} & fronted, flapped r, as in Spanish `torero' & \I{R} \\
  \B{s} & sharp s sound as in `soft' or `floss' & \I{s} \\
  \B{t} & unaspirated as the t in `stop' at the start of a syllable & \I{t} \\
             & unreleased at the end of a syllable & \I{t\textcorner} \\
  \B{tx} & pronounced t while holding breath & \I{t'} \\
  \B{ts} & as the ts in `tsunami' or `cats' & \I{\t{ts}} \\
  \B{v} & as in `very' or `have' & \I{v} \\
  \B{w} & as in `what' or `will' & \I{w} \\
  \B{y} & as in `yet' & \I{j} \\ 
  \B{z} & as in `zoom' or `Oz' & \I{z} \\
\end{tabular} \\

\subsubsection*{Stress}

Word stress in Na'vi is unpredictable. It must be learned. Once learned it has the advantage that the stress usually remains on the same syllable no matter whether affixes are added. With mono\hyp{}syllabic words the primary stress is clear and has not been marked in the IPA. The stress in poly\hyp{}syllabic words is indicated by color marking in the main entry, e.g. \B{\cp{ta}ron}. For those comfortable with the IPA the stress mark ( \B{\I{"}} ) is used, e.g. \I{["ta.Ron]}.

\subsubsection*{Sound Change -- Lenition}

Lenition is a phonological change from one sound into another. Prefixes or adpositions that cause lenition are marked with a \texttt{+}. In Na'vi the following consonants undergo lenition: 
\begin{center}
\begin{tabular}{ l c l || l c l || l c l || l c l }
   \B{'} & \ding{234} & (\emph{disappears}) & \B{kx} & \ding{234} & \B{k} & \B{px} & \ding{234} & \B{p} & \B{tx} & \ding{234} &  \B{t} \\ 
   \B{k} & \ding{234} & \B{h} & \B{p} & \ding{234} & \B{f} & \B{t} & \ding{234} & \B{s} & \B{ts} & \ding{234} & \B{s} \\
\end{tabular} 
\end{center}

This is important to know since sometimes a word in use could be lenited and therefore is not found under the subsequent entry but rather under its unlenited rootform. Be aware of that.

\subsection*{A Note on Dialects}

Although Forest Na'vi is chosen for this dictionary as a point from which to start, it is by no means preferred or superior to any other dialect spoken on Pandora. The known dialects so far are mutually comprehensible, which means that it should not be difficult to switch between one dialect and another.

\subsubsection*{Of Forest Na'vi (Lì'fya Na'rìngä)}

This dialect is spoken by the Omatikaya Clan. So far, it is the best studied and well\hyp{}known variety of the Na'vi language and is therefore chosen to represent Na'vi in this dictionary as a point of reference.  

\subsubsection*{Of Reef Na'vi (Lì'fya Wionä)} 

This dialect is spoken by the Reef Clans like the Metkayina. It differs from Forest Na'vi in the following aspects:

\begin{itemize}

	\item The consonant cluster sy- \I{[sj]} is pronounced \I{[S]}, like the sh in `shore' or `ash'. The words \B{syawm}, \emph{know}, \B{syaw}, \emph{call}, or \B{syeha}, \emph{breath}, are pronounced \I{[Sawm]}, \I{[Saw]}, and \I{["SE.ha]} respectively in Reef Na'vi.
	\item The consonant cluster tsy- \I{[\t{ts}j]} is pronounced \I{[\t{tS}]}, like the ch in `church'. Thus the words \B{tsyal}, \emph{wing}, or \B{tsyeytsyìp}, \emph{tiny bite}, are pronounced \I{[\t{tS}al]} and \I{["\t{tS}Ej.\t{tS}Ip\textcorner]} respectively in Reef Na'vi.
	\item The glottal stop is lost beween non\hyp{}identical vowels, e.g. \B{fra'u}, \emph{everything}, becomes \B{frau}; \B{Lo'ak}, \emph{the name Lo'ak}, becomes \B{Loak}.
	\item The glottal stop between identical vowels \emph{may be} left out, e.g. \B{rä'ä}, \emph{do not}, may become \B{rää} \I{[R\ae."{}\ae]}; \B{me'em}, \emph{harmonious}, may become \B{meem} \I{[mE."{}Em]}.
		\subitem Given that speakers of Reef Na'vi are used to hearing these sequences of identical vowels, forms like \B{apxaa} \I{[a."ba.a]}, \B{meeveng} \I{[mE."{}E.vEN]}, \B{seii} \I{[sE.i."{}i]}, and \B{peekxinùm} \I{[pE.E."gi.nUm]} remain as they are. 
	\item At the beginning of a syllable, the ejectives \B{px} \I{[p']}, \B{tx} \I{[t']}, and \B{kx} \I{[k']} are pronounced \I{[b]}, \I{[d]}, and \I{[g]} respectively in this dialect, e.g. \B{txon} \I{[don]}, \emph{night}; \B{holpxay} \I{[hol."baj]}, \emph{number}, or \B{kxitx} \I{[git']}, \emph{death}.
		\subitem If the word ends in an ejective, the syllable boundaries will be reanalysed when a suffix is added, e.g. \B{'awkx} \I{[Pawk']}, but \B{'awkxìl} \I{["Paw.gIl]}
		\subitem When two ejectives across syllable boundaries get in contact within a word, they undergo regressive assimilation, e.g. \B{atxkxe} and \B{ekxtxu} are \I{[ad."gE]} and \I{[Eg."du]} rather than *\I{[at'."gE]} and *\I{[Ek'."du]}.
	\item Reef Na'vi distinguishes between a tense u sound \I{[u]} as in `food', and a lax u sound \I{[U]} as in `good'. Analogous to i and ì in Forest Na'vi, the sounds are marked u and ù, e.g. \B{pum} \I{[pum]} \emph{or} \I{[pUm]} (\textsc{rn:} \B{pùm} \I{[pUm]}) \emph{``dummy noun''}. This can also lead to minimal pairs in RN, e.g. \B{tsun} \I{[\t{ts}un]}\textsuperscript{\textsc{rn}} \emph{heel} and \B{tsùn} \I{[\t{ts}Un]}\textsuperscript{\textsc{rn}} \emph{can, be able}. These differences in sound are only marked in the respective entries for words with the lax ù sound. Assume that all other u sounds are tense in Reef Na'vi.
	\item In unstressed syllables \B{ä} \I{[\ae]} tends to change to \B{e} \I{[E]} in Reef Na'vi, e.g. \B{nge\-yä} \I{["NE.j\ae]} and \B{tä\-txaw} \I{[t\ae."t'aw]} may be pronounced \I{["NE.jE]} and \I{[tE."daw]}, respectively.
	\item The patientive ending \B{-ti} is favored over its other variants and the topical case \B{-(ì)ri} may also appear sentence final. 

\end{itemize}

% \B{} \I{[]} \emph{}



% ===== section word classes =====

 \pdfbookmark[1]{Of Nouns, Verbs, and Adjectives}{Of Nouns, Verbs, and Adjectives}
\subsection*{Of Nouns}

\begin{itemize}

	\item It is assumed that the reader knows about word formation and how to form plurals in Na'vi. Thus, not all plural forms are given. However, with words where the prefix causes elision plurals are given, e.g. \B{'en} (\emph{pl.} \B{men}, \B{pxen}, \B{ayen\slash en}) or where there are irregular plurals, e.g. \B{'ll\-ngo} (\emph{pl.} \B{me\-'ll\-ngo}, \B{pxe\-'ll\-ngo}, \emph{always:} \B{ay\-'ll\-ngo}) or: \B{'u} (\emph{no short plural, always:} \B{ay\-u}). Where these elisions cause pairs of homonyms, a subsequent entry will refer to the main entry, e.g. \B{pxen}\textsuperscript{\on{2}}~:~\B{'en}.

	\item This project aims to become an exhaustive reference dictionary. Thus, proper names of people, places, flora and fauna and loan words have been incorporated into the main body of the dictionary.

\end{itemize}

\subsection*{Of Verbs}

\begin{itemize}

	\item It is assumed that the reader knows about word formation and the basic concept of how Na'vi creates tense, aspect and affect in verbs by use of infixes. Thus, these infix positions in regular verbs (mono\hyp{}syllabic and di\hyp{}syllabic verbs) as well as in \B{si}\hyp{}verbs have \emph{not been marked}. However, in irregular verbs, i.e. compound verbs with elements of other word classes or two verbs, the infix positions have been indicated with the system devised by Dr.~Frommer, i.e. \B{sre\-se\-'a}, \emph{inf.}\on{23}, which gives \B{sres}‹pre-1›‹1›\B{e'}‹2›\B{a}. Additionally they are given in the familiar way from the \emph{Taronyu} Dictionary via the raised dot in the IPA (\I{[sRE.s$\cdot$E."P$\cdot$a]}).

	\item The appearance of a verb can change quite drastically through the use of infixes. To familiarize the user\slash reader with that change, the example sentences always try to cover as many different forms as possible, starting from the core verb in use, to the ‹\B{iv}› form to more complicated constructs if they could be found.
	
\end{itemize}

\subsection*{Of Adjectives}

Adjectives undergo the least changes. In their basic function they only appear in their root form given in the dictionary as the predicative adjective (\B{tute lu 'ewan}, ``The person is young'') or with an attached \B{a} as the attributive adjective (\B{tute a'ewan lu tstunwi} \emph{or} \B{'ewana tute lu tstunwi}, ``The young person is kind''). It has been attempted to find and show at least one example of each of these uses.

\pdfbookmark[1]{A Note on Spelling}{A Note on Spelling and Hyphenation}
\subsection*{A Note on Spelling and Hyphenation}

Careful thought has been given to the fact that Na'vi is a spoken language and that rules of syllabification are handled differently than in English or other languages. Prefixes, infixes and suffixes can change the appearance of a word quite drastically. A vowel can make up a whole syllable. In general, the sequence VCV (where V stands for a vowel, pseudo vowel or diphthong and C for a consonant) will be syllabified V.CV rather than VC.V. However, in instances of case endings this rule has been slightly and carefully altered to isolate the suffix as well as the root word. The reasoning being that the root word in \B{tsafnetxonìri} (``concerning that kind of night'') would be more recognizable when given as \B{tsa-fne-txon-ìri} rather than \B{tsa-fne-txo-nì-ri}, as it would be spoken. For the same reason, some words---especially those with long vowel clusters---have been prevented from a line break to leave them recognizable, e.g. \B{tsaioiyä} (``of that adornment'') as \B{tsa-ioi-yä} rather than \B{tsa-i-o-i-yä}. 

In the same vein, \B{fayhilvanit} (``those rivers'' (\emph{patientive})) is syllabified as \B{fay-hil-van-it} rather than \B{fay-hil-va-nit} to retain the root noun \B{kilvan} (river) as visible as possible.

Agentive, patientive or dative single endings have not been treated that way to avoid illegal word boundaries, e.g. \B{ayyayol} (``the birds'' (\emph{agentive})) is still syllabified as \B{ay-ya-yol} rather than *\B{ay-ya-yo-l}.

This practice was not followed, however, with verbs affected by infixes. Infixes rearrange the core syllabification of a verb. Ignoring this would lead to illegal consonant clusters. Whereas \B{terkup} (``die'') is syllabified as \B{ter-kup}, adding infixes rearranges it---in the most extreme case---to e.g. \B{täpeykìyeverkeiup} (``cause oneself to die''), syllabified as \B{tä-pey-kì-ye-ver-ke-i-up}. 

Recognizing the root verb is a matter of practice. In extreme cases the translation should give a hint as to the core verb.


 							% input: fortlaufender Text beim Einfügen

% ===== section abbreviations =====

 \newpage
\pdfbookmark[0]{Abbreviations}{Abbreviations}
\section*{Abbreviations}

\begin{multicols}{2}
\begin{footnotesize}
\begin{tabular}{ l c l }
  \emph{adj.} & = & adjective \\
  \emph{adp.} & = & adposition \\
  \emph{adp.}\texttt{+} & = & adposition that causes lenition \\
  \emph{adv.} & = & adverb \\
  \emph{astr.} & = & astronomical \\
  \emph{coll.} & = & colloquial \\
  \emph{conj.} & = &  conjunction\\
  \emph{dem.} & = & demonstrative \\
  \emph{fig.} & = & figurative \\
  \textsc{fn} & = & \emph{Forest Na'vi} dialect specifically \\
  \emph{inf.} & = &  infix positions in irregular verbs\\
  \emph{inter.} & = & interrogative, question word \\
  \emph{intj.} & = & interjection \\
  \textsc{lw} & = & loan word (borrowed mostly from English) \\
  \emph{n.} & = & noun \\
  \emph{nfp.} & = & not for people \\
  \emph{num.} & = & numeral, number \\
  \emph{ofp.} & = & only for people \\
  \emph{ord. num.} & = & ordinal number \\
  \emph{part.} & = &  particle \\
  \emph{pl.} & = & plural (forms), i.e. dual, trial, and plural \\
  \emph{pn.} & = & pronoun \\
  \textsc{rn} & = & \emph{Reef Na'vi} dialect specifically \\
  \emph{slang} & = & slang \\
  s.o. & = & someone, somebody \\
  s.t. & = & something \\
  \emph{text.} & = & textile \\
  \emph{vin.} & = & intransitive verb \\
  \emph{vtr.} & = &  transitive verb\\
  \emph{vulg.} & = & vulgar, impolite, rude \\ [0.3cm]
References \\
  \textsc{\cb{a}}\slash \textsc{\cb{a\on{2}}} & = & line appears in the movie \emph{Avatar} (\on{2009})\slash \emph{Avatar: The Way of Water} (\on{2022}) \\
  \textsc{\cb{a\textsuperscript{c}}} & = & a line that was cut from the movie \\
  \textsc{\cb{b}} & = & link to \href{http://naviteri.org/}{\texttt{Na'viteri}} by Dr Frommer \\
  \textsc{\cb{dh}} & = & word or line appeared in \emph{Dark Horse Comics} (\on{2019}) \\
  \textsc{\cb{g}} & = & a line from the \emph{Avatar Game} (\on{2009}) \\
  \textsc{\cb{fop}} & = & a line from the Game \emph{Avatar: Frontiers of Pandora} (\on{2023}) \\
  \textsc{\cb{ld}} & = & link to \href{}{\texttt{Lemondrop}} \\
  \textsc{\cb{ll}} & = & link to \href{}{\texttt{The Language Log}} \\
  \textsc{\cb{ln}} & = & link to the \href{http://forum.learnnavi.org/}{\texttt{LearnNa'vi forums}} \\
  \textsc{\cb{pp}} & = & link to \href{}{\texttt{Pandorapedia}} \\
  \textsc{\cb{pr}} & = & link to the \href{}{\texttt{Public Reading Project}} \\
  \textsc{\cb{tn}} & = & link to \href{}{\texttt{Why is this night…?}} \\
  \textsc{\cb{wiki}} & = & link to the \href{}{\texttt{LearnNa'vi Wiki}} \\ 
  \textsc{\cb{yt}} & = & link to a \href{}{\texttt{YouTube video}} \\ [0.3cm]
Signs  \\
  \ding{96} & = & botany, flora \\
  \ding{214} & = & opposite of \dots \\
  \ding{43} & = & see \dots, compare with entry \dots  \\
  \ding{246} & = & military, combat, archery \\
  \texttt{+} & = & a prefix that causes lenition. \\
\end{tabular}
\end{footnotesize}
\end{multicols} 						% include: forciert einen Seitenumbruch beim Einfügen

% ===== section apostrophe ======

 \newpage
\pdfbookmark[0]{Apostrophe}{'}
\begin{center}
\section*{' \I{[P]} (\emph{tìftang}, Apostrophe, Glottal Stop)}
\end{center}

\begin{multicols}{2}

% 'a'aw
\hi 
\HT{'a'aw}{\B{\cp{'a}'aw}} \I{["Pa.Paw]} \emph{adj.} (with singular of the noun; only with countable nouns) a few, several. 
\B{Lu po\-ru 'ey\-lan a\-'a\-'aw.} He has several friends.~\LINK{http://naviteri.org/2010/07/vocabulary-update/}{b} 
\B{Oel tse\-'a 'a\-'aw\-a tu\-tet.} I see several people.~\LINK{http://naviteri.org/2010/07/vocabulary-update/}{b} 
(\ding{43}~\HL{hol}{\B{hol}})

% 'ak 
\hi
\HT{'ak}{\B{'ak}} \I{[Pak\textcorner]} \emph{intj.} ow, ouch.~\LINK{https://forum.learnnavi.org/language-updates/correspondence-with-karyu-pawl/}{ln}

% 'akra
\hi 
\HT{'akra}{\B{\cp{'ak}ra}} \I{["Pak\textcorner.Ra]} \emph{n.} (fertile) soil (in which plants can grow). 
\B{'ak\-ra a\-pay\-nga'}, moist soil; 
\B{'ak\-ra le\-pay}, watery, saturated soil.~\LINK{http://naviteri.org/2011/05/weather-part-2-and-a-bit-more-2/}{b}

%'aku
\hi 
\HT{'aku}{\B{\cp{'a}ku}} \I{["Pa.ku]} \emph{vtr.} (\emph{a.~of clothing}) take off, remove.  
\B{Ru\-txe me\-hawn\-ven\-it 'i\-va\-ku.} Please take off your shoes.~\LINK{http://naviteri.org/2011/08/new-vocabulary-clothing/}{b} 
(\emph{b.~fig.~with} \B{ta}) remove, take away.  \B{Pot 'a\-ku fì\-tseng\-ta!} Get him out of here!~\LINK{http://naviteri.org/2011/08/new-vocabulary-clothing/}{b} 
(\emph{c.~vulg.}) \B{'ä\-pi\-va\-ku}, get lost!~\LINK{http://forum.learnnavi.org/naviteri/new-vocabulary-clothing-19454/msg482699/\#msg482699}{ln} (\ding{214}~\HL{yemstokx}{\B{yem\-stokx}}, \HL{ioi si}{\B{\mbox{ioi} si}})

% 'al
\hi 
\HT{'al}{\B{'al}} \I{[Pal]} \emph{vtr.} waste.
\B{Rä\-'ä 'i\-val syu\-vet!} Don't waste food.~\LINK{http://naviteri.org/2014/07/tsawlultxamaw-a-ayliu-post-meetup-words/}{b} \\
\textsc{Derivations}: \ding{43} 
\on{1}.~\HL{tI'al}{\B{tì\-'al}}, wastefulness. 
\on{2}.~\HL{le'al}{\B{le\-'al}}, wasteful (not for people). 
\on{3}. \HL{nI'al}{\B{nì\-'al}}, wastefully.

%'ali'ä
\hi
\HT{'ali'A}{\B{'a\cp{li}'ä}} \I{[Pa."li.P\ae]} \emph{n.} choker, collar. 
\B{'Ali\-'ät vul\-sìn key\-kur za\-'u fì\-tseng.} Hang the choker on the branch and come here.~\LINK{http://naviteri.org/2013/04/wheres-the-bathroom-and-other-useful-things/}{b}

% 'ampi
\hi
\HT{'ampi}{\B{\cp{'am}pi}} \I{["Pam.pi]} \emph{vtr.} (\texttt{+}\emph{control}) touch.   
\B{Oe\-ti 'am\-pi rä\-'ä, ma \mbox{skxawng}!} Don't touch me, you moron!~\LINK{http://naviteri.org/2012/11/renu-ayinanfyaya-the-senses-paradigm/}{b}  
\B{Ze\-ne fko 'ivam\-pi prrnenit nì\-flrr fra\-krr.} One must always touch a baby gently.~\LINK{http://naviteri.org/2013/02/vospxi-ayol-posti-apup-short-post-for-a-short-month/}{b}
(\emph{a. idiom}) \B{(Na) lo\-re\-yu 'aw\-nam\-pi}, extreme shyness, [lit.: (like) a touched helicoradian] \B{Lu por mok\-ri a\-mik\-lor, slä lo\-re\-yu 'aw\-nam\-pi lu. Ke tsun ri\-vol eo su\-te.} She has a beautiful voice, but she's extremely shy. She can't sing in front of people.~\LINK{http://naviteri.org/2014/11/twenty-before-the-holidays/}{b} \\
\textsc{Derivations}: \ding{43} 
\on{1}. \HL{'ampirikx}{\B{'am\-pi\-rikx}}, leaf pitcher plant.

% 'ampirikx
\hi 
\HT{'ampirikx}{\B{\cp{'am}pirikx}} \I{["Pam.pi.Rik']} \emph{n.}, \ding{96} leaf pitcher plant, \emph{pseudocenia simplex}. 
(\ding{43}~\HL{'ampi}{\B{'am\-pi}}, \HL{rikx}{\B{rikx}})

% 'ana
\hi
\HT{'ana}{\B{\cp{'a}na}} \I{["Pa.na]} \emph{n.}, \ding{96} hanging vine.
(\emph{a. proverb}) \B{Spä skxawng sìn 'a\-na a\-flì.} ``A fool jumps onto a thin vine,'' i.e. \emph{don't engage in an unpromising and\slash or potentially risky cause}.~\LINK{http://naviteri.org/2021/08/ulte-ayyoratu-leiu-and-the-winners-are-2/}{b}
\B{Tsa\-ye\-rik te\-rul ne 'awkx. Spä skxawng sìn 'a\-na a\-flì.} That hexapede is running toward the cliff. Only a fool jumps onto a thin vine.~\LINK{http://naviteri.org/2021/08/ulte-ayyoratu-leiu-and-the-winners-are-2/}{b} \\
\textsc{Derivations}: \ding{43} 
\on{1}.~\HL{syep'an}{\B{syep\-'an}}, lift vine.

% 'anla
\hi
\HT{'anla}{\B{\cp{'an}la}} \I{["Pan.la]} \emph{vtr.} yearn for. 
\B{Tsa\-ye\-rik\-tsyìp\-ìl li 'an\-la sa'\-nok\-it.} That little hexapede is already yearning for its mother.~\LINK{http://naviteri.org/2013/01/awvea-posti-zisita-amip-first-post-of-the-new-year/}{b} \\  
\textsc{Derivations}: \ding{43} 
\on{1}.~\HL{sA'anla}{\B{sä\-'an\-la}}, yearning.

% 'angtsìk
\hi 
\HT{'angtsIk}{\B{\cp{'ang}tsìk}} \I{["PaN.\t{ts}Ik\textcorner]} \emph{n., fauna} hammerhead, \emph{rhinoquadruculus hammercephali}. 
\B{'Ang\-tsìk a kxap si lu le\-hrr\-ap, ma 'itan.} A threatening hammerhead is dangerous, son.~\LINK{http://naviteri.org/2015/03/kap-si-ayunil-saylahe-threats-dreams-and-other-things/}{b}
\B{Weyn\-flit\-it 'ang\-tsìk\-ìl sro\-lu' tspang.} Wainfleet was crushed and killed by a hammerhead.~\LINK{http://naviteri.org/2012/07/meetings-waterfalls-and-more/}{b}
\B{Äo fì\-ut\-ral lu tsmìm 'ang\-tsìk\-ä.} There's a hammerhead track under this tree.~\LINK{http://naviteri.org/2013/02/vospxi-ayol-posti-apup-short-post-for-a-short-month/}{b}

% 'ango
\hi
\HT{'ango}{\B{\cp{'a}ngo}} \I{["Pa.No]} \emph{adj.} (\emph{of sound}) soft. 
\B{Zawr ye\-rik\-ä lu 'a\-ngo.} The call of the hexapede is quiet.~\LINK{http://naviteri.org/2014/11/twenty-before-the-holidays/}{b}
(\ding{214}~\HL{wok}{\B{wok}})

% 'are
\hi
\HT{'are}{\B{\cp{'a}re}} \I{["Pa.RE]} \emph{n.} poncho, cape, shawl. 
\B{Mo\-'at \mbox{alu} Tsa\-hìk lu O\-ma\-ti\-ka\-ya\-ä le\-'awa ha\-pxì\-tu a \mbox{ioi} sä\-pi fa 'a\-re.} Mo'at, the Tsahik, is the only member of the Omatikaya who wears a poncho.~\LINK{http://naviteri.org/2011/08/new-vocabulary-clothing/}{b}

% 'asap / si
\hi 
\HT{'asap}{\B{\cp{'a}sap}} \I{["Pa.sap\textcorner]} \textsc{i}. \emph{n.} sudden shock. 
\B{Fwa tse\-'a pe\-yä tì\-fkey\-tok\-it le\-fkrr lo\-lä\-ngu oer 'a\-sap nì\-ngay.} It was a real shock to me to see him in his current condition. \\
\textsc{ii}. \B{\cp{'a}sap si} \I{["Pa.sap\textcorner{} si]} \emph{vin.} be shocked, be startled. 
\B{Oe 'a\-sap so\-li krr\-a tsa\-fmawn\-it stawm.} I was startled when I heard the news.~\LINK{http://naviteri.org/2019/06/50a-liu-amip-40-new-words/}{b} 

% 'avun
\hi 
\HT{'avun}{\B{\cp{'a}vun}} \I{["Pa.vun]} \emph{or} \I{[.vUn]} (\textsc{rn}: \B{\cp{'a}vùn} \I{["Pa.vUn]}) \emph{vtr.} save (\emph{time, food, etc.}). 
\B{Nga\-ri txo fì\-kem si\-vi fì\-fya, krr\-ti 'a\-vun.} If you do it like this, you'll save time.~\LINK{http://naviteri.org/2021/12/zolau-niprrte-ma-3746-welcome-2022/}{b} 
(\ding{43}~\HL{tarep}{\B{tarep}}; \HL{zong}{\B{zong}})

% 'aw
\hi
\HT{'aw}{\B{'aw}} \I{[Paw]} \emph{num., adj.} one. (\emph{base}: \ding{43}~\B{'aw-}, \emph{rest}: \ding{43}~\HL{-aw}{\B{\mbox{-aw}}})   
\B{'aw\-a tì\-pawm}, one question. \B{Ke tsun tu\-te a\-'aw tsa\-tskxe\-ti a\-ku'\-up kxi\-vel\-tek nì\-'aw\-tu.} One person alone can't lift that heavy rock.~\LINK{http://naviteri.org/2012/01/mipa-zisit-ayliu-amip-new-words-for-the-new-year/}{b} \\ 
\textsc{Derivations}: \ding{43} 
\on{1}.~\HL{'aw\-ve}{\B{'aw\-ve}}, first.   
\on{2}.~\HL{le'aw}{\B{le\-'aw}}, only. 
\on{3}.~\HL{'awlie}{\B{'aw\-lie}}, once (in the past). 
\on{4}.~\HL{'awlo}{\B{'aw\-lo}}, once, one time. 
\on{5}.~\HL{'awpo}{\B{'aw\-po}}, one (individual). 
\on{6}.~\HL{'awsiteng}{\B{'aw\-si\-teng}}, together. 
\on{7}.~\HL{'awstengyem}{\B{'aw\-steng\-yem}}, join (s.t. together).
\on{8}.~\HL{vospxI'aw}{\B{Vo\-spxì\-'aw}}, January.

% 'awkx
\hi
\HT{'awkx}{\B{'awkx}} \I{[Pawk']} \emph{n.} cliff. 
\B{Ik\-ran ya\-wo\-lo ftu txew 'awkx\-ä.} The banshee took to the air from the edge of a cliff.~\LINK{http://naviteri.org/2012/03/spring-vocabulary-part-1/}{b} 
\B{Ke tsun oe fì\-'awkx\-it tsyi\-vìl.} I can't scale this cliff.~\LINK{http://naviteri.org/2013/01/lifyari-po-peyi-how-good-is-her-navi/}{b} 

% 'awlie
\hi
\HT{'awlie}{\B{'aw\cp{li}e}} \I{[Paw."li.E]} \emph{adv.} (\emph{a.}) once (in the past) (\emph{marks an experience of something}). 
\B{Oel yo\-lom 'aw\-li\-e tey\-lut.} I once ate \emph{teylu}.~\LINK{http://naviteri.org/2013/01/awvea-posti-zisita-amip-first-post-of-the-new-year/}{b} 
(\emph{b. in questions}) ever. \B{Sra\-ke ko\-lä nga 'aw\-li\-e ne Nu Yor\-kì?} Have you ever been to New York?~\LINK{http://naviteri.org/2013/01/awvea-posti-zisita-amip-first-post-of-the-new-year/}{b}
(\ding{43}~\HL{'aw}{\B{'aw}}, \HL{lie}{\B{lie}})

% 'awlo
\hi
\HT{'awlo}{\B{\cp{'aw}lo}} \I{["Paw.lo]} \emph{adv.} once, one time (\emph{marks a single instance of something}).  
\B{Oel yo\-lom tey\-lut 'aw\-lo nì\-'aw.} I ate \emph{teylu} only once.~\LINK{http://naviteri.org/2013/01/awvea-posti-zisita-amip-first-post-of-the-new-year/}{b} 
(\ding{43}~\HL{'aw}{\B{'aw}}, \HL{alo}{\B{alo}}) \\
\textsc{Derivations}: \ding{43}
\on{1}.~\HL{kawlo}{\B{kaw\-lo}}, not once.

% 'awm
\hi
\HT{'awm}{\B{'awm}} \I{[Pawm]} \emph{n.} camp.
\B{Sra\-ne, 'awm\-it oel ta\-syok.} Yes, I intend to be at the camp.~\LINK{http://naviteri.org/2013/04/zun-zel-counterfactual-conditionals/\#comment-2363}{b} \\
\textsc{Derivations}: \ding{43} 
\on{1}. \HL{tI'awm}{\B{tì\-'awm}}, camping.

% 'awnìm
\hi
\HT{'awnIm}{\B{\cp{'aw}nìm}} \I{["P$\cdot$aw.nIm]} \emph{vtr., inf.}\on{11} avoid. \B{Ngal 'e\-raw\-nìm oe\-ti srak?} Are you avoiding me?~\LINK{http://naviteri.org/2011/09/miscellaneous-vocabulary/}{b}
(\ding{43}~\HL{'I'awn}{\B{'ì\-'awn}}, \HL{alIm}{\B{a\-lìm}})

% 'awpo
\hi
\HT{'awpo}{\B{\cp{'aw}po}} \I{["Paw.po]} \emph{pron.} one (individual\slash person). \B{Ha ftxey 'aw\-pot a ay\-nga\-kip.} So choose one among you.~\LINK{http://naviteri.org/2013/08/taronway-the-hunt-song/}{b} (\ding{43}~\HL{'aw}{\B{'aw}}, \HL{po}{\B{po}}, \HL{fko}{\B{fko}})

% 'awsiteng
\hi
\HT{'awsiteng}{\B{'awsi\cp{teng}}} \I{[Paw.si."tEN]} \emph{adv.} together.   
\B{Ze\-ne oe 'aw\-si\-teng tì\-kang\-kem si\-vi fo\-hu.} I must work together with them.~\LINK{http://forum.learnnavi.org/news-announcements/a-response-from-paul-frommer!/}{ln} 
\B{Ko\-ren a\-'aw\-ve tì\-ru\-sey\-ä 'aw\-si\-teng}, The First Rule for Living Together.~\LINK{http://wiki.learnnavi.org/index.php?title=Corpus\#Good_Morning_America}{wiki} 
(\emph{a. idiom}) \B{'Aw\-si\-teng lu me\-fo la\-nay'\-ka.} \emph{praising how well two people work together and complement each other} [lit.: they are a slinger (together)].~\LINK{http://naviteri.org/2021/08/ulte-ayyoratu-leiu-and-the-winners-are-2/}{b}
\B{To\-la\-ron me\-fol me\-sa\-li\-o\-a\-ngit! Tew\-ti, 'aw\-si\-teng lu me\-fo la\-nay'\-ka.} They hunted two sturmbeest? Wow, they work very well together.~\LINK{http://naviteri.org/2021/08/ulte-ayyoratu-leiu-and-the-winners-are-2/}{b}
(\ding{43}~\HL{'aw}{\B{'aw}}, \HL{teng}{\B{teng}}) \\
\textsc{Derivations}: \ding{43} 
\on{1}. \HL{'awstengyem}{\B{'aw\-steng\-yem}}, join.

% 'awstengyem
\hi
\HT{'awstengyem}{\B{'awsteng\cp{yem}}} \I{[Paw.stEN."j$\cdot$Em]}~\emph{vtr., inf.}\on{33} join (two or more things together).   
\B{Tseng\-a 'aw\-steng\-yä\-pem fì\-me\-kem\-yo lu mek\-tseng.} Where these two walls come together there's a gap.~\LINK{http://naviteri.org/2013/04/wheres-the-bathroom-and-other-useful-things/}{b} 
\B{'aw\-steng\-yä\-po\-le\-i\-em}, (we) have happily united ourselves.~\LINK{http://forum.learnnavi.org/2012-meetup/vote-on-your-favorite-t-shirt-design/msg547581/\#msg547581}{ln} 
(\ding{43}~\HL{'awsiteng}{\B{'aw\-si\-teng}}, \HL{yem}{\B{yem}})

% 'awve
\hi
\HT{'awve}{\B{\cp{'aw}ve}} \I{["Paw.vE]} \emph{ord.~num.} first. \B{Ko\-ren a\-'aw\-ve tì\-ru\-sey\-ä 'aw\-si\-teng}, The First Rule for Living Together.~\LINK{http://wiki.learnnavi.org/index.php?title=Corpus\#Good_Morning_America}{wiki} \B{'aw\-ve\-a po\-stì zì\-sìt\-ä a\-mip}, first post of the New Year.~\LINK{http://naviteri.org/2013/01/awvea-posti-zisita-amip-first-post-of-the-new-year/}{b} 
(\ding{43}~\HL{'aw}{\B{'aw}}, \HL{-ve}{\B{-ve}})

% 'ä'
\hi
\HT{'A'}{\B{'ä'}} \I{[P\ae{}P]} \emph{intj.} whoops! oops! 
\B{'Ä'! \mbox{Oel} tsngal\-it tì\-mung\-zup.} Whoops! I just dropped the cup.~\LINK{http://naviteri.org/2011/09/miscellaneous-vocabulary/}{b}  

% 'ä'o
\hi 
\HT{'A'o}{\B{\cp{'ä}'o}} \I{["P\ae.Po]} \emph{n.}, \ding{96} pitcher plant, \emph{plantapurpurea cordata}.

% 'älek
\hi 
\HT{'Alek}{\B{\cp{'ä}lek}} \I{["P\ae.lEk\textcorner]} \emph{adj.} determined. 
\B{Tì\-flä\-ri lo\-lu po 'ä\-lek.} She was determined to succeed.~\LINK{http://naviteri.org/2021/09/aawa-ayliu-amip-a-few-new-words/}{b} \\
\textsc{Derivations}: \ding{43} 
\on{1}.~\HL{tI'Alek}{\B{tì\-'ä\-lek}}, determination. 
\on{2}.~\HL{nI'Alek}{\B{nì\-'ä\-lek}}, determinedly, with determination.

% 'änsyem
\hi
\HT{'Ansyem}{\B{'än\cp{syem}}} \I{[P\ae{}n."{}sjEm]} \emph{adj.} complete. 

% 'ärìp
\hi
\HT{'ArIp}{\B{\cp{'ä}rìp}} \I{["P\ae.RIp\textcorner]} \emph{vtr.} move (s.t.). \\  
\textsc{Derivations}: \ding{43} 
\on{1}.~\HL{kA'ArIp}{\B{kä\-'ä\-rìp}}, push.  
\on{2}.~\HL{za'ArIp}{\B{za\-'ä\-rìp}}, pull.

% 'e'al
\hi
\HT{'e'al}{\B{\cp{'e}'al}} \I{["PE.Pal]} \emph{adj.} worst. 
\B{Ru\-txe, ay\-nga\-ri fì\-kem fe\-yä 'e\-'a\-la to\-pur ka\-ngay sì\-yi nì\-'aw.} Please, this will only confirm their worst fears about you.~\LINK{http://forum.learnnavi.org/language-updates/two-words-in-use-eal-and-kangay-si/msg564284/\#msg564284}{ln}  
(\ding{214}~\HL{swey}{\B{swey}})

% 'e'in
\hi
\HT{'e'in}{\B{'e\cp{'in}}} \I{[PE."Pin]} \emph{n.} pod, gourd. \\
\textsc{Derivations}: \ding{43} 
\on{1}. \HL{'e'insey}{\B{'e\-'in\-sey}}, drinking gourd.

% 'e'insey
\hi
\HT{'e'insey}{\B{'e\cp{'in}sey}} \I{[PE."Pin.sEj]} \emph{n.} drinking gourd.
(\ding{43}~\HL{'e'in}{\B{'e\-'in}}, \HL{sey}{\B{sey}})

% 'efu
\hi
\HT{'efu}{\B{\cp{'e}fu}} \I{["PE.fu]} (\emph{a. vtr.}) feel, sense, perceive. \B{Lì'\-fya\-ri le\-Na'\-vi \mbox{oel} 'e\-fu ay\-nge\-yä tì\-yawn\-it.} I feel your love for the Na'vi language.~\LINK{http://forum.learnnavi.org/news-announcements/a-response-from-paul-frommer!/}{ln} 
(\emph{b. vin.}) \emph{with predicate adj. for inner feelings} 
\B{'efu} \ding{43} \HL{nitram}{\B{nit\-ram}}, \HL{keftxo}{\B{ke\-ftxo}}, \HL{som}{\B{som}}, \HL{wew}{\B{wew}}, \HL{ohakx}{\B{o\-hakx}}, \HL{vAng}{\B{väng}}, \HL{ye}{\B{ye}}, \HL{yehakx}{\B{ye\-hakx}}, \HL{yevAng}{\B{ye\-väng}}, \HL{am'ake}{\B{am\-'a\-ke}}, \HL{tunu}{\B{tu\-nu}}, \HL{kxAl}{\B{kxäl}}, \HL{rayl}{\B{rayl}} be\slash feel happy, upset, hot, cold, hungry, thirsty, satisfied, full, quenched, confident\slash sure, romantic, apathetic, sorry (for). 
\B{Me\-nga 'efu väng srak?} Are you two thirsty?~\LINK{http://naviteri.org/2011/01/new-year-new-vocabulary/}{b} 
\B{Ngey\-ä tì\-kang\-kem\-ì\-ri 'e\-fe\-i\-u oe ye nì\-txan.} I'm very satisfied with your work.~\LINK{https://naviteri.org/2011/04/\%E2\%80\%99a\%E2\%80\%99awa-li\%E2\%80\%99fyavi-amip\%E2\%80\%94a-few-new-expressions/}{b} 
\B{Fu\-la ho\-ren sì ay\-lì'\-fya\-vi a\-mip le\-sar sa\-yi me\-ngar oe\-ti nit\-ram 'ey\-ke\-fu nì\-txan.} It makes me very happy that the new rules and words are useful to the two of you.~\LINK{http://naviteri.org/2012/06/spring-vocabulary-part-3/comment-page-1/\#comment-1515}{b}\\  
\textsc{Derivations}: \ding{43} 
\on{1}.~\HL{tI'efu}{\B{tì\-'e\-fu}}, feeling.
\on{2}.~\HL{txukxefu}{\B{txu\-kxe\-fu}}, care, be concerned about.

% 'eko
\hi
\HT{'eko}{\B{\cp{'e}ko}} \I{["{}E.ko]} \emph{vtr.} (\emph{a.} \ding{246}) attack. \B{Tsawl\-a pa\-lu\-lu\-kan\-ìl oe\-ti 'ìl\-me\-ko.} The big thanator has just attacked me.~\LINK{http://naviteri.org/2011/05/some-miscellaneous-vocabulary/}{b} 
(\emph{b. of the weather}) \emph{severe precipitation}. \B{Fì\-re\-won tom\-pa\-meyp zar\-mup, slä set 'ì\-me\-ko nì\-txan nang.} It was drizzling this morning, but it's really started coming down now!~\LINK{http://naviteri.org/2011/04/yafkeykiri-plltxe-frapo-everyone-talks-about-the-weather/}{b}

% 'ekong
\hi
\HT{'ekong}{\B{\cp{'e}kong}} \I{["PE.koN]} \emph{n.} (\emph{pl.} \B{me\-kong}, \B{pxe\-kong}, \B{ay\-e\-kong\slash e\-kong}) (rhythmic) beat. \B{'e\-kong te'\-lan\-ä}, the beat of the hearts.~\LINK{http://wiki.learnnavi.org/index.php/Weaving_Song}{wiki} \\
\textsc{Derivations}: \ding{43} 
\on{1}.~\HL{lI'kong}{\B{lì'\-kong}}, syllable. 
\on{2}.~\HL{txe'lankong}{\B{txe'\-lan\-kong}}, heartbeat.

% 'ekxin
\hi
\HT{'ekxin}{\B{'e\cp{kxin}}} \I{[PE."k'in]} \emph{adj.} tight. \B{Fì\-ra\-spu' 'e\-kxin lu nì\-hawng.} These leggings are too tight.~\LINK{http://naviteri.org/2012/07/meetings-waterfalls-and-more/}{b} (\ding{214}~\HL{um}{\B{um}}) 

% 'ekxinum
\hi
\HT{'ekxinum}{\B{'e\cp{kxi}num}} \I{[PE."k'i.num]} \emph{or} \I{[-nUm]} (\textsc{rn}: \B{'e\-\cp{kxi}\-nùm} \I{[PE."gi.nUm]}) \textsc{i.} \emph{n.} tightness\slash looseness. \\ 
\textsc{ii.} % 'ekxinumpe pekxinum
\B{'e\cp{kxi}numpe} \I{[PE."k'i.num.pE]} \emph{or} \B{pe\-\cp{kxi}\-num} \I{[pE."k'i.num]} \emph{inter.} how tight\slash loose? 
\B{Ngal mo\-lay' pxaw\-pxun\-it Lo\-ak\-ä a krr pe\-kxi\-num?} When you tried on Loak's armband, how tight was it?~\LINK{http://naviteri.org/2012/07/meetings-waterfalls-and-more/}{b}
(\ding{43}~\HL{'ekxin}{\B{'e\-kxin}}, \HL{um}{\B{um}}) 

% 'ele'wll
\hi 
\HT{'ele'wll}{\B{\cp{'e}le'wll}} \I{["PE.lEP.w\s{l}]} \emph{n.}, \ding{96} thorny paw (\emph{cactus-like Pandoran plant}), \emph{sigmoidia barbata}. %~\LINK{}{b}
(\ding{43}~\HL{'ewll}{\B{'e\-wll}}) 

% 'em
\hi
\HT{'em}{\B{'em}} \I{[PEm]} \emph{vtr.} cook. 
\B{Ke new oe mi\-vay' tsngan\-ti a 'o\-lem Ri\-nil.} I don't want to taste the meat that Rini cooked.~\LINK{http://naviteri.org/2012/11/renu-ayinanfyaya-the-senses-paradigm/}{b}
\B{Pey\-ral\-ìl zet wur\-a wu\-tsot a\-'aw\-nem pxel sngel.} Peyral treats a cool cooked meal like garbage.~\LINK{http://naviteri.org/2012/10/mipa-vospxi-mipa-ayliu-new-words-for-the-new-month/}{b} \\
\textsc{Derivations}: \ding{43} 
\on{1}.~\HL{'emyu}{\B{'e\-myu}}, cook. 
\on{2}.~\HL{tI'em}{\B{tì\-'em}}, the art of cooking.

% 'emyu
\hi
\HT{'emyu}{\B{\cp{'em}yu}} \I{["PEm.ju]} \emph{n.}~(\emph{pl.} \B{mem\-yu}, \B{pxem\-yu}, \B{ay\-em\-yu\slash emyu}) cook, someone who cooks. \B{Oe lu 'em\-yu.} I am a cook.~\LINK{http://naviteri.org/2010/09/getting-to-know-you-part-2/}{b} 
(\ding{43}~\HL{'em}{\B{'em}})

% 'en / si
\hi
\HT{'en}{\B{'en}} \I{[PEn]} \textsc{i.} \emph{n.} (\emph{pl.} \B{men}, \B{pxen}, \B{ay\-en\slash en}) (informed) guess, hunch, intuition. 
\B{Pe\-lun fì\-tseng\-ne nga zo\-la\-'u fte ti\-va\-ron? -- Lo\-lu 'en.} Why did you come here to hunt? -- It was a guess (hunch).~\LINK{http://naviteri.org/2011/01/new-year-new-vocabulary/}{b}
\B{Fì\-u\-van\-ì\-ri lu nga\-ru pxen nì\-'aw.} You only get three guesses in this game.~\LINK{http://naviteri.org/2011/01/new-year-new-vocabulary/}{b} \\ 
\textsc{ii.} \B{\cp{'en} si} \I{["PEn si]} \emph{vin.} make an informed guess.   
\B{'En si oe, Saw\-tu\-te\-ol tìl\-mok fì\-tseng\-it.} I guess some Skypeople were just here.~\LINK{http://naviteri.org/2011/01/new-year-new-vocabulary/}{b} \\  
\textsc{Derivations}: \ding{43} 
\on{1}.~\HL{le'en}{\B{le\-'en}}, speculative, intuitive.  
\on{2}.~\HL{nI'en}{\B{nì\-'en}}, making an informed guess.

% 'eng
\hi 
\HT{'eng}{\B{'eng}} \I{[PEN]} \emph{n., fauna} (\emph{pl.} \B{meng}, \B{pxeng}, \B{ay\-eng\slash eng}) beak of a bird or animal. 

% 'engeng
\hi
\HT{'engeng}{\B{\cp{'eng}eng}} \I{["PEN.EN]} \emph{adj.} level. \\
\textsc{Derivations}: \ding{43} 
\on{1}.~\HL{nI'eng}{\B{nì\-'eng}}, levelly, equally.

% 'eoio
\hi
\HT{'eoio}{\B{\cp{'e}oio}} \I{["PE.o.i.o]} \emph{adj.} ceremonious.   \B{Nì\-lun ay\-\mbox{ioi} a\-\mbox{'eoio} ay\-eyk\-tan\-ä lu lor fra\-to.} Of course the ceremonial wardrobes of the leaders are the most beautiful.~\LINK{http://naviteri.org/2011/08/new-vocabulary-clothing/}{b} \\  
\textsc{Derivations}: \ding{43} 
\on{1}.~\HL{sA'eoio}{\B{sä\-'e\-oio}}, ritual, rite.   
\on{2}.~\HL{nI'eoio}{\B{nì\-'e\-oio}}, ceremoniously.

% 'etnaw
\hi
\HT{'etnaw}{\B{\cp{'et}naw}} \I{["PEt\textcorner.naw]} \emph{n.} (\emph{pl.} \B{met\-naw}, \B{pxet\-naw}, \B{ay\-et\-naw\slash et\-naw}) shoulder. 
\B{Lo\-lu po\-ru txukx\-a skxir a ftu 'et\-naw a\-ski\-en ne me\-pun a\-ftär.} It had a deep wound from its right shoulder to its two left arms.~\LINK{http://naviteri.org/2020/04/awa-tskxekengtsyip-a-mikyunfpi-niul-one-more-little-listening-exercise/}{b}

% 'evan
\hi
\HT{'evan}{\B{\cp{'e}van}} \I{["PE.van]} \emph{n., coll.} (\emph{pl.} \B{me\-van}, \B{pxe\-van}, \B{ay\-e\-van\slash e\-van}) boy.    \B{'E\-van\-ìl \mbox{alo} a\-'aw\-ve nì\-'aw\-tu na'\-rìng\-it tar\-mok.} The boy was alone in the forest for the first time.~\LINK{http://naviteri.org/2011/09/miscellaneous-vocabulary/}{b} 
(\ding{43}~\HL{'evengan}{\B{'e\-veng\-an}})

% 'eve
\hi
\HT{'eve}{\B{\cp{'e}ve}} \I{["PE.vE]} \emph{n., coll.} (\emph{pl.} \B{me\-ve}, \B{pxe\-ve}, \B{ay\-e\-ve\slash e\-ve}) girl. \B{Se\-vin\-a tsa\-'e\-ve\-ru a\-hì\-'i sa'\-sem \mbox{ioi} so\-li fa mik\-tsang.} The parents put earrings on that pretty little girl.~\LINK{http://naviteri.org/2011/08/new-vocabulary-clothing/}{b} 
(\ding{43}~\HL{'evenge}{\B{'e\-veng\-e}})

% 'eveng
\hi
\HT{'eveng}{\B{\cp{'e}veng}} \I{["PE.vEN]} \emph{n.} (\emph{pl.} \B{me\-veng}, \B{pxe\-veng}, \B{ay\-e\-veng/e\-veng}) child. 
\B{Nì\-ran lu Lo\-ak mi 'e\-veng slä tsun ti\-va\-ron nì\-teng\-fya na fyeyn\-tu.} Loak is still really just a boy but he can hunt the same as an adult.~\LINK{http://naviteri.org/2012/10/mipa-vospxi-mipa-ayliu-new-words-for-the-new-month/}{b}  \\
\textsc{Derivations}: \ding{43} 
\on{1}.~\HL{'evengan}{\B{'e\-veng\-an}}, boy. 
\on{2}.~\HL{'evenge}{\B{'e\-veng\-e}}, girl. 
\on{3}.~\HL{'evi}{\B{'e\-vi}}, kid. 
\on{4}.~\HL{eywa'eveng}{\B{Ey\-wa\-'eveng}}, (the moon) Pandora. 
\on{5}.~\HL{kalweyaveng}{\B{kal\-wey\-a\-veng}}, `son of a bitch'.

% 'evengan
\hi
\HT{'evengan}{\B{\cp{'e}vengan}} \I{["PE.vEN.an]} \emph{n.} (\emph{pl.} \B{me\-ve\-ngan}, \B{pxe\-ve\-ngan}, \B{ay\-e\-ve\-ngan\slash e\-ve\-ngan}) boy. \B{Txe\-wì lu 'e\-veng\-an a so\-la\-lew zì\-sìt a\-vo\-mun.} Txewì is a ten\hyp{}year\hyp{}old boy.~\LINK{http://naviteri.org/2010/07/ting-mikyun-fte-tslivam-listening-comprehension-1-3/}{b} 
(\ding{43}~\HL{'eveng}{\B{'e\-veng}}, \HL{'evan}{\B{'e\-van}}) 

% 'evenge
\hi
\HT{'evenge}{\B{\cp{'e}venge}} \I{["PE.vEN.E]} \emph{n.} (\emph{pl.} \B{me\-veng\-e}, \B{pxe\-veng\-e}, \B{ay\-e\-veng\-e\slash e\-veng\-e}) girl. 
(\ding{43}~\HL{'eveng}{\B{'e\-veng}}, \HL{'eve}{\B{'e\-ve}})  

% 'evi
\hi
\HT{'evi}{\B{\cp{'e}vi}} \I{["PE.vi]} \emph{n.} (\emph{pl.} \B{me\-vi}, \B{pxe\-vi}, \B{ay\-e\-vi\slash e\-vi}) kid (\emph{affectionate form of `child'}). 
\B{Slä ay\-e\-vi nì\-wotx tso\-lun tsli\-vam (pot).} But all children could understand (her).~\LINK{http://naviteri.org/2010/07/ting-mikyun-fte-tslivam-listening-comprehension-1-3/}{b} 
(\ding{43}~\HL{'eveng}{\B{'e\-veng}})

% 'ewan
\hi
\HT{'ewan}{\B{\cp{'e}wan}} \I{["PE.wan]} \emph{adj.} young.   
\B{ay\-lì'\-u\-fa aw\-ngey\-ä 'ey\-lan\-ä a\-'e\-wan}, in the words of our young friend.~\LINK{http://forum.learnnavi.org/news-announcements/a-response-from-paul-frommer!/}{ln} 
\B{Ay\-hem\-ì\-ri 'e\-wan\-a tsa\-nan\-tang\-ur a\-hì\-'i tìng na\-ri.} Look at what that little young viperwolf is doing.~\LINK{http://naviteri.org/2012/01/more-additions-to-the-lexicon/}{b} 
(\ding{214}~\HL{koak}{\B{ko\-ak}}) \\
\textsc{Derivations}: \ding{43}
\on{1}.~\HL{kewan}{\B{ke\-wan}}, age.
\on{2}.~\HL{tI'ewan}{\B{tì\-'e\-wan}}, youth.

% 'ewll
\hi
\HT{'ewll}{\B{\cp{'e}wll}} \I{["PE.w\s{l}]} \emph{n.} \ding{96} (\emph{pl.} \B{me\-wll}, \B{pxe\-wll}, \B{ay\-e\-wll\slash e\-wll}) plant. 
\B{\mbox{Oel} vewng fu\-ta ay\-e\-veng ni\-vume te\-ri ay\-ewll na'\-rìng\-ä.} I see to it that the children learn about the forest plants.~\LINK{http://naviteri.org/2010/09/getting-to-know-you-part-2/}{b} 
\B{Lo\-ak\-ìl pä\-nu\-to\-lìng fu\-ta kar oe\-ru fya\-'ot a 'ìp fko nem\-fa \mbox{ewll}.} Loak promised he’d teach me how to vanish into the bushes.~\LINK{http://naviteri.org/2013/01/awvea-posti-zisita-amip-first-post-of-the-new-year/}{b} \\
\textsc{Derivations}: \ding{43} 
\on{1}.~\HL{prrwll}{\B{prr\-wll}}, moss. 
\on{2}.~\HL{paywll}{\B{pay\-wll}}, water plant. 
\on{3}.~\HL{pxiwll}{\B{pxi\-wll}}, hermit bud plant. 
\on{4}.~\HL{vAfewll}{\B{vä\-fe\-wll}}, Centipede. 
\on{5}.~\HL{fwAkIwll}{\B{fwä\-kì\-wll}}, mantis orchid. 
\on{6}.~\HL{'ele'wll}{\B{'e\-le'\-wll}}, thorny paw. 
\on{7}.~\HL{pa'liwll}{\B{pa'\-li\-wll}}, direhorse pitcher plant. 
\on{8}.~\HL{tompawll}{\B{tom\-pa\-wll}}, geode.
\on{9}.~\HL{hIrumwll}{\B{hì\-rum\-wll}}, puffer plant. 
\on{10}.~\HL{fkxakewll}{\B{fkxa\-ke\-wll}}, itch plant. 
\on{11}.~\HL{kllpxiwll}{\B{kll\-pxi\-wll}}, lionberry. 
\on{12}.~\HL{fngapsutxwll}{\B{fngap\-sutx\-wll}}, metalfollowing plant.
\on{13}.~\HL{yawrwll}{\B{yawr\-wll}}, feather blade.

% 'ewrang
\hi
\HT{'ewrang}{\B{\cp{'ew}rang}} \I{["PEw.RaN]} \emph{n., text.} (\emph{pl.} \B{mew\-rang}, \B{pxew\-rang}, \B{ay\-ew\-rang\slash ew\-rang}) loom. \\
\textsc{Derivations}: \ding{43} 
\on{1}.~\HL{sa'ewrang}{\B{sa\-'ew\-rang}}, great loom, mother loom.

% 'eylan
\hi
\HT{'eylan}{\B{\cp{'ey}lan}} \I{["PEj.lan]} \emph{n.} (\emph{pl.} \B{mey\-lan}, 
\B{pxey\-lan}, \B{ay\-ey\-lan\slash ey\-lan}) friend.  
\B{ay\-ey\-lan\-ur oe\-yä sì ey\-lan\-ur lì'\-fya\-yä le\-Na'\-vi nìwotx}, to my friends and all friends of the Na'vi language.~\LINK{http://forum.learnnavi.org/news-announcements/a-response-from-paul-frommer!/}{ln} 
\B{Me\-nga lu oe\-yä 'ey\-lan.} You two are my friends.~\LINK{http://naviteri.org/2011/07/number-in-na\%E2\%80\%99vi/}{b} 
% \B{aw\-nge\-yä 'ey\-lan\-hu mll\-te (oe)}, (I) agree with our friend.~\LINK{http://masempul.org/2010/07/\%E2\%80\%99ite-polaheiem/}{sempul} 
\B{Sra\-ke smon ngar oe\-yä mey\-lan alu En\-tu sì Ka\-mun?} Do you know my friends En\-tu and Ka\-mun?~\LINK{http://naviteri.org/2010/09/getting-to-know-you-part-1/}{b} 
\B{ay\-nga\-ru oe\-yä pxey\-lan\-ti}, let me introduce to you my three friends.~\LINK{http://naviteri.org/2012/07/tskxekengtsyip-a-mikyunfpi-a-little-listening-exercise/}{b} 
(\ding{214}~\HL{kxutu}{\B{kxu\-tu}}, \HL{wAtu}{\B{wä\-tu}}) \\  
\textsc{Derivations}: \ding{43} 
\on{1}.~\HL{tI'eylan}{\B{tì\-'ey\-lan}}, friendship.
\on{2}.~\HL{'eylanay}{\B{'ey\-la\-nay}}, acquaintance.
\on{3}.~\HL{'eylyong}{\B{'eyl\-yong}}, pet.
\on{4}.~\HL{let'eylan}{\B{let\-'ey\-lan}}, friendly (ofp.).
\on{5}.~\HL{txeylan}{\B{txey\-lan}}, best friend.

% 'eylanay
\hi
\HT{'eylanay}{\B{'eyla\cp{nay}}} \I{[PEj.la."{}naj]} \emph{n.} (\emph{pl.} \B{mey\-la\-nay}, \B{pxey\-la\-nay}, \B{ay\-ey\-la\-nay\slash ey\-la\-nay}) acquaintance (with the potential of becoming a friend).%~\LINK{}{b}  
(\ding{43}~\HL{'eylan}{\B{'ey\-lan}})

% 'eylyong
\hi 
\HT{'eylyong}{\B{\cp{'eyl}yong}} \I{["PEjl.joN]} \emph{n.} pet. 
\B{Tì\-kan fì\-pay\-oang\-ä ke lu fwa tsun fko pot yi\-vom; lu oe\-yä 'eyl\-yong.} This fish is not meant to be eaten; it's my pet.~\LINK{http://naviteri.org/2022/11/krr-ao-an-exciting-time/}{b}
(\ding{43}~\HL{'eylan}{\B{'ey\-lan}}, \HL{ioang}{\B{ioang}})

% 'eyng
\hi
\HT{'eyng}{\B{'eyng}} \I{[PEjN]} \emph{vin.} answer, respond. 
\B{'Aw\-a tì\-pawm\-ì\-ri 'i\-veyng oe set;  ay\-la\-ri zu\-saw\-krr 'ay\-eyng.} Let me answer one of the questions now; the others (I'll) answer in the future.~\LINK{http://forum.learnnavi.org/language-updates/email-from-frommer-re-numbers-\%28or-the-full-number-chart!\%29/}{ln}
\B{'Eyng oe\-ru set!} Answer me now!~\LINK{http://naviteri.org/2020/05/tipusawm-tiuseyng-si-okvur-a-eltur-titxen-si-asking-answering-and-an-interesting-story/}{b}
(\ding{214}~\HL{pawm}{\B{pawm}}, \HL{vin}{\B{vin}}) \\  
\textsc{Derivations}: \ding{43} 
\on{1}.~\HL{tI'eyng}{\B{tì\-'eyng}}, answer, respond.   
\on{2}.~\HL{nI'eyng}{\B{nì\-'eyng}}, in response.

% 'eyt
\hi
\HT{'eyt}{\B{'eyt}} \I{[PEjt\textcorner]} \emph{num.}, \textsc{lw} eight, (the symbol) \on{8} (\emph{only as Human decimal figure}). 
(\ding{43}~\HL{nayn}{\B{nayn}})

% 'i'a
\hi
\HT{'i'a}{\B{\cp{'i'}a}} \I{["PiP.a]} \emph{vin.} end, conclude. 
\B{ye'\-rìn 'ì\-yi\-'a sä\-nu\-me a tsa\-ri kll\-fro' oe}, soon the teaching I am responsible for will end.~\LINK{http://forum.learnnavi.org/language-updates/trr-rrtaya/}{ln}
\B{Ney\-ti\-ril nan\-tang\-it tso\-lukx fte pe\-yä tì\-sraw\-ti 'ey\-ki\-vi\-'a.} Neytiri stabbed the viperwolf to end its pain.~\LINK{http://naviteri.org/2017/12/zisit-amip-lefpom-ma-eylan-happy-new-year-friends/}{b}
(\ding{214}~\HL{sngA'i}{\B{sngä\-'i}})  \\  
\textsc{Derivations}: \ding{43} 
\on{1}.~\HL{tI'i'a}{\B{tì\-'i'\-a}}, ending, conclusion.  
\on{2}.~\HL{nI'i'a}{\B{nì\-'i'\-a}}, finally.

% 'ia
\hi
\HT{'ia}{\B{\cp{'i}a}} \I{["Pi.a]} \emph{vin.} lose oneself (\emph{spiritual sense}). \B{Eo Ey\-wa oe \mbox{'ia}.} I lose myself before Eywa.~\LINK{http://forum.learnnavi.org/language-updates/sieyng-a-ftu-naring-3-proverbs-what-we-asked-for-and-some-vocabulary/msg327513/\#msg327513}{a\slash ln}

% 'ipu
\hi
\HT{'ipu}{\B{\cp{'i}pu}} \I{["Pi.pu]} \emph{adj.} humorous, funny, amusing.   \B{Kaw\-krr ke lu peyä ay\-vur 'ipu kaw\-'it.} His stories are never a bit amusing.~\LINK{http://naviteri.org/2012/01/mipa-zisit-ayliu-amip-new-words-for-the-new-year/}{b} \\  
\textsc{Derivations}: \ding{43} 
\on{1}.~\HL{tI'ipu}{\B{tì\-'i\-pu}}, humor. 
\on{2}.~\HL{sA'ipu}{\B{sä\-'i\-pu}}, s.t.~humorous.

% 'it
\hi
\HT{'it}{\B{'it}} \I{[Pit\textcorner]} \emph{n.} bit, small amount.\\  
\textsc{Derivations}: \ding{43} 
\on{1}.~\HL{nI'it}{\B{nì\-'it}}, a bit. 
\on{2}.~\HL{kaw'it}{\B{kaw\-'it}}, (not) at all.

% 'itan
\hi
\HT{'itan}{\B{\cp{'i}tan}} \I{["Pi.tan]} \emph{n.} son. \B{Fra\-lo a ta\-ron, oe\-yä 'itan txo\-pu si nì\-nä\-nän.} Each time he hunts, my son becomes less and less afraid.~\LINK{http://naviteri.org/2012/02/trr-asawnung-lefpom-happy-leap-day/}{b} \B{Peyä 'itan\-ìri lu ho\-na nì\-txan a fì\-'u law lu fra\-por.} It's clear to everyone that his son is very cute.~\LINK{http://naviteri.org/2012/01/more-additions-to-the-lexicon/}{b} \\
\textsc{Derivations}: \ding{43} 
\on{1}.~\HL{maitan}{\B{mai\-tan}}, my son (address).

% 'ite
\hi
\HT{'ite}{\B{\cp{'i}te}} \I{["Pi.tE]} \emph{n.} daughter. \B{Za\-'u fì\-tseng, ma 'ite\-tsyìp.} Come here, little daughter.~\LINK{http://naviteri.org/2010/07/diminutives-conversational-expressions/}{b} \B{Nga\-ri leiu oe\-ru nrra nì\-txan, ma 'ite.} I'm very proud of you, daughter.~\LINK{http://naviteri.org/2012/07/fivospxiya-aylifyavi-amip-this-months-new-expressions-pt-2/}{b} \\
\textsc{Derivations}: \ding{43} 
\on{1}.~\HL{maite}{\B{mai\-te}}, my daughter (address).

% 'ì'awn
\hi
\HT{'I'awn}{\B{'ì\cp{'awn}}} \I{[PI."Pawn]} \emph{vin.} remain, stay. 
\B{Nga\-ri hu Ey\-wa sa\-lew ti\-re\-a, tokx 'ì'awn slu Na'\-vi\-yä ha\-pxì.} Your spirit goes with Eywa, your body remains to become part of the People.~\LINK{http://wiki.learnnavi.org/index.php?title=Corpus\#Behind_the_Scenes}{a\slash wiki} 
\B{'Ì'awn alìm!} Stay back!~\LINK{http://forum.learnnavi.org/language-updates/near-distant-and-irregular-adverbs/}{a\slash ln}
\B{Ke\-tsran tu\-te ni\-vew hi\-vum, po\-ru pll\-txe san ru\-txe 'i\-vì\-'awn.} No matter who wants to leave, tell them to please stay.~\LINK{http://naviteri.org/2013/03/whoever-whatever-whenever/}{b} \\
\textsc{Derivations}: \ding{43} 
\on{1}. \HL{'awnIm}{\B{'aw\-nìm}}, avoid.

% 'ìheyu
\hi
\HT{'Iheyu}{\B{'ì\cp{he}yu}} \I{[PI."hE.ju]} \emph{n.} spiral.   
\B{'ì\-he\-yu sì\-rey\-ä}, \emph{the spiral of the lives}.~\LINK{http://wiki.learnnavi.org/index.php/Weaving_Song}{wiki} 
\B{Ay\-wayl yìm ki\-fkey\-ä | 'Ì\-he\-yut a\-vo\-mrr}, \emph{The songs bind} | \emph{the thirteen spirals of the solid world.}~\LINK{http://www.pandorapedia.com/navi/music/ritual_music}{pp} \\
\textsc{Derivations}: \ding{43} 
\on{1}.~\HL{anIheyu}{\B{anì\-he\-yu}}, fibonacci.

% 'ìn
\hi
\HT{'In}{\B{'ìn}} \I{[PIn]} \emph{vin.} be busy, be occupied.   
\B{'Ìn nga fya\-pe nì\-fkrr?} What's been keeping you busy lately?~\LINK{http://naviteri.org/2010/09/getting-to-know-you-part-3/}{b} \\  
\textsc{Derivations}: \ding{43} 
\on{1}.~\HL{tIn}{\B{tìn}}, activity. 
\on{2}.~\HL{tIk'In}{\B{tìk\-'ìn}}, free time. 
\on{3}.~\HL{sulIn}{\B{su\-lìn}}, be busy (positive). 
\on{4}.~\HL{vrrIn}{\B{vrr\-ìn}}, be busy (negative). 
\on{5}.~\HL{kan'In}{\B{kan\-'ìn}}, focus on, specialize in.

% 'ìnglìsì
\hi
\HT{'InglIsI}{\B{\cp{'Ìng}lìsì}} \I{["PIN.lI.sI]} \emph{n.} English language. (\emph{mostly used in its derivations} \ding{43}~\B{le\-'Ìng\-lì\-sì} \emph{and} \ding{43}~\B{nì\-'Ìng\-lì\-sì}) 
\B{Krro lu 'Ìng\-lì\-sì txur nì\-hawng mì re\-'o.} Sometimes English is too strong in (my) head.~\LINK{http://naviteri.org/2011/09/\%E2\%80\%9Cby-the-way-what-are-you-reading\%E2\%80\%9D/comment-page-1/\#comment-1092}{b}
\B{Fu\-ri\-a fì\-lì'\-fya\-vit fkol ta 'Ìng\-lì\-sì mo\-lu\-nge, sìl\-pey oe tsnì ay\-ha\-pxì\-tu lì'\-fya\-o\-lo'\-ä aw\-nge\-yä ke stì\-ye\-vi.} That this expression is brought from English, will hopefully not anger the members of the language group.~\LINK{http://forum.learnnavi.org/language-updates/txelanit-hivawl/}{ln} \\
\textsc{Derivations}: \ding{43} 
\on{1}.~\HL{le'InglIsI}{\B{le\-'Ìng\-lì\-sì}}, English, regarding English.

% 'ìp
\hi
\HT{'Ip}{\B{'ìp}} \I{[PIp\textcorner]} \emph{vin.} disappear, vanish, recede from view. \B{Lo\-ak\-ìl pä\-nu\-to\-lìng fu\-ta kar oe\-ru fya\-'ot a 'ìp fko nem\-fa \mbox{ewll}.} Loak promised he'd teach me how to vanish into the bushes.~\LINK{http://naviteri.org/2013/01/awvea-posti-zisita-amip-first-post-of-the-new-year/}{b}  
\B{Oe\-yä tska\-lep 'o\-lìp ning\-yen.} My crossbow has mysteriously disappeared.~\LINK{http://naviteri.org/2013/01/awvea-posti-zisita-amip-first-post-of-the-new-year/}{b} 
(\emph{a. proverb}) \B{Nì\-win 'ìp tì\-'e\-wan, nì\-win pä\-hem tì\-ko\-ak.} Youth vanishes quickly, old age arrives quickly.~\LINK{https://naviteri.org/2025/04/mevosinga-liu-amip-twenty-new-words/}{b} 
(\ding{214}~\HL{srer}{\B{srer}})

% 'llngo
\hi
\HT{'llngo}{\B{\cp{'ll}ngo}} \I{["P\s{l}.No]} \emph{n.} (\emph{pl.} \B{me\-'ll\-ngo}, \B{pxe\-'ll\-ngo}, \B{ay\-'ll\-ngo}) hip. 
\B{Fu\-r\-i\-a 'll\-ngo nge\-yä zo\-slu, 'e\-fu oe nit\-ram nì\-'aw.} I'm just happy that your hip is getting better.~\LINK{http://naviteri.org/2013/07/pximaw-tsawlultxa-just-after-the-meet-up/comment-page-1/\#comment-2484}{b}

% 'o'
\hi
\HT{'o'}{\B{'o'}} \I{[PoP]} \emph{adj.} (bringing) fun, exciting.   
\B{Ay\-u\-van le\-tokx 'o' lu nì\-txan.} Sports are great fun.~\LINK{http://naviteri.org/2012/06/spring-vocabulary-part-3/}{b} 
\B{'O'a Ftxozä Hälowinä, ma frapo!} Happy Halloween, everyone!~\LINK{http://naviteri.org/2011/10/more-vocabulary-a-bit-of-grammar/}{b} \\  
\textsc{Derivations}: \ding{43} 
\on{1}.~\HL{tI'o'}{\B{tì\-'o'}}, fun, excitement. 
\on{2}.~\HL{nI'o'}{\B{nì\-'o'}}, enjoyably, fun\hyp{}ly.

% 'ok
\hi
\HT{'ok}{\B{'ok}} \I{[Pok\textcorner]} \emph{n.} memory, remembrance. \B{Teng\-krr pa\-lu\-lu\-kan moe\-ne kxll sar\-mi, pol\-txe Ney\-ti\-ril ay\-lì\-'ut a fra\-krr 'ok seyä la\-yu oer.} As the thanator was charging towards the two of us, Neytiri said something I will always remember.~\LINK{http://thereadinginpublicproject.blogspot.de/2010/05/paul-frommer.html}{pr} 
\B{'Ok oe\-yä lu ay\-ngar srak?} Do you all remember me?~\LINK{http://naviteri.org/2012/06/spring-vocabulary-part-3/}{b} 
(\ding{43}~\HL{zerok}{\B{ze\-rok}}) \\
\textsc{Derivations}: \ding{43} 
\on{1}.~\HL{'okrol}{\B{'ok\-rol}}, (ancient) history. 
\on{2}.~\HL{'okvur}{\B{'ok\-vur}}, (recent) history. 
\on{3}.~\HL{'oktrr}{\B{'ok\-trr}}, day of commemoration.

% 'okrol
\hi
\HT{'okrol}{\B{'ok\cp{rol}}} \I{[Pok\textcorner."Rol]} \emph{n.} (ancient) history. 
(\ding{43}~\HL{'ok}{\B{'ok}}, \HL{rol}{\B{rol}}; \HL{'okvur}{\B{'ok\-vur}})

% 'oktrr
\hi 
\HT{'oktrr}{\B{\cp{'ok}trr}} \I{["Pok\textcorner.t\s{r}]} \emph{n.} day of commemoration. 
(\ding{43}~\HL{'ok}{\B{'ok}}, \HL{trr}{\B{trr}})

% 'okvur
\hi
\HT{'okvur}{\B{'ok\cp{vur}}} \I{[Pok\textcorner."vuR]} \emph{or} \I{[."vUR]} \emph{n.} (non\hyp{}ancient) history. \B{oe\-yä so\-ai\-ä 'ok\-vur}, my family's history. \B{'ok\-vur Saw\-tu\-te\-yä mì Ey\-wa\-'e\-veng}, the history of the Sky People on Pandora.~\LINK{http://naviteri.org/2010/07/vocabulary-update/}{b} 
(\ding{214}~\HL{'ok}{\B{'ok}}, \HL{vur}{\B{vur}}; \HL{'okrol}{\B{'ok\-rol}})

% 'om
\hi
\HT{'om}{\B{'om}} \I{[Pom]} \emph{adj.} violet to purple to magenta.
\B{'om na mik\-yun}, the purplish color on the inside of a Na'vi ear.~\LINK{http://naviteri.org/2010/08/mipa-ayopin-mipa-ayliu-new-colors-new-words/}{b} \\
\textsc{Derivations}: \ding{43} 
\on{1}.~\HL{'ompin}{\B{'om\-pin}}, purple color.

% 'ompin
\hi
\HT{'ompin}{\B{\cp{'om}pin}} \I{["Pom.pin]} \emph{n.} the color \ding{43}~\HL{'om}{\B{'om}}. 
(\ding{43}~\HL{'opin}{\B{'o\-pin}})

% 'on / 'on si
\hi
\HT{'on}{\B{'on}} \I{[Pon]} \textsc{i}. \emph{n.} shape, form.
\B{Tsun fko ay\-on\-ti fì\-wopx\-ä ni\-vìn fte ya\-fkeyk\-it sre\-si\-ve\-'a.} The shapes of clouds can be used to predict the weather.~\LINK{http://naviteri.org/2013/09/on-si-salewfya-shapes-and-directions/}{b} 
\B{'On\-ìri tsko\-ä lu tì\-kin sä\-la\-tem\-ä a\-hì\-'i.} The form of the bow requires a small change.~\LINK{http://naviteri.org/2015/11/vomuna-liu-amip-ten-new-words/}{b} \\
\textsc{ii}. \B{\cp{'on} si} \I{["Pon si]} \emph{vin.} shape; give shape (to s.t.) (\emph{physical shaping or metaphorically}).
\B{Oe 'on si tskxe\-ru fte na ik\-ran li\-vam.} I shape a rock to look like an ik\-ran.~\LINK{http://naviteri.org/2018/07/ta-sulfatu-a-ayliu-niul-more-words-from-our-experts/}{b}
\B{Olo'\-ìri po\-an zu\-saw\-krr\-ur 'on so\-li.} He shaped the tribe's future.~\LINK{http://naviteri.org/2018/07/ta-sulfatu-a-ayliu-niul-more-words-from-our-experts/}{b} \\
\textsc{Derivations}: \ding{43} 
\on{1}.~\HL{ko'on}{\B{ko\-'on}}, ring, oval. 
\on{2}.~\HL{'ongop}{\B{'o\-ngop}}, design (\emph{for the larger aspects of design}).

% 'ong
\hi
\HT{'ong}{\B{'ong}} \I{[PoN]} \emph{vin.} unfold, blossom. (\emph{a.}) \B{Na se\-ze | A 'ong ne tsaw\-ke.} Like the blue\hyp{}flower | Which unfolds to the sun.~\LINK{http://www.pandorapedia.com/navi/music/ritual_music}{pp} 
\B{Lì'\-fya\-ri le\-Na'\-vi 'Rr\-ta\-mì, vay set 'al\-mong a fra\-'u ze\-ra\-'u ta ngrr\-po\-ngu.} Everything that has gone on with (blossomed regarding) Na'vi until now on Earth has come from a grassroots movement.~\LINK{http://wiki.learnnavi.org/index.php?title=Corpus\#Good_Morning_America}{wiki} 
(\emph{b. idiom})  \B{'Ivong Na'\-vi!} Let Na'vi bloom!~\LINK{http://forum.learnnavi.org/news-announcements/a-response-from-paul-frommer!/}{ln} 
\B{'ivong nìk\-'ong}, slow is fine. \B{Ke ze\-ne win sä\-pi\-vi. 'Ivong nìk\-'ong.} Take your time; don't rush. Slow is fine.~\LINK{http://naviteri.org/2010/09/quick-follow-up/}{b} \\  
\textsc{Derivations}: \ding{43} 
\on{1}.~\HL{tI'ong}{\B{tì\-'ong}}, unfolding, blooming. 
\on{2}.~\HL{'onglawn}{\B{'ong\-lawn}}, \emph{exhiliration of first bonding}.

% 'onglawn
\hi 
\HT{'onglawn}{\B{\cp{'ong}lawn}} \I{["PoN.lawn]} \emph{n.} \emph{exhiliration of first bonding} (\emph{used with} \ding{43}~\HL{'efu}{\B{'efu}}). 
\B{Kaw\-krr ke tswa\-ya' oel krr\-it a 'e\-fu 'ong\-lawn\-it.} I'll never forget the time I experienced \emph{'ong\-lawn}.~\LINK{http://naviteri.org/2020/11/vospxivopeya-ayliu-amip-novembers-new-words/}{b}
\B{Lu 'ong\-lawn tì\-'e\-fu a\-ko\-sman fra\-to mì hi\-fkey.} \emph{'Ong\-lawn} is the most wonderful feeling in the world.~\LINK{http://naviteri.org/2020/11/vospxivopeya-ayliu-amip-novembers-new-words/}{b}
(\ding{43}~\HL{'ong}{\B{'ong}}, \HL{lawnol}{\B{law\-nol}})

% 'ongokx
\hi
\HT{'ongokx}{\B{\cp{'o}ngokx}} \I{["Po.Nok']} \emph{vin.} be born.   
\B{oe 'o\-lo\-ngokx mì sray a txam\-pay\-ì\-ri sim}, I was born in a town near the ocean~\LINK{http://naviteri.org/2010/09/getting-to-know-you-part-2/}{b}
(\ding{43}~\HL{'ong}{\B{'ong}}, \HL{nokx}{\B{nokx}}) \\
\textsc{Derivations}: \ding{43} 
\on{1}.~\HL{tI'ongokx}{\B{tì\-'o\-ngokx}}, birth.

% 'ongop
\hi 
\HT{'ongop}{\B{\cp{'o}ngop}} \I{["Po.N$\cdot$op\textcorner]} \emph{vtr., inf.}\on{22} design (\emph{for the larger aspects of design; shape or structure}).
\B{Fì\-tsko\-ti a\-fyo\-lup 'o\-ngo\-lop oe\-yä sem\-pul\-ìl.} This exquisite bow was designed by my father.~\LINK{http://naviteri.org/2014/07/tsawlultxamaw-a-ayliu-post-meetup-words/}{b}
(\ding{43}~\HL{rengop}{\B{re\-ngop}}; \HL{'on}{\B{'on}}, \HL{ngop}{\B{ngop}}) \\
\textsc{Derivations}: \ding{43} 
\on{1}.~\HL{'ongopyu}{\B{'o\-ngop\-yu}}, designer.

% 'ongopyu
\hi 
\HT{'ongopyu}{\B{\cp{'o}ngopyu}} \I{["Po.Nop\textcorner.ju]} \emph{n.} designer 
(\ding{43}~\HL{'ongop}{\B{'o\-ngop}})

% 'opin
\hi
\HT{'opin}{\B{\cp{'o}pin}} \I{["Po.pin]} \emph{n.} color. \B{Tsa\-'o\-pin hek nì\-'it, slä su\-nu oer, ha ha'.} That color is bit odd, but I like it, so it’s a good fit for me.~\LINK{http://naviteri.org/2012/10/mipa-vospxi-mipa-ayliu-new-words-for-the-new-month/}{b} \\
\textsc{Derivations}: \ding{43} 
\on{1}.~\HL{'ompin}{\B{'om\-pin}}, the color violet\slash purple\slash megenta.
\on{2}.~\HL{eampin}{\B{eam\-pin}}, the color blue\slash green.
\on{3}.~\HL{kllvawmpin}{\B{kll\-vawm\-pin}}, the color brown.
\on{4}.~\HL{layompin}{\B{la\-yom\-pin}}, the color black.
\on{5}.~\HL{neympin}{\B{neym\-pin}}, the color value `light'.
\on{6}.~\HL{ngulpin}{\B{ngul\-pin}}, the color grey.
\on{7}.~\HL{rimpin}{\B{rim\-pin}}, the color yellow.
\on{8}.~\HL{teyrpin}{\B{teyr\-pin}}, the color white.
\on{9}.~\HL{tumpin}{\B{tum\-pin}}, the color red\slash orange.
\on{10}.~\HL{vawmpin}{\B{vawm\-pin}}, the color value `dark'.
\on{11}.~\HL{pxayopin}{\B{pxay\-opin}}, colorful. 
\on{12}.~\HL{'opinvultsyIp}{\B{'opin\-vul\-tsyìp}}, crayon. 
\on{13}.~\HL{'opinsung}{\B{'opin\-sung}}, color, color in.

% 'opinsung
\hi
\HT{'opinsung}{\B{\cp{'o}pinsung}} \I{["Po.pin.s$\cdot$uN]} \emph{or} \I{[.s$\cdot$UN]} (\textsc{rn}: \B{\cp{'o}pinsùng} \I{["Po.pin.s$\cdot$UN]}) \emph{inf.}\on{33} (\emph{a. vtr.}) color, color in; paint, apply paints\slash dyes.
\B{'Eveng\-ìl fì\-rel\-it 'opin\-so\-lung fa pin\-vul a\-rey\-pay\-tun.} The child colored (in) this picture with a red crayon.~\LINK{http://naviteri.org/2018/03/100a-liu-amip-64-new-words-part-1/\#comment-28355}{b}
(\emph{b. vin.}) \B{'Eveng 'opin\-sar\-mung mì fuk.} The child was coloring in a book.~\LINK{http://naviteri.org/2018/03/100a-liu-amip-64-new-words-part-1/\#comment-28355}{b}
(\ding{43}~\HL{'opin}{\B{'opin}}, \HL{sung}{\B{sung}})

% 'opinvultsyìp
\hi 
\HT{'opinvultsyIp}{\B{\cp{'o}pinvultsyìp}} \I{["Po.pin.vul.\t{ts}jIp\textcorner]} \emph{or} \I{[.vUl.]} \emph{n.} crayon (\emph{often shortened to} \ding{43}~\HL{pinvul}{\B{pinvul}}). 
(\ding{43}~\HL{'opin}{\B{'opin}}, \HL{vultsyIp}{\B{vul\-tsyìp}})

% 'ora
\hi
\HT{'ora}{\B{\cp{'o}ra}} \I{["Po.Ra]} \emph{n.} lake. \B{Fo kel\-ku si few 'o\-ra.} They live on the opposite side of the lake.~\LINK{http://naviteri.org/2011/10/more-vocabulary-a-bit-of-grammar/}{b} \\
\textsc{Derivations}: \ding{43} 
\on{1}.~\HL{'oratsyIp}{\B{'o\-ra\-tsyìp}}, pond, pool.

% 'oratsyìp
\hi
\HT{'oratsyIp}{\B{\cp{'o}ratsyìp}} \I{["Po.Ra.\t{ts}jIp\textcorner]} \emph{n.} pond, pool. \B{Ne 'o\-ra\-tsyìp po\-lä\-hem, yak si nì\-ftär.} When you arrive at the pond, turn to the left.~\LINK{http://naviteri.org/2013/09/on-si-salewfya-shapes-and-directions/}{b}
(\ding{43}~\HL{'ora}{\B{'o\-ra}})

% 'otxang
\hi
\HT{'otxang}{\B{'o\cp{txang}}} \I{[Po."t'aN]} \emph{n.} musical instrument (\emph{generic term}).   \B{Pam fì\-'o\-txang\-ä su\-nu oer nì\-ngay.} I really like the sound of this instrument.~\LINK{http://naviteri.org/2011/05/some-miscellaneous-vocabulary/}{b} 

% 'rrko
\hi
\HT{'rrko}{\B{\cp{'rr}ko}} \I{["P\s{r}.ko]} (\textsc{imperf} \B{'rr\-ko}) \emph{vin.} roll.
\B{Rum 'o\-lrr\-ko oe\-ne kll\-te\-sìn.} The ball rolled towards me on the ground.~\LINK{http://naviteri.org/2015/04/some-new-words-for-may-day/}{b}
\B{'Eveng\-ìl skxe\-vit 'ey\-krr\-ko sko uvan.} Children roll pebbles as a game.~\LINK{http://naviteri.org/2015/04/some-new-words-for-may-day/}{b}
\B{Tseyk 'ä\-pey\-ka\-mrr\-ko äo ut\-ral a zo\-lup fte hi\-vi\-fwo ftu ay\-sre' pa\-lu\-lu\-kan\-ä.} Jake rolled under the fallen tree to escape from the thanator's teeth.~\LINK{http://naviteri.org/2015/04/some-new-words-for-may-day/}{b}

% 'rrpxom
\hi
\HT{'rrpxom}{\B{\cp{'rr}pxom}} \I{["P\s{r}.p'om]} \emph{n.} thunder.%~\LINK{}{b}

% 'rrta
\hi
\HT{'Rrta}{\B{\cp{'Rr}ta}} \I{["P\s{r}.ta]} \emph{n.}, {\textsc{lw}} (\emph{planet}) Earth. \B{mì 'Rrta}, on Earth. 
\B{(Pol) Inan puk\-ot a te\-ri ay\-sam a 'Rr\-ta\-mì.} He's reading a book about the wars on Earth.~\LINK{http://naviteri.org/2011/09/\%E2\%80\%9Cby-the-way-what-are-you-reading\%E2\%80\%9D/}{ln} 
\B{Trr 'Rrtayä}, Earth Day.~\LINK{http://forum.learnnavi.org/language-updates/trr-rrtaya/}{ln} 
\B{Tsamsiyu a mì Saw 'Rrtayä}, \emph{Tsamsiyu} in the skies of Earth.~\LINK{http://naviteri.org/2013/03/tsamsiyu-a-mi-saw-rrtaya-warrior-in-our-sky/}{b} 

% 'u 
\hi
\HT{'u}{\B{'u}} \I{[Pu]} (\textsc{rn}: \B{'u} \I{[Pu]}) \textsc{i}.  \emph{n.} (\emph{no short plural, always:} \B{ayu}) thing (\emph{object, fact, abstraction}). 
\B{Kxawm set 'u a\-ngä\-zìk fra\-to lu la\-'a a\-yll.} Maybe the most difficult thing now is social distance.~\LINK{http://naviteri.org/2020/03/teri-tifkeytok-lefkrr-about-the-present-situation/}{b} 
\B{Krr\-ka tì\-rey a\-yol, pol mo\-lu'\-ni pxay\-a ay\-ut a tsran\-ten.} During her short life, she accomplished many important things.~\LINK{http://naviteri.org/2020/11/vospxivopeya-ayliu-amip-novembers-new-words/}{b}  \\
\textsc{ii}. \HL{'upe}{\B{\cp{'u}pe}} \I{["Pu.pE]} \emph{inter.} what (thing) (\emph{alternative form:} \B{pe\cp{u}}). \B{Tsa\-lì\-'u\-ri alu X, ral lu 'u\-pe?} What does X mean?~\LINK{http://naviteri.org/2010/09/getting-to-know-you-part-2/}{b} 
\B{Nga new yi\-vom 'u\-pet fì\-txon?} What do you want to eat tonight?~\LINK{http://naviteri.org/2013/03/whoever-whatever-whenever/}{b} 
(\ding{43}~\HL{zum}{\B{zum}}). \\
\textsc{Derivations}: \ding{43} 
\on{1}.~\HL{'uo}{\B{'uo}}, something. 
\on{2}.~\HL{fI'u}{\B{fì\-'u}}, this (thing). 
\on{3}.~\HL{tsa'u}{\B{tsa\-'u}}, that (thing). 
\on{4}.~\HL{fra'u}{\B{fra\-'u}}, everything. 
\on{5}.~\HL{ke'u}{\B{ke\-'u}}, nothing.
% (\emph{b}.) \B{tsa\cp{'u}} \I{[\t{ts}a."Pu]} \emph{pn., n.} (\emph{with case endings or enclitic adp.s often shortened to:} \B{tsa-} \emph{or} \B{tsaw-}, \emph{i.e.} \B{tsal, tsat(i), tsar(u), tseyä, tsari}) that (thing). \B{Ngal new a tsa\-'ut rä\-'ä wi\-vo, ma 'e\-vi.} Don't reach for what you want, child.~\LINK{http://naviteri.org/2012/01/mipa-zisit-ayliu-amip-new-words-for-the-new-year/}{b} \B{Fko zene sivung tsat tsatseng.} One can add that there.~\LINK{http://naviteri.org/2013/08/fmawn-a-fkol-kilmulat-this-just-in/comment-page-1/\#comment-2476}{b} \B{Fra\-'uä ran ngä\-pop fa fra\-syon tseyä.} The \emph{ran} of each thing arises from the totality of its attributes.~\LINK{http://naviteri.org/2012/10/mipa-vospxi-mipa-ayliu-new-words-for-the-new-month/}{b} \B{Tsa\-'u\-ri po pll\-ngay.} He admits it.~\LINK{http://naviteri.org/2011/07/txantsana-ultxa-mi-siatll-great-meeting-in-seattle/}{b}

% 'ul
\hi
\HT{'ul}{\B{'ul}} \I{[Pul]} \emph{or} \I{[PUl]} (\emph{a}.) \emph{vin.} increase. (\emph{b. to form correlative comparisons}) (i.) \B{'ul … 'ul}  the more … the more. \B{'Ul tskxe\-keng si, 'ul fnan.} The more you practice, the better you'll get.~\LINK{http://naviteri.org/2012/02/trr-asawnung-lefpom-happy-leap-day/}{b} 
\B{'Ul tu\-te, 'ul tì\-ngä\-zìk.} The more people, the more problems.~\LINK{http://naviteri.org/2012/02/trr-asawnung-lefpom-happy-leap-day/}{b} 
(ii.) \B{'ul … nän} the more … the less. (\ding{214}~\HL{nAn}{\B{nän}})\\  
\textsc{Derivations}: \ding{43} 
\on{1}.~\HL{tI.'ul}{\B{tì\-'ul}}, increase. 
\on{2}.~\HL{tsan'ul}{\B{tsan\-'ul}}, improve, get better. 
\on{3}.~\HL{fe'ul}{\B{fe\-'ul}}, worsen. 

% 'umtsa
\hi
\HT{'umtsa}{\B{\cp{'um}tsa}} \I{["{}um.\t{ts}a]} \emph{or} \I{["{}Um.]} (\textsc{rn}: \B{\cp{'ùm}tsa} \I{["PUm.\t{ts}a]}) \emph{n.} medicine (\HL{wA}{\B{wä}}, against).
\B{Ra\-lul ngo\-lop 'um\-tsat fa ay\-syu\-lang fwä\-kì\-wll\-ä.} Ralu made medicine from (\emph{or}: using the) flowers of the Mantis orchid.~\LINK{http://naviteri.org/2014/11/twenty-before-the-holidays/}{b} 
\B{Fì\-sä\-spxin\-ìri nge\-yä ke lä\-ngu kea 'um\-tsa.} Unfortunately there is no medicine for this disease of yours.~\LINK{http://naviteri.org/2014/11/twenty-before-the-holidays/}{b} 
\B{Wä fì\-sä\-spxin a 'um\-tsa le\-iu set tsu\-ka\-nom.} Medicine against this disease is happily now available.~\LINK{http://naviteri.org/2021/03/mipa-ayliu-si-aylifyavi-new-words-and-expressions/}{b} 

% 'uo
\hi
\HT{'uo}{\B{\cp{'u}o}} \I{["Pu.o]} \emph{pn., n.} something. \B{'Uo a fpi rey\-'eng Ey\-wa\-'e\-veng\-mì, 'Rr\-ta\-mì tsran\-ten nì\-txan aw\-nga\-ru nì\-wotx.} Something that matters a lot to all of us for the sake of The Balance of Life on both Pandora and Earth.~\LINK{http://forum.learnnavi.org/news-announcements/a-response-from-paul-frommer!/}{ln} 
\B{'Uol ik\-ran\-it txo\-pu sley\-ko\-la\-tsu, ta\-lun\-a po tsìk ya\-wo.} Something must have frightened the banshee, because it suddenly took to the air.~\LINK{http://naviteri.org/2012/01/mipa-zisit-ayliu-amip-new-words-for-the-new-year/}{b}
\B{Fì\-na\-er\-ìri ngal ew\-ku 'u\-ot a\-stxong srak?} Do you taste something strange in this drink?~\LINK{http://naviteri.org/2012/11/renu-ayinanfyaya-the-senses-paradigm/}{b}
(\ding{43}~\HL{'u}{\B{'u}})

% 'upe
\hi
\B{\cp{'u}pe}~:~\HL{'u}{\B{'u}}. 

% 'upxare
\hi
\HT{'upxare}{\B{'u\cp{pxa}re}} \I{[Pu."p'a.RE]} (\textsc{rn}: \B{'ù\cp{pxa}re} \I{[PU.ba.RE]}) \emph{n.} message.   
\B{Lu pì\-lok\-ur pxe\-sì\-kan sì pxe\-fne\-'upxa\-re.} The blog has three aims and three kinds of messages.~\LINK{http://naviteri.org/2010/06/first-post/}{b}
\B{Fpo\-le' ay\-ngal oer a 'upxa\-ret sto\-lawm oel.} I have heard the message you have sent me.~\LINK{http://forum.learnnavi.org/news-announcements/a-response-from-paul-frommer!/}{ln}
\B{Spä\-ngaw oel fu\-ta fì\-'upxa\-re\-yä tì\-ral\-peng ke lu eyawr.} Unfortunately, I don't believe the translation of this message is correct.~\LINK{http://naviteri.org/2015/03/kap-si-ayunil-saylahe-threats-dreams-and-other-things/}{b}

% 'ur
\hi
\HT{'ur}{\B{'ur}}\textsuperscript{\on{1}} \I{[PuR]} \emph{or} \I{[PUR]} (\textsc{rn}: \B{'ùr} \I{[PUR]}) \emph{n.} (\emph{a. sensation}) sight, look, appearance. 
(\emph{b. to form the middle voice}) \B{Nik\-re\-ri Ri\-ni\-yä 'ur fkan lor.} Rini's hair looks beautiful.~\LINK{http://naviteri.org/2012/11/renu-ayinanfyaya-the-senses-paradigm/}{b}

\hi
\B{'ur}\textsuperscript{\on{2}}~:~\HL{'u}{\B{'u}}.

\end{multicols}

% ====== Sektion A ======

 \pdfbookmark[0]{A}{A}
\begin{center}
\section*{A \I{[a]}}
\end{center}

\begin{multicols}{2}

% a
\hi
\HT{a}{\B{a}} \I{[a]} \emph{part.} \emph{used to subordinate an attributive sentence to a noun (relative clause)}. 
\B{Fpo\-le' ay\-ngal oer fì\-txan nì\-ftxa\-vang a 'u\-pxa\-ret sto\-lawm oel.} I have heard the message that you have sent me so passionately.~\LINK{http://forum.learnnavi.org/news-announcements/a-response-from-paul-frommer!/}{ln} 
\B{ke rey a tu\-te}, person who doesn't live.~\LINK{http://forum.learnnavi.org/language-updates/participles/}{ln} 
\B{fì\-txon yom sat a lu syä\-'ä nì\-'aw}, tonight (we) eat only those [vegetables] that are bitter.~\LINK{http://whyisthisnight.com/addl-more.html\#na\%27vi}{tn} 
\B{Tsap\-'a\-lu\-te, mì u\-pxa\-re a fpo\-le' oel ay\-nga\-ru trr\-am lu kxey\-ey}, Apologies, in the message that I sent you yesterday was a mistake.~\LINK{http://forum.learnnavi.org/language-updates/another-email-from-frommer/msg68133/\#msg68133}{ln} 
(with \emph{adp.s}) \B{Ma oeyä ey\-lan a\-yaw\-ne a lì'\-fya\-o\-lo'\-mì}, My beloved friends in the language clan.~\LINK{http://forum.learnnavi.org/news-announcements/happy-birthday-to-karyu-pawl/msg556899/\#msg556899}{ln} 
\B{Txan\-tsan\-a ul\-txa a mì Siätll so\-lu\-ne\-i\-u oer nì\-txan.} I liked the fantastic meeting in Seattle very much.~\LINK{http://naviteri.org/2011/07/attributive-\%E2\%80\%9Ca\%E2\%80\%9D-and-truncated-style/}{b} 
(with \emph{vin}.) \B{Kun\-sìp\-ì\-ri txan\-a tì\-meyp lu tsyal a mìn.} The gunship's main weakness is the rotor system. [lit.: the wing that turns.]~\LINK{http://forum.learnnavi.org/language-updates/ftarpa-and-skienpa-added-to-dict/msg570362/\#msg570362}{a\slash ln} 
\B{Lu fmawn a el\-tur tì\-txen si nì\-txan, ke\-fyak?} These are very interesting news, aren't they? [lit.: news that are interesting]~\LINK{http://naviteri.org/2011/09/miscellaneous-vocabulary/}{b} \\
(\emph{to form conj.s}): \ding{43}~\B{\HL{alu}{\B{a\-lu}}, \HL{krra}{krr\-a}, \HL{takrra}{ta\-krr\-a}, \HL{mawkrra}{maw\-krr\-a}, \HL{srekrra}{sre\-krr\-a}, \HL{fwa}{fwa}, \HL{fula}{fu\-la}, \HL{futa}{fu\-ta}, \HL{fura}{\B{fu\-ra}}, \HL{furia}{fu\-ri\-a}, \HL{fayluta}{fay\-lut\-a}, \HL{fmawnta}{fmawn\-ta}, \HL{kuma}{kum\-a}, \HL{teynga}{tey\-nga}, \HL{teyngla}{teyng\-la}, \HL{teyngta}{teyng\-ta}, \HL{teyngra}{\B{teyng\-ra}}, \HL{teyngA}{\B{tey\-ngä}}, \HL{teyngria}{\B{teyng\-ri\-a}}, \HL{tsenga}{tse\-nga}, \HL{tsonta}{tson\-ta}} 

% -a-
\hi
\HT{a-}{\B{a-}} \emph{or} \B{-a} \I{[a]} \emph{pref., suf.} (\emph{a. used to attribute an adj. to a noun} (\emph{attributive adj.})) 
\B{tsa\-krr pa\-ye\-'un swey\-a fya\-'ot}, then (I) will decide on the best way.~\LINK{http://forum.learnnavi.org/news-announcements/a-response-from-paul-frommer!/}{ln} 
\B{Ay\-lì\-'u\-fa aw\-nge\-yä 'ey\-lan\-ä a\-'e\-wan}, in the words of our young friend.~\LINK{http://forum.learnnavi.org/news-announcements/a-response-from-paul-frommer!/}{ln} 
(\emph{b. used to attribute active and passive participles to a noun}) 
\B{Pa\-lu\-lu\-kan a\-tu\-sa\-ron lu le\-hrr\-ap.} A hunting thanator is dangerous.~\LINK{http://forum.learnnavi.org/language-updates/participles/}{ln} 
\B{lì'\-fya a\-paw\-nll\-txe}, a spoken language.~\LINK{http://forum.learnnavi.org/language-updates/a-new-participle-infix-awn/msg140685/\#msg140685}{ln} 

% a fì'u
\hi
\B{a fì\cp{'u}}~:~\HL{fwa}{\B{fwa}}

% a fì'ul
\hi
\B{a fì\cp{'ul}}~:~\HL{fula}{\B{fula}}

% a fì'ur
\hi
\B{a fì\cp{'ur}}~:~\HL{fura}{\B{fura}}

% a fì'uri
\hi
\B{a fì\cp{'u}ri}~:~\HL{furia}{\B{furia}}

% a fì'ut
\hi
\B{a fì\cp{'ut}}~:~\HL{futa}{\B{futa}}

% a krr
\hi
\HT{a krr}{\B{a \cp{krr}}} \emph{or} \HL{krra}{\B{krr a}}, \HL{krra}{\B{\cp{krr}a}} \I{[a "k\s{r}]} \emph{or} \I{["k\s{r}.a]} \emph{conj.} when (\emph{to conclude a clause of time}). 
\B{tì\-'eyng\-it oel to\-lel a krr, ay\-nga\-ru pa\-yeng}, when I have recieved an answer, (I) will tell you.~\LINK{http://forum.learnnavi.org/news-announcements/a-response-from-paul-frommer!/}{ln} 
(\ding{43}~\HL{krr}{\B{krr}})

% a tsa'u
\hi
\B{a tsa\cp{'u}}~:~\HL{tsawa}{\B{tsawa}}

% a tsa'ul
\hi
\B{a tsa\cp{'ul}}~:~\HL{tsala}{\B{tsala}}

% a tsa'ur
\hi
\B{a tsa\cp{'ur}}~:~\HL{tsara}{\B{tsara}}

% a tsa'uri
\hi
\B{a tsa\cp{'u}ri}~:~\HL{tsaria}{\B{tsaria}}

% a tsa'ut
\hi
\B{a tsa\cp{'ut}}~:~\HL{tsata}{\B{tsata}}

% afpawng / si
\hi 
\HT{afpawng}{\B{a\cp{fpawng}}} \I{[a."fpawN]} \textsc{i}.~\emph{n.} grief. 
\B{Maw \mbox{kxitx} sem\-pul\-ä lar\-mä\-ngu Pey\-ral\-ä a\-fpawng txew\-luke.} After (her) father's death, Pey\-ral's grief was endless.~\LINK{http://naviteri.org/2018/03/100a-liu-amip-64-new-words-part-1/}{b} \\
\textsc{ii}. \B{a\cp{fpawng} si} \I{[a."fpawN si]} \emph{vin.} grieve. \\
\textsc{Derivatives}: \ding{43} 
\on{1}.~\HL{nafpawng}{\B{na\-fpawng}}, grievingly, with grief.

% Aha'ri
\hi
\B{A\cp{ha'}ri} \I{[a."haP.Ri]} \emph{n.} \emph{female name}.

% aho
\hi 
\HT{aho}{\B{a\cp{ho}}} \I{[a."ho]} \emph{vin.} pray. 
\B{Ey\-wa\-ru a\-ho, ma 'i\-tan, fte Nawm\-a Sa'\-nok\-ìl tì\-ye\-vìng ngar tì\-txur\-it.} Pray to Eywa, my son, that Great Mother will give you strength.~\LINK{http://naviteri.org/2018/03/100a-liu-amip-64-new-words-part-1/}{b} \\
\textsc{Derivatives}: \ding{43} 
\on{1}.~\HL{tIaho}{\B{tì\-aho}}, prayer (abstract). 
\on{2}.~\HL{saho}{\B{sa\-ho}}, prayer.

% akrrmaw
\hi
\B{a\cp{krr}maw}~:~\HL{mawkrra}{\B{mawkrra}}

% akrrta
\hi
\B{a\cp{krr}ta}~:~\HL{takrra}{\B{takrra}}

% akum
\hi
\B{a\cp{kum}}~:~\HL{kuma}{\B{kuma}}

% Akwey
\hi
\B{Ak\cp{wey}} \I{[ak\textcorner."wEj]} \emph{n.} \emph{male name}.
\B{Trr\-am\-ä ay\-u\-van\-ìri mak\-to Ak\-wey nì\-yew\-la ha snaytx.} Akwey rode disappointingly in yesterday's games so he lost.~\LINK{http://naviteri.org/2012/10/mipa-vospxi-mipa-ayliu-new-words-for-the-new-month/}{b}

% alaksi
\hi
\HT{alaksi}{\B{a\cp{lak}si}} \I{[a."lak\textcorner.si]} \emph{adj.} ready. 
\B{Alaksi srak, ma frapo?} Are you ready, everyone?~\LINK{http://naviteri.org/2012/07/tskxekengtsyip-a-mikyunfpi-a-little-listening-exercise/}{b}
\B{Ra\-lu rì\-kxi krr\-a srew, rì\-'ìr si pa\-lu\-kan\-ur a lu a\-lak\-si fte spi\-vä.} Ra\-lu does a shake while dancing, imitating a thanator that's ready to leap.~\LINK{http://naviteri.org/2018/03/100a-liu-amip-64-new-words-part-1/}{b}

% alìm
\hi
\HT{alIm}{\B{a\cp{lìm}}} \I{[a."lIm]} \emph{adv.} far away, at a distance. 
\B{'Ì'awn alìm!} Stay back!~\LINK{http://forum.learnnavi.org/language-updates/near-distant-and-irregular-adverbs/}{a\slash ln} 
(\ding{43}~\HL{lIm}{\B{lìm}}; \ding{214} \HL{asim}{\B{asim}} ) \\
\textsc{Derivatives}:\ding{43} \on{1}. \HL{'awnIm}{\B{'awnìm}}, avoid.

% ‹alm›
\hi
‹\HT{alm}{\B{alm}}› \I{[al.m]} \emph{inf. in} ‹1› \emph{to form the perfective past}. 
\B{na'\-rìng a ke lal\-mu tsar kea rìk}, the forest that had no leaves.~\LINK{http://forum.learnnavi.org/language-updates/for-\%28duration\%29-jesus-loan-word/}{ln} 
\B{Fkxi\-let a tsaw\-fa poe ioi sä\-pal\-mi ngo\-lop Va'\-rul.} Va'ru is the one who created the necklace she was wearing.~\LINK{http://naviteri.org/2011/08/new-vocabulary-clothing/}{b} 
\B{Lì'\-fya\-ri le\-Na'\-vi 'Rr\-ta\-mì, vay set 'al\-mong a fra\-'u ze\-ra\-'u ta ngrr\-po\-ngu.} Everything that has gone on with (blossomed regarding) Na'vi until now on Earth has come from a grassroots movement.~\LINK{http://wiki.learnnavi.org/index.php?title=Corpus\#Good_Morning_America}{wiki}

% alo
\hi
\HT{alo}{\B{\cp{a}lo}} \I{["{}a.lo]} \emph{n., adv.} time, turn, instance, \emph{one of a number of repeated or recurring actions}. 
\B{Ayupxareri angim nìhawng lu alo oeyä!} Now it's my turn for a message that's too long!~\LINK{http://forum.learnnavi.org/language-updates/txelanit-hivawl/}{ln}
\B{Fmi pi\-vll\-txe alo avol nì\-win!} Try to say (that) eight times fast!~\LINK{http://naviteri.org/2011/09/miscellaneous-vocabulary/}{b} 
\B{nì\-sìl\-pey alo ahay}, hopefully, next time.~\LINK{http://naviteri.org/2013/07/pximaw-tsawlultxa-just-after-the-meet-up/}{b} \\  
\textsc{Derivatives}:\ding{43}
\on{1}.~\HL{'awlo}{\B{'aw\-lo}}, once.  
\on{2}.~\HL{melo}{\B{me\-lo}}, twice. 
\on{3}.~\HL{pxelo}{\B{pxe\-lo}}, three times. 
\on{4}.~\HL{hayalo}{\B{ha\-ya\-lo}}, next time. 
\on{5}.~\HL{hamalo}{\B{ha\-ma\-lo}}, last time.

% alu
\hi
\HT{alu}{\B{\cp{a}lu}} \I{["{}a.lu]} \emph{conj.} (\emph{a. used for apposition}) that is, in other words \B{Tska\-lep\-it oel to\-lìng oe\-yä tsmu\-kan\-ur alu Ìstaw.} I gave the crossbow to my brother Istaw.~\LINK{http://naviteri.org/2010/07/vocabulary-update/}{b} 
\B{Txoa li\-vu, yaw\-ne lu oer So\-rewn … alu … ke tsun oeng mun\-txa sli\-vu} Sorry, but I love Sorewn \dots in other words, you and I cannot marry.~\LINK{http://naviteri.org/2010/07/vocabulary-update/}{b} 
\B{Srane, lì\-'ul alu mì fra\-krr ley\-ka\-tem lì\-'ut a\-hay tsa\-fya.} Yes, the word \emph{mì} always changes the next word like that.~\LINK{http://forum.learnnavi.org/language-updates/fmawno/msg293595/\#msg293595}{ln} 
(\emph{b. in contrastive comparisons}) \B{Fì\-kar\-yu alu fì\-po lu tsul\-fä\-tu; tsa\-kar\-yu alu tsa\-po lu \mbox{skxawng}.} \emph{This} teacher is a master; \emph{that} teacher is a fool.~\LINK{http://naviteri.org/2011/12/one-more-for-2011/}{b}
(\emph{c. to mean}) ``as, in the capacity of''. \B{Oe alu Tsa\-hìk ke tsun mì\-ftxe\-le tsngi\-vaw\-vìk; oe alu sa'\-nok tsun.} As Tsahik, I cannot weep over this matter; as a mother, I can.~\LINK{http://naviteri.org/2012/03/spring-vocabulary-part-2/}{b} 
(\ding{43}~\HL{sko}{\B{sko}})

% alunta
\hi
\B{a\cp{lun}ta}~:~\HL{talun1}{\B{taluna}}

% ‹aly›
\hi
‹\HT{aly}{\B{aly}}› \I{[al.j]} \emph{inf. in} ‹1› \emph{to form the perfective future}.

% ‹am›
\hi
‹\HT{am}{\B{am}}› \I{[a.m]} \emph{inf. in} ‹1› \emph{to form the past tense}. 
\B{Ay\-lì'\-fya yaw\-ne le\-ru oer ta\-krr\-a 'e\-veng la\-mu.} I've loved languages since I was a child.~\LINK{http://naviteri.org/2011/01/new-year-new-vocabulary/}{b}

% am'a
\hi
\HT{am'a}{\B{am\cp{'a}}} \I{[am."Pa]} \emph{n.} doubt. 
\B{Ke lu oer am'a.} I have no doubt.~\LINK{http://naviteri.org/2011/05/weather-part-2-and-a-bit-more-2/}{b} \\
\textsc{Derivatives}: \ding{43} \on{1}. \HL{am'ake}{\B{am'ake}}, sure, confident. 
\on{2}.~\HL{am'aluke}{\B{am'aluke}}, without a doubt.

% am'ake
\hi
\HT{am'ake}{\B{am\cp{'a}ke}} \I{[am."Pa.kE]} \emph{adj.} sure, confident. 
\B{Tsa\-ri\-a pol aw\-nga\-ti ke txa\-yìng oe 'e\-fu am\-'a\-ke nì\-wotx.} I'm entirely confident that he won't abandon us.~\LINK{http://naviteri.org/2012/06/spring-vocabulary-part-3/}{b}
(\ding{43}~\HL{am'a}{\B{am\-'a}})

% am'aluke
\hi
\HT{am'aluke}{\B{am\cp{'a}luke}} \I{[am."Pa.lu.kE]} \emph{adv.} without a doubt. 
\B{Am\-'a\-lu\-ke sna\-yaytx Saw\-tu\-te, ya\-yo\-ra' Na'\-vi.} Without a doubt the Sky People will lose and the People will win.~\LINK{http://naviteri.org/2012/06/spring-vocabulary-part-3/}{b}
(\ding{43}~\HL{am'a}{\B{am\-'a}}, \HL{luke}{\B{lu\-ke}})

% amay
\hi 
\HT{amay}{\B{a\cp{may}}} \I{[a."maj]} \emph{n.} brook. 

% Amerika
\hi
\B{Amerika} \I{[a."mE.Ri.ka]} \emph{n.}, \textsc{lw} America, North America.
\B{Trr le\-fpom, ma Su\-te Ame\-ri\-ka\-yä.} Good morning, people of America.~\LINK{http://wiki.learnnavi.org/index.php?title=Corpus\#Good_Morning_America}{wiki}

% Amhul
\hi
\B{Am\cp{hul}} \I{[am."hul]} \emph{or} \I{[."hUl]} \emph{n.} \emph{female name}.

% anìheyu
\hi
\HT{anIheyu}{\B{anì\cp{he}yu}} \I{[a.nI."hE.ju]} \emph{n.}~\ding{96} fibonacci, ``blue spiral plant'', \emph{orbis caeruleus}.
(\ding{43}~\HL{ean}{\B{e\-an}}, \HL{'Iheyu}{\B{'ì\-he\-yu}})

% Anuk
\hi
\B{Anuk} \I{[a.nuk\textcorner]} \emph{or} \I{[a.nUk\textcorner]} (\textsc{rn}: \B{A\-nùk} \I{[a.nUk\textcorner]})  \emph{n.} \emph{given name}.~\LINK{}{dh}

% Anurai
\hi
\B{\cp{A}nurai} \I{["{}a.nu.Ra.i]} \emph{n.} \emph{clan name, specialise in artisinal crafts}.
\B{Oe lu Anu\-rai\-yä syen\-a ha\-pxì a rey.} I am the last living member of the Anurai.~\LINK{http://naviteri.org/2015/05/cirque-du-soleils-toruk-the-first-flight-teaser-video-online/}{b}

% Aonung
\hi 
\HT{Aonung}{\B{\cp{A}onung}} \I{["{}a.o.nuN]} \emph{or} \I{[.nUN]} \emph{n. male name}. 

% apxa
\hi
\HT{apxa}{\B{a\cp{pxa}}} \I{[a."p'a]} \emph{adj} large. 
\B{Oeyk tsa\-tì\-snaytx\-ä lu a\-pxa sä\-nu\-i tì\-eyk\-tan\-ä.} That loss was caused by a great failure of leadership.~\LINK{http://naviteri.org/2013/04/wheres-the-bathroom-and-other-useful-things/}{b} 
\B{Mì sray apxa kelku si soaiahu.} (Txewì) lives in a large village with his family.~\LINK{http://naviteri.org/2010/07/ting-mikyun-fte-tslivam-listening-comprehension-1-3/}{b}  
\B{skxawng apxa}, prodigious moron.~\LINK{http://wiki.learnnavi.org/index.php?title=Corpus\#Good_Morning_America}{wiki} \\
\textsc{Derivatives}:\ding{43} \on{1}. \HL{apxangrr}{\B{apxangrr}}, delta tree.

% apxangrr
\hi
\HT{apxangrr}{\B{a\cp{pxa}ngrr}} \I{[a."p'a.N\s{r}]} \emph{n.} \ding{96} delta tree, ``large root'', \emph{magelum triangulare}.
(\ding{43}~\HL{apxa}{\B{a\-pxa}}, \HL{ngrr}{\B{ngrr}})

% ‹arm›
\hi
‹\HT{arm}{\B{arm}}› \I{[aR.m]} \emph{inf. in} ‹1› \emph{to form the imperfective past}. 
\B{Nga mal lar\-mu oer!} I trusted you!~\LINK{http://naviteri.org/2012/01/more-additions-to-the-lexicon/}{b} 
\B{Oel hu Txe\-wì trr\-am na'\-rìng\-it tar\-mok, tso\-le\-'a syep\-tu\-tet a\-tsawl fra\-to mì sì\-rey.} Yesterday, while I was with Txewì in the forest, we saw the biggest Trapper I've ever seen.~\LINK{http://wiki.learnnavi.org/index.php?title=Corpus\#New_York_Times}{ny\slash wiki}
\B{Po\-el oe\-ti lar\-mawk.} She was talking about me.~\LINK{http://naviteri.org/2010/09/getting-to-know-you-part-1/}{b} 
\B{Tì\-kang\-kem\-ì\-ri var\-mrr\-ìn oe nì\-wotx.} I was completely overwhelmed at work.~\LINK{http://naviteri.org/2010/09/getting-to-know-you-part-3/}{b} 
\B{\mbox{kxawm} oe har\-ma\-hä\-ngaw}, maybe I was sleeping.~\LINK{http://wiki.learnnavi.org/index.php?title=Canon\#Extracts_from_various_emails}{wiki}

% Artsut
\hi
\B{Artsut} \I{[aR.\t{ts}ut\textcorner]} \emph{or} \I{[.\t{ts}Ut\textcorner]} (\textsc{rn}: \B{Ar\-tsùt} \I{[aR.\t{ts}Ut\textcorner]}) \emph{n.} \emph{female name}.~\LINK{}{dh}

% Arvok
\hi
\B{Arvok} \I{[aR.vok\textcorner]} \emph{n.} \emph{male name}.~\LINK{}{dh}

% ‹ary›
\hi
‹\HT{ary}{\B{ary}}› \I{[aR.j]} \emph{inf.~in} ‹1› \emph{to form the imperfective future}.

% asim
\hi
\HT{asim}{\B{a\cp{sim}}} \I{[a."sim]} \emph{adv.} nearby, at close range. 
\B{New oe ri\-vun a\-sim tì\-fnu\-nga'\-a tseng\-it}, I want to find a quiet place nearby.~\LINK{http://naviteri.org/2013/01/awvea-posti-zisita-amip-first-post-of-the-new-year/}{b}
(\ding{43}~\HL{sim}{\B{sim}}, \ding{214}~\HL{alIm}{\B{alìm}})

% asip
\hi 
\HT{asip}{\B{\cp{a}sip}} \I{["{}a.sip\textcorner]} \emph{n.} (a.) tall thin mass or pile of s.t. (b.) tower. 
\B{Tu\-te\-ol a\-sip\-it ay\-tä\-re\-mä txo\-lu\-la mì na'\-rìng.} Someone built a tower of bones in the forest.~\LINK{http://naviteri.org/2023/09/aawa-ayliu-si-aylifyavi-amip-a-few-new-words-and-expressions/}{b}  

% ‹asy›
\hi
‹\HT{asy}{\B{asy}}› \I{[a.sj]} \emph{inf. in} ‹1› \emph{to form the intentional future}. 
\B{Ay\-oe ke wa\-syem.} We will not fight.~\LINK{http://forum.learnnavi.org/language-updates/ltasygt-with-ke-and-nga/}{g\slash ln} 
\B{Ke za\-syup lì\-'Ona ne kxu\-tu a mì\-fa fu a wrr\-pa.} The l'Ona will not perish to the enemy within or the enemy without.~\LINK{http://forum.learnnavi.org/language-updates/ltasygt-with-ke-and-nga/}{g\slash ln}

% atan 1
\hi
\HT{atan}{\B{a\cp{tan}}}\textsuperscript{\on{1}} \I{[a."tan]} \emph{n.} light. 
\B{Fra\-tseng lu atan, fnu fra\-po.} Everywhere is light, everyone is quiet.~\LINK{http://forum.learnnavi.org/news-announcements/christmas-song-ninavi-on-youtube/}{ln} \\ 
\textsc{Derivatives}: \ding{43} 
\on{1}.~\HL{atanvi}{\B{atanvi}}, ray (of light). \\  
\on{2}.~\HL{atanzaw}{\B{atanzaw}}, forked lightning. \\  
\on{3}.~\HL{syuratan}{\B{syuratan}}, bioluminescence. \\
\on{4}.~\HL{txanatan}{\B{txanatan}}, bright, vivid.

% Atan 2
\hi
\B{A\cp{tan}}\textsuperscript{\on{2}} \I{[a."tan]} \emph{n.} \emph{male name}.~\LINK{}{dh}

% atanvi
\hi
\HT{atanvi}{\B{a\cp{tan}vi}} \I{[a."tan.vi]} \emph{n.} ray (of light).
(\ding{43}~\HL{atan}{\B{a\-tan}})

% atanzaw
\hi
\HT{atanzaw}{\B{a\cp{tan}zaw}} \I{[a."tan.zaw]} \emph{n.} forked lightning.
(\ding{43}~\HL{atan}{\B{a\-tan}}, \HL{swizaw}{\B{swi\-zaw}}; \HL{rawm}{\B{rawm}})

% Ateyo
\hi
\B{A\cp{te}yo} \I{[a."tE.jo]} \emph{n.} \emph{male name}.
\B{Tsay\-sam\-si\-yu lu su\-pe? -- Fo lu Ka\-mun, Ra\-lu, Ìstaw, sì Ate\-yo.} They're Ka\-mun, Ra\-lu, Ìstaw, and Ate\-yo.~\LINK{http://naviteri.org/2011/07/number-in-na\%E2\%80\%99vi/}{b}

% atokirina'
\hi
\HT{atokirina'}{\B{atoki\cp{ri}na'}} \I{[a.to.ki."Ri.naP]} \emph{n.}, \ding{96} \emph{atokirina}, seeds of the great tree, woodsprite. 
\B{Pxay\-a a\-to\-ki\-ri\-na' ne fì\-ke\-tu\-wong zì\-ma\-'u.} Many \emph{atokirina} just came to this alien.~\LINK{http://forum.learnnavi.org/language-updates/pseudo-canon-from-bd-live-content/}{a\textsuperscript{c}\slash ln}
(\ding{43}~\HL{rina'}{\B{ri\-na'}})

% atxar
\hi
\HT{atxar}{\B{a\cp{txar}}} \I{[a."t'aR]} \emph{adj.} smell of living animals (\emph{as found around a watering hole or animal nest}). 

% atxkxe
\hi
\HT{atxkxe}{\B{atx\cp{kxe}}} \I{[at'."k'E]} \emph{n} land. 
\B{Tsun aw\-nga kel\-ku si\-vi nì\-'eng Saw\-tu\-te\-hu mì atx\-kxe aw\-nge\-yä.} We can share our land with the Skypeople.~\LINK{http://naviteri.org/2012/01/more-additions-to-the-lexicon/}{b} 
\B{Fì\-atx\-kxe\-ti ke ra\-ya\-'un ay\-oel kaw\-krr!} We will never give up this land.~\LINK{http://naviteri.org/2021/09/aawa-ayliu-amip-a-few-new-words/}{b} \\
\textsc{Derivatives}: \ding{43} 
\on{1}.~\HL{atxkxerel}{\B{atx\-kxe\-rel}}, map. 
\on{2}.~\HL{txanlokxe}{\B{txan\-lo\-kxe}}, tribal or clan domain; country.

% atxkxerel
\hi
\HT{atxkxerel}{\B{atx\cp{kxe}rel}} \I{[at'."k'E.REl]} \emph{n.} map.
(\ding{43}~\HL{atxkxe}{\B{atx\-kxe}}, \HL{rel}{\B{rel}})

% ‹ats›
\hi
‹\HT{ats}{\B{ats}}› \I{[a.\t{ts}]} \emph{inf. in} ‹2› \emph{to form the evidential, i.e. uncertainty or indirect knowledge}. 
\B{Fpìr\-mìl oel fu\-ta ay\-nga na\-tsew tsi\-ve\-'a fì\-'ut.} I was just thinking that you might want to see this.~\LINK{http://forum.learnnavi.org/language-updates/the-evidential-at-last/msg106055/\#msg106055}{ln} 
\B{Srung sa\-tsi fra\-por.} (That) might help everyone.~\LINK{http://naviteri.org/2011/08/reported-speech-reported-questions/comment-page-1/\#comment-1062}{b}
\B{Po\-an yaw\-ne la\-tsu poe\-ru nì\-lam.} Apparently she loves him.~\LINK{http://naviteri.org/2011/05/weather-part-2-and-a-bit-more-2/}{b}

% au
\hi
\HT{au}{\B{\cp{a}u}} \I{["{}a.u]} \emph{n.} drum (made of skin). 
\B{Po pam\-tse\-o si fa au nìl\-tsan nì\-txan.} She plays the drums very well.~\LINK{http://naviteri.org/2011/01/new-year-new-vocabulary/}{b}  \\  
\textsc{Derivatives}: \ding{43}
\on{1}.~\HL{wokau}{\B{wokau}}, pendulum drum.

% au
\hi
\HT{au2}{\B{a\cp{u}}} \I{[a."{}u]} \emph{intj.} \emph{exclamation of consternation}.~\LINK{https://forum.learnnavi.org/intermediate/my-dictionary/msg263252/\#msg263252}{ln} 

% aungia
\hi
\HT{aungia}{\B{a\cp{u}ngia}} \I{[a."{}u.Ni.a]} \emph{n.} sign, omen. 
\B{Tsa\-krr za\-'u a\-u\-ngi\-a ta Ey\-wa.} Then came a sign from Eywa.~\LINK{http://naviteri.org/2010/09/getting-to-know-you-part-1/comment-page-1/\#comment-296}{a\slash b} 
\B{Ze\-ne Sa'\-nok fì\-a\-u\-ngi\-a\-ti ral\-pi\-veng.} Mother has to interpret this omen.~\LINK{http://forum.learnnavi.org/language-updates/pseudo-canon-from-bd-live-content/}{ln\slash a\textsuperscript{c}}

% awaiei
\hi
\HT{awaiei}{\B{awai\cp{e}i}} \I{[a.wa.i."{}E.i]} \emph{n.} \ding{96} banshee of paradise, \emph{bansheba terrestre}.

% ‹ay›
\hi
‹\HT{ay1}{\B{ay}}› \I{[a.j]} \emph{inf. in} ‹1› (to form the future tense). 
\B{Pol\-txe Ney\-ti\-ril ay\-lì\-'ut a fra\-krr 'ok seyä la\-yu oer.} Neytiri said something that I will always remember.~\LINK{http://thereadinginpublicproject.blogspot.de/2010/05/paul-frommer.html}{pr}. 
\B{Tì\-'eyng\-it oel to\-lel a krr, ay\-nga\-ru pa\-yeng, tsa\-krr pa\-ye\-'un swey\-a fya\-'ot a za\-mi\-vu\-nge oel ay\-ngar ay\-lì\-'ut ho\-ren\-ti\-sì lì'\-fya\-yä le\-Na'\-vi.} When I've received an answer, I will tell you and then decide on the best way to bring you the words and rules of the Na'vi language.~\LINK{http://forum.learnnavi.org/news-announcements/a-response-from-paul-frommer!/}{ln}

\end{multicols}


% ====== section AW ======

 \pdfbookmark[0]{AW}{AW}
\begin{center}
\section*{AW \I{[aw]}}
\end{center}

\begin{multicols}{2}

% -aw
\hi
\HT{-aw}{\B{-aw}} \I{[.aw]} \emph{num., suf.} one (\emph{rest of} \ding{43}~\HL{'aw}{\B{'aw}} \emph{to form higher numbers}). \B{volaw}, nine (decimal), \on{11} (octal). \B{mevolaw}, \on{17} (decimal), \on{21} (octal); etc.

% ‹awn›
\hi
‹\HT{awn}{\B{awn}}› \I{[aw.n]} \emph{inf. in} ‹1› (to form the passive participle) 
(\emph{a. always used attributively}) \B{Ay\-lì\-'u a\-paw\-nll\-txe nìl\-tsan!} Words well-spoken.~\LINK{http://forum.learnnavi.org/language-updates/a-new-participle-infix-awn/msg140685/\#msg140685}{ln} 
\B{Ru\-txe oeyä tì\-o\-eyk\-tìng\-it a\-saw\-nung ni\-vìn.} Please look at my added explanation.~\LINK{http://naviteri.org/2011/03/word-order-and-case-marking-with-modals/comment-page-1/\#comment-581}{b} 
\B{Saw\-no'\-ha ioit ko\-la\-ne\-i\-om oel u\-van\-fa!} I got (won) the prized piece of jewelry in the game!~\LINK{http://naviteri.org/2012/06/spring-vocabulary-part-3/}{b}
(\emph{b. shortened in military jargon}) \B{Tì\-kan taw\-na\-tep!} Target lost!~\LINK{http://naviteri.org/2011/05/some-miscellaneous-vocabulary/}{g\slash b}
(\emph{c. shortened in proverbial expression}) \B{(Na) lo\-re\-yu 'aw\-nam\-pi}, extreme shyness, [lit.: (like) a touched helicoradian] \B{Lu por mok\-ri a\-mik\-lor, slä lo\-re\-yu 'aw\-nam\-pi lu. Ke tsun ri\-vol eo su\-te.} She has a beautiful voice, but she's extremely shy. She can't sing in front of people.~\LINK{http://naviteri.org/2014/11/twenty-before-the-holidays/}{b}
(\emph{d. in special adv. use}) \B{nìaw\cp{no}mum}, as you know, as is known. \B{nì\cp{zaw}nong}, safely.

% awnga
\hi
\HT{awnga}{\B{aw\cp{nga}}} \I{[aw."{}Na]} \emph{pn.} (\emph{shortened form of} \ding{43}~\HL{ayoeng}{\B{ay\-oeng}}) we all (\emph{inclusive}). 
\B{Aw\-nga sngi\-vä'i ko!} Let us begin!~\LINK{http://naviteri.org/2010/06/first-post/}{b} 
\B{aw\-ngal yom wu\-tsot teng\-krr he\-reyn nì\-pxim}, we eat dinner while sitting upright.~\LINK{http://whyisthisnight.com/addl-more.html\#na\%27vi}{tn} 
\B{Za\-ya\-'u Saw\-tu\-te fte aw\-nga\-ti ski\-va\-'a kay zì\-sìt a\-pxey!} The Sky People will come to destroy us three years from now!~\LINK{http://naviteri.org/2011/09/miscellaneous-vocabulary/}{b} 
\B{Po aw\-nga\-ru ka\-vuk so\-li.} He betrayed us.~\LINK{http://naviteri.org/2010/08/a-na\%E2\%80\%99vi-alphabet/comment-page-1/\#comment-189}{b}  
\B{Slo\-lu po tsul\-fä\-tu lì'\-fya\-yä aw\-nge\-yä.} He's become a master of our language.~\LINK{http://naviteri.org/2013/01/lifyari-po-peyi-how-good-is-her-navi/}{b}

\end{multicols}

 
% ===== Section AY =====

 \pdfbookmark[0]{AY}{AY}
\begin{center}
\section*{AY \I{[aj]}}
\end{center}

\begin{multicols}{2}

% ay+
\hi
\HT{ay}{\B{ay}}\texttt{+} \I{[aj.]} \emph{pref. to form the plural of a noun}. 
\B{ay\-nu\-me\-yu}, students; 
\B{ay\-zì\-sìt}, years; 
\B{ay\-sä\-o\-mum}, information; 
\B{ay\-ya\-yo}, birds. (can be left out when the noun undergoes lenition) \B{ay\-sa\-ron\-yu\slash sa\-ron\-yu}, hunters; 
\B{ay\-har\-yu\slash har\-yu}, teachers. \\
\textsc{Derivations:} \ding{43} \on{1}.~\HL{fay}{\B{fay}}\texttt{+}, those. 
\on{2}.~\HL{pay1}{\B{pay}}\texttt{+}, which (of them)? 
\on{3}.~\HL{fray}{\B{fray}}\texttt{+}, all (of them).

% ‹ay›
\hi
‹\B{ay}› \I{[a.j]} \emph{inf. in} ‹1› (to form the future tense). 
\B{Pol\-txe Ney\-ti\-ril ay\-lì\-'ut a fra\-krr 'ok seyä la\-yu oer.} Neytiri said something that I will always remember.~\LINK{http://thereadinginpublicproject.blogspot.de/2010/05/paul-frommer.html}{pr}. 
\B{Tì\-'eyng\-it oel to\-lel a krr, ay\-nga\-ru pa\-yeng, tsa\-krr pa\-ye\-'un swey\-a fya\-'ot a za\-mi\-vu\-nge oel ay\-ngar ay\-lì\-'ut ho\-ren\-ti\-sì lì'\-fya\-yä le\-Na'\-vi.} When I've received an answer, I will tell you and then decide on the best way to bring you the words and rules of the Na'vi language.~\LINK{http://forum.learnnavi.org/news-announcements/a-response-from-paul-frommer!/}{ln}

% ayfeyä
\hi
\B{ayfeyä}~:~\HL{ayfo}{\B{ayfo}}

% ayfo
\hi
\HT{ayfo}{\B{ay\cp{fo}}}, \emph{also} \HL{fo}{\B{fo}} \I{[aj."fo]} \emph{also} \I{[fo]} \emph{pn.} (\emph{gen.} \B{ay\-fe\-yä\slash feyä}) they (3\textsuperscript{rd} person plural). 
\B{Ay\-fo so\-lop ìlä hil\-van fa u\-ran.} They traveled along the river by boat.~\LINK{http://naviteri.org/2011/07/txantsana-ultxa-mi-siatll-great-meeting-in-seattle/}{b} 
\B{Ay\-oengal fte Fay\-saw\-tu\-tet ki\-vu\-rakx, ze\-ne nì\-'aw\-ve ay\-fot tsli\-vam.} We must understand these Sky People if we are to drive them out.~\LINK{http://forum.learnnavi.org/language-updates/confirmation-on-kurakx-eytukan-line-and-toitslan/}{ln} 
\B{Fo\-ru 'u\-pxa\-ret oel fpo\-le'.} I have sent them a message.~\LINK{http://forum.learnnavi.org/news-announcements/a-response-from-paul-frommer!/}{ln} \B{hol\-pxay ay\-zek\-wä\-yä fe\-yä}, the number of their fingers.~\LINK{http://forum.learnnavi.org/language-updates/talun-taweyk-tafral-lun-oeyk-holpxay-hol/}{ln} 
\B{Fo\-ri maw\-krr\-a fa ren\-ten ioi sä\-po\-li ho\-lum.} After they put on their goggles, they left.~\LINK{http://naviteri.org/2011/08/new-vocabulary-clothing/}{b}
(\ding{43}~\HL{po}{\B{po}}, \HL{ay}{\B{ay}}\texttt{+})

% ayla
\hi
\HT{ayla}{\B{ay\cp{la}}} \I{[aj."la]} : \HL{aylahe}{\B{aylahe}}. 

% aylahe
\hi
\HT{aylahe}{\B{ay\cp{la}he}} \I{[aj."la.hE]} \emph{pn.} the others, others (\emph{forms:}) \B{aylahe}, \B{aylahel}, \B{aylahet(i)}, \B{aylaheyä}, \B{aylaher(u)}, \B{aylaheri}). 
\B{Kem\-ìri a nga\-ru prr\-te' ke lu, tsa\-kem rä\-'ä si\-vi ay\-la\-he\-ru.} As for any action that you don't like, don't act that way to others.~\LINK{http://wiki.learnnavi.org/index.php?title=Corpus\#Good_Morning_America}{wiki} 
(\emph{shortened forms:} \B{ayla}, \B{aylal}, \B{aylat(i)}, \B{ayleyä}, \HL{aylaru}{\B{ay\-lar(u)}}, \B{aylari})
\B{Fì\-po\-ti oel tspì\-yang, fte tì\-ke\-nong li\-ye\-vu ay\-la\-ru.} I will kill this one as a lesson to the others.~\LINK{http://forum.learnnavi.org/language-updates/capturing-avatar/}{b}
\B{'Aw\-a tì\-pawm\-ì\-ri 'i\-veyng oe set;  ay\-la\-ri zu\-saw\-krr 'ay\-eyng.} Let me answer one of the questions now; the others (I'll) answer in the future.~\LINK{http://forum.learnnavi.org/language-updates/email-from-frommer-re-numbers-\%28or-the-full-number-chart!\%29/}{ln}
\ding{43}~\HL{lahe}{\B{la\-he}})

% aylal
\hi
\HT{aylal}{\B{ay\cp{lal}}} \I{[aj."lal]} : \HL{aylahe}{\B{aylahe}}. 

% aylari
\hi
\HT{aylari}{\B{ay\cp{la}ri}} \I{[aj."la.Ri]} : \HL{aylahe}{\B{aylahe}}. 

% aylaru
\hi
\HT{aylaru}{\B{ay\cp{la}ru}} \I{[aj."la.Ru]} : \HL{aylahe}{\B{aylahe}}. 

% aylat
\hi
\HT{aylat}{\B{ay\cp{lat}}} \I{[aj."lat\textcorner]} : \HL{aylahe}{\B{aylahe}}. 

% ayleyä
\hi
\HT{ayleyA}{\B{ay\cp{le}yä}} \I{[aj."lE.j\ae]} : \HL{aylahe}{\B{aylahe}}. 

% aynga
\hi
\HT{aynga}{\B{ay\cp{nga}}} \I{[aj."Na]} \emph{pn.} (\emph{gen.} \B{ay\-\cp{nge}\-yä}) you (all) (2\textsuperscript{nd} person plural). 
\B{Le\-i\-u ay\-nga tstu\-nwi nì\-txan nì\-wotx.} You are all very kind.~\LINK{http://forum.learnnavi.org/news-announcements/happy-birthday-to-karyu-pawl/msg556899/\#msg556899}{ln} 
\B{Ay\-nga ne\-to ri\-vikx!} Step back! [lit.: move away you all]~\LINK{http://wiki.learnnavi.org/index.php?title=Na\%27vi_from_Avatar_Movie\#Scene_in_Hometree}{a\slash wiki} 
\B{Txo tsi\-ve\-'a ay\-ngal key\-eyt, ru\-txe oe\-ru pi\-veng.} If you see mistakes, please tell me.~\LINK{http://naviteri.org/2010/06/first-post/}{b} 
\B{Oel ay\-nga\-ti ka\-me\-i\-e, ma oe\-yä ey\-lan.} I \textsc{see} you all, my friends.~\LINK{http://forum.learnnavi.org/news-announcements/a-response-from-paul-frommer!/}{ln} 
\B{Sìl\-pey oe, la\-yu oe\-ru ye'\-rìn sìl\-tsan\-a fmawn a tsun oe ay\-nga\-ru ti\-vìng.} I hope, soon I will have good news that I can give to you all.~\LINK{http://forum.learnnavi.org/news-announcements/a-response-from-paul-frommer!/}{ln} 
\B{Si\-vu\-nu ay\-ngar fì\-tskxe\-keng\-tsyìp!} May you all enjoy this little excercise!~\LINK{http://naviteri.org/2010/07/ting-mikyun-fte-tslivam-listening-comprehension-1-3/}{b} 
\B{Ay\-nge\-yä tì\-ftu\-si\-a 'o' li\-vu nì\-'aw!} May your studying be fun!~\LINK{http://naviteri.org/2012/10/sound-files-added-to-previous-post/}{b} \B{Sìl\-pey oe, ay\-nga\-ri fì\-tì\-päng\-kxo\-tsyìp el\-tur tì\-txen sil\-vi.} I hope, this little discussion has been interesting for you.~\LINK{http://naviteri.org/2010/07/diminutives-conversational-expressions/}{b} 
\B{Te\-rì\-ran ay\-oe ay\-nga\-ne.} We are walking to you.~\LINK{http://wiki.learnnavi.org/index.php/Hunting_Song}{wiki} 
(\ding{43}~\HL{nga}{\B{nga}})

% ayngenga
\hi
\HT{ayngenga}{\B{aynge\cp{nga}}} \I{[aj.NE."Na]} \emph{pn.} (\emph{gen.} \B{ay\-nge\cp{nge}\-yä}) you all (2\textsuperscript{nd} person plural; deferential or ceremonial form of \B{ay\-nga}). 
\B{Ay\-nge\-nga\-ru o\-he\-yä tsmuk\-it alu Ne\-wey te Tska\-ha So\-rewn\-'i\-te.} Allow me to introduce my sister, Newey te Tskaha Sorewn'ite.~\LINK{http://naviteri.org/2010/09/getting-to-know-you-part-1/}{b}
(\ding{43}~\HL{nga}{\B{nga}})

% ayoe
\hi
\HT{ayoe}{\B{ay\cp{o}e}} \I{[aj."{}o.e]} (\emph{with case endings or enclitic adp.s always pronounced as:} \I{[aj."wE-]}) \emph{pn.} we (1\textsuperscript{st} person plural, exclusive). 
\B{Te\-rì\-ran ay\-oe ay\-nga\-ne.} We are walking to you.~\LINK{http://wiki.learnnavi.org/index.php/Hunting_Song}{wiki} 
\B{Nì\-aw\-no\-mum ngo\-lop ay\-oel fì\-nu\-mul\-txa\-ti fpi sngä\-'i\-yu.} As you know, we've created this class for beginner.~\LINK{http://naviteri.org/2012/07/tskxekengtsyip-a-mikyunfpi-a-little-listening-exercise/}{b} 
\B{ul\-te kaw\-tu ke tsun ay\-oer tì\-ftang si\-vi}, and no one can stop us.~\LINK{http://forum.learnnavi.org/language-updates/stop!/}{ln} 
(\ding{43}~\HL{oe}{\B{oe}})

% ayohe
\hi
\HT{ayohe}{\B{ay\cp{o}he}} \I{[aj."{}o.hE]} \emph{pn.} (\emph{deferential or ceremonial form of}) we (1\textsuperscript{st} person plural exclusive).  The inclusive form is \B{ohe ay\-nge\-nga\-sì}, I and you, or \B{ohe pxe\-nge\-nga\-sì}, I and (the three of) you. 
(\ding{43}~\HL{ayoe}{\B{ayoe}})

% ayoeng
\hi
\HT{ayoeng}{\B{ay\cp{oeng}}} \I{[aj."wEN]} \emph{or} \HL{awnga}{\B{aw\cp{nga}}} \I{[aw."Na]} \emph{pn.} we (1\textsuperscript{st} person plural inclusive) (\emph{with case endings, the forms are:} \B{ayoe\cp{ngal}, ayoe\cp{ngat}(i), ayoe\cp{ngar}(u), ayoe\cp{nge}yä, ayoe\cp{nga}ri} \emph{with a stress shift to the} \B{nga}\emph{-part}) \I{[aj.wE."Nal]} etc. 
\B{Ay\-oengal fte Fay\-saw\-tu\-tet ki\-vu\-rakx, ze\-ne nì\-'aw\-ve ay\-fot tsli\-vam.} We must understand these Sky People if we are to drive them out.~\LINK{http://forum.learnnavi.org/language-updates/confirmation-on-kurakx-eytukan-line-and-toitslan/}{ln} 
(\ding{43}~\HL{awnga}{\B{aw\-nga}}, \HL{oeng}{\B{oeng}})

% ayram alusìng
\hi
\HT{ayram alusIng}{\B{Ay\cp{ram} Alu\cp{sìng}}} \emph{also} \B{ayRam aLusìng} \I{[aj."Ram a.lu."{}sIN]} \emph{n.} the Floating Mountains, the Hallelujah Mountains. 
\B{Aw\-nga kel\-ku si nu\-ä ay\-ram a\-lu\-sìng.} We live beyond the flying mountains.~\LINK{http://naviteri.org/2011/10/more-vocabulary-a-bit-of-grammar/}{b} 
\B{Fwa tsyìl Ay\-ram\-it A\-lu\-sìng lu sä\-'e\-o\-i\-o a ze\-ne fra\-po si\-vi fte sl\-ivu ta\-ron\-yu.} Climbing Iknimaya is a ritual that everyone has to perform to become a hunter.~\LINK{http://naviteri.org/2013/04/wheres-the-bathroom-and-other-useful-things/}{b}
(\ding{43}~\HL{ram}{\B{ram}}, \HL{lIng}{\B{lìng}})

% ayu
\hi
\B{ay\cp{u}} \I{[aj."{}u]}~:~\HL{'u}{\B{'u}}

% ayvitrayä ramunong
\hi
\HT{ayvitrayA ramunong}{\B{Ayvit\cp{ra}yä Ra\cp{mu}nong}} \I{[aj.vit\textcorner."Ra.j\ae{} Ra."mu.noN]} \emph{n.} the Well of Souls (\emph{the closest connection to Eywa on Pandora, a point of extreme spiritual significance to the Na'vi}).
\B{Oe tsway\-on io Ay\-vit\-ra\-yä Ra\-mu\-nong.} I fly over the Tree of Souls.~\LINK{http://forum.learnnavi.org/language-updates/over-under-and-quickly/}{ln}
(\ding{43}~\HL{vitra}{\B{vit\-ra}}, \HL{ramunong}{\B{ra\-mu\-nong}})

\end{multicols}

% ===== Section Ä =====

 \pdfbookmark[0]{AE}{AE}
\begin{center}
\section*{Ä \I{[\ae]}}
\end{center}

\begin{multicols}{2}

% -ä
\hi
\HT{-A}{\B{-ä}} \I{[.\ae]} \emph{suf. to mark the genitive case of a n. ending in a consonant, pseudovowel or u, o}. (\ding{43}~\HL{-yA}{\B{-yä}})

% äie
\hi
\HT{Aie}{\B{ä\cp{i}e}} \I{[\ae."{}i.E]} \emph{n.} vision (\emph{spiritual sense}). 
\B{Fì\-ta\-ron\-yur a\-pxan te\-swo\-ti\-vìng äiet a\-ngay.} Grant this worthy warrior true vision.~\LINK{https://wiki.learnnavi.org/Na\%27vi_from_Avatar_Movie\#Becoming_one_of_the_Na.27vi}{a\textsuperscript{c}}

% Änsìt
\hi
\B{\cp{Än}sìt} \I{["{}\ae{}n.sIt\textcorner]} \emph{n. male name}.

% ‹äng›
\hi
‹\HT{Ang}{\B{äng}}› \I{[\ae.N]} \emph{or} ‹\B{eng}› \I{[E.N]} (\emph{before i, ì}) \emph{inf. in} ‹2› (to indicate negative feelings of the speaker). 
\B{Slä vay set, ke pa\-mä\-hä\-ngem kea tì\-'eyng.} But unfortunately until now, no answer arrived.~\LINK{http://forum.learnnavi.org/news-announcements/a-response-from-paul-frommer!/}{ln}. 
\B{Fì\-skxawng\-ì\-ri tsap\-'a\-lu\-te se\-ngi oe.} I apologize for this moron.~\LINK{http://wiki.learnnavi.org/index.php?title=Corpus\#Times_Online}{wiki} (\ding{214}~‹\HL{ei}{\B{ei}}›)

% äo
\hi
\HT{Ao}{\B{\cp{ä}o}} \I{["{}\ae.o]} \emph{adp.} under, below. 
\B{Pol sä\-'o\-ti wo\-lan äo ay\-rìk.} He hid his tool under the leaves.~\LINK{http://naviteri.org/2011/02/new-vocabulary-ii\%E2\%80\%94part-1/}{b} 
\B{Tsa\-txon ay\-ut\-ral\-äo krr a pol oe\-ti no\-lìn nì\-wa\-we.} That night beneath the trees when she looked at me meaningfully.~\LINK{http://naviteri.org/2011/05/some-miscellaneous-vocabulary/}{b} (\ding{214}~\HL{io}{\B{io}})

% ‹äp› ‹äpeyk›
\hi
‹\HT{Ap}{\B{äp}}› \I{[\ae.p]} \emph{inf. in} ‹pre-1› \textsc{i}. \on{1}. (to form a reflexive statement; only possible with \emph{vtr.} and certain \B{si}-verbs). \B{Oe tsä\-pe\-'a.} I see myself.~\LINK{http://forum.learnnavi.org/language-updates/the-reflexive-\%28from-a-frommer-email!!!!\%29/}{ln} \B{oe ngar mu\-wä\-pi\-vìn\-txu}, I introduce myself to you.~\LINK{http://naviteri.org/2010/09/getting-to-know-you-part-1/}{b} \B{Ke ze\-ne win sä\-pi\-vi.} Take your time; don’t rush.~\LINK{http://naviteri.org/2010/09/quick-follow-up/}{b} \on{2}. (together with \B{fìtsap} to convey `each other') \B{Zì\-sìt\-o a\-vol ke tsä\-po\-le\-'a fo fì\-tsap.} They haven't seen each other in eight years.~\LINK{http://naviteri.org/2011/12/one-more-for-2011/}{b} \\
\textsc{ii}. ‹\B{äpeyk}› \I{[\ae.pEj.k]} \emph{inf. in} ‹pre-1› (to form a causative reflexive, i.e. `cause oneself to do s.t.') \B{po tä\-pey\-kì\-ye\-ver\-ke\-i\-up nì\-näk}, I am jazzed that he is apparently about to drink himself to death.~\LINK{http://forum.learnnavi.org/language-updates/sieyng-a-ftu-naring-8-participles-causatives-and-reflexives-oh-my!/}{ln}

% ätxäle / si
\hi
\HT{AtxAle}{\B{ä\cp{txä}le}} \I{[\ae{}."t'\ae{}.lE]} \textsc{i}. \emph{n.} request. \\
\textsc{ii}. \HT{AtxAle si}{\B{ä\cp{txä}le si}} \I{[\ae{}."t'\ae{}.lE si]} \emph{vin.} request (\emph{uses} \B{tsnì} \emph{to introduce the request}). 
\B{Ä\-txä\-le si tsnì li\-vu o\-he\-ru U\-nil\-ta\-ron.} I respectfully request the Dream Hunt.~\LINK{http://wiki.learnnavi.org/index.php?title=Corpus\#Behind_the_Scenes}{wiki} 
(\emph{a}.) \B{ä\-tx\-äle si oe pi\-vawm}, may I ask. 
\B{Ä\-txä\-le su\-yi o\-he pi\-vawm, pe\-o\-lo' lu\-yu pum nge\-nge\-yä?} May I ask what tribe you belong to?~\LINK{http://naviteri.org/2010/09/getting-to-know-you-part-2/}{b} 
(\emph{b. proverb}) \B{Ä\-txä\-le si pa\-lu\-lu\-kan\-ur tsnì smar\-it li\-vo\-nu.} \emph{or} \B{Ä\-txä\-le pa\-(lu)\-lu\-kan\-ur.} \emph{a futile gesture, an attempt to achieve something that might be desirable but will clearly not happen} [lit.: Ask a thanator to release its prey.]~\LINK{http://naviteri.org/2010/07/vocabulary-update/}{b}
(\emph{c. idiom}) \B{Tì\-ka\-ra lu ä\-txä\-le pa\-lu\-kan\-ur.} Resistance is futile [lit.: `reistance is a request to the thanator'].~\LINK{http://naviteri.org/2023/09/aawa-ayliu-si-aylifyavi-amip-a-few-new-words-and-expressions/}{b} \\
\textsc{Derivations}: \ding{43} 
\on{1}.~\HL{tsantxAl}{\B{tsan\-txäl}}, invitation.

% äzan / si
\hi 
\HT{Azan}{\B{ä\cp{zan}}} \I{[\ae."zan]} \textsc{i}. \emph{n.} force, compulsion. \\
\textsc{ii}. \B{ä\cp{zan} si} \I{[\ae."zan si]} \emph{vin.} force, compel (\emph{used with} \ding{43}~\HL{tsnI}{\B{tsnì}})
\B{Fo ä\-zan so\-li oer tsnì tsa\-kem si\-vi.} They forced me to do it.~\LINK{http://naviteri.org/2020/05/aawa-liu-amip-si-vurway-alor-a-few-new-words-and-a-beautiful-poem/}{b} \\
\textsc{Derivatives}: \ding{43} 
\on{1}.~\HL{Azanluke}{\B{ä\-zan\-lu\-ke}}, voluntarily, without force.
\on{2}.~\HL{Azantu}{\B{ä\-zan\-tu}}, domineering person.

% äzanluke
\hi 
\HT{Azanluke}{\B{ä\cp{zan}luke}} \I{[\ae."zan.lu.kE]} \emph{adv.} voluntarily, without force or compulsion. 
\B{Vin oel äzanluke futa nga kivä poehu.} I request, without compulsion, that you go with her (i.e., I'm asking you to go with her, but I won't force you.)~\LINK{http://naviteri.org/2020/05/aawa-liu-amip-si-vurway-alor-a-few-new-words-and-a-beautiful-poem/}{b}
\B{Vin oel futa nga kivä poehu äzanluke.} I request that you go with her voluntarily.~\LINK{http://naviteri.org/2020/05/aawa-liu-amip-si-vurway-alor-a-few-new-words-and-a-beautiful-poem/}{b}
(\ding{43}~\HL{Azan}{\B{äzan}}, \HL{luke}{\B{lu\-ke}})

% äzantu
\hi 
\HT{Azantu}{\B{ä\cp{zan}tu}} \I{[\ae."zan.tu]} \emph{n.} domineering person; one who is bossy, authoritarian, or dictatorial. 
(\ding{43}~\HL{Azan}{\B{äzan}}, \HL{-tu}{\B{-tu}})

\end{multicols}

% ===== Section E =====

 \textsc{Derivatives}: \pdfbookmark[0]{E}{E}
\begin{center}
\section*{E \I{[E]}}
\end{center}

\begin{multicols}{2}

% e
\hi
\HT{e}{\B{e}} \I{[E]} \emph{intj.} ah, oh (i.e., `ah, I understand').~\LINK{https://forum.learnnavi.org/language-updates/navi-equivalent-for-ah/msg698116/\#msg698116}{ln}

% eampin
\hi
\HT{eampin}{\B{\cp{e}ampin}} \I{["{}E.am.pin]} \emph{n.} the color \ding{43}~\HL{ean}{\B{ean}}. 
(\ding{43}~\HL{'opin}{\B{'o\-pin}})

% ean
\hi
\HT{ean}{\B{\cp{e}an}} \I{["{}E.an]} \emph{adj.} blue, green. 
\B{Fì\-syu\-lang lu ean na ta'\-leng.} This flower is skin\hyp{}blue.~\LINK{http://naviteri.org/2010/08/mipa-ayopin-mipa-ayliu-new-colors-new-words/}{b} 
\B{Fì\-syu\-lang a\-ean\hyp{}na\hyp{}ta'\-leng lor lu nì\-txan.\slash Ta'\-leng\-na\hyp{}ean\-a fì\-syu\-lang lor lu nì\-txan.} This skin\hyp{}blue flower is very beautiful.~\LINK{http://naviteri.org/2010/08/mipa-ayopin-mipa-ayliu-new-colors-new-words/}{b} \\
\textsc{Derivatives}: \ding{43} 
\on{1}.~\HL{eampin}{\B{eam\-pin}}, (the color) blue\slash green. 
\on{2}.~\HL{rIkean}{\B{rì\-ke\-an}}, green. 
\on{3}.~\HL{ta'lengean}{\B{ta'\-le\-nge\-an}}, blue. 
\on{4}.~\HL{eanean}{\B{ean\-ean}}, cheadle.
\on{5}.~\HL{eanfwopx}{\B{ean\-fwopx}}, mist bloom.

% eanean
\hi
\HT{eanean}{\B{\cp{e}an\cp{e}an}} \I{["{}E.an."{}E.an]} \emph{n.} \ding{96} cheadle, herbaceous plant, \emph{thylakoidia spiralis}. 
(\ding{43}~\HL{ean}{\B{ean}})

% eanfwopx
\hi 
\HT{eanfwopx}{\B{\cp{e}anfwopx}} \I{["{}E.an.fwop']} \emph{n.} \ding{96} mist bloom [lit.: `blue dust cloud']. 
(\ding{43}~\HL{ean}{\B{ean}}, \HL{fwopx}{\B{fwopx}})

% ehetx / si
\hi 
\HT{ehetx}{\B{e\cp{hetx}}} \I{[E."hEt']} \textsc{i}. \emph{n.} (\emph{pl.} \B{me\-hetx}, \B{pxe\-hetx}, \B{ay\-e\-hetx}) excuse. \\
\textsc{ii}. \B{e\cp{hetx} si} \I{[E."hEt' si]} \emph{vin.} make an excuse\slash make excuses.
\B{Fu\-ria nga ke tsan\-'ul, var nga e\-hetx si\-vi nì\-'aw.} Regarding your lack of improvement, you only keep making excuses.~\LINK{https://naviteri.org/2024/02/mipa-ayliu-si-aylifyavi-niul-more-new-words-and-expressions/}{b}

% ‹ei›
\hi
‹\HT{ei}{\B{ei}}› \I{[E.i]} \emph{or} ‹\B{eiy}› \I{[E.i.j]} (\emph{before i, ì, ll, rr})  \emph{inf. in} ‹2› (to indicate positive feelings of the speaker). 
\B{Pe\-yä tì\-kang\-kem\-ìl mi vey\-krr\-e\-i\-yìn pot nì\-wotx.} I'm glad, his work is still completely overwhelming him.~\LINK{http://naviteri.org/2010/09/getting-to-know-you-part-3/}{b} 
(\emph{a. idiom}) \B{oel nga\-ti ka\-me\-i\-e}, I \textsc{see} you.~\LINK{http://www.vanityfair.com/online/oscars/2009/12/brushing-up-on-navi-the-language-of-avatar}{vf} 
(\emph{b. idiom}) \B{sre\-fe\-re\-i\-ey nì\-prr\-te'}, looking forward. \B{Tsa\-ri\-a nga\-hu ye'\-rìn ul\-txa si nì\-mun, oe sre\-fe\-re\-i\-ey nì\-prr\-te'.} I'm looking forward to getting together with you again soon.~\LINK{http://naviteri.org/2011/02/new-vocabulary-part-2/}{b} 
(\ding{214}~‹\HL{Ang}{\B{äng}}›)

% ekxan / ekxan si
\hi
\HT{ekxan}{\B{e\cp{kxan}}} \I{[E."k'an]} \textsc{i}. \emph{n.} (\emph{pl.} \B{me\-kxan}, \B{pxe\-kxan}, \B{ay\-e\-kxan}) barricade, obstruction. 
\B{Rä\-'ä fmi\-vi li\-vok fu em\-ki\-vä ay\-ekxan\-it a fkol ngo\-lop fpi sì\-kxu\-ke ay\-frr\-tuä sì ay\-ioang\-ä.} Don't try to approach or cross the barricades that were created for the sake of the visitors' and animals' safety.~\LINK{http://naviteri.org/2017/06/ayliu-a-ta-eywaeveng-kifkey-uniltirantokxa-words-from-disneys-pandora-the-world-of-avatar/}{b} \\
\textsc{ii}. \B{e\cp{kxan} si} \I{[E."k'an si]} \emph{vin.} exclude, keep out, bar.
\B{Sra\-ke fì\-kxem\-yo tsun tsay\-ioang\-ur le\-hrr\-ap ekxan si\-vi?} Can this wall keep out those dangerous animals?~\LINK{http://naviteri.org/2017/09/zisikrr-amip-ayliu-amip-new-words-for-the-new-season/}{b} \\
\textsc{Derivatives}: \ding{43} 
\on{1}.~\HL{kawkxan}{\B{kaw\-kxan}}, free, unblocked, unobstructed, clear.
\on{2}.~\HL{lekxan}{\B{le\-kxan}}, blocked, obstructed; frustrated.

% ekxtxu
\hi
\HT{ekxtxu}{\B{ekx\cp{txu}}} \I{[Ek'."t'u]} \emph{adj.} rough.   
\B{Ta'\-leng prr\-nen\-ä lu faoi, pum ko\-ak\-tu\-ä ekx\-txu.} A baby's skin is smooth, an old person's is rough.~\LINK{http://naviteri.org/2012/01/mipa-zisit-ayliu-amip-new-words-for-the-new-year/}{b} 
\B{Kll\-te lu ekx\-txu.} The ground is rough.~\LINK{http://naviteri.org/2013/04/wheres-the-bathroom-and-other-useful-things/}{b} 
(\ding{214}~\HL{faoi}{\B{faoi}})

% eltu
\hi
\HT{eltu}{\B{\cp{el}tu}} \I{["{}El.tu]} \textsc{i}. \emph{n.}~(\emph{pl.} \B{mel\-tu}, \B{pxel\-tu}, \B{ay\-el\-tu}) 
(\emph{a.}) brain. 
(\emph{b.}) \B{\cp{el}tu le\cp{fngap}} \I{["{}El.tu lE."fNap\textcorner]} computer, PC [lit.: metallic brain]. \\
\textsc{ii}. \B{\cp{el}tu si} \I{["{}El.tu si]} \emph{vin.} pay attention, quit goofing off. 
\B{El\-tu si! Tsa\-tstal a\-fwem lu litx nì\-txan.} Pay attention! That blunt knife is very sharp.~\LINK{http://naviteri.org/2012/03/spring-vocabulary-part-1/}{b} \\
\textsc{iii}. \B{\cp{el}tur tì\cp{txen} si} \I{["{}El.tuR tI."t'En si]} \emph{vin.} be interesting, intriguing.   
\B{Tsa\-'u el\-tur tì\-txen si.} That's interesting.~\LINK{http://forum.learnnavi.org/language-updates/txelanit-hivawl/}{ln} 
\B{Tsa\-vur oe\-yä el\-tur tì\-txen sol\-i.} That story intrigued me.~\LINK{http://forum.learnnavi.org/language-updates/txelanit-hivawl/}{ln} 
\B{Tsa\-vur fe\-yä ay\-el\-tur tì\-txen so\-li.} They found that story interesting.~\LINK{http://forum.learnnavi.org/language-updates/txelanit-hivawl/}{ln} \B{Pxoe\-ri tsa\-vur el\-tur tì\-txen so\-li.} The three of us were interested in that story.~\LINK{http://forum.learnnavi.org/language-updates/txelanit-hivawl/}{ln} \\
\textsc{iv}. \B{\cp{el}tut hey\cp{ka}haw} \I{["{}El.tut\textcorner hEj."ka.haw]} \emph{or} \I{[.tUt\textcorner]} \emph{ph.} be boring. 
\B{Tsa\-sä\-ftxu\-lì\-'ul pe\-yä el\-tut hey\-ko\-la\-haw nì\-txan.} That speech of his was very boring.~\LINK{https://naviteri.org/2024/02/mipa-ayliu-si-aylifyavi-niul-more-new-words-and-expressions/}{b} 
(\ding{43}~\HL{skeykx}{\B{skeykx}}) \\ 
\textsc{Derivatives}: \ding{43} 
\on{1}.~\HL{eltungawng}{\B{el\-tu\-ngawng}}, brainworm.

% eltungawng
\hi
\HT{eltungawng}{\B{\cp{el}tungawng}} \I{["{}El.tu.NawN]} \emph{n., fauna}~(\emph{pl.} \B{mel\-tu\-ngawng}, \B{pxel\-tu\-ngawng}, \B{ay\-el\-tu\-ngawng}) brainworm. 
(\ding{43}~\HL{eltu}{\B{el\-tu}}, \HL{ngawng}{\B{ngawng}})

% emkä
\hi
\HT{emkA}{\B{em\cp{kä}}} \I{[Em."k$\cdot$\ae]} \emph{vtr.}\on{22} cross (s.t.). 
\B{Rä\-'ä fmi\-vi li\-vok fu em\-ki\-vä ay\-ekxan\-it a fkol ngo\-lop fpi sì\-kxuke ay\-frr\-tuä sì ay\-ioang\-ä.} Do not attempt to approach or cross any barriers designed for Guest and animal safety.~\LINK{http://naviteri.org/2017/06/ayliu-a-ta-eywaeveng-kifkey-uniltirantokxa-words-from-disneys-pandora-the-world-of-avatar/}{b}
(\ding{43}~\HL{kA}{\B{kä}}) \\
\textsc{Derivatives}: \ding{43} 
\on{1}.~\HL{semkA}{\B{sem\-kä}}, bridge.  
\on{2}.~\HL{emkAfya}{\B{em\-kä\-fya}}, ford, crossing. 

% emkäfya
\hi 
\HT{emkAfya}{\B{em\cp{kä}fya}} \I{[Em."k\ae.fja]} \emph{n.} ford, crossing. 
\B{Fì\-tseng pay\-fya vi\-rä ka ngip a\-reng, ha tsun aw\-nga tsat si\-var sko em\-kä\-fya.} Here the stream spreads over a shallow area, so we can use it as a ford.~\LINK{http://naviteri.org/2019/06/50a-liu-amip-40-new-words/}{b} 
(\ding{43}~\HL{emkA}{\B{em\-kä}}, \HL{fya'o}{\B{fya\-'o}}; \HL{semkA}{\B{sem\-kä}})

% emrey
\hi
\HT{emrey}{\B{em\cp{rey}}} \I{[Em."R$\cdot$Ej]} \emph{vin., inf.}\on{22} survive (a life-threatening situation). 
\B{Tsawl\-a pa\-lu\-lu\-kan\-ìl oe\-ti 'ìl\-me\-ko a krr, Nawm\-a Sa'\-nok lrr\-tok sei\-yi ul\-te oe em\-ro\-ley.} The Great Mother was looking after me when the big thanator attacked and I've survived.~\LINK{http://naviteri.org/2011/05/some-miscellaneous-vocabulary/}{b}
(\ding{43}~\HL{rey}{\B{rey}}) \\
\textsc{Derivatives}: \ding{43} 
\on{1}.~\HL{lemrey}{\B{lem\-rey}}, surviving. 
\on{2}.~\HL{temrey}{\B{tem\-rey}}, survival. 
\on{3}.~\HL{nemrey}{\B{nem\-rey}}, as if one's life were at stake.

% emza'u
\hi
\HT{emza'u}{\B{em\cp{za}'u}} \I{[Em."z$\cdot$a.P$\cdot$u]} \emph{vtr., inf.}\on{23} pass (a test). \B{Sì\-fme\-tok\-it em\-zo\-la\-'u ohel.} I have passed the tests.~\LINK{http://wiki.learnnavi.org/index.php?title=Corpus\#Behind_the_Scenes}{wiki} 
\B{Nän ftia, nän lu skxom a em\-za\-'u.} The less you study, the less chance you have of passing.~\LINK{http://naviteri.org/2012/02/trr-asawnung-lefpom-happy-leap-day/}{b}
(\ding{43}~\HL{za'u}{\B{za\-'u}})

% Entu
\hi
\B{\cp{En}tu} \I{["{}En.tu]} \emph{n. male name}. 
\B{Ma En\-tu, fì\-kem rä\-'ä si\-vi; lu ke\-fpom\-ro\-nga'.} Entu, don't do this; it's not healthy (mentally).~\LINK{http://naviteri.org/2014/10/tson-si-fpomron-obligation-and-mental-health/}{b}
\B{En\-tul kan\-'ìn tì\-wu\-sem\-it.} Entu specializes in fighting.~\LINK{http://naviteri.org/2010/09/getting-to-know-you-part-1/}{b}

% ‹eng›
\hi
‹\HT{eng}{\B{eng}}›~:~‹\HL{Ang}{\B{äng}}›

% eo
\hi
\HT{eo}{\B{\cp{e}o}} \I{["{}E.o]} \emph{adp.} before, in front of.
\B{Eo ay\-oeng lu txan\-a tì\-kawng.} A great evil is upon us.~\LINK{http://wiki.learnnavi.org/index.php?title=Corpus\#Behind_the_Scenes}{wiki} 
\B{Eo Ey\-wa oe \mbox{'ia}.} I lose myself before Eywa.~\LINK{http://forum.learnnavi.org/language-updates/sieyng-a-ftu-naring-3-proverbs-what-we-asked-for-and-some-vocabulary/msg327513/\#msg327513}{a\slash ln} 
\B{Ut\-ral a lu few pay\-fya a eo kel\-ku oeyä tsawl lu nì\-txan.} The tree on the other side of the stream in front of my house is very tall.~\LINK{http://naviteri.org/2011/01/new-year-new-vocabulary/}{b} 
(\emph{a. idiom}) \B{tokx eo tokx}, face to face, in person.~\LINK{http://forum.learnnavi.org/language-updates/fmawno/msg293595/\#msg293595}{ln} 
(\emph{b. idiom}) \B{na fkxen eo fkio}, to go to waste [lit.: `like vegetable food before a tetrapteron'] 
\B{Nge\-yä fì\-tì\-kang\-kem\-vi a\-txan\-tsan ke sla\-yu na fkxen eo fkio.} This excellent work of yours will not go to waste.~\LINK{https://naviteri.org/2024/12/odds-and-ends/}{b}
(\ding{214}~\HL{uo}{\B{uo}}) \\
\textsc{Derivatives}: \ding{43} 
\on{1}.~\HL{eolI'uvi}{\B{eo\-lì\-'u\-vi}}, prefix.

% eolì'uvi
\hi 
\HT{eolI'uvi}{\B{\cp{e}olì'uvi}} \I{["{}E.o.lI.Pu.vi]} \emph{n.} prefix. 
\B{Lu tsa\-lì\-'ur alu fì\-haw\-re'\-ti me\-lì\-'u\-vi alu 'aw\-a eo\-lì\-'u\-vi sì 'aw\-a uo\-lì\-'u\-vi.} The word \emph{fì\-haw\-re'\-ti} has two affixes--one prefix and one suffix.~\LINK{http://naviteri.org/2017/10/more-language-for-talking-about-language/}{b}
(\ding{43}~\HL{eo}{\B{eo}}, \HL{lI'uvi}{\B{lì\-'uvi}})

% ep'ang
\hi
\HT{ep'ang}{\B{ep\cp{'ang}}} \I{[Ep\textcorner."PaN]} \emph{adj.} complex. \B{Ul\-te ke stì\-yawm ay\-ngal fì\-trr kea ho\-ren\-it a\-ep\-'ang.} And you will not hear any complex rules today.~\LINK{http://naviteri.org/2012/07/tskxekengtsyip-a-mikyunfpi-a-little-listening-exercise/}{b} 
(\ding{214}~\HL{fyin}{\B{fyin}})\\
\textsc{Derivatives}: \ding{43} 
\on{1}. \HL{fyinep'ang}{\B{fyi\-nep\-'ang}}, complexity.

% epxang
\hi
\HT{epxang}{\B{e\cp{pxang}}} \I{[E."p'aN]} \emph{n.}~(\emph{pl.} \B{me\-pxang}, \B{pxe\-pxang}, \B{ay\-e\-pxang}) stone jar (\emph{to hold small toxic arachnid}).
\B{Epxang\-mì lu 'u\-pe?--Tì\-lo\-ho.} What's in the stone jar?--It's a surprise.~\LINK{http://http://naviteri.org/2019/06/50a-liu-amip-40-new-words}{b}

% ‹er›
\hi
‹\HT{er}{\B{er}}› \I{[E.R]} \emph{inf. in} ‹1› (to form imperfective aspect, on-going action). \B{Ney\-ti\-ri he\-ra\-haw.} Neytiri is sleeping.~\LINK{http://wiki.learnnavi.org/Corpus\#MSNBC}{wiki} \on{1}. (required with \B{tengkrr}) \B{aw\-ngal yom wu\-tsot teng\-krr he\-reyn nì\-pxim}, we eat dinner while sitting upright.~\LINK{http://whyisthisnight.com/addl-more.html\#na\%27vi}{tn} \on{2}. (can combine with tense infixes to form a frame narrative) \B{Oel hu Txe\-wì trr\-am na'\-rìng\-it tar\-mok, tso\-le'a syep\-tu\-tet a\-tsawl fra\-to mì sì\-rey.} Yesterday I was with Txewì in the forest and we saw the biggest Trapper I've ever seen.~\LINK{http://wiki.learnnavi.org/index.php?title=Corpus\#New_York_Times}{wiki} 
(\ding{43}~‹\HL{arm}{\B{arm}}›, ‹\HL{ary}{\B{ary}}›, ‹\HL{Irm}{\B{ìrm}}›, ‹\HL{Iry}{\B{ìry}}›; \ding{214}~‹\HL{ol}{\B{ol}}›)

% etrìp
\hi
\HT{etrIp}{\B{\cp{et}rìp}} \I{["{}Et\textcorner.RIp\textcorner]} \emph{adj.} favorable, auspicious. 
(\emph{a}) (idiom) 
\B{Et\-rìp\-a syay\-vi!} Good luck!~\LINK{http://naviteri.org/2010/07/vocabulary-update/}{b} 
(\ding{214}~\HL{ketrIp}{\B{ket\-rìp}}) \\
\textsc{Derivatives}: \ding{43} 
\on{1}.~\HL{netrIp}{\B{net\-rìp}}, luckily, happily.
\on{2}.~\HL{ketrIp}{\B{ket\-rìp}}, unfavorable, inauspicious.

% Europa
\hi
\B{Eu\cp{ro}pa} \I{[E.u."Ro.pa]} \emph{n.}, \textsc{lw} Europe. 

% eyawr
\hi
\HT{eyawr}{\B{e\cp{yawr}}} \I{[E."jawR]} \emph{adj.} correct, right. 
\B{Sä\-pll\-txe\-vi\-ri a\-la\-he, lu lì'\-fya\-vi nge\-yä e\-yawr, lu pum oe\-yä e\-yawr nì\-teng.} Concerning the other comment, your expression is correct, mine is correct as well.~\LINK{http://naviteri.org/2013/01/awvea-posti-zisita-amip-first-post-of-the-new-year/comment-page-1/\#comment-1903}{b} 
\B{Slä fma\-syi ri\-vun e\-yawr\-a tì\-'eyng\-it.} But I'll try to find the correct answer.~\LINK{http://naviteri.org/2011/08/new-vocabulary-clothing/comment-page-1/\#comment-920}{b} 
(\ding{214}~\HL{keyawr}{\B{ke\-yawr}}) \\
\textsc{Derivatives}: \ding{43} 
\on{1}.~\HL{tIyawr}{\B{tì\-yawr}}, correctness. 
\on{2}.~\HL{nIyawr}{\B{nì\-yawr}}, correctly. 
\on{3}.~\HL{keyawr}{\B{ke\-yawr}}, incorrect. 
\on{4}.~\HL{eyawrfya}{\B{e\-yawr\-fya}}, righ way, correct path.

% eyawrfya
\hi 
\HT{eyawrfya}{\B{e\cp{yawr}fya}} \I{[E."jawR.fja]} \emph{n.} right way (of doing s.t.), correct path. 
\B{Ney\-ti\-ril Tsyeyk\-ur ko\-lar eyawr\-fyat a fyep tsko\-ti.} Neytiri taught Jake the right way to hold a bow.~\LINK{http://naviteri.org/2016/06/mrrvola-lifyavi-amip-forty-new-expressions/}{b} 
\B{Eyawr\-fya\-ri ze\-ne tsli\-vam fya\-'ot a mìn ki\-fkey.} To know the right way, you have to understand how the world turns.~\LINK{http://naviteri.org/2016/06/mrrvola-lifyavi-amip-forty-new-expressions/}{b} 
(\ding{43}~\HL{eyawr}{\B{eyawr}}, \HL{fya'o}{\B{fya\-'o}})

% eyaye
\hi
\B{e\cp{ya}ye} \I{[E."ja.jE]} \emph{n.} \ding{96}~(\emph{pl.} \B{me\-ya\-ye}, \B{pxe\-ya\-ye}, \B{ay\-e\-ya\-ye}) warbonnet, \emph{bellicum pennatum}.  

\end{multicols}

% ===== Section EW =====

 \pdfbookmark[0]{EW}{EW}
\begin{center}
\section*{EW \I{[Ew]}}
\end{center}

\begin{multicols}{2}

% ewktswo
\hi
\HT{ewktswo}{\B{\cp{ewk}tswo}} \I{["{}Ewk\textcorner.\t{ts}wo]} \emph{n.} (sense of) taste, flavor.
(\ding{43}~\HL{ewku}{\B{ew\-ku}}, \HL{tsu'o}{\B{tsu\-'o}})

% ewku
\hi
\HT{ewku}{\B{\cp{ew}ku}} \I{["{}Ew.ku]} \emph{vtr.} (\texttt{-}\emph{control}) taste. \B{Fì\-na\-er\-ìri ngal ew\-ku 'u\-ot a\-stxong srak?} Do you taste something strange in this drink?~\LINK{http://naviteri.org/2012/11/renu-ayinanfyaya-the-senses-paradigm/}{b}
(\ding{43}~\HL{may'}{\B{may'}}, \HL{tIng ftxI}{\B{tìng ftxì}}) \\
\textsc{Derivations:} \ding{43} \on{1}.~\HL{ewktswo}{\B{ewktswo}}, (sense of) taste, flavor.

% ewro
\hi
\B{\cp{ew}ro} \I{["{}Ew.Ro]} \emph{n.}, \textsc{lw} euro (currency).
\B{Kì\-ma\-nom oel mip\-a el\-tut le\-fngap yoa ew\-ro.} I just bought a new computer.~\LINK{http://naviteri.org/2014/02/barter-and-exchange/}{b}

\end{multicols}

% ===== Section EY =====

 \pdfbookmark[0]{EY}{EY}
\begin{center}
\section*{EY \I{[Ej]}}
\end{center}

\begin{multicols}{2}

% eyk
\hi
\HT{eyk}{\B{eyk}} \I{[Ejk\textcorner]} \emph{vtr.} lead. \B{Tsam\-po\-ngut Tsu'\-tey\-ìl i\-veyk.} Tsu'tey will lead the war party.~\LINK{http://wiki.learnnavi.org/index.php?title=Na\%27vi_from_Avatar_Movie\#Preparing_the_war_party}{a\slash wiki} \B{Eyk Ka\-mun a fra\-lo lä\-ngu tsa\-sä\-ta\-ron vel\-ke nì\-wotx.} Every time Kamun is in charge, the hunt is a mess.~\LINK{http://naviteri.org/2012/10/mipa-vospxi-mipa-ayliu-new-words-for-the-new-month/}{b} \\
\textsc{Derivatives}: \ding{43} \on{1}.~\HL{eyktan}{\B{eyktan}}, leader. 
\on{2}.~\HL{eykyu}{\B{eykyu}}, leader (of a small group).

% ‹eyk›
\hi
‹\HT{eyk2}{\B{eyk}}› \I{[Ej.k]} \emph{inf. in} pre-‹1› (\emph{a. to make a causative statement, to make s.o. do s.t.}) 
\B{Ey\-tu\-kan\-ìl Ney\-ti\-rir ye\-rik\-it tey\-ko\-la\-ron.} Eytukan made Neytiri hunt a hexapede.~\LINK{http://wiki.learnnavi.org/index.php?title=Canon\#More_extracts_from_various_emails}{wiki} 
(\emph{b. alternatively with} \HL{fa}{\B{fa}} \emph{to introduce the caused person}) \B{Ey\-tu\-kan\-ìl fa Ney\-ti\-ri ye\-ri\-kit tey\-ko\-la\-ron.} Eytukan had a hexapede hunted by Neytiri.~\LINK{http://wiki.learnnavi.org/index.php?title=Canon\#Extracts_from_various_emails}{wiki}
(\emph{c. to make a \emph{vin.} into a \emph{vtr.}}) \B{A\-pxa tì\-väng\-ìl po\-ti stey\-ko\-li.} (His) great thirst made him angry [lit.: `great thirst caused him to be angry'].~\LINK{http://naviteri.org/2011/01/new-year-new-vocabulary/}{b} 
\B{oe\-ri U\-nil\-tì\-ran\-tokx\-ìl tì\-rey\-ti ley\-ko\-la\-te\-i\-em nì\-txan}, \emph{Avatar} has greatly changed my life for the better.~\LINK{http://forum.learnnavi.org/language-updates/pseudo-canon-from-bd-live-content/msg350756/\#msg350756}{ln} 
\B{Po\-el tì\-kang\-kem\-it sngey\-ko\-lä\-'i.} She began the work.~\LINK{http://naviteri.org/2010/06/first-post/}{b}
(\emph{d. with vtr. modals}) \B{Fì\-pam\-tseo\-l oe\-ru ney\-kew fu\-ta srew.} This music makes me want to dance.~\LINK{https://naviteri.org/2020/12/mrra-tipangkxotsyip-five-little-discussions/}{b}

% eykil
\hi
\HT{eykil}{\B{ey\cp{kil}}} \I{[Ej."k$\cdot$il]} \emph{vtr., inf.}\on{22} bend s.t., pulling things together.
\B{Me\-tew\-it fì\-swek\-ä ey\-ki\-vil.} Pull the two ends of this bar together.~\LINK{http://naviteri.org/2013/04/wheres-the-bathroom-and-other-useful-things/}{b} 
(\ding{43}~\HL{il}{\B{il}})

% eyktan
\hi
\HT{eyktan}{\B{\cp{eyk}tan}} \I{["{}Ejk\textcorner.tan]} \emph{n.}~(\emph{pl.} \B{meyk\-tan}, \B{pxeyk\-tan}, \B{ay\-eyk\-tan}) leader. \B{Na'viru set lu nawma eyktan amip a larmu Tawtute.} The Na'vi have a great leader now who was a Skyperson.~\LINK{http://naviteri.org/2010/07/ting-mikyun-fte-tslivam-listening-comprehension-1-3/}{b} \B{Nì\-lun ay\-ioi a\-'eoio ay\-eyk\-tan\-ä lu lor fra\-to.} Of course the ceremonial wardrobes of the leaders are the most beautiful.~\LINK{http://naviteri.org/2011/08/new-vocabulary-clothing/}{b} 
(\ding{43}~\HL{eyk}{\B{eyk}}) \\
\textsc{Derivatives}: \ding{43} \on{1}.~\HL{tIeyktan}{\B{tìeyktan}}, leadership. 
\on{2}.~\HL{eyktanay}{\B{eyktanay}}, deputy, general.

% eyktanay
\hi
\HT{eyktanay}{\B{eykta\cp{nay}}} \I{[Ejk\textcorner.ta."naj]} \emph{n.}~(\emph{pl.} \B{meyk\-ta\-nay}, \B{pxeyk\-ta\-nay}, \B{ay\-eyk\-ta\-nay}) deputy, general (s.o.~who is one step down in rank from leader). 
\B{Ul\-te ftxo\-ley pol ay\-eyk\-ta\-nay\-ti a lu stum nì\-ftxan kawng na po.} And he has chosen deputies who are almost as bad as he is.~\LINK{http://naviteri.org/2016/12/tengkrr-perahem-zisit-amip-as-the-new-year-arrives/}{b}
(\ding{43}~\HL{eyktan}{\B{eyk\-tan}})

% eykyu
\hi
\HT{eykyu}{\B{\cp{eyk}yu}} \I{["{}Ejk\textcorner.ju]} \emph{n.}~(\emph{pl.} \B{meyk\-yu}, \B{pxeyk\-yu}, \B{ay\-eyk\-yu}) leader (typically of a small group). 
\B{Vll eyk\-yu ne\-fä fte po\-ngu fä\-ki\-vä.} The leader signals the party to ascend by pointing upwards.~\LINK{http://naviteri.org/2012/06/spring-vocabulary-part-3/}{b} 
\B{Mì sa\-ngek a sä\-vll\-it ngo\-lop eyk\-yul tar\-po\-ngu\-ä.}
The sign on the tree trunk was made by the leader of the hunting party.~\LINK{http://naviteri.org/2017/12/zisit-amip-lefpom-ma-eylan-happy-new-year-friends/}{b}
(\ding{43}~\HL{eyk}{\B{eyk}})

% Eytukan
\hi
\B{\cp{Ey}tukan} \I{["{}Ej.tu.kan]} \emph{n. male name}. 
\B{Pol\-txe Ey\-tu\-kan san oe ka\-yä sìk, slä oel pot ke spaw.} Eytukan said he would come but I don't believe him.~\LINK{http://forum.learnnavi.org/language-updates/verb-for-write/msg175491/\#msg175491}{ln}
\B{Spaw oel fu\-ta Mo\-'at\-ìl tso\-le\-'a Ney\-ti\-rit. -- Ke\-he. Tso\-le\-'a Ney\-ti\-rit Ey\-tu\-kan\-ìl.} I believe Mo\-'at saw Ney\-ti\-ri. -- No. It's Ey\-tu\-kan who saw Ney\-ti\-ri.~\LINK{http://naviteri.org/2011/03/word-order-and-case-marking-with-modals/}{b}
\B{New Va'\-ru tskot Ey\-tu\-kan\-ä za\-sri\-vìn.} Va'ru wants to borrow Eytukan's bow.~\LINK{http://naviteri.org/2012/06/spring-vocabulary-part-3/}{b}

% Eywa
\hi
\HT{eywa}{\B{\cp{Ey}wa}} \I{["{}Ej.wa]} \emph{n.} Eywa, `the Great Mother' (the deity of the Na'vi who protects the balance of life). 
\B{Ey\-wal zey\-ki\-vo ngat nì\-win.} May Eywa heal you quickly.~\LINK{http://naviteri.org/2010/07/vocabulary-update/}{b}
\B{Spaw Na'\-vil fu\-ta tì\-yawn Ey\-wa\-yä lu txew\-lu\-ke.} The Na'vi believe that Eywa's love is boundless.~\LINK{http://naviteri.org/2014/11/twenty-before-the-holidays/}{b} 
(\emph{a. idiom}) \B{Ey\-wa nga\-hu} Goodbye! [lit.: Eywa (be) with you!]~\LINK{http://forum.learnnavi.org/news-announcements/a-response-from-paul-frommer!/}{ln} 
\B{Ey\-wa va\-yun.} Ey\-wa will provide.~\LINK{http://naviteri.org/2020/05/aawa-liu-amip-si-vurway-alor-a-few-new-words-and-a-beautiful-poem/}{b} \\
\textsc{Derivatives}: \ding{43} \on{1}. \HL{eywa'eveng}{\B{Eywa'eveng}}, (the moon) Pandora.

% Eywa'eveng
\hi
\HT{eywa'eveng}{\B{Eywa\cp{'e}veng}} \I{[""Ej.wa."PE.vEN]} \emph{n.} (the moon) Pandora [lit.: child of Eywa]. \B{Pol sla'\-tsu ay\-i\-o\-ang\-it a tse\-'a fkol mì Ey\-wa\-'e\-veng.} She describes animals seen on Pandora.~\LINK{http://naviteri.org/2011/01/new-year-new-vocabulary/}{b} 
(\ding{43}~\HL{eywa}{\B{Ey\-wa}}, \HL{'eveng}{\B{'e\-veng}}, \HL{eyweveng}{\B{Ey\-we\-veng}})

% Eyweveng
\hi
\HT{eyweveng}{\B{Ey\cp{we}veng}} \I{[Ej."wE.vEN]} \emph{n.}~(\emph{short form of}) \ding{43}~\HL{eywa'eveng}{\B{Ey\-wa\-'e\-veng}}. 
\B{Oel ron\-srel\-ngop fu\-ta Ey\-we\-veng\-it tok.} I imagine that I'm on Pandora.~\LINK{http://naviteri.org/2011/02/new-vocabulary-part-2/}{b} 
\B{Txon\-krr lu syu\-ra\-tan na'\-rìng\-ä Ey\-we\-veng\-ä lor nì\-txan.} At night, the bioluminescence of the Pandoran forest is very beautiful.~\LINK{http://naviteri.org/2012/03/spring-vocabulary-part-1/}{b} 
(\ding{43}~\HL{eywa}{\B{Ey\-wa}}, \HL{'eveng}{\B{'e\-veng}})

\end{multicols}

% ===== Section F =====

 \pdfbookmark[0]{F}{F}
\begin{center}
\section*{F \I{[f]} (\emph{Fä})}
\end{center}

\begin{multicols}{2}

% fa
\hi
\HT{fa}{\B{fa}} \I{[fa]} \emph{adp.} with, by means of. 
\B{Ay\-lì\-'u\-fa aw\-nge\-yä 'ey\-lan\-ä a\-'e\-wan.} In the words of our young friend.~\LINK{http://forum.learnnavi.org/news-announcements/a-response-from-paul-frommer!/}{ln} 
\B{O\-lo\-me\-i\-um oel fa key\-rel fu\-ta lu yaw\-ne oe po\-ru.} I knew by the expression on her face that she loves me.~\LINK{http://naviteri.org/2011/05/some-miscellaneous-vocabulary/}{b}
\B{Tsun fko tsa\-tse\-nge\-ne ki\-vä nì\-sle\-le fu fa fwa ik\-ran\-it mak\-to nì\-'aw.} You can only get there by swimming or riding an ikran.~\LINK{http://naviteri.org/2011/02/new-vocabulary-ii\%E2\%80\%94part-1/}{b}
(\emph{a. with certain verbs}) \ding{43}~\HL{ioi}{\B{ioi s(äp)i fa}}, adorn (oneself) with; \HL{pamtseo}{\B{pam\-tseo si fa}}, play an instrument; \HL{rawn}{\B{rawn fa}}, replace s.t.~with s.t.~else; \HL{tswayon}{\B{tsway\-on fa}}, fly with s.t.; \HL{pxek}{\B{pxek ki\-nam\-fa\slash pxun\-fa}}, kick, shove. \HL{rumtsyokx}{\B{ta\-kuk fa rum\-tsyokx}}, punch s.o.; \HL{rIktsyokx}{\B{ta\-kuk fa rìk\-tsyokx}}, slap s.o.
(\emph{b. alternative to inroduce the caused person in causative sentences}) \B{Ey\-tu\-kan\-ìl fa Ney\-ti\-ri ye\-ri\-kit tey\-ko\-la\-ron.} Eytukan had a hexapede hunted by Neytiri.~\LINK{http://wiki.learnnavi.org/index.php?title=Canon\#Extracts_from_various_emails}{wiki}
(\emph{c. proverb}) \B{Tì\-tok slu tìk\-tok fa pam\-tsyìp a\-'aw.} Presence becomes absence by means of one little sound. [i.e. `Small things can make a big difference.']~\LINK{https://naviteri.org/2024/06/ayhapxi-tokxa-si-ayliu-alahe-parts-of-the-body-and-more/}{b}

% fay+
\hi
\HT{fay}{\B{fay}}\texttt{+} \I{[faj.]} \emph{prod.~pref.~on nouns} (\emph{short form of} \HL{fI}{\B{fì}}\texttt{-} \emph{and} \HL{ay}{\B{ay}}\texttt{+}) these (\dots~here) (\emph{proximal plural demonstrative prefix}). 
\B{fay\-vrr\-tep fì\-tse\-nge lu kxa\-nì}, these demons are forbidden here.~\LINK{http://wiki.learnnavi.org/index.php?title=Corpus\#Times_Online}{wiki} 
\B{Sra\-ke tsun oe fay\-u\-pxa\-ret tsl\-ivam nì\-ftue?} Can I understand these messages easily?~\LINK{http://forum.learnnavi.org/language-updates/just-a-few-interesting-little-words/}{ln}
(\ding{43} \HL{tsay}{\B{tsay}\texttt{+}})

% fayluta
\hi
\HT{fayluta}{\B{fay\cp{lu}ta}} \I{[faj."lu.ta]} \emph{conj.} (\emph{short form of} \B{fay\-lì\-'ut a}) these words that\dots (to form reported speech) \B{Pol\-txe pol fay\-lu\-ta oe new ki\-vä.} She said, `I want to go'. \emph{or} She said she wanted to go.~\LINK{http://naviteri.org/2011/08/reported-speech-reported-questions/}{b} 

% fäkä
\hi
\HT{fAkA}{\B{fä\cp{kä}}} \I{[f\ae."k$\cdot$\ae]} \emph{vin., inf.}\on{22} go up, ascend. \B{Vll eyk\-yu ne\-fä fte po\-ngu fä\-ki\-vä.} The leader signals the party to ascend by pointing upwards.~\LINK{http://naviteri.org/2012/06/spring-vocabulary-part-3/}{b}
(\ding{43}~\HL{fApa}{\B{fä\-pa}}, \HL{kA}{\B{kä}}; \ding{214}~\HL{kllkA}{\B{kll\-kä}})

% fäpa
\hi
\HT{fApa}{\B{\cp{fä}pa}} \I{["f\ae.pa]} \emph{n.} top. \B{ro fäpa}, at the top. \LINK{http://forum.learnnavi.org/language-updates/all-adpositions/msg137792/\#msg137792}{ln}
(\ding{43}~\HL{pa'o}{\B{pa\-'o}}; \ding{214}~\HL{kllpa}{\B{kll\-pa}}) \\
\textsc{Derivations:} \ding{43} \on{1}.~\HL{fApyo}{\B{fäpyo}}, roof.

% fäpyo
\hi 
\HT{fApyo}{\B{\cp{fäp}yo}} \I{["f\ae{}p\textcorner.jo]} \emph{n.} roof. 
\B{Tom\-pa ze\-rup ul\-te pay lipx kxam\-lä fäp\-yo.} It's raining and water is dripping through the roof.~\LINK{http://naviteri.org/2020/04/aawa-u-amip-a-few-new-things/}{b} 
(\ding{43}~\HL{fApa}{\B{fäpa}}, \HL{yo}{\B{yo}})

% fäza'u
\hi
\HT{fAza'u}{\B{fä\cp{za}'u}} \I{[f\ae."z$\cdot$a.P$\cdot$u]} \emph{vin., inf.}\on{23} come up, ascend. (\emph{a.})  \B{Fä\-zi\-va\-'u ne tse\-nge a oel tok!} Come up to where I am!~\LINK{http://naviteri.org/2012/01/more-additions-to-the-lexicon/}{b} 
(\emph{b. of astronomical bodies rising}) \B{Fä\-za'u tsaw\-ke krr\-pe?} When will the sun come up?~\LINK{http://naviteri.org/2012/01/more-additions-to-the-lexicon/}{b} 
(\ding{43}~\HL{fApa}{\B{fä\-pa}}, \HL{za'u}{\B{za\-'u}}; \HL{fpxAkIm}{\B{fpxä\-kìm}}; \ding{214}~\HL{kllkA}{\B{kll\-kä}}, \HL{kllza'u}{\B{kll\-za\-'u}})

% fahew
\hi
\HT{fahew}{\B{fa\cp{hew}}} \I{[fa."hEw]} \emph{n.} (\emph{a. sensation}) smell. 
\B{Oe\-ri ta pe\-yä fa\-hew a\-ke\-wong on\-tu te\-ya lä\-ngu.} My nose is full of his alien smell.~\LINK{http://languagelog.ldc.upenn.edu/nll/?p=1977}{ll\slash a} 
\B{Fì\-syu\-lang\-it syam. Fa\-hew lor lu nì\-txan, ke\-fyak?} Smell this flower. Its fragrance is beautiful, isn't it?~\LINK{http://naviteri.org/2012/11/renu-ayinanfyaya-the-senses-paradigm/}{b} 
(\emph{b. to form the middle voice}) \B{Ra\-lu\-ri fa\-hew fkan oe\-ru na ye\-rik.} Ralu smells like a hexapede to me.~\LINK{http://naviteri.org/2012/11/renu-ayinanfyaya-the-senses-paradigm/}{b}

% faoi
\hi
\HT{faoi}{\B{\cp{fa}oi}} \I{["fa.o.i]} \emph{adj.} smooth. 
\B{Ta'\-leng prr\-nen\-ä lu faoi, pum ko\-ak\-tu\-ä ekx\-txu.} A baby's skin is smooth, an old person's is rough.~\LINK{http://naviteri.org/2012/01/mipa-zisit-ayliu-amip-new-words-for-the-new-year/}{b} 
(\ding{214}~\HL{ekxtxu}{\B{ekx\-txu}})

% fe'
\hi
\HT{fe'}{\B{fe'}} \I{[fEP]} \emph{adj., nfp.} bad. (\emph{a.}) \B{Peyä tsa\-tì\-pe\-'un a swey\-lu txo wi\-vem ayoeng Oma\-ti\-ka\-ya\-wä lu fe'.} His decision to fight against the Omaticaya was a bad one.~\LINK{http://naviteri.org/2011/07/txantsana-ultxa-mi-siatll-great-meeting-in-seattle/}{b} \B{tì\-len afe'}, bad event, something bad that happened.~\LINK{http://naviteri.org/2011/07/txantsana-ultxa-mi-siatll-great-meeting-in-seattle/}{b} 
(\emph{b.~proverb}) \B{Kem a\-muiä, kum a\-fe'.} Proper action, bad result. (\emph{something that should have turned out well didn't}).~\LINK{http://naviteri.org/2012/06/spring-vocabulary-part-3/}{b} 
(\ding{214}~\HL{sIltsan}{\B{sìl\-tsan}}) \\
\textsc{Derivations}: \ding{43} \on{1}.~\HL{nIfe'}{\B{nì\-fe'}}, badly, in a bad way. 
\on{2}.~\HL{fe'ul}{\B{fe\-'ul}}, worsen, get worse.
\on{3}.~\HL{fekem}{\B{fe\-kem}}, accident.
\on{4}.~\HL{fe'ran}{\B{fe'\-ran}}, flawed nature.
\on{5}.~\HL{fe'lup}{\B{fe'\-lup}}, tacky, in poor taste. 
\on{6}.~\HL{fekum}{\B{fe\-kum}}, disadvantage, drawback, downside.
\on{7}.~\HL{fe'pey}{\B{fe'\-pey}}, dread, expect something bad to happen.

% fe'lup
\hi
\HT{fe'lup}{\B{\cp{fe'}lup}} \I{["fEP.lup\textcorner]} \emph{or} \I{[.lUp\textcorner]} (\textsc{rn}: \B{\cp{fe'}lùp} \I{["fEP.lUp\textcorner]}) \emph{adj.} tacky, in poor taste.
(\ding{43}~\HL{fe'}{\B{fe'}}, \HL{lupra}{\B{lup\-ra}})

% fe'pey
\hi 
\HT{fe'pey}{\B{fe'\cp{pey}}} \I{[fE."p$\cdot$Ej]} \emph{vin., inf.}\on{22} dread, expect something bad to happen. (\emph{used with} \ding{43} \HL{tsnI}{\B{tsnì}})
\B{Krr\-a pä\-hem Saw\-tu\-te, pxay\-a Na'\-vi fe'\-par\-mey.} When the Sky People arrived, many Na'vi felt dread.
\B{Po fe'\-po\-ley tsnì 'i\-tan sne\-yä tì\-fme\-tok\-it ke em\-zì\-ye\-va\-'u.} He feared his son might not pass the test.~\LINK{http://naviteri.org/2019/06/50a-liu-amip-40-new-words/}{b}
(\ding{43}~\HL{fe'}{\B{fe'}}, \HL{pey}{\B{pey}}; \ding{214}~\HL{sIlpey}{\B{sìl\-pey}})

% fe'ran
\hi
\HT{fe'ran}{\B{\cp{fe'}ran}} \I{["fEP.Ran]} \emph{n.} flawed nature, something ill-conceived or inherently defective. \B{Fì\-tì\-hawl\-ì\-ri fe'\-ran law lä\-ngu fra\-por.} Unfortunately the flawed nature of this plan is obvious to everyone.~\LINK{http://naviteri.org/2012/10/mipa-vospxi-mipa-ayliu-new-words-for-the-new-month/}{b} 
(\ding{43}~\HL{fe'}{\B{fe'}}, \HL{ran}{\B{ran}}; \ding{214}~\HL{loran}{\B{lo\-ran}}) \\
\textsc{Derivations:} \ding{43} \on{1}. \HL{fe'ranvi}{\B{fe'ranvi}}, blemish, flawed feature.

% fe'ranvi
\hi
\HT{}{\B{\cp{fe'}ranvi}} \I{["fEP.Ran.vi]} \emph{n.} blemish, deformity, stain, flawed feature.  \B{Hu\-fwa lu fil\-ur Va'\-ru\-ä fne\-fe'\-ran\-vi, tsal\-su\-ngay fpìl fu\-ta say\-rìp lu nì\-txan.} Although Va'ru's facial stripes are rather uneven, I still think he's very handsome.~\LINK{http://naviteri.org/2012/10/mipa-vospxi-mipa-ayliu-new-words-for-the-new-month/}{b}
(\ding{43}~\HL{fe'ran}{\B{fe'\-ran}})

% fe'ul
\hi
\HT{fe'ul}{\B{\cp{fe}'ul}} \I{["fE.P$\cdot$ul]} \emph{vin., inf.}\on{22} worsen, get worse. \B{Ke tsun oe tsli\-vam teyng\-ta tì\-ru\-sol peyä lum\-pe fe\-'ul krra oe tìng mik\-yun.} I can't understand why her singing gets worse when I listen.~\LINK{http://naviteri.org/2013/03/tsanerul-feerul-getting-better-getting-worse/}{b}
\B{Nge\-yä tsay\-lì\-'ul tì\-fkey\-tok\-it fe\-'ey\-ko\-lul nì\-'aw.} Those words of yours have only made the situation worse.~\LINK{http://naviteri.org/2013/03/tsanerul-feerul-getting-better-getting-worse/}{b} 
(\ding{214}~\HL{tsan'ul}{\B{tsan\-'ul}}, \ding{43}~\HL{fe'}{\B{fe'}}, \HL{'ul}{\B{'ul}})  \\
\textsc{Derivations:} \ding{43} \on{1}.~\HL{tIfe'ul}{\B{tìfe'ul}}, worsening (abstract) 
\on{2}.~\HL{sAfe'ul}{\B{säfe'ul}}, worsening (specific)

% few
\hi
\HT{few}{\B{few}} \I{[fEw]} \emph{adp.} across, aiming for the other side. \B{Po spä few pay\-fya fte smar\-it si\-vutx.} He jumped across the stream to track his prey.~\LINK{http://naviteri.org/2011/01/new-year-new-vocabulary/}{b} \\
\textsc{Derivations:} \ding{43} \on{1}.~\HL{fewtusok}{\B{fewtusok}}, opposite (side).

% fewtusok
\hi
\HT{fewtusok}{\B{\cp{few}tusok}} \I{["fEw.tu.sok\textcorner]} \emph{or coll.} \I{["fEw.\t{ts}ok\textcorner]} (\emph{in rapid speech}) \emph{adj.} opposite, on the opposite side. 
\B{Oe kaw\-krr ne few\-tu\-sok\-a pa\-'o kil\-van\-ä ke ka\-mä.} I never went to the opposite side of the river.~\LINK{http://naviteri.org/2011/01/new-year-new-vocabulary/}{b}
(\ding{43}~\HL{few}{\B{few}}, \HL{tok}{\B{tok}})

% fekem
\hi
\HT{fekem}{\B{\cp{fe}kem}} \I{["fE.kEm]} \emph{n.} accident. \B{Na\-ri si fte kea fe\-kem ke li\-ven ngar!} Be careful you don't have an accident!~\LINK{http://naviteri.org/2011/07/txantsana-ultxa-mi-siatll-great-meeting-in-seattle/}{b}
(\ding{43}~\HL{fe'}{\B{fe'}}, \HL{kem}{\B{kem}})

% fekum
\hi
\HT{fekum}{\B{\cp{fe}kum}} \I{["fE.kum]} \emph{or} \I{[.kUm]} \emph{n.} disadvantage, drawback, downside. 
\B{Tì\-ta\-ron\-ìri lä\-ngu tì\-kak\-pam fe\-kum. Kin fkol fra\-inan\-fyat.} I'm sorry to say that deafness is a disadvantage for hunting. You need all your senses.~\LINK{http://naviteri.org/2015/08/aylifyavi-lereyfya-2-cultural-terms-2-and-more/}{b} 
(\ding{43}~\HL{fe'}{\B{fe'}}, \HL{kum}{\B{kum}}, \ding{214}~\HL{tsankum}{\B{tsan\-kum}}) \\
\textsc{Derivations:} \ding{43} \on{1}.~\HL{fekumnga'}{\B{fekumnga'}}, disadvantageous.

% fekumnga'
\hi
\HT{fekumnga'}{\B{\cp{fe}kumnga'}} \I{["fE.kum.NaP]} \emph{or} \I{[.kUm.]} \emph{adj.} disadvantageous. 
(\ding{43}~\HL{fekum}{\B{fe\-kum}}, \ding{214}~\HL{tsankumnga'}{\B{tsan\-kum\-nga'}})

% fewi
\hi
\HT{fewi}{\B{\cp{fe}wi}} \I{["fE.wi]} \emph{vtr.} chase. \B{Rum\-it a no\-lu\-i rä\-'ä fe\-wi.} Don't chase after a foul ball.~\LINK{http://naviteri.org/2013/04/wheres-the-bathroom-and-other-useful-things/}{b}

% feyä
\hi
\B{feyä}~:~\HL{ayfo}{\B{ayfo}}

% fil
\hi 
\HT{fil}{\B{fil}}\textsuperscript{\on{1}} \I{[fil]} \emph{n.} child's toy, plaything.

% 
\hi 
\B{fil}\textsuperscript{\on{2}} : \HL{pil}{\B{pil}}

% fì-
\hi
\HT{fI}{\B{fì-}} \I{[fI.]} \emph{prod. pref. on nouns} (\emph{a.}) this (here) (\emph{proximal singular demonstrative prefix}). \B{Fì\-lì\-'u\-ä ral lu tì\-me\-'em.} This word means harmony.~\LINK{http://forum.learnnavi.org/language-updates/trr-rrtaya/ }{ln} 
(\emph{b. with words describing time}) present, this \B{fì\-trr}, today; \B{fì\-kaym}, this evening; \B{fì\-txon}, tonight; \B{fì\-kin\-trr}, this week; etc. 
(\emph{c. with words describing time and} \B{-o}) duration \B{fì\-trr\-o}, all day (long). 
(\ding{43}~\HL{tsa}{\B{tsa-}}, \HL{fay}{\B{fay}}\texttt{+})

% fì'u
\hi
\HT{fI'u}{\B{fì\cp{'u}}} \I{[fI."Pu]} \emph{pn., n.} this (thing). 
\B{fay\-sä\-num\-vi\-ri ru\-txe fì\-'ut tsli\-vam}, concerning these lessons, please understand this.~\LINK{http://naviteri.org/2010/06/first-post/}{b} 
\B{Fì\-'u\-ti ni\-vìn, ma smuk.} Look at this, siblings.~\LINK{http://naviteri.org/2013/03/tsamsiyu-a-mi-saw-rrtaya-warrior-in-our-sky/}{b} 
(\emph{a. to form contrastive demonstratives}, !stress change!) \B{Fì\-fkxen alu \cp{fì}\-'u lu ftxì\-lor; tsa\-fkxen alu \cp{tsa}\-'u nga\-ti tspang.} \emph{This} vegetable is delicious; \emph{that} one will kill you.~\LINK{http://naviteri.org/2011/12/one-more-for-2011/}{b}
(\ding{43}~\HL{'u}{\B{'u}}) \\
\textsc{Derivations:} \ding{43} \on{1}.~\HL{fwa}{\B{fwa}}, that (\emph{subjective}). 
\on{2}.~\HL{fula}{\B{fula}}, that (\emph{agentive}). 
\on{3}.~\HL{futa}{\B{futa}}, that (\emph{patientive}). 
\on{4}.~\HL{furia}{\B{furia}}, that (\emph{topic}). 

% fìfya 
\hi
\HT{fIfya}{\B{fì\cp{fya}}} \I{[fI."fja]} \emph{adv.} this way, like this. 
\B{Oel nga\-ti stum ke stä\-ngawm krra nga fì\-fya tse\-rì\-syì.} I can barely hear you when you're whispering like this.~\LINK{http://naviteri.org/2011/09/miscellaneous-vocabulary/}{b} 
(\emph{a. idiom}) \B{fì\-fya tsa\-fya}, one way or another. 
\B{Slä fì\-fya tsa\-fya, nga\-ri ti\-re\-al tsa\-tse\-nget ta\-yok.} But one way or the other your spirit will be there.~\LINK{http://naviteri.org/2013/04/sin-asok-sin-zusawkrra-recent-and-upcoming-activities/comment-page-1/\#comment-2316}{b}
(\ding{43}~\HL{fya'o}{\B{fya\-'o}})

% fìkem
\hi
\HT{fIkem}{\B{fì\cp{kem}}} \I{[fI."kEm]} \emph{pn., n.} \textsc{i}.~this (action). 
\B{Ay\-nga\-ri fì\-kem fe\-yä 'e\-'al\-a to\-pur ka\-ngay sì\-yi nì\-'aw.} This will only confirm their worst fears about you.~\LINK{http://forum.learnnavi.org/language-updates/two-words-in-use-eal-and-kangay-si/msg564284/\#msg564284}{ln} \\
\textsc{ii}.~\HT{fIkem si}{\B{fìkem si}} \I{[fI."kEm si]} \emph{vin.} do this (action). \B{Swey lu fwa nga fì\-kem ke si\-vi nì\-sngum.} It's best that you not do this fretfully.~\LINK{http://naviteri.org/2011/01/new-year-new-vocabulary/}{b} 
(\ding{43}~\HL{kem}{\B{kem}})

% fìkintrr
\hi
\HT{fIkintrr}{\B{fì\cp{kin}trr}} \I{[fI."kin.t\s{r}]} \emph{adv.} this (present) week. 
(\ding{43}~\HL{kintrr}{\B{kin\-trr}})

% fìme+
\hi
\HT{fIme}{\B{fìme}}\texttt{+} \I{[fI.mE]} \emph{prod. pref. on n.s} these two (here) (\emph{proximal singular demonstrative prefix}).
\B{Tse\-nga 'aw\-steng\-yä\-pem fì\-me\-kem\-yo lu mek\-tseng a tsun fpxi\-vä\-kìm hì\-'ang tsaw\-ìlä.} Where these two walls come together there's a gap through which insects can get in.~\LINK{http://naviteri.org/2013/04/wheres-the-bathroom-and-other-useful-things/}{b}

% fìmuntrr
\hi
\HT{fImuntrr}{\B{fì\cp{mun}trr}} \I{[fI."mun.t\s{r}]}~\emph{or}~\I{[."mUn.]} \emph{adv.} this (present) weekend. 
\B{Sra\-ne, fì\-mun\-trr lo\-lu oer ma\-wey.} Yes, I had a quiet weekend.~\LINK{http://naviteri.org/2021/04/mipa-ayliu-mipa-sioeykting-new-words-new-explanations/\#comment-33458}{b}
(\ding{43}~\HL{muntrr}{\B{mun\-trr}})

% fìpo
\hi
\HT{fIpo}{\B{fì\cp{po}}} \I{[fI."po]} (\emph{gen.}~\B{fìpeyä}) \emph{pn.} this one (person or thing). 
\B{Fì\-po\-ti oel tspì\-yang, fte tì\-ke\-nong li\-ye\-vu ay\-la\-ru.} I will kill this one as a lesson to the others.~\LINK{http://forum.learnnavi.org/language-updates/capturing-avatar/}{b} 
\B{Fì\-por syaw fko Ìstaw.} This is Ìstaw.~\LINK{http://naviteri.org/2010/09/getting-to-know-you-part-1/}{b} (\emph{a. in contrastive demonstratives}, !stress change!) 
\B{Fì\-kar\-yu alu \cp{fì}\-po lu tsul\-fä\-tu; tsa\-kar\-yu alu \cp{tsa}\-po lu \mbox{skxawng}.} This teacher is a master; that teacher is a fool.~\LINK{http://naviteri.org/2011/12/one-more-for-2011/}{b} 
(\ding{43}~\HL{po}{\B{po}})

% fìpxe+
\hi
\HT{fIpxe}{\B{fìpxe}}\texttt{+} \I{[fI.p'E]} \emph{prod. pref. on n.s} these three (here) (\emph{proximal singular demonstrative prefix}).

% fìrewon
\hi
\HT{fIrewon}{\B{fì\cp{re}won}} \I{[fI."RE.won]} \emph{n.; adv.} this (present) morning. 
(\ding{43}~\HL{rewon}{\B{re\-won}})

% fìtrr
\hi
\HT{fItrr}{\B{fì\cp{trr}}} \I{[fI."t\s{r}]} \emph{n.; adv.} this (present) day, today. 
\B{fì\-trr\-mì le\-tsran\-ten}, on this important day.~\LINK{http://forum.learnnavi.org/language-updates/trr-rrtaya/}{ln} 
\B{Ul\-te fì\-trr\-ä ftxo\-zä\-ri, sìl\-pey oe, ay\-nga\-ru prr\-te' li\-vu.} And I hope you my have pleasure with today's clebration.~\LINK{http://forum.learnnavi.org/language-updates/trr-rrtaya/}{ln} 
(\ding{43}~\HL{trr}{\B{trr}})

% fìtxan
\hi
\HT{fItxan}{\B{fì\cp{txan}}} \I{[fI."t'an]} \emph{adv.} so, to such an extent. 
(\emph{a.}) \B{Po ä\-txä\-le so\-li nì\-ho\-na fì\-txan, ke tsun oe sti\-vo.} She asked so sweetly that I couldn't refuse.~\LINK{http://naviteri.org/2012/01/more-additions-to-the-lexicon/}{b} 
(\emph{b. in result clauses}) \B{Lu poe se\-vin fì\-txan kum\-a yaw\-ne slo\-lu oer.} She was so beautiful that I fell in love with her.~\LINK{http://naviteri.org/2012/06/spring-vocabulary-part-3/}{b}
(\ding{43}~\HL{txan}{\B{txan}})

% fìtxon
\hi
\HT{fItxon}{\B{fì\cp{txon}}} \I{[fI."t'on]} \emph{n., adv.} tonight, this (present) night. 
\B{Fì\-txon na ton a\-la\-he nì\-wotx pe\-lun ke lu teng?} Why is this night not the same like all other nights?~\LINK{http://whyisthisnight.com/addl-more.html\#na\%27vi}{tn} 
(\ding{43}~\HL{txon}{\B{txon}})

% fìtsap
\hi
\HT{fItsap}{\B{fì\cp{tsap}}} \I{[fI."\t{ts}ap\textcorner]} \emph{adv.} each other. (\emph{a. with vtr.}, ‹\B{äp}› \emph{is required}) 
\B{Me\-fo fì\-tsap mä\-po\-le\-yam teng\-krr tsngaw\-vìk.} The two of them hugged each other and wept.~\LINK{http://naviteri.org/2011/10/more-vocabulary-a-bit-of-grammar/}{b} 
(\emph{b. with vin.}) \B{Tse\-nu sì Lo\-ak fì\-tsap ke ha' kaw\-'it.} Tsenu and Loak are a terrible match for each other. (\emph{neither one is to blame for the mismatch})~\LINK{http://naviteri.org/2012/10/mipa-vospxi-mipa-ayliu-new-words-for-the-new-month/}{b} 
\B{Moe smon (moe\-ru) fì\-tsap.} We know each other.~\LINK{http://naviteri.org/2011/12/one-more-for-2011/}{b} 
\B{Me\-fo yaw\-ne lu (snor) fì\-tsap.} They love each other.~\LINK{http://naviteri.org/2011/12/one-more-for-2011/}{b}

% fìtseng
\hi
\HT{fItseng}{\B{fì\cp{tseng}}} \I{[fI."\t{ts}EN]} \emph{or} \B{fì\cp{tse}nge} \I{[fI."\t{ts}E.NE]} \emph{n.; adv.} here, this place. 
\B{Fay\-vrr\-tep fì\-tse\-nge lu kxa\-nì.} These demons are forbidden here.~\LINK{http://wiki.learnnavi.org/index.php?title=Corpus\#Times_Online}{wiki\slash a} 
\B{Nga za\-'u fì\-tseng pxìm srak?} Do you come here often?~\LINK{http://wiki.learnnavi.org/index.php?title=Corpus\#Lemondrop.com}{wiki\slash ld} 
\B{Nì\-'i'a tok oel nì\-mun fì\-tse\-nget!} Finally, I'm here again.~\LINK{http://naviteri.org/2011/07/txantsana-ultxa-mi-siatll-great-meeting-in-seattle/}{b} 
\B{'En si oe, Saw\-tu\-te\-ol tìl\-mok fì\-tseng\-it.} I guess, some Skypeople were just here.~\LINK{http://naviteri.org/2011/01/new-year-new-vocabulary/}{b} 
(\ding{43}~\HL{tsenge}{\B{tse\-nge}})

% fìvospxì
\hi
\HT{fIvospxI}{\B{fìvo\cp{spxì}}} \I{[fI.vo."{}sp'I]} \emph{n.; adv.} this (present) month. 
(\ding{43}~\HL{vospxI}{\B{vo\-spxì}})

% fkan
\hi
\HT{fkan}{\B{fkan}} \I{[fkan]} \emph{vin. resemble in a sensory modality, come to the senses as}, (used to form the `middle voice'). 
\B{Nik\-re\-ri Ri\-ni\-yä 'ur fkan lor.} Rini's hair looks beautiful.~\LINK{http://naviteri.org/2012/11/renu-ayinanfyaya-the-senses-paradigm/}{b} 
\B{Nik\-re\-ri Ri\-ni\-yä fkan lor.} Rini's hair is pleasant to the senses.~\LINK{http://naviteri.org/2012/11/renu-ayinanfyaya-the-senses-paradigm/}{b} 
\B{Ra\-lu\-ri fa\-hew fkan oe\-ru na ye\-rik.} Ralu smells like a hexapede to me.~\LINK{http://naviteri.org/2012/11/renu-ayinanfyaya-the-senses-paradigm/}{b} 
\B{Ra\-lu\-ri fkan na ye\-rik.} Ralu smells (looks\slash sounds) like a hexapede.~\LINK{http://naviteri.org/2012/11/renu-ayinanfyaya-the-senses-paradigm/}{b} \B{Fì\-na\-er\-ì\-ri sur fkan lor.} This drink tastes good.~\LINK{http://naviteri.org/2012/11/renu-ayinanfyaya-the-senses-paradigm/}{b} 
\B{Fì\-na\-er\-ì\-ri fkan ftxì\-lor.} This drink tastes good.~\LINK{http://naviteri.org/2012/11/renu-ayinanfyaya-the-senses-paradigm/}{b} 
\B{Syu\-ve\-ri ay\-ftxo\-zä\-yä ay\-nga\-ru fki\-van on\-lor ftxì\-lor\-sì nì\-wotx!} May the food of your holidays taste and smell good.~\LINK{http://naviteri.org/2012/11/renu-ayinanfyaya-the-senses-paradigm/}{b}

% fkarut
\hi
\HT{fkarut}{\B{\cp{fka}rut}} \I{["fka.Rut\textcorner]} \emph{or} \I{[.RUt\textcorner]} (\textsc{rn}: \B{\cp{fka}rùt} \I{["fka.RUt\textcorner]}) \emph{vtr.} peel. 

% fkay
\hi
\HT{fkay}{\B{fkay}} \I{[fkaj]} \emph{adj.} hateful. 

% fkew
\hi
\HT{fkew}{\B{fkew}} \I{[fkEw]} \emph{adj.} mighty. 
\B{na ay\-sa\-ngek a\-fkew}, like mighty trunks (of trees).~\LINK{http://www.pandorapedia.com/navi/music/ritual_music}{pp}

% -fkeyk
\hi
\HT{-fkeyk}{\B{-fkeyk}} \I{[.fkEjk\textcorner]} \emph{prod. suf. on nouns to mean} situation, state, status. 
\B{Kil\-van\-fkeyk lu fya\-pe fì\-trr?} What's the condition of the river today?~\LINK{http://naviteri.org/2011/04/yafkeykiri-plltxe-frapo-everyone-talks-about-the-weather/}{b} 
\B{Saw\-tu\-te\-ri ron\-sem\-fkeyk\-it ke tsun kaw\-tu tsli\-vam.} No one can understand the state of mind of the Sky People.~\LINK{http://naviteri.org/2011/04/yafkeykiri-plltxe-frapo-everyone-talks-about-the-weather/}{b} 
(\emph{a. conv. phrase}) \B{Nga\-fkeyk pe\-fya\slash fya\-pe?} `What’s your status?\slash How ya doin'?'~\LINK{http://naviteri.org/2018/02/negative-questions-in-navi/\#comment-28201}{b}
(\ding{43}~\HL{tIfkeytok}{\B{tì\-fkey\-tok}})

% fkeytok
\hi
\HT{fkeytok}{\B{\cp{fkey}tok}} \I{["fkEj.t$\cdot$ok\textcorner]} \emph{vin., inf.}\on{22} exist. \B{Ngal fwe\-rew a tu\-te ke fkey\-tok.} The person you're looking for doesn't exist.~\LINK{http://naviteri.org/2011/04/yafkeykiri-plltxe-frapo-everyone-talks-about-the-weather/}{b} 
(\ding{43}~\HL{kifkey}{\B{ki\-fkey}}, \HL{tok}{\B{tok}}) \\
\textsc{Derivations:} \ding{43} \on{1}. \HL{tIfkeytok}{\B{tìfkeytok}}, state, condition, situation.

% fkeyä
\hi
\B{\cp{fke}yä}~:~\HL{fko}{\B{fko}}

% fkio
\hi
\HT{fkio}{\B{\cp{fki}o}} \I{["fki.o]} \emph{n., fauna} tetrapteron. 

% fkip
\hi
\HT{fkip}{\B{fkip}} \I{[fkip\textcorner]} \emph{adp.} up among. 
(\emph{a. idiom}) \B{(na) fwam\-pop fkip fì\-wopx}, ``Fish out of water.'' [lit.: (like a) tapirus in the clouds].~\LINK{http://naviteri.org/2021/08/ulte-ayyoratu-leiu-and-the-winners-are-2/}{b}
(\ding{43}~\HL{kip}{\B{kip}}, \HL{takip}{\B{ta\-kip}}, \HL{tafkip}{\B{ta\-fkip}}) 

% fko
\hi
\HT{fko}{\B{fko}} \I{[fko]} (\emph{gen.}~\B{fkeyä}) \emph{pn.} (\emph{in the general sense}) one, you, they; (\emph{also used to convey passive voice}).  
\B{Oe\-ru syaw fko Ney\-ti\-ri.} My name is Neytiri [lit.: one calls (to) me Neytiri].~\LINK{http://forum.learnnavi.org/language-updates/the-reflexive-\%28from-a-frommer-email!!!!\%29/}{ln} 
\B{Ney\-ti\-ri\-ti fkol pä\-nu\-to\-lìng oe\-ru!} Neytiri was promised to me.~\LINK{http://forum.learnnavi.org/language-updates/panuting/}{ln\slash a\textsuperscript{c}} 
\B{Pa\-lu\-lu\-kan new fkot yi\-vom.} Thanator wants to eat you.~\LINK{http://naviteri.org/2012/10/mipa-vospxi-mipa-ayliu-new-words-for-the-new-month/}{b} 
\B{Hem nge\-yä zen\-ke fko\-ru li\-vu, Tsay\-hem a nga\-ru prr\-te' ke lu.} As for any action you don't like, don't do that action to others.~\LINK{http://wiki.learnnavi.org/index.php?title=Corpus\#Good_Morning_America}{wiki}
(\emph{a. proverb}) \B{Ke kur fko fa kxe\-tse.} ``One can't hang by a tail,'' i.e. \emph{don't rely on s.t.\slash s.o. untrustworthy or useless, just as a Na'vi tail can't be relied on to bear weight}.~\LINK{http://naviteri.org/2021/08/ulte-ayyoratu-leiu-and-the-winners-are-2/}{b}

% fkxake
\hi
\HT{fkxake}{\B{\cp{fkxa}ke}} \I{["fk'a.kE]} \emph{vin.} itch. 
\B{Sran, sran, oe\-ri pì\-tìk tsa\-tse\-nget a mì tal! Fkxa\-ke nì\-ftxan kum\-a ter\-kup.} Yeah, yeah, scratch that place on my back! It itches like crazy!~\LINK{http://naviteri.org/2020/07/vospxivol-lefpom-happy-august/}{b}
(\ding{214}~\HL{pItIk}{\B{pì\-tìk}}) \\
\textsc{Derivations:} \ding{43} \on{1}.~\HL{fkxakewll}{\B{fkxakewll}}, itch plant.

% fkxakewll
\hi
\B{\cp{fkxa}kewll} \I{["fk'a.kE.w\s{l}]} \emph{n.}, \ding{96} itch plant, \emph{cynaroidia glauca}. 
(\ding{43}~\HL{fkxake}{\B{fkxa\-ke}}, \HL{'ewll}{\B{'e\-wll}})

% fkxara
\hi 
\HT{fkxara}{\B{\cp{fkxa}ra}} \I{["fk'a.Ra]} \emph{n.} stress (\emph{mental or emotional feeling}). 
\B{Krra oe ftxu\-lì\-'u, pxìm 'efu fkxa\-rat nì\-txan.} When I give a speech, I often feel a lot of stress.~\LINK{http://naviteri.org/2018/03/100a-liu-amip-64-new-words-part-1/}{b} \\
\textsc{Derivations:} \ding{43} \on{1}.~\HL{fkxaranga'}{\B{fkxaranga'}}, stressful.

% fkxaranga'
\hi 
\HT{fkxaranga'}{\B{\cp{fkxa}ranga'}} \I{["fk'a.Ra.NaP]} \emph{adj.} stressful. 
\B{Nge\-yä fpom\-tokx\-ìri fì\-tì\-fkey\-tok a\-fkxa\-ra\-nga' lu le\-hrr\-ap.} This stressful situation is dangerous to your health.~\LINK{http://naviteri.org/2018/03/100a-liu-amip-64-new-words-part-1/}{b} 
(\ding{43}~\HL{fkxara}{\B{fkxa\-ra}})

% fkxen
\hi
\HT{fkxen}{\B{fkxen}} \I{[fk'En]} \emph{n.}~\ding{96} food of vegetable (as opposed to animal) origin. 
\B{aw\-ngal yom fkxen\-ti le\-rìk nì\-wotx}, we eat leafy vegetable-food.~\LINK{http://whyisthisnight.com/addl-more.html\#na\%27vi}{tn} 
\B{Fkxen\-ti pxìm yom fkol nì\-yrr.} Vegetables are often eaten raw.~\LINK{http://naviteri.org/2013/01/awvea-posti-zisita-amip-first-post-of-the-new-year/}{b}

% fkxile
\hi
\HT{fkxile}{\B{\cp{fkxi}le}} \I{["fk'i.lE]} \emph{n.} bib necklace. 
\B{Fkxi\-le pewn\-ta tu\-té\-yä kur.} The bib necklace hangs from the woman's neck.~\LINK{http://naviteri.org/2013/04/wheres-the-bathroom-and-other-useful-things/}{b}
\B{Pll\-txe nì\-tsleng! Tsa\-fkxi\-let ke to\-lìng ngar En\-tul!}Liar! En\-tu didn't give you that necklace!~\LINK{http://naviteri.org/2017/09/zisikrr-amip-ayliu-amip-new-words-for-the-new-season/}{b}
\B{Poe\-ri sä\-fmong lor\-a tsa\-fkxi\-le\-yä lo\-lu na ay\-skxe mì te'\-lan.} For her, the theft of that beautiful necklace was like stones in her heart.~\LINK{http://naviteri.org/2021/09/aawa-ayliu-amip-a-few-new-words/}{b} 

% flawkx
\hi
\HT{flawkx}{\B{flawkx}} \I{[flawk']} \emph{n.} leather.

% flawm
\hi
\HT{flawm}{\B{flawm}} \I{[flawm]} \emph{n.} cheek.

% flä
\hi
\HT{flA}{\B{flä}} \I{[fl\ae]} \emph{vin.} succeed. 
\B{Tì\-ste\-ftxaw nge\-yä flo\-lä nìl\-tsan!} Your examination succeeded well!~\LINK{http://wiki.learnnavi.org/index.php?title=Canon\#Extracts_from_various_emails}{wiki} 
\B{Tì\-tu\-sa\-ron\-ì\-ri fte fli\-vä, zene fko si\-vutx smar\-it ni\-nan nì\-no.} To succeed at hunting, you have to track your prey by reading (the forest) with attention to detail.~\LINK{http://naviteri.org/2011/09/\%E2\%80\%9Cby-the-way-what-are-you-reading\%E2\%80\%9D/}{b} 
(\emph{a. proverb}) \B{Flä ke flä, ley sä\-fmi.} Whether you succeed or not, the attempt has value.~\LINK{http://naviteri.org/2014/03/value-and-worth/}{b} 
(\ding{214}~\HL{nui}{\B{nui}}) \\
\textsc{Derivations:} \ding{43} \on{1}.~\HL{sAflA}{\B{sä\-flä}}, (instance of) success.
\on{2}.~\HL{tIflA}{\B{tì\-flä}}, success (in general). 
\on{3}.~\HL{nIflA}{\B{nì\-flä}}, successfully.

% flew
\hi
\HT{flew}{\B{flew}} \I{[flEw]} \emph{n.} throat. 

% flefle
\hi
\HT{flefle}{\B{\cp{fle}fle}} \I{["flE.flE]} \emph{n.}, \ding{96} Sol's delight, magnetotrophic plant, \emph{calamariphyllum elegans}. 

% flel
\hi 
\HT{flel}{\B{flel}} \I{[flEl]} \emph{vtr.} trick s.o., fool s.o. 
\B{En\-tul Pey\-ral\-it fì\-txan flo\-lel ku\-ma fpìl poe san oe yaw\-ne lu por.} En\-tu fooled Pey\-ral so mucht that she thought he loved her.~\LINK{https://naviteri.org/2024/03/fleltrra-ayliu-words-for-april-fools-day/}{b} \\
\textsc{Derivations}: \ding{43} \on{1}.~\HL{tIflel}{\B{tìflel}}, trickery (abstract).
\on{2}.~\HL{sAflel}{\B{sä\-flel}}, trick, hoax.
\on{3}.~\HL{fleltu}{\B{flel\-tu}}, fool, sucker.
\on{4}.~\HL{fleltrr}{\B{flel\-trr}}, April Fool's Day.

% fleltrr
\hi 
\HT{fleltrr}{\B{\cp{flel}trr}} \I{["flEl.t\s{r}]} \emph{n.} April Fool's Day. 
\B{Flel\-tu slu rä\-'ä! Fì\-trr lu Flel\-trr!} Don't be fooled! Today is April Fool's Day!.~\LINK{https://naviteri.org/2024/03/fleltrra-ayliu-words-for-april-fools-day/}{b}
(\ding{43}~\HL{flel}{\B{flel}}, \HL{trr}{\B{trr}}) 

% fleltu
\hi 
\HT{fleltu}{\B{\cp{flel}tu}} \I{["flEl.tu]} \emph{n.} fool, sucker, mark, s.o. easily tricked. 
\B{Flel\-tu slu rä\-'ä! Fì\-trr lu Flel\-trr!} Don't be fooled! Today is April Fool's Day!~\LINK{https://naviteri.org/2024/03/fleltrra-ayliu-words-for-april-fools-day/}{b}
(\ding{43}~\HL{flel}{\B{flel}})

% flì
\hi
\HT{flI}{\B{flì}} \I{[flI]} \emph{adj., nfp.} thin. 
\B{Krro krro, flì\-a vul a\-ru\-sey to nutx\-a pum a\-ke\-ru\-sey lu txur.} A thin living branch is sometimes stronger than a thick dead one.~\LINK{http://naviteri.org/2011/11/\%E2\%80\%99a\%E2\%80\%99awa-ayli\%E2\%80\%99u-amip-ni\%E2\%80\%99aw-\%E2\%80\%94-a-few-new-words-only/}{b} 
\B{Tskxe\-pay\-ri lu frir a\-flì sìn 'o\-ra.} There's a thin layer of ice on the lake.~\LINK{http://naviteri.org/2018/07/ta-sulfatu-a-ayliu-niul-more-words-from-our-experts/}{b}
(\emph{a. proverb}) \B{Spä skxawng sìn 'a\-na a\-flì.} ``A fool jumps onto a thin vine,'' i.e. \emph{don't engage in an unpromising and\slash or potentially risky cause}.~\LINK{http://naviteri.org/2021/08/ulte-ayyoratu-leiu-and-the-winners-are-2/}{b}
\B{Tsa\-ye\-rik te\-rul ne 'awkx. Spä skxawng sìn 'a\-na a\-flì.} That hexapede is running toward the cliff. Only a fool jumps onto a thin vine.~\LINK{http://naviteri.org/2021/08/ulte-ayyoratu-leiu-and-the-winners-are-2/}{b}
(\ding{214}~\HL{nutx}{\B{nutx}}; \ding{43}~\HL{litsi}{\B{li\-tsi}}) \\
\textsc{Derivations:} \ding{43} \on{1}.~\HL{flInutx}{\B{flìnutx}}, thickness.

% flìnutx
\hi
\HT{flInutx}{\B{flì\cp{nutx}}} \I{[flI."nut']} \emph{n.} thickness.   
\B{Sre fwa sìn tskxe\-pay tì\-ran, ze\-ne fko flì\-nutx\-it sti\-ve\-ftxaw.} It's necessary to check the thickness of the ice before walking on it.~\LINK{http://naviteri.org/2011/11/\%E2\%80\%99a\%E2\%80\%99awa-ayli\%E2\%80\%99u-amip-ni\%E2\%80\%99aw-\%E2\%80\%94-a-few-new-words-only/}{b}
(\ding{43}~\HL{flI}{\B{flì}}, \HL{nutx}{\B{nutx}})

% flrr
\hi
\HT{flrr}{\B{flrr}} \I{[fl\s{r}]} \emph{adj.} gentle, mild, tender (\emph{for both people and things}). 
\B{Keng tsam\-si\-yu ze\-ne flrr li\-vu ay\-e\-veng\-hu.} Even a warrior must be gentle with children.~\LINK{http://naviteri.org/2013/02/vospxi-ayol-posti-apup-short-post-for-a-short-month/}{b} 
\B{Flrr\-a tom\-pa ze\-rup.} A gentle rain is falling.~\LINK{http://naviteri.org/2013/02/vospxi-ayol-posti-apup-short-post-for-a-short-month/}{b} 
(\ding{214}~\HL{kakan}{\B{ka\-kan}}) \\
\textsc{Derivations:} \ding{43} \on{1}.~\HL{nIflrr}{\B{nìflrr}}, tenderly, mildly, gently. 
\on{2}.~\HL{tIflrr}{\B{tìflrr}}, gentleness, tenderness.
\on{3}.~\HL{flrrtsawl}{\B{flrr\-tsawl}}, sailfin goliath. 

% flrrtsawl
\hi 
\HT{flrrtsawl}{\B{\cp{flrr}tsawl}} \I{["fl\s{r}.\t{ts}awl]} \emph{n., fauna} sailfin goliath. 
(\ding{43}~\HL{flrr}{\B{flrr}}, \HL{tsawl}{\B{tsawl}})

% fmal
\hi
\HT{fmal}{\B{fmal}} \I{[fmal]} \emph{vtr.} sustain. \B{Te\-ri fwa fmal fì\-lì'\-fya\-ti a\-yaw\-ne.} On keeping this beloved language alive.~\LINK{http://naviteri.org/2011/12/ulte-yora\%E2\%80\%99tu-leiu-and-the-winner-is/}{b} \B{Lu\-ke pay, ke tsun ay\-oe tì\-rey\-ti fmi\-val.} Without water we cannot sustain life.~\LINK{http://forum.learnnavi.org/language-updates/for-\%28duration\%29-jesus-loan-word/msg372740/\#msg372740}{ln\slash g}

% fmawn
\hi
\HT{fmawn}{\B{fmawn}} \I{[fmawn]} \emph{n.} news, something to report.  
\B{Sìl\-pey oe, la\-yu oe\-ru ye'\-rìn sìl\-tsan\-a fmawn.} I hope that I will have good news soon.~\LINK{http://forum.learnnavi.org/news-announcements/a-response-from-paul-frommer!/}{ln} 
\B{Slä lu oe\-ru fmawn\-o a\-sìl\-tsan.} But I have some good news.~\LINK{http://forum.learnnavi.org/language-updates/trr-rrtaya/}{ln} 
\B{Nì\-sngä\-'i fmawn\-it fo nar\-mew wi\-van, slä nì\-'i'\-a fra\-por lo\-lo\-nu.} They originally wanted to hide the news, but in the end they revealed it everyone.~\LINK{http://naviteri.org/2011/10/more-vocabulary-a-bit-of-grammar/}{b} 
(\emph{a. idiom}) \B{Pe\-fmawn?} `What's the report?\slash What's the news?\slash What's new?'~\LINK{http://naviteri.org/2018/02/negative-questions-in-navi/\#comment-28201}{b} \\
\textsc{Derivations:} \ding{43} \on{1}.~\HL{fmawnta}{\B{fmawnta}}, (the news) that.

% fmawnta
\hi
\HT{fmawnta}{\B{\cp{fmawn}ta}} \I{["fmawn.ta]} \emph{conj.} (\emph{short form for} \B{fmawnit a}) \emph{to introduce reported speech of the verbs} \ding{43}~\HL{peng}{\B{peng}} \emph{and} \HL{stawm}{\B{stawm}}, (the news) that. 
\B{Sto\-lawm oel fmawn\-ta fo new hi\-vum.} I heard they want to leave.\LINK{http://naviteri.org/2011/08/reported-speech-reported-questions/}{b} 
\B{Ngal po\-leng oer fmawn\-ta po to\-ler\-kup.} You told me that he died.~\LINK{http://naviteri.org/2011/08/reported-speech-reported-questions/}{b} 
(\ding{43}~\HL{fmawn}{\B{fmawn}})

% fmetok
\hi
\HT{fmetok}{\B{\cp{fme}tok}} \I{["fmE.tok\textcorner]} \emph{vtr.} test (s.o.).
\B{Ney\-ti\-ril fma\-ye\-tok fì\-tsam\-si\-yut.} Neytiri will test this `warrior'.~\LINK{https://forum.learnnavi.org/language-updates/about-fmetok-and-pofo/}{a\textsuperscript{c}\slash ln} \\
\textsc{Derivations:} \ding{43} \on{1}.~\HL{tIfmetok}{\B{tìfmetok}}, test.

% fmi
\hi
\HT{fmi}{\B{fmi}} \I{[fmi]} \emph{vtr., modal} try, attempt. 
\B{Fmi pi\-vll\-txe alo avol nì\-win!} Try to say that eight times fast!~\LINK{http://naviteri.org/2011/09/miscellaneous-vocabulary/}{b} 
\B{Oe fmo\-li nu\-ä\-ngi.} I tried, but unfortunately I failed.~\LINK{http://naviteri.org/2013/04/wheres-the-bathroom-and-other-useful-things/}{b} 
\B{Fma\-yi oe 'i\-veyng ye'\-rìn.} I will try to answer soon.~\LINK{http://forum.learnnavi.org/language-updates/another-email-ta-frommer-no-answers-as-of-yet-but-a-few-interesting-points/msg196441/\#msg196441}{ln} (\ding{43}~\B{may'})\\
\textsc{Derivations:} \ding{43} \on{1}.~\HL{tIfmi}{\B{tìfmi}}, attempt. 
\on{2}.~\HL{sAfmi}{\B{säfmi}}, attempt.

% fmokx
\hi
\HT{fmokx}{\B{fmokx}} \I{[fmok']} \emph{n.} jealousy, envy (\emph{neutral connotation unless otherwise specified with} ‹\B{äng}› \emph{or} ‹\B{ei}›). 
\B{lu oe\-ru fmokx}, I'm jealous.~\LINK{http://naviteri.org/2011/10/more-vocabulary-a-bit-of-grammar/}{b} 
\B{Fu\-ri\-a fì\-txan fnä\-ngan Ul\-rey\-ìl tsko swi\-zaw\-it, lu oe\-ru fmokx.} I'm jealous of the fact that Ul\-rey is such an excellent archer.~\LINK{http://naviteri.org/2011/10/more-vocabulary-a-bit-of-grammar/}{b} 
\B{Fko pll\-txe\-i\-e san me\-nga mun\-txa slo\-lu sìk!~Sey\-kxel sì ni\-tram!~Slä lu oe\-ru fmokx nì\-'it.} I'm so happy to hear you got married! Congratulations! I am a little envious, though.~\LINK{http://naviteri.org/2011/10/more-vocabulary-a-bit-of-grammar/}{b} \\
\textsc{Derivations:} \ding{43} \on{1}.~\HL{nIfmokx}{\B{nìfmokx}}, jealously, enviously.

% fmong
\hi 
\HT{fmong}{\B{fmong}} \I{[fmoN]} \emph{vtr.} steal, rob. 
\B{Fo fmong, fo nekx, fo ska'a.} They steal, they burn, they destroy.~\LINK{http://naviteri.org/2020/06/mi-tanlokxe-oeya-srr-afpxamo-terrible-days-in-my-country/}{b} \\
\textsc{Derivations:} \ding{43} \on{1}.~\HL{tIfmong}{\B{tì\-fmong}}, theft.  
\on{2}.~\HL{fmongyu}{\B{fmong\-yu}}, thief.
\on{3}.~\HL{sAfmong}{\B{sä\-fmong}}, theft (\emph{particular instance}).

% fmongyu
\hi 
\HT{fmongyu}{\B{\cp{fmong}yu}} \I{["fmoN.ju]} \emph{n.} thief. 
(\ding{43}~\HL{fmong}{\B{fmong}}) 

% fnan
\hi
\HT{fnan}{\B{fnan}} \I{[fnan]} \emph{vtr.} be good at. 
\B{Sra\-ke fnan ngal lì'\-fya\-ti le\-Na'\-vi?} Are you good at Na'vi?~\LINK{http://naviteri.org/2010/09/getting-to-know-you-part-3/}{b} 
\B{Fol fnan fu\-ta 'em tey\-lut a fì\-'u sìn pe\-yì?} How good are they at cooking \emph{teylu}?~\LINK{http://naviteri.org/2013/01/lifyari-po-peyi-how-good-is-her-navi/}{b}
(\ding{214}~\HL{wAtx}{\B{wätx}})

% fnawe'
\hi
\HT{fnawe'}{\B{fna\cp{we'}}} \I{[fna."wEP]} \emph{adj.} cowardly. \\
\textsc{Derivations:} \ding{43} \on{1}.~\HL{fnawe'tu}{\B{fna\-we'\-tu}}, coward.
\on{2}.~\HL{tIfnawe'}{\B{tì\-fna\-we'}}, cowardice.

% fnawe'tu
\hi
\HT{fnawe'tu}{\B{fna\cp{we'}tu}} \I{[fna."wEP.tu]} \emph{n.} coward. 
\B{Txan\-tstew sä\-ro\-'a si, fna\-we'\-tu ke si.} A hero does great deeds, a coward does not.~\LINK{http://naviteri.org/2012/03/spring-vocabulary-part-1/}{b}
\B{Mu\-nge fna\-we'\-tul lewng\-it soaia\-ru sne\-yä.} A coward brings shame to his\slash her family.~\LINK{http://naviteri.org/2021/09/aawa-ayliu-amip-a-few-new-words/}{b}
(\emph{a. proverb}) \B{Kxìm utu\-ftu fna\-we'\-tu.} ``A coward commands from the canopy,'' i.e. \emph{a real leader will have `boots on the ground' and will help out, whereas a coward will only tell people (from afar) what to do}.~\LINK{http://naviteri.org/2021/08/ulte-ayyoratu-leiu-and-the-winners-are-2/}{b}
(\ding{43}~\HL{fnawe'}{\B{fna\-we'}})

% fne-
\hi
\HT{fne}{\B{fne-}} \I{[fnE]} \emph{pref. on nouns} kind (of), type (of). 
\B{Kel\-ut\-ral lu fne\-ut\-ral a\-zey.} Hometree is a special kind of tree.~\LINK{http://naviteri.org/2012/06/spring-vocabulary-part-3/}{b} 
\B{la\-ye\-i\-u oer krr nì\-'ul fte ngi\-vop ay\-lì\-'ut sì tsay\-fne\-sä\-num\-vit a tsun fra\-por srung si\-vi}, I will have more time to create words and those kinds of lessons that will help everyone.~\LINK{http://forum.learnnavi.org/language-updates/trr-rrtaya/}{ln} 
\B{Lu pì\-lok\-ur pxe\-sì\-kan sì pxe\-fne\-'u\-pxa\-re.} The blog has three aims and three kinds of messages.~\LINK{http://naviteri.org/2010/06/first-post/}{b} 
\B{Por war\-mo\-u fra\-fne\-i\-o\-ang nì\-tut, slä nì\-pxi pxe\-syì\-rä\-fì.} He just couldn't get over all the animals, but was especially taken with the 3 giraffe.~\LINK{http://forum.learnnavi.org/language-updates/txelanit-hivawl/}{ln} 
(\ding{43}~\HL{fnel}{\B{fnel}})

% fnel
\hi
\HT{fnel}{\B{fnel}} \I{[fnEl]} \emph{n.} kind, type. 
\B{tsa\-fnel syu\-lang\-ä}, that kind of flower.~\LINK{http://forum.learnnavi.org/language-updates/trr-rrtaya/}{ln} 
(\ding{43}~\HL{fne}{\B{fne-}}) \\
\textsc{Derivations:} \ding{43} \on{1}.~\HL{fnelan}{\B{fnelan}}, male. 
\on{2}.~\HL{fnele}{\B{fnele}}, female.
\on{3}.~\HL{lI'fyafnel}{\B{lì'\-fya\-fnel}}, dialect, variety of a language.

% fnelan
\hi 
\HT{fnelan}{\B{\cp{fne}lan}} \I{["fnE.lan]} \emph{n.} male.
(\ding{43}~\HL{fnel}{\B{fnel}}) \\
\textsc{Derivations:} \ding{43} \on{1}.~\HL{lefnelan}{\B{lefnelan}}, male.

% fnele
\hi 
\HT{fnele}{\B{\cp{fne}le}} \I{["fnE.lE]} \emph{n.} female.
(\ding{43}~\HL{fnel}{\B{fnel}}) \\
\textsc{Derivations:} \ding{43} \on{1}.~\HL{lefnele}{\B{lefnele}}, female.

% fnepe
\hi
\HT{fnepe}{\B{\cp{fne}pe}} \I{["fnE.pE]} \emph{inter.} which kind. 
(\ding{43}~\HL{fnel}{\B{fnel}}, \HL{pefnel}{\B{pe\-fnel}})

% fnetxum
\hi
\HT{fnetxum}{\B{fne\cp{txum}}} \I{[fnE."t'um]} \emph{or} \I{[."t'Um]} (\textsc{rn}: \B{fne\cp{txùm}} \I{[fne."dUm]}) \emph{idiom} 
\B{Oe\-ri lu swo\-a fne\-txum.} I'm allergic to alcohol [lit.: alcohol is a kind of poison for me].~\LINK{http://forum.learnnavi.org/language-updates/allergies/}{ln}
(\ding{43}~\HL{fne}{\B{fne-}}, \HL{txum}{\B{txum}})

% fnu 
\hi
\HT{fnu}{\B{fnu}} \I{[fnu]} \emph{vin.} be quiet, be still. 
\B{Fnu, ma smuk, fnu! Oel he\-fi ye\-rik\-it!} Quiet, everyone! I smell a hexapede!~\LINK{http://naviteri.org/2012/11/renu-ayinanfyaya-the-senses-paradigm/}{b} 
\B{Oe ä\-txä\-le so\-li tsnì fra\-po fni\-vu.} I request that everyone be quiet.~\LINK{http://naviteri.org/2011/08/reported-speech-reported-questions/comment-page-1/\#comment-1060}{b}\\
\textsc{Derivations:} \ding{43} \on{1}.~\HL{tIfnu}{\B{tìfnu}}, quiet, silence. 
\on{2}.~\HL{nIfnu}{\B{nìfnu}}, quietly.

% fngap
\hi
\HT{fngap}{\B{fngap}} \I{[fNap\textcorner]} \emph{n.} metal. \B{Fu\-ri\-a txu\-la tsa\-lew\-ti lu srä sìl\-tsan to fngap.} For constructing that cover, woven cloth is better than metal.~\LINK{http://naviteri.org/2013/07/pximaw-tsawlultxa-just-after-the-meet-up/}{b}\\
\textsc{Derivations:} \ding{43} \on{1}.~\HL{lefngap}{\B{lefngap}}, metallic. 
\on{2}.~\HL{fngapsutxwll}{\B{fngapsutxwll}}, metalfollowing plant.

% fngapsutxwll
\hi
\HT{fngapsutxwll}{\B{\cp{fngap}sutxwll}} \I{["fNap\textcorner.sut'.w\s{l}]} \emph{n.}, \ding{96} metalfollowing plant, \emph{pennanemone ciliare}.
(\ding{43}~\HL{fngap}{\B{fngap}}, \HL{sutx}{\B{sutx}}, \HL{'ewll}{\B{'e\-wll}})

% fngä' 
\hi
\HT{fngA'}{\B{fngä'}} \I{[fN\ae{}P]} \emph{vin.} relieve oneself. 
\B{Fko tsun fngi\-vä' pe\-seng?} Where is the bathroom? [lit.: Where can one relieve oneself?]~\LINK{http://naviteri.org/2013/04/wheres-the-bathroom-and-other-useful-things/}{b}

% fngä'tseng
\hi 
\HT{fngA'tseng}{\B{\cp{fngä'}tseng}} \I{["fN\ae{}P.\t{ts}EN]} \emph{n.} restroom (\emph{general term, any `place of elimination'}). 
(\ding{43}~\HL{fngA'}{\B{fngä'}}, \HL{tseng}{\B{tseng(e)}}; \HL{moafngA'}{\B{mo a fngä'}})

% fngo'
\hi
\HT{fngo'}{\B{fngo'}} \I{[fNoP]} \emph{vtr.} require, demand. 
\B{Fol fte ay\-spe\-'e\-tut li\-vo\-nu fngo' 'u\-pet?} What are they demanding in order for them to release the captives?~\LINK{http://naviteri.org/2012/01/mipa-zisit-ayliu-amip-new-words-for-the-new-year/}{b} 
\B{Fì\-fne\-tì\-fkey\-tok\-ìl fngo' fu\-ta kem si\-vi fko pxi\-ye'\-rìn.} This kind of situation requires immediate action.~\LINK{http://naviteri.org/2012/01/mipa-zisit-ayliu-amip-new-words-for-the-new-year/}{b} 
\B{Kar\-yul fngo\-lo' fu\-ta ay\-nu\-me\-yu pi\-va\-te ye'\-krr.} The teacher required the students to arrive early.~\LINK{http://naviteri.org/2012/01/mipa-zisit-ayliu-amip-new-words-for-the-new-year/}{b} \\  
\textsc{Derivations:} \ding{43} \on{1}.~\HL{sAfngo'}{\B{säfngo'}}, requirement, demand.

% fo
\hi
\HT{fo}{\B{fo}} : \HL{ayfo}{\B{ayfo}} 

% fpak
\hi
\HT{fpak}{\B{fpak}} \I{[fpak\textcorner]} \emph{vin.} hold off, suspend an action (that is already in progress). 

% fpap
\hi 
\HT{fpap}{\B{fpap}} \I{[fpap\textcorner]} \emph{vtr.} pound (\emph{striking heavily and repeatedly}). 
\B{Krr\-a sti nì\-txan, pol fyan\-yot fpap fa me\-syokx.} When he's angry, he pounds the table with his hand.~\LINK{http://naviteri.org/2020/11/vospxivopeya-ayliu-amip-novembers-new-words/}{b}

% fpe'
\hi
\HT{fpe'}{\B{fpe'}} \I{[fpEP]} \emph{vtr.} send. 
\B{Fpo\-le' ay\-ngal oer a 'u\-pxa\-re.} The message you have sent me.~\LINK{http://forum.learnnavi.org/news-announcements/a-response-from-paul-frommer!/}{ln}

% fpeio
\hi
\HT{fpeio}{\B{fpe\cp{i}o}} \I{[fpE."{}i.o]} \emph{n.} (ceremonial) challenge. 

% fpi
\hi
\HT{fpi}{\B{fpi}} \I{[fpi]} \emph{adp.}\texttt{+} for (the benefit or sake of). \B{Fay\-u\-pxa\-re la\-yu ay\-sngä\-'i\-yu\-fpi.} These messages will be for beginners.~\LINK{http://naviteri.org/2010/06/first-post/}{b} \B{Stxe\-lit fpo\-le' oel nga\-fpi.} I sent the gift for you (\emph{i.e. I sent it to someone else for your sake})~\LINK{http://forum.learnnavi.org/language-updates/i-have-something-to-say/}{ln} \B{Oe\-ri tì\-rey\-ti oel stxe\-nu\-to\-lìng fpi o\-lo' aw\-nge\-yä.} I offered my life for the sake of our clan.~\LINK{http://naviteri.org/2011/09/miscellaneous-vocabulary/}{b}

% fpivìl
\hi 
\HT{fpivIl}{\B{fpi\cp{vìl}}} \I{[fpi."vIl]} \emph{intj.} ``hmm, let's see, let me think''. 
\B{Fpi\-vìl … Kxawm nga\-ru tì\-yawr.} Hmm . . . Perhaps you're right.~\LINK{http://naviteri.org/2019/06/50a-liu-amip-40-new-words/}{b} 
(\ding{43}~\HL{fpIl}{\B{fpìl}})

% fpìl
\hi
\HT{fpIl}{\B{fpìl}} \I{[fpIl]} \emph{vtr.} think. 
\B{Fwa fpìl ngal fu\-ta fay\-lì'\-fya\-vi li\-vu le\-sar oe\-ru prr\-te' lu.} That you think that these expressions are useful pleases me.~\LINK{http://naviteri.org/2010/09/getting-to-know-you-part-1/comment-page-1/\#comment-293}{b} 
\B{Fra\-krr po fpìl nì\-tsim nì\-wotx.} Her thinking is always completely original.~\LINK{http://naviteri.org/2011/10/more-vocabulary-a-bit-of-grammar/}{b} 
\B{Fpìr\-mìl oel fu\-ta ay\-nga na\-tsew tsi\-ve\-'a fì\-'ut.} I was just thinking that you might want to see this.~\LINK{http://forum.learnnavi.org/language-updates/the-evidential-at-last/msg106055/\#msg106055}{ln} 
(\emph{a. with} \B{san … sìk}) \B{Lu su\-te a fpìl san oe\-ri tì\-rey tì\-sraw si, ha tì\-rey\-ti nge\-yä oel tì\-sraw sey\-ka\-syi nì\-teng … sìk.} There are people who think, `My life is hurting, so I will make your life hurt as well ….'~\LINK{http://naviteri.org/2020/06/mi-tanlokxe-oeya-srr-afpxamo-terrible-days-in-my-country/}{b} 
(\emph{b. conv.}) \B{fpìl oe …} `I think …' \B{Fpìl oe, fì\-vur za\-yaw\-prr\-te' ay\-nga\-ne!} I think, this story will be enjoyable for you all!~\LINK{http://naviteri.org/2020/04/keltrrtrra-tieylan-an-unusual-friendship/}{b} \\  
\textsc{Derivations:} \ding{43} \on{1}.~\HL{sAfpIl}{\B{säfpìl}}, idea, thought.
\on{2}.~\HL{fpIlfya}{\B{fpìlfya}}, thought pattern, way of thinking.
\on{3}.~\HL{kanfpIl}{\B{kanfpìl}}, concentrate, focus one's attention. 
\on{4}.~\HL{fpivIl}{\B{fpi\-vìl}}, ``hmmm \ldots, let me think''.
\on{5}.~\HL{txakrrfpIl}{\B{txa\-krr\-fpìl}}, consider, ponder.

% fpìlfya
\hi
\HT{fpIlfya}{\B{\cp{fpìl}fya}} \I{["fpIl.fja]} \emph{n.} thought pattern, way of thinking. \B{Slä sa'\-nok\-ì\-ri fpìl\-fya lu ke\-teng.} But mother's way of thinking is different.~\LINK{http://naviteri.org/2010/07/ting-mikyun-fte-tslivam-listening-comprehension-1-3/}{b} 
(\ding{43}~\HL{fpIl}{\B{fpìl}}, \HL{fya'o}{\B{fya\-'o}}) \\ 
\textsc{Derivations:} \ding{43} \on{1}.~\HL{snafpIlfya}{\B{snafpìlfya}}, philosophy.

% fpom
\hi
\HT{fpom}{\B{fpom}} \I{[fpom]} \emph{n.} well\hyp{}being, happiness; peace.
\B{Pxìm lu tì\-new le\-hawng kxu\-tu fpom\-ä.} Excessive desire is often the enemy of peace.~\LINK{http://naviteri.org/2014/11/twenty-before-the-holidays/}{b}
\B{Oey fpom\-it sweyn!} Leave me alone! [lit.: do not disturb my peace]~\LINK{http://naviteri.org/2021/09/aawa-ayliu-amip-a-few-new-words/}{b} 
(\emph{a. idiom}) \B{Nga\-ru lu fpom srak?} How are you? [lit.: Do you have well\hyp{}being?]~\LINK{http://wiki.learnnavi.org/index.php?title=Corpus\#Vanity_Fair}{wiki} \\
\textsc{Derivations:} \ding{43} \on{1}.~\HL{lefpom}{\B{lefpom}}, happy. 
\on{2}.~\HL{fpomtokx}{\B{fpomtokx}}, (physical) health. 
\on{3}.~\HL{fpomron}{\B{fpomron}}, (mental) health.

% fpomron
\hi
\HT{fpomron}{\B{fpom\cp{ron}}} \I{[fpom."Ron]} \emph{n.} (mental) health, well\hyp{}being.
\B{Po\-ri fpom\-tokx sì fpom\-ron yo'.} His physical and mental health are perfect.~\LINK{http://naviteri.org/2014/10/tson-si-fpomron-obligation-and-mental-health/}{b}
(\ding{43}~\B{fpom}, \B{ron\-sem}) \\
\textsc{Derivations:} \ding{43} \on{1}.~\HL{lefpomron}{\B{lefpomron}}, (mentally) healthy. 
\on{2}.~\HL{kelfpomron}{\B{kelfpomron}}, (mentally) unhealthy. 
\on{3}.~\HL{fpomronga'}{\B{fpomronga'}}, (mentally) healthful. 
\on{4}.~\HL{kefpomronga'}{\B{kefpomronga'}}, (mentally) unhealthful. 

% fpomronga'
\hi
\HT{fpomronga'}{\B{fpom\cp{ro}nga'}} \I{[fpom."Ro.NaP]} \emph{adj., nfp.} (mentally) healthful.
\B{Fwa lawk ay\-sì\-'efu\-ti ay\-ey\-lan\-kip lu fpom\-ro\-nga'.} It's healthy among friends to discuss feelings.~\LINK{http://naviteri.org/2014/10/tson-si-fpomron-obligation-and-mental-health/}{b}
(\ding{43}~\HL{fpomron}{\B{fpom\-ron}}, \ding{214}~\HL{kefpomronga'}{\B{ke\-fpom\-ro\-nga'}})

% fpomtokx
\hi
\HT{fpomtokx}{\B{fpom\cp{tokx}}} \I{[fpom."tok']} \emph{n.} (physical) health. 
\B{Po\-ri fpom\-tokx sì fpom\-ron yo'.} His physical and mental health are perfect.~\LINK{http://naviteri.org/2014/10/tson-si-fpomron-obligation-and-mental-health/}{b}
\B{Sìl\-pey oe tsnì fì\-zì\-sìt\-ìl ay\-nga\-ru za\-mì\-ye\-vu\-nge txan\-a fpom\-it sì fpom\-tokx\-it.} I hope this year will bring you much happiness and health.~\LINK{http://naviteri.org/2013/12/mipa-zisit-lefpom-happy-new-year/}{b}
\B{Nge\-yä fpom\-tokx\-ìri fì\-tì\-fkey\-tok a\-fkxa\-ra\-nga' lu le\-hrr\-ap.} This stressful situation is dangerous to your health.~\LINK{http://naviteri.org/2018/03/100a-liu-amip-64-new-words-part-1/}{b}
(\ding{43}~\B{fpom}, \B{tokx}) \\
\textsc{Derivations:} \ding{43} \on{1}.~\HL{lefpomtokx}{\B{lefpomtokx}}, (physically) healthy. 
\on{2}.~\HL{kelfpomtokx}{\B{kelfpomtokx}}, (physically) unhealthy. 
\on{3}.~\HL{fpomtokxnga'}{\B{fpomtokxnga'}}, (physically) healthful.  
\on{4}.~\HL{kefpomtokxnga'}{\B{kefpomtokxnga'}}, (physically) unhealthful.

% fpomtokxnga'
\hi
\HT{fpomtokxnga'}{\B{fpom\cp{tokx}nga'}} \I{[fpom."tok'.NaP]} \emph{adj., nfp.} (physically) healthful.
\B{Tsat rä\-'ä yi\-vom! Ke lu fpom\-tokx\-nga'.} Don't eat that. It's not healthful  (i.e., it will make you unhealthy).~\LINK{http://naviteri.org/2014/10/tson-si-fpomron-obligation-and-mental-health/}{b}
(\ding{43}~\HL{fpomtokx}{\B{fpom\-tokx}}, \ding{214}~\HL{kefpomtokxnga'}{\B{ke\-fpom\-tokx\-nga'}})

% fpxafaw
\hi
\HT{fpxafaw}{\B{\cp{fpxa}faw}} \I{["fp'a.faw]} \emph{n., fauna} medusa, \emph{electromedusa aerae}. 

% fpxamo
\hi 
\HT{fpxamo}{\B{\cp{fpxa}mo}} \I{["fp'a.mo]} \emph{adj.} terrible, horrible, awful. 
\B{Maw\-krr\-a fko lie so\-li tì\-len\-ur a\-fpxa\-mo fì\-txan, tì\-rey ke lu teng kaw\-krr.} After experiencing such a terrible event, life is never the same.~\LINK{http://naviteri.org/2018/08/fmawnti-stolawm-srak-have-you-heard-the-news/}{b} 
\B{Maw fpxa\-mo\-a kin\-trr a\-fpxa\-mo mì tan\-lo\-kxe oe\-yä …} After a truly horrible week in my country … \LINK{http://naviteri.org/2022/05/results-of-the-tttc/}{b}
(\ding{214}~\HL{kosman}{\B{kosman}}) \\
\textsc{Derivations:} \ding{43} \on{1}.~\HL{tIfpxamo}{\B{tìfpxamo}}, horror. 
\on{2}.~\HL{nIfpxamo}{\B{nìfpxamo}}, horribly, terribly, awfully.

% fpxäkìm
\hi
\HT{fpxAkIm}{\B{\cp{fpxä}kìm}} \I{["fp'\ae.kIm]} \emph{vin.} (\emph{a.}) enter (\B{ìlä}, through).   
\B{Lu mek\-tseng a tsun fpxi\-vä\-kìm hì\-'ang tsaw\-ìlä.}  There's a gap through which insects can get in.~\LINK{http://naviteri.org/2013/04/wheres-the-bathroom-and-other-useful-things/}{b} 
\B{Fpxo\-lä\-kìm fo ìlä rawng a\-hì\-'i.} They entered through a small doorway.~\LINK{http://naviteri.org/2016/06/mrrvola-lifyavi-amip-forty-new-expressions/}{b} 
(\emph{b. of astronomical bodies rising}) \B{Tsaw\-ke fpxe\-rä\-kìm nem\-fa taw.} \emph{or} \B{Tsaw\-ke fpxe\-rä\-kìm.} The sun is rising.~\LINK{http://naviteri.org/2012/01/more-additions-to-the-lexicon/}{b} (\ding{214}~\HL{hum}{\B{hum}})

% fra-
\hi
\HT{fra}{\B{fra-}} \I{[fRa.]} \emph{pref. on nouns} every, all. 
\B{Po\-ri wem\-tswo fra\-tsam\-si\-yur ro\-lo\-'a nì\-txan.} His ability to fight greatly impressed all the warriors.~\LINK{http://naviteri.org/2012/03/spring-vocabulary-part-2/}{b} 
\B{Fra\-syu\-ra\-ti fkol za\-sro\-lìn nì\-'aw ul\-te trr\-o ze\-ne tey\-ki\-vä\-txaw.} All energy is only borrowed, and one day it will have to be given back.~\LINK{http://naviteri.org/2012/03/spring-vocabulary-part-1/}{b}

% fra'u
\hi
\HT{fra'u}{\B{\cp{fra}'u}} \I{["fRa.Pu]} (\textsc{rn}: \B{frau} \I{["fRa.u]}) \emph{n.} everything. 
\B{Tse \dots za\-'u fra\-'u ne tu\-te lem\-wey\-pey, ke\-fyak?} Well, everything comes to the patient people, doesn't it?~\LINK{http://naviteri.org/2012/03/spring-vocabulary-part-2/}{b} 
\B{Sra\-ke tso\-lun fra\-'ut tsli\-vam kxey\-ey\-lu\-ke?} Could you understand everything without error?~\LINK{http://naviteri.org/2012/10/txon-eywaevenga-text-and-translation/}{b} 
\B{Fra\-'uä ran ngä\-pop fa fra\-syon tseyä.} The \emph{ran} of each thing arises from the totality of its attributes.~\LINK{http://naviteri.org/2012/10/mipa-vospxi-mipa-ayliu-new-words-for-the-new-month/}{b}
(\emph{a. coll. shortened to} \B{fraw}) \B{'U\-pe ke zo?--Fraw mì la\-'ang.} What's the matter?--Everything is screwed up.~\LINK{http://naviteri.org/2020/07/vospxivol-lefpom-happy-august/}{b}
(\ding{43}~\HL{'u}{\B{'u}}) \\
\textsc{Derivations:} \ding{43} \on{1}.~\HL{frawzo}{\B{frawzo}}, all is well.

% frakrr
\hi
\HT{frakrr}{\B{\cp{fra}krr}} \I{["fRa.k\s{r}]} \textsc{i}. \emph{n.} every time. \\
\textsc{ii}. \emph{adv.} always. \B{Pol\-txe Ney\-ti\-ril ay\-lì\-'ut a fra\-krr 'ok se\-yä la\-yu oer.} Neytiri said something that I will always remember.~\LINK{http://thereadinginpublicproject.blogspot.de/2010/05/paul-frommer.html}{pr}
(\ding{43}~\HL{krr}{\B{krr}})

% fralo
\hi
\HT{fralo}{\B{\cp{fra}lo}} \I{["fRa.lo]} \emph{adv.} every time, each time. 
\B{Fra\-lo poe pol\-txe san ke\-he.} Each time she said `no'.~\LINK{http://forum.learnnavi.org/language-updates/txelanit-hivawl/}{ln}
(\ding{43}~\HL{alo}{\B{alo}}) 

% franse
\hi
\HT{franse}{\B{\cp{Fran}se}} \I{["fRan.sE]} \emph{n.}, \textsc{lw} French, the French language. \\
\textsc{Derivations:} \ding{43} \on{1}.~\HL{nIfranse}{\B{nìFranse}}, in French.

% frapo
\hi
\HT{frapo}{\B{\cp{fra}po}} \I{["fRa.po]} (\emph{gen.}~\B{frapeyä}) \emph{pn.} everyone, everybody. 
\B{Fra\-po ne mì\-fa!} Everbody inside!~\LINK{http://naviteri.org/2012/07/meetings-waterfalls-and-more/}{b} 
\B{tslam fra\-pol fu\-ta slu po O\-lo'\-eyk\-tan a\-mip}, everyone understood that he had become the new Clan Leader.~\LINK{http://naviteri.org/2012/06/spring-vocabulary-part-3/}{b} 
\B{Po\-leng Ney\-ti\-ril hang\-vur\-it a fra\-pot hey\-kang\-ham.} Neytiri told a joke that made everyone laugh.~\LINK{http://naviteri.org/2012/01/mipa-zisit-ayliu-amip-new-words-for-the-new-year/}{b} 
\B{Fì\-tì\-hawl\-ì\-ri fe'\-ran law lä\-ngu fra\-por.} Unfortunately the flawed nature of this plan is obvious to everyone.~\LINK{http://naviteri.org/2012/10/mipa-vospxi-mipa-ayliu-new-words-for-the-new-month/}{b} 
(\ding{43}~\HL{po}{\B{po}})

% frato
\hi
\HT{frato}{\B{\cp{fra}to}} \I{["fRa.to]} \emph{part.} (\emph{superlative marker}) than all. 
(\emph{a. with (scalable) adjectives and adverbs}) 
\B{Oel hu Txe\-wì trr\-am na'\-rìng\-it tar\-mok, tso\-le\-'a syep\-tu\-tet a\-tsawl fra\-to mì sì\-rey.} Yesterday I was with Txewì in the forest and we saw the biggest Trapper I've ever seen.~\LINK{http://wiki.learnnavi.org/index.php?title=Corpus\#New_York_Times}{wiki} 
\B{New oe ti\-vìng ay\-ngar lì\-'ut a tì\-'e\-fu\-mì oe\-yä lu lor fra\-to mì lì'\-fya le\-Na'\-vi.} I want to give you the word that, in my opinion, is the most beautiful in the Na'vi language.~\LINK{http://forum.learnnavi.org/language-updates/trr-rrtaya/}{ln} 
\B{Pol ye\-rik\-it ta\-ron nìl\-tsan fra\-to.} He hunts yerik the best of all.~\LINK{http://forum.learnnavi.org/language-updates/confirmations-\%28comparisons-gerund\%29/}{ln} 
(\emph{b. with verbs}) \B{Tsran\-ten fra\-to a syon tsam\-si\-yu\-ä lu tì\-tstew.} The most important characteristic of a warrior is bravery.~\LINK{http://naviteri.org/2012/10/mipa-vospxi-mipa-ayliu-new-words-for-the-new-month/}{b} 
\B{Tì\-'o'\-ì\-ri pe\-u sunu ngar fra\-to?} What is your favorite way to have fun?~\LINK{http://naviteri.org/2010/09/getting-to-know-you-part-3/}{b} 
\B{Nga\-ru le\-i\-o\-a\-e si oe fra\-to, ma 'ey\-lan.} I respect you the most of all, friend.~\LINK{http://naviteri.org/2012/02/trr-asawnung-lefpom-happy-leap-day/}{b} 
\B{pi\-za\-yu Ney\-ti\-ri\-yä txan\-rar\-mo\-'a fra\-to}, Ney\-ti\-ri's ancestor was the most famous.~\LINK{http://naviteri.org/2012/03/spring-vocabulary-part-1/}{b}

% fratrr
\hi
\HT{fratrr}{\B{fra\cp{trr}}} \I{[fRa."t\s{r}]} \emph{adv.} every day.
\B{Lì'\-fya\-ri le\-Na'\-vi nga tsan\-'e\-re\-i\-ul fra\-trr.} I'm delighted that your Na'vi is improving every day.~\LINK{http://naviteri.org/2013/03/tsanerul-feerul-getting-better-getting-worse/}{b}
(\ding{43}~\HL{trr}{\B{trr}})

% fratxon
\hi
\HT{fratxon}{\B{fra\cp{txon}}} \I{[fRa."t'on]} \emph{adv.} every night. 
(\ding{43}~\HL{txon}{\B{txon}})

% fratseng
\hi
\HT{fratseng}{\B{\cp{fra}tseng}} \I{["fRa.\t{ts}EN]} \emph{adv.} everywhere. 
\B{Tsun Txil\-te pam\-rel\-it ivi\-nan; ta\-fral puk\-ot a\-nutx mu\-nge fra\-tseng.} Txilte knows how to read; therefore she brings a thick book wherever she goes.~\LINK{http://naviteri.org/2011/11/\%E2\%80\%99a\%E2\%80\%99awa-ayli\%E2\%80\%99u-amip-ni\%E2\%80\%99aw-\%E2\%80\%94-a-few-new-words-only/}{b}
(\ding{43}~\HL{tsenge}{\B{tse\-nge}})


% fraw
\hi
\HT{fraw}{\B{fraw}} \I{[fRaw]}~:~\HL{fra'u}{\B{fra'u}}

% frawzo
\hi
\HT{frawzo}{\B{fraw\cp{zo}}} \I{[fRaw."zo]} (\emph{idiom}) all is well, everything is fine, everything is \textsc{ok}. 
\B{Sìl\-pey Txe\-wì tsnì tsi\-vun set fko pi\-vll\-txe san fraw\-zo.} Txewì hopes that now one can say, `everything is fine'.~\LINK{http://naviteri.org/2010/07/ting-mikyun-fte-tslivam-listening-comprehension-1-3/}{b} 
\B{Nge\-yä tsa\-sä\-ngä\-'än\-tsyìp\-ì\-ri set fra\-wzo srak?} Have you recovered from being down for a while?~\LINK{http://naviteri.org/2013/02/vospxi-ayol-posti-apup-short-post-for-a-short-month/}{b}
(\ding{43}~\HL{fra'u}{\B{fra\-'u}}, \HL{zo}{\B{zo}})

% fray+
\hi
\HT{fray}{\B{fray}}\texttt{+} \I{[fRaj.]} (\emph{pref. to nouns}) every, all of \LINK{https://forum.learnnavi.org/language-updates/navi-details-from-avatarmeet-2013/}{ln}. 
\B{Kay\-mo zo\-la\-'u fray\-ul\-txa\-tu ne kel\-ku moe\-yä.} One evening all the members of the meeting came to our house.~\LINK{http://naviteri.org/2014/09/stxeli-alor-the-text/}{b} 
(\ding{43}~\HL{fra}{\B{fra-}}, \HL{ay}{\B{ay}}\texttt{+})

% frir
\hi 
\HT{frir}{\B{frir}} \I{[fRiR]} \emph{n.} layer. 
\B{Tskxe\-pay\-ri lu frir a\-flì sìn 'o\-ra.} There's a thin layer of ice on the lake.~\LINK{http://naviteri.org/2018/07/ta-sulfatu-a-ayliu-niul-more-words-from-our-experts/}{b}
\B{Fay\-frir le\-tskxe lor lu nì\-txan.} These stone layers are very beautiful.~\LINK{http://naviteri.org/2018/07/ta-sulfatu-a-ayliu-niul-more-words-from-our-experts/}{b} \\
\textsc{Derivations:} \ding{43} \on{1}.~\HL{lefrir}{\B{lefrir}}, layered. 
\on{2}.~\HL{nIfrir}{\B{nìfrir}}, in layers.

% frìp
\hi
\HT{frIp}{\B{frìp}} \I{[fRIp\textcorner]} \emph{vtr.} bite. \B{Frìp, frìp, ma Evi\-tsyìp, To\-ruk\-ä tsawl\-a kxa!} The big mouth of the Great Leonopteryx bites, bites, children.~\LINK{http://forum.learnnavi.org/language-updates/learnings-from-the-siatlla-song-translations/msg473566/\#msg473566}{ln} \\
\textsc{Derivations:} \ding{43} \on{1}.~\HL{sAfrIp}{\B{säfrìp}}, a bite.

% frrfen
\hi
\HT{frrfen}{\B{\cp{frr}fen}} \I{["f\s{r}.fEn]} (\emph{with} ‹\B{er}›: \B{frr\-fen}) \emph{vtr.} visit. 
\B{Oe\-ri so\-la\-lew wum zì\-sìt °a14 a krr, fo\-lrr\-fen spo\-not alo a\-'aw\-ve.} When I was about 12 years old I visited an island for the first time.~\LINK{http://naviteri.org/2012/10/mipa-vospxi-mipa-ayliu-new-words-for-the-new-month/}{b} \\
\textsc{Derivations:} \ding{43} \on{1}.~\HL{frrtu}{\B{frrtu}}, visitor.
\on{2}.~\HL{sAfrrfen}{\B{sä\-frr\-fen}}, a visit.

% frrtu
\hi
\HT{frrtu}{\B{\cp{frr}tu}} \I{["f\s{r}.tu]} (\textsc{rn}: \B{frrtu} \I{["f\s{r}.tu]}) \emph{n.} visitor, guest. 
\B{Sko frr\-tu, fra\-krr mì hel\-ku nge\-yä lu oe\-ru ho\-an nì\-txan, ma tsmuk.} I always feel very comfortable as a guest in your home, brother.~\LINK{http://naviteri.org/2015/08/aylifyavi-lereyfya-2-cultural-terms-2-and-more/}{b}
\B{Rä\-'ä fmi\-vi li\-vok fu em\-ki\-vä ay\-ekxan\-it a fkol ngo\-lop fpi sì\-kxu\-ke ay\-frr\-tu\-ä sì ay\-ioang\-ä.} Do not attempt to approach or cross any barriers designed for Guest and animal safety.~\LINK{http://naviteri.org/2017/06/ayliu-a-ta-eywaeveng-kifkey-uniltirantokxa-words-from-disneys-pandora-the-world-of-avatar/}{b}
(\ding{43}~\HL{frrfen}{\B{frr\-fen}})

% fta
\hi
\HT{fta}{\B{fta}}\textsuperscript{\on{1}} \I{[fta]} \textsc{i}.~\emph{n.} knot. \\
\textsc{ii}.~\HT{fta si}{\B{fta si}} \I{[fta si]} \emph{vin.} knot, make or tie a knot. \B{Txur\-tel\-mì fo fta so\-li fte\-ke ka tsyokx fwi\-vi.} They tied a knot in the rope so it wouldn't slip through their hands.~\LINK{http://naviteri.org/2013/04/wheres-the-bathroom-and-other-useful-things/}{b}

% fta
\hi
\B{fta}\textsuperscript{\on{2}} : \HL{futa}{\B{futa}}

% ftang
\hi
\HT{ftang}{\B{ftang}} \I{[ftang]} \emph{vin., modal} stop (to do s.t.\slash an action). \B{Ik\-ran\-ì\-ri krra hu tu\-te tsa\-heyl si, ftang li\-vu yrr.} When an ikran bonds with a person, it ceases to be wild.~\LINK{http://naviteri.org/2013/01/awvea-posti-zisita-amip-first-post-of-the-new-year/}{b} \B{Hay\-a yak\-ro fti\-vang.~Sa\-lew rä\-'ä.} Stop at the next fork.~Do not proceed further.~\LINK{http://naviteri.org/2013/09/on-si-salewfya-shapes-and-directions/}{b} \\
\textsc{Derivations:} \ding{43} \on{1}.~\HL{tIftang}{\B{tìftang}}, stopping, glottal stop. 
\on{2}.~\HL{nIlkeftang}{\B{nìlkeftang}}, continuously, incessantly. 
\on{3}.~\HL{ftanglen}{\B{ftanglen}}, prevent.

% ftanglen
\hi 
\HT{ftanglen}{\B{ftang\cp{len}}} \I{[ftaN."lEn]} \emph{vtr.} prevent. 
\B{Tsran\-ten nì\-txan fwa ftang\-len aw\-ngal fu\-ta fì\-sä\-spxin vi\-vi\-rä.} It's very important that we prevent this disease from spreading.~\LINK{http://naviteri.org/2020/05/aawa-liu-amip-si-vurway-alor-a-few-new-words-and-a-beautiful-poem/}{b} 
(\ding{43}~\HL{ftang}{\B{ftang}}, \HL{len}{\B{len}}) \\
\textsc{Derivations:} \ding{43} \on{1}.~\HL{tIftanglen}{\B{tìftanglen}}, prevention.

% ftawnemkrr
\hi
\HT{ftawnemkrr}{\B{ftaw\cp{nem}krr}} \I{[ftaw."nEm.k\s{r}]} \emph{n.} past. 
\B{txo tsa\-kem li\-ven mì ftaw\-nem\-krr}, if it happens in the past.~\LINK{http://naviteri.org/2013/01/awvea-posti-zisita-amip-first-post-of-the-new-year/comment-page-1/\#comment-1910}{b} 
\B{Ul\-te ftaw\-nem\-krr\-ìri, ‹imv› tam nì\-teng.} And concerning the past ‹imv› suffices as well.~\LINK{http://naviteri.org/2013/01/awvea-posti-zisita-amip-first-post-of-the-new-year/comment-page-1/\#comment-1910}{b}
(\ding{43}~\HL{ftem}{\B{ftem}}, \HL{krr}{\B{krr}})

% ftär
\hi
\HT{ftAr}{\B{ftär}} \I{[ft\ae{}R]} \emph{adj.} left (\emph{not right}). 
\B{Lo\-lu po\-ru txukx\-a skxir a ftu 'et\-naw a\-ski\-en ne me\-pun a\-ftär.} It had a deep wound from its right shoulder to its two left arms.~\LINK{http://naviteri.org/2020/04/awa-tskxekengtsyip-a-mikyunfpi-niul-one-more-little-listening-exercise/}{b}
(\ding{214}~\HL{skien}{\B{ski\-en}}) \\
\textsc{Derivations:} \ding{43} \on{1}.~\HL{ftArpa}{\B{ftärpa}}, left side.   
\on{2}.~\HL{nIftAr}{\B{nìftär}}, on the left.

% ftärpa
\hi
\HT{ftArpa}{\B{\cp{ftär}pa}} \I{["ft\ae{}R.pa]} \emph{n.} left side. \B{Mìn ne ftär\-pa.} Turn to the left.~\LINK{http://forum.learnnavi.org/language-updates/ftarpa-and-skienpa-added-to-dict/msg570362/\#msg570362}{b} 
(\ding{214}~\HL{skiempa}{\B{ski\-em\-pa}}, \ding{43}~\HL{ftAr}{\B{ftär}}, \HL{nIftAr}{\B{nì\-ftär}})

% fte
\hi
\HT{fte}{\B{fte}} \I{[ftE]} \emph{conj.} (\emph{requires} ‹\B{iv}›) so that, (in order) to. 
\B{Aw\-nga rol fte ki\-vame}, we sing to \textsc{see}.~\LINK{http://www.pandorapedia.com/navi/music/ritual_music}{pp}  
\B{He\-tu\-wong\-ìl aw\-nge\-yä swo\-tut ska\-'a, fte kll\-ki\-vu\-lat ke\-ru\-sey\-a tskxet.} The aliens destroy our sacred place to dig up dead rock.~\LINK{http://forum.learnnavi.org/language-updates/participles/}{{ln\slash a}\textsuperscript{c}} 
\B{New oel fu\-ta u\-pxa\-re oe\-yä lu\-ke key\-ey li\-vu nì\-wotx, fte eyawr\-a sì\-ke\-nong\-it ti\-vìng su\-te\-ru.} I want my messages to be without errors, in order to give correct examples to people.~\LINK{http://forum.learnnavi.org/language-updates/verb-for-write/msg175592/\#msg175592}{ln} 
(\ding{214}~\HL{fteke}{\B{fte\-ke}})

% fteke
\hi
\HT{fteke}{\B{\cp{fte}ke}} \I{["ftE.kE]} \emph{conj.} (\emph{requires} ‹\B{iv}›) so that not, lest. 
\B{Nari so\-li ay\-oe fte\-ke nì\-hawng li\-vok.} We were careful not to get too close.~\LINK{http://wiki.learnnavi.org/index.php?title=Corpus\#New_York_Times}{wiki} 
\B{Tì\-ran nì\-fnu txan\-txew\-vay fte\-ke ay\-ye\-rik\-ìl aw\-nga\-ti sti\-vawm.} Walk as quietly as possible so the hexapedes won't hear us.~\LINK{http://naviteri.org/2012/10/snewsyea-ftxoza-halowina-a-spooky-halloween/}{b} 
(\ding{214}~\HL{fte}{\B{fte}})

% ftem
\hi
\HT{ftem}{\B{ftem}} \I{[ftEm]} \emph{vtr.} pass by (s.t.). \\
\textsc{Derivations:} \ding{43} \on{1}.~\HL{ftawnemkrr}{\B{ftawnemkrr}}, (the) past.

% ftia
\hi
\HT{ftia}{\B{fti\cp{a}}} \I{[fti."{}a]} \emph{vtr.} study. 
\B{Zì\-sìt\-o a\-mrr fto\-li\-a o\-he, slä ze\-ne fko ni\-vume nì\-txan.} I studied for five years but there is much to learn.~\LINK{http://wiki.learnnavi.org/index.php?title=Na\%27vi_from_Avatar_Movie\#Grace_and_Norm_in_the_Link_room}{wiki\slash a} 
\B{Fti\-a oel lì'\-fya\-ti le\-Na'\-vi nì\-'o' nì\-wotx!} Studying Na'vi is a ton of fun for me.~\LINK{http://naviteri.org/2010/09/getting-to-know-you-part-3/}{b} \\
\textsc{Derivations:} \ding{43} \on{1}.~\HL{tIftia}{\B{tì\-ftia}}, studying.

% ftu
\hi
\HT{ftu}{\B{ftu}} \I{[ftu]} (\textsc{rn}: \B{ftu} \I{[ftu]}) \emph{adp.} from (\emph{direction}). \B{Ftu oe ne\-to rikx, ma \mbox{skxawng}!} Get away from me, moron!~\LINK{http://www.youtube.com/watch?v=ojr_BE8ZmVE}{yt} 
\B{Ftu fì\-tseng ze\-ne hi\-vum.} We gotta get out of here.~\LINK{http://forum.learnnavi.org/language-updates/txing-and-hum/}{ln} 
\B{Nga za\-'u pe\-senge\-ftu\slash tseng\-pe\-ftu?} Where are you from?~\LINK{http://naviteri.org/2010/09/getting-to-know-you-part-2/}{b} 
(\ding{214}~\HL{ne}{\B{ne}}) \\
\textsc{Derivations:} \ding{43} \on{1}.~\HL{ftumfa}{\B{ftum\-fa}}, out of, from inside. 
\on{2}.~\HL{ftuopa}{\B{ftuo\-pa}}, from behind.

% ftue
\hi
\HT{ftue}{\B{\cp{ftu}e}} \I{["ftu.E]} \emph{adj.} easy, simple.   
\B{Sì\-pawm nge\-yä ke lu ftu\-e}, your questions aren't easy.~\LINK{http://wiki.learnnavi.org/index.php?title=Canon\#More_extracts_from_various_emails}{wiki} 
\B{Fì\-pam\-tseo\-tu\-ri ke la\-yu ftue fwa run fkol sä\-rawn\-it a tam.} It won't be easy to find a satisfactory replacement for this musician.~\LINK{http://naviteri.org/2012/01/mipa-zisit-ayliu-amip-new-words-for-the-new-year/}{b}
(\ding{214}~\HL{ngAzIk}{\B{ngä\-zìk}}) \\
\textsc{Derivations:} \ding{43} \on{1}.~\HL{nIftue}{\B{nìftue}}, easily.

% ftumfa
\hi
\HT{ftumfa}{\B{\cp{ftum}fa}} \I{["ftum.fa]} \emph{adp.} out of, from inside. 
\B{Ri\-ti tswo\-lay\-on ftum\-fa slär.} The stingbat flew out of the cave.~\LINK{http://naviteri.org/2013/04/wheres-the-bathroom-and-other-useful-things/}{b} 
\B{Rey\-pay skxir\-ftum\-fa he\-rum.} Blood is coming out of the wound.~\LINK{http://naviteri.org/2013/04/wheres-the-bathroom-and-other-useful-things/}{b} 
(\ding{214}~\HL{nemfa}{\B{nem\-fa}}, \ding{43}~\HL{ftu}{\B{ftu}}, \HL{mIfa}{\B{mì\-fa}})

% ftuopa
\hi 
\HT{ftuopa}{\B{\cp{ftu}opa}} \I{["ftu.o.pa]} \emph{adp.} from behind. 
\B{Sro\-ler fwä\-kì ftu\-opa tskxe.} A mantis appeared from behind a rock.~\LINK{http://naviteri.org/2019/06/50a-liu-amip-40-new-words/}{b} 
(\ding{43}~\HL{ftu}{\B{ftu}}, \HL{uo}{\B{uo}}, \HL{pa'o}{\B{pa'o}})

% ftxavang
\hi
\HT{ftxavang}{\B{\cp{ftxa}vang}} \I{["ft'a.vaN]} \emph{adj.} passionate. \\
\textsc{Derivations:} \ding{43} \on{1}.~\HL{tIftxavang}{\B{tìftxavang}}, passion. 
\on{2}.~\HL{nIftxavang}{\B{nìftxavang}}, passionately. 

% ftxey
\hi
\HT{ftxey}{\B{ftxey}} \I{[ft'Ej]} \textsc{i}. \emph{vtr.} choose. 
\B{ftxey 'aw\-pot a ay\-nga\-kip} choose one among you.~\LINK{http://wiki.learnnavi.org/index.php/Hunting_Song}{wiki} 
\B{Ul\-te ftxo\-ley pol ay\-eyk\-ta\-nay\-ti a lu stum nì\-ftxan kawng na po.} And he has chosen deputies who are almost as bad as he is.~\LINK{http://naviteri.org/2016/12/tengkrr-perahem-zisit-amip-as-the-new-year-arrives/}{b}  \\
\textsc{ii}. \emph{conj.} (\emph{a.}) whether (or). 
\B{Fra\-po, ftxey sngä\-'i\-yu ftxey tsul\-fä\-tu, tsì\-ye\-vun fì\-tse\-nge ri\-vun 'u\-ot le\-sar.} Everyone, whether beginner or master, can find something useful here.~\LINK{http://naviteri.org/2010/06/first-post/}{b} 
(\emph{b.}) whether (or not). \B{Pi\-veng oer ftxey nga new ri\-vey fu\-ke.} Tell me whether you want to live (or not)?~\LINK{http://forum.learnnavi.org/language-updates/if-vs-whether-ftxey-fuke/msg156629/\#msg156629}{ln} 
\B{Ftxey nga za\-'u fu\-ke?} Are you coming (or not)?~\LINK{http://forum.learnnavi.org/language-updates/if-vs-whether-ftxey-fuke/msg156629/\#msg156629}{ln} \\
\textsc{Derivations:} \ding{43} \on{1}.~\HL{tIftxey}{\B{tìftxey}}, choice.

% ftxì
\hi
\HT{ftxI}{\B{ftxì}} \I{[ft'I]} \textsc{i}.~\emph{n.} tongue.
\B{Nga\-ri ftxì vi\-var le\-fpom\-tokx li\-vu.} May your tongue remain healthy.~\LINK{http://naviteri.org/2014/10/tson-si-fpomron-obligation-and-mental-health/comment-page-1/\#comment-10624}{b}  \\  
\textsc{ii}.~\B{tìng ftxì} \I{[tIN ft'I]} \emph{vin.} taste, take a taste. 
(\ding{43}~\HL{syam}{\B{syam}}) \\
\textsc{Derivations:} \ding{43} \on{1}.~\HL{ftxIlor}{\B{ftxìlor}}, delicious. 
\on{2}.~\HL{ftxIvA'}{\B{ftxìvä'}}, bad-tasting.

% ftxìlor 
\hi
\HT{ftxIlor}{\B{ftxì\cp{lor}}} \I{[ft'I."loR]} \emph{adj.} delicious, good\hyp{}tasting. 
\B{Tey\-lu a oel ye\-rom lu ftxì\-lor.} The teylu I'm eating is delicious.~\LINK{http://naviteri.org/2012/03/spring-vocabulary-part-1/}{b} 
\B{Fì\-na\-er\-ìri fkan ftxì\-lor.} This drink tastes good.~\LINK{http://naviteri.org/2012/11/renu-ayinanfyaya-the-senses-paradigm/}{b} 
(\ding{214}~\HL{ftxIvA'}{\B{ftxì\-vä'}})

% ftxìvä'
\hi
\HT{ftxIvA'}{\B{ftxì\cp{vä'}}} \I{[ft'I."v\ae{}P]} \emph{adj.} bad-tasting. 
\B{Fì\-na\-er ftxì\-vä' lu nì\-hawng, ha swey\-lu txo ngal tsat kxi\-vukx nì\-win.} This drink tastes horrible, so you'd better swallow it quickly.~\LINK{http://naviteri.org/2011/11/\%E2\%80\%99a\%E2\%80\%99awa-ayli\%E2\%80\%99u-amip-ni\%E2\%80\%99aw-\%E2\%80\%94-a-few-new-words-only/}{b} 
(\ding{214}~\HL{ftxIlor}{\B{ftxì\-lor}})

% ftxozä
\hi
\HT{ftxozA}{\B{ftxo\cp{zä}}} \I{[ft'o."z\ae]} \textsc{i}. \emph{n.} holiday, celebration, happy occasion.  
\B{Ay\-ftxo\-zä Le\-fpom!} Happy Holidays!~\LINK{http://wiki.learnnavi.org/index.php?title=Corpus\#Out_of_Office_AutoReply}{wiki} 
\B{Ftxo\-zä\-ri ay\-lrr\-tok ngar.} Smiles to you on your celebration.~\LINK{http://forum.learnnavi.org/language-updates/good-luck-with-this-one!/}{ln} \\
\textsc{ii}. \HT{ftxozA si}{\B{ftxo\cp{zä} si}} \I{[ft'o."z\ae{} si]} \emph{vin.} celebrate. 
\B{Teng\-krr ftxo\-zä se\-re\-i\-yi aw\-nga, ke tswi\-va' ay\-lom\-tu\-ti ko!} A toast: To absent friends. [lit.: While we are celebrating, let us not forget those who we wish could also be here (but can't).]~\LINK{http://naviteri.org/2011/07/txantsana-ultxa-mi-siatll-great-meeting-in-seattle/}{b} \\
\textsc{Derivations:} \ding{43} \on{1}.~\HL{leftxozA}{\B{leftxozä}}, celebratory.

% ftxulì'u
\hi
\HT{ftxulI'u}{\B{ftxu\cp{lì}'u}} \I{[ft'$\cdot$u."lI.Pu]} \emph{vin., inf.}\on{11} orate, give a speech. 
\B{Krra oe ftxu\-lì\-'u, pxìm 'efu fkxa\-rat nì\-txan.} When I give a speech, I often feel a lot of stress.~\LINK{http://naviteri.org/2018/03/100a-liu-amip-64-new-words-part-1/}{b}
\B{Maw\-krra Tsyeyk ftxo\-lu\-lì\-'u, tslam fra\-pol fu\-ta slu po O\-lo'\-eyk\-tan a\-mip.} After Jake spoke, everyone understood that he had become the new Clan Leader.~\LINK{http://naviteri.org/2012/06/spring-vocabulary-part-3/}{b} \\
\textsc{Derivations:} \ding{43} \on{1}.~\HL{ftxulI'uyu}{\B{ftxulì'uyu}}, orator.   
\on{2}.~\HL{tIftxulI'u}{\B{tìftxulì'u}}, speech-making, public speaking. 
\on{3}.~\HL{sAftxulI'u}{\B{säftxulì'u}}, speech, oration.

% ftxulì'uyu
\hi
\HT{ftxulI'uyu}{\B{ftxu\cp{lì}'uyu}} \I{[ft'u."lI.Pu.ju]} \emph{n.} orator, (public) speaker.
(\ding{43}~\HL{ftxulI'u}{\B{ftxu\-lì\-'u}})

% fu
\hi
\HT{fu1}{\B{fu}} \I{[fu]} \emph{conj.} (\emph{a.}) or. 
\B{Tsun fko tsa\-tse\-nge\-ne ki\-vä nì\-sle\-le fu fa fwa ik\-ran\-it mak\-to nì\-'aw.} You can only get there by swimming or riding an ikran.~\LINK{http://naviteri.org/2011/02/new-vocabulary-ii\%E2\%80\%94part-1/}{b} 
\B{Ke za\-syup lì\-'O\-na ne kxu\-tu a mì\-fa fu a wrr\-pa.} The l'Ona will not perish to the enemy within or the enemy without.~\LINK{http://forum.learnnavi.org/language-updates/ltasygt-with-ke-and-nga/}{ln\slash g} 
(\emph{b. choice statement}) 
\B{Nul\-new oel fì\-'ut fu tsa\-'ut.} I want this or that.~\LINK{http://naviteri.org/2019/09/choice-statements-vs-choice-questions-and-some-insults/}{b}
(\emph{c. choice question}) 
\B{Nul\-new ngal fu fì\-'ut fu tsa\-'ut?} Do you want this or that?~\LINK{http://naviteri.org/2019/09/choice-statements-vs-choice-questions-and-some-insults/}{b} 
\B{New nga fu yi\-vom fu ni\-väk?} Do you want to eat or drink?~\LINK{http://naviteri.org/2019/09/choice-statements-vs-choice-questions-and-some-insults/\#comment-30389}{b}
(\emph{d. non\hyp{}choice question}) 
\B{Sra\-ke nul\-new ngal fì\-'ut fu tsa\-'ut?} Do you want this or that?~\LINK{http://naviteri.org/2019/09/choice-statements-vs-choice-questions-and-some-insults/}{b} \\
\textsc{Derivations:} \ding{43} \on{1}.~\HL{fuke}{\B{fuke}}, or not.

% -fu
\hi
\HT{fu2}{\B{-fu}} \I{[.fu]} \emph{num., suf.}~six (\emph{rest form of} \ding{43}~\HL{pukap}{\B{pu\-kap}} \emph{to form higher numbers}). \B{vofu}, \on{14} (decimal), \on{16} (octal). \B{mevofu}, \on{22} (decimal), \on{26} (octal); etc.

% fuke
\hi
\HT{fuke}{\B{fu\cp{ke}}} \I {[fu."kE]} \emph{conj.} or not. \B{Oeng swey\-lu txo ki\-vä fu\-ke?} What do you think? Should we go or not?~\LINK{http://naviteri.org/2012/07/fivospxiya-aylifyavi-amip-this-months-new-expressions-pt-2/}{b} \B{Ftxey nga za\-'u fu\-ke?} Are you coming (or not)?~\LINK{http://forum.learnnavi.org/language-updates/if-vs-whether-ftxey-fuke/msg156629/\#msg156629}{ln} 
(\ding{43}~\HL{ftxey}{\B{ftxey}}, \HL{fu}{\B{fu}})

% fula
\hi
\HT{fula}{\B{\cp{fu}la}}, \B{fì\-\cp{'ul} a}, \emph{or} \B{a fì\-\cp{'ul}} \I{["fu.la]} \emph{conj.} that (\emph{to nominalize an agentive subordinate clause}) 
\B{Fu\-la ho\-ren sì ay\-lì'\-fya\-vi a\-mip le\-sar sa\-yi me\-ngar oe\-ti nit\-ram 'ey\-ke\-fu nì\-txan.} That the rules and new expressions are useful to you two makes me feel very happy.~\LINK{http://naviteri.org/2012/06/spring-vocabulary-part-3/comment-page-1/\#comment-1515}{b} 
\B{Vì\-ming\-kap oe\-ti fu\-la poe ke li ke pol\-txe san oe za\-sya\-'u.} It just occurred to me that she hasn't yet said she's coming.~\LINK{http://naviteri.org/2011/09/\%E2\%80\%9Cby-the-way-what-are-you-reading\%E2\%80\%9D/}{b} 
\B{Fay\-lì\-'u su\-nu ay\-ngar a fì\-'ul oe\-ti nit\-ram sley\-ke\-i\-u nì\-txan.} That you all like these words makes me very happy.~\LINK{http://naviteri.org/2012/03/spring-vocabulary-part-1/comment-page-1/\#comment-1443}{b} 
(\ding{43}~\HL{fI'u}{\B{fì\-'u}}, \emph{nominalizer of} \ding{43}~\HL{tsa'u}{\B{tsa\-'u}}, \emph{i.e.} \HL{tsala}{\B{tsa\-la}} \emph{exists but is rarely in use})

% fura
\hi
\HT{fura}{\B{\cp{fu}ra}}, \B{fì\-\cp{'ur} a}, \emph{or} \B{a fì\-\cp{'ur}} \I{["fu.Ra]} \emph{conj.} that (\emph{to introduce a dative subordinate clause}) 
\B{Fì\-'u\-pxa\-re ka\-ngay si fu\-ra po za\-ya\-'u.} This message confirms that he will come.~\LINK{http://naviteri.org/2020/04/aawa-u-amip-a-few-new-things/}{b}
(\ding{43}~\HL{fI'u}{\B{fì\-'u}}, \HL{tsara}{\B{tsa\-ra}} \emph{exists but is rarely in use})

% furia
\hi
\HT{furia}{\B{\cp{fu}ria}}, \B{fì\-\cp{'u}\-ri a}, \emph{or} \B{a fì\-\cp{'u}\-ri} \I{["fu.Ri.a]} \emph{conj.} that (\emph{to nominalize a topic subordinate clause}) 
\B{Fu\-ri\-a nì\-'Ìng\-lì\-sì pam\-rel si\-vi, oe\-ru txoa li\-vu.} Forgive me that I write in English.~\LINK{http://forum.learnnavi.org/language-updates/verb-for-write/msg139677/\#msg139677}{ln}
\B{Ngal oe\-yä 'u\-pxa\-ret ay\-su\-te\-ru fpo\-le' a fì\-'u\-ri, nga\-ru ir\-ayo sei\-yi oe nì\-txan!} Concerning the fact, that you sent my message to the people, I thank you very much!~\LINK{http://wiki.learnnavi.org/index.php?title=Canon\#Extracts_from_various_emails}{wiki} 
(\ding{43}~\HL{fI'u}{\B{fì\-'u}}, \emph{nominalizer of} \ding{43}~\HL{tsa'u}{\B{tsa\-'u}}, \emph{i.e.} \HL{tsaria}{\B{tsa\-ri\-a}} \emph{exists but is rarely in use})

% futa
\hi
\HT{futa}{\B{\cp{fu}ta}}, \B{fì\-\cp{'ut} a}, \emph{or} \B{a fì\-\cp{'ut}} \I{["fu.ta]} \emph{or coll.} \I{[fta]} \emph{conj.} that (\emph{to nominalize a patientive subordinate clause}) 
(\emph{a. with vtr.}) \B{Ul\-te omum oel fu\-ta tì\-fya\-wìn\-txu\-ri oe\-yä pe\-rey ay\-nga nì\-wotx.} And I know that you all are waiting for my guidance.~\LINK{http://forum.learnnavi.org/news-announcements/a-response-from-paul-frommer!/}{ln} 
(\emph{b. with transitive modals}, requires ‹\B{iv}› in the dependent verb) \B{\mbox{Oel} new fu\-ta ta\-ron\-yu ki\-vä.} I want the hunter to go.~\LINK{http://forum.learnnavi.org/language-updates/frommerian-email/}{b} 
\B{Po\-el sto\-la\-tso fu\-ta me\-fo ti\-va\-ron tsa\-ha'\-ngir.} She must have refused to hunt that afternoon.~\LINK{http://forum.learnnavi.org/language-updates/sto-has-the-same-syntax-as-new/msg567034/\#msg567034}{b} 
\B{(oel) nul\-new fu\-ta si\-vop oeng nì\-txi\-lu\-ke}, (I) prefer to travel leisurely.~\LINK{http://naviteri.org/2011/10/more-vocabulary-a-bit-of-grammar/}{b}
(\emph{c. rarely written as} \B{fta}) \B{Ro\-lun oel olo'\-ti a 'e\-fu fta oe man tsa\-tseng\-mì.} I've found a community where I feel I belong.~\LINK{http://naviteri.org/2022/06/solalew-maul-zisita-half-the-year-is-over/}{b} 
(\ding{43}~\HL{fI'u}{\B{fì\-'u}}, \emph{nominalizer of} \ding{43}~\HL{tsa'u}{\B{tsa\-'u}}, \emph{i.e.} \HL{tsata}{\B{tsa\-ta}} \emph{exists but is rarely in use})

% fwa
\hi
\HT{fwa}{\B{fwa}}, \B{fì\-\cp{'u} a}, \emph{or} \B{a fì\-\cp{'u}} \I{[fwa]} \emph{conj.} that (\emph{to nominalize a subordinate clause}). (\emph{a. with adj.}) \B{Tsun oe nga\-hu nì\-Na'\-vi pi\-väng\-kxo a fì\-'u oe\-ru prr\-te' lu.} It's a pleasure to be able to chat with you in Na'vi.~\LINK{http://wiki.learnnavi.org/index.php?title=Corpus\#Times_Online}{wiki} 
\B{Swey lu fwa nga fì\-kem ke si\-vi nì\-sngum.} It's best that you not do this fretfully.~\LINK{http://naviteri.org/2011/01/new-year-new-vocabulary/}{b} 
(\emph{b. with vin.}) \B{Su\-nu oe\-ru fwa fì\-tseng\-it te\-rok oel nì\-mun.} I like that I'm here again.~\LINK{http://naviteri.org/2010/08/a-na\%E2\%80\%99vi-alphabet/}{b} 
\B{Lam oer fwa po lu ka\-nu nì\-txan.} It seems to me that he is very clever.~\LINK{http://forum.learnnavi.org/language-updates/txelanit-hivawl/}{ln} 
\B{Toruk Makto po\-lä\-hem a fì\-'u ro\-lo\-'a nì\-txan Oma\-ti\-ka\-ya\-ru.} The arrival of \emph{Toruk Makto} made a great impression on the Omatikaya.~\LINK{http://naviteri.org/2012/03/spring-vocabulary-part-1/}{b}
(\emph{c. to introduce verb phrases}) \B{Oe\-yä txin\-tìn lu fwa stä'\-nì fay\-o\-ang\-it.} My central societal occupation is to catch fish.~\LINK{http://naviteri.org/2011/05/weather-part-2-and-a-bit-more-2/}{b}
(\emph{d. with} \B{to} \emph{to compare verb phrases}) \B{Ftu\-e lu fwa ta\-ron ngong\-a i\-o\-ang\-it to fwa ta\-ron pum\-it a lu wa\-lak sì win.} It's easier to hunt lethargic animals than to hunt perky, speedy ones.~\LINK{http://naviteri.org/2011/10/more-vocabulary-a-bit-of-grammar/}{b} 
(\emph{e. with adp.s to introduce a verb phrase}) \B{Sre fwa sìn tskxe\-pay tì\-ran, zene fko flì\-nutx\-it sti\-ve\-ftxaw.} It's necessary to check the thickness of the ice before walking on it.~\LINK{http://naviteri.org/2011/11/\%E2\%80\%99a\%E2\%80\%99awa-ayli\%E2\%80\%99u-amip-ni\%E2\%80\%99aw-\%E2\%80\%94-a-few-new-words-only/}{b} 
(\emph{f. with} \B{sweylu} \emph{to mean}) `should have' \B{Swey\-lu fwa nga ko\-lä.} You should have gone.~\LINK{http://naviteri.org/2011/04/\%E2\%80\%99a\%E2\%80\%99awa-li\%E2\%80\%99fyavi-amip\%E2\%80\%94a-few-new-expressions/}{b}
(\ding{43}~\HL{fI'u}{\B{fì\-'u}}, \emph{nominalizer of} \ding{43}~\HL{tsa'u}{\B{tsa\-'u}}, \emph{i.e.} \HL{tsawa}{\B{tsawa}} \emph{exists but is rarely in use})

% fwal
\hi 
\HT{fwal}{\B{fwal}} \I{[fwal]} \emph{vtr.} wipe. 
\B{Yo\-ti fwal ru\-txe. Mei slo\-lu maw tom\-pa.} Please wipe the table. It's gotten wet after the rain.~\LINK{http://naviteri.org/2020/07/vospxivol-lefpom-happy-august/}{b}

% fwampop
\hi
\HT{fwampop}{\B{\cp{fwam}pop}} \I{["fwam.pop\textcorner]} \emph{n., fauna} tapirus. \B{Oel tso\-le\-'a fwam\-pop\-it a\-win fra\-to.} I saw the fastest Tapirus.~\LINK{http://forum.learnnavi.org/language-updates/confirmations-\%28comparisons-gerund\%29/}{ln}
\B{Fwam\-pop\-ä tem\-rey tsa\-tseng\-mì ngä\-zìk lu nì\-wotx.} It's very hard for a tapirus to survive there.~\LINK{http://naviteri.org/2011/05/some-miscellaneous-vocabulary/}{b}
(\emph{a. idiom}) \B{(na) fwam\-pop fkip fì\-wopx}, ``Fish out of water.'' [lit.: (like a) tapirus in the clouds].~\LINK{http://naviteri.org/2021/08/ulte-ayyoratu-leiu-and-the-winners-are-2/}{b}

% fwang
\hi
\HT{fwang}{\B{fwang}} \I{[fwaN]} \emph{adj.} (\emph{of taste and smell}) savory, umami, rich.

% fwäkì
\hi
\HT{fwAkI}{\B{\cp{fwä}kì}} \I{["fw\ae.kI]} \emph{n., fauna} mantis.
\B{Sro\-ler fwä\-kì ftu\-opa tskxe.} A mantis appeared from behind a rock.~\LINK{http://naviteri.org/2019/06/50a-liu-amip-40-new-words/}{b} 
(\emph{a. proverb}) \B{Fwä\-kì ke fwe\-fwi.} It's against s.o.'s\slash s.t.'s nature [lit.: a mantis doesn't whistle].~\LINK{http://naviteri.org/2011/09/miscellaneous-vocabulary/comment-page-1/\#comment-1143}{b} \\
\textsc{Derivations:} \ding{43} \on{1}. \HL{fwAkIwll}{\B{fwäkìwll}}, mantic orchid.

% fwäkìwll
\hi
\HT{fwAkIwll}{\B{\cp{fwä}kìwll}} \I{["fw\ae.kI.w\s{l}]} \emph{n.}~\ding{96} mantis orchid, \emph{pandorchidus insectoralis}. 
\B{Ra\-lul ngo\-lop 'um\-tsat fa ay\-syu\-lang fwä\-kì\-wll\-ä.} Ralu made medicine from flowers of the Mantis orchid.~\LINK{http://naviteri.org/2014/11/twenty-before-the-holidays/}{b} 
(\ding{43}~\HL{fwAkI}{\B{fwä\-kì}}, \HL{'ewll}{\B{'e\-wll}})

% fwefwi
\hi
\HT{fwefwi}{\B{\cp{fwe}fwi}} \I{["fwE.fwi]} \emph{vin.} whistle.   
\B{Nga fwe\-fwi nìl\-ke\-ftang a fì\-'u lu sä\-srätx a\-txan oe\-ru.} It's a major annoyance to me that you whistle continuously.~\LINK{http://naviteri.org/2011/09/miscellaneous-vocabulary/}{b} \\
\textsc{Derivations:} \ding{43} \on{1}.~\HL{nIfwefwi}{\B{nìfwefwi}}, by whistling.

% fwew
\hi
\HT{fwew}{\B{fwew}} \I{[fwEw]} \emph{vtr.} look for, seek. \B{Ngal fwe\-rew a tu\-te ke fkey\-tok.} The person you're looking for doesn't exist.~\LINK{http://naviteri.org/2011/04/yafkeykiri-plltxe-frapo-everyone-talks-about-the-weather/}{b} 
\B{Tsa\-'u\-pxa\-ret fwo\-lew oel slä ke tsä\-ngun ri\-vun.} I looked for that message but I couldn't find it.~\LINK{http://naviteri.org/2017/06/ayliu-a-ta-eywaeveng-kifkey-uniltirantokxa-words-from-disneys-pandora-the-world-of-avatar/}{b}
(\ding{214}~\HL{run}{\B{run}})

% fwel
\hi
\HT{fwel}{\B{fwel}} \I{[fwEl]} \emph{adj.} broken. 

% fwem
\hi
\HT{fwem}{\B{fwem}} \I{[fwEm]} \emph{adj.} (\emph{of a point}) dull, blunt. 
\B{Ko\-ak\-tan\-ä ay\-sre' lä\-ngu fwem.} An old man's teeth are dull.~\LINK{http://naviteri.org/2012/03/spring-vocabulary-part-1/}{b} 
\B{Tsa\-tstal a\-fwem lu litx nì\-txan.} That blunt knife is very sharp.~\LINK{http://naviteri.org/2012/03/spring-vocabulary-part-1/}{b} 
(\ding{214}~\HL{pxi}{\B{pxi}}, \ding{43}~\HL{tete}{\B{te\-te}})

% fwep
\hi 
\HT{fwep}{\B{fwep}} \I{[fwEp\textcorner]} \emph{n.} dust (\emph{on a surface}).
\B{Slär\-ìl nga' fwep\-it a\-txan.} The cave is very dusty.~\LINK{http://naviteri.org/2014/07/tsawlultxamaw-a-ayliu-post-meetup-words/}{b}
\B{Txe\-wìl 'o\-la\-ku fwep\-it ftum\-fa kel\-ku sne\-yä.} Txewì removed the dust from inside his house.~\LINK{http://naviteri.org/2014/07/tsawlultxamaw-a-ayliu-post-meetup-words/}{b} \\
\textsc{Derivations:} \ding{43} \on{1}.~\HL{fwopx}{\B{fwopx}}, dust cloud, dust (in the air).

% fwi
\hi 
\HT{fwi}{\B{fwi}} \I{[fwi]} \emph{vin.} slip, slide.
\B{Na\-ri si! Kll\-te\-sìn lu pay a\-txan. Fwi rä\-'ä!} Be careful! There’s a lot of water on the ground. Don't slip!~\LINK{http://naviteri.org/2013/04/wheres-the-bathroom-and-other-useful-things/}{b}
\B{Txur\-tel\-mì fo fta so\-li fte\-ke ka tsyokx fwi\-vi.} They tied a knot in the rope so it wouldn't slip through their hands.~\LINK{http://naviteri.org/2013/04/wheres-the-bathroom-and-other-useful-things/}{b}
\B{Po\-ti fwey\-ko\-li ay\-oel ìlä rong.} We let him slide through the tunnel.~\LINK{http://naviteri.org/2013/04/wheres-the-bathroom-and-other-useful-things/}{b}

% fwìng / si
\hi 
\HT{fwIng}{\B{fwìng}} \I{[fwIN]} \textsc{i}. \emph{n.} humiliation, embarrassment, loss of face. 
\B{Ra\-lu\-ri fwa tì\-fme\-tok\-it ke em\-zo\-la\-'u lä\-ngu fwìng a\-txan.} Ra\-lu's not passing the test was a great humiliation (to him).~\LINK{http://naviteri.org/2018/04/100a-liu-amip-64-new-words-part-2/}{b} 
\B{Peyä hem\-ìl a\-ke\-muia\-nga' za\-mo\-lu\-nge fwìng\-it ay\-oer.} His dishonorable behavior (lit.: actions) humiliated us (lit.: brought humiliation to us).~\LINK{http://naviteri.org/2020/05/aawa-liu-amip-si-vurway-alor-a-few-new-words-and-a-beautiful-poem/}{b} \\
\textsc{ii}. \B{fwìng si} \I{[fwIN si]} \emph{vin.} humiliate. \\
\textsc{Derivations:} \ding{43} \on{1}.~\HL{txanfwIngtu}{\B{txanfwìngtu}}, bastard, loser (\emph{insult}).

% fwopx
\hi 
\HT{fwopx}{\B{fwopx}} \I{[fwop']} \emph{n.} dust (\emph{in the air}).
\B{Tsaw\-ke slo\-lu vawm ta\-lun fwopx.} The sun became dark due to dust in the air.~\LINK{http://naviteri.org/2014/07/tsawlultxamaw-a-ayliu-post-meetup-words/}{b}
(\ding{43}~\HL{fwep}{\B{fwep}}, \HL{pIwopx}{\B{pì\-wopx}}) \\
\textsc{Derivations:} \ding{43} \on{1}.~\HL{eanfwopx}{\B{ean\-fwopx}}, mist bloom.

% fwum
\hi 
\HT{fwum}{\B{fwum}} \I{[fwum]} \emph{or} \I{[fwUm]} \emph{vin.} float (\emph{on the surface of a liquid}). 
\B{Me\-rìk a\-lor pay\-sìn fwar\-mum.} Two beautiful leaves were floating on the water.~\LINK{http://naviteri.org/2021/12/zolau-niprrte-ma-3746-welcome-2022/}{b}
\B{Ke omum teyng\-ta fì\-u\-ran a\-ku'\-up fu fwa\-yum fu wa\-yapx.} I don't know (or: It's not known) whether this heavy boat will float or sink.~\LINK{http://naviteri.org/2021/12/zolau-niprrte-ma-3746-welcome-2022/}{b}

% fya'o
\hi
\HT{fya'o}{\B{\cp{fya}'o}} \I{["fja.Po]} \emph{n.} path, way; method, manner. \B{Fol lì'\-fya\-ti aw\-nge\-yä sar a fya\-'o lu ko\-sman.} It's wonderful the way they use our language.~\LINK{http://naviteri.org/2011/09/miscellaneous-vocabulary/}{b} \B{pa\-ye\-'un swey\-a fya\-'ot} I will decide the best way.~\LINK{http://forum.learnnavi.org/news-announcements/a-response-from-paul-frommer!/}{ln} \\  
\textsc{Derivations}: \ding{43} \on{1}.~\HL{fyape}{\B{fyape}\slash \B{pefya}}, how?   
\on{2}.~\HL{nIfya'o}{\B{nìfya'o}}, in a manner (that). 
\on{3}.~\HL{eyawrfya}{\B{eyawrfya}}, righ way, correct path. 
\on{4}.~\HL{fyalI'u}{\B{fya\-lì\-'u}}, adverb. 
\on{5}.~\HL{yo'kofya}{\B{yo'\-ko\-fya}}, cycle.

% fyalì'u
\hi 
\HT{fyalI'u}{\B{\cp{fya}lì'u}} \I{["fja.lI.Pu]} \emph{n.} adverb. 
(\ding{43}~\HL{fya'o}{\B{fya\-'o}}, \HL{lI'u}{\B{lì\-'u}})

% fyan
\hi
\HT{fyan}{\B{fyan}} \I{[fjan]} \emph{n.} \emph{constructed device for keeping s.t. off the ground and clean}. \\
\textsc{Derivations:} \ding{43} \on{1}.~\HL{fyanyo}{\B{fyanyo}}, table. 
\on{2}.~\HL{fyanyI}{\B{fyanyì}}, shelf. 
\on{3}.~\HL{kurfyan}{\B{kurfyan}}, hamper, suspended rack.

% fyanyì
\hi
\HT{fyanyI}{\B{\cp{fyan}yì}} \I{["fjan.jI]} \emph{n.} shelf.
(\ding{43}~\HL{fyan}{\B{fyan}}, \HL{yI}{\B{yì}})

% fyanyo
\hi
\HT{fyanyo}{\B{\cp{fyan}yo}} \I{["fjan.jo]} \emph{n.} table, elevated utilitarian surface; \emph{coll.}: \ding{43}~\HL{yo}{\B{yo}}.
\B{Oe\-yä tstal\-ìl tok pe\-se\-nget? -- Lu yo\-sìn.} Where's my knife? -- It's on the table.~\LINK{http://naviteri.org/2015/08/aylifyavi-lereyfya-2-cultural-terms-2-and-more/}{b}
(\ding{43}~\HL{fyan}{\B{fyan}}, \HL{yo}{\B{yo}})

% fyape
\hi
\HT{fyape}{\B{\cp{fya}pe}} \emph{or} \B{pe\cp{fya}} \I{["fja.pE]} \emph{or} \I{[pE."fja]} \emph{inter.} how?   
\B{Fya\-pe fko syaw ngar?} What's your name?~\LINK{http://forum.learnnavi.org/language-updates/the-reflexive-\%28from-a-frommer-email!!!!\%29/}{ln} 
\B{'Ìn nga fya\-pe nì\-fkrr?} What's been keeping you busy lately?~\LINK{http://naviteri.org/2010/09/getting-to-know-you-part-3/}{b} 
\B{Ya\-fkeyk za\-'u fya\-pe?} How's the weather?~\LINK{http://naviteri.org/2011/04/yafkeykiri-plltxe-frapo-everyone-talks-about-the-weather/}{b}
\B{Mak\-to fya\-pe?} ``How are things? How's everything going? How are you doing?'
\B{A: Mak\-to fya\-pe?---B: Zong. Nga\-ri tut?---A: Nìk\-sran. Oe\-ru lu fpom, slä oey 'i\-tan lu spxin.} A: How's everything?---B: Good. You?---A: So\hyp{}so. I'm fine, but my son is sick.~\LINK{http://naviteri.org/2020/05/tipusawm-tiuseyng-si-okvur-a-eltur-titxen-si-asking-answering-and-an-interesting-story/}{b} 

% fyawìntxu
\hi
\HT{fyawIntxu}{\B{fyawìn\cp{txu}}} \I{[fja.w$\cdot$In."t'$\cdot$u]} \emph{vin., inf.}\on{23} guide. \B{oe new nì\-txan ay\-nga\-ru fya\-wi\-vìn\-txu}, I want to guide you.~\LINK{http://forum.learnnavi.org/news-announcements/a-response-from-paul-frommer!/}{ln} 
(\ding{43}~\HL{fya'o}{\B{fya\-'o}}, \HL{wIntxu}{\B{wìn\-txu}}) 
\on{1}.~\HL{tIfyawIntxu}{\B{tìfyawìntxu}}, guidance. 
\on{2}.~\HL{fyawIntxuyu}{\B{fyawìntxuyu}}, guide.

% fyawìntxuyu
\hi 
\HT{fyawIntxuyu}{\B{fyawìn\cp{txu}yu}} \I{[fja.wIn."t'u.ju]} \emph{n.} guide (\emph{of a person}). 
(\ding{43}~\HL{fyawIntxu}{\B{fya\-wìn\-txu}}, \HL{-yu}{\B{-yu}})

% fyel
\hi
\HT{fyel}{\B{fyel}} \I{[fjEl]} \emph{vtr.} seal, cement, make impervious (\B{wä}, against) (\emph{making something tight and secure, so that it is unlikely (and not intended) to be broken}). 
\B{Txo fkol ke fyi\-vel u\-ran\-it pay\-wä, ze\-ne fko sli\-ve\-le.} If one does not seal a boat against water, one must swim.~\LINK{http://naviteri.org/2012/06/spring-vocabulary-part-3/}{b}

% fyeng
\hi 
\HT{fyeng}{\B{fyeng}} \I{[fjEN]} \emph{adj.} steep. 
\B{Na\-ri si! Fay\-ram\-tsyìp lu fyeng.} Be careful! These hills are steep.~\LINK{http://naviteri.org/2019/06/50a-liu-amip-40-new-words/}{b}

% fyengkekin
\hi 
\HT{fyengkekin}{\B{\cp{fyeng}kekin}} \I{["fjEN.kE.kin]} \emph{conv.} `no need to use honorific language' (\emph{from} \B{lì'fyeng kelkin}). 
(\ding{43}~\HL{lI'fyeng}{\B{lì'\-fyeng}}, \HL{kelkin}{\B{kel\-kin}}) 

% fyep
\hi
\HT{fyep}{\B{fyep}} \I{[fjEp\textcorner]} \emph{vtr.} (\emph{a.}) hold in the hand, grasp, grip. \B{Ngä\-zìk lu fwa var tsko\-ti fyi\-vep teng\-krr ut\-ral\-it tsye\-rìl.} It's hard to keep holding a bow while climbing a tree.~\LINK{http://naviteri.org/2012/03/spring-vocabulary-part-1/}{b} 
(\emph{b. general holding}) \B{Oel tstal\-it fyo\-lep fa ay\-sre'.} I held the knife in my teeth.~\LINK{http://naviteri.org/2012/03/spring-vocabulary-part-1/}{b} \\
\textsc{Derivations:} \ding{43} \on{1}.~\HL{sAfyep}{\B{säfyep}}, handle, grip.

% fyeyn
\hi
\HT{fyeyn}{\B{fyeyn}} \I{[fjEjn]} \emph{adj.} ripe, mature, adult. \\  
\on{1}. \HL{fyeyntu}{\B{fyeyntu}}, adult person.

% fyeyntu
\hi
\HT{fyeyntu}{\B{\cp{fyeyn}tu}} \I{[fjEjn.tu]} \emph{n.} adult person. 
\B{Le\-hrr\-ap lu fwa evi\-tsyìp sle\-le mì hil\-van lu\-ke fwa fyeyn\-tu te\-rìng nari.} It's dangerous for tiny ones to swim in the river without an adult watching.~\LINK{http://naviteri.org/2011/02/new-vocabulary-ii\%E2\%80\%94part-1/}{b} 
\B{Nì\-ran lu Lo\-ak mi 'e\-veng slä tsun ti\-va\-ron nì\-teng\-fya na fyeyn\-tu.} Loak is still really just a boy but he can hunt the same as an adult.~\LINK{http://naviteri.org/2012/10/mipa-vospxi-mipa-ayliu-new-words-for-the-new-month/}{b} 
(\ding{43}~\HL{fyeyn}{\B{fyeyn}}) \\
\textsc{Derivations:} \ding{43} \on{1}.~\HL{nIfyeyntu}{\B{nìfyeyntu}}, maturely.

% fyin
\hi
\HT{fyin}{\B{fyin}} \I{[fjin]} \emph{adj.} simple. 
\B{Ta\-fral kì\-syar a ay\-sä\-num\-vi sì ay\-lì'\-fya\-vi lu fyin.} That's why I'll teach lessons and expressions that are simple.~\LINK{http://naviteri.org/2012/07/tskxekengtsyip-a-mikyunfpi-a-little-listening-exercise/}{b} 
(\ding{214}~\HL{ep'ang}{\B{ep\-'ang}}) \\
\textsc{Derivations:} \ding{43} \on{1}.~\HL{fyinep'ang}{\B{fyinep'ang}}, complexity.

% fyinep'ang / fyinep'angpe
\hi
\HT{fyinep'ang}{\B{fyinep\cp{'ang}}} \I{[fjin.Ep\textcorner."PaN]} \textsc{i}. \emph{n.} degree of complexity. \\  
\textsc{ii}. \B{fyinep\cp{'ang}pe} \I{[fjin.Ep\textcorner."PaN]} \emph{inter.} how complex? 
\B{Ngal ke tslä\-ngam teyng\-ta fì\-tì\-ngä\-zìk\-ì\-ri pe\-fyin\-ep\-'ang.} Unfortunately you don't understand how complex this problem is.~\LINK{http://naviteri.org/2012/07/fivospxiya-aylifyavi-amip-this-months-new-expressions-pt-2/}{b}
(\ding{43}~\HL{fyin}{\B{fyin}}, \HL{ep'ang}{\B{ep\-'ang}})

% fyìp
\hi 
\HT{fyIp}{\B{fyìp}} \I{[fjIp\textcorner]} \emph{n.} tendril, tentacle. \\
\textsc{Derivations:} \ding{43} \on{1}.~\HL{fyIpmaut}{\B{fyìpmaut}}, squid fruit tree.

% fyìpmaut
\hi 
\HT{fyIpmaut}{\B{fyìpmaut}} \I{[fjIp\textcorner.ma.ut\textcorner]} \emph{n.} \ding{96} squid fruit tree, \emph{octocrus folliculus}.~\LINK{https://forum.learnnavi.org/language-updates/expansion-to-flora-fauna-and-places-from-the-valley-of-mo'ara/}{ln} 
(\ding{43}~\HL{fyIp}{\B{fyìp}}, \HL{mauti}{\B{mau\-ti}}, \HL{utral}{\B{ut\-ral}})

% fyole
\hi
\HT{fyole}{\B{\cp{fyo}le}} \I{["fjo.lE]} \emph{adj.} sublime, beyond perfection. 
\B{Ran tì\-ru\-sol\-ä pe\-yä lu fyo\-le.} The \emph{ran} of her singing is sublime.~\LINK{http://naviteri.org/2012/10/mipa-vospxi-mipa-ayliu-new-words-for-the-new-month/}{b} 
\B{Nge\-yä tì\-kang\-kem a\-fyo\-le me\-u\-i\-a le\-i\-u lì'\-fya\-o\-lo'\-ru, ma tsmu\-kan.} Your sublime work honors the language clan, brother.~\LINK{http://naviteri.org/2012/10/audio-and-video-learning-materials-for-navi-101/}{b} \\
\textsc{Derivations:} \ding{43} \on{1}.~\HL{fyolup}{\B{fyolup}}, exquisite, sublime in style.

% fyolup
\hi
\HT{fyolup}{\B{\cp{fyo}lup}} \I{["fjo.lup\textcorner]} \emph{or} \I{[.lUp\textcorner]} (\textsc{rn}: \B{\cp{fyo}lùp} \I{["fjo.lUp\textcorner]}) \emph{adj.} exquisite, sublime in style.
\B{Fì\-tsko\-ti a\-fyo\-lup 'o\-ngo\-lop oe\-yä sem\-pul\-ìl.} This exquisite bow was designed by my father.~\LINK{http://naviteri.org/2014/07/tsawlultxamaw-a-ayliu-post-meetup-words/}{b}
\B{Fyo\-lup\-a ay\-sä\-ngop ay\-nge\-yä oe\-ru te\-ya si nì\-ngay.} Your exquisite creations fill me greatly with joy.~\LINK{http://naviteri.org/2014/12/ayngaru-tsinga-yoratu-may-i-present-the-four-winners/}{b}
(\ding{43}~\HL{lupra}{\B{lup\-ra}}, \HL{fyole}{\B{fyo\-le}}) 

% fyuatx
\hi 
\HT{fyuatx}{\B{fyu\cp{atx}}} \I{[fju."{}at']} \emph{n., fauna} anemonoid.~\LINK{https://forum.learnnavi.org/language-updates/expansion-to-flora-fauna-and-places-from-the-valley-of-mo'ara/}{ln} 

% fyufye
\hi 
\HT{fyufye}{\B{\cp{fyu}fye}} \I{["fju.fjE]} (\textsc{rn}: \B{\cp{fyù}fye} \I{["fjU.fjE]}) \emph{vin.} splash. 
\B{Ra\-nu kil\-van\-mì fyar\-mu\-fye na tsawl\-a pa\-yo\-ang a\-pì\-saw.} Ra\-nu was splashing in the river like a big clumsy fish.~\LINK{http://naviteri.org/2023/09/aawa-ayliu-si-aylifyavi-amip-a-few-new-words-and-expressions/}{b}

\end{multicols}

% ===== section H ===== 

 \pdfbookmark[0]{H}{H}
\begin{center}
\section*{H \I{[h]} \emph{Hä}}
\end{center}

\begin{multicols}{2}

% ha'
\hi
\HT{ha'}{\B{ha'}} \I{[haP]} \emph{vin.} fit, suit, complement, inherently enhance (two entities (things, people, situations, \dots) fit or suit each other; they ``go together'' well).  
\B{Tse\-nu sì Lo\-ak fì\-tsap ke ha' kaw\-'it.} Tse\-nu and Lo\-ak are a terrible match for each other.~\LINK{http://naviteri.org/2012/10/mipa-vospxi-mipa-ayliu-new-words-for-the-new-month/}{b} 
\B{Tse\-nu Lo\-ak\-ur ke ha'.} Tse\-nu is a bad match for Lo\-ak.~\LINK{http://naviteri.org/2012/10/mipa-vospxi-mipa-ayliu-new-words-for-the-new-month/}{b} 
\B{Tse\-nur Lo\-ak ke hä\-nga'.} Lo\-ak isn't good for Tse\-nu.~\LINK{http://naviteri.org/2012/10/mipa-vospxi-mipa-ayliu-new-words-for-the-new-month/}{b} 
\B{Hu\-fwa nge\-yä tì\-hawl\-ìri ke lu kea kxey\-ey, tsal\-su\-ngay oe\-ru ke ha' nì\-tam.} Although there's nothing wrong with your plan, it just doesn't suit me.~\LINK{http://naviteri.org/2012/10/mipa-vospxi-mipa-ayliu-new-words-for-the-new-month/}{b} 
\B{Ngay. Tsa\-'o\-pin hek nì\-'it, slä su\-nu oer, ha ha'.} True. That color is bit odd, but I like it, so it's a good fit for me.~\LINK{http://naviteri.org/2012/10/mipa-vospxi-mipa-ayliu-new-words-for-the-new-month/}{b} \\
\textsc{Derivation}: \on{1}.~\HL{leha'}{\B{le\-ha'}}, suitable, appropriate, fitting.

% ha'ngir
\hi
\HT{ha'ngir}{\B{ha'\cp{ngir}}} \I{[haP."NiR]} \emph{n., adv.} afternoon. \B{Po\-el sto\-la\-tso fu\-ta me\-fo ti\-va\-ron tsa\-ha'\-ngir.} She must have refused their request to hunt that afternoon.~\LINK{http://forum.learnnavi.org/language-updates/sto-has-the-same-syntax-as-new/msg567034/\#msg567034}{ln} \\
(I.) \B{ha'ngi\cp{ram}} \I{[haP.Ni."Ram]} \emph{n., adv.} yesterday afternoon. \\ 
(II.) \B{ha'ngi\cp{ray}} \I{[haP.Ni."Raj]} \emph{n., adv.} tomorrow afternoon. 

% ha
\hi
\HT{ha}{\B{ha}} \I{[ha]} \emph{conj.} so, in that case. 
\B{Tom\-pa 'e\-ko nì\-hawng, ha ze\-ne aw\-nga kll\-wi\-vo.} The rain is too strong, so we must land.~\LINK{http://naviteri.org/2012/01/mipa-zisit-ayliu-amip-new-words-for-the-new-year/}{b}

% hafyon
\hi 
\HT{hafyon}{\B{ha\cp{fyon}}} \I{[ha."fjon]} \emph{n.} wisdom.
\B{Nge\-nge\-yä ha\-fyon\-ìri ira\-yo, ma le\-i\-o\-ko\-ak\-te.} Thank you for your wisdom, respected female elder.~\LINK{https://naviteri.org/2025/03/conlang-adventure-and-more/}{b} 
(\ding{43}~\HL{tItxantslusam}{\B{tìtxantslusam}}) \\
\textsc{Derivatives}: \ding{43} \on{1}.~\HL{lafyon}{\B{la\-fyon}}, wise. 
\on{2}.~\HL{hafyonga'}{\B{ha\-fyo\-nga'}}, wise.

% hafyonga'
\hi 
\HT{hafyonga'}{\B{ha\cp{fyo}nga'}} \I{[ha."fjo.NaP]} \emph{adj., nfp.} wise. 
\B{ay\-lì\-'u a\-ha\-fyo\-nga'}, wise words.~\LINK{http://naviteri.org/2018/03/100a-liu-amip-64-new-words-part-1/}{b}
(\ding{43}~\HL{hafyon}{\B{hafyon}}, \HL{lafyon}{\B{lafyon}})

% hahaw
\hi
\HT{hahaw}{\B{\cp{ha}haw}} \I{["ha.haw]} \emph{vin.} sleep. 
(\emph{a.}) \B{Ney\-ti\-ri he\-ra\-haw.} Ney\-ti\-ri is sleeping.~\LINK{http://wiki.learnnavi.org/index.php?title=Corpus\#MSNBC}{wiki}  
\B{Sem\-pul lì'\-fya\-yä har\-ma\-haw.} The father of the language was sleeping.~\LINK{http://forum.learnnavi.org/language-updates/a-long-answer-to-a-quick-question/msg155104/\#msg155104}{ln} 
\B{Lu na'\-rìng tsuk\-ha\-haw.} It's possible to sleep in the forest.~\LINK{http://naviteri.org/2011/03/\%E2\%80\%9Creceptive-ability\%E2\%80\%9D-and-hesitation/}{b} 
(\emph{b. idiom}) \B{Hi\-va\-haw nìm\-wey.} Goodnight (\emph{when it's bedtime}). 
(\emph{c. proverb}) \B{Ha\-haw nì\-'aw txo pa\-lu\-kan smi\-von ngar.} ``Only sleep if you are familiar with the Thanator,'' i.e. \emph{don't think you're safe unless you're aware of the danger}.~\LINK{http://naviteri.org/2021/08/ulte-ayyoratu-leiu-and-the-winners-are-2/}{b} \\
\textsc{Derivatives}: \ding{43} 
\on{1}.~\HL{hawtsyIp}{\B{haw\-tsyìp}}, nap. 
\on{2}.~\HL{tsukhahaw}{\B{tsuk\-ha\-haw}}, `sleepable', able to sleep in. 
\on{3}.~\HL{snomo}{\B{sno\-mo a ha\-haw}}, bedroom, area to sleep.
\on{4}.~\HL{haway}{\B{ha\-way}}, lullaby. 
\on{5}.~\HL{yehaw}{\B{ye\-haw}}, well\hyp{}rested.

% ham
\hi
\HT{ham}{\B{ham}} \I{[ham]} \emph{adj.} last, previous. 
(\ding{214}~\HL{hay}{\B{hay}}) \\
\textsc{Derivatives}: \ding{43} \on{1}.~\HL{hamalo}{\B{ha\-ma\-lo}}, last time.

% hamalo
\hi
\HT{hamalo}{\B{ha\cp{ma}lo}} \I{[ha."ma.lo]} \emph{adv.} last time. 
(\ding{214}~\HL{hayalo}{\B{hayalo}}, \ding{43}~\HL{}{\B{ham}}, \HL{alo}{\B{alo}}) 

% hametsì
% \hi
% \B{ha\cp{me}tsì} \I{[ha."mE.\t{ts}I]} \emph{n.}, \textsc{lw} Chametz (\emph{leavened bread}). \B{aw\-ngal yom ha\-me\-tsìt, yom ma\-tsat, ke tsran\-ten}, it doesn't matter whether we eat Chametz or Matzah.~\LINK{http://whyisthisnight.com/addl-more.html\#na\%27vi}{this night}

% han
\hi 
\HT{han}{\B{han}} \I{[han]} \emph{vtr.} lose (\emph{not having s.t. you once had}). \\
\textsc{Derivatives}: \ding{43} \on{1}.~\HL{tIhah}{\B{tì\-han}}, loss.

% hangham
\hi
\HT{hangham}{\B{\cp{hang}ham}} \I{["haN.ham]} \emph{vin.} laugh. 
\B{Oe ho\-lang\-ham nì\-teng.} I laughed too.~\LINK{http://masempul.org/2010/07/\%E2\%80\%99ite-polaheiem/}{sempul} 
\B{Po\-leng Ney\-ti\-ril hang\-vur\-it a fra\-pot hey\-kang\-ham.} Ney\-ti\-ri told a joke that made everyone laugh.~\LINK{http://naviteri.org/2012/01/mipa-zisit-ayliu-amip-new-words-for-the-new-year/}{b}\\
\textsc{Derivatives}: \ding{43} 
\on{1}.~\HL{hangvur}{\B{hang\-vur}}, joke, funny story.
\on{2}.~\HL{tIhangham}{\B{tì\-hang\-ham}}, laughter.
\on{3}.~\HL{lehangham}{\B{le\-hang\-ham}}, laughable, ridiculous, absurd.

% hangvur
\hi
\HT{hangvur}{\B{\cp{hang}vur}} \I{["haN.vuR]} \emph{or} \I{[.vUR]} \emph{n.} joke, funny story. 
\B{Po\-leng Ney\-ti\-ril hang\-vur\-it a fra\-pot hey\-kang\-ham.} Neytiri told a joke that made everyone laugh.~\LINK{http://naviteri.org/2012/01/mipa-zisit-ayliu-amip-new-words-for-the-new-year/}{b}
(\ding{43}~\HL{hangham}{\B{hang\-ham}}, \HL{vur}{\B{vur}})

% hapxì
\hi
\HT{hapxI}{\B{ha\cp{pxì}}} \I{[ha."p'I]} \emph{n.} part. 
\B{Nga\-ri hu Ey\-wa sa\-lew ti\-re\-a, tokx 'ì'awn slu Na'\-vi\-yä ha\-pxì.} Your spirit goes with Eywa, your body remains to become part of the People.~\LINK{http://wiki.learnnavi.org/index.php?title=Corpus\#Behind_the_Scenes}{a\slash wiki} 
\B{Pì\-lok\-ä fì\-ha\-pxì\-yä tì\-kan lu law.} The aim of this part of the blog is clear.~\LINK{http://naviteri.org/2010/06/first-post/}{b} \\
\textsc{Derivatives}: \ding{43} \on{1}.~\HL{hapxItu}{\B{ha\-pxì\-tu}}, member.   
\on{2}.~\HL{kApxI}{\B{kä\-pxì}}, rear (part\slash section).   
\on{3}.~\HL{zapxI}{\B{za\-pxì}}, front (part\slash section).   
\on{4}.~\HL{hImpxI}{\B{hìm\-pxì}}, minority. 
\on{5}.~\HL{txampxI}{\B{txam\-pxì}}, majority.   
\on{6}.~\HL{vospxI}{\B{vo\-spxì}}, month. 
\on{7}.~\HL{pxI}{\B{\mbox{-pxì}}}, \emph{to form fractions}. 
\on{8}.~\HL{pxImun'i}{\B{pxì\-mun\-'i}}, devide.
\on{9}.~\HL{swapxI}{\B{swa\-pxì}}, family member.
\on{10}.~\HL{napxI}{\B{na\-pxì}}, partially, in part.

% hapxìtu
\hi
\HT{hapxItu}{\B{ha\cp{pxì}tu}} \I{[ha."p'I.tu]} \emph{n.} member. 
\B{sìl\-pey oe tsnì ay\-ha\-pxì\-tu lì'\-fya\-o\-lo'\-ä aw\-nge\-yä ke stì\-ye\-vi.} I hope that the members of the language clan won't be angry.~\LINK{http://forum.learnnavi.org/language-updates/txelanit-hivawl/}{ln} 
\B{Me\-ha\-pxì\-tul lì'\-fya\-o\-lo'\-ä pxe\-stxe\-lit a\-lor to\-lìng aw\-nga\-ru.} Two members of the language clan have given us three beautiful gifts.~\LINK{http://naviteri.org/2010/12/beautiful-christmas-carols-sung-in-navi/}{b} 
\B{Ul\-te pi\-veng ru\-txe ay\-ha\-pxì\-tur a\-la\-he san se\-i\-yi i\-ra\-yo sìk nì\-teng.} And please tell the other members `thank you', too. 
\B{Va'\-rul zä\-nget ik\-ran\-it sn\-eyä pxel ha\-pxì\-tu so\-a\-i\-ä.} Va'\-ru treats his ikran like a member of the family.~\LINK{http://naviteri.org/2012/10/mipa-vospxi-mipa-ayliu-new-words-for-the-new-month/}{b} 
(\ding{43}~\HL{hapxI}{\B{hapxì}})

% hasa'
\hi 
\HT{hasa'}{\B{ha\cp{sa'}}} \I{[ha."{}saP]} \emph{n.} orphan. 
\B{Krr\-ka tsam slä\-ngu pxay\-a 'e\-veng ha\-sa'.} During war, sad to say, many children become orphans.~\LINK{https://naviteri.org/2024/09/pxevola-liu-amip-two-dozen-new-words/}{b} 

% hasey / si
\hi
\HT{hasey}{\B{ha\cp{sey}}} \I{[ha."{}sEj]} \textsc{i}.~\emph{adj.} done, finished. \B{Hu Saw\-tu\-te a tsam lu ha\-sey.} The war with the Skypeople is finished.~\LINK{http://naviteri.org/2010/07/ting-mikyun-fte-tslivam-listening-comprehension-1-3/}{b} \\
\textsc{ii}. \HT{hasey si}{\B{ha\cp{sey} si}} \I{[ha."{}sEj si]} \emph{vin.} accomplish, bring to a conclusion. 
\B{Nì\-'i'\-a po tsa\-tì\-kang\-kem\-vir ha\-sey so\-li.} She finally completed the project.~\LINK{http://naviteri.org/2011/04/\%E2\%80\%99a\%E2\%80\%99awa-li\%E2\%80\%99fyavi-amip\%E2\%80\%94a-few-new-expressions/}{b} 
\B{Sä\-tsway\-on a\-'aw\-ve tsa\-heyl\-ur ha\-sey si; ke tsun nga pi\-vey.} First flight seals the bond; you cannot wait.~\LINK{http://naviteri.org/2012/06/spring-vocabulary-part-3/}{b}
\B{Ha\-sey si fu\-ra yom!} Finish eating!~\LINK{http://naviteri.org/2020/11/vospxivopeya-ayliu-amip-novembers-new-words/}{b}
(\emph{a. idiom}) \B{Sey\-so\-nìl\-\cp{tsan}!} \I{[sEj.so.nIl."\t{ts}an]} Well done! (\emph{from} \B{ha\-sey so\-li nìl\-tsan})~\LINK{http://naviteri.org/2011/04/\%E2\%80\%99a\%E2\%80\%99awa-li\%E2\%80\%99fyavi-amip\%E2\%80\%94a-few-new-expressions/}{b} (\ding{43}~\HL{mu'ni}{\B{mu'\-ni}}) \\
\textsc{Derivatives}: \ding{43} \on{1}.~\HL{seysonIltsan}{\B{sey\-so\-nìl\-tsan}}, well done, bravo.

% haway
\hi 
\HT{haway}{\B{\cp{ha}way}} \I{["ha.waj]} \emph{n.} lullaby. 
(\ding{43}~\HL{hahaw}{\B{ha\-haw}}, \HL{way}{\B{way}})

% hawl
\hi
\HT{hawl}{\B{hawl}} \I{[hawl]} \emph{vtr.} prepare. 
\B{Ta\-fral ho\-lawl ay\-oel ay\-nga\-fpi 'a\-'aw\-a tì\-päng\-kxo\-tsyìp\-it.} That's why we've prepared several little dialogs for you.~\LINK{http://naviteri.org/2012/07/tskxekengtsyip-a-mikyunfpi-a-little-listening-exercise/}{b} \\
\textsc{Derivatives}: \ding{43} \on{1}.~\HL{tIhawl}{\B{tì\-hawl}}, preparation, plan.

% hawmpam
\hi
\HT{hawmpam}{\B{\cp{hawm}pam}} \I{["hawm.pam]} \emph{n.} noise (\emph{sound that is excessive, unnecessary, inappropriate, unexpected, or startling})
\B{Fì\-hawm\-pam pe\-lun, ma 'itan? Fni\-vu set!} Why all this noise, son? Be quiet now!~\LINK{http://naviteri.org/2014/11/twenty-before-the-holidays/}{b} 
(\ding{43}~\HL{vApam}{\B{vä\-pam}}, \HL{hawng}{\B{hawng}}, \HL{pam}{\B{pam}}) \\
\textsc{Derivatives}: \ding{43} \on{1}.~\HL{lehawmpam}{\B{le\-hawm\-pam}}, noisy. 
\on{2}.~\HL{nIhawmpam}{\B{nì\-hawm\-pam}}, noisily.

% hawntu
\hi 
\HT{hawntu}{\B{\cp{hawn}tu}} \I{["hawn.tu]} \emph{n.} one under s.o.'s protection, protected person. 
\B{Oey yawn\-tu lu oey hawn\-tu.} My beloved is under my protection.~\LINK{https://naviteri.org/2024/03/fleltrra-ayliu-words-for-april-fools-day/}{b}
(\ding{43}~\HL{hawnu}{\B{hawnu}})

% hawntsyokx
\hi
\HT{hawntsyokx}{\B{hawn\cp{tsyokx}}} \I{[hawn."\t{ts}jok']} \emph{n.} glove. (\emph{plurals in all but very careful or ceremonial conversation usually pronounced} \B{me\-hawn\-\cp{tsyokx}} \I{[mawn."\t{ts}jok']}, \B{pxe\-hawn\-\cp{tsyokx}} \I{[p'awn."\t{ts}jok']}, \B{ay\-hawn\-\cp{tsyokx}} \I{[aj.awn."\t{ts}jok']})
(\ding{43}~\HL{hawnu}{\B{haw\-nu}}, \HL{tsyokx}{\B{tsyokx}})

% hawnu
\hi
\HT{hawnu}{\B{\cp{haw}nu}} \I{["haw.nu]} \emph{vtr.} protect, shelter (\B{wä}, from). 
\B{Fì\-fne\-ya\-yol tsrul\-it txu\-la mì song\-ropx ut\-ral\-ä fte ay\-li\-nit hi\-vaw\-nu wä sar\-ni\-o\-ang.} This kind of bird builds its nest in a tree cavity so as to protect its young from predators.~\LINK{http://naviteri.org/2013/04/wheres-the-bathroom-and-other-useful-things/}{b}
(\ding{43}~\HL{tIhawnu}{\B{tì\-haw\-nu si}}) \\
\textsc{Derivatives}: \ding{43} \on{1}.~\HL{tIhawnu}{\B{tì\-haw\-nu}}, protection. 
\on{2}.~\HL{hawntsyokx}{\B{hawn\-tsyokx}}, glove. 
\on{3}.~\HL{hawnven}{\B{hawn\-ven}}, shoe. 
\on{4}.~\HL{hawre'}{\B{haw\-re'}}, hat.
\on{5}.~\HL{hawntu}{\B{hawn\-tu}}, protected person.

% hawnven
\hi
\HT{hawnven}{\B{hawn\cp{ven}}} \I{[hawn."vEn]} \emph{n.} shoe. (\emph{plurals in all but very careful or ceremonial conversation usually pronounced} \B{me\-hawn\-\cp{ven}} \I{[mawn."vEn]}, \B{pxe\-hawn\-\cp{ven}} \I{[p'awn."vEn]}, \B{ay\-hawn\-\cp{ven}} \I{[aj.awn."vEn]}) 
\B{Rutxe me\-hawn\-ven\-it 'i\-va\-ku.} Please take off your shoes.~\LINK{http://naviteri.org/2011/08/new-vocabulary-clothing/}{b} 
\B{Hawn\-ven\-ì\-ri lu oe\-ru mun\-sna a\-mrr.} I have five pairs of shoes.~\LINK{http://naviteri.org/2012/07/fivospxiya-aylifyavi-amip-this-months-new-expressions-pt-2/}{b}
(\ding{43}~\HL{hawnu}{\B{haw\-nu}}, \HL{venu}{\B{ve\-nu}})

% hawng
\hi
\HT{hawng}{\B{hawng}} \I{[hawN]} \emph{n.} overabundance. \\
\textsc{Derivatives}: \ding{43} \on{1}.~\HL{nIhawng}{\B{nìhawng}}, too, excessively. 
\on{2}.~\HL{lehawng}{\B{lehawng}}, excessive. 
\on{3}.~\HL{hawngkrr}{\B{hawngkrr}}, late (\emph{adv.}). 
\on{4}.~\HL{hawmpam}{\B{hawmpam}}, noise.

% hawngkrr
\hi
\HT{hawngkrr}{\B{\cp{hawng}krr}} \I{["hawN.k\s{r}]} \emph{adv.} late.   
\B{Hawng\-krr rä\-'ä zi\-va\-'u!} Don't come late!~\LINK{http://naviteri.org/2010/07/vocabulary-update/}{b} 
(\ding{214}~\HL{ye'krr}{\B{ye'krr}}) \\
\textsc{Derivatives}: \ding{43} \on{1}.~\HL{lehawngkrr}{\B{lehawngkrr}}, late (\emph{adj.}).

% hawre'
\hi
\HT{hawre'}{\B{haw\cp{re'}}} \I{[haw."REP]} \emph{n.} hat. (\emph{plurals in all but very careful or ceremonial conversation usually pronounced} \B{me\-haw\-\cp{re'}} \I{[maw."REP]}, \B{pxe\-haw\-\cp{re'}} \I{[p'aw."REP]}, \B{ay\-haw\-\cp{re'}} \I{[aj.aw."REP]}) 
\B{Su\-nu oer haw\-re' a ngal yo\-lem\-stokx.} I like the hat you’re wearing.~\LINK{http://naviteri.org/2011/08/new-vocabulary-clothing/}{b} 
\B{Fì\-re\-won ngal lum\-pe kea ha\-wre'\-it ke yo\-lem\-stokx?} Why didn’t you put on a hat this morning?~\LINK{http://naviteri.org/2011/08/new-vocabulary-clothing/}{b}
(\ding{43}~\HL{hawnu}{\B{haw\-nu}}, \HL{re'o}{\B{re\-'o}})

% hawtsyìp
\hi
\HT{hawtsyIp}{\B{\cp{haw}tsyìp}} \I{["haw.\t{ts}jIp\textcorner]} \emph{n.} nap. 
\B{Oel new fu\-ta li\-vu oer set haw\-tsyìp.} I want to take a nap now.~\LINK{http://naviteri.org/2011/07/txantsana-ultxa-mi-siatll-great-meeting-in-seattle/}{b} 
\B{Sa'\-nur le\-ru haw\-tsyìp.} Mommy is taking a nap.~\LINK{http://naviteri.org/2011/07/txantsana-ultxa-mi-siatll-great-meeting-in-seattle/}{b}
(\ding{43}~\HL{hahaw}{\B{ha\-haw}})

% hay
\hi
\HT{hay}{\B{hay}} \I{[haj]} \emph{adj.} next, following. 
\B{Sìl\-pey oe, tsun oeng ul\-txa si\-vi ro hay\-a ul\-txa!} I hope, we can meet at the next meeting.~\LINK{http://naviteri.org/2011/07/txantsana-ultxa-mi-siatll-great-meeting-in-seattle/comment-page-1/\#comment-844}{b}   
\B{Alo a\-hay oe za\-sya\-'u mì re\-won}, next time I'll come in the morning.~\LINK{http://forum.learnnavi.org/language-updates/txelanit-hivawl/}{ln} 
(\ding{214}~\HL{ham}{\B{ham}}) \\
\textsc{Derivatives}: \ding{43} \on{1}.~\HL{hayalo}{\B{ha\-ya\-lo}}, next time.

% hayalo
\hi
\HT{hayalo}{\B{ha\cp{ya}lo}} \I{[ha."ja.lo]} \emph{adv.} next time. (\emph{a. idiom}) \B{ha\-ya\-lo ta oe} \emph{or} \B{hay\-a\-lo oe\-ta}~\I{[ha."ja.lo "wE.ta]}. You're welcome [lit.: next time from me] (\emph{phrase for responding to thanks}).~\LINK{http://naviteri.org/2010/07/diminutives-conversational-expressions/}{b} 
(\emph{b. idiom}) \B{slä ha\-ya\-lo a\-la\-he} ``\ldots but that is for another time'' (\emph{set phrase in storytelling})~\LINK{http://naviteri.org/2020/04/aawa-u-amip-a-few-new-things/}{b}  \\
(i.) \B{ha\-\cp{ya}\-lo\-vay} \I{[ha."ja.lo.vaj]}, until next time.~\LINK{http://naviteri.org/2010/07/diminutives-conversational-expressions/}{b} 
(\ding{214}~\HL{hamalo}{\B{ha\-ma\-lo}}, \ding{43}~\HL{hay}{\B{hay}}, \HL{alo}{\B{alo}})

% hän
\hi
\HT{hAn}{\B{hän}} \I{[h\ae{}n]} \emph{n.} net; web.
\B{Tsun fko si\-var hän\-it fte pay\-oang\-it sti\-vä'\-nì.} One can use a net to catch a fish.~\LINK{http://naviteri.org/2015/04/some-new-words-for-may-day/}{b}
\B{Fì\-hì\-'ang\-ìl txu\-la hän\-ti fte smar\-it syi\-vep.} This insect constructs a web to trap its prey.~\LINK{http://naviteri.org/2015/04/some-new-words-for-may-day/}{b}

% he'a
\hi 
\HT{he'a}{\B{\cp{he}'a}} \I{["hE.Pa]} \emph{vin.} cough. 
\B{Nga\-ri krr\-a he\-'a, swey\-lu txo kxa\-ru lew si\-vi.} When you cough, it's best to cover your mouth.~\LINK{http://naviteri.org/2020/07/vospxivol-lefpom-happy-august/}{b} \\
\textsc{Derivatives}: \ding{43} \on{1}.~\HL{sAhe'a}{\B{sä\-he\-'a}}, a cough.

% hefi
\hi
\HT{hefi}{\B{\cp{he}fi}} \I{["hE.fi]} \emph{vtr.} smell (-\emph{control}). 
\B{Oel he\-fi ye\-rik\-it!} I smell a hexapede!~\LINK{http://naviteri.org/2012/11/renu-ayinanfyaya-the-senses-paradigm/}{b} 
(\ding{43}~\HL{syam}{\B{syam}}, \HL{ontu}{\B{tìng on\-tu}}) \\
\textsc{Derivatives}: \ding{43} \on{1}.~\HL{hefitswo}{\B{he\-fi\-tswo}}, sense of smell.

% hefitswo
\hi
\HT{hefitswo}{\B{\cp{he}fitswo}} \I{["hE.fi.\t{ts}wo]} \emph{n.} sense of smell. 
(\ding{43}~\HL{hefi}{\B{he\-fi}})

% hek
\hi
\HT{hek}{\B{hek}} \I{[hEk\textcorner]} \emph{vin.} be curious, odd, strange, unexpected. 
\B{Nge\-yä sä\-fpìl Saw\-tu\-te\-teri he\-i\-ek oer nì\-txan.} Your idea about the Sky People is very interesting to me (because it seems unusual).~\LINK{http://naviteri.org/2010/07/vocabulary-update/}{b} 
\B{Fo nì\-Na'\-vi pll\-txe nì\-fya\-'o a hek.} They speak Na'vi strangely.~\LINK{http://naviteri.org/2010/07/vocabulary-update/}{b} 
\B{hu\-fwa lì\-'u\-pam hek nì\-'it}, although the pronunciation is a bit strange.~\LINK{http://naviteri.org/2010/07/ting-mikyun-fte-tslivam-listening-comprehension-1-3/}{b} 
\B{Ay\-ron\-srel pe\-yä hä\-ngek nì\-txan.} His imaginings are (unpleasantly) weird.~\LINK{http://naviteri.org/2011/02/new-vocabulary-part-2/}{b} \\
\textsc{Derivatives}: \ding{43} \on{1}.~\HL{nIhek}{\B{nì\-hek}}, strangely, oddly.

% hena
\hi
\HT{hena}{\B{\cp{he}na}} \I{["hE.na]} \emph{vtr.} carry.
\B{Ru\-txe hi\-ve\-na fì\-e\-pxang\-it fpi oe.} Please carry this stone jar for me.~\LINK{http://naviteri.org/2014/05/mipa-ayliu-mipa-aysafpil-new-words-new-ideas/}{b} \\
\textsc{Derivatives}: \ding{43} 
\on{1}.~\HL{sAhena}{\B{sä\-he\-na}}, container, vessel, carrier. 
\on{2}.~\HL{kahena}{\B{ka\-he\-na}}, transport.

% henga / si
\hi 
\HT{henga}{\B{\cp{he}nga}} \I{["hE.Na]} \textsc{i}. \emph{n.} formally polite behavior. 
\B{He\-nga kel\-kin.} Formality isn't necessary.~\LINK{http://naviteri.org/2022/02/lifyengteri-concerning-honorific-language/}{b} \\
\textsc{ii}. \B{\cp{he}nga si} \I{["hE.Na si]} \emph{vin.} act in a formally polite or honorific way.
\B{Krr\-a ul\-txa si nga tsa\-txan\-ro'\-tu\-hu, ze\-rok fu\-ta ze\-ne he\-nga si\-vi, ma 'e\-veng.} When you meet that famous person, remember that you have to be formally polite, child.~\LINK{http://naviteri.org/2022/02/lifyengteri-concerning-honorific-language/}{b} 
\B{He\-nga rä\-'ä si, ma tsmuk.} Don't be so formal, bro\slash sis.~\LINK{http://naviteri.org/2022/02/lifyengteri-concerning-honorific-language/}{b} \\
\textsc{Derivatives}: \ding{43} \on{1}.~\HL{leheng}{\B{leheng}}, formally polite. 

% hermeyp
\hi
\HT{hermeyp}{\B{her\cp{meyp}}} \I{[hER."mEjp\textcorner]} \emph{n.} snow flurries. 
(\ding{43}~\HL{herwI}{\B{her\-wì}}, \HL{meyp}{\B{meyp}})

% hertxayo
\hi
\HT{hertxayo}{\B{\cp{her}txayo}} \I{["hER.t'a.jo]} \emph{n.} field of snow. 
(\ding{43}~\HL{herwI}{\B{her\-wì}}, \HL{txayo}{\B{txa\-yo}})

% herwì
\hi
\HT{herwI}{\B{\cp{her}wì}} \I{["hER.wI]} \emph{n.} snow. \B{Her\-wì ze\-re\-i\-up fì\-trr\-o nì\-wotx!} It's been snowing all day!~\LINK{http://naviteri.org/2011/04/yafkeykiri-plltxe-frapo-everyone-talks-about-the-weather/}{b} \\
\textsc{Derivatives}: \ding{43} 
\on{1}.~\HL{herwIva}{\B{her\-wì\-va}}, snowflake.  
\on{2}.~\HL{hermeyp}{\B{her\-meyp}}, snow flurries. 
\on{3}.~\HL{hertxayo}{\B{her\-txa\-yo}}, field of snow. 
\on{4}.~\HL{txanfwerwI}{\B{txan\-fwer\-wì}}, blizzard. 

% herwìva
\hi
\HT{herwIva}{\B{\cp{her}wìva}} \I{["hER.wI.va]} \emph{n.} snowflake.
(\ding{43}~\HL{herwI}{\B{her\-wì}}, \HL{Ilva}{\B{ìl\-va}}) 

% hewne
\hi
\HT{hewne}{\B{\cp{hew}ne}} \I{["hEw.nE]} \emph{adj.} soft (\emph{of an object}). 
(\ding{214}~\HL{txa'}{\B{txa'}}) 

% heyn
\hi
\HT{heyn}{\B{heyn}} \I{[hEjn]} \emph{vin.} sit. 
(\emph{a.}) \B{aw\-ngal yom wu\-tsot teng\-krr he\-reyn nì\-pxim}, we eat dinner while sitting upright.~\LINK{http://whyisthisnight.com/addl-more.html\#na\%27vi}{tn}
\B{Lä\-ngu fì\-seyn kel\-ho\-an nì\-ngay. Tsun oe hi\-veyn tsaw\-sìn 'a\-'aw\-a swaw\-tsyìp nì\-'aw.} This chair is really uncomfortable. I can only sit on it a few seconds.~\LINK{http://naviteri.org/2015/08/aylifyavi-lereyfya-2-cultural-terms-2-and-more/}{b}
(\emph{b. proverb}) \B{Txìm a\-'aw ke tsun hi\-veyn mì tal me\-fa'\-li\-yä.} You can't sit on the fence; you need to decide [lit.: One butt can't sit on the backs of two direhorses].~\LINK{http://naviteri.org/2013/02/vospxi-ayol-posti-apup-short-post-for-a-short-month/}{b} \\
\textsc{Derivatives}: \ding{43} 
\on{1}.~\HL{seyn}{\B{seyn}}, chair, stool, bench.
\on{2}.~\HL{heynyI}{\B{heyn\-yì}}, lap.

% heynyì
\hi 
\HT{heynyI}{\B{\cp{heyn}yì}} \I{["hEjn.jI]} \emph{n.} lap. 
\B{Za\-'u heyn sìn oey heyn\-yì, ma 'e\-vi.} Come sit on my lap, my child.~\LINK{https://naviteri.org/2024/06/ayhapxi-tokxa-si-ayliu-alahe-parts-of-the-body-and-more/}{b}
(\ding{43}~\HL{heyn}{\B{heyn}}, \HL{yI}{\B{yì}}) 

% heyr
\hi
\HT{heyr}{\B{heyr}} \I{[hEjR]} \emph{n.} chest (\emph{of human or animal body}), \emph{the area between the stomach and throat}. 
\B{Oe\-ri heyr tì\-sraw sä\-ngi ta\-lu\-na zi\-ze'\-ìl oet sngo\-lap tsa\-tseng.} My chest hurts because a hellfire wasp stung me there.~\LINK{http://naviteri.org/2015/03/kap-si-ayunil-saylahe-threats-dreams-and-other-things/}{b}
\B{Po\-an\-it tso\-lukx po\-el fa tstal nì\-txukx nem\-fa heyr.} She stabbed him deeply in the chest with a knife.~\LINK{http://naviteri.org/2019/06/50a-liu-amip-40-new-words/}{b}

% hifwo
\hi
\HT{hifwo}{\B{\cp{hi}fwo}} \I{["hi.fwo]} \emph{vin.} flee, escape (\B{ftu}, from). \B{Pa\-lu\-kan\-it tso\-le\-'a, ye\-rik lopx hi\-fwo kxam\-lä ze\-swa.} Spotting a thanator, the hexapede panicked and escaped through the grass.~\LINK{http://naviteri.org/2012/07/meetings-waterfalls-and-more/}{b}
\B{Tseyk 'ä\-pey\-ka\-mrr\-ko äo ut\-ral a zo\-lup fte hi\-vi\-fwo ftu ay\-sre' pa\-lu\-lu\-kan\-ä.} Jake rolled under the fallen tree to escape from the thanator's teeth.~\LINK{http://naviteri.org/2015/04/some-new-words-for-may-day/}{b}
(\emph{a. idiom}) \B{Hì\-pey ta\-ron\-yu, hi\-fwo ye\-rik.} ``He who hesitates is lost.'' [lit.: The hunter hesitates and the hexapede escapes]~\LINK{http://naviteri.org/2015/11/vomuna-liu-amip-ten-new-words/}{b}

% hilu
\hi 
\HT{hilu}{\B{\cp{hi}lu}} \I{["hi.lu]} \emph{n.} nipple. 

% -hin
\hi
\HT{hin}{\B{-hin}} \I{[.hin]} \emph{num., suf.} seven (\emph{rest of} \ding{43}~\HL{kinA}{\B{ki\-nä}} \emph{to form higher numbers}). \B{vo\-hin}, \on{15} (decimal), \on{17} (octal). \B{me\-vo\-hin}, \on{23} (decimal), \on{27} (octal); etc.

% hipx
\hi 
\HT{hipx}{\B{hipx}} \I{[hip']} \emph{vtr.} control. 
\B{Kar\-yu a\-sìl\-tsan ze\-ne tsi\-vun ay\-nu\-me\-yut hi\-vipx mì num\-tseng\-vi.} A good teacher has to be able to control (his/her) students in the classroom.~\LINK{http://naviteri.org/2018/04/100a-liu-amip-64-new-words-part-2/}{b} \\
\textsc{Derivatives}: \ding{43} \on{1}.~\HL{tIhipx}{\B{tì\-hipx}}, control.

% hiup
\hi 
\HT{hiup}{\B{\cp{hi}up}} \I{["hi.up\textcorner]} \emph{or} \I{[.Up\textcorner]} (\textsc{rn}: \B{\cp{hi}ùp} \I{["hi.Up\textcorner]}) (\emph{a. vin.}) spit.
\B{Tsa\-kem rä\-'ä si, ma 'itan. Fwa hi\-up fì\-tseng ke lu muiä.} Don't do that, son. It's not proper to spit here.~\LINK{http://naviteri.org/2016/06/mrrvola-lifyavi-amip-forty-new-expressions/}{b}
(\emph{b. vtr.}) spit out. \\
\textsc{Derivatives}: \ding{43} \on{1}.~\HL{hiupwopx}{\B{hi\-up\-wopx}}, cloud spitter.

% hiupwopx
\hi 
\HT{hiupwopx}{\B{\cp{hi}upwopx}} \I{["hi.up\textcorner.wop']} (\textsc{rn}: \B{\cp{hi}ùpwopx} \I{["hi.Up\textcorner.wop']}) \emph{n.} \ding{96} cloud spitter. 
(\ding{43}~\HL{hiup}{\B{hiup}}, \HL{pIwopx}{\B{pì\-wopx}})

% hiyìk
\hi
\HT{hiyIk}{\B{\cp{hi}yìk}} \I{["hi.jIk\textcorner]} \emph{adj.} funny, strange. 

% hì'ang
\hi
\HT{hI'ang}{\B{\cp{hì}'ang}} \I{["hI.PaN]} \emph{n.} insect. 
\B{Lu mek\-tseng a tsun fpxi\-vä\-kìm hì\-'ang tsaw\-ìlä.} There's a gap through which insects can get in.~\LINK{http://naviteri.org/2013/04/wheres-the-bathroom-and-other-useful-things/}{b}
(\ding{43}~\HL{hI'i}{\B{hì\-'i}}, \HL{ioang}{\B{ioang}})

% hì'i
\hi
\HT{hI'i}{\B{\cp{hì}'i}} \I{["hI.Pi]} \emph{adj.} little, small (\emph{in size}). 
\B{Fì\-syu\-lang a\-rim lu hì\-'i fra\-to.} This yellow flower is the smallest of all.~\LINK{http://naviteri.org/2010/08/mipa-ayopin-mipa-ayliu-new-colors-new-words/}{b} 
\B{su\-te tsun oe\-yä hì\-'i\-a tì\-ngop\-it si\-var}, people can use my strange/funny little creation.~\LINK{http://forum.learnnavi.org/language-updates/score!-\%281\%29-the-verb-to-use-\%282\%29-more-detail-on-teya-si/}{ln} 
\B{win\-a sä\-wä\-sul\-tsyìp a\-hì\-'i}, a quick little contest.~\LINK{http://naviteri.org/2012/10/wina-sawasultsyip-ahii-a-quick-little-contest/}{b} 
(\ding{214}~\HL{tsawl}{\B{tsawl}}) \\
\textsc{Derivatives}: \ding{43} \on{1}.~\HL{tsawlhI'}{\B{tsawl\-hì'}}, size. 
\on{2}.~\HL{hI'ang}{\B{hì\-'ang}}, insect.   
\on{3}.~\HL{hIkrr}{\B{hì\-krr}}, second, very short time. 
\on{4}.~\HL{hItxoa}{\B{hì\-txoa}}, excuse me. 
\on{5}.~\HL{tanhI}{\B{tan\-hì}}, star. 
\on{6}.~\HL{hIno}{\B{hì\-no}}, fine, detailed. 
\on{7}.~\HL{hIpey}{\B{hì\-pey}}, hesitate, hold back. 
\on{8}.~\HL{hIrumwll}{\B{hì\-rum\-wll}}, puffer plant.

% hìkrr
\hi
\HT{hIkrr}{\B{\cp{hì}krr}} \I{["hI.k\s{r}]} \textsc{i}. \emph{n.} second, a very short time, a moment. 
\B{Maw hì\-krr ay\-oe tì\-yä\-txaw.} We'll return after a short time.~\LINK{http://wiki.learnnavi.org/index.php?title=Corpus\#Good_Morning_America}{wiki} \\
\textsc{ii}.~\emph{adv.} for a moment, for a second. 
\B{Tsun mi\-vä\-kxu hì\-krr srak?} May I interrupt for a moment?~\LINK{http://naviteri.org/2010/09/getting-to-know-you-part-1/}{b} 
(\ding{43}~\HL{hI'i}{\B{hì\-'i}}, \HL{krr}{\B{krr}})

% hìm
\hi
\HT{hIm}{\B{hìm}} \I{[hIm]} \emph{adj.} small (\emph{in quantity}), a little. 
(\ding{214}~\HL{txan}{\B{txan}}) \\
\textsc{Derivatives}: \ding{43} \on{1}.~\HL{hImtxan}{\B{hìm\-txan}}, amount. 
\on{2}.~\HL{hImpxI}{\B{hìm\-pxì}}, minority. 
\on{3}.~\HL{hImtxew}{\B{hìm\-txew}}, minimum. 

% hìmpxì
\hi
\HT{hImpxI}{\B{hìm\cp{pxì}}} \I{[hIm."p'I]} \emph{n.} minority, least, small part. 
\B{Hìm\-pxì Saw\-tu\-te\-yä lu tstu\-nwi}, A minority of the Sky People are kind.~\LINK{http://naviteri.org/2011/11/\%E2\%80\%99a\%E2\%80\%99awa-ayli\%E2\%80\%99u-amip-ni\%E2\%80\%99aw-\%E2\%80\%94-a-few-new-words-only/}{b} 
(\ding{214}~\HL{txampxI}{\B{txam\-pxì}}, \ding{43}~\HL{hIm}{\B{hìm}}, \HL{hapxI}{\B{ha\-pxì}})

% hìmtxampe
\hi
\HT{hImtxampe}{\B{hìm\cp{txam}pe}} \I{[hIm."t'am.pE]} \emph{or} \B{pìm\cp{txan}} \I{[pIm.\-"t'an]} \emph{inter.} how much? 
\B{Fì\-nik\-roi ley pìm\-txan?} How valuable is this hair ornament?~\LINK{http://naviteri.org/2014/03/value-and-worth/}{b}
(\ding{43}~\HL{hImtxan}{\B{hìm\-txan}})

% hìmtxan
\hi
\HT{hImtxan}{\B{hìm\cp{txan}}} \I{[hIm."t'an]} \emph{n.} amount. \\
\textsc{Derivatives}: \ding{43} \on{1}.~\HL{hImtxampe}{\B{hìm\-txam\-pe}\slash \B{pìm\-txan}}, how much? 

% hìmtxew
\hi
\HT{hImtxew}{\B{hìm\cp{txew}}} \I{[hIm."t'Ew]} \emph{n.} minimum. 
(\ding{214}~\HL{txantxew}{\B{txan\-txew}}) \\
\textsc{Derivatives}: \ding{43} \on{1}.~\HL{hItxewvay}{\B{hìm\-txew\-vay}}, minimally.

% hìmtxewvay
\hi
\HT{hImtxewvay}{\B{hìm\cp{txew}vay}} \I{[hIm."t'Ew.vaj]} \emph{adv.} minimally. 
(\ding{43}~\HL{hImtxew}{\B{hìm\-txew}})

% hìno
\hi
\HT{hIno}{\B{hì\cp{no}}} \I{[hI."no]} \emph{adj.} fine, detailed, precise (\emph{of things}). 
\B{Pol ngop fra\-krr sì\-ke\-nong\-it a hì\-no lu nì\-hawng.} He always creates excessively detailed examples.~\LINK{http://naviteri.org/2010/09/getting-to-know-you-part-2/}{b} 
\B{Fì\-way\-ri hì\-no\-a re\-nut ngam\-pam\-ä ke tsä\-ngun oe tsl\-ivam.} I'm afraid I can't understand the intricate rhyme scheme of this poem.~\LINK{http://naviteri.org/2012/01/mipa-zisit-ayliu-amip-new-words-for-the-new-year/}{b} 
(\ding{43}~\HL{leno}{\B{le\-no}}, \HL{nIno}{\B{nì\-no}})

% hìpey
\hi
\HT{hIpey}{\B{\cp{hì}pey}} \I{["hI.p$\cdot$Ej]} \emph{vin., inf.}\on{22} hesitate, hold back for a short time; \emph{defer the start of an action, reluctance to begin or accomplish an action, for whatever reason}.
\B{Fu\-ria peng fmawn\-it Ey\-tu\-kan\-ur po hì\-po\-ley.} He hesitated to tell Eytukan the news.~\LINK{http://naviteri.org/2015/11/vomuna-liu-amip-ten-new-words/}{b}
(\emph{a. idiom}) \B{Hì\-pey ta\-ron\-yu, hi\-fwo ye\-rik.} ``He who hesitates is lost.'' [lit.: The hunter hesitates and the hexapede escapes]~\LINK{http://naviteri.org/2015/11/vomuna-liu-amip-ten-new-words/}{b} 
(\emph{b. proverb}) \B{'Uo\-ri hì\-pey, kxawm nga\-ru ke ley.} ``If you hesitate doing something, it might not be important to you,'' i.e. \emph{more meant\slash used as a motivation for s.o. hesitating, or even as a teasing to get s.o. into action}.~\LINK{http://naviteri.org/2021/08/ulte-ayyoratu-leiu-and-the-winners-are-2/}{b}
(\ding{43}~\HL{hI'i}{\B{hì\-'i}}, \HL{pey}{\B{pey}}) \\
\textsc{Derivatives}: \ding{43} \on{1}.~\HL{tIhIpey}{\B{tì\-hì\-pey}}, hesitation. 
\on{2}.~\HL{lehIpey}{\B{le\-hì\-pey}}, hesitant, in a state of hesitation. 
\on{3}.~\HL{nIhIpey}{\B{nì\-hì\-pey}}, hesitantly.

% hìrumwll
\hi
\HT{hIrumwll}{\B{hì\cp{rum}wll}} \I{[hI."Rum.w\s{l}]} \emph{or} \I{[."RUm.]} \emph{n.}~\ding{96} puffer plant. 
(\ding{43}~\HL{hI'i}{\B{hì\-'i}}, \HL{rum}{\B{rum}}, \HL{'ewll}{\B{'e\-wll}})

% hìtxoa
\hi
\HT{hItxoa}{\B{hì\cp{txo}a}} \I{[hI."t'o.a]} \emph{idiom} excuse me, pardon me. \B{Hì\-txo\-a, ke tslo\-lam.} Sorry, I didn't get that.~\LINK{http://naviteri.org/2010/09/getting-to-know-you-part-2/}{b} (\ding{43}~\HL{hI'i}{\B{hì\-'i}}, \HL{txoa}{\B{txoa}})

% hoan
\hi
\HT{hoan}{\B{\cp{ho}an}} \I{["ho.an]} \emph{n.} comfort.
\B{Sko frr\-tu, fra\-krr mì hel\-ku nge\-yä lu oe\-ru ho\-an nì\-txan, ma tsmuk.} I always feel very comfortable as a guest in your home, brother.~\LINK{http://naviteri.org/2015/08/aylifyavi-lereyfya-2-cultural-terms-2-and-more/}{b} \\
\textsc{Derivatives}: \ding{43} \on{1}.~\HL{lehoan}{\B{le\-ho\-an}}, comfortable. 
\on{2}.~\HL{nIhoan}{\B{nì\-ho\-an}}, comfortably.

% hoet
\hi
\HT{hoet}{\B{\cp{ho}et}} \I{["ho.Et\textcorner]} \emph{adj.} vast, broad, expansive. \B{'Rr\-ta\-mì a tam\-pay lu ho\-et.} The oceans on Earth are vast.~\LINK{http://naviteri.org/2012/06/spring-vocabulary-part-3/}{b} 
\B{Lì'\-fya le\-Na'\-vi lu teng\-wotx, ken ay\-reym\-sä\-tsa le\-ke\-teng\-syon Ey\-wa\-'e\-veng\-ä, ken ho\-et\-a ta\-yo a mì\-kam ay\-olo'.} The language of the Na'vi is consistent, despite the exo\hyp{}moon's varied landscapes and the vast expanse between clans.~\LINK{https://naviteri.org/2026/06/tskxekeng-a-mikyunfpi-pamrel-listening-exercise-text/}{b} \\
\textsc{Derivatives}: \ding{43} \on{1}.~\HL{nIhoet}{\B{nì\-ho\-et}}, widely, pervasively.

% hol
\hi
\HT{hol}{\B{hol}} \I{[hol]} \emph{adj.} few. 
\B{Tsay\-lì\-'u lu hol nì\-'aw.} Those words are few.~\LINK{http://naviteri.org/2011/09/\%E2\%80\%9Cby-the-way-what-are-you-reading\%E2\%80\%9D/comment-page-1/\#comment-1117}{b} 
\B{Oel tse\-'a hol\-a tu\-tet.} I see only a few people.~\LINK{http://naviteri.org/2010/07/vocabulary-update/}{b} 
(\ding{214}~\HL{pxay}{\B{pxay}}) \\
\textsc{Derivatives}: \ding{43} \on{1}.~\HL{holpxay}{\B{hol\-pxay}}, number.

% holpxay
\hi
\HT{holpxay}{\B{hol\cp{pxay}}} \I{[hol."p'aj]} \emph{n.} number. 
\B{hol\-pxay ay\-zek\-wä\-yä feyä}, the number of their fingers.~\LINK{http://forum.learnnavi.org/language-updates/talun-taweyk-tafral-lun-oeyk-holpxay-hol/}{ln} 
\B{Ko\-ren Hol\-pxay\-ä}, Number Principle.~\LINK{http://naviteri.org/2011/07/number-in-na\%E2\%80\%99vi/}{b} 
(\ding{43}~\HL{hol}{\B{hol}}, \HL{pxay}{\B{pxay}})

% holpxaype
\hi
\HT{holpxaype}{\B{hol\cp{pxay}pe}} \I{[hol."p'aj.pE]}~\emph{or}~\B{pol\cp{pxay}}~\I{[pol."p'aj]} \emph{inter.} how many? (behaves like an ordinary \emph{adj.}) \B{Nga\-ri so\-la\-lew pol\-pxay\-a zì\-sìt?} How old are you?~\LINK{http://naviteri.org/2010/09/getting-to-know-you-part-2/}{b} 
(\ding{43}~\HL{holpxay}{\B{hol\-pxay}})

% hona
\hi
\HT{hona}{\B{\cp{ho}na}} \I{["ho.na]} \emph{adj.} endearing, adorable, cute. 
\B{Pe\-yä 'i\-tan\-ìri lu ho\-na nì\-txan.} His son is very cute.~\LINK{http://naviteri.org/2012/01/more-additions-to-the-lexicon/}{b} 
\B{ho\-na pa\-lu\-kan\-tsyìp a\-hì\-'i}, sweet little cat.~\LINK{http://naviteri.org/2012/01/more-additions-to-the-lexicon/}{b} 
\B{Prr\-nen lrr\-tok si a krr, srer me\-song\-tsyìp a\-ho\-na.} When the baby smiles, two adorable dimples appear.~\LINK{http://naviteri.org/2012/03/spring-vocabulary-part-1/}{b} \\
\textsc{Derivatives}: \ding{43} \on{1}.~\HL{nIhona}{\B{nì\-ho\-na}}, sweetly, endearingly.  
\on{2}.~\HL{tIhona}{\B{tì\-ho\-na}}, cuteness, adorableness.

% horenleykekyu
\hi 
\HT{horenleykekyu}{\B{ho\cp{ren}ley\cp{kek}yu}} \I{[ho.""ren.lEj."kEk\textcorner.ju]} \emph{n.} law enforcer, (\emph{on Earth}) police officer. (\emph{often shortened to} \ding{43}~\HL{leykekyu}{\B{ley\-kek\-yu}}) 
\B{Kin\-trr\-am ay\-oe tar\-mìng na\-ri teng\-krr 'aw\-a ho\-ren\-ley\-kek\-yul tu\-tan\-it a\-'aw tspe\-rang--lun\-lu\-ke nì\-wotx, nìk\-'ong, nì\-ze\-vakx.} Last week we watched while one law enforcer killed a man--without any reason, slowly, cruelly.~\LINK{http://naviteri.org/2020/06/mi-tanlokxe-oeya-srr-afpxamo-terrible-days-in-my-country/}{b} 
(\ding{43}~\HL{leykekyu}{\B{leykekyu}}, \HL{koren}{\B{ko\-ren}})

% hrrap
\hi
\HT{hrrap}{\B{\cp{hrr}ap}} \I{["h\s{r}.ap\textcorner]} \emph{n.} danger. \\
\textsc{Derivatives}: \ding{43} 
\on{1}.~\HL{lehrrap}{\B{le\-hrr\-ap}}, dangerous.     
\on{2}.~\HL{yrrap}{\B{yrr\-ap}}, storm. 
\on{3}.~\HL{penghrrap}{\B{peng\-hrr\-ap}}, binary sunshine. 
\on{4}.~\HL{penghrr}{\B{peng\-hrr}}, warn.

% hu
\hi
\HT{hu}{\B{hu}} \I{[hu]} \emph{adp.} with (\emph{together with s.o., accompaniment}). \B{Nga\-ri hu Ey\-wa sa\-lew ti\-re\-a, tokx 'ì\-'awn slu Na'\-vi\-yä ha\-pxì.} Your spirit goes with Eywa, your body remains to become part of the People.~\LINK{http://wiki.learnnavi.org/index.php?title=Corpus\#Behind_the_Scenes}{a\slash wiki} \B{Lar\-mu tsa\-tsam\-si\-yu\-hu tì\-vawm a lu stxong ay\-oer.} That warrior carried with him a darkness unknown to us.~\LINK{http://naviteri.org/2010/07/vocabulary-update/}{b} 
(\emph{a. idiom}) \B{Ey\-wa nga\-hu.} Ey\-wa (be) with you. 
(\emph{b. with verbs}) \HL{pAngkxo}{\B{päng\-kxo hu}}, speak\slash chat with; 
\HL{kAteng}{\B{kä\-teng hu}}, spend time with; 
\HL{mllte}{\B{mll\-te hu}}, agree with; 
\HL{wem}{\B{wem hu}}, fight (together) with (against s.o.~else); 
\HL{ultxa}{\B{ul\-txa si hu}}, meet (intentionally) with; 
\HL{tsaheylu}{\B{tsa\-heyl si hu}}, bond with, establish the neural connection with; 
\HL{uvan}{\B{u\-van si hu}}, play with s.o., 
\HL{nawang}{\B{na\-wang hu}}, merge, become one with.
\HL{nalsteng}{\B{nal\-steng hu}}, feel empthay\slash compassion for s.o.\\
\textsc{Derivatives}: \ding{43} 
\on{1}.~\HL{huslew}{\B{hu\-slew}}, accompany. 

% hufwa
\hi
\HT{hufwa}{\B{hu\cp{fwa}}} \I{[hu."fwa]} \emph{conj.} although. \B{hu\-fwa lì\-'u\-pam hek nì\-'it}, although the pronunciation is a bit strange.~\LINK{http://naviteri.org/2010/07/ting-mikyun-fte-tslivam-listening-comprehension-1-3/}{b} \B{Fì\-ut\-ral\-ìri ta\-ngek\-ä zir fkan vawt, hu\-fwa ke rey.} The trunk of the tree feels solid, although it's dead.~\LINK{http://naviteri.org/2013/09/on-si-salewfya-shapes-and-directions/}{b} \B{Hu\-fwa ro\-lun oel 'a\-'aw\-a kxey\-ey\-ti, fì\-tì\-kang\-kem\-vi lu txan\-tsan ka wotx.} Although I found a few errors, this piece of work is generally excellent.~\LINK{http://naviteri.org/2012/07/fivospxiya-aylifyavi-amip-this-months-new-expressions-pt-2/}{b}

% hufwe
\hi
\HT{hufwe}{\B{hu\cp{fwe}}} \I{[hu."fwE]} \emph{n.} (\emph{a.}) wind. 
\B{ay\-nga\-ti spi\-vu\-le hu\-fwel}, the wind may propel you.~\LINK{http://forum.learnnavi.org/language-updates/naviya-the-other-vocative/}{a\slash ln} \B{Tì\-ran hu\-fwe.} It's breezy (but pleasantly so). 
\B{Tul hu\-fwe nì\-win.} It's very windy.~\LINK{http://naviteri.org/2011/05/weather-part-2-and-a-bit-more-2/}{b}
(\emph{b. expression}) \B{nì\-win na hu\-fwe}, ``as fast as the wind'' (\emph{any situation to express rapidity}). \B{na hu\-fwe}, fluent (not just rapid) use of language.
\B{Fte\-ria oel lì'\-fya\-ti le\-Na'\-vi, slä mi ke tsä\-ngun pi\-vll\-txe na hu\-fwe.} I'm studying Na'vi, but I'm afraid I still can't speak it fluently.~\LINK{http://naviteri.org/2015/08/aylifyavi-lereyfya-2-cultural-terms-2-and-more/}{b} \\
\textsc{Derivatives}: \ding{43} \on{1}.~\HL{hufwetsyIp}{\B{hu\-fwe\-tsyìp}}, breeze, light wind. 
\on{2}.~\HL{txanfwerwI}{\B{txan\-fwer\-wì}}, blizzard. 
\on{3}.~\HL{nawfwe}{\B{naw\-fwe}}, fluent (for speech).

% hufwetsyìp
\hi
\HT{hufwetsyIp}{\B{hu\cp{fwe}tsyìp}} \I{[hu."fwE.\t{ts}jIp\textcorner]} \emph{n.} breeze, light wind.
\B{Hu\-fwe\-tsyìp le\-fpom tar\-mì\-ran.} A pleasant little breeze was blowing.~\LINK{http://naviteri.org/2011/05/weather-part-2-and-a-bit-more-2/}{b} 
\B{Rìk a\-ukxo mì hu\-fwe\-tsyìp ka\-fi sar\-mi.} The dry leaf was sailing in the breeze.~\LINK{http://naviteri.org/2023/12/mipa-zisit-ayliu-amip-new-year-new-words/}{b}
(\ding{43}~\HL{hufwe}{\B{hu\-fwe}})

% hultstxem
\hi
\HT{hultstxem}{\B{hul\cp{tstxem}}} \I{[hul."\t{ts}t'Em]} \emph{or} \I{[hUl.]} (\textsc{rn}: \B{hùl\cp{tstxem}} \I{[hUl."\t{ts}t'Em]}) \emph{vtr.} (\emph{usually negative connotation}) hinder, be an obstacle (to a person or thing). \B{Ke new oel fu\-ta fì\-tì\-päng\-kxot ay\-nge\-yä hi\-vul\-tstxem, slä tsun mi\-vä\-kxu hì\-krr nì\-'aw srak?.} I don't want to derail your chat, but can I interrupt for just a moment?~\LINK{http://naviteri.org/2010/09/getting-to-know-you-part-1/}{b}
(\ding{43}~\HL{mAkxu}{\B{mä\-kxu}})

% hum
\hi
\HT{hum}{\B{hum}}\textsuperscript{\on{1}} \I{[hum]} \emph{or} \I{[hUm]} (\textsc{rn}: \B{hùm} \I{[hUm]}) \emph{vin.} (\emph{a.}) exit, leave, depart. 
\B{Hum ne wrr\-pa!} Get outside!~\LINK{http://forum.learnnavi.org/language-updates/txing-and-hum/}{ln} 
\B{Ftu fì\-tseng ze\-ne hi\-vum.} We gotta get outta here.~\LINK{http://forum.learnnavi.org/language-updates/txing-and-hum/}{ln} 
\B{Po\-lä\-hem Saw\-tu\-te kam zì\-sìt a\-mrr, hum me\-zì\-sìt maw\-krr.} The Sky People arrived five years ago and left two years later.~\LINK{http://naviteri.org/2011/09/miscellaneous-vocabulary/}{b}
\B{Rey\-pay skxir\-ftum\-fa he\-rum.} Blood is coming out of the wound.~\LINK{http://naviteri.org/2013/04/wheres-the-bathroom-and-other-useful-things/}{b} 
(\emph{b. for astronomical bodies setting}) \B{Tsay\-san\-hì ha\-yum ye'\-rìn taw\-ftu.}~\emph{or} \B{Tsay\-san\-hì ha\-yum ye'\-rìn.} Those stars will soon set.~\LINK{http://naviteri.org/2012/01/more-additions-to-the-lexicon/}{b} 
(\ding{214}~\HL{fpxAkIm}{\B{fpxä\-kìm}})

\hi
\B{hum}\textsuperscript{\on{2}} : \HL{kum}{\B{kum}}

% hupx
\hi
\HT{hupx}{\B{hupx}} \I{[hup']} \emph{or} \I{[hUp']} \emph{vtr.} miss, not hit a target. 
\B{Txe\-wìl ye\-rik\-it ko\-lan slä hä\-ngupx.} Txe\-wì aimed at the hexapede but unfortunately missed.~\LINK{http://naviteri.org/2019/12/tengkrr-zisit-leratem-as-the-year-changes/}{b}
(\ding{214}~\HL{takuk}{\B{ta\-kuk}})

% huru
\hi
\HT{huru}{\B{\cp{hu}ru}} \I{["hu.Ru]} \emph{n.} cooking pot. 

% huslew
\hi 
\HT{huslew}{\B{\cp{hu}slew}} \I{["hu.sl$\cdot$Ew]} \emph{vtr., inf.}\on{22} accompany. 
\B{Sra\-ke ni\-vew nga hu\-sli\-vew oe\-ti ftxo\-zä\-ne?} Would you like to accompany me to the celebration?~\LINK{https://naviteri.org/2026/04/tsivola-liu-amip-thirty-two-new-words/}{b}
(\ding{43}~\HL{hu}{\B{hu}}, \HL{salew}{\B{sa\-lew}}) \\
\textsc{Derivations}: \ding{43} 
\on{1}.~\HL{tIhuslew}{\B{tì\-hu\-slew}}, accompaniment.
\on{2}.~\HL{huslewtu}{\B{hu\-slew\-tu}}, accompanier, escort, companion.

% huslewtu
\hi 
\HT{huslewtu}{\B{\cp{hu}slewtu}} \I{["hu.slew.tu]} \emph{n.} accompanier, escort, companion. 
(\ding{43}~\HL{huslew}{\B{hu\-slew}}) 

% huta
\hi 
\HT{huta}{\B{\cp{hu}ta}} \I{["hu.ta]} \emph{adj.} unexpected (usually for positive outcomes). 
\B{tì\-'o\-ngokx a\-hu\-ta}, an unexpected birth.~\LINK{http://naviteri.org/2022/12/neytiriya-waytelem-neytiris-song-cord/}{b\slash a\on{2}} 

\end{multicols}
 
% ===== Section I =====

 \pdfbookmark[0]{I}{I}
\begin{center}
\section*{I \I{[i]}}
\end{center}

\begin{multicols}{2}

% i'en
\hi
\HT{i'en}{\B{\cp{i}'en}} \I{["{}i.PEn]} \emph{n.} \emph{stringed Pandoran instrument}. 
\B{Pe\-o\-txang\-fa nga pam\-tseo si?--Pam\-tse\-o si fa i\-'en.} What instrument do you play?--I play the \emph{i'en}.~\LINK{http://naviteri.org/2011/05/some-miscellaneous-vocabulary/}{b} 
\B{Tew\-ti, nga lu tsul\-fä\-tu i\-'en\-ä.} Wow, you are a master on the \emph{i'en}.~\LINK{http://naviteri.org/2011/01/new-year-new-vocabulary/}{b}

% iknimaya
\hi
\HT{iknimaya}{\B{ikni\cp{ma}ya}} \I{[ik\textcorner.ni."ma.ja]} \emph{n.} Stairway to Heaven (\emph{a rite of passage of the Na'vi}). 
\B{Tsyìl Ik\-ni\-ma\-yat ul\-te tsa\-heyl si ik\-ran\-hu a fì\-'u lu tì\-fme\-tok a ze\-ne fra\-ta\-ron\-yu a\-'e\-wan em\-zi\-va\-'u.} Scaling Iknimaya and bonding with a banshee is a test that every young hunter must pass.~\LINK{http://naviteri.org/2012/01/more-additions-to-the-lexicon/}{b}

% ikran
\hi
\HT{ikran}{\B{\cp{ik}ran}} \I{["{}ik\textcorner.Ran]} \emph{n., fauna} ikran, mountain banshee. 
\B{Oe\-yä ik\-ran sli\-vu nga, tsa\-krr oeng 'aw\-si\-teng mi\-vak\-to.} Be my \emph{ikran} and let's ride together.~\LINK{http://wiki.learnnavi.org/index.php?title=Corpus\#Lemondrop.com}{wiki} 
\B{Nge\-yä ik\-ran\-ìl ngip\-it le\-tam kin fte tsi\-vun kll\-pi\-vä.} Your ikran needs enough open space to be able to land.~\LINK{http://naviteri.org/2015/03/kap-si-ayunil-saylahe-threats-dreams-and-other-things/}{b}
\B{Va'\-rul zä\-nget ik\-ran\-it sn\-eyä pxel ha\-pxì\-tu so\-a\-i\-ä.} Va'ru treats his \emph{ikran} like a member of the family.~\LINK{http://naviteri.org/2012/10/mipa-vospxi-mipa-ayliu-new-words-for-the-new-month/}{b} 
\B{nan\-tang ke na\-mew fpi\-ve' ik\-ran\-ur upxa\-ret fa ngu\-way}, the viperwolf did not want to give messages to the banshee by howling.~\LINK{http://forum.learnnavi.org/language-updates/about-subordination/msg156378/\#msg156378}{ln} 
\B{Ik\-ran\-ìri krr\-a hu tu\-te tsa\-heyl si, ftang li\-vu yrr.} When an \emph{ikran} bonds with a person, it ceases to be wild.~\LINK{http://naviteri.org/2013/01/awvea-posti-zisita-amip-first-post-of-the-new-year/}{b} 
\B{Tsyeyk tswa\-mayon fa ik\-ran.} Jake flew with an \emph{ikran}.~\LINK{http://naviteri.org/2011/09/miscellaneous-vocabulary/}{b} 
(\emph{a.~special usage}) \B{\cp{ik}\-ran \cp{mak}\-to}, the act of riding an ikran.  \\  
\textsc{Derivatives}: \ding{43} \on{1}.~\HL{ikranay}{\B{ik\-ra\-nay}}, forest banshee.
\on{2}.~\HL{kxetsikran}{\B{kxe\-tsik\-ran}}, banshee's tail.

% ikranay
\hi
\HT{ikranay}{\B{ikra\cp{nay}}} \I{[ik\textcorner.Ra."naj]} \emph{n., fauna} forest banshee, lesser banshee (\emph{smaller cousin to the mountain banshee}). 
(\ding{43}~\HL{ikran}{\B{ik\-ran}})

% ikut
\hi
\HT{ikut}{\B{\cp{i}kut}} \I{["{}i.kut\textcorner]} \emph{or} \I{[.kUt\textcorner]} (\textsc{rn}: \B{\cp{i}kùt} \I{["{}i.kUt\textcorner]}) \emph{n.} large pestle (\emph{grinding tool}); meal\hyp{}mashing pole. 

% il
\hi
\HT{il}{\B{il}} \I{[il]} \emph{vin.} bend, \emph{s.t.~straight that bends or hinges at a joint}. 
\B{Txo vul i\-vil nì\-hawng kxì\-ye\-vakx.} If a branch bends too much, it might break.~\LINK{http://naviteri.org/2013/04/wheres-the-bathroom-and-other-useful-things/}{b} 
\B{Me\-tew\-it fì\-swek\-ä ey\-ki\-vil.} Pull the two ends of this bar together.~\LINK{http://naviteri.org/2013/04/wheres-the-bathroom-and-other-useful-things/}{b} \\
\textsc{Derivatives}: \ding{43} \on{1}.~\HL{til}{\B{til}}, joint, hinge. \\  
\on{2}.~\HL{eykil}{\B{ey\-kil}}, bend s.t., pulling things together.

% ilu
\hi 
\HT{ilu}{\B{\cp{i}lu}} \I{["{}i.lu]} \emph{n., fauna} ilu (\emph{a large plesiosaur\hyp{}like sea creature}). 

% ‹ilv›
\hi
‹\HT{ilv}{\B{ilv}}› \I{[il.v]} \emph{inf. in} ‹1› \emph{to form the perfective subjunctive}.
\B{Ul\-te sìl\-pey oe, fì\-'upxa\-ret inan a fì\-'u sil\-vu\-nu ay\-nga\-ru nì\-wotx!} And I hope, you all liked to read this message.~\LINK{http://naviteri.org/2011/09/\%E2\%80\%9Cby-the-way-what-are-you-reading\%E2\%80\%9D/}{b}

% ‹imv›
\hi
‹\HT{imv}{\B{imv}}› \I{[im.v]} \emph{inf. in} ‹1› \emph{to form the past subjunctive}.

% inan
\hi
\HT{inan}{\B{i\cp{nan}}} \I{[i."nan]} \emph{vtr.} (\emph{a.}) read (e.g. the forest), gain knowledge from sensory input. (\emph{with inf. stress changes to first syllable, e.g.}~\B{i\cp{vi}nan} \I{[i."vi.nan]}) 
\B{Tì\-o\-mum\-mì oe\-yä, pol na'\-rìng\-it i\-nan nìl\-tsan.} As far as I know, he reads the forest well.~\LINK{http://naviteri.org/2011/09/\%E2\%80\%9Cby-the-way-what-are-you-reading\%E2\%80\%9D/}{b} 
(\emph{b. colloquially}) `know what someone is about, know someone's ``deal'''.
\B{Pol ze\-ret oe\-ti pxel tu\-te a ke inan pot.} He's treating me like I don't know him.~\LINK{http://naviteri.org/2012/10/mipa-vospxi-mipa-ayliu-new-words-for-the-new-month/}{b} \\
\textsc{Derivatives}: \ding{43} \on{1}.~\HL{tinan}{\B{ti\-nan}}, reading. 
\on{2}.~\HL{ninan}{\B{ni\-nan}}, by reading. 
\on{3}.~\HL{inanfya}{\B{inan\-fya}}, sense (means of perception).

% inanfya
\hi
\HT{inanfya}{\B{i\cp{nan}fya}} \I{[i."nan.fja]} \emph{n.} sense (\emph{means of perception}).
\B{Tì\-ta\-ron\-ìri lä\-ngu tì\-kak\-pam fe\-kum. Kin fkol fra\-inan\-fyat.} I'm sorry to say that deafness is a disadvantage for hunting. You need all your senses.~\LINK{http://naviteri.org/2015/08/aylifyavi-lereyfya-2-cultural-terms-2-and-more/}{b} 
\B{re\-nu ay\-i\-nan\-fya\-yä}, sense paradigm.~\LINK{http://naviteri.org/2012/11/renu-ayinanfyaya-the-senses-paradigm/}{b} 
(\ding{43}~\HL{inan}{\B{inan}}, \HL{fya'o}{\B{fya\-'o}})

% ingyen
\hi
\HT{ingyen}{\B{\cp{ing}yen}} \I{["{}iN.jEn]} \emph{n.} feeling of mystery, feeling of noncomprehension.
\B{Lu oer ing\-yen a Ì\-staw nim lu fì\-txan kum\-a pxìm wä\-pan.} It's a mystery to me why Ìstaw is so shy that he frequently hides.~\LINK{http://naviteri.org/2013/01/awvea-posti-zisita-amip-first-post-of-the-new-year/}{b} \\
\textsc{Derivatives}: \ding{43} \on{1}.~\HL{ningyen}{\B{ning\-yen}}, mysteriously. 
\on{2}.~\HL{ingyentsim}{\B{ing\-yen\-tsim}}, mystery, riddle. 
\on{3}.~\HL{ingyentsyIp}{\B{ing\-yen\-tsyìp}}, trick, clever methodology. 
\on{4}.~\HL{ingyenga'}{\B{ing\-ye\-nga'}}, mysterious, puzzling. 

% ingyentsim
\hi
\HT{ingyentsim}{\B{\cp{ing}yentsim}} \I{["{}iN.jEn.\t{ts}im]} \emph{n.} mystery, riddle, enigma, conundrum.
\B{Fwa po ke zo\-la\-'u lu ing\-yen\-tsim.} It's a mystery that he didn't come.~\LINK{http://naviteri.org/2013/01/awvea-posti-zisita-amip-first-post-of-the-new-year/}{b} 
\B{Sra\-ke tsun nga fì\-ing\-yen\-tsim\-ur tì\-ke\-zin si\-vi?} Can you solve this riddle?~\LINK{http://naviteri.org/2021/04/mipa-ayliu-mipa-sioeykting-new-words-new-explanations/}{b} 
(\ding{43}~\HL{ingyen}{\B{ing\-yen}}, \HL{tsim}{\B{tsim}})

% ingyentsyìp
\hi
\HT{ingyentsyIp}{\B{\cp{ing}yentsyìp}} \I{["{}iN.jEn.\t{ts}jIp\textcorner]} \emph{n.} trick, slight of hand, clever methodology, special methodology.
\B{Pol\-txe po san lu ing\-yen\-tsyìp a\-zey.} He said there's a special trick to it.~\LINK{http://naviteri.org/2013/01/awvea-posti-zisita-amip-first-post-of-the-new-year/}{b} 
(\ding{43}~\HL{ingyen}{\B{ing\-yen}})

% ingyenga'
\hi
\HT{ingenga'}{\B{\cp{ing}yenga'}} \I{["{}iN.jE.NaP]} \emph{adj.} mysterious, puzzling, enigmatic.
\B{Pe\-yä ay\-lì\-'u a\-ing\-ye\-nga' lo\-lu sngum\-tsim ay\-oe\-ru nì\-wotx.} His mysterious words worried us all.~\LINK{http://naviteri.org/2013/01/awvea-posti-zisita-amip-first-post-of-the-new-year/}{b} 
(\ding{43}~\HL{ingyen}{\B{ing\-yen}})

% io
\hi
\HT{io}{\B{\cp{i}o}} \I{["{}i.o]} \emph{adp.} above, over. 
\B{Oe tsway\-on io Ay\-vit\-ra\-yä Ra\-mu\-nong.} I fly over the Tree of Souls.~\LINK{http://forum.learnnavi.org/language-updates/over-under-and-quickly/}{ln} 
(\ding{214}~\HL{Ao}{\B{äo}}) \\
\textsc{Derivatives}: \ding{43} \on{1}.~\HL{ionar}{\B{io\-nar}}, banshee rider visor.

% ionar
\hi
\HT{ionar}{\B{\cp{i}onar}} \I{["{}i.o.naR]} \emph{n.} banshee rider visor. 
(\ding{43}~\HL{io}{\B{io}}, \HL{nari}{\B{na\-ri}}; \HL{renten}{\B{ren\-ten}})

% ioang
\hi
\HT{ioang}{\B{i\cp{o}ang}} \I{[i."{}o.aN]} \emph{n., fauna} animal, beast. 
\B{Fì\-i\-o\-ang lu tsuk\-yom.} This animal is edible.~\LINK{http://naviteri.org/2011/03/\%E2\%80\%9Creceptive-ability\%E2\%80\%9D-and-hesitation/}{b} 
\B{Na'\-rìng\-ìl nga' pxay\-a ioang\-it.} The forest contains many animals.~\LINK{http://naviteri.org/2011/05/weather-part-2-and-a-bit-more-2/}{b} 
\B{Ioang\-ìri, tu\-te\-ri, tsun fko si\-var lì\-'ut alu tsmìm.} One can use the word `track' for animals or people.~\LINK{http://naviteri.org/2011/09/\%E2\%80\%9Cby-the-way-what-are-you-reading\%E2\%80\%9D/comment-page-1/\#comment-1118}{b} \\
\textsc{Derivatives}: \ding{43} 
\on{1}.~\HL{payoang}{\B{pay\-oang}}, fish.   
\on{2}.~\HL{hI'ang}{\B{hì\-'ang}}, insect.   
\on{3}.~\HL{yayo}{\B{ya\-yo}}, bird.  
\on{4}.~\HL{tarnioang}{\B{tar\-ni\-oang}}, predator animal. 
\on{5}.~\HL{talioang}{\B{ta\-li\-oang}}, sturmbeest. 
\on{6}.~\HL{yomioang}{\B{yom\-ioang}}, chalice plant. 
\on{7}.~\HL{yomhI'ang}{\B{yom\-hì\-'ang}}, dakteron. 
\on{8}.~\HL{tireaioang}{\B{Ti\-re\-a\-ioang}}, Spirit Animal.
\on{9}.~\HL{'eylyong}{\B{'eyl\-yong}}, pet.

% ioi / ioi si / ioi säpi
\hi
\HT{ioi}{\B{i\cp{o}i}} \I{[i."{}o.i]} \textsc{i}. \emph{n.} (item of) adornment or ceremonial apparel.
\B{Nì\-lun ay\-ioi a\-'eoio ay\-eyk\-tan\-ä lu lor fra\-to.} Of course the ceremonial wardrobes of the leaders are the most beautiful.~\LINK{http://naviteri.org/2011/08/new-vocabulary-clothing/}{b} 
\B{Saw\-no'\-ha ioit ko\-la\-ne\-i\-om oel u\-van\-fa!} I got (won) the prized piece of jewelry in the game!~\LINK{http://naviteri.org/2012/06/spring-vocabulary-part-3/}{b} (\ding{43}~\HL{pxen}{\B{pxen}}) \\
\textsc{ii}. \HT{ioi si}{\B{i\cp{o}i si}} \I{[i."{}o.i si]} \emph{vin.} adorn s.o. (with, \B{fa}), put on (adornment). 
\B{Se\-vin\-a tsa\-'e\-ve\-ru a\-hì\-'i me\-sa'\-sem ioi so\-li fa mik\-tsang.} The parents adorned that pretty little girl with earrings.~\LINK{http://naviteri.org/2011/08/new-vocabulary-clothing/}{b} \\
\textsc{iii}. \HT{ioi sApi}{\B{i\cp{o}i sä\cp{pi}}} \I{[i."{}o.i s\ae."pi]} \emph{or coll.} \I{[spi]} \emph{vin.} adorn oneself (\HL{fa}{\B{fa}}, with), put on (adornment). 
\B{Fo\-ri maw\-krr\-a fa ren\-ten ioi sä\-po\-li ho\-lum.} After they put on their goggles, they left.~\LINK{http://naviteri.org/2011/08/new-vocabulary-clothing/}{b}
(\ding{214}~\HL{'aku}{\B{'a\-ku}}, \ding{43}~\HL{yemstokx}{\B{yem\-stokx}})

% irayo / irayo si
\hi
\HT{irayo}{\B{i\cp{ra}yo}} \I{[i."Ra.jo]} \textsc{i}.~\emph{n.} thank; (\emph{a. idiom, intj.}) `Thanks,' `Thank you.' 
\B{Txa\-na ira\-yo.} Many thanks.~\LINK{http://naviteri.org/2011/12/ulte-yora\%E2\%80\%99tu-leiu-and-the-winner-is/}{b} \\
\textsc{ii}. \HT{irayo si}{\B{i\cp{ra}yo si}} \I{[i."Ra.jo si]} \emph{vin.} thank s.o. (\emph{in the dative}) for s.t. (\emph{with the topic}). 
(\emph{a.}) \B{Tsa\-krr\-vay, ay\-nge\-yä tìm\-wey\-pey\-ri ira\-yo sei\-yi oe.} Until then, I thank you all for your patience.~\LINK{http://forum.learnnavi.org/language-updates/trr-rrtaya/}{ln} 
\B{Pxe\-fe\-yä tì\-tstun\-wi\-ri aw\-nga ira\-yo si\-vi ko!} Let us all thank those three for their kindness.~\LINK{http://naviteri.org/2012/07/tskxekengtsyip-a-mikyunfpi-a-little-listening-exercise/}{b}
\B{Fu\-ri\-a ngal oe\-yä 'u\-pxa\-ret ay\-su\-te\-ru fpo\-le', nga\-ru ira\-yo sei\-yi oe nì\-txan!} I thank you very much for sending my message to people.~\LINK{http://wiki.learnnavi.org/index.php?title=Canon\#Extracts_from_various_emails}{wiki}
(\emph{b. word order can be broken}) \B{Oel ay\-nga\-ti ka\-me\-i\-e, ma oe\-yä ey\-lan, ul\-te ay\-nga\-ru sei\-yi ira\-yo.} I \textsc{see} you all, my friends, and thank you.~\LINK{http://forum.learnnavi.org/news-announcements/a-response-from-paul-frommer!/}{ln}

% ‹irv›
\hi
‹\HT{irv}{\B{irv}}› \I{[iR.v]} \emph{inf. in} ‹1› \emph{to form the imperfective subjunctive}.

% -it
\hi
\HT{-it}{\B{-it}} \I{[.it\textcorner]} \emph{suf. to mark the patientive case of a n. ending in a consonant, pseudovowel or diphthong}. 
(\ding{43}~\HL{-t}{\B{-t}}, \HL{-ti}{\B{-ti}})

% ‹iv›
\hi
‹\HT{iv}{\B{iv}}› \I{[i.v]} \emph{inf. in} ‹1›. (\emph{a.}) \emph{to form the subjunctive}. 
(\emph{b.}) \emph{required after} \B{fte}, \B{fteke} \emph{and all modal verbs}.
(\emph{c.}) \emph{optionally to indicate the imperative}.

% ‹iyev›
\hi
‹\HT{iyev}{\B{iyev}}› \I{[i.jE.v]} \emph{inf. in} ‹1› \emph{to form the (approximate) future subjunctive}. 
(\ding{43}~\HL{Iyev}{‹\B{ìyev}›})

\end{multicols}

% ===== Section ì =====

 \pdfbookmark[0]{SHORT I}{SHORT I}
\begin{center}
\section*{Ì \I{[I]}}
\end{center}

\begin{multicols}{2}

% ìì
\hi
\HT{II}{\B{ìì}} \I{[I:]} \emph{intj.} um, er (\emph{hesitation marker}). \B{Lu oe\-ru … ìì … tì\-ngä\-zìk a\-hì\-'i.} I have … um … a slight problem.~\LINK{http://naviteri.org/2011/03/\%E2\%80\%9Creceptive-ability\%E2\%80\%9D-and-hesitation/}{b}

% -ìl
\hi
\HT{Il}{\B{-ìl}} \I{[.Il]} \emph{suf. to mark the agentive case of a n. ending in a consonant, pseudovowel or diphthong}. 
(\ding{43}~\HL{l}{\B{-l}})

% ìlä
\hi
\HT{IlA}{\B{ì\cp{lä}}} \I{[I."l\ae]} \emph{or} \B{\cp{ì}lä} \I{["{}I.l\ae]} \emph{adp.}\texttt{+} 
(\emph{a.}) by, via, following; in. \B{Ay\-fo so\-lop ìlä hil\-van fa uran.} They traveled along the river by boat.~\LINK{http://naviteri.org/2011/07/txantsana-ultxa-mi-siatll-great-meeting-in-seattle/}{b}
\B{Po\-ti fwey\-ko\-li ay\-oel ìlä rong.} We let him slide through the tunnel.~\LINK{http://naviteri.org/2013/04/wheres-the-bathroom-and-other-useful-things/}{b}
\B{Na'\-vi ìlä ho\-'on kll\-kxo\-lem teng\-krr re\-rol.} The People were standing in a circle, singing.~\LINK{http://naviteri.org/2013/09/on-si-salewfya-shapes-and-directions/}{b}
(\emph{b.}) according to. \B{Ìlä Fey\-ral, mun\-txa so\-li Ra\-lu sì Ne\-wey nì\-wan me\-srr\-am.} According to Peyral, Ralu and Newey were secretly married the day before yesterday.~\LINK{http://naviteri.org/2012/07/meetings-waterfalls-and-more/}{b} (\ding{214}~\HL{wA}{\B{wä}})
(\emph{c. with} \HL{latem}{\B{la\-tem}}) depend on. \B{Tì\-flä la\-tem ìlä seyng\-a ftxey fkol sä\-nu\-met li\-vek fu\-ke.} Success depends on whether or not one follows instructions.~\LINK{http://naviteri.org/2012/07/meetings-waterfalls-and-more/}{b} \\
\textsc{Derivatives}: \ding{43} \on{1}.~\HL{kxamlA}{\B{kxam\-lä}}, through, via the middle of.

% ìley
\hi
\HT{Iley}{\B{ì\cp{ley}}} \I{[I."lEj]} \emph{intj. sound of a war cry}.

% ‹ìlm›
\hi
‹\HT{Ilm}{\B{ìlm}}› \I{[Il.m]} \emph{inf.~in} ‹1› \emph{to indicate perfective approximate past}.

% ìlva
\hi
\HT{Ilva}{\B{\cp{ìl}va}} \I{["{}Il.va]} \emph{n.} flake, drop, chip. \\  
\textsc{Derivatives}: \ding{43} \on{1}.~\HL{herwIva}{\B{her\-wì\-va}}, snowflake. 
\on{2}.~\HL{payIva}{\B{pay\-ìva}}, drop of water. 
\on{3}.~\HL{txepIva}{\B{txe\-pì\-va}}, ash, cinder. 
\on{4}.~\HL{tompIva}{\B{tom\-pì\-va}}, raindrop.

% ‹ìly›
\hi
‹\HT{Ily}{\B{ìly}}› \I{[Il.j]} \emph{inf.~in} ‹1› \emph{to indicate perfective approximate future}.

% ‹ìm›
\hi
‹\HT{Im}{\B{ìm}}› \I{[I.m]} \emph{inf.~in} ‹1› \emph{to indicate approximate past}.

% ìpxa
\hi
\HT{Ipxa}{\B{ì\cp{pxa}}} \I{[I."p'a]} \emph{n.}~\ding{96} fern. 
\B{Fnu, fnu, wä\-pan wä ku | Uo tsa\-ìpxa.} Quiet, quiet, hide yourself from harm | Behind that fern.~\LINK{http://forum.learnnavi.org/language-updates/learnings-from-the-siatlla-song-translations/msg473566/\#msg473566}{ln}

% -ìri
\hi
\HT{-Iri}{\B{-ìri}} \I{[.I.ri]} \emph{suf. to mark the topical case of a n. ending in a consonant or pseudovowel}.
(\ding{43}~\HL{-ri}{\B{\mbox{-ri}}})

% ‹ìrm›
\hi
‹\HT{Irm}{\B{ìrm}}› \I{[IR.m]} \emph{inf.~in} ‹1› \emph{to indicate imperfective approximate past}.

% ‹ìry›
\hi
‹\HT{Iry}{\B{ìry}}› \I{[IR.m]} \emph{inf.~in} ‹1› \emph{to indicate imperfective approximate future}.

% Ìstaw
\hi
\B{Ì\cp{staw}} \I{[I."staw]} \emph{n. male name}. 
\B{Fì\-po lu Ìstaw.} \emph{or} \B{Fì\-por syaw fko Ìstaw.} This is Ìstaw.~\LINK{http://naviteri.org/2010/09/getting-to-know-you-part-1/}{b}
\B{Wo\-leyn Ìstawl yey\-fyat mì hll\-te.} Ìstaw drew a line on the ground.~\LINK{http://naviteri.org/2013/09/on-si-salewfya-shapes-and-directions/}{b}

% ‹ìsy›
\hi
‹\HT{Isy}{\B{ìsy}}› \I{[I.sj]} \emph{inf.~in} ‹1› \emph{to indicate determinational approximate future}.
\B{Pì\-syeng oe ngar.} I intend to tell you.~\LINK{http://forum.learnnavi.org/language-updates/a-collection/}{ln}
\B{Ta\-fral ke lì\-syek oel nge\-yä ke\-ye\-'ung\-it.} Therefore I will not heed your insanity.~\LINK{http://forum.learnnavi.org/language-updates/ltasygt-with-ke-and-nga/}{g\slash ln}

% ‹ìy›
\hi
‹\HT{Iy}{\B{ìy}}› \I{[I.j]} \emph{inf.~in} ‹1› \emph{to indicate approximate future}. 
\B{Maw hì\-krr ay\-oe tì\-yä\-txaw.} We'll be right back after this break.~\LINK{http://wiki.learnnavi.org/index.php?title=Corpus\#Good_Morning_America}{wiki}
\B{ye'\-rìn 'ì\-yi\-'a sä\-nu\-me a tsa\-ri kll\-fro' oe}, soon the teaching I am responsible for will end.~\LINK{http://forum.learnnavi.org/language-updates/trr-rrtaya/}{ln}

% ‹ìyev›
\hi
‹\HT{Iyev}{\B{ìyev}}› \I{[I.jE.v]} \emph{inf.~in} ‹1› \emph{to indicate subjunctive future}.
\B{Sìl\-pey oe, fra\-po tsì\-ye\-vun fì\-tse\-nge ri\-vun 'uot le\-sar.} I hope that everyone may find something useful here.~\LINK{http://naviteri.org/2010/06/first-post/}{b}
\B{Lu oe\-ru sngum a sa\-ron\-yu ke tì\-ye\-vä\-txaw.} I'm worried that the hunters will not return (soon, as expected).~\LINK{http://naviteri.org/2011/01/new-year-new-vocabulary/}{b}
(\emph{a. idiom}) \B{Kì\-ye\-va\-me.} \textsc{See} you soon.
(\ding{43}~\HL{iyev}{‹\B{iyev}›})

\end{multicols}

% ===== section K =====

 \pdfbookmark[0]{K}{K}
\begin{center}
\section*{K \I{[k]} (\emph{KeK})}
\end{center}

\begin{multicols}{2}

% ka
\hi
\HT{ka}{\B{ka}}\textsuperscript{\on{1}} \I{[ka]} \emph{adp.} (\emph{a. locative}) across. 
\B{Txur\-tel\-mì fo fta so\-li fte\-ke ka tsyokx fwi\-vi.} They tied a knot in the rope so it wouldn't slip through their hands.~\LINK{http://naviteri.org/2013/04/wheres-the-bathroom-and-other-useful-things/}{b}
(\emph{b. temporal}) 
\B{Nge\-yä tì\-eyk\-tan\-ìri, tì\-slan\-ìri sì tsran\-ten fra\-to a tì\-'ey\-lan\-ìri a ka ay\-zì\-sìt nì\-wotx, sei\-yi oe ira\-yo nì\-txan.} Thank you so much for your leadership, your support, and most importantly your friendship throughout the years.~\LINK{http://naviteri.org/2012/03/spring-vocabulary-part-1/}{b}
(\emph{to disambiguate durative interpretation together with indefinite} \ding{43}~\B{-o}) 
\B{Pol lì'\-fya\-ti le\-Na'\-vi fto\-lia ka trr\-o (a\-'aw).} He studied Na'vi for a day.~\LINK{http://naviteri.org/2014/09/stxeli-alor-the-text/comment-page-1/\#comment-10619}{b}
(\emph{c. idiom}) \B{ka wotx}, generally, for the most part. 
\B{Hu\-fwa ro\-lun oel 'a\-'aw\-a kxey\-ey\-ti, fì\-tì\-kang\-kem\-vi lu txan\-tsan ka wotx.} Although I found a few errors, this piece of work is generally excellent.~\LINK{http://naviteri.org/2012/07/fivospxiya-aylifyavi-amip-this-months-new-expressions-pt-2/}{b} \\
\textsc{Derivatives}: \ding{43} \on{1}.~\HL{krrka}{\B{krr\-ka}}, during.
\on{2}.~\HL{katIng}{\B{ka\-tìng}}, distribute.
\on{3}.~\HL{kahena}{\B{ka\-he\-na}}, transport.

\hi
\B{ka}\textsuperscript{\on{2}} : \HL{kxa}{\B{kxa}}

% Ka'ani
\hi
\B{Ka'ani} \I{[ka.Pa.ni]} \emph{n.} \emph{male name}.~\LINK{}{dh}

% kafi / si
\hi 
\HT{kafi}{\B{\cp{ka}fi}} \I{["ka.fi]} \textsc{i}. \emph{n.} sail. \\
\textsc{ii}. \B{kafi si} \I{["ka.fi si]} \emph{vin.} to sail, move by means of a sail (\emph{lit. and fig.}).
\B{Rìk a\-ukxo mì hu\-fwe\-tsyìp ka\-fi sar\-mi.} The dry leaf was sailing in the breeze.~\LINK{http://naviteri.org/2023/12/mipa-zisit-ayliu-amip-new-year-new-words/}{b} \\
\textsc{Derivatives}: \ding{43} \on{1}.~\HL{kaifuran}{\B{ka\-fi\-u\-ran}}, sailboat. 

% kafiuran
\hi 
\HT{kafiuran}{\B{\cp{ka}fiuran}} \I{["ka.fi.u.Ran]} \emph{n.} sailboat. 
(\ding{43}~\HL{kafi}{\B{kafi}}, \HL{uran}{\B{uran}})

% kahena
\hi 
\HT{kahena}{\B{ka\cp{he}na}} \I{[ka."h$\cdot$E.n$\cdot$a]} \emph{vtr., inf.}\on{23} transport. 
\B{Fwa ka\-he\-na fì\-uran\-it a\-tsawl ftu tsray oe\-yä ne pum nge\-yä la\-yu ngä\-zìk.} It's going to be difficult to transport this large boat from my village to yours.~\LINK{}{} 
(\ding{43}~\HL{ka}{\B{ka}}, \HL{hena}{\B{he\-na}}) \\
\textsc{Derivatives}: \ding{43} \on{1}.~\HL{tIkanhena}{\B{tì\-ka\-hena}}, transportation (abstract). 
\on{2}.~\HL{sAkahena}{\B{sä\-ka\-he\-na}}, means of transport, vehicle.

% kakan
\hi 
\HT{kakan}{\B{\cp{ka}kan}} \I{["ka.kan]} \emph{adj.} rough (\emph{behavior; used for people and things}). 
\B{Ka\-kan\-a ay\-lì\-'u\-ri a pol\-txe oel nì\-sti, tsap\-'a\-lu\-te.} I apologize for the rough words that I spoke in anger.~\LINK{http://naviteri.org/2021/09/aawa-ayliu-amip-a-few-new-words/}{b} 
(\ding{214}~\HL{flrr}{\B{flrr}}; \ding{43}~\HL{ekxtxu}{\B{ekx\-txu}}) \\
\textsc{Derivatives}: \on{1}.~\HL{nIkakan}{\B{nì\-ka\-kan}}, roughly.

% kakmokri
\hi 
\HT{kakmokri}{\B{kak\cp{mok}ri}} \I{[kak\textcorner."mok\textcorner.Ri]} \emph{adj.} mute.
\B{Txo po lu kak\-mok\-ri, fya\-pe sä\-fpìl\-yewn?} If he's mute, how does he communicate?~\LINK{https://naviteri.org/2024/02/mipa-ayliu-si-aylifyavi-niul-more-new-words-and-expressions/}{b} 
(\ding{43}~\HL{mokri}{\B{mokri}}) \\
\textsc{Derivatives}: \ding{43} \on{1}.~\HL{tIkakmokri}{\B{tì\-kak\-mok\-ri}}, muteness.
\on{2}.~\HL{nIkakmokri}{\B{nì\-kak\-mok\-ri}}, mutely.

% kakpam
\hi
\HT{kakpam}{\B{kak\cp{pam}}} \I{[kak\textcorner."pam]} \emph{adj.} deaf.
\B{Hì\-krr\-o me\-fo kak\-pam lar\-mu maw\-krr\-a pxo\-lor kun\-sìp.} The two of them were deaf for a short time after the gunship exploded.~\LINK{http://naviteri.org/2014/05/mipa-ayliu-mipa-aysafpil-new-words-new-ideas/}{b} 
(\ding{43}~\HL{pam}{\B{pam}}) \\
\textsc{Derivatives}: \ding{43} \on{1}.~\HL{tIkakpam}{\B{tì\-kak\-pam}}, deafness.

% kakrel
\hi
\HT{kakrel}{\B{kak\cp{rel}}} \I{[kak\textcorner."REl]} \emph{adj.} blind. 
(\ding{43}~\HL{rel}{\B{rel}}) \\
\textsc{Derivatives}: \ding{43} \on{1}.~\HL{tIkakrel}{\B{tì\-kak\-rel}}, blindness.

% kakzir
\hi 
\HT{kakzir}{\B{kak\cp{zir}}} \I{[kak\textcorner."ziR]} \emph{adj.} (\emph{a. physically}) numb. 
\B{Po\-ri pxun kak\-zir la\-tsu.} His arm must be\slash appears to be numb.~\LINK{https://naviteri.org/2025/11/kaltxi-nimun-hello-again/}{b}
(\emph{b. emotionally}) numb. 
\B{Te\-ri kxitx pe\-yä stawm a krr, oe kak\-zir slo\-lu nì\-wotx.} When I heard about her death, I became completely numb.~\LINK{https://naviteri.org/2025/11/kaltxi-nimun-hello-again/}{b}
(\ding{43}~\HL{zir}{\B{zir}}) \\
\textsc{Derivations}: \ding{43} 
\on{1}.~\HL{tIkakzir}{\B{tì\-kak\-zir}}, numbness.

% kali'weya
\hi
\HT{kali'weya}{\B{kali'\cp{wey}a}} \I{[ka.liP."wEj.a]} \emph{n., fauna} species of poisonous arachnid, \emph{scorpiosista virosae}. \\
\textsc{Derivatives}: \ding{43} \on{1}.~\HL{kalweyaveng}{\B{kal\-we\-ya\-veng}}, `son of a bitch'.

% kalin
\hi
\HT{kalin}{\B{ka\cp{lin}}} \I{[ka."lin]} \emph{adj.} sweet. 
\B{Fì\-na\-er\-ìri sur fkan oe\-ru ka\-lin.} This drink tastes sweet to me.~\LINK{http://naviteri.org/2012/11/renu-ayinanfyaya-the-senses-paradigm/}{b} \\
\textsc{Derivatives}: \ding{43} \on{1}.~\HL{kalintu}{\B{ka\-lin\-tu}}, sweet person.

% kalintu
\hi 
\HT{kalintu}{\B{ka\cp{lin}tu}} \I{[ka."lin.tu]} \emph{n.} sweet person. 
\B{Lu nga ka\-lin\-tu nì\-ngay.} You're really a sweet person.~\LINK{https://naviteri.org/2024/09/pxevola-liu-amip-two-dozen-new-words/}{b} 
(\ding{43}~\HL{kalin}{\B{ka\-lin}}) 

% kaltxì
\hi
\HT{kaltxI}{\B{kal\cp{txì}}} \I{[kal."t'I]} \textsc{i}. \emph{idiom} Hello! Greetings! 
\B{Kal\-txì!} Hello!~\LINK{http://wiki.learnnavi.org/index.php?title=Corpus\#Vanity_Fair}{wiki} 
\B{Kal\-txì ay\-nga\-ru}, I greet you all.~\LINK{http://naviteri.org/2013/03/whoever-whatever-whenever/}{b} 
\B{Kal\-txì ta Ko\-pen\-han}, Hi from Copenhagen.~\LINK{http://naviteri.org/2010/09/kaltxi-ta-kopenhan-hi-from-copenhagen/}{b} \\
\textsc{ii}. \HT{kaltxI si}{\B{kal\cp{txì} si}} \I{[kal."t'I si]} \emph{vin.} say hello, greet. 
\B{Kal\-txì si\-vi, Ìstaw.} Say hello, Ìstaw.~\LINK{http://naviteri.org/2010/09/getting-to-know-you-part-1/}{b} 
\B{Za\-'u kal\-txì si ko!} Come and say hello!~\LINK{http://naviteri.org/2010/09/getting-to-know-you-part-3/}{b} 
\B{Lo\-ak kal\-txì so\-li oer.} Loak greeted me.~\LINK{http://naviteri.org/2010/09/getting-to-know-you-part-1/comment-page-1/\#comment-294}{b}
(\ding{43}~\HL{kxI}{\B{kxì}})

% kalweyaveng
\hi 
\HT{kalweyaveng}{\B{kal\cp{wey}aveng}} \I{[kal."wEj.a.vEN]} \emph{n.} `son of a bitch', \emph{lit.}: `child of a poisonous arachnid' (\emph{insult}). 
\B{Txo sran, tsa\-krr fwa syaw por ``kal\-wey\-a\-veng'' lu mu\-iä nì\-wotx!} If yes, then to call them `son of a bitch' is entirely justified.~\LINK{http://naviteri.org/2019/09/choice-statements-vs-choice-questions-and-some-insults/\#comment-30388}{b}
(\ding{43}~\HL{kali'weya}{\B{ka\-li'\-wey\-a}}, \HL{'eveng}{\B{'e\-veng}})

% kam 
\hi
\HT{kam}{\B{kam}}\textsuperscript{\on{1}} \I{[kam]} \emph{adp.} ago. 
\B{Tskot sngo\-lä\-'i po si\-var 'a\-'aw\-a trr\-kam.} He started to use the bow several days ago.~\LINK{http://naviteri.org/2011/09/miscellaneous-vocabulary/}{b} 
\B{Po\-lä\-hem Saw\-tu\-te kam zì\-sìt a\-mrr, hum me\-zì\-sìt maw\-krr.} The Sky People arrived five years ago and left two years later.~\LINK{http://naviteri.org/2011/09/miscellaneous-vocabulary/}{b}
(\ding{214}~\HL{kay}{\B{kay}}) 

\hi
\HT{kam}{\B{kam}}\textsuperscript{\on{2}} : \HL{kxam}{\B{kxam}}

% kame
\hi
\HT{kame}{\B{\cp{ka}me}} \I{["ka.mE]} \emph{vtr.} see (\emph{spiritual sense}), acknowledge, see into, understand, connect. (\emph{a}) \B{Oel nga\-ti ka\-me\-i\-e.} I \textsc{see} you.~\LINK{http://wiki.learnnavi.org/index.php?title=Corpus\#UGO_Movie_Blog}{wiki} 
(\emph{b, idiom}) \B{Kì\-ye\-vame.} Good-bye. \textsc{See} you soon.~\LINK{http://forum.learnnavi.org/news-announcements/a-response-from-paul-frommer!/}{ln}
(\ding{43}~\HL{tse'a}{\B{tse\-'a}})

% Kamun
\hi
\B{\cp{Ka}mun} \I{["ka.mun]}~\emph{or}~\I{[.mUn]} \emph{n. male name}. \B{Pol\-txe Ka\-mun san oe za\-sya\-'u sìk.} Kamun said, ``I will come.''~\LINK{http://naviteri.org/2010/08/a-na\%E2\%80\%99vi-alphabet/comment-page-1/\#comment-191}{b}
(\emph{shortened with} \B{-tsyìp}) \B{Kam\-tsyìp\-ìl wu\-tsot ye\-rom.} Little Kamun is having dinner.~\LINK{http://naviteri.org/2010/07/diminutives-conversational-expressions/}{b}

% kan
\hi
\HT{kan}{\B{kan}} \I{[kan]} (\emph{a. vtr.}) aim (s.t. at\slash against s.t.\slash s.o.; \emph{used with} \ding{43} \HL{ne}{\B{ne}} \emph{or} \HL{wA}{\B{w\"a}}).
\B{Ney\-ti\-ril tsko\-ti ke\-ran ne ye\-rik.} Neytiri is aiming her bow at a hexapede.~\LINK{http://naviteri.org/2015/11/vomuna-liu-amip-ten-new-words/comment-page-1/\#comment-22792}{g\slash b}
\B{Pol tsko\-ti ko\-lan wä ku\-tu.} He raised his bow against the enemy.~\LINK{http://naviteri.org/2015/11/vomuna-liu-amip-ten-new-words/comment-page-1/\#comment-22792}{g\slash b}
(\emph{b. vtr.}) aim (at). 
\B{Ta\-ron\-yul le\-hì\-pey kan smar\-it nìl\-ke\-ftang slä ke ta\-kuk kaw\-krr.} A hesitant hunter will aim at a prey forever but never hit it.~\LINK{http://naviteri.org/2015/11/vomuna-liu-amip-ten-new-words/}{b}
(\emph{c. vtr., modal}) intend (\emph{either with} ‹\B{iv}› \emph{or} \B{futa} \emph{and} ‹\B{iv}›). 
\B{Oe kan ki\-vä.} I intend to go.~\LINK{http://naviteri.org/2015/11/vomuna-liu-amip-ten-new-words/comment-page-1/\#comment-22792}{b} 
\B{Oel kan fu\-ta po ki\-vä.} I intend him to go.~\LINK{http://naviteri.org/2015/11/vomuna-liu-amip-ten-new-words/comment-page-1/\#comment-22792}{b} 
(\emph{d. conv., self\hyp{}correction as} \B{ko\cp{lan}}) ``I mean …, rather''.
\B{Oe\-ri tsyokx tì\-sraw si … ko\-lan ze\-kwä.} My hand hurts--I mean, my finger.~\LINK{http://naviteri.org/2014/05/mipa-ayliu-mipa-aysafpil-new-words-new-ideas/}{b}
\B{Fo ko\-lä tsa\-tseng fte ti\-va\-ron ye\-ri … ìì … ko\-lan ta\-li\-o\-ang\-it.} They went there to hunt hexa … um … I mean sturmbeast.~\LINK{http://naviteri.org/2014/05/mipa-ayliu-mipa-aysafpil-new-words-new-ideas/}{b} 
(\emph{e. proverb}) \B{Fwa kan ke tam; ze\-ne swi\-zaw\-it li\-vo\-nu.} Intent is not enough; it's action that counts. [lit.: To aim is not enough; one must release the arrow.]~\LINK{http://naviteri.org/2013/02/vospxi-ayol-posti-apup-short-post-for-a-short-month/}{b} \\
\textsc{Derivatives}: \ding{43} \on{1}.~\HL{tIkan}{\B{tì\-kan}}, aim, goal, purpose, target. 
\on{2}.~\HL{kan'In}{\B{kan\-'ìn}}, focus on, specialize in.  
\on{3}.~\HL{kanfpIl}{\B{kan\-fpìl}}, concentrate. 
\on{4}.~\HL{kantseng}{\B{kan\-tseng}}, destination.

% kan'ìn
\hi
\HT{kan'In}{\B{\cp{kan}'ìn}} \I{["kan.PIn]} \emph{vtr.} focus on, specialize in, be particularly interested in.  \B{En\-tul kan\-'ìn tì\-wu\-sem\-it.} Entu specializes in fighting.~\LINK{http://naviteri.org/2010/09/getting-to-know-you-part-1/}{b}
(\ding{43}~\HL{kan}{\B{kan}}, \HL{'In}{\B{'ìn}})

% kanfpìl
\hi
\HT{kanfpIl}{\B{\cp{kan}fpìl}} \I{["kan.fpIl]} \emph{vin.} concentrate, focus one's attention. \B{Fu\-ri\-a sne\-yä tsko\-ti ngop po kan\-fpìl.} He's concentrating on making his bow.~\LINK{http://naviteri.org/2011/10/more-vocabulary-a-bit-of-grammar/}{b}
\B{Txo new nga tsli\-vam, ze\-ne ki\-van\-fpìl.} If you want to understand, you have to concentrate.~\LINK{http://naviteri.org/2011/10/more-vocabulary-a-bit-of-grammar/}{b}
(\ding{43}~\HL{kan}{\B{kan}}, \HL{fpIl}{\B{fpìl}})

% kanom
\hi
\HT{kanom}{\B{\cp{ka}nom}} \I{["ka.nom]} \emph{vtr.} acquire, get.
\B{Oe\-yä tsmu\-kan\-ìl mip\-a tsko\-ti kì\-ma\-ne\-i\-om.} My brother just got a new bow, I'm happy to say.~\LINK{http://naviteri.org/2012/01/more-additions-to-the-lexicon/}{b} \\
\textsc{Derivatives}: \ding{43} \on{1}.~\HL{sAkanom}{\B{sä\-ka\-nom}}, s.t.~acquired, an acquisition, a possession.
\on{2}.~\HL{tsukanom}{tsu\-ka\-nom}, available, obtainable.

% kantseng
\hi 
\HT{kantseng}{\B{\cp{kan}tseng}} \I{["kan.\t{ts}EN]} \emph{n.} destination. 
\B{Nge\-yä fì\-tì\-sop\-ìri pe\-han\-tseng?} This journey of yours--what's its destination?~\LINK{http://naviteri.org/2019/08/aawa-lifya-si-lifyavi-amip-a-few-new-words-and-expressions/}{b} 
(\ding{43}~\HL{kan}{\B{kan}}, \HL{tseng}{\B{tseng(e)}})

% kanu
\hi
\HT{kanu}{\B{\cp{ka}nu}} \I{["ka.nu]} \emph{adj.} smart, intelligent (\emph{of a person}). 
\B{Lam oer fwa po lu ka\-nu nì\-txan.} It seems to me that he is very smart.~\LINK{http://forum.learnnavi.org/language-updates/txelanit-hivawl/}{ln} \\
\textsc{Derivatives}: \ding{43} \on{1}.~\HL{tIkanu}{\B{tì\-ka\-nu}}, intelligence.  
\on{2}.~\HL{nIkanu}{\B{nì\-ka\-nu}}, intelligently, in a smart way. 

% kangay / si
\hi
\HT{kangay}{\B{ka\cp{ngay}}} \I{[ka."Naj]} \textsc{i}. \emph{adj.} valid. \\
\textsc{ii}. \HT{kangay si}{\B{ka\cp{ngay} si}} \I{[ka."Naj si]} \emph{vin.} confirm, validate. 
\B{Ru\-txe, ay\-nga\-ri fì\-kem fe\-yä 'e\-'a\-la to\-pur ka\-ngay sì\-yi nì\-'aw.} Please, this will only confirm their worst fears about you.~\LINK{http://forum.learnnavi.org/language-updates/two-words-in-use-eal-and-kangay-si/msg564284/\#msg564284}{ln}
\B{Fì\-'u\-pxa\-re ka\-ngay si fu\-ra po za\-ya\-'u.} This message confirms that he will come.~\LINK{http://naviteri.org/2020/04/aawa-u-amip-a-few-new-things/}{b}

% kangkem
\hi 
\HT{kangkem}{\B{\cp{kang}kem}} \I{["kaN.kEm]} \emph{n.} work (\emph{coll. of} \ding{43}~\HL{tIkangkem}{\B{tì\-kang\-kem}}). 
\B{Kang\-kem val!} Work hard!~\LINK{http://naviteri.org/2022/12/trr-anawm-polaheiem-the-great-day-has-arrived/}{b} 
(\emph{a. idiom}) \B{Po\-ri kang\-kem lu kawng\-kem.} He's really lazy [lit.: `for him, work is a crime'].~\LINK{http://naviteri.org/2023/09/aawa-ayliu-si-aylifyavi-amip-a-few-new-words-and-expressions/}{b}

% Kapsavan
\hi 
\HT{kapsavan}{\B{Kapsavan}} \I{[kap\textcorner.sa.van]} \emph{n.} the river in the Valley of Mo\-'a\-ra along which the Shaman of Songs resides, the Na'\-vi River Journey river.~\LINK{https://forum.learnnavi.org/language-updates/expansion-to-flora-fauna-and-places-from-the-valley-of-mo'ara/}{ln} 

% kar
\hi
\HT{kar}{\B{kar}} \I{[kaR]} \emph{vtr.} teach. 
\B{Lo\-ak\-ìl pä\-nu\-to\-lìng fu\-ta kar oe\-ru fya\-'ot a 'ìp fko nem\-fa ewll.} Loak promised he'd teach me how to vanish into the bushes.~\LINK{http://naviteri.org/2013/01/awvea-posti-zisita-amip-first-post-of-the-new-year/}{b}
\B{Stxe\-nu\-to\-lìng oel fu\-ta lì'\-fyat le\-Na'\-vi po\-e\-ru ki\-var.} I offered to teach her Na'vi.~\LINK{http://naviteri.org/2011/09/miscellaneous-vocabulary/}{b}
\B{Ta\-fral kì\-syar a ay\-sä\-num\-vi sì ay\-lì'\-fya\-vi lu fyin.} That's why I'll teach lessons and expressions that are simple.~\LINK{http://naviteri.org/2012/07/tskxekengtsyip-a-mikyunfpi-a-little-listening-exercise/}{b} 
\B{Fo\-ri tsa\-fne\-tì\-ku\-sar ke lu le\-ha'.} That kind of teaching isn't appropriate for them.~\LINK{http://naviteri.org/2020/11/vospxivopeya-ayliu-amip-novembers-new-words/}{b}  \\
\textsc{Derivatives}: \ding{43} \on{1}.~\HL{karyu}{\B{ka\-ryu}}, teacher. 

% kara
\hi 
\HT{kara}{\B{ka\cp{ra}}} \I{[ka."Ra]} (\emph{a. vin.}) resist. 
\B{Fol nga\-ti spo\-le\-'e a krr, nga lum\-pe ke ka\-ra?} When they captured you, why didn't you resist?~\LINK{http://naviteri.org/2023/09/aawa-ayliu-si-aylifyavi-amip-a-few-new-words-and-expressions/}{b}
(\emph{b. vtr.}) resist s.t.\slash s.o.
\B{Ay\-sä\-la\-tem\-it a za\-mo\-lu\-nge Saw\-tu\-tel nga fmi ki\-va\-ra. Lä\-ngu kel\-tsun.} You try to resist the changes brought by the Sky People. Sadly, that's impossible.~\LINK{http://naviteri.org/2023/09/aawa-ayliu-si-aylifyavi-amip-a-few-new-words-and-expressions/}{b} \\
\textsc{Derivatives}: \ding{43} \on{1}.~\HL{tIkara}{\B{tì\-ka\-ra}}, resistance. 

% karyu
\hi
\HT{karyu}{\B{\cp{kar}yu}} \I{["kaR.ju]} \emph{n.} teacher.
\B{Lo\-lu tok\-tor Kì\-rey\-sì kar\-yu a\-sìl\-tsan.} Doctor Grace was a good teacher.~\LINK{http://naviteri.org/2010/07/ting-mikyun-fte-tslivam-listening-comprehension-1-3/}{b} 
\B{Kar\-yul fngo\-lo' fu\-ta ay\-nu\-me\-yu pi\-va\-te ye'\-krr.} The teacher required the students to arrive early.~\LINK{http://naviteri.org/2012/01/mipa-zisit-ayliu-amip-new-words-for-the-new-year/}{b}
\B{Nga ze\-ne ay\-lì\-'ut kar\-yu\-ä yi\-vu\-ne, ma 'evi.} You must listen to your teacher, my son.~\LINK{http://naviteri.org/2012/11/renu-ayinanfyaya-the-senses-paradigm/}{b}
\B{Fì\-sä\-tsan\-'ul\-ìri nge\-yä lu oe\-ru sko har\-yu nrra.} As your teacher, I'm proud of your improvement.~\LINK{http://naviteri.org/2013/03/tsanerul-feerul-getting-better-getting-worse/}{b}
(\ding{43}~\HL{kar}{\B{kar}}) \\
\textsc{Derivatives}: \ding{43} \on{1}.~\HL{karyunay}{\B{kar\-yu\-nay}}, apprentice teacher.

% karyunay
\hi
\HT{karyunay}{\B{karyu\cp{nay}}} \I{[kaR.ju."naj]} \emph{n.} apprendice teacher. 
(\ding{43}~\HL{karyu}{\B{kar\-yu}})

% katir
\hi 
\HT{katir}{\B{\cp{ka}tir}} \I{["ka.tiR]} \emph{n.} rainbow. 

% katìng
\hi 
\HT{katIng}{\B{\cp{ka}tìng}} \I{["ka.t$\cdot$IN]} \emph{vtr., inf.}\on{22} distribute. 
\B{Eyk\-yul ay\-swi\-zaw\-ti ka\-to\-lìng ay\-ha\-pxì\-tur tsam\-po\-ngu\-ä.} The leader distributed the arrows to the members of the war party.~\LINK{https://naviteri.org/2024/02/mipa-ayliu-si-aylifyavi-niul-more-new-words-and-expressions/}{b} 
(\ding{43}~\HL{ka}{\B{ka}}, \HL{tIng}{\B{tìng}}) \\
\textsc{Derivatives}: \ding{43} \on{1}.~\HL{tIkatIng}{\B{tì\-ka\-tìng}}, distribution. 

% kato
\hi
\HT{kato}{\B{\cp{ka}to}} \I{["ka.to]} \emph{n.} rhythm. 
\B{Tom\-pa\-yä ka\-to | Tsaw\-ke\-yä ka\-to | Ka\-tot tä\-ftxu oel | Nì\-ean nì\-rim.} The rhythm of the rain | The rhythm of the sun | I weave the rhythm | In blue and yellow.~\LINK{http://wiki.learnnavi.org/index.php/Weaving_Song}{wiki}

% kavan
\hi 
\HT{kavan}{\B{\cp{ka}van}} \I{["ka.van]} \emph{vtr.} support (\emph{physically}). 
\B{Fol kar\-ma\-van ko\-ak\-tet teng\-krr fme\-ri po ti\-vì\-ran.} They supported the old woman as she was trying to walk.~\LINK{http://naviteri.org/2019/06/50a-liu-amip-40-new-words/}{b}
(\ding{43}~\HL{slan}{\B{slan}})

% kavuk
\hi
\HT{kavuk}{\B{ka\cp{vuk}}} \I{[ka."vuk\textcorner]} \emph{or} \I{[."vUk\textcorner]} (\textsc{rn}: \B{ka\cp{vùk}} \I{[ka."vUk\textcorner]}) \textsc{i}.~\emph{n.} treachery, betrayal. 
\B{Lo\-lu ka\-vuk, slä Tse\-nul tì\-ngay\-it ko\-lu\-lat.} There was treachery, but Tsenu revealed the truth.~\LINK{http://naviteri.org/2011/10/more-vocabulary-a-bit-of-grammar/}{b} \\
\textsc{ii}.~\HL{kavuk si}{\B{kavuk si}} \I{[ka."vuk\textcorner{} si]} \emph{vin.} betray. 
\B{Po aw\-nga\-ru ka\-vuk so\-li.} He betrayed us.~\LINK{http://naviteri.org/2010/08/a-na\%E2\%80\%99vi-alphabet/comment-page-1/\#comment-189}{b}

% kaw'it
\hi
\HT{kaw'it}{\B{kaw\cp{'it}}} \I{[kaw."Pit\textcorner]} \emph{adv.} (\emph{used with} \ding{43}~\HL{ke}{\B{ke}} \emph{to mean}) (not) a bit, (not) at all.
\B{Txan\-krr ngal ke lawk pot kaw\-'it.} You haven't mentioned a thing about him in a long time.~\LINK{http://naviteri.org/2011/09/\%E2\%80\%9Cby-the-way-what-are-you-reading\%E2\%80\%9D/}{b} 
\B{Ke lu po tsam\-si\-yu kaw\-'it!} There's not a warrior's bone in his whole body!~\LINK{http://forum.learnnavi.org/language-updates/ke-kawit/msg174572/\#msg174572}{ln}
(\emph{a. conv.}) \B{Ke srätx kaw\-'it!\slash Kea sä\-srätx kaw\-'it!\slash Ke\-he kaw\-'it!} `Not at all!' [lit.: (it) doesn't bother (me) at all\slash no bother at all]~\LINK{http://naviteri.org/2020/11/vospxivopeya-ayliu-amip-novembers-new-words/}{b}
(\ding{43}~\HL{'it}{\B{'it}})

% kawkrr
\hi
\HT{kawkrr}{\B{\cp{kaw}krr}} \I{["kaw.k\s{r}]} \emph{n., adv.} no time; never (\emph{requires negation of the verb}). 
\B{Kaw\-krr ke sla\-yu nga Na'\-vi\-yä ha\-pxì!} You will never become part of the People.~\LINK{http://wiki.learnnavi.org/index.php?title=Na\%27vi_from_Avatar_Movie\#Jake.2C_honest_about_betrayal}{a\slash wiki}
\B{Fay\-lì\-'u a\-lor oe\-ru te\-ya si nì\-txan, ul\-te ke tswa\-ya' oel sat kaw\-krr.} These beautiful words fill me with joy, and I will never forget them.~\LINK{http://forum.learnnavi.org/language-updates/life-death-and-gerunds/}{ln}
(\ding{43}~\HL{krr}{\B{krr}})

% kawkxan
\hi 
\HT{kawkxan}{\B{kaw\cp{kxan}}} \I{[kaw."k'an]} \emph{adj.} free, unblocked, unobstructed, clear. 
\B{Nga tsun ki\-vä set. Fya\-'o lu kaw\-kxan.} You can go now. The way is clear.~\LINK{http://naviteri.org/2018/03/100a-liu-amip-64-new-words-part-1/}{b}
(\ding{43}~\HL{ekxan}{\B{e\-kxan}})

% kawl
\hi 
\HT{kawl}{\B{kawl}} \I{[kawl]} \emph{adv.} hard, diligently. 
\B{Mak\-to kawl, ma sam\-si\-yu, fte tsi\-vun pi\-vä\-hem nì\-win!} Ride hard, warriors, so you can get there fast!~\LINK{http://naviteri.org/2017/09/zisikrr-amip-ayliu-amip-new-words-for-the-new-season/}{b} 
(\ding{43}~\HL{val}{\B{val}}) 

% kawlo
\hi 
\HT{kawlo}{\B{\cp{kaw}lo}} \I{["kaw.lo]} \emph{adv.} not once (\emph{requires negation of the verb}). 
\B{Oel keng kawlo ke solar kea räptulì'fyat ngahu!} I have never even once used vulgar language with you!~\LINK{https://naviteri.org/2024/09/pxevola-liu-amip-two-dozen-new-words/}{b} 
\B{Oe kaw\-lo ke ko\-lä tsa\-tseng\-ne.} I have never once gone there.~\LINK{https://naviteri.org/2024/09/pxevola-liu-amip-two-dozen-new-words/\#comment-50641}{b}
(\ding{43}~\HL{ke}{\B{ke}}, \HL{'awlo}{\B{'aw\-lo}})

% kawnomum
\hi 
\HT{kawnomum}{\B{kaw\cp{no}mum}} \I{[kaw."no.mum]} \emph{or} \I{[.mUm]} (\textsc{rn}: \B{kaw\cp{no}nùm} \I{[kaw."no.nUm]}) \emph{adj.} unknown. 
(\ding{43}~\HL{ke}{\B{ke}}, \HL{awn}{‹\B{awn}›}, \HL{omum}{\B{omum}})

% kawng
\hi
\HT{kawng}{\B{kawng}} \I{[kawN]} \emph{adj.} bad, evil. 
\B{su\-te\-kip nì\-wotx, ftxey Na'\-vi ftxey Saw\-tu\-te, lu sìl\-tsan lu kawng}, among all people, whether Na'vi or Skypeople, there are good and bad (ones).~\LINK{http://naviteri.org/2010/07/ting-mikyun-fte-tslivam-listening-comprehension-1-3/}{b} 
\B{Lam set fwa Saw\-tu\-te a\-kawng ho\-lum.} Now it seems that the evil Skypeople have left.~\LINK{http://naviteri.org/2010/07/ting-mikyun-fte-tslivam-listening-comprehension-1-3/}{b} 
\B{Kawng\-a hem a\-na\-fì'u lo\-lä\-ngen alo a\-pxay mì ok\-vur ay\-oe\-yä.} Many such bad actions have happened many times in our history.~\LINK{http://naviteri.org/2020/06/mi-tanlokxe-oeya-srr-afpxamo-terrible-days-in-my-country/}{b} 
(\ding{214}~\HL{sIltsan}{\B{sìl\-tsan}}) \\
\textsc{Derivatives}: \ding{43} \on{1}.~\HL{tIkawng}{\B{tì\-kawng}}, evil. 
\on{2}.~\HL{kawnglan}{\B{kawng\-lan}}, malicious.   
\on{3}.~\HL{kawngsar}{\B{kawng\-sar}}, exploit.
\on{4}.~\HL{kawngtu}{\B{kawng\-tu}}, bad person.
\on{5}.~\HL{kawngkem}{\B{kawng\-kem}}, evil deed, crime.

% kawnglan
\hi
\HT{kawnglan}{\B{\cp{kawng}lan}} \I{["kawN.lan]} \emph{adj.} malicious, bad\hyp{}hearted. 
\B{Hìm\-pxì Saw\-tu\-te\-yä lu tstun\-wi, slä fe\-yä txam\-pxì lä\-ngu kawng\-lan.} A minority of the Sky People are kind, but the majority are malicious.~\LINK{http://naviteri.org/2011/11/\%E2\%80\%99a\%E2\%80\%99awa-ayli\%E2\%80\%99u-amip-ni\%E2\%80\%99aw-\%E2\%80\%94-a-few-new-words-only/}{b}
(\ding{43}~\HL{kawng}{\B{kawng}}, \HL{txe'lan}{\B{txe'\-lan}}) \\
\textsc{Derivatives}: \ding{43} \on{1}.~\HL{tstunlan}{\B{tstun\-lan}}, kind\hyp{}hearted.

% kawngkem
\hi 
\HT{kawngkem}{\B{\cp{kawng}kem}} \I{["kawN.kEm]} \emph{n.} evil deed, crime. 
\B{Tsa\-kawng\-kem\-ìri lu oe layl!} I am innocent of that crime!~\LINK{http://naviteri.org/2023/09/aawa-ayliu-si-aylifyavi-amip-a-few-new-words-and-expressions/}{b} 
(\emph{a. idiom}) \B{Po\-ri kang\-kem lu kawng\-kem.} He's really lazy [lit.: `for him, work is a crime'].~\LINK{http://naviteri.org/2023/09/aawa-ayliu-si-aylifyavi-amip-a-few-new-words-and-expressions/}{b} 
(\ding{43}~\HL{kawng}{\B{kawng}}, \HL{kem}{\B{kem}}) 

% kawngsar
\hi
\HT{kawngsar}{\B{\cp{kawng}sar}} \I{["kawN.s$\cdot$aR]} \emph{vtr., inf.}\on{22} exploit. 
\B{Ay\-nge\-yä ki\-fkey\-ti fol kawng\-sar nì\-tut, fì\-'u\-ri ke\-kem ke si ay\-nga.} They continuously exploit your world and you do nothing about it.~\LINK{http://naviteri.org/2011/07/txantsana-ultxa-mi-siatll-great-meeting-in-seattle/}{b}
(\ding{43}~\HL{kawng}{\B{kawng}}, \HL{sar}{\B{sar}})

% kawngtu
\hi 
\HT{kawngtu}{\B{\cp{kawng}tu}} \I{["kawN.tu]} \emph{n.} bad person, `bad guy'. 
\B{Lal\-a tsa\-rel\-mì a\-ru\-sikx, yem\-stokx tsan\-tul haw\-re'\-ti a\-teyr, kawng\-tul pum\-it a\-la\-yon.} In that old movie, the good guy wears a white hat, the bad guy a black one.~\LINK{http://naviteri.org/2021/12/zolau-niprrte-ma-3746-welcome-2022/}{b}
(\ding{214}~\HL{tsantu}{\B{tsantu}}; \ding{43}~\HL{kawng}{\B{kawng}}, \HL{tute1}{\B{tute}})

% kawtu
\hi
\HT{kawtu}{\B{\cp{kaw}tu}} \I{["kaw.tu]} \emph{pn.} no one, nobody (\emph{requires double negation}). 
\B{kaw\-tu ke tsun ay\-oer tì\-ftang si\-vi}, no one can stop us.~\LINK{http://forum.learnnavi.org/language-updates/stop!/}{ln} 
\B{En\-tu lu tu\-vom\-a ta\-ron\-yu. Kaw\-tut na po ke tso\-le\-'a oel mì sì\-rey.} Entu is an incredible hunter. I've never seen anyone like him before.~\LINK{http://naviteri.org/2015/04/some-new-words-for-may-day/}{b}
\B{Pe\-yä sä\-ftxu\-lì\-'u lo\-lä\-ngu ke\-sran ul\-te kaw\-tur slan\-ti\-re ke si.} Unfortunately his speech was only so-so and inspired no one.~\LINK{http://naviteri.org/2012/10/mipa-vospxi-mipa-ayliu-new-words-for-the-new-month/}{b}
(\ding{43}~\HL{tute1}{\B{tu\-te}})

% kawtseng
\hi 
\HT{kawtseng}{\B{\cp{kaw}tseng}} \I{["kaw.\t{ts}EN]} \emph{adv.} nowhere (\emph{requires negation of the verb}). 
\B{Nga ke tsun wä\-pi\-van; nga\-ri ke lu ke\-a u\-tu\-ru kaw\-tseng.} You cannot hide; there is no sanctuary for you anywhere.~\LINK{http://naviteri.org/2022/12/trr-anawm-polaheiem-the-great-day-has-arrived/}{b} 
(\ding{43}~\HL{tseng}{\B{tse\-nge}})

% kay
\hi
\HT{kay}{\B{kay}} \I{[kaj]} \emph{adp.} in (\emph{the future}), from now. 
\B{Za\-ya\-'u Saw\-tu\-te fte aw\-nga\-ti ski\-va\-'a kay zì\-sìt a\-pxey!} The Sky People will come to destroy us three years from now!~\LINK{http://naviteri.org/2011/09/miscellaneous-vocabulary/}{b}  
\B{Oe pä\-hem kay\slash maw trr\-pxì\-vi a\-mrr.} I'll be there in five minutes.~\LINK{https://naviteri.org/2024/05/krrteri-about-time/}{b} 
(\ding{214}~\HL{kam}{\B{kam}}) 

% kaym / kaymam / kaymay
\hi
\HT{kaym}{\B{kaym}} \I{[kajm]} \emph{n.} evening; \emph{adv.} in the evening. 
\B{Nän yom kxam\-trr, 'ul 'e\-fu o\-hakx kaym.} The less you eat at noon, the hungrier you'll feel in the evening.~\LINK{http://naviteri.org/2012/02/trr-asawnung-lefpom-happy-leap-day/}{b} 
\B{Kaym\-krr\-ka to\-lel mo\-el ta ay\-ul\-txa\-tu stxe\-lit a\-ko\-sman.} During the evening we received a wonderful gift from the meeting members.~\LINK{http://naviteri.org/2014/08/stxeli-alor-a-beautiful-gift/}{b} \\  
(i.) \B{kay\cp{mam}} \I{[kaj."mam]} \emph{n., adv.} last, yesterday evening. \\
(ii.) \B{kay\cp{may}} \I{[kaj."maj]} \emph{n., adv.} next, tomorrow evening. 

% kä-
\hi
\HT{kA-}{\B{kä-}} \I{[k\ae.]} \emph{non\hyp{}productive part., pref. to mean} away from, out, e.g. in \ding{43} \HL{kA'ArIp}{\B{kä\-'ä\-rìp}}, push; \HL{kAmakto}{\B{kä\-mak\-to}}, ride out; \HL{kApxI}{\B{kä\-pxì}}, rear; \HL{kAsrIn}{\B{kä\-srìn}}, lend; \HL{kAsatseng}{\B{kä\-sa\-tseng}}, out there; etc.

% kä
\hi
\HT{kA}{\B{kä}} \I{[k\ae]} \emph{vin.} go. 
\B{Tsa\-tse\-nge le\-hrr\-ap lu fì\-txan kum\-a tsa\-ne ke kä aw\-nga kaw\-krr.} That place is so dangerous we never go there.~\LINK{http://naviteri.org/2012/06/spring-vocabulary-part-3/}{b}  
\B{Ne\-to rä\-'ä ki\-vä!} Don't go away.~\LINK{http://wiki.learnnavi.org/index.php?title=Corpus\#Good_Morning_America}{wiki} 
\B{Oel new fu\-ta ta\-ron\-yu ki\-vä.} I want the hunter to go.~\LINK{http://forum.learnnavi.org/language-updates/frommerian-email/}{ln}
\B{Sra\-ke ko\-lä nga 'aw\-li\-e ne Nu Yor\-kì?} Have you ever been to New York?~\LINK{http://naviteri.org/2013/01/awvea-posti-zisita-amip-first-post-of-the-new-year/}{b} 
\B{Oe kaw\-krr ne few\-tu\-so\-ka pa\-'o kil\-van\-ä ke ka\-mä.} I never went to the opposite side of the river.~\LINK{http://naviteri.org/2011/01/new-year-new-vocabulary/}{b}
\B{Lam ngay oer, fo ka\-yä ìlä hil\-van tup na'\-rìng.} I bet they'll take the river rather than the forest.~\LINK{http://forum.learnnavi.org/language-updates/concerning-tup/}{ln} \\
\textsc{Derivatives}: \ding{43} \on{1}.~\HL{kAteng} {\B{kä\-teng}}, spend time with.     
\on{2}.~\HL{kllkA}{\B{kll\-kä}}, go down, descend. 
\on{3}.~\HL{fAkA}{\B{fä\-kä}}, go up, ascend. 
\on{4}.~\HL{emkA}{\B{em\-kä}}, cross (s.t.).

% kä'ärìp
\hi
\HT{kA'ArIp}{\B{kä\cp{'ä}rìp}} \I{[k\ae."P$\cdot$\ae.R$\cdot$Ip\textcorner]} \emph{vtr., inf.}\on{23} push (\emph{coll. also} \ding{43}~\HL{kArIp}{\B{kä\-rìp}}). 
(\ding{43}~\HL{'ArIp}{\B{'ärìp}}, \HL{kA-}{\B{kä-}}; \ding{214}~\HL{za'ArIp}{\B{za\-'ä\-rìp}}) 

% kämakto
\hi
\HT{kAmakto}{\B{kä\cp{mak}to}} \I{[k\ae."m$\cdot$ak\textcorner.t$\cdot$o]} \emph{vin., inf.}\on{23} ride out. 
\B{kä\-mak\-to nì\-win}, ride out quickly.~\LINK{http://forum.learnnavi.org/language-updates/naviya-the-other-vocative/}{a\slash ln}
(\ding{43}~\HL{makto}{\B{mak\-to}}, \HL{kA-}{\B{kä-}})

% kämunge
\hi 
\HT{kAmunge}{\B{kä\cp{mu}nge}} \I{[k\ae."m$\cdot$u.N$\cdot$E]} (\textsc{rn}: \B{kä\cp{mù}nge} \I{[k\ae."mU.NE]}) \emph{vtr., inf.}\on{23} take (\emph{away from the speaker}). 
(\ding{43}~\HL{munge}{\B{mu\-nge}}, \HL{kA-}{\B{kä-}}; \ding{214}~\HL{zamunge}{\B{za\-mu\-nge}})

% käpxì
\hi
\HT{kApxI}{\B{kä\cp{pxì}}} \I{[k\ae."p'I]} \emph{n.} rear (part or section). 
\B{Zo\-la\-'u ftu kä\-pxì num\-tseng\-vi\-yä sä\-mewn a\-wok.} From the back of the classroom came a loud sigh.~\LINK{https://naviteri.org/2025/04/mevosinga-liu-amip-twenty-new-words/}{b}
(\ding{214}~\HL{zapxI}{\B{za\-pxì}}) 

% kärìp
\hi 
\HT{kArIp}{\B{\cp{kä}rìp}} \I{["k\ae.RIp\textcorner]} \emph{vtr.} (\emph{coll. form of}) push. 
(\ding{43}~\HL{kA'ArIp}{\B{kä\-'ä\-rìp}}) 

% käsrìn
\hi
\HT{kAsrIn}{\B{kä\cp{srìn}}} \I{[k\ae."sR$\cdot$In]} \emph{vtr.}\on{22} lend.
\B{Sne\-yä ma\-sat\-it pol kä\-sro\-lìn oer.} He lent me his breastplate.~\LINK{http://naviteri.org/2012/01/more-additions-to-the-lexicon/}{b}
(\ding{43}~\HL{srIn}{\B{srìn}}, \HL{kA-}{\B{kä-}}; \ding{214}~\HL{zasrIn}{\B{za\-srìn}}) 

% käteng
\hi
\HT{kAteng}{\B{\cp{kä}teng}} \I{["k$\cdot$\ae.tEN]} \emph{vin., inf.}\on{11} spend time, hang out (with, \HL{hu}{\B{hu}}). 
\B{Nga pe\-su\-hu kä\-teng nì\-trr\-trr.} Who do you typically spend time with?~\LINK{http://naviteri.org/2010/09/getting-to-know-you-part-3/}{b} 
\B{Kä\-teng oe hu ey\-lan Per\-lin\-mì a mrr\-trr lu tsyeym a ke tsun tswi\-va' kaw\-krr.} The week I spent with my friends in Berlin is a treasure that I'll never forget.~\LINK{http://naviteri.org/2013/06/looking-back-looking-forward/}{b}
\B{Kaym\-o zo\-la\-'u fray\-ul\-txa\-tu ne kel\-ku moe\-yä, fte ki\-ve\-i\-ä\-teng nì\-'o'.} All the members came to our house for an evening to hang out in a fun way.~\LINK{http://naviteri.org/2014/08/stxeli-alor-a-beautiful-gift/}{b}
(\ding{43}~\HL{kA-}{\B{kä-}}, \HL{teng}{\B{teng}})

% käsatseng 
\hi
\HT{kAsatseng}{\B{kä\cp{sa}tseng}} \I{[k\ae."{}sa.\t{ts}EN]} \emph{adv.} out there. 
(\ding{43}~\HL{tsatseng}{\B{tsa\-tseng}}, \HL{kA-}{\B{kä-}})

% ke
\hi
\HT{ke}{\B{ke}} \I{[kE]} \emph{adv.} not.
(\emph{a. with verbs}) 
\B{Ke lu oer set krr atxan}, I don't have much time now.~\LINK{http://forum.learnnavi.org/language-updates/verb-for-write/msg139677/\#msg139677}{ln}
\B{oel pot ke spaw} I don't believe him.~\LINK{http://wiki.learnnavi.org/index.php?title=Canon\#Extracts_from_various_emails}{wiki} 
\B{ke rey a tu\-te}, person who doesn't live\slash isn't alive.~\LINK{http://wiki.learnnavi.org/index.php?title=Canon\#Participles}{wiki}
(\emph{b. with} \ding{43}~\HL{kea}{\B{kea}}, \HL{ke'u}{\B{ke\-'u}}, \HL{kawtu}{\B{kaw\-tu}}, \HL{kawkrr}{\B{kaw\-krr}}, \HL{kawtseng}{\B{kaw\-tseng}}, \HL{kaw'it}{\B{kaw\-'it}}, \HL{kekem}{\B{ke\-kem}} \emph{to form}) double negation.
\B{Vay set ke pa\-mä\-hä\-ngem kea tì\-'eyng.} Up to now no answer has arrived.~\LINK{http://forum.learnnavi.org/news-announcements/a-response-from-paul-frommer!/}{ln}
\B{Ke\-'u ke lu ngay.} Nothing is true.~\LINK{https://forum.learnnavi.org/language-updates/double-negatives-not-optional/}{ln}
\B{Ul\-te kaw\-tu ke tsun ay\-oer tì\-ftang si\-vi.} And no one can stop us.~\LINK{http://forum.learnnavi.org/language-updates/stop!/}{ln}
\B{Kaw\-krr ke sla\-yu nga Na'\-vi\-yä ha\-pxì!} You will never be one of The People!~\LINK{http://wiki.learnnavi.org/index.php?title=Na\%27vi_from_Avatar_Movie\#Jake.2C_honest_about_betrayal}{a\slash wiki}
\B{Ke lu po tsam\-si\-yu kaw\-'it!} He is not a worrior at all!~\LINK{http://forum.learnnavi.org/language-updates/ke-kawit/msg174572}{ln}
\B{Ay\-nge\-yä ki\-fkey\-ti fol kawng\-sar nì\-tut, fì\-'u\-ri ke\-kem ke si ay\-nga.} They continuously exploit your world and you do nothing about it.~\LINK{http://naviteri.org/2011/07/txantsana-ultxa-mi-siatll-great-meeting-in-seattle/}{b}
(\emph{c. with adj.s; idiom}) \B{Ke pxan.} `Thank you, I don't deserve your praise.' (\emph{in response to a praise or compliment}) [lit.: not worthy].
(\emph{d. with} \ding{43}~\HL{ki1}{\B{ki}} \emph{to mean}) `not … but rather'. 
\B{Nga pll\-txe ke nì\-fyeyn\-tu ki nì\-'e\-veng.} You speak not like an adult but a child.~\LINK{http://naviteri.org/2010/07/vocabulary-update/}{b}
(\emph{e. with adp.s}) \B{Kll\-kxa\-yem fì\-tì\-kang\-kem oe\-yä ro\-fa--ke io--pum fe\-yä.} My work will stand beside--not above--theirs.~\LINK{http://naviteri.org/2010/06/first-post/}{b} \\
\textsc{Derivatives}: \ding{43} \on{1}.~\HL{ke'aw}{\B{ke\-'aw}}, divided, torn apart. 
\on{2}.~\HL{ke'u}{\B{ke\-'u}}, nothing. 
\on{3}.~\HL{kefyak}{\B{ke\-fyak}}, … right? … isn't it? 
\on{4}.~\HL{kezemplltxe}{\B{ke\-zem\-pll\-txe}}, needless to say. 
\on{5}.~\HL{zenke}{\B{zen\-ke}}, must not. 
\on{6}.~\HL{kekem}{\B{ke\-kem}}, nothing (action).
\on{7}.~\HL{keve'o}{\B{ke\-ve\-'o}}, chaos.
\on{8}.~\HL{kawlo}{\B{kaw\-lo}}, not once.

% ke-
\hi
\HT{ke2}{\B{ke-}} \I{[kE.]} \emph{productive pref. on adj.s to form the opposite, unless there is a separate lexical item}.~\LINK{https://forum.learnnavi.org/language-updates/how-to-form-negativeopposite-adjectivesadverbs/}{ln} 

% ke'aw
\hi
\HT{ke'aw}{\B{ke\cp{'aw}}} \I{[kE."Paw]} \emph{adj.} divided, torn apart, strife\hyp{}ridden. 
(\ding{43}~\HL{ke}{\B{ke}}, \HL{'aw}{\B{'aw}})

% ke'u
\hi
\HT{ke'u}{\B{\cp{ke}'u}} \I{["kE.Pu]} \emph{pn., n.} nothing (\emph{requires negation of the verb}). 
(\emph{a.}) \B{Ke lä\-ngu fì\-po\-stì\-mì ke\-'u a lu mip.} I'm afraid there's nothing new in this post.~\LINK{http://naviteri.org/2013/08/taronway-the-hunt-song/}{b} 
\B{Fì\-ta\-ron\-yu\-tsyìp ke tsun ke\-'ut sti\-vä'\-nì.} This little hunter can't catch anything.~\LINK{http://naviteri.org/2010/07/diminutives-conversational-expressions/}{b} 
(\emph{b. idiom with the dative}, !stress change!) \B{Oeru ke\cp{'u}.} I don't care. [lit.: it's nothing to me]~\LINK{http://naviteri.org/2013/03/whoever-whatever-whenever/}{b}
\B{Nga new yi\-vom 'u\-pet fì\-txon? -- Ke\-tsran. Oe\-ru ke\-'u.} What do you want to eat tonight?  Whatever. I don't care.~\LINK{http://naviteri.org/2013/03/whoever-whatever-whenever/}{b}
\B{Oe tì\-kang\-kem si trr\-txon nì\-wotx, nga\-ru ke\-'u!} I work all day and all night, and you don't give a damn!~\LINK{http://naviteri.org/2013/03/whoever-whatever-whenever/}{b}
(\ding{43}~\HL{'u}{\B{'u}}, \HL{ke}{\B{ke}})

% kea
\hi
\HT{}{\B{\cp{ke}a}} \I{["kE.a]} \emph{adj.} no (… \emph{noun}) (\emph{requires negation of the verb}). 
\B{slä vay set ke pa\-mä\-hä\-ngem kea tì\-'eyng}, but until now no answer arrived.~\LINK{http://forum.learnnavi.org/news-announcements/a-response-from-paul-frommer!/}{ln} 
\B{Ke lu po\-ru kea smu\-ke.} He doesn't have any sisters.~\LINK{http://naviteri.org/2010/07/ting-mikyun-fte-tslivam-listening-comprehension-1-3/}{b}

% kefpomronga'
\hi
\HT{kefpomronga'}{\B{kefpom\cp{ro}nga'}} \I{[kE.fpom."Ro.NaP]} \emph{adj., nfp.} (mentally) unhealthful.
\B{Ma En\-tu, fì\-kem rä\-'ä si\-vi; lu ke\-fpom\-ro\-nga'.} Entu, don't do this; it's not healthy (mentally).~\LINK{http://naviteri.org/2014/10/tson-si-fpomron-obligation-and-mental-health/}{b}
(\ding{43}, \ding{214}~\HL{fpomronga'}{\B{fpom\-ro\-nga'}})

% kefpomtokxnga'
\hi
\HT{kefpomtokxnga'}{\B{kefpom\cp{tokx}nga'}} \I{[kE.fpom."tok'.NaP]} \emph{adj., nfp.} (physcially) unhealthful.
(\ding{43}, \ding{214}~\HL{fpomtokxnga'}{\B{fpom\-tokx\-nga'}})

% keftxo
\hi
\HT{keftxo}{\B{ke\cp{ftxo}}} \I{[kE."ft'o]} \emph{adj.} unhappy, upset; (\emph{for inner feelings used with} \B{'efu}) \B{Krr a sto\-lawm Txe\-wì te\-ri hem a poe\-ru lo\-len, 'o\-le\-fu ke\-ftxo nì\-txan.} When Txewì heard about what happened to her, he was very upset.~\LINK{http://naviteri.org/2010/07/ting-mikyun-fte-tslivam-listening-comprehension-1-3/}{b} 
\B{Po ho\-lum a\-krr\-maw, oe 'e\-fu ke\-ftxo.} After he left, I was unhappy.~\LINK{http://naviteri.org/2011/08/new-vocabulary-clothing/comment-page-1/\#comment-986}{b} 
(\emph{a. idiom}) \B{Ke\-ftxo.} `How sad!' \B{Tìng na\-ri. Tsa\-ye\-rik\-tsyìp\-ìl li 'an\-la sa'\-nok\-it. Ke\-ftxo!} Look! That little hexapede is already yearning for its mother. How sad!~\LINK{http://naviteri.org/2013/01/awvea-posti-zisita-amip-first-post-of-the-new-year/}{b} 
(\emph{b. idiom}) \B{Ngari keftxo (fwa) ke tok.} ``We missed you.\slash Sorry you couldn't make it.\slash Too bad you couldn't be there.''~\LINK{http://naviteri.org/2019/08/aawa-lifya-si-lifyavi-amip-a-few-new-words-and-expressions/}{b} \\
\textsc{Derivatives}: \ding{43} \on{1}.~\HL{tIkeftxo}{\B{tì\-ke\-ftxo}}, sadness. 
\on{2}.~\HL{nIkeftxo}{\B{nì\-ke\-ftxo}}, unfortunately, sadly.

% kefyak
\hi
\HT{kefyak}{\B{ke\cp{fyak}}} \I{[kE."fjak\textcorner]} \emph{idiom, tag question} … right? … isn't it? (\emph{from} \B{ke fì\-fya srak})
\B{Nga lu So\-rewn, ke\-fyak?} You're Sorewn, right?~\LINK{http://naviteri.org/2010/09/getting-to-know-you-part-1/}{b}

% kehe
\hi
\HT{kehe}{\B{\cp{ke}he}} \I{["kE.hE]} \emph{intj.} no. 
\B{Alo a\-mrr po\-an po\-lawm, slä fra\-lo poe pol\-txe san ke\-he.} He asked five times, but each time she said, `no.'~\LINK{http://forum.learnnavi.org/language-updates/txelanit-hivawl/}{ln}
\B{Sra\-ke fnan ngal lì'\-fya\-ti le\-Na'\-vi?~Ke\-he.} Are you good at Na'vi?~No.~\LINK{http://naviteri.org/2010/09/getting-to-know-you-part-3/}{b}
(\emph{a. conv.}) \B{Ke\-he kaw\-'it!} `Not at all!'~\LINK{http://naviteri.org/2020/11/vospxivopeya-ayliu-amip-novembers-new-words/}{b}
(\ding{214}~\HL{sran}{\B{sra\-ne}}) \\
\textsc{Derivatives}: \ding{43} \on{1}.~\HL{kesran}{\B{ke\-sran}}, mediocre. 
\on{2}.~\HL{srankehe}{\B{sran\-ke\-he}}, more or less.

% kekem
\hi
\HT{kekem}{\B{\cp{ke}kem}} \I{["kE.kEm]} \textsc{i}. \emph{adv.} no action or activity, nothing. \\
\textsc{ii}. \HT{kekem si}{\B{\cp{ke}kem ke si}} \I{["kE.kEm kE si]} \emph{vin.} do nothing.
\B{Ay\-nge\-yä ki\-fkey\-ti fol kawng\-sar nì\-tut, fì\-'u\-ri ke\-kem ke si ay\-nga.} They continuously exploit your world and you do nothing about it.~\LINK{http://naviteri.org/2011/07/txantsana-ultxa-mi-siatll-great-meeting-in-seattle/}{b}
(\ding{43}~\HL{ke}{\B{ke}}, \HL{kem}{\B{kem}})

% Kekunan
\hi
\B{\cp{Ke}kunan} \I{["kE.ku.nan]} \emph{n.} \emph{clan name, specialise in banshee riding}.

% kelemweypey
\hi
\HT{kelemweypey}{\B{kelem\cp{wey}pey}} \I{[kE.lEm."wEj.pEj]} \emph{adj.} impatient.
(\ding{214}, \ding{43}~\HL{lemweypey}{\B{lem\-wey\-pey}})

% kelfpomron
\hi
\HT{kelfpomron}{\B{kelfpom\cp{ron}}} \I{[kEl.fpom."Ron]} \emph{adj., ofp.} (mentally) unhealthy.
\B{Ke tsun nga tì\-kxìm si\-vi oer. Lu nga kel\-fpom\-ron!} You can't order me around. You're mentally unsound!~\LINK{http://naviteri.org/2014/10/tson-si-fpomron-obligation-and-mental-health/}{b}
(\ding{43}, \ding{214} \HL{lefpomron}{\B{le\-fpom\-ron}})

% kelfpomtokx
\hi
\HT{kelfpomtokx}{\B{kelfpom\cp{tokx}}} \I{[kEl.fpom."tok']} \emph{adj.} (physcially) unhealthy.
(\ding{43}, \ding{214} \HL{lefpomtokx}{\B{le\-fpom\-tokx}})

% kelha'
\hi
\HT{kelha'}{\B{kel\cp{ha'}}} \I{[kEl."haP]} \emph{adj.} inappropriate, unsuitable, unfitting.~\LINK{http://naviteri.org/2020/11/vospxivopeya-ayliu-amip-novembers-new-words/\#comment-31844}{b}
(\ding{43}~\HL{leha'}{\B{le\-ha'}})

% kelhoan
\hi
\HT{kelhoan}{\B{kel\cp{ho}an}} \I{[kEl."ho.an]} \emph{adj.} uncomfortable.
\B{Lä\-ngu fì\-seyn kel\-ho\-an nì\-ngay. Tsun oe hi\-veyn tsaw\-sìn 'a\-'aw\-a swaw\-tsyìp nì\-'aw.} This chair is really uncomfortable. I can only sit on it a few seconds.~\LINK{http://naviteri.org/2015/08/aylifyavi-lereyfya-2-cultural-terms-2-and-more/}{b}
(\ding{43}~\HL{hoan}{\B{ho\-an}}, \ding{214}~\HL{lehoan}{\B{le\-ho\-an}})

% kelkin
\hi
\HT{kelkin}{\B{kel\cp{kin}}} \I{[kEl."kin]} \emph{adj.} unnecessary.  (\ding{43}, \ding{214}~\HL{lekin}{\B{le\-kin}}) \\
\textsc{Derivatives}: \ding{43} 
\on{1}.~\HL{nIkelkin}{\B{nì\-kel\-kin}}, unnecessarily.

% kelku
\hi
\HT{kelku}{\B{\cp{kel}ku}} \I{["kEl.ku]} \textsc{i}. \emph{n.} home, house. 
\B{Lu me\-nge\-yä kel\-ku na'\-rìng\-ä lu\-ke kxu a\-txan a fì\-'u fmawn a\-sìl\-tsan lu nì\-ngay.} It's truly good news that your forest house is without much harm.\LINK{http://forum.learnnavi.org/language-updates/luke-kxu-without-harm/}{ln} 
\B{Aw\-ngal tok kel\-kut.} We're home.~\LINK{http://naviteri.org/2011/10/more-vocabulary-a-bit-of-grammar/}{b} 
\B{Ay\-nga\-ri li\-vu hel\-ku sang ul\-te te'\-lan le\-fpom.} May your homes be warm and your hearts be happy.~\LINK{http://naviteri.org/2013/09/on-si-salewfya-shapes-and-directions/}{b} 
\B{Ut\-ral a lu few pay\-fya a eo kel\-ku oe\-yä tsawl lu nì\-txan.} The tree on the other side of the stream in front of my house is very tall.~\LINK{http://naviteri.org/2011/01/new-year-new-vocabulary/}{b} \\
\textsc{ii}. \HT{kelku si}{\B{\cp{kel}ku si}} \emph{vin.} live, dwell. 
\B{Mì sray apxa kel\-ku si so\-a\-i\-a\-hu.} (He) lives with his family in a large village.~\LINK{http://naviteri.org/2010/07/ting-mikyun-fte-tslivam-listening-comprehension-1-3/}{b}
\B{Tsun aw\-nga kel\-ku si\-vi nì\-'eng Saw\-tu\-te\-hu mì atx\-kxe aw\-nge\-yä.} We can share our land with the Skypeople.~\LINK{http://naviteri.org/2012/01/more-additions-to-the-lexicon/}{b} \\
\textsc{Derivatives}: \ding{43} 
\on{1}.~\HL{kelutral}{\B{Ke\-lut\-ral}}, Hometree.  
\on{2}.~\HL{prrku}{\B{prr\-ku}}, womb.

% kelpxìmrun
\hi 
\HT{kelpxImrun}{\B{kel\cp{pxìm}run}} \I{[kEl."p'Im.Run]} \emph{or} \I{[.RUn]} \emph{adj.} rare. 
(\ding{43}~\HL{lepxImrun}{\B{le\-pxìm\-run}}) \\
\textsc{Derivatives}: \ding{43} 
\on{1}.~\HL{tIkelpxImrun}{\B{tì\-kel\-pxìm\-run}}, rarity, rareness. 

% kelsar
\hi
\HT{kelsar}{\B{kel\cp{sar}}} \I{[kEl."saR]} \emph{adj.} useless, in vain. 
(\ding{214}~\HL{lesar}{\B{le\-sar}})

% keltrrtrr
\hi 
\HT{keltrrtrr}{\B{kel\cp{trr}trr}} \I{[kEl."t\s{r}.t\s{r}]} \emph{adj.} unusual.  
\B{Kel\-trr\-trr\-a tì'eylan.} An unusual friendship.~\LINK{http://naviteri.org/2020/04/keltrrtrra-tieylan-an-unusual-friendship/}{b}
(\ding{43}~\HL{ke}{\B{ke}}, \HL{letrrtrr}{\B{le\-trr\-trr}}) \\
\textsc{Derivatives}: \ding{43} \on{1}.~\HL{txankeltrrtrr}{\B{txan\-kel\-trr\-trr}}, extraordinary.

% keltsun
\hi
\HT{keltsun}{\B{kel\cp{tsun}}} \I{[kEl."\t{ts}un]} \emph{or} \I{[."\t{ts}Un]} (\textsc{rn}: \B{kel\cp{tsùn}} \I{[kEl."\t{ts}Un]}) \emph{adj.} impossible. 
\B{Fwa fpìl nì\-fya\-'o a\-ve\-nga' fì\-tì\-zin\-mì lu kel\-tsun.} It's impossible to think clearly (i.e., in an organized manner) in this muddled situation.~\LINK{http://naviteri.org/2021/04/mipa-ayliu-mipa-sioeykting-new-words-new-explanations/\#comment-33483}{b}
(\ding{214}~\HL{letsunslu}{\B{le\-tsun\-slu}}) 

% Kelutral
\hi
\HT{kelutral}{\B{\cp{Kel}utral}} \I{["kEl.ut\textcorner.Ral]} \emph{n.} Hometree (\emph{home of the Omaticaya clan}). 
\B{Kel\-ut\-ral lu fne\-ut\-ral a\-zey.} Hometree is a special kind of tree.~\LINK{http://naviteri.org/2012/06/spring-vocabulary-part-3/}{b}
\B{Saw\-tu\-te ze\-ra\-'u fte fol Kel\-ut\-ral\-ti ski\-va\-'a.} The Skypeople are coming to destroy Hometree.~\LINK{http://wiki.learnnavi.org/index.php?title=Corpus\#Behind_the_Scenes}{wiki} 
\B{mì hil\-van a lok Kel\-ut\-ral}, in the river near Hometree.~\LINK{http://naviteri.org/2010/07/thoughts-on-ambiguity/}{b}
(\ding{43}~\HL{kelku}{\B{kel\-ku}}, \HL{utral}{\B{ut\-ral}})

% kem
\hi
\HT{kem}{\B{kem}} \I{[kEm]} \textsc{i}. \emph{n.} action, deed. 
\B{Kea kem le\-yew\-la rä\-'ä si, ru\-txe.} Please, don't let me down.~\LINK{http://naviteri.org/2012/10/mipa-vospxi-mipa-ayliu-new-words-for-the-new-month/}{b}
\B{Krr a sto\-lawm Txe\-wì te\-ri hem a poe\-ru lo\-len, 'o\-le\-fu ke\-ftxo nì\-txan.} When Txewì heard about what happened to her, he was very upset.~\LINK{http://naviteri.org/2010/07/ting-mikyun-fte-tslivam-listening-comprehension-1-3/}{b} 
\B{Hem le\-'en tsun le\-hrr\-ap li\-vu.} Speculative moves can be dangerous.~\LINK{http://naviteri.org/2011/01/new-year-new-vocabulary/}{b}
\B{Krr\-o len ay\-hem a\-na\-fì\-'u.} Sometimes such things happen.~\LINK{http://forum.learnnavi.org/language-updates/just-another-few-words/}{ln} 
\B{Saw\-tu\-te\-yä hem\-ìri lu aw\-nga\-ru ya\-yayr.} The Skypeople's actions confuse us.~\LINK{http://naviteri.org/2011/01/new-year-new-vocabulary/}{b}
(\emph{idiom}) \B{Kem\-ìri a nga\-ru prr\-te' ke lu, tsa\-kem rä\-'ä si\-vi ay\-la\-he\-ru.} As for any action that you don't like, don't act that way to others.~\LINK{http://wiki.learnnavi.org/index.php?title=Corpus\#Good_Morning_America}{wiki} 
(\emph{a. proverb}) \B{Hem a\-sru\-nga' nì\-'ul, hum a\-sngu\-nga' nì\-nän.}
More helpful actions lead to less troubling outcomes.~\LINK{http://naviteri.org/2015/03/kap-si-ayunil-saylahe-threats-dreams-and-other-things/}{b} \\
\textsc{ii}. \B{\cp{kem} si} \I{["kEm si]} \emph{vin.} do (an action).   \B{Kem si li trr\-ay\-sre.} Do it by tomorrow.~\LINK{http://naviteri.org/2011/02/new-vocabulary-part-2/}{b}
\B{Oel tse\-'a kem\-it a tsa\-tse\-nge nga so\-li.} I see what you did there.~\LINK{http://forum.learnnavi.org/language-updates/just-another-few-words/}{ln} \\
% \textsc{iii}. \B{fì\cp{kem}} \I{[fI."kEm]} \emph{pn., n.} this (action), it.
% \B{Ru\-txe, ay\-nga\-ri fì\-kem fe\-yä 'e\-'a\-la to\-pur ka\-ngay sì\-yi nì\-'aw.} Please, this will only confirm their worst fears about you.~\LINK{http://forum.learnnavi.org/language-updates/two-words-in-use-eal-and-kangay-si/msg564284/\#msg564284}{ln} \\
% \textsc{iv}. \B{\cp{tsa}kem} \I{["\t{ts}a.kEm]} \emph{pn., n.} that (action), it.
% \B{Tsa\-kem kop krr\-nekx.} That also takes time.~\LINK{http://forum.learnnavi.org/language-updates/just-a-few-interesting-little-words/}{ln} \\
% \textsc{v}. \B{\cp{kem}pe} \I{["kEm.pE]} \emph{inter.} what (action)?
% \B{Kem\-pe le\-ren?} What's happening?~\LINK{http://naviteri.org/2011/05/some-miscellaneous-vocabulary/}{b} \\
% \textsc{vi}. \B{\cp{kem}pe si} \I{["kEm.pE si]} \emph{inter.} do what (action)? 
% \B{Kem\-pe si nga, ma sa'\-nu\-tsyìp?} What are you doing, little mommy?~\LINK{http://naviteri.org/2010/07/diminutives-conversational-expressions/}{b} \\
\textsc{Derivatives}: \ding{43} 
\on{1}.~\HL{kempe}{\B{kem\-pe\slash pe\-hem}}, what (action). 
\on{2}.~\HL{fIkem}{\B{fì\-kem}}, this (action).
\on{3}.~\HL{tsakem}{\B{tsa\-kem}}, that (action).
\on{4}.~\HL{fekem}{\B{fe\-kem}}, accident. 
\on{5}.~\HL{tIkangkem}{\B{tì\-kang\-kem}}, work.
\on{6}.~\HL{kemlI'u}{\B{kem\-lì\-'u}}, verb.
\on{7}.~\HL{kawngkem}{\B{kawng\-kem}}, evil deed, crime.

% kemlì'u 
\hi
\HT{kemlI'u}{\B{\cp{kem}lì'u}} \I{["kEm.lI.Pu]} \emph{n.} verb. 
\B{Kem\-lì\-'ut alu var fkol sar pxel pum alu new.} You can use the verb `var' like (the verb) `new'.~\LINK{http://naviteri.org/2011/03/word-order-and-case-marking-with-modals/comment-page-1/\#comment-582}{b} 
(\ding{43}~\HL{kem}{\B{kem}}, \HL{lI'u}{\B{lì\-'u}}) \\
\textsc{Derivatives}: \ding{43} \on{1}.~\HL{kemlI'uvi}{\B{kem\-lì\-'u\-vi}}, infix.

% kemlì'uvi
\hi
\HT{kemlI'uvi}{\B{\cp{kem}lì'uvi}} \I{["kEm.lI.Pu.vi]} \emph{n.} verb infix. 
\B{Tsa\-lì\-'u\-ri alu zey\-ko lu kem\-lì\-'u\-vi (lì\-'u a\-txan\-tsan nang!) mì kam\-tseng.} Concerning the verb `zeyko', there is an infix (excellent word!) in the middle.~\LINK{http://naviteri.org/2010/07/vocabulary-update/comment-page-1/\#comment-95}{b}
(\ding{43}~\HL{kemlI'u}{\B{kem\-lì\-'u}})

% kempe / si
\hi
\HT{kempe}{\B{\cp{kem}pe}} \emph{or} \B{pe\cp{hem}} \I{["kEm.pE]} \emph{or} \I{[pE."hEm]} \textsc{i}. \emph{inter.} what (action)?
\B{Kem\-pe le\-ren?} What's happening?~\LINK{http://naviteri.org/2011/05/some-miscellaneous-vocabulary/}{b} \\
\textsc{ii}. \B{\cp{kem}pe si} \emph{or} \B{pe\cp{hem} si} \I{["kEm.pE si]} \emph{or} \I{[pE."hEm si]} \emph{vin., inter.} do what (action)? 
\B{Kem\-pe si nga, ma sa'\-nu\-tsyìp?} What are you doing, little mommy?~\LINK{http://naviteri.org/2010/07/diminutives-conversational-expressions/}{b}
\B{Pol po\-le\-'un fu\-ta pe\-hem si nì\-'en.} He decided what to do on a hunch.~\LINK{http://naviteri.org/2011/01/new-year-new-vocabulary/}{b}

% kemu'nitkan
\hi 
\HT{kemu'nitkan}{\B{ke\cp{mu'}nitkan}} \I{[kE."muP.nit\textcorner.kan]} \emph{or} \I{[."mUP.]} \emph{adj.} ineffective. 
(\ding{214}~\HL{mu'nitkan}{\B{mu'\-nit\-kan}}) 

% kemuia / si
\hi 
\HT{kemuia}{\B{ke\cp{mu}ia}} \I{[kE."mu.i.a]} \textsc{i}. \emph{n.} dishonor. \\
\textsc{ii}. \B{ke\cp{mu}ia si} \I{[kE."mu.i.a si]} \emph{vin.} dishonor.
\B{Nge\-yä tì\-fna\-we' ke\-muia so\-li fì\-soaia\-ru.} Your cowardice has dishonored this family.~\LINK{http://naviteri.org/2020/05/aawa-liu-amip-si-vurway-alor-a-few-new-words-and-a-beautiful-poem/}{b}
(\ding{214}~\HL{meuia}{\B{meuia}}; \ding{43}~\HL{meuia}{\B{meuia}}, \HL{ke}{\B{ke}}) \\
\textsc{Derivatives}: \ding{43} \on{1}.~\HL{kemuianga'}{\B{ke\-mu\-ia\-nga'}}, dishonorable.

% kemuianga'
\hi 
\HT{kemuianga'}{\B{ke\cp{mu}ianga'}} \I{[kE."mu.i.a.NaP]} \emph{adj.} dishonorable. 
\B{Peyä hem\-ìl a\-ke\-muia\-nga' za\-mo\-lu\-nge fwìng\-it ay\-oer.} His dishonorable behavior (lit.: actions) humiliated us (lit.: brought humiliation to us).~\LINK{http://naviteri.org/2020/05/aawa-liu-amip-si-vurway-alor-a-few-new-words-and-a-beautiful-poem/}{b}
(\ding{43}~\HL{kemuia}{\B{ke\-mu\-ia}}, \HL{-nga'}{\B{-nga'}})

% kemum
\hi 
\HT{kemum}{\B{ke\cp{mum}}} \I{[kE."mum]} \emph{or} \I{[."mUm]} (\textsc{rn}: \B{ke\cp{mùm}} \I{[kE."mUm]})~:~\HL{omum}{\B{omum}}

% kemwiä
\hi 
\HT{kemwiA}{\B{kem\cp{wi}ä}} \I{[kEm."wi.\ae]} \emph{adj.} improper, unfair, wrong, unjustified. 
\B{Lu kem\-wi\-ä!} It's unfair.~\LINK{http://naviteri.org/2017/09/zisikrr-amip-ayliu-amip-new-words-for-the-new-season/}{b} 
(\ding{214}~\HL{muiA}{\B{muiä}}) \\
\textsc{Derivatives}: \ding{43} \on{1}.~\HL{tIkemwiA}{\B{tì\-kem\-wiä}}, unfairness, injustice.
\on{2}.~\HL{leymkem}{\B{leym\-kem}}, protest.

% ken
\hi 
\HT{ken}{\B{ken}}\textsuperscript{\on{1}} \I{[kEn]} \emph{adp.} despite, in spite of. 
\B{Ken tì\-naw\-ri pe\-yä, ke flo\-lä.} Despite her talent, she didn't succeed.~\LINK{https://naviteri.org/2024/02/mipa-ayliu-si-aylifyavi-niul-more-new-words-and-expressions/}{b}
\B{ken fwa lu por 'aw\-a na\-ri nì\-'aw, lu Ma\-ti ta\-ron\-yu a\-swey.} In spite of having only one eye, Ma\-ti is the best hunter.~\LINK{https://naviteri.org/2024/02/mipa-ayliu-si-aylifyavi-niul-more-new-words-and-expressions/}{b}

% ken
\hi 
\HT{ken2}{\B{ken}}\textsuperscript{\on{2}} \I{[kEn]} \emph{vin.} behave assertively and with confidence. 
\B{Eyk\-tan a\-sìl\-tsan, ken ay\-skxe mì te'\-lan, zene ki\-ven fra\-krr.} A good leader, despite the stones in his heart, must always behave confidently.~\LINK{https://naviteri.org/2025/11/kaltxi-nimun-hello-again/}{b} 
(\emph{a. idiom}) \B{Kll\-kxem ki\-ven!} `Stand up straight!' [lit.: ``stand straight and act with confidence'']~\LINK{https://naviteri.org/2025/11/kaltxi-nimun-hello-again/}{b}\slash \LINK{https://naviteri.org/2026/01/follow-up-to-txep-si-txeva/}{a\on{3}}

% ken'aw
\hi
\HT{ken'aw}{\B{ken\cp{'aw}}} \I{[kEn."Paw]} \emph{adv.} not only (\emph{together with} \ding{43}~\HL{slAkop}{\B{slä\-kop}}).
\B{Nge\-yä tsmu\-ke lu ken\-'aw lor slä\-kop ka\-nu.} Your sister is not only beautiful but also intelligent.~\LINK{http://naviteri.org/2014/10/tson-si-fpomron-obligation-and-mental-health/}{b}
\B{Ay\-nga\-ri sìl\-pey oe tsnì ken\-'aw fpom\-tokx slä\-kop fpom\-ron yi\-vo'.} I hope that not only your physical but also your mental health will be perfect.~\LINK{http://naviteri.org/2014/10/tson-si-fpomron-obligation-and-mental-health/}{b}
(\ding{43}~\HL{ke}{\B{ke}}, \HL{nI'aw}{\B{nì\-'aw}})

% kenong
\hi
\HT{kenong}{\B{\cp{ke}nong}} \I{["kE.noN]} \emph{vtr.} model, represent, exemplify.
\B{Fay\-hem\-ìl pe\-yä ke ke\-nong tì\-sayt a fkol fngo' po\-ta.} These actions of his do not represent the loyalty that is required of him.~\LINK{http://naviteri.org/2022/01/aawa-tipangkxotsyip-a-teri-horen-lifyaya-a-few-little-discussions-about-grammar/}{b}
 \\
\textsc{Derivatives}: \ding{43} \on{1}.~\HL{tIkenong}{\B{tì\-ke\-nong}}, example.

% kenten
\hi
\HT{kenten}{\B{\cp{ken}ten}} \I{["kEn.tEn]} \emph{n., fauna} fan lizard, \emph{fanisaurus pennatus}.
\B{Mìn, mìn, ken\-ten mìn. | Pe\-lun ke 'efu spxin?} Turn, turn, the fanlizard turns. | Why doesn't it feel sick?~\LINK{http://forum.learnnavi.org/language-updates/learnings-from-the-siatlla-song-translations/msg473566/\#msg473566}{ln}
(\emph{a. idiom}) \B{(Na) ken\-ten mì kum\-pay}, being in an environment where one is prevented from acting naturally or doing what one wants to do, [lit.: (like) a fan lizard in gel] \B{Nar\-mew oe fo\-ru na'\-rìng\-ä tì\-lor\-it wi\-vìn\-txu, slä ke tsä\-ngun fo tsli\-vam. 'Efu oe na ken\-ten mì kum\-pay.} I wanted to show them the beauty of the forest, but sadly, they weren't able to understand. I felt completely stymied.~\LINK{http://naviteri.org/2014/11/twenty-before-the-holidays/}{b}

% kenzen
\hi
\HT{kenzen}{\B{ken\cp{zen}}} \I{[kEn."zEn]} \emph{adv.} not necessarily.
\B{Fay\-sä\-fpìl fay\-sì\-'e\-fu\-sì lu pum oe\-yä nì\-'aw, ken\-zen pum su\-te\-yä a za\-mo\-lu\-nge aw\-ngar rel\-it a\-ru\-sikx alu Unil\-tì\-ran\-tokx.} These thoughts and feelings are only mine, not necessarily of the people who brought us the movie \emph{Avatar}.~\LINK{http://naviteri.org/2016/12/tengkrr-perahem-zisit-amip-as-the-new-year-arrives/}{b}
(\ding{43}~\HL{nIzen}{\B{nì\-zen}})

% keng
\hi
\HT{keng}{\B{keng}} \I{[kEN]} \emph{adv.} even.
\B{Yom tey\-lut keng oel!} Even \emph{I} eat teylu!~\LINK{http://forum.learnnavi.org/language-updates/even-keng/}{ln} 
\B{Ngay\-txoa, krro tì\-kxey si keng kar\-yu.} Sorry, sometimes even the teacher messes up.~\LINK{http://naviteri.org/2012/02/trr-asawnung-lefpom-happy-leap-day/}{b}
\B{Ke su\-nu oer keng nì\-'it.} I don't like (them) one bit.~\LINK{http://naviteri.org/2013/01/awvea-posti-zisita-amip-first-post-of-the-new-year/}{b}

% keri
\hi 
\HT{keri}{\B{\cp{ke}ri}} \I{["kE.Ri]} \emph{adj.} fierce (\emph{of people or things}). 
\B{tsam\-si\-yu a\-ke\-ri}, fierce warrior or things.~\LINK{https://naviteri.org/2025/11/kaltxi-nimun-hello-again/}{b} 
\B{ay\-lì\-'u a\-ke\-ri}, fierce words.~\LINK{https://naviteri.org/2025/11/kaltxi-nimun-hello-again/}{b} 
(\emph{a. idiom}) \B{Ay\-lì\-'u a\-ke\-ri | Le\-kxu to tì\-tse\-ri.} Fierce words can be more harmful than you think.~\LINK{https://naviteri.org/2025/11/kaltxi-nimun-hello-again/}{b} (\emph{longer form:}) \B{Ay\-lì\-'u a\-ke\-ri tsun le\-kxu li\-vu nì\-'ul to tì\-tse\-ri.}

% kerìsmìsì
\hi 
\HT{kerIsmIsI}{\B{Ke\cp{rì}smìsì}} \I{[kE."RI.smI.sI]} (\emph{gen.} \B{Ke\-rì\-smì\-sä}) \textsc{lw}, \emph{n.} Christmas.
\B{Ftxo\-zä Ke\-rì\-smì\-sä Le\-fpom.} Merry Christmas.~\LINK{http://naviteri.org/2022/01/aawa-tipangkxotsyip-a-teri-horen-lifyaya-a-few-little-discussions-about-grammar/}{b}

% kerusey
\hi
\HT{kerusey}{\B{\cp{ke}rusey}} \I{["kE.Ru.sEj]} \emph{adj.} dead.
\B{He\-tu\-wong\-ìl aw\-nge\-yä swo\-tut ska\-'a, fte kll\-ki\-vu\-lat ke\-ru\-sey\-a tskxet.} The aliens destroy our sacred place to dig up dead rock.~\LINK{http://forum.learnnavi.org/language-updates/participles/}{ln}
\B{Krro krro, flì\-a vul a\-ru\-sey to nutx\-a pum a\-ke\-ru\-sey lu txur.} A thin living branch is sometimes stronger than a thick dead one.~\LINK{http://naviteri.org/2011/11/\%E2\%80\%99a\%E2\%80\%99awa-ayli\%E2\%80\%99u-amip-ni\%E2\%80\%99aw-\%E2\%80\%94-a-few-new-words-only/}{b} 
(\ding{214}~\HL{rusey}{\B{ru\-sey}}) 

% kesätarenga'
\hi 
\HT{kesAtarenga'}{\B{kesä\cp{ta}renga'}} \I{[kE.s\ae."ta.RE.NaP]} \emph{or coll.} \I{[kE."{}sta.RE.NaP]} \emph{adj.} irrelevant. 
(\ding{214}~\HL{sAtarenga'}{\B{sä\-ta\-re\-nga'}})

% kesran
\hi
\HT{kesran}{\B{ke\cp{sran}}} \I{[kE."{}sRan]} \emph{adj.} so\hyp{}so, mediocre (in quality).
\B{Pe\-yä sä\-ftxu\-lì\-'u lo\-lä\-ngu ke\-sran ul\-te kaw\-tur slan\-ti\-re ke si.} Unfortunately his speech was only so-so and inspired no one.~\LINK{http://naviteri.org/2012/10/mipa-vospxi-mipa-ayliu-new-words-for-the-new-month/}{b} 
\B{Fu\-ri\-a tä\-ftxu ngal tok yì\-pet?~Tok yìt a\-ke\-sran.} How's your weaving?~I'm so-so.~\LINK{http://naviteri.org/2013/01/lifyari-po-peyi-how-good-is-her-navi/}{b}
(\ding{43}~\HL{kehe}{\B{ke\-he}}, \HL{sran}{\B{sra\-ne}}) \\
\textsc{Derivatives}: \ding{43} \on{1}.~\HL{nIksran}{\B{nìk\-sran}}, in a mediocre manner.

% ketartu
\hi 
\HT{ketartu}{\B{ke\cp{tar}tu}} \I{[kE."taR.tu]} \emph{n.} outcast. 
(\ding{43}~\HL{tare}{\B{ta\-re}}) 

% keteng
\hi
\HT{keteng}{\B{\cp{ke}teng}} \I{["kE.tEN]} \emph{adj.} different. 
\B{Slä sa'\-nok\-ìri fpìl\-fya lu ke\-teng.} But mother's way of thinking is different.~\LINK{http://naviteri.org/2010/07/ting-mikyun-fte-tslivam-listening-comprehension-1-3/}{b} 
(\ding{214}~\HL{teng}{\B{teng}})  \\
\textsc{Derivatives}: \ding{43} 
\on{1}.~\HL{tIketeng}{\B{tì\-ke\-teng}}, difference.
\on{2}.~\HL{ketengsyon}{\B{ke\-teng\-syon}}, variety, diversity.

% ketengsyon
\hi 
\HT{ketengsyon}{\B{\cp{ke}tengsyon}} \I{["kE.tEN.sjon]} \emph{n.} variety, diversity.
\B{Naw\-num\-tseng\-ä ay\-nu\-me\-yu\-ri ley ay\-nga\-ru tì\-teng\-wotx, ay\-oe\-ru ke\-teng\-syon.} Regarding university students, you value uniformity, we value diversity.~\LINK{https://naviteri.org/2025/06/vospikin-lefpom-happy-july/}{b} 
(\ding{43}~\HL{keteng}{\B{keteng}}, \HL{syon}{\B{syon}}) \\
\textsc{Derivatives}: \ding{43} 
\on{1}.~\HL{leketengsyon}{\B{le\-ke\-teng\-syon}}, diverse, various.

% ketrìp
\hi 
\HT{ketrIp}{\B{\cp{ket}rìp}} \I{["kEt\textcorner.RIp\textcorner]} \emph{adj.} unfavorable, inauspicious. 
\B{Nga ket\-rìp\-a krr zo\-la\-'u; Ra\-lu set sti ul\-te ke new nga\-hu pi\-väng\-kxo.} You came at the wrong time; Ralu is angry and won't speak with you.~\LINK{http://naviteri.org/2019/06/50a-liu-amip-40-new-words/}{b}
(\ding{214}~\HL{etrIp}{\B{et\-rìp}})

% ketsran
\hi
\HT{ketsran}{\B{ke\cp{tsran}}} \I{[kE."\t{ts}Ran]} \textsc{i}. \emph{adj.} (\emph{to form compound relative pronouns, i.e.}) whatever, whoever, however. 
\B{Ke\-tsran\-a tu\-te a ni\-vew hi\-vum tsun tsa\-kem si\-vi.} Whoever wants to leave can do so.~\LINK{http://naviteri.org/2013/03/whoever-whatever-whenever/}{b}
\B{'U a\-ke\-tsran tsun ti\-vam.} Anything at all will be fine.~\LINK{http://naviteri.org/2013/03/whoever-whatever-whenever/}{b} \\
\textsc{ii}. \emph{conj.} 
\B{Ke\-tsran tseng\-ne nga ki\-vä, kä oe tsa\-tseng nì\-teng.} Wherever you go, I'll go there too.~\LINK{http://naviteri.org/2013/03/whoever-whatever-whenever/}{b}
\B{Ke\-tsran fya\-'o si\-vu\-nu ngar, kem si.} Do it however you'd like.~\LINK{http://naviteri.org/2013/03/whoever-whatever-whenever/}{b} \\
\textsc{iii}. (\emph{idiom}) \B{Nga new yi\-vom 'u\-pet fì\-txon?--Ke\-tsran. Oe\-ru ke\-'u.} What do you want to eat tonight?--Whatever. I don't care.~\LINK{http://naviteri.org/2013/03/whoever-whatever-whenever/}{b}
(\ding{43}~\HL{tsranten}{\B{tsran\-ten}})

% ketstun
\hi 
\HT{ketstun}{\B{ke\cp{tstun}}} \I{[kE."\t{ts}tun]} \emph{or} \I{[."\t{ts}tUn]} \emph{adj., ofp.} unkind, emotionally cold. 
(\emph{a. proverb}) \B{Tu\-te a\-ke\-tstun sli\-vu 'ey\-lan ke tsun.} An unkind person cannot be a friend.~\LINK{https://naviteri.org/2025/06/vospikin-lefpom-happy-july/}{b}  
(\ding{43}~\HL{tstunwi}{\B{tstun\-wi}}, \HL{tIktstunga'}{\B{tìk\-tstu\-nga'}}) \\
\textsc{Derivatives}: \ding{43} 
\on{1}.~\HL{tIktstun}{\B{tìk\-tstun}}, unkindness, emotional coldness.

% ketsuk-
\hi
\HT{ketsuk}{\B{ketsuk-}} \I{[kE.\t{ts}uk\textcorner.]} \emph{or} \I{[.\t{ts}Uk\textcorner.]} (\textsc{rn}: \B{ketsùk-} \I{[kE.\t{ts}Uk\textcorner.]}) \emph{prod. pref. on vtr. to form receptive ability, that s.t. is not capable of receiving the action of the verb}, ``un-…-able''. \B{ke\-tsuk\-tswa'}, unforgettable.~\LINK{http://naviteri.org/2011/03/\%E2\%80\%9Creceptive-ability\%E2\%80\%9D-and-hesitation/}{b} 
(\ding{214}~\HL{tsuk}{\B{tsuk-}})

% ketsukanom
\hi 
\HT{ketsukanom}{\B{ketsu\cp{ka}nom}} \I{[kE.\t{ts}u."ka.nom]} (\textsc{rn}: \B{ketsù\cp{ka}nom} \I{[kE.\t{ts}U."ka.nom]}) \emph{adj.} unavailable, unobtainable. 
(\ding{43}~\HL{ketsuk}{\B{ketsuk-}}, \HL{kanom}{\B{ka\-nom}}; \ding{214}~\HL{tsukanom}{\B{tsu\-ka\-nom}})

% ketsuklewn
\hi 
\HT{ketsuklewn}{\B{ketsuk\cp{lewn}}} \I{[kE.\t{ts}uk\textcorner."lEwn]} \emph{or} \I{[.\t{ts}Uk\textcorner.]} (\textsc{rn}: \B{ketsùk\cp{lewn}} \I{[kE.\t{ts}Uk\textcorner."lEwn]}) \emph{adj.} intolerable, unaccaptable, unbearable. 
\B{Tsa\-fne\-vo\-ìk lu ke\-tsuk\-lewn.} That kind of behavior is intolerable.~\LINK{http://naviteri.org/2021/03/mipa-ayliu-si-aylifyavi-new-words-and-expressions/}{b}
(\ding{43}~\HL{ketsuk}{\B{ke\-tsuk-}}, \HL{lewn}{\B{lewn}}) 

% ketsuktiam
\hi
\HT{ketsuktiam}{\B{ketsuk\cp{ti}am}} \I{[kE.\t{ts}uk\textcorner."ti.am]} \emph{or} \I{[.\t{ts}Uk\textcorner.]} (\textsc{rn}: \B{ketsùk\cp{ti}am} \I{[kE.\t{ts}Uk\textcorner]."ti.am}) \emph{adj.} uncountable (\emph{for countable nouns}).
\B{Hol\-pxay san\-hì\-yä a mì saw lu ke\-tsuk\-ti\-am; keng ke tsun fko tsi\-ve\-'a sat nì\-wotx.} The number of stars in the sky is infinite; it's not even possible to see them all.~\LINK{http://naviteri.org/2014/11/twenty-before-the-holidays/}{b} 
(\ding{43}~\HL{tiam}{\B{ti\-am}}, \HL{txewluke}{\B{txew\-lu\-ke}})

% ketsuktswa'
\hi
\HT{ketsuktswa'}{\B{ketsuk\cp{tswa'}}} \I{[kE.\t{ts}uk\textcorner."\t{ts}waP]} \emph{or} \I{[.\t{ts}Uk\textcorner.]} (\textsc{rn}: \B{ketsùk\cp{tswa'}} \I{[kE.\t{ts}Uk\textcorner."\t{ts}waP]}) \emph{adj.} unforgettable.
(\ding{43}~\HL{tswa'}{\B{tswa'}})

% ketuwong
\hi
\HT{ketuwong}{\B{\cp{ke}tuwong}} \I{["kE.tu.woN]} \emph{n.} alien (person, creature).
\B{Txo\-pu rä\-'ä si, lu ke\-tu\-wong\-o nì\-'aw.} Don't be afraid, it's only an alien.~\LINK{http://forum.learnnavi.org/language-updates/a-collection/}{ln\slash a\textsuperscript{c}}
\B{He\-tu\-wong\-ìl aw\-nge\-yä swo\-tut ska\-'a, fte kll\-ki\-vu\-lat ke\-ru\-sey\-a tskxet.} The aliens destroy our sacred place to dig up dead rock.~\LINK{http://forum.learnnavi.org/language-updates/participles/}{ln}
\B{Tsa\-swi\-rä\-ti lo\-nu! Ay\-nga ne\-to ri\-vikx! Fì\-ke\-tu\-wong\-ti oel stì\-ye\-ftxaw.} Release this creature! Step back! I will look at this alien.~\LINK{http://wiki.learnnavi.org/index.php?title=Na\%27vi_from_Avatar_Movie\#Scene_in_Hometree}{wiki\slash a}
(\ding{43}~\HL{kewong}{\B{ke\-wong}})

% kew
\hi
\HT{kew}{\B{kew}} \I{[kEw]} \emph{num.} zero, 0. 

% key
\hi
\HT{key}{\B{key}} \I{[kEj]} \emph{n.} face. 
\B{Oe\-ri key tun slo\-lu nì\-txan.} My face became very red.~\LINK{http://naviteri.org/2011/04/\%E2\%80\%99a\%E2\%80\%99awa-li\%E2\%80\%99fyavi-amip\%E2\%80\%94a-few-new-expressions/comment-page-1/\#comment-694}{b}
\B{Nga\-ri key ne\-re\-i\-yrr, ma 'ey\-lan.} Your face is glowing, friend.~\LINK{http://naviteri.org/2012/07/fivospxiya-aylifyavi-amip-this-months-new-expressions-pt-2/}{b} \\
\textsc{Derivatives}: \ding{43} \on{1}.~\HL{keyrel}{\B{key\-rel}}, facial expression.

% keykur
\hi
\HT{keykur}{\B{key\cp{kur}}} \I{[kEj."k$\cdot$uR]} \emph{or} \I{[."kUR]} (\textsc{rn}: \B{key\cp{kùr}} \I{[kEj."kUR]}) \emph{vtr., inf.}\on{22} hang (on, onto, \ding{43}~\HL{sIn}{\B{sìn}}), let hang.
\B{'A\-li\-'ät vul\-sìn key\-kur za\-'u fì\-tseng.} Hang the choker on the branch and come here.~\LINK{http://naviteri.org/2013/04/wheres-the-bathroom-and-other-useful-things/}{b}
(\ding{43}~\HL{kur}{\B{kur}})

% keyn
\hi 
\HT{keyn}{\B{keyn}} \I{[kEjn]} \emph{vtr.} put down. 
\B{Ngey tsko\-ti ki\-veyn. Li ye\-rik ho\-li\-fwo.} Put down your bow. The hexapede has already run away.~\LINK{http://naviteri.org/2018/03/100a-liu-amip-64-new-words-part-1/}{b} 
(\ding{214}~\HL{kxeltek}{\B{kxeltek}}) \\
\textsc{Derivatives}: \ding{43} \on{1}.~\HL{keynven}{\B{keyn\-ven}}, step.

% keynven
\hi 
\HT{keynven}{\B{keyn\cp{ven}}} \I{[k$\cdot$Ejn."vEn]} \emph{vin., inf.}\on{11} step. 
\B{Na\-ri si teng\-krr ke\-reyn\-ven fì\-tseng. Lu kll\-te ekx\-txu.} Step carefully here. The ground is rough.~\LINK{http://naviteri.org/2019/08/aawa-lifya-si-lifyavi-amip-a-few-new-words-and-expressions/}{b}
\B{Na\-ri si keyn\-ven!} Step carefully!~\LINK{http://naviteri.org/2019/08/aawa-lifya-si-lifyavi-amip-a-few-new-words-and-expressions/}{b}
(\emph{a. idiom}) \B{Po keyn\-ven sìn ketse.} He is socially awkward. (\emph{lit.}: `He steps on tails.')~\LINK{http://naviteri.org/2019/08/aawa-lifya-si-lifyavi-amip-a-few-new-words-and-expressions/}{b} 
(\ding{43}~\HL{keyn}{\B{keyn}}, \HL{venu}{\B{ve\-nu}}) \\
\textsc{Derivatives}: \ding{43}
\on{1}.~\HL{sAkeynven}{\B{sä\-keyn\-ven}}, step.

% keyrel
\hi
\HT{keyrel}{\B{\cp{key}rel}} \I{["kEj.rEl]} \emph{n.} facial expression.
\B{O\-lo\-me\-i\-um oel fa key\-rel fu\-ta lu yaw\-ne oe po\-ru.} I knew by the expression on her face that she loves me.~\LINK{http://naviteri.org/2011/05/some-miscellaneous-vocabulary/}{b}
(\ding{43}~\HL{key}{\B{key}}, \HL{rel}{\B{rel}})

% keve'o
\hi 
\HT{keve'vo}{\B{ke\cp{ve}'o}} \I{[kE."vE.Po]} \emph{n.} chaos, disorder. 
\B{Nge\-yä ay\-sä\-fpìl\-mì tse\-'ä\-nga oel ke\-ve\-'ot nì\-'aw.} In your thinking, I see only chaos.~\LINK{https://naviteri.org/2024/09/pxevola-liu-amip-two-dozen-new-words/}{b} 
(\ding{43}~\HL{ve'o}{\B{ve\-'o}}) \\
\textsc{Derivatives}: \ding{43} 
\on{1}.~\HL{lekve'o}{\B{lek\-ve\-'o}}, chaotic. 
\on{2}.~\HL{kekve'otu}{\B{ke\-ve\-'o\-tu}}, chaotic. 

% keve'otu
\hi 
\HT{keve'otu}{\B{ke\cp{ve}'otu}} \I{[kE."vE.Po.tu]} \emph{n.} troublemaker, \emph{one who habitually creates chaos and disorder}. 
\B{Ay\-sä\-leym\-fe'\-it pe\-yä yu\-ne rä\-'ä. Lu po ke\-ve\-'o\-tu nì\-'aw.} Don't listen to his complaints. He's just a troublemaker.~\LINK{https://naviteri.org/2025/03/conlang-adventure-and-more/}{b} 
(\ding{43}~\HL{keve'o}{\B{ke\-ve\-'o}})

% kewan
\hi 
\HT{kewan}{\B{\cp{ke}wan}} \I{["kE.wan]} \emph{n.} age. 
\B{Nge\-yä ke\-wan pìm\-txan?} How old are you? [formal and stiff]~\LINK{http://naviteri.org/2016/06/mrrvola-lifyavi-amip-forty-new-expressions/}{b} 
(\emph{a. proverb}) \B{Koakturi kewanti keyìl ke wan.} Some things can't be covered up. [lit.: `An old person's face doesn't hide their age.']~\LINK{http://naviteri.org/2016/06/mrrvola-lifyavi-amip-forty-new-expressions/}{b} 
(\ding{43}~\HL{koak}{\B{ko\-ak}}, \HL{'ewan}{\B{'e\-wan}})

% kewong
\hi
\HT{kewong}{\B{\cp{ke}wong}} \I{["kE.woN]} \emph{adj.} alien, unknown.
\B{Oe\-ri ta pe\-yä fa\-hew a\-ke\-wong on\-tu te\-ya lä\-ngu.} My nose is full of his alien smell.~\LINK{http://languagelog.ldc.upenn.edu/nll/?p=1977}{ll} \\
\textsc{Derivatives}: \ding{43} \on{1}.~\HL{ketuwong}{\B{ke\-tu\-wong}}, alien.

% keyawr
\hi
\HT{keyawr}{\B{ke\cp{yawr}}} \I{[kE."jawR]} \emph{adj.} wrong, incorrect. 
\B{Ul\-te tsa\-lì'\-fya\-vi nge\-yä ke lu ke\-yawr kaw\-'it.} And your expression isn't wrong at all.~\LINK{http://naviteri.org/2011/08/new-vocabulary-clothing/comment-page-1/\#comment-995}{b}
(\ding{214}~\HL{eyawr}{\B{eyawr}})

% keye'ung
\hi
\HT{keye'ung}{\B{ke\cp{ye}'ung}} \I{[kE."jE.PuN]} \emph{or} \I{[.PUN]} (\textsc{rn}: \B{ke\cp{ye}ùng} \I{[kE."jE.UN]}) \emph{n.} insanity. 
\B{Ta\-fral ke lì\-syek oel nge\-yä ke\-ye\-'ung\-it.} Therefore I will not heed your insanity.~\LINK{http://forum.learnnavi.org/language-updates/ltasygt-with-ke-and-nga/}{ln\slash g}

% kezemplltxe
\hi
\HT{kezemplltxe}{\B{kezempll\cp{txe}}} \I{[kE.zEm.p\s{l}."t'E]} \emph{adv.} of course, needless to say. 
\B{New Va'\-ru tskot Ey\-tu\-kan\-ä za\-sri\-vìn.~Ke\-zem\-pll\-txe pa\-yll\-txe san ke\-he.} Va'ru wants to borrow Eytukan's bow.~Of course he'll say no.~\LINK{http://naviteri.org/2012/06/spring-vocabulary-part-3/}{b}
(\ding{43}~\HL{ke}{\B{ke}}, \HL{zene}{\B{ze\-ne}}, \HL{plltxe}{\B{pll\-txe}})

% kezin
\hi 
\HT{kezin}{\B{\cp{ke}zin}} \I{["kE.zin]} \emph{adj.} untangled. 
(\ding{214}~\HL{zin}{\B{zin}}, \HL{ke}{\B{ke}}) \\
\textsc{Derivatives}: \ding{43}
\on{1}.~\HL{tIkezin}{\B{tì\-ke\-zin}}, s.t. in an untangled state; solution.

% ki
\hi
\HT{ki1}{\B{ki}} \I{[ki]} \emph{conj. with} \ding{43}~\HL{ke}{\B{ke}} \emph{to mean} not … but rather, not … but instead. 
\B{Nga pll\-txe ke nì\-fyeyn\-tu ki nì\-'e\-veng.} You speak not like an adult but a child.~\LINK{http://naviteri.org/2010/07/vocabulary-update/}{b} 
\B{Tsyokx\-ìri ke lu zek\-wä a\-tsìng ki a\-mrr.} The hand doesn't have four but five digits.~\LINK{http://naviteri.org/2010/07/ting-mikyun-fte-tslivam-listening-comprehension-1-3/}{b}

% ki-
\hi
\HT{ki2}{\B{ki-}} \I{[ki.]} \emph{num., pref.} seven (\emph{base to form higher numbers or fractions}). 
\B{\cp{ki}vol}, \on{56} (decimal), \on{70} (octal); 
\B{ki\cp{pxì}}, one-seventh, $\frac{1}{7}$; 
\B{ki\-pxì a\-tsìng}, four-seventh, $\frac{4}{7}$.
(\ding{43}~\HL{kinA}{\B{ki\-nä}})

% kian
\hi 
\HT{kian}{\B{ki\cp{an}}} \I{[ki."{}an]} \emph{vtr.} blame. 
\B{Oe\-ti kian rä\-'ä! Ke no\-lui oe!} Don't blame me! It wasn't my fault.~\LINK{http://naviteri.org/2021/04/mipa-ayliu-mipa-sioeykting-new-words-new-explanations/}{b} \\
\textsc{Derivatives}: \on{1}.~\HL{kiantu}{\B{ki\-an\-tu}}, blameworthy person.

% kiantu
\hi 
\HT{kiantu}{\B{ki\cp{an}tu}} \I{[ki."{}an.tu]} \emph{n.} blameworthy person. 
\B{Fì\-tì\-snaytx\-ìri lu nge\-yä tsmu\-kan ki\-an\-tu.} Your brother is to blame for this loss.~\LINK{http://naviteri.org/2021/04/mipa-ayliu-mipa-sioeykting-new-words-new-explanations/}{b}
(\ding{43}~\HL{kian}{\B{ki\-an}})

% kifkey
\hi
\HT{kifkey}{\B{ki\cp{fkey}}} \I{[ki."fkEj]} \emph{n.} the (physical, solid) world. 
\B{Pll\-txe Saw\-tu\-te san ki\-fkey\-ìl ay\-oe\-yä tok txan\-a ngip\-it a san\-hì\-kip.} The Sky People say that their world is in the great space among the stars.~\LINK{http://naviteri.org/2015/03/kap-si-ayunil-saylahe-threats-dreams-and-other-things/}{b}
\B{tì\-kang\-kem si\-vi fo\-hu a Unil\-tì\-ran\-tokx\-it sì ki\-fkey\-ti Ey\-wa\-'e\-veng\-ä za\-mo\-lu\-nge aw\-ngar}, working with those that brought us \emph{Avatar} and the World of Pandora.~\LINK{http://forum.learnnavi.org/news-announcements/a-response-from-paul-frommer!/}{ln}
\B{Ay\-wayl yìm ki\-fkey\-ä, 'ì\-he\-yut a\-vo\-mrr}, The Songs bind the thirteen spirals of the solid world.~\LINK{http://www.pandorapedia.com/navi/music/ritual_music}{pp}
\B{Oe new si\-vop ne tsa\-ki\-fkey le\-ron\-srel.} I want to journey to that imaginary world.~\LINK{http://naviteri.org/2011/02/new-vocabulary-part-2/}{b} 
(i.) \B{tì\-fti\-a ki\-fkey\-ä}, science [lit.: the study of the physical world].~\LINK{http://wiki.learnnavi.org/index.php?title=Corpus\#Science_Magazine}{wiki} 
(ii.) \B{tì\-fti\-a\-tu ki\-fkey\-ä}, scientist.~\LINK{http://naviteri.org/2015/03/kap-si-ayunil-saylahe-threats-dreams-and-other-things/}{b} \\
\textsc{Derivatives}: \ding{43} \on{1}.~\HL{fkeytok}{\B{fkey\-tok}}, exist.

% kilvan
\hi
\HT{kilvan}{\B{kil\cp{van}}} \I{[kil."van]} \emph{n.} river. 
\B{kil\-van a\-reng}, shallow river; \B{kil\-van a\-txukx}, deep river; \B{ngima kilvan}, long river. 
\B{Na\-ri si, ma 'itan. Kil\-van tswe\-sya si nì\-win nì\-txan.} Be careful, son. The river is flowing very swiftly.~\LINK{http://naviteri.org/2018/04/100a-liu-amip-64-new-words-part-2/}{b}
\B{Kil\-van\-ä tì\-fkey\-tok lu fya\-pe fì\-trr?} What's the condition of the river today?~\LINK{http://naviteri.org/2011/04/yafkeykiri-plltxe-frapo-everyone-talks-about-the-weather/}{b}
\B{Kil\-van\-ìri tsa\-tseng slo\-snep\-pe?} How wide is the river there?~\LINK{http://naviteri.org/2012/06/spring-vocabulary-part-3/}{b}

% kin
\hi
\HT{kin}{\B{kin}} \I{[kin]} \emph{vtr.} need. 
\B{Tì\-ta\-ron\-ìri lä\-ngu tì\-kak\-pam fe\-kum. Kin fkol fra\-inan\-fyat.} I'm sorry to say that deafness is a disadvantage for hunting. You need all your senses.~\LINK{http://naviteri.org/2015/08/aylifyavi-lereyfya-2-cultural-terms-2-and-more/}{b} 
\B{Po 'e\-fu ngeyn ul\-te kin tì\-rawn\-it nì\-txan.} He is tired and very much needs to be replaced.~\LINK{http://naviteri.org/2012/01/mipa-zisit-ayliu-amip-new-words-for-the-new-year/}{b} \\
\textsc{Derivatives}: \ding{43} \on{1}.~\HL{tIkin}{\B{tì\-kin}}, (the) need. 
\on{2}.~\HL{lekin}{\B{le\-kin}}, necessary. 

% kinam
\hi
\HT{kinam}{\B{ki\cp{nam}}} \I{[ki."nam]} \emph{n.} leg.
\B{Ut\-ral\-ä a\-nawm | Ay\-ri\-na' lu ay\-oeng | A pe\-yä tì\-txur mì hi\-nam aw\-nge\-yä.} We are all seeds | Of the Great Tree | Whose strength is in our legs.~\LINK{http://www.pandorapedia.com/navi/music/ritual_music}{pp} \\
\textsc{Derivatives}: \ding{43} \on{1}.~\HL{kinamtil}{\B{ki\-nam\-til}}, knee.

% kinamtil
\hi
\HT{kinamtil}{\B{ki\cp{nam}til}} \I{[ki."nam.til]} \emph{n.} knee. (\ding{43}~\HL{kinam}{\B{kinam}}, \HL{til}{\B{til}})

% kinä
\hi
\HT{kinA}{\B{\cp{ki}nä}} \I{["ki.n\ae]} \emph{num.} seven. (\emph{base}: \ding{43}~\HL{ki2}{\B{ki-}}, \emph{rest}: \ding{43}~\HL{hin}{\B{\mbox{-hin}}}) 
\B{Maria kxam\-lä na'\-rìng kä a zì\-sìt\-o a\-ki\-nä ke lal\-mu tsar kea rìk.} Maria went through the forest that had no leaf for seven years.~\LINK{http://forum.learnnavi.org/language-updates/for-\%28duration\%29-jesus-loan-word/}{ln} \\
\textsc{Derivatives}: \ding{43} \on{1}.~\HL{kintrr}{\B{kin\-trr}}, (\on{7}-day) week. 
\on{2}.~\HL{vospxIkin}{\B{Vo\-spxì\-kin}}, July.

% kintrr
\hi
\HT{kintrr}{\B{\cp{kin}trr}} \I{[kin.t\s{r}]} \emph{n., adv.} (\on{7}-day) week.
\B{Maw fpxa\-mo\-a kin\-trr a\-fpxa\-mo mì tan\-lo\-kxe oe\-yä …} After a truly horrible week in my country … \LINK{http://naviteri.org/2022/05/results-of-the-tttc/}{b} \\
\textsc{Derivatives}: \ding{43} \on{1}.~\B{kintrr\cp{am}} \emph{n., adv.}~last week. 
\B{Kin\-trr\-am fay\-fam\-rel\-si\-yu\-hu oe ul\-txa so\-le\-i\-yi.} Last week I met with those writers.~\LINK{http://naviteri.org/2013/08/fmawn-a-fkol-kilmulat-this-just-in/}{b}
\B{Kin\-trr\-am\-ä sä\-tsyìl lu le\-hrr\-ap slä 'o' nì\-txan.} Last week's climb was dangerous but very exciting.~\LINK{http://naviteri.org/2012/01/more-additions-to-the-lexicon/}{b} \\  
\on{2}.~\B{kintrr\cp{ay}}  \emph{n., adv.}~next week.
\B{Zun nga srung sil\-vi oer, zel ke kì\-ye\-vä oe ne Wa\-syìng\-ton kin\-trr\-ay.} If you had helped me, I wouldn’t be going to Washington next week.~\LINK{http://naviteri.org/2013/04/zun-zel-counterfactual-conditionals/}{b}
(\ding{43}~\HL{kinA}{\B{ki\-nä}}, \HL{trr}{\B{trr}})

% kinglor
\hi 
\HT{kinglor}{\B{\cp{king}lor}} \I{["kiN.loR]} \emph{n., fauna} kinglor. 
(\ding{43}~\HL{kIng}{\B{kìng}}, \HL{lor}{\B{lor}}) 

% kip
\hi
\HT{kip}{\B{kip}} \I{[kip\textcorner]} \emph{adp.} among.
\B{pi\-za\-yu Ney\-ti\-ri\-yä txan\-rar\-mo\-'a fra\-to kip ay\-ha\-pxì\-tu O\-ma\-ti\-ka\-ya\-ä}, Neytiri's ancestor was the most famous among the members of the Omaticaya.~\LINK{http://naviteri.org/2012/03/spring-vocabulary-part-1/}{b}   
\B{Tì\-'e\-fu\-mì oe\-yä su\-te\-kip nì\-wotx, ftxey Na'\-vi ftxey Saw\-tu\-te, lu sìl\-tsan lu kawng.} In my opinion among all people, whether Na'vi or Skypeople, there are good and bad.~\LINK{http://naviteri.org/2010/07/ting-mikyun-fte-tslivam-listening-comprehension-1-3/}{b} \\
\textsc{Derivatives}: \ding{43} \on{1}.~\HL{fkip}{\B{fkip}}, up among. 
\on{2}.~\HL{takip}{\B{ta\-kip}}, from among 

% kipxì
\hi
\HT{kipxI}{\B{ki\cp{pxì}}} \I{[ki."p'I]} \emph{num., n.} one\hyp{}seventh, $\frac{1}{7}$. 
(\ding{43}~\HL{ki2}{\B{ki-}}, \HL{pxI}{\B{-pxì}})

% kiri
\hi 
\HT{kiri}{\B{\cp{Ki}ri}} \I{["ki.Ri]} \emph{n. female name}.  

% kive
\hi
\HT{kive}{\B{\cp{ki}ve}} \I{["ki.vE]} \emph{ord. num., adj.} seventh. 

% kivol
\hi
\HT{kivol}{\B{\cp{ki}vol}} \I{["ki.vol]} \emph{num., adj.} \on{56} (decimal), \on{70} (octal). 
(\ding{43}~\HL{ki2}{\B{ki-}}, \HL{vol}{\B{vol}})

% kì'ong
\hi
\HT{kI'ong}{\B{kì\cp{'ong}}} \I{[kI."PoN]} \emph{adj.} slow. 
(\emph{a. idiom}) \B{Kxe\-tse kì\-'ong!} ``Don't get angry.\slash Calm down.'' [lit.: slow tail] (\emph{short for} \B{Nga\-ri kxe\-tse kì\-'ong li\-vu.}).~\LINK{http://naviteri.org/2021/08/ulte-ayyoratu-leiu-and-the-winners-are-2/}{b}
(\ding{214}~\HL{win}{\B{win}}) \\
\textsc{Derivatives}: \ding{43} \on{1}.~\HL{nIk'ong}{\B{nìk\-'ong}}, slowly.

% kìm
\hi
\HT{kIm}{\B{kìm}} \I{[kIm]} \emph{vtr.} spin. 
\B{Pol rum\-it ko\-lìm.} He spun the ball.~\LINK{http://forum.learnnavi.org/language-updates/ftarpa-and-skienpa-added-to-dict/msg570362/\#msg570362}{ln}

% kìmar
\hi
\HT{kImar}{\B{kì\cp{mar}}} \I{[kI."maR]} \emph{adj.} (\emph{a.~of foods, vegetable or animal}) in season. 
\B{Tey\-lu kì\-mar lì\-yu a fì\-'u oe\-ru te\-ya si.} It fills me with joy that \emph{teylu} is about to be in season.~\LINK{http://naviteri.org/2011/10/more-vocabulary-a-bit-of-grammar/}{b}
(\emph{b.~of events and situations that occur at a convenient or appropriate time}) \B{Kì\-mar lu fwa fo ke päng\-kxo oe\-hu.} It was opportune that they didn't chat with me.~\LINK{http://naviteri.org/2011/10/more-vocabulary-a-bit-of-grammar/}{b} \\
\textsc{Derivatives}: \ding{43} \on{1}.~\HL{nIkmar}{\B{nìk\-mar}}, convenient, in time.

% kìng
\hi
\HT{kIng}{\B{kìng}} \I{[kIN]} \emph{n.} thread. 
\B{na way\-te\-lem\-ä hìng}, like the threads of a song cord.~\LINK{http://www.pandorapedia.com/navi/music/ritual_music}{pp} 
\B{Lä\-ngu fay\-hìng zin nì\-wotx; ke tsun sat si\-var.} Unfortunately, these threads are all tangled up; they can't be used.~\LINK{http://naviteri.org/2021/04/mipa-ayliu-mipa-sioeykting-new-words-new-explanations/}{b} \\
\textsc{Derivatives}: \ding{43} \on{1}.~\HL{lI'ukIng}{\B{lì\-'u\-kìng}}, sentence.

% kìte'e
\hi
\HT{kIte'e}{\B{kìte\cp{'e}}} \I{[kI.tE."PE]} \textsc{i}. \emph{n.} service. \\
\textsc{ii}. \B{kìte\cp{'e} si} \I{[kI.tE."PE si]} \emph{vin.} serve, be of service. \B{Eo nge\-nga kll\-kxem ohe, a\-lak\-si (fte) Oma\-ti\-ka\-ya\-ru kì\-te\-'e si\-vi.} I stand before you, ready to serve the Omaticaya people.~\LINK{http://wiki.learnnavi.org/index.php?title=Na\%27vi_from_Avatar_Movie}{wiki\slash a}

% kìyevame
\hi
\HT{kIyevame}{\B{kìye\cp{va}me}} \I{[kI.jE."va.mE]} \emph{idiom} \textsc{See} you again soon! Goodbye! 
\B{Kì\-ye\-va\-me ul\-te Ey\-wa nga\-hu.} Goodbye for now, and may Eywa be with you.~\LINK{http://forum.learnnavi.org/news-announcements/a-response-from-paul-frommer!/}{ln} 
(\ding{43}~\HL{kame}{\B{ka\-me}})

% kllfro'
\hi
\HT{kllfro'}{\B{kll\cp{fro'}}} \I{[k\s{l}."fR$\cdot$oP]} \emph{vin., inf.}\on{22} be responsible for (\emph{with the topic}), having s.o. or s.t. as a job, duty, or area of concern. 
\B{Po\-ri ze\-ne kll\-fri\-vo' nga.} He is your responsibility [lit.: you are responsible for him].~\LINK{http://forum.learnnavi.org/language-updates/misc-answers/}{ln\slash a\textsuperscript{c}}
\B{Ye'\-rìn 'ì\-yi'\-a sä\-nu\-me a tsa\-ri kll\-fro' oe.} Soon the teaching that I am responsible for will end.~\LINK{http://forum.learnnavi.org/language-updates/trr-rrtaya/}{ln}

% kllkä
\hi
\HT{kllkA}{\B{kll\cp{kä}}} \I{[k\s{l}."k$\cdot$\ae]} \emph{vin., inf.}\on{22} (\emph{a.}) go down, descend. 
(\emph{b. with} \HL{ftu}{\B{ftu}}) get off, dismount.~\LINK{https://naviteri.org/2025/04/mevosinga-liu-amip-twenty-new-words/}{b}
(\emph{c. for astronomical bodies setting})
(\ding{214}~\HL{fAkA}{\B{fä\-kä}}, \ding{43}~\HL{kllte}{\B{kll\-te}}, \HL{kA}{\B{kä}}, \HL{hum}{\B{hum}}) 

% kllkulat
\hi
\HT{kllkulat}{\B{kll\cp{ku}lat}} \I{[k\s{l}."k$\cdot$u.l$\cdot$at\textcorner]} \emph{vtr., inf.}\on{23} dig up.  
\B{He\-tu\-wong\-ìl aw\-nge\-yä swo\-tut ska\-'a, fte kll\-ki\-vu\-lat ke\-ru\-sey\-a tskxet.} The aliens destroy our sacred place to dig up dead rock.~\LINK{http://forum.learnnavi.org/language-updates/participles/}{ln\slash a\textsuperscript{c}}
(\ding{214}~\HL{kllyem}{\B{kll\-yem}}, \ding{43}~\HL{kllte}{\B{kll\-te}}, \HL{kulat}{\B{ku\-lat}}) 

% kllkxem
\hi
\HT{kllkxem}{\B{kll\cp{kxem}}} \I{[k\s{l}."k'$\cdot$Em]} \emph{vin., inf.}\on{22} stand (\HL{mI}{\B{mì}}, in; \HL{sIn}{\B{sìn}}, on; \HL{IlA}{\B{ìlä}}, in shape\slash form of). 
(\emph{a. locally}) \B{Tsa\-ye\-rik kll\-kxem sìn yì a\-kxayl.} That \emph{yerik} is standing on a high ledge.~\LINK{http://naviteri.org/2013/01/lifyari-po-peyi-how-good-is-her-navi/}{b}
\B{Kll\-kxa\-yem fì\-tì\-kang\-kem oe\-yä ro\-fa--ke io--pum fe\-yä.} My work will stand beside--not above--theirs.~\LINK{http://naviteri.org/2010/06/first-post/}{b}
\B{oe\-yä ay\-syu\-lang set mì fay kll\-kxe\-rem}, my flowers are standing in water now.~\LINK{http://naviteri.org/2010/09/getting-to-know-you-part-1/comment-page-1/\#comment-296}{b}
\B{Na'\-vi ìlä ho\-'on kll\-kxo\-lem teng\-krr re\-rol.} The People were standing in a circle, singing.~\LINK{http://naviteri.org/2013/09/on-si-salewfya-shapes-and-directions/}{b}
(\emph{b. metaphorically with} \B{sìn yì} \emph{to describe anything scalable, e.g. water level, temperature, talent, anger, etc.}) \B{Lì'\-fya\-ri po kll\-kxem sìn pe\-yì?} How good is her Na'vi?~\LINK{http://naviteri.org/2013/01/lifyari-po-peyi-how-good-is-her-navi/}{b}
(\emph{c. idiom}) \B{Kll\-kxem ki\-ven!} `Stand up straight!' [lit.: ``stand straight and act with confidence'']~\LINK{https://naviteri.org/2025/11/kaltxi-nimun-hello-again/}{b}
(\ding{43}~\HL{kxem}{\B{kxem}})

% kllpa
\hi
\HT{kllpa}{\B{\cp{kll}pa}} \I{["k\s{l}.pa]} \emph{n.} bottom. 
(\ding{214}~\HL{fApa}{\B{fä\-pa}}, \ding{43}~\HL{kllte}{\B{kll\-te}}, \HL{pa'o}{\B{pa\-'o}}) 

% kllpä 
\hi
\HT{kllpA}{\B{kll\cp{pä}}} \I{[k\s{l}."p$\cdot$\ae]} \emph{vin., inf.}\on{22} land, arrive at the ground. 
\B{Maw sä\-tsway\-on a\-yol ay\-oe kll\-po\-lä mì ta\-yo a lu ro\-fa kil\-van.} After a short flight we landed in a field beside the river.~\LINK{http://naviteri.org/2012/01/mipa-zisit-ayliu-amip-new-words-for-the-new-year/}{b}
(\ding{43}~\B{kllwo}, \HL{kllte}{\B{kll\-te}}, \HL{pAhem}{\B{pä\-hem}})

% kllpxiwll
\hi
\HT{kllpxiwll}{\B{kll\cp{pxi}wll}} \I{[k\s{l}."p'i.w\s{l}]} \emph{n.}~\ding{96} lionberry, \emph{cynaroidia decumbens}.
(\ding{43}~\HL{kllte}{\B{kll\-te}}, \HL{pxi}{\B{pxi}}, \HL{'ewll}{\B{'e\-wll}})

% kllpxìltu
\hi
\HT{kllpxIltu}{\B{kll\cp{pxìl}tu}} \I{[k\s{l}."p'Il.tu]} \emph{n.} territory. 
\B{ul\-te kop su\-teo a\-la\-he a tok fra\-kll\-pxìl\-tut ki\-fkey\-ä} and (there are) also some others who are in all territories of the world.~\LINK{http://forum.learnnavi.org/language-updates/ke-kawit/msg174572/\#msg174572}{ln}

% kllrikx
\hi
\HT{kllrikx}{\B{kll\cp{rikx}}} \I{[k\s{l}."Rik']} \emph{n.} earthquake. 
\B{Txe\-wì pll\-txe san kll\-rikx txewm la\-mu sìk.} Txewì says that the earthquake was frightening.~\LINK{http://naviteri.org/2015/11/vomuna-liu-amip-ten-new-words/}{b}
(\ding{43}~\HL{kllte}{\B{kll\-te}}, \HL{rikx}{\B{rikx}})

% kllte
\hi
\HT{kllte}{\B{\cp{kll}te}} \I{["k\s{l}.tE]} \emph{n.} ground.  
\B{Kll\-te lu ekx\-txu.} The ground is rough.~\LINK{http://naviteri.org/2013/04/wheres-the-bathroom-and-other-useful-things/}{b}
\B{Fwa ya\-wo ftu kll\-te to fwa tsway\-on ftu 'awkx lu ngä\-zìk.} Taking off from the ground is harder than flying off a cliff.~\LINK{http://naviteri.org/2012/01/mipa-zisit-ayliu-amip-new-words-for-the-new-year/}{b}
(\emph{a. idiom}) \B{ne kllte}, get down! \\
\textsc{Derivatives}: \ding{43} \on{1}.~\HL{nekll}{\B{ne\-kll}}, see below.
\on{2}.~\HL{kllrikx}{\B{kll\-rikx}}, earthquake.
\on{3}.~\HL{kllkA}{\B{kll\-kä}}, go down.
\on{4}.~\HL{kllkulat}{\B{kll\-ku\-lat}}, dig up.
\on{5}.~\HL{kllkxem}{\B{kll\-kxem}}, stand.
\on{6}.~\HL{kllpa}{\B{kll\-pa}}, bottom.
\on{7}.~\HL{kllpA}{\B{kll\-pä}}, land, arrive at the ground.
\on{8}.~\HL{kllpxiwll}{\B{kll\-pxi\-wll}}, lionberry.
\on{9}.~\HL{klltxay}{\B{kll\-txay}}, lie on the ground.
\on{10}.~\HL{klltseng}{\B{kll\-tseng}}, position.
\on{11}.~\HL{kllvawm}{\B{kll\-vawm}}, brown.
\on{12}.~\HL{kllwo}{\B{kll\-wo}}, alight, land.
\on{13}.~\HL{kllyem}{\B{kll\-yem}}, bury.
\on{14}.~\HL{kllza'au}{\B{kll\-za\-'u}}, come down, descend.

% klltxay
\hi
\HT{klltxay}{\B{kll\cp{txay}}} \I{[k\s{l}."t'$\cdot$aj]} \emph{vin., inf.}\on{22} lie on the ground.
(\ding{43}~\HL{kllte}{\B{kll\-te}}, \HL{txay}{\B{txay}}) \\
\textsc{Derivatives}: \ding{43} \on{1}.~\B{kll\-txey\-kay}, lay (s.t.) on the ground.

% klltxeykay
\hi
\HT{klltxeykay}{\B{klltxey\cp{kay}}} \I{[k\s{l}.t'Ej."k$\cdot$aj]} \emph{vtr., inf.}\on{33} lay (s.t.) on the ground. 
(\ding{43}~\HL{klltxay}{\B{kll\-txay}})

% klltseng
\hi
\HT{klltseng}{\B{\cp{kll}tseng}} \I{["k\s{l}.\t{ts}EN]} \emph{n.} position. 
\B{Sra\-ne, tsun fko si\-var hem\-lì\-'u\-vit mì lì\-'u alu swey\-lu, ul\-te sra\-ne, kll\-tseng lu \on{2},\on{2}.} Yes, you can use infixes in the word \emph{sweylu}, and (their) position is \on{2},\on{2}.~\LINK{http://naviteri.org/2011/04/\%E2\%80\%99a\%E2\%80\%99awa-li\%E2\%80\%99fyavi-amip\%E2\%80\%94a-few-new-expressions/comment-page-1/\#comment-677}{b}
(\ding{43}~\HL{kllte}{\B{kll\-te}}, \HL{tsenge}{\B{tse\-nge}})

% kllvawm
\hi
\HT{kllvawm}{\B{\cp{kll}vawm}} \I{["k\s{l}.vawm]} \emph{adj.} brown. 
\B{Tsa\-tu\-tan\-ur a fkol tspo\-lang lo\-lu ta'\-leng a\-kll\-vawm; tspang\-yur lu pum a\-teyr.} That person that was killed had brown skin, the killer had white.~\LINK{http://naviteri.org/2020/06/mi-tanlokxe-oeya-srr-afpxamo-terrible-days-in-my-country/}{b}
(\ding{43}~\HL{kllte}{\B{kll\-te}}, \HL{vawm}{\B{vawm}}) \\
\textsc{Derivatives}: \ding{43} \on{1}.~\HL{kllvawmpin}{\B{kll\-vawm\-pin}}, (the) brown (color).

% kllvawmpin
\hi
\HT{kllvawmpin}{\B{\cp{kll}vawmpin}} \I{["k\s{l}.vawm.pin]} \emph{n.} the color \ding{43}~\HL{kllvawm}{\B{kll\-vawm}}.
(\ding{43}~\HL{'opin}{\B{'o\-pin}})

% kllwo 
\hi
\HT{kllwo}{\B{kll\cp{wo}}} \I{[k\s{l}."w$\cdot$o]} \emph{vin., inf.}\on{22} alight, land (\emph{process}), reach for the ground.  
\B{Tom\-pa 'e\-ko nì\-hawng, ha ze\-ne aw\-nga kll\-wi\-vo.} The rain is too strong, so we must land.~\LINK{http://naviteri.org/2012/01/mipa-zisit-ayliu-amip-new-words-for-the-new-year/}{b}
(\ding{214}~\HL{yawo}{\B{ya\-wo}}, \ding{43}~\HL{kllte}{\B{kll\-te}}, \HL{wo}{\B{wo}}, \HL{kllpA}{\B{kll\-pä}}) 

% kllyem
\hi
\HT{kllyem}{\B{kll\cp{yem}}} \I{[k\s{l}."j$\cdot$Em]} \emph{vtr., inf.}\on{22} bury. 
\B{Po\-ti kll\-yo\-lem ay\-oel äo ut\-ral\-o a\-lor a ro\-fa kil\-van.} We buried her under a beautiful tree beside the river.~\LINK{http://naviteri.org/2012/01/more-additions-to-the-lexicon/}{b}
(\ding{214}~\HL{kllkulat}{\B{kll\-ku\-lat}}, \ding{43}~\HL{kllte}{\B{kll\-te}}, \HL{yem}{\B{yem}}) 

% kllza'u
\hi
\HT{kllza'u}{\B{kll\cp{za}'u}} \I{[k\s{l}."z$\cdot$a.P$\cdot$u]} \emph{vin., inf.}\on{23} come down, descend. 
\B{Kll\-zo\-la\-'u Mo\-'at fa sna\-yì teng\-krr pe\-rll\-txe san Ay\-nga ne\-to ri\-vikx!} Mo'at came down the stairway saying, ``Get back, all of you!''~\LINK{http://naviteri.org/2013/01/lifyari-po-peyi-how-good-is-her-navi/}{b}
(\ding{214}~\HL{fAza'u}{\B{fä\-za\-'u}}, \ding{43}~\HL{kllte}{\B{kll\-te}}, \HL{za'u}{\B{za\-'u}}) 

% ko 
\hi
\HT{ko}{\B{ko}} \I{[ko]} \emph{sentence\hyp{}final part. for solicit agreement.} (\emph{a.}) ``Don't you think so?'' ``Wouldn't you agree?'' ``Let's \dots'' ``Why don't you \dots'' ``I'll do X, OK?''
\B{Aw\-nga sngi\-vä\-'i ko!} Let us begin!~\LINK{http://naviteri.org/2010/06/first-post/}{b}
\B{Var ni\-vu\-me ko!} Let's keep learning!~\LINK{http://forum.learnnavi.org/language-updates/as-x-as-y-keep-on-keepin-on/}{ln}
\B{Za\-'u kal\-txì si ko!} Come and say hello!~\LINK{http://naviteri.org/2010/09/getting-to-know-you-part-3/}{b}
\B{Teng\-krr ftxo\-zä se\-re\-i\-yi aw\-nga, ke tswi\-va’ ay\-lom\-tu\-ti ko!} A toast:~To absent friends. [lit.: While we are celebrating, let us not forget those who we wish could also be here]~\LINK{http://naviteri.org/2011/07/txantsana-ultxa-mi-siatll-great-meeting-in-seattle/}{b}
\B{Ul\-txa si\-vi oeng sìn ram\-tsyìp txon\-'ong\-ay.--Ye\-i\-o'!~Tsa\-krr\-vay ko!} Let's meet on the hill tomorrow at nightfall.--Perfect!~See you then.~\LINK{http://naviteri.org/2012/01/mipa-zisit-ayliu-amip-new-words-for-the-new-year/}{b}
\B{Nga lä\-pi\-vawk nì\-'it nì\-'ul ko.} Tell me a bit more about yourself.~\LINK{http://naviteri.org/2010/09/getting-to-know-you-part-2/}{b}
\B{Lu fo le\-hrr\-ap.~Tsun tu\-tet tspi\-vang ko.} Those things are dangerous. They can kill a person, you know.~\LINK{http://wiki.learnnavi.org/index.php?title=Corpus\#New_York_Times}{wiki}
\B{Fnu, ma 'e\-vi.~Sa'\-nur le\-ru haw\-tsyìp.~Tsi\-vu\-rokx ko.}
Quiet, young one.~Mommy is taking a nap.~Let her rest.~\LINK{http://naviteri.org/2011/07/txantsana-ultxa-mi-siatll-great-meeting-in-seattle/}{b}
(\emph{b. idiom}) \B{Li ko.} Well, get to it, then.\slash Let's get on it.~\LINK{http://naviteri.org/2011/02/new-vocabulary-part-2/}{b}
\B{May' ko.} Give it a try!\slash Have a go!~\LINK{http://naviteri.org/2011/05/some-miscellaneous-vocabulary/}{b}

% ko'on
\hi
\HT{ko'on}{\B{\cp{ko}'on}} \I{["ko.Pon]} \emph{n.} ring, oval, closed shaped (\emph{roughly circular}). 
\B{Na'\-vi ìlä ho\-'on kll\-kxo\-lem teng\-krr re\-rol.} The People were standing in a circle, singing.~\LINK{http://naviteri.org/2013/09/on-si-salewfya-shapes-and-directions/}{b}
(\ding{43}~\HL{koum}{\B{ko\-um}}, \HL{'on}{\B{'on}}; \HL{yo'ko}{\B{yo'\-ko}}) \\
\textsc{Derivatives}: \ding{43} \on{1}.~\HL{ko'onspul}{\B{ko\-'on\-spul}}, sunflower gigantus.

% ko'onspul
\hi 
\HT{ko'onspul}{\B{\cp{ko}'onspul}} \I{["ko.Pon.spul]} \emph{or} \I{[.spUl]} \emph{n.} \ding{96} sunflower gigantus. 
(\ding{43}~\HL{ko'on}{\B{ko'on}}, \HL{spule}{\B{spule}})

% koak
\hi
\HT{koak}{\B{\cp{ko}ak}} \I{["ko.ak\textcorner]} \emph{adj.} old, aged (\emph{for living things}). 
\B{Ko\-ak nì\-ngay, ke\-fyak?} Really old, right?~\LINK{http://forum.learnnavi.org/news-announcements/happy-birthday-to-karyu-pawl/msg556899/\#msg556899}{ln}
(\ding{214}~\HL{'ewan}{\B{'e\-wan}}, \ding{43}~\HL{spuwin}{\B{spu\-win}}) \\
\textsc{Derivatives}: \ding{43} \on{1}.~\HL{koaktu}{\B{ko\-ak\-tu}}, old person. 
\on{2}.~\HL{koaktan}{\B{ko\-ak\-tan}}, old man.   
\on{3}.~\HL{koakte}{\B{ko\-ak\-te}}, old woman. 
\on{4}.~\HL{kewan}{\B{ke\-wan}}, age.
\on{5}.~\HL{tIkoak}{\B{tì\-ko\-ak}}, old age.

% koaktan
\hi
\HT{koaktan}{\B{\cp{ko}aktan}} \I{["ko.ak\textcorner.tan]} \emph{n.} old man.  
\B{Ay\-sre' lom lu tsa\-ko\-ak\-tan\-ur.} The old man misses his teeth.~\LINK{http://naviteri.org/2011/07/txantsana-ultxa-mi-siatll-great-meeting-in-seattle/}{b}
\B{Ko\-ak\-tan\-ä ay\-sre' lä\-ngu fwem.} An old man's teeth are dull.~\LINK{http://naviteri.org/2012/03/spring-vocabulary-part-1/}{b}
(\ding{43}~\HL{koak}{\B{ko\-ak}})

% koakte
\hi
\HT{koakte}{\B{\cp{ko}akte}} \I{["ko.ak\textcorner.tE]} \emph{n.} old woman.
\B{Po\-leng ay\-oe\-ru ko\-ak\-tel vur\-it a\-txan\-lal.} The old woman told us an ancient story.~\LINK{http://naviteri.org/2015/08/aylifyavi-lereyfya-2-cultural-terms-2-and-more/}{b}    
\B{Tsa\-ko\-ak\-te\-ri stawm\-tswo lu fe'.} That old woman's hearing is poor.~\LINK{http://naviteri.org/2012/11/renu-ayinanfyaya-the-senses-paradigm/}{b}
(\ding{43}~\HL{koak}{\B{ko\-ak}})

% koaktu
\hi
\HT{koaktu}{\B{\cp{ko}aktu}} \I{["ko.ak\textcorner.tu]} \emph{n.} old person.
\B{Ul\-te txo smi\-von ngar ay\-ho\-ak\-tu fo\-ti pa\-lang fte tsi\-vun ivo\-mum teyng\-ta ftxey lu fo\-ru fpom fu\-ke.} And if you know elders contact them to know whether they are alright or not.~\LINK{http://naviteri.org/2020/03/teri-tifkeytok-lefkrr-about-the-present-situation/}{b}
\B{Ta'\-leng prr\-nen\-ä lu fa\-o\-i, pum ko\-ak\-tu\-ä ekx\-txu.} A baby's skin is smooth, an old person's is rough.~\LINK{http://naviteri.org/2012/01/mipa-zisit-ayliu-amip-new-words-for-the-new-year/}{b}
(\emph{a. proverb}) \B{Koakturi kewanti keyìl ke wan.} Some things can't be covered up. [lit.: `An old person's face doesn't hide their age.']~\LINK{http://naviteri.org/2016/06/mrrvola-lifyavi-amip-forty-new-expressions/}{b}
(\ding{43}~\HL{koak}{\B{ko\-ak}}, \HL{koaktan}{\B{ko\-ak\-tan}}, \HL{koakte}{\B{ko\-ak\-te}}) \\
\textsc{Derivatives}: \ding{43} \on{1}.~\HL{koaktutral}{\B{ko\-ak\-tut\-ral}}, goblin thistle.

% koaktutral
\hi 
\HT{koaktutral}{\B{ko\cp{ak}tutral}} \I{[ko."{}ak\textcorner.tut\textcorner.Ral]} \emph{or} \I{[.tUt\textcorner.]} \emph{n.} \ding{96} goblin thistle, \emph{cobalus carduus}.~\LINK{https://forum.learnnavi.org/language-updates/expansion-to-flora-fauna-and-places-from-the-valley-of-mo'ara/}{ln}  
(\ding{43}~\HL{koaktu}{\B{ko\-ak\-tu}}, \HL{utral}{\B{ut\-ral}})

% kolan
\hi
\HT{kolan}{\B{ko\cp{lan}}} \I{[ko."lan]} \emph{conv.} ``I mean \ldots, rather''.
(\ding{43}~\HL{kan}{\B{kan}})

% kom
\hi 
\HT{kom}{\B{kom}} \I{[kom]} \emph{vin., modal} dare. 
\B{Oe ke kom ki\-vä.} I don't dare to go.
\B{Nga kom pi\-vll\-txe oe\-hu tsa\-fya srak?} You dare to speak to me like that?~\LINK{http://naviteri.org/2019/06/50a-liu-amip-40-new-words/}{b}

% komum
\hi
\HT{komum}{\B{ko\cp{mum}}} \I{[ko."mum]} \emph{or} \I{[."mUm]} (\textsc{rn}: \B{ko\cp{mùm}} \I{[ko."mUm]})~:~\HL{omum}{\B{omum}}

% kop
\hi
\HT{kop}{\B{kop}} \I{[kop\textcorner]} \textsc{i}. \emph{adv.} also, too, as well, additionally. 
\B{Tsa\-kem kop krr\-nekx.} That also takes time.~\LINK{http://forum.learnnavi.org/language-updates/just-a-few-interesting-little-words/}{ln}
\B{Kop oe\-ru lo\-lä\-ngu lie a ha\-pxì\-tu soaiä ter\-kup.} Sadly, I too have had the experience of a family member dying.~\LINK{http://naviteri.org/2013/01/awvea-posti-zisita-amip-first-post-of-the-new-year/}{b}
(\emph{a. idiom}) \B{Kxe\-tse sì mik\-yun kop pll\-txe.} Body language is important [lit.: The tail and ears also speak].~\LINK{http://forum.learnnavi.org/language-updates/grid-tidbits/}{ln}  \\
\textsc{ii}. \emph{conj.} and also.
\B{Pe\-ioang tsun mì tam\-pay kop mì reym ri\-vey?} What animal can live both in the sea and on the land?~\LINK{http://naviteri.org/2015/04/some-new-words-for-may-day/}{b} \\
\textsc{Derivatives}: \ding{43} \on{1}.~\HL{slAkop}{\B{slä\-kop}}, but also.

% koren
\hi
\HT{koren}{\B{ko\cp{ren}}} \I{[ko."REn]} \emph{n.} rule, guideline.  
\B{Ko\-ren A\-'aw\-ve Tì\-ru\-sey\-ä 'Aw\-si\-teng}, the First Rule of Living Together.~\LINK{http://wiki.learnnavi.org/index.php?title=Corpus\#Good_Morning_America}{wiki}
\B{Oe\-ri tsa\-ko\-ren\-it pxìm tswä\-nga'.} I often forget that rule.~\LINK{http://naviteri.org/2011/09/\%E2\%80\%9Cby-the-way-what-are-you-reading\%E2\%80\%9D/comment-page-1/\#comment-1092}{b}
\B{za\-mi\-vu\-nge oel ay\-ngar ay\-lì\-'ut ho\-ren\-ti\-sì lì'\-fya\-yä le\-Na'\-vi}, I will bring you the words and rules of the Na'vi language.~\LINK{http://forum.learnnavi.org/news-announcements/a-response-from-paul-frommer!/}{ln}
\B{ke tsun fko ni\-vu\-me fte nì\-lì'\-fya a\-mip pi\-vll\-txe nìl\-tsan fa fwa tìng mik\-yun ho\-ren\-ur nì\-'aw}, you can't learn to speak a new language well just by listening rules.~\LINK{http://naviteri.org/2012/07/tskxekengtsyip-a-mikyunfpi-a-little-listening-exercise/}{b} \\
(i.) \B{koren ayll} law, societal rule.
\B{Pol aw\-nga\-ru ho\-ren\-it a\-yll ley\-kek.} He makes us obey the laws.~\LINK{http://naviteri.org/2020/06/mi-tanlokxe-oeya-srr-afpxamo-terrible-days-in-my-country/}{b} \\
\textsc{Derivatives}: \ding{43} \on{1}.~\HL{lekoren}{\B{le\-ko\-ren}}, lawlike, regarding rules. 
\on{2}.~\HL{horenleykekyu}{\B{ho\-ren\-ley\-kek\-yu}}, law enforcer.

% kosman
\hi
\HT{kosman}{\B{ko\cp{sman}}} \I{[ko."{}sman]} \emph{adj.} wonderful, terrific, fantastic.
\B{Fol lì'\-fya\-ti aw\-nge\-yä sar a fya\-'o lu ko\-sman.} It's wonderful the way they use our language.~\LINK{http://naviteri.org/2011/09/miscellaneous-vocabulary/}{b} 
\B{Nge\-yä sä\-mok\-ìri a\-ko\-sman se\-i\-yi oe ira\-yo.} Thanks for that excellent suggestion (of yours).~\LINK{http://naviteri.org/2012/01/more-additions-to-the-lexicon/}{b} \\
\textsc{Derivatives}: \ding{43} \on{1}.~\HL{nIksman}{\B{nìk\-sman}}, wonderfully.

% kot
\hi
\HT{kot}{\B{kot}} \I{[kot\textcorner]} \emph{n.}, \textsc{lw} coat.

% koum
\hi
\HT{koum}{\B{\cp{ko}um}} \I{["ko.um]} \emph{or} \I{[.Um]} (\textsc{rn}: \B{\cp{ko}ùm} \I{["ko.Um]}) \emph{adj.} rounded, curved.
\B{Fì\-tskxe\-ri fa\-'o lu yey; ke lu ko\-um.} This rock has straight sides; it's not rounded.~\LINK{http://naviteri.org/2013/09/on-si-salewfya-shapes-and-directions/}{b} \\
\textsc{Derivatives}: \ding{43} \on{1}.~\HL{ko'on}{\B{ko\-'on}}, ring, oval, closed shape roughly circular. 
\on{2}.~\HL{yo'ko}{\B{yo'\-ko}}, circle, perfectly circular ring.

% krr
\hi
\HT{krr}{\B{krr}} \I{[k\s{r}]} \emph{uncountable n.} time. 
\B{maw\-krr la\-ye\-i\-u oer krr nì\-'ul}, after that I will have more time.~\LINK{http://forum.learnnavi.org/language-updates/trr-rrtaya/}{ln}
\B{Ke lu oer set krr a\-txan.} I don't have much time now.~\LINK{http://forum.learnnavi.org/language-updates/verb-for-write/msg139677/\#msg139677}{ln} 
\B{Nga\-ri txo fì\-kem si\-vi fì\-fya, krr\-ti 'a\-vun.} If you do it like this, you'll save time.~\LINK{http://naviteri.org/2021/12/zolau-niprrte-ma-3746-welcome-2022/}{b} 
\B{Slä mì tam\-pxì krr\-ä, ke sar fkol kea hem\-lì\-'u\-vit fì\-lì\-'u\-mì.} But most of the time, one doesn't use any infixes in this word.~\LINK{http://naviteri.org/2011/04/\%E2\%80\%99a\%E2\%80\%99awa-li\%E2\%80\%99fyavi-amip\%E2\%80\%94a-few-new-expressions/comment-page-1/\#comment-677}{b} 
(\emph{a. idiom, poetic}) \B{tì\-'i'\-a\-vay krr\-ä}, forever, 'til the end of time. 
\B{To\-ler\-kup tu\-te; pam\-tseo pe\-yä tì\-'i'\-a\-vay krr\-ä ra\-yey.} The man has died; his music will live forever.~\LINK{http://naviteri.org/2015/06/na-ayskxe-mi-telan-sad-news/}{b} \\
\textsc{Derivatives}: \ding{43} \on{1}.~\HL{frakrr}{\B{fra\-krr}}, always. 
\on{2}.~\HL{ftawnemkrr}{\B{ftaw\-nem\-krr}}, past. 
\on{3}.~\HL{hIkrr}{\B{hì\-krr}}, short time, second. 
\on{4}.~\HL{kawkrr}{\B{kaw\-krr}}, never. 
\on{5}.~\HL{hawngkrr}{\B{hawng\-krr}}, late. 
\on{6}.~\HL{krra}{\B{krra}}\slash \HL{a krr}{\B{a krr}}, when, as. 
\on{7}.~\HL{krrnekx}{\B{krr\-nekx}}, take\slash consume time. 
\on{8}.~\HL{krro}{\B{krro}}, sometimes. 
\on{9}.~\HL{krro krro}{\B{krro krro}}, at times, on occasion. 
\on{10}.~\HL{lefkrr}{\B{le\-fkrr}}, current. 
\on{11}.~\HL{mawkrr}{\B{maw\-krr}}, after. 
\on{12}.~\HL{nIfkrr}{\B{nì\-fkrr}}, currently, lately. 
\on{13}.~\HL{nulkrr}{\B{nul\-krr}}, longer. 
\on{14}.~\HL{sekrr}{\B{se\-krr}}, present, now. 
\on{15}.~\HL{sngA'ikrr}{\B{sngä\-'i\-krr}}, beginning. 
\on{16}.~\HL{srekrr}{\B{sre\-krr}}, before. 
\on{17}.~\HL{tengkrr}{\B{teng\-krr}}, while, at the same time. 
\on{18}.~\HL{takrra}{\B{ta\-krra\slash akrr\-ta}}, since. 
\on{19}.~\HL{tsakrr}{\B{tsakrr}}, then, at that time. 
\on{20}.~\HL{txankrr}{\B{txan\-krr}}, (for a) long time. 
\on{21}.~\HL{txonkrr}{\B{txon\-krr}}, at night. 
\on{22}.~\HL{vaykrr}{\B{vay\-krr}}, until. 
\on{23}.~\HL{ye'krr}{\B{ye'\-krr}}, early. 
\on{24}.~\HL{zIsIkrr}{\B{zì\-sì\-krr}}, season. 
\on{25}.~\HL{zusawkrr}{\B{zu\-saw\-krr}}, future.  
\on{26}.~\HL{krrka}{\B{krr\-ka}}, during.
\on{27}.~\HL{txakrrfpIl}{\B{txa\-krr\-fpìl}}, consider, ponder.
\on{28}.~\HL{krrlik}{\B{krr\-lik}}, point in time.

% krra
\hi
\HT{krra}{\B{\cp{krr}a}}, \B{krr a} \emph{or} \HL{a krr}{\B{a krr}} \I{["k\s{r}.a]} \emph{conj.} when, at the time that. 
\B{Oel tsko\-ti nga\-ru ta\-syìng krra oeng ul\-txa si.} I'll give you the bow when we meet.~\LINK{http://naviteri.org/2012/03/spring-vocabulary-part-2/}{b}
\B{pam\-rel sa\-yi trr\-ay krr a skxom la\-tsu oe\-ru}, I'll write tomorrow when I have the chance.~\LINK{http://forum.learnnavi.org/language-updates/just-a-few-interesting-little-words/}{ln}
\B{Tì\-'eyng\-it oel to\-lel a krr, ay\-nga\-ru pa\-yeng.} When I've received an answer, I will tell you.~\LINK{http://forum.learnnavi.org/news-announcements/a-response-from-paul-frommer!/}{ln}
\B{Oel fo\-ru fì\-ay\-lì\-'ut to\-lìng a krr, kxawm oe har\-ma\-hä\-ngaw.} When I gave them these words I was probably sleeping.~\LINK{http://wiki.learnnavi.org/index.php?title=Canon\#Extracts_from_various_emails}{wiki} 
(\ding{43}~\HL{krr}{\B{krr}})

% krrka
\hi
\HT{krrka}{\B{\cp{krr}ka}} \I{["k\s{r}.ka]} \emph{adp.} during.
\B{Krr\-ka tsawl\-ul\-txa vo\-spxì\-am.} During the great meeting last month.~\LINK{http://naviteri.org/2014/08/stxeli-alor-a-beautiful-gift/}{b}
\B{Kaym\-krr\-ka to\-lel mo\-el ta ay\-ul\-txa\-tu stxe\-lit a\-ko\-sman.} During the evening we received a wonderful gift from the meeting members.~\LINK{http://naviteri.org/2014/08/stxeli-alor-a-beautiful-gift/}{b}
(\ding{43}~\HL{krr}{\B{krr}}, \HL{ka}{\B{ka}})

% krrlik
\hi 
\HT{krrlik}{\B{\cp{krr}lik}} \I{["k\s{r}.lik\textcorner]} \textsc{i}. \emph{n.} point in time. 
\B{Krr\-lik lu trr\-pxì a\-mrr\-ve.} The time is five o'clock.~\LINK{https://naviteri.org/2024/05/krrteri-about-time/}{b} \\
\textsc{ii}. \B{\cp{krr}likpe} \I{["k\s{r}.lik\textcorner.pE]} \emph{or} \B{pe\cp{hrr}lik} \I{[pE."h\s{r}.lik\textcorner]} \emph{inter.} what point in time? what hour?
(\ding{43}~\HL{krr}{\B{krr}}, \HL{lik}{\B{lik}}) 

% krro
\hi
\HT{krro}{\B{\cp{krr}o}} \I{["k\s{r}.o]} \emph{adv.} sometimes.
\B{Krro len ay\-hem a\-na\-fì\-'u.} Sometimes such things happen.~\LINK{http://forum.learnnavi.org/language-updates/just-another-few-words/}{ln}
\B{Krro lu 'Ìng\-lì\-sì txur nì\-hawng mì re\-'o.} Sometimes English is too strong in my head.~\LINK{http://naviteri.org/2011/09/\%E2\%80\%9Cby-the-way-what-are-you-reading\%E2\%80\%9D/comment-page-1/\#comment-1092}{b}
\B{Ngay\-txoa, krro tì\-kxey si keng kar\-yu.} I'm sorry, sometimes even the teacher messes up.~\LINK{http://naviteri.org/2012/02/trr-asawnung-lefpom-happy-leap-day/}{b}
(\ding{43}~\HL{krr}{\B{krr}})

% krro krro
\hi
\HT{krro krro}{\B{\cp{krr}o \cp{krr}o}} \I{["k\s{r}.o "k\s{r}.o]} \emph{adv.} at times, on occasion.
\B{Krro krro fì\-tse\-nge oe tì\-kxey sa\-yi.} At times, I will make mistakes here.~\LINK{http://naviteri.org/2010/06/first-post/}{b}
\B{Krro krro, flì\-a vul a\-ru\-sey to nutx\-a pum a\-ke\-ru\-sey lu txur.} A thin living branch is sometimes stronger than a thick dead one.~\LINK{http://naviteri.org/2011/11/\%E2\%80\%99a\%E2\%80\%99awa-ayli\%E2\%80\%99u-amip-ni\%E2\%80\%99aw-\%E2\%80\%94-a-few-new-words-only/}{b}
(\ding{43}~\HL{krr}{\B{krr}})

% krrnekx
\hi
\HT{krrnekx}{\B{krr\cp{nekx}}} \I{[k\s{r}."n$\cdot$Ek']} \emph{vin., inf.}\on{22} take\slash consume time. 
\B{Tsa\-kem kop krr\-nekx.} That also takes time.~\LINK{http://forum.learnnavi.org/language-updates/just-a-few-interesting-little-words/}{ln}
\B{fwa oel fì\-tì\-kang\-kem\-vit sngey\-ki\-vä\-'i krr\-no\-lekx nì\-txan}, me starting this project took a long time.~\LINK{http://naviteri.org/2010/06/first-post/}{b}
(\ding{43}~\HL{krr}{\B{krr}}, \HL{nekx}{\B{nekx}})

% krrpe
\hi
\HT{krrpe}{\B{\cp{krr}pe}} \emph{or} \HL{pehrr}{\B{pe\cp{hrr}}} \I{["k\s{r}.pE]} \emph{or} \I{[pE."h\s{r}]} \emph{inter.} when?
\B{Fä\-za\-'u tsaw\-ke krr\-pe?} When will the sun come up?~\LINK{http://naviteri.org/2012/01/more-additions-to-the-lexicon/}{b}
\B{Nga\-ri sno\-mot krr\-pe ve\-zey\-ko, ma 'itan?} When are you going to organize your room, son?~\LINK{http://naviteri.org/2012/10/mipa-vospxi-mipa-ayliu-new-words-for-the-new-month/}{b}
(\ding{43}~\HL{krr}{\B{krr}})

% ku'up
\hi
\HT{ku'up}{\B{\cp{ku'}up}} \I{["kuP.up\textcorner]} \emph{or} \I{["kUP.Up\textcorner]} (\textsc{rn}: \B{\cp{ku}ùp} \I{["ku.Up\textcorner]}) \emph{adj.} heavy (\emph{physical weight}). 
\B{Ay\-skxe a mì sa\-sna\-'o ku'\-up lu nì\-txan.} The rocks in that pile are very heavy.~\LINK{http://naviteri.org/2012/03/spring-vocabulary-part-2/}{b}
\B{Ke tsun tu\-te a\-'aw tsa\-tskxe\-ti a\-ku'\-up kxi\-vel\-tek nì\-'aw\-tu.} One person alone can't lift that heavy rock.~\LINK{http://naviteri.org/2012/01/mipa-zisit-ayliu-amip-new-words-for-the-new-year/}{b}
(\emph{a. idiom}) \B{Ngä\-zìk fwa fkol rum\-it a\-ku\-'up pey\-ku\-wup.} \emph{it's difficult to make s.o. stubborn or inept do what you want them to} [lit.: it's hard to bounce a heavy ball].~\LINK{http://naviteri.org/2023/12/mipa-zisit-ayliu-amip-new-year-new-words/}{b}
(\ding{214}~\HL{syo}{\B{syo}}) \\
\textsc{Derivatives}: \ding{43} \on{1}.~\HL{syokup}{\B{syo\-kup}}, weight.

% kulat
\hi
\HT{kulat}{\B{\cp{ku}lat}} \I{["ku.lat\textcorner]} \emph{vtr.} reveal, bring forth, uncover (\emph{literally and metaphorically}). 
\B{Maw txan\-tom\-pa, pxay\-a rìk\-äo la\-mu tska\-lep pe\-yä, ha tsat ku\-lat ay\-oel.} After the rainstorm, his crossbow was under a lot of leaves, so we uncovered it.~\LINK{http://naviteri.org/2011/10/more-vocabulary-a-bit-of-grammar/}{b}
\B{Lo\-lu ka\-vuk, slä Tse\-nul tì\-ngay\-it ko\-lu\-lat.} There was treachery, but Tsenu revealed the truth.~\LINK{http://naviteri.org/2011/10/more-vocabulary-a-bit-of-grammar/}{b} \\
\textsc{Derivatives}: \ding{43} \on{1}.~\HL{kllkulat}{\B{kll\-ku\-lat}}, dig out.

% kum
\hi
\HT{kum}{\B{kum}} \I{[kum]}~\emph{or}~\I{[kUm]} \emph{n.} result.
\B{Po\-ri tì\-kak\-rel tì\-kak\-pam\-sì kum tsam\-ä lu.} His blindness and deafness are a result of the war.~\LINK{http://naviteri.org/2014/05/mipa-ayliu-mipa-aysafpil-new-words-new-ideas/}{b}
\B{Mllte oe: am'aluke zamayunge tsatìpuseyìl humit apxan.} I agree: without a doubt the waiting will bring worthy results.~\LINK{http://naviteri.org/2013/08/fmawn-a-fkol-kilmulat-this-just-in/comment-page-1/\#comment-2478}{b}
(\emph{a. proverb}) \B{Kem a\-muiä, kum a\-fe'.} Proper action, bad result. (\emph{something that should have turned out well didn't}).~\LINK{http://naviteri.org/2012/06/spring-vocabulary-part-3/}{b}
\B{Hem a\-sru\-nga' nì\-'ul, hum a\-sngu\-nga' nì\-nän.}
More helpful actions lead to less troubling outcomes.~\LINK{http://naviteri.org/2015/03/kap-si-ayunil-saylahe-threats-dreams-and-other-things/}{b} \\
\textsc{Derivatives}: \ding{43} \on{1}.~\HL{kuma}{\B{ku\-ma\slash a\-kum}}, that (\emph{as a result}). 
\on{2}.~\HL{tsankum}{\B{tsan\-kum}}, advantage, benefit, upside, gain. 
\on{3}.~\HL{fekum}{\B{fe\-kum}}, disadvantage, drawback, downside.

% kuma
\hi
\HT{kuma}{\B{\cp{kum}a}} \emph{or} \B{a\cp{kum}} \I{["kum.a]} \emph{or} \I{[a."kum]} \emph{conj. with} \HL{fItxan}{\B{fì\-txan}} \emph{or} \HL{nIftxan}{\B{nì\-ftxan}} \emph{to form result clauses.} that (\emph{as a result}). 
\B{Lu poe se\-vin nì\-ftxan kum\-a yaw\-ne slo\-lu oer.}~\emph{or} \B{Poe yaw\-ne slo\-lu oer a\-kum, nì\-ftxan lu se\-vin.} She was so beautiful that I fell in love with her.~\LINK{http://naviteri.org/2012/06/spring-vocabulary-part-3/}{b}
\B{Lu oer ing\-yen a Ìstaw nim lu fì\-txan kum\-a pxìm wä\-pan.} It's a mystery to me why Ìstaw is so shy that he frequently hides.~\LINK{http://naviteri.org/2013/01/awvea-posti-zisita-amip-first-post-of-the-new-year/}{b}
(\emph{without} \B{fì\-txan} \emph{or} \B{nì\-ftxan}) \B{Pxe\-fo\-ru oe srung so\-li, kum\-a oe\-ru set pxe\-fo srung se\-ri.} I helped the three of them, so they're now helping me.~\LINK{http://naviteri.org/2012/06/spring-vocabulary-part-3/}{b}
(\emph{a. coll. idiom}) \B{[verb] nì\-ftxan kum\-a ter\-kup}, \emph{commenting on an action that is so intense that it results in dying}; ``like crazy, \ldots{} to death''. 
\B{Sran, sran, oe\-ri pì\-tìk tsa\-tse\-nget a mì tal! Fkxa\-ke nì\-ftxan kum\-a ter\-kup.} Yeah, yeah, scratch that place on my back! It itches like crazy!~\LINK{http://naviteri.org/2020/07/vospxivol-lefpom-happy-august/}{b}
(\ding{43}~\HL{kum}{\B{kum}})

% kunsìp
\hi
\HT{kunsIp}{\B{\cp{kun}sìp}} \I{["kun.sIp\textcorner]} \emph{or} \I{["kUn.]} (\textsc{rn}: \B{\cp{kùn}sìp} \I{["kUn.sIp\textcorner]}) \emph{n.} \textsc{lw} gunship. 
\B{Kun\-sìp\-ìri txan\-a tì\-meyp lu tsyal a mìn.} The gunship's main weakness is the rotor system.~\LINK{http://forum.learnnavi.org/language-updates/ftarpa-and-skienpa-added-to-dict/msg570362/\#msg570362}{ln\slash a\textsuperscript{c}}

% kung
\hi 
\HT{kung}{\B{kung}} \I{[kuN]} \emph{or} \I{[kUN]} (\textsc{rn}: \B{kùng} \I{[kUN]}) \emph{adj.} putrid, fetid, rotten (\emph{referred to rotten meat or a pile of dead and rotting animal matter}). 
\B{Tsa\-fa\-hew a\-on\-vä' ftu kung\-a ioang za\-'u.} That stinking smell comes from a rotten animal.~\LINK{http://naviteri.org/2019/09/choice-statements-vs-choice-questions-and-some-insults/}{b}
(\ding{43}~\HL{snAm}{\B{snäm}}) \\
\textsc{Derivatives}: \ding{43} \on{1}.~\HL{kurkung}{\B{kur\-kung}}, asshole (\emph{insult}).

% kur
\hi
\HT{kur}{\B{kur}} \I{[kuR]} \emph{or} \I{[kUR]} (\textsc{rn}: \B{kùr} \I{[kUR]}) \emph{vin.} hang (\HL{ta}{\B{ta}}, from).
\B{Fkxi\-le pewn\-ta tu\-té\-yä kur.} The bib necklace hangs from the woman's neck.~\LINK{http://naviteri.org/2013/04/wheres-the-bathroom-and-other-useful-things/}{b}
(\emph{a. proverb}) \B{Ke kur fko fa kxe\-tse.} ``One can't hang by a tail,'' i.e. \emph{don't rely on s.t.\slash s.o. untrustworthy or useless, just as a Na'vi tail can't be relied on to bear weight}.~\LINK{http://naviteri.org/2021/08/ulte-ayyoratu-leiu-and-the-winners-are-2/}{b} \\
\textsc{Derivatives}: \ding{43} \on{1}.~\HL{keykur}{\B{key\-kur}}, hang (s.t.), let hang. 
\on{2}.~\HL{kurfyan}{\B{kur\-fyan}}, hamper, suspended rack.

% kurakx
\hi
\HT{kurakx}{\B{ku\cp{rakx}}} \I{[ku."Rak']} \emph{vtr.} drive out. 
\B{Ay\-oe\-ngal fte Fay\-saw\-tu\-tet ki\-vu\-rakx, ze\-ne nì\-'aw\-ve ay\-fot tsli\-vam.} We must understand these Sky People if we are to drive them out.~\LINK{http://forum.learnnavi.org/language-updates/confirmation-on-kurakx-eytukan-line-and-toitslan/}{ln\slash a\textsuperscript{c}}

% kurfyan
\hi
\HT{kurfyan}{\B{\cp{kur}fyan}} \I{["kuR.fjan]} \emph{or} \I{["kUR.]} (\textsc{rn}: \B{\cp{kùr}fya} \I{["kUR.fja]}) \emph{n.} hamper, suspended rack.
(\ding{43}~\HL{kur}{\B{kur}}, \HL{fyan}{\B{fyan}}) \\
\textsc{Derivatives}: \ding{43} \on{1}.~\HL{snokfyan}{\B{snok\-fyan}}, personal belongings rack. 
\on{2}.~\HL{kurfyavi}{\B{kur\-fya\-vi}}, hook.

% kurfyavi
\hi
\HT{kurfyavi}{\B{\cp{kur}fyavi}} \I{["kuR.fja.vi]} \emph{or} \I{["kUR.]} (\textsc{rn}: \B{\cp{kùr}fyavi} \I{["kUR.fja.vi]}) \emph{n.} hook (\emph{for hanging or suspending an item}).
(\ding{43}~\HL{kurfyan}{\B{kur\-fyan}}) 

% kurkung
\hi 
\HT{kurkung}{\B{\cp{kur}kung}} \I{["kuR.kuN]} \emph{or} \I{["kUR.kUN]} (\textsc{rn}: \B{\cp{kur}kùng} \I{["kuR.kUN]}) \emph{n., vulg.} asshole (\emph{strong insult}). 
(\ding{43}~\HL{kuru}{\B{ku\-ru}}, \HL{kung}{\B{kung}})

% kuru
\hi
\HT{kuru}{\B{kuru}} \I{["ku.Ru]} \emph{n.} neural queue.
\ding{43}~\HL{tswin}{\B{tswin}} \\
\textsc{Derivatives}: \ding{43} \on{1}.~\HL{kurkung}{\B{kur\-kung}}, asshole (\emph{strong insult}). 

\end{multicols}

% ===== Section KX =====

 \pdfbookmark[0]{KX}{KX}
\begin{center}
\section*{KX \I{[k']} (\emph{KxeKx})}
\end{center}

\begin{multicols}{2}

% kxa
\hi
\HT{kxa}{\B{kxa}} \I{[k'a]} \emph{n.} mouth. 
\B{Kxa\-ri tstu\-tswo tsran\-ten krra ke lu kea sä\-fpìl le\-sar.} When one has no useful thoughts, the ability to close one's mouth is important.~\LINK{http://naviteri.org/2012/03/spring-vocabulary-part-2/}{b}

% kxakx
\hi
\HT{kxakx}{\B{kxakx}} \I{[k'ak']} \emph{vin.} break, snap in two.
\B{Txo vul ivil nì\-hawng kxì\-ye\-vakx.} If a branch bends too much, it might break.~\LINK{http://naviteri.org/2013/04/wheres-the-bathroom-and-other-useful-things/}{b}

% kxam
\hi
\HT{kxam}{\B{kxam}} \I{[k'am]} \emph{n.} middle, midpoint. \\ 
\textsc{Derivations}: \ding{43} \on{1}.~\HL{kxamlA}{\B{kxam\-lä}}, through.  
\on{2}.~\HL{kxamtrr}{\B{kxam\-trr}}, midday, noon. 
\on{3}.~\HL{kxamtxon}{\B{kxam\-txon}}, midnight.
\on{4}.~\HL{kxamtseng}{\B{kxam\-tseng}}, center.
\on{5}.~\HL{kxamyI}{\B{kxam\-yì}}, intermediate level. 
\on{6}.~\HL{mIkam}{\B{mì\-kam}}, between. 

% kxamlä
\hi
\HT{kxamlA}{\B{\cp{kxam}lä}} \I{["k'am.l\ae]} \emph{adp.} through, via the middle of. 
\B{Txon\-am teng\-krr tar\-mì\-ran oe kxam\-lä na'\-rìng, sro\-ler eo ut\-ral a\-tsawl txewm\-a vrr\-tep.} Last night as I was walking through the forest, a frightening demon appeared in front of a big tree.~\LINK{http://naviteri.org/2012/03/spring-vocabulary-part-1/}{b}
(\ding{43}~\HL{kxam}{\B{kxam}}, \HL{IlA}{\B{ìlä}})

% kxamtrr
\hi
\HT{kxamtrr}{\B{\cp{kxam}trr}} \I{["k'am.t\s{r}]} \emph{n.}; \emph{adv.} midday, noon.
\B{Nän yom kxam\-trr, 'ul 'e\-fu o\-hakx kaym.} The less you eat at noon, the hungrier you'll feel in the evening.~\LINK{http://naviteri.org/2012/02/trr-asawnung-lefpom-happy-leap-day/}{b}
\B{Kxam\-trr lam fwa san\-hì a mì saw 'o\-lìp nì\-wotx slä tsa\-krr ke tsun fko sat tsi\-ve\-'a nì\-'aw.} At mid-day it seems that the stars in the sky have all vanished but they just can't be seen then.~\LINK{http://naviteri.org/2012/03/spring-vocabulary-part-1/}{b}
(\ding{43}~\HL{kxam}{\B{kxam}}, \HL{trr}{\B{trr}}) \\
\textsc{Derivations}: \ding{43} \on{1}.~\HL{kxamtrrmaw}{\B{kxam\-trr\-maw}}, after midday. 
\on{2}.~\HL{srekamtrr}{\B{sre\-kam\-trr}}, before midday.

% kxamtrrmaw
\hi
\HT{kxamtrrmaw}{\B{\cp{kxam}trrmaw}} \I{["k'am.t\s{r}.maw]} \emph{n.}; \emph{adv.} after midday, after noon.
(\ding{43}~\HL{kxamtrr}{\B{kxam\-trr}}, \HL{maw}{\B{maw}})

% kxamtxomaw
\hi
\HT{kxamtxomaw}{\B{\cp{kxam}txomaw}} \I{["k'am.t'o.maw]} \emph{n.}; \emph{adv.} after the middle of the night.
(\ding{43}~\HL{kxamtxon}{\B{kxam\-txon}}, \HL{maw}{\B{maw}})

% kxamtxon
\hi
\HT{kxamtxon}{\B{\cp{kxam}txon}} \I{["k'am.t'on]} \emph{n.}; \emph{adv.} middle of the night, midnight.
(\ding{43}~\HL{kxam}{\B{kxam}}, \HL{txon}{\B{txon}}) \\
\textsc{Derivations}: \ding{43} \on{1}.~\HL{kxamtxomaw}{\B{kxam\-txo\-maw}}, after midnight. 
\on{2}.~\HL{srekamtxon}{\B{sre\-kam\-txon}}, before midnight.

% kxamtseng
\hi
\HT{kxamtseng}{\B{\cp{kxam}tseng}} \I{["k'am.\t{ts}EN]} \emph{n.} center, place in the middle. 
\B{Re\-rol teng\-krr ke\-rä | Ìlä fya\-'o a\-vol | Ne kxam\-tseng.} We sing our way, down the eight paths, to the center.~\LINK{http://www.pandorapedia.com/navi/music/ritual_music}{pp}
(\ding{43}~\HL{kxam}{\B{kxam}}, \HL{tseng}{\B{tse\-nge}})

% kxamyì
\hi
\HT{kxamyI}{\B{\cp{kxam}yì}} \I{["k'am.jI]} \emph{n.} intermediate level.
\B{Lì'\-fya\-ri tsa\-taw\-tu\-te pe\-yì?--Kxam\-yì.~Pll\-txe nìk\-sran, slä tsun fko pe\-yä ay\-lì\-'ut tsli\-vam.} How is that Sky Person's Na'vi?--It's intermediate.~His speech is mediocre, but you can understand him.~\LINK{http://naviteri.org/2013/01/lifyari-po-peyi-how-good-is-her-navi/}{b}
(\ding{43}~\HL{kxam}{\B{kxam}}, \HL{yI}{\B{yì}})

% kxanì
\hi
\HT{kxanI}{\B{\cp{kxa}nì}} \I{["k'a.nI]} \emph{adj.} forbidden.  
\B{Fay\-vrr\-tep fì\-tse\-nge lu kxa\-nì.} These demons are forbidden here.~\LINK{http://wiki.learnnavi.org/index.php?title=Corpus\#Times_Online}{wiki}

% kxangangang
\hi
\HT{kxangangang}{\B{\cp{kxang}angang}} \I{["k'aN.aN.aN]} \emph{onom.} boom. 

% kxange
\hi 
\HT{kxange}{\B{\cp{kxa}nge}} \I{["k'a.NE]} \emph{vin.} yawn (\emph{as a result of fatigue or boredom}). 
\B{Oe kxì\-ma\-nge ta\-lu\-na 'e\-fu ngeyn.} I just yawned because I feel tired.~\LINK{http://naviteri.org/2018/07/ta-sulfatu-a-ayliu-niul-more-words-from-our-experts/}{b}
\B{Keng krra sä\-num\-vi el\-tur tì\-txen ke si, nga swey\-lu txo ke kxi\-va\-nge mì num\-tseng.} Even when the lesson isn't interesting, you shouldn't yawn in school.~\LINK{http://naviteri.org/2018/07/ta-sulfatu-a-ayliu-niul-more-words-from-our-experts/}{b} \\
\textsc{Derivations}: \ding{43} \on{1}.~\HL{sAkxange}{\B{sä\-kxa\-nge}}, a yawn.

% kxap / si
\hi
\HT{kxap}{\B{kxap}} \I{[k'ap\textcorner]} \textsc{i}. \emph{n.} threat.
\B{Tì\-pä\-hem Saw\-tu\-te\-yä kxap a\-txan lar\-mu Na'\-vi\-ru.} The arrival of the Sky People was a great threat to the Na'\-vi.~\LINK{http://naviteri.org/2015/03/kap-si-ayunil-saylahe-threats-dreams-and-other-things/}{b}
\B{Kap sì ay\-u\-nil say\-la\-he.} Threats, dreams, and other things.~\LINK{http://naviteri.org/2015/03/kap-si-ayunil-saylahe-threats-dreams-and-other-things/}{b}
(\emph{a. with} \B{wä}) a threat to s.t.\slash s.o., threaten.
\B{Krr\-a kxap lar\-mu sì\-rey\-wä fe\-yä nì\-wotx.} When their lives were threatened.~\LINK{http://naviteri.org/2015/05/cirque-du-soleils-toruk-the-first-flight-teaser-video-online/}{b} \\
\textsc{ii}. \HT{kxap si}{\B{kxap si}} \I{[k'ap\textcorner{} si]} \emph{vin.} (\emph{a.}) threaten (\ding{43}~\HL{fko}{\B{fkor}}, s.o.). 
\B{Sra\-ke nga kxap si oer, ma \mbox{skxawng}?} Are you threatening me, you moron?~\LINK{http://naviteri.org/2015/03/kap-si-ayunil-saylahe-threats-dreams-and-other-things/}{b}
(\emph{b. of animals and people}) threatening. \B{'Ang\-tsìk a kxap si lu le\-hrr\-ap, ma 'itan.} A threatening hammerhead is dangerous, son.~\LINK{http://naviteri.org/2015/03/kap-si-ayunil-saylahe-threats-dreams-and-other-things/}{b} \\
\textsc{Derivations}: \ding{43} \on{1}.~\HL{kxapnga'}{\B{kxap\-nga'}}, threatening. 
\on{2}.~\HL{nIkxap}{\B{nì\-kxap}}, threateningly.

% kxapnga'
\hi
\HT{kxapnga'}{\B{\cp{kxap}nga'}} \I{[k'ap\textcorner.NaP]} \emph{adj.} threatening (\emph{not for people}).
\B{Nge\-yä ay\-lì\-'ul a\-kxap\-nga' txo\-pu ke sley\-ku oet kaw\-'it.} Your threatening words don't scare me one bit.~\LINK{http://naviteri.org/2015/03/kap-si-ayunil-saylahe-threats-dreams-and-other-things/}{b}
(\ding{43}~\HL{kxap si}{\B{kxap si}} (\emph{b.}), \HL{kxap}{\B{kxap}})

% kxawm
\hi
\HT{kxawm}{\B{kxawm}} \I{[k'awm]} \emph{adv.} perhaps, maybe.
\B{Oel fo\-ru fì\-ay\-lì\-'ut to\-lìng a krr, kxawm oe har\-ma\-hä\-ngaw.} When I gave them these words I was probably sleeping.~\LINK{http://wiki.learnnavi.org/index.php?title=Canon\#Extracts_from_various_emails}{wiki} 
\B{Kxawm swo\-at nì\-hawng nil\-vä\-ngäk oel txon\-am.} Maybe I drank too much last night.~\LINK{http://naviteri.org/2013/12/mipa-zisit-lefpom-happy-new-year/comment-page-1/\#comment-2585}{b}

% kxayl
\hi
\HT{kxayl}{\B{kxayl}} \I{[k'ajl]} \emph{adj.} high (\emph{scalable height}). 
\B{Ney\-ti\-ri\-ri lu yì tì\-sti\-yä kxayl nì\-ngay.} Neytiri is really angry [lit.: Neytiri's level of anger is truly high].~\LINK{http://naviteri.org/2013/01/lifyari-po-peyi-how-good-is-her-navi/}{b}
\B{Tsa\-ye\-rik kll\-kxem sìn yì a\-kxayl.} That \emph{yerik} is standing on a high ledge.~\LINK{http://naviteri.org/2013/01/lifyari-po-peyi-how-good-is-her-navi/}{b} \\ 
\textsc{Derivations}: \ding{43} \on{1}.~\HL{kxaylyI}{\B{kxayl\-yì}}, high level.

% kxaylte
\hi
\HT{kxaylte}{\B{\cp{kxayl}te}} \I{["k'ajl.tE]} \emph{n.}~\ding{96} cillaphant, \emph{clawenia purpurea}.

% kxaylyì
\hi
\HT{kxaylyI}{\B{\cp{kxayl}yì}} \I{["k'ajl.jI]} \emph{n.} high level.
\B{Lì'\-fya\-ri tsa\-taw\-tu\-te pe\-yì?--Kxayl\-yì. Ke tsun oe spi\-vaw. Slo\-lu po tsul\-fä\-tu lì'\-fya\-yä aw\-nge\-yä.} How is that Sky Person's Na'vi?--It's excellent. I can't believe it. He's become a master of our language.~\LINK{http://naviteri.org/2013/01/lifyari-po-peyi-how-good-is-her-navi/}{b}
(\ding{43}~\HL{kxayl}{\B{kxayl}}, \HL{yI}{\B{yì}})

% kxänäng
\hi
\HT{kxAnAng}{\B{\cp{kxä}näng}} \I{["k\ae.n\ae{}N]} \emph{adj.} putrid, \emph{smell of decaying animal\slash flesh}.
\B{Fì\-tsngan\-ur a sno\-läm lä\-ngu fa\-hew a\-kxä\-näng.} This rotten meat has a putrid smell.~\LINK{http://naviteri.org/2015/11/vomuna-liu-amip-ten-new-words/}{b}

% kxeltek
\hi
\HT{kxeltek}{\B{\cp{kxel}tek}} \I{["k'El.tEk\textcorner]} \emph{vtr.} pick up, lift.   
\B{Pxi\-set nge\-yä tska\-lep\-it kxel\-tek!} Pick up your crossbow right now!~\LINK{http://naviteri.org/2012/01/mipa-zisit-ayliu-amip-new-words-for-the-new-year/}{b} 
\B{Ke tsun tu\-te a\-'aw tsa\-tskxe\-ti a\-ku'\-up kxi\-vel\-tek nì\-'aw\-tu.} One person alone can't lift that heavy rock.~\LINK{http://naviteri.org/2012/01/mipa-zisit-ayliu-amip-new-words-for-the-new-year/}{b}
(\ding{214}~\HL{keyn}{\B{keyn}})

% kxem
\hi
\HT{kxem}{\B{kxem}} \I{[k'Em]} \emph{vin.} be vertical.   
\B{Fì\-ru\-mut lum\-pe ke kxem?} Why isn't this puffball tree vertical?~\LINK{http://naviteri.org/2012/03/spring-vocabulary-part-1/}{b} \\  
\textsc{Derivations}: \ding{43} \on{1}.~\HL{kllkxem}{\B{kllkxem}}, stand. 
\on{2}.~\HL{kxemyo}{\B{kxem\-yo}}, wall. 
\on{3}.~\HL{nIkxem}{\B{nì\-kxem}}, vertically.

% kxemyo
\hi
\HT{kxemyo}{\B{\cp{kxem}yo}} \I{["k'Em.jo]} \emph{n.} wall, vertical surface. 
\B{Rìk a\-'aw syar\-mar sìn kxem\-yo.} A leaf was sticking to the wall.~\LINK{http://naviteri.org/2020/11/vospxivopeya-ayliu-amip-novembers-new-words/}{b}
\B{Tseng\-a 'aw\-steng\-yä\-pem fì\-me\-kem\-yo lu mek\-tseng.} Where these two walls come together there's a gap.~\LINK{http://naviteri.org/2013/04/wheres-the-bathroom-and-other-useful-things/} {b}
(\ding{43}~\HL{kxem}{\B{kxem}}, \HL{yo}{\B{yo}})

% kxener
\hi
\HT{kxener}{\B{\cp{kxe}ner}} \I{["k'E.nER]} \emph{n.} smoke. 
\B{Kxe\-ner lu sä\-vll txep\-ä.} Smoke is a sign\slash an indication that there is fire.~\LINK{http://naviteri.org/2017/12/zisit-amip-lefpom-ma-eylan-happy-new-year-friends/}{b} \\
(i.) \B{tswìk \cp{kxe}nerit} \I{[\t{ts}wIk\textcorner{} "k'E.nE.Rit\textcorner]} \emph{vin.} smoke (\emph{a cigarette})
\B{Ke tung Na'\-vil fu\-ta tswìk kxe\-ner\-it Mo\-'a\-ra\-ka nì\-wotx.} The Na'vi have banned smoking in Mo'ara.~\LINK{http://naviteri.org/2017/06/ayliu-a-ta-eywaeveng-kifkey-uniltirantokxa-words-from-disneys-pandora-the-world-of-avatar/}{b}

% kxetse
\hi
\HT{kxetse}{\B{\cp{kxe}tse}} \I{["k'E.\t{ts}E]} \emph{n.} tail. 
(\emph{a.}) \B{Nge\-yä kxe\-tse lor lu nì\-txan.} Your tail is very beautiful.~\LINK{http://www.youtube.com/watch?v=ojr_BE8ZmVE}{yt}
\B{Ut\-ral\-ti a tsa\-uo Lo\-ak wä\-par\-man pol vol fa kxe\-tse.} She used her tail to point out the tree Loak was hiding behind.~\LINK{http://naviteri.org/2012/06/spring-vocabulary-part-3/}{b}
(\emph{b. idiom}) \B{Kxe\-tse sì mik\-yun kop pll\-txe.} Body language is important [lit.: The tail and ears also speak].~\LINK{http://forum.learnnavi.org/language-updates/grid-tidbits/}{ln} 
(\emph{c. idiom}) \B{Po keyn\-ven sìn ke\-tse.} He is socially awkward. (\emph{lit.}: `He steps on tails.')~\LINK{http://naviteri.org/2019/08/aawa-lifya-si-lifyavi-amip-a-few-new-words-and-expressions/}{b} 
(\emph{d. proverb}) \B{Ke kur fko fa kxe\-tse.} ``One can't hang by a tail,'' i.e. \emph{don't rely on s.t.\slash s.o. untrustworthy or useless, just as a Na'vi tail can't be relied on to bear weight}.~\LINK{http://naviteri.org/2021/08/ulte-ayyoratu-leiu-and-the-winners-are-2/}{b}
(\emph{e. idiom}) \B{Kxe\-tse kì\-'ong!} ``Don't get angry.\slash Calm down.'' [lit.: slow tail] (\emph{short for} \B{Nga\-ri kxe\-tse kì\-'ong li\-vu.}).~\LINK{http://naviteri.org/2021/08/ulte-ayyoratu-leiu-and-the-winners-are-2/}{b} \\
\textsc{Derivations}: \ding{43} \on{1}.~\HL{kxetsikran}{\B{kxe\-tsik\-ran}}, banshee's tail.

% kxetsikran
\hi 
\HT{kxetsikran}{\B{kxe\cp{tsik}ran}} \I{[k'E."\t{ts}ik\textcorner.Ran]} \emph{n.} \ding{96} banshee's tail. 
(\ding{43}~\HL{kxetse}{\B{kxetse}}, \HL{ikran}{\B{ikran}})

% kxeyey
\hi
\HT{kxeyey}{\B{\cp{kxey}ey}} \I{["k'Ej.Ej]} \emph{n.} error, mistake.
\B{Tsa\-kxey\-ey li zo\-slo\-lu.} That mistake has already been fixed.~\LINK{http://naviteri.org/2011/02/new-vocabulary-part-2/comment-page-1/\#comment-497}{b}
\B{Txo kxey\-ey\-ti ay\-ngal tsi\-ve\-'a, ru\-txe oe\-ru pi\-veng.} If you see a mistake, please tell me.~\LINK{http://naviteri.org/2011/05/some-miscellaneous-vocabulary/}{b}
\B{Ke ro\-lun oel kea ay\-key\-eyt!} I have found no mistakes!~\LINK{http://naviteri.org/2011/03/word-order-and-case-marking-with-modals/comment-page-1/\#comment-584}{b}
\B{Kxey\-ey\-ri tsap\-'a\-lu\-te se\-ngi oe.} I apologize for the mistake.~\LINK{http://forum.learnnavi.org/language-updates/frommerian-email/msg56711/\#msg56711}{ln}
\B{Slä krr a nì\-Na'\-vi pam\-rel si oe, new oel fu\-ta upxa\-re oe\-yä lu\-ke key\-ey li\-vu nì\-wotx, fte eyawr\-a sì\-ke\-nong\-it ti\-vìng su\-te\-ru.} But when I write Na'vi, I want my messages to be without any mistakes to give people correct examples.~\LINK{http://forum.learnnavi.org/language-updates/verb-for-write/msg175592/\#msg175592}{ln} \\  
\textsc{Derivations}: \ding{43} \on{1}.~\HL{tIkxey}{\B{tì\-kxey}}, incorrectness.

% kxitx
\hi
\HT{kxitx}{\B{kxitx}} \I{[k'it']} \emph{n.} death (\emph{the event of an individual's death}).
\B{Sa'\-sem\-ìri lu 'eveng\-ä \mbox{kxitx} yeng\-wal a\-tu\-vom.} For a parent, the greatest sorrow is the death of a child.~\LINK{http://naviteri.org/2015/04/some-new-words-for-may-day/}{b}
\B{Maw \mbox{kxitx} sem\-pul\-ä lar\-mä\-ngu Pey\-ral\-ä a\-fpawng txew\-luke.} After (her) father's death, Pey\-ral's grief was endless.~\LINK{http://naviteri.org/2018/03/100a-liu-amip-64-new-words-part-1/}{b}
(\ding{43}~\HL{tIterkup}{\B{tì\-ter\-kup}})

% kxì
\hi
\HT{kxI}{\B{kxì}} \I{[k'I]} \emph{intj.} hi, hiya (\emph{casual greeting}).
\B{Kxì, ma 'ey\-lan! Kem\-pe le\-ren?} Hey dude! What's happenin'?.~\LINK{http://naviteri.org/2018/02/negative-questions-in-navi/}{b}
(\ding{43}~\HL{kaltxI}{\B{kal\-txì}})

% kxìm
\hi
\HT{kxIm}{\B{kxìm}} \I{[k'Im]} \emph{vtr.} command, order, assign a task (\emph{with} \ding{43}~\HL{txonta}{\B{tson\-ta}}).
\B{Ay\-evengur kxo\-lìm sa'\-nok\-ìl fì\-tson\-it.} Mother imposed this task on the children.~\LINK{http://naviteri.org/2014/10/tson-si-fpomron-obligation-and-mental-health/}{b}
\B{Ay\-eveng\-ur kxo\-lìm sa'\-nok\-ìl tson\-ta pay\-it za\-mu\-nge.} The children were told by their mother to fetch water.~\LINK{http://naviteri.org/2014/10/tson-si-fpomron-obligation-and-mental-health/}{b}
(\emph{a. proverb}) \B{Kxìm utu\-ftu fna\-we'\-tu.} ``A coward commands from the canopy,'' i.e. \emph{a real leader will have `boots on the ground' and will help out, whereas a coward will only tell people (from afar) what to do}.~\LINK{http://naviteri.org/2021/08/ulte-ayyoratu-leiu-and-the-winners-are-2/}{b} \\
\textsc{Derivations}: \ding{43} \on{1}.~\HL{kxImyu}{\B{kxìm\-yu}}, commander, one with authority over another. 
\on{2}.~\HL{tIkxIm}{\B{tì\-kxìm}}, commanding, ordering, assigning tasks.

% kxìmyu
\hi
\HT{kxImyu}{\B{\cp{kxìm}yu}} \I{["k'Im.ju]} \emph{n.} commander, one with authority over another.
\B{Nge\-yä kxìm\-yu pe\-su?} Who's your boss?~\LINK{http://naviteri.org/2014/10/tson-si-fpomron-obligation-and-mental-health/}{b}
(\emph{a. idiom, politeness formula})
\B{Kxìm\-yu nga.} ``Please! Go ahead. You first.'' [lit.: (my) commander (is) you]~\LINK{http://naviteri.org/2014/10/tson-si-fpomron-obligation-and-mental-health/}{b}
(\ding{43}~\HL{kxIm}{\B{kxìm}})

% kxll / kxll si
\hi
\HT{kxll}{\B{kxll}} \I{[k'\s{l}]} \textsc{i}.~\emph{n.} charge. \\
\textsc{ii}.~\HT{kxll si}{\B{kxll si}} \I{[k'\s{l} si]} \emph{vin.} charge (toward, \B{ne}).   
\B{Teng\-krr pa\-lu\-lu\-kan moe\-ne kxll sar\-mi, pol\-txe Ney\-ti\-ril ay\-lì\-'ut a fra\-krr 'ok se\-yä la\-yu oer.} As the thanator was charging towards the two of us, Neytiri said something I will always remember.~\LINK{http://thereadinginpublicproject.blogspot.de/2010/05/paul-frommer.html}{pr}

% kxon
\hi 
\HT{kxon}{\B{kxon}} \I{[k'on]} \emph{n.} internal organ. 
\B{Lu tokx\-ur tu\-te\-yä pol\-pxay\-a kxon?} How many internal organs does a person's body have?~\LINK{https://naviteri.org/2024/06/ayhapxi-tokxa-si-ayliu-alahe-parts-of-the-body-and-more/}{b}

% kxor
\hi
\HT{kxor}{\B{kxor}} \I{[k'oR]} \emph{rarely countable n.} a wall or bank of powerful waterfalls noted for its deafening roar and deadly force. 

% kxu / kxu si
\hi
\HT{kxu}{\B{kxu}} \I{[k'u]} \textsc{i}. \emph{n.} harm.
\B{Lu me\-nge\-yä kel\-ku na'\-rìng\-ä lu\-ke kxu a\-txan a fì\-'u fmawn a\-sìl\-tsan lu nì\-ngay.} It's truly good news that your forest house is without much harm.\LINK{http://forum.learnnavi.org/language-updates/luke-kxu-without-harm/}{ln}
\B{Fnu, fnu, wä\-pan wä ku | Uo tsa\-ìpxa.} Quiet, quiet, hide yourself from harm | Behind that fern.~\LINK{http://forum.learnnavi.org/language-updates/learnings-from-the-siatlla-song-translations/msg473566/\#msg473566}{ln}  \\  
\textsc{ii}. \B{kxu si} \I{[k'u si]} \emph{vin.} harm.
\B{Po\-eru kxu rä\-'ä si!} Do not harm her!~\LINK{http://naviteri.org/2017/06/ayliu-a-ta-eywaeveng-kifkey-uniltirantokxa-words-from-disneys-pandora-the-world-of-avatar/}{b} 
\B{Rä\-'ä fmi\-vi kxu si\-vi ay\-ioang\-ur fu hel\-kur fe\-yä} Do not attempt to harm the animals or their habitats.~\LINK{http://naviteri.org/2017/06/ayliu-a-ta-eywaeveng-kifkey-uniltirantokxa-words-from-disneys-pandora-the-world-of-avatar/}{b} \\
\textsc{Derivations}: \ding{43} \on{1}.~\HL{kxutu}{\B{kxu\-tu}}, enemy. 
\on{2}.~\HL{kxuke}{\B{kxu\-ke}}, safe. 
\on{3}.~\HL{mAkxu}{\B{mä\-kxu}}, interrupt. 
\on{4}.~\HL{kxutslu}{\B{kxu\-tslu}}, risk.

% kxuke
\hi
\HT{kxuke}{\B{\cp{kxu}ke}} \I{["k'u.kE]} \emph{adj.} safe. 
(\ding{43}~\HL{kxu}{\B{kxu}}, \HL{luke}{\B{lu\-ke}}) \\
\textsc{Derivations}: \ding{43} \on{1}.~\HL{tIkxuke}{\B{tì\-kxu\-ke}}, safety.

% kxukx
\hi
\HT{kxukx}{\B{kxukx}} \I{[k'uk']} \emph{or} \I{[k'Uk']} \emph{vtr.} swallow.   
\B{Fì\-na\-er ftxì\-vä' lu nì\-hawng, ha swey\-lu txo ngal tsat kxi\-vukx nì\-win.} This drink tastes horrible, so you'd better swallow it quickly.~\LINK{http://naviteri.org/2011/11/\%E2\%80\%99a\%E2\%80\%99awa-ayli\%E2\%80\%99u-amip-ni\%E2\%80\%99aw-\%E2\%80\%94-a-few-new-words-only/}{b} 

% kxum
\hi
\HT{kxum}{\B{kxum}} \I{[k'um]} \emph{or} \I{[k'Um]} \emph{adj.} viscous, gelatinous, thick. \\
\textsc{Derivations}: \ding{43} \on{1}.~\HL{kxumpay}{\B{kxum\-pay}}, viscous liquid, gel.

% kxumpay
\hi
\HT{kxumpay}{\B{\cp{kxum}pay}} \I{["k'um.paj]} \emph{or} \I{["k'Um.]} (\textsc{rn}: \B{\cp{kxum}\-pay} \I{["gum.paj]}) \emph{n.} viscous liquid, gel. 
(\emph{a. idiom}) \B{(Na) ken\-ten mì kum\-pay}, being in an environment where one is prevented from acting naturally or doing what one wants to do, [lit.: (like) a fan lizard in gel] \B{Nar\-mew oe fo\-ru na'\-rìng\-ä tì\-lor\-it wi\-vìn\-txu, slä ke tsä\-ngun fo tsli\-vam. 'Efu oe na ken\-ten mì kum\-pay.} I wanted to show them the beauty of the forest, but sadly, they weren't able to understand. I felt completely stymied.~\LINK{http://naviteri.org/2014/11/twenty-before-the-holidays/}{b} 
(\ding{43}~\HL{kxum}{\B{kxum}}, \HL{pay}{\B{pay}}) \\
\textsc{Derivations}: \ding{43} \on{1}.~\HL{txekxumpay}{\B{txe\-kxum\-pay}}, magma, lava.
\on{2}.~\HL{kxumpaysyar}{\B{kxum\-pay\-syar}}, glue.

% kxumpaysyar
\hi 
\HT{kxumpaysyar}{\B{\cp{kxum}paysyar}} \I{["k'um.paj.sjaR]} \emph{or} \I{["k'Um.]} \emph{n.} glue. 
\B{Pol kxum\-pay\-syar\-it so\-lar syey\-kar rìk\-it sìn kxem\-yo.} She stuck the leaf onto the wall with glue. [lit.: ``She used glue (and then) stuck the leaf onto the wall.'']~\LINK{http://naviteri.org/2020/11/vospxivopeya-ayliu-amip-novembers-new-words/}{b}
(\ding{43}~\HL{kxumpay}{\B{kxum\-pay}}, \HL{syar}{\B{syar}})

% kxutu
\hi
\HT{kxutu}{\B{\cp{kxu}tu}} \I{["k'u.tu]} (\textsc{rn}: \B{\cp{kxu}tu} \I{["gu.tu]}) \emph{n.} enemy. 
\B{Ke za\-syup lì\-'O\-na ne kxu\-tu a mì\-fa fu a wrr\-pa.} The l'Ona will not perish to the enemy within or the enemy without.~\LINK{http://forum.learnnavi.org/language-updates/ltasygt-with-ke-and-nga/}{ln\slash g}
\B{Ay\-ha\-pxì\-tu po\-ngu\-ä txo\-pu si nì\-nän ta\-krr\-a Va'\-rul pxe\-ku\-tut lä\-txayn.} The members of the group are less afraid since Va’ru defeated three of the enemy.~\LINK{http://naviteri.org/2012/02/trr-asawnung-lefpom-happy-leap-day/}{b}
\B{Nì\-trr\-trr pxì\-mun\-'i sam\-si\-yul ay\-swi\-zaw\-it ku\-tu\-ä a\-law\-nä\-txayn sno\-kip nì\-'eng.} Warriors typically share the arrows of their defeated enemies among themselves.~\LINK{http://naviteri.org/2012/01/more-additions-to-the-lexicon/}{b}
(\ding{43}~\HL{kxu}{\B{kxu}})

% kxutslu
\hi 
\HT{kxutslu}{\B{\cp{kxu}tslu}} \I{["k'u.\t{ts}lu]} \emph{n.} risk. 
(\ding{43}~\HL{kxu}{\B{kxu}}, \HL{tItsunslu}{\B{tìtsunslu}})

\end{multicols}

% ===== Section L =====   

 \pdfbookmark[0]{L}{L}
\begin{center}
\section*{L \I{[l]} (\emph{LeL})}
\end{center}

\begin{multicols}{2}

% -l
\hi
\HT{l}{\B{-l}} \I{[l]} \emph{suf. to mark the agentive case of a n. ending in a vowel or diphthong}. 
(\ding{43}~\HL{Il}{\B{-ìl}})

% la'a
\hi 
\HT{la'a}{\B{\cp{la}'a}} \I{["la.Pa]} \emph{n.} physical separation, distance between two places or objects. 
\B{Kxawm set 'u a\-ngä\-zìk fra\-to lu la\-'a a\-yll.} Maybe the most difficult thing now is social distance.~\LINK{http://naviteri.org/2020/03/teri-tifkeytok-lefkrr-about-the-present-situation/}{b} \\
\textsc{Derivatives}: \ding{43} \on{1}.~\HL{pela'a}{\B{pe\-la\-'a}}, how near, how far, what distance.
\on{2}.~\HL{la'ape}{\B{la\-'a\-pe}}, how near, how far, what distance.

% la'ang
\hi 
\HT{la'ang}{\B{\cp{la}'ang}} \I{["la.Pang]} \emph{n.} pile of stinking, rotting animal matter. 
\B{'U\-pe ke zo?--Fraw mì la\-'ang.} What's the matter?--Everything is screwed up.~\LINK{http://naviteri.org/2020/07/vospxivol-lefpom-happy-august/}{b}
(\emph{a. expression}) \B{Pe\-la\-'ang?} `What the hell? What's going on?' [lit.: what pile of stinking, rotting animal matter?]~\LINK{http://naviteri.org/2020/07/vospxivol-lefpom-happy-august/\#comment-31420}{b}

% la'ape
\hi 
\HT{la'ape}{\B{\cp{la}'ape}} \I{["la.Pa.pE]} \emph{inter.}  how near? how far? what distance?
(\ding{43}~\HL{la'a}{\B{la'a}}, \HL{pe}{\B{pe}\texttt{+}}; \HL{pelImsim}{\B{pelìmsim\slash lìmsimpe}})

% la'um
\hi 
\HT{la'um}{\B{\cp{la}'um}} \I{["la.Pum]} \emph{or} \I{[.PUm]} (\textsc{rn}: \B{\cp{la}\-ùm} \I{["la.Um]}) \emph{vin.} pretend (\emph{used with} \ding{43}~\HL{tsnI}{\B{tsnì}}). 
\B{Pll\-txe po san nga yaw\-ne lu oer sìk, sl\"a la\-'um nì\-'aw.} He says he loves you, but he’s only pretending.~\LINK{http://naviteri.org/2020/02/some-words-for-leap-year-day/}{b}
\B{Lum\-pe nga la\-'um tsnì ke tsun sri\-vew?} Why are you pretending (that) you can't dance?~\LINK{http://naviteri.org/2020/02/some-words-for-leap-year-day/}{b} \\
\textsc{Derivatives}: \ding{43} \on{1}.~\HL{tIla'um}{\B{tì\-la\-'um}}, pretence.

% lafyon
\hi 
\HT{lafyon}{\B{la\cp{fyon}}} \I{[la."fjon]} \emph{adj., ofp.} wise. 
\B{eyk\-tan la\-fyon}, a wise leader.~\LINK{http://naviteri.org/2018/03/100a-liu-amip-64-new-words-part-1/}{b}
(\ding{43}~\HL{hafyon}{\B{hafyon}}, \HL{hafyonga'}{\B{hafyonga'}})

% lahe
\hi
\HT{lahe}{\B{\cp{la}he}} \I{["la.hE]} \textsc{i}.~\emph{adj.} other, else. 
\B{Fì\-txon na ton a\-la\-he nì\-wotx pe\-lun ke lu teng?} Why isn't this night the same like all the others?~\LINK{http://whyisthisnight.com/addl-more.html\#na\%27vi}{tn}
\B{Lu nga sì la\-he\-a ay\-ha\-pxì\-tu lì'\-fya\-o\-lo'\-ä aw\-nge\-yä ay\-aw\-po a ma\-wey\-pe\-rey.} It's you and the other members of the sprachbund who are being patient.~\LINK{http://forum.learnnavi.org/language-updates/txelanit-hivawl/}{ln}
\B{Lu law 'uo a\-la\-he, ma ey\-lan.} Something else is clear to me, my friends.~\LINK{http://naviteri.org/2010/06/first-post/}{b}
(\emph{c. idiom}) \B{slä ha\-ya\-lo a\-la\-he} ``\ldots but that is for another time'' (\emph{set phrase in storytelling})~\LINK{http://naviteri.org/2020/04/aawa-u-amip-a-few-new-things/}{b} \\
\textsc{ii}.~\emph{pn.} 
\B{tsa\-kem rä\-'ä si\-vi ay\-la\-he\-ru}, don't do that action to others.~\LINK{http://wiki.learnnavi.org/index.php?title=Corpus\#Good_Morning_America}{wiki} 
\B{Oel po\-ti tspì\-yang fte tì\-ke\-nong lì\-ye\-vu ay\-la\-he\-ru\slash aylaru.} I will kill him as an example to the others.~\LINK{http://forum.learnnavi.org/language-updates/pseudo-canon-from-bd-live-content/}{ln} \\
\textsc{Derivatives}: \ding{43} \on{1}.~\HL{lapo}{\B{la\-po}}, another (one). 
\on{2}.~\HL{saylahe}{\B{say\-la\-he}}, et cetera.

% lak
\hi 
\HT{lak}{\B{lak}} \I{[lak\textcorner]} \emph{n.} shell, \emph{hard covering of a plant or animal}. 
(\ding{43}~\HL{sum}{\B{sum}}) \\
\textsc{Derivatives}: \ding{43} \on{1}.~\HL{tslikxyutsawlak}{\B{tslikxyu tsawlak}}, scarab crawler.

% lal
\hi
\HT{lal}{\B{lal}} \I{[lal]}  \emph{adj.} old (\emph{not for people}).
\B{lal\-a sä\-fpìl}, an old idea.~\LINK{http://naviteri.org/2015/08/aylifyavi-lereyfya-2-cultural-terms-2-and-more/}{b} 
\B{Hum zì\-sìt a\-lal, pä\-hem pum a\-mip. Yo'\-ko\-fya a\-tì\-'i\-lu\-ke.} The old year leaves, the new one arrives. An endless cycle.~\LINK{http://naviteri.org/2017/12/zisit-amip-lefpom-ma-eylan-happy-new-year-friends/}{b}
(\ding{214}~\HL{mip}{\B{mip}}) \\
\textsc{Derivatives}: \ding{43} \on{1}.~\HL{txanlal}{\B{txan\-lal}}, ancient, very old.

% lam
\hi
\HT{lam}{\B{lam}} \I{[lam]}  \emph{vin.} seem, appear. (\emph{a.}) 
\B{fay\-he\-tu\-wong fmi Na'\-vi\-na li\-vam}, these aliens try to look like Na'vi.~\LINK{http://wiki.learnnavi.org/Na\%27vi_from_Avatar_Movie}{wiki\slash a\textsuperscript{c}}
\B{Poe\-ri unil\-tì\-ran\-tokx\-it tar\-mok a krr lam stum nì\-ay\-fo.} When she was in her avatar body she almost appear like them.~\LINK{http://naviteri.org/2010/07/ting-mikyun-fte-tslivam-listening-comprehension-1-3/}{b}
(\emph{b. idiom}) \B{lam oer (fwa)}, it seems\slash appears to me (that), seemingly.~\LINK{http://forum.learnnavi.org/language-updates/txelanit-hivawl/}{ln}
\B{lam fwa}, it seems that. \B{Lam ngay oer, fo ka\-yä ìlä hil\-van tup na'\-rìng.} I bet they'll take the river rather than the forest.~\LINK{http://forum.learnnavi.org/language-updates/concerning-tup/}{ln} \\  
\textsc{Derivatives}: \ding{43} \on{1}.~\HL{tIlam}{\B{tì\-lam}}, appearance.   
\on{2}.~\HL{nIlam}{\B{nì\-lam}}, apparently. 

% lamaytxa
\hi 
\HT{lamaytxa}{\B{lamay\cp{txa}}} \I{[la.maj."t'a]} \emph{n.} flood, powerful gathering of water. 
\B{La\-may\-txal a\-txan pxay\-a kel\-kut sko\-la\-'ä\-nga.} The great flood sadly destroyed many homes.~\LINK{http://naviteri.org/2018/08/fmawnti-stolawm-srak-have-you-heard-the-news/}{b}

% lan
\hi 
\HT{lan}{\B{lan}} \I{[lan]} \emph{n.} resin. \\
\textsc{Derivatives}: \ding{43} \on{1}.~\HL{lanutral}{\B{la\-nut\-ral}}, dandetiger.

% lanay'ka
\hi 
\HT{lanay'ka}{\B{la\cp{nay'}ka}} \I{[la."najP.ka]} \emph{n., fauna} slinger, \emph{acediacutus xenoterribili}. 
(\emph{a. idiom}) \B{'Aw\-si\-teng lu me\-fo la\-nay'\-ka.} \emph{praising how well two people work together and complement each other} [lit.: they are a slinger (together)].~\LINK{http://naviteri.org/2021/08/ulte-ayyoratu-leiu-and-the-winners-are-2/}{b}
\B{To\-la\-ron me\-fol me\-sa\-li\-o\-a\-ngit! Tew\-ti, 'aw\-si\-teng lu me\-fo la\-nay'\-ka.} They hunted two sturmbeest? Wow, they work very well together.~\LINK{http://naviteri.org/2021/08/ulte-ayyoratu-leiu-and-the-winners-are-2/}{b}
(\emph{b. idiom}) \B{(na) la\-nay'\-ka lu\-ke re\-'o}, ``completely lost,'' [lit.: Like a slinger without a head].~\LINK{http://naviteri.org/2021/08/ulte-ayyoratu-leiu-and-the-winners-are-2/}{b}
\B{Po maw kxitx mun\-txa\-tu\-ä 'a\-me\-fu na la\-nay'\-ka re'\-o\-lu\-ke.} After the death of his spouse, he felt completely lost.~\LINK{http://naviteri.org/2021/08/ulte-ayyoratu-leiu-and-the-winners-are-2/}{b}

% lante
\hi 
\HT{lante}{\B{\cp{lan}te}} \I{["lan.tE]} \emph{vin.} wander. 
\B{Fo ka na'\-rìng le\-ran\-te teng\-krr syu\-vet fwe\-rew.} They're wandering across the forest looking for food.~\LINK{http://naviteri.org/2023/05/reef-navi-part-2-morphology-syntax-lexicon-and-more/}{b}

% lanutral
\hi 
\HT{lanutral}{\B{\cp{la}nutral}} \I{["la.nut\textcorner.Ral]} \emph{or} \I{[.nUt\textcorner.]} \emph{n.} \ding{96} dandetiger, \emph{candea inflata}.~\LINK{https://forum.learnnavi.org/language-updates/expansion-to-flora-fauna-and-places-from-the-valley-of-mo'ara/}{ln} 
(\ding{43}~\HL{lan}{\B{lan}}, \HL{utral}{\B{ut\-ral}})

% lang
\hi
\HT{lang}{\B{lang}} \I{[laN]} \emph{vtr.} investigate, explore. 
\B{Lum\-pe le\-rang Kel\-ut\-ral\-it Saw\-tu\-tel?} Why are the humans exploring Hometree?~\LINK{http://naviteri.org/2014/11/twenty-before-the-holidays/}{b} \\
\textsc{Derivatives}: \ding{43} \on{1}.~\HL{tIlang}{\B{tì\-lang}}, exploration. 
\on{2}.~\HL{sAlang}{\B{sä\-lang}}, exploration, investigation.

% lapo
\hi
\HT{lapo}{\B{\cp{la}po}} \I{["la.po]} (\emph{gen.}~\B{lapeyä}) \emph{pn.} another one (person or thing). 
(\ding{43}~\HL{lahe}{\B{la\-he}}, \HL{po}{\B{po}}) 

% lapx
\hi 
\HT{lapx}{\B{lapx}} \I{[lap']} \emph{vtr.} regret s.t. 
\B{Kem\-it a oe so\-li oel lä\-ngapx.} I regret what I did.~\LINK{http://naviteri.org/2022/12/trr-anawm-polaheiem-the-great-day-has-arrived/}{b}  \\
\textsc{Derivatives}: \ding{43} \on{1}.~\HL{tIlapx}{\B{tì\-lapx}}, regret. 

% lare
\hi 
\HT{lare}{\B{\cp{la}re}} \I{["la.RE]} \emph{vin.} be watchful, alert (\emph{staying alert and watchful, keenly aware of your surroundings, particularly in dangerous circumstances}). 
\B{Tsun na'\-rìng le\-hrr\-ap li\-vu, ma 'e\-veng. Ze\-ne li\-va\-re fra\-krr.} The forest can be dangerous, child. You need to be alert at all times.~\LINK{https://naviteri.org/2025/03/conlang-adventure-and-more/}{b} 
(\ding{43}~\HL{nari}{\B{na\-ri si}}) 

% laro / si
\hi
\HT{laro}{\B{\cp{la}ro}} \I{["la.Ro]} \textsc{i}. \emph{adj.} clean. \\ 
\textsc{ii}. \B{\cp{la}ro si} \I{["la.Ro si]} \emph{vin.} clean, make free of dirt.
\B{Txo ke li\-vu pay, tsun me\-syokx\-ur la\-ro si\-vi fa srä.} If water isn't available, you can clean your hands with a cloth.~\LINK{http://naviteri.org/2016/06/mrrvola-lifyavi-amip-forty-new-expressions/}{b} 
(\ding{214}~\HL{tsewtx}{\B{tsewtx}}; \ding{43}~\HL{yur}{\B{yur}})

% latem
\hi
\HT{latem}{\B{\cp{la}tem}} \I{["la.tEm]}  \emph{vin.} (\emph{a.}) change, i.e. s.t. changes.
\B{fra\-'u la\-tem}, everything changes.
\B{Ke\-zem\-pll\-txe, ta\-lun tì\-vi\-rä fì\-sä\-spxin\-ä, lo\-la\-tä\-ngem ki\-fkey, lo\-la\-tä\-ngem tì\-rey.} Needless to say, because of the spread of this illness, the world and life changes.~\LINK{http://naviteri.org/2020/03/teri-tifkeytok-lefkrr-about-the-present-situation/}{b}
(\emph{b.~with} \HL{IlA}{\B{ìlä}}) depend on. \B{Tì\-flä la\-tem ìlä seyng\-a ftxey fkol sä\-nu\-met li\-vek fu\-ke.} Success depends on whether or not one follows instructions.~\LINK{http://naviteri.org/2012/07/meetings-waterfalls-and-more/}{b}
(\ding{43}~\HL{leykatem}{\B{ley\-ka\-tem}}) \\
\textsc{Derivatives}: \ding{43} \on{1}.~\HL{tIlatem}{\B{tì\-la\-tem}}, change (abstract concept). 
\on{2}.~\HL{sAlatem}{\B{sä\-la\-tem}}, (instance of) change.

% latsi
\hi
\HT{latsi}{\B{la\cp{tsi}}} \I{[la."\t{ts}i]}  \emph{vin.} keep up with s.o. 
\B{Slä pll\-txe rä\-'ä san Oe ke tsun li\-va\-tsi sìk!} But don't say, `I can't keep up!'~\LINK{http://naviteri.org/2020/05/aawa-liu-amip-si-vurway-alor-a-few-new-words-and-a-beautiful-poem/\#comment-31193}{b}

% law / si
\hi
\HT{law}{\B{law}} \I{[law]} \textsc{i}.~\emph{adj.} clear, certain. (\emph{a.}) \B{Pì\-lok\-ä fì\-ha\-pxì\-yä tì\-kan lu law.} The aim of this part of the blog is clear.~\LINK{http://naviteri.org/2010/06/first-post/}{b}
\B{Poe pol\-txe nì\-fya\-'o a\-law.} She spoke clearly [lit.: in a clear manner].~\LINK{http://forum.learnnavi.org/language-updates/why-is-this-night/msg154875/\#msg154875}{ln}
(\emph{b.~with the dative}) be clear to one. \B{Pe\-yä 'i\-tan\-ìri lu ho\-na nì\-txan a fì\-'u law lu fra\-por.} It's clear to everyone that his son is very cute.~\LINK{http://naviteri.org/2012/01/more-additions-to-the-lexicon/}{b} 
\B{Nge\-yä Ti\-re\-a\-i\-o\-ang zo\-la\-'u a krr, law la\-yu nga\-ru.} When your Spirit Animal comes, you will know [lit.: it will be clear to you].~\LINK{http://forum.learnnavi.org/language-updates/auxilary-verb-si-possessive-dative-krr/}{ln} \\
\textsc{ii}.~\HT{law si}{\B{law si}} \I{[law si]}  \emph{vin.} make clear (\emph{allows for double dative}).   
\B{Tsun nga law si\-vi nì\-'it srak?} Could you make that a bit clearer?~\LINK{http://naviteri.org/2010/09/getting-to-know-you-part-2/}{b} 
\B{Ra\-lur law so\-li fo oe\-ru.} They made the meaning clear to me.~\LINK{http://naviteri.org/2013/06/looking-back-looking-forward/}{b} \\  
\textsc{Derivatives}: \ding{43} \on{1}.~\HL{nIlaw}{\B{nì\-law}}, clearly.

% lawk
\hi
\HT{lawk}{\B{lawk}} \I{[lawk\textcorner]}  \emph{vtr.} discourse on, talk about, say s.t. concerning.   
\B{Txan\-krr ngal ke lawk pot kaw\-'it.} You haven't mentioned a thing about him in a long time.~\LINK{http://naviteri.org/2011/09/\%E2\%80\%9Cby-the-way-what-are-you-reading\%E2\%80\%9D/}{b}
\B{Po\-el oe\-ti lar\-mawk.} She was talking about me.~\LINK{http://naviteri.org/2010/09/getting-to-know-you-part-1/}{b}
(\emph{b. in the reflexive} \ding{43}~\HL{lApawk}{\B{lä\-pawk}}) introduce oneself. 

% lawnol
\hi
\HT{lawnol}{\B{\cp{law}nol}} \I{["law.nol]} \emph{n.} great joy.
\B{Ay\-zì\-sìt a\-zu\-sa\-'u ta law\-nol te\-ya li\-vu nì\-wotx.} May the coming years be full of great joy.~\LINK{http://naviteri.org/2014/02/barter-and-exchange/}{b} \\
\textsc{Derivatives}: \ding{43} \on{1}.~\HL{'onglawn}{\B{'ong\-lawn}}, \emph{exhiliration of first bonding}.
\on{2}.~\HL{lawnoltsim}{\B{law\-nol\-tsim}}, source of (great) joy.

% lawnoltsim
\hi 
\HT{lawnoltsim}{\B{\cp{law}noltsim}} \I{["law.nol.\t{ts}im]} \emph{or coll.} \I{["law.no.\t{ts}im]} \emph{n.} source of (great) joy. 
\B{Oe\-ri leiu fì\-lì'\-fya\-olo' law\-nol\-tsim.} This language community is a source of great joy to me.~\LINK{naviteri.org/2022/05/results-of-the-tttc/}{b} 
(\ding{43}~\HL{lawnol}{\B{law\-nol}}, \HL{tsim}{\B{tsim}})

% lawr
\hi
\HT{lawr}{\B{lawr}} \I{[lawR]} \textsc{i}.~\emph{n.} tune, melody. \\
\textsc{ii}.~\HT{tIng lawr}{\B{tìng lawr}} \I{[tIN lawR]} \emph{vin.} sing wordlessly, give out a tune or melody.
\B{'E\-van\-ìl alo a\-'aw\-ve nì\-'aw\-tu na'\-rìng\-it tar\-mok, ha to\-lìng lawr nì\-fwe\-fwi fte\-ke txo\-pu si\-vi.} The boy was alone in the forest for the first time, so he whistled a tune to calm his fears.~\LINK{http://naviteri.org/2011/09/miscellaneous-vocabulary/}{b}

% layl
\hi 
\HT{layl}{\B{layl}} \I{[lajl]} \emph{adj.} innocent. 
\B{Tsa\-kawng\-kem\-ìri lu oe layl!} I am innocent of that crime!~\LINK{http://naviteri.org/2023/09/aawa-ayliu-si-aylifyavi-amip-a-few-new-words-and-expressions/}{b}
(\ding{214}~\HL{tokat}{\B{to\-kat}}) \\
\textsc{Derivatives}: \ding{43} \on{1}.~\HL{tIlayl}{\B{tì\-layl}}, innocence. 
\on{2}.~\HL{nIlayl}{\B{nì\-layl}}, innocently. 

% layompin
\hi
\HT{layompin}{\B{la\cp{yom}pin}} \I{[la."jom.pin]}  \emph{n.} the color \ding{43}~\HL{layon}{\B{la\-yon}}. 
(\ding{43}~\HL{'opin}{\B{'opin}})

% layon
\hi
\HT{layon}{\B{la\cp{yon}}} \I{[la."jon]}  \emph{adj.} black.
\B{mrr\-a ya\-yo a\-tsawl sì la\-yon\slash mrr\-a ya\-yo a lu tsawl sì la\-yon}, five big, black birds.~\LINK{http://naviteri.org/2021/02/aysipawm-si-aysieyng-questions-and-answers/\#comment-32798}{b}
\B{tsawl\-a ya\-yo a\-la\-yon}, a big, black bird.~\LINK{http://naviteri.org/2021/02/aysipawm-si-aysieyng-questions-and-answers/\#comment-32798}{b}
(\ding{214}~\HL{teyr}{\B{teyr}}) \\
\textsc{Derivatives}: \ding{43} \on{1}.~\HL{layompin}{\B{la\-yom\-pin}}, (the color) black.

% layro
\hi 
\HT{layro}{\B{\cp{lay}ro}} \I{["laj.Ro]} \emph{adj.} free (\emph{from oppression}). \\
\textsc{Derivatives}: \ding{43} \on{1}.~\HL{tIlayro}{\B{tì\-lay\-ro}}, freedom. 

% läpawk
\hi
\HT{lApawk}{\B{lä\cp{pawk}}} \I{[l\ae."p$\cdot$awk\textcorner]} \emph{vin., inf.}\on{22} introduce oneself, talk about oneself.
\B{Ru\-txe lä\-pi\-vawk nì\-'it.} Please talk a little bit about yourself.~\LINK{http://naviteri.org/2010/09/getting-to-know-you-part-1/}{b}
\B{Po lä\-po\-lawk.} He talked about himself.~\LINK{http://naviteri.org/2010/09/getting-to-know-you-part-1/}{b}
\B{Nga lä\-pi\-vawk nì\-no ko.} Tell me all about yourself.~\LINK{http://naviteri.org/2010/09/getting-to-know-you-part-2/}{b}
(\emph{reflexive of} \ding{43}~\HL{lawk}{\B{lawk}})

% lätxayn
\hi
\HT{lAtxayn}{\B{lä\cp{txayn}}} \I{[l\ae."t'ajn]}  \emph{vtr.} defeat in battle, conquer.
\B{Fwa ay\-ioang a\-pxay fì\-txan Na'\-vi\-ru srung so\-li fte Saw\-tu\-tet li\-vä\-txayn lu pa\-rul nì\-ngay.} That so many animals helped the Na’vi defeat the Sky People was a genuine miracle.~\LINK{http://naviteri.org/2018/03/100a-liu-amip-64-new-words-part-1/}{b}
\B{Nì\-trr\-trr pxì\-mun\-'i sam\-si\-yul ay\-swi\-zaw\-it ku\-tu\-ä a\-law\-nä\-txayn sno\-kip nì\-'eng.} Warriors typically share the arrows of their defeated enemies among themselves.~\LINK{http://naviteri.org/2012/01/more-additions-to-the-lexicon/}{b}  \\  
\textsc{Derivatives}: \ding{43} \on{1}.~\HL{sAlAtxayn}{\B{sä\-lä\-txayn}}, defeat.

% le-
\hi
\HT{le}{\B{le-}} \I{[lE.]} \emph{unprod. pref. to form adj.s with unpredictable meaning}.

% le'aw
\hi
\HT{le'aw}{\B{le\cp{'aw}}} \I{[lE."Paw]}  \emph{adj.} only, sole.   
\B{Lu tsa\-tsam\-si\-yu le\-'aw\-a ha\-pxì\-tu tsam\-po\-ngu\-ä a mal lu moer.} That warrior is the only member of the war party that we both trust.~\LINK{http://naviteri.org/2012/01/more-additions-to-the-lexicon/}{b}
(\ding{43}~\HL{'aw}{\B{'aw}})

% le'awtu
\hi 
\HT{le'awtu}{\B{le\cp{'aw}tu}} \I{[lE."Paw.tu]} \emph{adj.} (\emph{a. neutral connotation}) alone, on one's own, lone, by oneself, solitary. 
\B{Oe 'e\-fu le\-'aw\-tu.} I feel alone.~\LINK{http://naviteri.org/2018/04/100a-liu-amip-64-new-words-part-2/}{b}
\B{Oe lu le\-'aw\-a tu\-te a tsun srung si\-vi, ul\-te 'e\-fu le\-'aw\-tu nì\-ngay.} I'm the only one who can help, and I feel really alone.~\LINK{http://naviteri.org/2018/04/100a-liu-amip-64-new-words-part-2/}{b} 
\B{Le\-'aw\-tu\-a ta\-li\-oang\-ìri lu ki\-fkey tse\-nge le\-hrr\-ap.} The world is a dangerous place for a lone sturmbeest.~\LINK{http://naviteri.org/2018/04/100a-liu-amip-64-new-words-part-2/}{b}
(\emph{b. often with} \ding{43} ‹\HL{Ang}{\B{äng}}›) lonely.
\B{Oe 'e\-fä\-ngu le\-'aw\-tu.} I feel lonely.~\LINK{http://naviteri.org/2018/04/100a-liu-amip-64-new-words-part-2/}{b}
\B{Oe lu le\-'aw\-a tu\-te a tsun srung si\-vi, ul\-te 'e\-fu\slash 'e\-fä\-ngu le\-'aw\-tu nì\-ngay.} I'm the only one who can help, and I feel really lonely.~\LINK{http://naviteri.org/2018/04/100a-liu-amip-64-new-words-part-2/}{b} 

% le'al
\hi 
\HT{le'al}{\B{le\cp{'al}}} \I{[lE."Pal]} \emph{n.} wasteful (\emph{not for people}).
\B{Fwa sar pay\-it fì\-txan lu le\-'al.} Using this much water is wasteful.~\LINK{http://naviteri.org/2014/07/tsawlultxamaw-a-ayliu-post-meetup-words/}{b} 
(\ding{43}~\HL{'al}{\B{'al}})

% le'en
\hi
\HT{le'en}{\B{le\cp{'en}}} \I{[lE."PEn]}  \emph{adj.} speculative, intuitive (\emph{of an action, not a person}).
\B{Hem le\-'en tsun le\-hrrap li\-vu.} Speculative moves can be dangerous.~\LINK{http://naviteri.org/2011/01/new-year-new-vocabulary/}{b}
(\ding{43}~\HL{'en}{\B{'en}})

% le'Ìnglìsì
\hi
\HT{le'InglIsI}{\B{le\cp{'Ìng}lìsì}} \I{[lE."PIN.lI.sI]} \emph{adj.}~\textsc{lw} English, regarding English. 
\B{Kxawm tsa\-lì\-'u le\-'Ìng\-lì\-sì alu ``No\"el'' mì ron\-sem lar\-mu.} Maybe I had the English word ``No\"el'' in my mind.~\LINK{http://naviteri.org/2011/01/new-year-new-vocabulary/comment-page-1/\#comment-499}{b}
(\ding{43}~\HL{'InglIsI}{\B{'Ìng\-lì\-sì}})

% lefkeytongay
\hi 
\HT{lefkeytongay}{\B{lefkeyto\cp{ngay}}} \I{[lE.fkEj.to."Naj]} \emph{adj.} real. 
\B{Yu\-ne oet! Ke lu fì\-vrr\-tep tu\-te le\-fkey\-to\-ngay!} Listen to me! This demon is not a real person!~\LINK{http://naviteri.org/2019/08/aawa-lifya-si-lifyavi-amip-a-few-new-words-and-expressions/}{b} 
(\ding{43}~\HL{tIfkeytongay}{\B{tì\-fkey\-to\-ngay}})

% lefkrr
\hi
\HT{lefkrr}{\B{le\cp{fkrr}}} \I{[lE."fk\s{r}]}  \emph{adj.} current. 
\B{Tì\-fkey\-tok le\-fkrr le\-hrr\-ap lu nì\-txan.} The current situation is very dangerous.~\LINK{http://naviteri.org/2011/04/yafkeykiri-plltxe-frapo-everyone-talks-about-the-weather/}{b}
\B{… fu\-la tsun le\-fkrr\-a sì\-len\-it tswi\-va' hì\-krr (…) oe\-ti 'ey\-ke\-fu nit\-ram nì\-txan.} … that we can forget current events for a short time, makes me very happy.~\LINK{http://naviteri.org/2022/05/results-of-the-tttc/}{b}
(\ding{43}~\HL{krr}{\B{krr}})

% lefnelan
\hi 
\HT{lefnelan}{\B{le\cp{fne}lan}} \I{[lE."fnE.lan]} \emph{adj.} male. 
(\ding{43}~\HL{fnelan}{\B{fne\-lan}})

% lefnele
\hi 
\HT{lefnele}{\B{le\cp{fne}le}} \I{[lE."fnE.lE]} \emph{adj.} female. 
(\ding{43}~\HL{fnele}{\B{fne\-le}})

% lefngap
\hi
\HT{lefngap}{\B{le\cp{fngap}}} \I{[lE."fNap\textcorner]}  \emph{adj.} metallic. \B{el\-tu le\-fngap}, computer, \textsc{pc} [lit.: metallic brain].~\LINK{http://forum.learnnavi.org/language-updates/frommerian-words-from-nytimes-6442/}{ln}
\B{Kì\-ma\-nom oel mip\-a el\-tut le\-fngap yoa ew\-ro.} I just bought a new computer.~\LINK{http://naviteri.org/2014/02/barter-and-exchange/}{b}
(\ding{43}~\HL{fngap}{\B{fngap}})

% lefpom
\hi
\HT{lefpom}{\B{le\cp{fpom}}} \I{[lE."fpom]}  \emph{adj., nfp.} happy, peaceful, joyous, pleasant. 
\B{trr le\-fpom}, good day, good morning (cf.~French \emph{bon jour}).
\B{Trr le\-fpom, ma Su\-te Ame\-ri\-ka\-yä.} Good morning, America.~\LINK{http://wiki.learnnavi.org/index.php?title=Corpus\#Good_Morning_America}{wiki}
\B{Txon le\-fpom (li\-vu ngar).} (May you have a) good night. 
\B{Ay\-ftxo\-zä Le\-fpom!}~Happy Holidays!
\B{Mipa Zìsìt Lefpom, ma frapo!}~Happy New Year, y'all!~\LINK{http://naviteri.org/2011/01/new-year-new-vocabulary/}{b}
\B{Txep\-ìl tìng le\-fpom\-a sä\-nrr\-ti.} The fire gives a pleasant glow.~\LINK{http://naviteri.org/2012/07/fivospxiya-aylifyavi-amip-this-months-new-expressions-pt-2/}{b}
(\ding{43}~\HL{fpom}{\B{fpom}}; \HL{nitram}{\B{nit\-ram}})

% lefpomron
\hi
\HT{lefpomron}{\B{lefpom\cp{ron}}} \I{[lE.fpom."Ron]} \emph{adj., ofp.} (mentally) healthy.
(\ding{43}~\HL{fpomron}{\B{fpom\-ron}}, \ding{214}~\HL{kelfpomron}{\B{kel\-fpom\-ron}})

% lefpomtokx
\hi
\HT{lefpomtokx}{\B{lefpom\cp{tokx}}} \I{[lE.fpom."tok']}  \emph{adj., ofp.} (physically) healthy.
\B{Nga\-ri ftxì vi\-var le\-fpom\-tokx li\-vu, ma tsmuk!} May your tongue remain healthy, sibling!~\LINK{http://naviteri.org/2014/10/tson-si-fpomron-obligation-and-mental-health/comment-page-1/\#comment-10624}{b}
(\ding{43}~\HL{fpomtokx}{\B{fpom\-tokx}}, \ding{214}~\HL{kelfpomtokx}{\B{kel\-fpom\-tokx}})

% lefrir
\hi 
\HT{lefrir}{\B{le\cp{frir}}} \I{[lE."fRiR]} \emph{adj.} layered. 
(\ding{43}~\HL{frir}{\B{frir}})

% leftxozä
\hi 
\HT{leftxozA}{\B{leftxo\cp{zä}}} \I{[lE.ft'o."z\ae]} \emph{adj.} celebratory. 
\B{Ay\-nge\-yä lor\-a fì\-'u\-pxa\-re le\-ftxo\-zä oe\-ru te\-ya so\-li nì\-ngay.} Your beautiful, celebratory message truly filled me with joy.~\LINK{https://forum.learnnavi.org/language-updates/new-words-leftxoza-(celebratory)-poston-(boston)-kelni-(koln-cologne)/}{ln} 
(\ding{43}~\HL{ftxozA}{\B{ftxo\-zä}})

% leha'
\hi 
\HT{leha'}{\B{le\cp{ha'}}} \I{[lE."haP]} \emph{adj.} appropriate, suitable, fitting. 
\B{Fo\-ri tsa\-fne\-tì\-ku\-sar ke lu le\-ha'.} That kind of teaching isn't appropriate for them.~\LINK{http://naviteri.org/2020/11/vospxivopeya-ayliu-amip-novembers-new-words/}{b} 
(\ding{214}~\HL{kelha'}{\B{kel\-ha'}} \ding{43}~\HL{ha'}{\B{ha'}})

% lehangham
\hi 
\HT{lehangham}{\B{le\cp{hang}ham}} \I{[lE."haN.ham]} \emph{adj.} laughable, ridiculous, absurd. 
\B{Fo\-le\-nge fra\-pol pe\-yä tì\-hawl\-it le\-hang\-ham.} Everyone ridiculed his ridiculous plan.~\LINK{naviteri.org/2025/07/ayu-leketengsyon-various-things/}{b}  
(\ding{43}~\HL{hangham}{\B{hang\-ham}})

% lehawmpam
\hi
\HT{lehawmpam}{\B{le\cp{hawm}pam}} \I{[lE."hawm.pam]}  \emph{adj.} noisy. 
\B{Ta\-ron\-yul le\-hawm\-pam ska\-'a sä\-ta\-ron\-it.} A noisy hunter destroys the hunt.~\LINK{http://naviteri.org/2014/11/twenty-before-the-holidays/}{b}
(\ding{43}~\HL{hawmpam}{\B{hawmpam}})

% lehawng
\hi
\HT{lehawng}{\B{le\cp{hawng}}} \I{[lE."hawN]}  \emph{adj.} excessive. 
\B{Pxìm lu tì\-new le\-hawng kxu\-tu fpom\-ä.} Excessive desire is often the enemy of peace.~\LINK{http://naviteri.org/2014/11/twenty-before-the-holidays/}{b}
(\ding{43}~\HL{hawng}{\B{hawng}})

% lehawngkrr
\hi
\HT{lehawngkrr}{\B{le\cp{hawng}krr}} \I{[lE."hawN.k\s{r}]}  \emph{adj.} late. 
\B{tì\-pä\-hem le\-hawng\-krr}, a late arrival.
(\ding{43}~\HL{hawngkrr}{\B{hawng\-krr}}, \ding{214}~\HL{leye'krr}{\B{le\-ye'\-krr}})

% leheng
\hi 
\HT{leheng}{\B{le\cp{heng}}} \I{[lE."hEN]} \emph{adj.} formally polite. \\
\textsc{Derivatives}: \ding{43} \on{1}.~\HL{lI'fyeng}{\B{lì'\-fyeng}}, honorific language. 

% lehìpey
\hi
\HT{lehIpey}{\B{le\cp{hì}pey}} \I{[lE."hI.pEj]} \emph{adj.} hesitant, in a state of hesitation.
\B{Ta\-ron\-yul le\-hì\-pey kan smar\-it nìl\-ke\-ftang slä ke ta\-kuk kaw\-krr.} A hesitant hunter will aim at a prey forever but never hit it.~\LINK{http://naviteri.org/2015/11/vomuna-liu-amip-ten-new-words/}{b}
(\ding{43}~\HL{hIpey}{\B{hì\-pey}})

% lehoan
\hi
\HT{lehoan}{\B{le\cp{ho}an}} \I{[lE."ho.an]} \emph{adj.} comfortable.
(\ding{43}~\HL{hoan}{\B{ho\-an}}, \ding{214}~\HL{kelhoan}{\B{kel\-ho\-an}})

% lehrrap
\hi
\HT{lehrrap}{\B{le\cp{hrr}ap}} \I{[lE."h\s{r}.ap\textcorner]}  \emph{adj.} dangerous. 
\B{Lu fo le\-hrr\-ap.} They are dangerous.~\LINK{http://wiki.learnnavi.org/index.php?title=Corpus\#New_York_Times}{wiki}
(\ding{43}~\HL{hrrap}{\B{hrrap}})

% leioae / si
\hi
\HT{leioae}{\B{leio\cp{a}e}} \I{[lE.i.o."{}a.E]} \textsc{i}.~\emph{n.} respect.   
\B{Lu\-ke le\-i\-o\-a\-e o\-lo'\-ä ke tsun kea eyk\-tan fli\-vä.} Without the respect of the clan, no leader can succeed.~\LINK{http://naviteri.org/2012/02/trr-asawnung-lefpom-happy-leap-day/}{b}
\B{le\-i\-o\-a\-e a\-mek,} feigned respect.~\LINK{http://naviteri.org/2012/02/trr-asawnung-lefpom-happy-leap-day/}{b} \\
\textsc{ii}.~\HT{leioae si}{\B{leio\cp{a}e si}} \I{[lE.i.o."{}a.E si]}  \emph{vin.} to respect (s.o.~with the dative). 
\B{Nga\-ru le\-i\-o\-a\-e si oe fra\-to, ma 'ey\-lan.} I respect you the most of all, friend.~\LINK{http://naviteri.org/2012/02/trr-asawnung-lefpom-happy-leap-day/}{b} \\
\textsc{Derivatives}: \ding{43} 
\on{1}.~\HL{leiokoaktu}{\B{le\-io\-ko\-ak\-tan}}, respected elder.

% leiokoaktan
\hi 
\HT{leiokoaktan}{\B{leio\cp{ko}aktan}} \I{[lE.i.o."ko.ak\textcorner.tan]} \emph{n} respected male elder. 
(\ding{214}~\HL{leiokoakte}{\B{le\-i\-o\-ko\-ak\-te}}, \ding{43}~\HL{leiokoaktu}{\B{le\-i\-o\-ko\-ak\-tu}})

% leiokoakte
\hi 
\HT{leiokoakte}{\B{leio\cp{ko}akte}} \I{[lE.i.o."ko.ak\textcorner.tE]} \emph{n.} respected female elder. 
\B{Nge\-nge\-yä ha\-fyon\-ìri ira\-yo, ma le\-i\-o\-ko\-ak\-te.} Thank you for your wisdom, respected female elder.~\LINK{https://naviteri.org/2025/03/conlang-adventure-and-more/}{b} 
(\ding{214}~\HL{leiokoaktan}{\B{le\-i\-o\-ko\-ak\-tan}}, \ding{43}~\HL{leiokoaktu}{\B{le\-i\-o\-ko\-ak\-tu}}) 

% leiokoaktu
\hi 
\HT{leiokoaktu}{\B{leio\cp{ko}aktu}} \I{[lE.i.o."ko.ak\textcorner.tu]} \emph{n.} respected elder. 
(\ding{43}~\HL{leioae}{\B{le\-i\-o\-a\-e}}, \HL{koaktu}{\B{ko\-ak\-tu}}) \\
\textsc{Derivatives}: \ding{43} 
\on{1}.~\HL{leiokoaktan}{\B{le\-io\-ko\-ak\-tan}}, respected male elder.
\on{2}.~\HL{leiokoakte}{\B{le\-i\-o\-ko\-ak\-te}}, respected female elder.

% lek
\hi
\HT{lek}{\B{lek}} \I{[lEk\textcorner]}  \emph{vtr.} heed, obey, follow (command, order, suggestion etc.). 
\B{Tì\-flä la\-tem ìlä seyng\-a ftxey fkol sä\-nu\-met li\-vek fu\-ke.} Success depends on whether or not one follows instructions.~\LINK{http://naviteri.org/2012/07/meetings-waterfalls-and-more/}{b}
\B{Ta\-fral ke lì\-syek oel nge\-yä ke\-ye\-'ung\-it.} Therefore I will not heed your insanity.~\LINK{http://forum.learnnavi.org/language-updates/ltasygt-with-ke-and-nga/}{ln\slash g} \\
\textsc{Derivatives}: \ding{43} \on{1}.~\HL{lek}{\B{ley\-kek}}, enforce.

% leketengsyon
\hi 
\HT{leketengsyon}{\B{le\cp{ke}tengsyon}} \I{[lE."kE.tEN.sjon]} \emph{adj.} diverse, various. 
\B{Pll\-txe Saw\-tu\-te fa ay\-lì'\-fya le\-ke\-teng\-syon.} Humans speak diverse (OR: a variety of) languages.~\LINK{https://naviteri.org/2025/06/vospikin-lefpom-happy-july/}{b} 
\B{Ay\-u le\-ke\-teng\-syon.} Various things.~\LINK{https://naviteri.org/2025/07/ayu-leketengsyon-various-things/}{b}
(\ding{43}~\HL{ketengsyon}{\B{ke\-teng\-syon}}) 

% lekin
\hi
\HT{lekin}{\B{le\cp{kin}}} \I{[lE."kin]}  \emph{adj.} necessary.
\B{Le\-kin lu tì\-txur, lu tì\-tstew.~Slä le\-tsran\-ten fra\-to lu tì\-mal.} Strength and courage are necessary.~But most important of all is trustworthiness.~\LINK{http://naviteri.org/2012/01/more-additions-to-the-lexicon/}{b}
(\ding{43}~\HL{kin}{\B{kin}}; \ding{214}~\HL{kelkin}{\B{kel\-kin}})

% lekoren
\hi
\HT{lekoren}{\B{leko\cp{ren}}} \I{[lE.ko."REn]}  \emph{adj.} lawlike, regarding rules.
\B{päng\-kxo\-yu le\-ko\-ren}, lawyer, one who discusses rules.~\LINK{http://forum.learnnavi.org/language-updates/frommerian-words-from-nytimes-6442/}{ln}
(\ding{43}~\HL{koren}{\B{ko\-ren}})

% lekve'o
\hi 
\HT{lekve'o}{\B{lek\cp{ve}'o}} \I{[lEk\textcorner."vE.Po]} \emph{adj.} chaotic. 
\B{Lam oe\-ru ki\-fkey lek\-ve\-'o.} The world seems chaotic to me.~\LINK{https://naviteri.org/2024/09/pxevola-liu-amip-two-dozen-new-words/}{b} 
(\ding{43}~\HL{ve'o}{\B{ve\-'o}}, \HL{velke}{\B{vel\-ke}}) 

% lekye'ung
\hi
\HT{lekye'ung}{\B{lek\cp{ye}'ung}} \I{[lEk\textcorner."jE.PuN]} \emph{or} \I{[.PUN]} (\textsc{rn}: \B{lek\cp{ye}ùng} \I{[lEk\textcorner.UN]}) \emph{adj.} insane, crazy. 
\B{Fì\-sä\-'an\-lal Ne\-we\-yä nga\-ti sley\-ka\-yu lek\-ye\-'ung!} This yearning for Newey is going to drive you crazy.~\LINK{http://naviteri.org/2013/01/awvea-posti-zisita-amip-first-post-of-the-new-year/}{b}
\B{Krr\-ka lek\-ye\-'ung\-a fay\-srr, sìl\-pey oe, li\-vu ay\-nga\-ru fpom nì\-wotx.} During these crazy days, I hope you are all well.~\LINK{http://naviteri.org/2020/11/vospxivopeya-ayliu-amip-novembers-new-words/}{b}
(\ding{43}~\HL{keye'ung}{\B{ke\-ye\-'ung}})

% lekxan
\hi 
\HT{lekxan}{\B{le\cp{kxan}}} \I{[lE."k'an]} \emph{adj.} (\emph{a. literally}) blocked, obstructed.
\B{Lu fì\-fya\-'o le\-kxan; ke tsun aw\-nga si\-va\-lew.} The path is obstructed; we can't proceed further.~\LINK{http://naviteri.org/2020/05/aawa-liu-amip-si-vurway-alor-a-few-new-words-and-a-beautiful-poem/}{b}
(\emph{b. metaphorically, with} \ding{43}~\HL{'efu}{\B{'efu}}) frustrated. 
\B{Oe pll\-txe, po ke tìng mik\-yun. 'E\-fu oe le\-kxan nì\-txan.} I talk, but he doesn't listen. I feel very frustrated.~\LINK{http://naviteri.org/2020/05/aawa-liu-amip-si-vurway-alor-a-few-new-words-and-a-beautiful-poem/}{b} 
(\ding{43}~\HL{ekxan}{\B{ekxan}}; \HL{tIkankxang'}{\B{tì\-kan\-kxa\-nga'}})

% lekxutslu
\hi 
\HT{lekxutslu}{\B{le\cp{kxu}tslu}} \I{[lE."k'u.\t{ts}lu]} \emph{adj.} risky. 
\B{Aw\-nga zen\-ke fì\-kem si\-vi. Lu le\-kxu\-tslu nì\-hawng.} We mustn't do this. It's too risky.~\LINK{http://naviteri.org/2020/04/aawa-u-amip-a-few-new-things/}{b}  
(\ding{43}~\HL{kxutslu}{\B{kxutslu}})

% lelewng
\hi 
\HT{lelewng}{\B{le\cp{lewng}}} \I{[lE."lEwN]} \emph{adj., ofp.} shameful. 
\B{tu\-te le\-lewng}, a shameful person.~\LINK{http://naviteri.org/2021/09/aawa-ayliu-amip-a-few-new-words/}{b}
(\ding{43}~\HL{lewng}{\B{lewng}}) 

% lelì'fya
\hi
\HT{lelI'fya}{\B{le\cp{lì'}fya}} \I{[lE."lIP.fja]}  \emph{adj.} concerning language.   
\B{tì\-päng\-kxo le\-lì'\-fya}, language discussion.~\LINK{http://naviteri.org/2010/06/first-post/}{b}
\B{sul\-fä\-tu le\-lì'\-fya}, language masters.~\LINK{http://naviteri.org/2010/08/a-na\%E2\%80\%99vi-alphabet/}{b}
(\ding{43}~\HL{lI'fya}{\B{lì'\-fya}})

% lemrey
\hi
\HT{lemrey}{\B{lem\cp{rey}}} \I{[lEm."Rej]}  \emph{adj.} surviving (e.g., of entities from a group some of whom have died). 
\B{Maw fwa fkol Kel\-ut\-ral\-it sko\-la\-'a, lem\-rey\-a ha\-pxì\-tul tsa\-soai\-ä txo\-lu\-la mip\-a kel\-ku\-ti mì ta\-yo.} After Hometree was destroyed, the surviving members of that family built a new home on the plains.~\LINK{http://naviteri.org/2011/05/some-miscellaneous-vocabulary/}{b}
(\ding{43}~\HL{emrey}{\B{em\-rey}})

% lemweypey
\hi
\HT{lemweypey}{\B{lem\cp{wey}pey}} \I{[lEm."wEj.pEj]}  \emph{adj.} patient.
\B{Za\-'u fra\-'u ne tu\-te lem\-wey\-pey, ke\-fyak?} Everything comes to the patient people, right?~\LINK{http://naviteri.org/2012/03/spring-vocabulary-part-2/}{b}
(\ding{43}~\HL{maweypey}{\B{ma\-wey\-pey}})

% len
\hi
\HT{len}{\B{len}} \I{[lEn]}  \emph{vin.} occur, happen. 
\B{Krro len ay\-hem a\-na\-fì\-'u.} Sometimes such actions happen.~\LINK{http://forum.learnnavi.org/language-updates/just-another-few-words/}{ln} 
\B{Krr a sto\-lawm Txe\-wì te\-ri hem a poe\-ru lo\-len, 'o\-le\-fu ke\-ftxo nì\-txan.} When Txewì heard about what happened to her, he felt very upset.~\LINK{http://naviteri.org/2010/07/ting-mikyun-fte-tslivam-listening-comprehension-1-3/}{b}
\B{Nari si fte kea fe\-kem ke li\-ven ngar!} Be careful you don't have an accident!~\LINK{http://naviteri.org/2011/07/txantsana-ultxa-mi-siatll-great-meeting-in-seattle/}{b} 
(\emph{a. idiom}) \B{Kem\-pe le\-ren?} `What's happenin'?' 
\B{Kxì, ma 'ey\-lan! Kem\-pe le\-ren?} Hey dude! What's happenin'?.~\LINK{http://naviteri.org/2018/02/negative-questions-in-navi/}{b} \\  
\textsc{Derivatives}: \ding{43} \on{1}.~\HL{tIlen}{\B{tì\-len}}, event. 
\on{2}.~\HL{ftanglen}{\B{ftang\-len}}, prevent.

% leNa'vi
\hi
\HT{lena'vi}{\B{le\cp{Na'}vi}} \I{[lE."naP.vi]} \emph{adj.} having to do with the Na'vi, regarding Na'vi. 
\B{Ay\-ey\-lan\-ur oe\-yä sì ey\-lan\-ur lì'\-fya\-yä le\-Na'\-vi nì\-wotx.} To my friends and all the frieands of the Na'vi language.~\LINK{http://forum.learnnavi.org/news-announcements/a-response-from-paul-frommer!/}{ln} 
\B{Sna\-fpìl\-fya\-ri le\-Na'\-vi krra smar\-it fkol tspang, tsran\-ten nì\-txan fwa po ke ngä\-'än nì\-kel\-kin.} It's important in Na'vi philosophy that the prey not suffer unnecessarily when it's killed.~\LINK{http://naviteri.org/2013/02/vospxi-ayol-posti-apup-short-post-for-a-short-month/}{b}
\B{Kin aw\-ngal le\-Na'\-vi\-a lì\-'ut a ral lu ``typo.''} We need a Na'vi word for `typo.'~\LINK{http://naviteri.org/2015/06/na-ayskxe-mi-telan-sad-news/comment-page-1/\#comment-20585}{b} 
(\ding{43}~\HL{na'vi}{\B{Na'vi}})

% lenalsteng
\hi 
\HT{lenalsteng}{\B{le\cp{nal}steng}} \I{[lE."nal.stEN]} \emph{adj., ofp.} compassionate. 
(\ding{43}~\HL{nalsteng}{\B{nal\-steng}}; \HL{tInalstenga'}{\B{tì\-nal\-ste\-nga'}})

% leno
\hi
\HT{leno}{\B{le\cp{no}}} \I{[lE."no]}  \emph{adj.} thorough, detail\hyp{}oriented (\emph{of a person}).
\B{Le\-no lu Lo\-ak nì\-txan.} Loak is very thorough.~\LINK{http://naviteri.org/2010/09/getting-to-know-you-part-2/}{b}

% lenomum
\hi
\HT{lenomum}{\B{le\cp{no}mum}} \I{[lE."no.mum]} \emph{or} \I{[.mUm]} (\textsc{rn}: \B{le\cp{no}mùm} \I{[lE."no.mUm]}) \emph{adj.} curious.
\B{Txo\-pu rä\-'ä si. Nga\-ti ke nìn pol nì\-kxap. Lu le\-no\-mum nì\-'aw.} Don't be scared. He's not looking at you threateningly. He's just curious.~\LINK{http://naviteri.org/2015/03/kap-si-ayunil-saylahe-threats-dreams-and-other-things/}{b}
(\ding{43}~\HL{newomum}{\B{new\-omum}})

% lenrra
\hi
\HT{lenrra}{\B{le\cp{nrr}a}} \I{[lE."n\s{r}.a]}  \emph{adj.} proud, \emph{feeling of pride one has in others}.
\B{Sa'\-nok le\-nrr\-a lrr\-tok so\-li krra prr\-nen alo a\-'aw\-ve pol\-txe.} The proud mother smiled when the baby spoke for the first time.~\LINK{http://naviteri.org/2012/07/fivospxiya-aylifyavi-amip-this-months-new-expressions-pt-2/}{b}
(\ding{43}~\HL{nrra}{\B{nrra}})

% lepay
\hi
\HT{lepay}{\B{le\cp{pay}}} \I{[lE."paj]}  \emph{adj.} watery, saturated. 
\B{'ak\-ra le\-pay}, watery, saturated soil.~\LINK{http://naviteri.org/2011/05/weather-part-2-and-a-bit-more-2/}{b}
(\ding{43}~\HL{pay}{\B{pay}}, \HL{paynga'}{\B{pay\-nga'}})

% lepwopx
\hi
\HT{lepwopx}{\B{lep\cp{wopx}}} \I{[lEp\textcorner."wop']}  \emph{adj.} cloudy.
\B{lep\-wopx nì\-hol}, lightly cloudy, just a few clouds.~\LINK{http://naviteri.org/2011/05/weather-part-2-and-a-bit-more-2/}{b}
\B{lep\-wopx nì\-pxay}, heavily cloud-covered, many clouds.~\LINK{http://naviteri.org/2011/05/weather-part-2-and-a-bit-more-2/}{b}
(\ding{43}~\HL{pIwopx}{\B{pì\-wopx}})

% lepxìmrun
\hi 
\HT{lepxImrun}{\B{le\cp{pxìm}run}} \I{[lE."p'Im.Run]} \emph{adj.} often\hyp{}found, common.
(\ding{43}~\HL{pxIm}{\B{pxìm}}, \HL{run}{\B{run}}; \ding{214}~\HL{kelpxImrun}{\B{kel\-pxìm\-run}})

% ler
\hi 
\HT{ler}{\B{ler}} \I{[lER]} \emph{adj.} steady, smooth (of motion).  \\
\textsc{Derivatives}: \ding{43} \on{1}.~\HL{nIler}{\B{nì\-ler}}, steadily.

% lereyfya
\hi
\HT{lereyfya}{\B{le\cp{rey}fya}} \I{[lE."REj.fja]}  \emph{adj.} cultural. 
\B{Ay\-lì'\-fya\-vi Le\-rey\-fya} Cultural Terms.~\LINK{http://naviteri.org/2015/07/aylifyavi-lereyfya-1-cultural-terms-1/}{b} 
(\ding{43}~\HL{reyfya}{\B{rey\-fya}}) 

% lerìk
\hi
\HT{lerIk}{\B{le\cp{rìk}}} \I{[lE."RIk\textcorner]}  \emph{adj.} leafy. 
\B{Aw\-ngal yom fkxen\-ti le\-rìk nì\-wotx.} We eat leafy vegetable.~\LINK{http://whyisthisnight.com/addl-more.html\#na\%27vi}{tn}
(\ding{43}~\HL{rIk}{\B{rìk}}) \\
\textsc{Derivatives}: \ding{43} \on{1}. \HL{yomyolerIk}{\B{yom\-yo le\-rìk}}, leaf plate.

% lerìn
\hi 
\HT{lerIn}{\B{le\cp{rìn}}} \I{[lE."RIn]} \emph{adj.} wooden, of wood. 
\B{tuk\-ru le\-rìn}, a spear made of wood.~\LINK{http://naviteri.org/2017/09/zisikrr-amip-ayliu-amip-new-words-for-the-new-season/}{b} 
(\ding{43}~\HL{rIn}{\B{rìn}})

% leronsrel
\hi
\HT{leronsrel}{\B{le\cp{ron}srel}} \I{[lE."Ron.sREl]} \emph{adj.} imaginary.
\B{Oe new si\-vop ne tsa\-ki\-fkey le\-ron\-srel.} I want to journey to that imaginary world.~\LINK{http://naviteri.org/2011/02/new-vocabulary-part-2/}{b}
(\ding{43}~\HL{ronsrel}{\B{ron\-srel}})

% lertu
\hi
\HT{lertu}{\B{\cp{ler}tu}} \I{["lER.tu]}  \emph{n.} colleague, co-worker.
\B{Nga\-ru oe\-yä ler\-tut.} \emph{or} \B{Oel mu\-wi\-vìn\-txu nga\-ru oe\-yä ler\-tut.} (\emph{formal}) Allow me to introduce my colleague.~\LINK{http://naviteri.org/2010/09/getting-to-know-you-part-1/}{b}

% lesar / si
\hi
\HT{lesar}{\B{le\cp{sar}}} \I{[lE."{}saR]} \textsc{i}.~\emph{adj.} useful.
\B{Tsuk\-yom\-a ioang lu le\-sar.} An edible animal is useful.~\LINK{http://naviteri.org/2011/03/\%E2\%80\%9Creceptive-ability\%E2\%80\%9D-and-hesitation/}{b}
\B{Kxa\-ri tstu\-tswo tsran\-ten krra ke lu kea sä\-fpìl le\-sar.} When one has no useful thoughts, the ability to close one's mouth is important.~\LINK{http://naviteri.org/2012/03/spring-vocabulary-part-2/}{b} \\
\textsc{ii}.~\B{le\cp{sar} si} \I{[lE."{}saR si]} \emph{vin.} be useful, be of use, come in handy.
\B{Fu\-la ho\-ren sì ay\-lì'\-fya\-vi a\-mip le\-sar sa\-yi me\-ngar oe\-ti nit\-ram 'ey\-ke\-fu nì\-txan.} It makes me very happy that the rules and new expressions will be useful to you both.~\LINK{http://naviteri.org/2012/06/spring-vocabulary-part-3/comment-page-1/\#comment-1515}{b}
(\ding{43}~\HL{sar}{\B{sar}})

% leskxir
\hi
\HT{leskxir}{\B{le\cp{skxir}}} \I{[lE."sk'iR]}  \emph{adj.} wounded.
\B{Oe\-ri ski\-en\-a tsyokx lu le\-skxir.} My right hand is wounded.~\LINK{http://naviteri.org/2014/05/mipa-ayliu-mipa-aysafpil-new-words-new-ideas/}{b}
(\ding{43}~\HL{skxir}{\B{skxir}})

% lesnonrra
\hi
\HT{lesnonrra}{\B{lesno\cp{nrr}a}} \I{[lE.sno."n\s{r}.a]}  \emph{adj.} full of self-pride, arrogant.
\B{Kea tsam\-si\-yu le\-sno\-nrr\-a ke tsun Na'\-vit iveyk.} No warrior full of self-pride can lead the People.~\LINK{http://naviteri.org/2012/07/fivospxiya-aylifyavi-amip-this-months-new-expressions-pt-2/}{b}
\B{Le\-sno\-nrr\-a tu\-te a swä\-pe\-ri fra\-krr slu eyk\-tan a\-fpxa\-mo.} An arrogant individual who always praises himself makes a terrible leader.~\LINK{naviteri.org/2025/07/ayu-leketengsyon-various-things/}{b} 
(\ding{43}~\HL{snonrra}{\B{sno\-nrr\-a}})

% lesngä'i
\hi
\HT{lesngA'i}{\B{le\cp{sngä}'i}} \I{[lE."{}sN\ae.Pi]}  \emph{adj.} original, existing at\slash{}from the start, first in a series.
\B{Swey\-lu txo ngal ke txi\-vìng sä\-fpìl\-it le\-sngä\-'i.} You shouldn't abandon your original idea.~\LINK{http://naviteri.org/2011/10/more-vocabulary-a-bit-of-grammar/}{b}
(\ding{43}~\HL{sngA'i}{\B{sngä\-'i}}, \HL{letsim}{\B{le\-tsim}})

% leso'ha
\hi
\HT{leso'ha}{\B{le\cp{so'}ha}} \I{[lE."{}soP.ha]}  \emph{adj.} enthusiastic, keen (\emph{of a person}).
\B{Ro\-le\-i\-un fu\-ta lu fo ka\-nu sì le\-so'\-ha nì\-wotx.} I found that they are intelligent and completely enthusiastic.~\LINK{http://naviteri.org/2013/08/fmawn-a-fkol-kilmulat-this-just-in/}{b}
(\ding{43}~\HL{so'ha}{\B{so'ha}})

% let'eylan
\hi 
\HT{let'eylan}{\B{let\cp{'ey}lan}} \I{[lEt\textcorner."PEj.lan]} \emph{adj., ofp.} friendly. 
\B{Ay\-ha\-pxì\-tu soaiä nge\-yä lu oe\-ru let\-'ey\-lan nì\-wotx.} Your family members are all friendly to me.~\LINK{https://naviteri.org/2024/09/pxevola-liu-amip-two-dozen-new-words/}{b} 
(\ding{43}~\HL{tI'eylanga'}{\B{tì\-'ey\-la\-nga'}}; \HL{tI'eylan}{\B{tì\-'ey\-lan}}) 

% letam
\hi
\HT{letam}{\B{le\cp{tam}}} \I{[lE."tam]}  \emph{adj.} sufficient, enough.
\B{Ngay\-txoa, nì\-aw\-no\-mum ke lo\-lu oer nì\-ke\-ftxo mì sok\-a srr ay\-skxom le\-tam fte lì'\-fya\-ri aw\-nge\-yä tì\-kang\-kem si\-vi.} My apologies: As you know, in recent days I have not had sufficient opportunity to work on our language.~\LINK{http://forum.learnnavi.org/language-updates/trr-rrtaya/}{ln}
(\ding{43}~\HL{tam}{\B{tam}})

% letokx
\hi
\HT{letokx}{\B{le\cp{tokx}}} \I{[lE."tok']}  \emph{adj.} bodily, physical.   
\B{Ay\-u\-van le\-tokx 'o' lu nì\-txan.} Sports are great fun.~\LINK{http://naviteri.org/2010/09/getting-to-know-you-part-3/}{b}
\B{Tì\-sraw le\-tokx sì tì\-ngu\-sä\-'än pxìm tä\-pa\-re fì\-tsap.} Physical pain and mental suffering are often interrelated.~\LINK{http://naviteri.org/2013/02/vospxi-ayol-posti-apup-short-post-for-a-short-month/}{b}
(\ding{43}~\HL{tokx}{\B{tokx}})

% letrr
\hi 
\HT{letrr}{\B{le\cp{trr}}} \I{[lE."t\s{r}]}  \emph{adj.} daily. 
\B{Oe new pi\-vll\-txe nì\-Na'\-vi mì oe\-yä le\-trr\-a tì\-rey.} I want to speak Na'vi in my daily life.~\LINK{http://forum.learnnavi.org/language-updates/frommerian-email/}{ln}
(\ding{43}~\HL{trr}{\B{trr}})

% letrrtrr
\hi
\HT{letrrtrr}{\B{le\cp{trr}trr}} \I{[lE."t\s{r}.t\s{r}]}  \emph{adj.} ordinary. 
\B{Ton\-ìri a\-la\-he, aw\-ngal yom wu\-tsot nì\-fya\-'o le\-trr\-trr; fì\-txon yom nì\-'eoio.} As for other nights, we eat dinner in an ordinary way; tonight we eat it ceremoniously.~\LINK{http://whyisthisnight.com/addl-more.html\#na\%27vi}{tn}
\B{Ta\-fral ho\-lawl ay\-oel ay\-nga\-fpi 'a\-'aw\-a tì\-päng\-kxo\-tsyìp\-it a tsun fko si\-var mì sì\-rey le\-trr\-trr.} That's why we prepared several little dialogues for you that one can use in ordinary life.~\LINK{http://naviteri.org/2012/07/tskxekengtsyip-a-mikyunfpi-a-little-listening-exercise/}{b}
\B{mo letrrtrr}, living room.~\LINK{http://naviteri.org/2012/10/mipa-vospxi-mipa-ayliu-new-words-for-the-new-month/}{b}
(\ding{43}~\HL{trr}{\B{trr}}) \\
\textsc{Derivatives}: \ding{43} \on{1}.~\HL{keltrrtrr}{\B{kel\-trr\-trr}}, unusual.

% letut
\hi 
\HT{letut}{\B{le\cp{tut}}} \I{[lE."tut\textcorner]} \emph{or} \I{[.tUt\textcorner]} (\textsc{rn}: \B{le\cp{tùt}} \I{[lE."tUt\textcorner]}) \emph{adj.} constant, continual. 
\B{Pe\-yä tì\-pu\-sll\-txel le\-tut oe\-ti srätx.} His constant talking annoys me.~\LINK{http://naviteri.org/2021/03/mipa-ayliu-si-aylifyavi-new-words-and-expressions/}{b}  
(\ding{43}~\HL{tut}{\B{tut}}, \HL{lukftang}{\B{luk\-ftang}})

% letwan
\hi
\HT{letwan}{\B{let\cp{wan}}} \I{[lEt\textcorner."wan]} \emph{adj.} dodgy, sneaky (\emph{of a person}).
(\ding{43}~\HL{tIwan}{\B{tì\-wan}})

% letsaktap
\hi 
\HT{letsaktap}{\B{le\cp{tsak}tap}} \I{[lE."\t{ts}ak\textcorner.tap\textcorner]} \emph{adj.} violent. 
(\ding{43}~\HL{tsaktap}{\B{tsak\-tap}})

% letsim
\hi
\HT{letsim}{\B{le\cp{tsim}}} \I{[lE."\t{ts}im]}  \emph{adj.} original, unique, not derived from another source.
\B{Swey\-lu txo ngal ke txi\-vìng sä\-fpìl\-it le\-tsim.} You shouldn't abandon your original idea.~\LINK{http://naviteri.org/2011/10/more-vocabulary-a-bit-of-grammar/}{b}
(\ding{43}~\HL{tsim}{\B{tsim}}, \HL{lesngA'i}{\B{le\-sngä\-'i}})

% letskxe
\hi 
\HT{letskxe}{\B{le\cp{tskxe}}} \I{[lE."\t{ts}k'E]} \emph{adj.} stony, of stone. 
\B{Fay\-frir le\-tskxe lor lu nì\-txan.} These stone layers are very beautiful.~\LINK{http://naviteri.org/2018/07/ta-sulfatu-a-ayliu-niul-more-words-from-our-experts/}{b}
(\ding{43}~\HL{tskxe}{\B{tskxe}})

% letsranten
\hi
\HT{letsranten}{\B{le\cp{tsran}ten}} \I{[lE."\t{ts}Ran.tEn]}  \emph{adj.} important.
\B{Fì\-trr\-mì le\-tsran\-ten--Trr 'Rr\-ta\-yä--new oe pi\-vll\-txe ay\-nga\-ru san kal\-txì sìk.} On this important day--Earth Day--I want to say hello to you.~\LINK{http://forum.learnnavi.org/language-updates/trr-rrtaya/}{ln}
\B{Mì\-ftxe\-le po\-ri lu oer le\-tsran\-ten\-a fmawn a new pi\-veng ngar.} By the way, I have some important news I want to tell you about him.~\LINK{http://naviteri.org/2011/09/\%E2\%80\%9Cby-the-way-what-are-you-reading\%E2\%80\%9D/}{b}
\B{Le\-kin lu tì\-txur, lu tì\-tstew. Slä le\-tsran\-ten fra\-to lu tì\-mal.} Strength and courage are necessary. But most important of all is trustworthiness.~\LINK{http://naviteri.org/2012/01/more-additions-to-the-lexicon/}{b}
(\ding{43}~\HL{tsranten}{\B{tsran\-ten}})

% letsunslu
\hi
\HT{letsunslu}{\B{le\cp{tsun}slu}} \I{[lE."\t{ts}un.slu]} \emph{or} \I{[."\t{ts}Un.]} (\textsc{rn}: \B{le\cp{tsùn}slu} \I{[lE."\t{ts}Un.slu]})  \emph{adj.} possible. 
(\ding{43}~\HL{tsunslu}{\B{tsun\-slu}}, \ding{214}~\HL{keltsun}{\B{kel\-tsun}})

% letswal
\hi 
\HT{letswal}{\B{le\cp{tswal}}} \I{[lE."\t{ts}wal]} \emph{adj., ofp.} powerful . 
(\ding{43}~\HL{tswal}{\B{tswal}}, \HL{tswalnga'}{\B{tswal\-nga'}})

% letxi
\hi
\HT{letxi}{\B{le\cp{txi}}} \I{[lE."t'i]}  \emph{adj.} hurried, frenzied.
(\ding{43}~\HL{txi}{\B{txi}}, \ding{214}~\HL{letxiluke}{\B{le\-txi\-lu\-ke}})

% letxiluke
\hi
\HT{letxiluke}{\B{le\cp{txi}luke}} \I{[lE."t'i.lu.kE]}  \emph{adj.} unhurried.
(\ding{43}~\HL{txi}{\B{txi}}, \HL{luke}{\B{lu\-ke}}; \ding{214}~\HL{letxi}{\B{le\-txi}})

% lew / si / säpi
\hi
\HT{lew}{\B{lew}} \I{[lEw]} \textsc{i}.~\emph{n.} cover, lid.
\B{Fu\-ri\-a txu\-la tsa\-lew\-ti lu srä sìl\-tsan to fngap.} For constructing that cover, woven cloth is better than metal.~\LINK{http://naviteri.org/2013/07/pximaw-tsawlultxa-just-after-the-meet-up/}{b} \\
\textsc{ii}.~\HT{lew si}{\B{lew si}} \I{[lEw si]}  \emph{vin.} cover. 
\B{Ne\-ni lew si fì\-txa\-yor vay txam\-pay nì\-wotx.} Sand covers this expanse all the way to the ocean.~\LINK{http://naviteri.org/2014/07/tsawlultxamaw-a-ayliu-post-meetup-words/}{b} \\
\textsc{iii}.~\HT{lew sApi}{\B{lew sä\cp{pi}}} \I{[lEw s\ae."p$\cdot$i]} \emph{vin.} cover oneself.
\B{Oe\-ru ke tsun li\-vam ke\-'u lor to Ey\-wa\-'e\-veng\-ä na'\-rìng a lew sä\-po\-li fa prr\-wll.} Nothing could be more beautiful to me than a Pandoran forest that has covered in moss.~\LINK{http://forum.learnnavi.org/language-updates/prrwll-moss-the-comparative-to-construct/}{ln}
\B{Kxawm mung\-wrr fì\-ki\-fkey a lew sä\-pì\-ye\-vi fa fpom sì lì'\-fya le\-Na'\-vi.} Except perhaps this world soon covering itself in peaceful well-being and the Na'vi language.~\LINK{http://forum.learnnavi.org/language-updates/prrwll-moss-the-comparative-to-construct/}{ln}

% lewn
\hi 
\HT{lewn}{\B{lewn}} \I{[lEwn]} \emph{vtr.} endure, stand, tolerate, bear. 
\B{Pe\-yä tì\-ru\-sol\-it ke tsun oe li\-vewn.} I can't stand her singing.~\LINK{http://naviteri.org/2021/03/mipa-ayliu-si-aylifyavi-new-words-and-expressions/}{b}
\B{Hu\-fwa tì\-sraw lu txan, tsun ay\-oe tsat li\-vewn.} Although the pain is great, we can endure it.~\LINK{http://naviteri.org/2021/03/mipa-ayliu-si-aylifyavi-new-words-and-expressions/}{b} \\
\textsc{Derivatives}: \on{1}.~\HL{ketsuklewn}{\B{ke\-tsuk\-lewn}}, intolerable, unacceptable, unbearable.

% lewng
\hi 
\HT{lewng}{\B{lewng}} \I{[lEwN]} \emph{n.} shame. 
\B{Mu\-nge fna\-we'\-tul lewng\-it soaia\-ru sne\-yä.} A coward brings shame to his\slash her family.~\LINK{http://naviteri.org/2021/09/aawa-ayliu-amip-a-few-new-words/}{b}
(\ding{214}~\HL{nrra}{\B{nrr\-a}}) \\
\textsc{Derivatives}: \on{1}.~\HL{lelewng}{\B{le\-lewng}}, shameful. \\
\on{2}.~\HL{lewnga'}{\B{lew\-nga'}}, shameful.

% lewnga'
\hi 
\HT{lewnga'}{\B{\cp{lew}nga'}} \I{["lEw.NaP]} \emph{adj., nfp.} shameful. 
\B{vo\-ìk a\-lew\-nga'}, shameful behavior.~\LINK{http://naviteri.org/2021/09/aawa-ayliu-amip-a-few-new-words/}{b} 
(\ding{43}~\HL{lewng}{\B{lewng}})

% ley
\hi
\HT{ley}{\B{ley}} \I{[lEj]} \emph{vin.} be of value, have some positive value, be worth s.t. 
\B{Fì\-'u ley.} This thing is valuable.~\LINK{http://naviteri.org/2014/03/value-and-worth/}{b} 
\B{Fì\-'u ke ley.} This thing is worthless.~\LINK{http://naviteri.org/2014/03/value-and-worth/}{b}
(\emph{a.~in comparisons}) \B{Oe\-ri tsaw ke ley fì\-'u\-to.} I don't value that over this.~\LINK{http://naviteri.org/2014/03/value-and-worth/}{b}
\B{Ma\-sat oeyä ley nì\-ftxan na pum nge\-yä.} My breastplate is as valuable as yours.~\LINK{http://naviteri.org/2014/03/value-and-worth/}{b}
(\emph{b.~with the dative}) \B{Sìl\-pey oe, fì\-po\-stì lil\-vey ay\-nga\-ru nì\-wotx!} I hope this post has been valuable to you.~\LINK{http://naviteri.org/2014/03/value-and-worth/}{b}
(\emph{c. proverb}) \B{Flä ke flä, ley sä\-fmi.} Whether you succeed or not, the attempt has value.~\LINK{http://naviteri.org/2014/03/value-and-worth/}{b} 
(\emph{d. proverb}) \B{'Uo\-ri hì\-pey, kxawm nga\-ru ke ley.} ``If you hesitate doing something, it might not be important to you,'' i.e. \emph{more meant\slash used as a motivation for s.o. hesitating, or even as a teasing to get s.o. into action}.~\LINK{http://naviteri.org/2021/08/ulte-ayyoratu-leiu-and-the-winners-are-2/}{b} \\
\textsc{Derivation}: \on{1}.~\HL{leytslam}{\B{ley\-tslam}}, appreciate.

% leykatem
\hi
\HT{leykatem}{\B{ley\cp{ka}tem}} \I{[lEj."k$\cdot$a.t$\cdot$Em]} \emph{vtr., inf.}\on{23} change s.t.
\B{Sra\-ne, lì\-'ul alu mì fra\-krr ley\-ka\-tem lì\-'ut a\-hay tsa\-fya.} Yes, the word \emph{in} always changes the following word like that.~\LINK{http://forum.learnnavi.org/language-updates/fmawno/msg293595/\#msg293595}{ln}
\B{Oe\-ri Unil\-tì\-ran\-tokx\-ìl tì\-rey\-ti ley\-ko\-la\-te\-i\-em nì\-txan.} \emph{Avatar} has greatly changed my life for the better.~\LINK{http://forum.learnnavi.org/language-updates/pseudo-canon-from-bd-live-content/msg350756/\#msg350756}{ln}
(\ding{43}~\HL{latem}{\B{la\-tem}})

% leykek
\hi 
\HT{leykek}{\B{ley\cp{kek}}} \I{[lEj."kEk\textcorner]} \emph{vtr.} enforce. 
\B{Pol ho\-ren\-it a\-yll ley\-kek.} He enforces the laws.~\LINK{http://naviteri.org/2020/06/mi-tanlokxe-oeya-srr-afpxamo-terrible-days-in-my-country/}{b} 
(\ding{43}~\HL{lek}{\B{lek}}) \\
\textsc{Derivatives}: \ding{43} \on{1}.~\HL{horenleykekyu}{\B{ho\-ren\-ley\-kek\-yu}}, law enforcer, police officer.

% leykekyu
\hi 
\HT{leykekyu}{\B{ley\cp{kek}yu}} \I{[lEj."kEk\textcorner.ju]} \emph{n.} law enforcer, (\emph{on Earth}) police officer. (\emph{shortened form of} \ding{43}~\HL{horenleykekyu}{\B{ho\-ren\-ley\-kek\-yu}}).
(\ding{43}~\HL{leykek}{\B{ley\-kek}}) 

% leym
\hi 
\HT{leym}{\B{leym}} \I{[lEjm]}  \emph{vin.} call out, cry out, exclaim.
\B{Zun Ey\-wa\-'e\-veng\-it oel ti\-vok, zel leym san Srung si ay\-oe\-ru, ma Ey\-wa sìk!} If I were on Pandora, I would call out, ``Help us, Eywa!''~\LINK{http://naviteri.org/2016/12/tengkrr-perahem-zisit-amip-as-the-new-year-arrives/}{b} \\
\textsc{Derivatives}: \ding{43} \on{1}.~\HL{tIleym}{\B{tì\-leym}}, call.
\on{2}.~\HL{leymfe'}{\B{leym\-fe'}}, complain.
\on{3}.~\HL{leymkem}{\B{leym\-kem}}, protest.

% leymfe'
\hi 
\HT{leymfe'}{\B{leym\cp{fe'}}} \I{[l$\cdot$Ejm."fEP]} \emph{vin., inf.}\on{11} complain. 
\B{Fo le\-reym\-fe' tsnì syu\-ve lu wew.} They're complaining that the food is cold.~\LINK{http://naviteri.org/2017/09/zisikrr-amip-ayliu-amip-new-words-for-the-new-season/}{b} 
(\ding{43}~\HL{leym}{\B{leym}}, \HL{fe'}{\B{fe'}}) \\
\textsc{Derivatives}: \ding{43}
\on{1}.~\HL{tIleymfe'}{\B{tì\-leym\-fe'}}, complaining. 
\on{2}.~\HL{sAleymfe'}{\B{sä\-leym\-fe'}}, complaint.

% leymkem
\hi 
\HT{leymkem}{\B{leym\cp{kem}}} \I{[l$\cdot$Ejm."kEm]}, \emph{coll. pronounced as} \I{[lejN."kEm]} \emph{vin., inf.}\on{11} protest (\B{tsnì}, that). 
\B{Oe leym\-kem! Fì\-tsam\-ìl Na'\-vit tì\-sraw sey\-ka\-yi nì\-'aw ul\-te ku\-tut ke lä\-txayn.} I protest! This war will only harm the People and not defeat the enemy.~\LINK{http://naviteri.org/2017/09/zisikrr-amip-ayliu-amip-new-words-for-the-new-season/}{b}
\B{Tsay\-hem\-ìri\slash Tsay\-hem\-te\-ri po lo\-leym\-kem.} She protested those actions.~\LINK{http://naviteri.org/2017/09/zisikrr-amip-ayliu-amip-new-words-for-the-new-season/}{b}
\B{Lo\-leym\-kem po tsnì fwa Ak\-wey slu olo'\-eyk\-tan lu kem\-wi\-ä.} He protested that it was unfair for Ak\-wey to become clan leader.~\LINK{http://naviteri.org/2017/09/zisikrr-amip-ayliu-amip-new-words-for-the-new-season/}{b} 
\B{Fpxa\-mo\-a fì\-kem\-ìl a\-fpxa\-mo ay\-oe\-ti tsngey\-ko\-law\-vìk, tsa\-krr ley\-ko\-leym\-kem.} This terrible, terrible action has made us cry, then protest.~\LINK{http://naviteri.org/2020/06/mi-tanlokxe-oeya-srr-afpxamo-terrible-days-in-my-country/}{b}
(\ding{43}~\HL{leym}{\B{leym}}, \HL{kemwiA}{\B{kem\-wi\-ä}}) \\
\textsc{Derivatives}: \ding{43} 
\on{1}.~\HL{tIleymkem}{\B{tì\-leym\-kem}}, protesting, protest. 
\on{2}.~\HL{sAleymkem}{\B{sä\-leym\-kem}}, protest, instance of protesting.

% leyn
\hi 
\HT{leyn}{\B{leyn}} \I{[lEjn]}  \emph{vtr.} repeat, do again.
\B{Ru\-txe li\-veyn.} Please repeat (it).~\LINK{http://naviteri.org/2010/09/getting-to-know-you-part-2/}{b}
\B{Rä\-'ä li\-veyn!} Don't do that again!~\LINK{http://naviteri.org/2010/09/getting-to-know-you-part-2/}{b} \\
\textsc{Derivatives}: \ding{43} 
\on{1}.~\HL{nIleyn}{\B{nì\-leyn}}, repeatedly.

% leyr
\hi 
\HT{leyr}{\B{leyr}} \I{[lEjR]} \textsc{i}. \emph{adj.} frozen. 
\B{Ke tsun ioang ri\-vun syu\-vet mì hll\-te a\-leyr.} Animals can't find food in the frozen ground.~\LINK{http://naviteri.org/2018/07/ta-sulfatu-a-ayliu-niul-more-words-from-our-experts/}{b} \\
\textsc{ii}. \B{leyr si} \I{[lEjR si]} \emph{vin.} freeze (s.t.).
\B{Txo awnga fì\-tsngan\-ur leyr si\-vi, tsun tsat yi\-vom kin\-trr\-ay.} If we freeze this meat, we can eat it next week.~\LINK{http://naviteri.org/2018/07/ta-sulfatu-a-ayliu-niul-more-words-from-our-experts/}{b} \\
\textsc{iii}. \B{slu leyr} \I{[slu lEjR]} \emph{vin.} s.t. freezes.
\B{Mì zì\-sì\-krr a\-txa\-wew slu ay\-ora leyr.} In the very cold season, the lakes freeze.~\LINK{http://naviteri.org/2018/07/ta-sulfatu-a-ayliu-niul-more-words-from-our-experts/}{b}

% leytslam
\hi 
\HT{leytslam}{\B{\cp{ley}tslam}} \I{["lEj.\t{ts}l$\cdot$am]} \emph{vtr., inf.}\on{22} appreciate. 
\B{Nge\-yä fay\-lì\-'ut a\-tì\-tstun\-wi\-nga' oel ley\-tslam, ma 'ey\-lan.} I appreciate your kind words, friend.~\LINK{http://naviteri.org/2020/11/vospxivopeya-ayliu-amip-novembers-new-words/}{b}
(\ding{43}~\HL{ley}{\B{ley}}, \HL{tslam}{\B{tslam}})

% leweopx
\hi 
\HT{leweopx}{\B{lewe\cp{opx}}} \I{[lE.wE."{}op']} \emph{adj.} wave\hyp{}like. 
\B{Tsay\-re\-nur le\-we\-opx a sìn ne\-ni tìng na\-ri.} Look at those wave\hyp{}like patterns in the sand.~\LINK{http://naviteri.org/2017/09/zisikrr-amip-ayliu-amip-new-words-for-the-new-season/}{b} 
(\ding{43}~\HL{weopx}{\B{we\-opx}})

% lewäte
\hi
\HT{lewAte}{\B{lewä\cp{te}}} \I{[lE.w\ae."tE]}  \emph{adj.} disagreeable, argumentative (\emph{of an agent}).
(\ding{43}~\HL{wAte}{\B{wä\-te}})

% leyapay
\hi
\HT{leyapay}{\B{le\cp{ya}pay}} \I{[lE."ja.paj]}  \emph{adj.} foggy, misty. 
(\ding{43}~\HL{yapay}{\B{ya\-pay}})

% leyayl
\hi 
\HT{leyayl}{\B{le\cp{yayl}}} \I{[lE."jajl]} \emph{adj.} nonsensical. 
(\ding{43}~\HL{yayl}{\B{yayl}})

% leye'krr
\hi
\HT{leye'krr}{\B{le\cp{ye'}krr}} \I{[lE."jEP.k\s{r}]}  \emph{adj.} early. 
\B{tì\-pä\-hem le\-ye'\-krr}, an early arrival.~\LINK{http://naviteri.org/2010/07/vocabulary-update/}{b}
(\ding{43}~\HL{ye'krr}{\B{ye'\-krr}}, \ding{214}~\HL{lehawngkrr}{\B{le\-hawng\-krr}})

% leyewla 
\hi
\HT{leyewla}{\B{le\cp{yew}la}} \I{[lE."jEw.la]}  \emph{adj.} disappointing. 
\B{Pe\-yä sä\-ngä\-'än\-ìri lu sä\-fe\-'ul le\-yew\-la nì\-txan.} The worsening of his depression is very disappointing.~\LINK{http://naviteri.org/2013/03/tsanerul-feerul-getting-better-getting-worse/}{b}
\B{Kea kem le\-yew\-la rä\-'ä si, ru\-txe.} Please don't let me down.~\LINK{http://naviteri.org/2012/10/mipa-vospxi-mipa-ayliu-new-words-for-the-new-month/}{b}
(\ding{43}~\HL{yewla}{\B{yew\-la}})

% lezeswa
\hi
\HT{lezeswa}{\B{le\cp{ze}swa}} \I{[lE."zE.swa]}  \emph{adj.} grassy.
(\ding{43}~\HL{zeswa}{\B{zeswa}}) 

% li
\hi
\HT{li}{\B{li}} \I{[li]}  \emph{adv.} (\emph{a.}) already.   
\B{Tì\-kang\-kem li ha\-sey lu.} The work is already finished.~\LINK{http://naviteri.org/2011/02/new-vocabulary-part-2/}{b}
(\emph{b. negated}) not yet.
\B{Vì\-ming\-kap oe\-ti fu\-la poe ke li ke pol\-txe san oe za\-sya\-'u.} It just occurred to me that she hasn't yet said she's coming.~\LINK{http://naviteri.org/2011/09/\%E2\%80\%9Cby-the-way-what-are-you-reading\%E2\%80\%9D/}{b}
\B{Fo li po\-lä\-hem srak?--Ke li.} Have they already arrived?--Not yet.~\LINK{http://naviteri.org/2011/02/new-vocabulary-part-2/}{b}
(\emph{c. in imperatives to convey strong urgency})
\B{Ngal mi fì\-tseng\-it te\-rok srak? Li \cp{kä}!} You're still here? Get going!~\LINK{http://naviteri.org/2011/02/new-vocabulary-part-2/}{b}
(\emph{d. set phrase})
\B{\cp{Li} ko.} Well, get to it, then.\slash Let's get on it.~\LINK{http://naviteri.org/2011/02/new-vocabulary-part-2/}{b} 
(\emph{e. hesitant or weak `yes'}) 
\B{ Nga mll\-te srak?--Li, slä\ldots} Do you agree?--Well, yes, I guess so, but\ldots{}~\LINK{http://naviteri.org/2011/02/new-vocabulary-part-2/}{b} 
(\emph{f. with} \ding{43}~\HL{sre}{\B{sre}}) by; before or up to, but not after. 
\B{Kem si li trr\-ay\-sre.} Do it by tomorrow.~\LINK{http://naviteri.org/2011/02/new-vocabulary-part-2/}{b} (\emph{also} \ding{43}~\HL{lisre}{\B{li\-sre}}) \\
\textsc{Derivatives}: \ding{43} \on{1}.~\HL{nIli}{\B{nì\-li}}, in advance. \\
\on{2}.~\HL{lisre}{\B{li\-sre}}, by; before but not after.

% lie / si
\hi
\HT{lie}{\B{\cp{li}e}} \I{["li.E]} \textsc{i}. \emph{n.} experience. 
\B{Kop oe\-ru lo\-lä\-ngu lie a ha\-pxì\-tu soaiä ter\-kup.} Sadly, I too have had the experience of a family member dying.~\LINK{http://naviteri.org/2013/01/awvea-posti-zisita-amip-first-post-of-the-new-year/}{b}
\B{Lu oe\-ru lie a yom tey\-lut.} I once ate \emph{teylu}.~\LINK{http://naviteri.org/2013/01/awvea-posti-zisita-amip-first-post-of-the-new-year/}{b} \\
\textsc{ii}. \B{\cp{li}e si} \I{["li.E si]} \emph{vin.} experience (\emph{an event, a feeling, a person}). 
\B{Tu\-te a ke\-ftxo fra\-to lu tsa\-po a tì\-yawn\-ur lie ke so\-li kaw\-krr.} The saddest person of all is the one who has never experienced love.~\LINK{http://naviteri.org/2018/04/100a-liu-amip-64-new-words-part-2/}{b} \\ 
\textsc{Derivatives}: \ding{43} \on{1}.~\HL{'awlie}{\B{'aw\-lie}}, once (\emph{experiential}).

% lik
\hi 
\HT{lik}{\B{lik}} \I{[lik\textcorner]} \emph{n.} point, spot, particular place or position in some area. 
\B{Tsko\-ti fyep fì\-lik\-ro.} Grasp the bow at this spot.~\LINK{https://naviteri.org/2024/05/krrteri-about-time/}{b}
\B{Ro sa\-lik a mì sä\-ftxu\-lì\-'u a\-tì\-'i\-lu\-ke pe\-yä, oe ho\-lum.} At that point in his endless speech, I left.~\LINK{https://naviteri.org/2024/05/krrteri-about-time/}{b} \\
\textsc{Derivatives}: \ding{43} \on{1}.~\HL{krrlik}{\B{krr\-lik}}, point in time.

% limang
\hi
\HT{limang}{\B{li\cp{mang}}} \I{[li."maN]}  \emph{n.} moot, celebratory gathering or event. 
\B{A\-ha'\-ri wrr\-ko\-lä fte li\-mang si\-vu\-nu por.} Aha'ri is out enjoying the moot.~\LINK{https://forum.learnnavi.org/language-updates/limangmoot-from-afop-dialogue/}{ln\slash fop}

% lini
\hi
\HT{lini}{\B{\cp{li}ni}} \I{["li.ni]}  \emph{countable n.} young (\emph{of an animal, bird, fish, insect}). 
\B{Fì\-fne\-ya\-yol tsrul\-it txu\-la mì song\-ropx ut\-ral\-ä fte ay\-li\-nit hi\-vaw\-nu wä sar\-ni\-o\-ang.} This kind of bird builds its nest in a tree cavity so as to protect its young from predators.~\LINK{http://naviteri.org/2013/04/wheres-the-bathroom-and-other-useful-things/}{b} \\
\textsc{Derivatives}: \ding{43} \on{1}.~\HL{nga'lini}{\B{nga'\-li\-ni}}, be pregnant (for animals).

% lipx
\hi 
\HT{lipx}{\B{lipx}} \I{[lip']} \emph{vin.} drip. 
\B{Tom\-pa ze\-rup ul\-te pay lipx kxam\-lä fäp\-yo.} It's raining and water is dripping through the roof.~\LINK{http://naviteri.org/2020/04/aawa-u-amip-a-few-new-things/}{b}
(\emph{a. idiom}) \B{Pay\-ìl a lipx tskxe\-ti ripx.} Dripping water pierces a stone. [\emph{persistent effort can accomplish unexpected and amazing things}]~\LINK{http://naviteri.org/2020/04/aawa-u-amip-a-few-new-things/}{b}  \\
\textsc{Derivatives}: \ding{43} \on{1}.~\HL{sAlipx}{\B{sä\-lipx}}, drip (of liquid).

% lisre
\hi
\HT{lisre}{\B{\cp{li}sre}} \I{["li.sRE]} \emph{adp.}\texttt{+} by, before, up to but not after.
\B{Kem si li\-sre srr\-ay.} Do it by tomorrow.~\LINK{http://naviteri.org/2011/02/new-vocabulary-part-2/}{b}
(\ding{43}~\HL{li}{\B{li}}, \HL{sre}{\B{sre}})

% liswa / si
\hi 
\HT{liswa}{\B{li\cp{swa}}} \I{[li."{}swa]} \textsc{i}. \emph{n.} nourishment. \\
\textsc{ii}. \B{li\cp{swa} si} \I{[li."{}swa si]} \emph{vin.} nourish, provide nourishment.
\B{Fì\-'e\-wll li\-swa si Na'\-vi\-ru.} This plant provides nourishment to the People.~\LINK{http://naviteri.org/2018/04/100a-liu-amip-64-new-words-part-2/}{b} 

% litx
\hi
\HT{litx}{\B{litx}} \I{[lit']}  \emph{adj.} sharp (as a blade).
\B{El\-tu si! Tsa\-tstal a\-fwem lu litx nì\-txan.} Pay attention! That blunt knife is very sharp.~\LINK{http://naviteri.org/2012/03/spring-vocabulary-part-1/}{b}
(\ding{214}~\HL{tete}{\B{te\-te}})

% litsi
\hi 
\HT{litsi}{\B{\cp{li}tsi}} \I{["li.\t{ts}i]} \emph{adj., ofp.} thin, lean, lithe. 
(\ding{214}~\HL{ompu}{\B{om\-pu}}; \ding{43}~\HL{flI}{\B{flì}}) 

% lì'fya
\hi
\HT{lI'fya}{\B{\cp{lì'}fya}} \I{["lIP.fja]}  \emph{n.} 
(\emph{a.}) language. 
\B{Ay\-lì'\-fya yaw\-ne le\-ru oer ta\-krr\-a 'e\-veng la\-mu.} I've loved languages since I was a child.~\LINK{http://naviteri.org/2011/01/new-year-new-vocabulary/}{b}
\B{Fol lì'\-fya\-ti aw\-nge\-yä sar a fya\-'o lu ko\-sman.} It's wonderful the way they use our language.~\LINK{http://naviteri.org/2011/09/miscellaneous-vocabulary/}{b}
(\emph{b.}) the language spoken on Pandora.
\B{Lì'\-fya\-ri pol tok pe\-yìt?} \emph{or} \B{Lì'\-fya\-ri po kll\-kxem sìn pe\-yì?}\slash \B{Lì'\-fya\-ri po sìn pe\-yì?}\slash \B{Lì'\-fya\-ri po pe\-yì?} How good is her Na'vi?~\LINK{http://naviteri.org/2013/01/lifyari-po-peyi-how-good-is-her-navi/}{b}
(\emph{c.}) \B{lì'\-fya le\-Na'\-vi}, the Na'vi language.
\B{Nge\-yä lì'\-fya le\-Na'\-vi txan\-tsan lu nì\-ngay.} Your Na'vi language is truly excellent.~\LINK{http://forum.learnnavi.org/language-updates/score!-\%281\%29-the-verb-to-use-\%282\%29-more-detail-on-teya-si/}{ln}
\B{Sra\-ke fnan ngal lì'\-fya\-ti le\-Na'\-vi?} Are you good at Na'vi?~\LINK{http://naviteri.org/2010/09/getting-to-know-you-part-3/}{b}
\B{Unil\-tì\-ran\-tokx\-ìri lì'\-fya\-yä le\-Na'\-vi wa\-wet fol tslam.} They understand the importance of the Na'vi language for \emph{Avatar}.~\LINK{http://naviteri.org/2013/08/fmawn-a-fkol-kilmulat-this-just-in/}{b}
\B{Lì'\-fya\-ri le\-Na'\-vi nga tsan\-'e\-re\-i\-ul fra\-trr.} I'm delighted that your Na'vi is improving every day.~\LINK{http://naviteri.org/2013/03/tsanerul-feerul-getting-better-getting-worse/}{b} \\  
\textsc{Derivatives}: \ding{43} \on{1}.~\HL{lI'fyavi}{\B{lì'\-fya\-vi}}, expression.   
\on{2}.~\HL{lI'fyaolo'}{\B{lì'\-fya\-o\-lo'}}, language group.
\on{3}.~\HL{lI'fyeng}{\B{lì'\-fyeng}}, honorific language.
\on{4}.~\HL{lI'fyafnel}{\B{lì'\-fya\-fnel}}, dialect, variety of a language.
\on{5}.~\HL{rAptulI'fya}{\B{räp\-tu\-lì'\-fya}}, coarse, vulgar language.

% lì'fyafnel
\hi 
\HT{lI'fyafnel}{\B{\cp{lì'}fyafnel}} \I{["lIP.fja.fnEl]} \emph{n.} dialect, variety of a language. 
\B{Lu mì `Fya\-'o Pay\-ä' mip\-a lì'\-fya\-fnel lì'\-fya\-yä le\-Na'\-vi.} There is a new dialect of the Na'vi language in \emph{Avatar: The Way of Water}.~\LINK{http://naviteri.org/2022/12/trr-anawm-polaheiem-the-great-day-has-arrived/}{b} 
(\ding{43}~\HL{lI'fya}{\B{lì'fya}}, \HL{fnel}{\B{fnel}})

% lì'fyaolo'
\hi
\HT{lI'fyaolo'}{\B{\cp{lì'}fyaolo'}} \I{["lIP.fja.o.loP]}  \emph{n.} language group, \emph{Sprachbund}.
\B{Nge\-yä tì\-kang\-kem a\-fyo\-le meuia leiu lì'\-fya\-o\-lo'\-ru, ma tsmu\-kan.} Your sublime work honors the language group, brother.~\LINK{http://naviteri.org/2012/10/audio-and-video-learning-materials-for-navi-101/}{b}
\B{Fu\-ri\-a fì\-lì'\-fya\-vit fkol ta 'Ìng\-lì\-sì mo\-lu\-nge, sìl\-pey oe tsnì ay\-ha\-pxì\-tu lì'\-fya\-o\-lo'\-ä aw\-nge\-yä ke stì\-ye\-vi.} That this expression is brought from English, will hopefully not anger the members of the language group.~\LINK{http://forum.learnnavi.org/language-updates/txelanit-hivawl/}{ln}
(\ding{43}~\HL{lI'fya}{\B{lì'\-fya}}, \HL{olo'}{\B{o\-lo'}}) \\
\textsc{Derivatives}: \ding{43}
\on{1}.~\HL{lI'fyaolo'tu}{\B{lì'\-fya\-o\-lo'\-tu}}, member of the language community.

% lì'fyaolo'tu
\hi
\HT{lI'fyaolo'}{\B{\cp{lì'}fyaolo'tu}} \I{["lIP.fja.o.loP.tu]}  \emph{n.} member of the language group\slash \emph{Sprachbund}.~\LINK{https://forum.learnnavi.org/language-updates/a-collection-of-questions-answered/msg679687/\#msg679687}{ln}
(\ding{43}~\HL{lI'fyaolo'}{\B{lì'\-fya\-olo'}}, \HL{-tu}{\B{-tu}})

% lì'fyavi
\hi
\HT{lI'fyavi}{\B{\cp{lì'}fyavi}} \I{["lIP.fja.vi]}  \emph{n.} expression, bit of language.
\B{Lì'\-fya\-vi le\-sar nì\-txan nang!} The expression is very useful.~\LINK{http://forum.learnnavi.org/language-updates/just-another-few-words/}{ln}
\B{\ldots slä nì\-sìl\-pey pum a nga' ay\-lì'\-fya\-vit le\-sar.} \ldots but hopefully one that contains useful expressions.~\LINK{http://naviteri.org/2013/02/vospxi-ayol-posti-apup-short-post-for-a-short-month/}{b}
\B{Tsa\-lì'\-fya\-vi\-ri fko ze\-ne si\-var lì\-'ut alu fì\-'u me\-lo.} One must use the word \emph{fì'u} twice in that expression.~\LINK{http://naviteri.org/2011/04/\%E2\%80\%99a\%E2\%80\%99awa-li\%E2\%80\%99fyavi-amip\%E2\%80\%94a-few-new-expressions/comment-page-1/\#comment-677}{b}
(\ding{43}~\HL{lI'fya}{\B{lì'\-fya}}, \HL{tIyewn}{\B{tì\-yewn}})

% lì'fyeng / si
\hi 
\HT{lI'fyeng}{\B{lì'\cp{fyeng}}} \I{[lIP."fjEN]} \textsc{i}. \emph{n.} honorific language. 
\B{Lì'\-fyeng kel\-kin.} No need to speak so formally.~\LINK{http://naviteri.org/2022/02/lifyengteri-concerning-honorific-language/}{b}  
\B{Lì'\-fyeng kel\-kin ko.} Let's not speak so formally with each other.~\LINK{http://naviteri.org/2022/02/lifyengteri-concerning-honorific-language/}{b} \\
\textsc{ii}. \B{lì'\cp{fyeng} si} \I{[lIP."fjEN si]} \emph{vin.} use honorific language.
\B{Lì'fyeng rä'ä si.} Don't use honorific language.~\LINK{http://naviteri.org/2022/02/lifyengteri-concerning-honorific-language/}{b} 
\B{Fwa lì'fyeng si lu kelkin.} It's not necessary to use honorific language.~\LINK{http://naviteri.org/2022/02/lifyengteri-concerning-honorific-language/}{b}
(\ding{43}~\HL{lI'fya}{\B{lì'\-fya}}, \HL{leheng}{\B{le\-heng}}) \\
\textsc{Derivatives}: \ding{43} \on{1}.~\HL{fyengkekin}{\B{fyeng\-ke\-kin}}, no need to use honorific language.

% lì'kong
\hi 
\HT{lI'kong}{\B{\cp{lì'}kong}} \I{["lIP.koN]} \emph{n.} syllable. 
\B{Lu tsa\-lì\-'ur alu fì\-haw\-re'\-ti tsìng\-a lì'\-kong.} The word \emph{fì\-haw\-re'\-ti} has four syllables.~\LINK{http://naviteri.org/2017/10/more-language-for-talking-about-language/}{b} 
(\ding{43}~\HL{lI'u}{\B{lì\-'u}}, \HL{'ekong}{\B{'e\-kong}})

% Lì'ona
\hi
\HT{lI'ona}{\B{Lì\cp{'o}na}} \I{[lI."Po.na]}  \emph{n. name of a clan}. l'Ona clan.
\B{Ke za\-syup lì\-'O\-na ne kxu\-tu a mì\-fa fu a wrr\-pa.} The l'Ona will not perish to the enemy within or the enemy without.~\LINK{http://forum.learnnavi.org/language-updates/ltasygt-with-ke-and-nga/}{g\slash ln}

% lì'u
\hi
\HT{lI'u}{\B{\cp{lì}'u}} \I{["lI.Pu]}  \emph{n.} word.
\B{Ma Ey\-tu\-kan, lu oe\-ru ay\-lì\-'u fra\-por.} Eytukan, I have something to say, to everyone.~\LINK{http://wiki.learnnavi.org/index.php?title=Corpus\#Behind_the_Scenes}{a\slash wiki}
\B{Nge\-yä tsay\-lì\-'ul tì\-fkey\-tok\-it fe\-'ey\-ko\-lul nì\-'aw.} Those words of yours have only made the situation worse.~\LINK{http://naviteri.org/2013/03/tsanerul-feerul-getting-better-getting-worse/}{b}
\B{Sra\-ne, lì\-'ul alu mì fra\-krr ley\-ka\-tem lì\-'ut a\-hay tsa\-fya.} Yes, the word \emph{in} always changes the following word like that.~\LINK{http://forum.learnnavi.org/language-updates/fmawno/msg293595/\#msg293595}{ln}
\B{Ro sngä\-'i\-tseng tsa\-lì\-'uä alu 'ey\-lan lu tì\-ftang.} At the beginning of the word \emph{'eylan} there's a glottal stop.~\LINK{http://naviteri.org/2012/07/fivospxiya-aylifyavi-amip-this-months-new-expressions-pt-2/}{b}
\B{Nge\-yä fay\-lì\-'u\-ri sei\-yi oe i\-ra\-yo.} I thank you for these words of yours.~\LINK{http://naviteri.org/2011/08/new-vocabulary-clothing/comment-page-1/\#comment-995}{b} \\ 
\textsc{Derivatives}: \ding{43} 
\on{1}.~\HL{lI'fya}{\B{lì'\-fya}}, language. 
\on{2}.~\HL{lI'upam}{\B{lì\-'u\-pam}}, pronunciation. 
\on{3}.~\HL{lI'upuk}{\B{lì\-'u\-puk}}, dictionary. 
\on{4}.~\HL{lI'uvan}{\B{lì\-'u\-van}}, pun, word play. 
\on{5}.~\HL{fayluta}{\B{fay\-lu\-ta}}, that (\emph{reported speech}). 
\on{6}.~\HL{ftxulI'u}{\B{ftxu\-lì\-'u}}, orate, give a speech. 
\on{7}.~\HL{kemlI'u}{\B{kem\-lì\-'u}}, verb. 
\on{8}.~\HL{tstxolI'u}{\B{tstxo\-lì\-'u}}, noun.
\on{9}.~\HL{lI'upe}{\B{lì\-'u\-pe\slash pe\-lì\-'u}}, what, which (word, utterance)?
\on{10}.~\HL{lI'ukIng}{\B{lì\-'u\-kìng}}, sentence.
\on{11}.~\HL{lI'kong}{\B{lì'\-kong}}, syllable.
\on{12}.~\HL{lI'uvi}{\B{lì\-'u\-vi}}, affix.
\on{13}.~\HL{syonlI'u}{\B{syon\-lì\-'u}}, adjective.
\on{14}.~\HL{fyalI'u}{\B{fya\-lì\-'u}}, adverb.
\on{15}.~\HL{starlI'u}{\B{star\-lì\-'u}}, adposition.
\on{16}.~\HL{tilI'u}{\B{ti\-lì\-'u}}, conjunction.
\on{17}.~\HL{rAptulI'u}{\B{räp\-tu\-lì\-'u}}, coarse or swear word.
\on{18}.~\HL{rawnlI'u}{\B{rawn\-lì\-'u}}, pronoun.
\on{19}.~\HL{tengralI'u}{\B{teng\-ra\-lì\-'u}}, synonym.
\on{20}.~\HL{wAralI'u}{\B{wä\-ra\-lì\-'u}}, antonym.

% lì'ukìng
\hi 
\HT{lI'ukIng}{\B{\cp{lì}'ukìng}} \I{["lI.Pu.kIN]} \emph{n.} sentence (of words), \emph{lit.}: word thread. 
(\ding{43}~\HL{lI'u}{\B{lì\-'u}}, \HL{kIng}{\B{kìng}}) \\
\textsc{Derivatives}: \ding{43}
\on{1}.~\HL{lI'ukIngvi}{\B{lì\-'u\-kìng\-vi}}, phrase.

% lì'ukìngvi
\hi 
\HT{lI'ukIngvi}{\B{\cp{lì}'ukìngvi}} \I{["lI.Pu.kIN.vi]} \emph{n.} phrase. 
(\ding{43}~\HL{lI'ukIng}{\B{lì\-'u\-kìng}})

% lì'upam
\hi
\HT{lI'upam}{\B{\cp{lì}'upam}} \I{["lI.Pu.pam]}  \emph{n.} pronunciation. 
\B{Poe tso\-lun nì\-Na'\-vi pi\-vll\-txe, hu\-fwa lì\-'u\-pam hek nì\-'it.} She can speak Na'vi although her pronunciation is a bit strange.~\LINK{http://naviteri.org/2010/07/ting-mikyun-fte-tslivam-listening-comprehension-1-3/}{b}
\B{Slä lì\-'u\-pam\-ìri pll\-txe fko san tì\-o\-mu\-mì sìk.} But as for pronunciation one says \emph{tìomumì}.~\LINK{http://naviteri.org/2011/09/\%E2\%80\%9Cby-the-way-what-are-you-reading\%E2\%80\%9D/comment-page-1/\#comment-1092}{b}
(\ding{43}~\HL{lI'u}{\B{lì'u}}, \HL{pam}{\B{pam}})

% lì'upe
\hi
\HT{lI'upe}{\B{\cp{lì}'upe}} \I{["lI.Pu.pE]} \emph{or} \B{pe\cp{lì}'u} \I{[pE."lI.Pu]} \emph{inter.} what, which (word, utterance)?
(\ding{43}~\HL{lI'u}{\B{lì\-'u}})

% lì'upuk
\hi
\HT{lI'upuk}{\B{\cp{lì}'upuk}} \I{["lI.Pu.puk\textcorner]} \emph{or} \I{[.pUk\textcorner]}  \emph{adj.} dictionary.
(\ding{43}~\HL{lI'u}{\B{lì'u}}, \HL{puk}{\B{puk}}) 

% lì'uvan
\hi
\HT{lI'uvan}{\B{\cp{lì}'uvan}} \I{["lI.Pu.van]}  \emph{n.} pun, word\hyp{}play.   
\B{Ay\-lì\-'u\-van a\-swey lu 'i\-pu, lu sìl\-ron\-sem.} The best puns are both funny and clever.~\LINK{http://naviteri.org/2012/01/mipa-zisit-ayliu-amip-new-words-for-the-new-year/}{b}
(\ding{43}~\HL{lI'u}{\B{lì\-'u}}, \HL{uvan}{\B{u\-van}})

% lì'uvi
\hi 
\HT{lI'uvi}{\B{\cp{lì}'uvi}} \I{["lI.Pu.vi]} \emph{n.} affix. 
\B{Lu tsa\-lì\-'ur alu fì\-haw\-re'\-ti me\-lì\-'u\-vi alu 'aw\-a eo\-lì\-'u\-vi sì 'aw\-a uo\-lì\-'u\-vi.} The word \emph{fì\-haw\-re'\-ti} has two affixes--one prefix and one suffix.~\LINK{http://naviteri.org/2017/10/more-language-for-talking-about-language/}{b}
(\ding{43}~\HL{lI'u}{\B{lì\-'u}}) \\
\textsc{Derivatives}: \ding{43} \on{1}.~\HL{eolI'uvi}{\B{eo\-lì\-'u\-vi}}, prefix. \\
\on{2}.~\HL{uolI'uvi}{\B{uo\-lì\-'u\-vi}}, suffix.

% lìktap
\hi
\HT{lIktap}{\B{\cp{lìk}tap}} \I{["lIk\textcorner.tap\textcorner]}  \emph{adj.} crooked.
\B{Ke tsun fko fì\-swi\-zaw\-ti si\-var--lu lìk\-tap.} This arrow can't be used--it's crooked.~\LINK{http://naviteri.org/2019/12/tengkrr-zisit-leratem-as-the-year-changes/}{b}
(\ding{214}~\HL{yey}{\B{yey}})

% lìlì
\hi 
\HT{lIlI}{\B{\cp{lì}lì}} \I{["lI.lI]} \emph{n.} breast. 

% lìm
\hi
\HT{lIm}{\B{lìm}} \I{[lIm]}  \emph{vin.} far, be far.
\B{'Ì\-'awn nì\-fya\-'o a lìm.} Stay away [lit.: remain in a way that is far].~\LINK{http://forum.learnnavi.org/language-updates/near-distant-and-irregular-adverbs/}{ln}
(\ding{43}~\HL{alIm}{\B{alìm}}; \ding{214}~\HL{sim}{\B{sim}})

% lìmsimpe 
\hi 
\HT{lImsimpe}{\B{\cp{lìm}simpe}} \I{["lIm.sim.pE]} \emph{inter.} how near? how far? what distance? (\emph{less common than} \ding{43}~\HL{la'ape}{\B{la'ape}}). 
(\ding{43}~\HL{pelImsim}{\B{pelìmsim}})

% lìng
\hi
\HT{lIng}{\B{lìng}} \I{[lIN]}  \emph{vin.} float in the air or water, hover. 
\ding{43}~\HL{ayram alusIng}{\B{Ay\-ram A\-lu\-sìng}}, the Floating Mountains. \\
\textsc{Derivatives}: \on{1}.~\HL{lIngtskxe}{\B{lìng\-tskxe}}, unobtanium.

% lìngtskxe
\hi 
\HT{lIngtskxe}{\B{\cp{lìng}tskxe}} \I{["lIN.\t{ts}k'E]} \emph{n.} unobtanium. 
(\ding{43}~\HL{lIng}{\B{lìng}}, \HL{tskxe}{\B{tskxe}})

% lo'a
\hi 
\HT{lo'a}{\B{\cp{lo}'a}} \I{["lo.Pa]} \emph{n.} totem; amulet. \\
\textsc{Derivatives}: \ding{43} \on{1}.~\HL{lo'akur}{\B{lo'akur}}, Toruk Makto Amulet.

% lo'ak
\hi 
\HT{lo'ak}{\B{\cp{Lo}'ak}} \I{["lo.Pak\textcorner]} \emph{n. male name}. 

% lo'akur
\hi 
\HT{lo'akur}{\B{lo'a\cp{kur}}} \I{[lo.Pa."kur]} \emph{or} \I{[."kUR]} (\textsc{rn}: \B{loa\cp{kùr}} \I{[lo.a."kUR]}) \emph{n.} Toruk Makto Amulet.
(\ding{43}~\HL{lo'a}{\B{lo'a}}, \HL{kur}{\B{kur}})

% Loak
\hi
\B{\cp{Lo}ak} \I{["lo.ak\textcorner]}  \emph{n., male name}. 
\B{Lo\-ak kal\-txì so\-li oer.} Loak greeted me.~\LINK{http://naviteri.org/2010/09/getting-to-know-you-part-1/comment-page-1/\#comment-294}{b}
\B{Lo\-ak\-ìl pä\-nu\-to\-lìng fu\-ta kar oe\-ru fya\-'ot a 'ìp fko nem\-fa ewll.} Loak promised he'd teach me how to vanish into the bushes.~\LINK{http://naviteri.org/2013/01/awvea-posti-zisita-amip-first-post-of-the-new-year/}{b}
\B{Tse\-nu Lo\-ak\-ur ke ha'.} Tsenu is a bad match for Loak.~\LINK{http://naviteri.org/2012/10/mipa-vospxi-mipa-ayliu-new-words-for-the-new-month/}{b}
\B{Ngal mo\-lay' pxaw\-pxun\-it Lo\-ak\-ä a krr pe\-kxin\-um?} When you tried on Loak's armband, how tight was it?~\LINK{http://naviteri.org/2012/07/meetings-waterfalls-and-more/}{b}

% loho
\hi 
\HT{loho}{\B{\cp{lo}ho}} \I{["lo.ho]} \emph{vin.} be surprising. 
\B{Tä\-ftxu\-tswo Ri\-ni\-yä lo\-ho oer nì\-txan.} Rini's ability to weave surprises me a lot.
\B{Fo tsìk sro\-ler a fì\-'u lo\-lo\-ho po\-an\-ur.\slash Lo\-lo\-ho po\-an\-ur fwa fo tsìk sro\-ler.} It surprised him that they suddenly appeared.~\LINK{http://naviteri.org/2019/06/50a-liu-amip-40-new-words/}{b} \\
\textsc{Derivatives}: \ding{43} \on{1}.~\HL{tIloho}{\B{tì\-lo\-ho}}, surprise. \\
\on{2}.~\HL{nIloho}{\B{nì\-lo\-ho}}, surprisingly. 

% loi
\hi
\HT{loi}{\B{\cp{lo}i}} \I{["lo.i]}  \emph{n.} egg.
\B{Ro\-lun ay\-oel tsrul\-mì hì\-'i\-a pxe\-lo\-it a\-teyr.} We found three little white eggs in the nest.~\LINK{http://naviteri.org/2013/06/looking-back-looking-forward/}{b}

% lok
\hi
\HT{lok1}{\B{lok}}\textsuperscript{\on{1}} \I{[lok\textcorner]}  \emph{adp.} close to. 
\B{(fwa sle\-le) mì hil\-van a lok Kel\-ut\-ral}, to swim in a river near Hometree.~\LINK{http://naviteri.org/2010/07/thoughts-on-ambiguity/}{b}

% lok
\hi
\HT{lok2}{\B{lok}}\textsuperscript{\on{2}} \I{[lok\textcorner]}  \emph{vtr.} get close to, approach.   
\B{Na\-ri so\-li ay\-oe fte\-ke nì\-hawng li\-vok.} We were careful not to get too close.~\LINK{http://wiki.learnnavi.org/index.php?title=Corpus\#New_York_Times}{wiki}
\B{Fra\-po ne mì\-fa! Le\-rok yrr\-ap apxa!} Everybody inside! A big storm is approaching!~\LINK{http://naviteri.org/2012/07/meetings-waterfalls-and-more/}{b}
\B{Sra\-ke pol la\-yok txew\-ti na'\-rìng\-ä?} Will he approach the edge of the forest?~\LINK{http://naviteri.org/2012/03/spring-vocabulary-part-1/}{b}
\B{Ke tsun aw\-nga pi\-vey nul\-krr -- txew lok.} We can't wait any longer -- time is almost up.~\LINK{http://naviteri.org/2012/03/spring-vocabulary-part-1/}{b}

% lom
\hi
\HT{lom}{\B{lom}} \I{[lom]}  \emph{adj.} missing, missed (\emph{as an absent person who is longed for; covers only s.t.~one once had but no longer does; a sense of emotional loss; not s.t.~one lacks but never had in the first place}).
\B{Nga lom lu oer.} I miss you.~\LINK{http://naviteri.org/2011/07/txantsana-ultxa-mi-siatll-great-meeting-in-seattle/}{b}
\B{Ay\-sre' lom lu tsa\-ko\-ak\-tan\-ur.} The old man misses his teeth.~\LINK{http://naviteri.org/2011/07/txantsana-ultxa-mi-siatll-great-meeting-in-seattle/}{b} \\
\textsc{Derivatives}: \ding{43} \on{1}.~\HL{lomtu}{\B{lom\-tu}}, missed person.

% lomtu
\hi
\HT{lomtu}{\B{\cp{lom}tu}} \I{["lom.tu]}  \emph{n.} (\emph{reserved for special circumstances}) missed person.
(\emph{a. idiom}) \B{Teng\-krr ftxo\-zä se\-rei\-yi aw\-nga, ke tswi\-va' ay\-lom\-tu\-ti ko!} A toast: To absent friends. [lit.: While we are celebrating, let us not forget those who we wish could also be here (but can't)]~\LINK{http://naviteri.org/2011/07/txantsana-ultxa-mi-siatll-great-meeting-in-seattle/}{b}
(\ding{43}~\HL{lom}{\B{lom}}) 

% lonu
\hi
\HT{lonu}{\B{lo\cp{nu}}} \I{[lo."nu]}  \emph{vtr.} release, let go. 
\B{Pot lo\-nu.} Release him.~\LINK{http://wiki.learnnavi.org/index.php?title=Na\%27vi_from_Avatar_Movie\#Scene_in_Hometree}{a\slash wiki}
\B{Nì\-sngä\-'i fmawn\-it fo nar\-mew wi\-van, slä nì\-'i'\-a fra\-por lo\-lo\-nu.} They originally wanted to hide the news, but in the end they revealed it everyone.~\LINK{http://naviteri.org/2011/10/more-vocabulary-a-bit-of-grammar/}{b}
(\emph{a. proverb}) \B{Fwa kan ke tam; ze\-ne swi\-zaw\-it li\-vo\-nu.} Intent is not enough; it's action that counts. [lit.: To aim is not enough; one must release the arrow.]~\LINK{http://naviteri.org/2013/02/vospxi-ayol-posti-apup-short-post-for-a-short-month/}{b} \\
\textsc{Derivatives}: \ding{43} \on{1}.~\HL{nIklonu}{\B{nìk\-lo\-nu}}, firmly, steadfastly, faithfully. \\
\on{2}.~\HL{stxenu}{\B{stxe\-nu}}, offer. \\
\on{3}.~\HL{lonusye}{\B{lo\-nu\-sye}}, exhale, blow.

% lonusye
\hi 
\HT{lonusye}{\B{lo\cp{nu}sye}} \I{[l$\cdot$o."n$\cdot$u.sjE]} \emph{vin., inf.}\on{12} exhale; blow (\ding{43}~\HL{ne}{\B{ne}}, on). 
\B{Txo syu\-ve som li\-vu nì\-hawng, lo\-nu\-sye tsa\-ne.} If the food is too hot, blow on it.~\LINK{http://naviteri.org/2020/07/vospxivol-lefpom-happy-august/}{b}
\B{Lo\-nu\-sye tong txep\-tsyìp\-it.} Blow out the flame.~\LINK{http://naviteri.org/2020/07/vospxivol-lefpom-happy-august/}{b}
(\ding{43}~\HL{lonu}{\B{lo\-nu}}, \HL{syeha}{\B{sye\-ha}}; \ding{214}~\HL{mungsye}{\B{mung\-sye}}) 

% lopx
\hi
\HT{lopx}{\B{lopx}} \I{[lop']}  \emph{vin.} panic.
\B{Pa\-lu\-kan\-it tso\-le\-'a, ye\-rik lopx hi\-fwo kxam\-lä ze\-swa.} Spotting a thanator, the hexapede panicked and escaped through the grass.~\LINK{http://naviteri.org/2012/07/meetings-waterfalls-and-more/}{b}

% lor
\hi
\HT{lor}{\B{lor}} \I{[loR]}  \emph{adj.} beautiful, pleasant to the senses.
\B{New oe ti\-vìng ay\-ngar lì\-'ut a tì\-'e\-fu\-mì oe\-yä lu lor fra\-to mì lì'\-fya le\-Na'\-vi.} I want to present to you the word that, in my opinion, is the most beautiful in the Na'vi language.~\LINK{http://forum.learnnavi.org/language-updates/trr-rrtaya/}{ln}
\B{Fay\-lì\-'u a\-lor oe\-ru te\-ya si nì\-txan, ul\-te ke tswa\-ya' oel sat kaw\-krr.} These beautiful words fill me with joy and I will never forget them.~\LINK{http://forum.learnnavi.org/language-updates/life-death-and-gerunds/}{ln} 
\B{Lor\-a fì\-tì\-kang\-kem me\-nge\-yä meuia lu\-yu aw\-nga\-ru nì\-wotx.} This beautiful work of theirs honors us very much.~\LINK{http://naviteri.org/2010/12/beautiful-christmas-carols-sung-in-navi/}{b}
\B{Lu po lor\-a tu\-té a\-lor.} She's an extremely beautiful woman. \emph{or} \B{Tsa\-tu\-té lu lor\-a pum a\-lor.} That woman is an extremely beautiful one.~\LINK{http://naviteri.org/2013/02/vospxi-ayol-posti-apup-short-post-for-a-short-month/}{b}
(\ding{214}~\HL{vA'}{\B{vä'}}) \\  
\textsc{Derivatives}: \on{1}.~\HL{tIlor}{\B{tì\-lor}}, beauty.   
\on{2}.~\HL{ftxIlor}{\B{ftxì\-lor}}, delicious. 
\on{3}.~\HL{onlor}{\B{on\-lor}}, good-smelling. 
\on{4}.~\HL{miklor}{\B{mik\-lor}}, pleasant sounding. 
\on{5}.~\HL{loran}{\B{lo\-ran}}, elegance. 
\on{6}.~\HL{loreyu}{\B{lo\-re\-yu}}, helicordian.
\on{7}.~\HL{lortsyal}{\B{lor\-tsyal}}, shimmyfly.
\on{8}.~\HL{kinglor}{\B{king\-lor}}, kinglor. 
\on{9}.~\HL{zeklor}{\B{zek\-lor}}, pleasant to the touch.
\on{10}.~\HL{lornusrr}{\B{lor\-nu\-srr}}, radiant.

% loran
\hi
\HT{loran}{\B{\cp{lo}ran}} \I{["lo.Ran]}  \emph{n.} elegance, grace.   
\B{Ya\-mì tsun fko tsi\-ve\-'a lo\-ran\-it re\-nu\-ä kil\-van\-ä slä kll\-te\-sìn wä\-pan.} From the air you can see the grace of the river's form but from the ground it's hidden.~\LINK{http://naviteri.org/2012/10/mipa-vospxi-mipa-ayliu-new-words-for-the-new-month/}{b} 
(\ding{43}~\HL{lor}{\B{lor}}, \HL{ran}{\B{ran}})

% loreyu
\hi
\HT{loreyu}{\B{lo\cp{re}yu}} \I{[lo."RE.ju]} \emph{n.} \ding{96} helicoradian, ``beautiful spiral'', \emph{helicoradium spirale}. 
\B{Mì na'\-rìng Tsyeyk\-ìl lo\-re\-yu\-ti 'o\-lam\-pi a tsa\-swaw wou oer nì\-wotx!} That instant when Jake touched the \emph{Helicoradiam Spirale} in the forest I was totally amazed.~\LINK{http://forum.learnnavi.org/language-updates/txelanit-hivawl/}{ln}
(\emph{a. idiom}) \B{(Na) lo\-re\-yu 'aw\-nam\-pi}, extreme shyness, [lit.: (like) a touched helicoradian] \B{Lu por mok\-ri a\-mik\-lor, slä lo\-re\-yu 'aw\-nam\-pi lu. Ke tsun ri\-vol eo su\-te.} She has a beautiful voice, but she's extremely shy. She can't sing in front of people.~\LINK{http://naviteri.org/2014/11/twenty-before-the-holidays/}{b}
(\ding{43}~\HL{lor}{\B{lor}}, \HL{'Iheyu}{\B{'ì\-he\-yu}})

% lornusrr
\hi 
\HT{lornusrr}{\B{\cp{lor}nusrr}} \I{["loR.nu.s\s{r}]} \emph{adj.} radiant. 
\B{Ey\-we\-veng lu lor\-nu\-srr.} Pandora is radiant.~\LINK{https://naviteri.org/2025/06/vospikin-lefpom-happy-july/}{b} 
(\ding{43}~\HL{lor}{\B{lor}}, \HL{nrr}{\B{nrr}}) 

% lortsyal
\hi 
\HT{lortsyal}{\B{\cp{lor}tsyal}} \I{["loR.\t{ts}jal]} \emph{n., fauna} shimmyfly.~\LINK{https://forum.learnnavi.org/language-updates/expansion-to-flora-fauna-and-places-from-the-valley-of-mo'ara/}{ln}
(\ding{43}~\HL{lor}{\B{lor}}, \HL{tsyal}{\B{tsyal}})

% lrrtok / si
\hi
\HT{lrrtok}{\B{\cp{lrr}tok}} \I{["l\s{r}.tok\textcorner]} \textsc{i}.~\emph{n.} smile. 
\B{Nge\-yä lrr\-tok lor lu nì\-txan.} Your smile is very beautiful.
\B{Ay\-nge\-yä ay\-u\-pxa\-re\-ri sì ay\-lrr\-tok\-ìri sei\-yi oe ira\-yo.} Thank you for all your messages and smiles.~\LINK{http://forum.learnnavi.org/news-announcements/happy-birthday-to-karyu-pawl/msg556899/\#msg556899}{ln}
(\emph{a.~idiom}) \B{Lrr\-tok ngar!\slash Ay\-lrr\-tok ngar (li\-vu)!} Good luck.\slash Good luck to you.~\LINK{http://forum.learnnavi.org/language-updates/good-luck-with-this-one!/}{ln} 
\B{Zì\-tsal\-trr\-ìri tì\-mun\-txa\-yä ay\-lrr\-tok!} Happy anniversary (of your marriage)!~\LINK{http://naviteri.org/2017/12/zisit-amip-lefpom-ma-eylan-happy-new-year-friends/}{b} \\
\textsc{ii}.~\HT{lrrtok si}{\B{\cp{lrr}tok si}} \I{["l\s{r}.tok\textcorner{} si]}  \emph{vin.}~(\emph{a.}) smile. 
\B{Prr\-nen lrr\-tok si a krr, srer me\-song\-tsyìp a\-ho\-na.} When the baby smiles, two adorable dimples appear.~\LINK{http://naviteri.org/2012/03/spring-vocabulary-part-1/}{b}
\B{Nì\-law Nawm\-a Sa'\-nok lrr\-tok so\-li me\-ngar.} Clearly the Great Mother smiled on you both.~\LINK{http://forum.learnnavi.org/language-updates/luke-kxu-without-harm/}{ln} 
(\emph{b. idiom}) shine (of the sun) \B{Tsaw\-ke lrr\-tok si.} The sun shines.~\LINK{http://forum.learnnavi.org/language-updates/some-news-from-berlin/}{ln}

% lu
\hi
\HT{lu}{\B{lu}} \I{[lu]}  \emph{vin.} \on{1}.~(\emph{copula}) be (am, is, are, was, were etc.).
\B{Rel\-tse\-o\-tu a\-txan\-tsan lu nga!} You are an excellent artist!~\LINK{http://forum.learnnavi.org/language-updates/art-related-vocab/}{ln}
\B{Taw\-sìp nge\-yä lu sngel\-tseng.} Your ship is a garbage scow.~\LINK{http://wiki.learnnavi.org/index.php?title=Corpus\#UGO_Movie_Blog_.28Twitter_Questions.29}{wiki\slash ugo}
\B{Lu nga win sì txur \slash{} Lu nga txan\-tslu\-sam}, You are fast and strong \slash{} You are wise.~\LINK{http://wiki.learnnavi.org/index.php/Hunting_Song}{wiki}
\B{Fì\-tskxe\-ri fa\-'o lu yey; ke lu ko\-um.} This rock has straight sides; it's not rounded.~\LINK{http://naviteri.org/2013/09/on-si-salewfya-shapes-and-directions/}{b} 
\on{2}.~(\emph{a. to indicate existence}) it is, it's. \B{Nga\-ru tì\-yawr, lu kxey\-ey.} You are right, it's a mistake.~\LINK{http://forum.learnnavi.org/language-updates/a-long-answer-to-a-quick-question/msg155104/\#msg155104}{ln} 
(\emph{b.}) there is\slash are.
\B{Äo fì\-ut\-ral lu tsmìm 'ang\-tsìk\-ä.} There's a hammerhead track under this tree.~\LINK{http://naviteri.org/2011/09/\%E2\%80\%9Cby-the-way-what-are-you-reading\%E2\%80\%9D/}{b}
\B{Lo\-lu txan\-tsan\-a ul\-txa mì Si\-ä\-tll.} There was a great meeting in Seattle.~\LINK{http://naviteri.org/2011/07/attributive-\%E2\%80\%9Ca\%E2\%80\%9D-and-truncated-style/}{b}
\B{Lo\-lu ka\-vuk, slä Tse\-nul tì\-ngay\-it ko\-lu\-lat.} There was treachery, but Tsenu revealed the truth.~\LINK{http://naviteri.org/2011/10/more-vocabulary-a-bit-of-grammar/}{b}
\B{Tskxe\-keng\-lu\-ke ke lu kea tì\-tsan\-'ul.} Without practice there is no improvement.~\LINK{http://naviteri.org/2013/03/tsanerul-feerul-getting-better-getting-worse/}{b}
(\emph{c. as parallel structure}) \B{Fra\-u\-van\-ìri lu yo\-ra'\-tu, lu snay\-tu.} For every game, there's a winner and a loser.~\LINK{http://naviteri.org/2011/12/a-note-on-the-word-\%E2\%80\%9Cyora\%E2\%80\%99tu\%E2\%80\%9D/}{b}
\B{Su\-te\-kip nì\-wotx, ftxey Na'\-vi, ftxey Saw\-tu\-te, lu sìl\-tsan, lu kawng.} There are good and bad among all people, whether Na'vi or Skypeople.~\LINK{http://naviteri.org/2010/07/ting-mikyun-fte-tslivam-listening-comprehension-1-3/}{b}
\on{3}.~(\emph{in combination with the dative to indicate possession}) have. (`have' constructions usually begin with the verb)
\B{Nga\-ru lu fpom srak?} How are you? [lit.: Do you have well-being?]~\LINK{http://naviteri.org/2010/07/vocabulary-update/}{b}
\B{La\-yu oe\-ru ye'\-rìn sìl\-tsan\-a fmawn.} Soon, I will have good news.~\LINK{http://forum.learnnavi.org/news-announcements/a-response-from-paul-frommer!/}{ln}
\B{Ke lu oer set krr a\-txan, ul\-te ke new tì\-kxey si\-vi.} I don't have much time now, and I don't want to make a mistake.~\LINK{http://forum.learnnavi.org/language-updates/verb-for-write/msg139677/\#msg139677}{ln}
(\emph{a. allows double dative})
\B{Lu oe\-ru ay\-lì\-'u fra\-por.} I have something to say to everyone.~\LINK{http://forum.learnnavi.org/language-updates/i-have-something-to-say/}{ln}
\B{Lu fra\-por ay\-lì\-'u oe\-ru.} Everyone has something to say to me.~\LINK{http://forum.learnnavi.org/language-updates/i-have-something-to-say/}{ln}
(\emph{b. idiomatic uses})
\B{lu fko\-ru ya\-yayr}, to be confused. \B{Saw\-tu\-te\-yä hem\-ìri lu aw\-nga\-ru ya\-yayr.} The Skypeople's actions confuse us.~\LINK{http://naviteri.org/2011/01/new-year-new-vocabulary/}{b}
\B{Slä 'en si oe, tsa\-me\-lì\-'u oe\-ru ya\-yayr lo\-la\-tsu.} But I guess, those two words would have confused me.~\LINK{http://naviteri.org/2011/04/\%E2\%80\%99a\%E2\%80\%99awa-li\%E2\%80\%99fyavi-amip\%E2\%80\%94a-few-new-expressions/comment-page-1/\#comment-677}{b}
\B{lu fko\-ru sngum}, to be worried. \B{Lu oe\-ru sngum a sa\-ron\-yu ke tì\-ye\-vä\-txaw.} I'm worried that the hunters will not return.~\LINK{http://naviteri.org/2011/01/new-year-new-vocabulary/}{b}
\B{lu fko\-ru ing\-yen}, to be puzzled. \B{Lu oer ing\-yen a Ìstaw nim lu fì\-txan kum\-a pxìm wä\-pan.} I'm puzzled that Ìstaw is so shy that he frequently hides.~\LINK{http://naviteri.org/2013/01/awvea-posti-zisita-amip-first-post-of-the-new-year/}{b}
\B{lu fko\-ru tì\-kin a} \texttt{+} \emph{verb}, need to do s.t. \B{Po\-ri aw\-nga\-ru lu tì\-kin a nu\-me nì\-'ul.} We need to learn more about him.~\LINK{http://wiki.learnnavi.org/index.php?title=Na\%27vi_from_Avatar_Movie\#Scene_in_Hometree}{a\slash wiki}
\B{lu fko\-ru haw\-tsyìp}, to take a nap. 
\B{Oel new fu\-ta li\-vu oer set haw\-tsyìp.} I want to take a nap now.~\LINK{http://naviteri.org/2011/07/txantsana-ultxa-mi-siatll-great-meeting-in-seattle/}{b} 
\B{Sa'\-nur le\-ru haw\-tsyìp.} Mommy is taking a nap.~\LINK{http://naviteri.org/2011/07/txantsana-ultxa-mi-siatll-great-meeting-in-seattle/}{b}
\B{lu fko\-ru txoa (a ta tu\-teo)}, to forgive s.o. \B{Oe\-ru txoa li\-vu.} Forgive me. [lit.: may I have forgiveness]~\LINK{http://forum.learnnavi.org/language-updates/verb-for-write/msg139677/\#msg139677}{ln} \B{Säm\-yam\-ìl po\-ru wa\-yìn\-txu fu\-ta nga\-ta lo\-lu li txoa.} A hug will show him that you've already forgiven him.~\LINK{http://naviteri.org/2011/10/more-vocabulary-a-bit-of-grammar/}{b}
\B{lu fko\-ru fmokx}, to be jealous (\emph{neutral unless otherwise specified with} ‹\B{äng}› \emph{or} ‹\B{ei}›). \B{Fu\-ria fì\-txan fnä\-ngan Ul\-rey\-ìl tsko swi\-zaw\-it, lu oe\-ru fmokx.} I'm jealous of the fact that Ulrey is such an excellent archer.~\LINK{http://naviteri.org/2011/10/more-vocabulary-a-bit-of-grammar/}{b}
\B{lu fko\-ru nrra}, to be proud of s.o. (with the topic). \B{Nga\-ri le\-iu oe\-ru nrra nì\-txan, ma 'ite.} I'm very proud of you, daughter.~\LINK{http://naviteri.org/2012/07/fivospxiya-aylifyavi-amip-this-months-new-expressions-pt-2/}{b}
\B{lu fko\-ru yew\-la}, to be disappointed. \B{Oer lu txan\-a yew\-la a ke tsun nga oe\-hu ki\-vä\-teng me\-srr\-ay.} I'm very disappointed you can't hang out with me the day after tomorrow.~\LINK{http://naviteri.org/2012/10/mipa-vospxi-mipa-ayliu-new-words-for-the-new-month/}{b}
\B{lu fko\-ru tì\-yawr}, to be right. \B{Nga\-ru tì\-yawr, lu kxey\-ey.} You are right, it's a mistake.~\LINK{http://forum.learnnavi.org/language-updates/a-long-answer-to-a-quick-question/msg155104/\#msg155104}{ln}
\B{lu fko\-ru tì\-kxey}, to be wrong. \B{Po ke tsun pi\-vll\-ngay san oe\-ru tì\-kxey.} He can't admit he's wrong.~\LINK{http://naviteri.org/2011/07/txantsana-ultxa-mi-siatll-great-meeting-in-seattle/}{b}
\B{lu fko\-ru ho\-an}, to feel comfortable.
\B{Sko frr\-tu, fra\-krr mì hel\-ku nge\-yä lu oe\-ru ho\-an nì\-txan, ma tsmuk.} I always feel very comfortable as a guest in your home, brother.~\LINK{http://naviteri.org/2015/08/aylifyavi-lereyfya-2-cultural-terms-2-and-more/}{b}
\B{lu fko\-ru tì\-hì\-pey}, to be hesitant.
\B{Sar tsa\-lì\-'ut a fì\-'u\-ri lu oe\-ru tì\-hì\-pey nì\-'it.} I'm a bit hesitant about using that word.~\LINK{http://naviteri.org/2015/11/vomuna-liu-amip-ten-new-words/}{b}
\on{4}.~(\emph{can be omitted in coll. conversation}) 
\B{Nga (lu) pe\-su?} Who are you?~\LINK{http://naviteri.org/2021/04/mipa-ayliu-mipa-sioeykting-new-words-new-explanations/}{b} \\
\textsc{Derivatives}: \ding{43} \on{1}.~\HL{luke}{\B{lu\-ke}}, without. 
\on{2}.~\HL{tIkelu}{\B{tì\-ke\-lu}}, lack.

% luan
\hi 
\HT{luan}{\B{\cp{lu}an}} \I{["lu.an]} (\textsc{rn}: \B{\cp{lu}an} \I{["lu.an]}) \emph{vtr.} owe (\emph{moral obligation to give s.t. to s.o.}). 
\B{Fol nge\-yä tsmu\-ke\-ru lu\-an tsko\-ti a\-mip} They owe your sister a new bow.~\LINK{https://naviteri.org/2024/02/mipa-ayliu-si-aylifyavi-niul-more-new-words-and-expressions/}{b}
\B{Oey vo\-ìk\-ìri a\-lew\-nga' lu\-an oel ngar tì\-oeyk\-tìng\-it.} I owe you an explanation for my shameful behavior.~\LINK{https://naviteri.org/2024/02/mipa-ayliu-si-aylifyavi-niul-more-new-words-and-expressions/}{b}

% luke
\hi
\HT{luke}{\B{\cp{lu}ke}} \I{["lu.kE]}  \emph{adp.} without.
\B{Lu me\-nge\-yä kel\-ku na'\-rìng\-ä lu\-ke kxu a\-txan a fì\-'u fmawn a\-sìl\-tsan lu nì\-ngay.} It's truly good news that your forest house is without great harm.~\LINK{http://forum.learnnavi.org/language-updates/luke-kxu-without-harm/}{ln} 
\B{Ke tsun oe tsaw\-te\-ri fpi\-vìl lu\-ke fwa sngä\-'i tsngi\-vaw\-vìk.} I can't think about it without starting to cry.~\LINK{http://forum.learnnavi.org/language-updates/txelanit-hivawl/}{ln}
\B{Sra\-ke tso\-lun fra\-'ut tsli\-vam kxey\-ey\-lu\-ke?} Could you understand everything without a mistake?~\LINK{http://naviteri.org/2012/10/txon-eywaevenga-text-and-translation/}{b} 
(\emph{a. idiom}) \B{(na) la\-nay'\-ka lu\-ke re\-'o}, ``completely lost,'' [lit.: Like a slinger without a head].~\LINK{http://naviteri.org/2021/08/ulte-ayyoratu-leiu-and-the-winners-are-2/}{b}
\B{Po maw kxitx mun\-txa\-tu\-ä 'a\-me\-fu na la\-nay'\-ka re'\-o\-lu\-ke.} After the death of his spouse, he felt completely lost.~\LINK{http://naviteri.org/2021/08/ulte-ayyoratu-leiu-and-the-winners-are-2/}{b} \\
\textsc{Derivatives}: \ding{43} \on{1}.~\HL{lukpen}{\B{luk\-pen}}, naked. 
\on{2}.~\HL{kxuke}{\B{kxu\-ke}}, safe. 
\on{3}.~\HL{velke}{\B{vel\-ke}}, chaotic. 
\on{4}.~\HL{am'aluke}{\B{am\-a'a\-lu\-ke}}, without a doubt. 
\on{5}.~\HL{letxiluke}{\B{le\-txi\-lu\-ke}}, unhurried. 
\on{6}.~\HL{nItxiluke}{\B{nì\-txi\-lu\-ke}}, unhurriedly, leisurely. 
\on{7}.~\HL{nItkanluke}{\B{nìt\-kan\-lu\-ke}}, accidentally. 
\on{8}.~\HL{renulke}{\B{re\-nul\-ke}}, irregular. 
\on{9}.~\HL{am'ake}{\B{am\-'a\-ke}}, sure, confident.
\on{10}.~\HL{txewluke}{\B{txew\-lu\-ke}}, endless, boundless.
\on{11}.~\HL{tI'iluke}{\B{tì\-'i\-lu\-ke}}, endless, never\hyp{}ending.
\on{12}.~\HL{lukftang}{\B{luk\-ftang}}, constant, continual.

% lukftang
\hi 
\HT{lukftang}{\B{luk\cp{ftang}}} \I{[luk\textcorner."ftaN]} \emph{or} \I{[lUk\textcorner.]} \emph{adj.} constant, continual. 
\B{Pe\-yä tì\-pu\-sll\-txel a\-luk\-ftang oe\-ti srätx.} His constant talking annoys me.~\LINK{http://naviteri.org/2021/03/mipa-ayliu-si-aylifyavi-new-words-and-expressions/}{b}
(\ding{43}~\HL{luke}{\B{lu\-ke}}, \HL{tIftang}{\B{tì\-ftang}}; \HL{letut}{\B{le\-tut}})

% lukpen
\hi
\HT{lukpen}{\B{luk\cp{pen}}} \I{[luk\textcorner."pEn]} \emph{or} \I{[lUk\textcorner.]}  \emph{adj.} naked, without clothing. 
(\ding{43}~\HL{luke}{\B{lu\-ke}}, \HL{pxen}{\B{pxen}})

% lumpe
\hi
\HT{lumpe}{\B{\cp{lum}pe}} \I{["lum.pE]} \emph{or} \I{["lUm.]} \emph{or} \HL{pelun}{\B{pe\cp{lun}}} \I{[pE."lun]} \emph{or} \I{[."lUn]} \emph{inter.} why? what reason?
\B{Oel lum\-pe tsat so\-lung mì upxa\-re?} Why did I add that in the message?~\LINK{http://naviteri.org/2010/07/vocabulary-update/comment-page-1/\#comment-95}{b}
\B{Tey\-nga lum\-pe fo ho\-lum ke lu law.} It's not clear why they left.~\LINK{http://naviteri.org/2011/08/reported-speech-reported-questions/}{b}
\B{Nge\-yä tsngal lum\-pe lu mek?} Why is your cup empty?~\LINK{http://naviteri.org/2012/02/trr-asawnung-lefpom-happy-leap-day/}{b}
(\ding{43}~\HL{lun}{\B{lun}})

% lun
\hi
\HT{lun}{\B{lun}} \I{[lun]} \emph{or} \I{[lUn]} \emph{n.} reason. 
(\ding{43}~\HL{oeyk}{\B{oeyk}}) \\  
\textsc{Derivatives}: \ding{43} \on{1}.~\HL{nIlun}{\B{nì\-lun}}, of course, logically.   
\on{2}.~\HL{talun1}{\B{ta\-lu\-na\slash a\-lun\-ta}}, because. 
\on{3}.~\HL{talun2}{\B{ta\-lun}}, because of. 
\on{4}.~\HL{lumpe}{\B{pe\-lun\slash lum\-pe}}, why? 
\on{5}.~\HL{velun}{\B{ve\-lun}}, logic.

% lupra
\hi
\HT{lupra}{\B{\cp{lup}ra}} \I{["lup\textcorner.Ra]} \emph{or} \I{["lUp\textcorner.]} (\textsc{rn}: \B{\cp{lùp}ra} \I{["lUp\textcorner].Ra}) \emph{n.} style.
\B{Pll\-txe fra\-po san fì\-fkxi\-le lor lu nì\-ngay sìk, slä oe\-ri ke su\-nu oer lup\-ra kaw\-'it.} Everybody says this necklace is really beautiful, but me, I don't like the style one bit.~\LINK{http://naviteri.org/2014/07/tsawlultxamaw-a-ayliu-post-meetup-words/}{b} \\
\textsc{Derivatives}: \ding{43} \on{1}.~\HL{fyolup}{\B{fyo\-lup}}, exquisite, sublime in taste. 
\on{2}.~\HL{fe'lup}{\B{fe'\-lup}}, tacky, in poor taste. 
\on{3}.~\HL{snolup}{\B{sno\-lup}}, personal style, aesthetic, presence.

\end{multicols}
 
% ===== Section M =====  derivations

 \pdfbookmark[0]{M}{M}
\begin{center}
\section*{M \I{[m]} (\emph{MeM})}
\end{center}

\begin{multicols}{2}

% ma
\hi
\HT{ma}{\B{ma}} \I{[ma]} \emph{part., vocative marker} `O, \emph{name}', `dear \emph{name}' (\emph{but usually not translated}).
\B{Tew\-ti, ma Prr\-ton!} Wow, Prrton!~\LINK{http://languagelog.ldc.upenn.edu/nll/?p=1977}{ll}
\B{Oel ay\-nga\-ti ka\-meie, ma oe\-yä ey\-lan.} I see you, my friends.~\LINK{http://forum.learnnavi.org/news-announcements/a-response-from-paul-frommer!/}{ln} 
\B{Ma smu\-kan, ma smu\-ke}, Brothers and sisters!~\LINK{https://forum.learnnavi.org/language-updates/naviya-the-other-vocative/}{ln\slash a}
(\ding{43}~\HL{-ya}{\B{-ya}}) \\
\textsc{Derivatives}: \ding{43} 
\on{1}.~\HL{maitan}{\B{ma\-i\-tan}}, my son (address). 
\on{2}.~\HL{maite}{\B{ma\-i\-te}}, my daughter (address).
\on{3}.~\HL{manawmtu}{\B{ma\-nawm\-tu}}, ``excuse me''. 
\on{4}.~\HL{matu}{\B{ma\-tu}}, ``excuse me, hey''.
\on{5}.~\HL{manga}{\B{ma\-nga}}, ``hey you''. 

% maitan
\hi 
\HT{maitan}{\B{ma\cp{i}tan}} \I{[ma."{}i.tan]} \emph{ph.} my son (\emph{form of address}). 
\B{Ma\-i\-tan za\-'u fì\-tseng.} Come here, son.~\LINK{http://naviteri.org/2019/06/50a-liu-amip-40-new-words/}{b} 
(\ding{43}~\HL{'itan}{\B{'itan}})

% maite
\hi 
\HT{maite}{\B{ma\cp{i}te}} \I{[ma."{}i.tE]} \emph{ph.} my daughter (\emph{form of address}).  
(\ding{43}~\HL{'ite}{\B{'ite}})

% makto
\hi
\HT{makto}{\B{\cp{mak}to}} \I{["mak\textcorner.to]} \emph{vtr.} ride.   
\B{Trr\-am\-ä ay\-u\-van\-ìri mak\-to Ak\-wey nì\-yew\-la ha snaytx.} Akwey rode disappointingly in yesterday's games so he lost.~\LINK{http://naviteri.org/2012/10/mipa-vospxi-mipa-ayliu-new-words-for-the-new-month/}{b}
\B{Tsun fko tsa\-tse\-nge\-ne ki\-vä nì\-sle\-le fu fa fwa ik\-ran\-it mak\-to nì\-'aw.} You can only get there by swimming or riding an ikran.~\LINK{http://naviteri.org/2011/02/new-vocabulary-ii\%E2\%80\%94part-1/}{b}
\B{Oe\-yä ik\-ran sli\-vu nga, tsa\-krr oeng 'aw\-si\-teng mi\-vak\-to.} Be my ikran and let's ride together.~\LINK{http://wiki.learnnavi.org/index.php?title=Corpus\#Lemondrop.com}{wiki} 
(\emph{a.~idiom}) \B{Mak\-to fya\-pe?} ``How are things? How's everything going? How are you doing?'
\B{A: Mak\-to fya\-pe?---B: Zong. Nga\-ri tut?---A: Nìk\-sran. Oe\-ru lu fpom, slä oey 'i\-tan lu spxin.} A: How's everything?---B: Good. You?---A: So\hyp{}so. I'm fine, but my son is sick.~\LINK{http://naviteri.org/2020/05/tipusawm-tiuseyng-si-okvur-a-eltur-titxen-si-asking-answering-and-an-interesting-story/}{b} 
(\emph{b.~idiom}) \B{Mak\-to zong.\slash Mi\-vak\-to nì\-zaw\-nong.} Take care.\slash Travel safely.~\LINK{http://forum.learnnavi.org/language-updates/makto-zong!/}{ln} 
(\emph{c.~special usage}) \B{Tor\-uk Mak\-to}, Rider of the Last Shadow, \emph{legendary hero of the Na'vi}. 
\B{Tor\-uk Mak\-to po\-lä\-hem a fì\-'u ro\-lo\-'a nì\-txan Oma\-ti\-ka\-ya\-ru.} The arrival of Toruk Makto made a great impression on the Omatikaya.~\LINK{http://naviteri.org/2012/03/spring-vocabulary-part-1/}{b} 
(\emph{d.~special usage}) \B{\cp{ik}\-ran \cp{mak}\-to}, the act of riding an ikran.~\textsc{\cb{dh}} 
(\emph{e.~special usage}) \B{\cp{pa'}\-li \cp{mak}\-to}, the act of riding a pa'\-li.~\textsc{\cb{dh}}  \\
\textsc{Derivatives}: \ding{43} 
\on{1}.~\HL{kAmakto}{\B{kämakto}}, ride out.
\on{2}.~\B{Toruk Makto}, Rider of the Last Shadow.

% mal 
\hi
\HT{mal}{\B{mal}} \I{[mal]} \emph{adj.} trustworthy, trust\hyp{}inspiring.  
\B{Fì\-tì\-kang\-kem\-vi\-ri le\-tsran\-ten ke new oe hu Ra\-lu tì\-kang\-kem si\-vi. Po ke lä\-ngu mal.} I don't want to work with Ralu on this important project. He's not trustworthy, unfortunately.~\LINK{http://naviteri.org/2012/01/more-additions-to-the-lexicon/}{b}
\B{Nga mal lar\-mu oer!} I trusted you!~\LINK{http://naviteri.org/2012/01/more-additions-to-the-lexicon/}{b} 
\B{Lu tsa\-tsam\-si\-yu le\-'aw\-a ha\-pxì\-tu tsam\-po\-ngu\-ä a mal lu moer.} That warrior is the only member of the war party that we both trust.~\LINK{http://naviteri.org/2012/01/more-additions-to-the-lexicon/}{b} \\
\textsc{Derivatives}: \ding{43} 
\on{1}.~\HL{tImal}{\B{tìmal}}, trustworthiness.  
\on{2}.~\HL{nImal}{\B{nìmal}}, trustingly, without hesitation.

% mam
\hi 
\HT{mam}{\B{mam}} \I{[mam]} \emph{vtr.} wrap. 
\B{Fì\-srä\-ti pxaw sey mi\-vam fte tsat hi\-vaw\-nu.} Wrap this cloth around the bowl to protect it.~\LINK{http://naviteri.org/2019/06/50a-liu-amip-40-new-words/}{b}

% man
\hi 
\HT{man}{\B{man}} \I{[man]} \emph{vin.} belong (\ding{43}~\HL{hu}{\B{hu}}, with; \HL{mI}{\B{mì}}, in) (\emph{fitting in; feeling comfortable as part of a group; being in a place, position, or relationship where one belongs}). 
\B{Man po fì\-po\-ngu\-mì nì\-law.} He clearly belongs in this group.~\LINK{http://naviteri.org/2022/06/solalew-maul-zisita-half-the-year-is-over/}{b}
\B{Nga man oe\-hu, ma pa\-ska\-lin.} You belong with me, honey.~\LINK{http://naviteri.org/2022/06/solalew-maul-zisita-half-the-year-is-over/}{b}
\B{Krro krro fpìl oel fu\-ta ke man oe kaw\-tseng.} Sometimes I think I don't belong anywhere.~\LINK{http://naviteri.org/2022/06/solalew-maul-zisita-half-the-year-is-over/}{b}
\B{Ro\-lun oel olo'\-ti a 'e\-fu fta oe man tsa\-tseng\-mì.} I've found a community where I feel I belong.~\LINK{http://naviteri.org/2022/06/solalew-maul-zisita-half-the-year-is-over/}{b} 

% manawmtu
\hi 
\HT{manawmtu}{\B{ma\cp{nawm}tu}} \I{[ma."nawm.tu]} \emph{intj.} ``excuse me sir, excuse me madam'' (\emph{attract s.o.'s attention politely\slash honorificly}). 
\B{Ma\-nawm\-tu, sra\-ke lu\-yu nge\-nga eyk\-tan fì\-olo'\-ä?} Excuse me, sir, are you the leader of this clan?~\LINK{http://naviteri.org/2021/09/aawa-ayliu-amip-a-few-new-words/}{b}
(\ding{43}~\HL{ma}{\B{ma}}, \HL{nawmtu}{\B{nawm\-tu}}; \HL{matu}{\B{ma\-tu}}, \HL{manga}{\B{ma\-nga}})

% mankuan
\hi 
\HT{mankuan}{\B{\cp{Man}kuan}} \I{["man.ku.an]} \emph{n., clan name} the Ash People.~\LINK{https://naviteri.org/2025/11/kaltxi-nimun-hello-again/}{b}

% manga
\hi 
\HT{manga}{\B{ma\cp{nga}}} \I{[ma."Na]} \emph{intj.} ``hey, hey you'' (\emph{attract s.o.'s attention one is close\slash superior to}). 
\B{Ma\-nga! Kem\-pe si?} Hey! What are you doing?~\LINK{http://naviteri.org/2021/09/aawa-ayliu-amip-a-few-new-words/}{b} 
(\ding{43}~\HL{ma}{\B{ma}}, \HL{nga}{\B{nga}}; \HL{manawmtu}{\B{ma\-nawm\-tu}}, \HL{matu}{\B{ma\-tu}})

% Maru
\hi
\B{Maru} \I{[ma.Ru]} \emph{n. female name}. \LINK{}{dh}
\B{Tì\-mun\-txa me\-fe\-yä zo\-law\-prr\-te' Ma\-ru\-ne nì\-layl. Ke lu por kea fmokx kaw\-'it.} Their marriage brought Ma\-ru pleasure. She felt no jealousy at all.~\LINK{http://naviteri.org/2023/09/aawa-ayliu-si-aylifyavi-amip-a-few-new-words-and-expressions/}{b}

% masat
\hi
\HT{masat}{\B{\cp{ma}sat}} \I{["ma.sat\textcorner]} \emph{n.}~\ding{246} breastplate (\emph{armor}).
\B{Ma\-sat oe\-yä ley nì\-ftxan na pum nge\-yä.} My breastplate is as valuable as yours.~\LINK{http://naviteri.org/2014/03/value-and-worth/}{b}
\B{Sne\-yä ma\-sat\-it pol kä\-sro\-lìn oer.} He lent me his breastplate.~\LINK{http://naviteri.org/2012/01/more-additions-to-the-lexicon/}{b}

% Mati
\hi 
\B{Mati} \I{[ma.ti]} \emph{name}.~\LINK{https://naviteri.org/2024/02/mipa-ayliu-si-aylifyavi-niul-more-new-words-and-expressions/}{b} 

% matu
\hi 
\HT{matu}{\B{ma\cp{tu}}} \I{[ma."tu]} \emph{intj.} ``excuse me, hey'' (\emph{attract s.o.'s attention in neutral language, not overly polite\slash familiar}). 
\B{Ma\-tu, ngal hawn\-tsyokx\-it tì\-mung\-zup.} Excuse me, you just dropped your glove.~\LINK{http://naviteri.org/2021/09/aawa-ayliu-amip-a-few-new-words/}{b}
(\ding{43}~\HL{ma}{\B{ma}}, \HL{tute1}{\B{tu\-te}}; \HL{manawmtu}{\B{ma\-nawm\-tu}}, \HL{manga}{\B{ma\-nga}}) 

% matsa
% \hi
% \B{\cp{ma}tsa} \I{["ma.\t{ts}a]} \emph{n.}, \textsc{lw} Matzah (\emph{unleavened bread}). \B{aw\-ngal yom ha\-me\-tsìt, yom ma\-tsat, ke tsran\-ten}, it doesn't matter whether we eat Chametz or Matzah.~\LINK{http://whyisthisnight.com/addl-more.html\#na\%27vi}{this night}

% maw
\hi
\HT{maw}{\B{maw}} \I{[maw]} \emph{adp.} after (\emph{temporal}).
\B{Maw hì\-krr ay\-oe tì\-yä\-txaw.} We'll be right back after this break.~\LINK{http://wiki.learnnavi.org/index.php?title=Corpus\#Good_Morning_America}{wiki} 
\B{Maw txan\-tom\-pa, pxay\-a rìk\-äo la\-mu tska\-lep pe\-yä, ha tsat ku\-lat ay\-oel.} After the rainstorm, his crossbow was under a lot of leaves, so we uncovered it.~\LINK{http://naviteri.org/2011/10/more-vocabulary-a-bit-of-grammar/}{b}
\B{Maw fwa fkol Kel\-ut\-ral\-it sko\-la\-'a, lem\-rey\-a ha\-pxì\-tul tsa\-soai\-ä txo\-lu\-la mip\-a kel\-ku\-ti mì ta\-yo.} After Hometree was destroyed, the surviving members of that family built a new home on the plains.~\LINK{http://naviteri.org/2011/05/some-miscellaneous-vocabulary/}{b}
(\emph{a. of telling time}) in.  
\B{Oe pä\-hem maw\slash kay trr\-pxì\-vi a\-mrr.} I'll be there in five minutes.~\LINK{https://naviteri.org/2024/05/krrteri-about-time/}{b} \\  
\textsc{Derivatives}: \ding{43} 
\on{1}.~\HL{mawkrr}{\B{maw\-krr}}, after, afterwards. 
\on{2}.~\HL{mawkrra}{\B{maw\-krr\-a\slash a\-krr\-maw}}, after (\emph{conj.}). 
\on{3}.~\HL{pximaw}{\B{pxi\-maw}}, right after. 
\on{4}.~\HL{txon'ongmaw}{\B{txon\-'ong\-maw}}, twilight, after sunset.  
\on{5}.~\HL{trr'ongmaw}{\B{trr\-'ong\-maw}}, dawn, time after sunrise.
\on{6}.~\HL{kxamtxomaw}{\B{kxam\-txo\-maw}}, after midnight.
\on{7}.~\HL{kxamtrrmaw}{\B{kxam\-trr\-maw}}, after noon.
\on{8}.~\HL{mawfwa}{\B{maw\-fwa}}, after (\emph{conj.}).
\on{9}.~\HL{mawftxele}{\B{maw\-ftxe\-le}}, belatedly, in hindsight.

% mawftxele
\hi 
\HT{mawftxele}{\B{maw\cp{ftxe}le}} \I{[maw."ft'E.lE]} \emph{adv.} belatedly, in hindsight, after the fact, as an afterthought. 
\B{Oel pe\-yä ftxo\-zä\-ti tswo\-lä\-nga', ha pol\-txe por san ftxo\-zä\-ri ay\-lrr\-tok nga\-ru sìk maw\-ftxe\-le.} Unfortunately I forgot his birthday, so I said ``Happy Birthday'' belatedly.~\LINK{http://naviteri.org/2017/09/zisikrr-amip-ayliu-amip-new-words-for-the-new-season/}{b}  
(\ding{43}~\HL{maw}{\B{maw}}, \HL{txele}{\B{txe\-le}})

% mawfwa
\hi 
\HT{mawfwa}{\B{\cp{maw}fwa}} \I{["maw.fwa]} \emph{conj.} after.
\B{Maw\-fwa zup tom\-pa, lu ngo\-a a\-txan fì\-ha\-pxì\-mì na'\-rìng\-ä.} After rain, there's a lot of mud in this part of the forest.~\LINK{http://naviteri.org/2014/07/tsawlultxamaw-a-ayliu-post-meetup-words/}{b}
(\ding{43}~\HL{maw}{\B{maw}}, \HL{fwa}{\B{fwa}})

% mawkrr
\hi
\HT{mawkrr}{\B{\cp{maw}krr}} \I{["maw.k\s{r}]} \emph{adv.} after (\emph{temporal}), afterwards, later.
\B{Ye'\-rìn 'ì\-yi'\-a sä\-nu\-me a tsa\-ri kll\-fro' oe; maw\-krr la\-ye\-i\-u oer krr nì\-'ul.} My teaching responsibilities will soon end; after that I will have more time.~\LINK{http://forum.learnnavi.org/language-updates/trr-rrtaya/}{ln} 
\B{Nì\-'aw\-ve oeng yo\-lom wu\-tsot; maw\-krr u\-van si.} First we had dinner; afterwards we played (a game).~\LINK{http://naviteri.org/2011/08/new-vocabulary-clothing/comment-page-1/\#comment-917}{b} 
\B{Po\-lä\-hem Saw\-tu\-te kam zì\-sìt a\-mrr, hum me\-zì\-sìt maw\-krr.} The Sky People arrived five years ago and left two years later.~\LINK{http://naviteri.org/2011/09/miscellaneous-vocabulary/}{b}
(\ding{43}~\HL{maw}{\B{maw}}, \HL{krr}{\B{krr}}, \HL{mawkrra}{\B{maw\-krr\-a}}; \ding{214}~\HL{srekrr}{\B{sre\-krr}})

% mawkrra
\hi
\HT{mawkrra}{\B{maw\cp{krr}a}} \I{[maw."k\s{r}.a]} \emph{or} \HL{mawkrra}{\B{a\-\cp{\-krr}\-maw}} \I{[a."k\s{r}. maw]} \emph{conj.} after (\emph{temporal}).
\B{Fo\-ri maw\-krr\-a fa ren\-ten ioi sä\-po\-li ho\-lum.} After they put on their goggles, they left.~\LINK{http://naviteri.org/2011/08/new-vocabulary-clothing/}{b}
\B{Maw\-krr\-a po ho\-lum, oe 'e\-fu ke\-ftxo.} I feel sad after he left.~\LINK{http://naviteri.org/2011/08/new-vocabulary-clothing/comment-page-1/\#comment-986}{b}
(\ding{43}~\HL{mawkrr}{\B{maw\-krr}}, \HL{mawfwa}{\B{maw\-fwa}})

% mawl
\hi
\HT{mawl}{\B{mawl}} \I{[mawl]} \emph{n., num.} half, one\hyp{}half, $\frac{1}{2}$.
\B{Po\-ri lu sno\-lup räp\-tum, ul\-te mawl ay\-lì\-'uä ke lu ngay.} His personal style is coarse and half of the words are not true.~\LINK{http://naviteri.org/2016/12/tengkrr-perahem-zisit-amip-as-the-new-year-arrives/}{b}
\B{Tsu'\-tey\-ìl to\-lìng oer mawl\-it smar\-ä.} Tsu'tey gave me a half of the prey.~\LINK{http://naviteri.org/2011/02/new-vocabulary-part-2/}{b}
(\emph{a. to tell time}) \B{trr\-pxì a\-mrr\-ve sì mawl} \emph{or coll.} \B{mrr\-ve sì mawl}, (it's) \on{5}:\on{30}.~\LINK{https://naviteri.org/2024/05/krrteri-about-time/}{b}

% mawey
\hi
\HT{mawey}{\B{ma\cp{wey}}} \I{[ma."wEj]} \emph{adj.} calm.
\B{Tok oel lo\-ra tsa\-mo\-ti a mì na'\-rìng a krr, 'e\-fu ma\-wey sì nit\-ram.} When I'm in that beautiful hollow in the forest, I feel calm and happy.~\LINK{http://naviteri.org/2012/10/mipa-vospxi-mipa-ayliu-new-words-for-the-new-month/}{b}
(\emph{a. idiom}) \B{Nga\-ri txe'\-lan ma\-wey li\-vu.} Don't worry (about it).~\LINK{http://forum.learnnavi.org/language-updates/fmawno/msg293595/\#msg293595}{ln} \\
\textsc{Derivatives}: \ding{43} 
\on{1}.~\HL{maweypey}{\B{ma\-wey\-pey}}, be patient. 

% maweypey
\hi
\HT{maweypey}{\B{ma\cp{wey}pey}} \I{[ma."wEj.p$\cdot$Ej]} \emph{vin., inf.}\on{33} be patient. 
\B{Lu nga sì la\-he\-a ay\-ha\-pxì\-tu lì'\-fya\-o\-lo'\-ä aw\-nge\-yä ay\-aw\-po a ma\-wey\-pe\-rey.} You and the other members of our \emph{Sprachbund} are the ones who are being patient.~\LINK{http://forum.learnnavi.org/language-updates/txelanit-hivawl/}{ln}
\B{Ze\-ne ma\-wey\-pi\-vey, ma ey\-lan.} You must be patient, friends.~\LINK{http://naviteri.org/2011/04/\%E2\%80\%99a\%E2\%80\%99awa-li\%E2\%80\%99fyavi-amip\%E2\%80\%94a-few-new-expressions/}{b} 
(\ding{43}~\HL{mawey}{\B{ma\-wey}}, \HL{pey}{\B{pey}}) \\  
\textsc{Derivatives}: \ding{43} 
\on{1}.~\HL{tImweypey}{\B{tìm\-wey\-pey}}, patience.   
\on{2}.~\HL{lemweypey}{\B{lem\-wey\-pey}}, patient.   
\on{3}.~\HL{nImweypey}{\B{nìm\-wey\-pey}}, patiently.  
\on{4}.~\HL{maweypeyyu}{\B{ma\-wey\-pey\-yu}}, patient person. 

% maweypeyyu 
\hi
\HT{maweypeyyu}{\B{ma\cp{wey}peyyu}} \I{[ma."wEj.pEj.ju]} \emph{n.} one who is patient, patient person. 
\B{Ay\-ma\-wey\-pey\-yu lä\-ngu nga sì la\-he\-a ay\-ha\-pxì\-tu lì'\-fya\-o\-lo'\-ä aw\-nge\-yä.} The one's who're being patient are you and the other \emph{Sprachbund} members.~\LINK{http://forum.learnnavi.org/language-updates/txelanit-hivawl/}{ln}
(\ding{43}~\HL{maweypey}{\B{ma\-wey\-pey}})

% mawup
\hi
\HT{mawup}{\B{\cp{ma}wup}} \I{["ma.wup\textcorner]} \emph{or} \I{[.wUp\textcorner]} \emph{n., fauna} turtapede. \\
\textsc{Derivatives}: \ding{43} 
\on{1}.~\HL{mawuptsyIp}{\B{ma\-wup\-tsyìp}}, turtle.

% mawuptsyìp
\hi 
\HT{mawuptsyIp}{\B{\cp{ma}wuptsyìp}} \I{["ma.wup\textcorner.\t{ts}jIp\textcorner]} \emph{or} \I{[.wUp\textcorner.]} \emph{n., fauna} turtle (Earth animal). 
(\ding{43}~\HL{mawup}{\B{ma\-wup}}) 

% may'
\hi
\HT{may'}{\B{may'}} \I{[majP]} \emph{vtr.} (\emph{a}.~\texttt{+}\emph{control}) taste, check out something by tasting. 
\B{Ke new oe mi\-vay' tsngan\-ti a 'o\-lem Ri\-nil.} I don't want to taste the meat that Rini cooked.~\LINK{http://naviteri.org/2012/11/renu-ayinanfyaya-the-senses-paradigm/}{b}
(\emph{b}.) try, sample, evaluate, check out, test\hyp{}drive.
\B{Ngal mo\-lay' pxaw\-pxun\-it Lo\-ak\-ä a krr pe\-kxin\-um?} When you tried on Loak’s armband, how tight was it?~\LINK{http://naviteri.org/2012/07/meetings-waterfalls-and-more/}{b}
(\emph{idiom}) \B{Mi\-vay'.\slash May' ko.} Give it a try.\slash Have a go.
(\emph{modal use}) \B{May' mi\-vak\-to pa'\-lit!} Try riding a direhorse!~\LINK{http://naviteri.org/2011/05/some-miscellaneous-vocabulary/}{b}

% mäkxu
\hi
\HT{mAkxu}{\B{mä\cp{kxu}}} \I{[m$\cdot$\ae."k'u]} \emph{vtr., inf.}\on{11} to interrupt, throw out of harmonious balance (\emph{only  for activities or established conditions})
\B{Tsun (oe) mi\-vä\-kxu hì\-krr srak?} May I interrupt a moment?~\LINK{http://naviteri.org/2010/09/getting-to-know-you-part-1/}{b}
\B{Pol mo\-lä\-kxu ul\-txa\-ti.} He interrupted the meeting.~\LINK{http://naviteri.org/2010/09/getting-to-know-you-part-1/}{b}

% mauti
\hi
\HT{mauti}{\B{\cp{ma}uti}} \I{["ma.u.ti]} \emph{n.}, \ding{96} fruit.
\B{Pe\-fne\-ma\-u\-ti su\-nu ngar fra\-to?} What's your favorite fruit?~\LINK{http://naviteri.org/2015/03/kap-si-ayunil-saylahe-threats-dreams-and-other-things/}{b} 
(\emph{a. idiom}) 
\B{Rey\-kì\-kxi ut\-ral\-ti, zup ma\-u\-ti.} If you shake the tree, the fruit will fall (i.e, ``actions have consequences'').~\LINK{http://naviteri.org/2018/03/100a-liu-amip-64-new-words-part-1/}{b} \\
\textsc{Derivatives}: \ding{43} 
\on{1}.~\HL{utumauti}{\B{u\-tu\-ma\-u\-ti}}, banana fruit. 
\on{2}.~\HL{paymaut}{\B{pay\-ma\-ut}}, fountain tree.
\on{3}.~\HL{tskxemauti}{\B{tskxe\-ma\-u\-ti}}, nut.

% me+
\hi
\HT{me+}{\B{me}\texttt{+}} \I{[mE.]} \emph{dual pref.} two (of s.t.).

% me'em
\hi
\HT{me'em}{\B{me\cp{'em}}} \I{[mE."PEm]} \emph{adj.} harmonious. \\
\textsc{Derivatives}: \ding{43} 
\on{1}.~\HL{tIme'em}{\B{tì\-me\-'em}}, (general) harmony.

% mewn
\hi 
\HT{mewn}{\B{mewn}} \I{[mEwn]} \emph{vin.} sigh. 
\B{Nga me\-rewn pe\-lun, ma pa\-ska\-lin?} Why are you sighing, sweetie?~\LINK{https://naviteri.org/2025/04/mevosinga-liu-amip-twenty-new-words/}{b} \\
\textsc{Derivatives}: \ding{43} 
\on{1}.~\HL{sAmewn}{\B{sä\-mewn}}, sigh, intance of sighing.

% mefan
\hi
\HT{mefan}{\B{me\cp{fan}}} \I{[mE."fan]} \emph{num., n.} two\hyp{}thirds, $\frac{2}{3}$.
(\ding{43}~\HL{pan}{\B{pan}})

% mefo
\hi
\HT{mefo}{\B{me\cp{fo}}} \I{[mE."fo]} (\emph{gen.} \B{me\-\cp{fe}\-yä}) \emph{pn.} they (two), the two of them (3\textsuperscript{rd} person dual).
\B{Me\-fo fì\-tsap mä\-po\-le\-yam teng\-krr tsngaw\-vìk.} The two of them hugged each other and wept.~\LINK{http://naviteri.org/2011/10/more-vocabulary-a-bit-of-grammar/}{b}
\B{'Aw\-po a me\-fo\-kip tsul\-fä\-tu le\-iu lì'\-fya\-yä le\-Na'\-vi.} One of the two is a master of the Na'vi language.~\LINK{http://naviteri.org/2014/02/barter-and-exchange/comment-page-1/\#comment-2619}{b}
(\emph{a. idiom}) \B{'Aw\-si\-teng lu me\-fo la\-nay'\-ka.} \emph{praising how well two people work together and complement each other} [lit.: they are a slinger (together)].~\LINK{http://naviteri.org/2021/08/ulte-ayyoratu-leiu-and-the-winners-are-2/}{b}
\B{To\-la\-ron me\-fol me\-sa\-li\-o\-a\-ngit! Tew\-ti, 'aw\-si\-teng lu me\-fo la\-nay'\-ka.} They hunted two sturmbeest? Wow, they work very well together.~\LINK{http://naviteri.org/2021/08/ulte-ayyoratu-leiu-and-the-winners-are-2/}{b}
(\ding{43}~\HL{po}{\B{po}})

% mei
\hi 
\HT{mei}{\B{\cp{me}i}} \I{["mE.i]} \emph{adj.} wet. 
\B{Kll\-te lu mei a krr, fwa fwi lu ftue.} When the ground is wet, it's easy to slip.~\LINK{http://naviteri.org/2020/04/aawa-u-amip-a-few-new-things/}{b}
(\ding{214}~\HL{ukxo}{\B{ukxo}}, \ding{43}~\HL{lepay}{\B{le\-pay}}, \HL{paynga'}{\B{pay\-nga'}}) \\
\textsc{Derivatives}: \ding{43} 
\on{1}.~\HL{meitayo}{\B{mei\-ta\-yo}}, wetlands.

% meitayo
\hi 
\HT{meitayo}{\B{mei\cp{ta}yo}} \I{[me.i."ta.jo]} \emph{or coll.} \I{[mej."ta.jo]} \emph{n.} wetlands. 
(\ding{43}~\HL{mei}{\B{mei}}, \HL{txayo}{\B{txa\-yo}})

% mek
\hi
\HT{mek}{\B{mek}} \I{[mEk\textcorner]} \emph{adj.}~(\emph{a}.) empty.   
\B{Nge\-yä tsngal lum\-pe lu mek?} Why is your cup empty?~\LINK{http://naviteri.org/2012/02/trr-asawnung-lefpom-happy-leap-day/}{b}
(\emph{b.~fig.}) `empty', having no valuable content. 
\B{mek\-a sä\-fpìl}, an empty\slash dumb idea~\LINK{http://naviteri.org/2012/02/trr-asawnung-lefpom-happy-leap-day/}{b}; 
\B{mek\-a sä\-pll\-txe\-vi}, an insipid\slash thoughtless comment~\LINK{http://naviteri.org/2012/02/trr-asawnung-lefpom-happy-leap-day/}{b}; 
\B{sä\-mok a\-mek}, a useless suggestion~\LINK{http://naviteri.org/2012/02/trr-asawnung-lefpom-happy-leap-day/}{b};
\B{\mbox{leioae} a\-mek}, feigned respect.~\LINK{http://naviteri.org/2012/02/trr-asawnung-lefpom-happy-leap-day/}{b} 
(\ding{214}~\HL{teya}{\B{te\-ya}}) \\
\textsc{Derivatives}: \ding{43} 
\on{1}.~\HL{mektseng}{\B{mek\-tseng}}, gap, breach.  
\on{2}.~\HL{momek}{\B{mo\-mek}}, hollow (not solid).

% mekong
\hi
\B{\cp{me}kong}~:~\HL{'ekong}{\B{'ekong}}

% mekre
\hi
\HT{mekre}{\B{\cp{mek}re}} \I{["mEk\textcorner.RE]} \emph{n.} supplies. 

% mektseng
\hi
\HT{mektseng}{\B{\cp{mek}tseng}} \I{["mEk\textcorner.\t{ts}EN]} \emph{n.} gap, breach.
\B{Lu mek\-tseng a tsun fpxi\-vä\-kìm hì\-'ang tsaw\-ìlä.}  There's a gap through which insects can get in.~\LINK{http://naviteri.org/2013/04/wheres-the-bathroom-and-other-useful-things/}{b}
\B{Fì\-po\-stì mek\-tseng\-ur te\-ya si.} This post fills a gap (in our understanding of Na'vi syntax).~\LINK{http://naviteri.org/2013/04/zun-zel-counterfactual-conditionals/}{b}
(\ding{43}~\HL{mek}{\B{mek}}, \HL{tseng}{\B{tseng}})

% mele'wll
\hi
\B{mele'wll}~:~\HL{'ele'wll}{\B{'ele'wll}}

% melo
\hi
\HT{melo}{\B{\cp{me}lo}} \I{["mE.lo]} \emph{adv.} twice.
\B{Tsa\-lì'\-fya\-vi\-ri fko ze\-ne si\-var lì\-'ut alu fì\-'u me\-lo.} One must use the word \emph{fì'u} twice in that expression.~\LINK{http://naviteri.org/2011/04/\%E2\%80\%99a\%E2\%80\%99awa-li\%E2\%80\%99fyavi-amip\%E2\%80\%94a-few-new-expressions/comment-page-1/\#comment-677}{b}
(\ding{43}~\HL{alo}{\B{alo}})

% memyu
\hi
\B{\cp{mem}yu}~:~\HL{'emyu}{\B{'emyu}}

% men
\hi
\B{men}~:~\HL{'en}{\B{'en}}

% menga
\hi
\HT{menga}{\B{me\cp{nga}}} \I{[mE."Na]} \emph{pn.} (\emph{gen.} \B{me\-\cp{nge}\-yä}) you (two), the two of you. (2\textsuperscript{nd} person dual). 
\B{Me\-nga 'e\-fu väng srak?} Are you two thirsty?~\LINK{http://naviteri.org/2011/01/new-year-new-vocabulary/}{b}
\B{Me\-nga lu kar\-yu.} You two are teachers.~\LINK{http://naviteri.org/2011/07/number-in-na\%E2\%80\%99vi/}{b}
\B{Ay\-oel me\-nga\-ti ka\-me\-i\-e.} We \textsc{see} you both.~\LINK{http://naviteri.org/2014/02/barter-and-exchange/}{b}
\B{Me\-nga\-ru tì\-yawr.} You are both right.~\LINK{http://naviteri.org/2010/09/getting-to-know-you-part-2/comment-page-1/\#comment-347}{b}
\B{Lor\-a fì\-tì\-kang\-kem me\-nge\-yä meuia lu\-yu aw\-nga\-ru nì\-wotx.} This beautiful work of theirs honors us very much.~\LINK{http://naviteri.org/2010/12/beautiful-christmas-carols-sung-in-navi/}{b}
(\ding{43}~\HL{nga}{\B{nga}})

% mengenga
\hi
\HT{mengenga}{\B{menge\cp{nga}}} \I{[mE.NE."Na]} \emph{pn.} (\emph{gen.} \B{me\-nge\-\cp{nge}\-yä}) you (two), the two of you. (2\textsuperscript{nd} person dual; deferential\slash ceremonial form).
(\ding{43}~\HL{ngenga}{\B{nge\-nga}})

% meoauniaea
\hi
\HT{meoauniaea}{\B{meoaunia\cp{e}a}} \I{[""mE.o.a.u.ni.a."{}E.a]} \emph{n.} harmony, living at one with nature.
\B{New oe ti\-vìng ay\-ngar lì\-'ut a tì\-'e\-fu\-mì oe\-yä lu lor fra\-to mì lì'\-fya le\-Na'\-vi: me\-o\-a\-u\-ni\-a\-e\-a.} I want to present to you the word that, in my opinion, is the most beautiful in the Na’vi language: \emph{me\-o\-a\-u\-ni\-a\-e\-a}.~\LINK{http://forum.learnnavi.org/language-updates/trr-rrtaya/}{ln}

% merki
\hi
\HT{merki}{\B{\cp{mer}ki}} \I{["mER.ki]} \emph{n.} ground rack (\emph{for smoking meats}).

% mesrram
\hi
\HT{mesrram}{\B{mesrr\cp{am}}} \I{[mE.s\s{r}."am]} \emph{n.; adv.} the day before yesterday, two days ago. 
\B{Ìlä Fey\-ral, mun\-txa so\-li Ra\-lu sì Ne\-wey nì\-wan me\-srr\-am.} According to Peyral, Ralu and Newey were secretly married the day before yesterday.~\LINK{http://naviteri.org/2012/07/meetings-waterfalls-and-more/}{b}
(\ding{43}~\HL{trram}{\B{trr\-am}})

% mesrray
\hi
\HT{mesrray}{\B{mesrr\cp{ay}}} \I{[mE.s\s{r}."aj]} \emph{n.; adv.} the day after tomorrow, two days from now.
\B{Oer lu txan\-a yew\-la a ke tsun nga oe\-hu ki\-vä\-teng me\-srr\-ay.} I'm very disappointed you can't hang out with me the day after tomorrow.~\LINK{http://naviteri.org/2012/10/mipa-vospxi-mipa-ayliu-new-words-for-the-new-month/}{b}
(\ding{43}~\HL{trray}{\B{trr\-ay}})

% metkayina
\hi
\B{Metka\cp{yi}na} \I{[mEt\textcorner.ka."ji.na]} \emph{n. clan name}, one clan of the reef people.

% metnaw
\hi
\B{\cp{met}naw}~:~\HL{'etnaw}{\B{'etnaw}}

% meuia / si
\hi
\HT{meuia}{\B{me\cp{u}ia}} \I{[mE."{}u.i.a]} \textsc{i}. \emph{n.} honor.
\B{Fì\-me\-u\-i\-a\-ri sei\-yi moe ira\-yo nì\-txan, ma smuk.} We both thank you very much for this honor, siblings.~\LINK{http://naviteri.org/2014/08/stxeli-alor-a-beautiful-gift/}{b}
(\emph{a.}) \B{Fay\-sul\-fä\-tu\-ä tì\-kang\-kem o\-he\-ru me\-u\-i\-a lu\-yu nì\-ngay.} The work of these experts truly honors me.~\LINK{http://naviteri.org/2010/06/first-post/}{b}
(\emph{b.~idiom}) \B{Oeru me\-u\-i\-a.} ``You're welcome.\slash The honor is mine.'' (in response to thanks) \\  
\textsc{ii}. \B{me\cp{u}ia si} \I{[me."{}u.i.a si]} \emph{vin.} honor. 
\B{Nge\-yä fay\-lì\-'u a\-tì\-tstun\-wi\-nga' oe\-ru meuia so\-li nì\-ngay.} These kind words of yours have honored me greatly.~\LINK{http://naviteri.org/2018/04/100a-liu-amip-64-new-words-part-2/}{b} 
(\ding{214}~\HL{kemuia}{\B{ke\-mu\-ia}}) \\
\textsc{Derivatives}: \ding{43} 
\on{1}.~\HL{meuianga'}{\B{me\-uia\-nga'}}, honorable. 
\on{2}.~\HL{kemuia}{\B{ke\-mu\-ia}}, dishonor.

% meuianga'
\hi
\HT{meuianga'}{\B{me\cp{u}ianga'}} \I{[mE."{}u.i.a.NaP]} \emph{adj.} honorable. 
(\ding{43}~\HL{meuia}{\B{me\-uia}})

% mevan
\hi
\B{\cp{me}van}~:~\HL{'evan}{\B{'evan}}

% meve
\hi
\B{\cp{me}ve}~:~\HL{'eve}{\B{'eve}}

% meveng
\hi
\B{\cp{me}veng}~:~\HL{'eveng}{\B{'eveng}}

% mevi
\hi
\B{\cp{me}vi}~:~\HL{'evi}{\B{'evi}}

% mevol
\hi
\HT{mevol}{\B{\cp{me}vol}} \I{["mE.vol]} \emph{num., adj.} \on{20} (octal), \on{16} (decimal).
\B{Po\-el haw\-re'\-tsyìp\-it ko\-la\-nom yoa txo\-lar a\-me\-vol.} She bought a little cap for \$\on{16}.~\LINK{http://naviteri.org/2014/02/barter-and-exchange/}{b}

% meyam
\hi
\HT{meyam}{\B{me\cp{yam}}} \I{[mE."jam]} \emph{vtr.} hug, embrace, hold in one's arms.   
\B{Ma sa'\-nu, oe txo\-pu si. Me\-yam oe\-ti!} Mommy, I'm scared. Hold me!~\LINK{http://naviteri.org/2011/10/more-vocabulary-a-bit-of-grammar/}{b} 
\B{Me\-fo fì\-tsap mä\-po\-le\-yam teng\-krr tsngaw\-vìk.} The two of them hugged each other and wept.~\LINK{http://naviteri.org/2011/10/more-vocabulary-a-bit-of-grammar/}{b} \\
\textsc{Derivatives}: \ding{43} 
\on{1}.~\HL{sAmyam}{\B{säm\-yam}}, hug, embrace.

% meylan
\hi
\B{\cp{mey}lan}~:~\HL{'eylan}{\B{'eylan}}

% meylanay
\hi
\B{meyla\cp{nay}}~:~\HL{'eylanay}{\B{'eylanay}}

% meylyong
\hi
\B{\cp{meyl}yong}~:~\HL{'eylyong}{\B{'eylyong}}

% meyp
\hi
\HT{meyp}{\B{meyp}} \I{[mEjp\textcorner]} \emph{adj.} (\emph{a.~of strength}) weak. 
(\emph{b.~fig.}) \B{Tsa\-sä\-pll\-txe\-ri sä\-wì\-ngay a to\-lìng ngal lu meyp.} The proof you gave of that statement is weak.~\LINK{http://naviteri.org/2015/04/some-new-words-for-may-day/}{b} \\  
\textsc{Derivatives}: \ding{43} 
\on{1}.~\HL{tImeyp}{\B{tì\-meyp}}, weakness. 
\on{2}.~\HL{nImeyp}{\B{nì\-meyp}}, weakly, loosely. 
\on{3}.~\HL{tompameyp}{\B{tom\-pa\-meyp}}, drizzle. 
\on{4}.~\HL{hermeyp}{\B{her\-meyp}}, snow flurry. 
\on{5}.~\HL{meyptu}{\B{meyp\-tu}}, weakling.

% meyptu
\hi 
\HT{meyptu}{\B{\cp{meyp}tu}} \I{["mEjp\textcorner.tu]} \emph{n.} weakling, \emph{a person of weak physique or character}. 
(\ding{43}~\HL{meyp}{\B{meyp}}, \HL{-tu}{\B{-tu}}, \ding{214}~\HL{txurtu}{\B{txur\-tu}})

% mi
\hi
\HT{mi}{\B{mi}} \I{[mi]} \emph{adv.} yet, still, as before.
\B{Pol trr\-am ye\-rik\-it tar\-ma\-ron.--Fì\-trr mi te\-ra\-ron.} He was hunting hexapede yesterday.--He's still hunting today.~\LINK{http://forum.learnnavi.org/language-updates/as-x-as-y-keep-on-keepin-on/}{ln}
\B{Tsa\-tsko lu spu\-win ul\-te ke lu mi txur, slä oe\-ri txan\-wa\-we le\-iu.} That bow is old and no longer strong, but it means a lot to me.~\LINK{http://naviteri.org/2011/05/some-miscellaneous-vocabulary/}{b} 

% miklor
\hi
\HT{miklor}{\B{mik\cp{lor}}} \I{[mik\textcorner."loR]}  \emph{adj.} pleasant sounding, beautiful sounding. 
\B{Lu por mok\-ri a\-mik\-lor, slä lo\-re\-yu 'aw\-nam\-pi lu.} She has a beautiful voice, but she's extremely shy.~\LINK{http://naviteri.org/2014/11/twenty-before-the-holidays/}{b}
(\ding{214}~\HL{mikvA'}{\B{mik\-vä'}}; \ding{43}~\HL{mikyun}{\B{mik\-yun}}, \HL{lor}{\B{lor}})

% miktsang
\hi
\HT{miktsang}{\B{\cp{mik}tsang}} \I{["mik\textcorner.\t{ts}aN]} \emph{n.} earring.
\B{Pey\-ral\-ä mik\-tsang\-it a\-mip ngal tso\-le\-ri srak?} Did you notice Pey\-ral's new earring?~\LINK{http://naviteri.org/2015/03/kap-si-ayunil-saylahe-threats-dreams-and-other-things/}{b}
\B{Se\-vin\-a tsa\-'e\-ve\-ru a\-hì\-'i me\-sa'\-sem ioi so\-li fa mik\-tsang.} The parents put earrings on that pretty little girl.~\LINK{http://naviteri.org/2011/08/new-vocabulary-clothing/}{b}
(\ding{43}~\HL{mikyun}{\B{mik\-yun}}, \HL{tsang}{\B{tsang}})

% mikvä'
\hi
\HT{mikvA'}{\B{mik\cp{vä'}}} \I{[mik\textcorner."v\ae{}P]}  \emph{adj.} bad\hyp{}sounding.
(\ding{214}~\HL{miklor}{\B{mik\-lor}}; \ding{43}~\HL{mikyun}{\B{mik\-yun}}, \HL{vA'}{\B{vä'}})

% mikyun
\hi
\HT{mikyun}{\B{\cp{mik}yun}} \I{["mik\textcorner.jun]} \emph{or} \I{[.jUn]} \textsc{i}.~\emph{n.} ear. 
\B{Lu Ney\-ti\-ri\-ru 'aw\-a mik\-yun a\-raw\-nipx nì\-'aw.} Neytiri has only one pierced ear.~\LINK{http://naviteri.org/2017/12/zisit-amip-lefpom-ma-eylan-happy-new-year-friends/}{b}
\B{'om na mik\-yun}, the purplish color on the inside of a Na'vi ear.~\LINK{http://naviteri.org/2010/08/mipa-ayopin-mipa-ayliu-new-colors-new-words/}{b}
\B{Tskxe\-keng\-tsyìp a Mik\-yun\-fpi}, A Little Listening Exercise.~\LINK{http://naviteri.org/2015/09/tskxekengtsyip-a-mikyunfpi-a-little-listening-exercise-2/}{b}
(\emph{idiom}) \B{Kxe\-tse sì mik\-yun kop pll\-txe.} Body language is important [lit.: the tail and ear also speak].~\LINK{http://forum.learnnavi.org/language-updates/grid-tidbits/}{ln} \\
\textsc{ii}. \HT{tIng mikyun}{\B{tìng \cp{mik}yun}} \I{[tIN "mik\textcorner.jun]} \emph{or coll.} \I{[tI."mik\textcorner.jun]} \emph{vin.} (\texttt{+}\emph{control}) listen, listen to. 
\B{Tìng mik\-yun fte tsli\-vam.} Listen in order to understand.~\LINK{http://naviteri.org/2010/07/ting-mikyun-fte-tslivam-listening-comprehension-1-3/}{b}
\B{Oe\-ru tìng mik\-yun, ma Ra\-lu.} Listen to me, Ralu.~\LINK{http://naviteri.org/2013/01/awvea-posti-zisita-amip-first-post-of-the-new-year/}{b}
\B{Po pa\-mll\-txe a krr, fra\-po tar\-mìng mik\-yun nì\-pxi, ta\-lu\-na mok\-ri lu sä\-tsì\-syì\-tsyìp.} When he spoke, everyone listened intently, because his voice was a tiny whisper.~\LINK{http://naviteri.org/2011/09/miscellaneous-vocabulary/}{b} \\
\textsc{Derivatives}: \ding{43} 
\on{1}.~\HL{miktsang}{\B{mik\-tsang}}, earring. 
\on{2}.~\HL{miklor}{\B{mik\-lor}}, pleasant\slash beautiful sounding. 
\on{3}.~\HL{mikvA'}{\B{mik\-vä'}}, bad\hyp{}sounding.

% mimu
\hi 
\HT{mimu}{\B{\cp{mi}mu}} \I{["mi.mu]} \emph{vtr.} develop. 
\B{Ze\-ne aw\-nga mi\-vi\-mu tì\-hawl\-it a lä\-txayn kxu\-tut.} We have to develop a plan to defeat the enemy.~\LINK{https://naviteri.org/2026/04/tsivola-liu-amip-thirty-two-new-words/}{b} \\
\textsc{Derivations}: \ding{43} 
\on{1}.~\HL{tImimu}{\B{tì\-mi\-mu}}, development.

% mip
\hi
\HT{mip}{\B{mip}} \I{[mip\textcorner]} \emph{adj.} new.
\B{Ke lä\-ngu fì\-po\-stì\-mì ke\-'u a lu mip.} Unfortunately, there is nothing in this post that is new.~\LINK{http://naviteri.org/2013/08/taronway-the-hunt-song/}{b}
\B{Mip\-a ay\-lì'\-fya\-vi, mip\-a ho\-ren}, New expressions, new rules.~\LINK{http://naviteri.org/2010/06/first-post/}{b}
\B{Zi\-va\-'u nì\-mun ye'\-rìn \ldots tsa\-krr ra\-yun ay\-ngal ay\-u\-pxa\-ret a\-mip.} Come again soon \ldots then you'll find new messages.~\LINK{http://naviteri.org/2010/06/first-post/}{b}
(\emph{a. idiom}) \B{Mip\-a zì\-sìt le\-fpom nga\-ru!} Happy New Year to you!~\LINK{http://naviteri.org/2011/01/new-year-new-vocabulary/}{b}
(\ding{214}~\HL{lal}{\B{lal}})

% mì
\hi
\HT{mI}{\B{mì}} \I{[mI]} \emph{adp.}\texttt{+} (\emph{a.~local})  in, on. 
\B{Mì sray a\-pxa kel\-ku si soaia\-hu.} (He) lives in a great village with his family.~\LINK{http://naviteri.org/2010/07/ting-mikyun-fte-tslivam-listening-comprehension-1-3/}{b}
\B{Nge\-yä tì\-rol\-tsyìp\-ìri, oe\-yä ay\-syu\-lang set mì fay kll\-kxe\-rem.} As for your little song, my flowers are standing in water now.~\LINK{http://naviteri.org/2010/09/getting-to-know-you-part-1/comment-page-1/\#comment-296}{b}
\B{Ul\-txa a mì na'\-rìng}, Meeting in the forest.~\LINK{http://naviteri.org/2010/12/catching-up/}{b}
\B{Pam\-tse\-ol ngop ay\-re\-nut | Mì ron\-sem\-ä tì\-fnu.} The music creates patterns | In the silence of the mind.~\LINK{http://www.pandorapedia.com/navi/music/ritual_music}{pp}
\B{Ay\-lì\-'u na ay\-skxe mì te'\-lan.} The words (are) like stones in (my) heart.~\LINK{http://wiki.learnnavi.org/index.php?title=Corpus\#Behind_the_Scenes}{wiki}
\B{Krro lu 'Ìng\-lì\-sì txur nì\-hawng mì re\-'o.} Sometimes English is too strong in my head.~\LINK{http://naviteri.org/2011/09/\%E2\%80\%9Cby-the-way-what-are-you-reading\%E2\%80\%9D/comment-page-1/\#comment-1092}{b}
\B{Wo\-leyn Ìstawl yey\-fyat mì hll\-te fte oeyk\-ti\-vìng fra\-po\-ru tì\-hawl\-te\-ri sne\-yä.} Ìstaw drew a line on the ground to explain his plan to everyone.~\LINK{http://naviteri.org/2013/09/on-si-salewfya-shapes-and-directions/}{b}
\B{Oe\-ri skxir a mì syokx tì\-sraw se\-ngi.} The wound on my hand hurts.~\LINK{http://naviteri.org/2011/05/some-miscellaneous-vocabulary/}{b}
\B{Txur\-tel\-mì fo fta so\-li fte\-ke ka tsyokx fwi\-vi.} They tied a knot in the rope so it wouldn't slip through their hands.~\LINK{http://naviteri.org/2013/04/wheres-the-bathroom-and-other-useful-things/}{b}
(\emph{b. local, fig.}) in. 
\B{Law lu oe\-ru fwa nga mì rel\-tseo no\-lu\-me nì\-txan!} I'm sure that you've learned a lot in music!~\LINK{http://forum.learnnavi.org/language-updates/art-related-vocab/}{ln}
\B{Ta\-fral ho\-lawl ay\-oel ay\-nga\-fpi 'a\-'aw\-a tì\-päng\-kxo\-tsyìp\-it a tsun fko si\-var mì sì\-rey le\-trr\-trr.} That's why we prepared several little dialogues for you that one can use in ordinary life.~\LINK{http://naviteri.org/2012/07/tskxekengtsyip-a-mikyunfpi-a-little-listening-exercise/}{b}
\B{Sìl\-pey oe, ay\-nga\-ri fra\-'u a mì sì\-rey vi\-var tsan\-'i\-vul fra\-fya, ma ey\-lan.} I hope, everything in your life continues to improve for you in every way, friends.~\LINK{http://naviteri.org/2013/03/tsanerul-feerul-getting-better-getting-worse/}{b}
\B{New oe ti\-vìng ay\-ngar lì\-'ut a tì\-'e\-fu\-mì oe\-yä lu lor fra\-to mì lì'\-fya le\-Na'\-vi: me\-o\-a\-u\-ni\-a\-e\-a.} I want to present to you the word that, in my opinion, is the most beautiful in the Na’vi language: \emph{me\-o\-a\-u\-ni\-a\-e\-a}.~\LINK{http://forum.learnnavi.org/language-updates/trr-rrtaya/}{ln}
\B{Sra\-ke tsun nga ri\-vun fì\-tì\-fkey\-tok\-mì a tì\-'i\-put?} Can you find the humor in this situation?~\LINK{http://naviteri.org/2012/01/mipa-zisit-ayliu-amip-new-words-for-the-new-year/}{b}
(\emph{c.~local, planetary surface}) on. 
\B{mì 'Rr\-ta\slash 'Rr\-ta\-mì}, on Earth.
\B{mì Ey\-wa\-'e\-veng\slash Ey\-wa\-'e\-veng\-mì}, on Pandora.
(\emph{d.~temporal, for range of time}) 
\B{Ke lo\-lu oer nì\-ke\-ftxo mì sok\-a srr ay\-skxom le\-tam fte lì'\-fya\-ri aw\-nge\-yä tì\-kang\-kem si\-vi.} In recent days I have not had sufficient opportunity to work on our language.~\LINK{http://forum.learnnavi.org/language-updates/trr-rrtaya/}{ln} 
\B{Fì\-trr\-mì le\-tsran\-ten}, on this important day.~\LINK{http://forum.learnnavi.org/language-updates/trr-rrtaya/}{ln}
\B{Alo a\-hay oe za\-sya\-'u mì re\-won.} Next time, I'll come in the morning.~\LINK{http://forum.learnnavi.org/language-updates/txelanit-hivawl/}{ln} 
\B{mì ftaw\-nem\-krr\slash zu\-saw\-krr}, in the past\slash future.
(\emph{e.~idiom})
\B{tì\-'e\-fu\-mì oe\-yä}, in my opinion.
\B{tì\-o\-mum\-mì oe\-yä} \I{[tI."{}o.mu.mI "wE.j\ae]}, to my knowledge, as far as I know.~\LINK{http://naviteri.org/2011/09/\%E2\%80\%9Cby-the-way-what-are-you-reading\%E2\%80\%9D/}{b}
\B{Ay\-te\-le a nge\-yä ha\-pxì\-mì ki\-fkey\-ä lu fya\-pe?} What's new in your part of the world?~\LINK{http://naviteri.org/2011/08/new-vocabulary-clothing/comment-page-1/\#comment-917}{b}
(\emph{f.~proverb}) 
\B{Txìm a\-'aw ke tsun hi\-veyn mì tal me\-fa'\-li\-yä.} You can't sit on the fence; you need to decide [lit.: One butt can't sit on the backs of two direhorses].~\LINK{http://naviteri.org/2013/02/vospxi-ayol-posti-apup-short-post-for-a-short-month/}{b} \\  
\textsc{Derivatives}: \ding{43} 
\on{1}.~\HL{mIfa}{\B{mìfa}}, inside.   
\on{2}.~\HL{mIftxele}{\B{mì\-ftxe\-le}}, by the way.   
\on{3}.~\HL{mIkam}{\B{mì\-kam}}, between. 
\on{4}.~\HL{mIso}{\B{mì\-so}}, away.

% mìfa
\hi
\HT{mIfa}{\B{\cp{mì}fa}} \I{["mI.fa]} \emph{or} \B{mì\cp{fa}} \I{[mI."fa]} \emph{n., adv.} inside.
\B{Fra\-po ne mì\-fa! Le\-rok yrr\-ap a\-pxa!} Everybody inside! A big storm is approaching!~\LINK{http://naviteri.org/2012/07/meetings-waterfalls-and-more/}{b}
\B{Ke za\-syup lì\-'O\-na ne kxu\-tu a mì\-fa fu a wrr\-pa.} The l'Ona will not perish to the enemy within or the enemy without.~\LINK{http://forum.learnnavi.org/language-updates/ltasygt-with-ke-and-nga/}{g\slash ln} 
(\ding{43}~\HL{mI}{\B{mì}}, \HL{pa'o}{\B{pa'o}}) \\  
\textsc{Derivatives}: \ding{43} 
\on{1}.~\HL{nemfa}{\B{nem\-fa}}, into. 
\on{2}.~\HL{ftumfa}{\B{ftum\-fa}}, out of. 
\on{3}.~\HL{tamIfa}{\B{ta\-mì\-fa}}, internal. 

% mìftxele
\hi
\HT{mIftxele}{\B{mì\cp{ftxe}le}} \I{[mI."ft'E.lE]} \emph{adv.} by the way, in this regard, related to this matter.
\B{Mì\-ftxe\-le po\-ri lu oer le\-tsran\-ten\-a fmawn a new pi\-veng ngar.} By the way, I have some important news I want to tell you about him.~\LINK{http://naviteri.org/2011/09/\%E2\%80\%9Cby-the-way-what-are-you-reading\%E2\%80\%9D/}{b}
\B{Sko Sa\-hìk ke tsun oe mì\-ftxe\-le tsngi\-vaw\-vìk.} As Tsahik, I cannot weep over this matter.~\LINK{http://naviteri.org/2012/03/spring-vocabulary-part-2/}{b}
(\ding{43}~\HL{nIvingkap}{\B{nì\-ving\-kap}})

% mìkam
\hi
\HT{mIkam}{\B{mì\cp{kam}}} \I{[mI."kam]} \emph{adp.} between. 
\B{Ta\-lu\-na fì\-tì\-'ey\-lan lu pum a mì\-kam tu\-tan sì syak\-syuk.} Because this friendship is between a man and a Prolemuris.~\LINK{http://naviteri.org/2020/04/awa-tskxekengtsyip-a-mikyunfpi-niul-one-more-little-listening-exercise/}{b}
(\ding{43}~\HL{mI}{\B{mì}}, \HL{kxam}{\B{kxam}}) 

% mìn
\hi
\HT{mIn}{\B{mìn}} \I{[mIn]} \emph{vin.} turn.
\B{Mìn ne ftär\-pa.} Turn to the left.~\LINK{http://forum.learnnavi.org/language-updates/ftarpa-and-skienpa-added-to-dict/msg570362/\#msg570362}{ln}
\B{Kun\-sìp\-ìri txan\-a tì\-meyp lu tsyal a mìn.} The gunship's main weakness is the rotor system.~\LINK{http://forum.learnnavi.org/language-updates/ftarpa-and-skienpa-added-to-dict/msg570362/\#msg570362}{ln}  \\  
\textsc{Derivatives}: \ding{43} 
\on{1}.~\HL{mInyu}{\B{mìn\-yu}}, turner; twisted lily.

% mìnyu
\hi
\HT{mInyu}{\B{\cp{mìn}yu}} \I{["mIn.ju]} \emph{n.} turner; \ding{96} twisted lily, \emph{pseudolilium contortum}.
(\ding{43}~\HL{mIn}{\B{mìn}})

% mìso
\hi
\HT{mIso}{\B{mì\cp{so}}} \I{[mI."{}so]} \emph{adv.} away (\emph{positional}). 
% \B{'Ì\-'awn mì\-so.} Stay away.
(\ding{43}~\HL{neto}{\B{ne\-to}})

% mll'an
\hi
\HT{mll'an}{\B{mll\cp{'an}}} \I{[m\s{l}."Pan]} (\textsc{perf.}~\B{mol'an}) \emph{vtr.} accept.   
\B{Mun\-txa\-tu\-ri So\-rewn\-ti ke tsun oe mi\-vll\-'an.} I can't accept Sorewn as my spouse.~\LINK{http://naviteri.org/2012/10/mipa-vospxi-mipa-ayliu-new-words-for-the-new-month/}{b}
\B{Stxe\-nu\-to\-lìng oel fu\-ta lì'\-fyat le\-Na'\-vi poe\-ru ki\-var. Mol\-'an nì\-prr\-te'.} I offered to teach her Na'vi. She accepted gladly.~\LINK{http://naviteri.org/2011/09/miscellaneous-vocabulary/}{b}
\B{Fu\-ta nga\-ta tel pxen\-yoa srät, mll\-'ei\-an oel.} I'm happy to take cloth from you for finished garments.~\LINK{http://naviteri.org/2014/02/barter-and-exchange/}{b} 
(\ding{214}~\HL{tsyAr}{\B{tsyär}}) \\  
\textsc{Derivatives}: \ding{43} 
\on{1}.~\HL{tImll'an}{\B{tì\-mll\-'an}}, acceptance.

% mllte
\hi
\HT{mllte}{\B{mll\cp{te}}} \I{[m\s{l}."tE]} (\textsc{perf.}~\B{mol\-te}) \emph{vin.} agree (\HL{hu}{\B{hu}}, with; \emph{with the thing agreed upon in the topic}). 
\B{Nga\-hu mll\-te oe.} I agree with you.~\LINK{http://forum.learnnavi.org/language-updates/txelanit-hivawl/}{ln}
\B{Txam\-pay lìm nì\-hawng a tsa\-ri mll\-te oe hu sem\-pul.} I agree with father that the ocean is too far away.~\LINK{http://forum.learnnavi.org/language-updates/txelanit-hivawl/}{ln}

% mo
\hi
\HT{mo}{\B{mo}} \I{[mo]} \emph{n.} space, hollow, enclosed open area.
\B{Tok oel lora tsamoti a mì na'rìng a krr, 'efu mawey sì nitram.} When I'm in that beautiful hollow in the forest, I feel calm and happy.~\LINK{http://naviteri.org/2012/10/mipa-vospxi-mipa-ayliu-new-words-for-the-new-month/}{b} \\
(i.) \HT{mo a hahaw}{\B{mo a \cp{ha}haw}}, bedroom, sleeping area. \\
(ii.) \HT{mo a yom}{\B{mo a yom}}, dining room, dining area. \\
(iii.) \HT{mo letrrtrr}{\B{mo le\cp{trr}trr}}, living room. \\  
(iv.) \HT{mo a fngA'}{\B{mo a fngä'}}, restroom (\emph{an enclosed private structure or room}). \\
\textsc{Derivatives}: \ding{43} 
\on{1}.~\HL{snomo}{\B{sno\-mo}}, private space that one can retreat to.  
\on{2}.~\HL{momek}{\B{mo\-mek}}, hollow, not solid.

% mo'ara
\hi
\HT{mo'ara}{\B{Mo\cp{'a}ra}} \I{[mo."Pa.Ra]} (\emph{a}.) \emph{n.} gathering place (\emph{generally for some specific purpose}).
(\emph{b}.) \emph{name of the Disney theme park}.

% mo'at
\hi
\HT{mo'at}{\B{\cp{Mo}'at}} \I{["mo.Pat\textcorner]} \emph{n., female name}.
\B{Mo\-'at \mbox{alu} Tsa\-hìk lu O\-ma\-ti\-ka\-ya\-ä le\-'awa ha\-pxì\-tu a \mbox{ioi} sä\-pi fa 'a\-re.} Mo'at, the Tsahik, is the only member of the Omatikaya who wears a poncho.~\LINK{http://naviteri.org/2011/08/new-vocabulary-clothing/}{b}
\B{Spaw oel fu\-ta Mo\-'at\-ìl tso\-le\-'a Ney\-ti\-rit.} I believe, Mo'at saw Neytiri.~\LINK{http://naviteri.org/2011/03/word-order-and-case-marking-with-modals/}{b}

% moe
\hi
\HT{moe}{\B{\cp{mo}e}} \I{["mo.E]} \emph{pn.} we (two, not you), the two of us (1\textsuperscript{st} person dual exclusive).
\B{Moe sìl\-pey nì\-teng tsnì pxo\-eng tsä\-pì\-ye\-ve\-'a fì\-tsap ye'\-rìn nì\-mun.} We also hope that we will see each other again soon.~\LINK{http://naviteri.org/2013/07/pximaw-tsawlultxa-just-after-the-meet-up/comment-page-1/\#comment-2483}{b}
\B{Moe smon (moe\-ru) fì\-tsap.} We know each other.~\LINK{http://naviteri.org/2011/12/one-more-for-2011/}{b}
\B{Teng\-krr pa\-lu\-lu\-kan moe\-ne kxll sar\-mi, pol\-txe Ney\-ti\-ril ay\-lì\-'ut a fra\-krr 'ok se\-yä la\-yu oer.} As the thanator was charging towards the two of us, Neytiri said something I will always remember.~\LINK{http://thereadinginpublicproject.blogspot.de/2010/05/paul-frommer.html}{pr}
(\ding{43}~\HL{oe}{\B{oe}}, \HL{me+}{\B{me}\texttt{+}}) \\
\textsc{Derivatives}: \ding{43}
\on{1}.~\HL{mohe}{mo\-he}, we two (deferential).

% mohe
\hi
\HT{mohe}{\B{\cp{mo}he}} \I{["mo.hE]} \emph{pn.} we (two, not you), the two of us (1\textsuperscript{st} person dual exclusive; deferential\slash ceremonial form). 
(\ding{43}~\HL{ohe}{\B{ohe}}, \HL{me+}{\B{me}\texttt{+}})

% mok
\hi
\HT{mok}{\B{mok}} \I{[mok\textcorner]} \emph{vtr.} suggest.
\B{Nge\-yä tì\-pawm\-ìri mok oel fu\-ta ngal si\-var kem\-lì\-'u\-vit alu ‹irv›.} As for your question, I suggest that you use the infix ‹irv›.~\LINK{http://naviteri.org/2013/01/awvea-posti-zisita-amip-first-post-of-the-new-year/comment-page-1/\#comment-1951}{b}
\B{Tsa\-tsmu\-kel alu Ri\-ni mo\-lok fu\-ta oe ngar mu\-wä\-pi\-vìn\-txu.} Sister Rini over there suggested that I introduce myself.~\LINK{http://naviteri.org/2010/09/getting-to-know-you-part-1/}{b} \\
\textsc{Derivatives}: \ding{43} 
\on{1}.~\HL{tImok}{\B{tì\-mok}}, suggestion, suggesting.   
\on{2}.~\HL{sAmok}{\B{sä\-mok}}, suggestion. 

% mokri
\hi
\HT{mokri}{\B{\cp{mok}ri}} \I{["mok\textcorner.Ri]} \emph{n.} voice. 
\B{Po pa\-mll\-txe a krr, fra\-po tar\-mìng mik\-yun nì\-pxi, ta\-lu\-na mok\-ri lu sä\-tsì\-syì\-tsyìp.} When he spoke, everyone listened intently, because his voice was a tiny whisper.~\LINK{http://naviteri.org/2011/09/miscellaneous-vocabulary/}{b}
\B{Fu\-la nga\-ri oel mok\-rit sta\-yawm, oe\-ti nit\-ram sley\-ku nì\-txan.} It makes me very happy that I will hear your voice.~\LINK{http://forum.learnnavi.org/language-updates/fula-short-form-of-fiul-a-and-tsal-oeti-steyki-nitram/}{ln}
(i.) \B{Ut\-ral Ay\-mok\-ri\-yä}, Tree of Voices (\emph{important spiritual site}).
\B{Ke lu kaw\-tu a nul\-ni\-vew oe po\-hu ti\-re\-a\-pi\-väng\-kxo äo Ut\-ral Ay\-mok\-ri\-yä.} There's nobody I'd rather commune with under the Tree of Voices.~\LINK{http://wiki.learnnavi.org/index.php?title=Corpus\#Lemondrop.com}{wiki} \\
\textsc{Derivatives}: \ding{43} 
\on{1}.~\HL{utral aymokriyA}{\B{Ut\-ral Ay\-mok\-ri\-yä}}, Tree of Voices. 
\on{2}.~\HL{spulmokri}{\B{spul\-mok\-ri}}, telephone. 
\on{3}.~\HL{kakmokri}{\B{kak\-mok\-ri}}, mute.

% momek
\hi
\HT{momek}{\B{\cp{mo}mek}} \I{["mo.mEk\textcorner]} \emph{adj.} hollow (\emph{not solid}).
\B{Tsa\-ta\-ngek\-ìri pam fkan mo\-mek.} That treetrunk sounds hollow.~\LINK{http://naviteri.org/2013/09/on-si-salewfya-shapes-and-directions/}{b}
(\ding{43}~\HL{mo}{\B{mo}}, \HL{mek}{\B{mek}})

% mong
\hi
\HT{mong}{\B{mong}} \I{[moN]} \emph{vtr.} depend on, rely on, trust for protection.
\B{Mong prr\-nen\-ìl fu\-ta sa'\-nul ve\-rewng nì\-tut.} The infant depends on her Mommy to look after her constantly.~\LINK{http://naviteri.org/2012/07/meetings-waterfalls-and-more/}{b}
\B{Tsun nga oet mi\-vong.} You can rely on me.~\LINK{http://naviteri.org/2012/07/meetings-waterfalls-and-more/}{b} \\
\textsc{Derivatives}: \ding{43} 
\on{1}.~\HL{tsukmong}{\B{tsukmong}}, reliable, dependable.

% mowan
\hi
\HT{mowan}{\B{mo\cp{wan}}} \I{[mo."wan]} \emph{adj.} (\emph{a}.) pleasing, enjoyable (\emph{implies physical or sensual pleasure, often has a sexual connotation}).
\B{Pll\-txe fko san nga\-ru lu mo\-wan Txil\-te ul\-te po\-ru nga.} I hear you like Txilte and vice versa.~\LINK{http://naviteri.org/2010/07/vocabulary-update/}{b} 
\B{Tì\-ru\-sol lu oe\-ru mo\-wan.} Singing is enjoyable to me.~\LINK{http://naviteri.org/2012/02/trr-asawnung-lefpom-happy-leap-day/}{b}
(\emph{b}.~\emph{slang}) like. \B{Tì\-tu\-sa\-ron mo\-wan lu oer nì\-ngay.} Hunting really turns me on.~\LINK{http://naviteri.org/2010/07/vocabulary-update/}{b}

% mowar / si
\hi
\HT{mowar}{\B{mo\cp{war}}} \I{[mo."waR]} \textsc{i}. \emph{n.} advice, bit or piece of advice.
\B{Nge\-yä mo\-war txan\-ley oer.} Your advice is invaluable to me.~\LINK{https://naviteri.org/2026/04/tsivola-liu-amip-thirty-two-new-words/}{b}
\B{Ma Ney\-ti\-ri, ay\-oel kin mo\-war\-it nge\-yä.} Neytiri, we need your advice.~\LINK{http://naviteri.org/2014/05/mipa-ayliu-mipa-aysafpil-new-words-new-ideas/}{b} \\
\textsc{ii}. \B{mo\cp{war si}} \I{[mo."waR si]} \emph{vin.} advise, give advice; (\B{tsnì}) advise s.o. to do s.t.
\B{Tsun oe mo\-war si\-vi ngar, slä ke tsun fya\-wi\-vìn\-txu.} I can advise you, but I can't guide you.~\LINK{http://naviteri.org/2014/05/mipa-ayliu-mipa-aysafpil-new-words-new-ideas/}{b}
\B{Poe mo\-war so\-li poan\-ur tsnì hi\-vum.} She advised him to leave.~\LINK{http://naviteri.org/2014/05/mipa-ayliu-mipa-aysafpil-new-words-new-ideas/}{b} \\
\textsc{Derivatives}: \ding{43} 
\on{1}.~\HL{mowarsiyu}{\B{mo\-war\-si\-yu}}, advisor.

% mowarsiyu
\hi
\HT{mowarsiyu}{\B{mo\cp{war}siyu}} \I{[mo."waR.si.ju]} \emph{n.} advisor.
\B{Lu fra\-eyk\-tan\-ur a\-sìl\-tsan txan\-tslu\-sam\-a ay\-mo\-war\-si\-yu.} Every good leader has wise advisors.~\LINK{http://naviteri.org/2014/05/mipa-ayliu-mipa-aysafpil-new-words-new-ideas/}{b}
\B{Ay\-sä\-mok\-ìri a\-txan\-tsan, oe\-yä ay\-mo\-war\-si\-yu\-ru ira\-yo nì\-txan!} I thank my advisors very much for the excellent suggestions!~\LINK{http://naviteri.org/2014/05/mipa-ayliu-mipa-aysafpil-new-words-new-ideas/}{b}
(\ding{43}~\HL{mowar}{\B{mo\-war}})

% mrr
\hi
\HT{mrr}{\B{mrr}} \I{[m\s{r}]} \emph{num., adj.} five.
\B{Alo a\-mrr po\-an po\-lawm, slä fra\-lo poe pol\-txe san ke\-he.} He asked five times but each time she said `no'.~\LINK{http://forum.learnnavi.org/language-updates/txelanit-hivawl/}{ln} 
\B{mrra riti}, five stingbats.~\LINK{http://naviteri.org/2011/07/number-in-na\%E2\%80\%99vi/}{b} \\
\textsc{Derivatives}: \ding{43} 
\on{1}.~\HL{mrrtrr}{\B{mrr\-trr}}, (\on{5}\hyp{}day) week. 
\on{2}.~\HL{vospxImrr}{\B{Vo\-spxì\-mrr}}, May.

% mrr-
\hi
\HT{mrr-}{\B{mrr-}} \I{[m\s{r}.]} \emph{num.~pref.} five (\emph{base to form higher numbers or fractions}). \B{mrr\-vol}, \on{40} (decimal), \on{50} (octal); \B{mrr\-pxì}, one-fifth, $\frac{1}{5}$; \B{mrr\-pxì a\-tsìng}, four-fifth, $\frac{4}{5}$.

% mrrpxì
\hi
\HT{mrrpxI}{\B{mrr\cp{pxì}}} \I{[m\s{r}."p'I]} \emph{n.} one-fifth, $\frac{1}{5}$.
(\ding{43}~\HL{mrr-}{\B{mrr-}}, \HL{hapxI}{\B{ha\-pxì}})

% mrrtrr / mrrtrram / mrrtrray
\hi
\HT{mrrtrr}{\B{\cp{mrr}trr}} \I{["m\s{r}.t\s{r}]} \emph{n., adv.} (\on{5}\hyp{}day) workweek.
\B{Kä\-teng oe hu ey\-lan Per\-lin\-mì a mrr\-trr lu tsyeym a ke tsun tswi\-va' kaw\-krr.} The (\on{5}-day) week I spent with my friends in Berlin is a treasure that I’ll never forget.~\LINK{http://naviteri.org/2013/06/looking-back-looking-forward/}{b} \\ 
(i.)~\B{mrr\-trr\-\cp{am}} \I{[m\s{r}.t\s{r}."{}am]} \emph{n., adv.} last (\on{5}\hyp{}day) workweek. \\
(ii.)~\B{mrr\-trr\-\cp{ay}} \I{[m\s{r}.t\s{r}"{}aj]} \emph{n., adv.} next (\on{5}\hyp{}day) workweek.

% mrrve
\hi
\HT{mrrve}{\B{\cp{mrr}ve}} \I{["m\s{r}.vE]} \emph{ord.~num., adj.} fifth. 

% mrrvol
\hi
\HT{mrrvol}{\B{\cp{mrr}vol}} \I{["m\s{r}.vol]} \emph{num., adj.} \on{40} (decimal), \on{50} (octal). 

% mu'ni
\hi 
\HT{mu'ni}{\B{\cp{mu'}ni}} \I{["muP.ni]} \emph{vtr.} accomplish, achieve (\emph{for achievements that are in some way significant}). 
\B{Krr\-ka tì\-rey a\-yol, pol mo\-lu'\-ni pxay\-a ay\-ut a tsran\-ten.} During her short life, she accomplished many important things.~\LINK{http://naviteri.org/2020/11/vospxivopeya-ayliu-amip-novembers-new-words/}{b} 
(\ding{43}~\HL{hasey si}{\B{ha\-sey si}}) \\
\textsc{Derivatives}: \ding{43} 
\on{1}.~\HL{tImu'ni}{\B{tì\-mu'\-ni}}, achievement, accomplishment.
\on{2}.~\HL{mu'nitkan}{\B{mu'\-nit\-kan}}, effective.

% mu'nitkan
\hi 
\HT{mu'nitkan}{\B{\cp{mu'}nitkan}} \I{["muP.nit\textcorner.kan]} \emph{or} \I{["mUP.]} \emph{adj.} effective. 
\B{Tsa\-ye\-rik\-ìl mu'\-nit\-kan\-a wan\-tseng\-it ro\-lun; ke tsun fko pot tsi\-ve\-'a kaw\-'it.} That hexapede has found an effective hiding place. He can't be seen at all.~\LINK{https://naviteri.org/2025/04/mevosinga-liu-amip-twenty-new-words/}{b}
(\ding{43}~\HL{mu'ni}{\B{mu'\-ni}}, \HL{tIkan}{\B{tì\-kan}}; \ding{214}~\HL{kemu'nitkan}{\B{ke\-mu'\-nit\-kan}}) \\
\textsc{Derivatives}: \ding{43} 
\on{1}.~\HL{tImu'nitkan}{\B{tì\-mu'\-nit\-kan}}, effectiveness.
\on{2}.~\HL{nImu'nitkan}{\B{nì\-mu'\-nit\-kan}}, effectively.

% muiä
\hi
\HT{muiA}{\B{mu\cp{i}ä}} \I{[mu."{}i.\ae]} \emph{adj.} (\emph{a}.) proper, fair, right, justified. 
\B{Tsay\-lì\-'u ke lu muiä!} Those words are improper!~\LINK{http://naviteri.org/2016/12/tengkrr-perahem-zisit-amip-as-the-new-year-arrives/}{b}
\B{Fra\-trr 'e\-rul a fay\-sä\-leym\-kem a\-yll muiä lu nì\-wotx.} The communal protests that are increasing every day are completely fair.~\LINK{http://naviteri.org/2020/06/mi-tanlokxe-oeya-srr-afpxamo-terrible-days-in-my-country/}{b}
\B{Nga sìl\-mi a tsa\-kem ke lu vo\-ìk a\-mu\-iä!} What you just did was not proper behavior!~\LINK{http://naviteri.org/2020/11/vospxivopeya-ayliu-amip-novembers-new-words/}{b} 
(\emph{b.~proverb}) \B{Kem a\-muiä, kum a\-fe'.} Proper action, bad result.\slash Well, their heart was in the right place (\emph{something that should have turned out well didn't}).~\LINK{http://naviteri.org/2012/06/spring-vocabulary-part-3/}{b} \\
\textsc{Derivatives}: \ding{43} 
\on{1}.~\HL{tImwiA}{\B{tìm\-wiä}}, fairness, justice.

% mulpxar
\hi
\HT{mulpxar}{\B{mul\cp{pxar}}} \I{[mul."p'aR]} \emph{or} \I{[mUl.]} \emph{n.}, \ding{96} roosterhead plant, \emph{alectophyllum molle}.

% -mun
\hi
\HT{-mun}{\B{-mun}} \I{[.mun]} \emph{or} \I{[.mUn]} \emph{num. suf.} two (\emph{to form higher numbers}). \B{vo\-mun}, \on{10} (decimal), \on{12} (octal). \B{me\-vo\-mun}, \on{18} (decimal), \on{22} (octal); etc. 
(\ding{43}~\HL{mune}{\B{mu\-ne}})

% mun'i
\hi
\HT{mun'i}{\B{mun\cp{'i}}} \I{[mun."Pi]} \emph{or} \I{[mUn.]} (\textsc{rn}: \B{mùn\cp{'i}} \I{[mUn."Pi]}) \emph{vtr.} cut. \\
\textsc{Derivatives}: \ding{43} 
\on{1}.~\HL{pxImun'i}{\B{pxì\-mun\-'i}}, divide, cut into parts. 
\on{2}.~\HL{syImawnun'i}{\B{syì\-maw\-nun\-'i}}, cut tie.

% mune
\hi
\HT{mune}{\B{\cp{mu}ne}} \I{["mu.nE]} \emph{num., adj.} two. (\emph{base}: \ding{43}~\B{me}\texttt{+}, \emph{rest}: \ding{43}~\HL{-mun}{\B{-mun}}, \emph{ordinal}: \ding{43}~\HL{muve}{\B{mu\-ve}})  \\
\textsc{Derivatives}: \ding{43} 
\on{1}.~\HL{nImun}{\B{nì\-mun}}, again. 
\on{2}.~\HL{muntrr}{\B{mun\-trr}}, weekend. 
\on{3}.~\HL{munsna}{\B{mun\-sna}}, pair.

% munsna
\hi
\HT{munsna}{\B{\cp{mun}sna}} \I{["mun.sna]} \emph{or} \I{["mUn.]} \emph{n.} pair.
\B{Hawn\-ven\-ìri lu oe\-ru mun\-sna a\-mrr.} I have five pairs of shoes.~\LINK{http://naviteri.org/2012/07/fivospxiya-aylifyavi-amip-this-months-new-expressions-pt-2/}{b}
\B{Hu\-fwa me\-fo le\-ru mun\-txa\-tu txan\-krr, mi lu mun\-snar ho\-na tì\-flrr a na pum me\-yaw\-ne\-tu\-ä a\-mip nì\-wotx.} Although the two of them have been mates a long time, they still have all the adorable tenderness of new sweethearts.~\LINK{http://naviteri.org/2013/02/vospxi-ayol-posti-apup-short-post-for-a-short-month/}{b}
(\ding{43}~\HL{mune}{\B{mu\-ne}}, \HL{sna'o}{\B{sna'o}})

% munsna-
\hi
\HT{munsna-}{\B{munsna-}} \I{[mun.sna.]} \emph{or} \I{[mUn.]} \emph{pref.~to mean} pair, couple. 
\B{Tì\-kang\-kem si fa mun\-sna\-tu\-te.} Work in pairs.~\LINK{http://naviteri.org/2012/07/fivospxiya-aylifyavi-amip-this-months-new-expressions-pt-2/}{b}
\B{Ta\-yel Tse\-nul pxe\-swi\-zaw\-ti yoa mun\-sna\-hawn\-ven.} Tsenu will receive three arrows in exchange for a pair of shoes.~\LINK{http://naviteri.org/2014/02/barter-and-exchange/}{b}

% muntrr / muntrram / muntrray
\hi
\HT{muntrr}{\B{\cp{mun}trr}} \I{["mun.t\s{r}]} \emph{or} \I{["mUn.]} \emph{n., adv.} weekend.  \\
(i.) \B{muntrr\cp{am}} \I{[mun.t\s{r}."{}am]} \emph{or} \I{[mUn.]} \emph{n., adv.} last weekend. 
\B{Mun\-trr\-am oe ko\-lä\-teng hu Ra\-lu ul\-te su\-nu oer tsa\-krr.} I spent last weekend with Ralu and had a good time [lit.: that time was pleasant to me].~\LINK{http://naviteri.org/2015/08/aylifyavi-lereyfya-2-cultural-terms-2-and-more/}{b} \\
(ii.) \B{muntrr\cp{ay}} \I{[mun.t\s{r}."{}aj]} \emph{or} \I{[mUn.]} \emph{n., adv.} next weekend. 
(\ding{43}~\HL{mune}{\B{mu\-ne}}, \HL{trr}{\B{trr}})

% muntxa / si / slu
\hi
\HT{muntxa}{\B{mun\cp{txa}}} \I{[mun."t'a]} \emph{or} \I{[mUn.]} \textsc{i}.~\emph{adj.} mated. \\
\textsc{ii}.a.~\B{mun\cp{txa} si} \I{[mun."t'a si]} \emph{or} \I{[mUn.]} \emph{vin.} mate (\ding{43}~\HL{hu}{\B{hu}}, with s.o.), marry s.o.
\B{Ìlä Fey\-ral, mun\-txa so\-li Ra\-lu sì Ne\-wey nì\-wan me\-srr\-am.} According to Peyral, Ralu and Newey were secretly married the day before yesterday.~\LINK{http://naviteri.org/2012/07/meetings-waterfalls-and-more/}{b}  
\B{Ne\-wey yaw\-ne lu oer ul\-te new oe mun\-txa si\-vi poe\-hu.} I love Ne\-wey and want to marry her.~\LINK{http://naviteri.org/2018/03/100a-liu-amip-64-new-words-part-1/}{b} \\
\textsc{ii}.b.~\B{mun\cp{txa} slu} I{[mun."t'a slu]} \emph{vin.} get married.
\B{Fko pll\-txe\-i\-e san me\-nga mun\-txa slo\-lu sìk!} I'm so happy to hear you got married!~\LINK{http://naviteri.org/2011/10/more-vocabulary-a-bit-of-grammar/}{b} \\
\textsc{Derivatives}: \ding{43} 
\on{1}.~\HL{muntxatu}{\B{mun\-txa\-tu}}, spouse.  
\on{2}.~\HL{tImuntxa}{\B{tì\-mun\-txa}}, marriage. 

% muntxatan
\hi
\HT{muntxatan}{\B{mun\cp{txa}tan}} \I{[mun."t'a.tan]} \emph{or} \I{[mUn.]} \emph{n.} husband, male spouse. 
(\ding{43}~\HL{muntxatu}{\B{mun\-txa\-tu}})

% muntxate
\hi
\HT{muntxate}{\B{mun\cp{txa}te}} \I{[mun."t'a.tE]} \emph{or} \I{[mUn.]} \emph{n.} wife, female spouse. 
(\ding{43}~\HL{muntxatu}{\B{mun\-txa\-tu}})

% muntxatu
\hi
\HT{muntxatu}{\B{mun\cp{txa}tu}} \I{[mun."t'a.tu]} \emph{or} \I{[mUn.]} \emph{n.} spouse. 
\B{Ma mun\-txa\-tu, oeng yaw\-ne lu (oe\-nga\-ru) fì\-tsap, ke\-fyak?} We love each other, don't we, my spouse?~\LINK{http://naviteri.org/2011/12/one-more-for-2011/}{b}
\B{Mip\-a me\-mun\-txa\-tu\-ru sey\-kxel sì ni\-tram.} Congratulations to the two newly-weds.~\LINK{http://naviteri.org/2014/02/barter-and-exchange/}{b}
\B{Mun\-txa\-tu\-ri So\-rewn\-ti ke tsun oe mi\-vll\-'an.} I can't accept Sorewn as my spouse.~\LINK{http://naviteri.org/2012/10/mipa-vospxi-mipa-ayliu-new-words-for-the-new-month/}{b}
(\ding{43}~\HL{muntxa}{\B{mun\-txa}}) \\
\textsc{Derivatives}: \ding{43} 
\on{1}.~\HL{muntxatan}{\B{mun\-txa\-tan}}, husband.   
\on{2}.~\HL{muntxate}{\B{mun\-txa\-te}}, wife.

% munge
\hi
\HT{munge}{\B{\cp{mu}nge}} \I{["mu.NE]} (\textsc{rn}: \B{\cp{mù}nge} \I{["mU.NE]}) \emph{vtr.} take, bring. 
\B{Tsun Txil\-te pam\-rel\-it ivi\-nan; ta\-fral puk\-ot a\-nutx mu\-nge fra\-tseng.} Txilte knows how to read; therefore she brings a thick book wherever she goes.~\LINK{http://naviteri.org/2011/11/\%E2\%80\%99a\%E2\%80\%99awa-ayli\%E2\%80\%99u-amip-ni\%E2\%80\%99aw-\%E2\%80\%94-a-few-new-words-only/}{b}
\B{Fu\-ri\-a fì\-lì'\-fya\-vit fkol ta 'Ìng\-lì\-sì mo\-lu\-nge, sìl\-pey oe tsnì ay\-ha\-pxì\-tu lì'\-fya\-o\-lo'\-ä aw\-nge\-yä ke stì\-ye\-vi.} That this expression is brought from English, will hopefully not anger the members of the language group.~\LINK{http://forum.learnnavi.org/language-updates/txelanit-hivawl/}{ln} \\
\textsc{Derivatives}: \ding{43} 
\on{1}.~\HL{sAmunge}{\B{sä\-mu\-nge}}, transportation device. 
\on{2}.~\HL{zamunge}{\B{za\-mu\-nge}}, bring. 
\on{3}.~\HL{mungwrr}{\B{mung\-wrr}}, except. 
\on{4}.~\HL{mungsye}{\B{mung\-sye}}, inhale.

% mungsye
\hi 
\HT{mungsye}{\B{\cp{mung}sye}} \I{["m$\cdot$uN.sjE]} \emph{or} \I{["m$\cdot$UN.]} (\textsc{rn}: \B{\cp{mùng}sye} \I{["m$\cdot$UN.SE]}) \emph{vin., inf.}\on{11} inhale. 
(\ding{43}~\HL{munge}{\B{mu\-nge}}, \HL{syeha}{\B{sye\-ha}}; \ding{214}~\HL{lonusye}{\B{lo\-nu\-sye}})

% mungwrr
\hi
\HT{mungwrr}{\B{mung\cp{wrr}}} \I{[muN."w\s{r}]} \emph{or} \I{[mUN.]} (\textsc{rn}: \B{mùng\cp{wrr}} \I{[mUN."w\s{r}]}) \emph{adp.} except.
\B{Me\-lì\-'u alu mung\-wrr sì nì\-fkrr ngam\-pam si.} The words \emph{mungwrr} and \emph{nìfkrr} rhyme.~\LINK{http://naviteri.org/2012/01/mipa-zisit-ayliu-amip-new-words-for-the-new-year/}{b}
\B{Oe\-ru ke tsun li\-vam ke\-'u lor to Ey\-wa\-'e\-veng\-ä na'\-rìng a lew sä\-po\-li fa prr\-wll, kxawm mung\-wrr fì\-ki\-fkey a lew sä\-pì\-ye\-vi fa fpom sì lì'\-fya le\-Na'\-vi.} Nothing could be more beautiful to me than a Pandoran forest that has covered in moss, except perhaps this world soon covering itself in peaceful well-being and the Na'vi language.~\LINK{http://forum.learnnavi.org/language-updates/prrwll-moss-the-comparative-to-construct/}{ln} 
(\emph{a. with} \HL{fwa}{\B{fwa}}) except that. 
\B{Po\-ru ke po\-leng oel ke\-'ut mung\-wrr fwa Ra\-lul ke tsa\-tse\-nget.} I told her nothing except that Ralu wasn't there.~\LINK{http://naviteri.org/2020/02/some-words-for-leap-year-day/}{b}
(\ding{43}~\HL{munge}{\B{mu\-nge}}) \\
\textsc{Derivatives}: \ding{43} 
\on{1}.~\HL{tImungwrr}{\B{tì\-mung\-wrr}}, exception. 
\on{2}.~\HL{mungwrrtxo}{\B{mung\-wrr\-txo}}, unless, except if.
\on{3}.~\HL{lemungwrr}{\B{le\-mung\-wrr}}, exceptional.

% mungwrrtxo
\hi 
\HT{mungwrrtxo}{\B{mung\cp{wrr}txo}} \I{[muN."w\s{r}.t'o]} \emph{or} \I{[mUN.]} \emph{or coll.} \I{[muN."w\s{r}.to]} (\textsc{rn}: \B{mùng\cp{wrr}txo} \I{[mUN."w\s{r}.do]}) \emph{conj.} unless, except if. 
\B{Tsak\-tap rä\-'ä si kaw\-krr mung\-wrr\-txo ke li\-vu kea fya\-'o a\-la\-he.} Never use violence unless there is no other way. ~\LINK{http://naviteri.org/2020/02/some-words-for-leap-year-day/}{b}
(\ding{43}~\HL{mungwrr}{\B{mung\-wrr}}, \HL{txo}{\B{txo}})

% muve
\hi
\HT{muve}{\B{\cp{mu}ve}} \I{["mu.vE]} \emph{ord. num., adj.} second. 
\B{Lì\-'ut a\-mu\-ve oel tswo\-la'.} I forgot the second word.~\LINK{http://naviteri.org/2011/04/\%E2\%80\%99a\%E2\%80\%99awa-li\%E2\%80\%99fyavi-amip\%E2\%80\%94a-few-new-expressions/comment-page-1/\#comment-677}{b}
\B{Ul\-te mu\-ve\-a sä\-pll\-txe le\-iu eyawr nì\-teng.} And the second statement is also correct, I'm happy to say.~\LINK{http://naviteri.org/2014/02/barter-and-exchange/comment-page-1/\#comment-2614}{b}
(\ding{43}~\HL{mune}{\B{mu\-ne}})

% muwìntxu
\hi
\HT{muwIntxu}{\B{muwìn\cp{txu}}} \I{[mu.w$\cdot$In."t'$\cdot$u]} \emph{vtr., inf.}\on{23} introduce, present (\emph{people or things}).   
\B{Oel mu\-wi\-vìn\-txu nga\-ru oe\-yä ler\-tut.} Let me present my colleague to you (\emph{formal}).\slash \B{Nga\-ru oe\-yä ler\-tut.} This is my colleague (\emph{coll.}).~\LINK{http://naviteri.org/2010/09/getting-to-know-you-part-1/comment-page-1/\#comment-297}{b}
\B{Tsa\-tsmu\-kel alu Ri\-ni mo\-lok fu\-ta oe ngar mu\-wä\-pi\-vìn\-txu.} Sister Rini over there suggested that I introduce myself.~\LINK{http://naviteri.org/2010/09/getting-to-know-you-part-1/}{b}
(\ding{43}~\HL{wIntxu}{\B{wìn\-txu}}) \\

\end{multicols}



% ===== Section N ===== 

 \pdfbookmark[0]{N}{N}
\begin{center}
\section*{N \I{[n]} (\emph{NeN})}
\end{center}

\begin{multicols}{2}

% na
\hi
\HT{na}{\B{na}} \I{[na]} \emph{adp.} (\emph{a.}) like, as.   
\B{Ay\-lì\-'u (lu) na ay\-skxe mì te'\-lan.} The words are like stones in my heart.~\LINK{http://wiki.learnnavi.org/index.php?title=Corpus\#Behind_the_Scenes}{a\slash wiki}
\B{Fì\-txon na ton a\-la\-he nì\-wotx pe\-lun ke lu teng?} Why is this night not the same like all other nights?~\LINK{http://whyisthisnight.com/addl-more.html\#na\%27vi}{tn}
(\emph{b. with colors to further subdivide the spectrum and name colors more specifically}) 
\B{Fì\-syu\-lang lu ta'\-leng\-na ean.} This flower is skin\hyp{}blue.~\LINK{http://naviteri.org/2010/08/mipa-ayopin-mipa-ayliu-new-colors-new-words/}{b}
\B{Fì\-syu\-lang a\-ean\hyp{}na\hyp{}ta'leng lor lu nì\-txan.} This skin\hyp{}blue flower is very beautiful.~\LINK{http://naviteri.org/2010/08/mipa-ayopin-mipa-ayliu-new-colors-new-words/}{b}
(\emph{c. in comparisons of equality}) as \emph{adj.}\slash \emph{adv.} as \emph{n.}\slash \emph{pn.} \B{Oe lu nì\-ftxan sìl\-tsan na nga.} I am as good as you.~\LINK{http://forum.learnnavi.org/language-updates/as-x-as-y-keep-on-keepin-on/}{ln}
\B{Oe tul nì\-ftxan nì\-win na nga.} I run as fast as you.~\LINK{http://forum.learnnavi.org/language-updates/confirmations-\%28comparisons-gerund\%29/}{ln}
\B{Oel ye\-rik\-it ta\-ron nì\-ftxan nìl\-tsan na nga.} I hunt \emph{yerik} as well as you.~\LINK{http://forum.learnnavi.org/language-updates/confirmations-\%28comparisons-gerund\%29/}{ln}
(\emph{d. idiom}) \B{nì\-win na hu\-fwe}, `as fast as the wind' (\emph{any situation to express rapidity}). \B{na hu\-fwe}, fluent (not just rapid) use of language.
\B{Fte\-ria oel lì'\-fya\-ti le\-Na'\-vi, slä mi ke tsä\-ngun pi\-vll\-txe na hu\-fwe.} I'm studying Na'vi, but I'm afraid I still can't speak it fluently.~\LINK{http://naviteri.org/2015/08/aylifyavi-lereyfya-2-cultural-terms-2-and-more/}{b}
(\emph{e. idiom}) \B{na fkxen eo fkio}, to go to waste [lit.: `like vegetable food before a tetrapteron'] 
\B{Nge\-yä fì\-tì\-kang\-kem\-vi a\-txan\-tsan ke sla\-yu na fkxen eo fkio.} This excellent work of yours will not go to waste.~\LINK{https://naviteri.org/2024/12/odds-and-ends/}{b}
(\ding{43}~\HL{pxel}{\B{pxel}})

% na'rìng
\hi
\HT{na'rIng}{\B{\cp{na'}rìng}} \I{["naP.RIN]} \emph{n.} forest. 
\B{Na'\-rìng sngä\-'i ni\-vrr txon\-krr syu\-ra\-tan\-fa.} The forest begins to glow at night with bioluminescence.~\LINK{http://naviteri.org/2012/07/fivospxiya-aylifyavi-amip-this-months-new-expressions-pt-2/}{b}
\B{Lu na'\-rìng tsuk\-ha\-haw.} One can sleep in the forest [lit.: the forest ist `sleepable'].~\LINK{http://naviteri.org/2011/03/\%E2\%80\%9Creceptive-ability\%E2\%80\%9D-and-hesitation/}{b}
\B{Na'\-rìng\-ìl nga' pxay\-a io\-ang\-it.} The forest contains many animals.~\LINK{http://naviteri.org/2011/05/weather-part-2-and-a-bit-more-2/}{b}
\B{Oel hu Txe\-wì trr\-am na'\-rìng\-it tar\-mok.} Yesterday I was with Txewì in the forest.~\LINK{http://wiki.learnnavi.org/index.php?title=Corpus\#New_York_Times}{wiki}
\B{Tì\-o\-mum\-mì oe\-yä, pol na'\-rìng\-it inan nìl\-tsan.} As far as I know, he reads the forest well.~\LINK{http://naviteri.org/2011/09/\%E2\%80\%9Cby-the-way-what-are-you-reading\%E2\%80\%9D/}{b}
\B{Sra\-ke pol la\-yok txew\-ti na'\-rìng\-ä?} Will he approach the edge of the forest?~\LINK{http://naviteri.org/2012/03/spring-vocabulary-part-1/}{b}
(\emph{a. idiom}) \B{zawr (a) mì na'\-rìng}, ``Old news.'' [lit.: an animal cry in the forest].~\LINK{http://naviteri.org/2021/08/ulte-ayyoratu-leiu-and-the-winners-are-2/}{b}
\B{Sra\-ke ngal sto\-lawm fu\-ta Tse\-nu Ra\-lur mo\-wan lu nì\-txan? -- Zawr mì na'\-rìng, ma tsmuk! Tsat o\-mum oel kin\-trr\-o.} Have you heard that Ralu has the hots for Tsenu? -- Old news, brother! I've known it for a week.~\LINK{http://naviteri.org/2021/08/ulte-ayyoratu-leiu-and-the-winners-are-2/}{b}

% Na'vi / nìNa'vi / leNa'vi
\hi
\HT{na'vi}{\B{\cp{Na'}vi}} \I{["naP.vi]} \emph{plural n.} (\emph{a.}) the Na'vi, the People, indigenous Pandorean sentient humanoids.  
\B{Tì\-'e\-fu\-mì oe\-yä su\-te\-kip nì\-wotx, ftxey Na'\-vi ftxey Saw\-tu\-te, lu sìl\-tsan lu kawng.} In my opinion, among all people, whether Na'vi or Skypeople, there are good and bad.~\LINK{http://naviteri.org/2010/07/ting-mikyun-fte-tslivam-listening-comprehension-1-3/}{b}
\B{Nì\-trr\-trr yom Na'\-vil wu\-tsot 'aw\-si\-teng pxaw yll\-txep.} The Na'vi regularly eat dinner together around a communal fire.~\LINK{http://naviteri.org/2012/07/fivospxiya-aylifyavi-amip-this-months-new-expressions-pt-2/}{b}
\B{Kea tsam\-si\-yu le\-sno\-nrr\-a ke tsun Na'\-vit iveyk.} No warrior full of self-pride can lead the People.~\LINK{http://naviteri.org/2012/07/fivospxiya-aylifyavi-amip-this-months-new-expressions-pt-2/}{b}
\B{Tsa\-sä\-lä\-txayn Na'\-vi\-ru srung so\-li nì\-'aw fte sli\-vu txur nì\-'ul.} That defeat only helped the People become stronger.~\LINK{http://naviteri.org/2012/01/more-additions-to-the-lexicon/}{b}
\B{Nga\-ri hu Ey\-wa sa\-lew ti\-re\-a, tokx 'ì\-'awn slu Na'\-vi\-yä ha\-pxì.} Your spirit goes with Eywa, your body remains to become part of the People.~\LINK{http://wiki.learnnavi.org/index.php?title=Corpus\#Behind_the_Scenes}{a\slash wiki}
\B{Na'\-vi\-ri txin\-a sä\-'o tì\-tu\-sa\-ron\-ä lu tsko swi\-zaw.} For the Na'vi the bow and arrow is the main hunting tool.~\LINK{http://naviteri.org/2011/01/new-year-new-vocabulary/}{b}
(\emph{b.}) \emph{shortened for} the Na'vi language. \B{'I\-vong Na'\-vi.} Let Na'vi bloom.\slash May the Na'vi language bloom.~\LINK{http://forum.learnnavi.org/news-announcements/a-response-from-paul-frommer!/}{ln} \\  
\textsc{Derivatives}: \ding{43} \on{1}.~\HL{lena'vi}{\B{le\-Na'\-vi}}, having to do with the Na'\-vi. 
\on{2}.~\HL{nIna'vi}{\B{nì\-Na'\-vi}}, like the Na'\-vi, Na'\-vi\-ly.

% naer
\hi
\HT{naer}{\B{na\cp{er}}} \I{[na."{}ER]} \emph{countable n.} drink, beverage.
\B{Fì\-na\-er ftxì\-vä' lu nì\-hawng, ha swey\-lu txo ngal tsat kxi\-vukx nì\-win.} This drink tastes horrible, so you'd better swallow it quickly.~\LINK{http://naviteri.org/2011/11/\%E2\%80\%99a\%E2\%80\%99awa-ayli\%E2\%80\%99u-amip-ni\%E2\%80\%99aw-\%E2\%80\%94-a-few-new-words-only/}{b}
\B{Fì\-na\-er\-ìri ngal ew\-ku 'u\-ot a\-stxong srak?} Do you taste something strange in this drink?~\LINK{http://naviteri.org/2012/11/renu-ayinanfyaya-the-senses-paradigm/}{b}
\B{Fì\-na\-er\-ìri sur fkan oe\-ru ka\-lin.} This drink tastes sweet to me.~\LINK{http://naviteri.org/2012/11/renu-ayinanfyaya-the-senses-paradigm/}{b}

% nafì'u
\hi
\HT{nafI'u}{\B{nafì\cp{'u}}} \I{[na.fI."Pu]} \emph{adj.} such.
\B{Na\-fì\-'u\-a mok\-ri lu sä\-kel\-pxìm\-run.} Such a voice is a rarity.~\LINK{https://naviteri.org/2026/04/tsivola-liu-amip-thirty-two-new-words/}{b}
\B{Krr\-o len ay\-hem a\-na\-fì\-'u.} Sometimes these things happen.~\LINK{http://forum.learnnavi.org/language-updates/just-another-few-words/}{ln}
(\ding{43}~\HL{na}{\B{na}}, \HL{fI'u}{\B{fì\-'u}})

% nafpawng
\hi 
\HT{nafpawng}{\B{na\cp{fpawng}}} \I{[na."fpawN]} \emph{adv.} grievingly, with grief. 
(\ding{43}~\HL{afpawng}{\B{afpawng}})

% nal
\hi 
\HT{nal}{\B{nal}} \I{[nal]} \emph{vin.} suffer. 
\B{Krr\-a tse\-'a oel tì\-fkey\-tok\-it txan\-lo\-kxe\-yä oey, nal nì\-txan mì te'\-lan.} When I see the condition of my country, I suffer greatly in my heart.~\LINK{naviteri.org/2025/07/ayu-leketengsyon-various-things/}{b} 
(\ding{43}~\HL{ngA'An}{\B{ngä\-'än}}) \\
\textsc{Derivatives}: \ding{43} 
\on{1}.~\HL{nalsteng}{\B{nal\-steng}}, feel empathy\slash compassion.

% nalsteng
\hi 
\HT{nalsteng}{\B{\cp{nal}steng}} \I{["n$\cdot$al.stEN]} \emph{vin., inf.}\on{11} feel empathy, feel compassion (for, \ding{43}~\HL{hu}{\B{hu}} \emph{or the topic}). 
\B{Po\-hu\slash Po\-ri  oe nal\-steng.} I feel compassion for him.~\LINK{naviteri.org/2025/07/ayu-leketengsyon-various-things/}{b} 
(\ding{43}~\HL{nal}{\B{nal}}, \HL{'awsiteng}{\B{'aw\-si\-teng}}) \\
\textsc{Derivatives}: \ding{43} 
\on{1}.~\HL{tInalsteng}{\B{tì\-nal\-steng}}, compassion, empathy.
\on{2}.~\HL{lenalsteng}{\B{le\-nal\-steng}}, compassionate (\emph{ofp.}).

% nalutsa
\hi 
\HT{nalutsa}{\B{na\cp{lu}tsa}} \I{[na."lu.\t{ts}a]} \emph{n., fauna} nalutsa.~\LINK{https://forum.learnnavi.org/language-updates/expansion-to-flora-fauna-and-places-from-the-valley-of-mo'ara/}{ln} 

% nam'ake
\hi
\HT{nam'ake}{\B{nam\cp{'a}ke}} \I{[nam."Pa.kE]} \emph{adv.} confidently. 
\B{Tu\-ke\-ru pol\-txe Ak\-wey nam\-'a\-ke, omum fu\-ta ke tsun poe sti\-vo.} Akwey spoke to Tuke confidently, knowing that she couldn't refuse.~\LINK{http://naviteri.org/2012/06/spring-vocabulary-part-3/}{b}
(\ding{43}~\HL{am'ake}{\B{am\-'a\-ke}})

% nantang
\hi
\HT{nantang}{\B{\cp{nan}tang}} \I{["nan.taN]} \emph{n., fauna} viperwolf, striped armored wolf, \emph{caniferratus costatus}.
\B{ay\-nan\-tang sì ay\-ri\-ti}, viperwolves and stingbats.~\LINK{http://naviteri.org/2010/07/thoughts-on-ambiguity/}{b}
\B{Nan\-tang\-ìl yom ye\-rik\-it.} Viperwolves eat hexapedes.~\LINK{http://naviteri.org/2011/07/number-in-na\%E2\%80\%99vi/}{b}
\B{Ay\-hem\-ìri 'e\-wan\-a tsa\-nan\-tang\-ur a\-hì\-'i tìng na\-ri.} Look at what that little young viperwolf is doing.~\LINK{http://naviteri.org/2012/01/more-additions-to-the-lexicon/}{b} \\  
\textsc{Derivatives}: \ding{43} \on{1}.~\HL{nantangtsyIp}{\B{nantangtsyìp}}, dog.

% nantangtsyìp
\hi 
\HT{nantangtsyIp}{\B{\cp{nan}tangtsyìp}} \I{["nan.taN.\t{ts}jIp\textcorner]} \emph{n.} dog (Earth animal). 
\B{Oey nan\-tang\-tsyìp o\-lu\-e' ta\-lu\-na yom nì\-hawng.} My dog vomited because it ate too much.~\LINK{http://naviteri.org/2016/06/mrrvola-lifyavi-amip-forty-new-expressions/}{b} 
(\ding{43}~\HL{nantang}{\B{nan\-tang}}, \HL{-tsyIp}{\B{-tsyìp}})

% nang
\hi
\HT{nang}{\B{nang}} \I{[naN]} \emph{part. for surprise, exclamation, encouragement (mirative); always sentence\hyp{}final; appears with a form of} \ding{43}~\HL{txan}{\B{txan}}.
\B{Lì'\-fya\-vi le\-sar nì\-txan nang!} A very useful expression, indeed!~\LINK{http://forum.learnnavi.org/language-updates/just-another-few-words/}{ln}
\B{kem\-lì\-'u\-vi (lì\-'u a\-txan\-tsan nang!)}, infix (a truly excellent word!)~\LINK{http://naviteri.org/2010/07/vocabulary-update/comment-page-1/\#comment-95}{b}

% napxì
\hi 
\HT{napxI}{\B{na\cp{pxì}}} \I{[na."p'I]} \emph{adv.} partially, in part. 
\B{Fì\-tì\-kang\-kem\-vi ha\-sey lu na\-pxì nì\-'aw.} This project is only partially complete.~\LINK{https://naviteri.org/2025/04/mevosinga-liu-amip-twenty-new-words/}{b}
(\ding{43}~\HL{hapxI}{\B{ha\-pxì}}) 

% naranawm
\hi 
\HT{naranawm}{\B{Nara\cp{nawm}}} \I{[naR.a."nawm]} \emph{n., astr.} Polyphemus [lit.: Great Eye]. 
\B{Na\-ra\-nawm\-ä mawl mi lu uo taw\-txew.} Half of Polyphemus is still behind the horizon.~\LINK{http://naviteri.org/2018/07/ta-sulfatu-a-ayliu-niul-more-words-from-our-experts/}{b}
(\ding{43}~\HL{nari}{\B{na\-ri}}, \HL{nawm}{\B{nawm}})

% nari / nari si
\hi
\HT{nari}{\B{\cp{na}ri}} \I{["na.Ri]} \textsc{i}.~\emph{n.} eye.
\B{Nge\-yä me\-na\-ri lor lu nì\-txan.} Your eyes are very beautiful.~\LINK{https://www.youtube.com/watch?v=ojr_BE8ZmVE}{yt}
\B{Oe\-ri me\-na\-ri\-ru pi\-ak si.} I open my eyes.~\LINK{http://forum.learnnavi.org/language-updates/open-close/}{ln}
\B{Po\-ri me\-na\-ri pi\-ak sä\-po\-li.} His eyes opened.~\LINK{http://forum.learnnavi.org/language-updates/open-close/}{ln} \\
\textsc{ii}.a.~\B{\cp{na}ri si} \I{["na.Ri si]} \emph{vin.} watch out, be careful. \B{Na\-ri so\-li ay\-oe fte\-ke nì\-hawng li\-vok.} We were careful not to get too close.~\LINK{http://wiki.learnnavi.org/index.php?title=Corpus\#New_York_Times}{wiki}
\B{Na\-ri si fte kea fe\-kem ke li\-ven ngar!} Be careful you don't have an accident!~\LINK{http://naviteri.org/2011/07/txantsana-ultxa-mi-siatll-great-meeting-in-seattle/}{b}
\B{Fì\-txe\-le\-ri ze\-ne oe na\-ri si\-vi!} I have to be careful about this matter!~\LINK{http://naviteri.org/2011/09/\%E2\%80\%9Cby-the-way-what-are-you-reading\%E2\%80\%9D/comment-page-1/\#comment-1093}{b}
\B{Na\-ri si txo\-ke\-fyaw tì\-ran nì\-nu.} Be careful or you'll trip.~\LINK{http://naviteri.org/2013/04/wheres-the-bathroom-and-other-useful-things/}{b} 
(\emph{to express `do X carefully'}) \B{nari si} \texttt{+} \emph{root verb} 
\B{Na\-ri si keyn\-ven!} Step carefully!~\LINK{http://naviteri.org/2019/08/aawa-lifya-si-lifyavi-amip-a-few-new-words-and-expressions/}{b}
\B{Na\-ri si lo\-nu swi\-zaw\-it.} Release the arrow carefully.~\LINK{http://naviteri.org/2019/08/aawa-lifya-si-lifyavi-amip-a-few-new-words-and-expressions/}{b} 
(\ding{43}~\HL{lare}{\B{la\-re}}) \\
\textsc{ii}.b.~\B{tìng \cp{na}ri} \I{[tIN."na.Ri]} \emph{or coll.} \I{[tI."na.Ri]} \emph{vin.} (\texttt{+}\emph{control}) look, look out, look at.
\B{Tìng na\-ri ne\-kll.} Look below.~\LINK{http://naviteri.org/2013/08/fmawn-a-fkol-kilmulat-this-just-in/}{b}
\B{Po\-ru tìng na\-ri!} Look at him.~\LINK{http://naviteri.org/2012/11/renu-ayinanfyaya-the-senses-paradigm/}{b}
\B{Le\-hrr\-ap lu fwa evi\-tsyìp sle\-le mì hil\-van lu\-ke fwa fyeyn\-tu te\-rìng na\-ri.} It's dangerous for tiny ones to swim in the river without an adult watching.~\LINK{http://naviteri.org/2011/02/new-vocabulary-ii\%E2\%80\%94part-1/}{b} \\  
\textsc{Derivatives}: \ding{43} \on{1}.~\HL{narlor}{\B{nar\-lor}}, (visually) beautiful. 
\on{2}.~\HL{narvA'}{\B{nar\-vä'}}, ugly, unsightly.
\on{3}.~\HL{naranawm}{\B{Na\-ra\-nawm}}, Polyphemus.
\on{4}.~\HL{naritxim}{\B{na\-ri\-txim}}, eyethorn.

% naritxim
\hi 
\HT{naritxim}{\B{\cp{na}ritxim}} \I{["na.Ri.t'im]} \emph{n.} \ding{96} eyethorn. 
(\ding{43}~\HL{nari}{\B{nari}}, \HL{txim}{\B{txim}}) 

% narlor
\hi
\HT{narlor}{\B{nar\cp{lor}}} \I{[naR."loR]}  \emph{adj.} beautiful visually. 
(\ding{214}~\HL{narvA'}{\B{nar\-vä'}}; \ding{43}~\HL{nari}{\B{na\-ri}}, \HL{lor}{\B{lor}})

% narvä'
\hi
\HT{narvA'}{\B{nar\cp{vä'}}} \I{[naR."v\ae{}P]}  \emph{adj.} ugly, unsightly. 
(\ding{214}~\HL{narlor}{\B{nar\-lor}}; \ding{43}~\HL{nari}{\B{na\-ri}}, \HL{vA'}{\B{vä'}})

% natkenong 
\hi
\HT{natkenong}{\B{nat\cp{ke}nong}} \I{[nat\textcorner."kE.noN]} \emph{conj.\slash adv.} for example.
\B{Slä lu 'a\-'aw\-a tì\-ke\-teng, nat\-ke\-nong, tsyokx\-ìri ke lu zek\-wä a\-tsìng ki a\-mrr.} But there are several differences, for example, the hand doesn't have four but five digits.~\LINK{http://naviteri.org/2010/07/ting-mikyun-fte-tslivam-listening-comprehension-1-3/}{b}
(\ding{43}~\HL{na}{\B{na}}, \HL{tIkenong}{\B{tì\-ke\-nong}})

% natxu
\hi 
\HT{natxu}{\B{na\cp{txu}}} \I{[na."t'u]} \emph{vtr.} disapprove of. 
\B{Oel nge\-yä tì\-hawl\-it na\-txu ulte tsa\-wä wa\-syem.} I disapprove of your plan and will oppose (fight against) it.~\LINK{http://naviteri.org/2020/02/some-words-for-leap-year-day/}{b} 
(\ding{214}~\HL{smaw}{\B{smaw}}) \\  
\textsc{Derivatives}: \ding{43} \on{1}.~\HL{tInatxu}{\B{tì\-na\-txu}}, disapproval.

% nawang
\hi 
\HT{nawang}{\B{\cp{na}wang}} \I{["na.waN]} \emph{vin.} merge, become one with (\HL{hu}{\B{hu}}, with). 
\B{Tì\-mun\-txa\-maw lam Ni\-nat\-ur fwa vit\-ra sne\-yä no\-la\-wang hu pum mun\-txa\-tu\-ä.} After her marriage, it seemed to Ninat that her soul had merged with that of her mate.~\LINK{http://naviteri.org/2018/08/fmawnti-stolawm-srak-have-you-heard-the-news/}{b} 

% nawfwe
\hi
\HT{nawfwe}{\B{naw\cp{fwe}}} \I{[naw."fwE]} \emph{adj.} fluent (for speech). (\emph{from:} \B{(nì\-win) na hu\-fwe})
\B{Toi\-tsye\-ri lu poe pll\-txe\-yu a\-naw\-fwe nì\-txan.} She's a very fluent speaker of German.~\LINK{http://naviteri.org/2015/08/aylifyavi-lereyfya-2-cultural-terms-2-and-more/}{b}

% nawkx 
\hi 
\HT{nawkx}{\B{nawkx}} \I{[nawk']} \emph{n., fauna} bone helm rhino.

% nawm
\hi
\HT{nawm}{\B{nawm}} \I{[nawm]} \emph{adj.} great, noble.
\B{Na'\-vi\-ru set lu nawm\-a eyk\-tan a\-mip a lar\-mu Taw\-tu\-te.} The Na'vi have a great new leader now who had been a Skyperson.~\LINK{http://naviteri.org/2010/07/ting-mikyun-fte-tslivam-listening-comprehension-1-3/}{b}
\B{Ut\-ral\-ä a\-nawm | Ay\-ri\-na' lu ay\-oeng.} We are all seeds | Of the Great Tree.~\LINK{http://www.pandorapedia.com/navi/music/ritual_music}{pp} \\
(i.) \B{Nawm\-a Sa'\-nok}, the Great Mother, Eywa. 
\B{Tsawl\-a pa\-lu\-lu\-kan\-ìl oe\-ti 'ìl\-me\-ko a krr, Nawm\-a Sa'\-nok lrr\-tok sei\-yi ul\-te oe em\-ro\-ley.} The Great Mother was looking after me when the big thanator attacked and I've survived.~\LINK{http://naviteri.org/2011/05/some-miscellaneous-vocabulary/}{b} \\  
\textsc{Derivatives}: \ding{43} \on{1}.~\HL{nawmtu}{\B{nawm\-tu}}, great, noble person. 
\on{2}.~\HL{naranawm}{\B{Na\-ra\-nawm}}, Polyphemus. 
\on{3}.~\HL{nawmtoruktek}{\B{nawm\-to\-ruk\-tek}}, Toruk Makto Totem.
\on{4}.~\HL{manawmtu}{\B{ma\-nawm\-tu}}, ``excuse me''.
\on{5}.~\HL{nawnumtseng}{\B{naw\-num\-tseng}}, university.

% nawmtoruktek
\hi
\HT{nawmtoruktek}{\B{nawm\cp{to}ruktek}} \I{[nawm."to.Ruk\textcorner.tek\textcorner]} \emph{or} \I{[.rUk.]} (\textsc{rn}: \B{nawm\cp{to}rùktek} \I{[nawm."to.RUk\textcorner.tEk\textcorner]}) \emph{n.} Toruk Makto Totem.
(\ding{43}~\HL{nawm}{\B{nawm}}, \HL{toruk}{\B{to\-ruk}}, \HL{tekre}{\B{tek\-re}})

% nawmtu
\hi
\HT{nawmtu}{\B{\cp{nawm}tu}} \I{["nawm.tu]} \emph{n.} great person. 
(\ding{43}~\HL{nawm}{\B{nawm}})
\textsc{Derivatives}: \ding{43} 
\on{1}.~\HL{manawmtu}{\B{ma\-nawm\-tu}}, ``excuse me''.

% nawnumtseng
\hi 
\HT{nawnumtseng}{\B{naw\cp{num}tseng}} \I{[naw."num.\t{ts}EN]} \emph{or} \I{["nUm.]} \emph{n.} university. 
\B{Naw\-num\-tseng\-ä ay\-nu\-me\-yu\-ri ley ay\-nga\-ru tì\-teng\-wotx, ay\-oe\-ru ke\-teng\-syon.} Regarding university students, you value uniformity, we value diversity.~\LINK{https://naviteri.org/2025/06/vospikin-lefpom-happy-july/}{b} 
(\ding{43}~\HL{nawm}{\B{nawm}}, \HL{numtseng}{\B{num\-tseng}})

% nawri
\hi
\HT{nawri}{\B{\cp{naw}ri}} \I{["naw.ri]} \emph{adj.} talented. 
\B{Nga lu rol\-yu a\-naw\-ri slä Ni\-nat lu pum a\-swey.} You're a talented singer but Ninat is the best (one).~\LINK{http://naviteri.org/2015/04/some-new-words-for-may-day/}{b} \\  
\textsc{Derivatives}: \ding{43} \on{1}.~\HL{tInawri}{\B{tì\-naw\-ri}}, talent.

% nayn
\hi
\HT{nayn}{\B{nayn}} \I{[najn]} \emph{num.}, \textsc{lw} nine, (the symbol) \on{9} (\emph{only as Human decimal figure}). 
(\ding{43}~\HL{'eyt}{\B{'eyt}})

% näk 
\hi
\HT{nAk}{\B{näk}} \I{[n\ae{}k\textcorner]} \emph{vtr.} drink.
\B{Nge\-yä tsngal lum\-pe lu mek? Näk nì\-'ul ko!} Why is your cup empty? Drink up!~\LINK{http://naviteri.org/2012/02/trr-asawnung-lefpom-happy-leap-day/}{b}
\B{Ngay\-txoa, ke tsun oe swo\-at ni\-väk.} Sorry, but I can't drink spirits.~\LINK{http://forum.learnnavi.org/language-updates/new-words-from-the-2012-avatarmeet-navi-class/msg551479/\#msg551479}{ln}
\B{Kxawm swo\-at nì\-hawng nil\-vä\-ngäk oel txo\-nam.} Perhaps I've had too much to drink last night.~\LINK{http://naviteri.org/2013/12/mipa-zisit-lefpom-happy-new-year/comment-page-1/\#comment-2585}{b} \\  
\textsc{Derivatives}: \ding{43} \on{1}.~\HL{nInAk}{\B{nì\-näk}}, by drinking, in a liquid way.

% nän
\hi
\HT{nAn}{\B{nän}} \I{[n\ae{}n]} \textsc{i}. \emph{vin.} decrease.
\B{Fra\-zì\-sì\-krr pay kil\-van\-ä nän nì\-'ul\-'ul.} Every season the river dries up more and more. \\
\textsc{ii}. \emph{conj. to form correlative comparisons} (\emph{a.}) \B{nän … nän}, the less … the less.   
\B{Nän ftia, nän lu skxom a em\-za\-'u.} The less you study, the less chance you have of passing.~\LINK{http://naviteri.org/2012/02/trr-asawnung-lefpom-happy-leap-day/}{b} 
(\emph{b.}) \B{nän … 'ul}, the less … the more. 
\B{Nän yom kxam\-trr, 'ul 'e\-fu ohakx kaym.} The less you eat at noon, the hungrier you'll feel in the evening.~\LINK{http://naviteri.org/2012/02/trr-asawnung-lefpom-happy-leap-day/}{b}
(\ding{214}~\HL{'ul}{\B{'ul}}) \\  
\textsc{Derivatives}: \ding{43} \on{1}.~\HL{nInAn}{\B{nì\-nän}}, less.

% ne 
\hi
\HT{ne}{\B{ne}} \I{[nE]} \emph{adp.} to, towards (\emph{direction}).   
\B{Re\-rol teng\-krr ke\-rä | Ìlä fya\-'o a\-vol | Ne kxam\-tseng.} We sing our way | Down the eight paths | To the center.~\LINK{http://www.pandorapedia.com/navi/music/ritual_music}{pp}
\B{Te\-rì\-ran ay\-oe ay\-nga\-ne.} We are walking your way.~\LINK{http://wiki.learnnavi.org/index.php/Hunting_Song}{wiki}
\B{Mìn ne ftär\-pa.} Turn to the left.~\LINK{http://forum.learnnavi.org/language-updates/ftarpa-and-skienpa-added-to-dict/msg570362/\#msg570362}{ln}
\B{Tswìl\-may\-on oe ftu Wa\-syìng\-ton ne kel\-ku.} I've just flown home from Washington.~\LINK{http://naviteri.org/2013/07/pximaw-tsawlultxa-just-after-the-meet-up/}{b}
(\emph{a. may be used to disambiguate the predicate of the verb} \ding{43}~\HL{slu}{\B{slu}}) \B{Ta\-ron\-yu slu ne tsam\-si\-yu.} The hunter becomes a warrior.~\LINK{http://wiki.learnnavi.org/Canon/2010/UltxaAyharyu\%C3\%A4\#Word_Order_with_Slu}{wiki}
(\emph{b.1 with} \ding{43}~\HL{za'u}{\B{za\-'u}}) 
\B{Zo\-la\-'u nì\-prr\-te' ne pì\-lok.} Welcome to the blog.~\LINK{http://naviteri.org/2010/06/first-post/}{b} 
(\emph{b.2 and} \ding{43}~\HL{nI'eng}{nì\-'eng}) \B{za\-'u nì\-'eng}, share an interest in common. 
\B{Tì\-ru\-sol za\-'u ne fo nì\-'eng.} They share an interest in singing.~\LINK{http://naviteri.org/2010/09/getting-to-know-you-part-1/}{b}
(\emph{b.3, an adv. and a sense organ}) `s.o. looks\slash sounds\slash tastes\slash smells good or bad. 
\B{Nga za\-'u nìl\-tsan ne na\-ri.} You look good.~\LINK{https://naviteri.org/2025/07/ayu-leketengsyon-various-things/}{b}
\B{Fì\-syu\-ve za\-'u nì\-fe' ne ftxì.} This food tastes bad.~\LINK{https://naviteri.org/2025/07/ayu-leketengsyon-various-things/}{b} 
\B{Nga za\-'u nì\-fpxa\-mo ne mik\-yun. Mok\-ri\-ri kem\-pe lo\-len?} You sound terrible. What happened to your voice?~\LINK{https://naviteri.org/2025/07/ayu-leketengsyon-various-things/}{b}
(\emph{c.}) \B{Ne kll\-te!} ``Get down!'' [lit.: to the ground] 
(\emph{d. with} \ding{43}~\HL{zup}{\B{zup}}) perish. \B{Ke za\-syup Lì\-'o\-na ne kxu\-tu.} The l'Ona will not perish to the enemy. [lit.: will not fall to the enemy]~\LINK{http://forum.learnnavi.org/language-updates/ltasygt-with-ke-and-nga/}{ln\slash g} 
(\emph{e. with} \ding{43}~\HL{kan}{\B{kan}}) aim s.t. at s.t.\slash s.o.
\B{Ney\-ti\-ril tsko\-ti ke\-ran ne ye\-rik.} Neytiri is aiming her bow at a hexapede.~\LINK{http://naviteri.org/2015/11/vomuna-liu-amip-ten-new-words/comment-page-1/\#comment-22792}{g\slash b}
(\emph{f. with} \ding{43}~\HL{lonusye}{\B{lo\-nu\-sye}}) blow on.
\B{Txo syu\-ve som li\-vu nì\-hawng, lo\-nu\-sye tsa\-ne.} If the food is too hot, blow on it.~\LINK{http://naviteri.org/2020/07/vospxivol-lefpom-happy-august/}{b}
(\ding{214}~\HL{ftu}{\B{ftu}}) \\  
\textsc{Derivatives}: \ding{43} \on{1}.~\HL{nefA}{\B{ne\-fä}}, up, upwards.   
\on{2}.~\HL{nekll}{\B{ne\-kll}}, down, below.   
\on{3}.~\HL{nemfa}{\B{nem\-fa}}, into, inside.   
\on{4}.~\HL{neto}{\B{ne\-to}}, away.

% ne'ìm
\hi
\HT{ne'Im}{\B{ne\cp{'ìm}}} \I{[nE."PIm]} \emph{adv.} back (\emph{direction}). 

% nefä
\hi
\HT{nefA}{\B{ne\cp{fä}}} \I{[nE."f\ae]} \emph{adv.} up (\emph{direction}), upwards. 
\B{Vll eyk\-yu ne\-fä fte po\-ngu fä\-ki\-vä.} The leader signals the party to ascend by pointing upwards.~\LINK{http://naviteri.org/2012/06/spring-vocabulary-part-3/}{b} 
(\ding{43}~\HL{ne}{\B{ne}}, \HL{fApa}{\B{fä\-pa}})

% Nekawn
\hi
\HT{nekawn}{\B{Nekawn}} \I{[nE.kawn]} \emph{n. female name}.~\LINK{}{dh}

% nekll
\hi
\HT{nekll}{\B{ne\cp{kll}}} \I{[nE."k\s{l}]} \emph{adv.} down (\emph{direction}).
\B{Tìng na\-ri ne\-kll.} Look below.~\LINK{http://naviteri.org/2013/08/fmawn-a-fkol-kilmulat-this-just-in/}{b}
\B{Po pxun\-fa pxek tsa\-krr me\-fo zup ne\-kll!} He shoved the two of them and they fell down.~\LINK{http://naviteri.org/2014/05/mipa-ayliu-mipa-aysafpil-new-words-new-ideas/}{b} 
(\ding{43}~\HL{ne}{\B{ne}}, \HL{kllte}{\B{kll\-te}})

% nekx
\hi
\HT{nekx}{\B{nekx}} \I{[nEk']} \emph{vtr.} burn, consume.
\B{Tsa\-srä a\-naw\-nekx ley nì\-'it nì\-'aw.} That burnt cloth is of little value.~\LINK{http://naviteri.org/2014/03/value-and-worth/}{b} 
(\ding{43}~\HL{palon}{\B{pa\-lon}}) \\  
\textsc{Derivatives}: \ding{43} \on{1}.~\HL{krrnekx}{\B{krr\-nekx}}, take time.

% nemfa
\hi
\HT{nemfa}{\B{\cp{nem}fa}} \I{["nEm.fa]} \emph{adp.} (\emph{a.}) into. 
\B{Lo\-ak\-ìl pä\-nu\-to\-lìng fu\-ta kar oe\-ru fya\-'ot a 'ìp fko nem\-fa ewll.} Loak promised he'd teach me how to vanish into the bushes.~\LINK{http://naviteri.org/2013/01/awvea-posti-zisita-amip-first-post-of-the-new-year/}{b}
(\emph{b. idiomatic use}) (i.) with \HL{fpxAkIm}{\B{fpxä\-kìm}} to mean ``born''.
\B{Fì\-tu\-te a\-hì\-'i a nem\-fa ki\-fkey fpxì\-mä\-kìm.} This little person that just entered into the world.~\LINK{http://masempul.org/2010/07/\%E2\%80\%99ite-polaheiem/}{sempul} (ii.) (for astronomical bodies rising) \B{Tsaw\-ke fpxe\-rä\-kìm nem\-fa taw.} The sun is rising.~\LINK{http://naviteri.org/2012/01/more-additions-to-the-lexicon/}{b}
(\ding{214}~\HL{ftumfa}{\B{ftum\-fa}}; \ding{43}~\HL{ne}{\B{ne}}, \HL{mIfa}{\B{mì\-fa}})

% nemrey
\hi
\HT{nemrey}{\B{nem\cp{rey}}} \I{[nEm."REj]} \emph{adv.} in a fashion as if one's life were at stake. 
\B{Kä li!--Slä\dots{}--Nem\-rey!} Get going!--But\dots{}--Run for your life!~\LINK{http://naviteri.org/2011/05/some-miscellaneous-vocabulary/}{b}
(\ding{43}~\HL{emrey}{\B{em\-rey}})

% neni
\hi
\HT{neni}{\B{\cp{ne}ni}} \I{["nE.ni]} \emph{n.} sand.
\B{Ne\-ni lew si fì\-txa\-yor vay txam\-pay nì\-wotx.} Sand covers this expanse all the way to the ocean.~\LINK{http://naviteri.org/2014/07/tsawlultxamaw-a-ayliu-post-meetup-words/}{b}
\B{Tsay\-re\-nur le\-we\-opx a sìn ne\-ni tìng na\-ri.} Look at those wave\hyp{}like patterns in the sand.~\LINK{http://naviteri.org/2017/09/zisikrr-amip-ayliu-amip-new-words-for-the-new-season/}{b} 

% neteyam
\hi 
\HT{neteyam}{\B{Ne\cp{te}yam}} \I{[nE."tE.jam]} \emph{n. male name}. 

% neto
\hi
\HT{neto}{\B{ne\cp{to}}} \I{[nE."to]} \emph{adv.} away (\emph{direction}). 
\B{Ay\-nga ne\-to ri\-vikx!} Step back! [lit.: move away you all]~\LINK{http://wiki.learnnavi.org/index.php?title=Na\%27vi_from_Avatar_Movie\#Scene_in_Hometree}{a\slash wiki}
\B{Ne\-to rä\-'ä ki\-vä!} Don't go away!~\LINK{http://wiki.learnnavi.org/index.php?title=Corpus\#Good_Morning_America}{wiki}
\B{Ftu oe ne\-to rikx!} Get away from me!~\LINK{http://naviteri.org/2011/08/new-vocabulary-clothing/comment-page-1/\#comment-994}{b}
(\ding{43}~\HL{mIso}{\B{mì\-so}})

% netrìp
\hi
\HT{netrIp}{\B{\cp{net}rìp}} \I{["nEt\textcorner.RIp\textcorner]} \emph{adv.} luckily, happily.
(\ding{43}~\HL{etrIp}{\B{et\-rìp}})

% new
\hi
\HT{new}{\B{new}} \I{[nEw]} \emph{vtr.} (\emph{a.}) want (s.t.).
\B{Nga\-tsyìp\-ìl new pe\-ut ta oe?} What does little you want from me?~\LINK{http://naviteri.org/2010/07/diminutives-conversational-expressions/}{b}
\B{Ngal new a tsa\-'ut rä\-'ä wi\-vo, ma 'e\-vi.} Don't reach for what you want, child.~\LINK{http://naviteri.org/2012/01/mipa-zisit-ayliu-amip-new-words-for-the-new-year/}{b}
(\emph{b. modal}) (\emph{either with} ‹\B{iv}› \emph{or} \B{futa} \emph{and} ‹\B{iv}›).
\B{Txo new nga ri\-vey, oe\-hu!} Come with me, if you want to live.~\LINK{http://wiki.learnnavi.org/index.php?title=Corpus\#UGO_Movie_Blog_.28Twitter_Questions.29}{wiki}
\B{Oe new nì\-txan ay\-nga\-ru fya\-wi\-vìn\-txu.} I want to guide you very much.~\LINK{http://forum.learnnavi.org/news-announcements/a-response-from-paul-frommer!/}{ln}
\B{Oe new yi\-vom tey\-lut.} I want to eat \emph{teylu}.~\LINK{http://wiki.learnnavi.org/index.php?title=Canon\#Extracts_from_various_emails}{wiki}
\B{Oe new ki\-vä.}\slash \B{Oel new fu\-ta ki\-vä.} I want to go.~\LINK{http://forum.learnnavi.org/language-updates/frommerian-email/}{ln}
\B{Oel new fu\-ta Ta\-ron\-yu ki\-vä.} I want Taronyu to go.~\LINK{http://forum.learnnavi.org/language-updates/frommerian-email/}{ln} 
\B{Fì\-pam\-tseo\-l oe\-ru ney\-kew fu\-ta srew.} This music makes me want to dance.~\LINK{https://naviteri.org/2020/12/mrra-tipangkxotsyip-five-little-discussions/}{b} \\  
\textsc{Derivatives}: \ding{43} 
\on{1}.~\HL{nulnew}{\B{nul\-new}}, prefer.   
\on{2}.~\HL{newomum}{\B{new\-omum}}, be curious.  
\on{3}.~\HL{txanew}{\B{txa\-new}}, greedy. 
\on{4}.~\HL{nInew}{\B{nì\-new}}, voluntarily, willingly, by desire.

% Newey
\hi
\HT{newey}{\B{Ne\cp{wey}}} \I{[nE."wEj]} \emph{n. female name}.
\B{Ìlä Fey\-ral, mun\-txa so\-li Ra\-lu sì Ne\-wey nì\-wan me\-srr\-am.} According to Pey\-ral, Ra\-lu and Ne\-wey were secretly married the day before yesterday.~\LINK{http://naviteri.org/2012/07/meetings-waterfalls-and-more/}{b}
\B{Fì\-sä\-'an\-lal Ne\-wey\-ä nga\-ti sley\-ka\-yu lek\-ye\-'ung!} This yearning for Newey is going to drive you crazy.~\LINK{http://naviteri.org/2013/01/awvea-posti-zisita-amip-first-post-of-the-new-year/}{b}

% newomum
\hi
\HT{newomum}{\B{new\cp{o}mum}} \I{[n$\cdot$Ew."{}o.mum]} \emph{or} \I{[.mUm]} (\textsc{rn}: \B{new\cp{o}mùm} \I{[n$\cdot$Ew."{}o.mUm]}) \emph{vin.}\on{11} be curious. 
(\ding{43}~\HL{new}{\B{new}}, \HL{omum}{\B{omum}}) \\  
\textsc{Derivatives}: \ding{43} \on{1}.~\HL{tInomum}{\B{tì\-no\-mum}}, curiosity.   
\on{2}.~\HL{lenomum}{\B{le\-no\-mum}}, curious.

% neympin
\hi
\HT{neympin}{\B{\cp{neym}pin}} \I{["nEjm.pin]} \emph{n.} the color \ding{43}~\HL{neyn}{\B{neyn}}.
(\ding{43}~\HL{'opin}{\B{'o\-pin}})

% neyn
\hi
\HT{neyn}{\B{neyn}} \I{[nEjn]} \emph{adj.} light colored, ``shades of white''. 
\B{neyn na txä\-rem}, the light color of bone.~\LINK{http://naviteri.org/2010/08/mipa-ayopin-mipa-ayliu-new-colors-new-words/}{b} 
\B{neyn na ya\-pay}, the light, nondescript color of mist or fog.~\LINK{http://naviteri.org/2010/08/mipa-ayopin-mipa-ayliu-new-colors-new-words/}{b} \\  
\textsc{Derivatives}: \ding{43} \on{1}.~\HL{neympin}{\B{neym\-pin}}, (the color) \B{neyn}.

% Neytep
\hi
\B{Neytep} \I{[nEj.tEp\textcorner]} \emph{n.} \emph{given name}.~\LINK{}{dh}

% Neytiri
\hi
\HT{neytiri}{\B{Ney\cp{ti}ri}} \I{[nEj."ti.ri]} \emph{n. female name}.
\B{Ney\-ti\-ri he\-ra\-haw.} Neytiri is sleeping.~\LINK{http://wiki.learnnavi.org/index.php?title=Corpus\#MSNBC}{wiki}
\B{Ney\-ti\-ril ye\-rik\-it to\-la\-ron.} Neytiri hunted a hexapede.~\LINK{http://wiki.learnnavi.org/index.php?title=Corpus\#MSNBC}{wiki}
\B{Ney\-ti\-ri\-ti fkol pä\-nu\-to\-lìng oe\-ru!} Neytiri was promised to me.~\LINK{http://forum.learnnavi.org/language-updates/panuting/}{ln\slash a}\textsuperscript{c}
\B{To\-ruk Mak\-to alu pi\-za\-yu Ney\-ti\-ri\-yä txan\-rar\-mo\-'a fra\-to kip ay\-ha\-pxì\-tu Oma\-ti\-ka\-yaä.} Neytiri's ancestor was the most famous Toruk Makto among the Omatikaya.~\LINK{http://naviteri.org/2012/03/spring-vocabulary-part-1/}{b}
\B{Ney\-ti\-ri\-ri lu yì tì\-sti\-yä kxayl nì\-ngay.} Neytiri is really angry.~\LINK{http://naviteri.org/2013/01/lifyari-po-peyi-how-good-is-her-navi/}{b}

% niä
\hi
\HT{niA}{\B{ni\cp{ä}}} \I{[ni"{}\ae]} \emph{vtr.} grab. 
\B{Oe mìn, to\-ruk, stum niä pol oe\-ti.} I turn, Toruk almost grabs me.~\LINK{https://wiki.learnnavi.org/Na\%27vi_from_Avatar_Movie\#Iknimaya}{a\textsuperscript{c}}

% nik
\hi
\HT{nik}{\B{nik}} \I{[nik\textcorner]} \emph{adj.} convenient (\emph{usable without much expenditure of effort}). \\  
\textsc{Derivatives}: \ding{43} \on{1}.~\HL{niktsyey}{\B{nik\-tsyey}}, food wrap.

% nikre 
\hi
\HT{nikre}{\B{\cp{nik}re}} \I{["nik\textcorner.RE]} \emph{n.} hair.
\B{vawm na nik\-re}, the dark color of Na'vi hair.~\LINK{http://naviteri.org/2010/08/mipa-ayopin-mipa-ayliu-new-colors-new-words/}{b}
\B{Ri\-ni\-ri nik\-re yä\-ngo' nì\-tut.} Rini's hair is always perfect.~\LINK{http://naviteri.org/2012/01/mipa-zisit-ayliu-amip-new-words-for-the-new-year/}{b}
\B{Nik\-re\-ri Ri\-ni\-yä 'ur fkan lor.} Rini's hair looks beautiful.~\LINK{http://naviteri.org/2012/11/renu-ayinanfyaya-the-senses-paradigm/}{b}
\B{Nik\-re\-ri Ri\-ni\-yä fkan lor.} Rini's hair is pleasant to the senses.~\LINK{http://naviteri.org/2012/11/renu-ayinanfyaya-the-senses-paradigm/}{b} 
\B{New sa'\-nok sli\-vayk nik\-ret 'eveng\-ä.} The mother wants to brush the child's hair.~\LINK{http://naviteri.org/2015/03/kap-si-ayunil-saylahe-threats-dreams-and-other-things/}{b} \\  
\textsc{Derivatives}: \ding{43} \on{1}.~\HL{nikroi}{\B{nik\-roi}}, hair ornament.

% nikroi 
\hi
\HT{nikroi}{\B{nik\cp{ro}i}} \I{[nik\textcorner."Ro.i]} \emph{n.} hair ornament. 
\B{Fì\-nik\-roi ley pìm\-txan?} How valuable is this hair ornament?~\LINK{http://naviteri.org/2014/03/value-and-worth/}{b}
\B{Kä\-sro\-lìn oel nik\-roit Pey\-ral\-ur yoa fwa po rol oer.} I loaned Peyral a hair ornament in exchange for her singing to me.~\LINK{http://naviteri.org/2014/02/barter-and-exchange/}{b}
(\ding{43}~\HL{nikre}{\B{nik\-re}}, \HL{ioi}{\B{ioi}})

% niktsyey
\hi
\HT{niktsyey}{\B{\cp{nik}tsyey}} \I{["nik\textcorner.\t{ts}jEj]} \emph{n.} food wrap (\emph{food items wrapped in edible leaves and vines}). 
(\ding{43}~\HL{nik}{\B{nik}}, \HL{tsyey}{\B{tsyey}})

% nim
\hi
\HT{nim}{\B{nim}} \I{[nim]} \emph{adj.} timid, shy.   
\B{Ye\-rik lu swi\-rä a\-nim nì\-txan.} The hexapede is a very timid creature.~\LINK{http://naviteri.org/2011/09/miscellaneous-vocabulary/}{b}
\B{Nim rä\-'ä lu!} Don't be shy!~\LINK{http://naviteri.org/2011/09/miscellaneous-vocabulary/}{b}

% ninan 
\hi
\HT{ninan}{\B{\cp{ni}nan}} \I{["ni.nan]} \emph{adv.} by reading. 
\B{Tì\-tu\-sa\-ron\-ìri fte fli\-vä, ze\-ne fko si\-vutx smar\-it ni\-nan nì\-no.} To succeed at hunting, you have to track your prey by reading (the forest) with attention to detail.~\LINK{http://naviteri.org/2011/09/\%E2\%80\%9Cby-the-way-what-are-you-reading\%E2\%80\%9D/}{b}
(\ding{43}~\HL{inan}{\B{inan}})

% Ninat
\hi
\B{Ni\cp{nat}} \I{[ni."nat\textcorner]} \emph{n. female name}.
\B{Nga lu rol\-yu a\-naw\-ri slä Ni\-nat lu pum a\-swey.} You're a talented singer but Ninat is the best (one).~\LINK{http://naviteri.org/2015/04/some-new-words-for-may-day/}{b}
\B{Ni\-nat\-ìri tì\-ru\-sol Txe\-wì\-yä lu vä\-pam.} To Ninat, Txewì's singing is noise.~\LINK{http://naviteri.org/2014/11/twenty-before-the-holidays/}{b}

% ningyen 
\hi
\HT{ningyen}{\B{\cp{ning}yen}} \I{["niN.jEn]} \emph{adv.} mysteriously, in a puzzling fashion.
\B{Oe\-yä tska\-lep 'o\-lìp ning\-yen.} My crossbow has mysteriously disappeared.~\LINK{http://naviteri.org/2013/01/awvea-posti-zisita-amip-first-post-of-the-new-year/}{b}
(\ding{43}~\HL{ingyen}{\B{ing\-yen}})

% nip
\hi
\HT{nip}{\B{nip}} \I{[nip\textcorner]} \emph{vin.} become stuck, get caught in s.t. 
\B{Ri\-ni fmar\-mi hi\-vi\-fwo slä ve\-nu no\-lip äo tskxe.} Ri\-ni tried to escape but her foot got caught under a rock.~\LINK{http://naviteri.org/2015/04/some-new-words-for-may-day/}{b}
\B{nip nìk\-lo\-nu\slash nì\-syep\slash nì\-meyp}, stuck tightly\slash stuck irremovably\slash weakly, loosely stuck.~\LINK{http://naviteri.org/2015/04/some-new-words-for-may-day/}{b} 

% nitram
\hi
\HT{nitram}{\B{nit\cp{ram}}} \I{[nit\textcorner."Ram]} \emph{adj.} happy (\emph{used with} \HL{'efu}{\B{'efu}}). 
(\emph{a. only for people})
\B{Saw\-tu\-te zo\-la\-'u a\-krr\-ta po ke 'ole\-fu nit\-ram.} Since the Skypeople came, she had not been happy.~\LINK{http://naviteri.org/2011/01/new-year-new-vocabulary/}{b}
\B{Fu\-la tsa\-yun oeng pi\-väng\-kxo ye'\-rìn oe\-ti nit\-ram sley\-ku nì\-txan.} It makes me very happy that we'll be able to chat soon.~\LINK{http://forum.learnnavi.org/language-updates/fula-short-form-of-fiul-a-and-tsal-oeti-steyki-nitram/}{ln}
\B{Fu\-la ho\-ren sì ay\-lì'\-fya\-vi a\-mip le\-sar sa\-yi me\-ngar oe\-ti nit\-ram 'ey\-ke\-fu nì\-txan.} It makes me feel happy that the rules and new expressions are useful to you both.~\LINK{http://naviteri.org/2012/06/spring-vocabulary-part-3/comment-page-1/\#comment-1515}{b}
(\emph{b. idioms}) \B{Sey\-kxel sì nit\-ram!} Congratulations!~\LINK{http://naviteri.org/2010/07/diminutives-conversational-expressions/}{b}
\B{Nit\-ram nì\-'aw!} Cheers!\slash Toast!~\LINK{http://naviteri.org/2012/07/tskxekengtsyip-a-mikyunfpi-a-little-listening-exercise/}{b}
(\ding{43}~\HL{tInitramnga'}{\B{tìnitramnga'}}) \\
\textsc{Derivatives}: \ding{43}
\on{1}.~\HL{tInitram}{\B{tì\-nit\-ram}}, happiness.

% nivi
\hi
\HT{nivi}{\B{\cp{ni}vi}} \I{["ni.vi]} \emph{n.} sleeping hammock (\emph{general term}).
\B{Tsa\-ni\-vi a\-txaw\-nu\-la nì\-fe' hu\-fwe\-mì pa\-lo\-la\-te.} That poorly constructed hammock fell apart in the wind.~\LINK{https://naviteri.org/2025/11/kaltxi-nimun-hello-again/}{b} \\  
\textsc{Derivatives}: \ding{43} \on{1}.~\HL{swaynivi}{\B{sway\-ni\-vi}}, family hammock. 
\on{2}.~\HL{snonivi}{\B{sno\-ni\-vi}}, single\hyp{}person hammock. 

% nì-
\hi
\B{nì}- \I{[nI.]} (\emph{a. prod. pref. to form adv.s from adj.s or ord. num.s}) 
\B{Ey\-wal zey\-ki\-vo ngat nì\-win.} May Eywa heal you quickly.~\LINK{http://naviteri.org/2010/07/vocabulary-update/}{b}
\B{Nì\-'aw\-ve oeng yo\-lom wu\-tsot; maw\-krr uvan si.} First we had dinner; afterwards we played.~\LINK{http://naviteri.org/2011/08/new-vocabulary-clothing/comment-page-1/\#comment-917}{b}
(\emph{b. prod. pref. on pron.s to mean}) like, as (s.o. does).
\B{Nì\-ay\-nga oe pe\-rey nì\-teng.} Like you I'm waiting too.~\LINK{http://forum.learnnavi.org/news-announcements/a-response-from-paul-frommer!/}{ln}
(\emph{c. unprod. pref. to form adv.s from other word classes}) e.g. \B{nì\-li}, in advance; \B{nì\-'en}, making an informed guess; \B{nì\-no}, in detail.

% nì'aw
\hi
\HT{nI'aw}{\B{nì\cp{'aw}}} \I{[nI."Paw]} \emph{adv.} only.
\B{Ru\-txe, ay\-nga\-ri fì\-kem fe\-yä 'e\-'a\-la to\-pur ka\-ngay sì\-yi nì\-'aw.} Please, this will only confirm their worst fears about you.~\LINK{http://forum.learnnavi.org/language-updates/two-words-in-use-eal-and-kangay-si/msg564284/\#msg564284}{ln}
\B{Soaia lu hì\-'i; nì\-'aw sem\-pul, sa'\-nok, tsmu\-kan sì Txe\-wì.} The family is small; only father, mother, a brother and Txewì.~\LINK{http://naviteri.org/2010/07/ting-mikyun-fte-tslivam-listening-comprehension-1-3/}{b}
(\ding{43}~\HL{'aw}{\B{'aw}}) \\
\textsc{Derivations}: \ding{43} 
\on{1}.~\HL{le'awnI'aw}{\B{le\-'aw\-nì\-'aw}}, unique.

% nì'awtu 
\hi
\HT{nI'awtu}{\B{nì\cp{'aw}tu}} \I{[nI."Paw.tu]} \emph{adv.} alone (\emph{as one person}). 
\B{Ke tsun tu\-te a\-'aw tsa\-tskxe\-ti a\-ku'\-up kxi\-vel\-tek nì\-'aw\-tu.} One person alone can't lift that heavy rock.~\LINK{http://naviteri.org/2012/01/mipa-zisit-ayliu-amip-new-words-for-the-new-year/}{b}
(\ding{43}~\HL{'aw}{\B{'aw}}, \HL{tute1}{\B{tu\-te}})

% nì'awve
\hi
\HT{nI'awve}{\B{nì\cp{'aw}ve}} \I{[nI."Paw.vE]} \emph{adv.} first, firstly. 
\B{Nì\-'aw\-ve oeng yo\-lom wu\-tsot; maw\-krr uvan si.} First we had dinner; afterwards we played.~\LINK{http://naviteri.org/2011/08/new-vocabulary-clothing/comment-page-1/\#comment-917}{b}
\B{Txo ni\-vew fko sä\-ro\-'a si\-vi, ze\-ne nì\-'aw\-ve ve\-nga' li\-vu.} If you want to accomplish great things, you first have to be organized.~\LINK{http://naviteri.org/2012/10/mipa-vospxi-mipa-ayliu-new-words-for-the-new-month/}{b}
(\ding{43}~\HL{'awve}{\B{'aw\-ve}})

% nì'älek
\hi 
\HT{nI'Alek}{\B{nì\cp{'ä}lek}} \I{[nI."P\ae.lEk\textcorner]} \emph{adv.} determinedly, with determination. 
(\ding{43}~\HL{'Alek}{\B{'ä\-lek}})

% nì'al
\hi 
\HT{nI'al}{\B{nì\cp{'al}}} \I{[tI."Pal]} \emph{adv.} wastefully. 
(\ding{43}~\HL{'al}{\B{'al}})

% nì'en
\hi
\HT{nI'en}{\B{nì\cp{'en}}} \I{[nI."PEn]} \emph{adv.} making an informed guess, acting on intuition.
\B{Pol po\-le\-'un fu\-ta pe\-hem si nì\-'en.} He decided what to do on a hunch.~\LINK{http://naviteri.org/2011/01/new-year-new-vocabulary/}{b}
(\ding{43}~\HL{'en}{\B{'en}})

% nì'eng
\hi
\HT{nI'eng}{\B{nì\cp{'eng}}} \I{[nI."PEN]} \emph{adv.} (\emph{a.}) levelly, equally.   
(\emph{b. idiom}) \B{za'u nì'eng}, share an interest in common (\HL{ne}{\B{ne}}, with). 
\B{Tì\-ru\-sol za\-'u ne fo nì\-'eng.} They share an interest in singing.~\LINK{http://naviteri.org/2010/09/getting-to-know-you-part-1/}{b}
(\emph{c. to express equal sharing}) \B{Fol tsngan\-it pxì\-mo\-lun\-'i nì\-'eng.} They shared the meat.~\LINK{http://naviteri.org/2012/01/more-additions-to-the-lexicon/}{b}
\B{Tsun aw\-nga kel\-ku si\-vi nì\-'eng Saw\-tu\-te\-hu mì atx\-kxe aw\-nge\-yä.} We can share our land with the Skypeople.~\LINK{http://naviteri.org/2012/01/more-additions-to-the-lexicon/}{b}
(\emph{d. to express value and worth; \emph{\B{nì'eng}} can be omitted in informal context}) \B{Ma\-sat oeyä ley nì\-ftxan na pum nge\-yä.} My breastplate is as valuable as yours.~\LINK{http://naviteri.org/2014/03/value-and-worth/}{b} 
(\emph{also monetary value on Earth}) \B{Oe\-yä el\-tu le\-fngap ley nì\-'eng na ew\-ro a\-za\-fu.} My computer is worth \on{70} euros.~\LINK{http://naviteri.org/2014/03/value-and-worth/}{b}
(\ding{43}~\HL{'engeng}{\B{'eng\-eng}})

% nì'eoio
\hi
\HT{nI'eoio}{\B{nì\cp{'e}oio}} \I{[nI."PE.o.i.o]} \emph{adv.} ceremoniously. 
\B{Aw\-ngal yom wu\-tsot nì\-fya\-'o le\-trr\-trr; fì\-txon yom nì\-'eoio.} We eat dinner in an ordinary manner; tonight we eat ceremoniously.~\LINK{http://whyisthisnight.com/addl-more.html\#na\%27vi}{tn}
(\ding{43}~\HL{'eoio}{\B{'eoio}})

% nì'eveng
\hi
\HT{nI'eveng}{\B{nì\cp{'e}veng}} \I{[nI."PE.vEN]} \emph{adv.} like a child, immaturely.
\B{Nga pll\-txe ke nì\-fyeyn\-tu ki nì\-'e\-veng.} You speak not like an adult but a child.~\LINK{http://naviteri.org/2010/07/vocabulary-update/}{b}
(\ding{43}~\HL{'eveng}{\B{'eveng}})

% nì'eyng
\hi
\HT{nI'eyng}{\B{nì\cp{'eyng}}} \I{[nI."EjN]} \emph{adv.} back, in response, in answer.
\B{Oel tel 'u\-pxa\-ret le\-Na'\-vi a krr, new oe nì\-teng\-fya pam\-rel si\-vi nì\-'eyng.} When I receive a Na'vi message, I want to write the same way in response.~\LINK{http://forum.learnnavi.org/language-updates/verb-for-write/msg175592/\#msg175592}{ln}
(\ding{43}~\HL{'eyng}{\B{'eyng}}) 

% nì'i'a
\hi
\HT{nI'i'a}{\B{nì\cp{'i'}a}} \I{[nI."PiP.a]} \emph{adv.} finally, at last. 
\B{Slä nì\-'i'a tsun oe pi\-vll\-txe san Zo\-la\-'u nì\-prr\-te' ne pì\-lok Na'\-vi\-te\-ri sìk!} But at last I can say ``Welcome to the blog Na'viteri!''~\LINK{http://naviteri.org/2010/06/first-post/}{b}
\B{Oe\-ri nì\-'i'a tsyokx zo\-slo\-lu.} My hand is finally healed.~\LINK{http://naviteri.org/2010/07/vocabulary-update/}{b}
(\ding{43}~\HL{'i'a}{\B{'i'a}})

% nì'it
\hi
\HT{nI'it}{\B{nì\cp{'it}}} \I{[nI."Pit\textcorner]} \emph{adv.} small amount, a bit. 
\B{Ay\-lì\-'u ngi\-an nì\-'it ske\-pek lu.} However the words are a bit formal.~\LINK{http://forum.learnnavi.org/language-updates/more-grammatical-tidbits/}{ln\slash a}
\B{Ru\-txe pi\-vll\-txe nì\-wok nì\-'it.} Please speak up a bit.~\LINK{http://naviteri.org/2011/09/miscellaneous-vocabulary/}{b}
\B{Tsa\-srä a\-naw\-nekx ley nì\-'it nì\-'aw.} That burnt cloth is of little value.~\LINK{http://naviteri.org/2014/03/value-and-worth/}{b}
(\ding{43}~\HL{'it}{\B{'it}})

% nì'Ìnglìsì
\hi
\HT{nI'InglIsI}{\B{nì\cp{'Ìng}lìsì}} \I{[nI."PIN.lI.sI]} \emph{adv.} English, in English.
\B{X-(ì)ri pe\-ral nì\-'Ìng\-lì\-sì?} What does X mean in English?~\LINK{http://naviteri.org/2010/09/getting-to-know-you-part-2/}{b}
\B{Tsun Txe\-wì kop pi\-vll\-txe nì\-'it nì\-'Ìng\-lì\-sì a no\-lu\-me mì num\-tseng tok\-tor Kì\-rey\-sä.} Txewì can also speak a bit of English that he has learned in doctor Grace's school.~\LINK{http://naviteri.org/2010/07/ting-mikyun-fte-tslivam-listening-comprehension-1-3/}{b}
(\ding{43}~\HL{'InglIsI}{\B{'Ìng\-lì\-sì}})

% nì'o'
\hi
\HT{nI'o'}{\B{nì\cp{'o'}}} \I{[nI."PoP]} \emph{adv.} in a manner that is very enjoyable,  `funly', enjoyably. 
\B{Ftia oel lì'\-fya\-ti le\-Na'\-vi nì\-'o' nì\-wotx!} Studying Na'vi is a ton of fun for me.~\LINK{http://naviteri.org/2010/09/getting-to-know-you-part-3/}{b}
\B{Tìng mik\-yun nì\-'o'!} Have fun listening!~\LINK{http://naviteri.org/2012/10/ting-mikyun-listen/}{b}
(\ding{43}~\HL{'o'}{\B{'o'}})

% nì'ul
\hi
\HT{nI'ul}{\B{nì\cp{'ul}}} \I{[nI."Pul]} \emph{or} \I{[."PUl]} \emph{adv.} more. (\emph{a.}) \B{Maw\-krr la\-ye\-iu oer krr nì\-'ul.} Afterwards I will have more time.~\LINK{http://forum.learnnavi.org/language-updates/trr-rrtaya/}{ln}
\B{Nga lä\-pi\-vawk nì\-'it nì\-'ul ko.} Tell me a bit more about yourself.~\LINK{http://naviteri.org/2010/09/getting-to-know-you-part-2/}{b}
\B{Nì\-'ul ye'\-rìn …} More soon …~\LINK{http://naviteri.org/2011/11/\%E2\%80\%99a\%E2\%80\%99awa-ayli\%E2\%80\%99u-amip-ni\%E2\%80\%99aw-\%E2\%80\%94-a-few-new-words-only/}{b}
(\emph{b. to form comparatives}) \B{Oe nga\-to yom nì\-'ul.} I eat more than you.~\LINK{http://forum.learnnavi.org/language-updates/prefer-x-to-y-24474/}{ln}
\B{Tsa\-sä\-lä\-txayn Na'\-vi\-ru srung so\-li nì\-'aw fte sli\-vu txur nì\-'ul.} That defeat only helped the People become stronger.~\LINK{http://naviteri.org/2012/01/more-additions-to-the-lexicon/}{b}
(\ding{43}~\HL{'ul}{\B{'ul}}) \\  
\textsc{Derivatives}: \ding{43} \on{1}.~\HL{nI'ul'ul}{\B{nì\-'ul\-'ul}}, more and more.

% nì'ul'ul
\hi
\HT{nI'ul'ul}{\B{nì\cp{'ul}'ul}} \I{[nI."Pul.Pul]} \emph{or} \I{[."PUl.PUl]} \emph{adv.} increasingly, more and more.
\B{Fra\-zì\-sì\-krr pay kil\-van\-ä nän nì\-'ul\-'ul.} Every season the river dries up more and more.~\LINK{http://naviteri.org/2012/02/trr-asawnung-lefpom-happy-leap-day/}{b}
\B{Sìn yì sngä\-'i\-yu\-ä, slä tsye\-rìl (hay\-a yì\-ne) nì\-'ul\-'ul.} They're at a beginner's level, but they're getting better and better.~\LINK{http://naviteri.org/2013/01/lifyari-po-peyi-how-good-is-her-navi/}{b}
(\ding{43}~\HL{nI'ul}{\B{nì\-'ul}})

% nìawnomum
\hi
\HT{nIawnomum}{\B{nìaw\cp{no}mum}} \I{[nI.aw."no.mum]} \emph{or} \I{[.mUm]}, \emph{in fast speech often} \I{[naw."no.mum]} (\textsc{rn}: \B{nìaw\cp{no}mùm} \I{[nI.aw."no.mUm]}) \emph{adv.} as you know, as is known, as everyone knows.
\B{Slä nì\-aw\-no\-mum, ze\-ne oe 'aw\-si\-teng tì\-kang\-kem si\-vi fo\-hu.} But as you know, I must work together with them.~\LINK{http://forum.learnnavi.org/news-announcements/a-response-from-paul-frommer!/}{ln}
\B{Nga nì\-aw\-no\-mum to oe\-tsyìp lu txur nì\-txan.} As everyone knows, you're a lot stronger than little old me.~\LINK{http://naviteri.org/2010/07/diminutives-conversational-expressions/}{b}
(\ding{43}~\HL{awn}{‹\B{awn}›}, \HL{omum}{\B{omum}})

% nìayoeng
\hi
\HT{nIayoeng}{\B{nìay\cp{oeng}}} \I{[naj."wEN]} \emph{adv.} like us, as we do. 
\B{Ma 'ite, aw\-nge\-yä fya\-'o\-ri ze\-ne nga sä\-nume si\-vi po\-ru fte tsi\-vun pil\-vll\-txe sì ti\-vì\-ran nì\-ay\-oeng.} My daughter, you will teach him our way to speak and walk as we do.~\LINK{http://wiki.learnnavi.org/index.php?title=Na\%27vi_from_Avatar_Movie\#Scene_in_Hometree}{a\slash wiki}
(\ding{43}~\HL{ayoeng}{\B{ay\-oeng}})

% nìfe'
\hi
\HT{nIfe'}{\B{nì\cp{fe'}}} \I{[nI."fEP]} \emph{adv.} badly.
\B{Oe pll\-ngay san mo\-lak\-to oe nì\-fe', ta\-fral sno\-laytx; wä\-tu lu oe\-to txur.} I acknowledge that I rode badly, so I lost; my opponent was stronger than I was.~\LINK{http://naviteri.org/2011/07/txantsana-ultxa-mi-siatll-great-meeting-in-seattle/}{b}
(\ding{43}~\HL{fe'}{\B{fe'}})

% nìfkeytongay
\hi
\HT{nIfkeytongay}{\B{nìfkeyto\cp{ngay}}} \I{[nI.fkEj.to."Naj]} \emph{adv.} actually, as a matter of fact.
\B{Tsat ngal tswo\-lä\-nga' nì\-lam. -- Nì\-fkey\-to\-ngay ke tswo\-la' kaw\-'it.} Apparently you've forgotten it. -- As a matter of fact, I haven't.~\LINK{http://naviteri.org/2014/06/teri-tsaliu-alu-nifkeytongay-about-nifkeytongay/}{b}
(\ding{43}~\HL{tIfkeytongay}{\B{tì\-fkey\-to\-ngay}})

% nìfkrr
\hi
\HT{nIfkrr}{\B{nì\cp{fkrr}}} \I{[nI."fk\s{r}]} \emph{adv.} lately. 
\B{'Ìn nga fya\-pe nì\-fkrr?} What's been keeping you busy lately?~\LINK{http://naviteri.org/2010/09/getting-to-know-you-part-3/}{b}
(\ding{43}~\HL{krr}{\B{krr}})

% nìflä
\hi 
\HT{nIflA}{\B{nì\cp{flä}}} \I{[nI."fl\ae]} \emph{adv.} successfully. 
\B{So\-leia! Ngal tì\-fme\-tok\-it em\-zo\-la\-'u nì\-flä! Sey\-kxel sì nit\-ram!} You rose to the challenge! You passed the test successfully! Congratulations!~\LINK{http://naviteri.org/2019/06/50a-liu-amip-40-new-words/}{b}
(\ding{43}~\HL{flA}{\B{flä}})

% nìflrr
\hi
\HT{nIflrr}{\B{nì\cp{flrr}}} \I{[nI."fl\s{r}]} \emph{adv.} gently, tenderly.
\B{Ze\-ne fko 'i\-vam\-pi prr\-nen\-it nì\-flrr fra\-krr.} One must always touch a baby gently.~\LINK{http://naviteri.org/2013/02/vospxi-ayol-posti-apup-short-post-for-a-short-month/}{b}
(\ding{43}~\HL{flrr}{\B{flrr}})

% nìfmokx
\hi
\HT{nIfmokx}{\B{nì\cp{fmokx}}} \I{[nI."fmok']} \emph{adv.} jealously, enviously.   
\B{Txe\-wìl tuk\-rut Lo\-ak\-ä nar\-mìn nì\-fmokx.} Txe\-wì was eyeing Lo\-ak's spear enviously.~\LINK{http://naviteri.org/2011/10/more-vocabulary-a-bit-of-grammar/}{b}
(\ding{43}~\HL{fmokx}{\B{fmokx}})

% nìfnu
\hi
\HT{nIfnu}{\B{nì\cp{fnu}}} \I{[nI."fnu]} \emph{adv.} silently.
\B{Tì\-ran nì\-fnu txan\-txew\-vay fte\-ke ay\-ye\-rik\-ìl aw\-nga\-ti sti\-vawm.} Walk as quietly as possible so the hexapedes won't hear us.~\LINK{http://naviteri.org/2012/10/snewsyea-ftxoza-halowina-a-spooky-halloween/}{b}
(\ding{43}~\HL{fnu}{\B{fnu}})

% nìfpxamo
\hi 
\HT{nIfpxamo}{\B{nì\cp{fpxa}mo}} \I{[nI."fp'a.mo]} \emph{adv.} horribly, terribly, awfully. 
\B{Fpä\-ngìl oe, txo\-nam oe ro\-lol nì\-fpxa\-mo.} Sadly, I think I sang terribly last night.~\LINK{http://naviteri.org/2018/08/fmawnti-stolawm-srak-have-you-heard-the-news/}{b} 
(\ding{43}~\HL{fpxamo}{\B{fpxa\-mo}})

% nìfrakrr
\hi
\HT{nIfrakrr}{\B{nì\cp{fra}krr}} \I{[nI."fRa.k\s{r}]} \emph{adv.} as always.
\B{Nì\-fra\-krr fol 'o\-lem a wu\-tso ftxì\-vä' lu nì\-ngay.} As always, the dinner they cooked tasted really terrible.~\LINK{http://naviteri.org/2011/05/weather-part-2-and-a-bit-more-2/}{b}
(\ding{43}~\HL{frakrr}{\B{fra\-krr}})

% nìFranse
\hi
\HT{nIfranse}{\B{nì\cp{Fran}se}} \I{[nI."fRan.sE]} \emph{adv.} French, in French.
\B{X-(ì)ri pe\-ral nì\-Fran\-se?} What does X mean in French?~\LINK{http://naviteri.org/2010/09/getting-to-know-you-part-2/}{b}

% nìfrir
\hi 
\HT{nIfrir}{\B{nì\cp{frir}}} \I{[nI."fRiR]} \emph{adv.} in layers. 
(\ding{43}~\HL{frir}{\B{frir}})

% nìftär
\hi
\HT{nIftAr}{\B{nì\cp{ftär}}} \I{[nE."ft\ae{}R]} \emph{adv.} to the left.
\B{Ne 'o\-ra\-tsyìp po\-lä\-hem, yak si nì\-ftär.} When you arrive at the pond, turn to the left.~\LINK{http://naviteri.org/2013/09/on-si-salewfya-shapes-and-directions/}{b}
(\ding{43}~\HL{ftAr}{\B{ftär}})

% nìftue
\hi
\HT{nIftue}{\B{nì\cp{ftu}e}} \I{[nI."ftu.E]} \emph{adv.} easily. 
\B{Sra\-ke tsun oe fay\-upxa\-ret tsli\-vam nì\-ftue?} Can I understand these messages easily?~\LINK{http://forum.learnnavi.org/language-updates/verb-for-write/msg175592/\#msg175592}{ln}
\B{Fay\-upxa\-re la\-yu ay\-sngä\-'i\-yu\-fpi, fte lì'\-fya\-ri aw\-nge\-yä fo tsì\-ye\-vun nì\-ftue nìl\-tsan\-sì ni\-vume.} These messages will be for beginners, so that they can learn about our language easily and well.~\LINK{http://naviteri.org/2010/06/first-post/}{b} 
(\ding{43}~\HL{ftue}{\B{ftue}})

% nìftxan
\hi
\HT{nIftxan}{\B{nì\cp{ftxan}}} \I{[nI."ft'an]} \emph{adv.} (\emph{a.}) so, to such an extent. 
\B{Fwa ay\-ngal fì\-tse\-nget tok nì\-pxay nì\-ftxan oe\-ru te\-ya si nì\-ngay.} It fills me with joy that you're here in such great number.~\LINK{http://naviteri.org/2012/07/tskxekengtsyip-a-mikyunfpi-a-little-listening-exercise/}{b}
(\emph{b. with adj. or adv. and} \HL{na}{\B{na}}\slash \HL{pxel}{\B{pxel}} \emph{to mean}) as \emph{adj.\slash adv.} as.
\B{Oe lu nì\-ftxan sìl\-tsan na nga.} I am as good as you.~\LINK{http://forum.learnnavi.org/language-updates/as-x-as-y-keep-on-keepin-on/}{ln} 
(\emph{A pronoun or definite point of comparison may take the topical}) 
\B{Nga\-ri lu oe nì\-ftxan sìl\-tsan.} I am as good as you.~\LINK{http://forum.learnnavi.org/language-updates/as-x-as-y-keep-on-keepin-on/}{ln}
\B{Oe tul nì\-ftxan nì\-win pxel nga.} I run as fast as you.~\LINK{https://forum.learnnavi.org/language-updates/eapressions-of-\%27like-we\%27/}{ln}
\B{Oel ye\-rik\-it ta\-ron nì\-ftxan nìl\-tsan na nga.} I hunt yerik as well as you.~\LINK{http://forum.learnnavi.org/language-updates/confirmations-\%28comparisons-gerund\%29/}{ln}
(\emph{c. to form result clauses with} \HL{kuma}{\B{kuma}}) so.
\B{Lu poe se\-vin nì\-ftxan kuma yaw\-ne slo\-lu oer.} She was so beautiful that I fell in love with her.~\LINK{http://naviteri.org/2012/06/spring-vocabulary-part-3/}{b}
(\emph{d. coll. idiom}) \B{[verb] nì\-ftxan kum\-a ter\-kup}, \emph{commenting on an action that is so intense that it results in dying}; ``like crazy, \ldots{} to death''. 
\B{Sran, sran, oe\-ri pì\-tìk tsa\-tse\-nget a mì tal! Fkxa\-ke nì\-ftxan kum\-a ter\-kup.} Yeah, yeah, scratch that place on my back! It itches like crazy!~\LINK{http://naviteri.org/2020/07/vospxivol-lefpom-happy-august/}{b}
(\ding{43}~\HL{fItxan}{\B{fì\-txan}})

% nìftxavang
\hi
\HT{nIftxavang}{\B{nì\cp{ftxa}vang}} \I{[nI."ft'a.vaN]} \emph{adv.} passionately, with all heart.
\B{Fpo\-le' ay\-ngal oer fì\-txan nì\-ftxa\-vang a 'upxa\-ret sto\-lawm oel.} I have heard the message you have sent me so passionately.~\LINK{http://forum.learnnavi.org/news-announcements/a-response-from-paul-frommer!/}{ln}
\B{Po\-ri mok\-ri lor lu nì\-txan ul\-te rol po nì\-ftxa\-vang.} Her voice is very beautiful and she sings passionately.~\LINK{http://naviteri.org/2014/06/teri-tsaliu-alu-nifkeytongay-about-nifkeytongay/}{b}
(\ding{43}~\HL{ftxavang}{\B{ftxa\-vang}})

% nìfwefwi
\hi
\HT{nIfwefwi}{\B{nì\cp{fwe}fwi}} \I{[nI."fwE.fwi]} \emph{adv.} by whistling, in a whistling manner.
\B{'E\-van\-ìl alo a\-'aw\-ve nì\-'aw\-tu na'\-rìng\-it tar\-mok, ha to\-lìng lawr nì\-fwe\-fwi fte\-ke txo\-pu si\-vi.} The boy was alone in the forest for the first time, so he whistled a tune to calm his fears.~\LINK{http://naviteri.org/2011/09/miscellaneous-vocabulary/}{b}
(\ding{43}~\HL{fwefwi}{\B{fwe\-fwi}})

% nìfya'o
\hi
\HT{nIfya'o}{\B{nì\cp{fya}'o}} \I{[nI."fja.Po]} \emph{adv.} in a manner. 
\B{Aw\-ngal yom wu\-tsot nì\-fya\-'o le\-trr\-trr; fì\-txon yom nì\-'eoio.} We eat dinner in an ordinary manner; tonight we eat ceremoniously.~\LINK{http://whyisthisnight.com/addl-more.html\#na\%27vi}{tn}
\B{Poe pol\-txe nì\-fya'o a\-law.} She spoke clearly.~\LINK{http://forum.learnnavi.org/language-updates/why-is-this-night/msg154875/\#msg154875}{ln}
\B{Fo nì\-Na'\-vi pll\-txe nì\-fya\-'o a hek.} They speak Na'vi strangely.~\LINK{http://naviteri.org/2010/07/vocabulary-update/}{b}
(\ding{43}~\HL{fya'o}{\B{fya\-'o}})

% nìfyeyntu
\hi
\HT{nIfyeyntu}{\B{nì\cp{fyeyn}tu}} \I{[nI."fjEjn.tu]} \emph{adv.} like an adult, maturely. 
\B{Nga pll\-txe ke nì\-fyeyn\-tu ki nì\-'e\-veng.} You speak not like an adult but a child.~\LINK{http://naviteri.org/2010/07/vocabulary-update/}{b}
(\ding{43}~\HL{fyeyntu}{\B{fyeyn\-tu}})

% nìhawmpam
\hi
\HT{nIhawmpam}{\B{nì\cp{hawm}pam}} \I{[nI."hawm.pam]} \emph{adj.} noisily. 
\B{Txo fko ti\-vul mì na'\-rìng nì\-hawm\-pam, stawm ay\-ioang.} If you run noisily in the forest, the animals will hear.~\LINK{http://naviteri.org/2014/11/twenty-before-the-holidays/}{b}
(\ding{43}~\HL{hawmpam}{\B{hawm\-pam}})

% nìhawng
\hi
\HT{nIhawng}{\B{nì\cp{hawng}}} \I{[nI."hawN]} \emph{adv.} too, excessively. 
\B{Na\-ri so\-li ay\-oe fte\-ke nì\-hawng li\-vok.} We were careful not to get too close.~\LINK{http://wiki.learnnavi.org/index.php?title=Corpus\#New_York_Times}{wiki}
\B{Ay\-upxa\-re\-ri a\-ngim nì\-hawng lu alo oe\-yä!} Now it's my turn for a message that's too long!~\LINK{http://forum.learnnavi.org/language-updates/txelanit-hivawl/}{ln}
\B{Fì\-na\-er ftxì\-vä' lu nì\-hawng, ha swey\-lu txo ngal tsat kxi\-vukx nì\-win.} This drink tastes horrible, so you'd better swallow it quickly.~\LINK{http://naviteri.org/2011/11/\%E2\%80\%99a\%E2\%80\%99awa-ayli\%E2\%80\%99u-amip-ni\%E2\%80\%99aw-\%E2\%80\%94-a-few-new-words-only/}{b}
(\ding{43}~\HL{hawng}{\B{hawng}})

% nìhay
\hi
\HT{nIhay}{\B{nì\cp{hay}}} \I{[nI."haj]} \emph{adv.} next (in line). 
\B{Nga kä pe\-seng nì\-hay?} Where are you going next?~\LINK{http://naviteri.org/2012/10/audio-and-video-learning-materials-for-navi-101/}{b}
(\ding{43}~\HL{hay}{\B{hay}})

% nìhek
\hi
\HT{nIhek}{\B{nì\cp{hek}}} \I{[nI."hEk\textcorner]} \emph{adv.} oddly, strangely   (\emph{sentence adverbial only}). 
\B{Nì\-hek fo nì\-Na'\-vi pll\-txe.} Strangely, they speak Na'\-vi.~\LINK{http://naviteri.org/2010/07/vocabulary-update/}{b}
(\ding{43}~\HL{hek}{\B{hek}})

% nìhìpey
\hi
\HT{nIhIpey}{\B{nì\cp{hì}pey}} \I{[nI."hI.pEj]} \emph{adv.} hesitantly.
(\ding{43}~\HL{hIpey}{\B{hì\-pey}})

% nìhoan
\hi
\HT{nIhoan}{\B{nì\cp{ho}an}} \I{[nI."ho.an]} \emph{adv.} comfortably.
(\ding{43}~\HL{hoan}{\B{ho\-an}})

% nìhoet
\hi
\HT{nIhoet}{\B{nì\cp{ho}et}} \I{[nI."ho.Et\textcorner]} \emph{adv.} widely, pervasively.   
\B{Run fkol tey\-lut nì\-ho\-et.} You find \emph{teylu} everywhere.~\LINK{http://naviteri.org/2012/06/spring-vocabulary-part-3/}{b}
(\ding{43}~\HL{hoet}{\B{ho\-et}})

% nìhol
\hi
\HT{nIhol}{\B{nì\cp{hol}}} \I{[nI."hol]} \emph{adv.} lightly, just a few.
\B{lep\-wopx nì\-hol}, lightly cloudy, just a few clouds.~\LINK{http://naviteri.org/2011/05/weather-part-2-and-a-bit-more-2/}{b}
(\ding{43}~\HL{hol}{\B{hol}})

% nìhona
\hi
\HT{nIhona}{\B{nì\cp{hon}a}} \I{[nI."ho.na]} \emph{adv.} endearingly, sweetly.
\B{Po ätxä\-le so\-li nì\-ho\-na fì\-txan, ke tsun oe sti\-vo.} She asked so sweetly that I couldn't refuse.~\LINK{http://naviteri.org/2012/01/more-additions-to-the-lexicon/}{b}
(\ding{43}~\HL{hona}{\B{ho\-na}})

% nìk'ärìp
\hi
\HT{nIk'ArIp}{\B{nìk\cp{'ä}rìp}} \I{[nIk\textcorner."P\ae.RIp\textcorner]} \emph{adv.} steadily (\emph{without letting it move}).
\B{fyep nìk\-'ä\-rìp}, hold steadily.~\LINK{http://naviteri.org/2012/03/spring-vocabulary-part-1/}{b}
(\ding{43}~\HL{'ArIp}{\B{'ä\-rìp}})

% nìk'ong
\hi
\HT{nIk'ong}{\B{nìk\cp{'ong}}} \I{[nIk\textcorner."PoN]} \emph{adv.} slowly. 
(\emph{a. idiom}) \B{Ke ze\-ne win sä\-pi\-vi. 'Ivong nìk\-'ong.} Take your time; don't rush. Slow is fine.~\LINK{http://naviteri.org/2010/09/quick-follow-up/}{b}
(\ding{43}~\HL{kI'ong}{\B{kì\-'ong}})

% nìkakan
\hi 
\HT{nIkakan}{\B{nì\cp{ka}kan}} \I{[nI."ka.kan]} \emph{adv.} roughly (\emph{for behavior, people and things}). 
\B{Sem\-pul\-ìl a\-sìl\-tsan sney eveng\-it ke txaw nì\-ka\-kan.} A good father doesn't punish his children roughly.~\LINK{http://naviteri.org/2021/09/aawa-ayliu-amip-a-few-new-words/}{b} 
(\ding{43}~\HL{kakan}{\B{ka\-kan}})

% nìkakmokri
\hi 
\HT{nIkakmokri}{\B{nìkak\cp{mok}ri}} \I{[nI.kak\textcorner."mok\textcorner.ri]} \emph{adv.} mutely. 
\B{Kll\-kxo\-lem fo nì\-kak\-mok\-ri lu\-ke fwa 'aw\-a lì\-'uti pll\-txe.} They stood there mutely without saying a word.~\LINK{https://naviteri.org/2024/02/mipa-ayliu-si-aylifyavi-niul-more-new-words-and-expressions/}{b} 
(\ding{43}~\HL{kakmokri}{\B{kak\-mok\-ri}}) 

% nìkanu
\hi
\HT{nIkanu}{\B{nì\cp{ka}nu}} \I{[nI."ka.nu]} \emph{adv.} intelligently, in a smart way. 
(\ding{43}~\HL{kanu}{\B{ka\-nu}})

% nìkeftxo
\hi
\HT{nIkeftxo}{\B{nìke\cp{ftxo}}} \I{[nI.kE."ft'o]} \emph{adv.} unfortunately, sadly. 
\B{Nì\-ke\-ftxo lu fì\-'e\-veng\-ur 'aw\-a sa'\-sem nì\-'aw.} Unfortunately this child has only one parent.~\LINK{http://naviteri.org/2011/08/new-vocabulary-clothing/comment-page-1/\#comment-917}{b}
\B{Nì\-ke\-ftxo oe\-ri ke tam.} Unfortunately, it doesn't work for me.~\LINK{http://naviteri.org/2011/09/\%E2\%80\%9Cby-the-way-what-are-you-reading\%E2\%80\%9D/comment-page-1/\#comment-1116}{b}
(\ding{43}~\HL{keftxo}{\B{ke\-ftxo}})

% nìkelkin
\hi
\HT{nIkelkin}{\B{nìkel\cp{kin}}} \I{[nI.kEl."kin]} \emph{adv.} unnecessarily.
\B{Sna\-fpìl\-fya\-ri le\-Na'\-vi krra smar\-it fkol tspang, tsran\-ten nì\-txan fwa po ke ngä\-'än nì\-kel\-kin.} It's important in Na'vi philosophy that the prey not suffer unnecessarily when it's killed.~\LINK{http://naviteri.org/2013/02/vospxi-ayol-posti-apup-short-post-for-a-short-month/}{b}
(\ding{43}~\HL{kelkin}{\B{kel\-kin}})

% nìkemweypey
\hi
\HT{nIkemweypey}{\B{nìkem\cp{wey}pey}} \I{[nI.kEm."wEj.pEj]} \emph{adv.} impatiently.
(\ding{43}~\HL{maweypey}{\B{ma\-wey\-pey}}, \HL{ke}{\B{ke}})

% nìklonu
\hi
\HT{nIklonu}{\B{nìklo\cp{nu}}} \I{[nIk\textcorner.lo."nu]} \emph{adv.} firmly, steadfastly, faithfully (\emph{without releasing it}).
\B{fyep nìk\-lo\-nu}, hold firmly.~\LINK{http://naviteri.org/2012/03/spring-vocabulary-part-1/}{b}
\B{nip nìk\-lo\-nu}, stuck tightly.~\LINK{http://naviteri.org/2015/04/some-new-words-for-may-day/}{b} 
(\ding{43}~\HL{lonu}{\B{lo\-nu}})

% nìkmar
\hi
\HT{nIkmar}{\B{nìk\cp{mar}}} \I{[nIk\textcorner."maR]} \emph{adv.} in the right season, opportunely.
\B{Po tsap\-'a\-lu\-te so\-li nìk\-mar.} He apologized at the right time.~\LINK{http://naviteri.org/2011/10/more-vocabulary-a-bit-of-grammar/}{b}
\B{Nìk\-mar zup tom\-pa set.} Now is a good time for it to be raining.~\LINK{http://naviteri.org/2011/10/more-vocabulary-a-bit-of-grammar/}{b}
\B{Fo ke pe\-räng\-kxo oe\-hu nìk\-mar.} They were chatting with me at a bad time.~\LINK{http://naviteri.org/2011/10/more-vocabulary-a-bit-of-grammar/}{b}
(\ding{43}~\HL{kImar}{\B{kì\-mar}}) 

% nìksman
\hi
\HT{nIksman}{\B{nìk\cp{sman}}} \I{[nIk\textcorner."sman]} \emph{adv.} wonderfully.
\B{Fol lì'\-fya\-ti aw\-nge\-yä sar nìk\-sman.} They use our language wonderfully.~\LINK{http://naviteri.org/2011/09/miscellaneous-vocabulary/}{b}
(\ding{43}~\HL{kosman}{\B{ko\-sman}})

% nìksran
\hi
\HT{nIksran}{\B{nìk\cp{sran}}} \I{[nIk\textcorner."sRan]} \emph{adv.} in a mediocre manner.
\B{Pll\-txe nìk\-sran, slä tsun fko pe\-yä ay\-lì\-'ut tsli\-vam.} His speech is mediocre, but you can understand him.~\LINK{http://naviteri.org/2013/01/lifyari-po-peyi-how-good-is-her-navi/}{b}
(\ding{43}~\HL{kesran}{\B{ke\-sran}})

% nìktungzup
\hi
\HT{nIktungzup}{\B{nìktung\cp{zup}}} \I{[nIk\textcorner.tuN."zup\textcorner]} \emph{or} \I{[.tUN."zUp\textcorner]} (\textsc{rn}: \B{nìktùng\cp{zup}} \I{[nIk\textcorner.tUN."zup\textcorner]}) \emph{adv.} carefully, firmly (\emph{without letting it fall}).
\B{fyep nìk\-tung\-zup}, hold firmly.~\LINK{http://naviteri.org/2012/03/spring-vocabulary-part-1/}{b}
(\ding{43}~\HL{tungzup}{\B{tung\-zup}})

% nìkx
\hi
\HT{nIkx}{\B{nìkx}} \I{[nIk']} \emph{n.} gravel.

% nìkxap
\hi
\HT{nIkxap}{\B{nì\cp{kxap}}} \I{[nI."k'ap\textcorner]} \emph{adv.} threateningly.
\B{Txo\-pu rä\-'ä si. Nga\-ti ke nìn pol nì\-kxap. Lu le\-no\-mum nì\-'aw.} Don't be scared. He's not looking at you threateningly. He's just curious.~\LINK{http://naviteri.org/2015/03/kap-si-ayunil-saylahe-threats-dreams-and-other-things/}{b}
(\ding{43}~\HL{kxap}{\B{kxap}})

% nìkxäl
\hi 
\HT{nIkxAl}{\B{nì\cp{kxäl}}} \I{[nI."k'\ae{}l]} \emph{adv.} apathetically. 
(\ding{43}~\HL{kxAl}{\B{kxäl}}) 

% nìkxem
\hi
\HT{nIkxem}{\B{nì\cp{kxem}}} \I{[nI."k'Em]} \emph{adv.} vertically.  
(\ding{43}~\HL{kxem}{\B{kxem}})

% nìlam
\hi
\HT{nIlam}{\B{nì\cp{lam}}} \I{[nI."lam]} \emph{adv.} apparently.
\B{Po\-an yaw\-ne la\-tsu poe\-ru nì\-lam.} Apparently she loves him.~\LINK{http://naviteri.org/2011/05/weather-part-2-and-a-bit-more-2/}{b}
\B{Tsat ngal tswo\-lä\-nga' nì\-lam.} Apparently you've forgotten it.~\LINK{http://naviteri.org/2014/06/teri-tsaliu-alu-nifkeytongay-about-nifkeytongay/}{b}
(\ding{43}~\HL{lam}{\B{lam}})

% nìlaw
\hi
\HT{nIlaw}{\B{nì\cp{law}}} \I{[nI."law]} \emph{adv.} clearly (\emph{sentence and manner adverbial}). 
\B{Poe pol\-txe nì\-law.} She spoke clearly.\slash Clearly, she spoke.~\LINK{http://forum.learnnavi.org/language-updates/why-is-this-night/msg154875/\#msg154875}{ln}
\B{Nì\-law po yaw\-ne lu ngar.} It's clear you love her.~\LINK{http://naviteri.org/2012/07/fivospxiya-aylifyavi-amip-this-months-new-expressions-pt-2/}{b}
(\ding{43}~\HL{law}{\B{law}})

% nìlayl
\hi 
\HT{nIlayl}{\B{nì\cp{layl}}} \I{[nI."lajl]} \emph{adv.} innocently. 
\B{zaw\-prr\-te' nì\-layl fko\-ne}, be pleasurable to one in an innocent way. 
\B{Tì\-mun\-txa me\-fe\-yä zo\-law\-prr\-te' Ma\-ru\-ne nì\-layl. Ke lu por kea fmokx kaw\-'it.} Their marriage brought Ma\-ru pleasure. She felt no jealousy at all.~\LINK{http://naviteri.org/2023/09/aawa-ayliu-si-aylifyavi-amip-a-few-new-words-and-expressions/}{b}
(\ding{43}~\HL{layl}{\B{layl}}; \ding{214}~\HL{nItokat}{\B{nì\-to\-kat}}) 

% nìleyn
\hi 
\HT{nIleyn}{\B{nì\cp{leyn}}} \I{[nI."lEjn]} \emph{adv.} repeatedly. 
(\ding{43}~\HL{leyn}{\B{leyn}}) 

% nìler
\hi 
\HT{nIler}{\B{nì\cp{ler}}} \I{[nI."lER]} \emph{adv.} (\emph{a. of motion}) steadily. 
\B{Nì\-sngä\-'i ke tsun Tsyeyk tswi\-vay\-on nì\-ler.} At first, Jake couldn't fly steadily.~\LINK{http://naviteri.org/2016/06/mrrvola-lifyavi-amip-forty-new-expressions/}{b}
(\emph{b. metaphorically of smooth, unbroken action}) 
\B{tì\-kang\-kem si nì\-ler}, to work steadily (without stopping).~\LINK{http://naviteri.org/2016/06/mrrvola-lifyavi-amip-forty-new-expressions/}{b} 
(\ding{43}~\HL{ler}{\B{ler}})

% nìli
\hi
\HT{nIli}{\B{nì\cp{li}}} \I{[nI."li]} \emph{adv.} in advance. 
\B{Nge\-yä stxe\-li\-ri a\-lor oe new nga\-ru pi\-vll\-txe san irayo sìk nì\-li.} I want to thank you in advance for your beautiful gift.~\LINK{http://naviteri.org/2011/05/weather-part-2-and-a-bit-more-2/}{b}
(\ding{43}~\HL{li}{\B{li}})

% nìlkeftang
\hi
\HT{nIlkeftang}{\B{nìlke\cp{ftang}}} \I{[nIl.kE."ftaN]} \emph{adv.} continuously, without stopping. 
\B{Nga fwe\-fwi nìl\-ke\-ftang a fì\-'u lu sä\-srätx a\-txan oe\-ru.} It's a major annoyance to me that you whistle continuously.~\LINK{http://naviteri.org/2011/09/miscellaneous-vocabulary/}{b}
(\ding{43}~\HL{ftang}{\B{ftang}}, \HL{luke}{\B{lu\-ke}})

% nìloho
\hi 
\HT{nIloho}{\B{nì\cp{lo}ho}} \I{[nI."lo.ho]} \emph{adv.} surprisingly. 
\B{Pol\-txe po nì\-lo\-ho san oe za\-sya\-'u.} Surprisingly, he said he would come.~\LINK{http://naviteri.org/2019/06/50a-liu-amip-40-new-words/}{b} 
(\ding{43}~\HL{loho}{\B{lo\-ho}})

% nìltsan
\hi
\HT{nIltsan}{\B{nìl\cp{tsan}}} \I{[nIl."\t{ts}an]} \emph{adv.} well.
\B{Pll\-txe nga nìl\-tsan!} You speak well!~\LINK{http://languagelog.ldc.upenn.edu/nll/?p=1977}{ll}
\B{Tì\-ste\-ftxaw nge\-yä flo\-lä nìl\-tsan!} Your examination succeeded well!~\LINK{http://wiki.learnnavi.org/index.php?title=Canon\#Extracts_from_various_emails}{wiki}
\B{Ay\-lì\-'u a\-paw\-nll\-txe nìl\-tsan!} Words well spoken!~\LINK{http://forum.learnnavi.org/language-updates/a-new-participle-infix-awn}{ln}
\B{Pol ye\-rik\-it ta\-ron nìl\-tsan fra\-to.} He hunts yerik better than anyone.~\LINK{http://forum.learnnavi.org/language-updates/confirmations-\%28comparisons-gerund\%29/}{ln}
(i.) \B{Sey\-so\-nìl\-\cp{tsan}!} \I{[sEj.so.nIl."\t{ts}an]} Well done! (\emph{from} \B{ha\-sey so\-li nìl\-tsan})~\LINK{http://naviteri.org/2011/04/\%E2\%80\%99a\%E2\%80\%99awa-li\%E2\%80\%99fyavi-amip\%E2\%80\%94a-few-new-expressions/}{b}
(\ding{43}~\HL{sIltsan}{\B{sìl\-tsan}})

% nìlun
\hi
\HT{nIlun}{\B{nì\cp{lun}}} \I{[nI."lun]} \emph{or} \I{[."lUn]} \emph{adv.} of course, logically, following common sense. 
\B{Nì\-lun ay\-ioi a\-'eoio ay\-eyk\-tan\-ä lu lor fra\-to.} Of course the ceremonial wardrobes of the leaders are the most beautiful.~\LINK{http://naviteri.org/2011/08/new-vocabulary-clothing/}{b}
(\ding{43}~\HL{lun}{\B{lun}})  

% nìmal
\hi
\HT{nImal}{\B{nì\cp{mal}}} \I{[nI."mal]} \emph{adv.} trustingly, without hesitation.  
\B{Ri\-ni tsa\-po\-hu ho\-lum nì\-mal nì\-wotx.} Rini left with that guy without thinking twice about it.~\LINK{http://naviteri.org/2012/01/more-additions-to-the-lexicon/}{b}
(\ding{43}~\HL{mal}{\B{mal}})

% nìmeyp
\hi
\HT{nImeyp}{\B{nì\cp{meyp}}} \I{[nI."mEjp\textcorner]} \emph{adv.} weakly, loosely.
\B{fyep nì\-meyp}, hold loosely.~\LINK{http://naviteri.org/2012/03/spring-vocabulary-part-1/}{b}
\B{nip nì\-meyp}, weakly, loosely stuck.~\LINK{http://naviteri.org/2015/04/some-new-words-for-may-day/}{b} 
(\ding{43}~\HL{meyp}{\B{meyp}})

% nìmu'nitkan
\hi 
\HT{nImu'nitkan}{\B{nì\cp{mu'}nitkan}} \I{[nI."muP.nit\textcorner.kan]} \emph{or} \I{[."mUP.]} \emph{adv.} effectively, in such a way as to reach a target. 
(\ding{43}~\HL{mu'nitkan}{\B{mu'\-nit\-kan}}) 

% nìmun
\hi
\HT{nImun}{\B{nì\cp{mun}}} \I{[nI."mun]} \emph{or} \I{[."mUn]} \emph{adv.} again.
\B{Nga\-ru irayo sei\-yi oe nì\-mun.} I thank you again.~\LINK{http://wiki.learnnavi.org/index.php?title=Canon\#Extracts_from_various_emails}{wiki} 
\B{Zi\-va\-'u nì\-mun ye'\-rìn.} Come again soon.~\LINK{http://naviteri.org/2010/06/first-post/}{b}
\B{Kal\-txì nì\-mun, ma oe\-yä ey\-lan.} Hello again, my friends.~\LINK{http://naviteri.org/2010/09/getting-to-know-you-part-3/}{b}
(\ding{43}~\HL{mune}{\B{mu\-ne}})

% nìmwey
\hi
\HT{nImwey}{\B{nìm\cp{wey}}} \I{[nIm."wEj]} \emph{adv.} calmly, peacefully. 
(\emph{a. idiom}) \B{hi\-va\-haw nìm\-wey} \I{[hi."va.haw nIm."wEj]} ``goodnight''.   (when it's bedtime).~\LINK{http://forum.learnnavi.org/language-updates/txelanit-hivawl/}{ln}
(\ding{43}~\HL{mawey}{\B{ma\-wey}})

% nìmweypey
\hi
\HT{nImweypey}{\B{nìm\cp{wey}pey}} \I{[nIm."wEj.pEj]} \emph{adv.} patiently.
(\ding{43}~\HL{maweypey}{\B{ma\-wey\-pey}})

% nìn
\hi
\HT{nIn}{\B{nìn}} \I{[nIn]} \emph{vtr.} look at (\texttt{+}\emph{control}). 
\B{Po\-ti nìn!} Look at him!~\LINK{http://naviteri.org/2012/11/renu-ayinanfyaya-the-senses-paradigm/}{b} 
\B{Nìn tsat!} Look at that!~\LINK{http://naviteri.org/2012/11/renu-ayinanfyaya-the-senses-paradigm/}{b}
\B{Ru\-txe oe\-yä tì\-oeyk\-tìng\-it a\-saw\-nung ni\-vìn.} Please look at my added explanation.~\LINK{http://naviteri.org/2011/03/word-order-and-case-marking-with-modals/comment-page-1/\#comment-581}{b}
\B{Txe\-wìl tuk\-rut Lo\-ak\-ä nar\-mìn nì\-fmokx.} Txewì was eyeing Loak's spear enviously.~\LINK{http://naviteri.org/2011/10/more-vocabulary-a-bit-of-grammar/}{b}
(\ding{43}~\HL{tIng}{\B{tìng na\-ri}})

% nìNa'vi
\hi
\HT{nIna'vi}{\B{nì\cp{Na'}vi}} \I{[nI."naP.vi]} \emph{adv.} like the Na'\-vi. 
\B{Tsun oe nga\-hu nì\-Na'\-vi pi\-väng\-kxo a fì\-'u oe\-ru prr\-te' lu.} It's a pleasure to be able to chat with you in Na'vi.~\LINK{http://wiki.learnnavi.org/index.php?title=Corpus\#Times_Online}{wiki}
\B{Nga nì\-Na'\-vi pll\-txe nìl\-tsan.} You speak Na'\-vi well.~\LINK{http://naviteri.org/2010/07/diminutives-conversational-expressions/}{b}
(\emph{a. idiom}) \B{X nì\-Na'\-vi slu lì\-'u\-pe\slash pe\-lì\-'u?} How do you say X in Na'\-vi?~\LINK{http://naviteri.org/2010/09/quick-follow-up/}{b}
(\ding{43}~\HL{na'vi}{\B{Na'\-vi}})

% nìnäk
\hi
\HT{nInAk}{\B{nì\cp{näk}}} \I{[nI."n\ae{}k\textcorner]} \emph{adv.} by drinking, in a liquid way. 
\B{Po tä\-pey\-kì\-ye\-ver\-kei\-up nì\-näk.} I'm so jazzed that he may be about to drink himself to death.~\LINK{http://naviteri.org/2010/12/catching-up/}{b}
(\ding{43}~\HL{nAk}{\B{näk}}) 

% nìnän
\hi
\HT{nInAn}{\B{nì\cp{nän}}} \I{[nI."n\ae{}n]} \emph{adv.} less.   
\B{Ru\-txe wi\-vem nì\-nän.} Please fight less.~\LINK{http://naviteri.org/2012/02/trr-asawnung-lefpom-happy-leap-day/}{b} 
\B{Ay\-ha\-pxì\-tu po\-ngu\-ä txo\-pu si nì\-nän ta\-krr\-a Va'\-rul pxe\-ku\-tut lä\-txayn.} The members of the group are less afraid since Va'ru defeated three of the enemy.~\LINK{http://naviteri.org/2012/02/trr-asawnung-lefpom-happy-leap-day/}{b}
(\ding{43}~\HL{nAn}{\B{nän}}, \ding{214}~\HL{nI'ul}{\B{nì\-'ul}})

% nìnänän
\hi
\HT{nInAnAn}{\B{nì\cp{nä}nän}} \I{[nI."n\ae.n\ae{}n]} \emph{adv.} decreasingly, less and less.
\B{Fra\-lo a ta\-ron, oe\-yä 'i\-tan txo\-pu si nì\-nä\-nän.} Each time he hunts, my son becomes less and less afraid.~\LINK{http://naviteri.org/2012/02/trr-asawnung-lefpom-happy-leap-day/}{b}
(\ding{43}~\HL{nInAn}{\B{nì\-nän}}, \ding{214}~\HL{nI'ul'ul}{\B{nì\-'ul\-'ul}})

% nìnew
\hi 
\HT{nInew}{\B{nì\cp{new}}} \I{[nI."nEw]} \emph{adv.} voluntarily, willingly, by desire. 
\B{Nga tsa\-kem so\-li nì\-new srak?!} You did that without being asked to?!~\LINK{http://naviteri.org/2018/07/ta-sulfatu-a-ayliu-niul-more-words-from-our-experts/}{b}
\B{Tsa\-sä\-num\-vit oel po\-ru ka\-yei\-ar nì\-new!} I'm happy to teach him that lesson!~\LINK{http://naviteri.org/2018/07/ta-sulfatu-a-ayliu-niul-more-words-from-our-experts/}{b}
\B{Oel pe\-lun ftxal\-mey nì\-new fu\-ta srung si skxawng\-ur a\-na\-fì\-'u?} Why did I choose, of my own free will, to help such a fool?~\LINK{http://naviteri.org/2018/07/ta-sulfatu-a-ayliu-niul-more-words-from-our-experts/}{b}
(\ding{43}~\HL{new}{\B{new}})

% nìno
\hi
\HT{nIno}{\B{nì\cp{no}}} \I{[nI."no]} \emph{adv.} in detail, thoroughly, with attention to detail. 
\B{Nga lä\-pi\-vawk nì\-no ko.} Tell me all about yourself. [lit.: introduce yourself in detail]~\LINK{http://naviteri.org/2010/09/getting-to-know-you-part-2/}{b}
\B{Tì\-tu\-sa\-ron\-ìri fte fli\-vä, ze\-ne fko si\-vutx smar\-it ni\-nan nì\-no.} To succeed at hunting, you have to track your prey by reading (the forest) with attention to detail.~\LINK{http://naviteri.org/2011/09/\%E2\%80\%9Cby-the-way-what-are-you-reading\%E2\%80\%9D/}{b}
(\ding{43}~\HL{leno}{\B{le\-no}}, \HL{hIno}{\B{hì\-no}}, \HL{tIno}{\B{tì\-no}})

% nìnrra
\hi
\HT{nInrra}{\B{nì\cp{nrr}a}} \I{[nI."n\s{r}.a]} \emph{adv.} proudly, with pride.   
\B{Sa'\-nok nì\-nrr\-a lrr\-tok so\-li krr\-a prr\-nen a\-lo a\-'aw\-ve pol\-txe.} The mother smiled with pride when the baby spoke for the first time.~\LINK{http://naviteri.org/2012/07/fivospxiya-aylifyavi-amip-this-months-new-expressions-pt-2/}{b}
(\ding{43}~\HL{nrra}{\B{nrra}})

% nìnu
\hi
\HT{nInu}{\B{nì\cp{nu}}} \I{[nI."nu]} \emph{adv.} failingly, falteringly, in vain, fruitlessly; \emph{not achieving the desired or expected end}.
\B{Oe\-ru txoa li\-vu. Pol\-txä\-nge nì\-nu.} Forgive me. I misspoke.~\LINK{http://naviteri.org/2013/04/wheres-the-bathroom-and-other-useful-things/}{b}
\B{Kll\-te lu ekx\-txu. Na\-ri si txo\-ke\-fyaw tì\-ran nì\-nu.} The ground is rough. Be careful or you'll trip.~\LINK{http://naviteri.org/2013/04/wheres-the-bathroom-and-other-useful-things/}{b}
(\ding{43}~\HL{nui}{\B{nui}})

% nìngay
\hi
\HT{nIngay}{\B{nì\cp{ngay}}} \I{[nI."Naj]} \emph{adv.} truly.   
\B{Oe\-yä swi\-zaw nì\-ngay ti\-va\-kuk.} Let my arrow strike true.~\LINK{http://wiki.learnnavi.org/index.php/Hunting_Song}{wiki}
\B{Ay\-lì\-'u nge\-yä lor lu nì\-ngay.} Your words are truly beautiful.~\LINK{http://forum.learnnavi.org/language-updates/combined-answers-1-\%28feb-16\%29/}{ln}
\B{Nì\-ngay pll\-txe nga.} You speak truth.~\LINK{http://forum.learnnavi.org/language-updates/confirmation-on-kurakx-eytukan-line-and-toitslan/}{ln\slash a\textsuperscript{c}}
\B{Zì\-sìt a\on{104}°! Ko\-ak nì\-ngay, ke\-fyak?} \on{68} years! Really old, right?~\LINK{http://forum.learnnavi.org/news-announcements/happy-birthday-to-karyu-pawl/msg556899/\#msg556899}{ln}
(\ding{43}~\HL{ngay}{\B{ngay}})

% nìngong
\hi
\HT{nIngong}{\B{nì\cp{ngong}}} \I{[nI."NoN]} \emph{adv.} lethargically, lazily.
(\ding{43}~\HL{ngong}{\B{ngong}})

% nìolo'
\hi
\HT{nIolo'}{\B{nìol\cp{o'}}} \I{[nI.o."loP]} \emph{adv.} (together) as members of a clan.
\B{Txil\-te Ri\-ni\-sì tä\-pa\-re fì\-tsap nì\-soaia, slä tsal\-su\-ngay ke nì\-olo' ta\-krr\-a Ri\-ni mun\-txa slo\-lu.} Txilte and Rini are related by blood, but nevertheless not by clan since the time Rini got married.~\LINK{http://naviteri.org/2012/06/spring-vocabulary-part-3/}{b}
(\ding{43}~\HL{olo'}{\B{olo'}}) 

% nìprrte'
\hi
\HT{nIprrte'}{\B{nì\cp{prr}te'}} \I{[nI."p\s{r}.tEP]} \emph{adv.} gladly, with pleasure.
\B{Fì\-'u\-pxa\-ret a\-syi\-nan oel nì\-prr\-te'.} I intend to read this message gladly.
(\emph{a. idiom}) \B{Zo\-la\-'u nì\-prr\-te'.} ``Welcome.'' (\HL{ne}{\B{ne}}, to)~\LINK{http://naviteri.org/2010/06/first-post/}{b}
(\emph{b.}) \B{Smon nì\-prr\-te'.} ``Nice to know you.''~\LINK{http://naviteri.org/2010/09/quick-follow-up/}{b}
(\emph{c.}) \B{To\-lä\-txaw nì\-prr\-te'.} ``Welcome back.''~\LINK{http://forum.learnnavi.org/language-updates/some-news-from-berlin/msg585072/\#msg585072}{ln}  
(\emph{d. responding to thanks}) \B{Nì\-prr\-te'.} ``You're welcome.''~\LINK{http://naviteri.org/2010/07/diminutives-conversational-expressions/}{b}
(\emph{e.}) \B{sre\-fe\-re\-i\-ey nì\-prr\-te'}, looking forward (to that). \B{Tsa\-ria nga\-hu ye'\-rìn ul\-txa si nì\-mun, oe sre\-fe\-re\-i\-ey nì\-prr\-te'.} I'm looking forward to getting together with you again soon.~\LINK{http://naviteri.org/2011/02/new-vocabulary-part-2/}{b}
(\ding{43}~\HL{prrte'}{\B{prr\-te'}}) \\  
\textsc{Derivatives}: \ding{43} 
\on{1}.~\HL{zawprrte'}{\B{zaw\-prr\-te'}}, be enjoyable.

% nìpxay
\hi
\HT{nIpxay}{\B{nì\cp{pxay}}} \I{[nI."p'aj]} \emph{adv.} many, heavily.
\B{Fwa ay\-ngal fì\-tse\-nget tok nì\-pxay nì\-ftxan oe\-ru te\-ya si nì\-ngay.} It fills me with joy that you are here in such a great number.~\LINK{http://naviteri.org/2012/07/tskxekengtsyip-a-mikyunfpi-a-little-listening-exercise/}{b}
(\emph{a. of weather}) \B{lep\-wopx nì\-pxay}, heavily cloud\hyp{}covered, many clouds.~\LINK{http://naviteri.org/2011/05/weather-part-2-and-a-bit-more-2/}{b}
(\ding{43}~\HL{pxay}{\B{pxay}})

% nìpxi
\hi
\HT{nIpxi}{\B{nì\cp{pxi}}} \I{[nI."p'i]} \emph{adv.} especially, pointedly, unambiguously. 
\B{Por war\-mou fra\-fne\-i\-o\-ang nì\-tut, slä nì\-pxi pxe\-syì\-rä\-fì.} He just couldn't get over all the animals, but was especially taken with the three giraffe.~\LINK{http://forum.learnnavi.org/language-updates/txelanit-hivawl/}{ln}
\B{Po pa\-mll\-txe a krr, fra\-po tar\-mìng mik\-yun nì\-pxi, ta\-lu\-na mok\-ri lu sä\-tsì\-syì\-tsyìp.} When he spoke, everyone listened intently, because his voice was a tiny whisper.~\LINK{http://naviteri.org/2011/09/miscellaneous-vocabulary/}{b}
(\ding{43}~\HL{pxi}{\B{pxi}})

% nìpxim
\hi
\HT{nIpxim}{\B{nì\cp{pxim}}} \I{[nI."p'im]} \emph{adv.} erectly, uprightly. 
\B{Ton\-ìri a\-la\-he, aw\-ngal yom wu\-tsot teng\-krr he\-reyn nì\-pxim.} As for other nights, we eat dinner while sitting erect.~\LINK{http://whyisthisnight.com/addl-more.html\#na\%27vi}{tn}
(\ding{43}~\HL{pxim}{\B{pxim}})

% nìpxul
\hi 
\HT{nIpxul}{\B{nì\cp{pxul}}} \I{[nI."p'ul]} \emph{or} \I{[.p'Ul]} \emph{adv.} formidably, imposingly.~\LINK{http://naviteri.org/2020/08/hiia-vur-a-teri-mefalukantsyip-a-little-story-about-two-cats/}{b}  
(\ding{43}~\HL{pxul}{\B{pxul}})

% nìran
\hi
\HT{nIran}{\B{nì\cp{ran}}} \I{[nI."Ran]} \emph{adv.} basically, fundamentally, in essence.
\B{Nì\-ran lu Lo\-ak mi 'e\-veng slä tsun ti\-va\-ron nì\-teng\-fya na fyeyn\-tu.} Loak is still really just a boy but he can hunt the same as an adult.~\LINK{http://naviteri.org/2012/10/mipa-vospxi-mipa-ayliu-new-words-for-the-new-month/}{b}
(\ding{43}~\HL{ran}{\B{ran}}) 

% nìrangal
\hi
\HT{nIrangal}{\B{nì\cp{rang}al}} \I{[nI."RaN.al]} \emph{adv.} ``I wish …'', ``oh that …'' (\emph{indicating counterfactuals}) 
(\emph{a. with the imperfective subjunctive} ‹\B{irv}› \emph{for a present counterfactual})
\B{Po\-el nì\-rang\-al tir\-vok fì\-tseng\-it.} I wish she were here. (But I know she's not.)~\LINK{http://naviteri.org/2013/01/awvea-posti-zisita-amip-first-post-of-the-new-year/}{b}
(\emph{b. with the perfective subjunctive} ‹\B{ilv}› \emph{or past subjunctive} ‹\B{imv}› \emph{for a past counterfactual})
\B{Poe nì\-rang\-al zim\-va\-'u\slash zil\-va\-'u trr\-am.} I wish she had come yesterday. (But she didn't.)~\LINK{http://naviteri.org/2013/01/awvea-posti-zisita-amip-first-post-of-the-new-year/}{b}
(\emph{c. with the future subjunctive} ‹\B{iyev}›\slash ‹\B{ìyev}› \emph{for a future counterfactual})
\B{Poe nì\-rang\-al zì\-ye\-va\-'u trr\-ay.} I wish she were coming tomorrow. (But I know she isn't.)~\LINK{http://naviteri.org/2013/01/awvea-posti-zisita-amip-first-post-of-the-new-year/}{b}
(\ding{43}~\HL{rangal}{\B{rang\-al}}, \HL{nIsIltsan}{\B{nì\-sìl\-pey}})

% nìräptum
\hi 
\HT{nIrAptum}{\B{nìräp\cp{tum}}} \I{[nI.R\ae{}p\textcorner."tum]} \emph{or} \I{[."tUm]} (\textsc{rn}: \B{nì\-räp\-tùm} \I{[nI.R\ae{}p\textcorner."tUm]}) \emph{adv.} coarsly, vulgarly. 
\B{Fya\-pe yaw\-ne lu fko\-ru tu\-te a fra\-krr vo\-ìk si fì\-txan nì\-räp\-tum?} How does one love a person who always behaves so coarsely?~\LINK{https://naviteri.org/2024/02/mipa-ayliu-si-aylifyavi-niul-more-new-words-and-expressions/}{b} 
(\ding{43}~\HL{rAptum}{\B{räp\-tum}})

% nìrim
\hi 
\HT{nIrim}{\B{nì\cp{rim}}} \I{[nI."Rim]} \emph{adv.} in yellow, `yellow-ly'.
\B{Tom\-pa\-yä ka\-to | Tsaw\-ke\-yä ka\-to | Ka\-tot tä\-ftxu oel | Nì\-ean nì\-rim.} The rhythm of the rain | The rhythm of the sun | I weave the rhythm | In blue and yellow.~\LINK{http://wiki.learnnavi.org/index.php/Weaving_Song}{wiki}
(\ding{43}~\HL{rim}{\B{rim}})

% nìrìkxi
\hi 
\HT{nIrIkxi}{\B{nìrì\cp{kxi}}} \I{[nI.RI."k'i]} \emph{adv.} shakily, tremblingly.
(\ding{43}~\HL{rIkxi}{\B{rì\-kxi}})

% nìronsrel
\hi
\HT{nIronsrel}{\B{nì\cp{ron}srel}} \I{[nI."Ron.sREl]} \emph{adv.} in\slash by imagination.
\B{Oe pxìm päng\-kxo nga\-hu nì\-ron\-srel.} I often talk with you in my imagination.\slash I often imagine I'm talking with you.~\LINK{http://naviteri.org/2011/02/new-vocabulary-part-2/}{b}
(\ding{43}~\HL{ronsrel}{\B{ron\-srel}})

% nìsay
\hi 
\HT{nIsay}{\B{nì\cp{say}}} \I{[nI."{}saj]} \emph{adv.} loyally.  
(\ding{43}~\HL{say}{\B{say}})

% nìsìlpey
\hi
\HT{nIsIlpey}{\B{nìsìl\cp{pey}}} \I{[nI.sIl."pEj]} \textsc{i}. \emph{sentence adv.} hopefully (\emph{expresses a hope that s.t. is true}) 
(\emph{a. with the subjunctive} ‹\B{iv}› \emph{for a present wish})
\B{Po\-el nì\-sìl\-pey ti\-vok fì\-tseng\-it.} I hope she's here.~\LINK{http://naviteri.org/2013/01/awvea-posti-zisita-amip-first-post-of-the-new-year/}{b}
(\emph{b. with the perfective subjunctive} ‹\B{ilv}› \emph{or past subjunctive} ‹\B{imv}› \emph{for a past wish})
\B{Poe nì\-sìl\-pey zim\-va\-'u\slash zil\-va\-'u trr\-am.} I hope she came yesterday.~\LINK{http://naviteri.org/2013/01/awvea-posti-zisita-amip-first-post-of-the-new-year/}{b}
(\emph{c. with the future subjunctive} ‹\B{iyev}›\slash ‹\B{ìyev}› \emph{for a future wish})
\B{Poe nì\-sìl\-pey zì\-ye\-va\-'u trr\-ay.} I hope she'll come tomorrow.~\LINK{http://naviteri.org/2013/01/awvea-posti-zisita-amip-first-post-of-the-new-year/}{b} \\
\textsc{ii}. \emph{manner adv.} in a hopeful way.
\B{Tsyeyk ä\-txä\-le so\-li nì\-sìl\-pey tsnì li\-vu por Unil\-ta\-ron.} Jake hopefully requested the Dream Hunt.~\LINK{http://naviteri.org/2013/01/awvea-posti-zisita-amip-first-post-of-the-new-year/}{b}
(\ding{43}~\HL{sIlpey}{\B{sìl\-pey}}, \HL{nIrangal}{\B{nì\-rang\-al}})

% nìskien
\hi
\HT{nIskien}{\B{nì\cp{ski}en}} \I{[nI."{}ski.En]} \emph{adv.} to the right.
\B{Yak si nì\-ski\-en.} Turn to the right.~\ding{43}~\LINK{http://naviteri.org/2013/09/on-si-salewfya-shapes-and-directions/}{b}
(\ding{43}~\HL{skien}{\B{ski\-en}}, \ding{214}~\HL{nIftAr}{\B{nì\-ftär}})

% nìslele
\hi
\HT{nIslele}{\B{nì\cp{sle}le}} \I{[nI."{}slE.lE]} \emph{adv.} by swimming.
\B{Tsun fko tsa\-tse\-nge\-ne ki\-vä nì\-sle\-le fu fa fwa ik\-ran\-it mak\-to nì\-'aw.} You can only get there by swimming or riding an ikran.~\LINK{http://naviteri.org/2011/02/new-vocabulary-ii\%E2\%80\%94part-1/}{b}
(\ding{43}~\HL{slele}{\B{sle\-le}})

% nìsngä'i
\hi
\HT{nIsngA'i}{\B{nì\cp{sngä}'i}} \I{[nI."{}sN\ae.Pi]} \emph{adv.} originally, at first.  
\B{Nì\-sngä\-'i fmawn\-it fo nar\-mew wi\-van, slä nì\-'i'\-a fra\-por lo\-lo\-nu.} They originally wanted to hide the news, but in the end they revealed it everyone.~\LINK{http://naviteri.org/2011/10/more-vocabulary-a-bit-of-grammar/}{b}
(\ding{43}~\HL{sngA'i}{\B{sngä\-'i}}, \HL{nItsim}{\B{nì\-tsim}})

% nìsngum 
\hi
\HT{nIsngum}{\B{nì\cp{sngum}}} \I{[nI."{}sNum]} \emph{or} \I{[."{}sNUm]} \emph{adv.} worryingly, fretfully. 
\B{Swey lu fwa nga fì\-kem ke si\-vi nì\-sngum.} It's best that you not do this fretfully.\slash It's best that you don't freak out about doing this.~\LINK{http://naviteri.org/2011/01/new-year-new-vocabulary/}{b}
(\ding{43}~\HL{sngum}{\B{sngum}})

% nìso'ha
\hi
\HT{nIso'ha}{\B{nì\cp{so'}ha}} \I{[nI."{}soP.ha]} \emph{adv.} enthusiastically.
(\ding{43}~\HL{so'ha}{\B{so'\-ha}})

% nìsoaia
\hi
\HT{nIsoaia}{\B{nìso\cp{a}ia}} \I{[nI.so."{}a.i.a]} \emph{adv.} (together) as members of a family. 
\B{Txil\-te Ri\-ni\-sì tä\-pa\-re fì\-tsap nì\-soaia, slä tsal\-su\-ngay ke nì\-olo' ta\-krr\-a Ri\-ni mun\-txa slo\-lu.} Txilte and Rini are related by blood, but nevertheless not by clan since the time Rini got married.~\LINK{http://naviteri.org/2012/06/spring-vocabulary-part-3/}{b}
(\ding{43}~\HL{soaia}{\B{soaia}})

% nìsok
\hi 
\HT{nIsok}{\B{nì\cp{sok}}} \I{[nI."{}sok\textcorner]} \emph{adv.} recently. 
(\ding{43}~\HL{sok}{\B{sok}})

% nìsteng
\hi
\HT{nIsteng}{\B{nì\cp{steng}}} \I{[nI."{}stEN]} \emph{adv.} similarly, in a similar fashion.
\B{Tseyk tswa\-may\-on fa ik\-ran sre\-krr; ta\-fral fmo\-li fì\-kem si\-vi fa to\-ruk nì\-steng.} Jake flew with an ikran before; therefore he tried to do it with a toruk in a similar fashion.~\LINK{http://naviteri.org/2011/09/miscellaneous-vocabulary/}{b} 
(\ding{43}~\HL{steng}{\B{steng}})

% nìsti
\hi
\HT{nIsti}{\B{nì\cp{sti}}} \I{[nI."{}sti]} \emph{adv.} angrily. 
(\ding{43}~\HL{sti}{\B{sti}})

% nìsung
\hi
\HT{nIsung}{\B{nì\cp{sung}}} \I{[nI."{}suN]} \emph{or} \I{[."{}sUN]} (\textsc{rn}: \B{nì\cp{sùng}} \I{[nI."{}sUN]}) \emph{adv.} besides, additionally, furthermore.
\B{Nì\-sung, txo tsa\-kem li\-ven mì zu\-saw\-krr, ze\-ne fko si\-var kem\-lì\-'u\-vit a\-lu ‹ìyev›.} Furthermore, if that happens in the future, one must use the infix ‹\B{ìyev}›.~\LINK{http://naviteri.org/2013/01/awvea-posti-zisita-amip-first-post-of-the-new-year/comment-page-1/\#comment-1910}{b}
(\ding{43}~\HL{sung}{\B{sung}})

% nìswey
\hi
\HT{nIswey}{\B{nì\cp{swey}}} \I{[nI."{}swEj]} \emph{adv.} optimally, best.
\B{Maw\-krr la\-ye\-iu oer krr nì\-'ul fte ngi\-vop ay\-lì\-'ut sì tsay\-fne\-sä\-num\-vit a tsun fra\-por srung si\-vi fte ni\-vu\-me sì zi\-ve\-rok nì\-swey.} After that I will have more time to create words and the kinds of lessons that can help everyone best learn and remember.~\LINK{http://forum.learnnavi.org/language-updates/trr-rrtaya/}{ln}
(\ding{43}~\HL{swey}{\B{swey}})

% nìsyayvi
\hi
\HT{nIsyayvi}{\B{nì\cp{syay}vi}} \I{[nI."{}sjaj.vi]} \emph{adv.} by chance, by coincidence, as luck would have it.
\B{Et\-rìp\-a syay\-vi!} Good luck!~\LINK{http://naviteri.org/2010/07/vocabulary-update/}{b}.
(\ding{43}~\HL{syayvi}{\B{syay\-vi}})

% nìsyen
\hi
\HT{nIsyen}{\B{nì\cp{syen}}} \I{[nI."{}sjEn]} \emph{adv.} lastly, finally.
\B{Ul\-te nì\-syen …}, And lastly …~\LINK{http://naviteri.org/2011/09/miscellaneous-vocabulary/}{b}
(\ding{43}~\HL{syen}{\B{syen}}, \HL{nI'i'a}{\B{nì'i'a}})

% nìsyep
\hi
\HT{nIsyep}{\B{nì\cp{syep}}} \I{[nI."{}sjEp\textcorner]} \emph{adv.} tightly, in an iron grip.
\B{fyep nì\-syep}, hold tightly.~\LINK{http://naviteri.org/2012/03/spring-vocabulary-part-1/}{b}
\B{nip nì\-syep}, stuck irremovably.~\LINK{http://naviteri.org/2015/04/some-new-words-for-may-day/}{b} 
(\ding{43}~\HL{syep}{\B{syep}})

% nìt'iluke
\hi 
\HT{nIt'iluke}{\B{nìt\cp{'i}luke}} \I{[nIt\textcorner."{}i.lu.kE]} \emph{adv.} never\hyp{}endingly, forever. 
(\ding{43}~\HL{tI'iluke}{\B{tì\-'i\-lu\-ke}})

% nìtam
\hi
\HT{nItam}{\B{nì\cp{tam}}} \I{[nI."tam]} \emph{adv.} enough.
\B{Hu\-fwa nge\-yä tì\-hawl\-ìri ke lu kea kxey\-ey, tsal\-su\-ngay oe\-ru ke ha' nì\-tam.} Although there's nothing wrong with your plan, it just doesn't suit me.~\LINK{http://naviteri.org/2012/10/mipa-vospxi-mipa-ayliu-new-words-for-the-new-month/}{b}
(\ding{43}~\HL{tam}{\B{tam}})

% nìteng
\hi
\HT{nIteng}{\B{nì\cp{teng}}} \I{[nI."tEN]} \emph{adv.} too, likewise, as well.   
\B{Nì\-ay\-nga oe pe\-rey nì\-teng.} Like you, I too am waiting.~\LINK{http://forum.learnnavi.org/news-announcements/a-response-from-paul-frommer!/}{ln}
\B{Lu poe lor, lu yu\-ey nì\-teng.} She's beautiful on the outside \emph{and} the inside.~\LINK{http://naviteri.org/2012/03/spring-vocabulary-part-1/}{b}
(\ding{43}~\HL{teng}{\B{teng}}, \HL{kop}{\B{kop}})

% nìtengfya
\hi
\HT{nItengfya}{\B{nì\cp{teng}fya}} \I{[nI."tEN.fja]} \emph{adv.} in the same way. 
\B{Oel tel 'u\-pxa\-ret le\-Na'\-vi a krr, new oe nì\-teng\-fya pam\-rel si\-vi nì\-'eyng.} When I receive a Na'vi message, I want to write the same way in response.~\LINK{http://forum.learnnavi.org/language-updates/verb-for-write/msg175592/\#msg175592}{ln}
(\ding{43}~\HL{tengfya}{\B{teng\-fya}})

% nìtkan
\hi
\HT{nItkan}{\B{nìt\cp{kan}}} \I{[nIt\textcorner."kan]} \emph{adv.} purposefully, deliberately.
(\ding{43}~\HL{tIkan}{\B{tì\-kan}})

% nìtkanluke
\hi
\HT{nItkanluke}{\B{nìt\cp{kan}luke}} \I{[nIt\textcorner."kan.lu.kE]} \emph{adv.} accidentally, unintentionally.
(\ding{43}~\HL{tIkan}{\B{tì\-kan}}, \HL{luke}{\B{lu\-ke}})

% nìToitsye
\hi
\HT{nItoitsye}{\B{nì\cp{To}itsye}} \I{[nI."to.i.\t{ts}jE]} \emph{adv.} German, in German.
\B{X-(ì)ri pe\-ral nì\-To\-i\-tsye?} What does X mean in German?~\LINK{http://naviteri.org/2010/09/getting-to-know-you-part-2/}{b}
(\ding{43}~\HL{toitsye}{\B{Toi\-tsye}})

% nìtokat
\hi 
\HT{nItokat}{\B{nì\cp{to}kat}} \I{[nI."to.kat\textcorner]} \emph{adv.} guiltily. 
\B{zaw\-prr\-te' nì\-to\-kat fko\-ne}, be pleasurable to one in a guilty way, \emph{schadenfreude, i.e. taking pleasure in s.o. else's misfortune}.
\B{En\-tul fot ve'\-kì ul\-te sä\-nui fe\-yä zo\-law\-prr\-te' nì\-to\-kat po\-ne.} En\-tu hates them and their failure brought him pleasure.~\LINK{http://naviteri.org/2023/09/aawa-ayliu-si-aylifyavi-amip-a-few-new-words-and-expressions/}{b}  
(\ding{43}~\HL{tokat}{\B{to\-kat}}; \ding{214}~\HL{nIlayl}{\B{nì\-layl}})

% nìtrea
\hi 
\HT{nItrea}{\B{nìt\cp{re}a}} \I{[nIt\textcorner."RE.a]} \emph{adv.} in spirit. 
\B{Ke tsä\-ngun Tsyìm zi\-va'\-u ftxo\-zä\-ne, slä tok nìt\-re\-a.} Sadly, Jim couldn't come to the celebration, but he was there in spirit.~\LINK{http://naviteri.org/2022/12/trr-anawm-polaheiem-the-great-day-has-arrived/}{b} 
(\ding{43}~\HL{tirea}{\B{ti\-re\-a}})

% nìtrrtrr
\hi
\HT{nItrrtrr}{\B{nì\cp{trr}trr}} \I{[nI."t\s{r}.t\s{r}]} \emph{adv.} daily, regularly, on a daily basis, typically.   
\B{Nga\-ru lu pe\-fne\-txin\-tìn nì\-trr\-trr?} What is your primary role (in society)?~\LINK{http://naviteri.org/2010/09/getting-to-know-you-part-2/}{b}
\B{Nga pe\-su\-hu kä\-teng nì\-trr\-trr?} Who do you usually hang out with?~\LINK{http://naviteri.org/2010/09/getting-to-know-you-part-3/}{b}
\B{Nì\-trr\-trr yom Na'\-vil wu\-tsot 'aw\-si\-teng pxaw yll\-txep.} The Na'vi regularly eat dinner together around a communal fire.~\LINK{http://naviteri.org/2012/07/fivospxiya-aylifyavi-amip-this-months-new-expressions-pt-2/}{b}
\B{Nì\-trr\-trr pxì\-mun\-'i sam\-si\-yul ay\-swi\-zaw\-it ku\-tu\-ä a\-law\-nä\-txayn sno\-kip nì\-'eng.} Warriors typically share the arrows of their defeated enemies among themselves.~\LINK{http://naviteri.org/2012/01/more-additions-to-the-lexicon/}{b}
(\ding{43}~\HL{letrrtrr}{\B{le\-trr\-trr}})

% nìtut
\hi
\HT{nItut}{\B{nì\cp{tut}}} \I{[nI."tut\textcorner]} \emph{or} \I{[."tUt\textcorner]} (\textsc{rn}: \B{nì\cp{tùt}} \I{[nI."tUt\textcorner]}) \emph{adv.} continually, continuously, always.
\B{Ay\-nge\-yä ki\-fkey\-ti fol kawng\-sar nì\-tut, fì\-'u\-ri ke\-kem ke si ay\-nga.} They continuously exploit your world and you do nothing about it.~\LINK{http://naviteri.org/2011/07/txantsana-ultxa-mi-siatll-great-meeting-in-seattle/}{b}
\B{Ri\-ni\-ri nik\-re yä\-ngo' nì\-tut.} Ri\-ni's hair is always perfect.~\LINK{http://naviteri.org/2012/01/mipa-zisit-ayliu-amip-new-words-for-the-new-year/}{b}
\B{Por war\-mou fra\-fne\-i\-o\-ang nì\-tut, slä nì\-pxi pxe\-syì\-rä\-fì.} He just couldn't get over all the animals, but was especially taken with the three giraffe.~\LINK{http://forum.learnnavi.org/language-updates/txelanit-hivawl/}{ln}
\B{Mong prr\-nen\-ìl fu\-ta sa'\-nul ve\-rewng nì\-tut.} The infant depends on her Mommy to look after her constantly.~\LINK{http://naviteri.org/2012/07/meetings-waterfalls-and-more/}{b}
\B{Pa\-lu\-lu\-kan\-ìri lu pxe\-syon a ze\-ne fko zi\-ve\-rok nì\-tut.} Three things about the thanator must always be kept in mind.~\LINK{http://naviteri.org/2012/10/mipa-vospxi-mipa-ayliu-new-words-for-the-new-month/}{b}
(\ding{43}~\HL{tut}{\B{tut}})

% nìtxan
\hi
\HT{nItxan}{\B{nì\cp{txan}}} \I{[nI."t'an]} \emph{adv.} (\emph{a.}) very much, very.
\B{Oe new nì\-txan ay\-nga\-ru fya\-wi\-vìn\-txu.} I want very much to guide you.~\LINK{http://forum.learnnavi.org/news-announcements/a-response-from-paul-frommer!/}{ln}
\B{Lam oer fwa po lu ka\-nu nì\-txan.} It seems to me that he is very intelligent.~\LINK{http://forum.learnnavi.org/language-updates/txelanit-hivawl/}{ln}
\B{Nga nì\-aw\-no\-mum to oe\-tsyìp lu txur nì\-txan.} As everyone knows, you're a lot stronger than little old me.~\LINK{http://naviteri.org/2010/07/diminutives-conversational-expressions/}{b}
\B{Ay\-nge\-yä ay\-sä\-pll\-txe\-vi\-ri sei\-yi oe ira\-yo nì\-txan ay\-nga\-ru nì\-wotx!} I thank you all very much for the comments!~\LINK{http://naviteri.org/2011/07/attributive-\%E2\%80\%9Ca\%E2\%80\%9D-and-truncated-style/}{b}
\B{Ira\-yo nì\-txan.} Thanks very much.~\LINK{http://naviteri.org/2010/06/irayo-thank-you-and-some-miscellaneous-thoughts/}{b}
(\emph{b. of time}) long.   
\B{Fwa oel fì\-tì\-kang\-kem\-vit sngey\-ki\-vä\-'i krr\-no\-lekx nì\-txan.} Me starting this project took a long time.~\LINK{http://naviteri.org/2010/06/first-post/}{b}
(\ding{43}~\HL{txan}{\B{txan}})

% nìtxankeltrrtrr
\hi 
\HT{nItxankeltrrtrr}{\B{nì\cp{txan}kel\cp{trr}trr}} \I{[nI."t'an.kEl."t\s{r}.t\s{r}]} \emph{adj.} extraordinarily. 
\B{Oey 'ey\-lan pll\-txe nì\-Na'\-vi na hu\-fwe nì\-txan\-kel\-trr\-trr.} My friend speaks Na'vi extraordinarily fluently.~\LINK{http://naviteri.org/2018/03/100a-liu-amip-64-new-words-part-1/}{b}
(\ding{43}~\HL{txankeltrrtrr}{\B{txan\-kel\-trr\-trr}})

% nìtxay
\hi
\HT{nItxay}{\B{nì\cp{txay}}} \I{[nI."t'aj]} \emph{adv.} horizontally.
(\ding{43}~\HL{txay}{\B{txay}})

% nìtxi
\hi
\HT{nItxi}{\B{nì\cp{txi}}} \I{[nI."t'i]} \emph{adv.} hurriedly, in a frenzied way. 
(\ding{43}~\HL{txi}{\B{txi}}; \ding{214}~\HL{nItxiluke}{\B{nì\-txi\-lu\-ke}})

% nìtxiluke
\hi
\HT{nItxiluke}{\B{nì\cp{txi}luke}} \I{[nI."t'i.lu.kE]} \emph{adv.} unhurriedly, leisurely   \B{Tsun oe nga\-hu tsa\-tse\-nge\-ne ki\-vä, slä nul\-new fu\-ta si\-vop oeng nì\-txi\-lu\-ke.} I can go there with you, but I prefer to travel leisurely.~\LINK{http://naviteri.org/2011/10/more-vocabulary-a-bit-of-grammar/}{b}
(\ding{43}~\HL{txi}{\B{txi}}, \HL{luke}{\B{lu\-ke}}; \ding{214}~\HL{nItxi}{\B{nì\-txi}})

% nìtxukx
\hi 
\HT{nItxukx}{\B{nì\cp{txukx}}} \I{[nI."t'uk']} \emph{or} \I{[."t'Uk']} (\textsc{rn}: \B{nìtxukx} \I{[nI."duk']}) \emph{adv.} deeply (\emph{also metaphorically}). 
\B{Po\-an\-it tso\-lukx po\-el fa tstal nì\-txukx nem\-fa heyr.} She stabbed him deeply in the chest with a knife.~\LINK{http://naviteri.org/2019/06/50a-liu-amip-40-new-words/}{b}
\B{Fì\-tì\-pawm\-te\-ri fpar\-mìl oe nì\-txukx, slä vay set ke ro\-lä\-ngun tì\-'eyng\-it.} I've thought about this question deeply, but I'm sad to say I haven't yet found the answer.~\LINK{http://naviteri.org/2019/06/50a-liu-amip-40-new-words/}{b}
(\ding{43}~\HL{txukx}{\B{txukx}})

% nìtsim
\hi
\HT{nItsim}{\B{nì\cp{tsim}}} \I{[nI."\t{ts}im]} \emph{adv.} originally, in an original way, with originality.   
\B{Fra\-krr po fpìl nì\-tsim nì\-wotx.} Her thinking is always completely original.~\LINK{http://naviteri.org/2011/10/more-vocabulary-a-bit-of-grammar/}{b}
(\ding{43}~\HL{tsim}{\B{tsim}}, \HL{nIsngA'i}{\B{nì\-sngä\-'i}})

% nìtsìsyì
\hi
\HT{nItsIsyI}{\B{nì\cp{tsì}syì}} \I{[nI."\t{ts}I.sjI]} \emph{adv.} by whispering, in a whisper.   
\B{Pol tstxo\-ti oe\-yä pol\-txe\-i\-e nì\-tsì\-syì.} He whispered my name, I'm happy to say.~\LINK{http://naviteri.org/2011/09/miscellaneous-vocabulary/}{b}
(\ding{43}~\HL{tsIsyI}{\B{tsìsyì}})

% nìtsleng
\hi 
\HT{nItsleng}{\B{nì\cp{tsleng}}} \I{[nI."\t{ts}lEN]} \emph{adv.} falsely. 
\B{Nì\-tsleng pll\-txe nga.} You speak falsely [i.e. you are a liar].~\LINK{http://naviteri.org/2017/09/zisikrr-amip-ayliu-amip-new-words-for-the-new-season/}{b} 
\B{Pll\-txe nì\-tsleng! Tsa\-fkxi\-let ke to\-lìng ngar En\-tul!}Liar! En\-tu didn't give you that necklace!~\LINK{http://naviteri.org/2017/09/zisikrr-amip-ayliu-amip-new-words-for-the-new-season/}{b}  
(\ding{43}~\HL{tsleng}{\B{tsleng}})

% nìtson
\hi
\HT{nItson}{\B{nì\cp{tson}}} \I{[nI."\t{ts}on]} \emph{adv.} dutifully, as an obligation.
\B{Pol vewng fra\-trr ay\-eveng\-it nì\-tson.} He observes his duty to care for the kids every day.~\LINK{http://naviteri.org/2014/10/tson-si-fpomron-obligation-and-mental-health/}{b}
(\ding{43}~\HL{tson}{\B{tson}})

% nìtstew
\hi 
\HT{nItstew}{\B{nì\cp{tstew}}} \I{[nI."\t{ts}tEw]} \emph{adv.} bravely, courageously.~\LINK{http://naviteri.org/2020/11/vospxivopeya-ayliu-amip-novembers-new-words/}{b} 
(\ding{43}~\HL{tstew}{\B{tstew}})

% nìtsyìl
\hi
\HT{nItsyIl}{\B{nì\cp{tsyìl}}} \I{[nI."\t{ts}jIl]} \emph{adv.} by climbing.
(\ding{43}~\HL{tsyIl}{\B{tsyìl}})

% nìTsyungwen
\hi
\HT{nItsyungwen}{\B{nì\cp{Tsyung}wen}} \I{[nI."\t{ts}juN.wEn]} \emph{or} \I{[."\t{ts}jUN.]} (\textsc{rn}: \B{nì\cp{Tsyùng}wen} \I{[nI."\t{tS}UN.wEn]}) \emph{adv.} like the Chinese, in Mandarin.
\B{X-(ì)ri pe\-ral nì\-Tsyung\-wen?} What does X mean in Chinese?~\LINK{http://naviteri.org/2010/09/getting-to-know-you-part-2/}{b}
(\ding{43}~\HL{tsyungwen}{\B{Tsyung\-wen}})

% nìvingkap
\hi
\HT{nIvingkap}{\B{nì\cp{ving}kap}} \I{[nI."viN.kap\textcorner]} \emph{adv.} by the way, incidentally (\emph{to change the subject and introduce something new into the conversation}). 
\B{Nì\-ving\-kap nge\-yä tsmu\-kan\-ur alu Ra\-lu lu fpom srak?} Oh by the way, how's your brother Ralu?~\LINK{http://naviteri.org/2011/09/\%E2\%80\%9Cby-the-way-what-are-you-reading\%E2\%80\%9D/}{b}
(\ding{43}~\HL{vingkap}{\B{ving\-kap}}, \HL{mIftxele}{\B{mì\-ftxe\-le}})

% nìwan
\hi
\HT{nIwan}{\B{nì\cp{wan}}} \I{[nI."wan]} \emph{adv.} secretly; in hiding, by hiding.  
\B{Sam\-si\-yu pe\-rey nì\-wan.} The hunters lie in wait, prepared to ambush.~\LINK{http://naviteri.org/2011/02/new-vocabulary-ii\%E2\%80\%94part-1/}{b}
\B{Ìlä Fey\-ral, mun\-txa so\-li Ra\-lu sì Ne\-wey nì\-wan me\-srr\-am.} According to Peyral, Ralu and Newey were secretly married the day before yesterday.~\LINK{http://naviteri.org/2012/07/meetings-waterfalls-and-more/}{b}
(\ding{43}~\HL{wan}{\B{wan}})

% nìwawe
\hi
\HT{nIwawe}{\B{nìwa\cp{we}}} \I{[nI.wa."wE]} \emph{adv.} meaningfully, significantly. 
\B{Tsa\-txon ay\-ut\-ral\-äo krr a pol oe\-ti no\-lìn nì\-wa\-we, olo\-me\-i\-um oel fa key\-rel fu\-ta lu yaw\-ne oe po\-ru.} That night beneath the trees when she looked at me meaningfully, I knew by the expression on her face that she loves me.~\LINK{http://naviteri.org/2011/05/some-miscellaneous-vocabulary/}{b}
(\ding{43}~\HL{wawe}{\B{wa\-we}})

% nìwä
\hi
\HT{nIwA}{\B{nì\cp{wä}}} \I{[nI."w\ae]} \emph{adv.} on the contrary, conversely. 
(\ding{43}~\HL{wA}{\B{wä}})

% nìwäte
\hi
\HT{nIwAte}{\B{nìwä\cp{te}}} \I{[nI.w\ae."tE]} \emph{adv.} disagreeably, begrudgingly.
(\ding{43}~\HL{wAte}{\B{wä\-te}}) 

% nìwin
\hi
\HT{nIwin}{\B{nì\cp{win}}} \I{[nI."win]} \emph{adv.} (\emph{a.}) fast, quickly. 
\B{kä\-mak\-to nì\-win}, ride out fast.~\LINK{http://forum.learnnavi.org/language-updates/naviya-the-other-vocative/}{ln\slash a}
\B{Ey\-wal zey\-ki\-vo ngat nì\-win.} May Eywa heal you quickly.~\LINK{http://naviteri.org/2010/07/vocabulary-update/}{b}
\B{Oe tul nì\-ftxan nì\-win na nga.} I run as fast as you.~\LINK{http://forum.learnnavi.org/language-updates/confirmations-\%28comparisons-gerund\%29/}{ln}
\B{Oe to nga tul nì\-win.} I run faster than you.~\LINK{http://forum.learnnavi.org/language-updates/confirmations-\%28comparisons-gerund\%29/}{ln}
\B{Po tul nì\-win fra\-to.} She is the fastest.~\LINK{http://forum.learnnavi.org/language-updates/confirmations-\%28comparisons-gerund\%29/}{ln}
(\emph{b. expression}) \B{nì\-win na hu\-fwe}, ``as fast as the wind'' (\emph{any situation to express rapidity}). \B{na hu\-fwe}, fluent (not just rapid) use of language.~\LINK{http://naviteri.org/2015/08/aylifyavi-lereyfya-2-cultural-terms-2-and-more/}{b}
(\ding{43}~\HL{win}{\B{win}})

% nìwok
\hi
\HT{nIwok}{\B{nì\cp{wok}}} \I{[nI."wok\textcorner]} \emph{adv.} loudly. 
\B{Ru\-txe pi\-vll\-txe nì\-wok nì\-'it.} Please speak up a bit.~\LINK{http://naviteri.org/2011/09/miscellaneous-vocabulary/}{b}
(\ding{43}~\HL{wok}{\B{wok}})

% nìwotx
\hi
\HT{nIwotx}{\B{nì\cp{wotx}}} \I{[nI."wot']} \emph{adv.} all (of), completely, \emph{in toto}.   
\B{Ay\-ftxo\-zä le\-fpom ay\-nga\-ru nì\-wotx!} Happy Holidays to You All!~\LINK{http://wiki.learnnavi.org/index.php?title=Corpus\#Out_of_Office_AutoReply}{wiki} 
\B{Ay\-ey\-lan\-ur oe\-yä sì ey\-lan\-ur lì'\-fya\-yä le\-Na'\-vi nì\-wotx.} To all my friends and friends of the Na'vi language.~\LINK{http://forum.learnnavi.org/news-announcements/a-response-from-paul-frommer!/}{ln} 
\B{Ri\-ni tsa\-po\-hu ho\-lum nì\-mal nì\-wotx.} Rini left with that guy without thinking twice about it.~\LINK{http://naviteri.org/2012/01/more-additions-to-the-lexicon/}{b}
\B{Nga\-ru tì\-yawr nì\-wotx!} You are completely right.~\LINK{http://naviteri.org/2010/08/a-na\%E2\%80\%99vi-alphabet/comment-page-1/\#comment-191}{b}
\B{Fti\-a oel lì'\-fya\-ti le\-Na'\-vi nì\-'o' nì\-wotx!} Studying Na'vi is a ton of fun for me.~\LINK{http://naviteri.org/2010/09/getting-to-know-you-part-3/}{b}
(\ding{43}~\HL{wotx}{\B{wotx}})

% nìyawr
\hi
\HT{nIyawr}{\B{nì\cp{yawr}}} \I{[nI."jawR]} \emph{adv.} correctly, rightly. 
(\ding{43}~\HL{eyawr}{\B{e\-yawr}})

% nìyewla
\hi
\HT{nIyewla}{\B{nì\cp{yew}la}} \I{[nI."jEw.la]} \emph{adv.} in a disappointing fashion, \emph{in a way failing to meet expectations}.   
\B{Trr\-am\-ä ay\-u\-van\-ìri mak\-to Ak\-wey nì\-yew\-la ha snaytx.} Akwey rode disappointingly in yesterday's games so he lost.~\LINK{http://naviteri.org/2012/10/mipa-vospxi-mipa-ayliu-new-words-for-the-new-month/}{b}
(\ding{43}~\HL{yewla}{\B{yew\-la}})

% nìyey
\hi
\HT{nIyey}{\B{nì\cp{yey}}} \I{[nI."jEj]} \emph{adv.} directly, straight to the point; just.
\B{Kem si nì\-yey!} Just do it!~\LINK{http://naviteri.org/2014/07/tsawlultxamaw-a-ayliu-post-meetup-words/}{b}
(\ding{43}~\HL{yey}{\B{yey}})

% nìyeyfya
\hi
\HT{nIyeyfya}{\B{nì\cp{yey}fya}} \I{[nI."jEj.fja]} \emph{adv.} straight ahead, in a straight line. 
\B{Sa\-lew nì\-yey\-fya.} Proceed straight ahead.~\LINK{http://naviteri.org/2013/09/on-si-salewfya-shapes-and-directions/}{b}
(\ding{43}~\HL{yey}{\B{yey}}, \HL{fya'o}{\B{fya\-'o}})

% nìyll
\hi
\HT{nIyll}{\B{nì\cp{yll}}} \I{[nI."j\s{l}]} \emph{adv.} communally, in a communal manner (\emph{to express sharing with the community}).
\B{Fol tsngan\-it pxì\-mo\-lun\-'i nì\-yll.} They shared the meat with the entire clan.~\LINK{http://naviteri.org/2012/01/more-additions-to-the-lexicon/}{b}
\B{Fì\-tey\-lu\-ri ke nar\-mew Va'\-ru yi\-vom nì\-yll.} Va'ru didn't want to share this \emph{teylu} with the Omatikaya.~\LINK{http://naviteri.org/2012/01/more-additions-to-the-lexicon/}{b}
(\ding{43}~\HL{yll}{\B{yll}}; \HL{nI'eng}{\B{nì\-'eng}})

% nìyo'
\hi
\HT{nIyo'}{\B{nì\cp{yo'}}} \I{[nI."joP]} \emph{adv.} perfectly, flawlessly.
(\emph{a. proverb}) \B{Txo ke nì\-yo' tsa\-krr nì\-yol.} `If you can't be flawless, at least be brief.'~\LINK{http://naviteri.org/2012/01/mipa-zisit-ayliu-amip-new-words-for-the-new-year/}{b}
(\ding{43}~\HL{yo'}{\B{yo'}}) 

% nìyol
\hi
\HT{nIyol}{\B{nì\cp{yol}}} \I{[nI."jol]} \emph{adv.}  briefly, shortly (\emph{of time}). 
(\emph{a. proverb}) \B{Txo ke nì\-yo' tsa\-krr nì\-yol.} `If you can't be flawless, at least be brief.'~\LINK{http://naviteri.org/2012/01/mipa-zisit-ayliu-amip-new-words-for-the-new-year/}{b}
(\ding{43}~\HL{yol}{\B{yol}})

% nìyrr
\hi
\HT{nIyrr}{\B{nì\cp{yrr}}} \I{[nI."j\s{r}]} \emph{adv.} naturally, without tampering with nature, without changing the nature of something. 
\B{Fkxen\-ti pxìm yom fkol nì\-yrr.} Vegetables are often eaten raw.~\LINK{http://naviteri.org/2013/01/awvea-posti-zisita-amip-first-post-of-the-new-year/}{b}
(\ding{43}~\HL{yrr}{\B{yrr}})

% nìzen
\hi
\HT{nIzen}{\B{nì\cp{zen}}} \I{[nI."zEn]} \emph{adv.} necessarily. 
(\ding{43}~\HL{zene}{\B{ze\-ne}}; \ding{214}~\HL{kenzen}{\B{ken\-zen}})

% nìzevakx
\hi 
\HT{nIzevakx}{\B{nì\cp{ze}vakx}} \I{[nI."zE.vak']} \emph{adv.} cruelly. 
\B{Kin\-trr\-am ay\-oe tar\-mìng na\-ri teng\-krr 'aw\-a ho\-ren\-ley\-kek\-yul tu\-tan\-it a\-'aw tspe\-rang--lun\-lu\-ke nì\-wotx, nìk\-'ong, nì\-ze\-vakx.} Last week we watched while one law enforcer killed a man--without any reason, slowly, cruelly.~\LINK{http://naviteri.org/2020/06/mi-tanlokxe-oeya-srr-afpxamo-terrible-days-in-my-country/}{b} 
(\ding{43}~\HL{zevakx}{\B{ze\-vakx}})

% nìzawnong
\hi
\HT{nIzawnong}{\B{nì\cp{zaw}nong}} \I{[nI."zaw.noN]} \emph{adv.} safely. 
\B{Si\-vop nì\-zaw\-nong, ma ay\-sop\-yu.} Travel safely, travellers.~\LINK{http://naviteri.org/2010/09/quick-follow-up/}{b}
\B{Ti\-va\-ron nì\-zaw\-nong, ma ey\-lan!} Hunt safely, friends!~\LINK{http://naviteri.org/2013/08/taronway-the-hunt-song/}{b}
(\emph{a. idiom}) \B{Mi\-vak\-to nì\-zaw\-nong.} `Take care.' `Travel safely.'~\LINK{http://forum.learnnavi.org/language-updates/makto-zong!/}{b}
(\emph{often shortened to} \HL{zong}{\B{zong}}) \emph{in the answer to the question} \B{Mak\-to fya\-pe?} ``How are things? How's everything going? How are you doing?'' \B{Nì\-zaw\-nong.\slash Zong.} ``Well.''~\LINK{http://naviteri.org/2020/05/tipusawm-tiuseyng-si-okvur-a-eltur-titxen-si-asking-answering-and-an-interesting-story/}{b}
(\ding{43}~\HL{zong}{\B{zong}})

% nokx
\hi
\HT{nokx}{\B{nokx}} \I{[nok']} \emph{vtr.} give birth to.
\B{Sa'\-nok\-ìl oe\-yä tsmu\-ket a\-mip no\-lokx trr\-am.} Mom gave birth to my new sister yesterday.~\LINK{http://naviteri.org/2010/09/getting-to-know-you-part-2/}{b} \\  
\textsc{Derivatives}: \ding{43} 
\on{1}. \HL{'ongokx}{\B{'o\-ngokx}}, be born.

% nong
\hi
\HT{nong}{\B{nong}} \I{[noN]} \emph{vtr.} follow, proceed after.
\B{Nong oet!} Follow me!~\LINK{http://naviteri.org/2022/01/aawa-tipangkxotsyip-a-teri-horen-lifyaya-a-few-little-discussions-about-grammar/}{b} \\
\textsc{Derivatives}: \ding{43} 
\on{1}.~\HL{nongspe'}{\B{nong\-spe'}}, persue with an intend to capture.

% nongspe'
\hi 
\HT{nongspe'}{\B{nong\cp{spe'}}} \I{[noN."{}spEP]} \emph{vtr.} persue with an intend to capture. 
\B{Ta\-ron\-yul ye\-rik\-it nar\-mong\-spe', slä tsun ye\-rik hi\-vi\-fwo.} The hunter was pursuing a hexapede, but the hexapede was able to escape.~\LINK{http://naviteri.org/2018/04/100a-liu-amip-64-new-words-part-2/}{b} 
(\ding{43}~\HL{nong}{\B{nong}}, \HL{spe'e}{\B{spe'e}})

% nopx
\hi
\HT{nopx}{\B{nopx}} \I{[nop']} \emph{vtr.} put away, store.
\B{Tsko swi\-zaw\-ti ni\-vopx, ma 'ite. Ke ta\-ron oeng fì\-trr.} Put away your bow and arrow, daughter. You and I are not hunting today.~\LINK{http://naviteri.org/2014/11/twenty-before-the-holidays/}{b}

% nrr
\hi
\HT{nrr}{\B{nrr}} \I{[n\s{r}]} \emph{vin.} glow, be luminous.   
\B{Na'\-rìng sngä\-'i ni\-vrr txon\-krr syu\-ra\-tan\-fa.} The forest begins to glow at night with bioluminescence.~\LINK{http://naviteri.org/2012/07/fivospxiya-aylifyavi-amip-this-months-new-expressions-pt-2/}{b}
\B{Nga\-ri key ne\-re\-i\-yrr, ma 'ey\-lan. Nì\-law po yaw\-ne lu ngar.} Your face is glowing, my friend (and I'm pleased to see it). It's clear you love her.~\LINK{http://naviteri.org/2012/07/fivospxiya-aylifyavi-amip-this-months-new-expressions-pt-2/}{b} \\  
\textsc{Derivatives}: \ding{43} 
\on{1}.~\HL{sAnrr}{\B{sä\-nrr}}, glow, an instance of glowing. 
\on{2}.~\HL{nrra}{\B{nrr\-a}}, pride, feeling of pride.
\on{3}.~\HL{lornusrr}{\B{lor\-nu\-srr}}, radiant.

% nrra
\hi
\HT{nrra}{\B{\cp{nrr}a}} \I{["n\s{r}.a]} \emph{n.} pride, feeling of pride.   
\B{Nga\-ri lu\slash leiu oe\-ru nrra nì\-txan, ma 'ite.} I'm very proud of you, daughter.~\LINK{http://naviteri.org/2012/07/fivospxiya-aylifyavi-amip-this-months-new-expressions-pt-2/}{b} 
(\ding{214}~\HL{lewng}{\B{lewng}}) \\  
\textsc{Derivatives}: \ding{43} 
\on{1}.~\HL{lenrra}{\B{le\-nrra}}, proud, feeling of pride in others. 
\on{2}.~\HL{nInrra}{\B{nì\-nrra}}, proudly, with pride. 
\on{3}.~\HL{snonrra}{\B{sno\-nrra}}, self\hyp{}pride.

% nuä
\hi
\HT{nuA}{\B{\cp{nu}ä}} \I{["nu.\ae]} \emph{adp.}\texttt{+} beyond.   
\B{Aw\-nga kel\-ku si nuä ay\-ram a\-lu\-sìng.} We live beyond the flying mountains.~\LINK{http://naviteri.org/2011/10/more-vocabulary-a-bit-of-grammar/}{b}
\B{Fo kel\-ku si nuä ora.} They live beyond the lake.~\LINK{http://naviteri.org/2011/10/more-vocabulary-a-bit-of-grammar/}{b} \\  
\textsc{Derivatives}: \ding{43} 
\on{1}.~\HL{nuAslew}{\B{nu\-ä\-slew}}, surpass, transcend. 

% nuäslew
\hi 
\HT{nuAslew}{\B{\cp{nu}äslew}} \I{["nu.\ae.sl$\cdot$Ew]} \emph{vtr., inf.}\on{33} surpass, transcend. 
\B{Nge\-yä way\-ìl nuä\-slew pum\-it oe\-yä.} Your song surpasses mine.~\LINK{https://naviteri.org/2026/04/tsivola-liu-amip-thirty-two-new-words/}{b} 
(\ding{43}~\HL{nuA}{\B{nuä}}, \HL{salew}{\B{sa\-lew}}) 

% nui
\hi
\HT{nui}{\B{\cp{nu}i}} \I{["nu.i]} \emph{vin.} (\emph{a.}) fail, falter, go astray (\emph{not obtain expected or desired result})  
\B{Oe fmo\-li nu\-ä\-ngi.} I tried, but unfortunately I failed.~\LINK{http://naviteri.org/2013/04/wheres-the-bathroom-and-other-useful-things/}{b}
\B{Rum\-it a no\-lu\-i rä\-'ä fe\-wi.} Don't chase after a foul ball.~\LINK{http://naviteri.org/2013/04/wheres-the-bathroom-and-other-useful-things/}{b}
(\emph{b.}) mess up, do wrong (\emph{used for placing blame}).
\B{No\-lu\-i \cp{nga}!} \emph{You} failed! It's \emph{your} fault! \emph{You're} the one who messed up.~\LINK{http://naviteri.org/2013/04/wheres-the-bathroom-and-other-useful-things/}{b}
(\ding{214}~\HL{flA}{\B{flä}}, \ding{43}~\HL{tIkxey}{\B{tì\-kxey si}}) \\  
\textsc{Derivatives}: \ding{43} 
\on{1}.~\HL{tInui}{\B{tì\-nui}}, failure (abstract).  
\on{2}.~\HL{sAnui}{\B{sä\-nui}}, failure (particular instance).   
\on{3}.~\HL{nInu}{\B{nì\-nu}}, failingly, falteringly, in vain. 

% nulkrr
\hi
\HT{nulkrr}{\B{nul\cp{krr}}} \I{[nul."k\s{r}]} \emph{or} \I{[nUl.]} \emph{adv.} longer (time). 
\B{Ke tsun aw\-nga pi\-vey nul\-krr--txew lok.} We can't wait any longer--time is almost up.~\LINK{http://naviteri.org/2012/03/spring-vocabulary-part-1/}{b}
\B{Ru\-txe ma\-wey\-pi\-vey nul\-krr nì\-'it, ma ey\-lan.} Please, wait a bit longer, friends.~\LINK{http://naviteri.org/2020/02/some-words-for-leap-year-day/}{b}
(\ding{43}~\HL{krr}{\B{krr}}, \HL{nI'ul}{\B{nì\-'ul}})

% nulnew
\hi
\HT{nulnew}{\B{nul\cp{new}}} \I{[nul."n$\cdot$Ew]} \emph{or} \I{[nUl.]} \emph{vtr.}\on{22} (\emph{a.}) prefer, would rather.
\B{Oel pay\-ti nul\-new to swo\-at.} I prefer water to alcohol.~\LINK{http://forum.learnnavi.org/language-updates/prefer-x-to-y-24474/}{ln}
(\emph{b. modal}) (\emph{either with} ‹\B{iv}› \emph{or} \B{futa} \emph{and} ‹\B{iv}›).
\B{Ke lu kaw\-tu a nul\-ni\-vew oe po\-hu ti\-rea\-pi\-väng\-kxo äo Ut\-ral Ay\-mok\-ri\-yä.} There's nobody I'd rather commune with under the Tree of Voices.~\LINK{http://wiki.learnnavi.org/Corpus\#Lemondrop.com}{ld\slash wiki}
\B{Tsun oe nga\-hu tsa\-tse\-nge\-ne ki\-vä, slä nul\-new fu\-ta si\-vop oeng nì\-txi\-lu\-ke.} I can go there with you, but I prefer to travel leisurely.~\LINK{http://naviteri.org/2011/10/more-vocabulary-a-bit-of-grammar/}{b}
(\ding{43}~\HL{new}{\B{new}}, \HL{nì'ul}{\B{nì\-'ul}})

% nume
\hi
\HT{nume}{\B{\cp{nu}me}} \I{["nu.mE]} \emph{vin.} learn, acquire knowledge or understanding (\HL{teri}{\B{teri}}, about).
\B{Tsko swi\-zaw\-fa a tì\-ta\-ron\-ìri ze\-ne fko ni\-vu\-me.} You must learn to hunt with a bow and arrow.~\LINK{http://forum.learnnavi.org/language-updates/ordinals-nume/}{ln\slash a\textsuperscript{c}}
\B{Maw\-krr la\-ye\-iu oer krr nì\-'ul fte ngi\-vop ay\-lì\-'ut sì tsay\-fne\-sä\-num\-vit a tsun fra\-por srung si\-vi fte ni\-vu\-me sì zi\-ve\-rok nì\-swey.} After that I will have more time to create words and the kinds of lessons that can help everyone best learn and remember.~\LINK{http://forum.learnnavi.org/language-updates/trr-rrtaya/}{ln}
\B{Var ni\-vu\-me ko!} Let's keep learning!~\LINK{http://forum.learnnavi.org/language-updates/as-x-as-y-keep-on-keepin-on/}{ln}
\B{Oel vewng fu\-ta ay\-e\-veng ni\-vu\-me te\-ri ay\-ewll na'\-rìng\-ä.} I see to it that the children learn about the forest plants.~\LINK{http://naviteri.org/2010/09/getting-to-know-you-part-2/}{b} \\
\textsc{Derivatives}: \ding{43}
\on{1}.~\HL{numeyu}{\B{nu\-me\-yu}}, student, learner. 
\on{2}.~\HL{numtseng}{\B{num\-tseng}}, school. 
\on{3}.~\HL{numtsengvi}{\B{num\-tseng\-vi}}, classroom. 
\on{4}.~\HL{numultxa}{\B{nu\-mul\-txa}}, class (for instruction). 
\on{5}.~\HL{numultxatu}{\B{nu\-mul\-txa\-tu}}, classmate. 
\on{6}.~\HL{snanumultxa}{\B{sna\-nu\-mul\-txa}}, course (for instruction).
\on{7}.~\HL{sAnume}{\B{sä\-nu\-me}}, instruction. 
\on{8}.~\HL{sAnumvi}{\B{sä\-num\-vi}}, lesson.

% numeyu
\hi
\HT{numeyu}{\B{\cp{nu}meyu}} \I{["nu.mE.ju]} \emph{n.} student, learner. 
\B{Oe lu nu\-me\-yu.} I am a student.~\LINK{http://naviteri.org/2010/09/getting-to-know-you-part-2/}{b}
\B{Kar\-yul fngo\-lo' fu\-ta ay\-nu\-me\-yu pi\-va\-te ye'\-krr.} The teacher required the students to arrive early.~\LINK{http://naviteri.org/2012/01/mipa-zisit-ayliu-amip-new-words-for-the-new-year/}{b}
(\ding{43}~\HL{nume}{\B{nu\-me}})

% numtseng
\hi
\HT{numtseng}{\B{\cp{num}tseng}} \I{["num.\t{ts}EN]} \emph{or} \I{["nUm.]} \emph{n.} school. 
\B{Tsun Txe\-wì kop pi\-vll\-txe nì\-'it nì\-'Ìng\-lì\-sì a no\-lu\-me mì num\-tseng tok\-tor Kì\-rey\-sä.} Txewì can also speak a bit of English that he has learned in doctor Grace's school.~\LINK{http://naviteri.org/2010/07/ting-mikyun-fte-tslivam-listening-comprehension-1-3/}{b}
(\ding{43}~\HL{nume}{\B{nu\-me}}, \HL{tseng}{\B{tse\-nge}}) \\  
\textsc{Derivatives}: \ding{43} 
\on{1}.~\HL{numtsengvi}{\B{num\-tseng\-vi}}, classroom.
\on{2}.~\HL{nawnumtseng}{\B{naw\-num\-tseng}}, university.

% numtsengvi
\hi
\HT{numtsengvi}{\B{\cp{num}tsengvi}} \I{["num.\t{ts}EN.vi]} \emph{or} \I{["nUm.]} \emph{n.} classroom, division of a school.
\B{num\-tseng\-vi ay\-sngä\-'i\-yu\-ä}, beginners' classroom.~\LINK{http://naviteri.org/2010/06/first-post/}{b}
\B{Zo\-la\-'u ftu kä\-pxì num\-tseng\-vi\-yä sä\-mewn a\-wok.} From the back of the classroom came a loud sigh.~\LINK{https://naviteri.org/2025/04/mevosinga-liu-amip-twenty-new-words/}{b}
\B{Kar\-yu a\-sìl\-tsan ze\-ne tsi\-vun ay\-nu\-me\-yut hi\-vipx mì num\-tseng\-vi.} A good teacher has to be able to control (his/her) students in the classroom.~\LINK{http://naviteri.org/2018/04/100a-liu-amip-64-new-words-part-2/}{b}
(\ding{43}~\HL{numtseng}{\B{num\-tseng}})

% numultxa
\hi
\HT{numultxa}{\B{numul\cp{txa}}} \I{[nu.mul."t'a]} \emph{or} \I{[.mUl.]} \emph{n.} (\textsc{rn}: \B{numùl\cp{txa}} \I{[nu.mUl."da]}) class (for instruction).   
\B{Te\-ri lì'\-fya le\-Na'\-vi a tsa\-nu\-mul\-txa lo\-le\-i\-u sä\-flä!} The class about the Na'vi language was a success.~\LINK{http://naviteri.org/2012/07/tskxekengtsyip-a-mikyunfpi-a-little-listening-exercise/}{b}
\B{Nì\-aw\-no\-mum ngo\-lop ay\-oel fì\-nu\-mul\-txa\-ti fpi sngä\-'i\-yu.} As you know, we created this class for beginners.~\LINK{http://naviteri.org/2012/07/tskxekengtsyip-a-mikyunfpi-a-little-listening-exercise/}{b}
\B{Mì Si\-ä\-tll a nu\-mul\-txa\-ri lì'\-fya\-yä le\-Na'\-vi, sìl\-pey oe slì\-ye\-vu nga nu\-mul\-txa\-tu!} I hope you'll participate in the upcoming Na'vi class in Seattle!~\LINK{http://naviteri.org/2012/07/meetings-waterfalls-and-more/}{b}
(\ding{43}~\HL{nume}{\B{nu\-me}}, \HL{ultxa}{\B{ul\-txa}}) \\  
\textsc{Derivatives}: \ding{43} 
\on{1}.~\HL{numultxatu}{\B{nu\-mul\-txa\-tu}}, classmate, member of a class.

% numultxatu
\hi
\HT{numultxatu}{\B{numul\cp{txa}tu}} \I{[nu.mul."t'a.tu]} \emph{or} \I{[.mUl.]} (\textsc{rn}: \B{numùl\cp{txa}tu} \I{[nu.mUl."da.tu]}) \emph{n.} classmate, member of a class.
\B{Mì Si\-ä\-tll a nu\-mul\-txa\-ri lì'\-fya\-yä le\-Na'\-vi, sìl\-pey oe slì\-ye\-vu nga nu\-mul\-txa\-tu!} I hope you'll participate in the upcoming Na'vi class in Seattle!~\LINK{http://naviteri.org/2012/07/meetings-waterfalls-and-more/}{b}
(\ding{43}~\HL{numultxa}{\B{nu\-mul\-txa}})

% nun
\hi 
\HT{nun}{\B{nun}} \I{[nun]} \emph{or} \I{[nUn]} (\textsc{rn}: \B{nun} \I{[nun]}) \emph{adj.} long (\emph{of time}). 
\B{Su\-nu oer Ra\-lu, slä fì\-sä\-frr\-fen pe\-yä lu nun nì\-hawng.} I like Ralu, but this visit of his is too long.~\LINK{http://naviteri.org/2024/02/mipa-ayliu-si-aylifyavi-niul-more-new-words-and-expressions/}{b} 
(\ding{214}~\HL{yol}{\B{yol}}) \\
\textsc{Derivatives}: \ding{43} 
\on{1}.~\HL{zIsItnun}{\B{zì\-sìt\-nun}}, leap year. 
\on{2}.~\HL{nunyol}{\B{nun\-yol}}, length (of time).

% nunyol
\hi 
\HT{nunyol}{\B{\cp{nun}yol}} \I{["nun.yol]} \emph{or} \I{["nUn.]} (\textsc{rn}: \B{\cp{nun}yol} \I{["nun.yol]}) \textsc{i}. \emph{n.} length (of time). \\
\textsc{ii}. \B{\cp{nun}yolpe} \I{["nun.yol.pE]} \emph{or} \B{pe\cp{nun}yol} \I{[pE."nun.yol]} \emph{inter.} what length? how long (of time)?
\B{Nga har\-ma\-haw pe\-nun\-yol?} How long were you sleeping?~\LINK{https://naviteri.org/2024/02/mipa-ayliu-si-aylifyavi-niul-more-new-words-and-expressions/}{b}
(\ding{43}~\HL{nun}{\B{nun}}, \HL{yol}{\B{yol}}; \HL{ngimpup}{\B{ngim\-pup}}) 

% nutx
\hi
\HT{nutx}{\B{nutx}} \I{[nut']} \emph{or} \I{[nUt']} \emph{adj., nfp.} thick.   
\B{Tsun Txil\-te pam\-rel\-it ivi\-nan; ta\-fral puk\-ot a\-nutx mu\-nge fra\-tseng.} Txilte knows how to read; therefore she brings a thick book wherever she goes.~\LINK{http://naviteri.org/2011/11/\%E2\%80\%99a\%E2\%80\%99awa-ayli\%E2\%80\%99u-amip-ni\%E2\%80\%99aw-\%E2\%80\%94-a-few-new-words-only/}{b}
\B{Krro krro, flì\-a vul a\-ru\-sey to nutx\-a pum a\-ke\-ru\-sey lu txur.} A thin living branch is sometimes stronger than a thick dead one.~\LINK{http://naviteri.org/2011/11/\%E2\%80\%99a\%E2\%80\%99awa-ayli\%E2\%80\%99u-amip-ni\%E2\%80\%99aw-\%E2\%80\%94-a-few-new-words-only/}{b}
(\ding{214}~\HL{flI}{\B{flì}}; \ding{43}~\HL{ompu}{\B{om\-pu}}) \\  
\textsc{Derivatives}: \ding{43} 
\on{1}.~\HL{flInutx}{\B{flì\-nutx}}, thickness.

\end{multicols}

% ===== Section NG ===== 

 \pdfbookmark[0]{NG}{NG}
\begin{center}
\section*{NG \I{[N]} (\emph{NgeNg})}
\end{center}

\begin{multicols}{2}

% nga
\hi
\HT{nga}{\B{nga}} \I{[Na]} \emph{pn.} (\emph{gen.} \B{\cp{nge}yä}) you (2\textsuperscript{nd} person singular). 
\B{Lu nga win sì txur | Lu nga txan\-tslu\-sam.} You are fast and strong | You are wise.~\LINK{http://naviteri.org/2013/08/taronway-the-hunt-song/}{b}
\B{Nga lu So\-rewn, ke\-fyak?} You're Sorewn, right?~\LINK{http://naviteri.org/2010/09/getting-to-know-you-part-1/}{b}
\B{Tew\-ti, nga lu tsul\-fä\-tu i\-'en\-ä.} Wow, you are a master on the \emph{i'en}.~\LINK{http://naviteri.org/2011/01/new-year-new-vocabulary/}{b}
\B{Pll\-txe nga nìl\-tsan!} You speak well!~\LINK{http://languagelog.ldc.upenn.edu/nll/?p=1977}{ll}
\B{Ngal fì\-'u\-pxa\-ret oe\-ru fpo\-le'.} You have sent me this message.~\LINK{http://forum.learnnavi.org/language-updates/a-few-small-things/msg299949/\#msg299949}{ln}
\B{Oel nga\-ti ka\-me\-i\-e.} I \textsc{see} you.~\LINK{http://wiki.learnnavi.org/index.php?title=Corpus\#UGO_Movie_Blog}{wiki}
\B{Ey\-wal zey\-ki\-vo ngat nì\-win.} May Eywa heal you quickly.~\LINK{http://naviteri.org/2010/07/vocabulary-update/}{b}
\B{Nga\-ru ira\-yo sei\-yi oe nì\-mun.} I (happily) thank you again.~\LINK{http://wiki.learnnavi.org/index.php?title=Canon\#Extracts_from_various_emails}{wiki}
\B{Oel mu\-wi\-vìn\-txu nga\-ru oe\-yä ler\-tut.} Let me present my colleague to you (\emph{formal}).\slash \B{Nga\-ru oe\-yä ler\-tut.} This is my colleague (\emph{coll.}).~\LINK{http://naviteri.org/2010/09/getting-to-know-you-part-1/comment-page-1/\#comment-297}{b}
\B{Fya\-pe fko syaw ngar?} What's your name? [lit.: how does one call you?]~\LINK{http://forum.learnnavi.org/language-updates/the-reflexive-\%28from-a-frommer-email!!!!\%29/}{ln}
\B{Tì\-ste\-ftxaw nge\-yä flo\-lä nìl\-tsan!} Your examination succeeded well!~\LINK{http://wiki.learnnavi.org/index.php?title=Canon\#Extracts_from_various_emails}{wiki} 
\B{Lu nge\-yä tì\-ftxey!} It's your choice!~\LINK{http://wiki.learnnavi.org/index.php?title=Canon\#Extracts_from_various_emails}{wiki} 
\B{Nga\-ri txe'\-lan ma\-wey li\-vu.} ``Don't worry (about it).''~\LINK{http://forum.learnnavi.org/language-updates/fmawno/}{ln}
\B{Nga\-ri so\-la\-lew pol\-pxay\-a zì\-sìt?} How old are you?~\LINK{http://naviteri.org/2010/09/getting-to-know-you-part-2/}{b}
\B{Tsun oe nga\-hu nì\-Na'\-vi pi\-väng\-kxo a fì\-'u oe\-ru prr\-te' lu.} It's a pleasure to be able to chat with you in Na'vi.~\LINK{http://wiki.learnnavi.org/index.php?title=Corpus\#Times_Online}{wiki}
\B{Ey\-wa nga\-hu.} Eywa be with you.~\LINK{http://forum.learnnavi.org/news-announcements/a-response-from-paul-frommer!/}{ln}
\B{Tul hu nga | Na'\-rìng\-ka nì\-'o'.} Run with you | Through the forest in a fun way.~\LINK{http://forum.learnnavi.org/language-updates/learnings-from-the-siatlla-song-translations/msg473566/\#msg473566}{ln}
(\emph{a. conv. phrase}) \B{Nga\-fkeyk pe\-fya\slash fya\-pe?} `What's your status?\slash How ya doin'?'~\LINK{http://naviteri.org/2018/02/negative-questions-in-navi/\#comment-28201}{b} \\
\textsc{Derivations}: \ding{43} 
\on{1}.~\HL{ngenga}{\B{nge\-nga}}, you (ceremonial).
\on{2}.~\HL{oeng}{\B{oeng}}, us two (inclusive).
\on{3}.~\HL{manga}{\B{ma\-nga}}, ``hey you''. 

% nga'
\hi
\HT{nga'}{\B{nga'}} \I{[NaP]} \emph{vtr.} (\emph{a.}) contain.   
\B{Na'\-rìng\-ìl nga' pxay\-a io\-ang\-it.} The forest contains many animals.~\LINK{http://naviteri.org/2011/05/weather-part-2-and-a-bit-more-2/}{b}
(\emph{b.}) be pregnant with. 
\B{Pol pxey\-a prr\-nen\-it nge\-i\-a'.} I'm delighted to say she's pregnant with triplets.~\LINK{http://naviteri.org/2024/02/mipa-ayliu-si-aylifyavi-niul-more-new-words-and-expressions/}{b}
\B{Krr\-a ngal oe\-ti ngar\-ma', 'e\-fu pe\-fya?} How did you feel when you were pregnant with me?~\LINK{http://naviteri.org/2024/02/mipa-ayliu-si-aylifyavi-niul-more-new-words-and-expressions/}{b} \\
\textsc{Derivations}: \ding{43} 
\on{1}.~\HL{nga'prrnen}{\B{nga'prrnen}}, be pregnant (for people).
\on{2}.~\HL{nga'lini}{\B{nga'lini}}, be pregnant (for animals).

% -nga'
\hi
\HT{-nga'}{\B{-nga'}} \I{[.NaP]} \emph{unprod. suf. to form adj.s from n.s with the meaning} containing, -ful; e.g., \B{meuia\-nga'}, honorable, honorably. \B{ral\-nga'}, meaningful, instructive. \B{tì\-tstew\-nga'}, courageous, brave, etc.

% nga'lini
\hi 
\HT{nga'lini}{\B{nga'\cp{li}ni}} \I{[NaP."li.ni]} \emph{vin.} be pregnant (\emph{for animals}). 
(\ding{43}~\HL{nga'}{\B{nga'}}, \HL{lini}{\B{li\-ni}}) \\
\textsc{Derivations}: \ding{43} 
\on{1}.~\HL{tInga'lini}{\B{tì\-nga'\-li\-ni}}, pregnancy (\emph{for animals}).

% nga'prrnen
\hi 
\HT{nga'prrnen}{\B{nga'\cp{prr}nen}} \I{[N$\cdot$aP."p\s{r}.nEn]} \emph{vin., ofp., inf.}\on{11} (\emph{a}.) be pregnant. 
\B{Zun ngal oey tsmu\-ket tsi\-ve\-'a, zel am\-'a\-lu\-ke i\-vo\-mum fu\-ta poe nga'\-prr\-nen.} If you saw my sister, you'd certainly know she was pregnant.~\LINK{http://naviteri.org/2024/02/mipa-ayliu-si-aylifyavi-niul-more-new-words-and-expressions/}{b}
(\emph{b.}) be pregnant \emph{with}, use \ding{43}~\HL{nga'}{\B{nga'}}, e.g. \B{Pol pxey\-a prr\-nen\-it nge\-i\-a'.} I'm delighted to say she's pregnant with triplets.~\LINK{http://naviteri.org/2024/02/mipa-ayliu-si-aylifyavi-niul-more-new-words-and-expressions/}{b}
(\ding{43}~\HL{nga'}{\B{nga'}}, \HL{prrnen}{\B{prr\-nen}}) \\
\textsc{Derivations}: \ding{43}
\on{1}.~\HL{tInga'prrnen}{\B{tì\-nga'\-prr\-nen}}, pregnancy (\emph{for people}).

% ngam
\hi
\HT{ngam}{\B{ngam}} \I{[Nam]} \emph{n.} echo.   
\B{Fì\-slär\-mì tsun fko sti\-vawm ngam\-it a\-pxay.} You can hear a lot of echoes in this cave.~\LINK{http://naviteri.org/2012/01/mipa-zisit-ayliu-amip-new-words-for-the-new-year/}{b} \\
\textsc{Derivations}: \ding{43}
\on{1}.~\HL{ngampam}{\B{ngam\-pam}}, rhyme.

% ngampam
\hi
\HT{ngampam}{\B{\cp{ngam}pam}} \I{["Nam.pam]} \textsc{i}. \emph{n.} rhyme.
\B{Fì\-way\-ri hì\-no\-a re\-nut ngam\-pam\-ä ke tsä\-ngun oe tsli\-vam.} I'm afraid I can't understand the intricate rhyme scheme of this poem.~\LINK{http://naviteri.org/2012/01/mipa-zisit-ayliu-amip-new-words-for-the-new-year/}{b} \\
\textsc{ii}. \B{\cp{ngam}pam si} \I{["Nam.pam si]} \emph{vin.} (\emph{a.}) rhyme.   
\B{Me\-lì\-'u alu mung\-wrr sì nì\-fkrr ngam\-pam si.} The words \emph{mung\-wrr} and \emph{nì\-fkrr} rhyme.~\LINK{http://naviteri.org/2012/01/mipa-zisit-ayliu-amip-new-words-for-the-new-year/}{b}
(\emph{b. fig.}) get along, fit together well. 
\B{New Ri\-ni sì Ra\-lu mun\-txa sli\-vu, slä tì\-'e\-fu\-mì oe\-yä, ngam\-pam ke si.} Rini and Ralu want to marry, but I feel they're not compatible.~\LINK{http://naviteri.org/2012/01/mipa-zisit-ayliu-amip-new-words-for-the-new-year/}{b}
(\ding{43}~\HL{ngam}{\B{ngam}}, \HL{pam}{\B{pam}})

% ngawng
\hi
\HT{ngawng}{\B{ngawng}} \I{[NawN]} \emph{n., fauna} worm. 
\B{Ma txan\-tslu\-sam\-a ngawng, a swok\-a ut\-ral\-ti yom.} O wise worm, who eats the Sacred Tree.~\LINK{https://wiki.learnnavi.org/Na\%27vi_from_Avatar_Movie\#Becoming_one_of_the_Na.27vi}{a\textsuperscript{c}}  \\
\textsc{Derivations}: \ding{43} 
\on{1}.~\HL{eltungawng}{\B{el\-tu\-ngawng}}, brainworm.

% ngay
\hi
\HT{ngay}{\B{ngay}} \I{[Naj]} \emph{adj.} true.   
\B{Oe\-yä txe'\-lan li\-vu ngay.} Let my heart be true.~\LINK{http://wiki.learnnavi.org/index.php/Hunting_Song}{wiki} 
\B{Ke\-'u ke lu ngay.} Nothing is true.~\LINK{http://forum.learnnavi.org/language-updates/double-negatives-not-optional/}{ln} 
\B{Lam ngay oer, fo ka\-yä ìlä hil\-van tup na'\-rìng.} I bet they'll take the river rather than the forest.~\LINK{http://forum.learnnavi.org/language-updates/concerning-tup/}{ln\slash g} 
\B{Tsyan\-ìri sì oe\-ri lo\-le\-iu tsa\-kaym tsyeym a\-ngay.} There was a true treasure that evening for John and I.~\LINK{http://naviteri.org/2014/08/stxeli-alor-a-beautiful-gift/}{b} \\  
\textsc{Derivations}: \ding{43}
\on{1}.~\HL{nIngay}{\B{nì\-ngay}}, truly.  
\on{2}.~\HL{tIngay}{\B{tì\-ngay}}, truth.   
\on{3}.~\HL{ngaytxoa}{\B{ngay\-txoa}}, sorry, my apologies.  
\on{4}.~\HL{pllngay}{\B{pll\-ngay}}, admit. 
\on{5}.~\HL{nIfkeytongay}{\B{nì\-fkey\-to\-ngay}}, actually, as a matter of fact. 
\on{6}.~\HL{wIngay}{\B{wì\-ngay}}, prove.
\on{7}.~\HL{pe'ngay}{\B{pe'\-ngay}}, judge, conclude.

% ngaytxoa
\hi
\HT{ngaytxoa}{\B{ngay\cp{txo}a}} \I{[Naj."t'o.a]} \emph{idiom} forgive me, sorry, my apologies (\emph{acknowledgment of guilt and regret}).
\B{Ngay\-txo\-a, nì\-aw\-no\-mum ke lo\-lu oer nì\-ke\-ftxo mì sok\-a srr ay\-skxom le\-tam fte lì'\-fya\-ri aw\-nge\-yä tì\-kang\-kem si\-vi.} My apologies: As you know, in recent days I have not had sufficient opportunity to work on our language.~\LINK{http://forum.learnnavi.org/language-updates/trr-rrtaya/}{ln}
\B{Kxey\-ey\-tsyìp\-ìri ngay\-txo\-a.} My opologies for the little mistake.~\LINK{http://naviteri.org/2013/02/vospxi-ayol-posti-apup-short-post-for-a-short-month/comment-page-1/\#comment-2108}{b}
(\ding{43}~\HL{ngay}{\B{ngay}}, \HL{txoa}{\B{txoa}})

% ngä'än 
\hi
\HT{ngA'An}{\B{ngä\cp{'än}}} \I{[N\ae."P\ae{}n]} \emph{vin.} suffer (mentally, emotionally), be miserable.
\B{Sna\-fpìl\-fya\-ri le\-Na'\-vi krra smar\-it fkol tspang, tsran\-ten nì\-txan fwa po ke ngä\-'än nì\-kel\-kin.} It's important in Na'vi philosophy that the prey not suffer unnecessarily when it's killed.~\LINK{http://naviteri.org/2013/02/vospxi-ayol-posti-apup-short-post-for-a-short-month/}{b} 
\B{Sra\-ne, skxir tì\-sraw si nì\-txan, slä ke nge\-rä\-'än oe kaw\-'it.} Yes, the wound is very painful, but I'm not in the least suffering emotionally.~\LINK{http://naviteri.org/2013/02/vospxi-ayol-posti-apup-short-post-for-a-short-month/}{b}
\B{Oe\-ri ke lu tìm\-wiä, ha nga\-ti oel ngey\-ka\-syä\-'än nì\-teng.} For me there isn't justice, so I'll make you be miserable as well.~\LINK{http://naviteri.org/2020/06/mi-tanlokxe-oeya-srr-afpxamo-terrible-days-in-my-country/}{b} \\  
\textsc{Derivations}: \ding{43}
\on{1}.~\HL{sAngA'An}{\B{sä\-ngä\-'än}}, bout of suffering, episode of drepression.

% ngäng
\hi
\HT{ngAng}{\B{ngäng}} \I{[N\ae{}N]} \emph{n.} stomach. 

% ngäzìk
\hi
\HT{ngAzIk}{\B{\cp{ngä}zìk}} \I{["N\ae.zIk\textcorner]} \emph{adj.} difficult, hard.
\B{Fwam\-pop\-ä tem\-rey tsa\-tseng\-mì ngä\-zìk lu nì\-wotx.} It's very hard for a tapirus to survive there.~\LINK{http://naviteri.org/2011/05/some-miscellaneous-vocabulary/}{b}
\B{Kxawm set 'u a\-ngä\-zìk fra\-to lu la\-'a a\-yll.} Maybe the most difficult thing now is social distance.~\LINK{http://naviteri.org/2020/03/teri-tifkeytok-lefkrr-about-the-present-situation/}{b}
(\ding{214}~\B{ftue}) \\  
\textsc{Derivations}: \ding{43}
\on{1}.~\HL{tIngAzIk}{\B{tìngäzìk}}, difficulty, problem.

% ngenga
\hi
\HT{ngenga}{\B{nge\cp{nga}}} \I{[NE."Na]} \emph{pn.} (\emph{gen.} \B{nge\-{\color{Plum}nge}\-yä}) you. (2\textsuperscript{nd} person singular; deferential\slash ceremonial form).
\B{Kal\-txì. Nge\-nga lu tu\-pe\slash pe\-su?--Oe lu Txe\-wì.} Hello. Who are you?--I'm Txewì.~\LINK{http://naviteri.org/2010/09/getting-to-know-you-part-3/}{b}
\B{Ä\-txä\-le su\-yi o\-he pi\-vawm, pe\-o\-lo' lu\-yu pum nge\-nge\-yä?} May I ask what tribe you belong to? (\emph{overly formal, stilted Na'vi})~\LINK{http://naviteri.org/2010/09/getting-to-know-you-part-2/}{b}
\B{Ey\-wa nge\-nga\-hu!--Ul\-te nge\-nga\-hu, ma 'E\-veng.} Eywa be with you!--And with you, 'Eveng.~\LINK{http://naviteri.org/2011/08/reported-speech-reported-questions/comment-page-1/\#comment-1069}{b}
(\ding{43}~\HL{nga}{\B{nga}})

% ngep
\hi
\HT{ngep}{\B{ngep}} \I{[NEp\textcorner]} \emph{n.} navel.

% nget
\hi
\HT{nget}{\B{nget}} \I{[NEt\textcorner]} \emph{adj.} smell of decaying wood and leaves.

% ngeyn
\hi
\HT{ngeyn}{\B{ngeyn}} \I{[NEjn]} \emph{adj.} tired. 
\B{Po 'e\-fu ngeyn ul\-te kin tì\-rawn\-it nì\-txan.} He is tired and very much needs to be replaced.~\LINK{http://naviteri.org/2012/01/mipa-zisit-ayliu-amip-new-words-for-the-new-year/}{b}
(\ding{214}~\HL{txen}{\B{txen}})

% ngian
\hi
\HT{ngian}{\B{ngi\cp{an}}} \I{[Ni."{}an]} \emph{adv.} however. 
\B{Ay\-lì\-'u ngi\-an nì\-'it ske\-pek lu.} However the words are a bit formal.~\LINK{http://forum.learnnavi.org/language-updates/more-grammatical-tidbits/}{ln\slash a}
\B{Fì\-kxey\-ey\-ri tsap\-'a\-lu\-te. Ngi\-an la\-he\-a lì\-'u alu Pol ke lu kxey\-ey.} Apologies for this mistake. However the other word `pol' is not a mistake.~\LINK{http://naviteri.org/2011/02/new-vocabulary-ii\%E2\%80\%94part-1/comment-page-1/\#comment-520}{b}

% ngim
\hi
\HT{ngim}{\B{ngim}} \I{[Nim]} \emph{adj.} long (\emph{physical length}).
\B{Ay\-upxa\-re\-ri a\-ngim nì\-hawng lu alo oe\-yä!} Now it's my turn for messages that are too long!~\LINK{http://forum.learnnavi.org/language-updates/txelanit-hivawl/}{ln}
(\ding{214}~\B{pup}) \\  
\textsc{Derivations}: \ding{43}
\on{1}.~\HL{ngimpup}{\B{ngim\-pup}}, length.

% ngimpup
\hi
\HT{ngimpup}{\B{ngim\cp{pup}}} \I{[Nim."pup\textcorner]} \emph{or} \I{[."pUp\textcorner]} (\textsc{rn}: \B{ngim\cp{pùp}} \I{[Nim."pUp\textcorner]}) \textsc{i}. \emph{n.} length (of physical extension). \\
\textsc{ii}. \B{\cp{ngim}puppe} \I{["Nim.pup\textcorner.pE]} \emph{or} \I{[."pUp\textcorner.]} \emph{or} \HL{pengimpup}{\B{pe\cp{ngim}pup}} \emph{or} \I{[."pUp\textcorner]} \emph{inter.} what length? how long (of physical extension)?
(\ding{43}~\HL{ngim}{\B{ngim}}, \HL{pup}{\B{pup}}; \HL{nunyol}{\B{nun\-yol}})

% ngip
\hi
\HT{ngip}{\B{ngip}} \I{[Nip\textcorner]} \emph{n.} space, open or borderless area.
\B{Nge\-yä ik\-ran\-ìl ngip\-it le\-tam kin fte tsi\-vun kll\-pi\-vä.} Your ikran needs enough open space to be able to land.~\LINK{http://naviteri.org/2015/03/kap-si-ayunil-saylahe-threats-dreams-and-other-things/}{b}
\B{Pll\-txe Saw\-tu\-te san ki\-fkey\-ìl ay\-oe\-yä tok txan\-a ngip\-it a san\-hì\-kip.} The Sky People say that their world is in the great space among the stars.~\LINK{http://naviteri.org/2015/03/kap-si-ayunil-saylahe-threats-dreams-and-other-things/}{b} \\
\textsc{Derivations}: \ding{43}
\on{1}.~\HL{pxawngip}{\B{pxawngip}}, environment.

% ngoa
\hi 
\HT{ngoa}{\B{\cp{ngo}a}} \I{["No.a]} \emph{n.} mud.
\B{Maw\-fwa zup tom\-pa, lu ngo\-a a\-txan fì\-ha\-pxì\-mì na'\-rìng\-ä.} After rain, there's a lot of mud in this part of the forest.~\LINK{http://naviteri.org/2014/07/tsawlultxamaw-a-ayliu-post-meetup-words/}{b}

% ngong
\hi
\HT{ngong}{\B{ngong}} \I{[NoN]} \emph{adj.} (\emph{a.}) lethargic, lacking sufficient energy.    
\B{Ftu\-e lu fwa ta\-ron ngong\-a io\-ang\-it to fwa ta\-ron pum\-it a lu wa\-lak sì win.} It's easier to hunt lethargic animals than to hunt perky, speedy ones.~\LINK{http://naviteri.org/2011/10/more-vocabulary-a-bit-of-grammar/}{b}
(\emph{b. of people with pejorative force}) lazy.
\B{Fì\-trr oe 'e\-fu ngong nì\-wotx.} I'm just not motivated to do anything today.~\LINK{http://naviteri.org/2011/10/more-vocabulary-a-bit-of-grammar/}{b}
(\ding{214}~\B{walak}) \\  
\textsc{Derivations}: \ding{43}
\on{1}.~\HL{nIngong}{\B{nì\-ngong}}, lethargically, lazily. 
\on{2}.~\HL{tIngong}{\B{tì\-ngong}}, lethargy, laziness.

% ngop
\hi
\HT{ngop}{\B{ngop}} \I{[Nop\textcorner]} \emph{vtr.} create.
\B{Pam\-tse\-ol ngop ay\-re\-nut}, Music creates patterns.~\LINK{http://www.pandorapedia.com/navi/music/ritual_music}{pp}
\B{Maw\-krr la\-ye\-iu oer krr nì\-'ul fte ngi\-vop ay\-lì\-'ut sì tsay\-fne\-sä\-num\-vit a tsun fra\-por srung si\-vi fte ni\-vu\-me sì zi\-ve\-rok nì\-swey.} After that I will have more time to create words and the kinds of lessons that can help everyone best learn and remember.~\LINK{http://forum.learnnavi.org/language-updates/trr-rrtaya/}{ln} 
\B{Nì\-aw\-no\-mum ngo\-lop ay\-oel fì\-nu\-mul\-txa\-ti fpi sngä\-'i\-yu.} As you know, we created this class for beginners.~\LINK{http://naviteri.org/2012/07/tskxekengtsyip-a-mikyunfpi-a-little-listening-exercise/}{b}
\B{Fra\-'u\-ä ran ngä\-pop fa fra\-syon tse\-yä.} The \emph{ran} of each thing arises from the totality of its attributes.~\LINK{http://naviteri.org/2012/10/mipa-vospxi-mipa-ayliu-new-words-for-the-new-month/}{b} \\  
\textsc{Derivations}: \ding{43}
\on{1}.~\HL{'ongop}{\B{'o\-ngop}}, design (form, shape). 
\on{2}.~\HL{ngopyu}{\B{ngop\-yu}}, creator.   
\on{3}.~\HL{tIngop}{\B{tì\-ngop}}, creation. 
\on{4}.~\HL{rengop}{\B{re\-ngop}}, design (finer detail).   
\on{5}.~\HL{ronsrelngop}{\B{ron\-srel\-ngop}}, imagine.

% ngopyu
\hi
\HT{ngopyu}{\B{\cp{ngop}yu}} \I{["Nop\textcorner.ju]} \emph{n.} creator. 
\B{Ul\-te set … mok\-ri ngop\-yu\-ä!} And now … the voice of the creator.~\LINK{http://naviteri.org/2012/10/ulte-sawasultsyipa-yoratu-leiu-and-the-winner-is/}{b} 
(\ding{43}~\B{ngop}) \\
\textsc{Derivations}: \ding{43}
\on{1}.~\HL{pamtseongopyu}{\B{pam\-tseo\-ngop\-yu}}, music creator, composer.

% ngrr
\hi
\HT{ngrr}{\B{ngrr}} \I{[N\s{r}]} \emph{n.} \ding{96} root. \\
\textsc{Derivations}: \ding{43}
\on{1}.~\HL{ngrrpongu}{\B{ngrr\-po\-ngu}}, grassroots movement.
\on{2}.~\HL{tsawlapxangrr}{\B{tsawl\-a\-pxa\-ngrr}}, unidelta tree.
\on{3}.~\HL{ngrrfpIl}{\B{ngrr\-fpìl}}, key assumption.

% ngrrfpìl 
\hi 
\HT{ngrrfpIl}{\B{\cp{ngrr}fpìl}} \I{["N\s{r}.fpIl]} \emph{n.} key assumption (\emph{a basic assumption that underlies a presentation or argument}). 
\B{Nì\-law lu pe\-yä ngrr\-fpìl fwa Saw\-tu\-te ke lu mal.} His assumption is clearly that the Skypeople can't be trusted.~\LINK{http://naviteri.org/2021/03/mipa-ayliu-si-aylifyavi-new-words-and-expressions/}{b}
(\ding{43}~\HL{ngrr}{\B{ngrr}}, \HL{sAfpIl}{\B{sä\-fpìl}}; \HL{txinfpIl}{\B{txin\-fpìl}})

% ngrrpongu
\hi
\HT{ngrrpongu}{\B{\cp{ngrr}pongu}} \I{["N\s{r}.po.Nu]} \emph{n.} grassroots movement, popular movement.
\B{Lì'\-fya\-ri le\-Na'\-vi 'Rr\-ta\-mì, vay set 'al\-mong a fra\-'u ze\-ra\-'u ta ngrr\-po\-ngu.} Everything that has gone on with (blossomed regarding) Na'vi until now on Earth has come from a grassroots movement.~\LINK{http://wiki.learnnavi.org/index.php?title=Corpus\#Good_Morning_America}{wiki}
(\ding{43}~\HL{ngrr}{\B{ngrr}}, \HL{pongu}{\B{po\-ngu}})

% ngul
\hi
\HT{ngul}{\B{ngul}} \I{[Nul]} \emph{or} \I{[NUl]} \emph{adj.} drab-colored, grey. 
\B{ngul na tskxe}, the drab color of stone.~\LINK{http://naviteri.org/2010/08/mipa-ayopin-mipa-ayliu-new-colors-new-words/}{b} \\ 
\textsc{Derivations}: \ding{43} 
\on{1}.~\HL{ngulpin}{\B{ngul\-pin}}, (the color) grey.

% ngulpin
\hi
\HT{ngulpin}{\B{\cp{ngul}pin}} \I{["Nul.pin]} \emph{or} \I{["NUl.]} \emph{n.} the color \ding{43}~\HL{ngul}{\B{ngul}}.
(\ding{43}~\HL{'opin}{\B{'opin}})

% ngungung
\hi 
\HT{ngungung}{\B{\cp{ngu}ngung}} \I{["Nu.NuN]} \emph{or} \I{[.NUN]} (\textsc{rn}: \B{\cp{ngù}ngùng} \I{["NU.NUN]}) \emph{vtr.} rub. 
\B{Nga\-ri pxun\-ti ngu\-ngung pe\-lun? Sra\-ke tì\-sraw si?} Why are you rubbing your arm? Does it hurt?~\LINK{http://naviteri.org/2020/07/vospxivol-lefpom-happy-august/}{b}

% nguway
\hi
\HT{nguway}{\B{\cp{ngu}way}} \I{["Nu.waj]} \emph{n.} howl, viperwolf cry. \\
\textsc{Derivations}: \ding{43}
\on{1}.~\HL{ronguway}{\B{ro\-ngu\-way}}, howl.

\end{multicols} 
 
% ===== Section O ===== 

 \pdfbookmark[0]{O}{O}
\begin{center}
\section*{O \I{[o]}}
\end{center}

\begin{multicols}{2}

% -o
\hi
\HT{-o}{\B{-o}} \I{[.o]} \emph{suf.} `indefinite -o'. 
(\emph{a. on regular nouns to show indefiniteness}) 
\B{Txo\-pu rä\-'ä si, lu ke\-tu\-wong\-o nì\-'aw.} Don't be afraid, it's just an alien.~\LINK{http://forum.learnnavi.org/language-updates/a-collection/}{ln\slash a} 
(\emph{b. on words relating to time of day or calender to refer to a point in time or duration; use} \ding{43}~\B{ro} \emph{or} \ding{43}~\B{ka} \emph{to disambiguate})
\B{Kaym\-o zo\-la\-'u fray\-ul\-txa\-tu ne kel\-ku moe\-yä fte yi\-vom wu\-tsot, fti\-via nì\-'it lì'\-fya\-ti le\-Na'\-vi, ul\-te ki\-vä\-teng nì\-'o'.} All the participants came to our house one evening to have dinner, study a little Na'vi, hang out together, and have fun.~\LINK{http://naviteri.org/2014/09/stxeli-alor-the-text/}{b}

% oare
\hi 
\HT{oare}{\B{o\cp{a}re}} \I{[o."{}a.RE]} \emph{n., astr.} moon (\emph{generic term}). 
\B{Pol\-pxay\-a oa\-ret tse\-'a ngal mì saw pxi\-set?} How many moons do you see in the sky right now?~\LINK{http://naviteri.org/2018/08/fmawnti-stolawm-srak-have-you-heard-the-news/}{b} 

% oe
\hi
\HT{oe}{\B{\cp{o}e}} \I{["{}o.E]} (\emph{with case endings or enclitic adp.s always pronounced as:} \I{["wE-]}) \emph{pn.} I (1\textsuperscript{st} person singular). 
(\emph{a. subjective}) 
\B{Li\-vu win sì txur oe ze\-ne | Ha nì\-'aw | Pxan li\-vu txo nì\-'aw oe nga\-ri | Tsa\-krr nga Na'\-vi\-ru yom\-tì\-yìng.} I must be fast and strong | So only | If I am worthy of you | Will you feed the People.~\LINK{http://naviteri.org/2013/08/taronway-the-hunt-song/}{b}
(\emph{b. agentive}) \B{oel} \I{[wEl]}
\B{Oel nga\-ti ka\-me\-i\-e.} I \textsc{see} you.~\LINK{http://wiki.learnnavi.org/index.php?title=Corpus\#UGO_Movie_Blog}{wiki}
\B{Ka\-tot tä\-ftxu oel | Nì\-ean nì\-rim}, I weave the rhythm | In yellow and blue.~\LINK{http://wiki.learnnavi.org/index.php/Weaving_Song}{wiki} 
\B{Pol\-txe Ey\-tu\-kan san oe ka\-yä sìk, slä oel pot ke spaw.} Eytukan said he would go, but I don't believe him.~\LINK{http://wiki.learnnavi.org/index.php?title=Canon\#Extracts_from_various_emails}{wiki}
(\emph{c. patientive}) \B{oet(i)} \I{[wEt\textcorner \slash "wE.ti]}
\B{Spi\-vaw oe\-ti ru\-txe, ma oe\-yä ey\-lan.} Please, believe me, friends.~\LINK{http://forum.learnnavi.org/news-announcements/a-response-from-paul-frommer!/}{ln}
\B{Fu\-la tsa\-yun oeng pi\-väng\-kxo ye'\-rìn ul\-te nga\-ri oel mok\-rit sta\-yawm, oe\-ti nit\-ram sley\-ku nì\-txan.} It makes me very happy that we'll be able to chat soon and that I'll hear your voice.~\LINK{http://forum.learnnavi.org/language-updates/fula-short-form-of-fiul-a-and-tsal-oeti-steyki-nitram/}{ln}
\B{Tsun nga oet mi\-vong.} You can rely on me.~\LINK{http://naviteri.org/2012/07/meetings-waterfalls-and-more/}{b}
(\emph{d. genitive}) \B{\cp{oe}yä} \I{["wE.j\ae]}
\B{Oe\-yä swi\-zaw nì\-ngay ti\-va\-kuk | Oe\-yä tuk\-rul txe'\-lan\-it ti\-va\-kuk.} Let my arrow strike true | Let my spear strike the heart.~\LINK{http://naviteri.org/2013/08/taronway-the-hunt-song/}{b}
\B{Ay\-ey\-lan\-ur oe\-yä sì ey\-lan\-ur lì'\-fya\-yä le\-Na'\-vi nì\-wotx.} To all my friends and friends of the Na'vi language.~\LINK{http://forum.learnnavi.org/news-announcements/a-response-from-paul-frommer!/}{ln}
(\emph{idiom}) \B{tì\-'e\-fu\-mì oe\-yä}, in my opinion. \B{tì\-o\-mum\-mì oe\-yä}, as far as I know, to my knowledge.
(\emph{e. dative}) \B{oer(u)} \I{[wER \slash "wE.Ru]}
\B{Tsun oe nga\-hu nì\-Na'\-vi pi\-väng\-kxo a fì\-'u oe\-ru prr\-te' lu.} It's a pleasure to be able to chat with you in Na'vi.~\LINK{http://wiki.learnnavi.org/index.php?title=Corpus\#Times_Online}{wiki}
\B{Oe\-ru syaw fko Txe\-wì.} My name is Txewì. [lit.: one calls (to) me Txewì]~\LINK{http://naviteri.org/2010/09/getting-to-know-you-part-3/}{b}
\B{Fpo\-le' ay\-ngal oer fì\-txan nì\-ftxa\-vang a 'upxa\-ret sto\-lawm oel.} I have heard the message you have sent me so passionately.~\LINK{http://forum.learnnavi.org/news-announcements/a-response-from-paul-frommer!/}{ln}
\B{Pi\-veng oer ftxey nga new ri\-vey fu\-ke.} Tell me whether you want to live or not.~\LINK{http://forum.learnnavi.org/language-updates/if-vs-whether-ftxey-fuke/msg156629/\#msg156629}{ln}
(\emph{f. topic}) \B{\cp{oe}ri} \I{["wE.Ri]}
\B{Oe\-ri lu swoa fne\-txum.} I'm allergic to intoxicating beverage.~\LINK{http://forum.learnnavi.org/language-updates/new-words-from-the-2012-avatarmeet-navi-class/msg551479/\#msg551479}{b}
\B{Oe\-ri so\-la\-lew zì\-sìt a\-pxe\-vol.} I'm \on{24} years old.~\LINK{http://naviteri.org/2010/09/getting-to-know-you-part-2/}{b}
(\emph{for inalienable possessions}) 
\B{Oe\-ri tì\-ngay\-ìl txe'\-lan\-it ti\-va\-kuk.} Let the truth strike my heart.~\LINK{http://naviteri.org/2013/08/taronway-the-hunt-song/}{b} 
\B{Oe\-ri ski\-en\-a tsyokx lu le\-skxir.} My right hand is wounded.~\LINK{http://naviteri.org/2014/05/mipa-ayliu-mipa-aysafpil-new-words-new-ideas/}{b}

% oeng
\hi
\HT{oeng}{\B{o\cp{eng}}} \I{[o."{}EN]} (\emph{with case endings or enclitic adp.s always pronounced as:} \I{[wE-]}) \emph{pn.} we (two), you and I, the two of us (1\textsuperscript{st} person dual inclusive).
(\emph{a. subjective}) 
\B{Nì\-'aw\-ve oeng yo\-lom wu\-tsot; maw\-krr uvan si.} First we had dinner; afterwards we played.~\LINK{http://naviteri.org/2011/08/new-vocabulary-clothing/comment-page-1/\#comment-917}{b}
(\emph{b. agentive}) \B{oe\cp{ngal}} \I{[wE."Nal]}
(\emph{c. patientive}) \B{oe\cp{ngat}(i)} \I{[wE."Nat\textcorner \slash wE."Na.ti]}
(\emph{d. genitive}) \B{oe\cp{nge}yä} \I{[wE."NE.j\ae]}
(\emph{e. dative}) \B{oe\cp{ngar}(u)} \I{[wE."NaR \slash wE."Na.Ru]} 
\B{Ma mun\-txa\-tu, oeng yaw\-ne lu oe\-nga\-ru fì\-tsap, ke\-fyak?} We love each other, don't we, my spouse?~\LINK{http://naviteri.org/2011/12/one-more-for-2011/}{b}
(\emph{f. topic}) \B{oe\cp{nga}ri} \I{[wE."Na.Ri]} 
(\ding{43}~\HL{oe}{\B{oe}}, \HL{nga}{\B{nga}})

% oeyk
\hi
\HT{oeyk}{\B{o\cp{eyk}}} \I{[o."{}Ejk\textcorner]} \emph{n.} cause.
\B{Oeyk tsa\-tì\-snaytx\-ä lu apxa sä\-nui tì\-eyk\-tan\-ä.} That loss was caused by a great failure of leadership.~\LINK{http://naviteri.org/2013/04/wheres-the-bathroom-and-other-useful-things/}{b} \\
\on{1}.~\HL{oeyktIng}{\B{oeyktìng}}, explain (why).

% oeyktìng
\hi
\HT{oeyktIng}{\B{o\cp{eyk}tìng}} \I{[o."{}Ejk\textcorner.t$\cdot$IN]} \emph{or} \I{["wEjk\textcorner.]} \emph{vtr.}\on{33} explain (why). 
\B{Ney\-ti\-ril Tsyey\-kur oeyk\-to\-lìng teyng\-ta fya\-pe ze\-ne vo\-ìk si\-vi tsa\-tì\-fkey\-tok\-mì.} Neytiri explained to Jake how to behave in that situation.~\LINK{http://naviteri.org/2020/11/vospxivopeya-ayliu-amip-novembers-new-words/}{b}
\B{Wo\-leyn Ìstawl yey\-fyat mì hll\-te fte oeyk\-ti\-vìng fra\-po\-ru tì\-hawl\-te\-ri sne\-yä.} Ìstaw drew a line on the ground to explain his plan to everyone.~\LINK{http://naviteri.org/2013/09/on-si-salewfya-shapes-and-directions/}{b} 
\B{Tsa\-lì'\-u\-ri alu nì\-fkey\-to\-ngay oe oeyk\-ta\-syìng ye'\-rìn mì fo\-stì a\-hay.} Concerning the word \emph{nì\-fkey\-to\-ngay} I'll explain (that) soon in a future post.~\LINK{http://naviteri.org/2014/05/mipa-ayliu-mipa-aysafpil-new-words-new-ideas/comment-page-1/\#comment-3841}{b}
\B{Tì\-yäkx ke lu sru\-nga', ma tsmuk. Nga txo sti, oeyk\-tìng teyng\-ta pe\-lun.} Snubbing isn't helpful, brother. If you're angry, explain why.~\LINK{http://naviteri.org/2015/03/kap-si-ayunil-saylahe-threats-dreams-and-other-things/}{b}
(\ding{43}~\HL{oeyk}{\B{oeyk}}) \\
\on{1}.~\HL{tIoeyktIng}{\B{tìoeyktìng}}, explanation.

% ohakx
\hi
\HT{ohakx}{\B{o\cp{hakx}}} \I{[o."hak']} \emph{adj.} hungry. 
\B{Nän yom kxam\-trr, 'ul 'e\-fu ohakx kaym.} The less you eat at noon, the hungrier you'll feel in the evening.~\LINK{http://naviteri.org/2012/02/trr-asawnung-lefpom-happy-leap-day/}{b}
\B{Sra\-ke nga 'e\-fu o\-hakx?--Sra\-ne, 'e\-fu nì\-txan. Yi\-vom ko!} Are you hungry?--Yes, quite. Let's eat.~\LINK{http://naviteri.org/2012/10/audio-and-video-learning-materials-for-navi-101/}{b}
(\ding{214}~\HL{vAng}{\B{väng}}) \\
\on{1}.~\HL{tIohakx}{\B{tìohakx}}, hunger.

% ohe
\hi
\HT{ohe}{\B{\cp{o}he}} \I{["{}o.hE]} \emph{pn.} I (1\textsuperscript{st} person singular; deferential\slash ceremonial form).
(\emph{a. subjective}) 
\B{Ä\-txä\-le su\-yi o\-he pi\-vawm, pe\-o\-lo' lu\-yu pum nge\-nge\-yä?} May I ask what tribe you belong to? (\emph{overly formal, stilted Na'vi})~\LINK{http://naviteri.org/2010/09/getting-to-know-you-part-2/}{b}
(\emph{b. agentive}) \B{\cp{o}hel} \I{["{}o.hEl]}
\B{Vu\-yin ohel Unil\-ta\-ron\-it.} I respectfully request the Dreamhunt.~\LINK{http://naviteri.org/2011/08/reported-speech-reported-questions/}{b}
(\emph{c. patientive}) \B{\cp{o}het(i)} \I{["{}o.hEt\textcorner \slash "{}o.hE.ti]}
(\emph{d. genitive}) \B{\cp{o}heyä} \I{["{}o.hE.j\ae]}
\B{Ay\-nge\-nga\-ru o\-he\-yä tsmuk\-it alu Ne\-wey te Tska\-ha So\-rewn\-'i\-te.} Allow me to introduce my sister, Newey te Tskaha Sorewn'ite.~\LINK{http://naviteri.org/2010/09/getting-to-know-you-part-1/}{b}
(\emph{e. dative}) \B{\cp{o}her(u)} \I{["{}o.hER \slash "{}o.hE.Ru]} 
\B{Ätxä\-le si tsnì li\-vu ohe\-ru Unil\-ta\-ron.} I respectfully request the Dreamhunt.~\LINK{http://naviteri.org/2011/08/reported-speech-reported-questions/}{b}
\B{Fay\-sul\-fä\-tu\-ä tì\-kang\-kem ohe\-ru meuia lu\-yu nì\-ngay.} The work of these masters truly honors me.~\LINK{http://naviteri.org/2010/06/first-post/}{b}
(\emph{f. topic}) \B{\cp{o}heri} \I{["{}o.hE.Ri]}
(\ding{43}~\HL{oe}{\B{oe}})

% oìsss
\hi
\HT{oIsss}{\B{o\cp{ìsss}}} \I{[o."{}Is:]} \textsc{i}.~\emph{intj.} watch it! (\emph{angry snarl}).  \\
\textsc{ii}.~\B{o\cp{ìsss} si} \I{[o."{}Is: si]} \emph{vin.} hiss (at s.o.)
\B{Nga lum\-pe oìsss so\-li por?} Why did you hiss at him?~\LINK{http://naviteri.org/2013/04/wheres-the-bathroom-and-other-useful-things/}{b}

% okup
\hi
\HT{okup}{\B{\cp{o}kup}} \I{["{}o.kup\textcorner]} \emph{or} \I{[.kUp\textcorner]} \emph{n.} milk.
\B{Sa'\-nok prr\-nen\-ur yom\-tìng fa okup sne\-yä.} A mother feeds an infant with her milk.~\LINK{http://naviteri.org/2013/06/looking-back-looking-forward/}{b}

% ‹ol›
\hi
‹\B{ol}› \I{[o.l]} \emph{inf. in} ‹1› \emph{to form the perfective aspect}.

% Olangi
\hi
\HT{olangi}{\B{Olangi}} \I{[o.la.Ni]} \emph{n.} \emph{clan name, clan of the plains and Direhorse riders}.

% olo'
\hi
\HT{olo'}{\B{o\cp{lo'}}} \I{[o."loP]} \emph{n.} clan, tribe; community.
\B{Ä\-txä\-le su\-yi o\-he pi\-vawm, pe\-o\-lo' lu\-yu pum nge\-nge\-yä?} May I ask what tribe you belong to? (\emph{overly formal, stilted Na'vi})~\LINK{http://naviteri.org/2010/09/getting-to-know-you-part-2/}{b}
\B{Tsa\-krr syay ay\-nge\-yä, syay o\-lo'\-ä oe\-yä la\-yu teng.} Then you will suffer the same fate as my clan.~\LINK{http://naviteri.org/2010/07/vocabulary-update/}{b} 
\B{Ay\-o\-lo'\-ru a\-la\-he peng zi\-va\-'u.} Tell the other clans (to) come.~\LINK{http://forum.learnnavi.org/language-updates/naviya-the-other-vocative/}{ln\slash a}
\B{Oe\-ri tì\-rey\-ti oel stxe\-nu\-to\-lìng fpi olo' aw\-nge\-yä.} I offered my life for the sake of our clan.~\LINK{http://naviteri.org/2011/09/miscellaneous-vocabulary/}{b} \\  
\on{1}.~\HL{nIolo'}{\B{nìo\-lo'}}, (together) as members of a clan. 
\on{2}.~\HL{olo'eyktan}{\B{olo'\-eyk\-tan}}, clan leader. 
\on{3}.~\HL{lI'fyaolo'}{\B{lì'\-fya\-o\-lo'}}, language clan, \emph{Sprachbund}. 
\on{4}.~\HL{uniltIrantokxolo'}{\B{Unil\-tì\-ran\-tokx\-o\-lo'}}, \emph{Avatar} (fan) community.
\on{5}.~\HL{txanlokxe}{\B{txan\-lo\-kxe}}, clan or tribal domain; country.

% olo'eyktan 1
\hi
\HT{olo'eyktan}{\B{olo\cp{'eyk}tan}}\textsuperscript{\on{1}} \I{[o.lo."PEjk\textcorner.tan]} \emph{n.} clan leader (\emph{generally gender\hyp{}neutral}). 
\B{Maw\-krr\-a Tsyeyk ftxo\-lu\-lì\-'u, tslam fra\-pol fu\-ta slu po Olo\-'eyk\-tan a\-mip.} After Jake spoke, everyone understood that he had become the new Clan Leader.~\LINK{http://naviteri.org/2012/06/spring-vocabulary-part-3/}{b}
\B{Zun li\-vu oe O\-lo'\-eyk\-tan, zel oe nga\-ru srung si\-vi set.} If I were Clan Leader, then I would help you now.~\LINK{http://naviteri.org/2013/04/zun-zel-counterfactual-conditionals/}{b}
(\ding{43}~\HL{olo'}{\B{olo'}}, \HL{eyktan}{\B{eyk\-tan}}; \HL{olo'eykte}{\B{olo'eykte}})

% olo'eyktan 2
\hi
\HT{olo'eyktan2}{\B{olo'eyk\cp{tan}}}\textsuperscript{\on{2}} \I{[o.lo.PEjk\textcorner."tan]} \emph{n.} male clan leader. 

% olo'eykte
\hi 
\HT{olo'eykte}{\B{olo'eyk\cp{te}}} \I{[o.loP.Ejk."tE]} \emph{n.} female clan leader. 
(\ding{43}~\HL{olo'eyktan}{\B{olo'eyktan}})

% Omatikaya
\hi
\HT{omatikaya}{\B{Omati\cp{ka}ya}} \I{[o.ma.ti."ka.ja]} (\emph{gen.} \B{O\-ma\-ti\-ka\-ya\-ä}) \emph{n.} the Omaticaya clan of the Na'vi.
\B{Mo\-'at \mbox{alu} Tsa\-hìk lu O\-ma\-ti\-ka\-ya\-ä le\-'awa ha\-pxì\-tu a \mbox{ioi} sä\-pi fa 'a\-re.} Mo'at, the Tsahik, is the only member of the Omaticaya who wears a poncho.~\LINK{http://naviteri.org/2011/08/new-vocabulary-clothing/}{b}
\B{To\-ruk Mak\-to po\-lä\-hem a fì\-'u ro\-lo\-'a nì\-txan Oma\-ti\-ka\-ya\-ru.} The arrival of Toruk Makto made a great impression on the Omaticaya.~\LINK{http://naviteri.org/2012/03/spring-vocabulary-part-1/}{b}
\B{Pe\-yä tsa\-tì\-pe\-'un a swey\-lu txo wi\-vem ay\-oeng Oma\-ti\-ka\-ya\-wä lu fe'.} His decision to fight against the Omaticaya was a bad one.~\LINK{http://naviteri.org/2011/07/txantsana-ultxa-mi-siatll-great-meeting-in-seattle/}{b}

% ompu
\hi 
\HT{ompu}{\B{\cp{om}pu}} \I{["{}om.pu]} \emph{adj., ofp.} fat, corpulent. 
(\ding{214}~\HL{litsi}{\B{li\-tsi}}; \ding{43}~\HL{nutx}{\B{nutx}}) 

% omum
\hi
\HT{omum}{\B{o\cp{mum}}} \I{[o."mum]} \emph{or} \I{[."mUm]} \emph{vtr.} know (a fact). (\emph{in any inflected or derived form, the stress shifts to the o, e.g.}~\B{i\cp{vo}mum} \I{[i."vo.mum]})    
\B{Ul\-te o\-mum oel fu\-ta tì\-fya\-wìn\-txu\-ri oe\-yä pe\-rey ay\-nga nì\-wotx.} And I know you are all waiting for my guidance.~\LINK{http://forum.learnnavi.org/news-announcements/a-response-from-paul-frommer!/}{ln}
\B{Pol ke tslam stum ke\-'ut, o\-mum lì\-'ut a\-vol nì\-'aw.} He understands almost nothing and only knows eight words.~\LINK{http://naviteri.org/2013/01/lifyari-po-peyi-how-good-is-her-navi/}{b}
\B{Tsa\-txon ay\-ut\-ral\-äo krr a pol oe\-ti no\-lìn nì\-wa\-we, olo\-me\-i\-um oel fa key\-rel fu\-ta lu yaw\-ne oe po\-ru.} That night beneath the trees when she looked at me meaningfully, I knew by the expression on her face that she loves me.~\LINK{http://naviteri.org/2011/05/some-miscellaneous-vocabulary/}{b} 
\B{Ke o\-mä\-ngum oel teyng\-ta ftu pe\-seng zo\-la\-'u tsa\-pam\-rel\-vi alu t.} I'm afraid I don't know where that letter t came from.~\LINK{http://naviteri.org/2014/02/barter-and-exchange/comment-page-1/\#comment-2612}{b}
\B{Ke o\-mum oel teyng\-ta fo kä pe\-seng\-ne.} I don't know where they're going.~\LINK{http://naviteri.org/2011/08/reported-speech-reported-questions/}{b}
(\emph{a. coll.}) \B{ko\cp{mum}} \I{[ko."mum]} \emph{or} \B{ke\cp{mum}} \I{[kE."mUm]} `I don't know \ldots' 
\B{Ko\-mum teyng\-ta kem\-pe lo\-len, ta\-lu\-na me\-sì\-ke\-nong\-mì ke lu kea key\-ey.} I don't know what happened because in the two examples there were no mistakes.~\LINK{http://naviteri.org/2021/03/mipa-ayliu-si-aylifyavi-new-words-and-expressions/\#comment-33161}{b}
(\ding{43}~\HL{syawm}{\B{syawm}}) \\
\textsc{Derivations}: \ding{43} \on{1}.~\HL{tIomum}{\B{tì\-o\-mum}}, knowledge. 
\on{2}.~\HL{sAomum}{\B{sä\-o\-um}}, information. 
\on{3}.~\HL{nIawnomum}{\B{nì\-aw\-no\-mum}}, as is known. 
\on{4}.~\HL{newomum}{\B{new\-o\-mum}}, be curious.

% onlor
\hi
\HT{onlor}{\B{on\cp{lor}}} \I{[on."loR]} \emph{adj.} good-smelling. 
\B{Syu\-ve\-ri ay\-ftxo\-zä\-yä ay\-nga\-ru fki\-van on\-lor ftxì\-lor\-sì nìwotx!} May the food of the celebrations smell and taste good to you.~\LINK{http://naviteri.org/2012/11/renu-ayinanfyaya-the-senses-paradigm/}{b}
(\ding{214}~\HL{onvA'}{\B{on\-vä'}}, \ding{43}~\HL{ontu}{\B{on\-tu}}, \HL{lor}{\B{lor}})

% ontu
\hi
\HT{ontu}{\B{\cp{on}tu}} \I{["{}on.tu]} \emph{n.} nose. 
\B{Oe\-ri ta pe\-yä fa\-hew a\-ke\-wong on\-tu te\-ya lä\-ngu.} My nose is full of his alien smell.~\LINK{http://languagelog.ldc.upenn.edu/nll/?p=1977}{ll} \\  
\on{1}.~\HL{onlor}{\B{onlor}}, good-smelling. \\  
\on{2}.~\HL{onvA'}{\B{onvä'}}, bad-smelling. \\
\on{3}.~\HL{ontsang}{\B{ontsang}}, nose ring. \\
\on{4}.~\HL{tIng}{\B{tìng ontu}}, (to) smell.

% ontsang
\hi
\HT{ontsang}{\B{\cp{on}tsang}} \I{["{}on.\t{ts}aN]} \emph{n.} nose ring.
(\ding{43}~\HL{ontu}{\B{on\-tu}}, \HL{tsang}{\B{tsang}})

% onvä'
\hi
\HT{onvA'}{\B{on\cp{vä'}}} \I{[on."v\ae{}P]} \emph{adj.} bad\hyp{}smelling.
\B{Tsa\-fa\-hew a\-on\-vä' ftu kung\-a ioang za\-'u.} That stinking smell comes from a rotten animal.~\LINK{http://naviteri.org/2019/09/choice-statements-vs-choice-questions-and-some-insults/}{b}
(\ding{214}~\HL{onlor}{\B{on\-lor}}, \ding{43}~\HL{ontu}{\B{on\-tu}}, \HL{vA'}{\B{vä'}}) \\
\on{1}.~\HL{vonvA'}{\B{vonvä'}}, asshole, dickhead.

% Otranyu
\hi
\HT{otranyu}{\B{Ot\cp{ran}yu}} \I{[ot\textcorner."Ran.ju]} \emph{n. name, probably clan}.

\end{multicols}

% ===== Section P =====

 \Chapter{P}{P}{P \I{[p]} (\emph{PeP})}{

  % pa'li
  \Headword{pa'li}{\B{\cp{pa'}li}}{
    \I{["paP.li]} \emph{n., fauna} direhorse, \emph{equidirus hoplites}.
    \B{Fmi mi\-vak\-to pa'\-lit!} Try to ride a direhorse!~\LINK{https://naviteri.org/2011/05/some-miscellaneous-vocabulary/}{b}
    \B{Lu pa'\-lir sngem a\-txur.} A direhorse has strong muscles.~\LINK{https://naviteri.org/2015/04/some-new-words-for-may-day/}{b}
    (\emph{proverb}) \B{Txìm a\-'aw ke tsun hi\-veyn mì tal me\-fa'\-li\-yä.} You can't take both positions or sit on the fence; you need to decide. [lit.: One butt can't sit on the backs of two direhorses]~\LINK{https://naviteri.org/2013/02/vospxi-ayol-posti-apup-short-post-for-a-short-month/}{b}
    (\emph{a.~special usage}) \B{\cp{pa'}\-li \cp{mak}\-to}, the act of riding a pa'\-li.  \\
    \textsc{Derivatives}: \ding{43}
    \on{1}.~\HL{pa'liwll}{\B{pa'liwll}}, direhorse pitcher plant.
  }

  % pa'liwll
  \Headword{pa'liwll}{\B{\cp{pa'}liwll}}{
    \I{["paP.li.w\s{l}]} \emph{n.}, \ding{96} direhorse pitcher plant, \emph{pseudocenia equina}.
  }

  % pa'o
  \Headword{pa'o}{\B{\cp{pa}'o}}{
    \I{["pa.Po]} \emph{n.} side.
    \B{Oe kaw\-krr ne few\-tu\-sok\-a pa\-'o kil\-van\-ä ke ka\-mä.} I never went to the opposite side of the river.~\LINK{https://naviteri.org/2011/01/new-year-new-vocabulary/}{b}
    \B{Fì\-tskxe\-ri fa\-'o lu yey; ke lu ko\-um.} This rock has straight sides; it's not rounded.~\LINK{https://naviteri.org/2013/09/on-si-salewfya-shapes-and-directions/}{b} \\
    \textsc{Derivatives}: \ding{43}
    \on{1}.~\HL{fApa}{\B{fä\-pa}}, top, topside.
    \on{2}.~\HL{ftArpa}{\B{ftär\-pa}}, left side.
    \on{3}.~\HL{kllpa}{\B{kll\-pa}}, bottom, bottom side.
    \on{4}.~\HL{pxawpa}{\B{pxaw\-pa}}, perimeter, border.
    \on{5}.~\HL{skiempa}{\B{ski\-em\-pa}}, right side.
    \on{6}.~\HL{wrrpa}{\B{wrr\-pa}}, outside.
    \on{7}.~\HL{rofa}{\B{ro\-fa}}, beside, alongside.
    \on{8}.~\HL{ftumfa}{\B{ftum\-fa}}, out of, from inside.
    \on{9}.~\HL{mIfa}{\B{mì\-fa}}, inside, inner side.
    \on{10}.~\HL{nemfa}{\B{nem\-fa}}, into, inside.
    \on{11}.~\HL{ftuopa}{\B{ftuo\-pa}}, from behind.
  }

  % pak
  \Headword{pak}{\B{pak}}{
    \I{[pak\textcorner]} \emph{part. for disparagement; always sentence\hyp{}final}.
    \B{Po\-an pak!? Ke lu po tsam\-si\-yu kaw\-'it!} Him!? There's not a warrior's bone in his whole body!~\LINK{https://forum.learnnavi.org/language-updates/ke-kawit/msg174572/\#msg174572}{ln}
    \B{Tsaw ke ley kaw\-'it pak!} That's not worth a thing!\slash That's just worthless junk!~\LINK{https://naviteri.org/2014/03/value-and-worth/}{b}
  }

  % palang
  \Headword{palang}{\B{\cp{pa}lang}}{
    \I{["pa.laN]} \emph{vtr.} contact (\emph{in a social sense}), communicate with.
    \B{Ul\-te txo smi\-von ngar ay\-ho\-ak\-tu, ftxey soaia\-mì ftxey sko ey\-lan, fo\-ti pa\-lang fte tsi\-vun ivo\-mum teyng\-ta ftxey lu fo\-ru fpom fu\-ke.} And if you know elders, whether in the family or as friends, contact them to know whether they are alright or not.~\LINK{https://naviteri.org/2020/03/mipa-sawasultsyip-a-new-contest/}{b}
    (\emph{a. expression}) \B{Pa\-lang ko!} Keep in touch!~\LINK{https://naviteri.org/2020/03/teri-tifkeytok-lefkrr-about-the-present-situation/}{b} \\
    \textsc{Derivatives}: \ding{43}
    \on{1}.~\HL{tIpalang}{\B{tìpalang}}, (social) contact.
  }

  % palate
  \Headword{palate}{\B{pa\cp{la}te}}{
    \I{[pa."l$\cdot$a.t$\cdot$E]} \emph{vin., inf.}\on{23} crumble, fall apart, disintegrate.
    \B{Tsa\-ni\-vi a\-txaw\-nu\-la nì\-fe' hu\-fwe\-mì pa\-lo\-la\-te.} That poorly constructed hammock fell apart in the wind.~\LINK{https://naviteri.org/2025/11/kaltxi-nimun-hello-again/}{b}
  }

  % palon
  \Headword{palon}{\B{\cp{pa}lon}}{
    \I{["pa.lon]} \emph{vin.} burn (\emph{of a fire}).
    \B{Txep a\-hì\-'i mì tep\-tseng par\-ma\-lon.} A little fire was burning in the fireplace.~\LINK{https://naviteri.org/2018/03/100a-liu-amip-64-new-words-part-1/}{b}
    \B{Txep pa\-ya\-lon, pa\-ya\-lon txep!} The fire will burn.~\LINK{https://naviteri.org/2026/01/follow-up-to-txep-si-txeva/}{b\slash a\on{3}}
    (\ding{43}~\HL{nekx}{\B{nekx}})
  }

  % palukan
  \Headword{palukan}{\B{pa\cp{lu}kan}}{
    \I{[pa."lu.kan]} \emph{n. short form of} \ding{43}~\HL{palulukan}{\B{pa\-lu\-lu\-kan}}. \emph{on Earth often used for} cat.
    \B{Ho\-na pa\-lu\-kan\-tsyìp a\-hì\-'i.} A sweet little cat.~\LINK{https://naviteri.org/2012/01/more-additions-to-the-lexicon/}{b}
    \B{Pa\-lu\-kan\-it tso\-le\-'a, ye\-rik lopx hi\-fwo kxam\-lä ze\-swa.} Spotting a thanator, the hexapede panicked and escaped through the grass.~\LINK{https://naviteri.org/2012/07/meetings-waterfalls-and-more/}{b}
    (\emph{a. idiom}) \B{Tì\-ka\-ra lu ä\-txä\-le pa\-lu\-kan\-ur.} Resistance is futile [lit.: `reistance is a request to the thanator'].~\LINK{https://naviteri.org/2023/09/aawa-ayliu-si-aylifyavi-amip-a-few-new-words-and-expressions/}{b}
  }

  % palukantsyìp
  \Headword{palukantsyIp}{\B{pa\cp{lu}kantsyìp}}{
    \I{[pa."lu.kan.\t{ts}jIp\textcorner]} \emph{n.} cat (Earth animal).
    \B{Teng\-krr hu pa\-lu\-kan\-tsyìp uvan se\-ri zo\-lìm oel mì sa'\-leng a 'u\-ot a lu txa' sì ekx\-txu.} While playing with my cat I felt something hard and rough on his skin.~\LINK{https://naviteri.org/2012/11/renu-ayinanfyaya-the-senses-paradigm/}{b}
    (\ding{43}~\HL{palukan}{\B{pa\-lu\-kan}})
  }

  % palulukan
  \Headword{palulukan}{\B{palu\cp{lu}kan}}{
    \I{[""pa.lu."lu.kan]} \emph{n., fauna} thanator, \emph{pestiapanthera ferox}.
    \B{Pa\-lu\-lu\-kan a\-tu\-sa\-ron lu le\-hrr\-ap.} A hunting thanator is dangerous.~\LINK{https://forum.learnnavi.org/language-updates/participles/}{ln}
    \B{Tsawl\-a pa\-lu\-lu\-kan\-ìl oe\-ti 'ìl\-me\-ko a krr, Nawm\-a Sa'\-nok lrr\-tok sei\-yi ul\-te oe em\-ro\-ley.} The Great Mother was looking after me when the big thanator attacked and I've survived.~\LINK{https://naviteri.org/2011/05/some-miscellaneous-vocabulary/}{b}
    \B{Oel tso\-le\-'a pa\-lu\-lu\-kan\-it a\-tsawl fra\-to mì sì\-rey.} I saw the biggest thanator I had ever seen.~\LINK{https://forum.learnnavi.org/language-updates/life-death-and-gerunds/}{ln}
    \B{Pa\-lu\-lu\-kan\-ìri lu pxe\-syon a ze\-ne fko zi\-ve\-rok nì\-tut.} Three things about the thanator must always be kept in mind.~\LINK{https://naviteri.org/2012/10/mipa-vospxi-mipa-ayliu-new-words-for-the-new-month/}{b}
    (\emph{a. proverb}) \B{Ätxä\-le si pa\-lu\-lu\-kan\-ur tsnì smar\-it li\-vo\-nu.\slash Ätxä\-le pa\-(lu)\-lu\-kan\-ur.} \emph{a futile gesture, an attempt to achieve something that might be desirable but will clearly not happen} [lit.: Ask a thanator to release its prey.]~\LINK{https://naviteri.org/2010/07/vocabulary-update/}{b}
    (\emph{b. proverb}) \B{Ha\-haw nì\-'aw txo pa\-lu\-kan smi\-von ngar.} ``Only sleep if you are familiar with the Thanator,'' i.e. \emph{don't think you're safe unless you're aware of the danger}.~\LINK{https://naviteri.org/2021/08/ulte-ayyoratu-leiu-and-the-winners-are-2/}{b}\\
    \textsc{Derivatives}: \ding{43}
    \on{1}.~\HL{palukantsyIp}{\B{palukantsyìp}}, cat.
  }

  % pam / pam si
  \Headword{pam}{\B{pam}}{
    \I{[pam]} \textsc{i}. \emph{n.} sound (\emph{sensation}).
    \B{Tsa\-ta\-ngek\-ìri pam fkan mo\-mek.} That (tree)\-trunk sounds hollow.~\LINK{https://naviteri.org/2013/09/on-si-salewfya-shapes-and-directions/}{b} \\
    \textsc{ii}. \B{pam si} \I{[pam si]} \emph{vin.} make a sound.
    \B{Fnu! Pam si rä\-'ä!} Be quiet! Don't make a sound!~\LINK{https://naviteri.org/2023/12/mipa-zisit-ayliu-amip-new-year-new-words/}{b}
    \B{Ut\-ral pam a\-wok so\-li krr\-a zup.} The tree made a loud sound when it fell.~\LINK{https://naviteri.org/2023/12/mipa-zisit-ayliu-amip-new-year-new-words/}{b} \\
    \textsc{Derivatives}: \ding{43}
    \on{1}.~\HL{lI'upam}{\B{lì\-'u\-pam}}, pronunciation.
    \on{2}.~\HL{ngampam}{\B{ngam\-pam}}, rhyme.
    \on{3}.~\HL{pamrel}{\B{pam\-rel}}, writing, s.t. written.
    \on{4}.~\HL{pamtseo}{\B{pam\-tseo}}, music.
    \on{5}.~\HL{pxorpam}{\B{pxor\-pam}}, ejective.
    \on{6}.~\HL{kakpam}{\B{kak\-pam}}, deaf.
    \on{7}.~\HL{vApam}{\B{vä\-pam}}, noise.
    \on{8}.~\HL{hawmpam}{\B{hawm\-pam}}, noise.
    \on{9}.~\HL{pamuvan}{\B{pam\-u\-van}}, soundplay.
    \on{10}.~\HL{tIranpam}{\B{tì\-ran\-pam}}, footstep (\emph{sound}).
    \on{11}.~\HL{stA'nIpam}{\B{stä'\-nì\-pam}}, recording.
    \on{12}.~\HL{plltxepam}{\B{pll\-txe\-pam}}, speech sound.
    \on{13}.~\HL{txampam}{\B{txam\-pam}}, soundblast colossus.
    \on{14}.~\HL{pamrIrlI'u}{\B{pam\-rìr\-lì'u}}, onomatopoetic word.
    \on{15}.~\HL{pamtsyIp}{\B{pam\-tsyìp}}, small\slash light sound.
  }

  % pamrel
  \Headword{pamrel}{\B{pam\cp{rel}}}{
    \I{[pam."REl]} \textsc{i}.~\emph{n.} writing, s.t. written; text.
    \B{Tok pe\-se\-nget pam\-rel\-ìl?} Where is the text?~\LINK{https://naviteri.org/2014/06/teri-tsaliu-alu-nifkeytongay-about-nifkeytongay/}{b}
    \B{Tsun Txil\-te pam\-rel\-it ivi\-nan; ta\-fral puk\-ot a\-nutx mu\-nge fra\-tseng.} Txilte knows how to read; therefore she brings a thick book wherever she goes.~\LINK{https://naviteri.org/2011/11/\%E2\%80\%99a\%E2\%80\%99awa-ayli\%E2\%80\%99u-amip-ni\%E2\%80\%99aw-\%E2\%80\%94-a-few-new-words-only/}{b}
    \B{Wìn\-txu fì\-'ul fu\-ta keng pam\-rel\-lu\-ke mì lì'\-fya, ke tswa' fkol ke\-'ut.} It shows that even without a written component to their language, nothing is forgotten.~\LINK{https://naviteri.org/2026/06/tskxekeng-a-mikyunfpi-pamrel-listening-exercise-text/}{b} \\
    \textsc{ii}.~\HT{pamrel si}{\B{pam\cp{rel} si}} \I{[pam."REl si]} \emph{vin.} write. (\emph{allows for double dative})
    \B{Oe 'upxa\-re\-ru tsmu\-kan\-ur oe\-yä pam\-rel so\-li.} I wrote my brother a message.~\LINK{https://naviteri.org/2013/06/looking-back-looking-forward/}{b}
    \B{Fu\-ri\-a nì\-'Ìng\-lì\-sì pam\-rel si\-vi, oe\-ru txoa li\-vu.} I'm sorry for writing in English.~\LINK{https://forum.learnnavi.org/language-updates/verb-for-write/}{ln}
    \B{Tsa\-lì\-'u\-ri fko pam\-rel si fya\-pe?} How is that word written?~\LINK{https://naviteri.org/2010/08/a-na\%E2\%80\%99vi-alphabet/}{b}
    \B{Pam\-rel sa\-yi trr\-ay krr a skxom la\-tsu oe\-ru.} I will write tomorrow when I have time.~\LINK{https://forum.learnnavi.org/language-updates/just-a-few-interesting-little-words/}{ln} \\
    \textsc{Derivatives}: \ding{43}
    \on{1}.~\HL{pamrelfya}{\B{pam\-rel\-fya}}, spelling. \\
    \on{2}.~\HL{pamrelsiyu}{\B{pam\-rel\-si\-yu}}, writer, author. \\
    \on{3}.~\HL{pamrelvi}{\B{pam\-rel\-vi}}, letter (\emph{phonetic symbol}). \\
    \on{4}.~\HL{pamrelvul}{\B{pam\-rel\-vul}}, pen, pencil.
  }

  % pamrelfya
  \Headword{pamrelfya}{\B{pam\cp{rel}fya}}{
    \I{[pam."REl.fja]} \emph{n.} spelling.
    \B{Pam\-rel\-fya\-ri fya\-pe?} How do you spell it?~\LINK{https://naviteri.org/2010/08/a-na\%E2\%80\%99vi-alphabet/}{b}
    \B{Lì\-'u\-ri alu tskxe pam\-rel fya\-pe?--Pam\-rel\-fya lu na Tsä, KxeKx, E.} How do you spell the word \emph{tskxe}?--It's spelled ts, kx, e.~\LINK{https://naviteri.org/2010/08/a-na\%E2\%80\%99vi-alphabet/}{b}
    (\ding{43}~\HL{pamrel}{\B{pam\-rel}})
  }

  % pamrelsiyu
  \Headword{pamrelsiyu}{\B{pam\cp{rel}siyu}}{
    \I{[pam."REl.si.ju]} \emph{n.} writer, author.
    \B{Kin\-trr\-am fay\-fam\-rel\-si\-yu\-hu oe ul\-txa so\-lei\-yi.} Last week I've met with these writers.~\LINK{https://naviteri.org/2013/08/fmawn-a-fkol-kilmulat-this-just-in/}{b}
    (\ding{43}~\HL{pamrel si}{\B{pam\-rel si}})
  }

  % pamrelvi
  \Headword{pamrelvi}{\B{pam\cp{rel}vi}}{
    \I{[pam."REl.vi]} \emph{n.} letter (\emph{phonetic symbol}).
    \B{Ke o\-mä\-ngum oel teyng\-ta ftu pe\-seng zo\-la\-'u tsa\-pam\-rel\-vi alu t.} I'm afraid I don't know where that letter t came from.~\LINK{https://naviteri.org/2014/02/barter-and-exchange/comment-page-1/\#comment-2612}{b}
    \B{U\-van Vo\-mun\-a Pam\-rel\-vi\-yä.} The Ten Letter Game.~\LINK{https://www.pbs.org/wgbh/nova/article/learn-navi-with-avatars-paul-frommer/}{nova}
    (\ding{43}~\HL{pamrel}{\B{pam\-rel}}) \\
    \textsc{Derivatives}: \ding{43}
    \on{1}.~\HL{snapamrelvi}{\B{snapamrelvi}}, alphabet.
  }

  % pamrelvul
  \Headword{pamrelvul}{\B{pam\cp{rel}vul}}{
    \I{[pam."REl.vul]} \emph{or} \I{[.vUl]} \emph{n.} pen, pencil, writing utensil (\emph{often shortened to} \ding{43}~\HL{relvul}{\B{rel\-vul}}).
    (i.) \B{pam\cp{rel}vul le\cp{rìn}}, pencil.
    (\ding{43}~\HL{pamrel}{\B{pam\-rel}}, \HL{vul}{\B{vul}}; \HL{pensIl}{\B{pen\-sìl}})
  }

  % pamrìrlì'u
  \Headword{pamrIrlI'u}{\B{pam\cp{rìr}lì'u}}{
    \I{[pam."RIR.lI.Pu]} \emph{n.} onomatopoetic word.
    (\ding{43}~\HL{pam}{\B{pam}}, \HL{rI'Ir}{\B{rì'ìr}}, \HL{lI'u}{\B{lì'u}})
  }

  % pamtseo
  \Headword{pamtseo}{\B{\cp{pam}tseo}}{
    \I{["pam.\t{ts}E.o]} \textsc{i}.~\emph{n.} music.
    \B{Tsa\-fne\-pam\-tseo ke zaw\-prr\-te' oe\-ne.} I don't enjoy that kind of music.~\LINK{https://naviteri.org/2015/04/some-new-words-for-may-day/}{b}
    \B{Pam\-tse\-ol ngop ay\-re\-nut | Mì ron\-sem\-ä tì\-fnu}, Music creates patterns | In the silence of the mind.~\LINK{https://web.archive.org/web/20160314132640/https://www.pandorapedia.com/navi/music/ritual\_music}{pp}
    \B{Pam\-tseo\-ri po su\-lä\-ngìn nì\-hawng.} He is overly engrossed in his music (and I'm displeased about it).~\LINK{https://naviteri.org/2010/09/getting-to-know-you-part-3/}{b} \\
    \textsc{ii}.~\HT{pamtseo si}{\B{\cp{pam}tseo si}} \I{["pam.\t{ts}E.o si]} \emph{vin.} play music, (\B{fa}) play an instrument.
    \B{Po pam\-tseo si fa au nìl\-tsan nì\-txan.} She plays the drums very well.~\LINK{https://naviteri.org/2011/01/new-year-new-vocabulary/}{b}
    (\ding{43}~\HL{pam}{\B{pam}}, \HL{tseo}{\B{tseo}}) \\
    \textsc{Derivatives}: \ding{43}
    \on{1}.~\HL{pamtseotu}{\B{pamtseotu}}, musician.
    \on{2}.~\HL{pamtseowll}{\B{pamtseowll}}, cat ear (plant).
    \on{3}.~\HL{pamtseongopyu}{\B{pamtseongopyu}}, music creator, composer.
    \on{4}.~\HL{pamtseovi}{\B{pamtseovi}}, musical piece.
  }

  % pamtseongopyu
  \Headword{pamtseongopyu}{\B{pamtseo\cp{ngop}yu}}{
    \I{[pam.\t{ts}E.o."Nop\textcorner.ju]} \emph{n.} music creator, composer.
    \B{Txo tì\-ke\-ftxo\-nga'\-a fmawn\-it ke stil\-vawm ay\-ngal, ze\-ne oe pi\-veng san Unil\-tì\-ran\-tokx\-ä pam\-tseo\-ngop\-yu a\-yaw\-ne to\-ler\-kä\-ngup.} If you haven't heard the sad news, I have to inform that the dear music creator of \emph{Avatar} has died.~\LINK{https://naviteri.org/2015/06/na-ayskxe-mi-telan-sad-news/}{b}
    (\ding{43}~\HL{pamtseo}{\B{pam\-tseo}}, \HL{ngopyu}{\B{ngop\-yu}})
  }

  % pamtseotu
  \Headword{pamtseotu}{\B{\cp{pam}tseotu}}{
    \I{["pam.\t{ts}E.o.tu]} \emph{n.} musician.
    \B{Fì\-pam\-tseo\-tu\-ri ke la\-yu ftue fwa run fkol sä\-rawn\-it a tam.} It won't be easy to find a satisfactory replacement for this musician.~\LINK{https://naviteri.org/2012/01/mipa-zisit-ayliu-amip-new-words-for-the-new-year/}{b}
    (\ding{43}~\HL{pamtseo}{\B{pam\-tseo}}, \HL{tute1}{\B{tu\-te}})
  }

  % pamtseovi
  \Headword{pamtseovi}{\B{\cp{pam}tseovi}}{
    \I{["pam.\t{ts}E.o.vi]} \emph{n.} musical piece.
    \B{Aw\-nga tìng mik\-yun ay\-lì\-'u\-lu\-ke a pam\-tseo\-vir ko!} Let's listen to some music without words.~\LINK{https://naviteri.org/2024/02/mipa-ayliu-si-aylifyavi-niul-more-new-words-and-expressions/}{b}
    (\ding{43}~\HL{pamtseo}{\B{pam\-tseo}})
  }

  % pamtseowll
  \Headword{pamtseowll}{\B{\cp{pam}tseowll}}{
    \I{["pam.\t{ts}E.o.w\s{l}]} \emph{n.}, \ding{96} cat ear, ``music plant'', \emph{felinafolia ferruginea}.
    (\ding{43}~\HL{pamtseo}{\B{pam\-tseo}}, \HL{'ewll}{\B{'ewll}})
  }

  % pamtsyìp
  \Headword{pamtsyIp}{\B{\cp{pam}tsyìp}}{
    \I{["pam.\t{ts}jIp\textcorner]} \emph{n.} small or light sound.
    (\emph{a. proverb}) \B{Tì\-tok slu tìk\-tok fa pam\-tsyìp a\-'aw.} Presence becomes absence by means of one little sound. [i.e. `Small things can make a big difference.']~\LINK{https://naviteri.org/2024/06/ayhapxi-tokxa-si-ayliu-alahe-parts-of-the-body-and-more/}{b}
    (\ding{43}~\HL{pam}{\B{pam}}, \HL{-tsyIp}{\B{-tsyìp}})
  }

  % pamuvan
  \Headword{pamuvan}{\B{\cp{pam}uvan}}{
    \I{["pam.u.van]} \emph{n.} sound play (\emph{playing with and enjoying the sounds of language}).
    (\ding{43}~\HL{pam}{\B{pam}}, \HL{uvan}{\B{uvan}})
  }

  % pan
  \Headword{pan}{\B{pan}}{
    \I{[pan]} (\emph{dual} \ding{43}~\HL{mefan}{\B{me\cp{fan}}}) \emph{n., num.} third, one third, $\frac{1}{3}$.
    \B{Tsu'\-tey\-ìl to\-lìng oer pan\-it smar\-ä.} Tsu'tey gave me a third of the prey.~\LINK{https://naviteri.org/2011/02/new-vocabulary-part-2/}{b}
  }

  % Pari
  \Headword{pari}{\B{Pa\cp{ri}}}{
    \I{[pa."Ri]} \emph{n.}, \textsc{lw} Paris (\emph{capital of France}).
    \B{Pa\-ri yaw\-ne lu oer nì\-ngay!} I truly love Paris.~\LINK{https://naviteri.org/2010/09/getting-to-know-you-part-1/}{b}
  }

  % parul
  \Headword{parul}{\B{pa\cp{rul}}}{
    \I{[pa."Rul]} \emph{or} \I{[."RUl]} (\textsc{rn}: \B{pa\cp{rùl}} \I{[pa."RUl]}) \emph{n.} miracle.
    \B{Fwa ay\-ioang a\-pxay fì\-txan Na'\-vi\-ru srung so\-li fte Saw\-tu\-tet li\-vä\-txayn lu pa\-rul nì\-ngay.} That so many animals helped the Na’vi defeat the Sky People was a genuine miracle.~\LINK{https://naviteri.org/2018/03/100a-liu-amip-64-new-words-part-1/}{b} \\
    \textsc{Derivatives}: \ding{43}
    \on{1}.~\HL{parulnga'}{\B{parulnga'}}, miraculous. \\
    \on{2}.~\HL{parultsyIp}{\B{parultsyìp}}, \emph{term of affection for children.}
  }

  % parulnga'
  \Headword{parulnga'}{\B{pa\cp{rul}nga'}}{
    \I{[pa."Rul.NaP]} \emph{or} \I{[."RUl.]} (\textsc{rn}: \B{pa\cp{rùl}nga'} \I{[pa."RUl.NaP]}) \emph{adj.} miraculous.
    (\ding{43}~\HL{parul}{\B{pa\-rul}})
  }

  % parultsyìp
  \Headword{parultsyIp}{\B{pa\cp{rul}tsyìp}}{
    \I{[pa."Rul.\t{ts}jIp\textcorner]} \emph{or} \I{[."RUl.]} (\textsc{rn}: \B{pa\cp{rùl}tsyìp} \I{[pa."RUl.\t{tS}Ip\textcorner]}) \emph{n. term of affection for children} (lit. `little wonder').
    \B{Txon le\-fpom, ma pa\-rul\-tsyìp. Hi\-va\-haw nìm\-wey.} Good night, my dear little one. Sleep peacefully.~\LINK{https://naviteri.org/2018/03/100a-liu-amip-64-new-words-part-1/}{b}
    (\ding{43}~\HL{parul}{\B{parul}})
  }

  % parwun
  \Headword{parwun}{\B{\cp{par}wun}}{
    \I{["paR.wun]} \emph{or} \I{[.wUn]} (\textsc{rn}: \B{\cp{par}wùn} \I{["paR.wUn]}) \emph{vtr.} inhabit.
    \B{Ay\-ioang\-pel fì\-spo\-not par\-wun?} What animals inhabit this island?~\LINK{https://naviteri.org/2026/04/tsivola-liu-amip-thirty-two-new-words/}{b}
    \B{Pxay\-a ioang\-ìl na'\-rìng\-it par\-wun.} Many animals inhabit the forest.~\LINK{https://naviteri.org/2026/04/tsivola-liu-amip-thirty-two-new-words/\#comment-68288}{b} \\
    \textsc{Derivations}: \ding{43}
    \on{1}.~\HL{parwuntu}{\B{par\-wun\-tu}}, inhabitant.
  }

  % parwuntu
  \Headword{parwuntu}{\B{\cp{par}wuntu}}{
    \I{["paR.wun.tu]} \emph{or} \I{[.wUn.]} (\textsc{rn}: \B{\cp{par}wùntu} \I{["paR.wUn.tu]}) \emph{n., ofp.} inhabitant.
    \B{Na fra\-ve\-'o ay\-ru\-sey\-ä, far\-wun\-tu Ey\-wa\-'ev\-eng\-ä sä\-fpìl\-yewn nì\-'aw\-nì\-'aw fa ay\-fya\-'o a ta\-re fey pxaw\-ngip\-it.} As with any ecosystem, the inhabitants of Pandora communicate in ways unique to their environment.~\LINK{https://naviteri.org/2026/06/tskxekeng-a-mikyunfpi-pamrel-listening-exercise-text/}{b}
    (\ding{43}~\HL{parwun}{\B{par\-wun}})
  }

  % paskalin
  \Headword{paskalin}{\B{paska\cp{lin}}}{
    \I{[pa.ska."lin]} \emph{n.} sweet berry (term of endearment).
    \B{Hivahaw nìmwey, ma paskalin.} Sleep well, honey.~\LINK{https://naviteri.org/2016/06/mrrvola-lifyavi-amip-forty-new-expressions/}{b}
    (\ding{43}~\HL{pasuk}{\B{pa\-suk}}, \HL{kalin}{\B{ka\-lin}})
  }

  % pasuk
  \Headword{pasuk}{\B{\cp{pa}suk}}{
    \I{["pa.suk\textcorner]} \emph{or} \I{[.sUk\textcorner]} (\textsc{rn}: \B{\cp{pa}sùk} \I{["pa.sUk\textcorner]}) \emph{n.}, \ding{96} berry.
    \B{Tsay\-fa\-suk tsu\-ka\-nom lu krr\-ka fì\-zì\-sì\-krr nì\-'aw.} Those berries are available during this season only.~\LINK{https://naviteri.org/2021/03/mipa-ayliu-si-aylifyavi-new-words-and-expressions/}{b} \\
    \textsc{Derivatives}: \ding{43}
    \on{1}.~\HL{vozampasukut}{\B{vozampasukut}}, grinch tree.
    \on{2}.~\HL{paskalin}{\B{paskalin}}, sweet berry (term of endearment).
  }

  % pate
  \Headword{pate}{\B{\cp{pa}te}}{
    \I{["pa.tE]} \emph{vin.} get to a place, arrive.
    \B{Kar\-yul fngo\-lo' fu\-ta ay\-nu\-me\-yu pi\-va\-te ye'\-krr.} The teacher required the students to arrive early.~\LINK{https://naviteri.org/2012/01/mipa-zisit-ayliu-amip-new-words-for-the-new-year/}{b}
    (\ding{43}~\HL{pAhem}{\B{pä\-hem}})
  }

  % paw
  \Headword{paw}{\B{paw}}{
    \I{[paw]} \emph{vin.} grow, germinate and develop (of a plant).
    \B{Fì\-ut\-ral paw kil\-van\-lok nì\-'aw. Tsawl slu nì\-win nì\-txan.} This tree only grows (i.e., germinates, develops) near a river. It grows (i.e., gets big) very quickly.~\LINK{https://naviteri.org/2018/10/aysrr-ayvospxi-ayzisikrr-days-months-seasons/}{b}
    \B{Fkol yo\-lem nem\-fa kll\-te ut\-ral\-it a\-kawng ul\-te tsat pey\-ko\-law.} A bad tree was planted and was made to grow.~\LINK{https://naviteri.org/2020/06/mi-tanlokxe-oeya-srr-afpxamo-terrible-days-in-my-country/}{b}  \\
    \textsc{Derivatives}: \ding{43}
    \on{1}.~\HL{tIpaw}{\B{tìpaw}}, growth.
    \on{2}.~\HL{zIskrrmipaw}{\B{zìskrrmipaw}}, spring (season).
  }

  % pawk
  \Headword{pawk}{\B{pawk}}{
    \I{[pawk\textcorner]} \emph{n.} horn, wind instrument. \\
    \textsc{Derivatives}: \ding{43}
    \on{1}.~\HL{txArpawk}{\B{txärpawk}}, Palulukan Bone Horn.
  }

  % pawm
  \Headword{pawm}{\B{pawm}}{
    \I{[pawm]} \emph{vtr.\slash vin.} ask.
    (\emph{a. transitive use})
    \B{Pol po\-lawm teng\-a tì\-pawm\-it a li oel pal\-mawm trr\-am.} He asked the same question that I had already asked yesterday.~\LINK{https://naviteri.org/2020/05/tipusawm-tiuseyng-si-okvur-a-eltur-titxen-si-asking-answering-and-an-interesting-story/}{b}
    \B{Oel po\-ru po\-lawm fì\-'ut.} I asked him this.~\LINK{https://forum.learnnavi.org/intermediate/question-about-pawm/msg323238/\#msg323238}{ln}
    (\emph{b. intransitive use})
    \B{Alo a\-mrr po\-an po\-lawm, slä fra\-lo poe pol\-txe san ke\-he.} He asked five times, but each time she said, ``no.''~\LINK{https://forum.learnnavi.org/language-updates/txelanit-hivawl/}{ln}
    (\emph{c. with direct speech})
    \B{Po\-lawm po san sra\-ke Sä\-li ho\-lum sìk.} He asked, ``Did Sally leave?''~\LINK{https://forum.learnnavi.org/language-updates/if-vs-whether-ftxey-fuke/}{ln}
    (\emph{d. optional} \B{san … sìk})
    \B{Po\-lawm po, Ney\-ti\-ri kä pe\-seng\-ne?} He asked where Neytiri was going.~\LINK{https://naviteri.org/2011/08/reported-speech-reported-questions/}{b}
    (\emph{e. with} \HL{ta}{\B{ta}} \emph{to form the indirect object})
    \B{Pol po\-lawm tì\-pawm\-it oe\-ta.} He asked me a question.~\LINK{https://naviteri.org/2020/05/tipusawm-tiuseyng-si-okvur-a-eltur-titxen-si-asking-answering-and-an-interesting-story/}{b}
    (\emph{f. phrase}) \B{ätxä\-le si oe pi\-vawm …}, may I ask …
    \B{Ä\-txä\-le su\-yi o\-he pi\-vawm, pe\-o\-lo' lu\-yu pum nge\-nge\-yä?} May I ask what tribe you belong to? (\emph{overly formal, stilted Na'vi})~\LINK{https://naviteri.org/2010/09/getting-to-know-you-part-2/}{b}
    (\ding{43}~\HL{vin}{\B{vin}}, \ding{214}~\HL{'eyng}{\B{'eyng}}) \\
    \textsc{Derivatives}: \ding{43}
    \on{1}.~\HL{tIpawm}{\B{tì\-pawm}}, question.
  }

  % pay+
  \Headword{pay+}{\B{pay}\texttt{+}}{
    \I{[paj.]} \emph{inter. pref. on nouns to mean} which (of them)? (\emph{plural})
    \B{Pol\-txe pol pay\-lì\-'ut?} What did she say? [lit.: what words did she say?]~\LINK{https://naviteri.org/2011/08/reported-speech-reported-questions/}{b}
    (\ding{43}~\HL{pe+}{\B{pe}}\texttt{+}, \HL{ay+}{\B{ay}}\texttt{+})
  }

  % pay
  \Headword{pay}{\B{pay}}{
    \I{[paj]} \emph{n.} water, liquid.
    \B{Fra\-zì\-sì\-krr pay kil\-van\-ä nän nì\-'ul\-'ul.} Every season the river dries up more and more.~\LINK{https://naviteri.org/2012/02/trr-asawnung-lefpom-happy-leap-day/}{b}
    \B{Na\-ri si! Kll\-te\-sìn lu pay a\-txan. Fwi rä\-'ä!} Be careful! There's a lot of water on the ground. Don't slip!~\LINK{https://naviteri.org/2013/04/wheres-the-bathroom-and-other-useful-things/}{b}
    \B{Oe\-ri pay\-ìl tìng rì\-'ìr\-it key\-ä.} My face is reflected in the water.~\LINK{https://naviteri.org/2011/09/miscellaneous-vocabulary/}{b}
    \B{Oel pay\-ti nul\-new to swo\-at.} I prefer water to alcohol.~\LINK{https://forum.learnnavi.org/language-updates/prefer-x-to-y-24474/}{ln}
    \B{Fwa sar pay\-it fì\-txan lu le\-'al.} Using this much water is wasteful.~\LINK{https://naviteri.org/2014/07/tsawlultxamaw-a-ayliu-post-meetup-words/}{b}
    \B{Pay\-ri som\-wew\-pe?} What's the temperature of the water?~\LINK{https://naviteri.org/2011/05/weather-part-2-and-a-bit-more-2/}{b}
    (\emph{a. idiom}) \B{Pay\-ìl a lipx tskxe\-ti ripx.} Dripping water pierces a stone. [{\em persistent effort can accomplish unexpected and amazing things}]~\LINK{https://naviteri.org/2020/04/aawa-u-amip-a-few-new-things/}{b} \\
    \textsc{Derivatives}: \ding{43}
    \on{1}.~\HL{payfya}{\B{pay\-fya}}, stream.
    \on{2}.~\HL{paynga'}{\B{pay\-nga'}}, moist, humid, damp.
    \on{3}.~\HL{lepay}{\B{le\-pay}}, watery.
    \on{4}.~\HL{yemfpay}{\B{yem\-fpay}}, dipping, immersion.
    \on{5}.~\HL{paywll}{\B{pay\-wll}}, dapophet.
    \on{6}.~\HL{payIva}{\B{pay\-ìva}}, drop of water.
    \on{7}.~\HL{yapay}{\B{ya\-pay}}, mist.
    \on{8}.~\HL{reypay}{\B{rey\-pay}}, blood.
    \on{9}.~\HL{payoang}{\B{pay\-oang}}, fish.
    \on{10}.~\HL{txampay}{\B{txam\-pay}}, sea, ocean.
    \on{11}.~\HL{tskxepay}{\B{tskxe\-pay}}, ice.
    \on{12}.~\HL{tsngawpay}{\B{tsngaw\-pay}}, tears.
    \on{13}.~\HL{txumpaywll}{\B{txum\-pay\-wll}}, scorpion thistle.
    \on{14}.~\HL{paysena}{\B{pay\-se\-na}}, water container.
    \on{15}.~\HL{kxumpay}{\B{kxum\-pay}}, gel.
    \on{16}.~\HL{paysmung}{\B{pay\-smung}}, water carrier.
    \on{17}.~\HL{paytxew}{\B{pay\-txew}}, shoreline, water's edge.
    \on{18}.~\HL{paysyul}{\B{pay\-syul}}, water lily.
    \on{19}.~\HL{paytsim}{\B{pay\-tsim}}, water fountain.
    \on{20}.~\HL{paymaut}{\B{pay\-ma\-ut}}, fountain tree.
    \on{21}.~\HL{kxapay}{\B{kxapay}}, spit, saliva.
  }

  % payakan
  \Headword{payakan}{\B{\cp{Pay}akan}}{
    \I{["paj.a.kan]} \emph{n., name of a rogue tul\-kun}.
  }

  % payfya
  \Headword{payfya}{\B{\cp{pay}fya}}{
    \I{["paj.fja]} \emph{n.} stream.
    \B{Po spä few pay\-fya fte smar\-it si\-vutx.} He jumped across the stream to track his prey.~\LINK{https://naviteri.org/2011/01/new-year-new-vocabulary/}{b}
    (\ding{43}~\HL{pay}{\B{pay}}, \HL{fya'o}{\B{fya'o}})
  }

  % payìva
  \Headword{payIva}{\B{\cp{pay}ìva}}{
    \I{["paj.I.va]} \emph{n.} drop of water.
    (\ding{43}~\HL{pay}{\B{pay}}, \HL{Ilva}{\B{ìlva}})
  }

  % paymaut
  \Headword{paymaut}{\B{\cp{pay}maut}}{
    \I{["paj.ma.ut\textcorner]} \emph{or} \I{[.Ut\textcorner]} \emph{n.} \ding{96} fountain tree [lit.: `liquid fruit'].
    (\ding{43}~\HL{pay}{\B{pay}}, \HL{mauti}{\B{mauti}})
  }

  % paynga'
  \Headword{paynga'}{\B{\cp{pay}nga'}}{
    \I{["paj.NaP]} \emph{adj.} moist, damp, humid.
    \B{Fì\-fne\-rìn ke ley krr\-a slu pay\-nga'.} This kind of wood is worthless when it gets damp.~\LINK{https://naviteri.org/2014/03/value-and-worth/}{b}
    \B{'ak\-ra a\-pay\-nga'}, moist soil.~\LINK{https://naviteri.org/2011/05/weather-part-2-and-a-bit-more-2/}{b}
    (\ding{43}~\HL{pay}{\B{pay}}, \HL{-nga'}{\B{-nga'}})
  }

  % payoang
  \Headword{payoang}{\B{pay\cp{o}ang}}{
    \I{[paj."{}o.aN]} \emph{countable n., fauna} fish.
    \B{Lu tsa\-fne\-pay\-oang ftxì\-lor fra\-to krr\-a lu yrr.} That kind of fish is most tasty when eaten as \emph{sashimi}.~\LINK{https://naviteri.org/2013/01/awvea-posti-zisita-amip-first-post-of-the-new-year/}{b}
    \B{Lu fay\-oang a\-lor mì hil\-van Ey\-wa\-'e\-veng\-ä.} There are beautiful fish in Pandora's rivers.~\LINK{https://naviteri.org/2010/07/thoughts-on-ambiguity/}{b}
    \B{Oe\-yä txin\-tìn lu fwa stä'\-nì fay\-oang\-it.} My central societal occupation is to catch fish.~\LINK{https://naviteri.org/2010/09/getting-to-know-you-part-2/}{b}
    \B{Tì\-kan fì\-pay\-oang\-ä ke lu fwa tsun fko pot yi\-vom; lu oe\-yä 'eyl\-yong.} This fish is not meant to be eaten; it's my pet.~\LINK{https://naviteri.org/2022/11/krr-ao-an-exciting-time/}{b}
    \B{Fol ko\-la\-nom po\-ta ay\-srok\-it fay\-oang\-yoa.} They acquired beads from him in exchange for fish.~\LINK{https://naviteri.org/2014/02/barter-and-exchange/}{b}
    (\ding{43}~\HL{pay}{\B{pay}}, \HL{ioang}{\B{ioang}})
  }

  % paysena
  \Headword{paysena}{\B{\cp{pay}sena}}{
    \I{["paj.sE.na]} \emph{n.} water container.
    \B{\mbox{Tsngal} lu fne\-pay\-se\-na.} A cup is a type of water container.~{\color{blue}\href{https://naviteri.org/2014/05/mipa-ayliu-mipa-aysafpil-new-words-new-ideas/}{\textsc{b}}}
    (\ding{43}~\HL{pay}{\B{pay}}, \HL{sAhena}{\B{sä\-he\-na}})
  }

  % paysmung
  \Headword{paysmung}{\B{\cp{pay}smung}}{
    \I{["paj.smuN]} \emph{or} \I{[.smUN]} (\textsc{rn}: \B{\cp{pay}smùng} \I{["paj.smUN]}) \emph{n.} water carrier.
    (\ding{43}~\HL{pay}{\B{pay}}, \HL{sAmunge}{\B{sä\-mu\-nge}})
  }

  % paystan
  \Headword{paystan}{\B{pay\cp{stan}}}{
    \I{[paj."{}stan]} \emph{inter.} who? what men? (\emph{alternate form:} \B{(ay)stampe})
    (\ding{43}~\HL{pay+}{\B{pay}}\texttt{+}, \HL{tutan}{\B{tu\-tan}})
  }

  % payste
  \Headword{payste}{\B{pay\cp{ste}}}{
    \I{[paj."{}stE]} \emph{inter.} who? what women? (\emph{alternate form:} \B{(ay)stepe})
    (\ding{43}~\HL{pay+}{\B{pay}}\texttt{+}, \HL{tute2}{\B{tu\-té}})
  }

  % paysu
  \Headword{paysu}{\B{pay\cp{su}}}{
    \I{[paj."{}su]} \emph{inter.} who? what people? (\emph{alternate form:} \B{(ay)supe})
    (\ding{43}~\HL{pay+}{\B{pay}}\texttt{+}, \HL{tute1}{\B{tu\-te}})
  }

  % paysyul
  \Headword{paysyul}{\B{\cp{pay}syul}}{
    \I{["paj.sjul]} \emph{or} \I{[.sjUl]} \emph{n.} \ding{96} water lily, \emph{inrigo lilliam}.~\LINK{https://forum.learnnavi.org/language-updates/expansion-to-flora-fauna-and-places-from-the-valley-of-mo'ara/}{ln}
    (\ding{43}~\HL{pay}{\B{pay}}, \HL{syulang}{\B{syu\-lang}})
  }

  % paytxew
  \Headword{paytxew}{\B{\cp{pay}txew}}{
    \I{["paj.t'Ew]} \emph{n.} shoreline, water's edge.
    (\ding{43}~\HL{pay}{\B{pay}}, \HL{txew}{\B{txew}})
  }

  % paytsim
  \Headword{paytsim}{\B{\cp{pay}tsim}}{
    \I{["paj.\t{ts}im]} \emph{n.} water fountain.~\LINK{https://forum.learnnavi.org/vocab-phrases/words-from-the-disney-cast-members-training-on-the-na'vi-language/}{ln}
    (\ding{43}~\HL{pay}{\B{pay}}, \HL{tsim}{\B{tsim}})
  }

  % paywll
  \Headword{paywll}{\B{\cp{pay}wll}}{
    \I{["paj.w\s{l}]} \emph{n.}, \ding{96}  dapophet, ``water plant'', \emph{aloeparilus succulentus}.
    (\ding{43}~\HL{pay}{\B{pay}}, \HL{'ewll}{\B{'ewll}})
  }

  % pähem
  \Headword{pAhem}{\B{\cp{pä}hem}}{
    \I{["p\ae.hEm]} \emph{vin.} arrive.
    \B{Fo\-ru 'upxa\-ret oel fpo\-le', slä vay set ke pa\-mä\-hä\-ngem kea tì\-'eyng.} I have sent them a message, but up to now no answer has arrived.~\LINK{https://forum.learnnavi.org/news-announcements/a-response-from-paul-frommer!/}{ln}
    \B{Fì\-tseo\-ri ke tsun kaw\-tu pi\-vä\-hem tì\-yo'\-ne; tsran\-ten tì\-pä\-hem\-ä tì\-fmi nì\-'aw.} In this art it's impossible to arrive at perfection; the only thing that matters is the attempt to arrive there.~\LINK{https://naviteri.org/2012/01/mipa-zisit-ayliu-amip-new-words-for-the-new-year/}{b}
    \B{Ne 'o\-ra\-tsyìp po\-lä\-hem, yak si nì\-ftär.} When you arrive at the pond, turn to the left.~\LINK{https://naviteri.org/2013/09/on-si-salewfya-shapes-and-directions/}{b}
    (\emph{a. proverb}) \B{Nì\-win 'ìp tì\-'e\-wan, nì\-win pä\-hem tì\-ko\-ak.} Youth vanishes quickly, old age arrives quickly.~\LINK{https://naviteri.org/2025/04/mevosinga-liu-amip-twenty-new-words/}{b}
    (\ding{43}~\HL{pate}{\B{pa\-te}}) \\
    \textsc{Derivatives}: \ding{43}
    \on{1}.~\HL{tIpAhem}{\B{tì\-pä\-hem}}, arrival.
    \on{2}.~\HL{kllpA}{\B{kll\-pä}}, land, reach the ground.
  }

  % pänu
  \Headword{pAnu}{\B{\cp{pä}nu}}{
    \I{["p\ae.nu]} \emph{n.} promise.
    \B{Oey fä\-nut oel fpal fra\-krr.} I always keep my promises.~\LINK{https://naviteri.org/2025/11/kaltxi-nimun-hello-again/}{b} \\
    \textsc{Derivatives}: \ding{43}
    \on{1}.~\HL{pAnutIng}{\B{pä\-nu\-tìng}}, promise.
  }

  % pänutìng
  \Headword{pAnutIng}{\B{\cp{pä}nutìng}}{
    \I{["p\ae.nu.t$\cdot$IN]} \emph{vtr.}\on{33} promise (s.t. to s.o.).
    \B{Ney\-ti\-ri\-ti fkol pä\-nu\-to\-lìng oe\-ru!} Neytiri was promised to me!~\LINK{https://forum.learnnavi.org/language-updates/panuting/}{ln\slash a\textsuperscript{c}}
    \B{Lo\-ak\-ìl pä\-nu\-to\-lìng fu\-ta kar oe\-ru fya\-'ot a 'ìp fko nem\-fa ewll.} Loak promised he'd teach me how to vanish into the bushes.~\LINK{https://naviteri.org/2013/01/awvea-posti-zisita-amip-first-post-of-the-new-year/}{b}
    (\ding{43}~\HL{pAnu}{\B{pä\-nu}}, \HL{tIng}{\B{tìng}})
  }

  % pängkxo
  \Headword{pAngkxo}{\B{päng\cp{kxo}}}{
    \I{[p\ae{}N."k'o]} \emph{vin.} chat, converse, have a conversation (\B{hu}, with s.o.).
    \B{Oe pxìm päng\-kxo nga\-hu nì\-ron\-srel.} I often talk with you in my imagination.\slash I often imagine I'm talking with you.~\LINK{https://naviteri.org/2011/02/new-vocabulary-part-2/}{b}
    \B{Tsun oe nga\-hu nì\-Na'\-vi pi\-väng\-kxo a fì\-'u oe\-ru prr\-te' lu.} It's a pleasure to be able to chat with you in Na'vi.~\LINK{https://wiki.learnnavi.org/index.php?title=Corpus\#Times\_Online}{wiki}
    \B{Fay\-upxa\-re\-mì oe pa\-yäng\-kxo te\-ri ho\-ren lì'\-fya\-yä le\-Na'\-vi.} In these messages I will chat about the rules of the Na'vi language.~\LINK{https://naviteri.org/2010/06/first-post/}{b}
    \B{Fo ke pe\-räng\-kxo oe\-hu nìk\-mar.} They were chatting with me at a bad time.~\LINK{https://naviteri.org/2011/10/more-vocabulary-a-bit-of-grammar/}{b}\\
    \textsc{Derivatives}: \ding{43}
    \on{1}.~\HL{tIpAngkxo}{\B{tì\-päng\-kxo}}, chat, conversation.
    \on{2}.~\HL{pAngkxoyu}{\B{päng\-kxo\-yu}}, chatter, converser.
    \on{3}.~\HL{tireapAngkxo}{\B{ti\-re\-a\-päng\-kxo}}, commune.
  }

  % pängkxoyu
  \Headword{pAngkxoyu}{\B{päng\cp{kxo}yu}}{
    \I{[p\ae{}N."k'o.ju]} \emph{n.} chatter, converser.
    (\ding{43}~\HL{pAngkxo}{\B{päng\-kxo}})
  }

  % pängkxoyu lekoren
  \Headword{pAngkxoyu lekoren}{\B{päng\cp{kxo}yu leko\cp{ren}}}{
    \I{[p\ae{}N."k'o.ju lE.ko."REn]} \emph{n.} lawyer, rule chatter.
  }

  % pätsì
  \Headword{pAtsI}{\B{\cp{pä}tsì}}{
    \I{["p\ae.\t{ts}I]} \emph{n.}, \textsc{lw} badge.
  }

  % pe+
  \Headword{pe+}{\B{pe}\texttt{+}}{
    \I{[pE.]} \emph{inter. pref. on nouns to mean} what? which?
    \B{Ä\-txä\-le su\-yi o\-he pi\-vawm, pe\-o\-lo' lu\-yu pum nge\-nge\-yä?} May I ask what tribe you belong to? (\emph{overly formal, stilted Na'vi})~\LINK{https://naviteri.org/2010/09/getting-to-know-you-part-2/}{b}
    \B{Nga\-ru lu pe\-fne\-txin\-tìn nì\-trr\-trr?} What is your primary role (in society)?~\LINK{https://naviteri.org/2010/09/getting-to-know-you-part-2/}{b}
    \B{Nga\-ri tstxo lu pe\-lì\-'u?} What's your name?~\LINK{https://naviteri.org/2010/09/getting-to-know-you-part-3/}{b} \\
    \textsc{Derivatives}: \ding{43}
    \on{1}.~\HL{pay+}{\B{pay}\texttt{+}}, which (of them)? (plural).
    \on{2}.~\HL{pem+}{\B{pem}\texttt{+}}, what two?
    \on{3}.~\HL{pep+}{\B{pep}\texttt{+}}, what three?
    \on{4}.~\HL{pefne-}{\B{pefne}\texttt{-}}, what kind?
  }

  % -pe
  \Headword{-pe}{\B{-pe}}{
    \I{[.pE]} \emph{inter. suf. on nouns to mean} what? which?
    \B{Ya\-ri som\-wew\-pe?} What's the temperature?~\LINK{https://naviteri.org/2011/05/weather-part-2-and-a-bit-more-2/}{b}
    \B{Kil\-van\-ìri tsa\-tseng slo\-snep\-pe?} How wide is the river there?~\LINK{https://naviteri.org/2012/06/spring-vocabulary-part-3/}{b}
  }

  % pe'ngay
  \Headword{pe'ngay}{\B{pe'\cp{ngay}}}{
    \I{[p$\cdot$EP."Naj]} \emph{vin., inf.}\on{11} judge, conclude (\emph{used with} \ding{43}~\HL{tsnI}{\B{tsnì}}).
    \B{Po\-ri key\-rel\-fa oe po\-le'\-ngay tsnì ke new zi\-va\-'u.} From her expression, I concluded that she didn’t want to come.~\LINK{https://naviteri.org/2021/03/mipa-ayliu-si-aylifyavi-new-words-and-expressions/}{b}
    (\ding{43}~\HL{pe'un}{\B{pe\-'un}}, \HL{ngay}{\B{ngay}}) \\
    \textsc{Derivatives}: \ding{43}
    \on{1}.~\HL{tIpe'ngay}{\B{tì\-pe'\-ngay}}, conclusion.
    \on{2}.~\HL{pe'ngayyu}{\B{pe'\-ngay\-yu}}, judge.
  }

  % pe'ngayyu
  \Headword{pe'ngayyu}{\B{pe'\cp{ngay}yu}}{
    \I{[pEP."Naj.ju]} \emph{n.} judge (\emph{as a person}).
    (\ding{43}~\HL{pe'ngay}{\B{pe'\-ngay}})
  }

  % pe'un
  \Headword{pe'un}{\B{\cp{pe}'un}}{
    \I{["pE.Pun]} \emph{or} \I{[.PUn]} (\textsc{rn}: \B{\cp{pe}ùn} \I{["pE.Un]}) \emph{vtr.} decide.
    \B{Tì\-'eyng\-it oel to\-lel a krr, tsa\-krr pa\-ye\-'un swey\-a fya\-'ot a za\-mi\-vu\-nge oel ay\-ngar ay\-lì\-'ut ho\-ren\-ti\-sì lì'\-fya\-yä le\-Na'\-vi.} When I receive an answer, I will then decide the best way to bring you the words and rules of Na'vi.~\LINK{https://forum.learnnavi.org/news-announcements/a-response-from-paul-frommer!/}{ln}
    \B{Pol po\-le\-'un fu\-ta pe\-hem si nì\-'en.} He decided what to do on a hunch.~\LINK{https://naviteri.org/2011/01/new-year-new-vocabulary/}{b} \\
    \textsc{Derivatives}: \ding{43}
    \on{1}.~\HL{tIpe'un}{\B{tì\-pe\-'un}}, decision.
    \on{2}.~\HL{pe'ngay}{\B{pe'\-ngay}}, judge, conclude.
  }

  % pefne-
  \Headword{pefne-}{\B{pefne-}}{
    \I{[pE.fnE.]} \emph{inter. pref. on nouns to mean} which kind?
    \B{Nga\-ru lu pe\-fne\-txin\-tìn nì\-trr\-trr?} What is your primary role (in society)?~\LINK{https://naviteri.org/2010/09/getting-to-know-you-part-2/}{b}
  }

  % pefnel
  \Headword{pefnel}{\B{pe\cp{fnel}}}{
    \I{[pE."fnEl]} \emph{inter.} which kind? (\emph{alternate of} \ding{43}~\HL{fnepe}{\B{fne\-pe}})
  }

  % pefya
  \Headword{pefya}{\B{pe\cp{fya}}}{
    \I{[pE."fja]} \emph{inter.} how? (\emph{alternate of} \ding{43}~\HL{fyape}{\B{fya\-pe}})
    \B{Ya\-fkeyk (za\-'u) pe\-fya?} How is the weather?~\LINK{https://naviteri.org/2011/04/yafkeykiri-plltxe-frapo-everyone-talks-about-the-weather/}{b}
    \B{Pe\-fya nga fpìl? Oeng swey\-lu txo ki\-vä fu\-ke?} What do you think? Should we go or not?~\LINK{https://naviteri.org/2012/07/fivospxiya-aylifyavi-amip-this-months-new-expressions-pt-2/}{b}
    (\ding{43}~\HL{pe+}{\B{pe}}\texttt{+}, \HL{fya'o}{\B{fya\-'o}})
  }

  % pefyinep'ang
  \Headword{pefyinep'ang}{\B{pefyinep\cp{'ang}}}{
    \I{[pE.fjin.Ep."PaN]} \emph{inter.} how complex?
    \B{Ngal ke tslä\-ngam teyng\-ta fì\-tì\-ngä\-zìk\-ìri pe\-fyin\-ep\-'ang.} Unfortunately you don't understand how complex this problem is.~\LINK{https://naviteri.org/2012/07/fivospxiya-aylifyavi-amip-this-months-new-expressions-pt-2/}{b}
    (\ding{43}~\HL{pe+}{\B{pe}}\texttt{+}, \HL{fyinep'ang}{\B{fyin\-ep\-'ang}})
  }

  % pehem
  \Headword{pehem}{\B{pe\cp{hem}}}{
    \I{[pE."hEm]} \textsc{i}. \emph{inter.} what (action)? (\emph{alternate of} \ding{43}~\HL{kempe}{\B{kem\-pe}})
    \B{Pol po\-le\-'un fu\-ta pe\-hem si nì\-'en.} He decided what to do on a hunch.~\LINK{https://naviteri.org/2011/01/new-year-new-vocabulary/}{b} \\
    \textsc{ii}. \B{pe\cp{hem} si} \I{[pE."hEm si]} \emph{vin., inter.} do what (action)?
    (\emph{a. idiom}) \B{Tsun pe\-hem?} (\emph{short for} \B{Tsun fko pe\-hem si\-vi?}) `What can one do?'
    \B{Ri\-ni yaw\-ne lu oer, slä oe yaw\-ne ke lä\-ngu por kaw\-'it.---Tsun pe\-hem?} I'm in love with Rini, but she doesn't love me one bit.---What are you gonna do? That's life.~\LINK{https://naviteri.org/2016/06/mrrvola-lifyavi-amip-forty-new-expressions/}{b}
    (\ding{43}~\HL{pe+}{\B{pe}}\texttt{+}, \HL{kem}{\B{kem}})
  }

  % pehrr
  \Headword{pehrr}{\B{pe\cp{hrr}}}{
    \I{[pE."h\s{r}]} \emph{inter.} when (time)? (\emph{alternate of} \ding{43}~\HL{krrpe}{\B{krr\-pe}}).
    (\ding{43}~\HL{pe+}{\B{pe}\texttt{+}}, \HL{krr}{\B{krr}})
  }

  % pehrrlik
  \Headword{pehrrlik}{\B{pe\cp{hrr}lik}}{
    \I{[pE."h\s{r}.lik\textcorner]} \emph{inter.} what point in time? what hour? (\emph{alternate form of} \ding{43} \HL{krrlik}{\B{krrlikpe}})
  }

  % pek
  \Headword{pek}{\B{pek}}{
    \I{[pEk\textcorner]} \emph{n.} fin (\emph{of an aquatic animal}).
  }

  % pekxinum
  \Headword{pekxinum}{\B{pe\cp{kxi}num}}{
    \I{[pE."k'i.num]} \emph{or} \I{[.nUm]} (\textsc{rn}: \B{pee\cp{kxi}nùm} \I{[pE.E"gi.nUm]}) \emph{inter.} how tight\slash loose? (\emph{alternate form:} \ding{43}~\HL{'ekxinum}{\B{'e\-kxin\-um\-pe}})
    \B{Ngal mo\-lay' pxaw\-pxun\-it Lo\-ak\-ä a krr pe\-kxin\-um?} When you tried on Loak's armband, how tight was it?~\LINK{https://naviteri.org/2012/07/meetings-waterfalls-and-more/}{b}
    (\ding{43}~\HL{pe+}{\B{pe}}\texttt{+}, \HL{'ekxinum}{\B{'e\-kxi\-num}})
  }

  % pela'a
  \Headword{pela'a}{\B{pe\cp{la}'a}}{
    \I{[pE."la.Pa]} \emph{inter.}  how near? how far? what distance?.
    \B{Ftu Kel\-ut\-ral ne Txin\-tseng Saw\-tu\-te\-yä pe\-la\-'a?} How far is Hometree from Hell's Gate?~\LINK{https://naviteri.org/2019/08/aawa-lifya-si-lifyavi-amip-a-few-new-words-and-expressions/}{b}
    (\ding{43}~\HL{la'a}{\B{la'a}}, \HL{pe+}{\B{pe}\texttt{+}}; \HL{pelImsim}{\B{pelìmsim\slash lìmsimpe}})
  }

  % pelìmsim
  \Headword{pelImsim}{\B{pe\cp{lìm}sim}}{
    \I{[pE."lIm.sim]} \emph{inter.} how near? how far? what distance? (\emph{less common than} \ding{43}~\HL{pela'a}{\B{pela'a}})
    (\ding{43}~\HL{lImsimpe}{\B{lìmsimpe}})
  }

  % pelun
  \Headword{pelun}{\B{pe\cp{lun}}}{
    \I{[pE."lun]} \emph{or} \I{[."lUn]} \emph{inter.} why? (\emph{alternative form:} \ding{43}~\HL{lumpe}{\B{lum\-pe}}).
    \B{Fì\-txon na ton a\-la\-he nì\-wotx pe\-lun ke lu teng?} Why is this night not the same as all the other nights?~\LINK{https://web.archive.org/web/20110718042916/http://whyisthisnight.com:80/addl-more.html}{tn}
    \B{Nga pe\-lun wä\-pe\-ran?} Why are you hiding?~\LINK{https://naviteri.org/2011/02/new-vocabulary-ii\%E2\%80\%94part-1/}{b}
    (\ding{43}~\HL{pe+}{\B{pe}}\texttt{+}, \HL{lun}{\B{lun}})
  }

  % pem+
  \Headword{pem+}{\B{pem}\texttt{+}}{
    \I{[pEm.]} \emph{inter. pref. on nouns to mean} what two? which two?
    (\ding{43}~\HL{pe+}{\B{pe}}\texttt{+}, \HL{me+}{\B{me}}\texttt{+})
  }

  % pemstan
  \Headword{pemstan}{\B{pem\cp{stan}}}{
    \I{[pEm."{}stan]} \emph{inter.} who? what two men?~\LINK{https://naviteri.org/2011/07/number-in-na\%E2\%80\%99vi/}{b} (\emph{alternate form:} \ding{43}~\HL{pemstan}{\B{me\-stam\-pe}})
    (\ding{43}~\HL{pem+}{\B{pem}}\texttt{+}, \HL{tutan}{\B{tu\-tan}})
  }

  % pemste
  \Headword{pemste}{\B{pem\cp{ste}}}{
    \I{[pEm."{}stE]} \emph{inter.} who? what two women?~\LINK{https://naviteri.org/2011/07/number-in-na\%E2\%80\%99vi/}{b} (\emph{alternate form:} \ding{43}~\B{me\-ste\-pe})
    (\ding{43}~\HL{pem+}{\B{pem}}\texttt{+}, \HL{tute2}{\B{tuté}})
  }

  % pemsu
  \Headword{pemsu}{\B{pem\cp{su}}}{
    \I{[pEm."{}su]} \emph{inter.} who? what two people?~\LINK{https://naviteri.org/2011/07/number-in-na\%E2\%80\%99vi/}{b} (\emph{alternate form:} \ding{43}~\HL{pemsu}{\B{me\-su\-pe}})
    (\ding{43}~\HL{pem+}{\B{pem}}\texttt{+}, \HL{tute1}{\B{tu\-te}})
  }

  % pensìl
  \Headword{pensIl}{\B{\cp{pen}sìl}}{
    \I{["pEn.sIl]} \emph{n.}, \textsc{lw} pencil.
    (\ding{43}~\HL{pamrelvul}{\B{pam\-rel\-vul}})
  }

  % penunyol
  \Headword{penunyol}{\B{pe\cp{nun}yol}}{
    \I{[pE."nun.yol]} \emph{or} \I{[."nUn.]}  \emph{inter.} what length? how long (of time)? (\emph{alternate form:} \ding{43}~\HL{nunyol}{\B{nun\-yol\-pe}}).
    \B{Nga har\-ma\-haw pe\-nun\-yol?} How long were you sleeping?~\LINK{https://naviteri.org/2024/02/mipa-ayliu-si-aylifyavi-niul-more-new-words-and-expressions/}{b}
  }

  % peng
  \Headword{peng}{\B{peng}}{
    \I{[pEN]} \emph{vtr.} tell, inform.
    \B{Lu poe tstun\-wi, ul\-te evi\-ru peng el\-tur tì\-txen si a ay\-vur\-it te\-ri 'Rr\-ta, alu tse\-nge a\-stxong nì\-ngay.} She is kind and tells the kids interesting stories about the truly strange place Earth.~\LINK{https://naviteri.org/2010/07/ting-mikyun-fte-tslivam-listening-comprehension-1-3/}{b}
    \B{Tì\-'eyng\-it oel to\-lel a krr, ay\-nga\-ru pa\-yeng, tsa\-krr pa\-ye\-'un swey\-a fya\-'ot a za\-mi\-vu\-nge oel ay\-ngar ay\-lì\-'ut ho\-ren\-ti\-sì lì'\-fya\-yä le\-Na'\-vi.} When I receive an answer, I will let you know, and I will then decide the best way to bring you the words and rules of Na'vi.~\LINK{https://forum.learnnavi.org/news-announcements/a-response-from-paul-frommer!/}{ln}
    \B{Pi\-veng oer ftxey nga new ri\-vey fu\-ke.} Tell me whether you want to live (or not).~\LINK{https://forum.learnnavi.org/language-updates/if-vs-whether-ftxey-fuke/}{ln}
    \B{New oe me\-nga\-ru pi\-veng fu\-ta tì\-kang\-kem\-ìri a\-txan\-tsan me\-nge\-yä fpi ``Trr Le\-fpom ma Ame\-ri\-ka'' oe\-ru te\-ya so\-li nì\-txan.} I want to tell you both that I am very much filled with joy for your excellent work for ``Good Morning Amerika.''~\LINK{https://forum.learnnavi.org/language-updates/verb-for-write/}{ln}
    (\emph{in indirect speech with} \HL{fmawnta}{\B{fmawn\-ta}})
    \B{Ngal po\-leng oer fmawn\-ta po to\-ler\-kup.} You told me that he died.~\LINK{https://naviteri.org/2011/08/reported-speech-reported-questions/}{b} \\
    \textsc{Derivatives}: \ding{43}
    \on{1}.~\HL{ralpeng}{\B{ralpeng}}, translate, interpret.
    \on{2}.~\HL{penghrrap}{\B{penghrrap}}, binary sunshine.
    \on{3}.~\HL{penghrr}{\B{penghrr}}, warn.
  }

  % penghrr
  \Headword{penghrr}{\B{peng\cp{hrr}}}{
    \I{[p$\cdot$EN."h\s{r}]} \emph{vin., inf.}\on{11} warn (\B{te\-ri}, about).
    \B{Tsyeyk Na'\-vi\-ru po\-leng\-hrr te\-ri Saw\-tu\-te.} Jake warned the Na'vi about the Skypeople.~\LINK{https://naviteri.org/2016/06/mrrvola-lifyavi-amip-forty-new-expressions/}{b}
    (\ding{43}~\HL{pllhrr}{\B{pll\-hrr}}; \HL{peng}{\B{peng}}, \HL{hrrap}{\B{hrr\-ap}}) \\
    \textsc{Derivatives}: \ding{43}
    \on{1}.~\HL{sApenghrr}{\B{säpenghrr}}, warning.
  }

  % penghrrap
  \Headword{penghrrap}{\B{peng\cp{hrr}ap}}{
    \I{[pEN."h\s{r}.ap\textcorner]} \emph{n.}, \ding{96} binary sunshine, fringed lamp, \emph{lucinaria fibriata}.
    (\ding{43}~\HL{peng}{\B{peng}}, \HL{hrrap}{\B{hrr\-ap}})
  }

  % pengimpup
  \Headword{pengimpup}{\B{pe\cp{ngim}pup}}{
    \I{[pE."Nim.pup\textcorner]} \emph{or} \I{[.pUp\textcorner]} (\textsc{rn}: \B{pe\cp{ngim}pùp} \I{[pE."Nim.pUp\textcorner]}) \emph{inter.} what length? how long (of physical extension)? (\emph{alternate form:} \ding{43}~\HL{ngimpup}{\B{ngim\-pup\-pe}})
  }

  % pep+
  \Headword{pep+}{\B{pep}\texttt{+}}{
    \I{[pEp\textcorner.]} \emph{inter. pref. on nouns to mean} what three? which three?
    (\ding{43}~\HL{pe+}{\B{pe}}\texttt{+}, \HL{pxe+}{\B{pxe}}\texttt{+})
  }

  % pepstan
  \Headword{pepstan}{\B{pep\cp{stan}}}{
    \I{[pEp\textcorner."{}stan]} \emph{inter.} who? what three men? \emph{alternate form:} \ding{43}~\HL{pepstan}{\B{pxe\-stam\-pe}}
    (\ding{43}~\HL{pep+}{\B{pep}}\texttt{+}, \HL{tutan}{\B{tu\-tan}})
  }

  % pepste
  \Headword{pepste}{\B{pep\cp{ste}}}{
    \I{[pEp\textcorner."{}stE]} \emph{inter.} who? what three women? \emph{alternate form:} \ding{43}~\HL{pepste}{\B{pxe\-ste\-pe}}
    (\ding{43}~\HL{pep+}{\B{pep}}\texttt{+}, \HL{tute2}{\B{tu\-té}})
  }

  % pepsu
  \Headword{pepsu}{\B{pep\cp{su}}}{
    \I{[pEp\textcorner."{}su]} \emph{inter.} who? what three people? \emph{alternate form:} \ding{43}~\HL{pepsu}{\B{pxe\-su\-pe}}
    \B{Pep\-sul tok fì\-tse\-nget?} What three people are here?~\LINK{https://naviteri.org/2011/07/number-in-na\%E2\%80\%99vi/}{b}
    (\ding{43}~\HL{pep+}{\B{pep}}\texttt{+}, \HL{tute1}{\B{tu\-te}})
  }

  % peseng
  \Headword{peseng}{\B{pe\cp{seng}}}{
    \I{[pE."{}sEN]} \emph{or} \B{pe\cp{se}nge} \I{[pE."{}sE.NE]} \emph{inter.} where? what place? (\emph{alternative form:} \ding{43}~\HL{tseng}{\B{tseng\-pe}})
    \B{Tok pe\-se\-nget pam\-rel\-ìl?} Where is the text?~\LINK{https://naviteri.org/2014/06/teri-tsaliu-alu-nifkeytongay-about-nifkeytongay/}{b}
    \B{Oe\-ta a tsa\-sä\-srìn\-ìl tok pe\-seng\-it?} Where's the thing (you) borrowed from me?~\LINK{https://naviteri.org/2012/01/more-additions-to-the-lexicon/}{b}
    \B{Nga zo\-la\-'u ftu pe\-seng?} Where are you from?~\LINK{https://naviteri.org/2010/09/getting-to-know-you-part-2/}{b}
    \B{Vo\-lin oel teyng\-ta Ney\-ti\-ri kä pe\-seng\-ne.} I asked where Neytiri was going.~\LINK{https://naviteri.org/2011/08/reported-speech-reported-questions/}{b}
    \B{Ay\-ngal syu\-vet sweyn pe\-seng fte\-ke ay\-i\-o\-ang tsi\-vun tsat ki\-va\-nom?} Where do you keep the food so that animals can’t get it?~\LINK{https://naviteri.org/2021/09/aawa-ayliu-amip-a-few-new-words/}{b}
    (\ding{43}~\HL{pe+}{\B{pe}}\texttt{+}, \HL{tseng}{\B{tse\-nge}})
  }

  % peslosnep
  \Headword{peslosnep}{\B{peslo\cp{snep}}}{
    \I{[pE.slo."{}snEp\textcorner]} \emph{inter.} what width? how wide? (\emph{alternate form:} \ding{43}~\HL{slosnep}{\B{slo\-snep\-pe}})
    (\ding{43}~\HL{pe+}{\B{pe}}\texttt{+}, \HL{slosnep}{\B{slo\-snep}})
  }

  % pesomwew
  \Headword{pesomwew}{\B{pesom\cp{wew}}}{
    \I{[pE.som."wEw]} \emph{inter.} what temperature? (\emph{alternate form:} \ding{43}~\HL{somwew}{\B{som\-wew\-pe}})
    (\ding{43}~\HL{pe+}{\B{pe}}\texttt{+}, \HL{somwew}{\B{som\-wew}})
  }

  % pesrrpxì
  \Headword{pesrrpxI}{\B{pesrr\cp{pxì}}}{
    \I{[pE.s\s{r}."p'I]} \emph{inter.} what part of the day? what hour? (\emph{alternate form:} \ding{43}~\HL{trrpxI}{\B{trr\-pxì\-pe}})
    \B{Fo pä\-hem pe\-srr\-pxì?---Sre\-kam\-trr.} What part of the day will they arrive?---Before noon.~\LINK{https://naviteri.org/2024/05/krrteri-about-time/}{b}
    (\ding{43}~\HL{krrlik}{\B{krrlikpe\slash pe\-hrr\-lik}})
  }

  % pestan
  \Headword{pestan}{\B{pe\cp{stan}} \emph{or} \B{tu\cp{tam}pe}}{
    \I{[pE."{}stan]} \emph{inter.} who? what man? \emph{alternate form:} \ding{43}~\HL{tutampe}{\B{tu\-tam\-pe}}
    (\ding{43}~\HL{pe+}{\B{pe}}\texttt{+}, \HL{tutan}{\B{tu\-tan}})
  }

  % peste
  \Headword{peste}{\B{pe\cp{ste}} \emph{or} \B{tu\cp{te}pe}}{
    \I{[pE."{}stE]} \emph{inter.} who? what woman? \emph{alternate form:} \ding{43}~\HL{tutepe}{\B{tu\-te\-pe}}
    (\ding{43}~\HL{pe+}{\B{pe}}\texttt{+}, \HL{tute2}{\B{tu\-té}})
  }

  % pesu
  \Headword{pesu}{\B{pe\cp{su}}}{
    \I{[pE."{}su]} \emph{inter.} who? what person? (\emph{alternate form:} \ding{43}~\HL{tupe}{\B{tu\-pe}})
    \B{Nge\-nga lu pe\-su?} Who are you?~\LINK{https://naviteri.org/2010/09/getting-to-know-you-part-3/}{b}
    \B{Fo lu pesu?} Who are they?~\LINK{https://naviteri.org/2011/07/number-in-na\%E2\%80\%99vi/}{b}
    \B{Nga pe\-su\-hu kä\-teng nì\-trr\-trr?} Who do you usually hang out with?~\LINK{https://naviteri.org/2010/09/getting-to-know-you-part-3/}{b}
    (\ding{43}~\HL{pe+}{\B{pe}}\texttt{+}, \HL{tute1}{\B{tu\-te}})
  }

  % peu
  \Headword{peu}{\B{pe\cp{u}}}{
    \I{[pE."{}u]} \emph{inter.} what (thing)? (\emph{alternate form:} \ding{43}~\HL{'u}{\B{'u\-pe}})
    \B{Tì\-'o'\-ìri pe\-u su\-nu ngar fra\-to?} What is your favorite way to have fun?~\LINK{https://naviteri.org/2010/09/getting-to-know-you-part-3/}{b}
    \B{Nga\-tsyìp\-ìl new pe\-ut ta oe?} What does little you want from me?~\LINK{https://naviteri.org/2010/07/diminutives-conversational-expressions/}{b}
    (\ding{43}~\HL{pe+}{\B{pe}}\texttt{+}, \HL{'u}{\B{'u}})
  }

  % pewn
  \Headword{pewn}{\B{pewn}}{
    \I{[pEwn]} \emph{n.} neck.
    \B{Fkxi\-le pewn\-ta tu\-té\-yä kur.} The bib necklace hangs from the woman's neck.~\LINK{https://naviteri.org/2013/04/wheres-the-bathroom-and-other-useful-things/}{b}
    (\emph{a. idiom})
    \B{Ke tsun fko tspi\-vang to\-ruk\-it fa fwa pewn\-ti snew.} \emph{Proverbial expression for a method that will not work.} [lit.: You can't kill a great leonopteryx by constricting its throat]~\LINK{https://naviteri.org/2012/10/snewsyea-ftxoza-halowina-a-spooky-halloween/}{b}
    (\emph{shortened to}) \B{pewn to\-ruk\-ä}. \B{Sra\-ke fpìl ay\-ngal fu\-ta aw\-nga tsun hu Saw\-tu\-te a tsam\-it 'i\-vaw\-nìm fa fwa päng\-kxo fo\-hu nì\-'aw, ma smuk? Pewn to\-ruk\-ä!} Do you think we can avoid a war with the Sky People solely by talking with them, my brothers? That won't work!~\LINK{https://naviteri.org/2012/10/snewsyea-ftxoza-halowina-a-spooky-halloween/comment-page-1/\#comment-1755}{b}
  }

  % pey
  \Headword{pey}{\B{pey}}{
    \I{[pEj]} \emph{vin.} wait, wait for (\emph{with the topic}).
    \B{Ul\-te o\-mum oel fu\-ta tì\-fya\-wìn\-txu\-ri oe\-yä pe\-rey ay\-nga nì\-wotx.} And I know you are all waiting for my guidance.~\LINK{https://forum.learnnavi.org/news-announcements/a-response-from-paul-frommer!/}{ln}
    \B{Nì\-ay\-nga oe pe\-rey nì\-teng.} Like you, I too am waiting.~\LINK{https://forum.learnnavi.org/news-announcements/a-response-from-paul-frommer!/}{ln}
    \B{Ke tsun aw\-nga pi\-vey nul\-krr--txew lok.} We can't wait any longer--time is almost up.~\LINK{https://naviteri.org/2012/03/spring-vocabulary-part-1/}{b} \\
    \textsc{Derivatives}: \ding{43}
    \on{1}.~\HL{sIlpey}{\B{sìl\-pey}}, hope.
    \on{2}.~\HL{maweypey}{\B{ma\-wey\-pey}}, be patient.
    \on{3}.~\HL{srefey}{\B{sre\-fey}}, expect.
    \on{4}.~\HL{fe'pey}{\B{fe'\-pey}}, feel dread, fear.
  }

  \Headword{-pey}{\B{-pey}}{
    \I{[.pEj]} \emph{num., suf.} three (\emph{rest form of} \ding{43}~\HL{pxey}{\B{pxey}} \emph{to form higher numbers}). \B{vopey}, \on{11} (decimal), \on{13} (octal). \B{me\-vo\-pey}, \on{19} (decimal), \on{23} (octal); etc.
  }

  % peyä
  \Redirect{\B{\cp{pe}yä}}{po}{\B{po}}.

  % Peyral
  \Headword{peyral}{\B{Pey\cp{ral}}}{
    \I{[pEj."Ral]} \emph{female name}.
    \B{Pll\-txe fra\-po san Pey\-ral ke tsun ri\-vol nìl\-tsan sìk.} They say that Peyral can't sing well.~\LINK{https://naviteri.org/2014/06/teri-tsaliu-alu-nifkeytongay-about-nifkeytongay/}{b}
    \B{Pey\-ral\-ìl zet wur\-a wu\-tsot a\-'aw\-nem pxel sngel.} Peyral treats a cool cooked meal like garbage.~\LINK{https://naviteri.org/2012/10/mipa-vospxi-mipa-ayliu-new-words-for-the-new-month/}{b}
    \B{Kä\-sro\-lìn oel nik\-ro\-it Pey\-ral\-ur yoa fwa po rol oer.} I loaned Peyral a hair ornament in exchange for her singing to me.~\LINK{https://naviteri.org/2014/02/barter-and-exchange/}{b}
    \B{Ìlä Fey\-ral, mun\-txa so\-li Ra\-lu sì Ne\-wey nì\-wan me\-srr\-am.} According to Peyral, Ralu and Newey were secretly married the day before yesterday.~\LINK{https://naviteri.org/2012/07/meetings-waterfalls-and-more/}{b}
  }

  % piak
  \Headword{piak}{\B{pi\cp{ak}}}{
    \I{[pi."{}ak\textcorner]} \textsc{i}. \emph{adj.} (\emph{a.}) open.
    (\emph{b. of the weather}) no clouds, completely clear.
    \B{Ya\-fkeyk fya\-pe set?--Taw lu pi\-ak.} What's the weather like now?--The sky is clear.~\LINK{https://naviteri.org/2012/10/audio-and-video-learning-materials-for-navi-101/}{b} \\
    \textsc{ii}. \HT{piak si}{\B{piak si}} \I{[pi."{}ak\textcorner{} si]} \emph{vin.} open.
    \B{Oe\-ri me\-na\-ri\-ru pi\-ak si.} I open my eyes.~\LINK{https://forum.learnnavi.org/language-updates/open-close/}{ln} \\
    \textsc{iii}. \HT{piak sApi}{\B{piak säpi}} \I{[pi."{}ak\textcorner{} s\ae."pi]} \emph{vin.} open by itself, open on its own.
    \B{Po\-ri me\-na\-ri pi\-ak sä\-po\-li.} His eyes opened.~\LINK{https://forum.learnnavi.org/language-updates/open-close/}{ln}
    (\ding{214}~\HL{tstu}{\B{tstu}})
  }

  % pil
  \Headword{pil}{\B{pil}}{
    \I{[pil]} \emph{n.} facial stripe.
    \B{ean na pil}, facial\hyp{}stripe blue.~\LINK{https://naviteri.org/2010/08/mipa-ayopin-mipa-ayliu-new-colors-new-words/}{b}
    \B{Hu\-fwa lu fil\-ur Va'\-ru\-ä fne\-fe'\-ran\-vi, tsal\-su\-ngay fpìl fu\-ta say\-rìp lu nì\-txan.} Although Va'ru's facial stripes are rather uneven, I still think he's very handsome.~\LINK{https://naviteri.org/2012/10/mipa-vospxi-mipa-ayliu-new-words-for-the-new-month/}{b}
    (\emph{a. proverbial expression}) \B{Po\-ri fra\-pil sì fra\-sil lor lu nì\-txan.} He\slash She is very attractive.~\LINK{https://naviteri.org/2024/06/ayhapxi-tokxa-si-ayliu-alahe-parts-of-the-body-and-more/}{b}
    \B{Fra\-pil fra\-sil!} `Wow! Check him\slash her out! He\slash she is gorgeous!'~\LINK{https://naviteri.org/2024/06/ayhapxi-tokxa-si-ayliu-alahe-parts-of-the-body-and-more/}{b} \\
    \textsc{Derivatives}: \ding{43}
    \on{1}.~\HL{teylupil}{\B{tey\-lu\-pil}}, \emph{teylu}\hyp{}face (\emph{childish insult}).
  }

  % pinvul
  \Headword{pinvul}{\B{\cp{pin}vul}}{
    \I{["pin.vul]} \emph{or} \I{[.vUl]} \emph{n.} crayon (\emph{short form of} \ding{43}~\HL{'opinvultsyIp}{\B{'opin\-vul\-tsyìp}}).
    \B{Pol rel\-it wo\-leyn fa pin\-vul.} He drew a picture with a crayon.~\LINK{https://naviteri.org/2018/03/100a-liu-amip-64-new-words-part-1/\#comment-28355}{b}
    (i.) \B{\cp{pin}vul le\cp{fngap}}, marker (\emph{for writing or drawing})~\LINK{https://naviteri.org/2018/03/100a-liu-amip-64-new-words-part-1/\#comment-28354}{b}
  }

  % pizayu
  \Headword{pizayu}{\B{\cp{pi}zayu}}{
    \I{["pi.za.ju]} \emph{n.} ancestor.
    \B{Vay fwa zo\-la\-'u Tsyeyk\-Su\-li, To\-ruk Mak\-to alu pi\-za\-yu Ney\-ti\-ri\-yä txan\-rar\-mo\-'a fra\-to kip ay\-ha\-pxì\-tu Oma\-ti\-ka\-yaä.} Until Jake Sully arrived, Neytiri's ancestor was the most famous Toruk Makto among the Omatikaya.~\LINK{https://naviteri.org/2012/03/spring-vocabulary-part-1/}{b}
    (\emph{a. in storytelling}) \emph{opening formula} \B{Ay\-fi\-za\-yu pll\-txu\-ye san trro la\-mu krr a …}, Once upon a time … [lit.: the ancestors say, some day there was a time that]
    \emph{closing formula} \B{Fì\-fya pll\-txu\-ye ay\-fi\-za\-yu.} Thus say the ancestors.~\LINK{https://forum.learnnavi.org/intermediate/translation-exercise-a-navajo-coyote-tale-in-navi/}{ln}
  }

  % pìlok
  \Headword{pIlok}{\B{pì\cp{lok}}}{
    \I{[pI."lok\textcorner]} \emph{n.}, \textsc{lw} blog.
    \B{Sìl\-pey oe, aw\-nge\-yä lì'\-fya\-olo'\-ìri fì\-pì\-lok lì\-ye\-vu pxan.} I hope, this blog will be worth for our language community.~\LINK{https://naviteri.org/2010/06/first-post/}{b}
    \B{Zo\-la\-'u nì\-prr\-te' ne pì\-lok Na'\-vi\-te\-ri.} Welcome to the \emph{Na'\-vi\-te\-ri blog}.~\LINK{https://naviteri.org/2010/06/first-post/}{b}
    \B{Lu pì\-lok\-ur pxe\-sì\-kan sì pxe\-fne\-'upxa\-re.} The blog has three aims and three kinds of messages.~\LINK{https://naviteri.org/2010/06/first-post/}{b}
    \B{Pì\-lok\-ä fì\-ha\-pxì\-yä tì\-kan lu law.} The aim of this part of the blog is clear.~\LINK{https://naviteri.org/2010/06/first-post/}{b}
  }

  % pìmtxan
  \Headword{pImtxan}{\B{pìm\cp{txan}}}{
    \I{[pIm."t'an]} \emph{inter., adj.} how much? (\emph{alternate form:} \ding{43}~\HL{hImtxampe}{\B{hìm\-txam\-pe}})
    \B{Fì\-nik\-roi ley pìm\-txan?} How valuable is this hair ornament?~\LINK{https://naviteri.org/2014/03/value-and-worth/}{b}
    \B{pìmtxana pay?} how much water?~\LINK{https://naviteri.org/2021/04/mipa-ayliu-mipa-sioeykting-new-words-new-explanations/}{b}
    \B{pìmtxana tìsom?} how hot?~\LINK{https://naviteri.org/2021/04/mipa-ayliu-mipa-sioeykting-new-words-new-explanations/}{b}
    (\ding{43}~\HL{pe+}{\B{pe}}\texttt{+}, \HL{hImtxan}{\B{hìm\-txan}})
  }

  % pìsaw
  \Headword{pIsaw}{\B{pì\cp{saw}}}{
    \I{[pI."{}saw]} (\emph{a. adj.}) clumsy, accident\hyp{}prone.
    \B{Po lu pì\-saw. Trr\-am tol\-tem ve\-nu\-ti sne\-yä nìt\-kan\-lu\-ke.} He is accident\hyp{}prone. Yesterday he unintentionally shot his own foot.~\LINK{https://naviteri.org/2017/09/zisikrr-amip-ayliu-amip-new-words-for-the-new-season/}{b}
    \B{Lu Saw\-tu\-te wok sì pì\-saw nì\-wotx, na prr\-nen.} The skypeople are all loud and clumsy, like a baby.~\LINK{https://naviteri.org/2017/09/zisikrr-amip-ayliu-amip-new-words-for-the-new-season/}{b}
    \B{Ra\-nu kil\-van\-mì fyar\-mu\-fye na tsawl\-a pa\-yo\-ang a\-pì\-saw.} Ra\-nu was splashing in the river like a big clumsy fish.~\LINK{https://naviteri.org/2023/09/aawa-ayliu-si-aylifyavi-amip-a-few-new-words-and-expressions/}{b}
    (\emph{b. intj.}) \emph{e.g. when you've acted clumsily with unintended negative results, often accompanied by a sharp intake of breath.} \\
    \textsc{Derivatives}: \ding{43}
    \on{1}.~\HL{tIpsaw}{\B{tìp\-saw}}, clumsiness.
  }

  % pìtìk
  \Headword{pItIk}{\B{\cp{pì}tìk}}{
    \I{["pI.tIk\textcorner]} \emph{vtr.} scratch non\hyp{}harmfully (\emph{as an itch}).
    \B{Sran, sran, oe\-ri pì\-tìk tsa\-tse\-nget a mì tal! Fkxa\-ke nì\-ftxan kum\-a ter\-kup.} Yeah, yeah, scratch that place on my back! It itches like crazy!~\LINK{https://naviteri.org/2020/07/vospxivol-lefpom-happy-august/}{b}
    (\ding{43}~\HL{tsupx}{\B{tsupx}})
  }

  % pìwopx
  \Headword{pIwopx}{\B{pì\cp{wopx}}}{
    \I{[pI."wop']} \emph{n.} cloud.
    \B{Tsun fko ay\-on\-ti fì\-wopx\-ä ni\-vìn fte ya\-fkeyk\-it sre\-si\-ve\-'a.} The shapes of clouds can be used to predict the weather.~\LINK{https://naviteri.org/2013/09/on-si-salewfya-shapes-and-directions/}{b}
    (\emph{a. idiom}) \B{(na) fwam\-pop fkip fì\-wopx}, ``Fish out of water.'' [lit.: (like a) tapirus in the clouds].~\LINK{https://naviteri.org/2021/08/ulte-ayyoratu-leiu-and-the-winners-are-2/}{b} \\
    \textsc{Derivatives}: \ding{43}
    \on{1}.~\HL{lepwopx}{\B{lepwopx}}, cloudy.
    \on{2}.~\HL{pIwopxtsyIp}{\B{pìwopxtsyìp}}, nebula.
    \on{3}.\HL{fwopx}{\B{fwopx}}, dust (in the air).
    \on{4}.~\HL{hiupwopx}{\B{hi\-up\-wopx}}, cloud spitter.
  }

  % pìwopxtsyìp
  \Headword{pIwopxtsyIp}{\B{pì\cp{wopx}tsyìp}}{
    \I{[pI."wop'.\t{ts}jIp\textcorner]} \emph{n.} nebula (\emph{astronomical}).
    (\ding{43}~\HL{pIwopx}{\B{pì\-wopx}}, -\HL{-tsyIp}{\B{tsyìp}})
  }

  % pllhrr
  \Headword{pllhrr}{\B{pll\cp{hrr}}}{
    \I{[p$\cdot$\s{l}."h\s{r}]} (\textsc{perf}, \B{pol\cp{hrr}})\emph{vin.}\on{11} warn (\B{teri}, about).
    \B{Tseyk Na'\-vi\-ru pol\-hrr te\-ri Saw\-tu\-te.} Jake warned the Na'vi about the Skypeople.~\LINK{https://naviteri.org/2017/01/meliuteri-alu-tung-si-pllhrr-about-tung-and-pllhrr/}{b}
    (\ding{43}~\HL{penghrr}{\B{peng\-hrr}}; \HL{plltxe}{\B{pll\-txe}}, \HL{hrrap}{\B{hrr\-ap}}) \\
    \textsc{Derivatives}: \ding{43}
    \on{1}.~\HL{sApllhrr}{\B{sä\-pll\-hrr}}, warning.
  }

  % pllngay
  \Headword{pllngay}{\B{pll\cp{ngay}}}{
    \I{[p$\cdot$\s{l}."Naj]} (\textsc{perf}, \B{pol\cp{ngay}}) \emph{vin.}\on{11} admit, acknowledge.
    (\emph{a. with direct speech}) \B{Po ke tsun pi\-vll\-ngay san oe\-ru tì\-kxey.} He can't admit he's wrong.~\LINK{https://naviteri.org/2011/07/txantsana-ultxa-mi-siatll-great-meeting-in-seattle/}{b}
    \B{Oe pll\-ngay san mo\-lak\-to oe nì\-fe', ta\-fral sno\-laytx; wä\-tu lu oe\-to txur.} I acknowledge that I rode badly, so I lost; my opponent was stronger than I was.~\LINK{https://naviteri.org/2011/07/txantsana-ultxa-mi-siatll-great-meeting-in-seattle/}{b}
    (\emph{b. with the topic}) \B{Tsa\-'u\-ri po pll\-ngay.} He admits it.~\LINK{https://naviteri.org/2011/07/txantsana-ultxa-mi-siatll-great-meeting-in-seattle/}{b}
    (\ding{43}~\HL{plltxe}{\B{pll\-txe}}, \HL{ngay}{\B{ngay}})
  }

  % plltxe
  \Headword{plltxe}{\B{pll\cp{txe}}}{
    \I{[p\s{l}."t'E]} (\textsc{perf}, \B{pol\cp{txe}}) \textsc{i}. \emph{vin.} speak, say.
    \B{Pll\-txe nga nìl\-tsan!} You speak well!~\LINK{https://languagelog.ldc.upenn.edu/nll/?p=1977}{ll}
    \B{Ke\-tsran tu\-te ni\-vew hi\-vum, po\-ru pll\-txe san ru\-txe 'i\-vì\-'awn.} No matter who wants to leave, tell them to please stay.~\LINK{https://naviteri.org/2013/03/whoever-whatever-whenever/}{b}
    \B{Ma 'ite, aw\-nge\-yä fya\-'o\-ri ze\-ne nga sä\-nume si\-vi po\-ru fte tsi\-vun pil\-vll\-txe sì ti\-vì\-ran nì\-ay\-oeng.} My daughter, you will teach him our way so that (he) can speak and walk as we do.~\LINK{https://wiki.learnnavi.org/index.php?title=Na\%27vi\_from\_Avatar\_Movie\#Scene\_in\_Hometree}{a\slash wiki}
    \B{New Va'\-ru tskot Ey\-tu\-kan\-ä za\-sri\-vìn. Ke\-zem\-pll\-txe pa\-yll\-txe san ke\-he.} Va'ru wants to borrow Eytukan's bow. Of course he'll say no.~\LINK{https://naviteri.org/2012/06/spring-vocabulary-part-3/}{b}
    \B{Sa'\-nok nì\-nrr\-a lrr\-tok so\-li krr\-a prr\-nen alo a\-'aw\-ve pol\-txe.} The mother smiled with pride when the baby spoke for the first time.~\LINK{https://naviteri.org/2012/07/fivospxiya-aylifyavi-amip-this-months-new-expressions-pt-2/}{b}
    \B{Fko pll\-txe\-i\-e san me\-nga mun\-txa slo\-lu sìk!} I'm so happy to hear you got married!~\LINK{https://naviteri.org/2011/10/more-vocabulary-a-bit-of-grammar/}{b} \\
    \textsc{ii}. \emph{vtr.}
    \B{Pol\-txe pol pay\-lì\-'ut?} What did she say? [lit.: What words did she say?]
    \B{Ke pol\-txe pol tsay\-lì\-'ut.} She didn't say that.
    \B{Pol\-txe pol fay\-lu\-ta oe new ki\-vä.} She said, `I want to go'. \emph{or} She said she wanted to go.~\LINK{https://naviteri.org/2011/08/reported-speech-reported-questions/}{b}
    (\emph{a. with} \ding{43}~\HL{fa}{\B{fa}}) speak a language.
    \B{Ka le\-ke\-teng\-syon\-a ay\-olo' Ey\-wa\-'e\-veng\-ä, pll\-txe fko 'aw\-a lì'\-fya\-fa nì\-'aw.} Throughout Pandora's various clans, only one language is spoken.~\LINK{https://naviteri.org/2026/06/tskxekeng-a-mikyunfpi-pamrel-listening-exercise-text/}{b}
    (\emph{b. with} \ding{43}~\HL{nInu}{\B{nì\-nu}}) misspeak, speak in vain.
    \B{Oe\-ru txoa li\-vu. Pol\-txä\-nge nì\-nu.} Forgive me. I misspoke.~\LINK{https://naviteri.org/2013/04/wheres-the-bathroom-and-other-useful-things/}{b}
    (\emph{c. compound}) \B{way a pll\-txe}, spoken poem.~\LINK{https://naviteri.org/2011/12/ulte-yora\%E2\%80\%99tu-leiu-and-the-winner-is/}{b}
    (\emph{d. idiom}) \B{Kxe\-tse sì mik\-yun kop pll\-txe.} Body language is important [lit.: the tail and ears also speak].~\LINK{https://forum.learnnavi.org/language-updates/grid-tidbits/}{ln}
    (\emph{d. idiom}) \B{pll\-txe räp\-tum}, speak vulgarly, use vulgar language, swear.\LINK{https://naviteri.org/2024/02/mipa-ayliu-si-aylifyavi-niul-more-new-words-and-expressions/}{b} \\
    \textsc{Derivatives}: \ding{43}
    \on{1}.~\HL{plltxeyu}{\B{pll\-txe\-yu}}, speaker.
    \on{2}.~\HL{sAplltxe}{\B{sä\-pll\-txe}}, statement.
    \on{3}.~\HL{kezemplltxe}{\B{ke\-zem\-pll\-txe}}, of course.
    \on{4}.~\HL{pllngay}{\B{pll\-ngay}}, admit, acknowledge.
    \on{5}.~\HL{pllhrr}{\B{pll\-hrr}}, warn.
    \on{6}.~\HL{wAte}{\B{wä\-te}}, argue.
    \on{7}.~\HL{plltxepam}{\B{pll\-txe\-pam}}, speech sound.
    \on{8}.~\HL{txurplltxe}{\B{txur\-pll\-txe}}, declare.
  }

  % plltxepam
  \Headword{plltxepam}{\B{pll\cp{txe}pam}}{
    \I{[p\s{l}."t'E.pam]} \emph{n.} speech sound; phone, phoneme.
    \B{Slä ke\-zem\-pll\-txe, tsa\-pll\-txe\-pam alu KX ke fkey\-tok nì\-Fran\-se.} But of course, the phoneme `kx' doesn't exist in French.~\LINK{https://naviteri.org/2020/05/aawa-liu-amip-si-vurway-alor-a-few-new-words-and-a-beautiful-poem/\#comment-31171}{b}
    (\ding{43}~\HL{plltxe}{\B{plltxe}}, \HL{pam}{\B{pam}})
  }

  % plltxeyu
  \Headword{plltxeyu}{\B{pll\cp{txe}yu}}{
    \I{[p\s{l}."t'E.ju]} \emph{n.} speaker.
    \B{Toi\-tsye\-ri lu poe pll\-txe\-yu a\-naw\-fwe nì\-txan.} She's a very fluent speaker of German.~\LINK{https://naviteri.org/2015/08/aylifyavi-lereyfya-2-cultural-terms-2-and-more/}{b}
    (\ding{43}~\HL{plltxe}{\B{pll\-txe}})
  }

  % po
  \Headword{po}{\B{po}}{
    \I{[po]} (\emph{gen.} \B{peyä}) \emph{pn.} he\slash she (3\textsuperscript{rd} person singular).
    (\emph{a. subjective}) \B{Lam oer fwa po lu ka\-nu nì\-txan.} It seems to me that he is very smart.~\LINK{https://forum.learnnavi.org/language-updates/txelanit-hivawl/}{ln}
    \B{Maw\-krr\-a Tsyeyk ftxo\-lu\-lì\-'u, tslam fra\-pol fu\-ta slu po Olo'\-eyk\-tan a\-mip.} After Jake spoke, everyone understood that he had become the new Clan Leader.~\LINK{https://naviteri.org/2012/06/spring-vocabulary-part-3/}{b}
    \B{Lu po lor\-a tu\-té a\-lor.} She's an extremely beautiful woman.~\LINK{https://naviteri.org/2013/02/vospxi-ayol-posti-apup-short-post-for-a-short-month/}{b}
    (\emph{b. agentive}) \B{pol} \I{[pol]}
    \B{Li pol te\-rok fì\-tseng\-it srak?} Is she already here?~\LINK{https://naviteri.org/2011/02/new-vocabulary-part-2/}{b}
    (\emph{c. patientive}) \B{pot(i)} \I{[pot\textcorner \slash "po.ti]}
    \B{Apxa tì\-väng\-ìl po\-ti stey\-ko\-li.} (His) great thirst made him angry.~\LINK{https://naviteri.org/2011/01/new-year-new-vocabulary/}{b}
    \B{Pol\-txe Ey\-tu\-kan san oe ka\-yä sìk, slä oel pot ke spaw.} Eytukan said he would come but I don't believe him.~\LINK{https://forum.learnnavi.org/language-updates/verb-for-write/msg175491/\#msg175491}{ln}
    (\emph{d. genitive}) \B{\cp{pe}yä} \I{["pE.j\ae]}
    \B{Ut\-ral\-ä A\-nawm | Ay\-ri\-na' lu ay\-oeng, | A pe\-yä tì\-txur mì hi\-nam aw\-nge\-yä.} We are all seeds | Of the Great Tree, | Whose strength is in our legs.~\LINK{https://web.archive.org/web/20160314132640/https://www.pandorapedia.com/navi/music/ritual\_music}{pp}
    \B{Ay\-ron\-srel pe\-yä hä\-ngek nì\-txan.} His imaginings are (unpleasantly) weird.~\LINK{https://naviteri.org/2011/02/new-vocabulary-part-2/}{b}
    \B{Pol 'o\-lem pe\-yä wu\-tsot.} He made his (i.e., someone else's) dinner.~\LINK{https://naviteri.org/2011/12/one-more-for-2011/}{b}
    (\ding{43}~\HL{sno}{\B{sne\-yä}})
    (\emph{e. dative}) \B{por(u)} \I{[poR \slash "po.Ru]}
    \B{Ma Ra\-lu, srung si por!} Ralu, help her!~\LINK{https://naviteri.org/2012/03/spring-vocabulary-part-2/}{b}
    \B{Ke su\-nu por fwa rol eo tu\-te a\-pxay.} She doesn't like to sing in front of many people.~\LINK{https://naviteri.org/2014/06/teri-tsaliu-alu-nifkeytongay-about-nifkeytongay/}{b}
    (\emph{f. topic}) \B{\cp{po}ri} \I{["po.Ri]}
    \B{Pori mokri lor lu nìtxan ulte rol po nìftxavang.} Her voice is very beautiful and she sings passionately.~\LINK{https://naviteri.org/2014/06/teri-tsaliu-alu-nifkeytongay-about-nifkeytongay/}{b}
    \B{Po\-ri ze\-ne kll\-fri\-vo' nga.} He is your responsibility [lit.: you are responsible for him].~\LINK{https://forum.learnnavi.org/language-updates/misc-answers/}{ln\slash a\textsuperscript{c}} \\
    \textsc{Derivatives}: \ding{43}
    \on{1}.~\HL{pan}{\B{po\-an}}, he.
    \on{2}.~\HL{poe}{\B{po\-e}}, she.
    \on{3}.~\HL{frapo}{\B{fra\-po}}, everyone, everybody.
    \on{4}.~\HL{fIpo}{\B{fì\-po}}, this one (person).
    \on{5}.~\HL{tsapo}{\B{tsa\-po}}, that one (person).
    \on{6}.~\HL{'awpo}{\B{'aw\-po}}, one (person).
    \on{7}.~\HL{lapo}{\B{la\-po}}, another one.
    \on{8}.~\HL{poho}{\B{po\-ho}}, he\slash she (deferential form).
  }

  % poan
  \Headword{poan}{\B{po\cp{an}}}{
    \I{[po."{}an]} \emph{pn.} he (3\textsuperscript{rd} person singular masculine).
    \B{Po\-an pak!? Ke lu po tsam\-si\-yu kaw\-'it!} Him!? There's not a warrior's bone in his whole body!~\LINK{https://forum.learnnavi.org/language-updates/ke-kawit/msg174572/\#msg174572}{ln}
    \B{Poe mo\-war so\-li poan\-ur tsnì hi\-vum.} She advised him to leave.~\LINK{https://naviteri.org/2014/05/mipa-ayliu-mipa-aysafpil-new-words-new-ideas/}{b}
    (\ding{43}~\HL{po}{\B{po}})  \\
    \textsc{Derivatives}: \ding{43}
    \on{1}.~\HL{pohan}{\B{po\-han}}, he (deferential).
  }

  % poe
  \Headword{poe}{\B{po\cp{e}}}{
    \I{[po."{}E]} \emph{pn.} she (3\textsuperscript{rd} person singular feminine).
    \B{Poe pol\-txe nì\-fya\-'o a\-law} She spoke clearly.~\LINK{https://forum.learnnavi.org/language-updates/why-is-this-night/msg154875/\#msg154875}{ln}
    \B{Po\-el oe\-ti lar\-mawk.} She was talking about me.~\LINK{https://naviteri.org/2010/09/getting-to-know-you-part-1/}{b}
    \B{Krr a sto\-lawm Txe\-wì te\-ri hem a poe\-ru lo\-len, 'o\-le\-fu ke\-ftxo nì\-txan.} When Txewì heard about what happened to her, he felt very upset.~\LINK{https://naviteri.org/2010/07/ting-mikyun-fte-tslivam-listening-comprehension-1-3/}{b}
    \B{Poe\-ri unil\-tì\-ran\-tokx\-it tar\-mok a krr lam stum nì\-ay\-fo.} When she is in her Avatar body, she appears almost like them.~\LINK{https://naviteri.org/2010/07/ting-mikyun-fte-tslivam-listening-comprehension-1-3/}{b}
    (\ding{43}~\HL{po}{\B{po}}) \\
    \textsc{Derivatives}: \ding{43}
    \on{1}.~\HL{pohe}{\B{po\-he}}, she (deferential).
  }

  % pohan
  \Headword{pohan}{\B{po\cp{han}}}{
    \I{[po."han]} \emph{pn.} he (3\textsuperscript{rd} person singular; deferential\slash ceremonial form).
    (\ding{43}~\HL{poho}{\B{po\-ho}}, \HL{poan}{\B{po\-an}})
  }

  % pohe
  \Headword{pohe}{\B{po\cp{he}}}{
    \I{[po."hE]} \emph{pn.} she (3\textsuperscript{rd} person singular; deferential\slash ceremonial form).
    (\ding{43}~\HL{poho}{\B{po\-ho}}, \HL{poe}{\B{po\-e}})
  }

  % poho
  \Headword{poho}{\B{\cp{po}ho}}{
    \I{["po.ho]} \emph{pn.} he\slash she (3\textsuperscript{rd} person singular; deferential\slash ceremonial form).
    \B{Sra\-ke lu\-yu po\-ho set ro hel\-ku?} Is he\slash she at home now?~\LINK{https://naviteri.org/2022/02/lifyengteri-concerning-honorific-language/}{b}
    (\ding{43}~\HL{po}{\B{po}}) \\
    \textsc{Derivatives}: \ding{43}
    \on{1}.~\HL{pohe}{\B{po\-he}}, she (deferential).
    \on{2}.~\HL{pohan}{\B{po\-han}}, he (deferential).
  }

  % polpxay
  \Headword{polpxay}{\B{pol\cp{pxay}}}{
    \I{[pol."p'aj]} \emph{inter., adj.} how many? (\emph{alternate form:} \ding{43}~\HL{holpxaype}{\B{hol\-pxay\-pe}})
    \B{Nga\-ri so\-la\-lew pol\-pxay\-a zì\-sìt?\slash Nga\-ri so\-la\-lew zì\-sìt a\-pol\-pxay?} How old are you?~\LINK{https://naviteri.org/2010/09/getting-to-know-you-part-2/}{b}
    (\ding{43}~\HL{pe+}{\B{pe}}\texttt{+}, \HL{holpxay}{\B{hol\-pxay}})
  }

  % poltxe
  \Redirect{\B{pol\cp{txe}} \I{[pol."t'E]}}{plltxe}{\B{plltxe}}.

  % pom
  \Headword{pom}{\B{pom}}{
    \I{[pom]} \emph{vtr.} kiss. \\
    \textsc{Derivatives}: \ding{43}
    \on{1}.~\HL{sApom}{\B{säpom}}, kiss.
  }

  % pon
  \Headword{pon}{\B{pon}}{
    \I{[pon]} \emph{vtr.} balance.
    \B{Fwa pon sey\-ti sìn ki\-nam\-til lu le\-hrr\-ap, ma 'itan. Tsun ne\-kll zi\-vup tsawng.} Balancing a cup on your knee is dangerous, son. It can fall to the ground and break.~\LINK{https://naviteri.org/2020/04/aawa-u-amip-a-few-new-things/}{b}
    \B{Nì\-sngä\-'i Tsyeyk lu pì\-saw ul\-te ke tsun vul\-sìn pä\-pi\-von.} At first Jake was clumsy and wasn't able to balance on a branch.~\LINK{https://naviteri.org/2020/04/aawa-u-amip-a-few-new-things/}{b}
  }

  % pongu
  \Headword{pongu}{\B{\cp{po}ngu}}{
    \I{["po.Nu]} \emph{n.} group or party of people.
    \B{Vll eyk\-yu ne\-fä fte po\-ngu fä\-ki\-vä.} The leader signals the party to ascend by pointing upwards.~\LINK{https://naviteri.org/2012/06/spring-vocabulary-part-3/}{b}
    \B{Ay\-ha\-pxì\-tu po\-ngu\-ä txo\-pu si nì\-nän ta\-krr\-a Va'\-rul pxe\-ku\-tut lä\-txayn.} The members of the group are less afraid since Va'ru defeated three of the enemy.~\LINK{https://naviteri.org/2012/02/trr-asawnung-lefpom-happy-leap-day/}{b} \\
    \textsc{Derivatives}: \ding{43}
    \on{1}.~\HL{ngrrpongu}{\B{ngrr\-po\-ngu}}, grassroots movement.
    \on{2}.~\HL{tsampongu}{\B{tsam\-po\-ngu}}, war party.
    \on{3}.~\HL{wempongu}{\B{wem\-po\-ngu}}, squad, battle party.
    \on{4}.~\HL{tarpongu}{\B{tar\-po\-ngu}}, hunting party.
  }

  % postì
  \Headword{postI}{\B{\cp{po}stì}}{
    \I{["po.stI]} \emph{n.}, \textsc{lw} (forum\slash blog) post.
    \B{'Aw\-ve\-a po\-stì zì\-sìt\-ä a\-mip.} First post of the new year.~\LINK{https://naviteri.org/2013/01/awvea-posti-zisita-amip-first-post-of-the-new-year/}{b}
    \B{Mì fo\-stì a kxey\-ey\-ti oel li zo\-sley\-ko\-lu.} I already fixed the mistake in the post.~\LINK{https://naviteri.org/2013/01/awvea-posti-zisita-amip-first-post-of-the-new-year/comment-page-1/\#comment-1910}{b}
  }

  % prrku
  \Headword{prrku}{\B{\cp{prr}ku}}{
    \I{["p\s{r}.ku]} \emph{n.} womb.
    (\ding{43}~\HL{prrnen}{\B{prr\-nen}}, \HL{kelku}{\B{kel\-ku}})
  }

  % prrnen
  \Headword{prrnen}{\B{\cp{prr}nen}}{
    \I{["p\s{r}.nEn]} \emph{n.} infant, baby.
    \B{Sa'\-nok nì\-nrr\-a lrr\-tok so\-li krr\-a prr\-nen a\-lo a\-'aw\-ve pol\-txe.} The mother smiled with pride when the baby spoke for the first time.~\LINK{https://naviteri.org/2012/07/fivospxiya-aylifyavi-amip-this-months-new-expressions-pt-2/}{b}
    \B{Mong prr\-nen\-ìl fu\-ta sa'\-nul ve\-rewng nì\-tut.} The infant depends on her Mommy to look after her constantly.~\LINK{https://naviteri.org/2012/07/meetings-waterfalls-and-more/}{b}
    \B{Ze\-ne fko 'i\-vam\-pi prr\-nen\-it nì\-flrr fra\-krr.} One must always touch a baby gently.~\LINK{https://naviteri.org/2013/02/vospxi-ayol-posti-apup-short-post-for-a-short-month/}{b}
    \B{Oel vewng frr\-nen\-it.} I look after the infants.~\LINK{https://naviteri.org/2010/09/getting-to-know-you-part-2/}{b}
    \B{Ta'\-leng prr\-nen\-ä lu faoi, pum ko\-ak\-tu\-ä ekx\-txu.} A baby's skin is smooth, an old person's is rough.~\LINK{https://naviteri.org/2012/01/mipa-zisit-ayliu-amip-new-words-for-the-new-year/}{b}
    \B{Sa'\-nok prr\-nen\-ur yom\-tìng fa okup sne\-yä.} A mother feeds an infant with her milk.~\LINK{https://naviteri.org/2013/06/looking-back-looking-forward/}{b} \\
    \textsc{Derivatives}: \ding{43}
    \on{1}.~\HL{prrku}{\B{prrku}}, womb.
    \on{2}.~\HL{prrnesyul}{\B{prrnesyul}}, bud (of a flower).
    \on{3}.~\HL{prrwll}{\B{prrwll}}, moss.
    \on{4}.~\HL{prrsmung}{\B{prrsmung}}, baby carrier.
    \on{5}.~\HL{nga'prrnen}{\B{nga'prrnen}}, be pregnant (for people).
  }

  % prrnesyul
  \Headword{prrnesyul}{\B{\cp{prr}nesyul}}{
    \I{["p\s{r}.nE.sjul]} \emph{or} \I{[.sjUl]} \emph{n.} bud (of a flower).
    \B{Ay\-ut\-ral\-ur lu frr\-ne\-syul, Ey\-wa tì\-haw\-nu si\-vi!} The trees have buds, may Eywa protect (them)!~\LINK{https://forum.learnnavi.org/news-announcements/christmas-song-ninavi-on-youtube/}{ln}
    (\ding{43}~\HL{prrnen}{\B{prr\-nen}}, \HL{syulang}{\B{syu\-lang}})
  }

  % prrsmung
  \Headword{prrsmung}{\B{\cp{prr}smung}}{
    \I{["p\s{r}.smuN]} \emph{or} \I{[.smUN]} (\textsc{rn}: \B{\cp{prr}smùng} \I{["p\s{r}.smUN]}) \emph{n.} baby carrier.
    (\ding{43}~\HL{prrnen}{\B{prr\-nen}}, \HL{sAmunge}{\B{sä\-mu\-nge}})
  }

  % prrte'
  \Headword{prrte'}{\B{\cp{prr}te'}}{
    \I{["p\s{r}.tEP]} \emph{adj.} pleasurable (of an activity).
    \B{Tsun oe nga\-hu nì\-Na'\-vi pi\-väng\-kxo a fì\-'u oe\-ru prr\-te' lu.} It's a pleasure to be able to chat with you in Na'vi.~\LINK{https://wiki.learnnavi.org/index.php?title=Corpus\#Times\_Online}{wiki}
    \B{Kem\-ìri a nga\-ru prr\-te' ke lu, tsa\-kem rä\-'ä si\-vi ay\-la\-he\-ru.} As for any action that you don't like, don't act that way to others.~\LINK{https://wiki.learnnavi.org/index.php?title=Corpus\#Good\_Morning\_America}{wiki}
    \B{Fì\-trr\-ä ftxo\-zä\-ri, sìl\-pey oe, ay\-nga\-ru prr\-te' li\-vu.} I hope you enjoy today’s celebration.~\LINK{https://forum.learnnavi.org/language-updates/trr-rrtaya/}{ln} \\
    \textsc{Derivatives}: \ding{43}
    \on{1}.~\HL{nIprrte'}{\B{nìprrte'}}, with pleasure.
    \on{2}.~\HL{tIprrte'}{\B{tìprrte'}}, pleasure.
  }

  % prrwll
  \Headword{prrwll}{\B{\cp{prr}wll}}{
    \I{["p\s{r}.w\s{l}]} \emph{n.} moss.
    \B{Oe\-ru ke tsun li\-vam ke\-'u lor to Ey\-wa\-'eveng\-ä na'\-rìng a lew sä\-po\-li fa prr\-wll.} Nothing could appear more beautiful to me than a Pandoran forest that has covered in moss.~\LINK{https://forum.learnnavi.org/language-updates/prrwll-moss-the-comparative-to-construct/}{ln}
    (\ding{43}~\HL{prrnen}{\B{prr\-nen}}, \HL{'ewll}{\B{'e\-wll}})
  }

  % pu-
  \Headword{pu-}{\B{pu-}}{
    \I{[pu.]} \emph{num. pref.} six (\emph{base form of \ding{43}~\HL{pukap}{\B{pu\-kap}} to form higher numbers or fractions}).
    \B{\cp{pu}vol}, \on{48} (decimal), \on{60} (octal);
    \B{pu\cp{pxì}}, one-sixth, $\frac{1}{6}$;
    \B{pu\-pxì a\-tsìng}, four\hyp{}sixth, $\frac{4}{6}$.
  }

  % puk
  \Headword{puk}{\B{puk}}{
    \I{[puk\textcorner]} \emph{or} \I{[pUk\textcorner]} (\textsc{rn}: \B{pùk} \I{[pUk\textcorner]}) \emph{n.}, \textsc{lw} book.
    \B{Lu oe\-ru puk.} I have a book.~\LINK{https://naviteri.org/2012/10/audio-and-video-learning-materials-for-navi-101/}{b}
    \B{Puk\-it a\-ke\-tsran ivi\-nan.} Read any book at all.~\LINK{https://naviteri.org/2013/03/whoever-whatever-whenever/}{b}
    \B{Tsun Txil\-te pam\-rel\-it ivi\-nan; ta\-fral puk\-ot a\-nutx mu\-nge fra\-tseng.} Txilte knows how to read; therefore she brings a thick book wherever she goes.~\LINK{https://naviteri.org/2011/11/\%E2\%80\%99a\%E2\%80\%99awa-ayli\%E2\%80\%99u-amip-ni\%E2\%80\%99aw-\%E2\%80\%94-a-few-new-words-only/}{b}
    \B{Kì\-rey\-sì\-ri lu tsa\-puk\-ä ay\-vur a te\-ri 'Rr\-ta txan\-wa\-we nì\-ngay, slä oe\-ri, hu\-fwa el\-tur tì\-txen si, lu ral\-nga' nì\-'aw.} For Grace, the stories in that book relating to Earth are personally meaningful, but for me, although interesting, they’re simply instructive.~\LINK{https://naviteri.org/2012/07/fivospxiya-aylifyavi-amip-this-months-new-expressions-pt-2/}{b} \\
    \textsc{Derivatives}: \ding{43}
    \on{1}.~\HL{puktsyIp}{\B{puk\-tsyìp}}, booklet.
    \on{2}.~\HL{lI'upuk}{\B{lì\-'u\-puk}}, dictionary.
  }

  % pukap
  \Headword{pukap}{\B{\cp{pu}kap}}{
    \I{["pu.kap\textcorner]} \emph{num.} six. (\emph{base}: \ding{43}~\HL{pu-}{\B{pu-}}, \emph{rest}: \ding{43}~\HL{-fu}{\B{\mbox{-fu}}}) \\
    \textsc{Derivatives}: \ding{43}
    \on{1}.~\HL{vospxIpuk}{\B{Vo\-spxì\-puk}}, June.
  }

  % puktsyìp
  \Headword{puktsyIp}{\B{\cp{puk}tsyìp}}{
    \I{["puk\textcorner.\t{ts}jIp\textcorner]} \emph{or} \I{["pUk\textcorner.]} (\textsc{rn}: \B{\cp{pùk}tsyìp} \I{["pUk\textcorner.\t{tS}Ip\textcorner]}) \emph{n.} booklet, pamphlet.
    (\ding{43}~\HL{puk}{\B{puk}}, \HL{-tsyIp}{\B{-tsyìp}})
  }

  % pum
  \Headword{pum}{\B{pum}}{
    \I{[pum]} \emph{or} \I{[pUm]} (\textsc{rn:} \B{pùm} \I{[pUm]})
    \emph{n.} (``dummy noun,'' \emph{not marked for number})
    (\emph{a.}) ``one'' (\emph{always used with a narrowing descriptor})
    \B{Lam set fwa Saw\-tu\-te a\-kawng ho\-lum, pum a\-sìl\-tsan 'ì\-'awn.} It now seems that the bad Skypeople have left, the good ones remain.~\LINK{https://naviteri.org/2010/07/ting-mikyun-fte-tslivam-listening-comprehension-1-3/}{b}
    \B{Ftu\-e lu fwa ta\-ron ngong\-a io\-ang\-it to fwa ta\-ron pum\-it a lu wa\-lak sì win.} It's easier to hunt lethargic animals than to hunt perky, speedy ones.~\LINK{https://naviteri.org/2011/10/more-vocabulary-a-bit-of-grammar/}{b}
    \B{Lu nga\-ru swi\-zaw nì\-hawng srak? Tìng oer pum\-it a\-mrr!} Do you have too many arrows? Give me five.~\LINK{https://naviteri.org/2025/01/tsaliu-alu-pum-aysaomum-asawnung-additional-information-about-pum/}{b}
    \B{Tìng oer pum\-it a\-'aw!} Give me one.~\LINK{https://naviteri.org/2025/01/tsaliu-alu-pum-aysaomum-asawnung-additional-information-about-pum/}{b}
    (\emph{a.1 optionally with} \B{fota} or \B{sawta}) ``of them''.
    \B{Lu nga\-ru 'e\-veng a\-pu\-kap, slä smon oer fo\-ta pum a\-mrr nì\-'aw.} You have six children, but I only know five of them.~\LINK{https://naviteri.org/2025/01/tsaliu-alu-pum-aysaomum-asawnung-additional-information-about-pum/}{b}
    (\emph{b. to convey emphasis})
    \B{Tsa\-tu\-té lu lor\-a pum a\-lor.} That woman is an extremely beautiful one.~\LINK{https://naviteri.org/2013/02/vospxi-ayol-posti-apup-short-post-for-a-short-month/}{b}
    (\emph{c. with possessive pronouns to form ``disjunctive possessives''}) ``possession, thing possessed.''
    \B{Kel\-ku nge\-yä lu tsawl; pum oe\-yä lu hì\-'i.} Your house is large; mine is small.~\LINK{https://naviteri.org/2010/06/first-post/}{b}
    \B{Ä\-txä\-le su\-yi o\-he pi\-vawm, pe\-o\-lo' lu\-yu pum nge\-nge\-yä?} May I ask what tribe you belong to? (\emph{overly formal, stilted Na'vi})~\LINK{https://naviteri.org/2010/09/getting-to-know-you-part-2/}{b}
    (\emph{d. idiom}) \B{Pum nge\-yä.} (\emph{in response to thanks}) ``You're welcome,'' i.e. The thanks belong to you.~\LINK{https://naviteri.org/2010/07/diminutives-conversational-expressions/}{b}
    (\emph{e. can be ommitted in coll. conversation})
    \B{Fì\-tsko lu (pum) oe\-yä.} This bow is mine.~\LINK{https://naviteri.org/2021/04/mipa-ayliu-mipa-sioeykting-new-words-new-explanations/}{b}
  }

  % pung
  \Headword{pung}{\B{pung}}{
    \I{[puN]} \emph{or} \I{[pUN]} \emph{vtr.} hurt, injure.
    \B{Ngal pe\-rung oet fì\-fya pe\-lun?} Why are you hurting me like this?~\LINK{https://naviteri.org/2021/09/aawa-ayliu-amip-a-few-new-words/}{b}
    \B{Te\-ya si oer fwa ngal paw\-nu\-nga ay\-ioang\-it zong.} It moves me that you save injured animals.~\LINK{https://naviteri.org/2021/09/aawa-ayliu-amip-a-few-new-words/}{b}
    (\ding{43}~\HL{tIsraw}{\B{tì\-sraw si}})
  }

  % pup
  \Headword{pup}{\B{pup}}{
    \I{[pup\textcorner]} \emph{or} \I{[pUp\textcorner]} (\textsc{rn}: \B{pùp} \I{[pUp\textcorner]}) \emph{adj.} short (\emph{physical length}).
    \B{Vo\-spxì A\-yol, Po\-stì A\-pup.} Short Post for a Short Month.~\LINK{https://naviteri.org/2013/02/vospxi-ayol-posti-apup-short-post-for-a-short-month/}{b}
    \B{Lu srey a\-pup ta Nga\-ri zì\-sìt a\-mip li\-vu le\-fpom.} It's a short version of \emph{Ngari zìsìt amip livu lefpom}.~\LINK{https://naviteri.org/2017/12/zisit-amip-lefpom-ma-eylan-happy-new-year-friends/\#comment-27964}{b}
    (\ding{214}~\HL{ngim}{\B{ngim}}) \\
    \textsc{Derivatives}: \ding{43}
    \on{1}.~\HL{ngimpup}{\B{ngimpup}}, length.
  }

  % pupxì
  \Headword{pupxI}{\B{pu\cp{pxì}}}{
    \I{[pu."p'I]} \emph{num.} one\hyp{}sixth, $\frac{1}{6}$.
  }

  % puve
  \Headword{puve}{\B{\cp{pu}ve}}{
    \I{["pu.vE]} \emph{adj., ord. num.} sixth. \\
    \textsc{Derivatives}: \ding{43}
    \on{1}.~\HL{trrpuve}{\B{trr\-pu\-ve}}, Friday.
  }

  % puvol
  \Headword{puvol}{\B{\cp{pu}vol}}{
    \I{["pu.vol]} \emph{num.} \on{48} (decimal), \on{60} (octal).
    (\ding{43}~\HL{pu-}{\B{pu-}}, \HL{vol}{\B{vol}})
  }

  % puwup
  \Headword{puwup}{\B{\cp{pu}wup}}{
    \I{["pu.wup\textcorner]} \emph{or} \I{[.wUp\textcorner]} (\textsc{rn}: \B{\cp{pu}wùp} \I{["pu.wUp\textcorner]}) \emph{vin.} bounce.
    \B{Rum 'aw\-lo po\-lu\-wup 'rr\-ko ne\-to.} The ball bounced once and rolled away.~\LINK{https://naviteri.org/2023/12/mipa-zisit-ayliu-amip-new-year-new-words/}{b}
    (\emph{a. idiom}) \B{Ngä\-zìk fwa fkol rum\-it a\-ku\-'up pey\-ku\-wup.} \emph{It's difficult to make s.o. stubborn or inept do what you want them to} [lit.: it's hard to bounce a heavy ball].~\LINK{https://naviteri.org/2023/12/mipa-zisit-ayliu-amip-new-year-new-words/}{b}
    \\

  }



}



% ===== Section PX =====

 \pdfbookmark[0]{PX}{PX}
\begin{center}
\section*{PX \I{[p']} (\emph{PxePx})}
\end{center}

\begin{multicols}{2}

% pxan
\hi
\HT{pxan}{\B{pxan}} \I{[p'an]} \emph{adj.} worthy, worthy of\slash worth for (\emph{with the topic}).   
\B{Pxan li\-vu txo nì\-'aw oe nga\-ri.} Only if I am worthy of you.~\LINK{http://wiki.learnnavi.org/index.php/Hunting_Song}{wiki}
\B{Sìl\-pey oe, aw\-nge\-yä lì'\-fya\-olo'\-ìri fì\-pì\-lok lì\-ye\-vu pxan.} I hope, this blog will be worth for our language community.~\LINK{http://naviteri.org/2010/06/first-post/}{b}
\B{Mll\-te oe: am\-'a\-lu\-ke za\-ma\-yu\-nge tsa\-tì\-pu\-sey\-ìl hum\-it a\-pxan.} I agree: without a doubt that hope will bring worthy results.~\LINK{http://naviteri.org/2013/08/fmawn-a-fkol-kilmulat-this-just-in/comment-page-1/\#comment-2478}{b}
(\emph{a. idiom}) \B{Ke pxan.} ``Don't mention it.\slash Thank you.'' (\emph{in response to a compliment}) [lit.: (I'm) not worthy].~\LINK{http://naviteri.org/2010/07/diminutives-conversational-expressions/}{b}

% pxasìk
\hi
\HT{pxasIk}{\B{\cp{pxa}sìk}} \I{["p'a.sIk\textcorner]} \emph{idiom, vulg.} Screw that! No way!~\LINK{http://www.lemondrop.com/2010/01/26/avatar-pick-up-lines-talk-dirty-in-Navi/}{ld}

% pxasul
\hi
\HT{pxasul}{\B{\cp{pxa}sul}} \I{["p'a.sul]} \emph{or} \I{[.sUl]} (\textsc{rn}: \B{\cp{pxa}sùl} \I{["ba.sUl]}) \emph{adj.} fresh, appealing as food.

% pxazang
\hi 
\HT{pxazang}{\B{\cp{pxa}zang}} \I{["p'a.zaN]} \emph{n., fauna} A\-ku\-la (\emph{shark\hyp{}like aquatic creature}).

% pxaw
\hi
\HT{pxaw}{\B{pxaw}} \I{[p'aw]} \emph{adp.} around. 
\B{Nì\-trr\-trr yom Na'\-vil wu\-tsot 'aw\-si\-teng pxaw yll\-txep.} The Na'vi regularly eat dinner together around a communal fire.~\LINK{http://naviteri.org/2012/07/fivospxiya-aylifyavi-amip-this-months-new-expressions-pt-2/}{b} \\
\textsc{Derivatives}: \ding{43} 
\on{1}.~\HL{pxawpa}{\B{pxaw\-pa}}, perimeter, border. 
\on{2}.~\HL{pxawpun}{\B{pxaw\-pxun}}, armband. 
\on{3}.~\HL{pxawngip}{\B{pxaw\-ngip}}, environment. 
\on{4}.~\HL{pxawtxap}{\B{pxaw\-txap}}, squeeze.
\on{5}.~\HL{pxawtok}{\B{pxaw\-tok}}, surround.

% pxawngip
\hi
\HT{pxawngip}{\B{\cp{pxaw}ngip}} \I{["p'aw.Nip\textcorner]} \emph{n.} environment. 
\B{Fwa 'e\-fu ma\-wey sì nit\-ram mì paw\-ngip a\-mip krr\-nekx.} It takes time (for them) to feel calm and happy in a new environment.~\LINK{http://naviteri.org/2015/09/tskxekengtsyip-a-mikyunfpi-a-little-listening-exercise-2/}{b}
\B{Na fra\-ve\-'o ay\-ru\-sey\-ä, far\-wun\-tu Ey\-wa\-'ev\-eng\-ä sä\-fpìl\-yewn nì\-'aw\-nì\-'aw fa ay\-fya\-'o a ta\-re fey pxaw\-ngip\-it.} As with any ecosystem, the inhabitants of Pandora communicate in ways unique to their environment.~\LINK{https://naviteri.org/2026/06/tskxekeng-a-mikyunfpi-pamrel-listening-exercise-text/}{b}
(\ding{43}~\HL{pxaw}{\B{pxaw}}, \HL{ngip}{\B{ngip}})

% pxawpa
\hi
\HT{pxawpa}{\B{\cp{pxaw}pa}} \I{["p'aw.pa]} \emph{n.} perimeter, circumference, border. 
\B{Tu\-tan a'\-ewan tar\-mìng na\-ri pxaw\-pa\-ro txan\-lo\-kxe\-yä.} A young man was watching at the border of his country.~\LINK{http://naviteri.org/2020/04/awa-tskxekengtsyip-a-mikyunfpi-niul-one-more-little-listening-exercise/}{b}
(\ding{43}~\HL{pxaw}{\B{pxaw}}, \HL{pa'o}{\B{pa'o}})

% pxawpxun
\hi
\HT{pxawpxun}{\B{\cp{pxaw}pxun}} \I{["p'aw.p'un]} \emph{or} \I{[.p'Un}] \emph{n.} armband.   
\B{Po\-ri pxun\-pxaw lu pxaw\-pxun.} Around his arm is an armband.~\LINK{http://naviteri.org/2011/08/new-vocabulary-clothing/}{b}
\B{Ngal mo\-lay' pxaw\-pxun\-it Lo\-ak\-ä a krr pe\-kxin\-um?} When you tried on Loak's armband, how tight was it?~\LINK{http://naviteri.org/2012/07/meetings-waterfalls-and-more/}{b}
(\ding{43}~\HL{pxaw}{\B{pxaw}}, \HL{pxun}{\B{pxun}})

% pxawtok
\hi 
\HT{pxawtok}{\B{\cp{pxaw}tok}} \I{["p'aw.t$\cdot$ok]} \emph{vtr., inf.}\on{22} surround. 
\B{Pxaw\-to\-lok sna\-nan\-tang\-ìl ye\-rik\-it.} The nan\-tang pack surrounded the ye\-rik.~\LINK{https://naviteri.org/2024/02/mipa-ayliu-si-aylifyavi-niul-more-new-words-and-expressions/}{b}  
(\ding{43}~\HL{pxaw}{\B{pxaw}}, \HL{tok}{\B{tok}})

% pxawtxap
\hi 
\HT{pxawtxap}{\B{pxaw\cp{txap}}} \I{[p'aw."t'$\cdot$ap\textcorner]} \emph{vtr., inf.}\on{22} squeeze. 
\B{Po\-el oey tsyokx\-it pxaw\-txo\-lap a krr, pol\-txe nì\-fnu san nga yaw\-ne lu oer.} When she squeezed my hand, she silently said she loved me.~\LINK{http://naviteri.org/2020/05/aawa-liu-amip-si-vurway-alor-a-few-new-words-and-a-beautiful-poem/}{b}  
(\ding{43}~\HL{pxaw}{\B{pxaw}}, \HL{txap}{\B{txap}})

% pxay
\hi
\HT{pxay}{\B{pxay}} \I{[p'aj]} \emph{adj.} many (\emph{for countable nouns}). 
(\emph{a. mainly used with the singular form of the noun}) 
\B{Na'\-rìng\-ìl nga' pxay\-a io\-ang\-it.} The forest contains many animals.~\LINK{http://naviteri.org/2011/05/weather-part-2-and-a-bit-more-2/}{b}
\B{Fì\-slär\-mì tsun fko sti\-vawm ngam\-it a\-pxay.} You can hear a lot of echoes in this cave.~\LINK{http://naviteri.org/2012/01/mipa-zisit-ayliu-amip-new-words-for-the-new-year/}{b}
(\emph{b. coll. with the plural form}) \B{Fu\-ri\-a fya\-pe fkol se\-rar lì'\-fya\-ti aw\-nge\-yä, le\-iu set nì\-law pxay\-a su\-te a oe\-to lu sìl\-tsan!} Concerning the use of our language, clearly many people are now better than me, I'm happy to say.~\LINK{http://naviteri.org/2010/08/a-na\%E2\%80\%99vi-alphabet/comment-page-1/\#comment-191}{b}
(\ding{214}~\HL{hol}{\B{hol}}, \ding{43}~\HL{'a'aw}{\B{'a'aw}}) \\
\textsc{Derivatives}: \ding{43}
\on{1}.~\HL{nIpxay}{\B{nì\-pxay}}, many, in great number. 
\on{2}.~\HL{holpxay}{\B{hol\-pxay}}, number. 
\on{3}.~\HL{pxayopin}{\B{pxay\-opin}}, colorful, multi\hyp{}colored. 
\on{4}.~\HL{pxayzekwA}{\B{pxay\-zek\-wä}}, spiny whips.

% pxayopin
\hi
\HT{pxayopin}{\B{\cp{pxay}opin}} \I{["p'aj.o.pin]} \emph{adj.} colorful, multi\hyp{}colored, variegated. 
\B{'En si oe, lor\-a tsa\-fkxi\-le\-ri a\-pxay\-opin so\-lar Tse\-nul srok\-it a\-vo\-zam.} I would guess that Tsenu used a thousand (lit. \on{512}) beads for that beautiful multi\hyp{}colored bib necklace.~\LINK{http://naviteri.org/2013/07/pximaw-tsawlultxa-just-after-the-meet-up/}{b}
(\ding{43}~\HL{pxay}{\B{pxay}}, \HL{'opin}{\B{'opin}})

% pxayzekwä
\hi
\HT{pxayzekwA}{\B{pxay\cp{zek}wä}} \I{[p'aj."zEk\textcorner.w\ae]} \emph{n.}, \ding{96} spiny whips, \emph{phytotinium poydactylum}. 
(\ding{43}~\HL{pxay}{\B{pxay}}, \HL{zekwA}{\B{zek\-wä}})

% pxe+
\hi
\HT{pxe}{\B{pxe}\texttt{+}} \I{[p'E.]} \emph{trial pref.} three (of something).

% pxefo
\hi
\HT{pxefo}{\B{pxe\cp{fo}}} \I{[p'E."fo]} (\emph{gen.} \B{pxe\-\cp{fe}\-yä}) \emph{pn.} they (three), the three of them, (3\textsuperscript{rd} person trial).
(\emph{a. subjective}) \B{Sra\-ke pxe\-fo li po\-lä\-ha\-tsem?} I wonder if the three of them have already arrived.~\LINK{http://naviteri.org/2011/10/more-vocabulary-a-bit-of-grammar/}{b}
(\emph{b. genitive}) \B{Pxe\-fe\-yä tì\-tstun\-wi\-ri aw\-nga ira\-yo si\-vi ko!} Let's thank the three of them for their kindness.~\LINK{http://naviteri.org/2012/07/tskxekengtsyip-a-mikyunfpi-a-little-listening-exercise/}{b}
(\emph{c. dative}) \B{Pxe\-fo\-ru oe srung so\-li, kuma oe\-ru set pxe\-fo srung se\-ri.} I helped the three of them, so they're now helping me.~\LINK{http://naviteri.org/2012/06/spring-vocabulary-part-3/}{b}
(\ding{43}~\HL{pxe}{\B{pxe}}\texttt{+}, \HL{po}{\B{po}})

% pxek
\hi
\HT{pxek}{\B{pxek}} \I{[p'Ek\textcorner]} \emph{vtr.} kick, shove (\emph{can be used with} \B{kinamfa} \emph{and} \B{pxunfa} \emph{to clear meaning}). 
\B{Ye\-rik\-ìl nan\-tang\-it pxo\-lek fte hi\-vi\-fwo.} The hexapede kicked the viperwolf in order to flee.~\LINK{http://naviteri.org/2014/05/mipa-ayliu-mipa-aysafpil-new-words-new-ideas/}{b}
\B{Po pxun\-fa pxek tsa\-krr me\-fo zup ne\-kll!} He shoved the two of them and they fell down.~\LINK{http://naviteri.org/2014/05/mipa-ayliu-mipa-aysafpil-new-words-new-ideas/}{b}

% pxel
\hi
\HT{pxel}{\B{pxel}} \I{[p'El]} \emph{adp.} like, as. 
\B{Kem\-lì\-'ut alu var fkol sar pxel pum alu new.} One uses the verb \emph{var} like \emph{new}.~\LINK{http://naviteri.org/2011/03/word-order-and-case-marking-with-modals/comment-page-1/\#comment-582}{b}
(\emph{a. required with} \HL{zet}{\B{zet}}) \B{Pey\-ral\-ìl zet wur\-a wu\-tsot a\-'aw\-nem pxel sngel.} Peyral treats a cool cooked meal like garbage.~\LINK{http://naviteri.org/2012/10/mipa-vospxi-mipa-ayliu-new-words-for-the-new-month/}{b}
(\emph{b. with} \HL{nIftxan}{\B{nìftxan}} \emph{and adj. or adv. to mean}) as \emph{adj.\slash adv.} as.
\B{Oe tul nì\-ftxan nì\-win pxel nga.} I run as fast as you.~\LINK{https://forum.learnnavi.org/language-updates/eapressions-of-\%27like-we\%27/}{ln}
(\ding{43}~\HL{na}{\B{na}})

% pxelo
\hi
\HT{pxelo}{\B{\cp{pxe}lo}} \I{["p'E.lo]} \emph{adv.} thrice, three times. 
(\ding{43}~\HL{pxe}{\B{pxe}}\texttt{+}, \HL{alo}{\B{alo}})

% pxen
\hi
\HT{pxen}{\B{pxen}}\textsuperscript{\on{1}} \I{[p'En]} \emph{n.} (item of) functional clothing.
\B{Pen\-it yem\-stokx!} Get dressed!~\LINK{http://naviteri.org/2011/08/new-vocabulary-clothing/}{b}
\B{Fu\-ta nga\-ta tel pxen\-yoa srät, mll\-'e\-i\-an oel.} I'm happy to take cloth from you for finished garments.~\LINK{http://naviteri.org/2014/02/barter-and-exchange/}{b}
(\ding{43}~\HL{ioi}{\B{ioi}}) \\
\textsc{Derivatives}: \ding{43}
\on{1}.~\HL{lukpen}{\B{luk\-pen}}, naked.

% pxen
\hi
\B{pxen}\textsuperscript{\on{2}} : \HL{'en}{\B{'en}}

% pxenga
\hi
\HT{pxenga}{\B{pxe\cp{nga}}} \I{[p'E."Na]} \emph{pn.} you (three), the three of you (2\textsuperscript{nd} person trial).
(\emph{a. genitive}) \B{Pxe\-nge\-yä tsul\-fä lì'\-fya\-yä le\-Na'\-vi oe\-ru te\-ya si.} Their mastery of the Na'vi language fills me with joy.~\LINK{http://naviteri.org/2010/12/navi-writing-contest-the-first-place-winners/}{b}
(\emph{b. dative}) \B{Ira\-yo nì\-txan pxe\-nga\-ru!} Thank you three very much!~\LINK{http://naviteri.org/2012/10/audio-and-video-learning-materials-for-navi-101/}{b}
(\ding{43}~\HL{pxe}{\B{pxe}}\texttt{+}, \HL{nga}{\B{nga}})

% pxengenga
\hi
\HT{pxengenga}{\B{pxenge\cp{nga}}} \I{[p'E.NE."Na]} \emph{pn.} you (three), the three of you (2\textsuperscript{nd} person trial; deferential or ceremonial form).
(\ding{43}~\HL{pxe}{\B{pxe}}\texttt{+}, \HL{ngenga}{\B{nge\-nga}})

% pxesrram
\hi
\HT{pxesrram}{\B{pxesrr\cp{am}}} \I{[p'E.s\s{r}."{}am]} \emph{n., adv.} three days ago. 
(\ding{43}~\HL{pxe}{\B{pxe}}\texttt{+}, \HL{trram}{\B{trr\-am}})

% pxesrray
\hi
\HT{pxesrray}{\B{pxesrr\cp{ay}}} \I{[p'E.s\s{r}."{}aj]} \emph{n., adv.} three days from now.
(\ding{43}~\HL{pxe}{\B{pxe}}\texttt{+}, \HL{trray}{\B{trr\-ay}})

% pxevol
\hi
\HT{pxevol}{\B{\cp{pxe}vol}} \I{["p'E.vol]} \emph{num., adj.} \on{30} (octal), \on{24} (decimal).
\B{Oe\-ri so\-la\-lew zì\-sìt a\-pxe\-vol.} I'm 24 years old.~\LINK{http://naviteri.org/2010/09/getting-to-know-you-part-2/}{b} 
(\ding{43}~\HL{pxe}{\B{pxe}}\texttt{+}, \HL{vol}{\B{vol}})

% pxevotrrpxì
\hi 
\HT{pxevotrrpxI}{\B{pxevotrr\cp{pxì}}} \I{[p'E.vo.t\s{r}."p'I]} \emph{n.} hour (in the \on{24}\hyp{}human cycle, \emph{often shortened to} \ding{43}~\HL{trrpxI}{\B{trrpxì}}).  
(\ding{43}~\HL{pxevol}{\B{pxevol}}, \HL{trrpxI}{\B{trrpxì}})

% pxey
\hi
\HT{pxey}{\B{pxey}} \I{[p'Ej]} \emph{num., adj.} three (\emph{base:} \B{pxey-}, \emph{rest:} \HL{-pey}{\B{-pey}}). 
\B{Za\-ya\-'u Saw\-tu\-te fte aw\-nga\-ti ski\-va\-'a kay zì\-sìt a\-pxey (fu: pxey\-a zì\-sìt\-kay)!} The Sky People will come to destroy us three years from now!~\LINK{http://naviteri.org/2011/09/miscellaneous-vocabulary/}{b} \\
\textsc{Derivatives}: \ding{43}
\on{1}.~\HL{vospxey}{\B{Vo\-spxey}}, March.

% pxeyve
\hi
\HT{pxeyve}{\B{\cp{pxey}ve}} \I{["p'Ej.vE]} \emph{num., adj.} third.
\B{ta\-ron\-yu a\-pxey\-ve\slash pxey\-ve\-a ta\-ron\-yu}, hunter number 3, the third hunter.~\LINK{http://forum.learnnavi.org/language-updates/update-on-counting/}{ln}

% pxi
\hi
\HT{pxi}{\B{pxi}} \I{[p'i]} \emph{adj.} sharp (\emph{as a point}). 
\B{Na\-ri si! Tsa\-ioang\-ur lu pxi\-a me\-slukx!} Careful! That animal has two sharp horns!~\LINK{http://naviteri.org/2016/06/mrrvola-lifyavi-amip-forty-new-expressions/}{b} 
(\ding{214}~\HL{fwem}{\B{fwem}}, \ding{43}~\HL{litx}{\B{litx}}) \\
\textsc{Derivatives}: \ding{43}
\on{1}.~\HL{nIpxi}{\B{nì\-pxi}}, especially, pointedly. 
\on{2}.~\HL{pximaw}{\B{pxi\-maw}}, right after. 
\on{3}.~\HL{pxiset}{\B{pxi\-set}}, right now. 
\on{4}.~\HL{pxisre}{\B{pxi\-sre}}, right before. 
\on{5}.~\HL{pxiswawam}{\B{pxi\-swaw\-am}}, just a moment ago. 
\on{6}.~\HL{pxiswaway}{\B{pxi\-swaw\-ay}}, in just a moment from now. 
\on{7}.~\HL{pxiut}{\B{pxi\-ut}}, razor palm. 
\on{8}.~\HL{pxiwll}{\B{pxi\-wll}}, hermit bud plant. 
\on{9}.~\HL{kllpxiwll}{\B{kll\-pxi\-wll}}, lionberry. 
\on{10}.~\HL{pxiye'rIn}{\B{pxi\-ye'\-rìn}}, immediately. 

% pxim
\hi
\HT{pxim}{\B{pxim}} \I{[p'im]} \emph{adj.} erect, upright. \\
\textsc{Derivatives}: \ding{43}
\on{1}.~\HL{nIpxim}{\B{nì\-pxim}}, errectly, rightly. 

% pximaw
\hi
\HT{pximaw}{\B{pxi\cp{maw}}} \I{[p'i."maw]} \emph{adp.} right after.
\B{Pxi\-maw Tsawl\-ul\-txa}, Just After the Meet\hyp{}up.~\LINK{http://naviteri.org/2013/07/pximaw-tsawlultxa-just-after-the-meet-up/}{b}
(\ding{43}~\HL{pxi}{\B{pxi}}, \HL{maw}{\B{maw}})

% pxir
\hi
\HT{pxir}{\B{pxir}} \I{[p'iR]} \emph{n.}, \textsc{lw} beer. 
\B{Su\-nu nga\-ru pxir srak?} Do you like beer?~\LINK{http://naviteri.org/2012/10/audio-and-video-learning-materials-for-navi-101/}{b}
\B{Pxir\-it fu swo\-at näk rä\-'ä; ni\-väk pay\-it nì\-'aw.} Don't drink beer or spirits; just drink water.~\LINK{http://naviteri.org/2014/07/tsawlultxamaw-a-ayliu-post-meetup-words/}{b}

% pxiset
\hi
\HT{pxiset}{\B{pxi\cp{set}}} \I{[p'i."sEt\textcorner]} \emph{adv.} right now.
\B{Pxi\-set nge\-yä tska\-lep\-it kxel\-tek!} Pick up your crossbow right now!~\LINK{http://naviteri.org/2012/01/mipa-zisit-ayliu-amip-new-words-for-the-new-year/}{b}
\B{Zun oe pxi\-set tir\-va\-ron, zel oe 'i\-ve\-fu nit\-ram.} If I were hunting right now, then I would be happy.~\LINK{http://naviteri.org/2013/04/zun-zel-counterfactual-conditionals/}{b}
(\ding{43}~\HL{pxi}{\B{pxi}}, \HL{set}{\B{set}})

% pxisre
\hi
\HT{pxisre}{\B{pxi\cp{sre}}} \I{[p'i."{}sRE]} \emph{adp.}\texttt{+} right before (\emph{temporal}).
(\ding{43}~\HL{pxi}{\B{pxi}}, \HL{sre}{\B{sre}}) 

% pxiswawam
\hi
\HT{pxiswawam}{\B{pxiswaw\cp{am}}} \I{[p'i.swaw."{}am]} \emph{adv.} just a moment ago.
(\ding{43}~\HL{pxi}{\B{pxi}}, \HL{swaw}{\B{swaw}})

% pxiswaway
\hi
\HT{pxiswaway}{\B{pxiswaw\cp{ay}}} \I{[p'i.swaw."{}aj]} \emph{adv.} in just a second from now.
(\ding{43}~\HL{pxi}{\B{pxi}}, \HL{swaw}{\B{swaw}}; \HL{pxiye'rIn}{\B{pxi\-ye'\-rìn}})

% pxiut
\hi
\HT{pxiut}{\B{\cp{pxi}ut}} \I{["p'i.ut\textcorner]} \emph{or} \I{[.Ut\textcorner]} \emph{n.}, \ding{96} razor palm, ``sharp tree'', \emph{saltcellar gracilis}. 
(\ding{43}~\HL{pxi}{\B{pxi}}, \HL{utral}{\B{utral}})

% pxiwll
\hi
\HT{pxiwll}{\B{\cp{pxi}wll}} \I{["p'i.w\s{l}]} \emph{n.}, \ding{96} hermit bud plant, \emph{cynaroidia discolor}. 
(\ding{43}~\HL{pxi}{\B{pxi}}, \HL{'ewll}{\B{'ewll}})

% pxiye'rìn
\hi
\HT{pxiye'rIn}{\B{pxi\cp{ye'}rìn}} \I{[p'i."jEP.RIn]} \emph{adv.} immediately (\emph{but not as soon as} \B{pxi\-swaw\-ay}).
\B{Fì\-fne\-tì\-fkey\-tok\-ìl fngo' fu\-ta kem si\-vi fko pxi\-ye'\-rìn.} This kind of situation requires immediate action.~\LINK{http://naviteri.org/2012/01/mipa-zisit-ayliu-amip-new-words-for-the-new-year/}{b}
(\ding{43}~\HL{pxi}{\B{pxi}}, \HL{ye'rIn}{\B{ye'rìn}}; \HL{pxiswaway}{\B{pxi\-swaw\-ay}})

% -pxì
\hi
\HT{pxI}{\B{-pxì}} \I{[.p'I]} \emph{num. suf. to form fractions with the base form of numbers}. 
\B{ki\-pxì}, one\hyp{}seventh, $\frac{1}{7}$. 
(\ding{43}~\HL{hapxI}{\B{ha\-pxì}})

% pxìm
\hi
\HT{pxIm}{\B{pxìm}} \I{[p'Im]} \emph{adv.} often.   
\B{Nga za\-'u fì\-tseng pxìm srak?} Do you come here often?~\LINK{http://www.lemondrop.com/2010/01/26/avatar-pick-up-lines-talk-dirty-in-Navi/}{ld}
\B{Oe pxìm päng\-kxo nga\-hu nì\-ron\-srel.} I often talk with you in my imagination.\slash I often imagine I'm talking with you.~\LINK{http://naviteri.org/2011/02/new-vocabulary-part-2/}{b}
\B{Sra\-ke tsun oe fay\-upxa\-ret tsli\-vam nì\-ftue? Tse … ze\-ne pi\-vll\-txe san pxìm tsa\-fya lu sìk.} Can I understand these messages easily? Well…  I must say ``Often that way''.~\LINK{http://forum.learnnavi.org/language-updates/verb-for-write/msg175592/\#msg175592}{ln} \\
\textsc{Derivatives}: \ding{43}
\on{1}.~\HL{lepxìmrun}{\B{le\-pxìm\-run}}, common, often found.

% pxìmun'i
\hi
\HT{pxImun'i}{\B{pxìmun\cp{'i}}} \I{[p'I.m$\cdot$un."P$\cdot$i]} \emph{or} \I{[.mUn.]} (\textsc{rn}: \B{pxìmùn\cp{'i}} \I{[bI.m$\cdot$Un."P$\cdot$i]}) \emph{vtr.}\on{23} divide, cut into parts.
\B{Nì\-trr\-trr pxì\-mun\-'i sam\-si\-yul ay\-swi\-zaw\-it ku\-tu\-ä a\-law\-nä\-txayn sno\-kip nì\-'eng.} Warriors typically share the arrows of their defeated enemies among themselves.~\LINK{http://naviteri.org/2012/01/more-additions-to-the-lexicon/}{b}
\B{Fol tsngan\-it pxì\-mo\-lun\-'i nì\-'eng.} They shared the meat.~\LINK{http://naviteri.org/2012/01/more-additions-to-the-lexicon/}{b}
\B{Fol tsngan\-it pxì\-mo\-lun\-'i nì\-yll.} They shared the meat with the entire clan.~\LINK{http://naviteri.org/2012/01/more-additions-to-the-lexicon/}{b} 
(\ding{43}~\HL{hapxI}{\B{hapxì}}, \HL{mun'i}{\B{mun'i}})

% pxoe
\hi
\HT{pxoe}{\B{\cp{pxo}e}} \I{["p'o.E]} \emph{pn.} we (three), the three of us (not you) (1\textsuperscript{st} person trial, exclusive).
(\emph{a. subjective}) \B{Fwa tso\-lun nga zi\-va'u moe\-yä kel\-ku\-ne fte tsi\-vun pxoe 'aw\-si\-teng ki\-vä\-teng nì\-mun oe\-ru te\-ya so\-lei\-yi nì\-ngay.} That you could come to our house, so that we three could spend time together again, filled me truly with joy.~\LINK{http://naviteri.org/2014/07/tsawlultxamaw-a-ayliu-post-meetup-words/comment-page-1/\#comment-7689}{b}
(\emph{b. topic}) \B{Pxoe\-ri tsa\-vur el\-tur tì\-txen so\-li.} The three of us were interested in that story.~\LINK{http://forum.learnnavi.org/language-updates/txelanit-hivawl/}{ln}
(\ding{43}~\HL{pxe}{\B{pxe}}\texttt{+}, \HL{oe}{\B{oe}})

% pxoeng
\hi
\HT{pxoeng}{\B{pxo\cp{eng}}} \I{[p'o."{}EN]} \emph{pn.} we (three), the three of us (1\textsuperscript{st} person trial, inclusive).
(\emph{a. subjective}) \B{Moe sìl\-pey nì\-teng tsnì pxo\-eng tsä\-pì\-ye\-ve\-'a fì\-tsap ye'\-rìn nì\-mun.} We both also hope that the three of us will see each other again soon.~\LINK{http://naviteri.org/2013/07/pximaw-tsawlultxa-just-after-the-meet-up/comment-page-1/\#comment-2483}{b}
(\ding{43}~\HL{pxe}{\B{pxe}}\texttt{+}, \HL{oeng}{\B{oeng}})

% pxohe
\hi
\HT{}{\B{\cp{pxo}he}} \I{["p'o.hE]} \emph{pn.} we (three), the three of us (not you) (1\textsuperscript{st} person trial, exclusive; deferential or ceremonial form)
% (\emph{a. subjective})
% (\emph{b. agentive})
% (\emph{c. patientive})
% (\emph{d. genitive})
% (\emph{e. dative})
% (\emph{f. topic})
(\ding{43}~\HL{pxe}{\B{pxe}}\texttt{+}, \HL{ohe}{\B{ohe}})

% pxor
\hi
\HT{pxor}{\B{pxor}} \I{[p'oR]} \emph{vin.} explode. 
\B{Txep\-ram pxor a krr, txan\-a txe\-kxum\-pay wrr\-za\-'u.} When a volcano erupts, a lot of lava comes out.~\LINK{http://naviteri.org/2015/03/kap-si-ayunil-saylahe-threats-dreams-and-other-things/}{b} 
\B{Hì\-krr\-o me\-fo kak\-pam lar\-mu maw\-krr\-a pxo\-lor kun\-sìp.} The two of them were deaf for a short time after the gunship exploded.~\LINK{http://naviteri.org/2014/05/mipa-ayliu-mipa-aysafpil-new-words-new-ideas/}{b} \\
\textsc{Derivatives}: \ding{43}
\on{1}.~\HL{pxorpam}{\B{pxor\-pam}}, ejective. 
\on{2}.~\HL{pxorna'}{\B{pxor\-na'}}, episoth. 
\on{3}.~\HL{pxorna'lor}{\B{pxor\-na'\-lor}}, sari.
\on{4}.~\HL{sApxor}{\B{sä\-pxor}}, explosion.

% pxorna'
\hi
\HT{pxorna'}{\B{\cp{pxor}na'}} \I{["p'oR.naP]} \emph{n.}, \ding{96} episoth, ``exploding seed'', \emph{croquembouche columnare}. 
(\ding{43}~\HL{pxor}{\B{pxor}}, \HL{rina'}{\B{ri\-na'}})

% pxorna'lor
\hi
\HT{pxorna'lor}{\B{\cp{pxor}na'lor}} \I{["p'oR.naP.""loR]} \emph{n.}, \ding{96} sari, ``beautiful exploding seed'', \emph{pandargonium cyanellum}. 
(\ding{43}~\HL{pxorna'}{\B{pxor\-na'}}, \HL{lor}{\B{lor}})

% pxorpam
\hi
\HT{pxorpam}{\B{\cp{pxor}pam}} \I{["p'oR.pam]} \emph{n.} ejective consonant. 
(\ding{43}~\HL{pxor}{\B{pxor}}, \HL{pam}{\B{pam}})

% pxul
\hi 
\HT{pxul}{\B{pxul}} \I{[p'ul]} \emph{or} \I{[p'Ul]} \emph{adj.} formidable, imposing. 
\B{Lu Saw\-tu\-te wä\-tu a\-pxul.} The Skypeople are formidable opponents.~\LINK{http://naviteri.org/2020/08/hiia-vur-a-teri-mefalukantsyip-a-little-story-about-two-cats/}{b} \\
\textsc{Derivatives}: \ding{43}
\on{1}.~\HL{tIpxul}{\B{tì\-pxul}}, formidableness, imposingness. 
\on{2}.~\HL{nIpxul}{\B{nì\-pxul}}, formidably, imposingly.

% pxun
\hi
\HT{pxun}{\B{pxun}} \I{[p'un]} \emph{or} \I{[p'Un]} \emph{n.} arm. 
\B{A pe\-yä tì\-txur … Mì pun | Na ay\-vul a\-hu\-saw\-nu.} Whose strength is in our arms | As sheltering branches.~\LINK{http://www.pandorapedia.com/navi/music/ritual_music}{pp} 
\B{Me\-pun\-ìri oe 'efu ngeyn nì\-ngay!} My arms are really tired!~\LINK{http://naviteri.org/2013/07/pximaw-tsawlultxa-just-after-the-meet-up/}{b}
\B{Po pxun\-fa pxek tsa\-krr me\-fo zup ne\-kll!} He shoved the two of them and they fell down.~\LINK{http://naviteri.org/2014/05/mipa-ayliu-mipa-aysafpil-new-words-new-ideas/}{b}
\B{Po\-ri pxun\-pxaw lu pxaw\-pxun.} Around his arm is an armband.~\LINK{http://naviteri.org/2011/08/new-vocabulary-clothing/}{b} \\
\textsc{Derivatives}: \ding{43}
\on{1}.~\HL{pxuntil}{\B{pxun\-til}}, elbow. 
\on{2}.~\HL{pxawpxun}{\B{pxaw\-pxun}}, armband.

% pxuntil
\hi
\HT{pxuntil}{\B{\cp{pxun}til}} \I{["p'un.til]} \emph{or} \I{["p'Un.]} \emph{n.} elbow. 
(\ding{43}~\HL{pxun}{\B{pxun}}, \HL{til}{\B{til}})

\end{multicols}


% ===== Section R =====

 \pdfbookmark[0]{R}{R}
\begin{center}
\section*{R \I{[R]} (\emph{ReR})}
\end{center}

\begin{multicols}{2}

% -r
\hi
\HT{-r}{\B{-r}} \I{[R]} \emph{suf. to mark the dative case of a n. ending in a vowel or the diphthongs aw and ew}.
(\ding{43}~\HL{-ru}{\B{-ru}}, \HL{-ur}{\B{-ur}})

% ra'un
\hi 
\HT{ra'un}{\B{\cp{ra}'un}} \I{["Ra.Pun]} \emph{or} \I{[.PUn]} (\textsc{rn}: \B{\cp{ra}ùn} \I{[Ra.Un]}) \emph{vtr.} surrender s.t., relinquish s.t., give up s.t. 
\B{Fì\-atx\-kxe\-ti ke ra\-ya\-'un ay\-oel kaw\-krr!} We will never give up this land.~\LINK{http://naviteri.org/2021/09/aawa-ayliu-amip-a-few-new-words/}{b} 
(\ding{43}~\HL{velek}{\B{ve\-lek}}) \\
\textsc{Derivatives}: \ding{43} 
\on{1}.~\HL{tIra'un}{\B{tì\-ra\-'un}}, surrender, relinquishment.

% rayl
\hi 
\HT{rayl}{\B{rayl}} \I{[Rajl]} \emph{adj.} feeling pity and sorrow for s.o.'s misfortune (\emph{used with} \ding{43}~\HL{'efu}{\B{'e\-fu}}). 
\B{Po\-ri sa'\-nok tì\-mer\-kup; 'e\-fu oe rayl.} His mother just died; I feel bad for him.~\LINK{https://naviteri.org/2025/06/vospikin-lefpom-happy-july/}{b}

% ral
\hi
\HT{ral}{\B{ral}} \I{[Ral]} \emph{n.} meaning.
\B{Fì\-lì\-'u\-ä ral lu tì\-me\-'em sì tì\-ru\-sey mì hi\-fkey na Nawm\-a Sa'\-nok\-ä ha\-pxì.} The meaning of this word is ``harmony, living in the world as part of the Great Mother.''~\LINK{http://forum.learnnavi.org/language-updates/trr-rrtaya/}{ln}
\B{Ral\-it fay\-lì\-'u\-ä oel ke tsli\-vam.} I don't understand the meaning of these words.~\LINK{http://naviteri.org/2011/05/some-miscellaneous-vocabulary/}{b}
\B{Ra\-lur law so\-li fo oe\-ru.} They made the meaning clear to me.~\LINK{http://naviteri.org/2013/06/looking-back-looking-forward/}{b}
\B{Sra\-ke tsun nga ral\-ur 'en si\-vi?} Can you guess the meaning?~\LINK{http://naviteri.org/2014/03/value-and-worth/comment-page-1/\#comment-2669}{b}
\B{Ral\-ìri ke 'efu oel kea tì\-ke\-teng\-it.} I don't feel any difference concerning the meaning.~\LINK{http://naviteri.org/2013/01/awvea-posti-zisita-amip-first-post-of-the-new-year/comment-page-1/\#comment-1903}{b} \\
\textsc{Derivatives}: \ding{43} 
\on{1}.~\HL{ralpeng}{\B{ral\-peng}}, interpret. 
\on{2}.~\HL{tafral}{\B{ta\-fral}}, therefore. 
\on{3}.~\HL{ralnga'}{\B{ral\-nga'}}, meaningful, instructive. 
\on{4}.~\HL{ralke}{\B{ral\-ke}}, meaningless, devoid of content.
\on{5}.~\HL{tengralI'u}{\B{teng\-ra\-lì\-'u}}, synonym.
\on{6}.~\HL{wAralI'u}{\B{wä\-ra\-lì\-'u}}, antonym.

% ralke
\hi 
\HT{ralke}{\B{\cp{ral}ke}} \I{["Ral.kE]} \emph{adj.} meaningless, devoid of content. 
\B{Txe\-wì ka trro nì\-wotx ftxo\-lu\-lì\-'u, slä ay\-lì\-'u peyä lä\-ngu ral\-ke.} Txewì spoke for an entire day, but sadly, his words were meaningless.~\LINK{http://naviteri.org/2019/08/aawa-lifya-si-lifyavi-amip-a-few-new-words-and-expressions/}{b} 
(\ding{43}~\HL{ral}{\B{ral}}, \HL{luke}{\B{lu\-ke}})

% ralnga'
\hi
\HT{ralnga'}{\B{\cp{ral}nga'}} \I{["Ral.NaP]} \emph{adj.} meaningful, instructive, s.t. from which a lesson can be learned.   
\B{Kì\-rey\-sì\-ri lu tsa\-puk\-ä ay\-vur a te\-ri 'Rr\-ta txan\-wa\-we nì\-ngay, slä oe\-ri, hu\-fwa el\-tur tì\-txen si, lu ral\-nga' nì\-'aw.} For Grace, the stories in that book relating to Earth are personally meaningful, but for me, although interesting, they’re simply instructive.~\LINK{http://naviteri.org/2012/07/fivospxiya-aylifyavi-amip-this-months-new-expressions-pt-2/}{b}
(\ding{43}~\HL{ral}{\B{ral}}, \HL{-nga'}{\B{-nga'}}; \ding{214}~\HL{ralke}{\B{ral\-ke}}) 

% ralpeng
\hi
\HT{ralpeng}{\B{ral\cp{peng}}} \I{[Ral."p$\cdot$EN]} \emph{vtr.}\on{22} interpret, translate. 
\B{Wa\-we\-ti ke tsun fko ral\-pi\-veng.} One can't explain (or: put into words) \emph{wawe}.~\LINK{http://naviteri.org/2011/05/some-miscellaneous-vocabulary/}{b}
\B{Ze\-ne Sa'\-nok fì\-au\-ngia\-ti ral\-pi\-veng.} Mother must interpret this sign.~\LINK{http://forum.learnnavi.org/language-updates/pseudo-canon-from-bd-live-content/}{ln\slash a\textsuperscript{c}}
(\ding{43}~\HL{ral}{\B{ral}}, \HL{peng}{\B{peng}}) \\
\textsc{Derivatives}: \ding{43} 
\on{1}.~\HL{tIralpeng}{\B{tì\-ral\-peng}}, translation, interpretation.

% Ralu
\hi
\B{\cp{Ra}lu} \I{["Ra.lu]} \emph{n.} \emph{male name}. 
\B{Ra\-lu to\-lì\-ran äo vul apxa.} Ralu walked under the large branch.~\LINK{http://forum.learnnavi.org/language-updates/over-under-and-quickly/}{ln}
\B{Ra\-lu\-ri fa\-hew fkan oe\-ru na ye\-rik.} Ralu smells like a hexapede to me.~\LINK{http://naviteri.org/2012/11/renu-ayinanfyaya-the-senses-paradigm/}{b}

% ram
\hi
\HT{ram}{\B{ram}} \I{[Ram]} \emph{n.} mountain.
\B{Ay\-ram A\-lu\-sìng}, the Floating Mountains. \\
\textsc{Derivatives}: \ding{43} 
\on{1}.~\HL{ramtsyIp}{\B{ram\-tsyìp}}, hill. \\
\on{2}.~\HL{txepram}{\B{txep\-ram}}, volcano.

% ramtsyìp
\hi
\HT{ramtsyIp}{\B{\cp{ram}tsyìp}} \I{["Ram.\t{ts}jIp\textcorner]} \emph{n.} hill.
\B{Ul\-txa si\-vi oeng sìn ram\-tsyìp txon\-'ong\-ay.} Let's meet on the hill tomorrow at nightfall.~\LINK{http://naviteri.org/2012/01/mipa-zisit-ayliu-amip-new-words-for-the-new-year/}{b}
(\ding{43}~\HL{ram}{\B{ram}}, \HL{-tsyIp}{\B{-tsyìp}}) 

% ramu
\hi 
\HT{ramu}{\B{Ramu}} \I{[Ra.mu]} \emph{male name}. 
\B{Ra\-mu ke lu txur nì\-txan, slä u\-van\-it yo\-lo\-ra', ti\-'a.} Ramu isn't very strong, but surprisingly, he won the game.~\LINK{http://naviteri.org/2022/11/krr-ao-an-exciting-time/}{b} 

% ramunong
\hi
\HT{ramunong}{\B{ra\cp{mu}nong}} \I{[Ra."mu.noN]} \emph{n.} well.
\ding{43}~\HL{ayvitrayA ramunong}{\B{Ayvitrayä Ramunong}}, the Well of Souls.

% ran
\hi
\HT{ran}{\B{ran}} \I{[Ran]} \emph{n.} (\emph{a.}) intrinsic character or nature, essence, constitution, (\emph{basic nature of s.t. resulting from the totality of its properties}).
\B{Fra\-'u\-ä ran ngä\-pop fa fra\-syon tse\-yä.} The \emph{ran} of each thing arises from the totality of its attributes.~\LINK{http://naviteri.org/2012/10/mipa-vospxi-mipa-ayliu-new-words-for-the-new-month/}{b}
\B{Ran tì\-ru\-sol\-ä pe\-yä lu fyo\-le.} The \emph{ran} of her singing is sublime.~\LINK{http://naviteri.org/2012/10/mipa-vospxi-mipa-ayliu-new-words-for-the-new-month/}{b}
(\emph{b. for people}) personality.
\B{Mun\-txa\-tu\-ri So\-rewn\-ti ke tsun oe mi\-vll\-'an. Ran pe\-yä oe\-ru ke ha'.} I can't accept Sorewn as my spouse. Her personality doesn't suit me.~\LINK{http://naviteri.org/2012/10/mipa-vospxi-mipa-ayliu-new-words-for-the-new-month/}{b} \\
\textsc{Derivatives}: \ding{43} 
\on{1}.~\HL{nIran}{\B{nì\-ran}}, basically, fundamentally, in essence. 
\on{2}.~\HL{loran}{\B{lo\-ran}}, elegance, grace. 
\on{3}.~\HL{fe'ran}{\B{fe'\-ran}}, flawed nature.

% Ranu
\hi 
\HT{ranu}{\B{Ranu}} \I{[Ra.nu]} \emph{n. given name}.

% rangal
\hi
\HT{rangal}{\B{\cp{rang}al}} \I{["RaN.al]} \emph{vin.} wish (\emph{used with} \HL{tsnI}{\B{tsnì}}). \\
\textsc{Derivatives}: \ding{43} 
\on{1}.~\HL{nIrangal}{\B{nì\-ra\-ngal}}, ``I wish that …''.
\on{2}.~\HL{sArangal}{\B{sä\-ra\-ngal}}, a wish.

% raspu'
\hi
\HT{raspu'}{\B{ra\cp{spu'}}} \I{[Ra."{}spuP]} \emph{or} \I{[."{}spUP]} \emph{n.} \ding{246} leggings (\emph{used in war}).
\B{Fì\-ra\-spu' 'e\-kxin lu nì\-hawng. Ke tsun oe yi\-vem\-stokx.} These leggings are too tight. I can't wear them.~\LINK{http://naviteri.org/2012/07/meetings-waterfalls-and-more/}{b}

% raw
\hi 
\HT{raw}{\B{raw}} \I{[Raw]} \emph{adp.} down to. 
\B{Ko\-lä oe raw kil\-van fte iva\-ho.} I went down to the river to pray.
\B{Kll\-za\-'u yì\-raw a\-mu\-ve.} Descend to the second level.
\B{Ti\-am ta vo\-mrr raw pxey.} Count down from thirteen to three. ~\LINK{http://naviteri.org/2019/06/50a-liu-amip-40-new-words/}{b} 
(\ding{214}~\HL{vay}{\B{vay}})

% rawke
\hi
\HT{rawke}{\B{\cp{raw}ke}} \I{["Raw.kE]} \emph{intj.} alarm cry, call to defense.

% rawm
\hi
\HT{rawm}{\B{rawm}} \I{[Rawm]} \emph{n.} (\emph{of weather}) lightning (\emph{general term}).
(\ding{43}~\HL{atanzaw}{\B{atan\-zaw}})  \\
\textsc{Derivatives}: \ding{43} 
\on{1}.~\HL{rawmpxom}{\B{rawm\-pxom}}, thunder and lightning.

% rawmpxom
\hi
\HT{rawmpxom}{\B{\cp{rawm}pxom}} \I{["Rawm.p'om]} \emph{n.} thunder and lightning.
(\ding{43}~\HL{rawm}{\B{rawm}}, \HL{'rrpxom}{\B{'rr\-pxom}}) 

% rawn
\hi
\HT{rawn}{\B{rawn}} \I{[Rawn]} \emph{vtr.} replace, substitute. 
\B{Ro\-lawn oel pa'\-lit fa ik\-ran.} I replaced my direhorse with a banshee.~\LINK{http://naviteri.org/2012/01/mipa-zisit-ayliu-amip-new-words-for-the-new-year/}{b} \\
\textsc{Derivatives}: \ding{43} 
\on{1}.~\HL{tIrawn}{\B{tì\-rawn}}, replacement, act of replacing. 
\on{2}.~\HL{sArawn}{\B{sä\-rawn}}, replacement, substitute.
\on{3}.~\HL{rawnlI'u}{\B{rawn\-lì\-'u}}, pronoun.

% rawnlì'u
\hi 
\HT{rawnlI'u}{\B{\cp{rawn}lì'u}} \I{["Rawn.lI.Pu]} \emph{n.} pronoun. 
(\ding{43}~\HL{rawn}{\B{rawn}}, \HL{lI'u}{\B{lì'u}})

% rawng
\hi 
\HT{rawng}{\B{rawng}} \I{[RawN]} \emph{n.} entrance, doorway. 
\B{Fpxo\-lä\-kìm fo ìlä rawng a\-hì\-'i.} They entered through a small doorway.~\LINK{http://naviteri.org/2016/06/mrrvola-lifyavi-amip-forty-new-expressions/}{b} 

% rawp
\hi
\HT{rawp}{\B{rawp}} \I{[Rawp\textcorner]} \emph{n.}, \ding{96} bladder polyp, \emph{polypiferus brevisimus}. 

% rä'ä
\hi
\HT{rA'A}{\B{rä\cp{'ä}}} \I{[R\ae."P\ae]} \emph{part. to form a negative imperative}, ``don't (do s.t.)!''.
(\emph{a. usually before the verb})
\B{Nim rä\-'ä lu!} Don't be shy!~\LINK{http://naviteri.org/2011/09/miscellaneous-vocabulary/}{b}
\B{Hawng\-krr rä\-'ä zi\-va\-'u!} Don't come late!~\LINK{http://naviteri.org/2010/07/vocabulary-update/}{b}
\B{Txo\-pu rä\-'ä si, lu ke\-tu\-wong\-o nì\-'aw.} Don't be afraid, it's just an alien.~\LINK{http://forum.learnnavi.org/language-updates/a-collection/}{ln\slash a} 
(\emph{b. after the verb for special emphasis}) \B{Oe\-ti 'am\-pi rä\-'ä, ma skxawng!} Don't touch me, you moron!~\LINK{http://naviteri.org/2012/11/renu-ayinanfyaya-the-senses-paradigm/}{b} 
\B{Hay\-a yak\-ro fti\-vang. Sa\-lew rä\-'ä.} Stop at the next fork. Do not proceed further.~\LINK{http://naviteri.org/2013/09/on-si-salewfya-shapes-and-directions/}{b}
\B{Na\-ri si! Kll\-te\-sìn lu pay a\-txan. Fwi rä\-'ä!} Be careful! There’s a lot of water on the ground. Don't slip!~\LINK{http://naviteri.org/2013/04/wheres-the-bathroom-and-other-useful-things/}{b}
(\emph{c. proverb}) \B{Kem\-ìri a nga\-ru prr\-te' ke lu, tsa\-kem rä\-'ä si\-vi ay\-la\-he\-ru.} As for any action that you don't like, don't act that way to others.~\LINK{http://wiki.learnnavi.org/index.php?title=Corpus\#Good_Morning_America}{wiki}

% räptulì'fya
\hi 
\HT{rAptulI'fya}{\B{räptu\cp{lì'}fya}} \I{[R\ae{}p\textcorner.tu."lIP.fja]} (\textsc{rn}: \B{räptù\cp{lì'}fya} \I{[REp\textcorner.tU."lIP.fja]}) \emph{n.} coarse, vulgar language. 
(\ding{43}~\HL{rAptum}{\B{räp\-tum}}, \HL{lI'fya}{\B{lì'\-fya}}) 

% räptulì'u
\hi 
\HT{rAptulI'u}{\B{räptu\cp{lì}'u}} \I{[R\ae{}p\textcorner.tu."lI.Pu]} (\textsc{rn}: \B{räptù\cp{lì}u} \I{[REp\textcorner.tU."lI.u]}) \emph{n.} coarse or swear word. 
(\ding{43}~\HL{rAptum}{\B{räp\-tum}}, \HL{lI'u}{\B{lì\-'u}}) 

% räptum
\hi
\HT{rAptum}{\B{räp\cp{tum}}} \I{[R\ae{}p\textcorner."tum]} \emph{or} \I{[."tUm]} (\textsc{rn}: \B{räp\cp{tùm}} \I{[REp\textcorner."tUm]}) \emph{adj.} coarse, vulgar, socially unacceptable.
\B{Po\-ri lu sno\-lup räp\-tum, ul\-te mawl ay\-lì\-'uä ke lu ngay.} His personal style is coarse and half of the words are not true.~\LINK{http://naviteri.org/2016/12/tengkrr-perahem-zisit-amip-as-the-new-year-arrives/}{b}
(\emph{a. idiom}) \B{Rä\-'ä räp\-tum!} Don't be impolite.~\LINK{http://naviteri.org/2016/12/tengkrr-perahem-zisit-amip-as-the-new-year-arrives/}{b}
\B{Ngal new pe\-ut? -- Rä\-'ä räp\-tum, ma 'eveng! Tsay\-lì\-'u ke lu muiä!} What do you want? -- Don't be impolite, child! Those words are improper!~\LINK{http://naviteri.org/2016/12/tengkrr-perahem-zisit-amip-as-the-new-year-arrives/}{b}
(\emph{b. idiom}) \B{pll\-txe räp\-tum}, speak vulgarly, use vulgar language, swear.\LINK{https://naviteri.org/2024/02/mipa-ayliu-si-aylifyavi-niul-more-new-words-and-expressions/}{b} \\
\textsc{Derivatives}: \ding{43} 
\on{1}.~\HL{nIrAptum}{\B{nì\-räp\-tum}}, coarsly, vulgarly.  
\on{2}.~\HL{rAptulI'u}{\B{räp\-tu\-lì\-'u}}, coarse or swear word.
\on{3}.~\HL{rAptulI'fya}{\B{räp\-tu\-lì'\-fya}}, coarse, vulgar language.

% räzekx
\hi 
\HT{rAzekx}{\B{\cp{rä}zekx}} \I{["R\ae.zEk']} \emph{n.} curse. \\
(i.) \B{tìng \cp{rä}zekx} \emph{vin.}, to curse (s.o.), wishing s.o. ill.
\B{Oe nga\-ru tìng rä\-zekx, ma kal\-wey\-a\-veng!} I curse you, (OR: Damn you!), you son of a bitch.~\LINK{https://naviteri.org/2025/06/vospikin-lefpom-happy-july/}{b} 
\B{Rä\-zekx ngar!} Damn you! [lit.: ``(I) curse you!'']~\LINK{https://naviteri.org/2025/06/vospikin-lefpom-happy-july/\#comment-54817}{b}
(\ding{214}~\HL{syawn}{\B{syawn}}) 

% re'o
\hi
\HT{re'o}{\B{\cp{re}'o}} \I{["RE.Po]} \emph{n.} head. 
\B{Oe\-ri ay\-som\-pì\-va sìn re\-'o var zi\-vup.} Raindrops keep falling on my head.~\LINK{http://naviteri.org/2011/04/yafkeykiri-plltxe-frapo-everyone-talks-about-the-weather/}{b}
\B{Krro lu 'Ìng\-lì\-sì txur nì\-hawng mì re\-'o.} Sometimes English is too strong in (my) head.~\LINK{http://naviteri.org/2011/09/\%E2\%80\%9Cby-the-way-what-are-you-reading\%E2\%80\%9D/comment-page-1/\#comment-1092}{b} 
(\emph{a. idiom}) \B{(na) la\-nay'\-ka lu\-ke re\-'o}, ``completely lost,'' [lit.: Like a slinger without a head].~\LINK{http://naviteri.org/2021/08/ulte-ayyoratu-leiu-and-the-winners-are-2/}{b}
\B{Po maw kxitx mun\-txa\-tu\-ä 'a\-me\-fu na la\-nay'\-ka re'\-o\-lu\-ke.} After the death of his spouse, he felt completely lost.~\LINK{http://naviteri.org/2021/08/ulte-ayyoratu-leiu-and-the-winners-are-2/}{b} \\
\textsc{Derivatives}: \ding{43} 
\on{1}.~\HL{hawre'}{\B{haw\-re'}}, hat. 
\on{2}.~\HL{zare'}{\B{za\-re'}}, forehead, brow. 

% rel
\hi
\HT{rel}{\B{rel}} \I{[REl]} \emph{n.} image, picture.
\B{Rel oe\-yä na unil\-tì\-ran\-yu lor lu nì\-ngay!} My picture as a Dreamwalker is truly beautiful.~\LINK{http://forum.learnnavi.org/language-updates/art-related-vocab/}{ln} 
\B{Lu tsay\-rel ko\-sman nì\-wotx!} Those pictures are wonderful!~\LINK{http://naviteri.org/2013/07/pximaw-tsawlultxa-just-after-the-meet-up/comment-page-1/\#comment-2485}{b}
\B{Fì\-me\-rel\-it 'o\-ngo\-lop aw\-nge\-yä tsul\-fä\-tul rel\-tseo\-ä alu Älìn.} The two portraits were created by our Master of Visual Arts, Alan.~\LINK{http://naviteri.org/2014/09/stxeli-alor-the-text/}{b}
\B{Ean\-a rel\-ìri lei\-am nga say\-rìp nì\-txan nang!} It's nice for me to see you look extremely handsome in your blue photo.~\LINK{http://forum.learnnavi.org/language-updates/eana-reliri-leiam-nga-sayrip-nitxan-nang!/}{ln}  \\ 
\textsc{Derivatives}: \ding{43} 
\on{1}.~\HL{reltseo}{\B{rel\-tseo}}, visual art.  
\on{2}.~\HL{rel arusikx}{\B{rel a\-ru\-sikx}}, movie, film, video. 
\on{3}.~\HL{pamrel}{\B{pam\-rel}}, writing, s.t. written. 
\on{4}.~\HL{keyrel}{\B{key\-rel}}, facial expression. 
\on{5}.~\HL{ronsrelngop}{\B{ron\-srel\-ngop}}, imagine. 
\on{6}.~\HL{ronsrel}{\B{ron\-srel}}, s.t. imagined. 
\on{7}.~\HL{syeprel}{\B{syep\-rel}}, camera. 
\on{8}.~\HL{kakrel}{\B{kak\-rel}}, blind. 
\on{9}.~\HL{atxkxerel}{\B{atx\-kxe\-rel}}, map.
\on{10}.~\HL{relvul}{\B{rel\-vul}}, pen, pencil.

% rel arusikx
\hi
\HT{rel arusikx}{\B{rel aru\cp{sikx}}} \I{[REl a.Ru."{}sik']} \emph{n.} movie, film, video [lit.: moving picture].
\B{Am\-'a\-lu\-ke hay\-a pxe\-rel a\-ru\-sikx wa\-you!} Doubtless the next three movies will be amazing!~\LINK{http://naviteri.org/2013/08/fmawn-a-fkol-kilmulat-this-just-in/}{b}
(\ding{43}~\HL{rel}{\B{rel}}, \HL{rikx}{\B{rikx}})

% reltseo
\hi
\HT{reltseo}{\B{\cp{rel}tseo}} \I{["REl.\t{ts}E.o]} \emph{n.} visual art.
\B{Law lu oe\-ru fwa nga mì rel\-tseo no\-lu\-me nì\-txan!} It's clear to me that you learned a lot in visual art.~\LINK{http://forum.learnnavi.org/language-updates/art-related-vocab/}{ln}
\B{Fì\-me\-rel\-it 'o\-ngo\-lop aw\-nge\-yä tsul\-fä\-tul rel\-tseo\-ä alu Älìn.} The two portraits were created by our Master of Visual Arts, Alan.~\LINK{http://naviteri.org/2014/09/stxeli-alor-the-text/}{b}
(\ding{43}~\HL{rel}{\B{rel}}, \HL{tseo}{\B{tseo}}) \\  
\textsc{Derivatives}: \ding{43} 
\on{1}.~\HL{reltseotu}{\B{rel\-tseo\-tu}}, artist.

% reltseotu
\hi
\HT{reltseotu}{\B{\cp{rel}tseotu}} \I{["REl.\t{ts}E.o.tu]} \emph{n.} artist.
\B{Rel\-tseo\-tu a\-txan\-tsan lu nga!} You are an excellent artist.~\LINK{http://forum.learnnavi.org/language-updates/art-related-vocab/}{ln}
(\ding{43}~\HL{reltseo}{\B{rel\-tseo}}, \HL{tute1}{\B{tu\-te}}) 

% relvul
\hi
\HT{relvul}{\B{\cp{rel}vul}} \I{["REl.vul]} \emph{or} \I{[.vUl]} \emph{n.} pen, pencil (\emph{short form of} \ding{43}~\HL{pamrelvul}{\B{pam\-rel\-vul}}).

% rem
\hi 
\HT{rem}{\B{rem}} \I{[REm]} \emph{n.} fuel. 
\B{Na'\-vi\-ri lu fì\-ut\-ral\-ä rìn rem le\-tsran\-ten.} The wood of this tree is an important fuel for the Na'vi.~\LINK{http://naviteri.org/2018/03/100a-liu-amip-64-new-words-part-1/}{b} 

% renten
\hi
\HT{renten}{\B{\cp{ren}ten}} \I{["REn.tEn]} \emph{n.} goggles (\emph{made from insect wings, carved wood etc.}).
\B{Fo\-ri maw\-krr\-a fa ren\-ten ioi sä\-po\-li ho\-lum.} After they put on their goggles, they left.~\LINK{http://naviteri.org/2011/08/new-vocabulary-clothing/}{b}

% renu
\hi
\HT{renu}{\B{\cp{re}nu}} \I{["RE.nu]} \emph{n.} pattern; paradigm. 
\B{Renu Ayinanfyayä}, the Senses Paradigm.~\LINK{http://naviteri.org/2012/11/renu-ayinanfyaya-the-senses-paradigm/}{b}
\B{Sìl\-pey oe, fì\-re\-nu a\-mip nga\-ru le\-sar lì\-ye\-vu.} I hope, this new (grammatical) pattern will be useful for you.~\LINK{http://naviteri.org/2013/03/whoever-whatever-whenever/comment-page-1/\#comment-2237}{b}
\B{Pam\-tse\-ol ngop ay\-re\-nut}, Music creates patterns.~\LINK{http://www.pandorapedia.com/navi/music/ritual_music}{pp}
\B{Tsay\-re\-nur le\-we\-opx a sìn ne\-ni tìng na\-ri.} Look at those wave\hyp{}like patterns in the sand.~\LINK{http://naviteri.org/2017/09/zisikrr-amip-ayliu-amip-new-words-for-the-new-season/}{b} 
\B{Ya\-mì tsun fko tsi\-ve\-'a lo\-ran\-it re\-nu\-ä kil\-van\-ä slä kll\-te\-sìn wä\-pan.} From the air you can see the grace of the river's form but from the ground it's hidden.~\LINK{http://naviteri.org/2012/10/mipa-vospxi-mipa-ayliu-new-words-for-the-new-month/}{b} \\  
\textsc{Derivatives}: \ding{43} 
\on{1}.~\HL{renulke}{\B{re\-nul\-ke}}, irregular, random. 
\on{2}.~\HL{rengop}{\B{re\-ngop}}, design (finer detail). 
\on{3}.~\HL{renu ngampamA}{\B{re\-nu ngam\-pa\-mä}}, rhyme scheme.

% renu ngampamä
\hi
\HT{renu ngampamA}{\B{\cp{re}nu \cp{ngam}pamä}} \I{["RE.nu "Nam.pa.m\ae]} \emph{n.} rhyme scheme.   
\B{Fì\-way\-ri hì\-no\-a re\-nut ngam\-pam\-ä ke tsä\-ngun oe tsli\-vam.} I'm afraid I can't understand the intricate rhyme scheme of this poem.~\LINK{http://naviteri.org/2012/01/mipa-zisit-ayliu-amip-new-words-for-the-new-year/}{b}
(\ding{43}~\HL{renu}{\B{re\-nu}}, \HL{ngampam}{\B{ngam\-pam}})

% renulke
\hi
\HT{renulke}{\B{\cp{re}nulke}} \I{["RE.nul.kE]} \emph{or} \I{[.nUl.]} \emph{adj.} irregular, random. 
\B{Eo ay\-fo a fya\-'o la\-mu ay\-skxe\-ta te\-ya sì re\-nul\-ke.} The path ahead of them was full of rocks and irregular.~\LINK{http://naviteri.org/2013/09/on-si-salewfya-shapes-and-directions/}{b}
(\ding{43}~\HL{renu}{\B{re\-nu}}, \HL{luke}{\B{lu\-ke}})

% reng
\hi
\HT{reng}{\B{reng}} \I{[REN]} \emph{adj.} (\emph{a. physical}) shallow. 
\B{kil\-van a\-reng}, a shallow river.~\LINK{http://naviteri.org/2010/07/vocabulary-update/}{b}
(\emph{b. sometimes applied metaphorically to thoughts, ideas, analyses, etc.}) shallow.
\B{Tì\-eyk\-tan\-ìri ay\-sä\-fpìl pe\-yä lu reng, ay\-sì\-hawl lu fe'\-ran.} Concerning leadership his thoughts are shallow, his plans are flawed.~\LINK{http://naviteri.org/2016/12/tengkrr-perahem-zisit-amip-as-the-new-year-arrives/}{b}
(\ding{214}~\HL{txukx}{\B{txukx}})

% rengop
\hi
\HT{rengop}{\B{\cp{re}ngop}} \I{["RE.N$\cdot$op\textcorner]} \emph{vtr., inf.}\on{22} design (\emph{for finer details, e.g. ornament design}).
\B{Nge\-yä tsa\-fkxi\-let tu\-pel re\-ngo\-lop?} Who designed that necklace of yours?~\LINK{http://naviteri.org/2013/01/awvea-posti-zisita-amip-first-post-of-the-new-year/}{b}
(\ding{43}~\HL{renu}{\B{re\-nu}}, \HL{ngop}{\B{ngop}}; \HL{'ongop}{\B{'o\-ngop}}) \\
\textsc{Derivatives}: \ding{43} 
\on{1}.~\HL{sArengop}{\B{sä\-re\-ngop}}, design (particular instance). 
\on{2}.~\HL{tIrengop}{\B{tì\-re\-ngop}}, design (abstract concept).

% rewon
\hi
\HT{rewon}{\B{\cp{re}won}} \I{["RE.won]} \emph{n.; adv.} morning.
\B{Fì\-re\-won tom\-pa\-meyp zar\-mup, slä set 'ì\-me\-ko nì\-txan nang.} It was drizzling this morning, but it's really started coming down now!~\LINK{http://naviteri.org/2011/04/yafkeykiri-plltxe-frapo-everyone-talks-about-the-weather/}{b}
\B{Alo a\-hay oe za\-sya\-'u mì re\-won.} Next time I'll come in the morning.~\LINK{http://forum.learnnavi.org/language-updates/txelanit-hivawl/}{ln} \\ 
(i.) \B{rewo\cp{nay}} \I{[RE.wo."naj]} \emph{n.; adv.} tomorrow morning.   
\B{Oeng re\-wo\-nay 'aw\-si\-teng ti\-va\-ron ko.} Let's you and I go hunting together tomorrow morning.~\LINK{http://naviteri.org/2011/02/new-vocabulary-part-2/}{b} \\  
(ii.) \B{rewo\cp{nam}} \I{[RE.wo."nam]} \emph{n.; adv.} yesterday morning. 

% rey
\hi
\HT{rey}{\B{rey}} \I{[REj]} \emph{vin.} live.  
\B{ke rey a tu\-te}, person who doesn't live\slash isn't alive.~\LINK{http://forum.learnnavi.org/language-updates/participles/}{ln}
\B{Txo new nga ri\-vey, oe\-hu!} Come with me, if you want to live!~\LINK{http://wiki.learnnavi.org/index.php?title=Corpus\#UGO_Movie_Blog_.28Twitter_Questions.29}{wiki} 
\B{Pi\-veng oer ftxey nga new ri\-vey fu\-ke.} Tell me whether you want to live (or not).~\LINK{http://forum.learnnavi.org/language-updates/if-vs-whether-ftxey-fuke/}{ln} 
\B{Fì\-lì\-'u\-ä ral lu tì\-me\-'em sì tì\-ru\-sey mì hi\-fkey na Nawm\-a Sa'\-nok\-ä ha\-pxì.} The meaning of this word is ``harmony, living in the world as part of the Great Mother.''~\LINK{http://forum.learnnavi.org/language-updates/trr-rrtaya/}{ln}
\B{Ko\-ren A\-'aw\-ve Tì\-ru\-sey\-ä 'Aw\-si\-teng}, the First Rule of Living Together.~\LINK{http://wiki.learnnavi.org/index.php?title=Corpus\#Good_Morning_America}{wiki} \\
\textsc{Derivatives}: \ding{43} 
\on{1}.~\HL{tIrey}{\B{tì\-rey}}, life. 
\on{2}.~\HL{rusey}{\B{ru\-sey}}, living thing; alive, living. 
\on{3}.~\HL{rey'eng}{\B{rey\-'eng}}, balance of life. 
\on{4}.~\HL{reyfya}{\B{rey\-fya}}, way of living, culture. 
\on{5}.~\HL{reypay}{\B{rey\-pay}}, blood. 
\on{6}.~\HL{emrey}{\B{em\-rey}}, survive.

% rey'eng
\hi
\HT{rey'eng}{\B{rey\cp{'eng}}} \I{[REj."PEN]} \emph{n.} balance of life.
\B{'Uo a fpi rey\-'eng Ey\-wa\-'e\-veng\-mì, 'Rr\-ta\-mì tsran\-ten nì\-txan aw\-nga\-ru nì\-wotx.} Something that matters a lot to all of us for the sake of The Balance of Life on both Pandora and Earth.~\LINK{http://forum.learnnavi.org/news-announcements/a-response-from-paul-frommer!/}{ln} 
(\ding{43}~\HL{rey}{\B{rey}}, \HL{'engeng}{\B{'eng\-eng}})

% Reya
\hi
\HT{Reya}{\B{Reya}} \I{[RE.ja]} \emph{n., female name} (\emph{short form of:} \ding{43}~\HL{Tsireya}{\B{Tsi\-re\-ya}}).

% reyfya
\hi
\HT{reyfya}{\B{\cp{rey}fya}} \I{["REj.fja]} \emph{n.} way of living, culture.
\B{'Rr\-ta\-mì a rey\-fya Saw\-tu\-te\-yä la\-tsu fe'\-ran nì\-ngay.} The Skypeople's culture on earth must truly be flawed.~\LINK{http://naviteri.org/2012/10/mipa-vospxi-mipa-ayliu-new-words-for-the-new-month/}{b}
(\ding{43}~\HL{rey}{\B{rey}}, \HL{fya'o}{\B{fya\-'o}}) \\
\textsc{Derivatives}: \ding{43} 
\on{1}.~\HL{lereyfya}{\B{le\-rey\-fya}}, cultural.

% reykol
\hi
\HT{reykol}{\B{rey\cp{kol}}} \I{[REj."k$\cdot$ol]} \emph{vtr., inf.}\on{22} play music (\emph{poetic}).  
\B{Tew\-ti, nga lu tsul\-fä\-tu i\-'en\-ä. Ngal tsat rey\-kì\-mol!} Wow, you are a master on the \emph{i'en}. You just made it sing!~\LINK{http://naviteri.org/2011/01/new-year-new-vocabulary/}{b}
(\ding{43}~\HL{rol}{\B{rol}}, ‹\HL{eyk2}{\B{eyk}}›)

% reym
\hi
\HT{reym}{\B{reym}} \I{[REjm]} \emph{n.} dry land (\emph{as distinct from water}). 
\B{Pe\-ioang tsun mì tam\-pay kop mì reym ri\-vey?} What animal can live both in the sea and on the land?~\LINK{http://naviteri.org/2015/04/some-new-words-for-may-day/}{b} 

% reypay
\hi
\HT{reypay}{\B{\cp{rey}pay}} \I{["REj.paj]} \emph{n.} blood. 
\B{Rey\-pay skxir\-ftum\-fa he\-rum.} Blood is coming out of the wound.~\LINK{http://naviteri.org/2013/04/wheres-the-bathroom-and-other-useful-things/}{b} 
(\ding{43}~\HL{rey}{\B{rey}}, \HL{pay}{\B{pay}}) \\
\textsc{Derivatives}: \ding{43} 
\on{1}.~\HL{reypaytun}{\B{rey\-pay\-tun}}, red. 
\on{2}.~\HL{reyptswIk}{\B{reyp\-tswìk}}, wolf tick.

% reypaytun
\hi 
\HT{reypaytun}{\B{\cp{rey}paytun}} \I{["REj.paj.tun]} \emph{or} \I{[.tUn]} \emph{adj.} red.
\B{'Eveng\-ìl fì\-rel\-it 'opin\-so\-lung fa pin\-vul a\-rey\-pay\-tun.} The child colored (in) this picture with a red crayon.~ \LINK{http://naviteri.org/2018/03/100a-liu-amip-64-new-words-part-1/\#comment-28355}{b}
(\ding{43}~\HL{reypay}{\B{rey\-pay}}, \HL{tun}{\B{tun}})

% reyptswìk
\hi 
\HT{reyptswIk}{\B{\cp{reyp}tswìk}} \I{["REjp\textcorner.\t{ts}wIk\textcorner]} \emph{n., fauna} wolf tick. 
(\ding{43}~\HL{reypay}{\B{rey\-pay}}, \HL{tswIk}{\B{tswìk}}) 

% -ri
\hi
\HT{-ri}{\B{-ri}} \I{[.Ri]} \emph{suf. to mark the topic case of a n. ending in a vowel or diphthong}.
(\ding{43}~\HL{-Iri}{\B{-ìri}})

% rikx
\hi
\HT{rikx}{\B{rikx}} \I{[Rik']} \emph{vin.} move, shift position. 
\B{Ftu oe ne\-to rikx!} Get away from me!~\LINK{http://naviteri.org/2011/08/new-vocabulary-clothing/comment-page-1/\#comment-994}{b}
\B{Tsun kxam\-lä na'\-rìng ri\-vikx nì\-fnu nì\-wotx.} (The thanator) can move silently through the forest.~\LINK{http://naviteri.org/2012/10/mipa-vospxi-mipa-ayliu-new-words-for-the-new-month/}{b} \\
\textsc{Derivatives}: \ding{43} 
\on{1}.~\HL{'ampirikx}{\B{'am\-pi\-rikx}}, leaf pitcher plant. 
\on{2}.~\HL{kllrikx}{\B{kll\-rikx}}, earthquake. 
\on{3}.~\HL{rel arusikx}{\B{rel a\-ru\-sikx}}, movie, film, video.

% rim
\hi
\HT{rim}{\B{rim}} \I{[Rim]} \emph{adj.} yellow.   
\B{Fì\-syu\-lang lu rim.} This flower is yellow.\LINK{http://naviteri.org/2010/08/mipa-ayopin-mipa-ayliu-new-colors-new-words/}{b}
\B{Fì\-syu\-lang a\-rim lu hì\-'i fra\-to.} This yellow flower is the smallest of all.\LINK{http://naviteri.org/2010/08/mipa-ayopin-mipa-ayliu-new-colors-new-words/}{b}
\B{rim na na\-ri}, the yellow of a Na'vi eye.~\LINK{http://naviteri.org/2010/08/mipa-ayopin-mipa-ayliu-new-colors-new-words/}{b} \\
\textsc{Derivatives}: \ding{43} 
\on{1}.~\HL{nIrim}{\B{nì\-rim}}, in yellow. 
\on{2}.~\HL{rimpin}{\B{rim\-pin}}, the color yellow. 

% rimpin
\hi
\HT{rimpin}{\B{\cp{rim}pin}} \I{["Rim.pin]} \emph{n.} the color \ding{43}~\HL{rim}{\B{rim}}.
\B{Ke su\-nu oe\-ru rim\-pin.} I don't like the color yellow.~\LINK{http://naviteri.org/2010/08/mipa-ayopin-mipa-ayliu-new-colors-new-words/}{b}
(\ding{43}~\HL{'opin}{\B{'opin}})

% rina'
\hi
\HT{rina'}{\B{ri\cp{na'}}} \I{[Ri."naP]} \emph{n.}, \ding{96} seed. 
\B{Ut\-ral\-ä A\-nawm | Ay\-ri\-na' lu ay\-oeng.} We are all seeds | Of the Great Tree.~\LINK{http://www.pandorapedia.com/navi/music/ritual_music}{pp}
\B{Ay\-fol za\-mo\-lu\-nge aw\-ngar ay\-ri\-na'\-it fa sä\-he\-na a\-pxa.} They brought us the seeds in a large container.~\LINK{http://naviteri.org/2014/05/mipa-ayliu-mipa-aysafpil-new-words-new-ideas/}{b} \\  
\textsc{Derivatives}: \ding{43} 
\on{1}.~\HL{atokirina'}{\B{ato\-ki\-ri\-na'}}, seeds of the great tree, woodsprite.  
\on{2}.~\HL{tsyorina'wll}{\B{tsyo\-ri\-na'\-wll}}, cycad.

% Rini
\hi
\B{\cp{Ri}ni} \I{["Ri.ni]} \emph{n.} \emph{female name}. 
\B{Ri\-ni tsa\-po\-hu ho\-lum nì\-mal nì\-wotx.} Rini left with that guy without thinking twice about it.~\LINK{http://naviteri.org/2012/01/more-additions-to-the-lexicon/}{b}
\B{Ke new oe mi\-vay' tsngan\-ti a 'o\-lem Ri\-nil.} I don't want to taste the meat that Rini cooked.~\LINK{http://naviteri.org/2012/11/renu-ayinanfyaya-the-senses-paradigm/}{b}
\B{Nik\-re\-ri Ri\-ni\-yä 'ur fkan lor.} Rini's hair looks beautiful.~\LINK{http://naviteri.org/2012/11/renu-ayinanfyaya-the-senses-paradigm/}{b}
\B{Ri\-ni\-ri nik\-re yä\-ngo' nì\-tut.} Rini's hair is always perfect.~\LINK{http://naviteri.org/2012/01/mipa-zisit-ayliu-amip-new-words-for-the-new-year/}{b}

% ripx
\hi 
\HT{ripx}{\B{ripx}} \I{[Rip']} \emph{vtr.} pierce. 
\B{Lu Ney\-ti\-ri\-ru 'aw\-a mik\-yun a\-raw\-nipx nì\-'aw.} Neytiri has only one pierced ear.~\LINK{http://naviteri.org/2017/12/zisit-amip-lefpom-ma-eylan-happy-new-year-friends/}{b}
(\emph{a. idiom}) \B{Pay\-ìl a lipx tskxe\-ti ripx.} Dripping water pierces a stone. [{\em persistent effort can accomplish unexpected and amazing things}]~\LINK{http://naviteri.org/2020/04/aawa-u-amip-a-few-new-things/}{b}

% riti
\hi
\HT{riti}{\B{\cp{ri}ti}} \I{["Ri.ti]} \emph{n., fauna} stingbat, \emph{scorpiobattus volansii}. 
\B{Ri\-ti tswo\-lay\-on ftum\-fa slär.} The stingbat flew out of the cave.~\LINK{http://naviteri.org/2013/04/wheres-the-bathroom-and-other-useful-things/}{b} 
\B{mrr\-a ri\-ti}, five stingbats.~\LINK{http://naviteri.org/2011/07/number-in-na\%E2\%80\%99vi/}{b}

% rì'ìr
\hi
\HT{rI'Ir}{\B{\cp{rì}'ìr}} \I{["RI.PIR]} \textsc{i}. \emph{n.} reflection.
\B{Oe\-ri pay\-ìl tìng rì\-'ìr\-it key\-ä.} My face is reflected in the water.~\LINK{http://naviteri.org/2011/09/miscellaneous-vocabulary/}{b} \\
\textsc{ii}. \HL{rI'Irsi}{\B{\cp{rì}'ìr si}} \I{["RI.PIR si]} \emph{vin.} reflect, imitate.   
\B{Rì\-'ìr rä\-'ä si\-vi tsmuk\-tur!} Don't imitate your sibling!~\LINK{http://naviteri.org/2011/09/miscellaneous-vocabulary/}{b}
\B{Por lu sno\-lup a new fra\-po rì\-'ìr si\-vi.} He\slash She has a personal style that everyone wants to emulate.~\LINK{http://naviteri.org/2014/07/tsawlultxamaw-a-ayliu-post-meetup-words/}{b}

% rìk
\hi
\HT{rIk}{\B{rìk}} \I{[RIk\textcorner]} \emph{n.} (\emph{a.} \ding{96}) leaf.
\B{Na'\-rìng a zì\-sìt\-o a\-ki\-nä ke lal\-mu tsar kea rìk.} A forest that had no leaf for seven years.~\LINK{http://forum.learnnavi.org/language-updates/for-\%28duration\%29-jesus-loan-word/}{ln}
\B{Pol sä\-'o\-ti wo\-lan äo ay\-rìk.} He hid his tool under the leaves.~\LINK{http://naviteri.org/2011/02/new-vocabulary-ii\%E2\%80\%94part-1/}{b}
\B{ean na rìk}, leaf\hyp{}green.~\LINK{http://naviteri.org/2010/08/mipa-ayopin-mipa-ayliu-new-colors-new-words/}{b}
(\emph{b. coll. form of} \ding{43}~\HL{yomyolerIk}{\B{yom\-yo le\-rìk}}) leaf plate. \\  
\textsc{Derivatives}: \ding{43} 
\on{1}.~\HL{lerIk}{\B{le\-rìk}}, leafy. 
\on{2}.~\HL{rIkean}{\B{rìk\-ean}}, leaf\hyp{}green.
\on{3}.~\HL{rIktsyokx}{\B{rìk\-tsyokx}}, flat hand.

% rìkean
\hi
\HT{rIkean}{\B{\cp{rì}kean}} \I{["RI.kE.an]} \emph{adj.} leaf\hyp{}green, i.e. green.
(\ding{43}~\HL{rIk}{\B{rìk}}, \HL{ean}{\B{ean}}; \HL{ta'lengean}{\B{ta'\-leng\-ean}})

% rìktsyokx
\hi 
\HT{rIktsyokx}{\B{\cp{rìk}tsyokx}} \I{["RIk\textcorner.\t{ts}jok']} \emph{n.} flat hand. 
\B{A\-te\-yol En\-tut to\-la\-kuk fa rìk\-tsyokx.} Ateyo slapped Entu (i.e. struck him with his flat hand).~\LINK{https://naviteri.org/2024/06/ayhapxi-tokxa-si-ayliu-alahe-parts-of-the-body-and-more/}{b} 
(\ding{43}~\HL{rIk}{\B{rìk}}, \HL{tsyokx}{\B{tsyokx}}) 

% rìkxi
\hi 
\HT{rIkxi}{\B{rì\cp{kxi}}} \I{[RI."k'i]} \emph{vin.} tremble, shake, shiver. 
\B{Po\-ri me\-syokx rì\-kxi, ha ke tsa\-yun ye\-rik\-it ti\-va\-kuk.} His hands tremble, so he will not be able to hit the hexapede.~\LINK{http://naviteri.org/2018/03/100a-liu-amip-64-new-words-part-1/}{b}
\B{Ra\-lu rì\-kxi krr\-a srew, rì\-'ìr si pa\-lu\-kan\-ur a lu a\-lak\-si fte spi\-vä.} Ra\-lu does a shake while dancing, imitating a thanator that's ready to leap.~\LINK{http://naviteri.org/2018/03/100a-liu-amip-64-new-words-part-1/}{b} 
(\emph{a. idiom}) 
\B{Rey\-kì\-kxi ut\-ral\-ti, zup ma\-u\-ti.} If you shake the tree, the fruit will fall (i.e, ``actions have consequences'').~\LINK{http://naviteri.org/2018/03/100a-liu-amip-64-new-words-part-1/}{b} \\
\textsc{Derivatives}: \ding{43} 
\on{1}.~\HL{nIrIkxi}{\B{nì\-rìkxi}}, shakily, tremblingly.

% rìn
\hi
\HT{rIn}{\B{rìn}} \I{[RIn]} \emph{n.} \ding{96} wood.
\B{Na'\-vi\-ri lu fì\-ut\-ral\-ä rìn rem le\-tsran\-ten.} The wood of this tree is an important fuel for the Na'vi.~\LINK{http://naviteri.org/2018/03/100a-liu-amip-64-new-words-part-1/}{b} 
\B{Fì\-fne\-rìn ke ley krr\-a slu pay\-nga'.} This kind of wood is worthless when it gets damp.~\LINK{http://naviteri.org/2014/03/value-and-worth/}{b}
\B{Txep\-ìri, yrr\-a rìn\-ti rä\-'ä sar, ma skxawng.} Don't use that green wood for a fire, you fool.~\LINK{http://naviteri.org/2013/01/awvea-posti-zisita-amip-first-post-of-the-new-year/}{b} \\
\textsc{Derivatives}: \ding{43} 
\on{1}.~\HL{lerIn}{\B{le\-rìn}}, wooden, of wood.

% ro
\hi
\HT{ro}{\B{ro}} \I{[Ro]} \emph{adp.}\texttt{+} (\emph{a. locative}) at.   
\B{ro fä\-pa}, at the top.~\LINK{http://forum.learnnavi.org/language-updates/all-adpositions/msg137792/\#msg137792}{ln}
\B{ro hel\-ku}, at home.~\LINK{http://forum.learnnavi.org/language-updates/all-adpositions/msg137792/\#msg137792}{ln}
\B{Sìl\-pey oe, tsun oeng ul\-txa si\-vi ro hay\-a ul\-txa!} I hope we can meet at the next meeting.~\LINK{http://naviteri.org/2011/07/txantsana-ultxa-mi-siatll-great-meeting-in-seattle/comment-page-1/\#comment-844}{b}
\B{Ro sngä\-'i\-tseng tsa\-lì\-'u\-ä alu 'ey\-lan lu tì\-ftang.} At the beginning of the word \emph{'eylan} there's a glottal stop.~\LINK{http://naviteri.org/2012/07/fivospxiya-aylifyavi-amip-this-months-new-expressions-pt-2/}{b}
\B{Hay\-a yak\-ro fti\-vang. Sa\-lew rä\-'ä.} Stop at the next fork. Do not proceed further.~\LINK{http://naviteri.org/2013/09/on-si-salewfya-shapes-and-directions/}{b}
\B{New oe ri\-vun a\-sim tì\-fnu\-nga'\-a tseng\-it a tsa\-ro tsun syi\-vor tsi\-vu\-rokx fte spä\-pi\-veng.} I want to find a quiet place nearby where I can chill out and rest to get my head back on straight.~\LINK{http://naviteri.org/2013/01/awvea-posti-zisita-amip-first-post-of-the-new-year/}{b}
(\emph{b. temporal to disambiguate point\hyp{}in\hyp{}time interpretation together with indefinite} \ding{43}~\B{-o}) 
\B{Ro srr\-o Ra\-lu zo\-la\-'u fte oe\-hu ul\-txa si\-vi.} One day Ralu came to meet with me.~\LINK{http://naviteri.org/2014/09/stxeli-alor-the-text/comment-page-1/\#comment-10619}{b} \\
\textsc{Derivatives}: \ding{43} 
\on{1}.~\HL{rofa}{\B{ro\-fa}}, beside, alongside.

% ro'a
\hi
\HT{ro'a}{\B{\cp{ro}'a}} \I{["Ro.Pa]} \emph{vin.} be impressive, inspire awe or respect.  
\B{To\-ruk Mak\-to po\-lä\-hem a fì\-'u ro\-lo\-'a nì\-txan Oma\-ti\-ka\-ya\-ru.} The arrival of Toruk Makto made a great impression on the Omaticaya.~\LINK{http://naviteri.org/2012/03/spring-vocabulary-part-1/}{b}
\B{Po\-ri wem\-tswo fra\-tsam\-si\-yur ro\-lo\-'a nì\-txan.} His ability to fight greatly impressed all the warriors.~\LINK{http://naviteri.org/2012/03/spring-vocabulary-part-2/}{b} \\  
\textsc{Derivatives}: \ding{43} 
\on{1}.~\HL{sAro'a}{\B{sä\-ro\-'a}}, feat, accomplishment, great deed.  
\on{2}.~\HL{txanro'a}{\B{txan\-ro\-'a}}, be famous.

% rofa
\hi
\HT{rofa}{\B{\cp{ro}fa}} \I{["Ro.fa]} \emph{adp.} beside, alongside. 
\B{Kll\-kxa\-yem fì\-tì\-kang\-kem oe\-yä ro\-fa--ke io--pum fe\-yä.} My work will stand beside--not above--theirs.~\LINK{http://naviteri.org/2010/06/first-post/}{b}
\B{Maw sä\-tsway\-on a\-yol ay\-oe kll\-po\-lä mì ta\-yo a lu ro\-fa kil\-van.} After a short flight we landed in a field beside the river.~\LINK{http://naviteri.org/2012/01/mipa-zisit-ayliu-amip-new-words-for-the-new-year/}{b}
(\ding{43}~\HL{ro}{\B{ro}}, \HL{pa'o}{\B{pa'o}})

% rol
\hi
\HT{rol}{\B{rol}} \I{[Rol]} \emph{vin.} sing.   
\B{Kä\-sro\-lìn oel nik\-ro\-it Pey\-ral\-ur yoa fwa po rol oer.} I loaned Peyral a hair ornament in exchange for her singing to me.~\LINK{http://naviteri.org/2014/02/barter-and-exchange/}{b}
\B{Pll\-txe fra\-po san Pey\-ral ke tsun ri\-vol nìl\-tsan sìk.} They say that Peyral can't sing well.~\LINK{http://naviteri.org/2014/06/teri-tsaliu-alu-nifkeytongay-about-nifkeytongay/}{b}
\B{Re\-rol ay\-oe ay\-nga\-ne}, We are singing your way.~\LINK{http://wiki.learnnavi.org/index.php/Hunting_Song}{wiki}
\B{Na'\-vi ìlä ho\-'on kll\-kxo\-lem teng\-krr re\-rol.} The People were standing in a circle, singing.~\LINK{http://naviteri.org/2013/09/on-si-salewfya-shapes-and-directions/}{b}
\B{Tì\-ru\-sol za\-'u ne fo nì\-'eng.} They share an interest in singing.~\LINK{http://naviteri.org/2010/09/getting-to-know-you-part-1/}{b} 
\B{Tì\-ru\-sol lu oe\-ru mo\-wan.} Singing is enjoyable to me.~\LINK{http://naviteri.org/2012/02/trr-asawnung-lefpom-happy-leap-day/}{b} \\
\textsc{Derivatives}: \ding{43} 
\on{1}.~\HL{rolyu}{\B{rol\-yu}}, singer.
\on{2}.~\HL{tIrol}{\B{tì\-rol}}, song. 
\on{3}.~\HL{ronguway}{\B{ro\-ngu\-way}}, howl.

% rolyu
\hi
\HT{rolyu}{\B{\cp{rol}yu}} \I{["Rol.ju]} \emph{n.} singer.
\B{Nì\-fkey\-to\-ngay lu po rol\-yu a\-txan\-tsan.} As a matter of fact, she is an excellent singer.~\LINK{http://naviteri.org/2014/06/teri-tsaliu-alu-nifkeytongay-about-nifkeytongay/}{b}
(\ding{43}~\HL{rol}{\B{rol}})

% Ronal
\hi
\HT{Ronal}{\B{\cp{Ro}nal}} \I{["Ro.nal]} \emph{n., female name}.

% ronsem
\hi
\HT{ronsem}{\B{\cp{ron}sem}} \I{["Ron.sEm]} \emph{n.} mind. 
\B{Pam\-tse\-ol ngop ay\-re\-nut | Mì ron\-sem\-ä tì\-fnu}, Music creates patterns | In the silence of the mind.~\LINK{http://www.pandorapedia.com/navi/music/ritual_music}{pp}
\B{Kxawm tsa\-lì\-'u le\-'Ìng\-lì\-sì alu ``No\"el'' mì ron\-sem lar\-mu.} Maybe I had the English word ``No\"el'' in my mind.~\LINK{http://naviteri.org/2011/01/new-year-new-vocabulary/comment-page-1/\#comment-499}{b}
\B{Saw\-tu\-te\-ri ron\-sem\-fkeyk\-it ke tsun kaw\-tu tsl\-ivam.} No one can understand the state of mind of the Sky People.~\LINK{http://naviteri.org/2011/04/yafkeykiri-plltxe-frapo-everyone-talks-about-the-weather/}{b} \\  
\textsc{Derivatives}: \ding{43} 
\on{1}.~\HL{ronsrel}{\B{ron\-srel}}, s.t. imagined. 
\on{2}.~\HL{fpomron}{\B{fpom\-ron}}, (mental) health. 
\on{3}.~\HL{sIlronsem}{\B{sìl\-ron\-sem}}, clever, smart (of things).

% ronsrel
\hi
\HT{ronsrel}{\B{\cp{ron}srel}} \I{["Ron.sREl]} \emph{n.} s.t. imagined, imagining.
\B{Ay\-ron\-srel pe\-yä hä\-ngek nì\-txan.} His imaginings are (unpleasantly) weird.~\LINK{http://naviteri.org/2011/02/new-vocabulary-part-2/}{b} 
(\ding{43}~\HL{ronsem}{\B{ron\-sem}}, \HL{rel}{\B{rel}}) \\
\textsc{Derivatives}: \ding{43} 
\on{1}.~\HL{ronsrelngop}{\B{ron\-srel\-gop}} 
\on{2}.~\HL{tIronsrel}{\B{tì\-ron\-srel}}, imagination. 
\on{3}.~\HL{leronsrel}{\B{le\-ron\-srel}}, imaginary. 
\on{4}.~\HL{nIronsrel}{\B{nì\-ron\-srel}}, in\slash by imagination.

% ronsrelngop
\hi
\HT{ronsrelngop}{\B{\cp{ron}srelngop}} \I{["Ron.sREl.N$\cdot$op\textcorner]} \emph{or coll.} \I{["Ron.sREw.Nop\textcorner]} \emph{vtr., inf.}\on{33} imagine, envision.
\B{Oel ron\-srel\-ngop fu\-ta Ey\-we\-veng\-it tok.} I imagine that I'm on Pandora.~\LINK{http://naviteri.org/2011/02/new-vocabulary-part-2/}{b}
\B{Tsat ke tsun oe ron\-srel\-ngi\-vop.} I can't imagine that.~\LINK{http://naviteri.org/2011/02/new-vocabulary-part-2/}{b}
(\ding{43}~\HL{ronsrel}{\B{ron\-srel}}, \HL{ngop}{\B{ngop}})

% rong
\hi
\HT{rong}{\B{rong}} \I{[RoN]} \emph{n.} tunnel.
\B{Po\-ti fwey\-ko\-li ay\-oel ìlä rong.} We let him slide through the tunnel.~\LINK{http://naviteri.org/2013/04/wheres-the-bathroom-and-other-useful-things/}{b}

% Rongloa
\hi
\B{\cp{Rong}loa} \I{["RoN.lo.a]} \emph{n.} \emph{clan name}.~\LINK{http://forum.learnnavi.org/language-updates/stxo/}{ln} 

% ronguway
\hi
\HT{ronguway}{\B{ro\cp{ngu}way}} \I{[R$\cdot$o."Nu.waj]} \emph{vin., inf.}\on{11} howl.
\B{Nan\-tang ro\-ngu\-way.} The viper wolf howls.~\LINK{http://forum.learnnavi.org/projects/kaltxi-ma-prrnen!-animal-learning-book/msg612075/\#msg612075}{ln}
(\ding{43}~\HL{rol}{\B{rol}}, \HL{nguway}{\B{ngu\-way}}) \\
\textsc{Derivatives}: \ding{43} 
\on{1}.~\HL{ronguwayyu}{\B{ro\-ngu\-way\-yu}}, howler.

% ronguwayyu
\hi
\HT{ronguwayyu}{\B{ro\cp{ngu}wayyu}} \I{[Ro."Nu.waj.ju]} \emph{n.} howler, one who howls.
(\ding{43}~\HL{ronguway}{\B{ro\-ngu\-way}})

% ropx
\hi
\HT{ropx}{\B{ropx}} \I{[Rop']} \emph{n.} hole (\emph{going clear through an object}). 
\B{Fì\-fne\-ya\-yol tsrul\-it txu\-la mì song\-ropx ut\-ral\-ä fte ay\-li\-nit hi\-vaw\-nu wä sar\-ni\-o\-ang. Fì\-tì\-kan\-ìri ropx ke ha'.} This kind of bird builds its nest in a tree cavity so as to protect its young from predators. For this purpose, a hole going right through the tree trunk isn't suitable.~\LINK{http://naviteri.org/2013/04/wheres-the-bathroom-and-other-useful-things/}{b} \\
\textsc{Derivatives}: \ding{43} 
\on{1}.~\HL{tsongropx}{\B{tsong\-ropx}}, hole, cavity.

% rotxo
\hi 
\HT{rotxo}{\B{Ro\cp{txo}}} \I{[Ro."t'o]} \emph{n. male name}.

% rou
\hi
\HT{rou}{\B{\cp{ro}u}} \I{["Ro.u]} \emph{vin.} be drunk, get drunk.   
\B{Sra\-ne, fì\-txon tsun nga ni\-väk swo\-at nì\-'it, slä rä\-'ä rou!} Yes, you can have a little alcohol tonight, but don't get drunk!~\LINK{http://naviteri.org/2012/10/audio-and-video-learning-materials-for-navi-101/}{b}
\B{Nga tsun ni\-väk wew\-a pay\-it. Ul\-te ke ra\-you!} You can drink cold water. And not get drung!~\LINK{http://naviteri.org/2012/10/audio-and-video-learning-materials-for-navi-101/}{b}

% -ru
\hi
\HT{-ru}{\B{-ru}} \I{[.Ru]} \emph{suf. to mark the dative case of a n. ending in a vowel, diphthong or glottal stop}.
(\ding{43}~\HL{-r}{\B{-r}}, \HL{-ur}{\B{-ur}})

% rum
\hi
\HT{rum}{\B{rum}} \I{[Rum]} \emph{or} \I{[RUm]} \emph{n.} ball. 
\B{Rum 'o\-lrr\-ko oe\-ne kll\-te\-sìn.} The ball rolled towards me on the ground.~\LINK{http://naviteri.org/2015/04/some-new-words-for-may-day/}{b}
\B{Pol rum\-it ko\-lìm.} He spun the ball.~\LINK{http://forum.learnnavi.org/language-updates/ftarpa-and-skienpa-added-to-dict/msg570362/\#msg570362}{ln}
\B{Rum\-it a no\-lu\-i rä\-'ä fe\-wi.} Don't chase after a foul ball.~\LINK{http://naviteri.org/2013/04/wheres-the-bathroom-and-other-useful-things/}{b} 
(\emph{a. idiom}) \B{Ngä\-zìk fwa fkol rum\-it a\-ku\-'up pey\-ku\-wup.} \emph{it's difficult to make s.o. stubborn or inept do what you want them to} [lit.: it's hard to bounce a heavy ball].~\LINK{http://naviteri.org/2023/12/mipa-zisit-ayliu-amip-new-year-new-words/}{b} \\  
\textsc{Derivatives}: \ding{43} 
\on{1}.~\HL{rumut}{\B{ru\-mut}}, puffball tree. 
\on{2}.~\HL{rumaut}{\B{ru\-ma\-ut}}, cannonball fruit tree.
\on{3}.~\HL{rumtsyokx}{\B{rum\-tsyokx}}, fist.

% rumaut
\hi 
\HT{rumaut}{\B{\cp{ru}maut}} \I{["Ru.ma.ut\textcorner]} \emph{or coll.} \I{["Ru.mawt\textcorner]} \emph{n.} \ding{96} cannonball fruit tree, \emph{ecdurus putamen pomus}.~\LINK{https://forum.learnnavi.org/language-updates/expansion-to-flora-fauna-and-places-from-the-valley-of-mo'ara/}{ln}  
(\ding{43}~\HL{rum}{\B{rum}}, \HL{mauti}{\B{mauti}})

% rumtsyokx
\hi 
\HT{rumtsyokx}{\B{\cp{rum}tsyokx}} \I{["Rum.\t{ts}jok']} \emph{or} \I{["RUm.]} \emph{n.} fist. 
\B{A\-te\-yol En\-tut to\-la\-kuk fa rum\-tsyokx.} Ateyo punched Entu (i.e. struck him with his fist).~\LINK{https://naviteri.org/2024/06/ayhapxi-tokxa-si-ayliu-alahe-parts-of-the-body-and-more/}{b} 
(\ding{43}~\HL{rum}{\B{rum}}, \HL{tsyokx}{\B{tsyokx}})

% rumut
\hi
\HT{rumut}{\B{\cp{ru}mut}} \I{["Ru.mut\textcorner]} \emph{or} \I{[.mUt\textcorner]} \emph{n.}, \ding{96} puffball tree, \emph{obesus rotundus}. 
\B{Fì\-rum\-ut lum\-pe ke kxem?} Why isn't this puffball tree vertical?~\LINK{http://naviteri.org/2012/03/spring-vocabulary-part-1/}{b}
(\ding{43}~\HL{rum}{\B{rum}}, \HL{utral}{\B{ut\-ral}})  \\
\textsc{Derivatives}: \ding{43} 
\on{1}.~\HL{vArumut}{\B{vä\-ru\-mut}}, vein pod.

% run
\hi
\HT{run}{\B{run}} \I{[Run]} \emph{or} \I{[RUn]} (\textsc{rn}: \B{run} \I{[Run]}) \emph{vtr.} (\emph{a.}) find, discover.
\B{Run fkol tey\-lut nì\-ho\-et.} You find \emph{teylu} everywhere.~\LINK{http://naviteri.org/2012/06/spring-vocabulary-part-3/}{b}  
\B{Slä fma\-syi ri\-vun e\-yawr\-a tì\-'eyng\-it.} But I'll try to find the correct answer.~\LINK{http://naviteri.org/2011/08/new-vocabulary-clothing/comment-page-1/\#comment-920}{b}
\B{Ira\-yo 'ey\-lan\-ur aw\-nge\-yä a kxey\-ey\-ti ro\-lun} Thanks to our friend who found the mistake.~\LINK{http://forum.learnnavi.org/language-updates/another-email-from-frommer/msg68133/\#msg68133}{ln} 
\B{Zi\-va\-'u nì\-mun ye'\-rìn \ldots tsa\-krr ra\-yun ay\-ngal ay\-u\-pxa\-ret a\-mip.} Come again soon \ldots then you'll find new messages.~\LINK{http://naviteri.org/2010/06/first-post/}{b}
\B{Ro\-le\-i\-un fu\-ta lu fo ka\-nu sì le\-so'\-ha nì\-wotx.} I found that they are intelligent and completely enthusiastic.~\LINK{http://naviteri.org/2013/08/fmawn-a-fkol-kilmulat-this-just-in/}{b}
(\emph{b. conv. idiom}) \B{Ro\cp{lun}!} ``Eureka! I found it!''~\LINK{http://naviteri.org/2010/07/vocabulary-update/}{b} 
(\ding{214}~\B{fwew}) \\  
\textsc{Derivatives}: \ding{43}
\on{1}.\HL{ultxarun}{\B{ul\-txa\-run}}, encounter, meet by chance. 
\on{2}.~\HL{lepxìmrun}{\B{le\-pxìm\-run}}, common, often found.
\on{3}.~\HL{tIrun}{\B{tì\-run}}, discovery.

% rurur
\hi
\HT{rurur}{\B{ru\cp{rur}}} \I{[Ru."RuR]} \emph{or} \I{[."RUR]} (\textsc{rn}: \B{rù\cp{rùr}} \I{[RU."RUR]}) \emph{uncountable n.} kind of waterfall, \emph{water that is aerated while flowing among the rocks of a very gradually sloping stream}.

% rusey
\hi
\HT{rusey}{\B{ru\cp{sey}}} \I{[Ru."{}sEj]} \textsc{i}. \emph{adj.} alive, living.
\B{Krro krro, flì\-a vul a\-ru\-sey to nutx\-a pum a\-ke\-ru\-sey lu txur.} A thin living branch is sometimes stronger than a thick dead one.~\LINK{http://naviteri.org/2011/11/\%E2\%80\%99a\%E2\%80\%99awa-ayli\%E2\%80\%99u-amip-ni\%E2\%80\%99aw-\%E2\%80\%94-a-few-new-words-only/}{b} 
(\ding{214}~\HL{kerusey}{\B{ke\-ru\-sey}}) \\ 
\textsc{ii}. \emph{n.} living thing.~\LINK{http://forum.learnnavi.org/language-updates/conceptual-adjectives/}{ln}
(\ding{43}~\HL{rey}{\B{rey}})

% rutxe
\hi
\HT{rutxe}{\B{ru\cp{txe}}} \I{[Ru."t'E]} \emph{intj.} please.
\B{Spi\-vaw oe\-ti ru\-txe, ma oe\-yä ey\-lan.} Please believe me, my friends.~\LINK{http://forum.learnnavi.org/news-announcements/a-response-from-paul-frommer!/}{ln}
\B{Ru\-txe ti\-vìng mik\-yun, ma fra\-po.} Your attention, please, everyone!~\LINK{http://naviteri.org/2010/09/getting-to-know-you-part-1/}{b}
\B{Syaw oer ru\-txe trr\-ay fa spul\-mok\-ri.} Please phone me tomorrow.~\LINK{http://naviteri.org/2012/07/fivospxiya-aylifyavi-amip-this-months-new-expressions-pt-2/}{b}

\end{multicols}

% ===== Section S ===== 

 \pdfbookmark[0]{S}{S}
\begin{center}
\section*{S \I{[s]} (\emph{Sä})}
\end{center}

\begin{multicols}{2}

% sa
\hi 
\HT{sa}{\B{sa}} \I{[sa]} \emph{vin.} (\emph{a.}) rise to a challenge. 
(\emph{b. idioms}) \B{Si\-\cp{va} ko!} Rise to the challenge! Courage! You can do it!~\LINK{http://naviteri.org/2016/06/mrrvola-lifyavi-amip-forty-new-expressions/}{b} 
\B{So\-lei\cp{a}}! \I{[so.lE.i."{}a]} \emph{or} \I{[so.lEj."{}a]} Congratulations! Nice going! You met the challenge! You did it!~\LINK{http://naviteri.org/2016/06/mrrvola-lifyavi-amip-forty-new-expressions/}{b}
\B{Sa\cp{sya}!} I'll rise to the challenge! I can do it!~\LINK{http://naviteri.org/2016/06/mrrvola-lifyavi-amip-forty-new-expressions/}{b}
\B{Ätxä\-le si, oer syi\-vaw trr\-ay ha'\-ngir fa Skxayp?---Sa\-sya!} Could you call me on Skype tomorrow afternoon?---Sure will!\slash Will do!~\LINK{http://naviteri.org/2016/06/mrrvola-lifyavi-amip-forty-new-expressions/}{b}

% sa'ewrang
\hi
\HT{sa'ewrang}{\B{sa\cp{'ew}rang}} \I{[sa."PEw.RaN]} \emph{n.} great loom, mother loom.
(\ding{43}~\HL{sa'nok}{\B{sa'\-nok}}, \HL{'ewrang}{\B{'ew\-rang}})

% sa'nok
\hi
\HT{sa'nok}{\B{\cp{sa'}nok}} \I{["{}saP.nok\textcorner]} \emph{n.} mother. 
\B{Trr\-am to\-ler\-kä\-ngup sa'\-nok a\-yaw\-ne.} My dear mother died yesterday.~\LINK{http://naviteri.org/2012/01/more-additions-to-the-lexicon/}{b}
\B{Sa'\-nok\-ìl oe\-yä tsmu\-ket a\-mip no\-lokx trr\-am.} Mom gave birth to my new sister yesterday.~\LINK{http://naviteri.org/2010/09/getting-to-know-you-part-2/}{b}
\B{Tsa\-ye\-rik\-tsyìp\-ìl li 'an\-la sa'\-nok\-it.} That little hexapede is already yearning for its mother.~\LINK{http://naviteri.org/2013/01/awvea-posti-zisita-amip-first-post-of-the-new-year/}{b}
\B{Slä sa'\-nok\-ìri fpìl\-fya lu ke\-teng.} But mothers way of thinking is different.~\LINK{http://naviteri.org/2010/07/ting-mikyun-fte-tslivam-listening-comprehension-1-3/}{b}
(i.) \B{Nawm\-a Sa'\-nok}, the Great Mother, Eywa.
(\emph{a. idiom}) \B{(Tsa\-ri) Nga\-ru lrr\-tok (si\-vi (Nawm\-a Sa'\-nok))!} Good luck with that! [lit.: (concerning that) may the great mother smile upon you]~\LINK{http://forum.learnnavi.org/language-updates/good-luck-with-this-one!/}{ln}
\B{Fì\-lì\-'u\-ä ral lu tì\-me'\-em sì tì\-ru\-sey mì hi\-fkey na Nawm\-a Sa'\-nok\-ä ha\-pxì.} The meaning of this word is ``harmony, living in the world as part of the Great Mother''.~\LINK{http://forum.learnnavi.org/language-updates/trr-rrtaya/}{ln} \\
\textsc{Derivations}: \ding{43}
\on{1}.~\HL{sa'nu}{\B{sa'\-nu}}, mommy. 
\on{2}.~\HL{sa'sem}{\B{sa'\-sem}}, parents, a parent. 
\on{3}.~\HL{sa'ewrang}{\B{sa'\-ewrang}}, great loom, mother loom.

% sa'nu
\hi
\HT{sa'nu}{\B{\cp{sa'}nu}} \I{["{}saP.nu]} \emph{n.} mommy. 
\B{Mong prr\-nen\-ìl fu\-ta sa'\-nul ve\-rewng nì\-tut.} The infant depends on her Mommy to look after her constantly.~\LINK{http://naviteri.org/2012/07/meetings-waterfalls-and-more/}{b}
\B{Fnu, ma 'evi. Sa'\-nur le\-ru haw\-tsyìp. Tsi\-vu\-rokx ko.} Quiet, young one. Mommy is taking a nap. Let her rest.~\LINK{http://naviteri.org/2011/07/txantsana-ultxa-mi-siatll-great-meeting-in-seattle/}{b}
\B{Nì\-trr\-trr oel unil\-tsa sa'\-nu\-ä tey\-lut.} I regularly dream of my mom's \emph{teylu}.~\LINK{http://naviteri.org/2015/03/kap-si-ayunil-saylahe-threats-dreams-and-other-things/}{b}
(\ding{43}~\HL{sa'nok}{\B{sa'\-nok}})

% sa'sem
\hi
\HT{sa'sem}{\B{\cp{sa'}sem}} \I{["{}saP.sEm]} \emph{n.} parents, one parent.
\B{Nì\-ke\-ftxo lu fì\-'e\-veng\-ur 'aw\-a sa'\-sem nì\-'aw.} Unfortunately this child has only one parent.~\LINK{http://naviteri.org/2011/08/new-vocabulary-clothing/comment-page-1/\#comment-917}{b}
\B{Se\-vin\-a tsa\-'eve\-ru a\-hì\-'i me\-sa'\-sem ioi so\-li fa mik\-tsang.} The parents put earrings on that pretty little girl.~\LINK{http://naviteri.org/2011/08/new-vocabulary-clothing/}{b}
\B{Sa'\-sem\-ìri lu 'eveng\-ä kxitx yeng\-wal a\-tu\-vom.} For a parent, the greatest sorrow is the death of a child.~\LINK{http://naviteri.org/2015/04/some-new-words-for-may-day/}{b}
(\ding{43}~\HL{sa'nok}{\B{sa'\-nok}}, \HL{sempul}{\B{sem\-pul}})

% sa'u
\hi
\HT{sa'u}{\B{\cp{sa}'u}} \I{["{}sa.Pu]}~:~\B{'u} (short plural of \HL{tsa'u}{\B{tsa'u}}).

% saa
\hi
\HT{saa}{\B{saa}} \I{[sa:]} \emph{intj.} \emph{threatening cry}.

% Saeyla
\hi
\B{Saeyla} \I{[sa.Ej.la]} \emph{n.} \emph{female name}.~\LINK{}{dh}

% saho
\hi 
\HT{saho}{\B{sa\cp{ho}}} \I{[sa."ho]} \emph{n.} prayer (\emph{concrete concept}). 
(\ding{43}~\HL{aho}{\B{aho}})

% sal
\hi
\HT{sal}{\B{sal}} \I{[sal]}~:~\B{'u} (\emph{shortened agentive form of the plural of} \HL{tsa'u}{\B{tsa\-'u}}).

% salew 
\hi
\HT{salew}{\B{sa\cp{lew}}} \I{[sa."lEw]} \emph{vin.} (\emph{a.}) proceed, go.
\B{Nga\-ri hu Ey\-wa sa\-lew ti\-rea, tokx 'ì\-'awn, slu Na'\-vi\-yä ha\-pxì.} Your spirit goes with Eywa, your body stays behind to become part of the People.~\LINK{http://wiki.learnnavi.org/index.php?title=Na\%27vi_from_Avatar_Movie\#Jake.27s_final_kill}{a\slash wiki}
\B{Ay\-nge\-yä ay\-srr tì\-flä\-ta te\-ya li\-vu ul\-te si\-va\-lew nì\-'o' nì\-'aw.} May your days be full of success and proceed only in a fun way.~\LINK{http://naviteri.org/2013/12/mipa-zisit-lefpom-happy-new-year/}{b}
\B{Sa\-lew nì\-yey\-fya.} Proceed straight ahead.~\LINK{http://naviteri.org/2013/09/on-si-salewfya-shapes-and-directions/}{b}
(\emph{b. used for stating age}) 
\B{Nga\-ri so\-la\-lew pol\-pxay\-a zì\-sìt?\slash Nga\-ri so\-la\-lew zì\-sìt a\-pol\-pxay?} How old are you?~\LINK{http://naviteri.org/2010/09/getting-to-know-you-part-2/}{b} 
\B{Oe\-ri so\-la\-lew (so\-lew) zì\-sìt a\-pxe\-vol.} I'm 24 years old.~\LINK{http://naviteri.org/2010/09/getting-to-know-you-part-2/}{b} \\
\textsc{Derivations}: \ding{43}
\on{1}.~\HL{salewfya}{\B{salewfya}}, direction, course. 
\on{2}.~\HL{zIsItsaltrr}{\B{zìsìtsaltrr}}, (yearly) anniversary.

% salewfya
\hi
\HT{salewfya}{\B{sa\cp{lew}fya}} \I{[sa."lEw.fja]} \emph{n.} direction, course.
\B{Swey\-lu set txo aw\-nga ki\-vä pe\-sa\-lew\-fya\slash pe\-sa\-lew\-fya\-ìlä?} What direction should we go in now?~\LINK{http://naviteri.org/2013/09/on-si-salewfya-shapes-and-directions/}{b}
\B{Aw\-nga ze\-ne vi\-var ìlä fì\-sa\-lew\-fya. Zen\-ke yak si\-vi.} We must continue in this direction. We must not go astray.~\LINK{http://naviteri.org/2013/09/on-si-salewfya-shapes-and-directions/}{b}
(\emph{a. idiom}) \B{Set pe\-sa\-lew\-fya?} What do we do now?~\LINK{http://naviteri.org/2013/09/on-si-salewfya-shapes-and-directions/}{b}
(\ding{43}~\HL{salew}{\B{sa\-lew}}, \HL{fya'o}{\B{fya\-'o}})

% san
\hi
\HT{san}{\B{san}} \I{[san]} \emph{part. to indicate the beginning of a quote with verbs of speech} \HL{plltxe}{\B{pll\-txe}}, \HL{peng}{\B{peng}}, \HL{pllngay}{\B{pll\-ngay}}, \HL{pawm}{\B{pawm}}, \emph{thought} \HL{fpIl}{\B{fpìl}}, \emph{and writing} \HL{pamrel}{\B{pam\-rel si}}; \emph{used together with} \ding{43}~\HL{sIk}{\B{sìk}} \emph{when the speaker continues to say s.t. after the cited quote.}
\B{Pol\-txe Ey\-tu\-kan san oe ka\-yä sìk, slä oel pot ke spaw.} Eytukan said he would go but I don't believe him.~\LINK{http://wiki.learnnavi.org/index.php?title=Canon\#Extracts_from_various_emails}{wiki}
\B{New oe pi\-vll\-txe ay\-nga\-ru san kal\-txì sìk ul\-te ti\-vìng ay\-ngar lì\-'ut a tì\-'e\-fu\-mì oe\-yä lu lor fra\-to mì lì'\-fya le\-Na'\-vi.} I want to say hello to you and present to you the word that, in my opinion, is the most beautiful in the Na'vi language.~\LINK{http://forum.learnnavi.org/language-updates/trr-rrtaya/}{ln}
\B{Po\-lawm po san sra\-ke Sä\-li ho\-lum sìk.} He asked, ``Did Sally leave?''~\LINK{http://forum.learnnavi.org/language-updates/if-vs-whether-ftxey-fuke/}{ln} 
\B{Alo a\-mrr po\-an po\-lawm, slä fra\-lo poe pol\-txe san ke\-he.} He asked five times, but each time she said, `no.'~\LINK{http://forum.learnnavi.org/language-updates/txelanit-hivawl/}{ln}
\B{Po ke tsun pi\-vll\-ngay san oe\-ru tì\-kxey.} He can't admit he's wrong.~\LINK{http://naviteri.org/2011/07/txantsana-ultxa-mi-siatll-great-meeting-in-seattle/}{b}
\B{Oe pll\-ngay san mo\-lak\-to oe nì\-fe', ta\-fral sno\-laytx; wä\-tu lu oe\-to txur.} I acknowledge that I rode badly, so I lost; my opponent was stronger than I was.~\LINK{http://naviteri.org/2011/07/txantsana-ultxa-mi-siatll-great-meeting-in-seattle/}{b}
\B{Oel ho\-ren\-it 'Ìng\-lì\-sì\-yä so\-lar, nì\-teng\-fya fko pam\-rel si\-vi san \ldots} I used the rules of English, the same way as you write, ``\ldots''.~\LINK{http://forum.learnnavi.org/language-updates/fmawno/}{ln}
\B{Oe pam\-rel so\-li nga\-ru 'upxa\-rer san Oe za\-ya\-'u trr\-ay.} I wrote a message to you, saying ``I'll come tomorrow''.~\LINK{http://forum.learnnavi.org/language-updates/some-news-from-berlin/}{ln}

% sang
\hi
\HT{sang}{\B{sang}}\textsuperscript{\on{1}} \I{[saN]} \emph{adj.} warm. 
\B{Ya\-ri som\-wew\-pe (set)? Ya lu sang.} What's the temperature? It's warm.~\LINK{http://naviteri.org/2012/10/audio-and-video-learning-materials-for-navi-101/}{b}
\B{Ay\-nga\-ri teng\-krr ya wur sle\-ru nì\-'ul, sìl\-pey oe, li\-vu hel\-ku sang ul\-te te'\-lan le\-fpom.} As for you all, while it's getting cooler, I hope your homes may be warm and your hearts happy.~\LINK{http://naviteri.org/2013/09/on-si-salewfya-shapes-and-directions/}{b}
\B{Sìl\-pey oe tsnì fpom li\-vu ay\-nga\-ru nì\-wotx ul\-te snge\-rä\-'i a zì\-sì\-krr a\-sang zi\-vaw\-prr\-te' ay\-nga\-ne.} I hope that you are all well and that you are enjoying the beginning warm season.~\LINK{http://naviteri.org/2015/04/some-new-words-for-may-day/}{b}
(\ding{214}~\HL{wur}{\B{wur}})

\hi
\B{sang}\textsuperscript{\on{2}} : \HL{tsang}{\B{tsang}}.

% sar
\hi
\HT{sar}{\B{sar}}\textsuperscript{\on{1}} \I{[saR]} \emph{vtr.} use.
\B{Tsat sa\-yar oel nìl\-tsan.} I will use it well.~\LINK{http://forum.learnnavi.org/language-updates/score!-\%281\%29-the-verb-to-use-\%282\%29-more-detail-on-teya-si/}{ln}
\B{Kem\-lì\-'ut alu var fkol sar pxel pum alu new.} One uses the verb \emph{var} like \emph{new}.~\LINK{http://naviteri.org/2011/03/word-order-and-case-marking-with-modals/comment-page-1/\#comment-582}{b}
\B{Fol lì'\-fya\-ti aw\-nge\-yä sar a fya\-'o lu ko\-sman.} It's wonderful the way they use our language.~\LINK{http://naviteri.org/2011/09/miscellaneous-vocabulary/}{b}
\B{Txep\-ìri, yrr\-a rìn\-ti rä\-'ä sar, ma skxawng.} Don't use that green wood for a fire, you fool.~\LINK{http://naviteri.org/2013/01/awvea-posti-zisita-amip-first-post-of-the-new-year/}{b}
\B{Tskot sngo\-lä\-'i po si\-var 'a\-'aw\-a trr\-kam.} He started to use the bow several days ago.~\LINK{http://naviteri.org/2011/09/miscellaneous-vocabulary/}{b}
\B{'En si oe, lor\-a tsa\-fkxi\-le\-ri a\-pxay\-opin so\-lar Tse\-nul srok\-it a\-vo\-zam.} I would guess that Tsenu used a thousand (lit. \on{512}) beads for that beautiful multi\hyp{}colored bib necklace.~\LINK{http://naviteri.org/2013/07/pximaw-tsawlultxa-just-after-the-meet-up/}{b}

\hi
\B{sar}\textsuperscript{\on{2}} : \B{'u} (\emph{shortened dative form of the plural of} \HL{tsa'u}{\B{tsa\-'u}}).

% sari
\hi
\B{\cp{sa}ri} \I{["{}sa.Ri]}~:~\B{'u} (\emph{shortened topic form of the plural of} \HL{tsa'u}{\B{tsa\-'u}}).

% sat
\hi
\B{sat} \I{[sat\textcorner]}~:~\B{'u} (\emph{shortened patientive form of the plural of} \HL{tsa'u}{\B{tsa\-'u}}).

% satu'li
\hi
\HT{satu'li}{\B{sa\cp{tu'}li}} \I{[sa."tuP.li]} \emph{n.} heritage.

% sau
\hi
\HT{sau}{\B{sa\cp{u}}} \I{[sa."{}u]} \emph{intj.} \emph{exclamation upon exertion}.

% sawtute
% \hi
% \HT{sawtute}{\B{\cp{saw}tute}} \I{["{}saw.tu.tE]}~:~\B{tawtute}. 

% say
\hi 
\HT{say}{\B{say}} \I{[saj]} \emph{adj.} loyal. 
\B{Le\-iu po ken\-'aw say\-rìp slä\-kop say.} He's not only handsome but also, I'm happy to say, loyal.~\LINK{http://naviteri.org/2021/12/zolau-niprrte-ma-3746-welcome-2022/}{b} \\
\textsc{Derivations}: \ding{43}
\on{1}.~\HL{tIsay}{\B{tì\-say}}, loyalty. 
\on{2}.~\HL{nIsay}{\B{nì\-say}}, loyally.

% saylahe
\hi
\HT{saylahe}{\B{say\cp{la}he}} \I{[saj."la.hE]} \emph{adv.} et cetera (etc.). (\emph{from} \B{sì ay\-la\-he}, \emph{in writing on Earth may be shortened to} \B{sl.}) 
\B{Kap sì ay\-u\-nil say\-la\-he.} Threats, dreams, and other things.~\LINK{http://naviteri.org/2015/03/kap-si-ayunil-saylahe-threats-dreams-and-other-things/}{b}

% sayrìp
\hi
\HT{sayrIp}{\B{\cp{say}rìp}} \I{["{}saj.RIp\textcorner]} \emph{adj.} handsome, good\hyp{}looking (\emph{primarily of males}).
\B{Hu\-fwa lu fil\-ur Va'\-ru\-ä fne\-fe'\-ran\-vi, tsal\-su\-ngay fpìl fu\-ta say\-rìp lu nì\-txan.} Although Va'ru's facial stripes are rather uneven, I still think he's very handsome.~\LINK{http://naviteri.org/2012/10/mipa-vospxi-mipa-ayliu-new-words-for-the-new-month/}{b}
(\ding{43}~\HL{sevin}{\B{se\-vin}}, \HL{lor}{\B{lor}})

% sä-
\hi
\HT{sA-}{\B{sä-}} \I{[s\ae.]} \emph{unprod. pref. to create n.s from v.s and adj.s with the meaning of (1) an instrument or tool, or the means by which s.t. is achieved, or of (2) a particular, concrete instance of a general action.}
(\ding{43}~\HL{tI-}{\B{tì-}})

% sä'anla
\hi
\HT{sA'anla}{\B{sä\cp{'an}la}} \I{[s\ae."Pan.la]} \emph{n.} yearning.
\B{Fì\-sä\-'an\-lal Ne\-we\-yä nga\-ti sley\-ka\-yu lek\-ye\-'ung!} This yearning for Newey is going to drive you crazy.~\LINK{http://naviteri.org/2013/01/awvea-posti-zisita-amip-first-post-of-the-new-year/}{b}
(\ding{43}~\HL{'anla}{\B{'an\-la}})

% sä'eoio
\hi
\HT{sA'eoio}{\B{sä\cp{'e}oio}} \I{[s\ae."PE.o.i.o]} \textsc{i}. \emph{n.} ceremony, ritual, rite. \\
 \textsc{ii}. \B{sä\cp{'e}oio si} \I{[s\ae."PE.o.i.o si]} \emph{vin.} take part in a ceremony, perform a ritual.
\B{Fwa tsyìl Ay\-ram\-it A\-lu\-sìng lu sä\-'e\-o\-i\-o a ze\-ne fra\-po si\-vi fte sl\-ivu ta\-ron\-yu.} Climbing Iknimaya is a ritual that everyone has to perform to become a hunter.~\LINK{http://naviteri.org/2013/04/wheres-the-bathroom-and-other-useful-things/}{b}
(\ding{43}~\HL{'eoio}{\B{'eoio}})

% sä'ipu
\hi
\HT{sA'ipu}{\B{sä\cp{'i}pu}} \I{[s\ae."Pi.pu]} \emph{n.} s.t. humorous, a particular instance of being humorous, e.g. a joke.
\B{Oe\-ru txoa li\-vu, ma 'ey\-lan. Rä\-'ä sti\-vi. Lu hì\-'i\-a sä\-'i\-pu nì\-'aw.} I'm sorry, friend. Don't be angry. It was just a small bit of humor.~\LINK{http://naviteri.org/2012/01/mipa-zisit-ayliu-amip-new-words-for-the-new-year/}{b}
(\ding{43}~\HL{'ipu}{\B{'i\-pu}}, \HL{tI'ipu}{\B{tì\-'i\-pu}})

% sä'o
\hi
\HT{sA'o}{\B{\cp{sä}'o}} \I{["{}s\ae.Po]} \emph{n.} tool, utensil.   
\B{Na'\-vi\-ri txin\-a sä\-'o tì\-tu\-sa\-ron\-ä lu tsko swi\-zaw.} For the Na'vi the bow and arrow is the main hunting tool.~\LINK{http://naviteri.org/2011/01/new-year-new-vocabulary/}{b}
\B{Pol sä\-'o\-ti wo\-lan äo ay\-rìk.} He hid his tool under the leaves.~\LINK{http://naviteri.org/2011/02/new-vocabulary-ii\%E2\%80\%94part-1/}{b} \\
\textsc{Derivations}: \ding{43}
\on{1}.~\HL{tsamsA'o}{\B{tsam\-sä\-'o}}, weapon of war.

% säfe'ul
\hi
\HT{sAfe'ul}{\B{sä\cp{fe}'ul}} \I{[s\ae."fE.Pul]} \emph{or} \I{[.PUl]} \emph{n.} worsening (\emph{specific instance}). 
\B{Pe\-yä sä\-ngä\-'än\-ìri lu sä\-fe\-'ul le\-yew\-la nì\-txan.} The worsening of his depression is very disappointing.~\LINK{http://naviteri.org/2013/03/tsanerul-feerul-getting-better-getting-worse/}{b}
(\ding{43}~\HL{fe'ul}{\B{fe\-'ul}}, \HL{tIfe'ul}{\B{tì\-fe\-'ul}})

% säflel
\hi 
\HT{sAflel}{\B{sä\cp{flel}}} \I{[s\ae."flEl]} \emph{n.} trick, hoax, dishonest act or scheme. 
\B{Pot spaw rä\-'ä! Lu fì\-'u sä\-flel!} Don't believe him! It's a trick!~\LINK{https://naviteri.org/2024/03/fleltrra-ayliu-words-for-april-fools-day/}{b}
(\ding{43}~\HL{flel}{\B{flel}}, \HL{tIflel}{\B{tì\-flel}}) \\
\textsc{Derivations}: \ding{43} 
\on{1}.~\HL{sAfleltsyIp}{\B{sä\-flel\-tsyìp}}, practical joke.

% säfleltsyìp
\hi 
\HT{sAfleltsyIp}{\B{sä\cp{flel}tsyìp}} \I{[s\ae."flEl.\t{ts}jIp\textcorner]} \emph{n.} practical joke. 
(\ding{43}~\HL{sAflel}{\B{säflel}})

% säflä
\hi
\HT{sAflA}{\B{sä\cp{flä}}} \I{[s\ae."fl\ae]} \emph{n.} success, an instance of succeeding.
\B{Te\-ri lì'\-fya le\-Na'\-vi a tsa\-nu\-mul\-txa lo\-le\-i\-u sä\-flä!} The class about the Na'vi language was a success.~\LINK{http://naviteri.org/2012/07/tskxekengtsyip-a-mikyunfpi-a-little-listening-exercise/}{b}
(\ding{43}~\HL{flA}{\B{flä}}, \HL{tIflA}{\B{tì\-flä}})

% säfmi
\hi
\HT{sAfmi}{\B{sä\cp{fmi}}} \I{[s\ae."fmi]} \emph{n.} attempt. 
(\emph{a. proverb}) \B{Flä ke flä, ley sä\-fmi.} Whether you succeed or not, the attempt has value.~\LINK{http://naviteri.org/2014/03/value-and-worth/}{b}
(\ding{43}~\HL{fmi}{\B{fmi}}, \HL{tIfmi}{\B{tì\-fmi}})

% säfmong
\hi 
\HT{sAfmong}{\B{sä\cp{fmong}}} \I{[s\ae."fmoN]} \emph{n.} theft (\emph{particular instance}). 
\B{Poe\-ri sä\-fmong lor\-a tsa\-fkxi\-le\-yä lo\-lu na ay\-skxe mì te'\-lan.} For her, the theft of that beautiful necklace was like stones in her heart.~\LINK{http://naviteri.org/2021/09/aawa-ayliu-amip-a-few-new-words/}{b} 
(\ding{43}~\HL{fmong}{\B{fmong}})

% säfngo'
\hi
\HT{sAfngo'}{\B{sä\cp{fngo'}}} \I{[s\ae."fNoP]} \emph{n.} requirement, demand.   
\B{Nge\-yä fay\-sä\-fngo'\-ìl nì\-wotx stey\-ke\-rä\-ngi oe\-ti nì\-hawng.} All these demands of yours are making me exceedingly angry.~\LINK{http://naviteri.org/2012/01/mipa-zisit-ayliu-amip-new-words-for-the-new-year/}{b}
(\ding{43}~\HL{fngo'}{\B{fngo'}})

% säfpìl
\hi
\HT{sAfpIl}{\B{sä\cp{fpìl}}} \I{[s\ae."fpIl]} \emph{n.} idea, thought.   
\B{Nge\-yä sä\-fpìl Saw\-tu\-te\-te\-ri hei\-ek oer nì\-txan.} Your idea about the Sky People is very interesting to me (because it seems unusual).~\LINK{http://naviteri.org/2010/07/vocabulary-update/}{b}
\B{Lor\-a ay\-lì\-'u, lor\-a ay\-sä\-fpìl.} Beautiful words and thoughts.~\LINK{http://masempul.org/2010/07/\%E2\%80\%99ite-polaheiem/}{sempul} 
\B{Saw\-tu\-te lu ay\-vrr\-tep nì\-wotx a sä\-fpìl\-it oel wä\-te.} I dispute the idea that the Skypeople are all demons.~\LINK{http://naviteri.org/2011/01/new-year-new-vocabulary/}{b}
\B{Pol sä\-fpìl\-it ve\-rar wi\-van.} He's keeping his idea a secret.~\LINK{http://naviteri.org/2011/02/new-vocabulary-ii\%E2\%80\%94part-1/}{b}
(\emph{a. metaphorically}) \B{ay\-sä\-fpìl a\-txukx}, deep thoughts.~\LINK{http://naviteri.org/2010/07/vocabulary-update/}{b}
\B{mek\-a sä\-fpìl}, an empty\slash dumb idea.~\LINK{http://naviteri.org/2012/02/trr-asawnung-lefpom-happy-leap-day/}{b}
(\emph{b. proverb}) \B{Sä\-fpìl a\-steng, tì\-kan a\-teng.} Great minds think alike [lit.: similar thought, same aim].~\LINK{http://naviteri.org/2011/09/miscellaneous-vocabulary/}{b}
(\ding{43}~\HL{fpIl}{\B{fpìl}}) \\
\textsc{Derivations}: \ding{43}
\on{1}.~\HL{txinfpIl}{\B{txin\-fpìl}}, main point.
\on{2}.~\HL{ngrrfpIl}{\B{ngrr\-fpìl}}, key assumption.
\on{3}.~\HL{sAfpIlyewn}{\B{sä\-fpìl\-yewn}}, communicate.

% säfpìlyewn
\hi 
\HT{sAfpIlyewn}{\B{sä\cp{fpìl}yewn}} \I{[s\ae."fpIl.j$\cdot$Ewn]} \emph{vin., inf.}\on{33} communicate. 
\B{Txo po lu kak\-mok\-ri, fya\-pe sä\-fpìl\-yewn?} If he's mute, how does he communicate?~\LINK{https://naviteri.org/2024/02/mipa-ayliu-si-aylifyavi-niul-more-new-words-and-expressions/}{b} 
(\ding{43}~\HL{sAfpIl}{\B{sä\-fpìl}}, \HL{yewn}{\B{yewn}}) \\
\textsc{Derivations}: \ding{43} 
\on{1}.~\HL{tIsAfpIlyewn}{\B{tì\-sä\-fpìl\-yewn}}, communication.

% säfrìp
\hi
\HT{sAfrIp}{\B{sä\cp{frìp}}} \I{[s\ae."fRIp\textcorner]} \emph{n.} a bite.
\B{Lam fwa skxir za\-'u ta sä\-frìp nan\-tang\-ä.} It seemed the wound came from a bite of a viperwolf.~\LINK{http://naviteri.org/2020/04/awa-tskxekengtsyip-a-mikyunfpi-niul-one-more-little-listening-exercise/}{b}
(\ding{43}~\HL{frIp}{\B{frìp}})

% säfrrfen
\hi 
\HT{sAfrrfen}{\B{sä\cp{frr}fen}} \I{[s\ae."f\s{r}.fEn]} \emph{n.} a visit, an instance of visiting. 
\B{Fu\-ria nga zo\-la\-'u ira\-yo; ngey sä\-frr\-fen txa\-su\-nu oer.} Thank you for coming; I enjoyed your visit very much.~\LINK{http://naviteri.org/2021/12/zolau-niprrte-ma-3746-welcome-2022/}{b} 
(\ding{43}~\HL{frrfen}{\B{frr\-fen}})

% säftxulì'u
\hi
\HT{sAftxulI'u}{\B{säftxu\cp{lì}'u}} \I{[s\ae.ft'u."lI.Pu]} \emph{n.} speech, oration.   
\B{Sä\-ftxu\-lì\-'u Tsyeyk\-ä Na'\-vi\-ru slan\-ti\-re so\-li nì\-wotx.} Jake's speech inspired all the People.~\LINK{http://naviteri.org/2012/06/spring-vocabulary-part-3/}{b}
\B{Fni\-vu! Sä\-ftxu\-lì\-'u\-ri ke tsun oe sti\-vawm.} Hush! I can't hear the speech.~\LINK{http://naviteri.org/2012/06/spring-vocabulary-part-3/}{b}
\B{Su\-nu oe\-ru nì\-txan ay\-sä\-ftxu\-lì\-'u pe\-yä.} I like his speeches a lot.~\LINK{http://naviteri.org/2014/11/twenty-before-the-holidays/}{b}
(\ding{43}~\HL{ftxulI'u}{\B{ftxu\-lì\-'u}})

% säfyep
\hi
\HT{sAfyep}{\B{sä\cp{fyep}}} \I{[s\ae."fjEp\textcorner]} \emph{n.} handle (\emph{the part of an object where one holds it, e.g. the handle of a cup}).
(\ding{43}~\HL{fyep}{\B{fyep}})

% sähe'a
\hi 
\HT{sAhe'a}{\B{sä\cp{he}'a}} \I{[s\ae."hE.Pa]} \emph{n.} a cough, an instance of coughing. 
(\ding{43}~\HL{he'a}{\B{he\-'a}})

% sähena
\hi
\HT{sAhena}{\B{sä\cp{he}na}} \I{[s\ae."hE.na]} \emph{n.} container, vessel, carrier.
\B{Ay\-fol za\-mo\-lu\-nge aw\-ngar ay\-ri\-na'\-it fa sä\-he\-na a\-pxa.} They brought us the seeds in a large container.~\LINK{http://naviteri.org/2014/05/mipa-ayliu-mipa-aysafpil-new-words-new-ideas/}{b}
(\ding{43}~\HL{hena}{\B{he\-na}}) \\  
\on{1}.~\HL{paysena}{\B{pay\-se\-na}}, water container. 
\on{2}.~\HL{tutsena}{\B{tut\-se\-na}}, stretcher. 
\on{3}.~\HL{tstalsena}{\B{tstal\-se\-na}}, knife sheath. 
\on{4}.~\HL{swizawsena}{\B{swi\-zaw\-se\-na\slash zaw\-se\-na}}, quiver.

% säkahena
\hi 
\HT{sAkahena}{\B{säka\cp{he}na}} \I{[s\ae.ka."hE.na]} \emph{or coll.} \I{[ska."hE.na]} \emph{n.} means of transport, transporation device, vehicle (\emph{typically s.t. that can move large things, including people}).  
(\ding{43}~\HL{kahena}{\B{ka\-he\-na}}, \HL{tIkahena}{\B{tì\-ka\-he\-na}}; \HL{sAmunge}{\B{sä\-mu\-nge}})

% säkanom
\hi
\HT{sAkanom}{\B{sä\cp{ka}nom}} \I{[s\ae."ka.nom]} \emph{or coll.} \I{["{}ska.nom]} \emph{n.} something acquired, an acquisition, a possession.   
\B{Tì\-'efu\-mì oe\-yä, nge\-yä fì\-sä\-ka\-nom lu le\-hrr\-ap ul\-te tsun nga\-ti tì\-sraw sey\-ki\-vi.} In my opinion, this acquisition of yours is dangerous and can hurt you.~\LINK{http://naviteri.org/2012/01/more-additions-to-the-lexicon/}{b}
(\ding{43}~\HL{kanom}{\B{ka\-nom}})

% säkeynven
\hi 
\HT{sAkeynven}{\B{säkeyn\cp{ven}}} \I{[s\ae.kEjn."vEn]} \emph{or coll.} \I{[skEjn."vEn]} \emph{n.} (\emph{a. physical}) step. 
\B{Rawng tsa\-slär\-ä ftu fì\-tsenge lu vol\-a skeyn\-ven nì\-'aw.} The entrance to the cave is only eight steps from here.~\LINK{https://naviteri.org/2024/09/pxevola-liu-amip-two-dozen-new-words/}{b} 
(\emph{b. metaphorical}) step in a procedure.
(\ding{43}~\HL{keynven}{\B{keyn\-ven}}) 

% säkxange
\hi 
\HT{sAkxange}{\B{sä\cp{kxa}nge}} \I{[s\ae."k'a.NE]} \emph{or} \I{["{}sk'a.NE]} \emph{n.} a yawn. 
(\ding{43}~\HL{kxange}{\B{kxa\-nge}})

% sälang
\hi
\HT{sAlang}{\B{sä\cp{lang}}} \I{[s\ae."laN]} \emph{or coll.} \I{[slaN]} \emph{n.} expedition, investigation.
\B{Fì\-txe\-lel fngo' sä\-lang\-it a\-ta\-wrr\-pa.} This matter requires an external investigation.~\LINK{http://naviteri.org/2021/12/zolau-niprrte-ma-3746-welcome-2022/}{b}
\B{Kum sä\-lang\-ä le\-yew\-la lä\-ngu. Ke ro\-lun aw\-ngal ke\-'ut.} The result of the investigation was, I'm sorry to say, disappointing. We found nothing.~\LINK{http://naviteri.org/2014/11/twenty-before-the-holidays/}{b}
(\ding{43}~\HL{lang}{\B{lang}}, \HL{tIlang}{\B{tì\-lang}})

% sälatem
\hi
\HT{sAlatem}{\B{sä\cp{la}tem}} \I{[s\ae."la.tEm]} \emph{or coll.} \I{["{}sla.tEm]} \emph{n.} (instance of) change, edit, modification.
\B{Ay\-sä\-la\-tem\-it a za\-mo\-lu\-nge Saw\-tu\-tel nga fmi ki\-va\-ra. Lä\-ngu kel\-tsun.} You try to resist the changes brought by the Sky People. Sadly, that's impossible.~\LINK{http://naviteri.org/2023/09/aawa-ayliu-si-aylifyavi-amip-a-few-new-words-and-expressions/}{b}
\B{'On\-ìri tsko\-ä lu tì\-kin sä\-la\-tem\-ä a\-hì\-'i.} The form of the bow requires a small change.~\LINK{http://naviteri.org/2015/11/vomuna-liu-amip-ten-new-words/}{b}
(\ding{43}~\HL{latem}{\B{la\-tem}}, \HL{tIlatem}{\B{tì\-la\-tem}})

% sälätxayn
\hi
\HT{sAlAtxayn}{\B{sälä\cp{txayn}}} \I{[s\ae.l\ae."t'ajn]} \emph{or coll.} \I{[sl\ae."t'ajn]} \emph{n.} defeat, an instance of defeat.
\B{Tsa\-sä\-lä\-txayn Na'\-vi\-ru srung so\-li nì\-'aw fte sli\-vu txur nì\-'ul.} That defeat only helped the People become stronger.~\LINK{http://naviteri.org/2012/01/more-additions-to-the-lexicon/}{b}
(\ding{43}~\HL{lAtxayn}{\B{lä\-txayn}})

% säleymfe'
\hi 
\HT{sAleymfe'}{\B{säleym\cp{fe'}}} \I{[s\ae.lEjm."fEP]} \emph{or coll.} \I{[slEjm."fEP]} \emph{n.} complaint. 
\B{Ay\-sä\-leym\-fe'\-it pe\-yä yu\-ne rä\-'ä. Lu po ke\-ve\-'o\-tu nì\-'aw.} Don't listen to his complaints. He's just a troublemaker.~\LINK{https://naviteri.org/2025/03/conlang-adventure-and-more/}{b} 
(\ding{43}~\HL{leymkem}{\B{leym\-kem}})

% säleymkem
\hi 
\HT{sAleymkem}{\B{säleym\cp{kem}}} \I{[s\ae.lEjm.\-"kEm]} \emph{or coll.} \I{[slEjm.\-"kEm]} \emph{n.} protest, instance of protesting. 
\B{Eyk\-tan\-ìl nge\-yä sä\-leym\-kem\-it sto\-lawm ul\-te pa\-ye\-'un teyng\-ta ze\-ne fko pe\-hem si\-vi.} The leader has heard your protest and will decide what must be done.~\LINK{http://naviteri.org/2017/09/zisikrr-amip-ayliu-amip-new-words-for-the-new-season/}{b} 
(\ding{43}~\HL{leymkem}{\B{leym\-kem}})

% sälipx
\hi 
\HT{sAlipx}{\B{sä\cp{lipx}}} \I{[s\ae."lip']} \emph{or coll.} \I{[slip']} \emph{n.} drop ({\em of a liquid}). 
(\ding{43}~\HL{lipx}{\B{lipx}})

% sämewn
\hi 
\HT{sAmewn}{\B{sä\cp{mewn}}} \I{[s\ae."mEwn]} \emph{or coll.} \I{[smEwn]} \emph{n.} sigh, an instance of sighing. 
\B{Zo\-la\-'u ftu kä\-pxì num\-tseng\-vi\-yä sä\-mewn a\-wok.} From the back of the classroom came a loud sigh.~\LINK{https://naviteri.org/2025/04/mevosinga-liu-amip-twenty-new-words/}{b}
(\ding{43}~\HL{mewn}{\B{mewn}}) 

% sämok
\hi
\HT{sAmok}{\B{sä\cp{mok}}} \I{[s\ae."mok\textcorner]} \emph{or coll.} \I{[smok\textcorner]} \emph{n.} a suggestion, a concrete instance of suggesting.
\B{Fe\-yä ay\-sä\-mok lu fe' nì\-wotx.} All of their suggestions are bad.~\LINK{http://naviteri.org/2012/02/trr-asawnung-lefpom-happy-leap-day/}{b}
\B{Oel sä\-mok\-ti nge\-yä txa\-krr\-fpo\-lìl.} I have considered your suggestion.~\LINK{http://naviteri.org/2023/12/mipa-zisit-ayliu-amip-new-year-new-words/}{b} 
\B{Nge\-yä sä\-mok\-ìri a\-ko\-sman sei\-yi oe ira\-yo.} Thanks for that excellent suggestion (of yours).~\LINK{http://naviteri.org/2012/01/more-additions-to-the-lexicon/}{b}
(\emph{a. fig.}) \B{sä\-mok a\-mek}, a useless suggestion.~\LINK{http://naviteri.org/2012/02/trr-asawnung-lefpom-happy-leap-day/}{b}
(\ding{43}~\HL{mok}{\B{mok}}, \HL{tImok}{\B{tì\-mok}})

% sämunge
\hi
\HT{sAmunge}{\B{sä\cp{mu}nge}} \I{[s\ae."mu.NE]} \emph{or coll.} \I{["{}smu.NE]} (\textsc{rn}: \B{sä\cp{mù}nge} \I{[sE."mU.NE]}) \emph{n.} transportation tool or device (\emph{general term for any object used to carry or transport something else}).
\B{Sìn skxir yem 'um\-tsat a ho\-le\-na li pol mì sä\-mu\-nge.} He put medicine, which he had already carried in his pouch, on the wound.~\LINK{http://naviteri.org/2020/04/awa-tskxekengtsyip-a-mikyunfpi-niul-one-more-little-listening-exercise/}{b}
(\ding{43}~\B{mu\-nge}, \B{sä\-he\-na}) \\
\textsc{Derivations}: \ding{43}
\on{1}.~\HL{syusmung}{\B{syu\-smung}}, tray. 
\on{2}.~\HL{paysmung}{\B{pay\-smung}}, water carrier. 
\on{3}.~\HL{prrsmung}{\B{prr\-smung}}, baby carrier.

% sämyam
\hi
\HT{sAmyam}{\B{säm\cp{yam}}} \I{[s\ae{}m."jam]} \emph{n.} hug, embrace.   
\B{Säm\-yam\-ìl po\-ru wa\-yìn\-txu fu\-ta nga\-ta lo\-lu li txo\-a.} A hug will show him that you've already forgiven him.~\LINK{http://naviteri.org/2011/10/more-vocabulary-a-bit-of-grammar/}{b}
(\ding{43}~\HL{meyam}{\B{me\-yam}})

% sänrr
\hi
\HT{sAnrr}{\B{sä\cp{nrr}}} \I{[s\ae."n\s{r}]} \emph{or coll.} \I{[sn\s{r}]} \emph{n.} (\emph{a.}) glow, an instance of glowing.   
\B{Txep\-ìl tìng le\-fpom\-a sä\-nrr\-ti.} The fire gives a pleasant glow.~\LINK{http://naviteri.org/2012/07/fivospxiya-aylifyavi-amip-this-months-new-expressions-pt-2/}{b} 
(\emph{b.}) light source; lamp.
(\ding{43}~\HL{nrr}{\B{nrr}}) \\
\textsc{Derivations}: \ding{43}
\on{1}.~\HL{tsmisnrr}{\B{tsmi\-snrr}}, bladder lantern, nectar lantern. 
\on{2}.~\HL{tawsnrrtsyIp}{\B{Taw\-snrr\-tsyìp}}, Alpha Centauri C.

% sänui
\hi
\HT{sAnui}{\B{sä\cp{nu}i}} \I{[s\ae."nu.i]} \emph{or coll.} \I{["{}snu.i]} \emph{n.} failure (particular instance of failure).
\B{Oeyk tsa\-tì\-snaytx\-ä lu apxa sä\-nui tì\-eyk\-tan\-ä.} That loss was caused by a great failure of leadership.~\LINK{http://naviteri.org/2013/04/wheres-the-bathroom-and-other-useful-things/}{b}
\B{Fì\-sä\-nu\-it ze\-ne nga tswi\-va', ma 'i\-tan.} You must forget about this failure, my son.~\LINK{http://naviteri.org/2013/04/wheres-the-bathroom-and-other-useful-things/}{b}
(\ding{43}~\HL{nui}{\B{nui}}, \HL{tInui}{\B{tì\-nui}})

% sänume / sänume si
\hi
\HT{sAnume}{\B{sä\cp{nu}me}} \I{[s\ae."nu.mE]} \emph{or coll.} \I{["{}snu.mE]} \textsc{i}. \emph{n.} teaching, instruction. 
\B{Ye'\-rìn 'ì\-yi\-'a sä\-nu\-me a tsa\-ri kll\-fro' oe.} My teaching responsibilities will soon end.~\LINK{http://forum.learnnavi.org/language-updates/trr-rrtaya/}{ln}
\B{Tì\-flä la\-tem ìlä seyng\-a ftxey fkol sä\-nu\-met li\-vek fu\-ke.} Success depends on whether or not one follows instructions.~\LINK{http://naviteri.org/2012/07/meetings-waterfalls-and-more/}{b}
\B{Nge\-yä te\-ri fay\-te\-le a ay\-sä\-nu\-me\-ri ngar ira\-yo sei\-yi ay\-oe nì\-wotx.} We all thank you for your teachings about these topics.~\LINK{http://forum.learnnavi.org/language-updates/combined-answers-1-\%28feb-16\%29/}{ln} \\
\textsc{ii}. \B{sä\cp{nu}me si} \I{[s\ae{}."nu.mE si]} \emph{or coll.} \I{["{}snu.mE si]} \emph{vin.} teach.
\B{Ma 'ite, aw\-nge\-yä fya\-'o\-ri ze\-ne nga sä\-nu\-me si\-vi po\-ru fte tsi\-vun pil\-vll\-txe sì ti\-vì\-ran nì\-ay\-oeng.} My daughter, you will teach him our way to speak and walk as we do.~\LINK{http://wiki.learnnavi.org/index.php?title=Na\%27vi_from_Avatar_Movie\#Scene_in_Hometree}{a\slash wiki}
(\ding{43}~\HL{nume}{\B{nu\-me}}) \\
\textsc{Derivations}: \ding{43}
\on{1}.~\HL{sAnumvi}{\B{sä\-num\-vi}}, lesson.

% sänumvi
\hi
\HT{sAnumvi}{\B{sä\cp{num}vi}} \I{[s\ae."num.vi]} \emph{or} \I{[."nUm.]} \emph{or coll.} \I{["{}snum.vi]} \emph{n.} lesson.
\B{Ta\-fral kì\-syar a ay\-sä\-num\-vi sì ay\-lì'\-fya\-vi lu fyin.} That's why I'll teach lessons and expressions that are simple.~\LINK{http://naviteri.org/2012/07/tskxekengtsyip-a-mikyunfpi-a-little-listening-exercise/}{b}
\B{Maw\-krr la\-ye\-iu oer krr nì\-'ul fte ngi\-vop ay\-lì\-'ut sì tsay\-fne\-sä\-num\-vit a tsun fra\-por srung si\-vi fte ni\-vu\-me sì zi\-ve\-rok nì\-swey.} After that I will have more time to create words and the kinds of lessons that can help everyone best learn and remember.~\LINK{http://forum.learnnavi.org/language-updates/trr-rrtaya/}{ln}
(\ding{43}~\HL{sAnume}{\B{sä\-nu\-me}})

% sängä'än
\hi
\HT{sAngA'An}{\B{sängä\cp{'än}}} \I{[s\ae.N\ae."P\ae{}n]} \emph{or coll.} \I{[sN\ae."P\ae{}n]} \emph{n.} bout of suffering, episode of depression.
\B{Nge\-yä tsa\-sä\-ngä\-'än\-tsyìp\-ì\-ri set fra\-wzo srak?} Have you recovered from being down for a while?~\LINK{http://naviteri.org/2013/02/vospxi-ayol-posti-apup-short-post-for-a-short-month/}{b}
\B{Pe\-yä sä\-ngä\-'än\-ìri lu sä\-fe\-'ul le\-yew\-la nì\-txan.} The worsening of his depression is very disappointing.~\LINK{http://naviteri.org/2013/03/tsanerul-feerul-getting-better-getting-worse/}{b}
(\ding{43}~\HL{ngA'An}{\B{ngä\-'än}})

% sängop
\hi
\HT{sAngop}{\B{sä\cp{ngop}}} \I{[s\ae."Nop\textcorner]} \emph{or coll.} \I{[sngop\textcorner]} \emph{n.} creation.
\B{Fyo\-lup\-a ay\-sä\-ngop ay\-nge\-yä oe\-ru te\-ya si nì\-ngay.} Your exquisite creations fill me greatly with joy.~\LINK{http://naviteri.org/2014/12/ayngaru-tsinga-yoratu-may-i-present-the-four-winners/}{b}
(\ding{43}~\HL{ngop}{\B{ngop}}, \HL{tIngop}{\B{tì\-ngop}})

% säomum
\hi
\HT{sAomum}{\B{sä\cp{o}mum}} \I{[s\ae."{}o.mum]} \emph{or} \I{[.mUm]} (\textsc{rn}: \B{sä\cp{o}mùm} \I{[sE."{}o.mUm]}) \emph{countable n.} (piece of) information.   
\B{Fì\-sä\-omum txan\-krr lo\-lu oer.} I had this information for a long time.~\LINK{http://naviteri.org/2013/08/fmawn-a-fkol-kilmulat-this-just-in/}{b}
\B{Ay\-sä\-omum\-ìri le\-sar sei\-yi oe nga\-ru ira\-yo nì\-txan.} Thanks very much for (all) this useful information.~\LINK{http://naviteri.org/2011/09/miscellaneous-vocabulary/}{b} 
(\ding{43}~\HL{omum}{\B{o\-mum}})

% säpenghrr
\hi 
\HT{sApenghrr}{\B{säpeng\cp{hrr}}} \I{[s\ae.pEN."h\s{r}]} \emph{or coll.} \I{[spEN."h\s{r}]} \emph{n.} warning. 
\B{Som\-wew\-ä tì\-'ul a ka 'Rr\-ta sä\-peng\-hrr lu aw\-nga\-ru nì\-wotx.} Global warming is a warning to us all.~\LINK{http://naviteri.org/2016/06/mrrvola-lifyavi-amip-forty-new-expressions/}{b} 
(\ding{43}~\HL{penghrr}{\B{peng\-hrr}}, \HL{sApllhrr}{\B{sä\-pll\-hrr}})

% säpllhrr
\hi 
\HT{sApllhrr}{\B{säpll\cp{hrr}}} \I{[s\ae.p\s{l}."h\s{r}]} \emph{or coll.} \I{[sp\s{l}."h\s{r}]} \emph{n.} warning. 
(\ding{43}~\HL{pllhrr}{\B{pll\-hrr}}; \HL{sApenghrr}{\B{sä\-peng\-hrr}})

% säplltxe
\hi
\HT{sAplltxe}{\B{säpll\cp{txe}}} \I{[s\ae.p\s{l}."t'E]} \emph{or coll.} \I{[sp\s{l}."t'E]} \emph{n.} statement.
\B{'Aw\-ve\-a sä\-pll\-txe nge\-yä lä\-ngu e\-yawr.} Your first statement is unfortunately correct.~\LINK{http://naviteri.org/2014/02/barter-and-exchange/comment-page-1/\#comment-2614}{b}
\B{Sä\-pll\-txel kar\-yu\-ä ke to\-la\-rä\-nge tì\-pawm\-it kaw\-'it.} The teacher's statement in no way pertained to the question. Drat!~\LINK{http://naviteri.org/2012/06/spring-vocabulary-part-3/}{b}
(\ding{43}~\HL{plltxe}{\B{pll\-txe}}) \\
\textsc{Derivations}: \ding{43}
\on{1}.~\HL{sAplltxevi}{\B{sä\-pll\-txe\-vi}}, comment.

% säplltxevi
\hi
\HT{sAplltxevi}{\B{säpll\cp{txe}vi}} \I{[s\ae.p\s{l}."t'E.vi]} \emph{or coll.} \I{[sp\s{l}."t'E.vi]} \emph{n.} comment.
\B{Tsa\-lì\-'u\-ri alu yo\-ra'\-tu, nge\-yä sä\-pll\-txe\-vi lu txan\-tsan.} Your comment concerning the word \emph{yora'tu} is excellent.~\LINK{http://naviteri.org/2011/12/ulte-yora\%E2\%80\%99tu-leiu-and-the-winner-is/comment-page-1/\#comment-1267}{b}
\B{Ay\-nge\-yä ay\-sä\-pll\-txe\-vi\-ri sei\-yi oe ira\-yo nì\-txan ay\-nga\-ru nì\-wotx!} I thank you all very much for the comments!~\LINK{http://naviteri.org/2011/07/attributive-\%E2\%80\%9Ca\%E2\%80\%9D-and-truncated-style/}{b}
(\emph{a. fig.}) \B{mek\-a sä\-pll\-txe\-vi}, an insipid\slash thoughtless comment.~\LINK{http://naviteri.org/2012/02/trr-asawnung-lefpom-happy-leap-day/}{b}
(\ding{43}~\HL{sAplltxe}{\B{sä\-pll\-txe}})

% säpom
\hi
\HT{sApom}{\B{sä\cp{pom}}} \I{[s\ae."pom]} \emph{or coll.} \I{[spom]} \emph{n.} kiss.
(\ding{43}~\HL{pom}{\B{pom}})

% säpxor
\hi 
\HT{sApxor}{\B{sä\cp{pxor}}} \I{[s\ae."p'oR]} \emph{or coll.} \I{[sp'oR]}  \emph{n.} explosion. 
\B{Stì\-mawm sä\-pxor\-it! Tul!} (I) just heard an explosion! Run!~\LINK{https://naviteri.org/2024/09/pxevola-liu-amip-two-dozen-new-words/}{b} 
(\ding{43}~\HL{pxor}{\B{pxor}}) 

% särangal
\hi 
\HT{sArangal}{\B{sä\cp{ra}ngal}} \I{[s\ae."Ra.Nal]} \emph{or coll.} \I{["{}sRa.Nal]} \emph{n.} wish. 
\B{Fay\-sä\-ra\-ngal\-ìri lu oe nga\-hu zam mì zam.} Concerning these wishes I'm a hundred percent with you.~\LINK{http://naviteri.org/2021/12/zolau-niprrte-ma-3746-welcome-2022/\#comment-39760}{b}
(\ding{43}~\HL{rangal}{\B{ra\-ngal}}) 

% särawn
\hi
\HT{sArawn}{\B{sä\cp{rawn}}} \I{[s\ae."Rawn]} \emph{or coll.} \I{[sRawn]} \emph{n.} replacement, substitute, something that replaces something else.
\B{Fì\-pam\-tseo\-tu\-ri ke la\-yu ftue fwa run fkol sä\-rawn\-it a tam.} It won't be easy to find a satisfactory replacement for this musician.~\LINK{http://naviteri.org/2012/01/mipa-zisit-ayliu-amip-new-words-for-the-new-year/}{b}
(\ding{43}~\HL{rawn}{\B{rawn}})

% särengop
\hi
\HT{sArengop}{\B{sä\cp{re}ngop}} \I{[s\ae."RE.Nop\textcorner]} \emph{or coll.} \I{["{}sRE.Nop\textcorner]} \emph{n.} design, a particular instance of designing.
\B{Fay\-sä\-re\-ngop\-it a\-vä' oe\-ru rä\-'ä wìn\-txu nì\-mun, ru\-txe. Ke su\-nu oer keng nì\-'it.} Please don't show me these ugly designs again. I don't like them one bit.~\LINK{http://naviteri.org/2013/01/awvea-posti-zisita-amip-first-post-of-the-new-year/}{b}
(\ding{43}~\HL{rengop}{\B{re\-ngop}}, \HL{tIrengop}{\B{tì\-re\-ngop}})

% säro'a / säro'a si
\hi
\HT{sAro'a}{\B{sä\cp{ro}'a}} \I{[s\ae."Ro.Pa]} \emph{or coll.} \I{["{}sRo.Pa]} \textsc{i}. \emph{n.} feat, accomplishment, great deed. \\
\textsc{ii}. \B{sä\cp{ro}'a si} \I{[s\ae."Ro.Pa si]} \emph{or coll.} \I{["{}sRo.Pa si]} \emph{vin.} do great deeds.  
\B{Txan\-tstew sä\-ro\-'a si, fna\-we'\-tu ke si.} A hero does great deeds, a coward does not.~\LINK{http://naviteri.org/2012/03/spring-vocabulary-part-1/}{b}
\B{Txo ni\-vew fko sä\-ro\-'a si\-vi, ze\-ne nì\-'aw\-ve ve\-nga' li\-vu.} If you want to accomplish great things, you first have to be organized.~\LINK{http://naviteri.org/2012/10/mipa-vospxi-mipa-ayliu-new-words-for-the-new-month/}{b}
(\ding{43}~\HL{ro'a}{\B{ro\-'a}})

% säseyto
\hi
\HT{sAseyto}{\B{säsey\cp{to}}} \I{[s\ae.sEj."to]} \emph{n.} butchering tool.
(\ding{43}~\HL{seyto}{\B{sey\-to}})

% säsìlpey
\hi
\HT{sAsIlpey}{\B{säsìl\cp{pey}}} \I{[s\ae.sIl."pEj]} \emph{n.} hope, a particular instance of hoping.
\B{Krr\-a wä\-tut tse\-'a, pe\-yä sä\-sìl\-pey a yo\-ra' 'o\-lìp.} When he saw his opponent, his hope of winning vanished.~\LINK{http://naviteri.org/2013/01/awvea-posti-zisita-amip-first-post-of-the-new-year/}{b}
(\ding{43}~\HL{sIlpey}{\B{sìl\-pey}}, \HL{tsIlpey}{\B{tsìl\-pey}})

% säspaw
\hi 
\HT{sAspaw}{\B{sä\cp{spaw}}} \I{[s\ae."{}spaw]} \emph{n.} belief (\emph{particular instance}). 
\B{Tsa\-sä\-spaw a\-tsleng lu le\-hrr\-ap.} That false belief is dangerous.~\LINK{http://naviteri.org/2023/12/mipa-zisit-ayliu-amip-new-year-new-words/}{b} 
(\ding{43}~\HL{spaw}{\B{spaw}}; \HL{tIspaw}{\B{tì\-spaw}}) 

% säspxin
\hi
\HT{sAspxin}{\B{sä\cp{spxin}}} \I{[s\ae."{}sp'in]} \emph{n.} disease, sickness.
\B{Fu\-ria ke tsun tì\-kang\-kem si\-vi, pe\-yä sä\-spxin lu tì\-la\-'um nì\-'aw.} As for not being able to work, his illness is only a pretence.~\LINK{http://naviteri.org/2020/02/some-words-for-leap-year-day/}{b} 
\B{Fì\-sä\-spxin\-ìri nge\-yä ke lä\-ngu kea 'um\-tsa.} Unfortunately there is no medicine for this disease of yours.~\LINK{http://naviteri.org/2014/11/twenty-before-the-holidays/}{b} 
(\ding{43}~\HL{spxin}{\B{spxin}}) \\
\textsc{Derivations}: \ding{43}
\on{1}.~\HL{sAspxintsyIp}{\B{sä\-spxin\-tsyìp}}, minor ailment.

% säspxintsyìp
\hi
\HT{sAspxintsyIp}{\B{sä\cp{spxin}tsyìp}} \I{[s\ae."{}sp'in.\t{ts}jIp\textcorner]} \emph{n.} minor ailment. 
(\ding{43}~\HL{sAspxin}{\B{sä\-spxin}})

% säsrätx
\hi
\HT{sAsrAtx}{\B{sä\cp{srätx}}} \I{[s\ae."{}sR\ae{}t']} \emph{n.} annoyance.
\B{Nga fwe\-fwi nìl\-ke\-ftang a fì\-'u lu sä\-srätx a\-txan oe\-ru.} It's a major annoyance to me that you whistle continuously.~\LINK{http://naviteri.org/2011/09/miscellaneous-vocabulary/}{b}
(\emph{a. expression}) \B{Sra\-ke srätx (ngat) txo \ldots?} (\emph{with} ‹iv›) `Would you mind if \ldots?'
\B{Kea sä\-srätx kaw\-'it!} `Not at all!'~\LINK{http://naviteri.org/2020/11/vospxivopeya-ayliu-amip-novembers-new-words/}{b}
(\ding{43}~\HL{srAtx}{\B{srätx}})

% säsrese'a
\hi 
\HT{sAsrese'a}{\B{säsrese\cp{'a}}} \I{[s\ae.sRE.sE."Pa]} \emph{n.} prediction. 
\B{Ay\-nge\-yä sä\-sre\-se\-'a lo\-lu e\-yawr!} Your prediction has been correct.~\LINK{http://naviteri.org/2019/08/aawa-lifya-si-lifyavi-amip-a-few-new-words-and-expressions/\#comment-30279}{b} 
(\ding{43}~\HL{srese'a}{\B{sre\-se'a}})

% säsrìn
\hi
\HT{sAsrIn}{\B{sä\cp{srìn}}} \I{[s\ae."{}sRIn]} \emph{n.} lent or borrowed thing.
\B{Oe\-ta a tsa\-sä\-srìn\-ìl tok pe\-seng\-it?} Where's the thing (you) borrowed from me?~\LINK{http://naviteri.org/2012/01/more-additions-to-the-lexicon/}{b}
(\ding{43}~\HL{srIn}{\B{srìn}})

% säsrung
\hi 
\HT{sAsrung}{\B{sä\cp{srung}}} \I{[s\ae."{}sRuN]} \emph{or} \I{[."{}sRUN]} \emph{n.} helper (\emph{inanimate}), s.t. that helps. 
\B{Oe\-yä tì\-tslam\-ìri tì\-oeyk\-tìng nge\-yä lo\-lu sä\-srung.} Your explanation helped my understanding.~\LINK{https://naviteri.org/2024/03/fleltrra-ayliu-words-for-april-fools-day/}{b}
(\ding{43}~\HL{srung}{\B{srung}}, \HL{srungsiyu}{\B{srung\-si\-yu}}) \\
\textsc{Derivations}: \ding{43} 
\on{1}.~\HL{txawnulsrung}{\B{txaw\-nul\-srung}}, machine.

% sästarsìm
\hi
\HT{sAstarsIm}{\B{sä\cp{star}sìm}} \I{[s\ae."{}staR.sIm]} \emph{n.} collection (put together intentionally by a person). 
(i.) \B{sä\-star\-sìm syu\-lang\-ä}, bouquet (\emph{collection of flowers selected and put together intentionally by a person}).~\LINK{http://naviteri.org/2012/03/spring-vocabulary-part-2/}{b}
(\ding{43}~\HL{starsIm}{\B{star\-sìm}})

% säsulìn
\hi
\HT{sAsulIn}{\B{sä\cp{su}lìn}} \I{[s\ae."{}su.lIn]} \emph{n.} hobby.   
\B{Oe tskxe\-keng si sä\-su\-lìn\-ur alu tsko swi\-zaw.} I practice my hobby, which is archery.~\LINK{http://naviteri.org/2010/09/getting-to-know-you-part-3/}{b}
(\ding{43}~\HL{sulIn}{\B{su\-lìn}})

% säsyep
\hi
\HT{sAsyep}{\B{sä\cp{syep}}} \I{[s\ae."{}sjEp\textcorner]} \emph{n.} trap.
(\ding{43}~\HL{syep}{\B{syep}})

% sätare
\hi
\HT{sAtare}{\B{sä\cp{ta}re}} \I{[s\ae."ta.RE]} \emph{or coll.} \I{["{}sta.RE]} \emph{n.} connection, relationship.   
\B{Nì\-ngay le\-iu oer sì sem\-pul\-ur sä\-ta\-re a\-sìl\-tsan.} Father and I really have a good relationship; it's nice.~\LINK{http://naviteri.org/2012/06/spring-vocabulary-part-3/}{b}
(\ding{43}~\HL{tare}{\B{ta\-re}}) \\
\textsc{Derivations}: \ding{43}
\on{1}.~\HL{starlI'u}{\B{star\-lì\-'u}}, adposition.
\on{2}.~\HL{sAtarenga'}{\B{sä\-ta\-re\-nga'}}, relevant, pertinent.

% sätarenga'
\hi 
\HT{sAtarenga'}{\B{sä\cp{ta}renga'}} \I{[s\ae."ta.RE.NaP]} \emph{or coll.} \I{["{}sta.RE.NaP]} \emph{adj.} relevant, pertinent. 
\B{Tsa\-sä\-pll\-txe\-vi\-ri a\-sä\-ta\-re\-nga' ira\-yo.} Thanks for that pertinent comment.~\LINK{http://naviteri.org/2021/03/mipa-ayliu-si-aylifyavi-new-words-and-expressions/}{b} 
(\ding{43}~\HL{sAtare}{\B{sä\-ta\-re}}, \HL{-nga'}{\B{-nga'}}; \ding{214}~\HL{kesAtarenga'}{\B{ke\-sä\-ta\-re\-nga'}})

% sätarep
\hi 
\HT{sAtarep}{\B{sä\cp{ta}rep}} \I{[s\ae."ta.REp\textcorner]} \emph{or coll.} \I{["{}sta.REp\textcorner]} \emph{n.} a rescue, an incident of rescuing.  
(\ding{43}~\HL{tarep}{\B{ta\-rep}})

% sätaron
\hi
\HT{sAtaron}{\B{sä\cp{ta}ron}} \I{[s\ae."ta.Ron]} \emph{or coll.} \I{["{}sta.Ron]} \emph{n.} hunt.   
\B{Eyk Ka\-mun a fra\-lo lä\-ngu tsa\-sä\-ta\-ron vel\-ke nì\-wotx. Ta\-ron\-yut yom smar\-ìl!} Every time Kamun is in charge, the hunt is a mess. Everything goes wrong that can!~\LINK{http://naviteri.org/2012/10/mipa-vospxi-mipa-ayliu-new-words-for-the-new-month/}{b}
\B{Ta\-ron\-yul le\-hawm\-pam ska\-'a sä\-ta\-ron\-it.} A noisy hunter destroys the hunt.~\LINK{http://naviteri.org/2014/11/twenty-before-the-holidays/}{b}
(\ding{43}~\HL{taron}{\B{ta\-ron}})

% sätsan'ul
\hi
\HT{sAtsan'ul}{\B{sä\cp{tsan}'ul}} \I{[s\ae."\t{ts}an.Pul]} \emph{or} \I{[.PUl]} \emph{n.} improvement, a specific instance of improving.
\B{Set pll\-txe nga nìl\-tsan nì\-ngay. Fì\-sä\-tsan\-'ul\-ìri nge\-yä lu oe\-ru sko har\-yu nrr\-a.} You speak really well now. As your teacher, I'm proud of your improvement.~\LINK{http://naviteri.org/2013/03/tsanerul-feerul-getting-better-getting-worse/}{b}
(\ding{43}~\HL{tsan'ul}{\B{tsan\-'ul}}, \HL{tItsan'ul}{\B{tì\-tsan\-'ul}})

% sätsawn
\hi 
\HT{sAtsawn}{\B{sä\cp{tsawn}}} \I{[s\ae."\t{ts}awn]} \emph{n.} harvest (\emph{particular instance}). 
\B{Fì\-zì\-sìt\-ä sä\-tsawn txan\-tsan lei\-u.} I'm pleased to say that this year's harvest was excellent.~\LINK{http://naviteri.org/2022/11/krr-ao-an-exciting-time/}{b} 
(\ding{43}~\HL{tsawn}{\B{tsawn}}) 

% sätseri
\hi 
\HT{sAtseri}{\B{sä\cp{tse}ri}} \I{[s\ae."\t{ts}E.Ri]} \emph{n.} observation, something noticed. 
\B{Sìl\-tsan\-a sä\-tse\-ri, sìl\-tsan\-a tì\-pawm.} Good observation, good question.~\LINK{http://naviteri.org/2020/08/hiia-vur-a-teri-mefalukantsyip-a-little-story-about-two-cats/\#comment-31473}{b} 
\B{Nge\-yä tsa\-sä\-tse\-ri\-ri a el\-tur tì\-txen si ira\-yo.} Thank you for that interesting observation of yours.~\LINK{http://naviteri.org/2020/11/vospxivopeya-ayliu-amip-novembers-new-words/}{b} 
(\ding{43}~\HL{tseri}{\B{tse\-ri}})

% sätsìsyì
\hi
\HT{sAtsIsyI}{\B{sä\cp{tsì}syì}} \I{[s\ae."\t{ts}I.sjI]} \emph{n.} a whisper.
\B{Po pa\-mll\-txe a krr, fra\-po tar\-mìng mik\-yun nì\-pxi, ta\-lu\-na mok\-ri lu sä\-tsì\-syì\-tsyìp.} When he spoke, everyone listened intently, because his voice was a tiny whisper.~\LINK{http://naviteri.org/2011/09/miscellaneous-vocabulary/}{b}
(\ding{43}~\HL{tsIsyI}{\B{tsì\-syì}})

% sätswayon
\hi
\HT{sAtswayon}{\B{sä\cp{tsway}on}} \I{[s\ae."\t{ts}waj.on]} \emph{n.} flight, an instance of flying.
\B{Sä\-tsway\-on a\-'aw\-ve tsa\-heyl\-ur ha\-sey si; ke tsun nga pi\-vey.} First flight seals the bond; you cannot wait.~\LINK{http://naviteri.org/2012/06/spring-vocabulary-part-3/}{b}
\B{Maw sä\-tsway\-on a\-yol ay\-oe kll\-po\-lä mì ta\-yo a lu ro\-fa kil\-van.} After a short flight we landed in a field beside the river.~\LINK{http://naviteri.org/2012/01/mipa-zisit-ayliu-amip-new-words-for-the-new-year/}{b}
(\ding{43}~\HL{tswayon}{\B{tsway\-on}})

% sätsyìl
\hi
\HT{sAtsyIl}{\B{sä\cp{tsyìl}}} \I{[s\ae."\t{ts}jIl]} \emph{n.} climbing event, a climb.   
\B{Kin\-trr\-am\-ä sä\-tsyìl lu le\-hrr\-ap slä 'o' nì\-txan.} Last week's climb was dangerous but very exciting.~\LINK{http://naviteri.org/2012/01/more-additions-to-the-lexicon/}{b} 
\B{Tsa\-sä\-tsyìl lo\-lu ngä\-zìk nì\-ngay!} That was a really hard climb!~\LINK{http://naviteri.org/2012/02/trr-asawnung-lefpom-happy-leap-day/}{b}
(\ding{43}~\HL{tsyIl}{\B{tsyìl}})

% sävll
\hi 
\HT{sAvll}{\B{sä\cp{vll}}} \I{[s\ae."v\s{l}]} \emph{n.} sign, signal, indication. 
\B{Kxe\-ner lu sä\-vll txep\-ä.} Smoke is a sign\slash an indication that there is fire.~\LINK{http://naviteri.org/2017/12/zisit-amip-lefpom-ma-eylan-happy-new-year-friends/}{b}
\B{Mì sa\-ngek a sä\-vll\-it ngo\-lop eyk\-yul tar\-po\-ngu\-ä.} The sign on the tree trunk was made by the leader of the hunting party.~\LINK{http://naviteri.org/2017/12/zisit-amip-lefpom-ma-eylan-happy-new-year-friends/}{b}
(\ding{43}~\HL{vll}{\B{vll}})

% säwäsul
\hi
\HT{sAwAsul}{\B{sä\cp{wä}sul}} \I{[s\ae."w\ae.sul]} \emph{or} \I{[.sUl]} \emph{or coll.} \I{["{}sw\ae.sul]} \emph{n.} a competition, a particular instance of competing.
(i.) \B{sä\-wä\-sul a tul}, foot race.~\LINK{http://naviteri.org/2012/10/wina-sawasultsyip-ahii-a-quick-little-contest/}{b}
(\ding{43}~\B{wä\-sul}) \\
\textsc{Derivations}: \ding{43}
\on{1}.~\HL{sAwAsultsyIp}{\B{sä\-wä\-sul\-tsyìp}}, contest.

% säwäsultsyìp
\hi
\HT{sAwAsultsyIp}{\B{sä\cp{wä}sultsyìp}} \I{[s\ae."w\ae.sul.\t{ts}jIp\textcorner]} \emph{or} \I{[.sUl.]} \emph{or coll.} \I{["{}sw\ae.sul.\t{ts}jIp\textcorner]} \emph{n.} contest.
\B{Sìl\-pey oe, fì\-sä\-wä\-sul\-tsyìp 'o' lì\-ye\-vu ay\-nga\-ru!} I hope, this little contest will be fun for you!~\LINK{http://naviteri.org/2012/10/wina-sawasultsyip-ahii-a-quick-little-contest/}{b}
\B{Win\-a Sä\-wä\-sul\-tsyìp A\-hì\-'i.} A Quick Little Contest.~\LINK{http://naviteri.org/2012/10/wina-sawasultsyip-ahii-a-quick-little-contest/}{b}
(\ding{43}~\HL{sAwAsul}{\B{sä\-wä\-sul}})

% säwäte
\hi
\HT{sAwAte}{\B{säwä\cp{te}}} \I{[s\ae.w\ae."tE]} \emph{or coll.} \I{[sw\ae."tE]} \emph{n.} point of contention, source of argument, thing disputed. 
(\ding{43}~\HL{wAte}{\B{wä\-te}}, \HL{tIwAte}{\B{tì\-wä\-te}})

% säwem
\hi 
\HT{sAwem}{\B{sä\cp{wem}}} \I{[s\ae."wEm]} \emph{or coll.} \I{[swEm]} \emph{n.} fight. 
\B{Fì\-sä\-wem\-ìri ze\-ne aw\-nga kawl hä\-pi\-vawl. Lu Saw\-tu\-te wä\-tu a\-pxul.} We must prepare diligently for this fight. The Skypeople are formidable opponents.~\LINK{http://naviteri.org/2020/08/hiia-vur-a-teri-mefalukantsyip-a-little-story-about-two-cats/}{b}
(\ding{43}~\HL{wem}{\B{wem}})

% säwìntxu
\hi
\HT{sAwIntxu}{\B{säwìn\cp{txu}}} \I{[s\ae.wIn."t'u]} \emph{or coll.} \I{[swIn."t'u]} \emph{n.} a showing, an exhibition.
(\ding{43}~\HL{wIntxu}{\B{wìn\-txu}})

% säwìngay
\hi
\HT{sAwIngay}{\B{säwì\cp{ngay}}} \I{[s\ae.wI."Naj]} \emph{or coll.} \I{[swI."Naj]} \emph{n.} proof (\emph{particular instance}).
\B{Tsa\-sä\-pll\-txe\-ri sä\-wì\-ngay a to\-lìng ngal lu meyp.} The proof you gave of that statement is weak.~\LINK{http://naviteri.org/2015/04/some-new-words-for-may-day/}{b}
(\ding{43}~\HL{wIngay}{\B{wì\-ngay}})

% säyäkx
\hi
\HT{sAyAkx}{\B{sä\cp{yäkx}}} \I{[s\ae."j\ae{}k']} \emph{or coll.} \I{[sj\ae.k']} \emph{n.} snub.
\B{Fì\-sä\-yäkx\-it ay\-oel ke tswa\-ya'.} We will not forget this snub.~\LINK{http://naviteri.org/2015/03/kap-si-ayunil-saylahe-threats-dreams-and-other-things/}{b}
(\ding{43}~\HL{yAkx}{\B{yäkx}})

% säyìm
\hi 
\HT{sAyIm}{\B{sä\cp{yìm}}} \I{[s\ae."jIm]} \emph{or coll.} \I{[sjIm]} \emph{n.} tie, something used for binding. 
(\ding{43}~\HL{yIm}{\B{yìm}}) \\
\textsc{Derivations}: \ding{43} 
\on{1}.~\HL{syImawnun'i}{\B{syì\-maw\-nun\-'i}}, cut tie. 

% säzärìp
\hi
\HT{sAzArIp}{\B{sä\cp{zä}rìp}} \I{[s\ae."z\ae.RIp\textcorner]} \emph{n.} direhorse lead, rein.
(\ding{43}~\HL{za'ArIp}{\B{za\-'ä\-rìp}}) 

% se'ayl
\hi
\HT{se'ayl}{\B{se\cp{'ayl}}} \I{[sE."Pajl]} \emph{countable n.} kind of waterfall, \emph{an individual tall, thin waterfall that pours down a sheer high cliff or off of a floating mountain}.

% sekrr
\hi
\HT{sekrr}{\B{\cp{se}krr}} \I{["{}sE.k\s{r}]} \emph{n.} present (\emph{temporal}). 
(\ding{43}~\HL{set}{\B{set}}, \HL{krr}{\B{krr}})

% semkä
\hi 
\HT{semkA}{\B{sem\cp{kä}}} \I{[sEm."k\ae]} \emph{n.} bridge.  
(\ding{43}~\HL{sA'o}{\B{sä\-'o}}, \HL{emkA}{\B{em\-kä}}; \HL{emkAfya}{\B{em\-kä\-fya}})

% sempu
\hi
\HT{sempu}{\B{\cp{sem}pu}} \I{["{}sEm.pu]} \emph{n.} daddy. 
\B{Ma sem\-pu, oey yom\-yo tso\-lawng!} Daddy, my plate broke!~\LINK{http://naviteri.org/2020/04/aawa-u-amip-a-few-new-things/}{b}
\B{Sem\-pu\-ti rä\-'ä srätx, ma yawn\-tu\-tsyìp. Tì\-kang\-kem se\-ri.} Don't bother daddy, little one. He's working.~\LINK{http://naviteri.org/2017/09/zisikrr-amip-ayliu-amip-new-words-for-the-new-season/}{b}
(\ding{43}~\HL{sempul}{\B{sem\-pul}})

% sempul
\hi
\HT{sempul}{\B{\cp{sem}pul}} \I{["{}sEm.pul]} \emph{or} \I{[.pUl]} \emph{n.} father. 
\B{Ma Sem\-pul, pxay\-a a\-to\-ki\-ri\-na' ne fì\-ke\-tu\-wong zì\-ma\-'u. Ze\-ne Sa'\-nok fì\-au\-ngia\-ti ral\-pi\-veng.} Father, many \emph{atokirina} just came to this alien. Mother must interpret this sign.~\LINK{http://forum.learnnavi.org/language-updates/pseudo-canon-from-bd-live-content/}{a\textsuperscript{c}\slash ln}
\B{Tì\-re\-ngop\-ìri ioi\-yä lu sem\-pul pe\-yä tsul\-fä\-tu.} Her father is a master designer of ceremonial adornments.~\LINK{http://naviteri.org/2013/01/awvea-posti-zisita-amip-first-post-of-the-new-year/}{b}
\B{Fì\-tsko\-ti a\-fyo\-lup 'o\-ngo\-lop oe\-yä sem\-pul\-ìl.} This exquisite bow was designed by my father.~\LINK{http://naviteri.org/2014/07/tsawlultxamaw-a-ayliu-post-meetup-words/}{b}
\B{Nì\-ngay le\-iu oer sì sem\-pul\-ur sä\-ta\-re a\-sìl\-tsan.} Father and I really have a good relationship; it's nice.~\LINK{http://naviteri.org/2012/06/spring-vocabulary-part-3/}{b}
\B{Txe\-wì\-yä sem\-pul\-ìri sngä\-'i\-krr Saw\-tu\-tet ve\-re'\-kì.} In the beginning Txewì's father hated the Sky People.~\LINK{http://naviteri.org/2010/07/ting-mikyun-fte-tslivam-listening-comprehension-1-3/}{b}
\B{Tsa\-ri\-a fkol po\-le\-'un fu\-ta Lo\-ak slu ta\-ron\-yu, sem\-pul 'efu ye.} Father is content that it's been decided Loak will be a hunter.~\LINK{http://naviteri.org/2011/04/\%E2\%80\%99a\%E2\%80\%99awa-li\%E2\%80\%99fyavi-amip\%E2\%80\%94a-few-new-expressions/}{b}
\B{Na\-ri si!  Tsa\-tsko lu spu\-win ul\-te ke lu mi txur, slä oe\-ri txan\-wa\-we le\-iu. Lu stxe\-li a sem\-pul\-ta.} Careful! That bow is old and no longer strong, but it means a lot to me. It was a gift from my father.~\LINK{http://naviteri.org/2011/05/some-miscellaneous-vocabulary/}{b}
(i.) \B{Sem\-pul Lì'\-fya\-yä}, father of the language, Dr.~Paul Frommer.~\LINK{http://forum.learnnavi.org/language-updates/a-long-answer-to-a-quick-question/msg155104/\#msg155104}{b} \\
\textsc{Derivations}: \ding{43}
\on{1}.~\HL{sempul}{\B{sem\-pu}}, daddy. 
\on{2}.~\HL{sa'sem}{\B{sa'\-sem}}, parent, parents.

% set
\hi
\HT{set}{\B{set}} \I{[sEt\textcorner]} \emph{adv.} now.   
\B{'Aw\-pot set ftxey ay\-ngal a lu ay\-nga\-kip.} Choose one among you now.~\LINK{http://naviteri.org/2013/08/taronway-the-hunt-song/}{b}
\B{'Aw\-a tì\-pawm\-ì\-ri 'i\-veyng oe set;  ay\-la\-ri zu\-saw\-krr 'ay\-eyng.} Let me answer one of the questions now; the others (I'll) answer in the future.~\LINK{http://forum.learnnavi.org/language-updates/email-from-frommer-re-numbers-\%28or-the-full-number-chart!\%29/}{ln}
\B{Fo\-ru 'upxa\-ret oel fpo\-le', slä vay set ke pa\-mä\-hä\-ngem kea tì\-'eyng.} I have sent them a message, but up to now no answer has arrived.~\LINK{http://forum.learnnavi.org/news-announcements/a-response-from-paul-frommer!/}{ln}
\B{Lì'\-fya\-ri le\-Na'\-vi 'Rr\-ta\-mì, vay set 'al\-mong a fra\-'u ze\-ra\-'u ta ngrr\-po\-ngu.} Everything that has gone on with (blossomed regarding) Na'vi until now on Earth has come from a grassroots movement.~\LINK{http://wiki.learnnavi.org/index.php?title=Corpus\#Good_Morning_America}{wiki} \\
\textsc{Derivations}: \ding{43}
\on{1}.~\HL{sekrr}{\B{se\-krr}}, present, present time. 
\on{2}.~\HL{pxiset}{\B{pxi\-set}}, right now.

% sevin
\hi
\HT{sevin}{\B{se\cp{vin}}} \I{[sE."vin]} \emph{adj.} pretty (\emph{primarily of females}).
\B{Lu poe se\-vin nì\-ftxan kum\-a yaw\-ne slo\-lu oer.} She was so beautiful that I fell in love with her.~\LINK{http://naviteri.org/2012/06/spring-vocabulary-part-3/}{b}
\B{Se\-vin\-a tsa\-'eve\-ru a\-hì\-'i me\-sa'\-sem ioi so\-li fa mik\-tsang.} The parents put earrings on that pretty little girl.~\LINK{http://naviteri.org/2011/08/new-vocabulary-clothing/}{b}
\B{Nga\-ri tswin\-tsyìp se\-vin nì\-txan lu nang!} What a pretty little queue you have!~\LINK{http://naviteri.org/2010/07/diminutives-conversational-expressions/}{b}
(\ding{43}~\HL{sayrIp}{\B{say\-rìp}}, \HL{lor}{\B{lor}})

% sey
\hi
\HT{sey}{\B{sey}} \I{[sEj]} \emph{n.} cup, bowl (\emph{minimally modified from naturally occurring resources}).
\B{Fì\-srä\-ti pxaw sey mi\-vam fte tsat hi\-vaw\-nu.} Wrap this cloth around the bowl to protect it.~\LINK{http://naviteri.org/2019/06/50a-liu-amip-40-new-words/}{b}
\B{Fwa pon sey\-ti sìn ki\-nam\-til lu le\-hrr\-ap, ma 'itan.} Balancing a cup on your knee is dangerous, son.~\LINK{http://naviteri.org/2020/04/aawa-u-amip-a-few-new-things/}{b}  \\
\textsc{Derivations}: \ding{43}
\on{1}.~\HL{'e'insey}{\B{'e\-'in\-sey}}, drinking gourd. 
\on{2}.~\HL{sumsey}{\B{sum\-sey}}, drinking vessel (from shell). 
\on{3}.~\HL{swoasey}{\B{swoa\-sey}}, kava bowl.

% seykxel
\hi
\HT{seykxel}{\B{sey\cp{kxel}}} \I{[sEj."k'El]} \emph{adj.} strong: inner strength, a quiet feeling of confidence in one's own capability.   
(\emph{a. conv.}) \B{Sey\-kxel sì nit\-ram!} ``Congratulations!'' [lit.: strong and happy] (\emph{for congratulating s.o. on their good fortune}).~\LINK{http://naviteri.org/2010/07/diminutives-conversational-expressions/}{b}

% seyn
\hi
\HT{seyn}{\B{seyn}} \I{[sEjn]} \emph{n.} chair, stool, bench; \emph{any tool or device to facilitate sitting}. 
\B{Lä\-ngu fì\-seyn kel\-ho\-an nì\-ngay. Tsun oe hi\-veyn tsaw\-sìn 'a\-'aw\-a swaw\-tsyìp nì\-'aw.} This chair is really uncomfortable. I can only sit on it a few seconds.~\LINK{http://naviteri.org/2015/08/aylifyavi-lereyfya-2-cultural-terms-2-and-more/}{b} 
(\ding{43}~\HL{sA}{\B{sä-}}, \HL{heyn}{\B{heyn}})

% seyri
\hi
\HT{seyri}{\B{\cp{sey}ri}} \I{["{}sEj.Ri]} \emph{n.} lip. 
\B{Nge\-yä me\-sey\-ri lor lu nì\-txan.} Your lips are very beautiful.~\LINK{https://www.youtube.com/watch?v=ojr_BE8ZmVE}{ted}

% seysonìltsan
\hi
\HT{seysonIltsan}{\B{seysonìl\cp{tsan}}} \I{[sEj.so.nIl."\t{ts}an]} \emph{conv.} ``Well done!'' (\emph{from} \B{ha\-sey so\-li nìl\-tsan})
\B{Nge\-yä tì\-kang\-kem\-ìri 'e\-fe\-iu oe ye nì\-txan. Sey\-so\-nìl\-tsan!} I'm very satisfied with your work. Well done!~\LINK{http://naviteri.org/2011/04/\%E2\%80\%99a\%E2\%80\%99awa-li\%E2\%80\%99fyavi-amip\%E2\%80\%94a-few-new-expressions/}{b}
(\ding{43}~\HL{hasey}{\B{ha\-sey si}}, \HL{nIltsan}{\B{nìl\-tsan}})

% seyto
\hi
\HT{seyto}{\B{sey\cp{to}}} \I{[sEj."to]} \emph{vtr.} butcher (\emph{i.e. separating or processing the carcass of a dead animal}).
\B{Aw\-ngal fì\-ye\-rik\-it nì\-win si\-vey\-to ko.} Let's cut up this hexapede quickly.~\LINK{http://naviteri.org/2015/08/aylifyavi-lereyfya-2-cultural-terms-2-and-more/}{b} \\
\textsc{Derivations}: \ding{43}
\on{1}.~\HL{sAseyto}{\B{sä\-sey\-to}}, butchering tool.

% seyä
\hi
\HT{seyA}{\B{\cp{se}yä}} \I{["{}sE.j\ae]}~:~\HL{'u}{\B{'u}} (\emph{shortened genitive form of the plural of} \HL{tsa'u}{\B{tsa\-'u})}.

% seze
\hi
\HT{seze}{\B{\cp{se}ze}} \I{["{}sE.zE]} \emph{n.} (\emph{a. \ding{96}}) blue flower. 
\B{Na se\-ze | A 'ong ne tsaw\-ke.} Like the blue\hyp{}flower | Which unfolds to the sun.~\LINK{http://www.pandorapedia.com/navi/music/ritual_music}{pp} 
(\emph{b. proper name}) \emph{name of Neytiri's banshee}.

% si
\hi
\HT{si}{\B{si}} \I{[si]} \emph{unprod. auxiliary v. to turn a n. or adj. into a v.:} \HL{tIkangkem}{\B{tì\-kang\-kem si}}, to work. \HL{win}{\B{win si}}, to hurry. \emph{Grammatically vin., the logical direct object takes the dative case:}  
\B{Nì\-'i'\-a po tsa\-tì\-kang\-kem\-vir ha\-sey so\-li.} She finally completed the project.~\LINK{http://naviteri.org/2011/04/\%E2\%80\%99a\%E2\%80\%99awa-li\%E2\%80\%99fyavi-amip\%E2\%80\%94a-few-new-expressions/}{b} 
\emph{Follows the non\hyp{}verbal element; with} \B{ira\-yo si} \emph{word order can be broken:} 
\B{Oel ay\-nga\-ti ka\-me\-i\-e, ma oe\-yä ey\-lan, ul\-te ay\-nga\-ru sei\-yi ira\-yo.} I \textsc{see} you all, my friends, and thank you.~\LINK{http://forum.learnnavi.org/news-announcements/a-response-from-paul-frommer!/}{ln} 
\emph{Can stand alone when the non\hyp{}verbal element has already been mentioned and is understood from context:}
\B{Nga tsap\-'alu\-te so\-li srak? -- Soli.} Did you apologize? -- (Yes,) I have.~\LINK{http://forum.learnnavi.org/language-updates/auxilary-verb-si-possessive-dative-krr/}{ln}
\emph{Negation precedes the auxiliary v.:} 
\B{Swey lu fwa nga fì\-kem ke si\-vi nì\-sngum.} It's best that you not do this fretfully.~\LINK{http://naviteri.org/2011/01/new-year-new-vocabulary/}{b}
\B{Txo\-pu rä\-'ä si, lu ke\-tu\-wong\-o nì\-'aw.} Don't be afraid, it's only an alien.~\LINK{http://forum.learnnavi.org/language-updates/a-collection/}{ln\slash a\textsuperscript{c}} \\
\emph{Special forms in} ‹\HL{Ap}{\B{äp}}› \emph{and} ‹\HL{eyk2}{\B{eyk}}› \emph{exist, e.g.:} \ding{43} \HL{ioi}{\B{ioi sä\-pi}}, adorn oneself. \HL{win}{\B{win sä\-pi}}, hurry oneself. \HL{tIsraw}{\B{tì\-sraw sey\-ki}}, hurt s.o. \B{sä\-pi} \emph{and its derivations is often shortened to} \I{[spi]}.

% sil
\hi 
\HT{sil}{\B{sil}}\textsuperscript{1} \I{[sil]} \emph{n.} body stripe (\emph{non\hyp{}facial}). 
(\emph{a. proverbial expression}) \B{Po\-ri fra\-pil sì fra\-sil lor lu nì\-txan.} He\slash She is very attractive.~\LINK{https://naviteri.org/2024/06/ayhapxi-tokxa-si-ayliu-alahe-parts-of-the-body-and-more/}{b} 
\B{Fra\-pil fra\-sil!} `Wow! Check him\slash her out! He\slash she is gorgeous!'~\LINK{https://naviteri.org/2024/06/ayhapxi-tokxa-si-ayliu-alahe-parts-of-the-body-and-more/}{b}
(\ding{43}~\HL{pil}{\B{pil}})

\hi
\B{sil}\textsuperscript{2} \I{[sil]}~:~\HL{til}{\B{til}}

% sim
\hi
\HT{sim}{\B{sim}}\textsuperscript{\on{1}} \I{[sim]} \emph{vin.} near, be near. 
\B{Oe 'o\-long\-okx mì sray a txam\-pay\-ìri sim.} I was born in a town near the ocean.~\LINK{http://naviteri.org/2010/09/getting-to-know-you-part-2/}{b} \\
\textsc{Derivations}: \ding{43}
\on{1}.~\HL{asim}{\B{asim}}, near by, at close range.

\hi
\B{sim}\textsuperscript{\on{2}} : \HL{tsim}{\B{tsim}}

% sì
\hi
\HT{sI}{\B{sì}} \I{[sI]} \emph{conj., suf.} and (\emph{for enumeration}).
\B{Ay\-ey\-lan\-ur oe\-yä sì ey\-lan\-ur lì'\-fya\-yä le\-Na'\-vi nì\-wotx.} To all my friends and friends of the Na'vi language.~\LINK{http://forum.learnnavi.org/news-announcements/a-response-from-paul-frommer!/}{ln}
\B{Lu nga win sì txur.} You are fast and strong.~\LINK{http://wiki.learnnavi.org/index.php/Hunting_Song}{wiki}
\B{Tsay\-fne\-sä\-num\-vi a tsun fra\-por srung si\-vi fte ni\-vu\-me sì zi\-ve\-rok nì\-swey.} The kinds of lessons that can help everyone best learn and remember.~\LINK{http://forum.learnnavi.org/language-updates/trr-rrtaya/}{ln}
(\emph{a. can be used as an enclitic})
\B{Ta 'ey\-lan kar\-yu\-sì ay\-nge\-yä}. From your friend and teacher.~\LINK{http://forum.learnnavi.org/news-announcements/a-response-from-paul-frommer!/}{ln}
\B{Tsa\-krr pa\-ye\-'un swey\-a fya\-'ot a za\-mi\-vu\-nge oel ay\-ngar ay\-lì\-'ut ho\-ren\-ti\-sì lì'\-fya\-yä le\-Na'\-vi.} I will then decide the best way to bring you the words and rules of Na'vi.~\LINK{http://forum.learnnavi.org/news-announcements/a-response-from-paul-frommer!/}{ln}
\B{Fay\-upxa\-re la\-yu ay\-sngä\-'i\-yu\-fpi, fte lì'\-fya\-ri aw\-nge\-yä fo tsì\-ye\-vun nì\-ftue nìl\-tsan\-sì ni\-vume.} These messages will be for beginners, so that they can learn about our language easily and well.~\LINK{http://naviteri.org/2010/06/first-post/}{b} 
(\emph{b. often merges in coll. pronunciation}) \B{ay\-nan\-tang sì ay\-ri\-ti} \emph{pronounced:} \I{[aj."nan.taN saj."Ri.ti]}, viperwolfs and stingbats.~\LINK{http://naviteri.org/2010/07/thoughts-on-ambiguity/}{b}
(\emph{c. idiom}) \B{Kxe\-tse sì mik\-yun kop pll\-txe.} Body language is important [lit.: The tail and ears also speak].~\LINK{http://forum.learnnavi.org/language-updates/grid-tidbits/}{ln}
(\ding{43}~\HL{ulte}{\B{ul\-te}})

% sìk
\hi
\HT{sIk}{\B{sìk}} \I{[sIk\textcorner]} \emph{part. to indicate the end of a quote with verbs of speech} \HL{plltxe}{\B{pll\-txe}}, \HL{peng}{\B{peng}}, \HL{pllngay}{\B{pll\-ngay}}, \HL{pawm}{\B{pawm}}, \emph{and writing} \HL{pamrel}{\B{pam\-rel si}} \emph{when the speaker continues to say s.t. after the cited quote; used together with} \ding{43}~\HL{san}{\B{san}}.
\B{Po\-lawm po san sra\-ke Sä\-li ho\-lum sìk.} He asked, ``Did Sally leave?''~\LINK{http://forum.learnnavi.org/language-updates/if-vs-whether-ftxey-fuke/}{ln}
\B{Pol\-txe Ey\-tu\-kan san oe ka\-yä sìk, slä oel pot ke spaw.} Eytukan said he would go but I don't believe him.~\LINK{http://wiki.learnnavi.org/index.php?title=Canon\#Extracts_from_various_emails}{wiki}

% sìlpey
\hi
\HT{sIlpey}{\B{sìl\cp{pey}}} \I{[sIl."p$\cdot$Ej]} \emph{vin., inf.}\on{22} hope (\emph{used optionally with} \HL{tsnI}{\B{tsnì}}, \emph{often with} ‹\B{iv}›), \B{sìl\-pey oe}, I hope, hopefully.~\LINK{http://forum.learnnavi.org/language-updates/txelanit-hivawl/}{ln}
\B{Sìl\-pey oe, la\-yu oe\-ru ye'\-rìn sìl\-tsan\-a fmawn.} I hope, I will soon have good news.~\LINK{http://forum.learnnavi.org/news-announcements/a-response-from-paul-frommer!/}{ln} 
\B{Ul\-te sìl\-pey oe, fay\-sì\-oeyk\-tìng li\-vu law nì\-wotx!} And I hope, all these explanations are clear.~\LINK{http://naviteri.org/2011/08/reported-speech-reported-questions/}{b}
\B{Fu\-ri\-a fì\-lì'\-fya\-vit fkol ta 'Ìng\-lì\-sì mo\-lu\-nge, sìl\-pey oe tsnì ay\-ha\-pxì\-tu lì'\-fya\-o\-lo'\-ä aw\-nge\-yä ke stì\-ye\-vi.} That this expression is brought from English, will hopefully not anger the members of the language group.~\LINK{http://forum.learnnavi.org/language-updates/txelanit-hivawl/}{ln}
(\ding{43}~\HL{sIltsan}{\B{sìl\-tsan}}, \HL{pey}{\B{pey}}; \ding{214}~\HL{fe'\-pey}{\B{fe'\-pey}}) \\
\textsc{Derivations}: \ding{43}
\on{1}.~\HL{nIsIlpey}{\B{nì\-sìl\-pey}}, hopefully. 
\on{2}.~\HL{tsIlpey}{\B{tsìl\-pey}}, hope (abstract). 
\on{3}.~\HL{sAsIlpey}{\B{sä\-sìl\-pey}}, a hope.

% sìlronsem
\hi
\HT{sIlronsem}{\B{sìl\cp{ron}sem}} \I{[sIl."Ron.sEm]} \emph{adj.} (\emph{of things}) clever, smart. 
\B{Ay\-lì\-'u\-van a\-swey lu 'i\-pu, lu sìl\-ron\-sem.} The best puns are both funny and clever.~\LINK{http://naviteri.org/2012/01/mipa-zisit-ayliu-amip-new-words-for-the-new-year/}{b}
(\ding{43}~\HL{sIltsan}{\B{sìl\-tsan}}, \HL{ronsem}{\B{ron\-sem}})

% sìltsan
\hi
\HT{sIltsan}{\B{sìl\cp{tsan}}} \I{[sIl."\t{ts}an]} \emph{adj.} good.   
\B{Tì\-'e\-fu\-mì oe\-yä su\-te\-kip nì\-wotx, ftxey Na'\-vi ftxey Saw\-tu\-te, lu sìl\-tsan lu kawng.} In my opinion among all people, whether Na'vi or Skypeople, there are good and bad.~\LINK{http://naviteri.org/2010/07/ting-mikyun-fte-tslivam-listening-comprehension-1-3/}{b}
\B{Sìl\-pey oe, la\-yu oe\-ru ye'\-rìn sìl\-tsan\-a fmawn.} I hope, I will soon have good news.~\LINK{http://forum.learnnavi.org/news-announcements/a-response-from-paul-frommer!/}{ln} 
\B{Slä lu oe\-ru fmawn\-o a\-sìl\-tsan.} But I have some good news.~\LINK{http://forum.learnnavi.org/language-updates/trr-rrtaya/}{ln}
\B{Po to oe lu sìl\-tsan.} S\slash he is better than me.~\LINK{http://forum.learnnavi.org/language-updates/prrwll-moss-the-comparative-to-construct/}{ln} 
(\ding{43}~\HL{swey}{\B{swey}}; \ding{214}~\HL{fe'}{\B{fe'}}, \HL{kawng}{\B{kawng}}) \\
\textsc{Derivations}: \ding{43}
\on{1}.~\HL{nIltsan}{\B{nìl\-tsan}}, well, in a good way. 
\on{2}.~\HL{tsIltsan}{\B{tsìl\-tsan}}, good, goodness.
\on{3}.~\HL{tsantu}{\B{tsan\-tu}}, good person.
\on{4}.~\HL{tsantxAl}{\B{tsan\-txäl}}, invitation.

% sìn
\hi
\HT{sIn}{\B{sìn}}\textsuperscript{\on{1}} \I{[sIn]} \emph{adp.} (\emph{a. locally}) on, onto.
\B{Ay\-wayl yìm ki\-fkey\-ä | 'Ì\-he\-yut a\-vo\-mrr | Sìn ti\-rea\-fya\-'o a\-vol | Na way\-te\-lem\-ä hìng.} The songs bind | the thirteen spirals of the solid world | To the eight spirit paths | Like the threads of a Songcord.~\LINK{http://www.pandorapedia.com/navi/music/ritual_music}{pp}
\B{Sre fwa sìn tskxe\-pay tì\-ran, ze\-ne fko flì\-nutx\-it sti\-ve\-ftxaw.} It's necessary to check the thickness of the ice before walking on it.~\LINK{http://naviteri.org/2011/11/\%E2\%80\%99a\%E2\%80\%99awa-ayli\%E2\%80\%99u-amip-ni\%E2\%80\%99aw-\%E2\%80\%94-a-few-new-words-only/}{b}
\B{Ul\-txa si\-vi oeng sìn ram\-tsyìp txon\-'ong\-ay.} Let's meet on the hill tomorrow at nightfall.~\LINK{http://naviteri.org/2012/01/mipa-zisit-ayliu-amip-new-words-for-the-new-year/}{b}
\B{Ya\-mì tsun fko tsi\-ve\-'a lo\-ran\-it re\-nu\-ä kil\-van\-ä slä kll\-te\-sìn wä\-pan.} From the air you can see the grace of the river's form but from the ground it's hidden.~\LINK{http://naviteri.org/2012/10/mipa-vospxi-mipa-ayliu-new-words-for-the-new-month/}{b} 
\B{Ngey tsko\-ti yem tsa\-yì\-sìn tsa\-krr za\-'u fì\-tseng!} Put your bow on that ledge and come here!~\LINK{http://naviteri.org/2013/01/lifyari-po-peyi-how-good-is-her-navi/}{b}
\B{Oe\-ri ay\-som\-pì\-va sìn re\-'o var zi\-vup.} Raindrops keep falling on my head.~\LINK{http://naviteri.org/2011/04/yafkeykiri-plltxe-frapo-everyone-talks-about-the-weather/}{b}
(\emph{b. with} \HL{fAkA}{\B{fäkä}}) get on, mount. \B{Po sìn pa'\-li fä\-ko\-lä mak\-to ne\-to.} She got on her direhorse and rode away.~\LINK{https://naviteri.org/2025/04/mevosinga-liu-amip-twenty-new-words/}{b}
(\emph{c. metaphorically with} \HL{kllkxem}{\B{kll\-kxem}} \emph{and} \HL{yI}{\B{yì}} \emph{to refer to the level of anything scalable: water level, temperature, talent, anger, etc.}) 
\B{Lì'\-fya\-ri po (kll\-kxem) sìn pe\-yì?} [\emph{coll.}: \I{[sIm.pE."jI]}]\slash \B{Lì'\-fyari po pe\-yì?} How good is her Na'vi?~\LINK{http://naviteri.org/2013/01/lifyari-po-peyi-how-good-is-her-navi/}{b}
\B{Sìn yì sngä\-'i\-yu\-ä, slä tsye\-rìl (hay\-a yì\-ne) nì\-'ul\-'ul.} They're at a beginner's level, but they're getting better and better.~\LINK{http://naviteri.org/2013/01/lifyari-po-peyi-how-good-is-her-navi/}{b}
\B{Sìn yì a ke tsun kaw\-tu spi\-vaw, nì\-win fra\-to.} He runs at an incredible level, faster than anyone else.~\LINK{http://naviteri.org/2013/01/lifyari-po-peyi-how-good-is-her-navi/}{b}

\hi
\B{sìn}\textsuperscript{\on{2}} : \B{tìn}

% -sìng
\hi
\HT{sIng}{\B{-sìng}} \I{[.sIN]} \emph{num., suf.} four (\emph{rest form of \ding{43}~\HL{tsIng}{\B{tsìng}} to form higher numbers}). \B{vo\-sìng}, \on{12} (decimal), \on{14} (octal). \B{me\-vo\-sìng}, \on{20} (decimal), \on{24} (octal); etc.

% ska'a
\hi
\HT{ska'a}{\B{ska\cp{'a}}} \I{[ska."Pa]} \emph{vtr.} destroy.
\B{He\-tu\-wong\-ìl aw\-nge\-yä swo\-tut ska\-'a, fte kll\-ki\-vu\-lat ke\-ru\-sey\-a tskxet.} The aliens destroy our sacred place to dig up dead rock.~\LINK{http://forum.learnnavi.org/language-updates/participles/}{ln\slash a\textsuperscript{c}}
\B{Zo\-la\-'u saw\-tu\-te fì\-tseng fte aw\-nga\-ti ski\-va\-'a.} The Skypeople have come to destroy us.~\LINK{http://naviteri.org/2010/07/ting-mikyun-fte-tslivam-listening-comprehension-1-3/}{b}
\B{Maw fwa fkol Kel\-ut\-ral\-it sko\-la\-'a, lem\-rey\-a ha\-pxì\-tul tsa\-soai\-ä txo\-lu\-la mip\-a kel\-ku\-ti mì ta\-yo.} After Hometree was destroyed, the surviving members of that family built a new home on the plains.~\LINK{http://naviteri.org/2011/05/some-miscellaneous-vocabulary/}{b}
\B{Lu oer sngum\-tsim a pol Ame\-ri\-kat ski\-ye\-va\-'a.} My source of worry is that he will destroy America.~\LINK{http://naviteri.org/2016/12/tengkrr-perahem-zisit-amip-as-the-new-year-arrives/}{b} \\
\textsc{Derivations}: \ding{43}
\on{1}.~\HL{tIska'a}{\B{tì\-ska\-'a}}, destruction.

% skepek
\hi
\HT{skepek}{\B{\cp{ske}pek}} \I{["{}skE.pEk\textcorner]} \emph{adj.} formal. 
\B{Ay\-lì\-'u ngi\-an nì\-'it ske\-pek lu.} However the words are a bit formal.~\LINK{http://forum.learnnavi.org/language-updates/more-grammatical-tidbits/}{ln\slash a}

% skeykx
\hi 
\HT{skeykx}{\B{skeykx}} \I{[skEjk']} \emph{adj.} bored. 
\B{Oe 'e\-fu skeykx ul\-te new ti\-vä\-txaw ne kel\-ku.} I'm bored and I want to go home.~\LINK{http://naviteri.org/2024/02/mipa-ayliu-si-aylifyavi-niul-more-new-words-and-expressions/}{b}
(\ding{214}~\HL{eltu}{\B{eltu si}})

% skiempa
\hi
\HT{skiempa}{\B{\cp{ski}empa}} \I{["{}ski.Em.pa]} \emph{n.} right side.
\B{Mìn ne ski\-em\-pa.} Turn to the right.~\ding{43}~\LINK{http://forum.learnnavi.org/language-updates/ftarpa-and-skienpa-added-to-dict/msg570362/\#msg570362}{ln}
(\ding{43}~\HL{skien}{\B{ski\-en}}, \HL{pa'o}{\B{pa\-'o}}; \ding{214}~\HL{ftArpa}{\B{ftär\-pa}})

% skien
\hi
\HT{skien}{\B{\cp{ski}en}} \I{["{}ski.En]} \emph{adj.} right (\emph{not left, direction}). 
\B{Oe\-ri ski\-en\-a tsyokx lu le\-skxir.} My right hand is wounded.~\LINK{http://naviteri.org/2014/05/mipa-ayliu-mipa-aysafpil-new-words-new-ideas/}{b}
(\ding{214}~\B{ftär}) \\
\textsc{Derivations}: \ding{43}
\on{1}.~\HL{skiempa}{\B{ski\-em\-pa}}, right side. 
\on{2}.~\HL{nIskien}{\B{nì\-ski\-en}}, to the right.

% sko
\hi
\HT{sko}{\B{sko}}\textsuperscript{\on{1}} \I{[sko]} \emph{adp.}\texttt{+} (\emph{a}) as, in the capacity of, in the role of.   
\B{Sko Sa\-hìk ke tsun oe mì\-ftxe\-le tsngi\-vaw\-vìk, sko sa'\-nok tsun.} As Tsahik, I cannot weep over this matter, as mother I can.~\LINK{http://naviteri.org/2012/03/spring-vocabulary-part-2/}{b}
\B{Set pll\-txe nga nìl\-tsan nì\-ngay. Fì\-sä\-tsan\-'ul\-ìri nge\-yä lu oe\-ru sko har\-yu nrr\-a.} You speak really well now. As your teacher, I'm proud of your improvement.~\LINK{http://naviteri.org/2013/03/tsanerul-feerul-getting-better-getting-worse/}{b}
\B{'Eveng\-ìl skxe\-vit 'ey\-krr\-ko sko uvan.} Children roll pebbles as a game.~\LINK{http://naviteri.org/2015/04/some-new-words-for-may-day/}{b}
(\emph{b}) \emph{to convey the meaning of} `I myself, you yourself' etc. 
\B{Nga sko nga pol\-txe san rä\-'ä ki\-vä!\slash Nga nga\-sko pol\-txe …} You yourself said don't go!~\LINK{https://naviteri.org/2024/03/fleltrra-ayliu-words-for-april-fools-day/}{b}

% sko
\hi
\B{sko}\textsuperscript{\on{2}} : \B{tsko}.

% skuka
\hi 
\HT{skuka}{\B{\cp{sku}ka}} \I{["{}sku.ka]} \emph{n., fauna} Sagittaria, \emph{cephalopod\hyp{}like aquatic creature with a hard shell and long tentacles}.~\LINK{https://forum.learnnavi.org/language-updates/expansion-to-flora-fauna-and-places-from-the-valley-of-mo'ara/}{ln}

% skxakep
\hi
\HT{skxakep}{\B{\cp{skxa}kep}} \I{["{}sk'a.kEp\textcorner]} \textsc{i}. \emph{adj.} probable. \\ 
\textsc{ii}. \emph{adv.} probably. 

% skxawng
\hi
\HT{skxawng}{\B{skxawng}} \I{[sk'awN]} \emph{n.} moron, idiot.
\B{Fì\-kar\-yu alu fì\-po lu tsul\-fä\-tu; tsa\-kar\-yu alu tsa\-po lu skxawng.} \emph{This} teacher is a master; \emph{that} teacher is a fool.~\LINK{http://naviteri.org/2011/12/one-more-for-2011/}{b}
\B{Oe\-ti 'am\-pi rä\-'ä, ma \mbox{skxawng}!} Don't touch me, you moron!~\LINK{http://naviteri.org/2012/11/renu-ayinanfyaya-the-senses-paradigm/}{b}
\B{Fì\-skxawng\-ìri tsap\-'a\-lu\-te se\-ngi oe.} I apologize for this moron.~\LINK{http://wiki.learnnavi.org/index.php?title=Corpus\#Times_Online}{wiki}
\B{Nem\-fa na'\-rìng, ma skxawng\-tsyìp! Kä li!} Into the forest, you moron! Get going! (\emph{endearing})~\LINK{http://naviteri.org/2011/05/some-miscellaneous-vocabulary/}{b}
(\emph{a. proverb}) \B{Spä skxawng sìn 'a\-na a\-flì.} ``A fool jumps onto a thin vine,'' i.e. \emph{don't engage in an unpromising and\slash or potentially risky cause}.~\LINK{http://naviteri.org/2021/08/ulte-ayyoratu-leiu-and-the-winners-are-2/}{b}
\B{Tsa\-ye\-rik te\-rul ne 'awkx. Spä skxawng sìn 'a\-na a\-flì.} That hexapede is running toward the cliff. Only a fool jumps onto a thin vine.~\LINK{http://naviteri.org/2021/08/ulte-ayyoratu-leiu-and-the-winners-are-2/}{b}

% skxir / si
\hi
\HT{skxir}{\B{skxir}} \I{[sk'iR]} \textsc{i}. \emph{n.} wound.  
\B{Oe\-ri skxir a mì syokx tì\-sraw se\-ngi.} The wound on my hand hurts.~\LINK{http://naviteri.org/2011/05/some-miscellaneous-vocabulary/}{b}
\B{Ze\-ne nga yi\-vur pxìm fì\-skxir\-it.} You must clean the wound frequently.~\LINK{http://naviteri.org/2011/05/some-miscellaneous-vocabulary/}{b}
\B{Rey\-pay skxir\-ftum\-fa he\-rum.} Blood is coming out of the wound.~\LINK{http://naviteri.org/2013/04/wheres-the-bathroom-and-other-useful-things/}{b} \\
\textsc{ii}. \B{skxir si} \I{[sk'iR si]} \emph{vin.} wound.
\B{Ta\-ron\-yu ye\-rik\-ur skxir so\-li.} The hunter wounded the hexapede.~\LINK{http://naviteri.org/2011/05/some-miscellaneous-vocabulary/}{b} \\
\textsc{Derivations}: \ding{43}
\on{1}.~\HL{leskxir}{\B{le\-skxir}}, wounded. \\
\on{2}.~\HL{skxirtsyIp}{\B{skxir\-tsyìp}}, cut, bruise, minor wound.

% skxirtsyìp
\hi
\HT{skxirtsyIp}{\B{\cp{skxir}tsyìp}} \I{["sk'iR.\t{ts}jIp\textcorner]} \emph{n.} cut, bruise, minor wound. 
(\ding{43}~\HL{skxir}{\B{skxir}}, \HL{tsyIp}{\B{-tsyìp}})

% skxom
\hi
\HT{skxom}{\B{skxom}} \I{[sk'om]} \emph{n.} chance, opportunity. 
\B{Ngay\-txo\-a, nì\-aw\-no\-mum ke lo\-lu oer nì\-ke\-ftxo mì sok\-a srr ay\-skxom le\-tam fte lì'\-fya\-ri aw\-nge\-yä tì\-kang\-kem si\-vi.} My apologies: As you know, in recent days I have not had sufficient opportunity to work on our language.~\LINK{http://forum.learnnavi.org/language-updates/trr-rrtaya/}{ln}
\B{Nän ftia, nän lu skxom a em\-za\-'u.} The less you study, the less chance you have of passing.~\LINK{http://naviteri.org/2012/02/trr-asawnung-lefpom-happy-leap-day/}{b}

% sl.
\hi
\B{sl.}~:~\HL{saylahe}{\B{saylahe}}.

% sla'tsu
\hi
\HT{sla'tsu}{\B{\cp{sla'}tsu}} \I{["{}slaP.\t{ts}u]} \emph{vtr.} describe.   
\B{Pol sla'\-tsu ay\-ioang\-it a tse\-'a fko mì Ey\-wa\-'e\-veng.} She describes animals seen on Pandora.~\LINK{http://naviteri.org/2011/01/new-year-new-vocabulary/}{b} \\
\textsc{Derivations}: \ding{43}
\on{1}.~\HL{tIsla'tsu}{\B{tì\-sla'\-tsu}}, description.

% slayk
\hi
\HT{slayk}{\B{slayk}} \I{[slajk\textcorner]} \emph{vtr.} brush, comb.
\B{New sa'\-nok sli\-vayk nik\-ret 'eveng\-ä.} The mother wants to brush the child's hair.~\LINK{http://naviteri.org/2015/03/kap-si-ayunil-saylahe-threats-dreams-and-other-things/}{b}

% slan
\hi
\HT{slan}{\B{slan}} \I{[slan]} \emph{vtr.} support (\emph{emotionally, socially, or personally}). 
\B{Tì\-wä\-te\-ri ngal oe\-ti pe\-lun ke slan kaw\-krr?} Why don't you ever support me in an argument?~\LINK{http://naviteri.org/2012/03/spring-vocabulary-part-1/}{b} 
\B{Ul\-te wä sì\-kawng a fì\-tì\-wu\-sem\-ìri, ze\-ne aw\-nga nì\-wotx fì\-tsap slä\-pi\-van.} And in this fighting against this evil we must all support each other.~\LINK{http://naviteri.org/2016/12/tengkrr-perahem-zisit-amip-as-the-new-year-arrives/}{b} 
(\ding{43}~\HL{kavan}{\B{ka\-van}}) \\
\textsc{Derivations}: \ding{43}
\on{1}.~\HL{tIslan}{\B{tì\-slan}}, support. 
\on{2}.~\HL{slantire}{\B{slan\-ti\-re}}, inspiration.

% slantire / si
\hi
\HT{slantire}{\B{slanti\cp{re}}} \I{[slan.ti."RE]} \textsc{i}. \emph{n.} inspiration.
\B{Mll\-te oe nì\-wotx. Fì\-slan\-ti\-re\-ri irayo nì\-txan, ma Tsyi\-li!} I agree completely. Thank you for this inspiration, Tsyi\-li.~\LINK{http://naviteri.org/2019/09/choice-statements-vs-choice-questions-and-some-insults/\#comment-30391}{b} \\
\textsc{ii}. \B{slanti\cp{re} si} \I{[slan.ti."RE si]} \emph{vin.} inspire.
\B{Sä\-ftxu\-lì\-'u Tsyeyk\-ä Na'\-vi\-ru slan\-ti\-re so\-li nì\-wotx.} Jake's speech inspired all the People.~\LINK{http://naviteri.org/2012/06/spring-vocabulary-part-3/}{b}
(\ding{43}~\B{slan}, \B{tirea})

% slä
\hi
\HT{slA}{\B{slä}} \I{[sl\ae]} \emph{conj.} but.   
\B{Fo\-ru 'upxa\-ret oel fpo\-le', slä vay set ke pa\-mä\-hä\-ngem kea tì\-'eyng.} I have sent them a message, but up to now no answer has arrived.~\LINK{http://forum.learnnavi.org/news-announcements/a-response-from-paul-frommer!/}{ln}
\B{Pol\-txe Ey\-tu\-kan san oe ka\-yä sìk, slä oel pot ke spaw.} Eytukan said he would go but I don't believe him.~\LINK{http://wiki.learnnavi.org/index.php?title=Canon\#Extracts_from_various_emails}{wiki} 
(\ding{43}~\HL{ki}{\B{ki}}) \\
\textsc{Derivations}: \ding{43}
\on{1}.~\HL{slAkop}{\B{slä\-kop}}, but also.

% släkop
\hi
\HT{slAkop}{\B{\cp{slä}kop}} \I{["{}sl\ae.kop\textcorner]} \emph{adv.} but also. 
(\emph{a. together with} \ding{43}~\HL{ken'aw}{\B{ken\-'aw}})
\B{Nge\-yä tsmu\-ke lu ken\-'aw lor slä\-kop ka\-nu.} Your sister is not only beautiful but also intelligent.~\LINK{http://naviteri.org/2014/10/tson-si-fpomron-obligation-and-mental-health/}{b}
\B{Ay\-nga\-ri sìl\-pey oe tsnì ken\-'aw fpom\-tokx slä\-kop fpom\-ron yi\-vo'.} I hope that not only your physical but also your mental health will be perfect.~\LINK{http://naviteri.org/2014/10/tson-si-fpomron-obligation-and-mental-health/}{b}
(\emph{b. on its own}) \B{Fra\-krr lu nge\-yä sì\-pawm ngä\-zìk slä\-kop le\-tsran\-ten.} Your questions are always difficult but also important.~\LINK{http://naviteri.org/2014/10/tson-si-fpomron-obligation-and-mental-health/}{b}
(\ding{43}~\HL{slA}{\B{slä}}, \HL{kop}{\B{kop}})

% slär
\hi
\HT{slAr}{\B{slär}} \I{[sl\ae{}R]} \emph{n.} cave.
\B{Slär\-ìl nga' fwep\-it a\-txan.} The cave is very dusty.~\LINK{http://naviteri.org/2014/07/tsawlultxamaw-a-ayliu-post-meetup-words/}{b}
\B{Fì\-slär\-mì tsun fko sti\-vawm ngam\-it a\-pxay.} You can hear a lot of echoes in this cave.~\LINK{http://naviteri.org/2012/01/mipa-zisit-ayliu-amip-new-words-for-the-new-year/}{b}
\B{Ri\-ti tswo\-lay\-on ftum\-fa slär.} The stingbat flew out of the cave.~\LINK{http://naviteri.org/2013/04/wheres-the-bathroom-and-other-useful-things/}{b} 

% slele
\hi
\HT{slele}{\B{\cp{sle}le}} \I{["{}slE.lE]} \emph{vin.} swim.   
\B{Le\-hrr\-ap lu fwa evi\-tsyìp sle\-le mì hil\-van lu\-ke fwa fyeyn\-tu te\-rìng nari-} It's dangerous for tiny ones to swim in the river without an adult watching.~\LINK{http://naviteri.org/2011/02/new-vocabulary-ii\%E2\%80\%94part-1/}{b} 
\B{Txo fkol ke fyi\-vel u\-ran\-it pay\-wä, ze\-ne fko sli\-ve\-le.} If one does not seal a boat against water, one must swim.~\LINK{http://naviteri.org/2012/06/spring-vocabulary-part-3/}{b} \\
\textsc{Derivations}: \ding{43}
\on{1}.~\HL{nIslele}{\B{nì\-sle\-le}}, by swimming.

% sleyku
\hi
\HT{sleyku}{\B{sley\cp{ku}}} \I{[slEj."k$\cdot$u]} \emph{vtr., inf.}\on{22} (\emph{a.}) produce, make.
\B{Nga pam\-rel so\-li san sley\-ki\-vu ut\-ral\-ìl lì'\-fya\-yä le\-Na'\-vi pxay\-a ay\-vur\-it a\-txur sìk.} You wrote, ``May the Na'vi language tree produce many strong stories.''~\LINK{http://forum.learnnavi.org/language-updates/life-death-and-gerunds/}{ln}
\B{Ta\-kuk set tsaw\-ke. Sley\-ku fra\-'ut u\-kxo.} Now the sun bears down. (It) makes everything dry.~\LINK{http://forum.learnnavi.org/language-updates/learnings-from-the-siatlla-song-translations/msg475008/\#msg475008}{ln}
(\emph{b. for emotional states}) ``make''. 
\B{Fu\-la tsa\-yun oeng pi\-väng\-kxo ye'\-rìn oe\-ti nit\-ram sley\-ku nì\-txan.} It makes me very happy that we'll be able to chat soon.~\LINK{http://forum.learnnavi.org/language-updates/fula-short-form-of-fiul-a-and-tsal-oeti-steyki-nitram/}{ln}
\B{Fì\-sä\-'an\-lal Ne\-we\-yä nga\-ti sley\-ka\-yu lek\-ye\-'ung!} This yearning for Newey is going to drive you crazy.~\LINK{http://naviteri.org/2013/01/awvea-posti-zisita-amip-first-post-of-the-new-year/}{b}
(i.) \B{txo\-pu sley\-ku}, frighten, scare, produce fear. \B{'Uol ik\-ran\-it txo\-pu sley\-ko\-la\-tsu, ta\-lun\-a po tsìk ya\-wo.} Something must have frightened the banshee, because it suddenly took to the air.~\LINK{http://naviteri.org/2012/01/mipa-zisit-ayliu-amip-new-words-for-the-new-year/}{b}
(\ding{43}~\HL{slu}{\B{slu}})

% sloa
\hi
\HT{sloa}{\B{\cp{slo}a}} \I{["{}slo.a]} \emph{adj.} wide.
\B{Tsa\-ut\-ral\-ìri ta\-ngek lu sloa nì\-txan; 'evi ke tsun tsyi\-vìl.} The trunk of that tree is very wide; the kid cannot climb it.~\LINK{http://naviteri.org/2012/06/spring-vocabulary-part-3/}{b}
(\ding{214}~\HL{snep}{\B{snep}}) \\
\textsc{Derivations}: \ding{43} 
\on{1}.~\HL{slosnep}{\B{slo\-snep}}, width.
\on{2}.~\HL{slotsyal}{\B{slo\-tsyal}}, stormglider.

% sloan
\hi
\HT{sloan}{\B{slo\cp{an}}} \I{[slo."{}an]} \emph{vtr.} pour (s.t.).
\B{Ru\-txe sli\-vo\-an ngal pay\-it oe\-fpi; 'efu väng nì\-txan.} Please pour me some water; I'm very thirsty.~\LINK{http://naviteri.org/2011/05/some-miscellaneous-vocabulary/}{b}

% slosnep / slosneppe
\hi
\HT{slosnep}{\B{slo\cp{snep}}} \I{[slo."{}snEp\textcorner]} \textsc{i}. \emph{n.} width. \\
\textsc{ii}. \B{slo\cp{snep}pe} \I{[slo."{}snEp\textcorner.pE]} \emph{or} \B{pe\-slo\cp{snep}} \I{[pE.slo."{}snEp\textcorner]} \emph{inter.} what width? how wide? 
\B{Kil\-van\-ìri tsa\-tseng slo\-snep\-pe?} How wide is the river there?~\LINK{http://naviteri.org/2012/06/spring-vocabulary-part-3/}{b}
(\ding{43}~\HL{sloa}{\B{sloa}}, \HL{snep}{\B{snep}})

% slotsyal
\hi 
\HT{slotsyal}{\B{\cp{slo}tsyal}} \I{["{}slo.\t{ts}jal]} \emph{n., fauna} stormglider [lit.: `wide wing']. 
(\ding{43}~\HL{sloa}{\B{sloa}}, \HL{tsyal}{\B{tsyal}}) 

% slu
\hi
\HT{slu}{\B{slu}} \I{[slu]} \emph{vin.}  become. 
(\emph{a. with n.s}) \B{Nga\-ri hu Ey\-wa sa\-lew ti\-rea, tokx 'ì\-'awn, slu Na'\-vi\-yä ha\-pxì.} Your spirit goes with Eywa, your body stays behind to become part of the People.~\LINK{http://wiki.learnnavi.org/index.php?title=Na\%27vi_from_Avatar_Movie\#Jake.27s_final_kill}{a\slash wiki}
\B{Tsa\-ri\-a fkol po\-le\-'un fu\-ta Lo\-ak slu ta\-ron\-yu, sem\-pul 'efu ye.} Father is content that it's been decided Loak will be a hunter.~\LINK{http://naviteri.org/2011/04/\%E2\%80\%99a\%E2\%80\%99awa-li\%E2\%80\%99fyavi-amip\%E2\%80\%94a-few-new-expressions/}{b}
\B{Mì Si\-ä\-tll a nu\-mul\-txa\-ri lì'\-fya\-yä le\-Na'\-vi, sìl\-pey oe slì\-ye\-vu nga nu\-mul\-txa\-tu!} I hope you'll participate in the upcoming Na'vi class in Seattle!~\LINK{http://naviteri.org/2012/07/meetings-waterfalls-and-more/}{b}
\B{Oe\-yä ik\-ran sli\-vu nga, tsa\-krr oeng 'aw\-si\-teng mi\-vak\-to.} Be my ikran and let's ride together.~\LINK{http://wiki.learnnavi.org/index.php?title=Corpus\#Lemondrop.com}{ld}
(\emph{b. optionally with} \B{ne} \emph{to disambiguate the predicate of a sentence}) 
\B{Ta\-ron\-yu slu ne tsam\-si\-yu.} The hunter becomes a warrior.~\LINK{http://wiki.learnnavi.org/Canon/2010/UltxaAyharyu\%C3\%A4\#Word_Order_with_Slu}{wiki}
(\emph{c. with adj.s}) 
\B{Fì\-fne\-rìn ke ley krr\-a slu pay\-nga'.} This kind of wood is worthless when it gets damp.~\LINK{http://naviteri.org/2014/03/value-and-worth/}{b} 
\B{Ay\-nga\-ri teng\-krr ya wur sle\-ru nì\-'ul, sìl\-pey oe, li\-vu hel\-ku sang ul\-te te'\-lan le\-fpom.} As for you all, while it's getting cooler, I hope your homes may be warm and your hearts happy.~\LINK{http://naviteri.org/2013/09/on-si-salewfya-shapes-and-directions/}{b}
\B{Oeri key tun slolu nìtxan.} I'm so embarrassed [lit.: my face has become very red].~\LINK{http://naviteri.org/2011/04/\%E2\%80\%99a\%E2\%80\%99awa-li\%E2\%80\%99fyavi-amip\%E2\%80\%94a-few-new-expressions/comment-page-1/\#comment-694}{b}
(\emph{d. idiom}) \B{X nì\-Na'\-vi slu lì\-'u\-pe\slash pe\-lì\-'u?} How do you say X in Na'vi?~\LINK{http://naviteri.org/2010/09/quick-follow-up/}{b} 
(\emph{e. proverb}) \B{Tì\-tok slu tìk\-tok fa pam\-tsyìp a\-'aw.} Presence becomes absence by means of one little sound. [i.e. `Small things can make a big difference.']~\LINK{https://naviteri.org/2024/06/ayhapxi-tokxa-si-ayliu-alahe-parts-of-the-body-and-more/}{b} \\
(i.) \B{tsawl slu}, grow. \\
(ii.) \B{yawne slu}, fall in love. \B{Lu poe se\-vin nì\-ftxan kuma yaw\-ne slo\-lu oer.} She was so beautiful that I fell in love with her.~\LINK{http://naviteri.org/2012/06/spring-vocabulary-part-3/}{b} \\
(iii.) \B{muntxa slu}, get married. \B{Fko pll\-txe\-i\-e san me\-nga mun\-txa slo\-lu sìk!} I'm so happy to hear you got married!~\LINK{http://naviteri.org/2011/10/more-vocabulary-a-bit-of-grammar/}{b}  \\
(iv.) \B{slu leyr}, freeze. \B{Mì zì\-sì\-krr a\-txa\-wew slu ay\-ora leyr.} In the very cold season, the lakes freeze.~\LINK{http://naviteri.org/2018/07/ta-sulfatu-a-ayliu-niul-more-words-from-our-experts/}{b} \\
\textsc{Derivations}: \ding{43}
\on{1}.~\HL{sleyku}{\B{sley\-ku}}, produce, make. 
\on{2}.~\HL{zoslu}{\B{zo\-slu}}, heal, become well, get fixed. 
\on{3}.~\HL{tsunslu}{\B{tsun\-slu}}, may, be possible.

% slukx
\hi 
\HT{slukx}{\B{slukx}} \I{[sluk']} \emph{or} \I{[slUk']} \emph{n.} horn (\emph{of an animal}). 
\B{Na\-ri si! Tsa\-ioang\-ur lu pxi\-a me\-slukx!} Careful! That animal has two sharp horns!~\LINK{http://naviteri.org/2016/06/mrrvola-lifyavi-amip-forty-new-expressions/}{b} 

% smaw
\hi 
\HT{smaw}{\B{smaw}} \I{[smaw]} \emph{vtr.} approve of. 
\B{Fay\-hemit oel smaw nì\-wotx.} I completely approve of these actions.~\LINK{http://naviteri.org/2020/02/some-words-for-leap-year-day/}{b}
(\ding{214}~\HL{natxu}{\B{na\-txu}}) \\
\textsc{Derivations}: \ding{43}
\on{1}.~\HL{tIsmaw}{\B{tì\-smaw}}, approval.

% smaoe
\hi
\HT{smaoe}{\B{sma\cp{o}e}} \I{[sma."{}o.E]} \emph{n.}, \ding{96} phalanxia, heavy thorned plant, \emph{phalanxia ferox}. 

% smar
\hi
\HT{smar}{\B{smar}} \I{[smaR]} \emph{n.} prey, a hunted thing.
\B{Tì\-tu\-sa\-ron\-ìri fte fli\-vä, ze\-ne fko si\-vutx smar\-it ni\-nan nì\-no.} To succeed at hunting, you have to track your prey by reading (the forest) with attention to detail.~\LINK{http://naviteri.org/2011/09/\%E2\%80\%9Cby-the-way-what-are-you-reading\%E2\%80\%9D/}{b}
\B{Tsu'\-tey\-ìl to\-lìng oer mawl\-it smar\-ä.} Tsu'tey gave me a half of the prey.~\LINK{http://naviteri.org/2011/02/new-vocabulary-part-2/}{b}
(\emph{a. proverb}) (i.) \B{Ätxä\-le si pa\-lu\-lu\-kan\-ur tsnì smar\-it li\-vo\-nu.} \emph{a futile gesture, an attempt to achieve something that might be desirable but will clearly not happen} [lit.: ask a thanator to release its prey.]~\LINK{http://naviteri.org/2010/07/vocabulary-update/}{b}
(ii.) \B{Ta\-ron\-yut yom smar\-ìl.} \emph{A totally botched situation, chaos} [lit.: the prey eats the hunter].
\B{Eyk Ka\-mun a fra\-lo lä\-ngu tsa\-sä\-ta\-ron vel\-ke nì\-wotx. Ta\-ron\-yut yom smar\-ìl!} Every time Kamun is in charge, the hunt is a mess. Everything goes wrong that can!~\LINK{http://naviteri.org/2012/10/mipa-vospxi-mipa-ayliu-new-words-for-the-new-month/}{b}

% smon
\hi
\HT{smon}{\B{smon}} \I{[smon]} \emph{vin.} be familiar to, ``know s.o.''.
\B{Po smon oer.} I know him.~\LINK{http://naviteri.org/2010/09/getting-to-know-you-part-1/}{b}
\B{Sra\-ke smon ngar oe\-yä mey\-lan alu En\-tu sì Ka\-mun?} Do you know my friends En\-tu and Ka\-mun?~\LINK{http://naviteri.org/2010/09/getting-to-know-you-part-1/}{b} 
\B{Moe smon (moe\-ru) fì\-tsap.} We know each other.~\LINK{http://naviteri.org/2011/12/one-more-for-2011/}{b}
(\emph{a. conv.}) \B{Smon nì\-prr\-te'.} ``Nice to know you.''~\LINK{http://naviteri.org/2010/09/quick-follow-up/}{b}
(\emph{b. proverb}) \B{Ha\-haw nì\-'aw txo pa\-lu\-kan smi\-von ngar.} ``Only sleep if you are familiar with the Thanator,'' i.e. \emph{don't think you're safe unless you're aware of the danger}.~\LINK{http://naviteri.org/2021/08/ulte-ayyoratu-leiu-and-the-winners-are-2/}{b}

% sna-
\hi
\HT{sna}{\B{sna-}} \I{[sna.]} \on{1}. \emph{prod. pref. to describe a naturally occurring group or collection of animals or plants.} 
\B{sna\-ta\-li\-o\-ang}, a herd of sturmbeest. 
\on{2}. \emph{unprod. pref. to describe a collection of n.s denoting non\hyp{}living things.} \B{sna\-txä\-rem}, skeleton.
(\ding{43}~\HL{sna'o}{\B{sna\-'o}})

% sna'o
\hi
\HT{sna'o}{\B{\cp{sna}'o}} \I{["{}sna.Po]} \emph{n.} set, group, pile, clump, stand.   
\B{Ay\-skxe a mì sa\-sna\-'o ku'\-up lu nì\-txan.} The rocks in that pile are very heavy.~\LINK{http://naviteri.org/2012/03/spring-vocabulary-part-2/}{b}
\B{Hay\-ku\-ri sna\-'o ay\-nge\-yä lor nì\-txan lu nang!} What a beautiful collection of haiku!~\LINK{http://naviteri.org/2012/10/ulte-sawasultsyipa-yoratu-leiu-and-the-winner-is/}{b}

% snafpìlfya
\hi
\HT{snafpIlfya}{\B{sna\cp{fpìl}fya}} \I{[sna."fpIl.fja]} \emph{n.} philosophy.
\B{Sna\-fpìl\-fya\-ri le\-Na'\-vi krra smar\-it fkol tspang, tsran\-ten nì\-txan fwa po ke ngä\-'än nì\-kel\-kin.} It's important in Na'vi philosophy that the prey not suffer unnecessarily when it's killed.~\LINK{http://naviteri.org/2013/02/vospxi-ayol-posti-apup-short-post-for-a-short-month/}{b}
(\ding{43}~\HL{sna}{\B{sna-}}, \HL{fpIlfya}{\B{fpìl\-fya}})

% snanumultxa
\hi
\HT{snanumltxa}{\B{snanumul\cp{txa}}} \I{[sna.nu.mul."t'a]} \emph{or} \I{[.mUl.]} (\textsc{rn}: \B{snanumùl\cp{txa}} \I{[sna.nu.mUl."da]}) \emph{n.} course (\emph{of lectures}). 
(\ding{43}~\HL{sna}{\B{sna-}}, \HL{numultxa}{\B{nu\-mul\-txa}})

% snapamrelvi
\hi
\HT{snapamrelvi}{\B{snapam\cp{rel}vi}} \I{[sna.pam."REl.vi]} \emph{n.} alphabet. 
\B{Sna\-pam\-rel\-vi sì Lì\-'u\-pam}, Alphabet and Pronunciation.~\LINK{http://naviteri.org/2012/10/audio-and-video-learning-materials-for-navi-101/}{b}
(\ding{43}~\HL{sna}{\B{sna-}}, \HL{pamrelvi}{\B{pam\-rel\-vi}})

% snatanhì
\hi
\HT{snatanhI}{\B{snatan\cp{hì}}} \I{[sna.tan."hI]} \emph{n., astr.} constellation (\emph{of stars}).
(\ding{43}~\HL{sna}{\B{sna-}}, \HL{tanhI}{\B{tan\-hì}})

% snatanhìtsyìp
\hi
\HT{snatanhItsyIp}{\B{snatan\cp{hì}tsyìp}} \I{[sna.tan."hI.\t{ts}jIp\textcorner]} \emph{n., astr.} star cluster. 
(\ding{43}~\HL{snatanhI}{\B{sna\-tan\-hì}}, \HL{tsyIp}{\B{-tsyìp}})

% snatxärem
\hi
\HT{snatxArem}{\B{sna\cp{txä}rem}} \I{[sna."t'\ae.REm]} \emph{n.} skeleton.
(\ding{43}~\HL{sna}{\mbox{\B{sna-}}}, \HL{txArem}{\B{txä\-rem}})

% snaytu
\hi
\HT{snaytu}{\B{\cp{snay}tu}} \I{["{}snaj.tu]} \emph{n.} loser.   
\B{Fra\-u\-van\-ìri lu yo\-ra'\-tu, lu snay\-tu.} For every game, there's a winner and a loser.~\LINK{http://naviteri.org/2011/12/a-note-on-the-word-\%E2\%80\%9Cyora\%E2\%80\%99tu\%E2\%80\%9D/}{b}
(\ding{43}~\HL{snaytx}{\B{snaytx}})

% snaytx
\hi
\HT{snaytx}{\B{snaytx}} \I{[snajt']} \emph{vtr.} lose (\emph{a game}).   
\B{Trr\-am\-ä ay\-u\-van\-ìri mak\-to Ak\-wey nì\-yew\-la ha snaytx.} Ak\-wey rode disappointingly in yesterday's games, so he lost.~\LINK{http://naviteri.org/2012/10/mipa-vospxi-mipa-ayliu-new-words-for-the-new-month/}{b}
\B{Ay\-oe sno\-laytx fì\-trr ta\-lun\-a oel rum\-it to\-lung\-zup.} We lost today because I dropped the ball.~\LINK{http://naviteri.org/2011/07/txantsana-ultxa-mi-siatll-great-meeting-in-seattle/}{b} 
\B{Oe pll\-ngay san mo\-lak\-to oe nì\-fe', ta\-fral sno\-laytx; wä\-tu lu oe\-to txur.} I acknowledge that I rode badly, so I lost; my opponent was stronger than I was.~\LINK{http://naviteri.org/2011/07/txantsana-ultxa-mi-siatll-great-meeting-in-seattle/}{b} \\
\textsc{Derivations}: \ding{43}
\on{1}.~\HL{snaytu}{\B{snay\-tu}}, loser.

% snayì
\hi
\HT{snayI}{\B{sna\cp{yì}}} \I{[sna."jI]} \emph{n.} (natural) staircase, series of step\hyp{}like levels. 
\B{Kll\-zo\-la\-'u Mo\-'at fa sna\-yì teng\-krr pe\-rll\-txe san Ay\-nga ne\-to ri\-vikx!} Mo'at came down the stairway saying, ``Get back, all of you!''~\LINK{http://naviteri.org/2013/01/lifyari-po-peyi-how-good-is-her-navi/}{b}
(\ding{43}~\HL{sna}{\B{sna-}}, \HL{yI}{\B{yì}})

% snäm
\hi
\HT{snAm}{\B{snäm}} \I{[sn\ae{}m]} \emph{vin.} rot, decay, degrade over time (\emph{of an object or s.t. abstract}). 
\B{Fì\-tsngan\-ur a sno\-läm lä\-ngu fa\-hew a\-kxä\-näng.} This rotten meat has a putrid smell.~\LINK{http://naviteri.org/2015/11/vomuna-liu-amip-ten-new-words/}{b}
\B{Ze\-ne fko tsko swi\-zaw\-it si\-var nì\-trr\-trr fte\-ke fì\-tsu\-'o sni\-väm.} One must use a bow and arrow regularly to prevent this ability degrading over time.~\LINK{http://naviteri.org/2015/11/vomuna-liu-amip-ten-new-words/}{b}

% snep
\hi
\HT{snep}{\B{snep}} \I{[snEp\textcorner]} \emph{adj.} narrow.
(\ding{214}~\HL{sloa}{\B{sloa}}) \\
\textsc{Derivations}: \ding{43}
\on{1}.~\HL{slosnep}{\B{slo\-snep}}, width.

% snew
\hi
\HT{snew}{\B{snew}} \I{[snEw]} \emph{vtr.} constrict, tighten. 
(\emph{a. idiom})
\B{Ke tsun fko tspi\-vang to\-ruk\-it fa fwa pewn\-ti snew.} \emph{Proverbial expression for a method that will not work.} [lit.: you can't kill a great leonopteryx by constricting its throat]~\LINK{http://naviteri.org/2012/10/snewsyea-ftxoza-halowina-a-spooky-halloween/}{b} 
(\emph{shortened to}) \B{pewn to\-ruk\-ä}. \\
\textsc{Derivations}: \ding{43}
\on{1}.~\HL{snewsye}{\B{snew\-sye}}, weird, spooky.

% snewsye
\hi
\HT{snewsye}{\B{\cp{snew}sye}} \I{["{}snEw.sjE]} \emph{adj.} weird, spooky.   
\B{Ay\-nga\-ri ftxo\-zä Hä\-lo\-win\-ä li\-vu 'o' sì snew\-sye txan\-txew\-vay.} May your Halloween be as fun and spooky as possible.~\LINK{http://naviteri.org/2012/10/snewsyea-ftxoza-halowina-a-spooky-halloween/}{b}
\B{Oe\-ri lu fì\-tseng snew\-sye nì\-hawng. Hi\-vum ko.} This place is too spooky for me. Let's get out of here.~\LINK{http://naviteri.org/2012/10/snewsyea-ftxoza-halowina-a-spooky-halloween/}{b}
(\ding{43}~\HL{snew}{\B{snew}}, \HL{syeha}{\B{sye\-ha}})

% sno
\hi
\HT{sno}{\B{sno}} \I{[sno]} (\emph{gen.} \ding{43}~\HL{sneyA}{\B{\cp{sne}yä}}) \emph{pn. 3\textsuperscript{rd} person to mean} ``his\slash her\slash its\slash their self'' or ``his\slash her\slash their own''; \emph{indifferent to number; most often used in its genitive form}.
(\emph{a. subjective})
\B{Nì\-trr\-trr pxì\-mun\-'i sam\-si\-yul ay\-swi\-zaw\-it ku\-tu\-ä a\-law\-nä\-txayn sno\-kip nì\-'eng.} Warriors typically share the arrows of their defeated enemies among themselves.~\LINK{http://naviteri.org/2012/01/more-additions-to-the-lexicon/}{b}
(\emph{b. genitive})
\B{Pol 'o\-lem sne\-yä wu\-tsot.} He made his (own) dinner.~\LINK{http://naviteri.org/2011/12/one-more-for-2011/}{b}
\B{Sa'\-nok prr\-nen\-ur yom\-tìng fa okup sne\-yä.} A mother feeds an infant with her milk.~\LINK{http://naviteri.org/2013/06/looking-back-looking-forward/}{b}
(\emph{c. dative})
\B{Po yaw\-ne lu snor.} He loves himself.~\LINK{http://naviteri.org/2011/12/one-more-for-2011/}{b}
\B{Me\-fo yaw\-ne lu (snor) fì\-tsap.} They (= those two) love each other.
\B{Fo smon (sno\-ru) fì\-tsap nì\-wotx.} They all know each other.~\LINK{http://naviteri.org/2011/12/one-more-for-2011/}{b} \\
\textsc{Derivations}: \ding{43}
\on{1}.~\HL{snolup}{\B{sno\-lup}}, personal style. 
\on{2}.~\HL{snomo}{\B{sno\-mo}}, private space, room.
\on{3}.~\HL{snonrra}{\B{sno\-nrra}}, self-pride. 
\on{4}.~\HL{snokfyan}{\B{snok\-fyan}}, personal belongings rack.
\on{5}.~\HL{snonivi}{\B{sno\-ni\-vi}}, single\hyp{}person hammock.
\on{6}.~\HL{snotipx}{\B{sno\-tipx}}, self\hyp{}control.

% snokfyan
\hi
\HT{snokfyan}{\B{\cp{snok}fyan}} \I{["snok\textcorner.fjan]} \emph{n.} personal belongings rack.
(\ding{43}~\HL{sno}{\B{sno}}, \HL{kurfyan}{\B{kur\-fyan}})

% snolup
\hi
\HT{snolup}{\B{\cp{sno}lup}} \I{["sno.lup\textcorner]} \emph{or} \I{[.lUp\textcorner]} (\textsc{rn}: \B{\cp{sno}lùp} \I{["sno.lUp\textcorner]}) \emph{n.} personal style or aesthetic, presence.
\B{Por lu sno\-lup a new fra\-po rì\-'ìr si\-vi.} He\slash She has a personal style that everyone wants to emulate.~\LINK{http://naviteri.org/2014/07/tsawlultxamaw-a-ayliu-post-meetup-words/}{b}
(\ding{43}~\HL{sno}{\B{sno}}, \HL{lupra}{\B{lup\-ra}})

% snomo
\hi
\HT{snomo}{\B{\cp{sno}mo}} \I{["{}sno.mo]} \emph{n.} private space that one can retreat to. 
\B{Nga\-ri sno\-mot krr\-pe ve\-zey\-ko, ma 'itan?} When are you going to organize your room, son?~\LINK{http://naviteri.org/2012/10/mipa-vospxi-mipa-ayliu-new-words-for-the-new-month/}{b}
(i.) \B{(sno)\-mo a ha\-haw}, bedroom.~\LINK{http://naviteri.org/2012/10/mipa-vospxi-mipa-ayliu-new-words-for-the-new-month/}{b}
(\ding{43}~\HL{sno}{\B{sno}}, \HL{mo}{\B{mo}})

% snonivi
\hi
\HT{snonivi}{\B{\cp{sno}nivi}} \I{["{}sno.ni.vi]} \emph{n.} single\hyp{}person hammock.
\B{Ze\-ne aw\-nga tì\-ngä\-zìk\-ur sno\-ni\-vi\-yä tì\-ke\-zin si\-vi nì\-win.} We need to solve the hammock problem quickly.~\LINK{http://naviteri.org/2021/04/mipa-ayliu-mipa-sioeykting-new-words-new-explanations/\#comment-33483}{b}
(\ding{43}~\HL{sno}{\B{sno}}, \HL{nivi}{\B{ni\-vi}})

% snonrra
\hi
\HT{snonrra}{\B{sno\cp{nrr}a}} \I{[sno."n\s{r}.a]} \emph{n.} self\hyp{}pride (\emph{negative connotation}), arrogance.
(\ding{43}~\B{sno}, \B{nrr\-a}) \\
\textsc{Derivations}: \ding{43}
\on{1}.~\HL{lesnonrra}{\B{le\-sno\-nrra}}, full of self\hyp{}pride.

% snotipx
\hi 
\HT{snotipx}{\B{sno\cp{tipx}}} \I{[sno."tip']} \emph{n.} self\hyp{}control. 
\B{Ke fkey\-tok tì\-eyk\-tan a\-tì\-flä\-nga' lu\-ke sno\-tipx.} Successful leadership does not exist without self-control.~\LINK{http://naviteri.org/2018/04/100a-liu-amip-64-new-words-part-2/}{b} 
(\ding{43}~\HL{sno}{\B{sno}}, \HL{tIhipx}{\B{tì\-hipx}})

% Snrrtsyìp
\hi 
\HT{snrrtsyIp}{\B{\cp{Snrr}tsyìp}} \I{["s\s{r}.\t{ts}jIp\textcorner]} \emph{n., astr., short form of} \ding{43}~\HL{tawsnrrtsyIp}{\B{Tawsnrrtsyìp}}.

% snumìna
\hi
\HT{snumIna}{\B{\cp{snu}mìna}} \I{["snu.mI.na]} (\textsc{rn}: \B{\cp{snù}mìna} \I{["snU.mI.na]}) \emph{adj.} dim (\emph{of a person}), stupid. 
\B{Oe\-ru po snu\-mìn\-a la\-tsam.} He seems dim to me.~\LINK{https://wiki.learnnavi.org/Na\%27vi_from_Avatar_Movie\#Wutso_in_Hometree}{a\textsuperscript{c}}
\B{snu\-mì\-na skxawng\slash skxawng a\-snu\-mì\-na}, stupid moron.~\LINK{http://wiki.learnnavi.org/index.php?title=Corpus\#Good_Morning_America}{wiki}

% sngap
\hi
\HT{sngap}{\B{sngap}} \I{[sNap\textcorner]} \emph{vtr.} sting. 
(\emph{a. idiom}) \B{sre fwa sngap zi\-ze'}, ``as quickly as possible'' [lit: before the hellfire wasp stings].~\LINK{http://naviteri.org/2012/07/fivospxiya-aylifyavi-amip-this-months-new-expressions-pt-2/}{b}

% sngä'i
\hi
\HT{sngA'i}{\B{\cp{sngä}'i}} \I{["{}sN\ae.Pi]} (\emph{a.}) \emph{vin.} begin, start, s.t. begins\slash starts.   
\B{Ye'\-rìn sngä\-'i tom\-pa.} Soon the rain starts.~\LINK{http://forum.learnnavi.org/language-updates/learnings-from-the-siatlla-song-translations/msg475008/\#msg475008}{ln}
\B{Aw\-nga sngi\-vä\-'i ko!} Let's begin!~\LINK{http://naviteri.org/2010/06/first-post/}{b}
\B{Tì\-kang\-kem sngo\-lä\-'i.} The work began.~\LINK{http://naviteri.org/2010/06/first-post/}{b}
(\emph{b. modal}) start, begin to do s.t. 
\B{Na'\-rìng sngä\-'i ni\-vrr txon\-krr syu\-ra\-tan\-fa.} The forest begins to glow at night with bioluminescence.~\LINK{http://naviteri.org/2012/07/fivospxiya-aylifyavi-amip-this-months-new-expressions-pt-2/}{b}
\B{Tskot sngo\-lä\-'i po si\-var 'a\-'aw\-a trr\-kam.} He started to use the bow several days ago.~\LINK{http://naviteri.org/2011/09/miscellaneous-vocabulary/}{b}
(\ding{43}~\HL{tsyul}{\B{tsyul}}, \ding{214}~\HL{'i'a}{\B{'i'\-a}}) \\  
\textsc{Derivations}:
\on{1}.~\HL{sngA'ikrr}{\B{sngä\-'i\-krr}}, beginning (time).  
\on{2}.~\HL{sngA'itseng}{\B{sngä\.'i\-tseng}}, beginning, starting position.  
\on{3}.~\HL{sngA'iyu}{\B{sngä\-'i\-yu}}, beginner. 
\on{4}.~\HL{nIsngA'i}{\B{nì\-sngä\-'i}}, originally, at first. 
\on{5}.~\HL{lesngA'i}{\B{le\-sngä\-'i}}, original, first in a series.

% sngä'ikrr
\hi
\HT{sngA'ikrr}{\B{\cp{sngä}'ikrr}} \I{["{}sN\ae.Pi.k\s{r}]} \textsc{i}. \emph{n.} beginning, start time. \\
\textsc{ii}. \emph{adv.} in the beginning, at the start.
\B{Txe\-wì\-yä sem\-pul\-ìri sngä\-'i\-krr Saw\-tu\-tet ve\-re'\-kì.} In the beginning, Txe\-wì's father hated the Skypeople.~\LINK{http://naviteri.org/2010/07/ting-mikyun-fte-tslivam-listening-comprehension-1-3/}{b}
\B{Tsa\-lì'\-fya\-vi\-ri, sngä\-'i\-krr lo\-lu am\-'a\-lu\-ke aw\-nge\-yä 'ey\-lan alu Tsyeyk tsa\-fne\-tu\-te.} As for that expression, our friend Jake was doubtless that kind of person in the beginning.~\LINK{http://naviteri.org/2019/08/aawa-lifya-si-lifyavi-amip-a-few-new-words-and-expressions/\#comment-30289}{b}
(\ding{43}~\HL{sngA'i}{\B{sngä\-'i}}, \HL{krr}{\B{krr}}, \ding{214}~\HL{tI'i'a}{\B{tì\-'i'a}})

% sngä'itseng
\hi
\HT{sngA'itseng}{\B{\cp{sngä}'itseng}} \I{["{}sN\ae.Pi.\t{ts}EN]} \emph{n.} beginning, starting position, initial location.
\B{Ro sngä\-'i\-tseng tsa\-lì\-'u\-ä alu 'ey\-lan lu tì\-ftang.} At the beginning of the word \emph{'eylan} there's a glottal stop.~\LINK{http://naviteri.org/2012/07/fivospxiya-aylifyavi-amip-this-months-new-expressions-pt-2/}{b}
(\ding{43}~\HL{sngA'i}{\B{sngä\-'i}}, \HL{tseng}{\B{tse\-nge}})

% sngä'iyu
\hi
\HT{sngA'iyu}{\B{\cp{sngä}'iyu}} \I{["{}sN\ae.Pi.ju]} \emph{n.} beginner.   
\B{Fra\-po, ftxey sngä\-'i\-yu ftxey tsul\-fä\-tu, tsì\-ye\-vun fì\-tse\-nge ri\-vun 'u\-ot le\-sar.} Everyone, whether beginner or master, can find something useful here.~\LINK{http://naviteri.org/2010/06/first-post/}{b} 
\B{num\-tseng\-vi ay\-sngä\-'i\-yu\-ä}, beginners' classroom.~\LINK{http://naviteri.org/2010/06/first-post/}{b}
\B{Sìn yì sngä\-'i\-yu\-ä, slä tsye\-rìl (hay\-a yì\-ne) nì\-'ul\-'ul.} They're at a beginner's level, but they're getting better and better.~\LINK{http://naviteri.org/2013/01/lifyari-po-peyi-how-good-is-her-navi/}{b}
\B{Nì\-aw\-no\-mum ngo\-lop ay\-oel fì\-nu\-mul\-txa\-ti fpi sngä\-'i\-yu.} As you know, we created this class for the beginner.~\LINK{http://naviteri.org/2012/07/tskxekengtsyip-a-mikyunfpi-a-little-listening-exercise/}{b}
(\ding{43}~\HL{sngA'i}{\B{sngä\-'i}}, \ding{214}~\HL{tsulfAtu}{\B{tsul\-fä\-tu}})

% sngeykä'i
\hi
\HT{sngeykA'i}{\B{sngey\cp{kä}'i}} \I{[sNEj."k$\cdot$\ae.P$\cdot$i]} \emph{vtr.} begin, start s.t.
\B{Poel tì\-kang\-kem\-it sngey\-ko\-lä\-'i.} She began the work.~\LINK{http://naviteri.org/2010/06/first-post/}{b} 
\B{Fwa oel fì\-tì\-kang\-kem\-vit sngey\-ki\-vä\-'i krr\-no\-lekx nì\-txan.} Me starting this project took a long time.~\LINK{http://naviteri.org/2010/06/first-post/}{b} 
(\ding{43}~\HL{sngA'i}{\B{sngä\-'i}}; \HL{tsyul}{\B{tsyul}})

% sngel
\hi
\HT{sngel}{\B{sngel}} \I{[sNEl]} \emph{countable n.} garbage, trash. 
\B{Oel vewng ay\-sngel\-it.} I tend to the refuse.~\LINK{http://naviteri.org/2010/09/getting-to-know-you-part-2/}{b}
\B{Ngal kaw\-krr sngel\-it ke mu\-nge wrr\-pa\-ne.} You never take the trash out.~\LINK{http://forum.learnnavi.org/language-updates/txelanit-hivawl/}{ln}
\B{Pey\-ral\-ìl zet wur\-a wu\-tsot a\-'aw\-nem pxel sngel.} Peyral treats a cool cooked meal like garbage.~\LINK{http://naviteri.org/2012/10/mipa-vospxi-mipa-ayliu-new-words-for-the-new-month/}{b} \\
\textsc{Derivations}: \ding{43}
\on{1}.~\HL{sngeltseng}{\B{sngel\-tseng}}, garbage dump.

% sngeltseng
\hi
\HT{sngeltseng}{\B{\cp{sngel}tseng}} \I{["{}sNEl.\t{ts}EN]} \emph{n.} garbage dump.
\B{Taw\-sìp nge\-yä lu sngel\-tseng.} Your ship is a garbage scow.~\LINK{http://wiki.learnnavi.org/index.php?title=Corpus\#UGO_Movie_Blog_.28Twitter_Questions.29}{ugo\slash wiki}
(\ding{43}~\HL{sngel}{\B{sngel}}, \HL{tseng}{\B{tse\-nge}})

% sngukx
\hi
\HT{sngukx}{\B{sngukx}} \I{[sNuk']} \emph{or} \I{[sNUk']} \emph{n.}, \ding{96} grub plant, \emph{limacina erecta}. 

% sngum
\hi
\HT{sngum}{\B{sngum}} \I{[sNum]} \emph{or} \I{[sNUm]} \textsc{i}. \emph{n.} worry.   
\B{Lu oe\-ru sngum a sa\-ron\-yu ke tì\-ye\-vä\-txaw.} I'm worried that the hunters will not return (soon, as expected).~\LINK{http://naviteri.org/2011/01/new-year-new-vocabulary/}{b} \\
\textsc{ii}. \B{sngum si} \I{[sNum si]} \emph{vin.} worry.
\B{Sngum rä\-'ä si.} Don't worry.~\LINK{http://naviteri.org/2014/09/stxeli-alor-the-text/comment-page-1/\#comment-10619}{b}
\B{Fu\-ri\-a po ke za\-ya\-'u oe sngum si.} I'm worried that he won't come.~\LINK{http://naviteri.org/2014/09/stxeli-alor-the-text/comment-page-1/\#comment-10734}{b} \\
\textsc{Derivations}: \ding{43}
\on{1}.~\HL{sngumtsim}{\B{sngum\-tsim}}, worrisome matter. 
\on{2}.~\HL{sngunga'}{\B{sngu\-nga'}}, worrisome, troubling. 
\on{3}.~\HL{txansngum}{\B{txan\-sngum}}, desperation.
 
% sngumtsim
\hi
\HT{sngumtsim}{\B{\cp{sngum}tsim}} \I{["{}sNum.\t{ts}im]} \emph{or} \I{["{}sNUm.]} \emph{n.} worrisome matter, source of worry.
\B{Lu oer sngum\-tsim fwa po ke zo\-la\-'u.} I have a feeling of worry that he didn't come.~\LINK{http://naviteri.org/2013/01/awvea-posti-zisita-amip-first-post-of-the-new-year/}{b}
\B{Pe\-yä ay\-lì\-'u a\-ing\-ye\-nga' lo\-lu sngum\-tsim ay\-oe\-ru nì\-wotx.} His mysterious words worried us all.~\LINK{http://naviteri.org/2013/01/awvea-posti-zisita-amip-first-post-of-the-new-year/}{b}
\B{Syor, ma 'ey\-lan, syor. Ke lu kea sngum\-tsim.} Relax, friend, relax. There's nothing to worry about.~\LINK{http://naviteri.org/2013/01/awvea-posti-zisita-amip-first-post-of-the-new-year/}{b}
(\ding{43}~\HL{sngum}{\B{sgnum}}, \HL{tsim}{\B{tsim}})

% sngunga'
\hi
\HT{sngunga'}{\B{\cp{sngu}nga'}} \I{["{}sNu.NaP]} \emph{adj.} worrisome, troubling.  
(\emph{a. proverb}) \B{Hem a\-sru\-nga' nì\-'ul, hum a\-sngu\-nga' nì\-nän.}
More helpful actions lead to less troubling outcomes.~\LINK{http://naviteri.org/2015/03/kap-si-ayunil-saylahe-threats-dreams-and-other-things/}{b}
(\ding{43}~\HL{sngum}{\B{sngum}}, \HL{-nga'}{\B{-nga'}})
 
% so'ha
\hi
\HT{so'ha}{\B{\cp{so'}ha}} \I{["{}soP.ha]} \emph{vtr.} (\emph{a.}) be enthusiastic about, show enthusiasm for, be excited about. 
\B{Oel so'\-ha fu\-ta trr\-ay nga\-hu kä.} I'm excited about going with you tomorrow.~\LINK{http://naviteri.org/2012/06/spring-vocabulary-part-3/}{b}
\B{Tse\-nul so'\-ha tey\-lut nì\-hawng nì\-'it, ke\-fyak?} Don't you think Tsenu is a bit too into \emph{teylu}?~\LINK{http://naviteri.org/2012/06/spring-vocabulary-part-3/}{b}
\B{Saw\-no'\-ha ioit ko\-la\-ne\-i\-om oel u\-van\-fa!} I got (won) the prized piece of jewelry in the game!~\LINK{http://naviteri.org/2012/06/spring-vocabulary-part-3/}{b}
(\emph{b. conv., intj.}) \B{Le\-iam fwa Txe\-wì Ri\-ni\-sì mun\-txa slì\-yu.---So'\-ha!} It looks like Txewì and Rini are about to get married.---That's great!~\LINK{http://naviteri.org/2012/06/spring-vocabulary-part-3/}{b} \\
\textsc{Derivations}: \ding{43}
\on{1}.~\HL{tIso'ha}{\B{tì\-so'\-ha}}, enthusiasm, having a good attitude. 
\on{2}.~\HL{leso'ha}{\B{le\-so'\-ha}}, enthusiastic, keen. 
\on{3}.~\HL{nIso'ha}{\B{nì\-so'\-ha}}, enthusiastically.
\on{4}.~\HL{so'hayu}{\B{so'\-ha\-yu}}, enthusiast.

% so'hayu
\hi
\HT{so'hayu}{\B{\cp{so'}hayu}} \I{["{}soP.ha.ju]} \emph{n.} enthusiast.
(\ding{43}~\HL{so'ha}{\B{so'\-ha}}) \\
\textsc{Derivations}: \ding{43}
\on{1}.~\HL{txanso'hayu}{\B{txanso'hayu}}, fan, enthusiast.

% so'yu
\hi 
\HT{so'yu}{\B{\cp{so'}yu}} \I{["{}soP.ju]} \emph{n.} fan, enthusiast; (\emph{coll. form of} \ding{43}~\HL{txanso'hayu}{\B{txan\-so'\-ha\-yu}}). 

% soaia
\hi
\HT{soaia}{\B{so\cp{a}ia}} \I{[so."{}a.i.a]} (\emph{gen.} \B{soaiä}) \emph{n.} family.
\B{Soaia lu hì\-'i -- nì\-'aw sem\-pul, sa'\-nok, tsmu\-kan sì Txe\-wì.} The family is small -- only father, mother, a brother and Txewì.~\LINK{http://naviteri.org/2010/07/ting-mikyun-fte-tslivam-listening-comprehension-1-3/}{b}
\B{Mu\-nge fna\-we'\-tul lewng\-it soaia\-ru sne\-yä.} A coward brings shame to his\slash her family.~\LINK{http://naviteri.org/2021/09/aawa-ayliu-amip-a-few-new-words/}{b}
\B{Va'\-rul zä\-nget ik\-ran\-it sn\-eyä pxel ha\-pxì\-tu so\-a\-i\-ä.} Va'ru treats his ikran like a member of the family.~\LINK{http://naviteri.org/2012/10/mipa-vospxi-mipa-ayliu-new-words-for-the-new-month/}{b} 
\B{Mì sray apxa kel\-ku si so\-a\-i\-a\-hu.} (He) lives with his family in a large village.~\LINK{http://naviteri.org/2010/07/ting-mikyun-fte-tslivam-listening-comprehension-1-3/}{b} \\
\textsc{Derivations}: \ding{43}
\on{1}.~\HL{nIsoaia}{\B{nì\-soaia}}, (together) as members of a family. 
\on{2}.~\HL{swaynivi}{\B{sway\-ni\-vi}}, family hammock.
\on{3}.~\HL{swapxI}{\B{swa\-pxì}}, family member.

% sok
\hi
\HT{sok}{\B{sok}} \I{[sok\textcorner]} \emph{adj.} recent.
\B{Ngay\-txo\-a, nì\-aw\-no\-mum ke lo\-lu oer nì\-ke\-ftxo mì sok\-a srr ay\-skxom le\-tam fte lì'\-fya\-ri aw\-nge\-yä tì\-kang\-kem si\-vi.} My apologies: As you know, in recent days I have not had sufficient opportunity to work on our language.~\LINK{http://forum.learnnavi.org/language-updates/trr-rrtaya/}{ln}
\B{Sìn A\-sok, Sìn Zu\-saw\-krr\-ä}, Recent and Upcoming Activities [lit.: recent activities, future's activities].~\LINK{http://naviteri.org/2013/04/sin-asok-sin-zusawkrra-recent-and-upcoming-activities/}{b}
\on{1}.~\HL{nIsok}{\B{nì\-sok}}, recently.

% solew
\hi
\HT{solew}{\B{so\cp{lew}}} \I{[so."lEw]}~:~\HL{salew}{\B{salew}}.

% som
\hi
\HT{som}{\B{som}} \I{[som]} \emph{adj.} hot. 
\B{Ya\-ri som\-wew\-pe (set)?---Ya lu som.} How's the temperature (now)?---It's hot.~\LINK{http://naviteri.org/2012/10/audio-and-video-learning-materials-for-navi-101/}{b}
(\ding{214}~\B{wew}) \\
\textsc{Derivations}: \ding{43}
\on{1}.~\HL{tIsom}{\B{tì\-som}}, heat. 
\on{2}.~\HL{txasom}{\B{txa\-som}}, very hot. 
\on{3}.~\HL{somwew}{\B{som\-wew}}, temperature. 
\on{4}.~\HL{somtIlor}{\B{som\-tì\-lor}}, popsicle, ``hot beauty''. 
\on{5}.~\HL{zIskrrsom}{\B{zì\-skrr\-som}}, summer.

% somtìlor
\hi
\HT{somtIlor}{\B{\cp{som}tìlor}} \I{["{}som.tI.loR]} \emph{n.}, \ding{96} popsicle, ``hot beauty'',
\emph{capsulatum virgatum}.
(\ding{43}~\HL{som}{\B{som}}, \HL{tIlor}{\B{tì\-lor}})

% somwew
\hi
\HT{somwew}{\B{som\cp{wew}}} \I{[som."wEw]} \textsc{i}. \emph{n.} temperature. 
\B{Som\-wew\-ä tì\-'ul a ka 'Rr\-ta sä\-peng\-hrr lu aw\-nga\-ru nì\-wotx.} Global warming (lit.: the increase in temperature across the Earth) is a warning to us all.~\LINK{http://naviteri.org/2016/06/mrrvola-lifyavi-amip-forty-new-expressions/}{b} \\
\textsc{ii}. \B{som\cp{wew}pe} \I{[som."{}wEw.pE]} \emph{or} \B{pe\-som\-\cp{wew}} \I{[pE.som."{}wEw]} \emph{inter.} what temperature? ``How cold\slash warm?''
\B{Ya\-ri som\-wew\-pe (set)?---Ya lu sang.} What's the temperature?---It's warm.~\LINK{http://naviteri.org/2012/10/audio-and-video-learning-materials-for-navi-101/}{b}
\B{Pay\-ri som\-wew\-pe?} What's the temperature of the water?~\LINK{http://naviteri.org/2011/05/weather-part-2-and-a-bit-more-2/}{b}
(\ding{43}~\HL{som}{\B{som}}, \HL{wew}{\B{wew}})

% sop
\hi
\HT{sop}{\B{sop}} \I{[sop\textcorner]} \emph{vin.} travel. 
\B{Tsun oe nga\-hu tsa\-tse\-nge\-ne ki\-vä, slä nul\-new fu\-ta si\-vop oeng nì\-txi\-lu\-ke.} I can go there with you, but I prefer to travel leisurely.~\LINK{http://naviteri.org/2011/10/more-vocabulary-a-bit-of-grammar/}{b}
\B{Ay\-fo so\-lop ìlä hil\-van fa u\-ran.} They traveled along (up, down) the river by boat.~\LINK{http://naviteri.org/2011/07/txantsana-ultxa-mi-siatll-great-meeting-in-seattle/}{b}
(\emph{a. idiom}) \B{Si\-vop nì\-zaw\-nong.} Travel safely.~\LINK{http://naviteri.org/2012/10/audio-and-video-learning-materials-for-navi-101/}{b} \\
\textsc{Derivations}: \ding{43}
\on{1}.~\HL{tIsop}{\B{tì\-sop}}, journey. 
\on{2}.~\HL{sopyu}{\B{sop\-yu}}, traveller.

% sopyu
\hi
\HT{sopyu}{\B{\cp{sop}yu}} \I{["sop\textcorner.ju]} \emph{n.} traveller.
\B{Si\-vop nì\-zaw\-nong, ma ay\-sop\-yu.} Travel safely, travellers.~\LINK{http://naviteri.org/2010/09/quick-follow-up/}{b}
(\ding{43}~\HL{sop}{\B{sop}})

% Sorewn
\hi
\B{So\cp{rewn}} \I{[so."REwn]} \emph{n.} \emph{female name}. 
\B{Nga lu So\-rewn, ke\-fyak?} You're Sorewn, right?~\LINK{http://naviteri.org/2010/09/getting-to-know-you-part-1/}{b}
\B{So\-rewn\-ìl kan\-'ìn tì\-'em\-it nì\-txan ul\-te kxawm tsa\-txe\-le me\-nga\-ne za\-'a\-tsu nì\-'eng.} Sorewn is really into cooking and you two might share that in common.~\LINK{http://naviteri.org/2010/09/getting-to-know-you-part-1/}{b}
\B{Mun\-txa\-tu\-ri So\-rewn\-ti ke tsun oe mi\-vll\-'an.} I can't accept Sorewn as my spouse.~\LINK{http://naviteri.org/2012/10/mipa-vospxi-mipa-ayliu-new-words-for-the-new-month/}{b}

% sosul
\hi
\HT{sosul}{\B{so\cp{sul}}} \I{[so."sul]} \emph{or} \I{[."sUl]} (\textsc{rn}: \B{so\cp{sùl}} \I{[so."sUl]}) \emph{adj.} pleasant smell (\emph{of nearby running water, rain, moist vegetation}).

% spakat
\hi 
\HT{spakat}{\B{\cp{spa}kat}} \I{["spa.kat\textcorner]} \emph{vin.} drown. 
\B{Fì\-ioang ke tsun sli\-ve\-le; nem\-fa pay zup tìk spa\-kat.} This animal cannot swim; if it falls into the water, it immediately drowns.~\LINK{http://naviteri.org/2021/12/zolau-niprrte-ma-3746-welcome-2022/}{b}

% spaw
\hi
\HT{spaw}{\B{spaw}} \I{[spaw]} \emph{vtr.} believe (s.o.).
\B{Spi\-vaw oe\-ti ru\-txe, ma oe\-yä ey\-lan.} Please believe me, my friends.~\LINK{http://forum.learnnavi.org/news-announcements/a-response-from-paul-frommer!/}{ln}
\B{Pol\-txe Ey\-tu\-kan san oe ka\-yä sìk, slä oel pot ke spaw.} Eytukan said he would come but I don't believe him.~\LINK{http://forum.learnnavi.org/language-updates/verb-for-write/msg175491/\#msg175491}{ln}
(\emph{a. conv.}) \B{spaw oe …\slash spaw oel fu\-ta …}, I believe (that).
\B{Spaw oel fu\-ta Mo\-'at\-ìl tso\-le\-'a Ney\-ti\-rit.} I believe Mo'at saw Neytiri.~\LINK{http://naviteri.org/2011/03/word-order-and-case-marking-with-modals/}{b}
\B{Spaw oe, fwa po ko\-lä lä\-ngu kxey\-ey.} I believe it was a mistake for him to have gone.~\LINK{http://naviteri.org/2011/04/\%E2\%80\%99a\%E2\%80\%99awa-li\%E2\%80\%99fyavi-amip\%E2\%80\%94a-few-new-expressions/}{b} \\
\textsc{Derivations}: \ding{43} 
\on{1}.~\HL{tIspaw}{\B{tì\-spaw}}, belief (abstract).
\on{2}.~\HL{sAspaw}{\B{sä\-spaw}}, belief (particular).

% spä
\hi
\HT{spA}{\B{spä}} \I{[sp\ae]} \emph{vin.} jump.
\B{Po spä few pay\-fya fte smar\-it si\-vutx.} He jumped across the stream to track his prey.~\LINK{http://naviteri.org/2011/01/new-year-new-vocabulary/}{b}
\B{Maw\-krr\-a pa\-lu\-lu\-kan po\-sìn spo\-lä, ke tsä\-ngun fko pot ti\-va\-rep.} Sadly, once the thanator had jumped on her, she could not be rescued.~\LINK{http://naviteri.org/2021/12/zolau-niprrte-ma-3746-welcome-2022/}{b}
(\emph{a. proverb}) \B{Spä skxawng sìn 'a\-na a\-flì.} ``A fool jumps onto a thin vine,'' i.e. \emph{don't engage in an unpromising and\slash or potentially risky cause}.~\LINK{http://naviteri.org/2021/08/ulte-ayyoratu-leiu-and-the-winners-are-2/}{b}
\B{Tsa\-ye\-rik te\-rul ne 'awkx. Spä skxawng sìn 'a\-na a\-flì.} That hexapede is running toward the cliff. Only a fool jumps onto a thin vine.~\LINK{http://naviteri.org/2021/08/ulte-ayyoratu-leiu-and-the-winners-are-2/}{b}

% spe'e
\hi
\HT{spe'e}{\B{spe\cp{'e}}} \I{[spE."PE]} \emph{vtr.} capture (s.o.\slash s.t.). 
\B{Fol nga\-ti spo\-le\-'e a krr, nga lum\-pe ke ka\-ra?} When they captured you, why didn't you resist?~\LINK{http://naviteri.org/2023/09/aawa-ayliu-si-aylifyavi-amip-a-few-new-words-and-expressions/}{b} \\
\textsc{Derivations}: \ding{43} 
\on{1}.~\HL{tIspe'e}{\B{tì\-spe\-'e}}, capture. 
\on{2}.~\HL{spe'etu}{\B{spe\-'e\-tu}}, captive. 
\on{3}.~\HL{nongspe'}{\B{nong\-spe'}}, pursue with an intent to capture.

% spe'etu
\hi
\HT{spe'etu}{\B{spe\cp{'e}tu}} \I{[spE."PE.tu]} \emph{n.} captive. 
\B{Po spe\-'e\-tu lu oe\-yä.} He is \emph{my} captive.~\LINK{http://forum.learnnavi.org/language-updates/pseudo-canon-from-bd-live-content/}{ln\slash a\textsuperscript{c}}
\B{Fol fte ay\-spe\-'e\-tut li\-vo\-nu fngo' 'u\-pet?} What are they demanding in order for them to release the captives?~\LINK{http://naviteri.org/2012/01/mipa-zisit-ayliu-amip-new-words-for-the-new-year/}{b}
(\ding{43}~\HL{spe'e}{\B{spe\-'e}})

% speng
\hi
\HT{speng}{\B{speng}} \I{[spEN]} \emph{vtr.} (\emph{a.}) restore. 
(\emph{b. metaphorically as} \B{spä\-peng}) restore oneself, get one's head back on straight.
\B{New oe ri\-vun a\-sim tì\-fnu\-nga'\-a tseng\-it a tsa\-ro tsun syi\-vor tsi\-vu\-rokx fte spä\-pi\-veng.} I want to find a quiet place nearby where I can chill out and rest to get my head back on straight.~\LINK{http://naviteri.org/2013/01/awvea-posti-zisita-amip-first-post-of-the-new-year/}{b}

% spono
\hi
\HT{spono}{\B{\cp{spo}no}} \I{["{}spo.no]} \emph{n.} island.   
\B{Ay\-oel ro\-lun mip\-a spo\-not mì hil\-van.} We found a new island in the river.~\LINK{http://naviteri.org/2012/03/spring-vocabulary-part-1/}{b}
\B{Oe\-ri so\-la\-lew wum zì\-sìt °a14 a krr, fo\-lrr\-fen spo\-not alo a\-'aw\-ve.} When I was about 12 years old I visited an island for the first time.~\LINK{http://naviteri.org/2012/10/mipa-vospxi-mipa-ayliu-new-words-for-the-new-month/}{b}

% spule
\hi
\HT{spule}{\B{\cp{spu}le}} \I{["{}spu.lE]} \emph{vtr.} propel.   
\B{Kä\-mak\-to nì\-win, ay\-nga\-ti spi\-vu\-le hu\-fwel.} You ride out as fast as the wind can carry you [lit.: ride out fast, may the wind propel you].~\LINK{http://forum.learnnavi.org/language-updates/naviya-the-other-vocative/}{ln\slash a} \\
\textsc{Derivations}: \ding{43}
\on{1}.~\HL{spulmokri}{\B{spul\-mok\-ri}}, telephone. 
\on{2}.~\HL{spulyaney}{\B{spul\-ya\-ney}}, canoe paddle.
\on{3}.~\HL{yIspul}{\B{yì\-spul}}, mermaid tail.
\on{4}.~\HL{ko'onspul}{\B{ko\-'on\-spul}} sunflower gigantus. 

% spulmokri
\hi
\HT{spulmokri}{\B{spul\cp{mok}ri}} \I{[spul."mok\textcorner.Ri]} \emph{or} \I{[spUl.]} \emph{n.} telephone.
\B{Syaw oer ru\-txe trr\-ay fa spul\-mok\-ri.} Please phone me tomorrow.~\LINK{http://naviteri.org/2012/07/fivospxiya-aylifyavi-amip-this-months-new-expressions-pt-2/}{b}
(\ding{43}~\HL{spule}{\B{spu\-le}}, \HL{mokri}{\B{mok\-ri}})

% spulyaney
\hi
\HT{spulyaney}{\B{spulya\cp{ney}}} \I{[spul.ja."nEj]} \emph{or} \I{[spUl.]} \emph{n.} canoe paddle.
(\ding{43}~\HL{spule}{\B{spu\-le}}, \HL{yaney}{\B{ya\-ney}})

% spuwin
\hi
\HT{spuwin}{\B{\cp{spu}win}} \I{["{}spu.win]} \emph{adj.} old (\emph{not for people}), former. 
\B{Tsa\-tsko lu spu\-win ul\-te ke lu mi txur, slä oe\-ri txan\-wa\-we le\-iu.} That bow is old and no longer strong, but it means a lot to me.~\LINK{http://naviteri.org/2011/05/some-miscellaneous-vocabulary/}{b}
(\ding{43}~\HL{lal}{\B{lal}})

% spxam
\hi
\HT{spxam}{\B{spxam}} \I{[sp'am]} \emph{n., funga} fungus. 
\textsc{Derivations}: \ding{43} 
\on{1}.~\HL{torukspxam}{\B{to\-ruk\-spxam}}, octoshroom.
\on{2}.~\HL{txepvispxam}{\B{txep\-vi\-spxam}}, sparkle pod. 

% spxin
\hi
\HT{spxin}{\B{spxin}} \I{[sp'in]} \emph{adj.} sick.
\B{Mìn, mìn, ken\-ten mìn. | Pe\-lun ke 'efu spxin?} Turn, turn, the fanlizard turns. | Why doesn't it feel sick?~\LINK{http://forum.learnnavi.org/language-updates/learnings-from-the-siatlla-song-translations/msg473566/\#msg473566}{ln}
\B{Oe\-yä 'ey\-lan ke 'e\-fu spxin. Slä ke tslä\-ngam pol fu\-ta nì\-fkey\-to\-ngay lu spxin nì\-txan.} My friend doesn't feel sick. But he doesn't understand that he is in fact very sick.~\LINK{http://naviteri.org/2020/05/tipusawm-tiuseyng-si-okvur-a-eltur-titxen-si-asking-answering-and-an-interesting-story/\#comment-31207}{b}
(\ding{214}~\B{le\-fpom\-tokx}) \\
\textsc{Derivations}: \ding{43}
\on{1}.~\HL{tIspxin}{\B{tì\-spxin}}, sickness, state of being sick. 
\on{2}.~\HL{sAspxin}{\B{sä\-spxin}}, sickness, disease, illness.
\on{3}.~\HL{spxintu}{\B{spxin\-tu}}, sick or ill person.

% spxintu
\hi 
\HT{spxintu}{\B{\cp{spxin}tu}} \I{["{}sp'in.tu]} \emph{n.} sick or ill person. 
\B{Lä\-ngu zam\-a spxin\-tu mì sray oe\-yä.} I'm sorry to say there are sixty\hyp{}four sick people in my village.~\LINK{https://naviteri.org/2025/06/vospikin-lefpom-happy-july/}{b} 
(\ding{43}~\HL{spxin}{\B{spxin}})

% srak / srake
\hi
\HT{srak}{\B{srak}} \I{[sRak\textcorner]} \emph{or} \B{\cp{sra}ke} \I{["{}sRa.kE]} \emph{part. to indicate a yes\slash no\hyp{}question.}   
(\emph{a. at the end}) \B{Nga\-ru lu fpom srak?} How are you? [lit.: do you have well\hyp{}being?]~\LINK{http://wiki.learnnavi.org/index.php?title=Corpus\#Vanity_Fair}{wiki}
\B{Nga zo srak?} Are you well? (\emph{e.g. have you recovered from your illness\slash fall, etc.?})~\LINK{http://naviteri.org/2010/07/vocabulary-update/}{b}
\B{Nga za\-'u fì\-tseng pxìm srak?} Do you come here often?~\LINK{http://wiki.learnnavi.org/index.php?title=Corpus\#Lemondrop.com}{wiki\slash ld}
\B{Ngal 'e\-raw\-nìm oe\-ti srak?} Are you avoiding me?~\LINK{http://naviteri.org/2011/09/miscellaneous-vocabulary/}{b}
\B{Tsun (oe) mi\-vä\-kxu hì\-krr srak?} May I interrupt a moment?~\LINK{http://naviteri.org/2010/09/getting-to-know-you-part-1/}{b}
(\emph{b. at the beginning}) \B{Sra\-ke nga za\-'u?} Are you coming?~\LINK{http://forum.learnnavi.org/language-updates/if-vs-whether-ftxey-fuke/}{ln}
\B{Sra\-ke tsun oe fay\-upxa\-ret tsli\-vam nì\-ftue?} Can I understand these messages easily?~\LINK{http://forum.learnnavi.org/language-updates/verb-for-write/msg175592/\#msg175592}{ln}
(\ding{43}~\HL{sran}{\B{sra\-ne}}, \HL{kehe}{\B{ke\-he}})

% srakat
\hi 
\HT{srakat}{\B{\cp{sra}kat}} \I{["{}sRa.kat\textcorner]} \emph{n., fauna} dinicthoid, \emph{large piranha\hyp{}like aquatic creature that lives in Pandora's lakes and murky lowland drainages}.~\LINK{https://forum.learnnavi.org/language-updates/expansion-to-flora-fauna-and-places-from-the-valley-of-mo'ara/}{ln}

% sran / srane
\hi
\HT{sran}{\B{sran}} \I{[sRan]} \emph{or} \B{\cp{sra}ne} \I{["{}sRa.nE]} \emph{intj.} yes, (\emph{coll.}) yeah.
\B{Sra\-ne, fì\-txon tsun nga ni\-väk swo\-at nì\-'it, slä rä\-'ä rou!} Yes, you can have a little alcohol tonight, but don't get drunk!~\LINK{http://naviteri.org/2012/10/audio-and-video-learning-materials-for-navi-101/}{b}
\B{Nga 'e\-fu väng srak?---Sran, 'e\-fu nì\-txan.} Are you thirsty?---Yep, I'm parched.~\LINK{http://naviteri.org/2012/10/audio-and-video-learning-materials-for-navi-101/}{b}
(\ding{214}~\HL{kehe}{\B{ke\-he}}; \ding{43}~\HL{li}{\B{li}}) \\
\textsc{Derivations}: \ding{43}
\on{1}.~\HL{kesran}{\B{ke\-sran}}, so\hyp{}so, mediocre. 
\on{2}.~\HL{srak}{\B{srak\slash sra\-ke}}, \emph{yes\slash no particle}. 
\on{3}.~\HL{srankehe}{\B{sran\-ke\-he}}, more or less.

% srankehe
\hi 
\HT{srankehe}{\B{sran\cp{ke}he}} \I{[sRan."kE.hE]} \emph{or} \I{[sRaN."kE.hE]} \emph{part., intj.} more or less, somewhat, yes and no, kind of (\emph{an equivocal response to a yes\slash no question, so as not to commit oneself}). 
\B{Sra\-ke fay\-sä\-fpìl lu pum ngey nì\-wotx?--Sran\-ke\-he.} Are all these ideas your own?--More or less.~\LINK{http://naviteri.org/2018/03/100a-liu-amip-64-new-words-part-1/}{b}
(\ding{43}~\HL{sran}{\B{sran(e)}}, \HL{kehe}{\B{ke\-he}})

% sraw
\hi
\HT{sraw}{\B{sraw}} \I{[sRaw]} \emph{adj.} painful. \\
\textsc{Derivations}: \ding{43}
\on{1}.~\HL{tIsraw}{\B{tì\-sraw}}, pain.

% srä
\hi
\HT{srA}{\B{srä}} \I{[sR\ae]} \emph{n.} cloth (\emph{a piece of cloth woven on a loom}).
\B{Tsa\-srä a\-naw\-nekx ley nì\-'it nì\-'aw.} That burnt cloth is of little value.~\LINK{http://naviteri.org/2014/03/value-and-worth/}{b}
\B{Fu\-ri\-a txu\-la tsa\-lew\-ti lu srä sìl\-tsan to fngap.} For constructing that cover, woven cloth is better than metal.~\LINK{http://naviteri.org/2013/07/pximaw-tsawlultxa-just-after-the-meet-up/}{b}
\B{Fu\-ta nga\-ta tel pxen\-yoa srät, mll\-'ei\-an oel.} I'm happy to take cloth from you for finished garments.~\LINK{http://naviteri.org/2014/02/barter-and-exchange/}{b} 

% srätx
\hi
\HT{srAtx}{\B{srätx}} \I{[sR\ae{}t']} \emph{vtr.} bother, annoy.   
\B{Oe\-ti rä\-'ä srätx.} Don't bother me.~\LINK{http://naviteri.org/2011/09/miscellaneous-vocabulary/}{b} 
(\emph{a. expression}) \B{Sra\-ke srätx (ngat) txo \ldots?} (\emph{with} ‹iv›) `Would you mind if \ldots?'~\LINK{http://naviteri.org/2020/11/vospxivopeya-ayliu-amip-novembers-new-words/}{b}
\B{Ke srätx (oet) txo \ldots} `I don't mind if \ldots'~\LINK{http://naviteri.org/2020/11/vospxivopeya-ayliu-amip-novembers-new-words/}{b}       
\B{Ke srätx kaw\-'it!\slash Kea sä\-srätx kaw\-'it!\slash Ke\-he kaw\-'it!} `Not at all!'~\LINK{http://naviteri.org/2020/11/vospxivopeya-ayliu-amip-novembers-new-words/}{b} \\
\textsc{Derivations}: \ding{43}
\on{1}.~\HL{sAsrAtx}{\B{sä\-srätx}}, annoyance.

% sre
\hi
\HT{sre}{\B{sre}} \I{[sRE]} \emph{adp.}\texttt{+} before (\emph{temporal}). (\emph{a.}) \B{Sre fwa sìn tskxe\-pay tì\-ran, ze\-ne fko flì\-nutx\-it sti\-ve\-ftxaw.} It's necessary to check the thickness of the ice before walking on it.~\LINK{http://naviteri.org/2011/11/\%E2\%80\%99a\%E2\%80\%99awa-ayli\%E2\%80\%99u-amip-ni\%E2\%80\%99aw-\%E2\%80\%94-a-few-new-words-only/}{b}
(\emph{b. of telling time}) \B{trr\-pxì\-vi a\-vo\-mun sre srr\-pxì a\-mrr\-ve} \emph{or coll.} \B{vo\-mun sre mrr\-ve}, (it's) \on{4}:\on{50}, or `ten to five'.~\LINK{https://naviteri.org/2024/05/krrteri-about-time/}{b}
(\emph{c. with} \B{li}) by; before or up to, but not after. 
\B{Kem si li trr\-ay\-sre.} Do it by tomorrow. (\emph{also} \ding{43}~\HL{lisre}{\B{li\-sre}})
(\emph{d. idiom}) \B{sre fwa sngap zi\-ze'}, ``as quickly as possible'' [lit: before the hellfire wasp stings].~\LINK{http://naviteri.org/2012/07/fivospxiya-aylifyavi-amip-this-months-new-expressions-pt-2/}{b}
(\ding{214}~\HL{maw}{\B{maw}}) \\
\textsc{Derivations}:
\on{1}.~\HL{lisre}{\B{li\-sre}}, by; up to, but not after. 
\on{2}.~\HL{srefey}{\B{sre\-fey}}, expect. 
\on{3}.~\HL{srefpIl}{\B{sre\-fpìl}}, assume, presume. 
\on{4}.~\HL{srese'a}{\B{sre\-se\-'a}}, prophesize, predict. 
\on{5}.~\HL{srefwa}{\B{sre\-fwa}}, before (\emph{conj.}). 
\on{6}.~\HL{srekrr}{\B{sre\-krr}}, before (\emph{adv.}). 
\on{7}.~\HL{srekamtrr}{\B{sre\-kam\-trr}}, time before noon.
\on{8}.~\HL{srekamtxon}{\B{sre\-kam\-txon}}, time before midnight. 
\on{9}.~\HL{sresrr'ong}{\B{sre\-srr\-'ong}}, before dawn. 
\on{10}.~\HL{sreton'ong}{\B{sre\-ton\-'ong}}, dusk.

% sre'
\hi
\HT{sre'}{\B{sre'}} \I{[sREP]} \emph{n.} tooth. 
\B{Ay\-sre' lom lu tsa\-ko\-ak\-tan\-ur.} The old man misses his teeth.~\LINK{http://naviteri.org/2011/07/txantsana-ultxa-mi-siatll-great-meeting-in-seattle/}{b}
\B{Oel tstal\-it fyo\-lep fa aysre'.} I held the knife in my teeth.~\LINK{http://naviteri.org/2012/03/spring-vocabulary-part-1/}{b}

% srefey
\hi
\HT{srefey}{\B{sre\cp{fey}}} \I{[sRE."f$\cdot$Ej]} \emph{vtr.}\slash \emph{vin.}\on{22} expect (\emph{takes} \HL{tsnI}{\B{tsnì}} \emph{when used as vin.}).
\B{Set sre\-fey oel fu\-ta tsam\-po\-ngu tä\-txaw maw txon\-'ong.\slash Set sre\-fey oe tsnì tsam\-po\-ngu …} I'm currently expecting the war party back after nightfall.~\LINK{http://naviteri.org/2011/02/new-vocabulary-part-2/}{b}
(\emph{a. idiom}) \B{sre\-fe\-re\-i\-ey nì\-prr\-te'}, looking forward (to s.t.).~\LINK{http://naviteri.org/2011/02/new-vocabulary-part-2/}{b} \B{Tsa\-ria nga\-hu ye'\-rìn ul\-txa si nì\-mun, oe sre\-fe\-re\-i\-ey nì\-prr\-te'.} I'm looking forward to getting together with you again soon.~\LINK{http://naviteri.org/2011/02/new-vocabulary-part-2/}{b}
(\ding{43}~\HL{sre}{\B{sre}}, \HL{pey}{\B{pey}})

% srefpìl
\hi
\HT{srefpIl}{\B{sre\cp{fpìl}}} \I{[sRE."fp$\cdot$Il]} \emph{vtr.}\slash \emph{vin., inf.}\on{22} assume, presume.
(\emph{a. vtr.}) \B{Sre\-fpìl oel fu\-ta nga lu tok\-tor Lì\-vìng\-sì\-ton.} Doctor Livingstone, I presume.~\LINK{http://naviteri.org/2014/10/tson-si-fpomron-obligation-and-mental-health/}{b}
(\emph{a. vin. used with} \HL{tsnI}{\B{tsnì}}) \B{Sre\-fpìl Oma\-ti\-ka\-ya tsnì Tsyeyk kaw\-krr ke ta\-yä\-txaw maw ka\-vuk sne\-yä.} The Omaticaya assumed that Jake would never return after his treachery.~\LINK{http://naviteri.org/2014/10/tson-si-fpomron-obligation-and-mental-health/}{b}
(\ding{43}~\HL{sre}{\B{sre}}, \HL{fpIl}{\B{fpìl}})

% srefwa
\hi
\HT{srefwa}{\B{\cp{sre}fwa}} \I{["{}sRE.fwa]} \emph{conj.} before.
\B{Sre\-fwa oe hum, new pi\-vll\-txe.} Before I leave, I want to speak.~\LINK{http://naviteri.org/2014/10/tson-si-fpomron-obligation-and-mental-health/}{b}
(\ding{43}~\HL{sre}{\B{sre}}, \HL{fwa}{\B{fwa}}, \ding{214}~\HL{mawfwa}{\B{maw\-fwa}})

% srekamtrr
\hi
\HT{srekamtrr}{\B{sre\cp{kam}trr}} \I{[sRE."kam.t\s{r}]} \emph{n.}; \emph{adv.} before midday, before noon.
\B{Fo pä\-hem trr\-pxì\-pe?--Sre\-kam\-trr.} What part of the day will they arrive?--Before noon.~\LINK{https://naviteri.org/2024/05/krrteri-about-time/}{b}
(\ding{43}~\HL{sre}{\B{sre}}, \HL{kxamtrr}{\B{kxam\-trr}})

% srekamtxon
\hi
\HT{srekamtxon}{\B{sre\cp{kam}txon}} \I{[sRE."kam.t'on]} \emph{n.}; \emph{adv.} before midnight.
(\ding{43}~\HL{sre}{\B{sre}}, \HL{kxamtxon}{\B{kxam\-txon}})

% srekrr
\hi
\HT{srekrr}{\B{sre\cp{krr}}} \I{[sRE."k\s{r}]} \emph{adv.} before (\emph{temporal}), beforehand.  
\B{Sre\-krr 'a\-me\-fu väng, set ye\-väng.} Before, I was thirsty; now my thirst has been quenched.~\LINK{http://naviteri.org/2011/04/\%E2\%80\%99a\%E2\%80\%99awa-li\%E2\%80\%99fyavi-amip\%E2\%80\%94a-few-new-expressions/}{b}
\B{Fay\-upxa\-re\-mì oe pa\-yäng\-kxo te\-ri ho\-ren lì'\-fya\-yä le\-Na'\-vi fpi su\-te a tsun sre\-krr tsat si\-var.} In these messages I will chat about the rules of the Na'vi language.for people who can already use it.~\LINK{http://naviteri.org/2010/06/irayo-thank-you-and-some-miscellaneous-thoughts/}{b}
\B{Tseyk tswa\-may\-on fa ik\-ran sre\-krr; ta\-fral fmo\-li fì\-kem si\-vi fa to\-ruk nì\-steng.} Jake flew with an ikran before; therefore he tried to do it with a \emph{toruk} in a similar fashion.~\LINK{http://naviteri.org/2011/09/miscellaneous-vocabulary/}{b} 
(\ding{43}~\HL{sre}{\B{sre}}, \HL{krr}{\B{krr}})

% srer
\hi
\HT{srer}{\B{srer}} \I{[sRER]} \emph{vin.} appear, materialize, come into view.
\B{Prr\-nen lrr\-tok si a krr, srer me\-song\-tsyìp a\-ho\-na.} When the baby smiles, two adorable dimples appear.~\LINK{http://naviteri.org/2012/03/spring-vocabulary-part-1/}{b}
\B{Txon\-am teng\-krr tar\-mì\-ran oe kxam\-lä na'\-rìng, sro\-ler eo ut\-ral a\-tsawl txewm\-a vrr\-tep.} Last night as I was walking through the forest, a frightening demon appeared in front of a big tree.~\LINK{http://naviteri.org/2012/03/spring-vocabulary-part-1/}{b}
(\ding{214}~\HL{'Ip}{\B{'ìp}})

% srese'a
\hi
\HT{srese'a}{\B{srese\cp{'a}}} \I{[sRE.s$\cdot$E."P$\cdot$a]} \emph{vtr., inf.}\on{23} prophesy, predict. 
\B{Tsun fko ay\-on\-ti fì\-wopx\-ä ni\-vìn fte ya\-fkeyk\-it sre\-si\-ve\-'a.} The shapes of clouds can be used to predict the weather.~\LINK{http://naviteri.org/2013/09/on-si-salewfya-shapes-and-directions/}{b}
(\ding{43}~\HL{sre}{\B{sre}}, \HL{tse'a}{\B{tse\-'a}}) \\
\textsc{Derivations}: \ding{43}
\on{1}.~\HL{sAsrese'a}{\B{sä\-sre\-se\-'a}}, prediction.

% sresrr'ong
\hi
\HT{sresrr'ong}{\B{sresrr\cp{'ong}}} \I{[sRE.s\s{r}."PoN]} \emph{n.}; \emph{adv.} before daybreak, before dawn.
\B{Sa'\-nok 'eveng\-ur tì\-txen si sre\-srr\-'ong.} Mother woke the child in the early morning.~\LINK{http://forum.learnnavi.org/language-updates/little-confirmation-titxen-si/}{ln}
(\ding{43}~\HL{sre}{\B{sre}}, \HL{trr'ong}{\B{trr\-'ong}})

% sreton'ong
\hi
\HT{sreton'ong}{\B{sreton\cp{'ong}}} \I{[sRE.ton."oN]} \emph{n.}; \emph{adv.} before darkness falls, before night unfolds.
(\ding{43}~\HL{sre}{\B{sre}}, \HL{txon'ong}{\B{txon\-'ong}})

% srew
\hi
\HT{srew}{\B{srew}} \I{[sREw]} \emph{vin.} dance. 
\B{Zun tom\-pa zì\-ye\-vup trr\-ay, zel fo srì\-ye\-vew.\slash … zel fo srew.} If it rained tomorrow, they'd do a dance.~\LINK{http://naviteri.org/2013/04/zun-zel-counterfactual-conditionals/}{b}

% srey
\hi
\HT{srey}{\B{srey}} \I{[sREj]} \emph{n.} version.
\B{Fì\-srey puk\-ä lu swey nì\-law.} This version of the book is clearly the best.~\LINK{http://naviteri.org/2014/07/tsawlultxamaw-a-ayliu-post-meetup-words/}{b}
\B{Lu srey a\-pup ta Nga\-ri zì\-sìt a\-mip li\-vu le\-fpom.} It's a short version of \emph{Ngari zìsìt amip livu lefpom}.~\LINK{http://naviteri.org/2017/12/zisit-amip-lefpom-ma-eylan-happy-new-year-friends/\#comment-27964}{b}

% srìn
\hi
\HT{srIn}{\B{srìn}} \I{[sRIn]} \emph{vtr.} temporarily transfer from one to another (\emph{mainly used in its derivational forms}). \\
\textsc{Derivations}: \ding{43}
\on{1}.~\HL{kAsrIn}{\B{kä\-srìn}}, lend. 
\on{2}.~\HL{zasrIn}{\B{za\-srìn}}, borrow. 
\on{3}.~\HL{sAsrIn}{\B{sä\-srìn}}, lent\slash borrowed thing.

% srok
\hi
\HT{srok}{\B{srok}} \I{[sRok\textcorner]} \emph{n.} bead (\emph{decorative}).
\B{'En si oe, lor\-a tsa\-fkxi\-le\-ri a\-pxay\-opin so\-lar Tse\-nul srok\-it a\-vo\-zam.} I would guess that Tsenu used a thousand (lit. \on{512}) beads for that beautiful multi\hyp{}colored bib necklace.~\LINK{http://naviteri.org/2013/07/pximaw-tsawlultxa-just-after-the-meet-up/}{b}
\B{Fol ko\-la\-nom po\-ta ay\-srok\-it fay\-oang\-yoa.} They acquired beads from him in exchange for fish.~\LINK{http://naviteri.org/2014/02/barter-and-exchange/}{b}

% sru'
\hi
\HT{sru'}{\B{sru'}} \I{[sRuP]} \emph{or} \I{[sRUP]} \emph{vtr.} crush, trample.   \B{Weyn\-flit\-it 'ang\-tsìk\-ìl sro\-lu' tspang.} Wainfleet was crushed and killed by a hammerhead.~\LINK{http://naviteri.org/2012/07/meetings-waterfalls-and-more/}{b}

% srung / srung si
\hi
\HT{srung}{\B{srung}} \I{[sRuN]} \emph{or} \I{[sRUN]} (\textsc{rn}: \B{srùng} \I{[sRUN]}) \textsc{i}. \emph{n.} help, assistance. \\
\textsc{ii}. \HT{srungsi}{\B{srung si}} \I{[sRuN si]} (\textsc{rn}: \B{srùng si} \I{[sRUN si]}) \emph{vin.} (\emph{a.}) help (s.o.). 
\B{Ma Ra\-lu, srung si por!} Ralu, help her!~\LINK{http://naviteri.org/2012/03/spring-vocabulary-part-2/}{b}
\B{Txan\-a srung so\-li po oer.} He helped me a lot.~\LINK{http://forum.learnnavi.org/language-updates/si-verbs-modification/}{ln}
\B{Tsa\-tì\-oeyk\-tìng a\-ya\-yayr\-nga' srung ke so\-li oer fte tsli\-vam teyng\-ta kem\-pe ze\-ne si\-vi.} That confusing explanation didn't help me understand what I have to do.~\LINK{http://naviteri.org/2013/01/awvea-posti-zisita-amip-first-post-of-the-new-year/}{b}
(\emph{b. of people}) helpful. \B{tu\-te a srung si\slash srung si a tu\-te}, a helpful person.~\LINK{http://naviteri.org/2015/03/kap-si-ayunil-saylahe-threats-dreams-and-other-things/}{b} \\  
\textsc{Derivations}: \ding{43} 
\on{1}.~\HL{srungsiyu}{\B{srung\-si\-yu}}, assistant, helper. 
\on{2}.~\HL{srungtsyIp}{\B{srung\-tsyìp}}, helpful hint, tip. 
\on{3}.~\HL{srunga'}{\B{sru\-nga'}}, helpful.
\on{4}.~\HL{sAsrungsiyu}{\B{srung\-si\-yu}}, assistant.
\on{5}.~\HL{sAsrung}{\B{sä\-srung}}, helper (inanimate).

% srunga' 
\hi
\HT{srunga'}{\B{\cp{sru}nga'}} \I{["sRu.NaP]} (\textsc{rn}: \B{\cp{srù}nga'} \I{["sRU.NaP]}) \emph{adj., nfp.} helpful.
\B{Tì\-yäkx ke lu sru\-nga', ma tsmuk. Nga txo sti, oeyk\-tìng teyng\-ta pe\-lun.} Snubbing isn't helpful, brother. If you're angry, explain why.~\LINK{http://naviteri.org/2015/03/kap-si-ayunil-saylahe-threats-dreams-and-other-things/}{b}
(\emph{a. proverb}) \B{Hem a\-sru\-nga' nì\-'ul, hum a\-sngu\-nga' nì\-nän.}
More helpful actions lead to less troubling outcomes.~\LINK{http://naviteri.org/2015/03/kap-si-ayunil-saylahe-threats-dreams-and-other-things/}{b} 
(\ding{43}~\HL{srung}{\B{srung}}, \HL{-nga'}{\B{-nga'}}, \HL{srung}{\B{srung si} (\emph{b. of people})})

% srungsiyu
\hi
\HT{srungsiyu}{\B{\cp{srung}siyu}} \I{["sRuN.si.ju]} \emph{or} \I{["sRUN.]} \emph{n.} assistant, helper, one who helps.
\B{Hu\-fwa lu oe fì\-trr kar\-yu a\-txin, le\-iu oer kop pxe\-srung\-si\-yu a fì\-lì'\-fya\-ri slo\-lu tsul\-fä\-tu.} Although I'm the main teacher today, I'm happy to say I also have three assistants who have become masters of the language.~\LINK{http://naviteri.org/2012/07/tskxekengtsyip-a-mikyunfpi-a-little-listening-exercise/}{b}
(\ding{43}~\HL{srung}{\B{srung si}}, \HL{-yu}{\B{-yu}})

% srungtsyìp
\hi
\HT{srungtsyIp}{\B{\cp{srung}tsyìp}} \I{["{}sRuN.\t{ts}jIp\textcorner]} (\textsc{rn}: \B{\cp{srùng}tsyìp} \I{["sRUN.\t{tS}Ip\textcorner]})  \emph{n.} helpful hint, tip.
\B{Srung\-tsyìp a\-1ve}, 1st tip.~\LINK{http://naviteri.org/2012/10/audio-and-video-learning-materials-for-navi-101/}{b}
(\ding{43}~\HL{srung}{\B{srung}}, \HL{-tsyIp}{\B{-tsyìp}})

% stang
\hi 
\HT{stang}{\B{stang}} \I{[staN]} \emph{n.} torso. 

% starlì'u
\hi 
\HT{starlI'u}{\B{\cp{star}lì'u}} \I{["{}staR.lI.Pu]} \emph{n.} adposition. 
(\ding{43}~\HL{sAtare}{\B{sä\-ta\-re}}, \HL{lI'u}{\B{lì\-'u}})

% starsìm
\hi
\HT{starsIm}{\B{\cp{star}sìm}} \I{["{}staR.sIm]} \emph{vtr.} gather, collect.
\B{Aw\-nga tsong\-ne ki\-vä fte sti\-var\-sìm tey\-lut.} Let's go to the valley to gather beetle larvae.~\LINK{http://naviteri.org/2012/03/spring-vocabulary-part-1/}{b} \\
\textsc{Derivations}: \ding{43}
\on{1}.~\HL{sAstarsIm}{\B{sä\-star\-sìm}}, collection.

% stawm
\hi
\HT{stawm}{\B{stawm}} \I{[stawm]} \emph{vtr.} (\texttt{-}\emph{control}) hear (s.t.), hear about s.t. (\B{te\-ri}).
\B{Txo fko ti\-vul mì na'\-rìng nì\-hawm\-pam, stawm ay\-ioang.} If you run noisily in the forest, the animals will hear.~\LINK{http://naviteri.org/2014/11/twenty-before-the-holidays/}{b}
\B{Fì\-slär\-mì tsun fko sti\-vawm ngam\-it a\-pxay.} You can hear a lot of echoes in this cave.~\LINK{http://naviteri.org/2012/01/mipa-zisit-ayliu-amip-new-words-for-the-new-year/}{b}
\B{Fpo\-le' ay\-ngal oer fì\-txan nì\-ftxa\-vang a 'upxa\-ret sto\-lawm oel.} I have heard the message you have sent me so passionately.~\LINK{http://forum.learnnavi.org/news-announcements/a-response-from-paul-frommer!/}{ln}
\B{Fol oe\-yä tì\-pawm\-it ke sto\-lä\-ngawm.} Unfortunately they didn't hear my question.~\LINK{http://naviteri.org/2012/11/renu-ayinanfyaya-the-senses-paradigm/}{b}
\B{Krr a sto\-lawm Txe\-wì te\-ri hem a poe\-ru lo\-len, 'o\-le\-fu ke\-ftxo nì\-txan.} When Txewì heard about what happened to her, he felt very upset.~\LINK{http://naviteri.org/2010/07/ting-mikyun-fte-tslivam-listening-comprehension-1-3/}{b}
(\ding{43}~\HL{yune}{\B{yu\-ne}}) \\
\textsc{Derivations}: \ding{43}
\on{1}.~\HL{stawmtswo}{\B{stawm\-tswo}}, (sense of) hearing.

% stawmtswo
\hi
\HT{stawmtswo}{\B{\cp{stawm}tswo}} \I{["{}stawm.\t{ts}wo]} \emph{n.} (sense of) hearing. 
\B{Tsa\-ko\-ak\-te\-ri stawm\-tswo lu fe'. Po\-hu a tì\-päng\-kxo ngä\-zìk lu nì\-txan.} That old woman's hearing is poor. A conversation with her is very difficult.~\LINK{http://naviteri.org/2012/11/renu-ayinanfyaya-the-senses-paradigm/}{b}
(\ding{43}~\HL{stawm}{\B{stawm}}, \HL{-tswo}{\B{-tswo}})

% stä'nì
\hi
\HT{stA'nI}{\B{\cp{stä'}nì}} \I{["{}st\ae{}P.nI]} \emph{vtr.} catch. 
\B{Oe\-yä txin\-tìn lu fwa stä'\-nì fay\-oang\-it.} My central societal occupation is to catch fish.~\LINK{http://naviteri.org/2010/09/getting-to-know-you-part-2/}{b}
\B{Fì\-ta\-ron\-yu\-tsyìp ke tsun ke\-'ut sti\-vä'\-nì.} This little hunter can't catch anything.~\LINK{http://naviteri.org/2010/07/diminutives-conversational-expressions/}{b} \\
\textsc{Derivations}: \ding{43}
\on{1}.~\HL{stA'nIpam}{\B{stä'\-nì\-pam}}, recording.

% stä'nìpam
\hi 
\HT{stA'nIpam}{\B{\cp{stä'}nìpam}} \I{["st\ae{}P.nI.pam]} \textsc{i}. \emph{n.} recording. 
\B{Fì\-vur\-way\-ri a\-lor sì stä'\-nì\-pam\-ìri tse\-yä ira\-yo nì\-txan, ma mesmuk!} Thank you two very much for this beautiful narravite poem and its recording!
(\ding{43}~\HL{stA'nI}{\B{stä'\-nì}}, \HL{pam}{\B{pam}}) \\
\textsc{ii}. \B{\cp{stä'}nìpam si} \I{["st\ae{}P.nI.pam si]} \emph{vin.} record, make a recording.~\LINK{http://naviteri.org/2019/10/vurway-alor-a-beautiful-narrative-poem-2/}{b}
\B{Sä\-ftxu\-lì\-'u a\-tì\-txur\-nga' nì\-txan nang! Fu\-ria tsa\-ru nga stä'\-nì\-pam so\-li, ira\-yo.} What a powerful speech! Thank you for recording it.~\LINK{http://naviteri.org/2019/10/vurway-alor-a-beautiful-narrative-poem-2/}{b}

% steftxaw
\hi
\HT{steftxaw}{\B{ste\cp{ftxaw}}} \I{[stE."ft'aw]} \emph{vtr.} examine, check. 
\B{Sre fwa sìn tskxe\-pay tì\-ran, ze\-ne fko flì\-nutx\-it sti\-ve\-ftxaw.} It's necessary to check the thickness of the ice before walking on it.~\LINK{http://naviteri.org/2011/11/\%E2\%80\%99a\%E2\%80\%99awa-ayli\%E2\%80\%99u-amip-ni\%E2\%80\%99aw-\%E2\%80\%94-a-few-new-words-only/}{b} 
\B{Tsa\-swi\-rä\-ti lo\-nu! Ay\-nga ne\-to ri\-vikx! Fì\-ke\-tu\-wong\-ti oel stì\-ye\-ftxaw.} Release this creature! Step back! I will look at this alien.~\LINK{http://wiki.learnnavi.org/index.php?title=Na\%27vi_from_Avatar_Movie\#Scene_in_Hometree}{wiki\slash a}
(\ding{43}~\HL{lang}{\B{lang}}) \\
\textsc{Derivations}: \ding{43}
\on{1}.~\HL{tIsteftxaw}{\B{tì\-ste\-ftxaw}}, examination.

% steng
\hi
\HT{steng}{\B{steng}} \I{[stEN]} \emph{adj.} similar.   
\B{Me\-unil\-tì\-ran\-tokx Tok\-tor Kì\-rey\-sä sì Tsyeyk\-ä ke lu teng ki steng.} The avatars of Grace and Jake aren't the same, but they are similar.~\LINK{http://naviteri.org/2011/09/miscellaneous-vocabulary/}{b}
\B{Nga\-ru tì\-yawr. Lo\-lu oer steng\-a sä\-fpìl.} You're right. I had a similar idea.~\LINK{http://naviteri.org/2011/09/miscellaneous-vocabulary/}{b}
(\emph{a. proverb}) \B{Sä\-fpìl a\-steng, tì\-kan a\-teng.} Great minds think alike [lit.: similar thought, same aim].~\LINK{http://naviteri.org/2011/09/miscellaneous-vocabulary/}{b} \\
\textsc{Derivations}: \ding{43}
\on{1}.~\HL{nIsteng}{\B{nì\-steng}}, similarly.

% steyki
\hi
\HT{steyki}{\B{stey\cp{ki}}} \I{[stEj."k$\cdot$i]} \emph{vtr., inf.}\on{22} anger (s.o.), make (s.o.) angry. 
\B{Ngal kaw\-krr sngel\-it ke mu\-nge wrr\-pa\-ne ul\-te nì\-tut ze\-ne ätxä\-le sir\-vä\-ngi ngar oe, stey\-ki fì\-txan.} Your never taking out the trash and my having to constantly ask you to do it makes me so angry.~\LINK{http://forum.learnnavi.org/language-updates/txelanit-hivawl/}{ln}
\B{Apxa tì\-väng\-ìl po\-ti stey\-ko\-li.} (His) great thirst made him angry.~\LINK{http://naviteri.org/2011/01/new-year-new-vocabulary/}{b}
\B{Nge\-yä fay\-sä\-fngo'\-ìl nì\-wotx stey\-ke\-rä\-ngi oe\-ti nì\-hawng.} All these demands of yours are making me exceedingly angry.~\LINK{http://naviteri.org/2012/01/mipa-zisit-ayliu-amip-new-words-for-the-new-year/}{b}
(\ding{43}~\HL{sti}{\B{sti}})

% sti
\hi
\HT{sti}{\B{sti}} \I{[sti]} \emph{vin.} angry, be angry. 
\B{Tì\-yäkx ke lu sru\-nga', ma tsmuk. Nga txo sti, oeyk\-tìng teyng\-ta pe\-lun.} Snubbing isn't helpful, brother. If you're angry, explain why.~\LINK{http://naviteri.org/2015/03/kap-si-ayunil-saylahe-threats-dreams-and-other-things/}{b}
\B{Oe\-ru txoa li\-vu, ma 'ey\-lan. Rä\-'ä sti\-vi. Lu hì\-'i\-a sä\-'i\-pu nì\-'aw.} I'm sorry, friend. Don't be angry. It was just a small bit of humor.~\LINK{http://naviteri.org/2012/01/mipa-zisit-ayliu-amip-new-words-for-the-new-year/}{b}
\B{Fu\-ri\-a fì\-lì'\-fya\-vit fkol ta 'Ìng\-lì\-sì mo\-lu\-nge, sìl\-pey oe tsnì ay\-ha\-pxì\-tu lì'\-fya\-o\-lo'\-ä aw\-nge\-yä ke stì\-ye\-vi.} That this expression is brought from English, will hopefully not anger the members of the language group.~\LINK{http://forum.learnnavi.org/language-updates/txelanit-hivawl/}{ln} \\
\textsc{Derivations}: \ding{43}
\on{1}.~\HL{steyki}{\B{stey\-ki}}, make (s.o.) angry. 
\on{2}.~\HL{tIsti}{\B{tì\-sti}}, anger. 
\on{3}.~\HL{nIsti}{\B{nì\-sti}}, angrily.

% stiwi
\hi 
\HT{stiwi}{\B{\cp{sti}wi}} \I{["{}sti.wi]} \textsc{i}. \emph{n.} mischief. \\
\textsc{ii}. \B{\cp{sti}wi si} \I{["{}sti.wi si]} \emph{vin.} be naughty, do mischief.
\B{Sti\-wi rä\-'ä si, ma 'e\-veng! U\-van si mì se\-ngo a\-la\-he.} Don't be naughty, child! Play somewhere else.~\LINK{http://naviteri.org/2016/06/mrrvola-lifyavi-amip-forty-new-expressions/}{b} \\
\textsc{Derivations}: \ding{43}
\on{1}.~\HL{stiwinga'}{\B{sti\-wi\-nga'}}, mischievous. 
\on{2}.~\HL{stiwiyu}{\B{sti\-wi\-yu}}, mischief\hyp{}maker.

% stiwinga'
\hi 
\HT{stiwinga'}{\B{\cp{sti}winga'}} \I{["{}sti.wi.NaP]} \emph{adj.} mischievous.
(\ding{43}~\HL{stiwi}{\B{sti\-wi}})

% stiwisiyu
\hi 
\HT{stiwisiyu}{\B{\cp{sti}wisiyu}} \I{["{}sti.wi.si.ju]} \emph{n.} mieschief\hyp{}maker. 
(\ding{43}~\HL{stiwi}{\B{sti\-wi}})

% sto
\hi
\HT{sto}{\B{sto}} \I{[sto]} \emph{vtr.} (\emph{a.}) refuse (s.t.).
\B{Po ätxä\-le so\-li nì\-ho\-na fì\-txan, ke tsun oe sti\-vo.} She asked so sweetly that I couldn't refuse.~\LINK{http://naviteri.org/2012/01/more-additions-to-the-lexicon/}{b}
\B{Sto\-lo oel stxe\-nut pe\-yä.} I refused his offer.~\LINK{http://forum.learnnavi.org/language-updates/sto-has-the-same-syntax-as-new/}{ln}
(\emph{b. modal; either with} ‹\B{iv}› \emph{or} \HL{futa}{\B{fu\-ta}} \emph{and} ‹\B{iv}›) refuse (to do s.t.). 
\B{Sto\-lo po hi\-vum fo\-hu.} She refused to leave with them.~\LINK{http://forum.learnnavi.org/language-updates/sto-has-the-same-syntax-as-new/}{ln}
\B{Po\-el sto\-la\-tso fu\-ta me\-fo ti\-va\-ron tsa\-ha'\-ngir.} She must have refused their request to hunt that afternoon.~\LINK{http://forum.learnnavi.org/language-updates/sto-has-the-same-syntax-as-new/}{ln}

% stum
\hi
\HT{stum}{\B{stum}} \I{[stum]} \emph{or} \I{[stUm]} (\textsc{rn}: \B{stùm} \I{[stUm]})  \emph{adv.} (\emph{a.}) almost. 
\B{Poe\-ri unil\-tì\-ran\-tokx\-it tar\-mok a krr lam stum nì\-ay\-fo.} When she is in her Avatar body, she appears almost like them.~\LINK{http://naviteri.org/2010/07/ting-mikyun-fte-tslivam-listening-comprehension-1-3/}{b}
\B{Lì'\-fya\-ri tsa\-taw\-tu\-te pe\-yì?---Tìm\-yì. Pol ke tslam stum ke\-'ut, omum lì\-'ut a\-vol nì\-'aw.} How is that Sky Person's Na'vi?---Pretty bad. He understands almost nothing and only knows eight words.~\LINK{http://naviteri.org/2013/01/lifyari-po-peyi-how-good-is-her-navi/}{b}
(\emph{b. negated}) hardly, barely, scarcely.
\B{Ru\-txe pi\-vll\-txe nì\-wok nì\-'it, oel nga\-ti stum ke stä\-ngawm krr\-a nga fì\-fya tsye\-rì\-syì!} Please speak up a bit, I can barely hear you when you're whispering like this!~\LINK{http://naviteri.org/2011/09/miscellaneous-vocabulary/}{b}

% stxang
\hi 
\HT{stxang}{\B{stxang}} \I{[st'aN]} \emph{n.} axe, hatchet, tomahawk. 

% stxeli
\hi
\HT{stxeli}{\B{\cp{stxe}li}} \I{["{}st'E.li]} \emph{n.} gift, present.   
\B{Ftxo\-zä\-ri fì\-tì\-ngop a\-lor lu swey\-a stxe\-li a to\-lel oel.} This beautiful creation is the best present I have received for (my birthday) celebration.~\LINK{http://forum.learnnavi.org/community-calendar/karyu-pawls-birthday/msg312733/\#msg312733}{ln}
\B{Tsa\-tsko lu spu\-win ul\-te ke lu mi txur, slä oe\-ri txan\-wa\-we le\-iu. Lu stxe\-li a sem\-pul\-ta.} That bow is old and no longer strong, but it means a lot to me. It was a gift from my father.~\LINK{http://naviteri.org/2011/05/some-miscellaneous-vocabulary/}{b}
\B{Fì\-me\-stxe\-li a\-lor kur set ta kxem\-yo a mì hel\-ku moe\-yä.} These two beautiful gifts are now hanging on a wall in our home.~\LINK{http://naviteri.org/2014/09/stxeli-alor-the-text/}{b}
\B{Stxe\-lit fpo\-le' oel nga\-ru.} I sent the gift to you.~\LINK{http://forum.learnnavi.org/language-updates/i-have-something-to-say/}{ln}
\B{Nge\-yä stxe\-li\-ri a\-lor oe new nga\-ru pi\-vll\-txe san ira\-yo sìk nì\-li.} I want to thank you in advance for your beautiful gift.~\LINK{http://naviteri.org/2011/05/weather-part-2-and-a-bit-more-2/}{b}
 \\
\textsc{Derivations}: \ding{43}
\on{1}.~\HL{stxenu}{\B{stxe\-nu}}, offer.

% stxenu
\hi
\HT{stxenu}{\B{\cp{stxe}nu}} \I{["{}stE.nu]} \emph{n.} offer. 
\B{Ngal lum\-pe oe\-yä stxe\-nut tsyo\-lä\-ngär?} Why did you reject my offer?~\LINK{http://naviteri.org/2011/09/miscellaneous-vocabulary/}{b}
\B{Po pol\-txe san tì\-tstew\-nga'\-a stxe\-nu\-ri ira\-yo sei\-yi oe nga\-ru nì\-txan.} She said, ``I thank you very much for your courageous offer.''~\LINK{http://naviteri.org/2011/09/miscellaneous-vocabulary/}{b}
(\ding{43}~\HL{stxeli}{\B{stxe\-li}}, \HL{lonu}{\B{lo\-nu}}) \\
\textsc{Derivations}: \ding{43}
\on{1}.~\HL{stxenutIng}{\B{stxe\-nu\-tìng}}, offer (s.t.).

% stxenutìng
\hi
\HT{stxenutIng}{\B{\cp{stxe}nutìng}} \I{["{}stE.nu.t$\cdot$IN]} \emph{vtr., inf.}\on{33} offer (s.t.).
\B{Oe\-ri tì\-rey\-ti oel stxe\-nu\-to\-lìng fpi olo' aw\-nge\-yä.} I offered my life for the sake of our clan.~\LINK{http://naviteri.org/2011/09/miscellaneous-vocabulary/}{b}
\B{Stxe\-nu\-to\-lìng oel fu\-ta lì'\-fyat le\-Na'\-vi poe\-ru ki\-var. Mol\-'an nì\-prr\-te'.} I offered to teach her Na’vi. She accepted gladly.~\LINK{http://naviteri.org/2011/09/miscellaneous-vocabulary/}{b}
(\ding{43}~\HL{stxenu}{\B{stxe\-nu}}, \HL{tIng}{\B{tìng}})

% stxong
\hi
\HT{stxong}{\B{stxong}} \I{[st'oN]} \emph{adj.} strange, unfamiliar, unknown (\emph{usually negative connotation, potentially threatening or dangerous}).   
\B{Lar\-mu tsa\-tsam\-si\-yu\-hu tì\-vawm a lu stxong ay\-oer.} That warrior carried with him a darkness unknown to us.~\LINK{http://naviteri.org/2010/07/vocabulary-update/}{b}
\B{Fì\-na\-er\-ìri ngal ew\-ku 'u\-ot a\-stxong srak?} Do you taste something strange in this drink?~\LINK{http://naviteri.org/2012/11/renu-ayinanfyaya-the-senses-paradigm/}{b}

% sulìn
\hi
\HT{sulIn}{\B{\cp{su}lìn}} \I{["{}su.lIn]} \emph{vin.} be busy (\emph{positive sense}), \emph{be engrossed in s.t. one finds especially pleasant and energizing}.
\B{Pam\-tseo\-ri po su\-lä\-ngìn nì\-hawng.} He is overly engrossed in his music (and I'm displeased about it).~\LINK{http://naviteri.org/2010/09/getting-to-know-you-part-3/}{b}
(\ding{43}~\HL{'In}{\B{'ìn}}, \ding{214}~\HL{vrrIn}{\B{vrr\-ìn}}) \\
\textsc{Derivations}: \ding{43}
\on{1}.~\HL{sAsulIn}{\B{sä\-su\-lìn}}, hobby.

% sum
\hi
\HT{sum}{\B{sum}} \I{[sum]} \emph{or} \I{[sUm]} \emph{n.} shell (\emph{from the ocean}). 
(\ding{43}~\HL{lak}{\B{lak}}) \\
\textsc{Derivations}: \ding{43}
\on{1}.~\HL{sumsey}{\B{sum\-sey}}, drinking vessel (made of shell).

% sumsey
\hi
\HT{sumsey}{\B{\cp{sum}sey}} \I{["{}sum.sEj]} \emph{or} \I{["{}sUm.]} \emph{n.} drinking vessel (\emph{made of shell}).
(\ding{43}~\HL{sum}{\B{sum}}, \HL{sey}{\B{sey}})

% sunkesun
\hi 
\HT{sunkesun}{\B{\cp{sun}kesun}} \I{["{}sun.kE.sun]} \emph{or} \I{["{}sUn.kE.sUn]} \emph{adv.} ``like it or not'' (\emph{default addressee is ``you''}). 
\B{Sun\-ke\-sun po sla\-yu olo'\-eyk\-tan.} Whether you like it or not, he's going to become chief.~\LINK{http://naviteri.org/2019/06/50a-liu-amip-40-new-words/}{b} 
(\ding{43}~\HL{sunu}{\B{su\-nu}}, \HL{ke}{\B{ke}})

% sunu
\hi
\HT{sunu}{\B{\cp{su}nu}} \I{["{}su.nu]} \emph{vin.} be pleasing or likeable, bring enjoyment.   
\B{Su\-nu oe\-ru tey\-lu.} I like \emph{teylu}.~\LINK{http://naviteri.org/2010/07/vocabulary-update/}{b}
\B{Fra\-por a ke su\-nu por zop\-lo si.} He insults everyone he doesn't like.~\LINK{http://naviteri.org/2016/12/tengkrr-perahem-zisit-amip-as-the-new-year-arrives/}{b}
\B{Su\-nu oe\-ru fwa fì\-tseng\-it te\-rok oel nì\-mun.} I'm happy to be here again.~\LINK{http://naviteri.org/2010/08/a-na\%E2\%80\%99vi-alphabet/}{b}
\B{Tì\-'o'\-ìri pe\-u su\-nu ngar fra\-to?} What is your favorite way to have fun?~\LINK{http://naviteri.org/2010/09/getting-to-know-you-part-3/}{b}
\B{Ke\-tsran fya\-'o si\-vu\-nu ngar, kem si.} Do it however you'd like.~\LINK{http://naviteri.org/2013/03/whoever-whatever-whenever/}{b}
\B{Ul\-te sìl\-pey oe, fì\-'upxa\-ret inan a fì\-'u sil\-vu\-nu ay\-nga\-ru nì\-wotx!} And I hope, you all liked to read this message.~\LINK{http://naviteri.org/2011/09/\%E2\%80\%9Cby-the-way-what-are-you-reading\%E2\%80\%9D/}{b} \\
\textsc{Derivations}: \ding{43}
\on{1}.~\HL{txasunu}{\B{txa\-su\-nu}}, love greatly, enjoy tremendously. 
\on{2}.~\HL{sunkesun}{\B{sun\-ke\-sun}}, ``like it or not''.

% sung 
\hi
\HT{sung}{\B{sung}} \I{[suN]} \emph{or} \I{[sUN]} (\textsc{rn}: \B{sùng} \I{[sUN]}) \emph{vtr.} add. 
\B{Fko ze\-ne si\-vung tsat tsa\-tseng.} One must add that there.~\LINK{http://naviteri.org/2013/08/fmawn-a-fkol-kilmulat-this-just-in/comment-page-1/\#comment-2476}{b}
\B{Oel lum\-pe tsat so\-lung mì upxa\-re?} Why did I add that in the message?~\LINK{http://naviteri.org/2010/07/vocabulary-update/comment-page-1/\#comment-95}{b}
\B{Ru\-txe oe\-yä tì\-oeyk\-tìng\-it a\-saw\-nung ni\-vìn.} Please look at my added explanation.~\LINK{http://naviteri.org/2011/03/word-order-and-case-marking-with-modals/comment-page-1/\#comment-581}{b}
(i.) \B{Trr A\-saw\-nung}, Leap Day.~\LINK{http://naviteri.org/2012/02/trr-asawnung-lefpom-happy-leap-day/}{b} \\
\textsc{Derivations}: \ding{43}
\on{1}.~\HL{nIsung}{\B{nì\-sung}}, besides, additionally, furthermore. 
\on{2}.~\HL{tIsung}{\B{tì\-sung}}, addition, post script. 
\on{3}.~\HL{'opinsung}{\B{'o\-pin\-sung}}, color, color in.

% sur
\hi
\HT{sur}{\B{sur}} \I{[suR]} \emph{or} \I{[sUR]} \emph{n.} (\emph{a. sensation}) taste, flavor. 
(\emph{b. to form the middle voice}) \B{Fì\-na\-er\-ìri sur fkan oe\-ru ka\-lin.} This drink tastes sweet to me.~\LINK{http://naviteri.org/2012/11/renu-ayinanfyaya-the-senses-paradigm/}{b}
\B{Fì\-na\-er\-ìri sur fkan lor.} This drink tastes good.~\LINK{http://naviteri.org/2012/11/renu-ayinanfyaya-the-senses-paradigm/}{b}

% susyang
\hi
\HT{susyang}{\B{su\cp{syang}}} \I{[su."sjaN]} \emph{adj.} fragile, delicate. 
\B{Lu fì\-vul su\-syang nì\-txan. Txo fko ti\-vìng zek\-wä kxakx.} This branch is very fragile. If you touch it, it'll break.~\LINK{http://naviteri.org/2015/04/some-new-words-for-may-day/}{b}
\B{Rey\-'eng lu su\-syang.} The balance of life is fragile.~\LINK{http://naviteri.org/2015/04/some-new-words-for-may-day/}{b}

% sutx
\hi
\HT{sutx}{\B{sutx}} \I{[sut']} \emph{or} \I{[sUt']} \emph{vtr.} track, follow. 
\B{Po spä few pay\-fya fte smar\-it si\-vutx.} He jumped across the stream to track his prey.~\LINK{http://naviteri.org/2011/01/new-year-new-vocabulary/}{b}
\B{Tì\-tu\-sa\-ron\-ìri fte fli\-vä, ze\-ne fko si\-vutx smar\-it ni\-nan nì\-no.} To succeed at hunting, you have to track your prey by reading (the forest) with attention to detail.~\LINK{http://naviteri.org/2011/09/\%E2\%80\%9Cby-the-way-what-are-you-reading\%E2\%80\%9D/}{b} \\
\textsc{Derivations}: \ding{43}
\on{1}.~\HL{fngapsutxwll}{\B{fngap\-sutx\-wll}}, metalfollowing plant.

% swapxì
\hi 
\HT{swapxI}{\B{swa\cp{pxì}}} \I{[swa."p'I]} \emph{n.} family member. 
\B{Ay\-swapx\-ìl oe\-yä tok fì\-tseng\-it nì\-wotx.} All the members of my family are here.~\LINK{http://naviteri.org/2021/09/aawa-ayliu-amip-a-few-new-words/}{b}
(\ding{43}~\HL{soaia}{\B{soaia}}, \HL{hapxì}{\B{ha\-pxì}}) 

% swaran
\hi 
\HT{swaran}{\B{\cp{swa}ran}} \I{["{s}wa.Ran]} \emph{adj.} humble, modest, self\hyp{}effacing. 
\B{Tsam\-si\-yu a\-sìl\-tsan lu tstew slä\-kop swa\-ran.} A good warrior is courageous but also humble.~\LINK{http://naviteri.org/2020/11/vospxivopeya-ayliu-amip-novembers-new-words/}{b} \\
\textsc{Derivations}: \ding{43}
\on{1}.~\HL{tIswaran}{\B{tì\-swa\-ran}}, humility, humbleness.

% swaw
\hi
\HT{swaw}{\B{swaw}} \I{[swaw]} \emph{n.} moment. 
\B{Mì na'\-rìng Tsyeyk\-ìl lo\-re\-yu\-ti 'o\-lam\-pi a tsa\-swaw wou oer nì\-wotx!} That instant when Jake touched the \emph{Helicoradiam Spirale} in the forest I was totally amazed.~\LINK{http://forum.learnnavi.org/language-updates/txelanit-hivawl/}{ln} \\
\textsc{Derivations}: \ding{43}
\on{1}.~\HL{pxisawam}{\B{pxi\-swaw\-am}}, just a moment ago. 
\on{2}.~\HL{pxiswaway}{\B{pxi\-swaw\-ay}}, in just second (from now). 
\on{3}.~\HL{swawtsyIp}{\B{swaw\-tsyìp}}, little moment. 

% swawtsyìp 
\hi
\HT{swawtsyIp}{\B{\cp{swaw}tsyìp}} \I{["{}swaw.\t{ts}jIp\textcorner]} \emph{n.} little moment, second.
\B{'Aw\-a swaw\-tsyìp. Oe fpe\-rìl.} Just a moment. I'm thinking.~\LINK{http://naviteri.org/2010/09/getting-to-know-you-part-3/}{b}
\B{'Aw\-a swaw\-tsyìp li\-vu oer.} May I have a tiny moment.~\LINK{http://naviteri.org/2010/09/getting-to-know-you-part-3/}{b}
(\ding{43}~\HL{swaw}{\B{swaw}}, \HL{-tsyIp}{\B{-tsyìp}})

% swaynivi
\hi
\HT{swaynivi}{\B{\cp{sway}nivi}} \I{["{}swaj.ni.vi]} \emph{n.} family hammock. 
(\ding{43}~\HL{soaia}{\B{soaia}}, \HL{nivi}{\B{ni\-vi}})

% swek
\hi
\HT{swek}{\B{swek}} \I{[swEk\textcorner]} \emph{n.} bar, rod, pole.
\B{Me\-tew\-it fì\-swek\-ä ey\-ki\-vil.} Pull the two ends of this bar together.~\LINK{http://naviteri.org/2013/04/wheres-the-bathroom-and-other-useful-things/}{b}

% swey
\hi
\HT{swey}{\B{swey}} \I{[swEj]} \emph{adj.} best.
\B{Fì\-srey puk\-ä lu swey nì\-law.} This version of the book is clearly the best.~\LINK{http://naviteri.org/2014/07/tsawlultxamaw-a-ayliu-post-meetup-words/}{b}
\B{Tì\-'eyng\-it oel to\-lel a krr, tsa\-krr pa\-ye\-'un swey\-a fya\-'ot a za\-mi\-vu\-nge oel ay\-ngar ay\-lì\-'ut ho\-ren\-ti\-sì lì'\-fya\-yä le\-Na'\-vi.} When I receive an answer, I will then decide the best way to bring you the words and rules of Na'vi.~\LINK{http://forum.learnnavi.org/news-announcements/a-response-from-paul-frommer!/}{ln}
\B{Ay\-lì\-'u\-van a\-swey lu 'i\-pu, lu sìl\-ron\-sem.} The best puns are both funny and clever.~\LINK{http://naviteri.org/2012/01/mipa-zisit-ayliu-amip-new-words-for-the-new-year/}{b}
\B{Swey lu fwa nga fì\-kem ke si\-vi nì\-sngum.} It's best that you not do this fretfully.\slash It's best that you don't freak out about doing this.~\LINK{http://naviteri.org/2011/01/new-year-new-vocabulary/}{b}
(\ding{214}~\HL{'e'al}{\B{'e'al}}) \\
\textsc{Derivations}: \ding{43}
\on{1}.~\HL{nIswey}{\B{nì\-swey}}, optimally. 
\on{2}.~\HL{sweylu}{\B{swey\-lu}}, should.

% sweylu
\hi
\HT{sweylu}{\B{\cp{swey}lu}} \I{["{}swEj.lu]} \emph{vin.} should, ``it's best (that)''.
(\emph{a. for s.t. that hasn't yet happened with} \HL{txo}{\B{txo}} \emph{and} ‹\B{iv}›)
\B{Fì\-na\-er ftxì\-vä' lu nì\-hawng, ha swey\-lu txo ngal tsat kxi\-vukx nì\-win.} This drink tastes horrible, so you'd better swallow it quickly.~\LINK{http://naviteri.org/2011/11/\%E2\%80\%99a\%E2\%80\%99awa-ayli\%E2\%80\%99u-amip-ni\%E2\%80\%99aw-\%E2\%80\%94-a-few-new-words-only/}{b}
\B{Swey\-lu txo nga ki\-vä.\slash Nga swey\-lu txo ki\-vä.} You should go.~\LINK{http://naviteri.org/2011/04/\%E2\%80\%99a\%E2\%80\%99awa-li\%E2\%80\%99fyavi-amip\%E2\%80\%94a-few-new-expressions/}{b}
\B{Swey\-lu txo nga ke ki\-vä.\slash Nga swey\-lu txo ke ki\-vä.} You shouldn't go.~\LINK{http://naviteri.org/2011/04/\%E2\%80\%99a\%E2\%80\%99awa-li\%E2\%80\%99fyavi-amip\%E2\%80\%94a-few-new-expressions/}{b}
\B{Swey\-lu txo ki\-vä.} I\slash you\slash she\slash one\slash etc. should go.~\LINK{http://naviteri.org/2011/04/\%E2\%80\%99a\%E2\%80\%99awa-li\%E2\%80\%99fyavi-amip\%E2\%80\%94a-few-new-expressions/}{b}
(\emph{b. for s.t. that's already happened with} \HL{fwa}{\B{fwa}} \emph{or} \HL{tsawa}{\B{tsawa}} \emph{and the past or perfective indicative})
\B{Swey\-lu fwa nga ko\-lä.} You should have gone.~\LINK{http://naviteri.org/2011/04/\%E2\%80\%99a\%E2\%80\%99awa-li\%E2\%80\%99fyavi-amip\%E2\%80\%94a-few-new-expressions/}{b}
(\ding{43}~\HL{swey}{\B{swey}}, \HL{lu}{\B{lu}})

% sweyn
\hi 
\HT{sweyn}{\B{sweyn}} \I{[swEjn]} \emph{vtr.} (\emph{a.}) keep, preserve. 
\B{Ay\-ngal syu\-vet sweyn pe\-seng fte\-ke ay\-ioang tsi\-vun tsat ki\-va\-nom?} Where do you keep the food so that animals can't get it?~\LINK{http://naviteri.org/2021/09/aawa-ayliu-amip-a-few-new-words/}{b}  
(\emph{b.}) leave alone, not disturb.
\B{Tsa\-ya\-yo\-tsrul\-it sweyn, ma 'i\-tan.} Don't disturb that bird's nest, son. ~\LINK{http://naviteri.org/2021/09/aawa-ayliu-amip-a-few-new-words/}{b} 
\B{Oey fpom\-it sweyn!} Leave me alone!\slash Do not disturb my peace!~\LINK{http://naviteri.org/2021/09/aawa-ayliu-amip-a-few-new-words/}{b} 

% Sweriye
% \hi
% \B{Sweriye} \I{[]} \emph{n.} Sweden. 

% swirä
\hi
\HT{swirA}{\B{swi\cp{rä}}} \I{[swi."R\ae]} \emph{n.} creature. 
\B{Ye\-rik lu swi\-rä a\-nim nì\-txan.} The hexapede is a very timid creature.~\LINK{http://naviteri.org/2011/09/miscellaneous-vocabulary/}{b}
\B{Tsa\-swi\-rä\-ti lo\-nu! Ay\-nga ne\-to ri\-vikx! Fì\-ke\-tu\-wong\-ti oel stì\-ye\-ftxaw.} Release this creature! Step back! I will look at this alien.~\LINK{http://wiki.learnnavi.org/index.php?title=Na\%27vi_from_Avatar_Movie\#Scene_in_Hometree}{wiki\slash a}

% swizaw
\hi
\HT{swizaw}{\B{swi\cp{zaw}}} \I{[swi."zaw]} \emph{n.} arrow.
\B{Oe\-yä swi\-zaw nì\-ngay ti\-va\-kuk.} Let my arrow strike true.~\LINK{http://wiki.learnnavi.org/index.php/Hunting_Song}{wiki}
\B{Nì\-trr\-trr pxì\-mun\-'i sam\-si\-yul ay\-swi\-zaw\-it ku\-tu\-ä a\-law\-nä\-txayn sno\-kip nì\-'eng.} Warriors typically share the arrows of their defeated enemies among themselves.~\LINK{http://naviteri.org/2012/01/more-additions-to-the-lexicon/}{b}
\B{Ta\-yel Tse\-nul pxe\-swi\-zaw\-ti yoa mun\-sna\-hawn\-ven.} Tsenu will receive three arrows in exchange for a pair of shoes.~\LINK{http://naviteri.org/2014/02/barter-and-exchange/}{b}
(\emph{a. with} \HL{tsweykayon}{\B{tswey\-ka\-yon}}) shoot an arrow.
\B{Swi\-zaw\-ti tswey\-kay\-on ne\-fä ne taw, tse\-nga zup ke lu law.} I shot an arrow into the air, | It fell to earth, I know not where.~\LINK{http://naviteri.org/2017/03/titusemteri-concerning-shooting/}{b} 
(i.) \B{tsko swizaw}, bow and arrow, archery.
\B{Na'\-vi\-ri txin\-a sä\-'o tì\-tu\-sa\-ron\-ä lu tsko swi\-zaw.} For the Na'vi the bow and arrow is the main hunting tool.~\LINK{http://naviteri.org/2011/01/new-year-new-vocabulary/}{b}
\B{Tsko swi\-zaw\-fa a tì\-ta\-ron\-ìri ze\-ne fko ni\-vu\-me.} You must learn to hunt with a bow and arrow.~\LINK{http://forum.learnnavi.org/language-updates/ordinals-nume/}{ln\slash a\textsuperscript{c}}
\B{Oe tskxe\-keng si sä\-su\-lìn\-ur alu tsko swi\-zaw.} I practice my hobby, which is archery.~\LINK{http://naviteri.org/2010/09/getting-to-know-you-part-3/}{b}
(\emph{a. proverb}) \B{Fwa kan ke tam; ze\-ne swi\-zaw\-it li\-vo\-nu.} Intent is not enough; it's action that counts. [lit.: To aim is not enough; one must release the arrow.]~\LINK{http://naviteri.org/2013/02/vospxi-ayol-posti-apup-short-post-for-a-short-month/}{b} \\
\textsc{Derivations}: \ding{43} 
\on{1}.~\HL{swizawsena}{\B{swi\-zaw\-se\-na}}, quiver. 
\on{2}.~\HL{tsko}{\B{tsko swi\-zaw}}, bow and arrow, archery. 
\on{3}.~\HL{atanzaw}{\B{atan\-zaw}}, forked lightning.
\on{4}.~\HL{winzaw}{\B{win\-zaw}}, arrow deer.

% swizawsena
\hi
\HT{swizawsena}{\B{swi\cp{zaw}sena}} \I{[swi."zaw.sE.na]} \emph{n.} quiver (\emph{attached to the ikran's saddle}), \emph{coll. often shortened to} \HL{zawsena}{\B{zaw\-se\-na}}.
(\ding{43}~\HL{swizaw}{\B{swi\-zaw}}, \HL{sAhena}{\B{sä\-he\-na}})

% swoa
\hi
\HT{swoa}{\B{\cp{swo}a}} \I{["{}swo.a]} \emph{n.} intoxicating beverage.
\B{Oe\-ri lu swoa fne\-txum.} I'm allergic to intoxicating beverage.~\LINK{http://forum.learnnavi.org/language-updates/new-words-from-the-2012-avatarmeet-navi-class/msg551479/\#msg551479}{b}   
\B{Sra\-ne, fì\-txon tsun nga ni\-väk swo\-at nì\-'it, slä rä\-'ä rou!} Yes, you can have a little alcohol tonight, but don't get drunk!~\LINK{http://naviteri.org/2012/10/audio-and-video-learning-materials-for-navi-101/}{b} \\
\textsc{Derivations}: \ding{43}
\on{1}.~\HL{swoasey}{\B{swoa\-sey}}, kava bowl.

% swoasey
\hi
\HT{swoasey}{\B{\cp{swo}asey}} \I{["{}swo.a.sEj]} \emph{n.} (hand\hyp{}sized) kava bowl (\emph{constructed from seed pods, used for drinking intoxicating beverages}). \\
(i.) \B{swoasey ayll} \emph{n.} large, social kava bowl. 
(\ding{43}~\HL{swoa}{\B{swoa}}, \HL{sey}{\B{sey}})

% swok
\hi
\HT{swok}{\B{swok}} \I{[swok\textcorner]} \emph{adj.} sacred. 
\B{Ma txan\-tslu\-sam\-a ngawng, a swok\-a ut\-ral\-ti yom.} O wise worm, who eats the Sacred Tree.~\LINK{https://wiki.learnnavi.org/Na\%27vi_from_Avatar_Movie\#Becoming_one_of_the_Na.27vi}{a\textsuperscript{c}}

% swotu
\hi
\HT{swotu}{\B{\cp{swo}tu}} \I{["{}swo.tu]} \emph{n.} sacred place.
\B{He\-tu\-wong\-ìl aw\-nge\-yä swo\-tut ska\-'a, fte kll\-ki\-vu\-lat ke\-ru\-sey\-a tskxet.} The aliens destroy our sacred place to dig up dead rock.~\LINK{http://forum.learnnavi.org/language-updates/participles/}{ln\slash a\textsuperscript{c}}

% Swotu Wayä
\hi 
\HT{swotu wayA}{\B{\cp{Swo}tu \cp{Way}ä}} \I{["{}swo.tu waj.\ae]} \emph{n.} Sacred Place of Song, \emph{situated in the Valley of Mo'ara}. 
(\ding{43}~\HL{swotu}{\B{swo\-tu}}, \HL{way}{\B{way}})

% swotulu
\hi
\HT{swotulu}{\B{\cp{swo}tulu}} \I{["{}swo.tu.lu]} \emph{n.} \emph{name of a sacred river}. 

% syaw
\hi
\HT{syaw}{\B{syaw}} \I{[sjaw]} \emph{vin.} call, name. 
\B{Fya\-pe fko syaw ngar? -- Oe\-ru syaw (fko) Ney\-ti\-ri. Nga\-ru tut? -- Oe\-ru syaw Tseyk.} What's your name? -- I'm Neytiri. How about you? -- I'm Jake.~\LINK{http://forum.learnnavi.org/language-updates/the-reflexive-\%28from-a-frommer-email!!!!\%29/}{ln}
\B{Fì\-por syaw fko Ìstaw.} This is Ìstaw.~\LINK{http://naviteri.org/2010/09/getting-to-know-you-part-1/}{b}
\B{Syaw oer ru\-txe trr\-ay fa spul\-mok\-ri.} Please phone me tomorrow.~\LINK{http://naviteri.org/2012/07/fivospxiya-aylifyavi-amip-this-months-new-expressions-pt-2/}{b}

% syawm
\hi 
\HT{syawm}{\B{syawm}} \I{[sjawm]} \emph{vtr.}, \textsc{rn} know. 
% \B{Ayoe ke syawm.} We don't know.~\LINK{}{a\on{2}} 
(\ding{43}~\HL{omum}{\B{o\-mum}}) 

% syawn
\hi 
\HT{syawn}{\B{syawn}} \I{[sjawn]} \textsc{i}.~\emph{n.} blessing. 
\B{Ne\-wey yaw\-ne lu oer ul\-te new oe mun\-txa si\-vi poe\-hu. Ru\-txe, ma sem\-pul, tìng mo\-er nge\-yä syawn\-it.} I love Ne\-wey and want to marry her. Please, father, give us your blessing.~\LINK{http://naviteri.org/2018/03/100a-liu-amip-64-new-words-part-1/}{b}
\B{Eywal tivìng ngar sneyä syawnit.} May Ey\-wa bless you. \\
\textsc{ii}.~\B{tìng syawn} \I{[t$\cdot$IN syawn]} \emph{vin.} bless.
\B{Eywa ti\-vìng syawn ngar, ma 'i\-te.} May Ey\-wa bless you, my daughter.~\LINK{http://naviteri.org/2018/03/100a-liu-amip-64-new-words-part-1/}{b} 
(\ding{214}~\HL{rAzekx}{\B{rä\-zekx}}) 

% syay
\hi
\HT{syay}{\B{syay}} \I{[sjaj]} \emph{n.} fate, destiny, the arc of one's life. 
\B{Tsa\-krr syay ay\-nge\-yä, syay o\-lo'\-ä oe\-yä la\-yu teng.} Then you will suffer the same fate as my clan.~\LINK{http://naviteri.org/2010/07/vocabulary-update/}{b} \\
\textsc{Derivations}: \ding{43}
\on{1}.~\HL{syayvi}{\B{syay\-vi}}, luck, chance.

% syayvi
\hi
\HT{syayvi}{\B{\cp{syay}vi}} \I{["{}sjaj.vi]} \emph{n.} luck, chance.
\B{Et\-rìp\-a syay\-vi!} Good luck!~\LINK{http://naviteri.org/2010/07/vocabulary-update/}{b} \\
\on{1.} \B{nìsyayvi}, by chance, by coincidence.
(\ding{43}~\HL{syay}{\B{syay}}, \HL{-vi}{\B{-vi}})

% syaksyuk
\hi
\HT{syaksyuk}{\B{syak\cp{syuk}}} \I{[sjak\textcorner."{}sjuk\textcorner]} \emph{or} \I{[."{}sjUk\textcorner]} (\textsc{rn}: \B{syak\cp{syùk}} \I{[Sak\textcorner."SUk\textcorner]}) \emph{n., fauna} prolemuris, \emph{prolemuris noctis}.

% syam
\hi
\HT{syam}{\B{syam}} \I{[sjam]} \emph{vtr.} smell (\texttt{+}\emph{control}). 
\B{Fì\-syu\-lang\-it syam. Fa\-hew lor lu nì\-txan, ke\-fyak?} Smell this flower. Its fragrance is beautiful, isn't it?~\LINK{http://naviteri.org/2012/11/renu-ayinanfyaya-the-senses-paradigm/}{b}
(\ding{43}~\HL{hefi}{\B{he\-fi}})

% syanan
\hi
\HT{syanan}{\B{\cp{sya}nan}} \I{["{}sja.nan]} \emph{countable n.} a type of waterfall, \emph{a single drop or series of smaller falls occurring sequentially along a stream or series of pools}.

% syar
\hi 
\HT{syar}{\B{syar}} \I{[sjar]} \emph{vin.} stick, stick to, adhere. 
\B{Rìk a\-'aw syar\-mar sìn kxem\-yo.} A leaf was sticking to the wall.~\LINK{http://naviteri.org/2020/11/vospxivopeya-ayliu-amip-novembers-new-words/}{b}
\B{Pol kxum\-pay\-syar\-it so\-lar syey\-kar rìk\-it sìn kxem\-yo.} She stuck the leaf onto the wall with glue. [lit.: ``She used glue (and then) stuck the leaf onto the wall.'']~\LINK{http://naviteri.org/2020/11/vospxivopeya-ayliu-amip-novembers-new-words/}{b} \\
\textsc{Derivations}: \ding{43}
\on{1}.~\HL{kxumpaysyar}{\B{kxumpaysyar}}, glue.

% syä'ä
\hi
\HT{syA'A}{\B{\cp{syä}'ä}} \I{["{}sj\ae.P\ae]} \emph{adj.}(of taste) bitter. 
\B{Ton\-ìri a\-la\-he, aw\-ngal yom fkxen\-ti le\-rìk nì\-wotx; fì\-txon yom sat a lu syä\-'ä nì\-'aw.} As for other nights, we eat all leafy vegetable; this night we eat only those that are bitter.~\LINK{http://whyisthisnight.com/addl-more.html\#na\%27vi}{tn} 
\B{To\-lìng ay\-oer ma\-u\-tit a\-syä\-'ä.} (It) has given us bitter fruit.~\LINK{http://naviteri.org/2020/06/mi-tanlokxe-oeya-srr-afpxamo-terrible-days-in-my-country/}{b}

% sye'otxang
\hi 
\HT{sye'otxang}{\B{\cp{sye}'otxang}} \I{["{}sjE.Po.t'aN]} \emph{n.} wind instrument. 
(\ding{43}~\HL{syeha}{\B{sye\-ha}}, \HL{'otxang}{\B{'o\-txang}})

% syeha
\hi
\HT{syeha}{\B{\cp{sye}ha}} \I{["{}sjE.ha]} \textsc{i}. \emph{n.} breath.  \\
\textsc{ii}. \B{\cp{sye}ha si} \I{["{}sjE.ha si]} \emph{vin.} breathe.   
\B{Ma Ra\-lu, srung si por! Nì\-win! Sye\-ha ke si!} Ralu, help her! Quick! She's not breathing!~\LINK{http://naviteri.org/2012/03/spring-vocabulary-part-2/}{b} \\
\textsc{Derivations}: \ding{43}
\on{1}.~\HL{snewsye}{\B{snew\-sye}}, weird, spooky. 
\on{2}.~\HL{sye'otxang}{\B{sye\-'o\-txang}}, wind instrument.
\on{3}.~\HL{mungsye}{\B{mung\-sye}}, inhale.
\on{4}.~\HL{lonusye}{\B{lo\-nu\-sye}}, exhale, blow.

% syen
\hi
\HT{syen}{\B{syen}} \I{[sjEn]} \emph{adj.} last, final (\emph{in a series}). 
\B{tì\-pawm a\-syen}, last question.~\LINK{http://naviteri.org/2010/08/mipa-ayopin-mipa-ayliu-new-colors-new-words/}{b}
(\ding{43}~\HL{tor}{\B{tor}}) \\
\textsc{Derivations}: \ding{43}
\on{1}.~\HL{nIsyen}{\B{nì\-syen}}, lastly, finally.

% syep
\hi
\HT{syep}{\B{syep}} \I{[sjEp\textcorner]} \emph{vtr.} trap (s.o.\slash s.t.).
\B{Fì\-hì\-'ang\-ìl txu\-la hän\-ti fte smar\-it syi\-vep.} This insect constructs a web to trap its prey.~\LINK{http://naviteri.org/2015/04/some-new-words-for-may-day/}{b} \\
\textsc{Derivation}: \ding{43} \on{1}.~\HL{sAsyep}{\B{sä\-syep}}, a trap. 
\on{2}.~\HL{nIsyep}{\B{nì\-syep}}, tightly, in an iron grip. 
\on{3}.~\HL{syeptute}{\B{syep\-tu\-te}}, hyneman. 
\on{4}.~\HL{syeprel}{\B{syep\-rel}}, camera.
\on{5}.~\HL{syep'an}{\B{syep\-'an}}, lift vine.

% syep'an
\hi 
\HT{syep'an}{\B{\cp{syep}'an}} \I{["{}sjEp\textcorner.Pan]} \emph{n.} \ding{96} lift vine [lit.: `trap vine']. 
(\ding{43}~\HL{syep}{\B{syep}}, \HL{'ana}{\B{'a\-na}})

% syeprel
\hi
\HT{syeprel}{\B{syep\cp{rel}}} \I{[sjEp\textcorner."REl]} \emph{n.} camera.
\B{Lu oe\-ru syep\-rel.} I have a camera.~\LINK{http://naviteri.org/2012/10/audio-and-video-learning-materials-for-navi-101/}{b}
(\ding{43}~\HL{syep}{\B{syep}}, \HL{rel}{\B{rel}})

% syeptute
\hi
\HT{syeptute}{\B{\cp{syep}tute}} \I{["{}sjEp\textcorner.tu.tE]} \emph{n.}, \ding{96} hyneman, ``person trapper'' (\emph{type of carnivorous plant}), \emph{pandoratonia myopora}. 
\B{Oel hu Txe\-wì trr\-am na'\-rìng\-it tar\-mok, tso\-le\-'a syep\-tu\-tet a\-tsawl fra\-to mì sì\-rey.} Yesterday I was with Txewì in the forest and we saw the biggest Trapper I've ever seen.~\LINK{http://wiki.learnnavi.org/index.php?title=Corpus\#New_York_Times}{wiki}
(\ding{43}~\HL{syep}{\B{syep}}, \HL{tute}{\B{tu\-te}})

% syewe
\hi 
\HT{syewe}{\B{\cp{sye}we}} \I{["{}sjE.wE]} \emph{n.} fat (substance in meat). \\
\textsc{Derivations}: \ding{43}
\on{1}.~\HL{syewenga'}{\B{sye\-we\-nga'}}, fatty.

% syewenga'
\hi 
\HT{syewenga'}{\B{\cp{sye}wenga'}} \I{["{}sjE.wE.NaP]} \emph{adj.} fatty. 
\B{Po\-an\-ur su\-nu tsngan a\-sye\-we\-nga'; poe\-ru ke su\-nu kaw\-'it.} He likes fatty meat; she doesn't like it at all.~\LINK{http://naviteri.org/2016/06/mrrvola-lifyavi-amip-forty-new-expressions/}{b} 
(\ding{43}~\HL{syewe}{\B{sye\-we}})

% syìmawnun'i
\hi 
\HT{syImawnun'i}{\B{syìmawnun\cp{'i}}} \I{[sjI.maw.nun."Pi]} \emph{or} \I{[.nUn.]} \emph{n.} cut tie, dissolved connection. 
\B{Oe\-nga\-ri sä\-ta\-re syì\-maw\-nun\-'i slo\-lä\-ngu.} As for the two of us, I'm sorry to say our relationship is dissolved.~\LINK{http://naviteri.org/2023/07/trr-tsyimawnuniya-lefpom-happy-independence-day/}{b} 
(\ding{43}~\HL{sAyIm}{\B{sä\-yìm}}, \HL{mun'i}{\B{mun\-'i}}) \\
\textsc{Derivations}: \ding{43} 
\on{1}.~\HL{tIsyImawnun'i}{\B{tì\-syì\-maw\-nun\-'i}}, independence. 

% syo
\hi
\HT{syo}{\B{syo}}\textsuperscript{\on{1}} \I{[sjo]} \emph{adj.} light, lightweight.
(\ding{214}~\HL{ku'up}{\B{ku'\-up}}) \\
\textsc{Derivations}: \ding{43}
\on{1}.~\HL{syokup}{\B{syo\-kup}}, weight.

\hi
\B{syo}\textsuperscript{\on{2}} : \HL{tsyo}{\B{tsyo}}

% syokup
\hi 
\HT{syokup}{\B{syo\cp{kup}}} \I{[sjo."kup\textcorner]} \emph{or} \I{[."kUp\textcorner]} (\textsc{rn}: \B{syo\cp{kùp}} \I{[So."kUp\textcorner]}) \emph{n.} weight (physical). 
(\ding{43}~\HL{syo}{\B{syo}}, \HL{ku'up}{\B{ku'\-up}})

% syon
\hi
\HT{syon}{\B{syon}} \I{[sjon]} \emph{n.} feature, trait, attribute, characteristic, point, aspect, facet, property. 
\B{Tsran\-ten fra\-to a syon tsam\-si\-yu\-ä lu tì\-tstew.} The most important characteristic of a warrior is bravery.~\LINK{http://naviteri.org/2012/10/mipa-vospxi-mipa-ayliu-new-words-for-the-new-month/}{b}  
\B{Pa\-lu\-lu\-kan\-ìri lu pxe\-syon a ze\-ne fko zi\-ve\-rok nì\-tut.} Three things about the thanator must always be kept in mind.~\LINK{http://naviteri.org/2012/10/mipa-vospxi-mipa-ayliu-new-words-for-the-new-month/}{b}
\B{Fra\-'u\-ä ran ngä\-pop fa fra\-syon tse\-yä.} The \emph{ran} of each thing arises from the totality of its attributes.~\LINK{http://naviteri.org/2012/10/mipa-vospxi-mipa-ayliu-new-words-for-the-new-month/}{b} \\
\textsc{Derivations}: \ding{43}
\on{1}.~\HL{syonlI'u}{\B{syon\-lì\-'u}}, adjective.
\on{2}.~\HL{ketengsyon}{\B{ke\-teng\-syon}}, variety, diversity.

% syonlì'u
\hi 
\HT{syonlI'u}{\B{\cp{syon}lì'u}} \I{["{}sjon.lI.Pu]} \emph{n.} adjective. 
\B{Kehe, ke tsun\-slu fwa yem me\-syon\-lì\-'ut uo tstxo\-lì\-'u.} No, it's not possible to put two adjectives after a noun.~\LINK{http://naviteri.org/2017/12/zisit-amip-lefpom-ma-eylan-happy-new-year-friends/\#comment-27964}{b} 
(\ding{43}~\HL{syon}{\B{syon}}, \HL{lI'u}{\B{lì\-'u}})

% syor
\hi
\HT{syor}{\B{syor}} \I{[sjoR]} \emph{vin.} relax, chill out.
\B{Syor, ma 'ey\-lan, syor. Ke lu kea sngum\-tsim.} Relax, friend, relax. There's nothing to worry about.~\LINK{http://naviteri.org/2013/01/awvea-posti-zisita-amip-first-post-of-the-new-year/}{b}
\B{New oe ri\-vun a\-sim tì\-fnu\-nga'\-a tseng\-it a tsa\-ro tsun syi\-vor tsi\-vu\-rokx fte spä\-pi\-veng.} I want to find a quiet place nearby where I can chill out and rest to get my head back on straight.~\LINK{http://naviteri.org/2013/01/awvea-posti-zisita-amip-first-post-of-the-new-year/}{b}
\B{Fì\-u\-van\-ìl oe\-ti syey\-kor nì\-txan.} I find this game very relaxing.~\LINK{http://naviteri.org/2013/01/awvea-posti-zisita-amip-first-post-of-the-new-year/}{b} \\
\textsc{Derivations}: \ding{43}
\on{1}.~\HL{tIsyor}{\B{tì\-syor}}, relaxation.

% syulang
\hi
\HT{syulang}{\B{\cp{syu}lang}} \I{["{}sju.lang]} \emph{n.}, \ding{96} flower. 
\B{Fì\-syu\-lang a\-rim lu hì\-'i fra\-to.} This yellow flower is the smallest of all.~\LINK{http://naviteri.org/2010/08/mipa-ayopin-mipa-ayliu-new-colors-new-words/}{b}
\B{Oe\-yä ay\-syu\-lang set mì fay kll\-kxe\-rem.} My flowers are standing in water now.~\LINK{http://naviteri.org/2010/09/getting-to-know-you-part-1/comment-page-1/\#comment-296}{b}
\B{Ra\-lul ngo\-lop 'um\-tsat fa ay\-syu\-lang fwä\-kì\-wll\-ä.} Ralu made medicine from (\emph{or}: using the) flowers of the Mantis orchid.~\LINK{http://naviteri.org/2014/11/twenty-before-the-holidays/}{b} 
\B{Fì\-syu\-lang\-it syam. Fa\-hew lor lu nì\-txan, ke\-fyak?} Smell this flower. Its fragrance is beautiful, isn't it?~\LINK{http://naviteri.org/2012/11/renu-ayinanfyaya-the-senses-paradigm/}{b}
(i.) \B{sä\-star\-sìm syu\-lang\-ä}, bouquet (\emph{collection of flowers selected and put together intentionally by a person}).~\LINK{http://naviteri.org/2012/03/spring-vocabulary-part-2/}{b} \\
\textsc{Derivations}: \ding{43}
\on{1}.~\HL{prrnesyul}{\B{prr\-ne\-syul}}, bud (of a flower).
\on{2}.~\HL{paysyul}{\B{pay\-syul}}, water lily.
\on{3}.~\HL{tsawksyul}{\B{tsawk\-syul}}, sun lily.
\on{4}.~\HL{tarsyu}{\B{tar\-syu}}, tarsyu.

% syura
\hi
\HT{syura}{\B{syu\cp{ra}}} \I{[sju."Ra]} \emph{n.} energy (\emph{physical and spiritual}).
\B{Fra\-syu\-ra\-ti fkol za\-sro\-lìn nì\-'aw ul\-te trr\-o ze\-ne tey\-ki\-vä\-txaw.} All energy is only borrowed, and one day it will have to be given back.~\LINK{http://naviteri.org/2012/03/spring-vocabulary-part-1/}{b} \\
\textsc{Derivations}: \ding{43}
\on{1}.~\HL{syuratan}{\B{syuratan}}, bioluminescence.

% syuratan
\hi
\HT{syuratan}{\B{syu\cp{ra}tan}} \I{[sju."Ra.tan]} \emph{n.} bioluminescence.   
\B{Txon\-krr lu syu\-ra\-tan na'\-rìng\-ä Ey\-we\-veng\-ä lor nì\-txan.} At night, the bioluminescence of the Pandoran forest is very beautiful.~\LINK{http://naviteri.org/2012/03/spring-vocabulary-part-1/}{b}
\B{Na'\-rìng sngä\-'i ni\-vrr txon\-krr syu\-ra\-tan\-fa.} The forest begins to glow at night with bioluminescence.~\LINK{http://naviteri.org/2012/07/fivospxiya-aylifyavi-amip-this-months-new-expressions-pt-2/}{b}
(\ding{43}~\HL{syura}{\B{syu\-ra}}, \HL{atan}{\B{a\-tan}})

% syusmung
\hi
\HT{syusmung}{\B{\cp{syu}smung}} \I{["{}sju.smuN]} \emph{or} \I{[.smUN]} (\textsc{rn}: \B{\cp{syu}smùng} \I{["Su.smUN]}) \emph{n.} tray. 
(\ding{43}~\HL{syuve}{\B{syu\-ve}}, \HL{sAmunge}{\B{sä\-mu\-nge}})

% syuve
\hi
\HT{syuve}{\B{\cp{syu}ve}} \I{["{}sju.vE]} \emph{uncountable n.} food. 
\B{Lu oe\-ru syu\-ve.} I have food.~\LINK{http://naviteri.org/2012/10/audio-and-video-learning-materials-for-navi-101/}{b}
\B{Rä\-'ä 'i\-val syu\-vet!} Don't waste food.~\LINK{http://naviteri.org/2014/07/tsawlultxamaw-a-ayliu-post-meetup-words/}{b}
\B{Zun ngal tsa\-fne\-syu\-vet tim\-vìng oer, zel li\-vu oe txur fì\-trr.} If you had given me that kind of food, I would be strong today.~\LINK{http://naviteri.org/2013/04/zun-zel-counterfactual-conditionals/}{b}
\B{Syu\-ve\-ri ay\-ftxo\-zä\-yä ay\-nga\-ru fki\-van on\-lor ftxì\-lor\-sì nìwotx!} May the food of the celebrations smell and taste good to you.~\LINK{http://naviteri.org/2012/11/renu-ayinanfyaya-the-senses-paradigm/}{b} \\
\textsc{Derivations}: \ding{43} 
\on{1}.~\HL{syuvekel}{\B{syu\-ve\-kel}}, famine. 
\on{2}.~\HL{syusmung}{\B{syu\-smung}}, tray. 
\on{3}.~\HL{tsyosyu}{\B{tsyo\-syu}}, food made from flour.
\on{4}.~\HL{okupsyu}{\B{o\-kup\-syu}}, dairy product.

% syuvekel
\hi
\HT{syuvekel}{\B{\cp{syu}vekel}} \I{["{}sju.vE.kEl]} \emph{n.} famine.
(\ding{43}~\HL{syuve}{\B{syu\-ve}}, \HL{tIkelu}{\B{tì\-ke\-lu}}) \\

\end{multicols} 


 
% ===== Section T ===== 

 \pdfbookmark[0]{T}{T}
\begin{center}
\section*{T \I{[t]} (\emph{TeT})}
\end{center}

\begin{multicols}{2}

% -t
\hi
\HT{-t}{\B{-t}} \I{[t\textcorner]} \emph{suf. to mark the patientive case of a n. ending in a vowel or diphthong}. 
(\ding{43}~\HL{-it}{\B{-it}}, \HL{-ti}{\B{-ti}})

% ta
\hi
\HT{ta}{\B{ta}} \I{[ta]} \emph{adp.} from.   
(\emph{a. origin}) from, of. 
\B{Oe\-ri ta pe\-yä fa\-hew a\-ke\-wong on\-tu te\-ya lä\-ngu.} My nose is full of his alien smell.~\LINK{http://languagelog.ldc.upenn.edu/nll/?p=1977}{ll\slash a}  
\B{Ta 'ey\-lan kar\-yu\-sì ay\-nge\-yä}. [This letter comes] From your friend and teacher.~\LINK{http://forum.learnnavi.org/news-announcements/a-response-from-paul-frommer!/}{ln}
\B{Fu\-ri\-a fì\-lì'\-fya\-vit fkol ta 'Ìng\-lì\-sì mo\-lu\-nge, sìl\-pey oe tsnì ay\-ha\-pxì\-tu lì'\-fya\-o\-lo'\-ä aw\-nge\-yä ke stì\-ye\-vi.} That this expression is brought from English, will hopefully not anger the members of the language group.~\LINK{http://forum.learnnavi.org/language-updates/txelanit-hivawl/}{ln}
\B{Kal\-txì ta Ko\-pen\-han.} Hi from Copenhagen.~\LINK{http://naviteri.org/2010/09/kaltxi-ta-kopenhan-hi-from-copenhagen/}{b}
\B{Pot 'a\-ku fì\-tseng\-ta!} Get him out of here!~\LINK{http://naviteri.org/2011/08/new-vocabulary-clothing/}{b}
\B{Säm\-yam\-ìl po\-ru wa\-yìn\-txu fu\-ta nga\-ta lo\-lu li txo\-a.} A hug will show him that you've already forgiven him.~\LINK{http://naviteri.org/2011/10/more-vocabulary-a-bit-of-grammar/}{b}
\B{Fkxi\-le pewn\-ta tu\-té\-yä kur.} The bib necklace hangs from the woman's neck.~\LINK{http://naviteri.org/2013/04/wheres-the-bathroom-and-other-useful-things/}{b}
\B{Za\-'u tsa\-tson ta Ey\-wa.} That obligation comes from Eywa.~\LINK{http://naviteri.org/2014/10/tson-si-fpomron-obligation-and-mental-health/}{b}
\B{Lu Ney\-ti\-rir ta Mo\-'at a tson a kar Tsyeyk\-ur ay\-fya\-'ot Na'\-vi\-yä.} Neytiri is under obligation by\slash from Mo'at to teach Jake the ways of the Na'vi.~\LINK{http://naviteri.org/2014/10/tson-si-fpomron-obligation-and-mental-health/}{b}
\B{tuk\-ru a txo\-lu\-la fkol ta rìn\slash tuk\-ru a ta rìn}, a spear made of wood.~\LINK{http://naviteri.org/2017/09/zisikrr-amip-ayliu-amip-new-words-for-the-new-season/}{b}
(\emph{b. in counting}) from. 
\B{Ti\-am ta pxey vay vol.} Count from three to eight.~\LINK{https://forum.learnnavi.org/language-updates/'count-to'-wa-and-sequential-genitive-(discussion)/}{ln}
(\emph{c. temporal}) from, since. 
\B{Trr\-'ong\-ta Txon\-'ong\-vay po to\-lì\-ran.} He walked from dawn until dusk.~\LINK{http://naviteri.org/2011/01/new-year-new-vocabulary/}{b}
(\emph{d. with} \HL{pawm}{\B{pawm}} \emph{to form the indirect object})
\B{Pol po\-lawm tì\-pawm\-it oe\-ta.} He asked me a question.~\LINK{http://naviteri.org/2020/05/tipusawm-tiuseyng-si-okvur-a-eltur-titxen-si-asking-answering-and-an-interesting-story/}{b}
(\emph{e. conv. phrase for responding to thanks}) \B{ha\-ya\-lo ta oe} \emph{or} \B{hay\-a\-lo oe\-ta}~\I{[ha."ja.lo "wE.ta]}. You're welcome [lit.: next time from me].~\LINK{http://naviteri.org/2010/07/diminutives-conversational-expressions/}{b}
(\ding{43}~\HL{ftu}{\B{ftu}}) \\
\textsc{Derivatives}: \ding{43}
\on{1}.~\HL{ta'em}{\B{ta\-'em}}, from above. 
\on{2}.~\HL{tafkip}{\B{ta\-fkip}}, from up among. 
\on{3}.~\HL{tafral}{\B{ta\-fral}}, therefore, because of that. 
\on{4}.~\HL{takip}{\B{ta\-kip}},  from among. 
\on{5}.~\HL{takrra}{\B{ta\-krra\slash akrr\-ta}}, since. 
\on{6}.~\HL{talun1}{\B{ta\-lun(a)\slash alun\-ta}}, because, because of. 
\on{7}.~\HL{tatlam}{\B{tat\-lam}}, apparently. 
\on{8}.~\HL{taweyk}{\B{ta\-weyk(a)}}, because.
\on{9}.~\HL{tamIfa}{\B{ta\-mì\-fa}}, internal. 
\on{10}.~\HL{tawrrpa}{\B{ta\-wrr\-pa}}, external.

% ta'em
\hi
\HT{ta'em}{\B{ta\cp{'em}}} \I{[ta."PEm]} \emph{adv.} from above. 
\B{Ik\-ran\-ti mak\-to, 'eko ta\-'em.} Take the Ikran, attack from above.~\LINK{http://wiki.learnnavi.org/index.php?title=Na\%27vi_from_Avatar_Movie\#Attack_on_hometree}{wiki\slash a}

% ta'leng
\hi
\HT{ta'leng}{\B{\cp{ta'}leng}} \I{["taP.lEN]} \emph{n.} skin. 
\B{Ta'\-leng prr\-nen\-ä lu faoi, pum ko\-ak\-tu\-ä ekx\-txu.} A baby's skin is smooth, an old person's is rough.~\LINK{http://naviteri.org/2012/01/mipa-zisit-ayliu-amip-new-words-for-the-new-year/}{b}
\B{Fì\-syu\-lang lu ean na ta'\-leng.} This flower is skin\hyp{}blue.~\LINK{http://naviteri.org/2010/08/mipa-ayopin-mipa-ayliu-new-colors-new-words/}{b}
\B{Teng\-krr hu pa\-lu\-kan\-tsyìp uvan se\-ri zo\-lìm oel mì sa'\-leng a 'u\-ot a lu txa' sì ekx\-txu.} While playing with my cat I felt something hard and rough on his skin.~\LINK{http://naviteri.org/2012/11/renu-ayinanfyaya-the-senses-paradigm/}{b} \\
\textsc{Derivatives}: \ding{43}
\on{1}.~\HL{ta'lengean}{\B{ta'\-le\-nge\-an}}, skin\hyp{}blue.

% ta'lengean
\hi
\HT{ta'lengean}{\B{\cp{ta'}lengean}} \I{["taP.lEN.E.an]} \emph{adj.} skin\hyp{}blue.
(\ding{43}~\HL{ta'leng}{\B{ta'\-leng}}, \HL{ean}{\B{ean}})

% ta'unui
\hi
\B{Ta'unui} \I{[ta.Pu.nu.i]} \emph{n. a clan name}, one clan of the reef people.

% tafkip
\hi
\HT{tafkip}{\B{ta\cp{fkip}}} \I{[ta."fkip\textcorner]} \emph{adp.} from up among. 
(\ding{43}~\HL{ta}{\B{ta}}, \HL{fkip}{\B{fkip}})

% tafral
\hi
\HT{tafral}{\B{ta\cp{fral}}} \I{[ta."fRal]} \emph{adv.} therefore, because of that.
\B{Ta\-fral ke lì\-syek oel nge\-yä ke\-ye\-'ung\-it.} Therefore I will not heed your insanity.~\LINK{http://forum.learnnavi.org/language-updates/ltasygt-with-ke-and-nga/}{ln\slash g}
\B{Oe pll\-ngay san mo\-lak\-to oe nì\-fe', ta\-fral sno\-laytx; wä\-tu lu oe\-to txur.} I acknowledge that I rode badly, so I lost; my opponent was stronger than I was.~\LINK{http://naviteri.org/2011/07/txantsana-ultxa-mi-siatll-great-meeting-in-seattle/}{b} 
(\ding{43}~\HL{ta}{\B{ta}}, \HL{ral}{\B{ral}})

% takip
\hi
\HT{takip}{\B{ta\cp{kip}}} \I{[ta."kip\textcorner]} \emph{adp.} from among. 
(\ding{43}~\HL{ta}{\B{ta}}, \HL{kip}{\B{kip}})

% takrra
\hi
\HT{takrra}{\B{ta\cp{krr}a}} \emph{or} \B{a\cp{krr}ta} \I{[ta."k\s{r}.a]} \emph{conj.} since.   
\B{Ay\-lì'\-fya yaw\-ne le\-ru oer ta\-krr\-a 'e\-veng la\-mu.} I've loved languages since I was a child.~\LINK{http://naviteri.org/2011/01/new-year-new-vocabulary/}{b} 
\B{Ay\-ha\-pxì\-tu po\-ngu\-ä txo\-pu si nì\-nän ta\-krr\-a Va'\-rul pxe\-ku\-tut lä\-txayn.} The members of the group are less afraid since Va'ru defeated three of the enemy.~\LINK{http://naviteri.org/2012/02/trr-asawnung-lefpom-happy-leap-day/}{b}
(\ding{43}~\HL{ta}{\B{ta}}, \HL{krr}{\B{krr}}, \HL{a}{\B{a}})

% taksyokx
\hi
\HT{taksyokx}{\B{tak\cp{syokx}}} \I{[ta$\cdot$k\textcorner."{}sjok']} \emph{vin., inf.}\on{11} clap hands.  
\B{Tak\-syokx, ma fra\-po!} Let's give (him\slash her\slash them) a hand, everybody!~\LINK{http://naviteri.org/2017/06/ayliu-a-ta-eywaeveng-kifkey-uniltirantokxa-words-from-disneys-pandora-the-world-of-avatar/}{b}
(\ding{43}~\HL{takuk}{\B{ta\-kuk}}, \HL{tsyokx}{\B{tsyokx}})

% takuk 1
\hi
\HT{takuk}{\B{\cp{ta}kuk}}\textsuperscript{\on{1}} \I{["ta.kuk\textcorner]} \emph{or} \I{[.kUk\textcorner]} (\textsc{rn}: \B{\cp{ta}kùk} \I{["ta.kUk\textcorner]}) \emph{vtr.} (\emph{a.}) strike, hit one's target, slap, beat.   
\B{Oe\-yä swi\-zaw nì\-ngay ti\-va\-kuk.} Let my arrow strike true.~\LINK{http://wiki.learnnavi.org/index.php/Hunting_Song}{wiki}
\B{Ta\-ron\-yul le\-hì\-pey kan smar\-it nìl\-ke\-ftang slä ke ta\-kuk kaw\-krr.} A hesitant hunter will aim at a prey forever but never hit it.~\LINK{http://naviteri.org/2015/11/vomuna-liu-amip-ten-new-words/}{b}
\B{Oe\-ti ta\-kuk rä\-'ä!} Don't hit me!~\LINK{https://forum.learnnavi.org/language-updates/little-confirmation-on-takuk/}{ln}
(\emph{b. with} \HL{fa}{\B{fa}}) \HL{rumtsyokx}{\B{ta\-kuk fa rum\-tsyokx}}, punch s.o. \B{A\-te\-yol En\-tut to\-la\-kuk fa rum\-tsyokx.} Ateyo punched Entu.~\LINK{https://naviteri.org/2024/06/ayhapxi-tokxa-si-ayliu-alahe-parts-of-the-body-and-more/}{b}; \HL{rIktsyokx}{\B{ta\-kuk fa rìk\-tsyokx}}, slap s.o. \B{A\-te\-yol En\-tut to\-la\-kuk fa rìk\-tsyokx.} Ateyo slapped Entu.~\LINK{https://naviteri.org/2024/06/ayhapxi-tokxa-si-ayliu-alahe-parts-of-the-body-and-more/}{b}
(\ding{214}~\HL{hupx}{\B{hupx}})
(\emph{c. of weather}) \B{Ta\-kuk set tsaw\-ke. Sley\-ku fra\-'ut u\-kxo.} Now the sun bears down. (It) makes everything dry.~\LINK{http://forum.learnnavi.org/language-updates/learnings-from-the-siatlla-song-translations/msg475008/\#msg475008}{ln} \\
\textsc{Derivatives}: \ding{43}
\on{1}.~\HL{tsamkuk}{\B{tsam\-kuk}}, war drum.

% Takuk 2
\hi
\B{\cp{Ta}kuk}\textsuperscript{\on{2}} \I{["ta.kuk\textcorner]} \emph{or} \I{[.kUk\textcorner]} \emph{n.} \emph{male name}.~\LINK{}{dh}

% talioang
\hi
\HT{talioang}{\B{\cp{ta}lioang}} \I{["tal.i.""{}o.aN]} \emph{n., fauna} sturmbeest, \emph{bovindicum monocerii}. 
\B{Fo ko\-lä tsa\-tseng fte ti\-va\-ron ye\-ri … ìì … ko\-lan ta\-li\-o\-ang\-it.} They went there to hunt hexa … um … I mean sturmbeast.~\LINK{http://naviteri.org/2014/05/mipa-ayliu-mipa-aysafpil-new-words-new-ideas/}{b} 

% talun / taluna
\hi
\HT{talun1}{\B{ta\cp{lun}}}\textsuperscript{\on{1}} \emph{or} \B{ta\cp{lu}na} \I{[ta."lun]} \emph{or} \I{[ta."lu.na]} \emph{conj.} because, since.
\B{Ay\-oe sno\-laytx fì\-trr ta\-lun\-a oel rum\-it to\-lung\-zup.} We lost today because I dropped the ball.~\LINK{http://naviteri.org/2011/05/some-miscellaneous-vocabulary/}{b}
\B{Tsa\-ik\-ran\-ìri ta\-lun\-a new nga\-ti tspi\-vang, law lu fwa mi ke lu zä\-fi.} Since that ikran wants to kill you, it's clear it's still not tame.~\LINK{http://naviteri.org/2013/01/awvea-posti-zisita-amip-first-post-of-the-new-year/}{b}
(\ding{43}~\HL{ta}{\B{ta}}, \HL{lun}{\B{lun}}, \HL{a}{\B{a}})

\hi
\HT{talun2}{\B{ta\cp{lun}}}\textsuperscript{\on{2}} \I{[ta."lun]} \emph{or} \I{[."lUn]} \emph{adp.} because of, due to.
\B{Tì\-'i'\-a\-ri tsam\-ä ze\-ne Saw\-tu\-te vi\-ve\-lek ta\-lun tì\-txur Ey\-wa\-yä.} At the end of the war, the Sky People had to give up due to the power of Eywa.~\LINK{http://naviteri.org/2012/03/spring-vocabulary-part-1/}{b}
\B{Tsaw\-ke slo\-lu vawm ta\-lun fwopx.} The sun became dark due to dust in the air.~\LINK{http://naviteri.org/2014/07/tsawlultxamaw-a-ayliu-post-meetup-words/}{b}
(\emph{a. to form a causative meaning with vin. modals}) \B{Tom\-pa\-kel\-ta\-lun ze\-ne tu\-te Kä\-lì\-for\-ni\-a\-yä pay\-it si\-var nì\-nän.} The drought causes Californians to use less water.~\LINK{http://naviteri.org/2015/04/some-new-words-for-may-day/}{b}
(\ding{43}~\HL{ta}{\B{ta}}, \HL{lun}{\B{lun}})

% tam
\hi
\HT{tam}{\B{tam}} \I{[tam]} \emph{vin.} suffice, be enough, ``do''. 
\B{Fì\-pam\-tseo\-tu\-ri ke la\-yu ftue fwa run fkol sä\-rawn\-it a tam.} It won't be easy to find a satisfactory replacement for this musician.~\LINK{http://naviteri.org/2012/01/mipa-zisit-ayliu-amip-new-words-for-the-new-year/}{b}
\B{Nì\-ke\-ftxo oe\-ri ke tam.} Unfortunately, it doesn't work for me.~\LINK{http://naviteri.org/2011/09/\%E2\%80\%9Cby-the-way-what-are-you-reading\%E2\%80\%9D/comment-page-1/\#comment-1116}{b}
\B{Pe\-yä 'itan\-ìri lu ho\-na nì\-txan a fì\-'u law lu fra\-por. Slä tì\-ho\-na nì\-'aw ke tam.} It's clear to everyone that his son is very cute. But cuteness alone isn't enough.~\LINK{http://naviteri.org/2012/01/more-additions-to-the-lexicon/}{b}
(\emph{a. idiom, conv.}) \B{Tam.} \textsc{ok}! 
\B{Tsun ti\-vam.} It's ok.~\LINK{http://forum.learnnavi.org/language-updates/ok!-we-got-ok/}{ln} 
\B{'U a\-ke\-tsran tsun ti\-vam.} Anything at all will be fine.~\LINK{http://naviteri.org/2013/03/whoever-whatever-whenever/}{b}
(\emph{b. idiom, conv.}) \B{Tam ke tam}. `Meh\slash so\hyp{}so.'~\LINK{http://naviteri.org/2012/10/mipa-vospxi-mipa-ayliu-new-words-for-the-new-month/comment-page-1/\#comment-1621}{b}
\B{Nga\-ru lu fpom srak?--Tam ke tam.} How are you?--Meh, so\hyp{}so.
(\emph{c. proverb}) \B{Fwa kan ke tam; ze\-ne swi\-zaw\-it li\-vo\-nu.} Intent is not enough; it's action that counts. [lit.: To aim is not enough; one must release the arrow.]~\LINK{http://naviteri.org/2013/02/vospxi-ayol-posti-apup-short-post-for-a-short-month/}{b}\\
\textsc{Derivatives}: \ding{43}
\on{1}.~\HL{letam}{\B{le\-tam}}, sufficient, enough. 
\on{2}.~\HL{nItam}{\B{nì\-tam}}, enough.

% tamìfa
\hi 
\HT{tamIfa}{\B{ta\cp{mì}fa}} \I{[ta."mI.fa]} \emph{adj.} internal. 
(\ding{214}~\HL{tawrrpa}{\B{ta\-wrr\-pa}}, \ding{43}~\HL{ta}{\B{ta}}, \HL{mIfa}{\B{mì\-fa}}) 

% tanhì
\hi
\HT{tanhI}{\B{tan\cp{hì}}} \I{[tan."hI]} \emph{n.} (\emph{a. astr.}) star.  
\B{Tsay\-san\-hì ha\-yum ye'\-rìn (taw\-ftu).} Those stars will soon set.~\LINK{http://naviteri.org/2012/01/more-additions-to-the-lexicon/}{b}
\B{Kxam\-trr lam fwa san\-hì a mì saw 'o\-lìp nì\-wotx slä tsa\-krr ke tsun fko sat tsi\-ve\-'a nì\-'aw.} At mid-day it seems that the stars in the sky have all vanished but they just can't be seen then.~\LINK{http://naviteri.org/2012/03/spring-vocabulary-part-1/}{b}
\B{Hol\-pxay san\-hì\-yä a mì saw lu ke\-tsuk\-ti\-am; keng ke tsun fko tsi\-ve\-'a sat nì\-wotx.} The number of stars in the sky is infinite; it's not even possible to see them all.~\LINK{http://naviteri.org/2014/11/twenty-before-the-holidays/}{b}
\B{Pll\-txe Saw\-tu\-te san ki\-fkey\-ìl ay\-oe\-yä tok txan\-a ngip\-it a san\-hì\-kip.} The Sky People say that their world is in the great space among the stars.~\LINK{http://naviteri.org/2015/03/kap-si-ayunil-saylahe-threats-dreams-and-other-things/}{b}
(\emph{b. of the body}) bioluminescent freckle. 
(\ding{43}~\HL{atan}{\B{a\-tan}}, \HL{hI'i}{\B{hì\-'i}})

% tanleng
\hi 
\HT{tanleng}{\B{\cp{tan}leng}} \I{["tan.lEN]} \emph{n.}, \ding{96} bark (of a tree). 
(\ding{43}~\HL{tangek}{\B{ta\-ngek}}, \HL{ta'leng}{\B{ta'\-leng}})

% tangek
\hi
\HT{tangek}{\B{\cp{ta}ngek}} \I{["taN.Ek\textcorner]} \emph{n.} trunk (\emph{of a tree}). 
\B{Tsa\-ut\-ral\-ìri ta\-ngek lu sloa nì\-txan; 'evi ke tsun tsyi\-vìl.} The trunk of that tree is very wide; the kid cannot climb it.~\LINK{http://naviteri.org/2012/06/spring-vocabulary-part-3/}{b}
\B{Fì\-ut\-ral\-ìri ta\-ngek\-ä zir fkan vawt, hu\-fwa ke rey.} The trunk of the tree feels solid, although it's dead.~\LINK{http://naviteri.org/2013/09/on-si-salewfya-shapes-and-directions/}{b}
\B{Tsa\-ta\-ngek\-ìri pam fkan mo\-mek.} That (tree)trunk sounds hollow.~\LINK{http://naviteri.org/2013/09/on-si-salewfya-shapes-and-directions/}{b}
\B{Na ay\-sa\-ngek a\-fkew.} Like mighty trunks (of trees).~\LINK{http://www.pandorapedia.com/navi/music/ritual_music}{pp} \\
\textsc{Derivatives}: \ding{43}
\on{1}.~\HL{tanleng}{\B{tan\-leng}}, bark.

% tare
\hi
\HT{tare}{\B{\cp{ta}re}} \I{["ta.RE]} \emph{vtr.} connect, relate to, have a relationship with.   
\B{Sä\-pll\-txel kar\-yu\-ä ke to\-la\-rä\-nge tì\-pawm\-it kaw\-'it.} The teacher's statement in no way pertained to the question. Drat!~\LINK{http://naviteri.org/2012/06/spring-vocabulary-part-3/}{b}
\B{Txil\-te Ri\-ni\-sì tä\-pa\-re fì\-tsap nì\-soaia, slä tsal\-su\-ngay ke nì\-olo' ta\-krr\-a Ri\-ni mun\-txa slo\-lu.} Txilte and Rini are related by blood, but nevertheless not by clan since the time Rini got married.~\LINK{http://naviteri.org/2012/06/spring-vocabulary-part-3/}{b}
(\emph{a. conv.}) \B{Ke ta\-re.} It's irrelevant.\slash It doesn't matter.~\LINK{http://naviteri.org/2012/06/spring-vocabulary-part-3/}{b} \\
\textsc{Derivatives}: \ding{43}
\on{1}.~\HL{sAtare}{\B{sä\-ta\-re}}, connection, relationship. 
\on{2}.~\HL{ketartu}{\B{ke\-tar\-tu}}, outcast.
\on{3}.~\HL{tarsyu}{\B{tar\-syu}}, tarsyu.

% tarep
\hi 
\HT{tarep}{\B{\cp{ta}rep}} \I{["ta.REp\textcorner]} \emph{vtr.} save, rescue (\emph{from a dangerous or distressing situation}). 
\B{Maw\-krr\-a pa\-lu\-lu\-kan po\-sìn spo\-lä, ke tsä\-ngun fko pot ti\-va\-rep.} Sadly, once the thanator had jumped on her, she could not be rescued.~\LINK{http://naviteri.org/2021/12/zolau-niprrte-ma-3746-welcome-2022/}{b}
(\emph{a. idiom}) \B{Ta\-rep!} ``Help!''~\LINK{http://naviteri.org/2021/12/zolau-niprrte-ma-3746-welcome-2022/}{b} 
(\ding{43}~\HL{'avun}{\B{'a\-vun}}; \HL{zong}{\B{zong}}) \\
\textsc{Derivatives}: \ding{43} 
\on{1}.~\HL{tareptu}{\B{ta\-rep\-tu}}, rescuee, rescued person. 
\on{2}.~\HL{sAtarep}{\B{sä\-ta\-rep}}, a rescue. 
\on{3}.~\HL{tarepyu}{\B{ta\-rep\-yu}}, rescuer.

% tareptu
\hi 
\HT{tareptu}{\B{\cp{ta}reptu}} \I{["ta.REp\textcorner.tu]} \emph{n.} rescuee, s.o. who has been rescued or saved, saved person, rescued person. 
\B{Maw tsa\-fe\-kem, new ta\-rep\-tu sne\-yä ta\-rep\-yur ira\-yo si\-vi.} After the accident, the rescuee wanted to thank his rescuer.~\LINK{http://naviteri.org/2021/12/zolau-niprrte-ma-3746-welcome-2022/}{b} 
(\ding{43}~\HL{tarep}{\B{ta\-rep}})

% tarepyu
\hi 
\HT{tarepyu}{\B{\cp{ta}repyu}} \I{["ta.REp\textcorner.ju]} \emph{n.} rescuer. 
\B{Maw tsa\-fe\-kem, new ta\-rep\-tu sne\-yä ta\-rep\-yur ira\-yo si\-vi.} After the accident, the rescuee wanted to thank his rescuer.~\LINK{http://naviteri.org/2021/12/zolau-niprrte-ma-3746-welcome-2022/}{b} 
(\ding{43}~\HL{tarep}{\B{ta\-rep}})

% tarnioang
\hi
\HT{tarnioang}{\B{\cp{tar}nioang}} \I{["taR.ni.o.aN]} \emph{n.} predator animal.
\B{Fì\-fne\-ya\-yol tsrul\-it txu\-la mì song\-ropx ut\-ral\-ä fte ay\-li\-nit hi\-vaw\-nu wä sar\-ni\-o\-ang. Fì\-tì\-kan\-ìri ropx ke ha'.} This kind of bird builds its nest in a tree cavity so as to protect its young from predators. For this purpose, a hole going right through the tree trunk isn't suitable.~\LINK{http://naviteri.org/2013/04/wheres-the-bathroom-and-other-useful-things/}{b}
(\ding{43}~\HL{taron}{\B{ta\-ron}}, \HL{ioang}{\B{ioang}})

% taron
\hi
\HT{taron}{\B{\cp{ta}ron}} \I{["ta.Ron]} \emph{vtr.} hunt. 
\B{Ney\-ti\-ril ye\-rik\-it to\-la\-ron.} Neytiri hunted a hexapede.~\LINK{http://wiki.learnnavi.org/index.php?title=Corpus\#MSNBC}{wiki}
\B{Pol ye\-rik\-it ta\-ron nìl\-tsan fra\-to.} He hunts yerik the best of all.~\LINK{http://forum.learnnavi.org/language-updates/confirmations-\%28comparisons-gerund\%29/}{ln} 
\B{Ftu\-e lu fwa ta\-ron ngong\-a io\-ang\-it to fwa ta\-ron pum\-it a lu wa\-lak sì win.} It's easier to hunt lethargic animals than to hunt perky, speedy ones.~\LINK{http://naviteri.org/2011/10/more-vocabulary-a-bit-of-grammar/}{b}
\B{Pe\-lun fì\-tseng\-ne nga zo\-la\-'u fte ti\-va\-ron? -- Lo\-lu 'en.} Why did you come here to hunt? -- It was a guess (hunch).~\LINK{http://naviteri.org/2011/01/new-year-new-vocabulary/}{b}
\B{Pa\-lu\-lu\-kan a te\-ra\-ron lu le\-hrr\-ap.} A thanator that's hunting is dangerous.~\LINK{http://forum.learnnavi.org/language-updates/participles/}{ln}
\B{Pa\-lu\-lu\-kan a\-tu\-sa\-ron lu le\-hrr\-ap.} A hunting thanator is dangerous.~\LINK{http://forum.learnnavi.org/language-updates/participles/}{ln}
\B{Pol trr\-am ye\-rik\-it tar\-ma\-ron. -- Fì\-trr mi te\-ra\-ron.} He was hunting hexapede yesterday. -- He's still hunting today.~\LINK{http://forum.learnnavi.org/language-updates/as-x-as-y-keep-on-keepin-on/}{ln} \\
\textsc{Derivatives}: \ding{43}
\on{1}.~\HL{taronyu}{\B{ta\-ron\-yu}}, hunter. 
\on{2}.~\HL{tItaron}{\B{tì\-ta\-ron}}, hunt. 
\on{3}.~\HL{sAtaron}{\B{sä\-ta\-ron}}, hunt. 
\on{4}.~\HL{tarnioang}{\B{tar\-ni\-oang}}, predator animal. 
\on{5}.~\HL{uniltaron}{\B{unil\-ta\-ron}}, dream hunt. 
\on{6}.~\HL{taronway}{\B{ta\-ron\-way}}, the Hunt Song.
\on{7}.~\HL{tarpongu}{\B{tar\-po\-ngu}}, hunting party.

% taronway
\hi
\HT{taronway}{\B{\cp{ta}ronway}} \I{["ta.Ron.waj]} \emph{n.} the Hunt Song.~\LINK{http://naviteri.org/2013/08/taronway-the-hunt-song/}{b}
(\ding{43}~\HL{taron}{\B{ta\-ron}}, \HL{way}{\B{way}})

% taronyu
\hi
\HT{taronyu}{\B{\cp{ta}ronyu}} \I{["ta.Ron.ju]} \emph{n.} hunter. 
\B{Oe lu ta\-ron\-yu.} I am a hunter.~\LINK{http://naviteri.org/2010/09/getting-to-know-you-part-2/}{b}
\B{Tsa\-ri\-a fkol po\-le\-'un fu\-ta Lo\-ak slu ta\-ron\-yu, sem\-pul 'efu ye.} Father is content that it's been decided Loak will be a hunter.~\LINK{http://naviteri.org/2011/04/\%E2\%80\%99a\%E2\%80\%99awa-li\%E2\%80\%99fyavi-amip\%E2\%80\%94a-few-new-expressions/}{b}
\B{Fì\-ta\-ron\-yu\-tsyìp ke tsun ke\-'ut sti\-vä'\-nì.} This little hunter can't catch anything.~\LINK{http://naviteri.org/2010/07/diminutives-conversational-expressions/}{b}
\B{Ta\-ron\-yul le\-hawm\-pam ska\-'a sä\-ta\-ron\-it.} A noisy hunter destroys the hunt.~\LINK{http://naviteri.org/2014/11/twenty-before-the-holidays/}{b}
\B{Lu oe\-ru sngum a sa\-ron\-yu ke tì\-ye\-vä\-txaw.} I'm worried that the hunters will not return (soon, as expected).~\LINK{http://naviteri.org/2011/01/new-year-new-vocabulary/}{b}
(\emph{a. proverb}) \B{Ta\-ron\-yut yom smar\-ìl.} \emph{A totally botched situation, chaos} [lit.: the prey eats the hunter].
\B{Eyk Ka\-mun a fra\-lo lä\-ngu tsa\-sä\-ta\-ron vel\-ke nì\-wotx. Ta\-ron\-yut yom smar\-ìl!} Every time Kamun is in charge, the hunt is a mess. Everything goes wrong that can!~\LINK{http://naviteri.org/2012/10/mipa-vospxi-mipa-ayliu-new-words-for-the-new-month/}{b}
\B{Hì\-pey ta\-ron\-yu, hi\-fwo ye\-rik.} ``He who hesitates is lost.'' [lit.: The hunter hesitates and the hexapede escapes]~\LINK{http://naviteri.org/2015/11/vomuna-liu-amip-ten-new-words/}{b}
(\ding{43}~\HL{taron}{\B{ta\-ron}}, \HL{-yu}{\B{-yu}})

% tarpongu
\hi
\HT{tarpongu}{\B{\cp{tar}pongu}} \I{["taR.po.Nu]} \emph{n.} hunting party. 
\B{Ni\-nat tì\-mung\-wrr sìl\-mi fte Ra\-lu tsi\-vun ki\-vä hu tar\-po\-ngu.} Ninat just made an exception so that Ralu can go with the hunting party.~\LINK{http://naviteri.org/2015/08/aylifyavi-lereyfya-2-cultural-terms-2-and-more/}{b}
(\ding{43}~\HL{taron}{\B{ta\-ron}}, \HL{pongu}{\B{po\-ngu}})

% tarsyu
\hi 
\HT{tarsyu}{\B{\cp{tar}syu}} \I{["taR.sju]} \emph{n.} \ding{96} tarsyu. 
(\ding{43}~\HL{tare}{\B{ta\-re}}, \HL{syulang}{\B{syu\-lang}})

% tatep
\hi
\HT{tatep}{\B{\cp{ta}tep}} \I{["ta.tEp\textcorner]} \emph{vtr.} lose, lose track\slash  awareness of s.t. 
\B{Tì\-kan taw\-na\-tep!} (\emph{mil.}) Target lost!~\LINK{http://naviteri.org/2011/05/some-miscellaneous-vocabulary/}{b\slash g}

% tatlam
\hi
\HT{tatlam}{\B{tat\cp{lam}}} \I{[tat\textcorner."lam]} \emph{adv.} apparently. 
(\ding{43}~\HL{ta}{\B{ta}}, \HL{tIlam}{\B{tì\-lam}})

% tautral
\hi
\HT{tautral}{\B{\cp{ta}utral}} \I{["ta.ut\textcorner.Ral]} \emph{n.}, \ding{96} sky tree.
(\ding{43}~\HL{taw}{\B{taw}}, \HL{utral}{\B{ut\-ral}})

% tayng
\hi 
\HT{tayng}{\B{tayng}} \I{[tajN]} \emph{n.}, \ding{96} thistle\hyp{}like plant (\emph{generic term}).  \\
\textsc{Derivatives}: \ding{43} 
\on{1}.~\HL{tompatayng}{\B{tom\-pa\-tayng}}, rain thistle. 

% Tayrangi
\hi
\HT{tayrangi}{\B{Tayrangi}} \I{[taj.Ra.Ni]} \emph{n., clan name, specialise in banshee fishing}.

% taweyk
\hi
\HT{taweyk}{\B{ta\cp{weyk}}} \emph{or} \B{ta\cp{wey}ka} \I{[ta."wEjk\textcorner]}  \I{[ta."wEj.ka]} \emph{conj.} because.
(\ding{43}~\HL{ta}{\B{ta}}, \HL{oeyk}{\B{o\-eyk}}, \HL{talun1}{\B{ta\-lun(a)}}) 

% tawrrpa
\hi 
\HT{tawrrpa}{\B{ta\cp{wrr}pa}} \I{[ta."w\s{r}.pa]} \emph{adj.} external. 
\B{Fì\-txe\-lel fngo' sä\-lang\-it a\-ta\-wrr\-pa.} This matter requires an external investigation.~\LINK{http://naviteri.org/2021/12/zolau-niprrte-ma-3746-welcome-2022/}{b}
(\ding{214}~\HL{tamIfa}{\B{ta\-mì\-fa}}, \ding{43}~\HL{ta}{\B{ta}}, \HL{wrrpa}{\B{wrr\-pa}}) 

% taw
\hi
\HT{taw}{\B{taw}} \I{[taw]} \emph{n.} sky.
\B{Kxam\-trr lam fwa san\-hì a mì saw 'o\-lìp nì\-wotx slä tsa\-krr ke tsun fko sat tsi\-ve\-'a nì\-'aw.} At mid\hyp{}day it seems that the stars in the sky have all vanished but they just can't be seen then.~\LINK{http://naviteri.org/2012/03/spring-vocabulary-part-1/}{b} 
\B{Hol\-pxay san\-hì\-yä a mì saw lu ke\-tsuk\-ti\-am; keng ke tsun fko tsi\-ve\-'a sat nì\-wotx.} The number of stars in the sky is infinite; it's not even possible to see them all.~\LINK{http://naviteri.org/2014/11/twenty-before-the-holidays/}{b} 
(\emph{a. of weather}) \B{Ya\-fkeyk fya\-pe set? -- Taw lu pi\-ak\slash lep\-wopx\slash tstu.} What's the weather like now? -- The sky is clear\slash cloudy\slash overcast.~\LINK{http://naviteri.org/2012/10/audio-and-video-learning-materials-for-navi-101/}{b} 
\B{Taw\-ri fya\-pe?} What's the sky like?~\LINK{http://naviteri.org/2011/05/weather-part-2-and-a-bit-more-2/}{b}
(\emph{b. for celestial bodies rising and setting})
\B{Tsaw\-ke fpxe\-rä\-kìm (nem\-fa taw).} The sun is rising.~\LINK{http://naviteri.org/2012/01/more-additions-to-the-lexicon/}{b}
\B{Tsay\-san\-hì ha\-yum ye'\-rìn (taw\-ftu).} Those stars will soon set.~\LINK{http://naviteri.org/2012/01/more-additions-to-the-lexicon/}{b} \\
\textsc{Derivatives}: \ding{43}
\on{1}.~\HL{tawtute}{\B{taw\-tu\-te}}, human, sky\hyp{}person. 
\on{2}.~\HL{tawsyuratan}{\B{taw\-syu\-ra\-tan}}, aurora. 
\on{3}.~\HL{tawtsngal}{\B{taw\-tsngal}}, panopyra. 
\on{4}.~\HL{tautral}{\B{ta\-ut\-ral}}, sky\hyp{}tree.
\on{5}.~\HL{tawsIp}{\B{taw\-sìp}}, spaceship.
\on{6}.~\HL{tawtxew}{\B{tawtxew}}, horizon, skyline.

% Tawkami
\hi
\HT{Tawkami}{\B{Taw\cp{ka}mi}} \I{[taw."ka.mi]} \emph{n.} \emph{clan name, experts in botany}.

% tawng
\hi
\HT{tawng}{\B{tawng}} \I{[tawN]} \emph{vin.} dive, duck (\emph{into water}).
\B{tawng nem\-fa pay}, dive into water.~\LINK{http://naviteri.org/2021/02/aysipawm-si-aysieyng-questions-and-answers/}{b} 

% tawsìp
\hi
\HT{tawsIp}{\B{\cp{taw}sìp}} \I{["taw.sIp\textcorner]} \emph{n.} spaceship, \emph{compound of Na'vi} \HL{taw}{\B{taw}} ``sky'' \emph{and English} \textsc{lw} ``ship''.
\B{Taw\-sìp nge\-yä lu sngel\-tseng.} Your ship is a garbage scow.~\LINK{http://wiki.learnnavi.org/index.php?title=Corpus\#UGO_Movie_Blog_.28Twitter_Questions.29}{ugo\slash wiki}

% Tawsnrrtsyìp
\hi 
\HT{tawsnrrtsyIp}{\B{Taw\cp{snrr}tsyìp}} \I{[taw."{}sn\s{r}.\t{ts}jIp\textcorner]} \emph{n., astr.} Alpha Centauri C\slash Proxima Centauri; \emph{often shortened to} \ding{43}~\HL{snrrtsyIp}{\B{Snrrtsyìp}} 
(\ding{43}~\HL{taw}{\B{taw}}, \HL{sAnrr}{\B{sä\-nrr}}, \HL{-tsyIp}{\B{-tsyìp}})

% tawsyuratan
\hi
\HT{tawsyuratan}{\B{tawsyu\cp{ra}tan}} \I{[taw.sju."Ra.tan]} \emph{n.} aurora.
(\ding{43}~\HL{taw}{\B{taw}}, \HL{syuratan}{\B{syu\-ra\-tan}}) 

% tawtute
\hi
\HT{tawtute}{\B{\cp{taw}tute}} \I{["taw.tu.tE]} \emph{n.} human, ``skyperson''.
\B{Na'\-vi\-ru set lu nawm\-a eyk\-tan a\-mip a lar\-mu Taw\-tu\-te.} The Na'vi have a great new leader now who had been a Skyperson.~\LINK{http://naviteri.org/2010/07/ting-mikyun-fte-tslivam-listening-comprehension-1-3/}{b}
\B{Lì'\-fya\-ri tsa\-taw\-tu\-te pe\-yì?} How is that Sky Person's Na'vi?.~\LINK{http://naviteri.org/2013/01/lifyari-po-peyi-how-good-is-her-navi/}{b}
\B{Tì\-'e\-fu\-mì oe\-yä su\-te\-kip nì\-wotx, ftxey Na'\-vi ftxey Saw\-tu\-te, lu sìl\-tsan lu kawng.} In my opinion among all people, whether Na'vi or Skypeople, there are good and bad.~\LINK{http://naviteri.org/2010/07/ting-mikyun-fte-tslivam-listening-comprehension-1-3/}{b}
\B{Lum\-pe le\-rang Kel\-ut\-ral\-it Saw\-tu\-tel?} Why are the humans exploring Hometree?~\LINK{http://naviteri.org/2014/11/twenty-before-the-holidays/}{b}
\B{Ay\-oe\-ngal fte Fay\-saw\-tu\-tet ki\-vu\-rakx, ze\-ne nì\-'aw\-ve ay\-fot tsli\-vam.} We must understand these Sky People if we are to drive them out.~\LINK{http://forum.learnnavi.org/language-updates/confirmation-on-kurakx-eytukan-line-and-toitslan/}{ln\slash a\textsuperscript{c}}
\B{Sra\-ne, su\-nu Saw\-tu\-te\-ru tì\-lang, slä ke omum fol teyng\-ta kem\-pe ze\-ne si\-vi maw\-krr\-a 'uo\-ti ro\-lun.} Yes, the Skypeople love exploration, but they don't know what to do once they find something.~\LINK{http://naviteri.org/2014/11/twenty-before-the-holidays/}{b} 
\B{Saw\-tu\-te\-yä hem\-ìri lu aw\-nga\-ru ya\-yayr.} The Skypeople's actions confuse us.~\LINK{http://naviteri.org/2011/01/new-year-new-vocabulary/}{b}
\B{Saw\-tu\-te\-ri tì\-fkey\-tok ke tsan\-'o\-lul kaw\-'it.} The situation with the Skypeople hasn't improved one bit.~\LINK{http://naviteri.org/2013/03/tsanerul-feerul-getting-better-getting-worse/}{b}
(\ding{43}~\HL{taw}{\B{taw}}, \HL{tute1}{\B{tu\-te}}) 

% tawtxew
\hi 
\HT{tawtxew}{\B{\cp{taw}txew}} \I{["taw.t'Ew]} \emph{n.} horizon, skyline (\B{sìn}\slash \B{ro}, on; \B{eo}, in front of; \B{uo}, behind). 
\B{Lu ay\-ram sìn taw\-txew.} There are mountains on the horizon.~\LINK{http://naviteri.org/2018/07/ta-sulfatu-a-ayliu-niul-more-words-from-our-experts/}{b}
\B{Na\-ra\-nawm\-ä mawl mi lu uo taw\-txew.} Half of Polyphemus is still behind the horizon.~\LINK{http://naviteri.org/2018/07/ta-sulfatu-a-ayliu-niul-more-words-from-our-experts/}{b}
(\ding{43}~\HL{taw}{\B{taw}}, \HL{txew}{\B{txew}})

% tawtsngal
\hi
\HT{tawtsngal}{\B{\cp{taw}tsngal}} \I{["taw.\t{ts}Nal]} \emph{n.}, \ding{96} panopyra, ``sky cup'', \emph{panopyra\slash maraca aerii}.
(\ding{43}~\HL{taw}{\B{taw}}, \HL{tsngal}{\B{tsngal}}) 

% täftxu
\hi
\HT{tAftxu}{\B{tä\cp{ftxu}}} \I{[t\ae."ft'u]} \emph{vtr.} weave.   
\B{Ka\-tot tä\-ftxu oel | Nì\-ean nì\-rim.} I weave the rhythm | In yellow and blue.~\LINK{http://wiki.learnnavi.org/index.php/Weaving_Song}{wiki}
\B{Fu\-ri\-a tä\-ftxu ngal tok yì\-pet? -- Tok yìt a\-ke\-sran.} How's your weaving? -- I'm so\hyp{}so.~\LINK{http://naviteri.org/2013/01/lifyari-po-peyi-how-good-is-her-navi/}{b} \\
\textsc{Derivatives}: \ding{43}
\on{1}.~\HL{tAftxuyu}{\B{tä\-ftxu\-yu}}, weaver. 
\on{2}.~\HL{ftxulI'u}{\B{ftxu\-lì\-'u}}, orate, give a speech.
\on{3}.~\HL{sAtAftxu}{\B{sä\-tä\-ftxu}}, piece of weaving.

% täftxuyu
\hi
\HT{tAftxuyu}{\B{tä\cp{ftxu}yu}} \I{[t\ae."ft'u.ju]} \emph{n.} weaver, one who weaves. 
\B{Pam\-tseol ngop ay\-re\-nut | Mì ron\-sem\-ä tì\-fnu | Teng\-fya ngop sä\-ftxu\-yul | Mì hi\-fkey.} Music creates patterns | In the silence of the mind | As weavers do | In the physical word.~\LINK{https://www.pandorapedia.com/navi/music/ritual_music.html}{pp}
(\ding{43}~\HL{tAftxu}{\B{tä\-ftxu}}, \HL{-yu}{\B{-yu}}) 

% tätxaw
\hi
\HT{tAtxaw}{\B{tä\cp{txaw}}} \I{[t\ae."t'aw]} \emph{vin.} return (\B{ne}, to).   
\B{Set sre\-fey oel fu\-ta tsam\-po\-ngu tä\-txaw maw txon\-'ong.} I'm currently expecting the war party back after nightfall.~\LINK{http://naviteri.org/2011/02/new-vocabulary-part-2/}{b}
\B{Krr\-pe nga tä\-txaw ne kel\-ku?} When will you return home?~\LINK{http://naviteri.org/2012/10/audio-and-video-learning-materials-for-navi-101/}{b}
\B{Maw\-krr\-a Saw\-tu\-te\-yä txam\-pxì ho\-lum, tä\-txaw Na'\-vi\-ne nì\-'i'\-a ve\-'o.} After most of the Sky People left, order finally returned to the People.~\LINK{http://naviteri.org/2012/10/mipa-vospxi-mipa-ayliu-new-words-for-the-new-month/}{b}
\B{Maw hì\-krr ay\-oe tì\-yä\-txaw.} We'll return after a short time.~\LINK{http://wiki.learnnavi.org/index.php?title=Corpus\#Good_Morning_America}{wiki}
\B{Sre\-fpìl Oma\-ti\-ka\-ya tsnì Tsyeyk kaw\-krr ke ta\-yä\-txaw maw ka\-vuk sne\-yä.} The Omaticaya assumed that Jake would never return after his treachery.~\LINK{http://naviteri.org/2014/10/tson-si-fpomron-obligation-and-mental-health/}{b}
(\emph{a. idiom}) \B{To\-lä\-txaw nì\-prr\-te'.} ``Welcome back.''~\LINK{http://forum.learnnavi.org/language-updates/some-news-from-berlin/msg585072/\#msg585072}{ln} 

% te
\hi
\HT{te}{\B{te}} \I{[tE]} \emph{part. used in full names}. 
\B{Ney\-ti\-ri te Tska\-ha Mo\-'at\-'ite}, Neytiri of the Tskaha, daughter of Mo'at.~\LINK{http://wiki.learnnavi.org/index.php?title=Na\%27vi_from_Avatar_Movie\#Wutso_in_Hometree}{wiki\slash a\textsuperscript{c}} 
\B{Ay\-nge\-nga\-ru ohe\-yä tsmuk\-it alu Ne\-wey te Tska\-ha So\-rewn\-'ite.} Allow me to introduce my sister, Ne\-wey te Tska\-ha So\-rewn\-'ite.~\LINK{http://naviteri.org/2010/09/getting-to-know-you-part-1/}{b}
\B{Tsu'\-tey te Rong\-loa, A\-te\-yo\-'itan}, Tsu'tey of the Rongloa, son of Ateyo.~\LINK{http://wiki.learnnavi.org/index.php?title=Na\%27vi_from_Avatar_Movie\#Toruk_Makto}{wiki\slash a}

% tekre
\hi
\HT{tekre}{\B{\cp{tek}re}} \I{["tEk\textcorner.rE]} \emph{n.} skull. \\
\textsc{Derivatives}: \ding{43}
\on{1}.~\HL{nawmtoruktek}{\B{nawm\-to\-ruk\-tek}}, Toruk Makto Totem.

% tel
\hi
\HT{tel}{\B{tel}} \I{[tEl]} \emph{vtr.} receive (\B{ta}, from, \B{yoa}, in exchange for).   
\B{Oel tel 'u\-pxa\-ret le\-Na'\-vi a krr, new oe nì\-teng\-fya pam\-rel si\-vi nì\-'eyng.} When I receive a Na'vi message, I want to write the same way in response.~\LINK{http://forum.learnnavi.org/language-updates/verb-for-write/msg175592/\#msg175592}{ln}
\B{Fu\-ta nga\-ta tel pxen\-yoa srät, mll\-'e\-i\-an oel.} I'm happy to take cloth from you for finished garments.~\LINK{http://naviteri.org/2014/02/barter-and-exchange/}{b}
\B{Tì\-'eyng\-it oel to\-lel a krr, ay\-nga\-ru pa\-yeng.} When I receive an answer, I will let you know.~\LINK{http://forum.learnnavi.org/news-announcements/a-response-from-paul-frommer!/}{ln}
(\emph{a. conv. idiom}) \B{To\-\cp{lel}}! ``Eureka! I got it! I understand!''~\LINK{http://naviteri.org/2010/07/vocabulary-update/}{b}

% telem
\hi
\HT{telem}{\B{te\cp{lem}}} \I{[tE."lEm]} \emph{n.} cord. \\
\textsc{Derivatives}: \ding{43}
\on{1}.~\HL{txurtel}{\B{txur\-tel}}, rope. 
\on{2}.~\HL{waytelem}{\B{way\-te\-lem}}, song cord.

% telisi
\hi 
\HT{telisi}{\B{teli\cp{si}}} \I{[tE.li."{}si]} \emph{n.} whirlwind. 

% tem
\hi
\HT{tem}{\B{tem}} \I{[tEm]} \emph{vin.} shoot (\B{ne}, at).
\B{Tem rä\-'ä!} Don't shoot!~\LINK{http://naviteri.org/2017/03/titusemteri-concerning-shooting/}{b} 
\B{Oe\-ne fko te\-rem!} Someone is shooting at me!~\LINK{http://naviteri.org/2017/03/titusemteri-concerning-shooting/}{b} 
\B{Txo nga ze\-ne ti\-vem, tsa\-tì\-tu\-sem li\-vu muiä.} If you must shoot, let it be justified.~\LINK{http://naviteri.org/2017/03/titusemteri-concerning-shooting/}{b} 
(\ding{43}~\HL{toltem}{\B{tol\-tem}}; \HL{tsweykayon}{\B{tswa\-yon} (\emph{b. \ding{246}})})

% temrey
\hi
\HT{temrey}{\B{tem\cp{rey}}} \I{[tEm."REj]} \emph{n.} survival (\emph{in the face of danger}). 
\B{Fwam\-pop\-ä tem\-rey tsa\-tseng\-mì ngä\-zìk lu nì\-wotx.} It's very hard for a tapirus to survive there.~\LINK{http://naviteri.org/2011/05/some-miscellaneous-vocabulary/}{b}
(\ding{43}~\HL{emrey}{\B{em\-rey}})

% teng
\hi
\HT{teng}{\B{teng}} \I{[tEN]} \emph{adj.} equal, same. 
\B{Fì\-txon na ton a\-la\-he nì\-wotx pe\-lun ke lu teng?} Why is this night not the same as all the other nights?~\LINK{http://whyisthisnight.com/addl-more.html\#na\%27vi}{tn} 
\B{Tsa\-krr syay ay\-nge\-yä, syay o\-lo'\-ä oe\-yä la\-yu teng.} Then you will suffer the same fate as my clan.~\LINK{http://naviteri.org/2010/07/vocabulary-update/}{b} 
\B{Me\-unil\-tì\-ran\-tokx Tok\-tor Kì\-rey\-sä sì Tsyeyk\-ä ke lu teng ki steng.} The avatars of Grace and Jake aren't the same, but they are similar.~\LINK{http://naviteri.org/2011/09/miscellaneous-vocabulary/}{b}
(\emph{a. proverb}) \B{Sä\-fpìl a\-steng, tì\-kan a\-teng.} Great minds think alike [lit.: similar thought, same aim].~\LINK{http://naviteri.org/2011/09/miscellaneous-vocabulary/}{b}
(\ding{214}~\HL{keteng}{\B{ke\-teng}}) \\
\textsc{Derivatives}: \ding{43} 
\on{1}.~\HL{nIteng}{\B{nì\-teng}}, too, also, likewise. 
\on{2}.~\HL{tengfya}{\B{teng\-fya}}, as, same way as (conj.). 
\on{3}.~\HL{tengkrr}{\B{teng\-krr}}, while (conj.). 
\on{4}.~\HL{'awsiteng}{\B{'aw\-si\-teng}}, together. 
\on{5}.~\HL{keteng}{\B{ke\-teng}}, different. 
\on{6}.~\HL{kAteng}{\B{kä\-teng}}, spend time with, hang out with. 
\on{7}.~\HL{'awstengyem}{\B{'aw\-steng\-yem}}, join (two or more things) together.
\on{8}.~\HL{tengralI'u}{\B{teng\-ra\-lì\-'u}}, synonym.
\on{9}.~\HL{tengwotx}{\B{teng\-wotx}}, uniform, non\hyp{}diverse.

% tengfya
\hi
\HT{tengfya}{\B{\cp{teng}fya}} \I{["tEN.fja]} \emph{conj.} as. 
\B{Pam\-tseol ngop ay\-re\-nut | Mì ron\-sem\-ä tì\-fnu | Teng\-fya ngop sä\-ftxu\-yul | Mì hi\-fkey.} Music creates patterns | In the silence of the mind | As weavers do | In the physical word.~\LINK{http://www.pandorapedia.com/navi/music/ritual_music}{pp}
\B{Teng\-fya tsun tsi\-ve\-'a, lup\-ra el\-tur tì\-txen si nì\-txan.} As you can see, the style is very interesting.~\LINK{http://naviteri.org/2014/09/stxeli-alor-the-text/}{b}
(\ding{43}~\B{teng}, \B{fya\-'o}) \\
\textsc{Derivatives}: \ding{43}
\on{1}.~\HL{nItengfya}{\B{nì\-teng\-fya}}, in the same way.

% tengkrr
\hi
\HT{tengkrr}{\B{teng\cp{krr}}} \I{[tEN."k\s{r}]} \emph{conj.} while (\emph{doing s.t.}). \emph{Usually with a form of} ‹\B{er}› \emph{in the following verb.}
\B{Re\-rol teng\-krr ke\-rä | Ìlä fya\-'o a\-vol | Ne kxam\-tseng.} We sing our way | Down the eight paths | To the center.~\LINK{http://www.pandorapedia.com/navi/music/ritual_music}{pp}
\B{Na'\-vi ìlä ho\-'on kll\-kxo\-lem teng\-krr re\-rol.} The People were standing in a circle, singing.~\LINK{http://naviteri.org/2013/09/on-si-salewfya-shapes-and-directions/}{b}
\B{Ay\-nga\-ri teng\-krr ya wur sle\-ru nì\-'ul, sìl\-pey oe, li\-vu hel\-ku sang ul\-te te'\-lan le\-fpom.} As for you all, while it's getting cooler, I hope your homes may be warm and your hearts happy.~\LINK{http://naviteri.org/2013/09/on-si-salewfya-shapes-and-directions/}{b}
\B{Txon\-am teng\-krr tar\-mì\-ran oe kxam\-lä na'\-rìng, sro\-ler eo ut\-ral a\-tsawl txewm\-a vrr\-tep.} Last night as I was walking through the forest, a frightening demon appeared in front of a big tree.~\LINK{http://naviteri.org/2012/03/spring-vocabulary-part-1/}{b}
(\emph{a. idiom}) \B{Teng\-krr ftxo\-zä se\-rei\-yi aw\-nga, ke tswi\-va' ay\-lom\-tu\-ti ko!} A toast: To absent friends. [lit.: While we are celebrating, let us not forget those who we wish could also be here (but can't)]~\LINK{http://naviteri.org/2011/07/txantsana-ultxa-mi-siatll-great-meeting-in-seattle/}{b}
(\ding{43}~\HL{teng}{\B{teng}}, \HL{krr}{\B{krr}})

% tengralì'u
\hi 
\HT{tengralI'u}{\B{\cp{teng}ralì'u}} \I{["tEN.Ra.lI.Pu]} \emph{n.} synonym. 
\B{Me\-lì\-'u alu pxel sì na lu teng\-ra\-lì\-'u.} The two words \emph{pxel} and \emph{na} are synonyms.~\LINK{https://naviteri.org/2024/09/pxevola-liu-amip-two-dozen-new-words/}{b} 
(\ding{43}~\HL{teng}{\B{teng}}, \HL{ral}{\B{ral}}, \HL{lI'u}{\B{lì\-'u}}; \ding{214}~\HL{wAralI'u}{\B{wä\-ra\-lì\-'u}}) 

% tengwotx
\hi 
\HT{tengwotx}{\B{\cp{teng}wotx}} \I{["tEN.wot']} \emph{adj.} uniform, non\hyp{}diverse.
\B{Lì'\-fya le\-Na'\-vi lu teng\-wotx, ken ay\-reym\-sä\-tsa le\-ke\-teng\-syon Ey\-wa\-'e\-veng\-ä, ken ho\-et\-a ta\-yo a mì\-kam ay\-olo'.} The language of the Na'vi is consistent, despite the exo\hyp{}moon's varied landscapes and the vast expanse between clans.~\LINK{https://naviteri.org/2026/06/tskxekeng-a-mikyunfpi-pamrel-listening-exercise-text/}{b}
(\ding{43}~\HL{teng}{\B{teng}}, \HL{nIwotx}{\B{nì\-wotx}}) \\
\textsc{Derivatives}: \ding{43} 
\on{1}.~\HL{tItengwotx}{\B{tì\-teng\-wotx}}, uniformity.

% teri
\hi
\HT{teri}{\B{te\cp{ri}}} \I{[tE."Ri]} \emph{adp.} about, concerning.
(\emph{a.})
\B{Lu fra\-'u te\-ri tì\-ftxa\-vang.} Everything is about passion.~\LINK{http://forum.learnnavi.org/news-announcements/paul-frommer-at-grid-2010/}{ln}
\B{Nge\-yä te\-ri fay\-te\-le a ay\-sä\-nu\-me\-ri ngar ira\-yo sei\-yi ay\-oe nì\-wotx.} We all thank you for your teachings about these topics.~\LINK{http://forum.learnnavi.org/language-updates/combined-answers-1-\%28feb-16\%29/}{ln}
(\emph{b. with the verbs} \ding{43}~\B{\HL{pAngkxo}{päng\-kxo}, \HL{peng}{peng}, \HL{stawm}{stawm}, \HL{nume}{nu\-me}, \HL{oeyktIng}{oeyk\-tìng} te\-ri}) chat, tell, hear, learn, explain about.
\B{Fay\-upxa\-re\-mì oe pa\-yäng\-kxo te\-ri ho\-ren lì'\-fya\-yä le\-Na'\-vi.} In these messages I will chat about the rules of the Na'vi language.~\LINK{http://naviteri.org/2010/06/first-post/}{b}
\B{Oe pi\-veng ay\-nga\-ru nì\-'it te\-ri Txe\-wì alu Na'\-vi\-yä 'e\-veng.} I'd like to tell you a bit about Txewì, a Na'vi child.~\LINK{http://naviteri.org/2010/07/ting-mikyun-fte-tslivam-listening-comprehension-1-3/}{b}
\B{Krr a sto\-lawm Txe\-wì te\-ri hem a poe\-ru lo\-len, 'o\-le\-fu ke\-ftxo nì\-txan.} When Txewì heard about what happened to her, he felt very upset.~\LINK{http://naviteri.org/2010/07/ting-mikyun-fte-tslivam-listening-comprehension-1-3/}{b}
\B{Oel vewng fu\-ta ay\-e\-veng ni\-vu\-me te\-ri ay\-ewll na'\-rìng\-ä.} I see to it that the children learn about the forest plants.~\LINK{http://naviteri.org/2010/09/getting-to-know-you-part-2/}{b}
\B{Wo\-leyn Ìstawl yey\-fyat mì hll\-te fte oeyk\-ti\-vìng fra\-po\-ru tì\-hawl\-te\-ri sne\-yä.} Ìstaw drew a line on the ground to explain his plan to everyone.~\LINK{http://naviteri.org/2013/09/on-si-salewfya-shapes-and-directions/}{b}

% terkup
\hi
\HT{terkup}{\B{\cp{ter}kup}} \I{["tER.kup\textcorner]} \emph{or} \I{[.kUp\textcorner]} (\textsc{rn}: \B{\cp{ter}kùp} \I{["tER.kUP\textcorner]}) \emph{vin.} die. 
\B{Kop oe\-ru lo\-lä\-ngu lie a ha\-pxì\-tu soaiä ter\-kup.} Sadly, I too have had the experience of a family member dying.~\LINK{http://naviteri.org/2013/01/awvea-posti-zisita-amip-first-post-of-the-new-year/}{b}
\B{Ngal po\-leng oer fmawn\-ta po to\-ler\-kup.} You told me that he died.~\LINK{http://naviteri.org/2011/08/reported-speech-reported-questions/}{b}
\B{Trr\-am to\-ler\-kä\-ngup sa'\-nok a\-yaw\-ne.} My dear mother died yesterday.~\LINK{http://naviteri.org/2012/01/more-additions-to-the-lexicon/}{b}
\B{Po tä\-pey\-kì\-ye\-ver\-kei\-up nì\-näk.} I'm so jazzed that he may be about to drink himself to death.~\LINK{http://naviteri.org/2010/12/catching-up/}{b}
(\emph{a. coll. idiom}) \B{[verb] nì\-ftxan kum\-a ter\-kup}, \emph{commenting on an action that is so intense that it results in dying}; ``like crazy, \ldots{} to death''. 
\B{Sran, sran, oe\-ri pì\-tìk tsa\-tse\-nget a mì tal! Fkxa\-ke nì\-ftxan kum\-a ter\-kup.} Yeah, yeah, scratch that place on my back! It itches like crazy!~\LINK{http://naviteri.org/2020/07/vospxivol-lefpom-happy-august/}{b} \\
\textsc{Derivatives}: \ding{43}
\on{1}.~\HL{tIterkup}{\B{tì\-ter\-kup}}, death, dying.

% teswotìng 
\hi
\HT{teswotIng}{\B{te\cp{swo}tìng}} \I{[tE."{}swo.tIN]} \emph{vtr.} grant. 
\B{Fì\-ta\-ron\-yur a\-pxan te\-swo\-ti\-vìng äiet a\-ngay.} Grant this worthy warrior true vision.~\LINK{https://wiki.learnnavi.org/Na\%27vi_from_Avatar_Movie\#Becoming_one_of_the_Na.27vi}{a\textsuperscript{c}}

% tete
\hi
\HT{tete}{\B{\cp{te}te}} \I{["tE.tE]} \emph{adj.} dull (\emph{of a blade}).
(\ding{214}~\HL{litx}{\B{litx}}, \ding{43}~\HL{fwem}{\B{fwem}}) 

% tewng
\hi
\HT{tewng}{\B{tewng}} \I{[tEwN]} \emph{n., text.} loincloth.

% tewti
\hi
\HT{tewti}{\B{\cp{tew}ti}} \I{["tEw.ti]} \emph{intj. expression of (pleasant) surprise}, ``Wow! Whoa!''.
\B{Tew\-ti, nga lu tsul\-fä\-tu i\-'en\-ä.} Wow, you are a master on the \emph{i'en}.~\LINK{http://naviteri.org/2011/01/new-year-new-vocabulary/}{b}

% teya / si  
\hi
\HT{teya}{\B{te\cp{ya}}} \I{[tE."ja]} \emph{adj.} full (\HL{ta}{\B{ta}}, of).
\B{Oe\-ri ta pe\-yä fa\-hew a\-ke\-wong on\-tu te\-ya lä\-ngu.} My nose is full of his alien smell.~\LINK{http://languagelog.ldc.upenn.edu/nll/?p=1977}{ll\slash a}
\B{Eo ay\-fo a fya\-'o la\-mu ay\-skxe\-ta te\-ya sì re\-nul\-ke.} The path ahead of them was full of rocks and irregular.~\LINK{http://naviteri.org/2013/09/on-si-salewfya-shapes-and-directions/}{b}
\B{Ay\-nge\-yä ay\-srr tì\-flä\-ta te\-ya li\-vu ul\-te si\-va\-lew nì\-'o' nì\-'aw.} May your days be full of success and proceed only in a fun way.~\LINK{http://naviteri.org/2013/12/mipa-zisit-lefpom-happy-new-year/}{b} \\
\textsc{ii}. \B{te\cp{ya} si} \I{[tE."ja si]} \emph{vin.} (\emph{a.}) fill, make full. 
\B{Fì\-po\-stì mek\-tseng\-ur te\-ya si.} This post fills a gap (in our understanding of Na'vi syntax).~\LINK{http://naviteri.org/2013/04/zun-zel-counterfactual-conditionals/}{b}  
(\emph{b. idiom}) \B{oe\-ru te\-ya si}, ``s.t. fills me with satisfaction\slash joy''. 
\B{Sì ay\-zì\-sìt\-ä ka\-to | Sì ekong te'\-lan\-ä | Te'\-lan le\-Na'\-vi | Oe\-ru te\-ya si.} And the rhythm of the years | And the beats of the hearts | The hearts of the People | Fill me.~\LINK{http://wiki.learnnavi.org/index.php/Weaving_Song}{wiki}
\B{Fay\-lì\-'u a\-lor oe\-ru te\-ya si nì\-txan, ul\-te ke tswa\-ya' oel sat kaw\-krr.} These beautiful words fill me with joy and I'll never forget them.~\LINK{http://forum.learnnavi.org/language-updates/life-death-and-gerunds/}{ln}
\B{Tì\-tstun\-wi ay\-nge\-yä oe\-ru te\-ya so\-li.} Your kindness have filled me with joy.~\LINK{http://forum.learnnavi.org/community-calendar/karyu-pawls-birthday/msg312733/\#msg312733}{ln}
\B{Moe\-ru te\-ya si nì\-ngay.} We're greatly touched.~\LINK{http://naviteri.org/2014/09/stxeli-alor-the-text/}{b}
(\ding{214}~\HL{mek}{\B{mek}})

% teylu
\hi
\HT{teylu}{\B{\cp{tey}lu}} \I{["tEj.lu]} \emph{n., fauna} beetle larva(e).   
\B{Tey\-lu kì\-mar lì\-yu a fì\-'u oe\-ru te\-ya si.} It fills me with joy that \emph{teylu} is about to be in season.~\LINK{http://naviteri.org/2011/10/more-vocabulary-a-bit-of-grammar/}{b}
\B{Oe new yi\-vom tey\-lut.} I want to eat \emph{teylu}.~\LINK{http://naviteri.org/2011/03/word-order-and-case-marking-with-modals/}{b}
\B{Fol fnan fu\-ta 'em tey\-lut a fì\-'u sìn pe\-yì?} How good are they at cooking \emph{teylu}?~\LINK{http://naviteri.org/2013/01/lifyari-po-peyi-how-good-is-her-navi/}{b}
\B{Fì\-tey\-lu\-ri ke nar\-mew Va'\-ru yi\-vom nì\-yll.} Va'ru didn't want to share this \emph{teylu} with the Omatikaya.~\LINK{http://naviteri.org/2012/01/more-additions-to-the-lexicon/}{b} \\
\textsc{Derivatives}: \ding{43}
\on{1}.~\HL{teylupil}{\B{tey\-lu\-pil}}, \emph{teylu}\hyp{}face (\emph{childish insult}). 

% teylupil
\hi 
\HT{teylupil}{\B{\cp{tey}lupil}} \I{["tEj.lu.pil]} \emph{n.} \emph{teylu}\hyp{}face (\emph{a somewhat childish insult}). 
(\ding{43}~\HL{teylu}{\B{tey\-lu}}, \HL{pil}{\B{pil}})

% teynga 
\hi
\HT{teynga}{\B{\cp{tey}nga}} \I{["tEj.Na]} \emph{conj. to form indirect questions; shortened form of} \ding{43}~\HL{tI'eyng}{\B{tì\-'eyng a}}.   
\B{Tey\-nga lum\-pe fo ho\-lum ke lu law.} It's not clear why they left.~\LINK{http://naviteri.org/2011/08/reported-speech-reported-questions/}{b}
\B{Tey\-nga pe\-su ni\-vew hi\-vum ke tsran\-ten, po\-ru pll\-txe san ru\-txe 'i\-vì\-'awn.} No matter who wants to leave, tell them to please stay.~\LINK{http://naviteri.org/2013/03/whoever-whatever-whenever/}{b}
(\emph{a. with adp.s})
\B{Ay\-oe pe\-räng\-kxo te\-ri tey\-nga lum\-pe po ho\-lum.} We're chatting about why he left.~\LINK{https://naviteri.org/2024/12/odds-and-ends/}{b}
\B{Tì\-flä la\-tem ìlä seyng\-a ftxey fkol sä\-nu\-met li\-vek fu\-ke.} Success depends on whether or not one follows instructions.~\LINK{http://naviteri.org/2012/07/meetings-waterfalls-and-more/}{b}
(\ding{43}~\HL{teyngla}{\B{teyng\-la}}, \HL{teyngta}{\B{teyng\-ta}}, \HL{teyngra}{\B{teyng\-ra}}, \HL{teyngA}{\B{tey\-ngä}}, \HL{teyngria}{\B{teyng\-ri\-a}})

% teyngä
\hi
\HT{teyngA}{\B{\cp{tey}ngä}} \I{["tEj.N\ae]} \emph{conj. to form indirect questions; shortened form of} \ding{43}~\HL{tI'eyng}{\B{tì\-'ey\-ngä a}}.   
\B{Tì\-run\-ìl tey\-ngä lum\-pe po ho\-lum oe\-ti ke\-ftxo 'ey\-ko\-le\-fu.} The discovery of why he left saddened me.~\LINK{https://naviteri.org/2024/12/odds-and-ends/}{b}
(\ding{43}~\HL{teynga}{\B{tey\-nga}}, \HL{teyngla}{\B{teyng\-la}}, \HL{teyngta}{\B{teyng\-ta}}, \HL{teyngra}{\B{teyng\-ra}}, \HL{teyngria}{\B{teyng\-ri\-a}})

% teyngla
\hi
\HT{teyngla}{\B{\cp{tey}ngla}} \I{["tEjN.la]} \emph{conj. to form indirect questions; shortened form of} \ding{43}~\HL{tI'eyng}{\B{tì\-'eyng\-ìl a}}.
\B{Teyng\-la lum\-pe po ho\-lum oe\-ti hey\-ko\-lang\-ham.} Why he left made me laugh.~\LINK{https://naviteri.org/2024/12/odds-and-ends/}{b}
(\ding{43}~\HL{teynga}{\B{tey\-nga}}, \HL{teyngta}{\B{teyng\-ta}}, \HL{teyngra}{\B{teyng\-ra}}, \HL{teyngA}{\B{tey\-ngä}}, \HL{teyngria}{\B{teyng\-ri\-a}})

% teyngra
\hi
\HT{teyngra}{\B{\cp{teyng}ra}} \I{["tEjN.Ra]} \emph{conj. to form indirect questions; shortened form of} \ding{43}~\HL{tI'eyng}{\B{tì\-'eyng\-ur a}}.
\B{Ru\-txe law si\-vi teyng\-ra lum\-pe po ho\-lum.} Please clarify why he left.~\LINK{https://naviteri.org/2024/12/odds-and-ends/}{b}
(\ding{43}~\HL{teynga}{\B{tey\-nga}}, \HL{teyngla}{\B{teyng\-la}}, \HL{teyngta}{\B{teyng\-ta}}, \HL{teyngA}{\B{tey\-ngä}}, \HL{teyngria}{\B{teyng\-ri\-a}})

% teyngria
\hi
\HT{teyngria}{\B{\cp{teyng}ria}} \I{["tEjN.Ri.a]} \emph{conj. to form indirect questions; shortened form of} \ding{43}~\HL{tI'eyng}{\B{tì\-'eyng\-ri a}}.   
\B{Teyng\-ria lum\-pe po ho\-lum oel ke tslam ke\-'ut.} I understand nothing about why he left.~\LINK{https://naviteri.org/2024/12/odds-and-ends/}{b}
(\ding{43}~\HL{teynga}{\B{tey\-nga}}, \HL{teyngla}{\B{teyng\-la}}, \HL{teyngta}{\B{teyng\-ta}}, \HL{teyngra}{\B{teyng\-ra}}, \HL{teyngA}{\B{tey\-ngä}})

% teyngta
\hi
\HT{teyngta}{\B{\cp{tey}ngta}} \I{["tEjN.ta]} \emph{conj. to form indirect questions; shortened form of} \ding{43}~\HL{tI'eyng}{\B{tì\-'eyng\-it a}}.
\B{Vo\-lin oel teyng\-ta Ney\-ti\-ri kä pe\-seng\-ne.} I asked where Neytiri was going.~\LINK{http://naviteri.org/2011/08/reported-speech-reported-questions/}{b} 
\B{Ngal ke tslä\-ngam teyng\-ta fì\-tì\-ngä\-zìk\-ìri pe\-fyin\-ep\-'ang.} Unfortunately you don't understand how complex this problem is.~\LINK{http://naviteri.org/2012/07/fivospxiya-aylifyavi-amip-this-months-new-expressions-pt-2/}{b}
\B{Ke o\-mä\-ngum oel teyng\-ta ftu pe\-seng zo\-la\-'u tsa\-pam\-rel\-vi alu t.} I'm afraid I don't know where that letter t came from.~\LINK{http://naviteri.org/2014/02/barter-and-exchange/comment-page-1/\#comment-2612}{b} 
(\ding{43}~\HL{teynga}{\B{tey\-nga}}, \HL{teyngla}{\B{teyng\-la}}, \HL{teyngra}{\B{teyng\-ra}}, \HL{teyngA}{\B{tey\-ngä}}, \HL{teyngria}{\B{teyng\-ri\-a}})

% teyr
\hi
\HT{teyr}{\B{teyr}} \I{[tEjR]} \emph{adj.} white. 
\B{Ro\-lun ay\-oel tsrul\-mì hì\-'i\-a pxe\-lo\-it a\-teyr.} We found three little white eggs in the nest.~\LINK{http://naviteri.org/2013/06/looking-back-looking-forward/}{b}
(\ding{214}~\HL{layon}{\B{la\-yon}}) \\
\textsc{Derivatives}: \ding{43}
\on{1}.~\HL{teyrpin}{\B{teyr\-pin}}, the color white.

% teyrpin
\hi
\HT{teyrpin}{\B{\cp{teyr}pin}} \I{["tEjR.pin]} \emph{n.} the color \ding{43}~\HL{teyr}{\B{teyr}}.
(\ding{43}~\HL{'opin}{\B{'o\-pin}})

% -ti 
\hi
\HT{-ti}{\B{-ti}} \I{[.ti]} \emph{suf. to mark the patientive case of a n. ending in a vowel, pseudo\hyp{}vowel, diphthong, or consonant, preferred ending in} \textsc{rn}.
(\ding{43}~\HL{-it}{\B{-it}}, \HL{-t}{\B{-t}}) 

% ti'a
\hi 
\HT{ti'a}{\B{\cp{ti}'a}} \I{["ti.Pa]} \emph{adv.} surprisingly (\emph{for unexpectedly positive outcomes}). 
\B{Ra\-mu ke lu txur nìtxan, slä u\-van\-it yo\-lo\-ra', ti\-'a.} Ra\-mu isn't very strong, but surprisingly, he won the game.~\LINK{http://naviteri.org/2022/11/krr-ao-an-exciting-time/}{b} 
(\ding{214}~\HL{um'a}{\B{um\-'a}}) 

% tiam
\hi
\HT{tiam}{\B{\cp{ti}am}} \I{["ti.am]} \emph{vtr.} count (\HL{ta}{\B{ta}}, from; \HL{vay}{\B{vay}}, (up) to; \HL{raw}{\B{raw}}, (down) to).
\B{Ru\-txe ti\-vi\-am ay\-srok\-it tsa\-krr hol\-pxay\-ti pi\-veng oer.} Please count the beads and tell me the number.~\LINK{http://naviteri.org/2014/11/twenty-before-the-holidays/}{b} 
\B{Ti\-am (ta pxey) vay vol.} Count (from three) to eight.~\LINK{https://forum.learnnavi.org/language-updates/'count-to'-wa-and-sequential-genitive-(discussion)/}{ln}
\B{Ti\-am ta vo\-mrr raw pxey.} Count down from thirteen to three.~\LINK{http://naviteri.org/2019/06/50a-liu-amip-40-new-words\%ef\%bb\%bf/}{b} \\
\textsc{Derivatives}: \ding{43}
\on{1}.~\HL{ketsuktiam}{\B{ke\-tsuk\-ti\-am}}, uncountable, infinite.

% til
\hi
\HT{til}{\B{til}} \I{[til]} \emph{n.} joint, hinge. 
(\ding{43}~\HL{il}{\B{il}}) \\
\textsc{Derivatives}: \ding{43}
\on{1}.~\HL{kinamtil}{\B{ki\-nam\-til}}, knee.
\on{2}.~\HL{pxuntil}{\B{pxun\-til}}, elbow.
\on{3}.~\HL{ventil}{\B{ven\-til}}, ankle.
\on{4}.~\HL{tilI'u}{\B{ti\-lì\-'u}}, conjunction.
\on{5}.~\HL{zektil}{\B{zek\-til}}, knuckle.

% tilì'u
\hi 
\HT{tilI'u}{\B{\cp{ti}lì'u}} \I{["ti.lI.Pu]} \emph{n.} conjunction. 
(\ding{43}~\HL{til}{\B{til}}, \HL{lI'u}{\B{lì\-'u}})

% tinan 
\hi
\HT{tinan}{\B{\cp{ti}nan}} \I{["ti.nan]} \emph{n.} reading.   
\B{Tsmìm\-ìri wätx fol ti\-nan\-it nì\-ngay.} They're really bad at reading animal tracks.~\LINK{http://naviteri.org/2011/09/\%E2\%80\%9Cby-the-way-what-are-you-reading\%E2\%80\%9D/}{b}
(\ding{43}~\HL{inan}{\B{inan}})

% Tipani
\hi
\HT{Tipani}{\B{\cp{Ti}pani}} \I{["ti.pa.ni]} \emph{n.} \emph{clan name, expert warriors and hunters}.

% tirea
\hi
\HT{tirea}{\B{ti\cp{re}a}} \I{[ti."RE.a]} \emph{n.} spirit (\emph{of a being}). 
\B{Nga\-ri hu Ey\-wa sa\-lew ti\-rea, tokx 'ì\-'awn, slu Na'\-vi\-yä ha\-pxì.} Your spirit goes with Eywa, your body stays behind to become part of the People.~\LINK{http://wiki.learnnavi.org/index.php?title=Na\%27vi_from_Avatar_Movie\#Jake.27s_final_kill}{a\slash wiki}
\B{Slä am\-'a\-lu\-ke nge\-yä ti\-re\-al tsa\-tse\-nget ta\-ye\-i\-ok.} But without a doubt your spirit will be there.~\LINK{http://naviteri.org/2012/07/fivospxiya-aylifyavi-amip-this-months-new-expressions-pt-2/comment-page-1/\#comment-1569}{b}
\B{nrr a ti\-re\-a}, glowing spirit.~\LINK{http://naviteri.org/2012/07/fivospxiya-aylifyavi-amip-this-months-new-expressions-pt-2/}{b} \\
\textsc{Derivatives}: \ding{43}
\on{1}.~\HL{tireafya'o}{\B{tirea\-fya\-'o}}, spirit path. 
\on{2}.~\HL{tireaioang}{\B{tirea\-ioang}}, spirit animal. 
\on{3}.~\HL{tireapAngkxo}{\B{tirea\-päng\-kxo}}, commune.
\on{4}.~\HL{slantire}{\B{slan\-ti\-re}}, inspiration. 
\on{5}.~\HL{tiretu}{\B{ti\-re\-tu}}, shaman.
\on{6}.~\HL{nItrea}{\B{nìt\-re\-a}}, in spirit.

% tireafya'o
\hi
\HT{tireafya'o}{\B{ti\cp{re}afya'o}} \I{[ti."RE.a.""fja.Po]} \emph{n.} spirit path. 
\B{Ay\-wayl yìm ki\-fkey\-ä | 'Ì\-he\-yut a\-vo\-mrr | Sìn ti\-rea\-fya\-'o a\-vol | Na way\-te\-lem\-ä hìng.} The songs bind | the thirteen spirals of the solid world | To the eight spirit paths | Like the threads of a Songcord.~\LINK{http://www.pandorapedia.com/navi/music/ritual_music}{pp}
(\ding{43}~\HL{tirea}{\B{ti\-rea}}, \HL{fya'o}{\B{fya\-'o}})

% tireaioang
\hi
\HT{tireaioang}{\B{ti\cp{re}aioang}} \I{[ti."RE.a.i.""{}o.aN]} \emph{n.} spirit animal.   
\B{Nge\-yä Ti\-re\-a\-i\-o\-ang zo\-la\-'u a krr, law la\-yu nga\-ru.} When your Spirit Animal comes, you will know.~\LINK{http://forum.learnnavi.org/language-updates/auxilary-verb-si-possessive-dative-krr/}{ln}
(\ding{43}~\HL{tirea}{\B{ti\-rea}}, \HL{ioang}{\B{io\-ang}})

% tireapängkxo
\hi
\HT{tireapAngkxo}{\B{ti\cp{re}apängkxo}} \I{[ti."RE.a.p$\cdot$\ae{}N.""k'$\cdot$o]} \emph{vin., inf.}\on{45} commune (\emph{about intimate contact}).   
\B{Ke lu kaw\-tu a nul\-ni\-vew oe po\-hu ti\-rea\-pi\-väng\-kxo äo Ut\-ral Ay\-mok\-ri\-yä.} There's nobody I'd rather commune with under the Tree of Voices.~\LINK{http://wiki.learnnavi.org/Corpus\#Lemondrop.com}{ld\slash wiki}
(\ding{43}~\HL{tirea}{\B{ti\-rea}}, \HL{pAngkxo}{\B{päng\-kxo}})

% tiretu
\hi 
\HT{tiretu}{\B{ti\cp{re}tu}} \I{[ti."RE.tu]} \emph{n.} shaman.
(\ding{43}~\HL{tirea}{\B{ti\-re\-a}}, \HL{tute1}{\B{tu\-te}})

% tì-
\hi
\HT{tI-}{\B{tì-}} \I{[tI.]} \emph{unprod. pref. to turn adjectives, verbs or nouns into abstract nouns.}
(\ding{43}~\HL{sA-}{\B{sä-}})

% tì- us
\hi
\B{tì-} + ‹\B{us}› \emph{prod. process to form the gerund of a verb}.
\B{Ay\-nge\-yä tì\-ftu\-si\-a 'o' li\-vu nì\-'aw!} May your studying be fun.~\LINK{http://naviteri.org/2012/10/sound-files-added-to-previous-post/}{b}
\B{Tì\-tu\-sa\-ron mo\-wan lu oer nì\-ngay.} Hunting really turns me on.~\LINK{http://naviteri.org/2010/07/vocabulary-update/}{b}
\B{Oel kan\-'ìn tì\-'u\-sem\-it.} I specialize in cooking.~\LINK{http://naviteri.org/2010/09/getting-to-know-you-part-1/comment-page-1/\#comment-296}{b}
\B{Ko\-ren A\-'aw\-ve Tì\-ru\-sey\-ä 'Aw\-si\-teng}, the First Rule of Living Together.~\LINK{http://wiki.learnnavi.org/index.php?title=Corpus\#Good_Morning_America}{wiki}

% tì'al
\hi 
\HT{tI'al}{\B{tì\cp{'al}}} \I{[tI."Pal]} \emph{n.} wastefulness.
\B{Tì\-'al lu zop\-lo a tsa\-ri ke tsun txoa li\-vu nì\-ftue.} Wastlefulness is an offense that cannot easily be forgiven.~\LINK{http://naviteri.org/2014/07/tsawlultxamaw-a-ayliu-post-meetup-words/}{b} 
(\ding{43}~\HL{'al}{\B{'al}})

% tì'awm
\hi
\HT{tI'awm}{\B{tì\cp{'awm}}} \I{[tI."Pawm]} \textsc{i}. \emph{n.} camping. \\
\textsc{ii}. \B{tì\cp{'awm} si} \I{[tI."Pawm si]} \emph{vin.} to camp. 
(\ding{43}~\HL{'awm}{\B{'awm}})

% tì'awnì'aw
\hi 
\HT{tI'awnI'aw}{\B{tì\cp{'aw}nì'aw}} \I{[tI."Paw.nI.Paw]} \emph{n.} uniqueness.  
(\ding{43}~\HL{le'awnI'aw}{\B{le\-'aw\-nì\-'aw}})

% tì'awpo
\hi
\HT{tI'awpo}{\B{tì\cp{'aw}po}} \I{[tI."Paw.po]} \emph{n.} individuality (\emph{negative connotation}); selfishness. 
(\ding{43}~\HL{'awpo}{\B{'aw\-po}})

% tì'älek
\hi 
\HT{tI'Alek}{\B{tì\cp{'ä}lek}} \I{[tI."P\ae.lEk\textcorner]} \emph{n.} determination. 
\B{Pe\-yä tì\-'ä\-lek oe\-ru ro\-lo\-'a nì\-txan.} His determination impressed me greatly.~\LINK{http://naviteri.org/2021/09/aawa-ayliu-amip-a-few-new-words/}{b} 
(\ding{43}~\HL{'Alek}{\B{'ä\-lek}})

% tì'efu
\hi
\HT{tI'efu}{\B{tì\cp{'e}fu}} \I{[tI."PE.fu]} \emph{n.} feeling.
\B{Fwa lawk ay\-sì\-'efu\-ti ay\-ey\-lan\-kip lu fpom\-ro\-nga'.} It's healthy among friends to discuss feelings.~\LINK{http://naviteri.org/2014/10/tson-si-fpomron-obligation-and-mental-health/}{b}
(\emph{a. idiom}) \B{tì\-'e\-fu\-mì oe\-yä}, ``in my opinion''. \B{New oe ti\-vìng ay\-ngar lì\-'ut a tì\-'e\-fu\-mì oe\-yä lu lor fra\-to mì lì'\-fya le\-Na'\-vi.} I want to give you the word that, in my opinion, is the most beautiful in the Na'vi language.~\LINK{http://forum.learnnavi.org/language-updates/trr-rrtaya/}{ln}
\B{tì\-'e\-fu\-wä oe\-yä}, ``contrary to my opinion''.~\LINK{https://forum.learnnavi.org/language-updates/'count-to'-wa-and-sequential-genitive-(discussion)/}{ln}
\B{Tì\-'efu\-wä oe\-yä, fpìl Pey\-ral\-ìl fu\-ta ke ze\-ne ay\-oeng ki\-vä.} I think we have to go, but Peyral doesn't.~\LINK{https://forum.learnnavi.org/language-updates/'count-to'-wa-and-sequential-genitive-(discussion)/}{ln}
(\ding{43}~\HL{'efu}{\B{'efu}})

% tì'em
\hi
\HT{tI'em}{\B{tì\cp{'em}}} \I{[tI."PEm]} \emph{n.} the art of cooking.
\B{So\-rewn\-ìl kan\-'ìn tì\-'em\-it nì\-txan ul\-te kxawm tsa\-txe\-le me\-nga\-ne za\-'a\-tsu nì\-'eng.} Sorewn is really into cooking and you two might share that in common.\LINK{http://naviteri.org/2010/09/getting-to-know-you-part-1/}{b}
(\ding{43}~\HL{'em}{\B{'em}})

% tì'eylan
\hi
\HT{tI'eylan}{\B{tì\cp{'ey}lan}} \I{[tI."PEj.lan]} \emph{n.} friendship.
\B{Nge\-yä tì\-eyk\-tan\-ìri, tì\-slan\-ìri sì tsran\-ten fra\-to a tì\-'ey\-lan\-ìri a ka ay\-zì\-sìt nì\-wotx, sei\-yi oe ira\-yo nì\-txan.} Thank you so much for your leadership, your support, and most importantly your friendship throughout the years.~\LINK{http://naviteri.org/2012/03/spring-vocabulary-part-1/}{b}
(\ding{43}~\HL{'eylan}{\B{'eylan}}) \\
\textsc{Derivatives}: \ding{43}
\on{1}.~\HL{tI'eylanga'}{\B{tì\-'ey\-la\-nga'}}, friendly (nfp.).

% tì'eylanga'
\hi 
\HT{tI'eylanga'}{\B{tì\cp{'ey}langa'}} \I{[tI."PEj.la.NaP]} \emph{adj., nfp.} friendly. 
\B{Nge\-yä ay\-lì\-'u\-ri a\-tì\-'ey\-la\-nga' sei\-yi ira\-yo.} Thank you for your friendly words.~\LINK{https://naviteri.org/2024/09/pxevola-liu-amip-two-dozen-new-words/}{b}
(\ding{43}~\HL{let'eylan}{\B{let\-'ey\-lan}}; \HL{tI'eylan}{\B{tì\-'ey\-lan}}) 

% tì'eyng
\hi
\HT{tI'eyng}{\B{tì\cp{'eyng}}} \I{[tI."PEjN]} \emph{n.} answer, response.   
\B{Fo\-ru 'upxa\-ret oel fpo\-le', slä vay set ke pa\-mä\-hä\-ngem kea tì\-'eyng.} I have sent them a message, but up to now no answer has arrived.~\LINK{http://forum.learnnavi.org/news-announcements/a-response-from-paul-frommer!/}{ln}
\B{Tì\-'eyng\-it oel to\-lel a krr, tsa\-krr pa\-ye\-'un swey\-a fya\-'ot a za\-mi\-vu\-nge oel ay\-ngar ay\-lì\-'ut ho\-ren\-ti\-sì lì'\-fya\-yä le\-Na'\-vi.} When I receive an answer, I will then decide the best way to bring you the words and rules of Na'vi.~\LINK{http://forum.learnnavi.org/news-announcements/a-response-from-paul-frommer!/}{ln}
\B{Ay\-nge\-yä sì\-pawm\-ìri kop fma\-yi fì\-tse\-nge ti\-vìng sì\-'eyng\-it.} As for your questions, I will also try to give answers here.~\LINK{http://naviteri.org/2010/06/first-post/}{b}
(\emph{a. to form indirect questions}) \B{Vo\-lin oel tì\-'eyng\-it a Ney\-ti\-ri kä pe\-seng\-ne.} I asked where Neytiri was going.~\LINK{http://naviteri.org/2011/08/reported-speech-reported-questions/}{b} 
(\emph{b. contracted forms for indirect questions:}) \HL{teynga}{\B{tey\-nga}} \HL{teyngla}{\B{teyng\-la}}, \HL{teyngta}{\B{teyng\-ta}}, \HL{teyngra}{\B{teyng\-ra}}, \HL{teyngA}{\B{tey\-ngä}}, \HL{teyngria}{\B{teyng\-ri\-a}}
(\ding{43}~\HL{'eyng}{\B{'eyng}})

% tì'ewan
\hi 
\HT{tI'ewan}{\B{tì\cp{'e}wan}} \I{[tI."PE.wan]} \emph{n.} youth (\emph{time of life}). 
(\emph{a. proverb}) \B{Nì\-win 'ìp tì\-'e\-wan, nì\-win pä\-hem tì\-ko\-ak.} Youth vanishes quickly, old age arrives quickly.~\LINK{https://naviteri.org/2025/04/mevosinga-liu-amip-twenty-new-words/}{b}
(\ding{43}~\HL{'ewan}{\B{'e\-wan}}; \ding{214}~\HL{tIkoak}{\B{tì\-ko\-vak}})

% tì'i'a
\hi
\HT{tI'i'a}{\B{tì\cp{'i'}a}} \I{[tI."PiP.a]} \emph{n.} ending, conclusion. 
\B{Tì\-'i'\-a\-ri tsam\-ä ze\-ne Saw\-tu\-te vi\-ve\-lek ta\-lun tì\-txur Ey\-wa\-yä.} At the end of the war, the Sky People had to give up due to the power of Eywa.~\LINK{http://naviteri.org/2012/03/spring-vocabulary-part-1/}{b}
(\emph{a. idiom, poetic}) \B{tì\-'i'\-a\-vay krr\-ä}, forever, 'til the end of time.~\LINK{http://forum.learnnavi.org/language-updates/karyuru-frakrr-irayo-si-ko!/}{ln}
\B{To\-ler\-kup tu\-te; pam\-tseo pe\-yä tì\-'i'\-a\-vay krr\-ä ra\-yey.} The man has died; his music will live forever.~\LINK{http://naviteri.org/2015/06/na-ayskxe-mi-telan-sad-news/}{b}
(\emph{b. to tell the full hour}) \B{tì\-'i'\-a trr\-pxì\-yä a\-mrr\-ve}, `five o'clock', \emph{usually abbreviated to} \B{trr\-pxì a\-mrr\-ve}, `five o'clock'  \emph{or even} \B{Lu mrr\-ve.} `It's five.'~\LINK{https://naviteri.org/2024/05/krrteri-about-time/}{b}
(\ding{214}~\HL{sngA'itseng}{\B{sngä\-'i\-tseng}}, \HL{sngA'ikrr}{\B{sngä\-'i\-krr}} \ding{43}~\HL{'i'a}{\B{'i'a}}) \\
\textsc{Derivatives}: \ding{43}
\on{1}.~\HL{tI'iluke}{\B{tì\-'i\-lu\-ke}}, endless, never\hyp{}ending.

% tì'iluke
\hi 
\HT{tI'iluke}{\B{tì\cp{'i}luke}} \I{[tI."Pi.lu.kE]} \emph{adj.} endless, never\hyp{}ending. 
\B{Hum zì\-sìt a\-lal, pä\-hem pum a\-mip. Yo'\-ko\-fya a\-tì\-'i\-lu\-ke.} The old year leaves, the new one arrives. An endless cycle.~\LINK{http://naviteri.org/2017/12/zisit-amip-lefpom-ma-eylan-happy-new-year-friends/}{b}
(\ding{43}~\HL{tI'i'a}{\B{tì\-'i\-'a}}, \HL{luke}{\B{lu\-ke}}; \HL{txewluke}{\B{txew\-lu\-ke}}) \\
\textsc{Derivatives}: \ding{43}
\on{1}.~\HL{nIt'iluke}{\B{nìt\-'i\-lu\-ke}}, never\hyp{}endingly, forever. 

% tì'ipu
\hi
\HT{tI'ipu}{\B{tì\cp{'i}pu}} \I{[tI."Pi.pu]} \emph{n.} the abstract concept of being humorous, that is, humor in general.
\B{Sra\-ke tsun nga ri\-vun fì\-tì\-fkey\-tok\-mì a tì\-'i\-put?} Can you find the humor in this situation?~\LINK{http://naviteri.org/2012/01/mipa-zisit-ayliu-amip-new-words-for-the-new-year/}{b}
(\ding{43}~\HL{'ipu}{\B{'ipu}}, \HL{sA'ipu}{\B{sä\-'ipu}})

% tì'o'
\hi
\HT{tI'o'}{\B{tì\cp{'o'}}} \I{[tI."PoP]} \emph{n.} fun, excitement.
\B{Mì zì\-sìt a\-mip lì\-ye\-vu ay\-nga\-ru nì\-wotx txan\-a fpom\-tokx, fpom\-ron, tì\-yawn, sì tì\-'o'.} May you have much physical health, mental health, love and fun in the new year.~\LINK{http://naviteri.org/2014/12/ayngaru-tsinga-yoratu-may-i-present-the-four-winners/}{b}
\B{Tì\-'o'\-ìri pe\-u su\-nu ngar fra\-to?} What is your favorite way to have fun?~\LINK{http://naviteri.org/2010/09/getting-to-know-you-part-3/}{b}
(\ding{43}~\HL{'o'}{\B{'o'}})

% tì'ong
\hi
\HT{tI'ong}{\B{tì\cp{'ong}}} \I{[tI."PoN]} \emph{n.} blooming, unfolding.
\B{tì\-'ong lì'\-fya\-yä le\-Na'\-vi}, unfolding\slash blossoming of the Na'vi language.~\LINK{http://naviteri.org/2010/09/getting-to-know-you-part-1/}{b}
(\ding{43}~\HL{'ong}{\B{'ong}}) 

% tì'ongokx
\hi 
\HT{tI'ongokx}{\B{tì\cp{'o}ngokx}} \I{[tI."Po.Nok']} \emph{n.} birth. 
\B{Tì\-'o\-ngokx\-ìri nge\-yä 'itan\-ä sey\-kxel sì nit\-ram!} Congratulations on the birth of your son!~\LINK{http://naviteri.org/2018/03/100a-liu-amip-64-new-words-part-1/}{b} 
(\ding{43}~\HL{'ongokx}{\B{'o\-ngokx}})

% tì'ul
\hi 
\HT{tI'ul}{\B{tì\cp{'ul}}} \I{[tI."{}ul]} \emph{or} \I{[."{}Ul]} \emph{n.} increase. 
\B{Som\-wew\-ä tì\-'ul a ka 'Rr\-ta sä\-peng\-hrr lu aw\-nga\-ru nì\-wotx.} Global warming (lit.: the increase in temperature across the Earth) is a warning to us all.~\LINK{http://naviteri.org/2016/06/mrrvola-lifyavi-amip-forty-new-expressions/}{b} 
(\ding{43}~\HL{'ul}{\B{'ul}})

% tìaho
\hi 
\HT{tIaho}{\B{tìa\cp{ho}}} \I{[tI.a."ho]} \emph{n.} prayer (\emph{in general, abstract idea}). 
(\ding{43}~\HL{aho}{\B{aho}})

% tìeyktan
\hi
\HT{tIeyktan}{\B{tì\cp{eyk}tan}} \I{[tI."Ejk\textcorner.tan]} \emph{n.} leadership.
\B{Ke fkey\-tok tì\-eyk\-tan a\-tì\-flä\-nga' lu\-ke sno\-tipx.} Successful leadership does not exist without self-control.~\LINK{http://naviteri.org/2018/04/100a-liu-amip-64-new-words-part-2/}{b}
\B{Oeyk tsa\-tì\-snaytx\-ä lu apxa sä\-nui tì\-eyk\-tan\-ä.} That loss was caused by a great failure of leadership.~\LINK{http://naviteri.org/2013/04/wheres-the-bathroom-and-other-useful-things/}{b}
\B{Nge\-yä tì\-eyk\-tan\-ìri, tì\-slan\-ìri sì tsran\-ten fra\-to a tì\-'ey\-lan\-ìri a ka ay\-zì\-sìt nì\-wotx, sei\-yi oe ira\-yo nì\-txan.} Thank you so much for your leadership, your support, and most importantly your friendship throughout the years.~\LINK{http://naviteri.org/2012/03/spring-vocabulary-part-1/}{b}
(\ding{43}~\HL{eyktan}{\B{eyk\-tan}})

% tìfe'ul
\hi
\HT{tIfe'ul}{\B{tì\cp{fe}'ul}} \I{[tI."fE.Pul]} \emph{n.} worsening.
(\ding{214}~\HL{tItsan'ul}{\B{tì\-tsan\-'ul}}, \ding{43}~\HL{fe'ul}{\B{fe\-'ul}}, \HL{sAfe'ul}{\B{sä\-fe\-'ul}})

% tìfenge
\hi 
\HT{tIfenge}{\B{tì\cp{fe}nge}} \I{[tI."fE.NE]} \emph{n.} ridicule.
(\ding{43}~\HL{fenge}{\B{fe\-nge}})

% tìfnawe'
\hi 
\HT{tIfnawe'}{\B{tìfna\cp{we'}}} \I{[tI.fna."wEP]} \emph{n.} cowardice. 
\B{Nge\-yä tì\-fna\-we' ke\-muia so\-li fì\-soaia\-ru.} Your cowardice has dishonored this family.~\LINK{http://naviteri.org/2020/05/aawa-liu-amip-si-vurway-alor-a-few-new-words-and-a-beautiful-poem/}{b}
(\ding{43}~\HL{fnawe'}{\B{fna\-we'}})

% tìfkeytok
\hi
\HT{tIfkeytok}{\B{tì\cp{fkey}tok}} \I{[tI."fkEj.tok\textcorner]} \emph{n.} state, condition, situation.
\B{Tì\-fkey\-tok le\-fkrr le\-hrr\-ap lu nì\-txan.} The current situation is very dangerous.~\LINK{http://naviteri.org/2011/04/yafkeykiri-plltxe-frapo-everyone-talks-about-the-weather/}{b}
\B{Fì\-fne\-tì\-fkey\-tok\-ìl fngo' fu\-ta kem si\-vi fko pxi\-ye'\-rìn.} This kind of situation requires immediate action.~\LINK{http://naviteri.org/2012/01/mipa-zisit-ayliu-amip-new-words-for-the-new-year/}{b}
\B{Nge\-yä tsay\-lì\-'ul tì\-fkey\-tok\-it fe\-'ey\-ko\-lul nì\-'aw.} Those words of yours have only made the situation worse.~\LINK{http://naviteri.org/2013/03/tsanerul-feerul-getting-better-getting-worse/}{b}
\B{Sra\-ke tsun nga ri\-vun fì\-tì\-fkey\-tok\-mì a tì\-'i\-put?} Can you find the humor in this situation?~\LINK{http://naviteri.org/2012/01/mipa-zisit-ayliu-amip-new-words-for-the-new-year/}{b}
(\emph{a. conv. phrase}) \B{Nga\-ri\slash Nge\-yä tì\-fkey\-tok (lu) fya\-pe?} How are you doing?~\LINK{http://naviteri.org/2011/04/yafkeykiri-plltxe-frapo-everyone-talks-about-the-weather/comment-page-1/\#comment-646}{b}
\B{Tì\-fkey\-tok pe\-fya\slash fya\-pe?} `What's the existing situation?\slash How are things?'~\LINK{http://naviteri.org/2018/02/negative-questions-in-navi/\#comment-28201}{b}
(\ding{43}~\HL{fkeytok}{\B{fkey\-tok}}, \HL{-fkeyk}{\B{-fkeyk}}) \\
\textsc{Derivatives}: \ding{43}
\on{1}.~\HL{tIfkeytongay}{\B{tì\-fkey\-to\-ngay}}, reality.

% tìfkeytongay
\hi 
\HT{tIfkeytongay}{\B{tìfkeyto\cp{ngay}}} \I{[tI.fkEj.to."Naj]} \emph{n.} reality. 
\B{Ay\-unil nge\-yä lu lor, slä fì\-txe\-le\-ri lu tì\-fkey\-to\-ngay ke\-teng.} Your dreams are beautiful, but the reality of this situation is different.~\LINK{http://naviteri.org/2019/08/aawa-lifya-si-lifyavi-amip-a-few-new-words-and-expressions/}{b} 
(\ding{43}~\HL{tIfkeytok}{\B{tì\-fkey\-tok}}, \HL{ngay}{\B{ngay}}) \\
\textsc{Derivatives}: \ding{43}
\on{1}.~\HL{nIfkeytongay}{\B{nì\-fkey\-to\-ngay}}, actually, as a matter of fact.

% tìflä
\hi
\HT{tIflA}{\B{tì\cp{flä}}} \I{[tI."fl\ae]} \emph{n.} success (in general).
\B{Tì\-flä la\-tem ìlä seyng\-a ftxey fkol sä\-nu\-met li\-vek fu\-ke.} Success depends on whether or not one follows instructions.~\LINK{http://naviteri.org/2012/07/meetings-waterfalls-and-more/}{b}
\B{Tì\-flä\-ri lo\-lu po 'ä\-lek.} She was determined to succeed.~\LINK{http://naviteri.org/2021/09/aawa-ayliu-amip-a-few-new-words/}{b}
\B{Ay\-nge\-yä ay\-srr tì\-flä\-ta te\-ya li\-vu ul\-te si\-va\-lew nì\-'o' nì\-'aw.} May your days be full of success and proceed only in a fun way.~\LINK{http://naviteri.org/2013/12/mipa-zisit-lefpom-happy-new-year/}{b}

(\ding{214}~\HL{tInui}{\B{tì\-nui}}, \ding{43}~\HL{flA}{\B{flä}}, \HL{sAflA}{\B{sä\-flä}}) \\
\textsc{Derivatives}: \ding{43}
\on{1}.~\HL{tIflAnga'}{\B{tì\-flä\-nga'}}, successful.

% tìflänga'
\hi 
\HT{tIflAnga'}{\B{tì\cp{flä}nga'}} \I{[tI."fl\ae.NaP]} \emph{adj., nfp.} successful. 
\B{tì\-hawl a\-tì\-flä\-nga'}, a successful plan.~\LINK{http://naviteri.org/2017/09/zisikrr-amip-ayliu-amip-new-words-for-the-new-season/}{b}  
(\ding{43}~\HL{tIflA}{\B{tì\-flä}})

% tìflel
\hi 
\HT{tIflel}{\B{tì\cp{flel}}} \I{[tI."flEl]} \emph{n.} trickery (\emph{abstract concept}). 
(\ding{43}~\HL{flel}{\B{flel}}, \HL{sAflel}{\B{sä\-flel}})

% tìflrr
\hi
\HT{tIflrr}{\B{tì\cp{flrr}}} \I{[tI."fl\s{r}]} \emph{n.} gentleness, tenderness.
\B{Hu\-fwa me\-fo le\-ru mun\-txa\-tu txan\-krr, mi lu mun\-snar ho\-na tì\-flrr a na pum me\-yaw\-ne\-tu\-ä a\-mip nì\-wotx.} Although the two of them have been mates a long time, they still have all the adorable tenderness of new sweethearts.~\LINK{http://naviteri.org/2013/02/vospxi-ayol-posti-apup-short-post-for-a-short-month/}{b}
(\ding{43}~\HL{flrr}{\B{flrr}})

% tìfmetok
\hi
\HT{tIfmetok}{\B{tì\cp{fme}tok}} \I{[tI."fmE.tok\textcorner]} \emph{n.} test. 
\B{Tsyìl Ik\-ni\-ma\-yat ul\-te tsa\-heyl si ik\-ran\-hu a fì\-'u lu tì\-fme\-tok a ze\-ne fra\-ta\-ron\-yu a\-'e\-wan em\-zi\-va\-'u.} Scaling Iknimaya and bonding with a banshee is a test that every young hunter must pass.~\LINK{http://naviteri.org/2012/01/more-additions-to-the-lexicon/}{b}
\B{Sì\-fme\-tok\-it em\-zo\-la\-'u o\-hel.} I have passed the tests.~\LINK{http://wiki.learnnavi.org/index.php?title=Corpus\#Behind_the_Scenes}{wiki\slash a\textsuperscript{c}}
(\ding{43}~\HL{fmetok}{\B{fme\-tok}})

% tìfmi
\hi
\HT{tIfmi}{\B{tì\cp{fmi}}} \I{[tI."fmi]} \emph{n.} attempt.
\B{Fì\-tseo\-ri ke tsun kaw\-tu pi\-vä\-hem tì\-yo'\-ne; tsran\-ten tì\-pä\-hem\-ä tì\-fmi nì\-'aw.} In this art it's impossible to arrive at perfection; the only thing that matters is the attempt to arrive there.~\LINK{http://naviteri.org/2012/01/mipa-zisit-ayliu-amip-new-words-for-the-new-year/}{b}
(\ding{43}~\HL{fmi}{\B{fmi}}, \HL{sAfmi}{\B{sä\-fmi}})

% tìfmong
\hi 
\HT{tIfmong}{\B{tì\cp{fmong}}} \I{[tI."fmoN]} \emph{n.} theft. 
(\ding{43}~\HL{fmong}{\B{fmong}})

% tìfnu
\hi
\HT{tIfnu}{\B{tì\cp{fnu}}} \I{[tI."fnu]} \emph{n.} quiet, silence. 
\B{Pam\-tse\-ol ngop ay\-re\-nut | Mì ron\-sem\-ä tì\-fnu}, Music creates patterns | In the silence of the mind.~\LINK{http://www.pandorapedia.com/navi/music/ritual_music}{pp}
\B{Oel vo\-lin tì\-fnut.\slash Oe ätxä\-le so\-li tsnì li\-vu tì\-fnu.} I ask for silence.~\LINK{http://naviteri.org/2011/08/reported-speech-reported-questions/comment-page-1/\#comment-1060}{b}
(\ding{43}~\HL{fnu}{\B{fnu}}) \\
\textsc{Derivatives}: \ding{43}
\on{1}.~\HL{tIfnunga'}{\B{tì\-fnu\-nga'}}, quiet.

% tìfnunga'
\hi
\HT{tIfnunga'}{\B{tì\cp{fnu}nga'}} \I{[tI."fnu.NaP]} \emph{adj.} quiet.
\B{New oe ri\-vun a\-sim tì\-fnu\-nga'\-a tseng\-it a tsa\-ro tsun syi\-vor tsi\-vu\-rokx fte spä\-pi\-veng.} I want to find a quiet place nearby where I can chill out and rest to get my head back on straight.~\LINK{http://naviteri.org/2013/01/awvea-posti-zisita-amip-first-post-of-the-new-year/}{b}
(\ding{43}~\HL{tIfnu}{\B{tì\-fnu}}, \HL{-nga'}{\B{-nga'}})

% tìfpxamo
\hi 
\HT{tIfpxamo}{\B{tì\cp{fpxa}mo}} \I{[tI."fp'a.mo]} \emph{n.} horror. 
(\ding{43}~\HL{fpxamo}{\B{fpxa\-mo}})

% tìftang
\hi
\HT{tIftang}{\B{tì\cp{ftang}}} \I{[tI."ftaN]} \textsc{i}. \emph{n.} (\emph{a.}) stopping.
(\emph{b. linguistic}) glottal stop, the letter \B{'}, the sound \I{[P]}. 
\B{Ro sngä\-'i\-tseng tsa\-lì\-'u\-ä alu 'ey\-lan lu tì\-ftang.} At the beginning of the word \emph{'eylan} there's a glottal stop.~\LINK{http://naviteri.org/2012/07/fivospxiya-aylifyavi-amip-this-months-new-expressions-pt-2/}{b} \\
\textsc{ii}. \B{tì\cp{ftang} si} \I{[tI."ftaN si]} \emph{vin.} stop (s.o.\slash s.t.).   
\B{Ul\-te kaw\-tu ke tsun ay\-oer tì\-ftang si\-vi.} And no one can stop us.~\LINK{http://forum.learnnavi.org/language-updates/stop!/}{ln} 
\B{Oe po\-ru tì\-ftang so\-li.} I stopped him.~\LINK{http://forum.learnnavi.org/language-updates/stop!/}{ln}
(\ding{43}~\HL{ftang}{\B{ftang}}) \\
\textsc{Derivatives}: \ding{43}
\on{1}.~\HL{lukftang}{\B{luk\-ftang}}, constant, continual.

% tìftanglen
\hi 
\HT{tIftanglen}{\B{tìftang\cp{len}}} \I{[tI.ftaN."lEn]} \emph{n.} prevention. 
\B{Tì\-kan la\-'a\-yä a\-yll lu tì\-ftang\-len tì\-spxin\-ä.} The goal of social distancing is the prevention of disease.~\LINK{http://naviteri.org/2020/05/aawa-liu-amip-si-vurway-alor-a-few-new-words-and-a-beautiful-poem/}{b} 
(\ding{43}~\HL{ftanglen}{\B{ftang\-len}})

% tìftia
\hi
\HT{tIftia}{\B{tìfti\cp{a}}} \I{[tI.fti."a]} \emph{n.} study. 
(i.) \B{tì\-ftia ki\-fkey\-ä}, science [lit.: the study of the physical world]. 
(\ding{43}~\HL{ftia}{\B{ftia}}) \\
\textsc{Derivatives}: \ding{43}
\on{1}.~\HL{tIftiatu}{\B{tì\-fti\-a\-tu}}, researcher, scientist.

% tìftiatu
\hi
\HT{tIftiatu}{\B{tìfti\cp{a}tu}} \I{[tI.fti."a.tu]} \emph{n.} researcher, scientist. 
(i.) \B{tì\-fti\-a\-tu ki\-fkey\-ä}, sientist. 
\B{Saw\-tu\-te\-ri sì\-fti\-a\-tu ki\-fkey\-ä var fmi\-vi Ey\-we\-veng\-it tsli\-vam, slä kaw\-krr ke fla\-yä.} The scientists among the Sky People keep trying to understand Pandora, but they will never succeed.~\LINK{http://naviteri.org/2015/03/kap-si-ayunil-saylahe-threats-dreams-and-other-things/}{b}
(\ding{43}~\HL{tIftia}{\B{tì\-ftia}})

% tìftxavang
\hi
\HT{tIftxavang}{\B{tì\cp{ftxa}vang}} \I{[tI."ft'a.vaN]} \emph{n.} passion.
\B{Lu fra\-'u te\-ri tì\-ftxa\-vang.} Everything is about passion.~\LINK{http://forum.learnnavi.org/news-announcements/paul-frommer-at-grid-2010/}{ln}
(\ding{43}~\HL{ftxavang}{\B{ftxa\-vang}})

% tìftxey
\hi
\HT{tIftxey}{\B{tì\cp{ftxey}}} \I{[tI."ft'Ej]} \emph{n.} choice, option.
\B{Lu nge\-yä tì\-ftxey!} It's your choice!~\LINK{http://wiki.learnnavi.org/index.php?title=Canon\#Extracts_from_various_emails}{wiki} 
\B{Fì\-txe\-le\-ri tì\-nui ke lu tì\-ftxey.} In this matter, failure is not an option.~\LINK{http://naviteri.org/2013/04/wheres-the-bathroom-and-other-useful-things/}{b}
(\ding{43}~\HL{ftxey}{\B{ftxey}})

% tìftxulì'u
\hi
\HT{tIftxulI'u}{\B{tìftxu\cp{lì}'u}} \I{[tI.ft'u."lI.Pu]} \emph{n.} speech\hyp{}making, public speaking.   
\B{Oe\-ri lu tì\-ftxu\-lì\-'u ngä\-zìk nì\-txan. Wätx nì\-'aw.} Public speaking is very difficult for me. I'm hopelessly bad at it.~\LINK{http://naviteri.org/2012/06/spring-vocabulary-part-3/}{b}
(\ding{43}~\HL{ftxulI'u}{\B{ftxu\-lì\-'u}}, \HL{sAftxulI'u}{\B{sä\-ftxu\-lì\-'u}}) 

% tìfyawìntxu
\hi
\HT{tIfyawIntxu}{\B{tìfyawìn\cp{txu}}} \I{[tI.fja.wIn."t'u]} \emph{n.} guidance.
\B{Ul\-te o\-mum oel fu\-ta tì\-fya\-wìn\-txu\-ri oe\-yä pe\-rey ay\-nga nì\-wotx.} And I know you are all waiting for my guidance.~\LINK{http://forum.learnnavi.org/news-announcements/a-response-from-paul-frommer!/}{ln}
(\ding{43}~\HL{fyawIntxu}{\B{fya\-wìn\-txu}}) 

% tìfyeyn
\hi
\HT{tIfyeyn}{\B{tì\cp{fyeyn}}} \I{[tI."fjEjn]} \emph{n.} ripeness, maturity, full fruition.
(\ding{43}~\HL{fyeyn}{\B{fyeyn}}) 

% tìhawl
\hi
\HT{tIhawl}{\B{tì\cp{hawl}}} \I{[tI."hawl]} \emph{n.} preparations, plan. 
\B{Tì\-hawl le\-sngä\-'i lu tì\-kang\-kem\-vi skxawng\-ä, slä pum alu fì\-'u yo' nì\-'aw.} The original plan was the work of an idiot, but this one is just perfect.~\LINK{http://naviteri.org/2011/10/more-vocabulary-a-bit-of-grammar/}{b}
\B{Hu\-fwa nge\-yä tì\-hawl\-ìri ke lu kea kxey\-ey, tsal\-su\-ngay oe\-ru ke ha' nì\-tam.} Although there’s nothing wrong with your plan, it just doesn't suit me.~\LINK{http://naviteri.org/2012/10/mipa-vospxi-mipa-ayliu-new-words-for-the-new-month/}{b}
\B{Wo\-leyn Ìstawl yey\-fyat mì hll\-te fte oeyk\-ti\-vìng fra\-po\-ru tì\-hawl\-te\-ri sne\-yä.} Ìstaw drew a line on the ground to explain his plan to everyone.~\LINK{http://naviteri.org/2013/09/on-si-salewfya-shapes-and-directions/}{b}
(\ding{43}~\HL{hawl}{\B{hawl}})

% tìhawnu / si
\hi
\HT{tIhawnu}{\B{tì\cp{haw}nu}} \I{[tI."haw.nu]} \textsc{i}. \emph{n.} protection.  \\
\textsc{ii}. \B{tì\cp{haw}nu si} \I{[tI."haw.nu si]} \emph{vin.} protect.
\B{Oma\-ti\-ka\-ya\-ru tì\-haw\-nu si\-vi.} Protect the people.~\LINK{http://wiki.learnnavi.org/index.php?title=Na\%27vi_from_Avatar_Movie\#Eytukan.27s_Death}{wiki\slash a}
(\ding{43}~\HL{hawnu}{\B{haw\-nu}}) 

% tìhawnuwll
\hi
\HT{tIhawnuwll}{\B{tì\cp{haw}nuwll}} \I{[tI."haw.nu.w\s{l}]} \emph{n.} spartan, ``protection plant'', \emph{kuchenium polyphyllum}. 
(\ding{43}~\HL{tIhawnu}{\B{tì\-haw\-nu}}, \HL{'ewll}{\B{'ewll}})

% tìhan
\hi 
\HT{tIhan}{\B{tì\cp{han}}} \I{[tI."han]} \emph{n.} loss. 
\B{Maw tì\-han sa'\-nok\-ä, Txe\-wì a\-fpawng sar\-mä\-ngi zì\-sì\-to a\-pxay.} Sadly, after the loss of his mother, Txewì grieved for many years.~\LINK{http://naviteri.org/2018/04/100a-liu-amip-64-new-words-part-2/}{b} 
\B{Nge\-yä tì\-han\-ìri ke\-ftxo.} Sorry for your loss.~\LINK{http://naviteri.org/2019/08/aawa-lifya-si-lifyavi-amip-a-few-new-words-and-expressions/\#comment-30281}{b}
(\ding{43}~\HL{han}{\B{han}})

% tìhangham
\hi 
\HT{tIhangham}{\B{tì\cp{hang}ham}} \I{[tI."haN.ham]} \emph{n.} laughter. 
\B{Txa\-su\-nu oer fwa stawm nge\-yä tì\-hang\-ham\-it.} I love to hear your laughter.~\LINK{http://naviteri.org/2022/12/trr-anawm-polaheiem-the-great-day-has-arrived/}{b}  
(\ding{43}~\HL{hangham}{\B{hang\-ham}})

% tìhipx
\hi 
\HT{tIhipx}{\B{tì\cp{hipx}}} \I{[tI."hip']} \emph{n.} control. 
(\ding{43}~\HL{hipx}{\B{hipx}}) \\
\textsc{Derivatives}: \ding{43}
\on{1}.~\HL{snotipx}{\B{sno\-tipx}}, self\hyp{}control. 

% tìhìpey
\hi
\HT{tIhIpey}{\B{tì\cp{hì}pey}} \I{[tI."hI.pEj]} \emph{n.} hesitation.
\B{Tì\-hì\-pey tsun krro krro le\-sar li\-vu.} Hesitation can sometimes be useful.~\LINK{http://naviteri.org/2015/11/vomuna-liu-amip-ten-new-words/}{b}
\B{Sar tsa\-lì\-'ut a fì\-'u\-ri lu oe\-ru tì\-hì\-pey nì\-'it.} I'm a bit hesitant about using that word.~\LINK{http://naviteri.org/2015/11/vomuna-liu-amip-ten-new-words/}{b}
(\ding{43}~\HL{hIpey}{\B{hì\-pey}})

% tìhona
\hi
\HT{tIhona}{\B{tì\cp{ho}na}} \I{[tI."ho.na]} \emph{n.} cuteness, adorableness. 
\B{Pe\-yä 'itan\-ìri lu ho\-na nì\-txan a fì\-'u law lu fra\-por. Slä tì\-ho\-na nì\-'aw ke tam.} It's clear to everyone that his son is very cute. But cuteness alone isn't enough.~\LINK{http://naviteri.org/2012/01/more-additions-to-the-lexicon/}{b}
(\ding{43}~\HL{hona}{\B{ho\-na}})

% tìhuslew
\hi 
\HT{tIhuslew}{\B{tì\cp{hu}slew}} \I{[tI."hu.slEw]} \emph{n.} accompaniment. 
(\ding{43}~\HL{huslew}{\B{hu\-slew}}) 

% tìk
\hi 
\HT{tIk}{\B{tìk}} \I{[tIk\textcorner]} \textsc{i}. \emph{adv.} immediately, without delay. 
\B{Tsa\-kem si tìk!} Do it immediately!~\LINK{http://naviteri.org/2021/12/zolau-niprrte-ma-3746-welcome-2022/}{b} \\
\textsc{ii}. \emph{conj., indicating that a second action immediately follows a first}.
\B{Fì\-ioang ke tsun sli\-ve\-le; nem\-fa pay zup tìk spa\-kat.} This animal cannot swim; if it falls into the water, it immediately drowns.~\LINK{http://naviteri.org/2021/12/zolau-niprrte-ma-3746-welcome-2022/}{b}
(\emph{a. idiom}) \B{Tse\-'a tìk yaw\-ne.} `Love at first sight.'~\LINK{http://naviteri.org/2021/12/zolau-niprrte-ma-3746-welcome-2022/}{b}
(\ding{43}~\HL{pxiye'rIn}{\B{pxi\-ye'\-rìn}})

% tìk'ìn
\hi
\HT{tIk'In}{\B{tìk\cp{'ìn}}} \I{[tIk\textcorner."PIn]} \emph{n.} free time, absence of being busy.  
\B{Tìk\-'ìn\-ìri kem\-pe si nga?} What do you do in your free time?~\LINK{http://naviteri.org/2010/09/getting-to-know-you-part-3/}{b}
(\ding{43}~\HL{'In}{\B{'ìn}}) 

% tìkoak
\hi 
\HT{tIkoak}{\B{tì\cp{ko}ak}} \I{[tI."ko.ak\textcorner]} \emph{n.} old age. 
(\emph{a. proverb}) \B{Nì\-win 'ìp tì\-'e\-wan, nì\-win pä\-hem tì\-ko\-ak.} Youth vanishes quickly, old age arrives quickly.~\LINK{https://naviteri.org/2025/04/mevosinga-liu-amip-twenty-new-words/}{b} 
(\ding{43}~\HL{koak}{\B{ko\-ak}}; \ding{214}~\HL{tI'ewan}{\B{tì\-'e\-wan}}) 

% tìktok
\hi 
\HT{tIktok}{\B{tìk\cp{tok}}} \I{[tIk\textcorner."tok\textcorner]} \emph{n.} absence. 
\B{Tìk\-tok\-it pe\-yä fkol tso\-le\-ri.} His absence was noted.~\LINK{https://naviteri.org/2024/06/ayhapxi-tokxa-si-ayliu-alahe-parts-of-the-body-and-more/}{b}
(\emph{a. proverb}) \B{Tì\-tok slu tìk\-tok fa pam\-tsyìp a\-'aw.} Presence becomes absence by means of one little sound. [i.e. `Small things can make a big difference.']~\LINK{https://naviteri.org/2024/06/ayhapxi-tokxa-si-ayliu-alahe-parts-of-the-body-and-more/}{b}
(\ding{43}~\HL{ke}{\B{ke}}, \HL{tok}{\B{tok}}, \ding{214}~\HL{tItok}{\B{tì\-tok}}) 

% tìktseri
\hi
\HT{tIktseri}{\B{tìk\cp{tse}ri}} \I{[tIk\textcorner."\t{ts}E.Ri]} \emph{n.} unawareness, lack of notice.
(\emph{a. proverb}) \B{Tìk\-tse\-ri lu tì\-meyp.} Lack of awareness is (a form of) weakness.~\LINK{http://naviteri.org/2015/03/kap-si-ayunil-saylahe-threats-dreams-and-other-things/}{b}
(\ding{43}~\HL{tseri}{\B{tse\-ri}})

% tìktstun
\hi 
\HT{tIktstun}{\B{tìk\cp{tstun}}} \I{[tIk\textcorner."\t{ts}tun]} \emph{or} \I{[."\t{ts}tUn]} \emph{n.} unkindness, emotional coldness.
(\ding{43}~\HL{ketstun}{\B{ke\-tstun}}) \\
\textsc{Derivatives}: \ding{43}
\on{1}.~\HL{tIktstunga'}{\B{tìk\-tstu\-nga'}}, unkind, emotional coldness.

% tìktstunga'
\hi 
\HT{tIktstunga'}{\B{tìk\cp{tstu}nga'}} \I{[tIk\textcorner."\t{ts}tu.NaP]} \emph{adj., nfp.} unkind, emotional coldness.
(\ding{43}~\HL{tIktstun}{\B{tìk\-tstun}}, \HL{-nga'}{\B{-nga'}}; \HL{ketstun}{\B{ke\-tstun}}; \ding{214}~\HL{tItstunwinga'}{\B{tì\-tstun\-wi\-nga'}}) 

% tìkahena
\hi 
\HT{tIkahena}{\B{tìka\cp{he}na}} \I{[tI.ka."hE.na]} \emph{n.} transportation (\emph{abstract concept}).  
(\ding{43}~\HL{kahena}{\B{ka\-he\-na}}, \HL{sAkahena}{\B{sä\-ka\-he\-na}})

% tìkakmokri
\hi 
\HT{tIkakmokri}{\B{tìkak\cp{mok}ri}} \I{[tI.kak\textcorner."mok\textcorner.ri]} \emph{n.} muteness. 
(\ding{43}~\HL{kakmokri}{\B{kak\-mok\-ri}})

% tìkakpam
\hi
\HT{tIkakpam}{\B{tìkak\cp{pam}}} \I{[tI.kak\textcorner."pam]} \emph{n.} deafness. 
\B{Tì\-ta\-ron\-ìri lä\-ngu tì\-kak\-pam fe\-kum. Kin fkol fra\-inan\-fyat.} I'm sorry to say that deafness is a disadvantage for hunting. You need all your senses.~\LINK{http://naviteri.org/2015/08/aylifyavi-lereyfya-2-cultural-terms-2-and-more/}{b} 
\B{Po\-ri tì\-kak\-rel tì\-kak\-pam\-sì kum tsam\-ä lu.} His blindness and deafness are a result of the war.~\LINK{http://naviteri.org/2014/05/mipa-ayliu-mipa-aysafpil-new-words-new-ideas/}{b}
(\ding{43}~\HL{kakpam}{\B{kak\-pam}})

% tìkakrel
\hi
\HT{tIkakrel}{\B{tìkak\cp{rel}}} \I{[tI.kak\textcorner."REl]} \emph{n.} blindness. 
\B{Po\-ri tì\-kak\-rel tì\-kak\-pam\-sì kum tsam\-ä lu.} His blindness and deafness are a result of the war.~\LINK{http://naviteri.org/2014/05/mipa-ayliu-mipa-aysafpil-new-words-new-ideas/}{b}
(\ding{43}~\HL{kakrel}{\B{kak\-rel}})

% tìkakzir
\hi 
\HT{tIkakzir}{\B{tìkak\cp{zir}}} \I{[tI.kak\textcorner."zir]} \emph{n.} numbness (\emph{physical and emotional}). 
(\ding{43}~\HL{kakzir}{\B{kak\-zir}})

% tìkan
\hi
\HT{tIkan}{\B{tì\cp{kan}}} \I{[tI."kan]} \emph{n.} aim, goal, purpose, target. 
\B{Pì\-lok\-ä fì\-ha\-pxì\-yä tì\-kan lu law.} The aim of this part of the blog is clear.~\LINK{http://naviteri.org/2010/06/first-post/}{b}
\B{Tì\-kan taw\-na\-tep!} (\emph{mil.}) Target lost!~\LINK{http://naviteri.org/2011/05/some-miscellaneous-vocabulary/}{b\slash g}
\B{Lu pì\-lok\-ur pxe\-sì\-kan sì pxe\-fne\-'upxa\-re.} The blog has three aims and three kinds of messages.~\LINK{http://naviteri.org/2010/06/first-post/}{b}
\B{Fu\-ria tsun fko nì\-'i'\-a tì\-kan\-it tsi\-ve\-'a, 'e\-fu oe nit\-ram nì\-txan nì\-'aw.} That one can finally see the aim, I feel very happy.~\LINK{http://naviteri.org/2013/08/fmawn-a-fkol-kilmulat-this-just-in/comment-page-1/\#comment-2479}{b}
\B{Fì\-tì\-kan\-ìri ropx ke ha'.} For this purpose, a hole going right through the tree trunk isn't suitable.~\LINK{http://naviteri.org/2013/04/wheres-the-bathroom-and-other-useful-things/}{b}
(\emph{a. proverb}) \B{Sä\-fpìl a\-steng, tì\-kan a\-teng.} Great minds think alike [lit.: similar thought, same aim].~\LINK{http://naviteri.org/2011/09/miscellaneous-vocabulary/}{b} 
(\ding{43}~\HL{kan}{\B{kan}}) \\
\textsc{Derivatives}: \ding{43}
\on{1}.~\HL{tIkangkem}{\B{tì\-kang\-kem}}, work.
\on{2}.~\HL{nItkan}{\B{nìt\-kan}}, intentionally. 
\on{3}.~\HL{tIkankxan}{\B{tì\-kan\-kxan}}, barrier to one's goals, source of frustration.
\on{4}.~\HL{mu'nitkan}{\B{mu'\-nit\-kan}}, effective.

% tìkankxan
\hi 
\HT{tIkankxan}{\B{tìkan\cp{kxan}}} \I{[tI.kan."k'an]} \emph{n.} barrier to one's goals, source of frustration. 
\B{Oe\-ri mì\-ftxe\-le, fwa pol oe\-ti ke slo\-lan lu tì\-kan\-kxan apxa.} In this regard, his not supporting me was a large barrier to (achieving) my goal.~\LINK{http://naviteri.org/2020/05/aawa-liu-amip-si-vurway-alor-a-few-new-words-and-a-beautiful-poem/}{b}
(\ding{43}~\HL{tIkan}{\B{tì\-kan}}, \HL{ekxan}{\B{e\-kxan}})

% tìkankxanga'
\hi 
\HT{tIkankxanga'}{\B{tìkan\cp{kxa}nga'}} \I{[tI.kan."k'a.NaP]} \emph{adj.} frustrating. 
\B{Fì\-tì\-fkey\-tok\-ìl a\-tì\-kan\-kxa\-nga' 'ey\-ke\-re\-fu oet lek\-ye\-'ung.} This frustrating situation is making me crazy.~\LINK{http://naviteri.org/2020/05/aawa-liu-amip-si-vurway-alor-a-few-new-words-and-a-beautiful-poem/}{b}
(\ding{43}~\HL{tIkankxan}{\B{tìkankxan}}, \HL{-nga'}{\B{-nga}}; \HL{lekxan}{\B{le\-kxan}})

% tìkanu
\hi
\HT{tIkanu}{\B{tì\cp{ka}nu}} \I{[tI."ka.nu]} \emph{n.} intelligence.
\B{Nge\-yä ke ftang a tì\-kang\-kem\-ìri sì tì\-ka\-nu\-ri sì fì\-lì'\-fya\-yä tì\-yawn\-ìri sei\-yi oe ira\-yo nì\-txan.} I thank you very much for your never\hyp{}ending work, intelligence and love for the language.~\LINK{http://naviteri.org/2012/07/meetings-waterfalls-and-more/comment-page-1/\#comment-1547}{b}
(\ding{43}~\HL{kanu}{\B{ka\-nu}})

% tìkangkem / si
\hi
\HT{tIkangkem}{\B{tì\cp{kang}kem}} \I{[tI."kaN.kEm]} \textsc{i}. \emph{n.} work (\emph{coll. shortened to} \ding{43}~\HL{kangkem}{\B{kang\-kem}}). 
\B{Fay\-sul\-fä\-tu\-ä tì\-kang\-kem ohe\-ru meuia lu\-yu nì\-ngay.} The work of these masters truly honors me.~\LINK{http://naviteri.org/2010/06/first-post/}{b}
\B{Pe\-yä tì\-kang\-kem\-ìl mi vey\-krr\-ei\-yìn pot nì\-wotx.} His work is still completely overwhelming him (and I'm glad).~\LINK{http://naviteri.org/2010/09/getting-to-know-you-part-3/}{b}
\B{Tì\-kang\-kem\-ìri var\-mrr\-ìn oe nì\-wotx.} I was completely swamped (overwhelmed) at work.~\LINK{http://naviteri.org/2010/09/getting-to-know-you-part-3/}{b}
\B{Poel tì\-kang\-kem\-it sngey\-ko\-lä\-'i.} She began the work.~\LINK{http://naviteri.org/2010/06/first-post/}{b}  \\
\textsc{ii}. \B{tì\cp{kang}kem si} \I{[tI."kaN.kEm si]} \emph{vin.} work.
\B{Slä nì\-aw\-no\-mum, ze\-ne oe 'aw\-si\-teng tì\-kang\-kem si\-vi fo\-hu.} But as you know, I must work together with them.~\LINK{http://forum.learnnavi.org/news-announcements/a-response-from-paul-frommer!/}{ln}
\B{Tì\-kang\-kem si val!} Work hard!~\LINK{http://naviteri.org/2022/12/trr-anawm-polaheiem-the-great-day-has-arrived/}{b}
(\emph{with the topic}) work at\slash on s.t. \B{Fì\-tì\-kang\-kem\-vi\-ri le\-tsran\-ten ke new oe hu Ra\-lu tì\-kang\-kem si\-vi.} I don't want to work with Ralu on this important project.~\LINK{http://naviteri.org/2012/01/more-additions-to-the-lexicon/}{b}
(\ding{43}~\HL{tIkan}{\B{tì\-kan}}, \HL{kem}{\B{kem}}) \\
\textsc{Derivatives}: \ding{43}
\on{1}.~\HL{tIkangkemvi}{\B{tì\-kang\-kem\-vi}}, project.

% tìkangkemvi
\hi
\HT{tIkangkemvi}{\B{tì\cp{kang}kemvi}} \I{[tI."kaN.kEm.vi]} \emph{n.} project, piece of work.
\B{Hu\-fwa ro\-lun oel 'a\-'aw\-a kxey\-ey\-ti, fì\-tì\-kang\-kem\-vi lu txan\-tsan ka wotx.} Although I found a few errors, this piece of work is generally excellent.~\LINK{http://naviteri.org/2012/07/fivospxiya-aylifyavi-amip-this-months-new-expressions-pt-2/}{b}
\B{Fwa oel fì\-tì\-kang\-kem\-vit sngey\-ki\-vä\-'i krr\-no\-lekx nì\-txan.} Me starting this project took a long time.~\LINK{http://naviteri.org/2010/06/first-post/}{b}
\B{Nì\-'i'\-a po tsa\-tì\-kang\-kem\-vir ha\-sey so\-li.} She finally completed the project.~\LINK{http://naviteri.org/2011/04/\%E2\%80\%99a\%E2\%80\%99awa-li\%E2\%80\%99fyavi-amip\%E2\%80\%94a-few-new-expressions/}{b} 
\B{Fì\-tì\-kang\-kem\-vi\-ri le\-tsran\-ten ke new oe hu Ra\-lu tì\-kang\-kem si\-vi.} I don't want to work with Ralu on this important project.~\LINK{http://naviteri.org/2012/01/more-additions-to-the-lexicon/}{b}
\B{Fì\-tì\-kang\-kem\-vi\-ri oel kin 'aw\-a trr\-ti nì\-'aw.} For this project I only need one day.~\LINK{http://naviteri.org/2014/09/stxeli-alor-the-text/}{b}
(\ding{43}~\HL{tIkangkem}{\B{tì\-kang\-kem}})

% tìkara
\hi 
\HT{tIkara}{\B{tìka\cp{ra}}} \I{[tI.ka."Ra]} \emph{n.} resistance. 
(\emph{a. idiom}) \B{Tì\-ka\-ra lu ä\-txä\-le pa\-lu\-kan\-ur.} Resistance is futile [lit.: `reistance is a request to the thanator'].~\LINK{http://naviteri.org/2023/09/aawa-ayliu-si-aylifyavi-amip-a-few-new-words-and-expressions/}{b}
(\ding{43}~\HL{kara}{\B{ka\-ra}})

% tìkatìng
\hi 
\HT{tIkatIng}{\B{tì\cp{ka}tìng}} \I{[tI."ka.tIN]} \emph{n.} distribution.  
(\ding{43}~\HL{katIng}{\B{ka\-tìng}})

% tìkawng
\hi
\HT{tIkawng}{\B{tì\cp{kawng}}} \I{[tI."kawN]} \emph{n.} evil. 
\B{Eo ay\-oeng lu txan\-a tì\-kawng.} A great evil is upon us.~\LINK{http://wiki.learnnavi.org/index.php?title=Na\%27vi_from_Avatar_Movie\#Jake_back_at_hometree}{wiki\slash a}
(\ding{43}~\HL{kawng}{\B{kawng}})

% tìkeftxo
\hi
\HT{tIkeftxo}{\B{tìke\cp{ftxo}}} \I{[tI.kE."ft'o]} \emph{n.} sadness.
(\ding{43}~\HL{keftxo}{\B{ke\-ftxo}}) \\
\textsc{Derivatives}: \ding{43}
\on{1}.~\HL{tIkeftxonga'}{\B{tì\-ke\-ftxo\-nga'}}, sad (\emph{nfp.}).

% tìkeftxonga'
\hi
\HT{tIkeftxonga'}{\B{tìke\cp{ftxo}nga'}} \I{[tI.kE."ft'o.NaP]} \emph{adj., nfp.} sad.
\B{Txo tì\-ke\-ftxo\-nga'\-a fmawn\-it ke stil\-vawm ay\-ngal, ze\-ne oe pi\-veng san Unil\-tì\-ran\-tokx\-ä pam\-tseo\-ngop\-yu a\-yaw\-ne to\-ler\-kä\-ngup.} If you haven't heard the sad news, I have to inform that the dear music creator of \emph{Avatar} has died.~\LINK{http://naviteri.org/2015/06/na-ayskxe-mi-telan-sad-news/}{b}
(\ding{43}~\HL{tIkeftxo}{\B{tì\-ke\-ftxo}})

% tìkelpxìmrun
\hi 
\HT{tIkelpxImrun}{\B{tìkel\cp{pxìm}run}} \I{[tI.kEl."p'Im.Run]} \emph{or} \I{[.RUn]} \emph{n.} rarity, rareness. 
\B{Tì\-kel\-pxìm\-run\-it na\-fì\-'u\-a mok\-ri\-yä ke tsun fko wä\-ti\-ve.} The rarity of such a voice can't be disputed.~\LINK{https://naviteri.org/2026/04/tsivola-liu-amip-thirty-two-new-words/}{b}  
(\ding{43}~\HL{kelpxImrun}{\B{kel\-pxìm\-run}}; \HL{sAkelpxImrun}{\B{sä\-kel\-pxìm\-run}}) 

% tìkelu
\hi
\HT{tIkelu}{\B{tì\cp{ke}lu}} \I{[tI."kE.lu]} \emph{n.} lack. 
\B{Tì\-ke\-lu tì\-eyk\-tan\-ä a\-sìl\-tsan lä\-ngu mì olo' aw\-nge\-yä tì\-ngä\-zìk.} Unfortunately, the lack of good leadership in our clan is a problem.~\LINK{http://naviteri.org/2015/04/some-new-words-for-may-day/}{b} \\
\textsc{Derivatives}: \ding{43}
\on{1}.~\HL{tompakel}{\B{tom\-pa\-kel}}, drought.
\on{2}.~\HL{syuvekel}{\B{syu\-ve\-kel}}, famine.

% tìkemwiä
\hi 
\HT{tIkemwiA}{\B{tìkem\cp{wi}ä}} \I{[tI.kEm."wi.\ae]} \emph{n.} unfairness, injustice. 
(\ding{43}~\HL{kemwiA}{\B{kem\-wi\-ä}})

% tìkenong
\hi
\HT{tIkenong}{\B{tì\cp{ke}nong}} \I{[tI."kE.noN]} \emph{n.} example.   
\B{Fì\-po\-ti oel tspì\-yang, fte tì\-ke\-nong li\-ye\-vu ay\-la\-ru.} I will kill this one as a lesson to the others.~\LINK{http://forum.learnnavi.org/language-updates/capturing-avatar/}{ln\slash a\textsuperscript{c}}
\B{Su\-nu oer fay\-sì\-ke\-nong nge\-yä.} I like these examples of yours.~\LINK{http://naviteri.org/2014/07/tsawlultxamaw-a-ayliu-post-meetup-words/comment-page-1/\#comment-7697}{b}
\B{Pol ngop fra\-krr sì\-ke\-nong\-it a hì\-no lu nì\-hawng.} He always creates excessively detailed examples.~\LINK{http://naviteri.org/2010/09/getting-to-know-you-part-2/}{b}
\B{Tsa\-tì\-ke\-nong\-ìri oel kxey\-eyt zo\-sley\-ko\-lu.} As for that example, I got the mistake fixed.~\LINK{http://naviteri.org/2013/01/awvea-posti-zisita-amip-first-post-of-the-new-year/comment-page-1/\#comment-1903}{b}
(\ding{43}~\HL{kenong}{\B{ke\-nong}}) \\
\textsc{Derivatives}: \ding{43}
\on{1}.~\HL{natkenong}{\B{nat\-ke\-nong}}, for example.

% tìketeng
\hi
\HT{tIketeng}{\B{tì\cp{ke}teng}} \I{[tI."kE.tEN]} \emph{n.} difference.
\B{Tì\-ke\-teng lu 'u\-pe?} What's the difference?~\LINK{http://naviteri.org/2014/08/stxeli-alor-a-beautiful-gift/comment-page-1/\#comment-9347}{b}
\B{Ral\-ìri ke 'efu oel kea tì\-ke\-teng\-it.} I don't feel any difference concerning the meaning.~\LINK{http://naviteri.org/2013/01/awvea-posti-zisita-amip-first-post-of-the-new-year/comment-page-1/\#comment-1903}{b}
(\ding{43}~\HL{keteng}{\B{ke\-teng}})

% tìkezin (si)
\hi 
\HT{tIkezin}{\B{tì\cp{ke}zin}} \I{[tI."kE.zin]} \textsc{i}. \emph{n.} (\emph{a.}) s.t. in an untangled state, ``untanglement''. (\emph{b.}) solution. 
\B{Tsa\-tì\-ngä\-zìk\-ìri tì\-ke\-zin lu fyin.} The solution to that problem is simple.~\LINK{http://naviteri.org/2021/04/mipa-ayliu-mipa-sioeykting-new-words-new-explanations/}{b}  
(\ding{43}~\HL{kezin}{\B{ke\-zin}}) \\
\textsc{ii}. \B{tì\cp{ke}zin si} \I{[tI."kE.zin si]} \emph{vin.} (\emph{a.}) untangle. (\emph{b.}) solve. 
\B{Sra\-ke tsun nga fì\-ing\-yen\-tsim\-ur tì\-ke\-zin si\-vi?} Can you solve this riddle?~\LINK{http://naviteri.org/2021/04/mipa-ayliu-mipa-sioeykting-new-words-new-explanations/}{b} 
\B{Ze\-ne aw\-nga tì\-ngä\-zìk\-ur sno\-ni\-vi\-yä tì\-ke\-zin si\-vi nì\-win.} We need to solve the hammock problem quickly.~\LINK{http://naviteri.org/2021/04/mipa-ayliu-mipa-sioeykting-new-words-new-explanations/\#comment-33483}{b}
(\ding{43}~\HL{kezin}{\B{ke\-zin}})

% tìkin
\hi
\HT{tIkin}{\B{tì\cp{kin}}} \I{[tI."kin]} \emph{n.} need (\emph{can refer either to the general state or concept or to a specific instance}). 
\B{'On\-ìri tsko\-ä lu tì\-kin sä\-la\-tem\-ä a\-hì\-'i.} The form of the bow requires a small change.~\LINK{http://naviteri.org/2015/11/vomuna-liu-amip-ten-new-words/}{b}
(\emph{a. with the dative to express}) (have the) need to (do s.t.) \B{Po\-ri aw\-nga\-ru lu tì\-kin a nu\-me nì\-'ul.} We need to learn more about him.~\LINK{http://wiki.learnnavi.org/index.php?title=Na\%27vi_from_Avatar_Movie\#Scene_in_Hometree}{wiki\slash a}  
(\emph{b. conv. phrase for responding to thanks}) \B{Kea tì\-kin.} ``You're welcome.'' [lit.: (there is) no need]
(\ding{43}~\HL{kin}{\B{kin}})

% tìkxal
\hi 
\HT{tIkxal}{\B{tì\cp{kxal}}} \I{[tI."k'al]} \emph{n.} strictness, discipline. 
(\emph{a. proverb}) \B{Tì\-kxal\-ìl tì\-kxäl\-it zey\-ko.} Discipline cures apathy.~\LINK{https://naviteri.org/2025/04/mevosinga-liu-amip-twenty-new-words/}{b}  
(\ding{43}~\HL{kxal}{\B{kxal}}) 

% tìkxäl 
\hi 
\HT{tIkxAl}{\B{tì\cp{kxäl}}} \I{[tI."k'\ae{}l]} \emph{n.} apathy. 
\B{Tì\-kxäl lu pe\-u? Lu tskxe\-pay a mì te'\-lan sì kxitx a mì el\-tu.} What is apathy? It's ice in the heart and a death in the brain.~\LINK{https://naviteri.org/2025/04/mevosinga-liu-amip-twenty-new-words/}{b} 
(\emph{a. proverb}) \B{Tì\-kxal\-ìl tì\-kxäl\-it zey\-ko.} Discipline cures apathy.~\LINK{https://naviteri.org/2025/04/mevosinga-liu-amip-twenty-new-words/}{b}  
(\ding{43}~\HL{kxAl}{\B{kxäl}}) 

% tìkxey / si
\hi
\HT{tIkxey}{\B{tì\cp{kxey}}} \I{[tI."k'Ej]} \textsc{i}. \emph{n.} incorrectness, mistakeness.
\B{Tì\-kxey zo\-slo\-lu.} The incorrectness has been fixed.~\LINK{http://naviteri.org/2013/12/mipa-zisit-lefpom-happy-new-year/comment-page-1/\#comment-2585}{b}
(\emph{a. idiom}) \B{Nga\-ru tì\-kxey.} You're wrong.
\B{Nì\-ngay lu ngar tì\-yawr, lu ngar tì\-kxey.} Truly, you are both right and wrong.~\LINK{http://naviteri.org/2010/07/vocabulary-update/comment-page-1/\#comment-95}{b} 
\B{Po ke tsun pi\-vll\-ngay san oe\-ru tì\-kxey.} He can't admit he's wrong.~\LINK{http://naviteri.org/2011/07/txantsana-ultxa-mi-siatll-great-meeting-in-seattle/}{b} \\
\textsc{ii}. \B{tì\cp{kxey} si} \I{[tI."k'Ej si]} \emph{vin.} mess up, foul, do wrong, make a mistake.
\B{Krr\-o krr\-o fì\-tse\-nge oe tì\-kxey sa\-yi.} From time to time I will mess up here.~\LINK{http://naviteri.org/2010/06/first-post/}{b}
\B{Ngay\-txoa, krro tì\-kxey si keng kar\-yu.} Sorry, sometimes even the teacher messes up.~\LINK{http://naviteri.org/2012/02/trr-asawnung-lefpom-happy-leap-day/}{b}
(\ding{214}~\HL{tIyawr}{\B{tì\-yawr}}, \ding{43}~\HL{kxeyey}{\B{kxey\-ey}})

% tìkxìm
\hi
\HT{tIkxIm}{\B{tì\cp{kxìm}}} \I{[tI."k'Im]} \textsc{i}. \emph{n.} commanding, ordering, assigning tasks.
\B{Sìl\-tsan\-a eyk\-tan ze\-ne fni\-van tì\-kxìm\-ti.} A good leader must be skilled at assigning tasks.~\LINK{http://naviteri.org/2014/10/tson-si-fpomron-obligation-and-mental-health/}{b} \\
\textsc{ii}. \B{tì\cp{kxìm} si} \I{[tI."k'Im si]} \emph{vin.} be above s.o. in a hierarchy, be s.o.'s superior.~\LINK{http://naviteri.org/2014/10/tson-si-fpomron-obligation-and-mental-health/}{b}
\B{Po tì\-kxìm si oer.} He is above me (in some relevant hierarchy).\slash I am under him.\slash He has authority over me.\slash He is my boss.~\LINK{http://naviteri.org/2014/10/tson-si-fpomron-obligation-and-mental-health/}{b}
\B{Pe\-su tì\-kxìm si ngar?} Who's your boss?~\LINK{http://naviteri.org/2014/10/tson-si-fpomron-obligation-and-mental-health/}{b} 
(\ding{43}~\HL{kxIm}{\B{kxìm}})

% tìkxuke
\hi
\HT{tIkxuke}{\B{tì\cp{kxu}ke}} \I{[tI."ku.kE]} \emph{n.} safety.
\B{Fmal tì\-kxu\-ket nge\-yä sì Mo\-'a\-ra\-yä.} Keep yourself and Mo\-'ara safe.~\LINK{http://naviteri.org/2017/06/ayliu-a-ta-eywaeveng-kifkey-uniltirantokxa-words-from-disneys-pandora-the-world-of-avatar/}{b}
(\ding{43}~\HL{kxuke}{\B{kxu\-ke}})

% tìla'um
\hi 
\HT{tIla'um}{\B{tì\cp{la}'um}} \I{[tI."la.Pum]} \emph{or} \I{[.PUm]} (\textsc{rn}: \B{tì\cp{la}ùm} \I{[tI."la.Um]}) \emph{n.} pretence. 
\B{Fu\-ria ke tsun tì\-kang\-kem si\-vi, pe\-yä sä\-spxin lu tì\-la\-'um nì\-'aw.} As for not being able to work, his illness is only a pretence.~\LINK{http://naviteri.org/2020/02/some-words-for-leap-year-day/}{b} 
(\ding{43}~\HL{la'um}{\B{la\-'um}})

% tìlätek
\hi 
\HT{tIlAtek}{\B{tì\cp{lä}tek}} \I{[tI."l\ae.tEk\textcorner]} \emph{n.} recognition. 
(\ding{43}~\HL{lAtek}{\B{lä\-tek}}) 

% tìlayl
\hi 
\HT{tIlayl}{\B{tì\cp{layl}}} \I{[tI."lajl]} \emph{n.} innocence.  
(\ding{43}~\HL{layl}{\B{layl}}; \ding{214}~\HL{tItokat}{\B{tì\-to\-kat}}) 

% tìlayro
\hi 
\HT{tIlayro}{\B{tì\cp{lay}ro}} \I{[tI."laj.ro]} \emph{n.} freedom. 
\B{Ay\-su\-tel nì\-wotx new tì\-lay\-rot.} All people want freedom.~\LINK{http://naviteri.org/2022/11/krr-ao-an-exciting-time/}{b} 
(\ding{43}~\HL{layro}{\B{lay\-ro}}) 

% tìlalim
\hi 
\HT{tIlalim}{\B{Tìla\cp{lim}}} \I{[tI.la."lim]} \emph{n., clan name} the Windtrader clan.~\LINK{https://naviteri.org/2025/11/kaltxi-nimun-hello-again/}{b} 
\B{Way Tì\-la\-lim\-ä}, The Windtrader Song.~\LINK{https://naviteri.org/2026/02/way-tilalima-the-windtrader-song/}{b}

% tìlam
\hi
\HT{tIlam}{\B{tì\cp{lam}}} \I{[tI."lam]} \emph{n.} appearance. 
\B{Ul\-te kxawm ke lu tì\-fkey\-tok nì\-ftxan fe' na tì\-lam.} And maybe the situation is not as bad as it seems.~\LINK{http://naviteri.org/2016/12/tengkrr-perahem-zisit-amip-as-the-new-year-arrives/}{b}
(\ding{43}~\HL{lam}{\B{lam}}) \\
\textsc{Derivatives}: \ding{43}
\on{1}.~\HL{tatlam}{\B{tat\-lam}}, apparently.

% tìlang
\hi
\HT{tIlang}{\B{tì\cp{lang}}} \I{[tI."laN]} \emph{n.} exploration (general sense). 
\B{Sra\-ne, su\-nu Saw\-tu\-te\-ru tì\-lang, slä ke omum fol teyng\-ta kem\-pe ze\-ne si\-vi maw\-krr\-a 'uo\-ti ro\-lun.} Yes, the Skypeople love exploration, but they don't know what to do once they find something.~\LINK{http://naviteri.org/2014/11/twenty-before-the-holidays/}{b} 
(\ding{43}~\HL{lang}{\B{lang}}, \HL{sAlang}{\B{sä\-lang}})

% tìlapx
\hi 
\HT{tIlapx}{\B{tì\cp{lapx}}} \I{[tI."lap']} \emph{n.} regret. 
\B{Tsa\-tì\-pe\-'un\-ìri ke lu oe\-ru ke\-a tì\-lapx.} I have no reget(s) about that decision.~\LINK{http://naviteri.org/2022/12/trr-anawm-polaheiem-the-great-day-has-arrived/}{b}  
(\ding{43}~\HL{lapx}{\B{lapx}})

% tìlatem
\hi
\HT{tIlatem}{\B{tì\cp{la}tem}} \I{[tI."la.tEm]} \emph{n.} change (abstract concept).
\B{Pxay\-a su\-te\-ri, tì\-la\-tem lu ngä\-zìk.} For many people, change is difficult.~\LINK{http://naviteri.org/2015/11/vomuna-liu-amip-ten-new-words/}{b}
(\ding{43}~\HL{latem}{\B{la\-tem}}, \HL{sAlatem}{\B{sä\-la\-tem}})

% tìlen
\hi
\HT{tIlen}{\B{tì\cp{len}}} \I{[tI."lEn]} \emph{n.} event, happening. 
\B{tì\-len a\-fe'}, a bad event.~\LINK{http://naviteri.org/2011/07/txantsana-ultxa-mi-siatll-great-meeting-in-seattle/}{b}
\B{Maw\-krr\-a fko lie so\-li tì\-len\-ur a\-fpxa\-mo fì\-txan, tì\-rey ke lu teng kaw\-krr.} After experiencing such a terrible event, life is never the same.~\LINK{http://naviteri.org/2018/08/fmawnti-stolawm-srak-have-you-heard-the-news/}{b} 
\B{… fu\-la tsun le\-fkrr\-a sì\-len\-it tswi\-va' hì\-krr (…) oe\-ti 'ey\-ke\-fu nit\-ram nì\-txan.} … that we can forget current events for a short time, makes me very happy.~\LINK{http://naviteri.org/2022/05/results-of-the-tttc/}{b}
(\ding{43}~\HL{len}{\B{len}})

% tìleym
\hi 
\HT{tIleym}{\B{tì\cp{leym}}} \I{[tI."lEjm]} \emph{n.} call. 
\B{Ey\-wal tì\-leym\-it aw\-nge\-yä sto\-lei\-awm!} Eywa has heard our call!~\LINK{http://naviteri.org/2017/09/zisikrr-amip-ayliu-amip-new-words-for-the-new-season/}{b} 
(\ding{43}~\HL{leym}{\B{leym}})

% tìleymfe'
\hi 
\HT{tIleymfe'}{\B{tìleym\cp{fe'}}} \I{[tI.lEjm."fEP]} \emph{n.} complaining (\emph{abstract concept}). 
(\ding{43}~\HL{leymfe'}{\B{leym\-fe'}})

% tìleymkem
\hi 
\HT{tIleymkem}{\B{tìleym\cp{kem}}} \I{[tI.lEjm."kEm]} \emph{n.} protesting, protest (\emph{abstract concept}).
\B{Mì ftaw\-nem\-krr, sì\-leym\-kem a tì\-kem\-wiä\-wä maw kin\-trr a\-'aw to\-ler\-kup; set ke te\-rer\-kup ki 'e\-rul.} In the past, the protests against injustice have died after one week; now they are not dying but rather increase.~\LINK{http://naviteri.org/2020/06/mi-tanlokxe-oeya-srr-afpxamo-terrible-days-in-my-country/\#comment-31259}{b}
(\ding{43}~\HL{leymkem}{\B{leym\-kem}})

% tìloho
\hi 
\HT{tIloho}{\B{tì\cp{lo}ho}} \I{[tI."lo.ho]} \emph{n.} surprise. 
\B{Epxang\-mì lu 'u\-pe?--Tì\-lo\-ho.} What's in the stone jar?--It's a surprise.~\LINK{http://naviteri.org/2019/06/50a-liu-amip-40-new-words}{b}
(\ding{43}~\HL{loho}{\B{lo\-ho}})

% tìlor
\hi
\HT{tIlor}{\B{tì\cp{lor}}} \I{[tI."loR]} \emph{n.} beauty.
\B{Nar\-mew oe fo\-ru na'\-rìng\-ä tì\-lor\-it wi\-vìn\-txu, slä ke tsä\-ngun fo tsli\-vam.} I wanted to show them the beauty of the forest, but sadly, they weren't able to understand.~\LINK{http://naviteri.org/2014/11/twenty-before-the-holidays/}{b} 
(\ding{43}~\HL{lor}{\B{lor}}) \\
\textsc{Derivatives}: \ding{43}
\on{1}.~\HL{somtIlor}{\B{som\-tì\-lor}}, popsicle.

% tìm
\hi
\HT{tIm}{\B{tìm}}\textsuperscript{\on{1}} \I{[tIm]} \emph{adj.} low. 
\B{Ta\-lun\-a vo\-spxì\-o a\-mrr ke zo\-lup tom\-pa, lä\-ngu yì kil\-van\-ä tìm nì\-txan.} Because it hasn't rained for five months, the river level is very low.~\LINK{http://naviteri.org/2013/01/lifyari-po-peyi-how-good-is-her-navi/}{b}
(\ding{214}~\HL{kxayl}{\B{kxayl}}) \\
\textsc{Derivatives}: \ding{43}
\on{1}.~\HL{tImyI}{\B{tìm\-yì}}, low level.

% tìm
\hi
\B{tìm}\textsuperscript{\on{2}} : \HL{txIm}{\B{txìm}}. 

% tìmal
\hi
\HT{tImal}{\B{tì\cp{mal}}} \I{[tI."mal]} \emph{n.} trustworthiness.   
\B{Le\-kin lu tì\-txur, lu tì\-tstew. Slä le\-tsran\-ten fra\-to lu tì\-mal.} Strength and courage are necessary. But most important of all is trustworthiness.~\LINK{http://naviteri.org/2012/01/more-additions-to-the-lexicon/}{b}
(\ding{43}~\HL{mal}{\B{mal}})

% tìme'em
\hi
\HT{tIme'em}{\B{tì\cp{me}'em}} \I{[tI."mE.PEm]} \emph{n.} harmony (general sense).
\B{Fì\-lì\-'u\-ä ral lu tì\-me\-'em sì tì\-ru\-sey mì hi\-fkey na Nawm\-a Sa'\-nok\-ä ha\-pxì.} The meaning of this word is ``harmony, living in the world as part of the Great Mother.''~\LINK{http://forum.learnnavi.org/language-updates/trr-rrtaya/}{ln} 
(\ding{43}~\HL{me'em}{\B{me\-'em}}, \HL{meoauniaea}{\B{me\-o\-a\-u\-ni\-a\-e\-a}})

% tìmeyp
\hi
\HT{tImeyp}{\B{tì\cp{meyp}}} \I{[tI."mEjp\textcorner]} \emph{n.} weakness.   
\B{Kun\-sìp\-ìri txan\-a tì\-meyp lu tsyal a mìn.} The gunship's main weakness is the rotor system.~\LINK{http://forum.learnnavi.org/language-updates/ftarpa-and-skienpa-added-to-dict/msg570362/\#msg570362}{ln\slash a\textsuperscript{c}}
(\emph{a. proverb}) \B{Tìk\-tse\-ri lu tì\-meyp.} Lack of awareness is (a form of) weakness.~\LINK{http://naviteri.org/2015/03/kap-si-ayunil-saylahe-threats-dreams-and-other-things/}{b}
(\ding{43}~\HL{meyp}{\B{meyp}})

% tìmimu
\hi 
\HT{tImimu}{\B{tì\cp{mi}mu}} \I{[tI."mi.mu]} \emph{n.} development. 
(\ding{43}~\HL{mimu}{\B{mi\-mu}}) 

% tìmll'an
\hi
\HT{tImll'an}{\B{tìmll\cp{'an}}} \I{[tI.m\s{l}."Pan]} \emph{n.} acceptance. 
(\ding{214}~\HL{tItsyAr}{\B{tì\-tsyär}}, \ding{43}~\HL{mll'an}{\B{mll\-'an}})

% tìmok
\hi
\HT{tImok}{\B{tì\cp{mok}}} \I{[tI."mok\textcorner]} \emph{n.} suggestion, the abstract idea of suggesting.
\B{Fì\-txe\-le\-ri tì\-mok ke tam; ze\-ne fko fngi\-vo'.} In this matter, suggesting won't cut it; you need to demand.~\LINK{http://naviteri.org/2012/02/trr-asawnung-lefpom-happy-leap-day/}{b}
(\ding{43}~\HL{mok}{\B{mok}}, \HL{sAmok}{\B{sä\-mok}})

% tìmu'ni
\hi 
\HT{tImu'ni}{\B{tì\cp{mu'}ni}} \I{[tI."muP.ni]} \emph{or} \I{[."mUP.]} \emph{n.}  achievement, accomplishment.~\LINK{http://naviteri.org/2020/11/vospxivopeya-ayliu-amip-novembers-new-words/}{b}
(\ding{43}~\HL{mu'ni}{\B{mu'\-ni}}) 

% tìmu'nitkan
\hi 
\HT{tImu'nitkan}{\B{tì\cp{mu'}nitkan}} \I{[tI."muP.nit\textcorner.kan]} \emph{or} \I{[."mUP.]} \emph{n.} effectiveness. 
(\ding{43}~\HL{mu'nitkan}{\B{mu'\-nit\-kan}}) 

% tìmuntxa
\hi
\HT{tImuntxa}{\B{tìmun\cp{txa}}} \I{[tI.mun."t'a]} \emph{or} \I{[.mUn.]} \emph{n.}  mating, Na'vi marriage. 
\B{Zì\-tsal\-trr\-ìri tì\-mun\-txa\-yä ay\-lrr\-tok!} Happy anniversary (of your marriage)!~\LINK{http://naviteri.org/2017/12/zisit-amip-lefpom-ma-eylan-happy-new-year-friends/}{b} 
(\ding{43}~\HL{muntxa}{\B{mun\-txa}})

% tìmungwrr / si
\hi
\HT{tImungwrr}{\B{tìmung\cp{wrr}}} \I{[tI.muN."w\s{r}]} \emph{or} \I{[.mUN.]} (\textsc{rn}: \B{tìmùng\cp{wrr}} \I{[tI.mUN."w\s{r}]}) \textsc{i}. \emph{n.} exception.
\B{Fwa taw\-tu\-te slu Na'\-vi\-yä ha\-pxì lo\-lu tì\-mung\-wrr apxa.} It was a great exception for a human to become part of the People.~\LINK{http://naviteri.org/2015/08/aylifyavi-lereyfya-2-cultural-terms-2-and-more/}{b} \\
\textsc{ii}. \B{tìmung\cp{wrr} si} \I{[tI.muN."w\s{r} si]} \emph{vin.} make an exception.
\B{Ni\-nat tì\-mung\-wrr sìl\-mi fte Ra\-lu tsi\-vun ki\-vä hu tar\-po\-ngu.} Ninat just made an exception so that Ralu can go with the hunting party.~\LINK{http://naviteri.org/2015/08/aylifyavi-lereyfya-2-cultural-terms-2-and-more/}{b}
(\ding{43}~\HL{mungwrr}{\B{mung\-wrr}})

% tìmweypey
\hi
\HT{tImweypey}{\B{tìm\cp{wey}pey}} \I{[tIm."wEj.pEj]} \emph{n.} patience.
\B{Tsa\-krr\-vay, ay\-nge\-yä tìm\-wey\-pey\-ri ira\-yo sei\-yi oe, ul\-te fì\-trr\-ä ftxo\-zä\-ri, sìl\-pey oe, ay\-nga\-ru prr\-te' li\-vu.} In the meantime, I thank you for your patience, and I hope you enjoy today's celebration.~\LINK{http://forum.learnnavi.org/language-updates/trr-rrtaya/}{ln}
(\ding{43}~\HL{maweypey}{\B{ma\-wey\-pey}})

% tìmwiä
\hi
\HT{tImwiA}{\B{tìm\cp{wi}ä}} \I{[tIm."wi.\ae]} \emph{n.} fairness, justice.   
\B{tì\-ngay sì tìm\-wi\-ä}, truth and justice.~\LINK{http://forum.learnnavi.org/language-updates/tirey-lu-muia/}{ln}
\B{Oe\-ri ke lu tìm\-wiä, ha nga\-ti oel ngey\-ka\-syä\-'än nì\-teng.} For me there isn't justice, so I'll make you be miserable as well.~\LINK{http://naviteri.org/2020/06/mi-tanlokxe-oeya-srr-afpxamo-terrible-days-in-my-country/}{b}
(\ding{43}~\HL{muiA}{\B{mu\-i\-ä}})

% tìmyì
\hi
\HT{tImyI}{\B{\cp{tìm}yì}} \I{["tIm.jI]} \emph{n.} low level.
\B{Lì'\-fya\-ri tsa\-taw\-tu\-te pe\-yì? -- Tìm\-yì. Pol ke tslam stum ke\-'ut, omum lì\-'ut a\-vol nì\-'aw.} How is that Sky Person's Na'vi? -- Pretty bad. He understands almost nothing and only knows eight words.~\LINK{http://naviteri.org/2013/01/lifyari-po-peyi-how-good-is-her-navi/}{b}
(\ding{43}~\HL{tIm}{\B{tìm}}, \HL{yI}{\B{yì}}, \HL{kxaylyI}{\B{kxayl\-yì}})

% tìn
\hi
\HT{tIn}{\B{tìn}} \I{[tIn]} \emph{n.} activity that keeps one busy. 
\B{Sìn A\-sok, Sìn Zu\-saw\-krr\-ä}, Recent and Upcoming Activities [lit.: recent activities, future's activities].~\LINK{http://naviteri.org/2013/04/sin-asok-sin-zusawkrra-recent-and-upcoming-activities/}{b}
(\ding{43}~\B{'ìn}) \\
\textsc{Derivatives}: \ding{43}
\on{1}.~\HL{tInvi}{\B{tìn\-vi}}, task, errand, step.
\on{2}.~\HL{txintIn}{\B{txin\-tìn}}, primary role in society.

% tìnalsteng
\hi 
\HT{tInalsteng}{\B{tì\cp{nal}steng}} \I{[tI."nal.stEN]} \emph{n.} compassion, empathy. 
(\ding{43}~\HL{nalsteng}{\B{nal\-steng}}) \\
\textsc{Derivatives}: \ding{43}  
\on{1}.~\HL{tInalstenga'}{\B{tì\-nal\-ste\-nga'}}, compassionate (\emph{nfp.}).

% tìnalstenga'
\hi 
\HT{tInalstenga'}{\B{tì\cp{nal}stenga'}} \I{[tI."nal.stE.NaP]} \emph{adj., nfp.} compassionate. 
(\ding{43}~\HL{tInalsteng}{\B{tì\-nal\-steng}}, \HL{-nga'}{\B{-nga'}}; \HL{lenalsteng}{\B{le\-nal\-steng}})

% tìnatxu
\hi 
\HT{tInatxu}{\B{tìna\cp{txu}}} \I{[tI.na."t'u]} \emph{n.} disapproval. 
(\ding{43}~\HL{natxu}{\B{na\-txu}}, \ding{214}~\HL{tIsmaw}{\B{tì\-smaw}})

% tìnawri
\hi
\HT{tInawri}{\B{tì\cp{naw}ri}} \I{[tI."naw.ri]} \emph{n.} talent. 
\B{Tì\-ru\-sol\-ìri ke lu oe\-ru kea tì\-naw\-ri kaw\-'it.} I have absolutely no talent for singing.~\LINK{http://naviteri.org/2015/04/some-new-words-for-may-day/}{b} 
(\ding{43}~\HL{nawri}{\B{naw\-ri}})

% tìnew
\hi
\HT{tInew}{\B{tì\cp{new}}} \I{[tI."nEw]} \emph{n.} desire (\emph{can refer either to the general state or concept or to a specific instance}).
\B{Tsam\-si\-yu\-ri lu tì\-yo\-ra'\-ä tì\-new le\-kin.} A warrior must have the desire for victory.~\LINK{http://naviteri.org/2014/11/twenty-before-the-holidays/}{b}
\B{Lu oer tì\-new a tse\-'a txam\-pay\-it.} I have a desire to see the ocean.~\LINK{http://naviteri.org/2014/11/twenty-before-the-holidays/}{b}
\B{Pxìm lu tì\-new le\-hawng kxu\-tu fpom\-ä.} Excessive desire is often the enemy of peace.~\LINK{http://naviteri.org/2014/11/twenty-before-the-holidays/}{b}
(\ding{43}~\HL{new}{\B{new}})

% tìnitram
\hi 
\HT{tInitram}{\B{tìnit\cp{ram}}} \I{[tI.nit\textcorner."Ram]} \emph{n.} happiness (\emph{usually resulting from a particular situation or event}). 
\B{Lu oe\-ru tì\-nit\-ram.} I feel happy.~\LINK{https://naviteri.org/2024/09/pxevola-liu-amip-two-dozen-new-words/}{b} 
(\ding{43}~\HL{nitram}{\B{nit\-ram}}) \\
\textsc{Derivatives}: \ding{43}
\on{1}.~\HL{tInitramnga'}{\B{tì\-nit\-ram\-nga'}}, happy (nfp.).

% tìnitramnga'
\hi 
\HT{tInitramnga'}{\B{tìnit\cp{ram}nga'}} \I{[tI.nit\textcorner."Ram.NaP]} \emph{or coll.} \I{[tI.nit\textcorner."Ra.NaP]} \emph{adj., nfp.} happy. 
\B{tì\-len a\-tì\-nit\-ram\-nga'}, happy event.~\LINK{https://naviteri.org/2024/09/pxevola-liu-amip-two-dozen-new-words/\#comment-50634}{b} 
(\ding{43}~\HL{tInitram}{\B{tì\-nit\-ram}}, \HL{nitram}{\B{nit\-ram}})

% tìno
\hi
\HT{tIno}{\B{tì\cp{no}}} \I{[tI."no]} \emph{n.} thoroughness, attention to detail. 
\B{Lä\-ngu tì\-kang\-kem fe\-yä lu\-ke tì\-no.} Unfortunately there is no attention to detail in their work.~\LINK{http://naviteri.org/2010/09/getting-to-know-you-part-2/}{b}

% tìnomum
\hi
\HT{tInomum}{\B{tì\cp{no}mum}} \I{[tI."no.mum]} \emph{or} \I{[.mUm]} (\textsc{rn}: \B{tì\cp{no}mùm} \I{[tI."no.mUm]}) \emph{n.} curiosity.
(\ding{43}~\HL{newomum}{\B{new\-o\-mum}})

% tìnui
\hi
\HT{tInui}{\B{tì\cp{nu}i}} \I{[tI."nu.i]} \emph{n.} failure.
\B{Fì\-txe\-le\-ri tì\-nui ke lu tì\-ftxey.} In this matter, failure is not an option.~\LINK{http://naviteri.org/2013/04/wheres-the-bathroom-and-other-useful-things/}{b}
(\ding{43}~\HL{nui}{\B{nui}}, \HL{sAnui}{\B{sä\-nui}})

% tìnvi / si
\hi 
\HT{tInvi}{\B{\cp{tìn}vi}} \I{["tIn.vi]} \textsc{i}. \emph{n.} task, errand, step (in an instruction). \\
\textsc{ii}. \B{\cp{tìn}vi si} \I{["tIn.vi si]} \emph{vin.} perform a task, run an errand.
\B{Oer txoa li\-vu, ke tsun oe ki\-vä nga\-hu. Ze\-ne pxay\-a tìn\-vi si\-vi.} Sorry, I can't go with you. I have a lot of errands to run.~\LINK{http://naviteri.org/2016/06/mrrvola-lifyavi-amip-forty-new-expressions/}{b} 
(\ding{43}~\HL{tIn}{\B{tìn}})

% tìng
\hi
\HT{tIng}{\B{tìng}} \I{[tIN]} \emph{vtr.} give (s.t. to s.o.).   
\B{Sìl\-pey oe, la\-yu oe\-ru ye'\-rìn sìl\-tsan\-a fmawn a tsun oe ay\-nga\-ru ti\-vìng.} I hope I will soon have good news to give you.~\LINK{http://forum.learnnavi.org/news-announcements/a-response-from-paul-frommer!/}{ln}
\B{Txep\-ìl tìng le\-fpom\-a sä\-nrr\-ti.} The fire gives a pleasant glow.~\LINK{http://naviteri.org/2012/07/fivospxiya-aylifyavi-amip-this-months-new-expressions-pt-2/}{b}
\B{Zun ngal tsa\-fne\-syu\-vet tim\-vìng oer, zel li\-vu oe txur fì\-trr.} If you had given me that kind of food, I would be strong today.~\LINK{http://naviteri.org/2013/04/zun-zel-counterfactual-conditionals/}{b}
\B{Tì\-oeyk\-tìng\-it oel tì\-syìng mì upxa\-re a\-hay.} I intend to give an explanation in the next message.~\LINK{http://naviteri.org/2011/12/ulte-yora\%E2\%80\%99tu-leiu-and-the-winner-is/comment-page-1/\#comment-1267}{b}
(\emph{a. with} \B{rì'ìr}) reflect, give reflection. \B{Oe\-ri pay\-ìl tìng rì\-'ìr\-it key\-ä.} My face is reflected in the water.~\LINK{http://naviteri.org/2011/09/miscellaneous-vocabulary/}{b} 
(\emph{b. with} \B{yoa} \emph{to convey trade}) \B{Oel to\-lìng nga\-ru tsngan\-it yoa fkxen.} I traded you meat for vegetables.~\LINK{http://naviteri.org/2014/02/barter-and-exchange/}{b} 
(\emph{c. special uses}) \\
(i.) \B{tìng ftxì} \I{[tIN ft'I]} \emph{vin.} (\texttt{+}\emph{control}) taste. (\ding{43}~\B{may'}, \B{ew\-ku}) \\
(ii.) \B{tìng \cp{mik}yun} \I{[tIm "mik\textcorner.jun]} \emph{vin.} (\texttt{+}\emph{control}) listen.
\B{Tìng mik\-yun fte tsli\-vam.} Listen to understand.~\LINK{http://naviteri.org/2010/07/ting-mikyun-fte-tslivam-listening-comprehension-1-3/}{b}
\B{Oe\-ru tìng mik\-yun, ma Ra\-lu.} Listen to me, Ralu.~\LINK{http://naviteri.org/2013/01/awvea-posti-zisita-amip-first-post-of-the-new-year/}{b}
\B{Ru\-txe ti\-vìng mik\-yun, ma fra\-po.} Your attention, please, everyone!~\LINK{http://naviteri.org/2010/09/getting-to-know-you-part-1/}{b}
\B{Po pa\-mll\-txe a krr, fra\-po tar\-mìng mik\-yun nì\-pxi, ta\-lu\-na mok\-ri lu sä\-tsì\-syì\-tsyìp.} When he spoke, everyone listened intently, because his voice was a tiny whisper.~\LINK{http://naviteri.org/2011/09/miscellaneous-vocabulary/}{b} (\ding{43}~\B{stawm}, \B{yu\-ne}) \\
(iii.) \B{tìng \cp{na}ri} \I{[tIn "na.Ri]} \emph{vin.} (\texttt{+}\emph{control}) look.
\B{Po\-ru tìng na\-ri.} Look at him.~\LINK{http://wiki.learnnavi.org/index.php?title=Na\%27vi_from_Avatar_Movie\#With_Direhorse}{wiki\slash a} 
\B{Ru\-txe ti\-vìng na\-ri ne\-kll.} Please, see below.~\LINK{http://naviteri.org/2011/07/txantsana-ultxa-mi-siatll-great-meeting-in-seattle/comment-page-1/\#comment-847}{b} 
\B{Le\-hrr\-ap lu fwa evi\-tsyìp sle\-le mì hil\-van lu\-ke fwa fyeyn\-tu te\-rìng nari-} It's dangerous for tiny ones to swim in the river without an adult watching.~\LINK{http://naviteri.org/2011/02/new-vocabulary-ii\%E2\%80\%94part-1/}{b} (\ding{43}~\B{tse\-'a}, \B{nìn}) \\ 
(iv.) \B{tìng \cp{on}tu} \I{[tIN "{}on.tu]} \emph{vin.} (\texttt{+}\emph{control}) smell. (\ding{43}~\B{he\-fi}, \B{syam}) \\
(v.) \B{tìng tseng} \I{[tIN \t{ts}EN]} \emph{vin.} back down. \\ 
(vi.) \B{tìng lawr} \I{[tIN lawR]} \emph{vin.} sing wordlessly, give out a tune or melody. 
\B{'E\-van\-ìl alo a\-'aw\-ve nì\-'aw\-tu na'\-rìng\-it tar\-mok, ha to\-lìng lawr nì\-fwe\-fwi fte\-ke txo\-pu si\-vi.} The boy was alone in the forest for the first time, so he whistled a tune to calm his fears.~\LINK{http://naviteri.org/2011/09/miscellaneous-vocabulary/}{b} \\ 
(vii.) \B{tìng zekwä} \I{[tIN "zEk\textcorner.w\ae]} \emph{vin.} (\texttt{+}\emph{control}) touch (intentionally). 
\B{Rä\-'ä tìng zek\-wä oer!} Don't touch me!~\LINK{http://naviteri.org/2015/04/some-new-words-for-may-day/}{b}
\B{Lu fì\-vul su\-syang nì\-txan. Txo fko ti\-vìng zek\-wä kxakx.} This branch is very fragile. If you touch it, it'll break.~\LINK{http://naviteri.org/2015/04/some-new-words-for-may-day/}{b} 
(\ding{43}~\B{zìm}, \B{'am\-pi}) 
(viii.) \B{tìng syawn} \I{[tIN sjawn]} \emph{vin.} bless, give blessing. 
\B{Eywa ti\-vìng syawn ngar, ma 'i\-te.} May Ey\-wa bless you, my daughter.~\LINK{http://naviteri.org/2018/03/100a-liu-amip-64-new-words-part-1/}{b} \\ 
(ix.) \B{tìng \cp{rä}zekx}  \I{[tIN "R\ae.zEk']} \emph{vin.} curse (s.o.), wishing s.o. ill.
\B{Oe nga\-ru tìng rä\-zekx, ma kal\-wey\-a\-veng!} I curse you, (OR: Damn you!), you son of a bitch.~\LINK{https://naviteri.org/2025/06/vospikin-lefpom-happy-july/}{b} 
\B{Rä\-zekx ngar!} Damn you! [lit.: ``(I) curse you!'']~\LINK{https://naviteri.org/2025/06/vospikin-lefpom-happy-july/\#comment-54817}{b} \\ 
\textsc{Derivatives}: \ding{43}
\on{1}.~\HL{oeyktIng}{\B{oeyk\-tìng}}, explain.
\on{2}.~\HL{pAnutIng}{\B{pä\-nu\-tìng}}, promise. 
\on{3}.~\HL{stxenutIng}{\B{stxe\-nu\-tìng}}, offer. 
\on{4}.~\HL{teswotIng}{\B{te\-swo\-tìng}}, grant.  
\on{5}.~\HL{yomtIng}{\B{yom\-tìng}}, feed.

% tìnga'lini
\hi 
\HT{tInga'lini}{\B{tìnga'\cp{li}ni}} \I{[tI.NaP."li.ni]} \emph{n.} pregnancy (\emph{for animals}). 
(\ding{43}~\HL{nga'lini}{\B{nga'\-li\-ni}}, \HL{tInga'prrnen}{\B{tì\-nga'\-prr\-nen}}) 

% tìnga'prrnen
\hi 
\HT{tInga'prrnen}{\B{tìnga'\cp{prr}nen}} \I{[tI.NaP."p\s{r}.nEn]} \emph{n.} pregnancy (\emph{for people}). 
(\ding{43}~\HL{nga'prrnen}{\B{nga'\-prr\-nen}}, \HL{tInga'lini}{\B{tì\-nga'\-li\-ni}})

% tìngay
\hi
\HT{tIngay}{\B{tì\cp{ngay}}} \I{[tI."Naj]} \emph{n.} truth.   
\B{tì\-ngay sì tìm\-wi\-ä}, truth and justice.~\LINK{http://forum.learnnavi.org/language-updates/tirey-lu-muia/}{ln}
\B{Oe\-ri tì\-ngay\-ìl txe'\-lan\-it ti\-va\-kuk.} Let the truth strike my heart.~\LINK{http://naviteri.org/2013/08/taronway-the-hunt-song/}{b} 
\B{Lo\-lu ka\-vuk, slä Tse\-nul tì\-ngay\-it ko\-lu\-lat.} There was treachery, but Tsenu revealed the truth.~\LINK{http://naviteri.org/2011/10/more-vocabulary-a-bit-of-grammar/}{b}
(\ding{43}~\HL{ngay}{\B{ngay}})

% tìngäzìk
\hi
\HT{tIngAzIk}{\B{tì\cp{ngä}zìk}} \I{[tI."N\ae.zIk\textcorner]} \emph{n.} difficulty, problem.
\B{Lu oe\-ru … ìì … tì\-ngä\-zìk a\-hì\-'i.} I have … um … a slight problem.~\LINK{http://naviteri.org/2011/03/\%E2\%80\%9Creceptive-ability\%E2\%80\%9D-and-hesitation/}{b}
\B{'Ul tu\-te, 'ul tì\-ngä\-zìk.} The more people, the more problems.~\LINK{http://naviteri.org/2012/02/trr-asawnung-lefpom-happy-leap-day/}{b}
\B{Pol zo\-let oe\-yä sä\-fpìl\-it pxel pum a tì\-ngä\-zìk\-it ngop.} He treated my idea as though it created a problem.~\LINK{http://naviteri.org/2012/10/mipa-vospxi-mipa-ayliu-new-words-for-the-new-month/}{b}
\B{Ngal ke tslä\-ngam teyng\-ta fì\-tì\-ngä\-zìk\-ìri pe\-fyin\-ep\-'ang.} Unfortunately you don't understand how complex this problem is.~\LINK{http://naviteri.org/2012/07/fivospxiya-aylifyavi-amip-this-months-new-expressions-pt-2/}{b} 
(\ding{43}~\HL{ngAzIk}{\B{ngä\-zìk}})

% tìngong
\hi
\HT{tIngong}{\B{tì\cp{ngong}}} \I{[tI."NoN]} \emph{n.} lethargy, laziness.
\B{Fwa Ìstaw\-hu 'aw\-si\-teng tì\-kang\-kem si ke su\-nu oer; tì\-ngong\-ìri ke lu kaw\-tu na po.} I don't like to work with Ìstaw; he's famous for his laziness.~\LINK{http://naviteri.org/2011/10/more-vocabulary-a-bit-of-grammar/}{b}
(\ding{43}~\HL{ngong}{\B{ngong}})

% tìngop
\hi
\HT{tIngop}{\B{tì\cp{ngop}}} \I{[tI."Nop\textcorner]} \emph{n.} creation.
\B{Ftxo\-zä\-ri fì\-tì\-ngop a\-lor lu swey\-a stxe\-li a to\-lel oel.} This beautiful creation is the best present I have received for (my birthday) celebration.~\LINK{http://forum.learnnavi.org/community-calendar/karyu-pawls-birthday/msg312733/\#msg312733}{ln}
(\ding{43}~\HL{ngop}{\B{ngop}}, \HL{sAngop}{\B{sä\-ngop}})

% tìoeyktìng
\hi
\HT{tIoeyktIng}{\B{tìo\cp{eyk}tìng}} \I{[tI.o."{}Ejk\textcorner.tIN]} \emph{or} \I{[tI."wEjk\textcorner.tIN]} \emph{n.} explanation. 
\B{Sìl\-pey oe tsnì fì\-tì\-oeyk\-tìng law li\-vu nga\-ru set.} I hope that this explanation is clear to you now.~\LINK{http://forum.learnnavi.org/language-updates/a-long-answer-to-a-quick-question/}{ln}
\B{Tì\-oeyk\-tìng za\-'u ye'\-rìn.} An explanation comes soon.~\LINK{http://naviteri.org/2011/02/new-vocabulary-ii\%E2\%80\%94part-1/comment-page-1/\#comment-520}{b}
\B{Ul\-te sìl\-pey oe, fay\-sì\-oeyk\-tìng li\-vu law nì\-wotx!} And I hope, all these explanations are clear.~\LINK{http://naviteri.org/2011/08/reported-speech-reported-questions/}{b}
\B{Ru\-txe oe\-yä tì\-oeyk\-tìng\-it a\-saw\-nung ni\-vìn.} Please look at my added explanation.~\LINK{http://naviteri.org/2011/03/word-order-and-case-marking-with-modals/comment-page-1/\#comment-581}{b}
\B{Fì\-tì\-oeyk\-tìng\-ìri ira\-yo, ma tsmuk.} Thanks for this explanation, brother.~\LINK{http://naviteri.org/2014/10/tson-si-fpomron-obligation-and-mental-health/comment-page-1/\#comment-10675}{b}
(\ding{43}~\HL{oeyktIng}{\B{oeyk\-tìng}})

% tìohakx
\hi
\HT{tIohakx}{\B{tìo\cp{hakx}}} \I{[tI.o."hak']} \emph{n.} hunger. 
(\ding{43}~\HL{ohakx}{\B{o\-hakx}})

% tìomum
\hi
\HT{tIomum}{\B{tì\cp{o}mum}} \I{[tI."{}o.mum]} \emph{or} \I{[.mUm]} (\textsc{rn}: \B{tì\cp{o}mùm} \I{[tI."{}o.mUm]}) \emph{n.} knowledge.
(\emph{idiom}) \B{tì\-o\-mum\-mì oe\-yä}, to my knowledge, as far as I know.
\B{Tì\-o\-mum\-mì oe\-yä, pol na'\-rìng\-it i\-nan nìl\-tsan.} As far as I know, he reads the forest well.~\LINK{http://naviteri.org/2011/09/\%E2\%80\%9Cby-the-way-what-are-you-reading\%E2\%80\%9D/}{b}
(\ding{43}~\HL{omum}{\B{omum}}, \HL{sAomum}{\B{sä\-omum}}) 

% tìpalang
\hi 
\HT{tIpalang}{\B{tì\cp{pa}lang}} \I{[tI."pa.laN]} \emph{n.} (social) contact. 
\B{Po\-hu ke lu oe\-ru kea tì\-pa\-lang kaw\-'it.} I have no contact with him whatsoever.~\LINK{http://naviteri.org/2020/03/teri-tifkeytok-lefkrr-about-the-present-situation/}{b} 
(\ding{43}~\HL{palang}{\B{pa\-lang}})

% tìpaw
\hi 
\HT{tIpaw}{\B{tì\cp{paw}}} \I{[tI."paw]} \emph{n.} growth. 
(\ding{43}~\HL{paw}{\B{paw}}) \\
\textsc{Derivatives}: \ding{43}
\on{1}.~\HL{zIskrrmipaw}{\B{zì\-skrr\-mi\-paw}}, spring (season).

% tìpawm
\hi
\HT{tIpawm}{\B{tì\cp{pawm}}} \I{[tI."pawm]} \emph{n.} question.   
\B{tì\-pawm a\-syen}, last question.~\LINK{http://naviteri.org/2010/08/mipa-ayopin-mipa-ayliu-new-colors-new-words/}{b} 
\B{tì\-pawm a\-tor}, final question.~\LINK{http://naviteri.org/2010/08/mipa-ayopin-mipa-ayliu-new-colors-new-words/}{b}
\B{Fra\-krr lu nge\-yä sì\-pawm ngä\-zìk slä\-kop le\-tsran\-ten.} Your questions are always difficult but also important.~\LINK{http://naviteri.org/2014/10/tson-si-fpomron-obligation-and-mental-health/}{b}
\B{Me\-sì\-pawm a\-sìl\-tsan nì\-txan nang!} Two very good questions!~\LINK{http://forum.learnnavi.org/language-updates/txelanit-hivawl/}{ln}
\B{Fol oe\-yä tì\-pawm\-it ke sto\-lä\-ngawm.} Unfortunately they didn't hear my question.~\LINK{http://naviteri.org/2012/11/renu-ayinanfyaya-the-senses-paradigm/}{b}
\B{Sä\-pll\-txel kar\-yu\-ä ke to\-la\-rä\-nge tì\-pawm\-it kaw\-'it.} The teacher's statement in no way pertained to the question. Drat!~\LINK{http://naviteri.org/2012/06/spring-vocabulary-part-3/}{b}
\B{Nge\-yä txan\-tsan\-a tì\-pawm\-ìri nga\-ru sei\-yi oe ira\-yo.} I thank you for your excellent question.~\LINK{http://forum.learnnavi.org/language-updates/another-email-from-frommer/}{ln}
(\ding{43}~\HL{pawm}{\B{pawm}})

% tìpähem
\hi
\HT{tIpAhem}{\B{tì\cp{pä}hem}} \I{[tI."p\ae.hEm]} \emph{n.} arrival.
\B{tì\-pä\-hem le\-ye'\-krr}, an early arrival.~\LINK{http://naviteri.org/2010/07/vocabulary-update/}{b}
\B{Fì\-tseo\-ri ke tsun kaw\-tu pi\-vä\-hem tì\-yo'\-ne; tsran\-ten tì\-pä\-hem\-ä tì\-fmi nì\-'aw.} In this art it's impossible to arrive at perfection; the only thing that matters is the attempt to arrive there.~\LINK{http://naviteri.org/2012/01/mipa-zisit-ayliu-amip-new-words-for-the-new-year/}{b}
(\ding{43}~\HL{pAhem}{\B{pä\-hem}})

% tìpängkxo
\hi
\HT{tIpAngkxo}{\B{tìpäng\cp{kxo}}} \I{[tI.p\ae{}N."k'o]} \emph{n.} conversation, discussion; dialog.  
\B{tì\-päng\-kxo le\-lì'\-fya}, language discussion.~\LINK{http://naviteri.org/2010/06/first-post/}{b}
\B{Tsa\-ko\-ak\-te\-ri stawm\-tswo lu fe'. Po\-hu a tì\-päng\-kxo ngä\-zìk lu nì\-txan.} That old woman's hearing is poor. A conversation with her is very difficult.~\LINK{http://naviteri.org/2012/11/renu-ayinanfyaya-the-senses-paradigm/}{b}
\B{Ke new oel fu\-ta fì\-tì\-päng\-kxot ay\-nge\-yä hi\-vul\-tstxem, slä tsun mi\-vä\-kxu hì\-krr nì\-'aw srak?.} I don't want to derail your chat, but can I interrupt for just a moment?~\LINK{http://naviteri.org/2010/09/getting-to-know-you-part-1/}{b}
(\ding{43}~\HL{pAngkxo}{\B{päng\-kxo}}) \\
\textsc{Derivatives}: \ding{43}
\on{1}.~\HL{tIpAngkxotsyIp}{\B{tì\-päng\-kxo\-tsyìp}}, dialog.

% tìpängkxotsyìp
\hi
\HT{tIpAngkxotsyIp}{\B{tìpäng\cp{kxo}tsyìp}} \I{[tI.p\ae{}N."k'o.\t{ts}Ip\textcorner]} \emph{n.} little conversation, dialog.
\B{Sìl\-pey oe, ay\-nga\-ri fì\-tì\-päng\-kxo\-tsyìp el\-tur tì\-txen sil\-vi.} I hope, you have found this dialog interesting.~\LINK{http://naviteri.org/2010/07/diminutives-conversational-expressions/}{b}
\B{Ta\-fral ho\-lawl ay\-oel ay\-nga\-fpi 'a\-'aw\-a tì\-päng\-kxo\-tsyìp\-it.} That's why we've prepared several little dialogs for you.~\LINK{http://naviteri.org/2012/07/tskxekengtsyip-a-mikyunfpi-a-little-listening-exercise/}{b}
(\ding{43}~\HL{tIpAngkxo}{\B{tì\-päng\-kxo}})

% tìpe'ngay
\hi 
\HT{tIpe'ngay}{\B{tìpe'\cp{ngay}}} \I{[tI.pEP."Naj]} \emph{n.} conclusion (\emph{referring to a determination}). \\
\textsc{i}. \B{tì\-pe'\-ngayt wrr\-zä\-rìp} \emph{vtr.} infer.
\B{Ngey ay\-lì\-'u\-ftu wrr\-zä\-rìp oel tì\-pe'\-ngayt a lu ngar yew\-la.} From your words, I infer that you're disappointed.~\LINK{http://naviteri.org/2021/03/mipa-ayliu-si-aylifyavi-new-words-and-expressions/}{b} \\
\textsc{ii}. \B{tì\-pe'\-ngayt wrr\-zey\-kä\-rìp} \emph{vtr.} imply.
\B{Ngey ay\-lì\-'ul wrr\-zey\-kä\-rìp tì\-pe'\-ngayt a lu ngar yew\-la.} Your words imply that you're disappointed.~\LINK{http://naviteri.org/2021/03/mipa-ayliu-si-aylifyavi-new-words-and-expressions/}{b}
(\ding{43}~\HL{pe'ngay}{\B{pe'\-ngay}}) 

% tìpe'un
\hi
\HT{tIpe'un}{\B{tì\cp{pe}'un}} \I{[tI."pE.Pun]} (\textsc{rn}: \B{tì\cp{pe}ùn} \I{[tI."pE.Un]}) \emph{n.} decision. 
\B{Peyä tsa\-tì\-pe\-'un a swey\-lu txo wi\-vem ayoeng Oma\-ti\-ka\-ya\-wä lu fe'.} His decision to fight against the Omaticaya was a bad one.~\LINK{http://naviteri.org/2011/07/txantsana-ultxa-mi-siatll-great-meeting-in-seattle/}{b}
\B{Tsa\-tì\-pe\-'un\-ìri ke lu oe\-ru ke\-a tì\-lapx.} I have no reget(s) about that decision.~\LINK{http://naviteri.org/2022/12/trr-anawm-polaheiem-the-great-day-has-arrived/}{b}
(\ding{43}~\HL{pe'un}{\B{pe\-'un}})

% tìprrte'
\hi
\HT{tIprrte'}{\B{tì\cp{prr}te'}} \I{[tI."p\s{r}.tEP]} \emph{n.} pleasure. 
(\emph{a. conv. phrase}) \B{Oe\-ru tì\-prr\-te'.} My pleasure.~\LINK{http://naviteri.org/2012/11/renu-ayinanfyaya-the-senses-paradigm/comment-page-1/\#comment-1851}{b}
(\ding{43}~\HL{prrte'}{\B{prr\-te'}})

% tìpsaw
\hi 
\HT{tIpsaw}{\B{tìp\cp{saw}}} \I{[tIp\textcorner."{}saw]} \emph{n.} clumsiness. 
\B{Po\-eyä tìp\-saw\-ìl txo\-pu sley\-ko\-lu ye\-rik\-it ha po hi\-fwo.} Her clumsiness scared the hexapede, so he escaped.~\LINK{http://naviteri.org/2017/09/zisikrr-amip-ayliu-amip-new-words-for-the-new-season/}{b} 
(\ding{43}~\HL{pIsaw}{\B{pì\-saw}})

% tìpxul
\hi 
\HT{tIpxul}{\B{tì\cp{pxul}}} \I{[tI."p'ul]} \emph{or} \I{[.p'Ul]} \emph{n.} formidableness, imposingness.~\LINK{http://naviteri.org/2020/08/hiia-vur-a-teri-mefalukantsyip-a-little-story-about-two-cats/}{b} 
(\ding{43}~\HL{pxul}{\B{pxul}}) 

% tìra'un
\hi 
\HT{tIra'un}{\B{tì\cp{ra}'un}} \I{[tI."Ra.Pun]} \emph{or} \I{[.PUn]} \emph{n.} surrender, relinquishment. 
\B{Pe\-yä tì\-ra\-'un tì\-eyk\-tan\-ä le\-yew\-la lu nì\-txan.} His surrender of leadership is very disappointing.~\LINK{http://naviteri.org/2021/09/aawa-ayliu-amip-a-few-new-words/}{b} 
(\ding{43}~\HL{ra'un}{\B{ra\-'un}}) 

% tìralpeng
\hi
\HT{tIralpeng}{\B{tìral\cp{peng}}} \I{[tI.Ral."pEN]} \emph{n.} translation, interpretation.
\B{Spä\-ngaw oel fu\-ta fì\-'upxa\-re\-yä tì\-ral\-peng ke lu eyawr.} Unfortunately, I don't believe the translation of this message is correct.~\LINK{http://naviteri.org/2015/03/kap-si-ayunil-saylahe-threats-dreams-and-other-things/}{b} 
(\ding{43}~\HL{ralpeng}{\B{ral\-peng}})

% tìran
\hi
\HT{tIran}{\B{tì\cp{ran}}} \I{[tI."Ran]} \emph{vin.} (\emph{a.}) walk. 
\B{Ma 'ite, aw\-nge\-yä fya\-'o\-ri ze\-ne nga sä\-nume si\-vi po\-ru fte tsi\-vun pil\-vll\-txe sì ti\-vì\-ran nì\-ay\-oeng.} My daughter, you will teach him our way to speak and walk as we do.~\LINK{http://wiki.learnnavi.org/index.php?title=Na\%27vi_from_Avatar_Movie\#Scene_in_Hometree}{a\slash wiki}
\B{Sre fwa sìn tskxe\-pay tì\-ran, ze\-ne fko flì\-nutx\-it sti\-ve\-ftxaw.} It's necessary to check the thickness of the ice before walking on it.~\LINK{http://naviteri.org/2011/11/\%E2\%80\%99a\%E2\%80\%99awa-ayli\%E2\%80\%99u-amip-ni\%E2\%80\%99aw-\%E2\%80\%94-a-few-new-words-only/}{b}
\B{Trr\-'ong\-ta Txon\-'ong\-vay po to\-lì\-ran.} He walked from dawn until dusk.~\LINK{http://naviteri.org/2011/01/new-year-new-vocabulary/}{b}
\B{Txon\-am teng\-krr tar\-mì\-ran oe kxam\-lä na'\-rìng, sro\-ler eo ut\-ral a\-tsawl txewm\-a vrr\-tep.} Last night as I was walking through the forest, a frightening demon appeared in front of a big tree.~\LINK{http://naviteri.org/2012/03/spring-vocabulary-part-1/}{b}
(\emph{b. with} \B{nì\-nu}) trip, ``walk failingly''.
\B{Kll\-te lu ekx\-txu. Na\-ri si txo\-ke\-fyaw tì\-ran nì\-nu.} The ground is rough. Be careful or you'll trip.~\LINK{http://naviteri.org/2013/04/wheres-the-bathroom-and-other-useful-things/}{b}
(\emph{c. of the weather}) \B{tì\-ran hu\-fwe}, it's breezy (but pleasantly so).~\LINK{http://naviteri.org/2011/05/weather-part-2-and-a-bit-more-2/}{b} 
\B{Hu\-fwe\-tsyìp le\-fpom tar\-mì\-ran.} A pleasant little breeze was blowing.~\LINK{http://naviteri.org/2011/05/weather-part-2-and-a-bit-more-2/}{b} 
(\emph{d. rhyming expression}) \B{Tslikx, tì\-ran, tul; | Ftu yì ne yì tsan\-'ul.} \emph{when learning s.t. new, you have to proceed from step to step: baby steps first, then bigger ones} [lit.: `Crawl, walk, run; from level to level get better.']~\LINK{http://naviteri.org/2023/09/aawa-ayliu-si-aylifyavi-amip-a-few-new-words-and-expressions/}{b} \\
\textsc{Derivatives}: \ding{43}
\on{1}.~\HL{uniltIranyu}{\B{u\-nil\-tì\-ran\-yu}}, dreamwalker. 
\on{2}.~\HL{uniltIrantokx}{\B{u\-nil\-tì\-ran\-tokx}}, avatar, dreamwalker body; (the movie) \emph{Avatar} (\on{2009}). 
\on{3}.~\HL{tIranpam}{\B{tì\-ran\-pam}}, footstep (\emph{sound}).

% tìranpam
\hi 
\HT{tIranpam}{\B{tì\cp{ran}pam}} \I{[tI."Ran.pam]} \emph{n.} footstep (\emph{sound}). 
\B{Oel stawm sì\-ran\-pam\-it! Le\-rok tu\-teo!} I hear footsteps! Someone is coming!~\LINK{http://naviteri.org/2019/06/50a-liu-amip-40-new-words}{b} 
(\ding{43}~\HL{tIran}{\B{tì\-ran}}, \HL{pam}{\B{pam}})

% tìranyu
\hi
\HT{tIranyu}{\B{tì\cp{ran}yu}} \I{[tI."Ran.ju]} \emph{n.} walker, one who walks.
(\ding{43}~\HL{tIran}{\B{tì\-ran}}) \\
\textsc{Derivatives}: \ding{43}
\on{1}.~\HL{uniltIranyu}{\B{u\-nil\-tì\-ran\-yu}}, dreamwalker.

% tìrawn
\hi
\HT{tIrawn}{\B{tì\cp{rawn}}} \I{[tI."Rawn]} \emph{n.} replacement, act of replacing.
\B{Po 'e\-fu ngeyn ul\-te kin tì\-rawn\-it nì\-txan.} He is tired and very much needs to be replaced.~\LINK{http://naviteri.org/2012/01/mipa-zisit-ayliu-amip-new-words-for-the-new-year/}{b}
(\ding{43}~\HL{rawn}{\B{rawn}})

% tìrengop
\hi
\HT{tIrengop}{\B{tì\cp{re}ngop}} \I{[tI."RE.Nop\textcorner]} \emph{n.} (abstract) design.
\B{Tì\-re\-ngop\-ìri ioi\-yä lu sem\-pul pe\-yä tsul\-fä\-tu.} Her father is a master designer of ceremonial adornments.~\LINK{http://naviteri.org/2013/01/awvea-posti-zisita-amip-first-post-of-the-new-year/}{b}
(\ding{43}~\HL{rengop}{\B{re\-ngop}}, \HL{sArengop}{\B{sä\-re\-ngop}})

% tìrey
\hi
\HT{tIrey}{\B{tì\cp{rey}}} \I{[tI."REj]} \emph{n.} life.  
\B{'Ìhe\-yu sì\-rey\-ä | Sì\-rey le\-Na'\-vi | Oe\-ru te\-ya si.} The spiral of the lives | the lives of the people | Fill me.~\LINK{http://wiki.learnnavi.org/index.php/Weaving_Song}{wiki} 
\B{mì oe\-yä le\-trr\-a tì\-rey}, in my daily life.~\LINK{http://forum.learnnavi.org/language-updates/frommerian-email/}{b} 
\B{Me\-nge\-yä tì\-rey\-ìl 'aw\-si\-teng me\-nga\-ru za\-mì\-ye\-vu\-nge txan\-a fpom\-it sì fpom\-tokx\-it.} May your life together bring you both much peace and health.~\LINK{http://naviteri.org/2014/02/barter-and-exchange/}{b}
\B{Oe\-ri Unil\-tì\-ran\-tokx\-ìl tì\-rey\-ti ley\-ko\-la\-te\-i\-em nì\-txan.} \emph{Avatar} has greatly changed my life for the better.~\LINK{http://forum.learnnavi.org/language-updates/pseudo-canon-from-bd-live-content/msg350756/\#msg350756}{ln}
\B{Nawm\-a sa'\-nok tì\-haw\-nu si\-vi | Aw\-nge\-yä tì\-rey\-ur ki\-fkey\-ur\-sì.} May the Great Mother protect | Our life and world.~\LINK{http://forum.learnnavi.org/news-announcements/christmas-song-ninavi-on-youtube/}{ln}
\B{No\-lu\-me oe nì\-txan te\-ri tì\-tsun\-slu tì\-rey\-ä a hi\-fkey\-mì a\-la\-he.} I learned a lot about the possibility of life on other worlds.~\LINK{http://naviteri.org/2012/07/meetings-waterfalls-and-more/}{b}
\B{Tsei\-un oe pi\-vll\-txe san oe\-ri lu tì\-rey sìl\-tsan nì\-ngay sìk.} I can happily say, `my life is truly good.'~\LINK{http://naviteri.org/2020/05/tipusawm-tiuseyng-si-okvur-a-eltur-titxen-si-asking-answering-and-an-interesting-story/\#comment-31216}{b}
(\ding{43}~\HL{rey}{\B{rey}})

% tìreyn
\hi
\HT{tIreyn}{\B{tì\cp{reyn}}} \I{[tI."REjn]} \emph{n.}, \textsc{lw} train.
\B{Su\-nu Tsyan\-ur tì\-reyn nì\-txan.} John likes trains a lot.~\LINK{http://naviteri.org/2014/08/stxeli-alor-a-beautiful-gift/}{b}

% tìrol
\hi
\HT{tIrol}{\B{tì\cp{rol}}} \I{[tI."Rol]} \emph{n.} song.
\B{Ul\-te lu yaw\-ne oer fì\-tì\-rol ay\-nge\-yä nì\-ngay!} And I truly love this song of yours!~\LINK{http://forum.learnnavi.org/community-calendar/karyu-pawls-birthday/msg312733/\#msg312733}{ln}
\B{Nge\-yä tì\-rol\-tsyìp\-ìri, oe\-yä ay\-syu\-lang set mì fay kll\-kxe\-rem.} As for your little song, my flowers are standing in water now.~\LINK{http://naviteri.org/2010/09/getting-to-know-you-part-1/comment-page-1/\#comment-296}{b}
(\ding{43}~\HL{rol}{\B{rol}})

% tìronsrel
\hi
\HT{tIronsrel}{\B{tì\cp{ron}srel}} \I{[tI."Ron.sREl]} \emph{n.} imagination.   
\B{Lu po\-ru tì\-ron\-srel a\-txa\-na\-tan.} She has a vivid imagination.~\LINK{http://naviteri.org/2011/02/new-vocabulary-part-2/}{b} 
(\ding{43}~\HL{ronsrel}{\B{ron\-srel}})

% tìrun
\hi
\HT{tIrun}{\B{tì\cp{run}}} \I{[tI."Run]} \emph{or} \I{[."RUn]} \emph{n.} discovery.
\B{Tì\-run\-ìl tey\-ngä lum\-pe po ho\-lum oe\-ti ke\-ftxo 'ey\-ko\-le\-fu.} The discovery of why he left saddened me.~\LINK{https://naviteri.org/2024/12/odds-and-ends/}{b}
(\ding{43}~\HL{run}{\B{run}})

% tìsay
\hi 
\HT{tIsay}{\B{tì\cp{say}}} \I{[tI."{}saj]} \emph{n.} loyalty. 
\B{Tì\-'ey\-lan\-ìri tsran\-ten fra\-to tì\-say.} What matters most in friendship is loyalty.~\LINK{http://naviteri.org/2021/12/zolau-niprrte-ma-3746-welcome-2022/}{b} 
(\ding{43}~\HL{say}{\B{say}}) 

% tìsäfpìlyewn
\hi 
\HT{tIsAfpIlyewn}{\B{tìsä\cp{fpìl}yewn}} \I{[tI.s\ae."fpIl.jEwn]} \emph{or coll.}  \I{[\t{ts}\ae."fpIl.jEwn]} \emph{n.} communication. 
\B{Tìsäfpìlyewn ke lu aylì'u nì'aw.} Communication is far more than a matter of words.~\LINK{https://naviteri.org/2026/06/tskxekeng-a-mikyunfpi-pamrel-listening-exercise-text/}{b}
(\ding{43}~\HL{sAfpIlyewn}{\B{sä\-fpìl\-yewn}}) 

% tìska'a
\hi
\HT{tIska'a}{\B{tìska\cp{'a}}} \I{[tI.ska."Pa]} \emph{n.} destruction. 
\B{Tì\-ska\-'a\-ri Kel\-ut\-ral\-ä Na'\-vi wa\-ya\-lew pe\-fya?} With the destruction of Hometree, how will the Na’vi ever move on?~\LINK{http://naviteri.org/2018/07/ta-sulfatu-a-ayliu-niul-more-words-from-our-experts/}{b}
(\ding{43}~\HL{ska'a}{\B{ska\-'a}})

% tìsla'tsu
\hi
\HT{tIsla'tsu}{\B{tì\cp{sla'}tsu}} \I{[tI."{}slaP.\t{ts}u]} \emph{n.} description.
(\ding{43}~\HL{sla'tsu}{\B{sla'\-tsu}}) 

% tìslan
\hi
\HT{tIslan}{\B{tì\cp{slan}}} \I{[tI."{}slan]} \emph{n.} support.
\B{Nge\-yä tì\-eyk\-tan\-ìri, tì\-slan\-ìri sì tsran\-ten fra\-to a tì\-'ey\-lan\-ìri a ka ay\-zì\-sìt nì\-wotx, sei\-yi oe ira\-yo nì\-txan.} Thank you so much for your leadership, your support, and most importantly your friendship throughout the years.~\LINK{http://naviteri.org/2012/03/spring-vocabulary-part-1/}{b}
(\ding{43}~\HL{slan}{\B{slan}})

% tìsmaw
\hi 
\HT{tIsmaw}{\B{tì\cp{smaw}}} \I{[tI."smaw]} \emph{n.} approval. 
\B{Moe\-y\"a tì\-mun\-txa\-ri tì\-smaw nge\-y\"a oe\-ru te\-ya si.} Your approval of our marriage fills me (with joy).~\LINK{http://naviteri.org/2020/02/some-words-for-leap-year-day/}{b}
(\ding{43}~\HL{smaw}{\B{smaw}}, \ding{214}~\HL{tInatxu}{\B{tì\-na\-txu}})

% tìsnaytx
\hi
\HT{tIsnaytx}{\B{tì\cp{snaytx}}} \I{[tI."{}snajt']} \emph{n.} loss (\emph{of a contest, game or fight}). 
\B{Oeyk tsa\-tì\-snaytx\-ä lu apxa sä\-nui tì\-eyk\-tan\-ä.} That loss was caused by a great failure of leadership.~\LINK{http://naviteri.org/2013/04/wheres-the-bathroom-and-other-useful-things/}{b}
(\ding{214}~\HL{tIyora'}{\B{tì\-yo\-ra'}}, \ding{43}~\HL{snaytx}{\B{snaytx}})

% tìso'ha
\hi
\HT{tIso'ha}{\B{tì\cp{so'}ha}} \I{[tI."{}soP.ha]} \emph{n.} enthusiasm; having a good attitude.
(\ding{43}~\HL{so'ha}{\B{so'\-ha}})

% tìsom
\hi
\HT{tIsom}{\B{tì\cp{som}}} \I{[tI."{}som]} \emph{n.} heat.
(\ding{214}~\HL{tIwew}{\B{tì\-wew}}, \ding{43}~\HL{som}{\B{som}})

% tìsop
\hi
\HT{tIsop}{\B{tì\cp{sop}}} \I{[tI."{}sop\textcorner]} \emph{n.} journey. 
\B{Nge\-yä fì\-tì\-sop\-ìri pe\-han\-tseng?} This journey of yours--what's its destination?~\LINK{http://naviteri.org/2019/08/aawa-lifya-si-lifyavi-amip-a-few-new-words-and-expressions/}{b} 
(\ding{43}~\HL{sop}{\B{sop}})

% tìspaw
\hi 
\HT{tIspaw}{\B{tì\cp{spaw}}} \I{[tI."{}spaw]} \emph{n.} belief (\emph{abstract concept}). 
\B{Tsran\-ten tì\-spaw, tsran\-ten nì\-'ul tì\-fkey\-to\-ngay.} Belief is important, but reality is more important.~\LINK{http://naviteri.org/2023/12/mipa-zisit-ayliu-amip-new-year-new-words/}{b} 
(\ding{43}~\HL{spaw}{\B{spaw}}; \HL{sAspaw}{\B{sä\-spaw}})

% tìspe'e
\hi
\HT{tIspe'e}{\B{tìspe\cp{'e}}} \I{[tI.spE."PE]} \emph{n.} capture.
(\ding{43}~\HL{spe'e}{\B{spe\-'e}})

% tìspxin
\hi
\HT{tIspxin}{\B{tì\cp{spxin}}} \I{[tI."{}sp'in]} \emph{n.} sickness, the state of being ill. 
(\ding{43}~\HL{spxin}{\B{spxin}}, \HL{sAspxin}{\B{sä\-spxin}})

% tìsraw / si / seyki
\hi
\HT{tIsraw}{\B{tì\cp{sraw}}} \I{[tI."{}sRaw]} \textsc{i}. \emph{n.} pain.
\B{Tì\-sraw le\-tokx sì tì\-ngu\-sä\-'än pxìm tä\-pa\-re fì\-tsap.} Physical pain and mental suffering are often interrelated.~\LINK{http://naviteri.org/2013/02/vospxi-ayol-posti-apup-short-post-for-a-short-month/}{b} 
\B{Tsam\-si\-yu ze\-ne tsi\-vun yi\-väkx sne\-yä tì\-sraw\-it.} Warriors must be able to ignore their own pain.~\LINK{http://naviteri.org/2015/03/kap-si-ayunil-saylahe-threats-dreams-and-other-things/}{b}  \\
\textsc{ii}. \B{tì\cp{sraw} si} \I{[tI."{}sRaw si]} \emph{vin.} hurt, be painful.
\B{Sra\-ne, skxir tì\-sraw si nì\-txan, slä ke nge\-rä\-'än oe kaw\-'it.} Yes, the wound is very painful, but I'm not in the least suffering emotionally.~\LINK{http://naviteri.org/2013/02/vospxi-ayol-posti-apup-short-post-for-a-short-month/}{b}
\B{Oe\-ri skxir a mì syokx tì\-sraw se\-ngi.} The wound on my hand hurts.~\LINK{http://naviteri.org/2011/05/some-miscellaneous-vocabulary/}{b} \\
\textsc{iii}. \B{tì\cp{sraw} sey\cp{ki}} \I{[tI."{}sRaw sEj."k$\cdot$i]} \emph{vtr., inf.}\on{22} hurt (s.o.\slash s.t.).
\B{Tì\-'efu\-mì oe\-yä, nge\-yä fì\-sä\-ka\-nom lu le\-hrr\-ap ul\-te tsun nga\-ti tì\-sraw sey\-ki\-vi.} In my opinion, this acquisition of yours is dangerous and can hurt you.~\LINK{http://naviteri.org/2012/01/more-additions-to-the-lexicon/}{b}
\B{Oe\-ri teng\-krr te\-rul mì na'\-rìng, tsun\-it a\-ski\-en tì\-sraw sey\-ko\-li.} While I was running in the forest, I hurt my right heel.~\LINK{http://naviteri.org/2022/12/trr-anawm-polaheiem-the-great-day-has-arrived/}{b} 
(\ding{43}~\HL{sraw}{\B{sraw}}, \HL{pung}{\B{pung}})

% tìsrese'a
\hi
\HT{tIsrese'a}{\B{tì\cp{sre}se'a}} \I{[tI."{}sRE.sE.Pa]} \emph{n.} prophecy. 
(\ding{43}~\HL{srese'a}{\B{sre\-se\-'a}})

% tìsteftxaw
\hi
\HT{tIsteftxaw}{\B{tìste\cp{ftxaw}}} \I{[tI.stE."ft'aw]} \emph{n.} examination.
\B{Tì\-ste\-ftxaw nge\-yä flo\-lä nìl\-tsan!} Your examination succeeded well!~\LINK{http://wiki.learnnavi.org/index.php?title=Canon\#Extracts_from_various_emails}{wiki}
(\ding{43}~\HL{steftxaw}{\B{ste\-ftxaw}})

% tìsti
\hi
\HT{tIsti}{\B{tì\cp{sti}}} \I{[tI."{}sti]} \emph{n.} anger. 
\B{Ney\-ti\-ri\-ri lu yì tì\-sti\-yä kxayl nì\-ngay.} Neytiri is really angry [lit.: Neytiri's level of anger is truly high].~\LINK{http://naviteri.org/2013/01/lifyari-po-peyi-how-good-is-her-navi/}{b}
(\ding{43}~\HL{sti}{\B{sti}})

% tìsung
\hi
\HT{tIsung}{\B{tì\cp{sung}}} \I{[tI."{}suN]} (\textsc{rn}: \B{tì\cp{sùng}} \I{[tI."sUN]}) \emph{n.} addition; post\hyp{}scriptum.
(\ding{43}~\HL{sung}{\B{sung}}, \HL{nIsung}{\B{nì\-sung}}) 

% tìswaran
\hi 
\HT{tIswaran}{\B{tì\cp{swa}ran}} \I{[tI."{s}wa.Ran]} \emph{n.} humility, humbleness.~\LINK{http://naviteri.org/2020/11/vospxivopeya-ayliu-amip-novembers-new-words/}{b}
(\ding{43}~\HL{swaran}{\B{swa\-ran}})

% tìswäper
\hi 
\HT{tIswAper}{\B{tìswä\cp{per}}} \I{[tI.sw\ae."pER]} \emph{n.} self\hyp{}praise. 
(\ding{43}~\HL{sweri}{\B{swe\-ri}})

% tìsweri
\hi 
\HT{tIsweri}{\B{tì\cp{swe}ri}} \I{[tI."{}swE.Ri]} \emph{n.} praise (general sense).
(\ding{43}~\HL{sweri}{\B{swe\-ri}})

% tìsyìmawnun'i
\hi 
\HT{tIsyImawnun'i}{\B{tìsyìmawnun\cp{'i}}} \I{[tI.sjI.maw.nun."Pi]} \emph{or} \I{[.nUn.]} (\textsc{rn}: \B{tìsyìmawnùn\cp{'i}} \I{[tI.SI.maw.nUn."Pi]}) \emph{n.} independence, freedom from a pre\hyp{}existing relationship. (\emph{often shortened and spelled} \ding{43}~\HL{tsyImawnun'i}{\B{tsyì\-maw\-nun\-'i}})
(\ding{43}~\HL{syImawnun'i}{\B{syì\-maw\-nun\-'i}}) 

% tìsyor
\hi
\HT{tIsyor}{\B{tì\cp{syor}}} \I{[tI."{}sjoR]} \emph{n.} relaxation.
\B{Krr\-a fko ta\-ron ke lu kea skxom tì\-syor\-ä.} When one hunts there's no opportunity for relaxation.~\LINK{http://naviteri.org/2013/01/awvea-posti-zisita-amip-first-post-of-the-new-year/}{b}
(\ding{43}~\HL{syor}{\B{syor}}) \\
\textsc{Derivatives}: \ding{43}
\on{1}.~\HL{tIsyortsyIp}{\B{tì\-syor\-tsyìp}}, break, small rest or relaxation.

% tìsyortsyìp
\hi 
\HT{tIsyortsyIp}{\B{tì\cp{syor}tsyìp}} \I{[tI."{}sjoR.\t{ts}jIp\textcorner]} \emph{n.} break, small rest or relaxation. 
\B{Tì\-kang\-kem so\-li oe kawl nì\-txan, 'efu ngeyn, ul\-te kin oel tì\-syor\-tsyìp\-it.} I've worked hard, I'm tired, and I need a break.~\LINK{http://naviteri.org/2018/04/100a-liu-amip-64-new-words-part-2/}{b} 
(\ding{43}~\HL{tIsyor}{\B{tì\-syor}}, \HL{-tsyIp}{\B{-tsyìp}})

% tìtaron
\hi
\HT{tItaron}{\B{tì\cp{ta}ron}} \I{[tI."ta.Ron]} \emph{n.} hunting (\emph{abstract}).
\B{Tsko swi\-zaw\-fa a tì\-ta\-ron\-ìri ze\-ne fko ni\-vu\-me.} You must learn to hunt with a bow and arrow.~\LINK{http://forum.learnnavi.org/language-updates/ordinals-nume/}{ln\slash a\textsuperscript{c}}
(\ding{43}~\HL{taron}{\B{ta\-ron}}, \HL{sAtaron}{\B{sä\-ta\-ron}})

% tìtengwotx
\hi 
\HT{tItengwotx}{\B{tì\cp{teng}wotx}} \I{[tI."tEN.wot']} \emph{n.} uniformity. 
\B{Naw\-num\-tseng\-ä ay\-nu\-me\-yu\-ri ley ay\-nga\-ru tì\-teng\-wotx, ay\-oe\-ru ke\-teng\-syon.} Regarding university students, you value uniformity, we value diversity.~\LINK{https://naviteri.org/2025/06/vospikin-lefpom-happy-july/}{b}
(\ding{43}~\HL{tengwotx}{\B{teng\-wotx}}) 

% tìterkup
\hi
\HT{tIterkup}{\B{tì\cp{ter}kup}} \I{[tI."tER.kup\textcorner]} (\textsc{rn}: \B{tì\cp{ter}kùp} \I{[tI."tER.kUp\textcorner]}) \emph{n.} death. 
(\ding{43}~\HL{terkup}{\B{ter\-kup}}, \HL{kxitx}{\B{kxitx}})

% tìtok
\hi 
\HT{tItok}{\B{tì\cp{tok}}} \I{[tI."tok\textcorner]} \emph{n.} presence. 
\B{Tì\-tok\-ìl nge\-yä ngop tì\-ngä\-zìk\-it fra\-po\-ru.} Your presence creates a difficulty for everyone.~\LINK{https://naviteri.org/2024/06/ayhapxi-tokxa-si-ayliu-alahe-parts-of-the-body-and-more/}{b}
(\emph{a. proverb}) \B{Tì\-tok slu tìk\-tok fa pam\-tsyìp a\-'aw.} Presence becomes absence by means of one little sound. [i.e. `Small things can make a big difference.']~\LINK{https://naviteri.org/2024/06/ayhapxi-tokxa-si-ayliu-alahe-parts-of-the-body-and-more/}{b}
(\ding{43}~\HL{tok}{\B{tok}}, \ding{214}~\HL{tIktok}{\B{tìktok}}) 

% tìtokat
\hi 
\HT{tItokat}{\B{tì\cp{to}kat}} \I{[tI."to.kat\textcorner]} \emph{n.} guilt. 
(\ding{43}~\HL{tokat}{\B{to\-kat}}; \ding{214}~\HL{tIlayl}{\B{tì\-layl}}) 

% tìtunu
\hi 
\HT{tItunu}{\B{tì\cp{tu}nu}} \I{[tI."tu.nu]} \emph{n.} romance. 
\B{Aw\-nga zen\-ke ti\-vung fu\-ta fì\-tì\-tu\-nu vi\-var.} We must not allow this romance to continue.~\LINK{http://naviteri.org/2021/04/mipa-ayliu-mipa-sioeykting-new-words-new-explanations/}{b} 
(\ding{43}~\HL{tunu}{\B{tu\-nu}})

% tìtxaw
\hi 
\HT{tItxaw}{\B{tì\cp{txaw}}} \I{[tI."t'aw]} \emph{n.} punishment. 
(\ding{43}~\HL{txaw}{\B{txaw}})  

% tìtxanew
\hi
\HT{tItxanew}{\B{tì\cp{txa}new}} \I{[tI."t'a.nEw]} \emph{n.} greed. 
(\ding{43}~\HL{txanew}{\B{txa\-new}})

% tìtxanro'a
\hi 
\HT{tItxanro'a}{\B{tìtxan\cp{ro}'a}} \I{[tI.t'an."Ro.Pa]} \emph{n.} fame, glory. 
\B{Tsa\-tu a new tì\-txan\-ro\-'at nì\-'aw ke sla\-yu eyk\-tan a\-txan\-tsan kaw\-krr.} The person who wants only fame will never become an excellent leader.~\LINK{https://naviteri.org/2024/06/ayhapxi-tokxa-si-ayliu-alahe-parts-of-the-body-and-more/}{b} 
(\ding{43}~\HL{txanro'a}{\B{txan\-ro\-'a}})

% tìtxantslusam
\hi 
\HT{tItxantslusam}{\B{tì\cp{txan}tslusam}} \I{[tI."t'an.\t{ts}lu.sam]} \emph{n.} wisdom (\emph{seldom used}). 
(\ding{43}~\HL{txantslusam}{\B{txan\-tslu\-sam}}, \HL{hafyon}{\B{ha\-fyon}})

% tìtxap
\hi 
\HT{tItxap}{\B{tì\cp{txap}}} \I{[tI."t'ap\textcorner]} \emph{n.} pressure (\emph{physical}). 
(\ding{43}~\HL{txap}{\B{txap}})

% tìtxeym
\hi 
\HT{tItxeym}{\B{tì\cp{txeym}}} \I{[tI."t'Ejm]} \emph{uncountable n.} interest. 
\B{Fe\-yä vur\-ìri lo\-lu txan\-a tì\-txeym.} There was great interest in their story.~\LINK{https://naviteri.org/2024/09/pxevola-liu-amip-two-dozen-new-words/}{b} 
(\ding{43}~\HL{txeym}{\B{txeym}}) \\
\textsc{Derivatives}: \ding{43}
\on{1}.~\HL{tItxeymtsim}{\B{tì\-txeym\-tsim}}, interest (\emph{countable}).

% tìtxeymtsim
\hi 
\HT{tItxeymtsim}{\B{tì\cp{txeym}tsim}} \I{[tI."t'Ejm.\t{ts}im]} \emph{countable n.} interest. 
\B{Lu oe\-ru pxay\-a tì\-txeym\-tsim.} I have a lot of interests.~\LINK{https://naviteri.org/2024/09/pxevola-liu-amip-two-dozen-new-words/\#comment-50634}{b} 
(\ding{43}~\HL{tItxeym}{\B{tì\-txey}}, \HL{tsim}{\B{tsim}}) 

% tìtxukxefu
\hi 
\HT{tItxukxefu}{\B{tìtxu\cp{kxe}fu}} \I{[tI.t'u."k'E.fu]} \emph{n.} care, concern. 
(\ding{43}~\HL{txukxefu}{\B{txu\-kxe\-fu}})  

% tìtxula
\hi
\HT{tItxula}{\B{tì\cp{txu}la}} \I{[tI."t'u.la]} (\textsc{rn}: \B{tì\cp{txù}la} \I{[tI."dU.la]}) \emph{n.} construction, constructed thing.
(\ding{43}~\HL{txula}{\B{txu\-la}}) 

% tìtxur
\hi
\HT{tItxur}{\B{tì\cp{txur}}} \I{[tI."t'uR]} \emph{n.} strength, power. 
\B{Ut\-ral\-ä A\-nawm | Ay\-ri\-na' lu ay\-oeng, | A pe\-yä tì\-txur mì hi\-nam aw\-nge\-yä.} We are all seeds | Of the Great Tree, | Whose strength is in our legs.~\LINK{http://www.pandorapedia.com/navi/music/ritual_music}{pp}
\B{Tì\-'i'\-a\-ri tsam\-ä ze\-ne Saw\-tu\-te vi\-ve\-lek ta\-lun tì\-txur Ey\-wa\-yä.} At the end of the war, the Sky People had to give up due to the power of Eywa.~\LINK{http://naviteri.org/2012/03/spring-vocabulary-part-1/}{b}
\B{Tu\-te\-ri tì\-txur lu txew\-nga'.} There are limits to a person's strength.~\LINK{http://naviteri.org/2012/03/spring-vocabulary-part-1/}{b}
\B{Ey\-wa\-ru a\-ho, ma 'i\-tan, fte Nawm\-a Sa'\-nok\-ìl tì\-ye\-vìng ngar tì\-txur\-it.} Pray to Eywa, my son, that Great Mother will give you strength.~\LINK{http://naviteri.org/2018/03/100a-liu-amip-64-new-words-part-1/}{b}
(\ding{43}~\HL{txur}{\B{txur}}) \\
\textsc{Derivatives}: \ding{43}
\on{1}.~\HL{tItxurnga'}{\B{tì\-txur\-nga'}}, powerful.

% tìtxurnga'
\hi
\HT{tItxurnga'}{\B{tì\cp{txur}nga'}} \I{[tI."t'uR.NaP]} \emph{adj., nfp.} powerful. 
\B{Po\-ri ay\-sam\-sä\-'o a\-tì\-txur\-nga' fra\-to mì hi\-fkey la\-yu mì syokx.} He will have the most powerful weapons of war in the world in his hands.~\LINK{http://naviteri.org/2016/12/tengkrr-perahem-zisit-amip-as-the-new-year-arrives/}{b}
(\ding{43}~\HL{tItxur}{\B{tì\-txur}})

% tìtxen / si
\hi
\HT{tItxen}{\B{tì\cp{txen}}} \textsc{i}. \emph{n.} awakeness, the state of waking. \\ 
\textsc{ii}. \B{tì\cp{txen} si} \I{[tI."t'En si]} \emph{vin.} to wake, to waken, wake up (\emph{with s.o.\slash s.t. in the dative}).
\B{Tì\-txen si, ma 'ite! Unil sar\-mi nga teng\-krr ze\-rawng. Lu fpom srak?} Wake up, daughter! You were dreaming and screaming. Are you okay?~\LINK{http://naviteri.org/2015/03/kap-si-ayunil-saylahe-threats-dreams-and-other-things/}{b}
\B{Sa'\-nok 'eveng\-ur tì\-txen si sre\-srr\-'ong.} Mother woke the child in the early morning.~\LINK{http://forum.learnnavi.org/language-updates/little-confirmation-titxen-si/}{ln}
(\emph{a. idiom}) \B{el\-tur tì\-txen si}, be interesting, intriguing [lit.: (s.t.) wakes the brain].   
\B{Tsa\-'u el\-tur tì\-txen si.} That's interesting.~\LINK{http://forum.learnnavi.org/language-updates/txelanit-hivawl/}{ln} 
\B{Tsa\-vur oe\-yä el\-tur tì\-txen sol\-i.} That story intrigued me.~\LINK{http://forum.learnnavi.org/language-updates/txelanit-hivawl/}{ln} 
\B{Tsa\-vur fe\-yä ay\-el\-tur tì\-txen so\-li.} They found that story interesting.~\LINK{http://forum.learnnavi.org/language-updates/txelanit-hivawl/}{ln} 
\B{Pxoe\-ri tsa\-vur el\-tur tì\-txen so\-li.} The three of us were interested in that story.~\LINK{http://forum.learnnavi.org/language-updates/txelanit-hivawl/}{ln}
\B{Lu poe tstun\-wi, ul\-te evi\-ru peng el\-tur tì\-txen si a ay\-vur\-it te\-ri 'Rr\-ta, alu tse\-nge a\-stxong nì\-ngay.} She is kind, and tells the kids interesting stories about Earth which is a truly strange place.~\LINK{http://naviteri.org/2010/07/ting-mikyun-fte-tslivam-listening-comprehension-1-3/}{b}
(\ding{43}~\HL{txen}{\B{txen}})

% tìtsan'ul
\hi
\HT{tItsan'ul}{\B{tì\cp{tsan}'ul}} \I{[tI."\t{ts}an.Pul]} \emph{or} \I{[.PUl]} \emph{n.} improvement.
\B{Tskxe\-keng\-luke ke lu kea tì\-tsan\-'ul.} Without practice there is no improvement.~\LINK{http://naviteri.org/2013/03/tsanerul-feerul-getting-better-getting-worse/}{b}
\B{Tì\-tsan\-'ul\-ìri lu nga\-ru ay\-sä\-mok srak?} Do you have any suggestions for improvement?~\LINK{http://naviteri.org/2013/03/tsanerul-feerul-getting-better-getting-worse/}{b}
(\ding{214}~\HL{tIfe'ul}{\B{tì\-fe\-'ul}}, \ding{43}~\HL{tsan'ul}{\B{tsan\-'ul}}, \HL{sAtsan'ul}{\B{sä\-tsan\-'ul}})

% tìtseri
\hi
\HT{tItseri}{\B{tì\cp{tse}ri}} \I{[tI."\t{ts}E.Ri]} \emph{n.} awareness, notice.
(\emph{a. idiom}) \B{to tì\cp{tse}ri} \I{[to tI."\t{ts}E.Ri]} than is apparent, than you are aware of. 
\B{Lu po\-e na nga nì\-'ul to tì\-tse\-ri.} She's more like you than you think (\emph{or:} than you know).~\LINK{http://naviteri.org/2018/04/100a-liu-amip-64-new-words-part-2/}{b} 
(\ding{43}~\HL{tseri}{\B{tse\-ri}})

% tìtslam
\hi
\HT{tItslam}{\B{tì\cp{tslam}}} \I{[tI."\t{ts}lam]} \emph{n.} understanding, intelligence.
\B{Slä tsal\-su\-ngay, txo tì\-tslam for, lu txa\-yo na'\-rìng\-to sìl\-tsan.} But even so, if they're smart they'll take open terrain over bush.~\LINK{http://naviteri.org/2011/09/\%E2\%80\%9Cby-the-way-what-are-you-reading\%E2\%80\%9D/}{b}
\B{Sìl\-pey oe, nge\-yä tì\-tslam\-ur srung so\-li nì\-'it.} Hopefully, (it) has helped your understanding a bit.~\LINK{http://forum.learnnavi.org/language-updates/minor-confirmation-uvan-si/}{ln}
(\ding{43}~\HL{tslam}{\B{tslam}}) 

% tìtsleng
\hi 
\HT{tItsleng}{\B{tì\cp{tsleng}}} \I{[tI."\t{ts}lEN]} \emph{n.} falsehood, falsity. 
(\ding{43}~\HL{tsleng}{\B{tsleng}}; \HL{fleyul}{\B{fley\-ul}})

% tìtstew
\hi
\HT{tItstew}{\B{tì\cp{tstew}}} \I{[tI."\t{ts}tew]} \emph{n.} courage, bravery.
\B{Le\-kin lu tì\-txur, lu tì\-tstew. Slä le\-tsran\-ten fra\-to lu tì\-mal.} Strength and courage are necessary. But most important of all is trustworthiness.~\LINK{http://naviteri.org/2012/01/more-additions-to-the-lexicon/}{b}
(\ding{43}~\HL{tstew}{\B{tstew}}) \\
\textsc{Derivatives}: \ding{43}
\on{1}.~\HL{tItstewnga'}{\B{tì\-tstew\-nga'}}, couragous, brave.

% tìtstewnga'
\hi
\HT{tItstewnga'}{\B{tì\cp{tstew}nga'}} \I{[tI."\t{ts}tew.NaP]} \emph{adj.} (\emph{of things}) courageous, brave.
\B{Po pol\-txe san tì\-tstew\-nga'\-a stxe\-nu\-ri ira\-yo sei\-yi oe nga\-ru nì\-txan.} She said, ``I thank you very much for your courageous offer.''~\LINK{http://naviteri.org/2011/09/miscellaneous-vocabulary/}{b}
\B{Kem a\-tì\-tstew\-nga' si tu\-te a\-tstew.} A brave person does brave deeds.~\LINK{http://naviteri.org/2011/09/miscellaneous-vocabulary/}{b}
(\ding{43}~\HL{tItstew}{\B{tì\-tstew}}, \HL{tstew}{\B{tstew}})

% tìtsukanom
\hi 
\HT{tItsukanom}{\B{tìtsu\cp{ka}nom}} \I{[tI.\t{ts}u."ka.nom]} (\textsc{rn}: \B{tìtsù\cp{ka}nom} \I{[tI.\t{ts}U."ka.nom]}) \emph{n.} availability. 
(\ding{43}~\HL{tsukanom}{\B{tsu\-ka\-nom}})

% tìtstunwi
\hi
\HT{tItstunwi}{\B{tì\cp{tstun}wi}} \I{[tI."\t{ts}tun.wi]} \emph{or} \I{[."\t{ts}tUn.]} \emph{n.} kindness.
\B{Tì\-tstun\-wi ay\-nge\-yä oe\-ru te\-ya so\-li.} Your kindness have filled me with joy.~\LINK{http://forum.learnnavi.org/community-calendar/karyu-pawls-birthday/msg312733/\#msg312733}{ln}
\B{Pxe\-fe\-yä tì\-tstun\-wi\-ri aw\-nga ira\-yo si\-vi ko!} Let's thank the three of them for their kindness.~\LINK{http://naviteri.org/2012/07/tskxekengtsyip-a-mikyunfpi-a-little-listening-exercise/}{b}
(\ding{43}~\HL{tstunwi}{\B{tstun\-wi}}) \\
\textsc{Derivatives}: \ding{43}
\on{1}.~\HL{tItstunwinga'}{\B{tì\-tstun\-wi\-nga'}}, kind.

% tìtstunwinga'
\hi 
\HT{tItstunwinga'}{\B{tì\cp{tstun}winga'}} \I{[tI."\t{ts}tun.wi.NaP]} \emph{or} \I{[."\t{ts}tUn.]} \emph{adj., nfp.} kind. 
\B{ay\-lì\-'u a\-tì\-tstun\-wi\-nga'}, kind words.~\LINK{http://naviteri.org/2017/09/zisikrr-amip-ayliu-amip-new-words-for-the-new-season/}{b} 
\B{Nge\-yä fay\-lì\-'u a\-tì\-tstun\-wi\-nga' oe\-ru meuia so\-li nì\-ngay.} These kind words of yours have honored me greatly.~\LINK{http://naviteri.org/2018/04/100a-liu-amip-64-new-words-part-2/}{b}
(\ding{43}~\HL{tItstunwi}{\B{tì\-tstun\-wi}}, \HL{-nga'}{\B{-nga'}}; \HL{tstunwi}{\B{tstun\-wi}} \ding{214}~\HL{tIktstunga'}{\B{tìk\-tstu\-nga'}})

% tìtsunslu
\hi
\HT{tItsunslu}{\B{tì\cp{tsun}slu}} \I{[tI."\t{ts}un.slu]} \emph{or} \I{["\t{ts}Un.]} (\textsc{rn}: \B{tì\cp{tsùn}slu} \I{[tI."\t{ts}Un.slu]}) \emph{n.} possibility.
\B{No\-lu\-me oe nì\-txan te\-ri tì\-tsun\-slu tì\-rey\-ä a hi\-fkey\-mì a\-la\-he.} I learned a lot about the possibility of life on other worlds.~\LINK{http://naviteri.org/2012/07/meetings-waterfalls-and-more/}{b}
(\ding{43}~\HL{tsunslu}{\B{tsun\-slu}}) \\
\textsc{Derivatives}: \ding{43}
\on{1}.~\HL{kxutslu}{\B{kxu\-tslu}}, risk.

% tìtsyär
\hi
\HT{tItsyAr}{\B{tì\cp{tsyär}}} \I{[tI."\t{ts}j\ae{}R]} \emph{n.} rejection. 
(\ding{214}~\HL{tImll'an}{\B{tì\-mll\-'an}}, \ding{43}~\HL{tsyAr}{\B{tsyär}})

% tìtsyul
\hi 
\HT{tItsyul}{\B{tì\cp{tsyul}}} \I{[tI."\t{ts}jul]} \emph{or} \I{[."\t{ts}jUl]} \emph{n.} beginning, start. 
\B{Sì\-tsyul nì\-wotx lu ngä\-zìk.} All beginnings are difficult.~\LINK{http://naviteri.org/2018/08/fmawnti-stolawm-srak-have-you-heard-the-news/}{b}
(\ding{43}~\HL{tsyul}{\B{tsyul}})

% tìvawm
\hi
\HT{tIvawm}{\B{tì\cp{vawm}}} \I{[tI."vawm]} \emph{n.} darkness. 
\B{Lar\-mu tsa\-tsam\-si\-yu\-hu tì\-vawm a lu stxong ay\-oer.} That warrior carried with him a darkness unknown to us.~\LINK{http://naviteri.org/2010/07/vocabulary-update/}{b}
(\ding{43}~\HL{vawm}{\B{vawm}})

% tìväng
\hi
\HT{tIvAng}{\B{tì\cp{väng}}} \I{[tI."v\ae{}N]} \emph{n.} thirst.
\B{Apxa tì\-väng\-ìl po\-ti stey\-ko\-li.} (His) great thirst made him angry.~\LINK{http://naviteri.org/2011/01/new-year-new-vocabulary/}{b}
(\ding{43}~\HL{vAng}{\B{väng}})

% tìve'kì
\hi
\HT{tIve'kI}{\B{tìve'\cp{kì}}} \I{[tI.vEP."kI]} \emph{n.} hatred. 
(\ding{214}~\HL{tIyawn}{\B{tì\-yawn}}, \ding{43}~\HL{ve'kI}{\B{ve'\-kì}})

% tìvirä
\hi 
\HT{tIvirA}{\B{tìvi\cp{rä}}} \I{[tI.vi."r\ae]} \emph{n.} spread, proliferation. 
\B{Ke\-zem\-pll\-txe, ta\-lun tì\-vi\-rä fì\-sä\-spxin\-ä, lo\-la\-tä\-ngem ki\-fkey, lo\-la\-tä\-ngem tì\-rey.} Needless to say, because of the spread of this illness, the world and life changes.~\LINK{http://naviteri.org/2020/03/teri-tifkeytok-lefkrr-about-the-present-situation/}{b} 
(\ding{43}~\HL{virA}{\B{virä}})

% tìwan
\hi
\HT{tIwan}{\B{tì\cp{wan}}} \I{[tI."wan]} \emph{n.} obfuscation, cover\hyp{}up.
(\ding{43}~\HL{wan}{\B{wan}}) \\
\textsc{Derivatives}: \ding{43}
\on{1}.~\HL{letwan}{\B{let\-wan}}, dodgy, sneaky

% tìwäsul
\hi
\HT{tIwAsul}{\B{tì\cp{wä}sul}} \I{[tI."w\ae.sul]} \emph{or} \I{[.sUl]}  \emph{n.} competition, the abstract idea of competition.
(\ding{43}~\HL{wAsul}{\B{wä\-sul}}, \HL{sAwAsul}{\B{sä\-wä\-sul}})

% tìwäte
\hi
\HT{tIwAte}{\B{tìwä\cp{te}}} \I{[tI.w\ae."tE]} \emph{n.} dispute, argument.
\B{Tì\-wä\-te\-ri ngal oe\-ti pe\-lun ke slan kaw\-krr?} Why don't you ever support me in an argument?~\LINK{http://naviteri.org/2012/03/spring-vocabulary-part-1/}{b}
(\ding{43}~\HL{wAte}{\B{wä\-te}}, \HL{sAwAte}{\B{sä\-wä\-te}}) \\
\textsc{Derivatives}: \ding{43}
\on{1}.~\HL{tIwAtenga'}{\B{tì\-wä\-te\-nga'}}, controversial.

% tìwätenga'
\hi 
\HT{tIwAtenga'}{\B{tìwä\cp{te}nga'}} \I{[tI.w\ae."tE.NaP]} \emph{adj., nfp.} controversial. 
\B{Fì\-puk\-mì a ay\-sä\-fpìl lu tì\-wä\-te\-nga' nì\-txan.} The ideas in this book are very controversial.~\LINK{https://naviteri.org/2024/09/pxevola-liu-amip-two-dozen-new-words/}{b} 
(\ding{43}~\HL{tIwAte}{\B{tì\-wä\-te}})

% tìwew
\hi
\HT{tIwew}{\B{tì\cp{wew}}} \I{[tI."wEw]} \emph{n.} coldness, cold.
(\ding{214}~\HL{tIsom}{\B{tì\-som}}, \ding{43}~\HL{wew}{\B{wew}})

% tìwesek
\hi 
\HT{tIwesek}{\B{tì\cp{we}sek}} \I{[tI."wE.sEk\textcorner]} \emph{n.} subtlety.
\B{Ke lo\-lu kea tì\-we\-sek mì sä\-ftxu\-lì\-'u a\-räp\-tum pe\-yä.} There was no subtlety in his coarse speech.~\LINK{https://naviteri.org/2026/04/tsivola-liu-amip-thirty-two-new-words/}{b} 
(\ding{43}~\HL{wesek}{\B{we\-sek}})

% tìwìngay
\hi
\HT{tIwIngay}{\B{tìwì\cp{ngay}}} \I{[tI.wI."Naj]} \emph{n.} proof, proving (\emph{abstract}).
\B{Txo new ngal fu\-ta su\-tel nge\-yä ay\-lì\-'u\-ti spi\-vaw, tsran\-ten tì\-wì\-ngay.} If you want people to believe you, proof is important.~\LINK{http://naviteri.org/2015/04/some-new-words-for-may-day/}{b}
(\ding{43}~\HL{wIngay}{\B{wì\-ngay}}, \HL{sAwIngay}{\B{sä\-wì\-ngay}})

% tìyawn
\hi
\HT{tIyawn}{\B{tì\cp{yawn}}} \I{[tI."jawn]} \emph{n.} love.   
\B{Spaw Na'\-vil fu\-ta tì\-yawn Ey\-wa\-yä lu txew\-lu\-ke.} The Na'vi believe that Eywa's love is boundless.~\LINK{http://naviteri.org/2014/11/twenty-before-the-holidays/}{b} 
\B{Lì'\-fya\-ri le\-Na'\-vi oel 'e\-fu ay\-nge\-yä tì\-yawn\-it.} I feel your love for the Na'vi language.~\LINK{http://forum.learnnavi.org/news-announcements/a-response-from-paul-frommer!/}{ln} 
\B{Nge\-yä ke ftang a tì\-kang\-kem\-ìri sì tì\-ka\-nu\-ri sì fì\-lì'\-fya\-yä tì\-yawn\-ìri sei\-yi oe ira\-yo nì\-txan.} I thank you very much for your never\hyp{}ending work, intelligence and love for the language.~\LINK{http://naviteri.org/2012/07/meetings-waterfalls-and-more/comment-page-1/\#comment-1547}{b}
(\emph{a. idiom}) \B{tì\-yewn tì\-yawn\-ä} `an expression of love'~\LINK{http://naviteri.org/2020/11/vospxivopeya-ayliu-amip-novembers-new-words/}{b}
(\ding{214}~\HL{tIve'kI}{\B{tì\-ve'\-kì}}, \ding{43}~\HL{yawne}{\B{yaw\-ne}})

% tìyawr
\hi
\HT{tIyawr}{\B{tì\cp{yawr}}} \I{[tI."jawR]} \emph{n.} correctness. 
(\emph{a. idiom}) \B{Nga\-ru tì\-yawr.} You're right. 
\B{Nga\-ru tì\-yawr nì\-wotx!} You are completely right.~\LINK{http://naviteri.org/2010/08/a-na\%E2\%80\%99vi-alphabet/comment-page-1/\#comment-191}{b} 
\B{Nì\-ngay lu ngar tì\-yawr, lu ngar tì\-kxey.} Truly, you are both right and wrong.~\LINK{http://naviteri.org/2010/07/vocabulary-update/comment-page-1/\#comment-95}{b} 
\B{Oe\-ru tì\-yawr srak?} Am I right?~\LINK{http://naviteri.org/2012/03/spring-vocabulary-part-1/comment-page-1/\#comment-1452}{b}
(\ding{214}~\HL{tIkxey}{\B{tì\-kxey}}, \ding{43}~\HL{eyawr}{\B{eyawr}})

% tìyäkx
\hi
\HT{tIyAkx}{\B{tì\cp{yäkx}}} \I{[tI."j\ae{}k']} \emph{n.} lack of notice; snubbing.
\B{Tì\-yäkx ke lu sru\-nga', ma tsmuk. Nga txo sti, oeyk\-tìng teyng\-ta pe\-lun.} Snubbing isn't helpful, brother. If you're angry, explain why.~\LINK{http://naviteri.org/2015/03/kap-si-ayunil-saylahe-threats-dreams-and-other-things/}{b} 
(\ding{43}~\HL{yAkx}{\B{yäkx}}, \HL{sAyAkx}{\B{sä\-yäkx}})

% tìyewn
\hi 
\HT{tIyewn}{\B{tì\cp{yewn}}} \I{[tI."jEwn]} \emph{n.} expression (\emph{of a thought or feeling}). 
(\emph{a. idiom}) \B{tì\-yewn tì\-yawn\-ä} `an expression of love'~\LINK{http://naviteri.org/2020/11/vospxivopeya-ayliu-amip-novembers-new-words/}{b}
(\ding{43}~\HL{yewn}{\B{yewn}}, \HL{lI'fyavi}{\B{lì'\-fya\-vi}})

% tìyo'
\hi
\HT{tIyo'}{\B{tì\cp{yo'}}} \I{[tI."joP]} \emph{n.} perfection.
\B{Fì\-tseo\-ri ke tsun kaw\-tu pi\-vä\-hem tì\-yo'\-ne; tsran\-ten tì\-pä\-hem\-ä tì\-fmi nì\-'aw.} In this art it's impossible to arrive at perfection; the only thing that matters is the attempt to arrive there.~\LINK{http://naviteri.org/2012/01/mipa-zisit-ayliu-amip-new-words-for-the-new-year/}{b}
(\ding{43}~\HL{yo'}{\B{yo'}})

% tìyora'
\hi
\HT{tIyora'}{\B{tìyo\cp{ra'}}} \I{[tI.jo."RaP]} \emph{n.} victory, a win. 
\B{Ul\-te kxawm zì\-ye\-va\-'u nga\-ne tì\-yo\-ra' zì\-sìt\-ay!} And maybe victory might come to you next year!~\LINK{http://naviteri.org/2011/12/ulte-yora\%E2\%80\%99tu-leiu-and-the-winner-is/}{b}
\B{Tsam\-si\-yu\-ri lu tì\-yo\-ra'\-ä tì\-new le\-kin.} A warrior must have the desire for victory.~\LINK{http://naviteri.org/2014/11/twenty-before-the-holidays/}{b}
(\ding{214}~\HL{tIsnaytx}{\B{tì\-snaytx}}, \ding{43}~\HL{yora'}{\B{yo\-ra'}})

% tìzevakx
\hi 
\HT{tIzevakx}{\B{tì\cp{ze}vakx}} \I{[tI."zE.vak']} \emph{n.} cruelty. 
(\ding{43}~\HL{zevakx}{\B{zevakx}}) \\
\textsc{Derivatives}: \ding{43}
\on{1}.~\HL{tIzevakxnga'}{\B{tì\-ze\-vakx\-nga'}}, cruel.

% tìzevakxnga'
\hi 
\HT{tIzevakxnga'}{\B{tì\cp{ze}vakxnga'}} \I{[tI."zE.vak'.NaP]} \emph{adj., nfp.} cruel. 
\B{ay\-lì\-'u a\-tì\-ze\-vakx\-nga'}, cruel words.~\LINK{http://naviteri.org/2020/06/mi-tanlokxe-oeya-srr-afpxamo-terrible-days-in-my-country/}{b}
(\ding{43}~\HL{zevakx}{\B{zevakx}}; \HL{tIzevakx}{\B{tì\-ze\-vakx}}, \HL{-nga'}{\B{-nga'}})

% tìzin (si)
\hi 
\HT{tIzin}{\B{tì\cp{zin}}} \I{[tI."zin]} \textsc{i}. \emph{n.} a tangle, tanglement; \emph{mass of s.t. twisted together}. 
\B{Fwa fpìl nì\-fya\-'o a\-ve\-nga' fì\-tì\-zin\-mì lu kel\-tsun.} It's impossible to think clearly (i.e., in an organized manner) in this muddled situation.~\LINK{http://naviteri.org/2021/04/mipa-ayliu-mipa-sioeykting-new-words-new-explanations/\#comment-33483}{b}
(\ding{43}~\HL{zin}{\B{zin}}) \\
\textsc{ii}. \B{tì\cp{zin} si} \I{[tI."zin si]} \emph{vin.} tangle, tangle up.
\B{Na\-ri si fte\-ke ay\-tur\-tel\-ur tì\-zin si\-vi!} Be careful not to tangle the ropes.~\LINK{http://naviteri.org/2021/04/mipa-ayliu-mipa-sioeykting-new-words-new-explanations/}{b} 

% to
\hi
\HT{to}{\B{to}} \I{[to]} \emph{part.} (\emph{comparative marker}) than.   
\B{Po to oe lu sìl\-tsan.} S\slash he is better than me.~\LINK{http://forum.learnnavi.org/language-updates/prrwll-moss-the-comparative-to-construct/}{ln}
\B{Nga nì\-aw\-no\-mum to oe\-tsyìp lu txur nì\-txan.} As everyone knows, you're a lot stronger than little old me.~\LINK{http://naviteri.org/2010/07/diminutives-conversational-expressions/}{b}
\B{Oe to nga tul nì\-win.} I run faster than you.~\LINK{http://forum.learnnavi.org/language-updates/confirmations-\%28comparisons-gerund\%29/}{ln}
\B{Oel to ngal ye\-rik\-it ta\-ron nìl\-tsan.} I hunt yerik better than you.~\LINK{http://forum.learnnavi.org/language-updates/confirmations-\%28comparisons-gerund\%29/}{ln}
\B{Oel pay\-ti nul\-new to swo\-at.} I prefer water to alcohol.~\LINK{http://forum.learnnavi.org/language-updates/prefer-x-to-y-24474/}{ln}
(\emph{can be attached as a suffix}) \B{Oe nga\-to yom nì\-'ul.} I eat more than you.~\LINK{http://forum.learnnavi.org/language-updates/prefer-x-to-y-24474/}{ln}
\B{Po\-to oe lu ko\-ak.} I am older than he.~\LINK{http://naviteri.org/2010/07/thoughts-on-ambiguity/}{b}
\B{Oe\-ri tsaw ke ley fì\-'u\-to.} I don't value that over this.~\LINK{http://naviteri.org/2014/03/value-and-worth/}{b}
(\emph{can refer to whole sentences}) \B{Ftu\-e lu fwa ta\-ron ngong\-a io\-ang\-it to fwa ta\-ron pum\-it a lu wa\-lak sì win.} It's easier to hunt lethargic animals than to hunt perky, speedy ones.~\LINK{http://naviteri.org/2011/10/more-vocabulary-a-bit-of-grammar/}{b}
(\emph{a. idiom}) \B{to tì\cp{tse}ri} \I{[to tI."\t{ts}E.Ri]} than is apparent, than you are aware of. 
\B{Lu po\-e na nga nì\-'ul to tì\-tse\-ri.} She's more like you than you think (\emph{or:} than you know).~\LINK{http://naviteri.org/2018/04/100a-liu-amip-64-new-words-part-2/}{b} 
(\ding{43}~\HL{frato}{\B{fra\-to}})

% Toitslan
% \hi
% \B{\cp{To}itslan} \I{["to.i.\t{ts}lan]} \emph{n.}, \textsc{lw} Germany.

% toitsye
% \hi
% \B{\cp{To}itsye} \I{"to.i.\t{ts}jE} \emph{adj.} German, German language.

% tok
\hi
\HT{tok}{\B{tok}} \I{[tok\textcorner]} \emph{vtr.} (\emph{a. local}) be at, occupy a space. 
\B{Nì\-'i'\-a tok oel nì\-mun fì\-tse\-nget!} Finally, I'm here again!~\LINK{http://naviteri.org/2011/07/txantsana-ultxa-mi-siatll-great-meeting-in-seattle/}{b}
\B{Pep\-sul tok fì\-tse\-nget?} What three people are here?~\LINK{http://naviteri.org/2011/07/number-in-na\%E2\%80\%99vi/}{b}
\B{Aw\-ngal tok kel\-kut.} We're home.~\LINK{http://naviteri.org/2011/10/more-vocabulary-a-bit-of-grammar/}{b}
\B{Tsìl\-pey\-ìl tok txe'\-lan\-it.} Hope lives within the heart.~\LINK{http://naviteri.org/2013/01/awvea-posti-zisita-amip-first-post-of-the-new-year/}{b}
\B{Li pol te\-rok fì\-tseng\-it srak?} Is she already here?~\LINK{http://naviteri.org/2011/02/new-vocabulary-part-2/}{b}
\B{Oel hu Txe\-wì trr\-am na'\-rìng\-it tar\-mok.} Yesterday I was with Txewì in the forest.~\LINK{http://wiki.learnnavi.org/index.php?title=Corpus\#New_York_Times}{wiki}
(\emph{b. can be omitted in casual speech when easily understood})
\B{Pol fì\-tse\-nget.} He's here.~\LINK{http://naviteri.org/2020/02/some-words-for-leap-year-day/}{b}
\B{Pol ke fì\-tse\-nget.} He's not here.~\LINK{http://naviteri.org/2020/02/some-words-for-leap-year-day/}{b}
\B{Po\-ru ke po\-leng oel ke\-'ut mung\-wrr fwa Ra\-lul ke tsa\-tse\-nget.} I told her nothing except that Ralu wasn't there.~\LINK{http://naviteri.org/2020/02/some-words-for-leap-year-day/}{b}
(\emph{c. metaphorically with} \B{yì}) 
\B{Lì'\-fya\-ri pol tok pe\-yìt?} How good is her Na'vi?~\LINK{http://naviteri.org/2013/01/lifyari-po-peyi-how-good-is-her-navi/}{b}
\B{Fu\-ri\-a tä\-ftxu ngal tok yì\-pet? -- Tok yìt a\-ke\-sran.} How's your weaving? -- I'm so\hyp{}so.~\LINK{http://naviteri.org/2013/01/lifyari-po-peyi-how-good-is-her-navi/}{b}
(\emph{d. idiom}) \B{Nga\-ri ke\-ftxo (fwa) ke tok.} ``We missed you.\slash Sorry you couldn't make it.\slash Too bad you couldn't be there.''~\LINK{http://naviteri.org/2019/08/aawa-lifya-si-lifyavi-amip-a-few-new-words-and-expressions/}{b} \\
\textsc{Derivatives}: \ding{43}
\on{1}.~\HL{fkeytok}{\B{fkey\-tok}}, exist. 
\on{2}.~\HL{fewtusok}{\B{few\-tu\-sok}}, opposite.
\on{3}.~\HL{tItok}{\B{tì\-tok}}, presence.
\on{4}.~\HL{pxawtok}{\B{pxaw\-tok}}, surround.

% tokat
\hi 
\HT{tokat}{\B{\cp{to}kat}} \I{["to.kat\textcorner]} \emph{adj.} guilty. 
(\ding{214}~\HL{layl}{\B{layl}}) \\
\textsc{Derivatives}: \ding{43} 
\on{1}.~\HL{tItokat}{\B{tì\-to\-kat}}, guilt. 
\on{2}.~\HL{nItokat}{\B{nì\-to\-kat}}, guiltily. 

% toktor
\hi
\HT{toktor}{\B{\cp{tok}tor}} \I{["tok\textcorner.toR]} \emph{n.}, \textsc{lw} doctor (\emph{academic degree}). 
\B{Lo\-lu tok\-tor Kì\-rey\-sì kar\-yu a\-sìl\-tsan.} Doctor Grace was a good teacher.~\LINK{http://naviteri.org/2010/07/ting-mikyun-fte-tslivam-listening-comprehension-1-3/}{b} 

% tokx
\hi
\HT{tokx}{\B{tokx}} \I{[tok']} \emph{n.} body.
\B{Nga\-ri hu Ey\-wa sa\-lew ti\-rea, tokx 'ì\-'awn, slu Na'\-vi\-yä ha\-pxì.} Your spirit goes with Eywa, your body stays behind to become part of the People.~\LINK{http://wiki.learnnavi.org/index.php?title=Na\%27vi_from_Avatar_Movie\#Jake.27s_final_kill}{a\slash wiki} 
\B{Fì\-po lu vrr\-tep a mì sokx a\-tsleng!} This one is a demon in a false body!~\LINK{http://wiki.learnnavi.org/index.php?title=Na\%27vi_from_Avatar_Movie\#Tsu.27tey_attacks_Jake}{wiki\slash a}
(\emph{a. idiom}) \B{tokx eo tokx}, face to face, in person.~\LINK{http://forum.learnnavi.org/language-updates/fmawno/}{ln} \\
\textsc{Derivatives}: \ding{43}
\on{1}.~\HL{letokx}{\B{le\-tokx}}, bodily, physical. 
\on{2}.~\HL{fpomtokx}{\B{fpom\-tokx}}, (physical) health. 
\on{3}.~\HL{uniltIrantokx}{\B{u\-nil\-tì\-ran\-tokx}}, avatar, ``dreamwalker body''. 
\on{4}.~\HL{yemstokx}{\B{yem\-stokx}}, put on (clothing). 

% toltem
\hi
\HT{toltem}{\B{tol\cp{tem}}} \I{[tol."tEm]} \emph{vtr.} shoot (s.o.\slash s.t.).
\B{Pll\-txe Ra\-lu san oe new ti\-vol\-tem ye\-rik\-it.} Ralu says he wants to shoot a hexapede.~\LINK{http://naviteri.org/2017/03/titusemteri-concerning-shooting/}{b}
(\ding{43}~\HL{tem}{\B{tem}}; \HL{tsweykayon}{\B{tsway\-on} (\emph{b. \ding{246}})})

% tompa
\hi
\HT{tompa}{\B{\cp{tom}pa}} \I{["tom.pa]} \emph{n.} rain.   
\B{Tom\-pa\-yä ka\-to | Tsaw\-ke\-yä ka\-to.} The rhythm of the rain | The rhythm of the sun.~\LINK{http://wiki.learnnavi.org/index.php/Weaving_Song}{wiki}
\B{Ye'\-rìn sngä\-'i tom\-pa.} Soon the rain starts.~\LINK{http://forum.learnnavi.org/language-updates/learnings-from-the-siatlla-song-translations/msg475008/\#msg475008}{ln}
(\emph{a. precipitation, with} \B{zup})
\B{Ze\-rup tom\-pa.} It's raining.~\LINK{http://naviteri.org/2011/04/yafkeykiri-plltxe-frapo-everyone-talks-about-the-weather/}{b}
\B{Nìk\-mar zup tom\-pa set.} Now is a good time for it to be raining.~\LINK{http://naviteri.org/2011/10/more-vocabulary-a-bit-of-grammar/}{b}
(\emph{with} \B{'eko}) \B{Tom\-pa 'e\-ko nì\-hawng, ha ze\-ne aw\-nga kll\-wi\-vo.} The rain is too strong, so we must land.~\LINK{http://naviteri.org/2012/01/mipa-zisit-ayliu-amip-new-words-for-the-new-year/}{b} \\
\textsc{Derivation}: \ding{43}
\on{1}.~\HL{tompameyp}{\B{tom\-pa\-meyp}}, drizzle. 
\on{2}.~\HL{tompawll}{\B{tom\-pa\-wll}}, geode. 
\on{3}.~\HL{txantompa}{\B{txan\-tom\-pa}}, rainstorm. 
\on{4}.~\HL{tomperwI}{\B{tom\-per\-wì}}, sleet. 
\on{5}.~\HL{tompIva}{\B{tom\-pì\-va}}, raindrop. 
\on{6}.~\HL{tompakel}{\B{tom\-pa\-kel}}, drought.
\on{7}.~\HL{tompatayng}{\B{tom\-pa\-tayng}}, rain thistle.

% tompakel
\hi
\HT{tompakel}{\B{\cp{tom}pakel}} \I{["tom.pa.kEl]} \emph{n.} drought.   
\B{Tom\-pa\-kel\-ta\-lun ze\-ne tu\-te Kä\-lì\-for\-ni\-a\-yä pay\-it si\-var nì\-nän.} Because of the drought, Californians have to use less water.~\LINK{http://naviteri.org/2015/04/some-new-words-for-may-day/}{b}
(\ding{43}~\HL{tompa}{\B{tom\-pa}}, \HL{tIkelu}{\B{tì\-ke\-lu}})

% tompameyp
\hi
\HT{tompameyp}{\B{tompa\cp{meyp}}} \I{[tom.pa."mEjp\textcorner]} \emph{n.} drizzle.
\B{Fì\-re\-won tom\-pa\-meyp zar\-mup, slä set 'ì\-me\-ko nì\-txan nang.} It was drizzling this morning, but it's really started coming down now!~\LINK{http://naviteri.org/2011/04/yafkeykiri-plltxe-frapo-everyone-talks-about-the-weather/}{b}
(\ding{43}~\HL{tompa}{\B{tom\-pa}}, \HL{meyp}{\B{meyp}})

% tompatayng
\hi 
\HT{tompatayng}{\B{\cp{tom}patayng}} \I{["tom.pa.tajN]} \emph{n.} \ding{96} rain thistle. 
(\ding{43}~\HL{tompa}{\B{tompa}}, \HL{tayng}{\B{tayng}})

% tompawll
\hi
\HT{tompawll}{\B{\cp{tom}pawll}} \I{["tom.pa.w\s{l}]} \emph{n.}, \ding{96} geode, ``rain plant'', \emph{tubineus azureus}.
(\ding{43}~\HL{tompa}{\B{tom\-pa}}, \HL{'ewll}{\B{'ewll}})

% tomperwì
\hi
\HT{tomperwI}{\B{\cp{tom}perwì}} \I{["tom.pER.wI]} \emph{n.} sleet.
(\ding{43}~\HL{tompa}{\B{tom\-pa}}, \HL{herwI}{\B{her\-wì}})

% tompìva
\hi
\HT{tompIva}{\B{\cp{tom}pìva}} \I{["tom.pI.va]} \emph{n.} raindrop.
\B{Oe\-ri ay\-som\-pì\-va sìn re\-'o var zi\-vup.} Raindrops keep falling on my head.~\LINK{http://naviteri.org/2011/04/yafkeykiri-plltxe-frapo-everyone-talks-about-the-weather/}{b}
(\ding{43}~\HL{tompa}{\B{tom\-pa}}, \HL{Ilva}{\B{ìl\-va}})

% tonowari
\hi
\HT{Tonowari}{\B{Tono\cp{wa}ri}} \I{[to.no."wa.ri]} \emph{n. male name}.

% tong
\hi 
\HT{tong}{\B{tong}} \I{[toN]} \emph{vtr.} put out, quench, exstinguish (\emph{of a fire}). 
\B{Maw\-krr\-a ngal txep\-it to\-long tsun hi\-vum.} After you've put out the fire you can leave.~\LINK{http://naviteri.org/2018/03/100a-liu-amip-64-new-words-part-1/}{b}
\B{Tong txep\-tsyìp\-it.} Extinguish the flame.~\LINK{http://naviteri.org/2020/07/vospxivol-lefpom-happy-august/}{b}
\B{Lo\-nu\-sye tong txep\-tsyìp\-it.} Blow out the flame.~\LINK{http://naviteri.org/2020/07/vospxivol-lefpom-happy-august/}{b}

% tor
\hi
\HT{tor}{\B{tor}} \I{[toR]} \emph{adj.} last, ultimate, terminal, bring about finality.
\B{tì\-pawm a\-tor}, ultimate question.~\LINK{http://naviteri.org/2010/08/mipa-ayopin-mipa-ayliu-new-colors-new-words/}{b} 
(\ding{43}~\HL{syen}{\B{syen}}) \\
\textsc{Derivatives}: \ding{43}
\on{1}.~\HL{toruk}{\B{to\-ruk}}, great leonopteryx.

% toruk
\hi
\HT{toruk}{\B{\cp{to}ruk}} \I{["to.Ruk\textcorner]} \emph{or} \I{[.RUk\textcorner]} (\textsc{rn}: \B{\cp{to}rùk} \I{["to.RUk\textcorner]}) \emph{n.} (\emph{a. fauna}) great leonopteryx, ``last shadow'', \emph{leonopteryx rex}. 
\B{Tseyk tswa\-may\-on fa ik\-ran sre\-krr; ta\-fral fmo\-li fì\-kem si\-vi fa to\-ruk nì\-steng.} Jake flew with an ikran before; therefore he tried to do it with a \emph{toruk} in a similar fashion.~\LINK{http://naviteri.org/2011/09/miscellaneous-vocabulary/}{b}
\B{ean na tsuk\-sìm to\-ruk\-ä}, blue like the chin of a great leonopteryx.~\LINK{http://naviteri.org/2010/08/mipa-ayopin-mipa-ayliu-new-colors-new-words/}{b}
(\emph{b. idiom})
\B{Ke tsun fko tspi\-vang to\-ruk\-it fa fwa pewn\-ti snew.} \emph{Proverbial expression for a method that will not work.} [lit.: You can't kill a great leonopteryx by constricting its throat]~\LINK{http://naviteri.org/2012/10/snewsyea-ftxoza-halowina-a-spooky-halloween/}{b} 
(\emph{shortened to}) \B{pewn to\-ruk\-ä}. \B{Sra\-ke fpìl ay\-ngal fu\-ta aw\-nga tsun hu Saw\-tu\-te a tsam\-it 'i\-vaw\-nìm fa fwa päng\-kxo fo\-hu nì\-'aw, ma smuk? Pewn to\-ruk\-ä!} Do you think we can avoid a war with the Sky People solely by talking with them, my brothers? That won't work!~\LINK{http://naviteri.org/2012/10/snewsyea-ftxoza-halowina-a-spooky-halloween/comment-page-1/\#comment-1755}{b} \\
(i.) \B{\cp{to}ruk \cp{mak}to} \I{["to.Ruk\textcorner{} "mak\textcorner.to]} \emph{n.} Rider of the Last Shadow (\emph{legendary hero of the Na'vi}).
\B{To\-ruk Mak\-to po\-lä\-hem a fì\-'u ro\-lo\-'a nì\-txan Oma\-ti\-ka\-ya\-ru.} The arrival of Toruk Makto made a great impression on the Omaticaya.~\LINK{http://naviteri.org/2012/03/spring-vocabulary-part-1/}{b}
(\ding{43}~\HL{tor}{\B{tor}}, \HL{uk}{\B{uk}}) \\
\textsc{Derivatives}: \ding{43}
\on{1}.~\HL{torukspxam}{\B{to\-ruk\-spxam}}, octoshroom. 
\on{2}.~\HL{nawmtoruktek}{\B{nawm\-to\-ruk\-tek}}, Toruk Makto Totem.

% torukspxam
\hi
\HT{torukspxam}{\B{\cp{to}rukspxam}} \I{["to.Ruk\textcorner.sp'am]} \emph{or} \I{[.RUk\textcorner.]} (\textsc{rn}: \B{\cp{to}rùkspxam} \I{["to.RUk\textcorner.sp'am]}) \emph{n., funga} octoshroom, ``great leonopteryx fungus'', \emph{fungimonium giganteum}.
(\ding{43}~\HL{toruk}{\B{to\-ruk}}, \HL{spxam}{\B{spxam}})

% trr
\hi
\HT{trr}{\B{trr}} \I{[t\s{r}]} \emph{n.} day.   
\B{Tom\-pa\-yä ka\-to | Tsaw\-ke\-yä ka\-to | Trr\-ä sì txon\-ä.} The rhythm of the rain | The rhythm of the sun | Of day and night.~\LINK{http://wiki.learnnavi.org/index.php/Weaving_Song}{wiki}
\B{Tskot sngo\-lä\-'i po si\-var kam trr 'a\-'aw\-a.} He started to use the bow several days ago.~\LINK{http://naviteri.org/2011/09/miscellaneous-vocabulary/}{b}
\B{Fì\-trr lu trr\-pe\-ve?} What day is it today?~\LINK{http://forum.learnnavi.org/language-updates/trr-rrtaya/}{ln}
\B{Ay\-nge\-yä ay\-srr tì\-flä\-ta te\-ya li\-vu ul\-te si\-va\-lew nì\-'o' nì\-'aw.} May your days be full of success and proceed only in a fun way.~\LINK{http://naviteri.org/2013/12/mipa-zisit-lefpom-happy-new-year/}{b}
\B{Fì\-tì\-kang\-kem\-vi\-ri oel kin 'aw\-a trr\-ti nì\-'aw.} For this project I only need one day.~\LINK{http://naviteri.org/2014/09/stxeli-alor-the-text/}{b}
(\emph{idiom}) \B{trr le\-fpom}, good day, good morning (cf. French \emph{bon jour}).
\B{Trr le\-fpom, ma Su\-te Ame\-ri\-ka\-yä.} Good morning, America.~\LINK{http://wiki.learnnavi.org/index.php?title=Corpus\#Good_Morning_America}{wiki} \\
(i.) \HT{trram}{\B{trr\cp{am}}} \I{[t\s{r}."{}am]} \emph{adv., n.} yesterday. 
\B{Oel hu Txe\-wì trr\-am na'\-rìng\-it tar\-mok, tso\-le\-'a syep\-tu\-tet a\-tsawl fra\-to mì sì\-rey.} Yesterday I was with Txewì in the forest and we saw the biggest Trapper I've ever seen.~\LINK{http://wiki.learnnavi.org/index.php?title=Corpus\#New_York_Times}{wiki}
\B{Trr\-am\-ä ay\-u\-van\-ìri mak\-to Ak\-wey nì\-yew\-la ha snaytx.} Akwey rode disappointingly in yesterday's games so he lost.~\LINK{http://naviteri.org/2012/10/mipa-vospxi-mipa-ayliu-new-words-for-the-new-month/}{b} \\
(ii.) \HT{trray}{\B{trr\cp{ay}}} \I{[t\s{r}."{}aj]} \emph{adv., n.} tomorrow. 
\B{Pam\-rel sa\-yi trr\-ay krr a skxom la\-tsu oe\-ru.} I will write tomorrow when I have the chance.~\LINK{http://forum.learnnavi.org/language-updates/just-a-few-interesting-little-words/}{ln} 
\B{Trr\-ay\-ri li\-vu nì\-sìl\-pey ay\-lrr\-tok a\-txan ta Ey\-wa aw\-nga\-ru nì\-wotx.} As for tomorrow, hopefully much smiles from Eywa to you all.~\LINK{http://forum.learnnavi.org/language-updates/luke-kxu-without-harm/msg132308/\#msg132308}{ln} \\
\textsc{Derivatives}: \ding{43}
\on{1}.~\HL{fratrr}{\B{fra\-trr}}, every day. 
\on{2}.~\HL{fItrr}{\B{fì\-trr}}, today. 
\on{3}.~\HL{letrr}{\B{le\-trr}}, daily. 
\on{4}.~\HL{nItrrtrr}{\B{nì\-trr\-trr}}, ordinarily. 
\on{5}.~\HL{trrtxon}{\B{trr\-txon}}, day and night cycle. 
\on{6}.~\HL{trr'ong}{\B{trr\-'ong}}, dawn, sunrise. 
\on{7}.~\HL{kxamtrr}{\B{kxam\-trr}}, midday, noon. 
\on{8}.~\HL{mrrtrr}{\B{mrr\-trr}}, (five\hyp{}day) work week. 
\on{9}.~\HL{kintrr}{\B{kin\-trr}}, seven\hyp{}day week. 
\on{10}.~\HL{muntrr}{\B{mun\-trr}}, weekend. 
\on{11}.~\HL{trr'awve}{\B{Trr\-'aw\-ve}}, Sunday. 
\on{12}.~\HL{trrmuve}{\B{Trr\-mu\-ve}}, Monday. 
\on{13}.~\HL{trrpxeyve}{\B{Trr\-pxey\-ve}}, Tuesday. 
\on{14}.~\HL{trrtsIve}{\B{Trr\-tsì\-ve}}, Wednesday. 
\on{15}.~\HL{trrmrrve}{\B{Trr\-mrr\-ve}}, Thursday. 
\on{16}.~\HL{trrpuve}{\B{Trr\-pu\-ve}}, Friday. 
\on{17}.~\HL{trrkive}{\B{Trr\-ki\-ve}}, Saturday. 
\on{18}.~\HL{zIsItsaltrr}{\B{zì\-sìt\-sal\-trr}}, (yearly) anniversary.
\on{19}.~\HL{trrpxI}{\B{trr\-pxì}}, part of day; hour.
\on{20}.~\HL{fleltrr}{\B{flel\-trr}}, April Fool's Day.

% trr'awve
\hi
\HT{trr'awve}{\B{Trr\cp{'aw}ve}} \I{[t\s{r}."Paw.vE]} \emph{n.} Sunday. 
(\ding{43}~\HL{trr}{\B{trr}}, \HL{'awve}{\B{'aw\-ve}})

% trr'ong
\hi
\HT{trr'ong}{\B{trr\cp{'ong}}} \I{[t\s{r}."PoN]} \emph{n.}; \emph{adv.} daybreak, dawn.
\B{Trr\-'ong\-ta Txon\-'ong\-vay po to\-lì\-ran.} He walked from dawn until dusk.~\LINK{http://naviteri.org/2011/01/new-year-new-vocabulary/}{b}
(\ding{43}~\HL{trr}{\B{trr}}, \HL{'ong}{\B{'ong}}) \\
\textsc{Derivatives}: \ding{43}
\on{1}.~\HL{sresrr'ong}{\B{sre\-srr\-'ong}}, before daybreak. 
\on{2}.~\HL{trr'ongmaw}{\B{trr\-'ong\-maw}}, after daybreak.  

% trr'ongmaw
\hi
\HT{trr'ongmaw}{\B{trr\cp{'ong}maw}} \I{[t\s{r}."PoN.maw]} \emph{n.}; \emph{adv.} after daybreak, after dawn.
(\ding{43}~\HL{trr'ong}{\B{trr\-'ong}}, \HL{maw}{\B{maw}})

% trrkive
\hi
\HT{trrkive}{\B{Trr\cp{ki}ve}} \I{[t\s{r}."ki.vE]} \emph{n.} Saturday. 
(\ding{43}~\HL{trr}{\B{trr}}, \HL{kive}{\B{ki\-ve}})

% trrmuve
\hi
\HT{trrmuve}{\B{Trr\cp{mu}ve}} \I{[t\s{r}."mu.vE]} \emph{n.} Monday. 
(\ding{43}~\HL{trr}{\B{trr}}, \HL{muve}{\B{mu\-ve}})

% trrmrrve
\hi
\HT{trrmrrve}{\B{Trr\cp{mrr}ve}} \I{[t\s{r}."m\s{r}.vE]} \emph{n.} Thursday. 
(\ding{43}~\HL{trr}{\B{trr}}, \HL{mrrve}{\B{mrr\-ve}})

% trrpuve
\hi
\HT{trrpuve}{\B{Trr\cp{pu}ve}} \I{[t\s{r}."pu.vE]} \emph{n.} Friday. 
(\ding{43}~\HL{trr}{\B{trr}}, \HL{puve}{\B{pu\-ve}})

% trrpxeyve
\hi
\HT{trrpxeyve}{\B{Trr\cp{pxey}ve}} \I{[t\s{r}."p'Ej.vE]} \emph{n.} Tuesday. 
(\ding{43}~\HL{trr}{\B{trr}}, \HL{pxeyve}{\B{pxey\-ve}})

% trrpxì
\hi 
\HT{trrpxI}{\B{trr\cp{pxì}}} \I{[t\s{r}."p'I]} \textsc{i}. \emph{n.} (\emph{a. on Pandora}) part of day. 
(\emph{b. on Earth}) hour (\emph{shortened from:} \B{trr\-pxì Saw\-tu\-te\-yä} \emph{and a variant of:} \ding{43}~\HL{pxevotrrpxI}{\B{pxe\-vo\-trr\-pxì}})
\B{Po tì\-kang\-kem so\-li (ka) trr\-pxì\-o a\-mrr.} She has worked for five hours.~\LINK{https://naviteri.org/2024/05/krrteri-about-time/}{b}
(\emph{b.1 to tell the full hour}) \B{tì\-'i\-'a trr\-pxì\-yä a\-mrr\-ve}, `five o'clock' [lit.: end of the fifth hour], \emph{usually abbreviated to} \B{trr\-pxì a\-mrr\-ve}, `five o'clock' [lit.: fifth hour] \emph{or even} \B{Lu mrr\-ve.} `It's five.' [lit.: `it's the fifth']~\LINK{https://naviteri.org/2024/05/krrteri-about-time/}{b}
(\emph{b.2 to tell minutes}) \B{trr\-pxì a\-mrr\-ve sì mawl} \emph{or coll.} \B{mrr\-ve sì mawl}, (it's) \on{5}:\on{30}.~\LINK{https://naviteri.org/2024/05/krrteri-about-time/}{b}
\B{trr\-pxì a\-mrr\-ve sì trr\-pxì\-vi a\-pxe\-vol} \emph{or coll.} \B{mrr\-ve sì pxe\-vol}, \on{5}:\on{24}.~\LINK{https://naviteri.org/2024/05/krrteri-about-time/}{b}
\B{trr\-pxì\-vi a\-vo\-mun sre srr\-pxì a\-mrr\-ve} \emph{or coll.} \B{vo\-mun sre mrr\-ve}, (it's) \on{4}:\on{50}, or `ten to five'.~\LINK{https://naviteri.org/2024/05/krrteri-about-time/}{b} \\
\textsc{ii.} \B{trr\cp{pxì}pe} \I{[t\s{r}."p'I.pE]} \emph{or} \B{pe\-srr\-\cp{pxì}} \I{[pE.s\s{r}."p'I]} \emph{inter.} what part of the day? what hour? 
\B{Fo pä\-hem trr\-pxì\-pe?--Sre\-kam\-trr.} What part of the day will they arrive?--Before noon.~\LINK{https://naviteri.org/2024/05/krrteri-about-time/}{b} 
(\ding{43}~\HL{trr}{\B{trr}}, \HL{hapxI}{\B{hapxì}}) \\
\textsc{Derivatives}: \ding{43} 
\on{1}.~\HL{trrpxIvi}{\B{trr\-pxì\-vi}}, minute.

% trrpxìvi
\hi 
\HT{trrpxIvi}{\B{trr\cp{pxì}vi}} \I{[t\s{r}."p'I.vi]} \emph{n.} minute (on Earth). 
\B{Oe pä\-hem maw\slash kay trr\-pxì\-vi a\-mrr.} I'll be there in five minutes.~\LINK{https://naviteri.org/2024/05/krrteri-about-time/}{b} 
(\emph{a. to tell time}) \B{trr\-pxì a\-mrr\-ve sì mawl} \emph{or coll.} \B{mrr\-ve sì mawl}, (it's) \on{5}:\on{30}.~\LINK{https://naviteri.org/2024/05/krrteri-about-time/}{b}
\B{trr\-pxì a\-mrr\-ve sì trr\-pxì\-vi a\-pxe\-vol} \emph{or coll.} \B{mrr\-ve sì pxe\-vol}, (it's) \on{5}:\on{24}.~\LINK{https://naviteri.org/2024/05/krrteri-about-time/}{b}
\B{trr\-pxì\-vi a\-vo\-mun sre srr\-pxì a\-mrr\-ve} \emph{or coll.} \B{vo\-mun sre mrr\-ve}, (it's) \on{4}:\on{50}, or `ten to five'.~\LINK{https://naviteri.org/2024/05/krrteri-about-time/}{b} \\
(\ding{43}~\HL{trrpxI}{\B{trrpxì}}) \\
\textsc{Derivatives}: \ding{43} 
\on{1}.~\HL{trrpxIvitsyIp}{\B{trr\-pxì\-vi\-tsyìp}}, second. 

% trrpxìvitsyìp
\hi 
\HT{trrpxIvitsyIp}{\B{trr\cp{pxì}vitsyìp}} \I{[t\s{r}."p'I.vi.\t{ts}jIp\textcorner]} \emph{n.} second (on Earth). 
(\ding{43}~\HL{trrpxIvi}{\B{trr\-pxì\-vi}}) 

% trrtsìve
\hi
\HT{trrtsIve}{\B{Trr\cp{tsì}ve}} \I{[t\s{r}."\t{ts}I.vE]} \emph{n.} Wednesday. 
(\ding{43}~\HL{trr}{\B{trr}}, \HL{tsIve}{\B{tsì\-ve}})

% trrtxon
\hi
\HT{trrtxon}{\B{trr\cp{txon}}} \I{[t\s{r}."t'on]} \emph{n.} day night cycle.
\B{Oe tì\-kang\-kem si trr\-txon nì\-wotx, nga\-ru ke\-'u!} I work all day and all night, and you don't give a damn!~\LINK{http://naviteri.org/2013/03/whoever-whatever-whenever/}{b}
(\ding{43}~\HL{trr}{\B{trr}}, \HL{txon}{\B{txon}})

% -tu
\hi
\HT{-tu}{\B{-tu}} \I{[tu]} \emph{suf. non\hyp{}productive to mainly turn word classes other than v.s into an agent, e.g., of n.s:} \B{pam\-tseo\-tu}, musician; \emph{of adj.s:} \B{nawm\-tu}, great person; \emph{other word classes:} \B{wä\-tu}, opponent; \emph{etc.}
(\ding{43}~\HL{-yu}{\B{-yu}})

% Tuk
\hi
\HT{Tuk}{\B{Tuk}} \I{[tuk\textcorner]} \emph{or} \I{[tUk\textcorner]} \emph{n., female name} (\emph{short form of:} \ding{43}~\HL{Tuktirey}{\B{Tuk\-ti\-rey}}).

% Tuke
\hi
\HT{Tuke}{\B{\cp{Tu}ke}} \I{["tu.kE]} \emph{n., female name}.
\B{Tu\-ke\-ru pol\-txe Ak\-wey nam\-'a\-ke, omum fu\-ta ke tsun poe sti\-vo.} Akwey spoke to Tuke confidently, knowing that she couldn't refuse.~\LINK{http://naviteri.org/2012/06/spring-vocabulary-part-3/}{b}

% tukru
\hi
\HT{tukru}{\B{tuk\cp{ru}}} \I{[tuk\textcorner."Ru]} \emph{or} \I{[tUk\textcorner.]} (\textsc{rn}: \B{tùk\cp{ru}} \I{[tUk\textcorner."Ru]}) \emph{n.} spear.
\B{Oe\-yä swi\-zaw nì\-ngay ti\-va\-kuk | Oe\-yä tuk\-rul txe'\-lan\-it ti\-va\-kuk.} Let my arrow strike true | Let my spear strike the heart.~\LINK{http://wiki.learnnavi.org/index.php/Hunting_Song}{wiki}
\B{Txe\-wìl tuk\-rut Lo\-ak\-ä nar\-mìn nì\-fmokx.} Txewì was eyeing Loak's spear enviously.~\LINK{http://naviteri.org/2011/10/more-vocabulary-a-bit-of-grammar/}{b}
\B{tuk\-ru a txo\-lu\-la fkol ta rìn\slash tuk\-ru a ta rìn\slash tuk\-ru le\-rìn}, a spear made of wood.~\LINK{http://naviteri.org/2017/09/zisikrr-amip-ayliu-amip-new-words-for-the-new-season/}{b}

% Tuktirey
\hi 
\HT{Tuktirey}{\B{Tuktirey}} \I{[tuk\textcorner.ti.REj]} \emph{n. female name} (\emph{short form:} \ding{43}~\HL{Tuk}{\B{Tuk}}).

% tul
\hi
\HT{tul}{\B{tul}} \I{[tul]} \emph{or} \I{[tUl]} \emph{vin.} (\emph{a.}) run. 
\B{Hum ne wrr\-pa! Tul!} Get outside! Run!~\LINK{http://forum.learnnavi.org/language-updates/txing-and-hum/}{ln\slash a\textsuperscript{c}} 
\B{Oe tul nì\-ftxan nì\-win na nga.} I run as fast as you.~\LINK{http://forum.learnnavi.org/language-updates/confirmations-\%28comparisons-gerund\%29/}{ln}
\B{Oe to nga tul nì\-win.} I run faster than you.~\LINK{http://forum.learnnavi.org/language-updates/confirmations-\%28comparisons-gerund\%29/}{ln}
\B{Po tul nì\-win fra\-to.} She is the fastest.~\LINK{http://forum.learnnavi.org/language-updates/confirmations-\%28comparisons-gerund\%29/}{ln}
\B{Txo fko ti\-vul mì na'\-rìng nì\-hawm\-pam, stawm ay\-ioang.} If you run noisily in the forest, the animals will hear.~\LINK{http://naviteri.org/2014/11/twenty-before-the-holidays/}{b}
(\emph{b. of the weather}) \B{Tul hu\-fwe nì\-win.} It's very windy.~\LINK{http://naviteri.org/2011/05/weather-part-2-and-a-bit-more-2/}{b}
(\emph{c. rhyming expression}) \B{Tslikx, tì\-ran, tul; | Ftu yì ne yì tsan\-'ul.} \emph{when learning s.t. new, you have to proceed from step to step: baby steps first, then bigger ones} [lit.: `Crawl, walk, run; from level to level get better.']~\LINK{http://naviteri.org/2023/09/aawa-ayliu-si-aylifyavi-amip-a-few-new-words-and-expressions/}{b}
(i.) \B{sä\-wä\-sul a tul}, foot race.~\LINK{http://naviteri.org/2012/10/wina-sawasultsyip-ahii-a-quick-little-contest/}{b}

% tulkun
\hi
\HT{tulkun}{\B{\cp{tul}kun}} \I{["tul.kun]} \emph{or} \I{["tUl.kUn]} (\textsc{rn}: \B{\cp{tùl}kun} \I{["tUl.kun]}) \emph{n., fauna} Tulkun (\emph{sentient whale\hyp{}like Pandoran sea creature}). 

% tumpasuk
\hi 
\HT{tumpasuk}{\B{\cp{tum}pasuk}} \I{["tum.pa.suk\textcorner]} \emph{or} \I{["tUm.pa.sUk\textcorner]} (\textsc{rn}: \B{\cp{tum}pasùk} \I{["tum.pa.sUk\textcorner]}) \emph{n.} \ding{96} celia fruit tree, \emph{pampinus bacca acinum}.~\LINK{https://forum.learnnavi.org/language-updates/expansion-to-flora-fauna-and-places-from-the-valley-of-mo'ara/}{ln} 

% tumpin
\hi
\HT{tumpin}{\B{\cp{tum}pin}} \I{["tum.pin]} \emph{or} \I{["tUm.]}  \emph{adj.} the color \ding{43}~\HL{tun}{\B{tun}}.
(\ding{43}~\HL{'opin}{\B{'o\-pin}})

% tun
\hi
\HT{tun}{\B{tun}} \I{[tun]} \emph{or} \I{[tUn]}  \emph{adj.} red, orange.
\B{Oe\-ri key tun slo\-lu nì\-txan.} I'm so embarrassed [lit.: my face has become very red].~\LINK{http://naviteri.org/2011/04/\%E2\%80\%99a\%E2\%80\%99awa-li\%E2\%80\%99fyavi-amip\%E2\%80\%94a-few-new-expressions/comment-page-1/\#comment-694}{b}
\B{tun na eyk\-tan}, leader red (\emph{the reddish color that distinguishes the dress of Na'vi leaders}).~\LINK{http://naviteri.org/2010/08/mipa-ayopin-mipa-ayliu-new-colors-new-words/}{b} \\
\textsc{Derivatives}: \ding{43}
\on{1}.~\HL{tumpin}{\B{tum\-pin}}, the color red\slash orange. 
\on{2}.~\HL{reypaytun}{\B{rey\-pay\-tun}}, red. 
\on{3}.~\HL{txeptun}{\B{txep\-tun}}, orange. 

% tunu
\hi 
\HT{tunu}{\B{\cp{tu}nu}} \I{["tu.nu]} \emph{adj.} romantic. 
\B{Nga\-ri 'e\-fu oe tu\-nu.} I feel romantic towards you; I have romantic feelings for you.~\LINK{http://naviteri.org/2021/04/mipa-ayliu-mipa-sioeykting-new-words-new-explanations/}{b}  
\B{Po yaw\-ne lu oer, slä po\-ri ke 'e\-fu oe tu\-nu.} I love him, but I don't have romantic feelings for him.~\LINK{http://naviteri.org/2021/04/mipa-ayliu-mipa-sioeykting-new-words-new-explanations/}{b} \\
\textsc{Derivatives}: \ding{43}
\on{1}.~\HL{tItunu}{\B{tì\-tu\-nu}}, romance.
\on{2}.~\HL{tunutu}{\B{tu\-nu\-tu}}, object of desire; ``crush''.

% tunutu
\hi 
\HT{tunutu}{\B{\cp{tu}nutu}} \I{["tu.nu.tu]} \emph{n.} object of desire; ``crush''. 
(\ding{43}~\HL{tunu}{\B{tu\-nu}})

% tung
\hi
\HT{tung}{\B{tung}} \I{[tuN]} \emph{or} \I{[tUN]} (\textsc{rn}: \B{tùng} \I{[tUN]}) \emph{vtr.} allow, let, permit. 
\B{Ke tung fkol tì\-wu\-sem\-it fì\-tseng.} Fighting isn't allowed here.~\LINK{http://naviteri.org/2017/01/meliuteri-alu-tung-si-pllhrr-about-tung-and-pllhrr/}{b}
\B{Pol to\-lung oe\-ru fu\-ta ki\-vä.\slash Pol to\-lung fu\-ta oe ki\-vä.} He allowed me to go.~\LINK{http://naviteri.org/2017/01/meliuteri-alu-tung-si-pllhrr-about-tung-and-pllhrr/}{b}
\B{Tung oer fu\-ta kä!} Let me go!~\LINK{http://naviteri.org/2017/01/meliuteri-alu-tung-si-pllhrr-about-tung-and-pllhrr/}{b}
\B{Ke tung Na'\-vil fu\-ta tswìk kxe\-ner\-it Mo\-'a\-ra\-ka nì\-wotx.} The Na'vi have banned smoking in Mo\-'ara.~\LINK{http://naviteri.org/2017/06/ayliu-a-ta-eywaeveng-kifkey-uniltirantokxa-words-from-disneys-pandora-the-world-of-avatar/}{b} \\
\textsc{Derivatives}: \ding{43}
\on{1}.~\HL{tungzup}{\B{tung\-zup}}, drop, let fall.

% tungzup
\hi
\HT{tungzup}{\B{tung\cp{zup}}} \I{[tuN."zup\textcorner]} \emph{or} \I{[tUN."zUp\textcorner]} (\textsc{rn}: \B{tùng\cp{zup}} \I{[tUN."zup\textcorner]}) \emph{vtr.} drop, let fall (\emph{accidentally, inadvertently}). 
\B{Ay\-oe sno\-laytx fì\-trr ta\-lun\-a oel rum\-it to\-lung\-zup.} We lost today because I dropped the ball.~\LINK{http://naviteri.org/2011/05/some-miscellaneous-vocabulary/}{b}
(\ding{43}~\HL{tung}{\B{tung}}, \HL{zup}{\B{zup}}, \HL{zeykup}{\B{zey\-kup}}) \\
\textsc{Derivatives}: \ding{43}
\on{1}.~\HL{nIktungzup}{\B{nìk\-tung\-zup}}, (hold) carefully, firmly.

% tup
\hi
\HT{tup}{\B{tup}} \I{[tup\textcorner]} \emph{or} \I{[tUp\textcorner]} (\textsc{rn}: \B{tùp} \I{[tUp\textcorner]}) \emph{conj.} instead of, rather than. 
\B{Lam ngay oer, fo ka\-yä ìlä hil\-van tup na'\-rìng.} I bet they'll take the river rather than the forest.~\LINK{http://forum.learnnavi.org/language-updates/concerning-tup/}{ln\slash g}

% tupe
\hi
\HT{tupe}{\B{\cp{tu}pe}} \emph{or} \HL{pesu}{\B{pe\cp{su}}} \I{["tu.pE]} \emph{or} \I{[pE."{}su]} \emph{inter.} who? what person?
\B{Kal\-txì. Nge\-nga lu tu\-pe? -- Oe lu Txe\-wì.} Hello. Who are you? -- I'm Txewì.~\LINK{http://naviteri.org/2010/09/getting-to-know-you-part-3/}{b}
\B{Nge\-yä tsa\-fkxi\-let tu\-pel re\-ngo\-lop?} Who designed that necklace of yours?~\LINK{http://naviteri.org/2013/01/awvea-posti-zisita-amip-first-post-of-the-new-year/}{b}
\B{Tsay\-sam\-si\-yu lu su\-pe? --  (Fo) lu Ka\-mun, Ra\-lu, Ìstaw, sì Ate\-yo.} Who are those warriors? [i.e. identify them] -- They're Ka\-mun, Ra\-lu, Ìstaw, and Ate\-yo.~\LINK{http://naviteri.org/2011/07/number-in-na\%E2\%80\%99vi/}{b}
(\ding{43}~\HL{-pe}{\B{pe}}\texttt{+}, \HL{tute1}{\B{tu\-te}}) 

% tut
\hi
\HT{tut}{\B{tut}} \I{[tut\textcorner]} \emph{or} \I{[tUt\textcorner]} (\textsc{rn}: \B{tùt} \I{[tUt\textcorner]}) \emph{part. of continuation,} ``And what about \ldots{}?'' 
\B{Oe\-ru syaw fko Ney\-ti\-ri. Nga\-ru tut?} My name is Neytiri. And how about you?~\LINK{http://naviteri.org/2012/10/audio-and-video-learning-materials-for-navi-101/}{b}
\B{Nga\-ru lu fpom srak? -- Lu fpom. Nga\-ru tut?} Are you well? -- I'm well. And you?~\LINK{http://naviteri.org/2012/10/audio-and-video-learning-materials-for-navi-101/}{b}
\B{Oe zo\-la\-'u ftu Wa\-sying\-ton. Nga\-ri tut?} I'm from Washington. And you?~\LINK{http://naviteri.org/2012/10/audio-and-video-learning-materials-for-navi-101/}{b}
\B{Su\-nu nga\-ru fkxen srak? -- Ke\-zem\-pll\-txe, su\-nu oe\-ru nì\-txan. -- Vey tut? -- Ngay\-txoa, ke tsun oe yi\-vom vey\-ti.} Do you like veggie stuff? -- Of course, I like it a lot. -- How about meat? -- Sorry, I can't eat meat.~\LINK{http://naviteri.org/2012/10/audio-and-video-learning-materials-for-navi-101/}{b} \\
\textsc{Derivatives}: \ding{43}
\on{1}.~\HL{nItut}{\B{nì\-tut}}, continually.
\on{2}.~\HL{letut}{\B{le\-tut}}, constant, continual.

% tutampe
\hi
\HT{tutampe}{\B{tu\cp{tam}pe}} \emph{or} \HL{pestan}{\B{pe\cp{stan}}} \I{[tu."tam.pe]} \emph{inter.} who? what man?~\LINK{http://naviteri.org/2011/07/number-in-na\%E2\%80\%99vi/}{b}
(\ding{43}~\HL{-pe}{\B{pe}}\texttt{+}, \HL{tutan}{\B{tu\-tan}}) 

% tutan
\hi
\HT{tutan}{\B{tu\cp{tan}}} \I{[tu."tan]} \emph{n.} male person, man. 
\B{Slä tsran\-ten fra\-to, tsa\-tu\-tan la\-yä\-ngu ye'\-rìn eyk\-tan.} But most important, that man will soon be leader.\LINK{http://naviteri.org/2016/12/tengkrr-perahem-zisit-amip-as-the-new-year-arrives/}{b}
\B{Ki\-fkey\-ri fì\-tu\-tan\-ìl a\-yay\-mak ke tslam stum ke\-'ut, ul\-te ke new ni\-vu\-me nì\-'ul.} This ignorant man understands almost nothing and doesn't want to learn more about the world.\LINK{http://naviteri.org/2016/12/tengkrr-perahem-zisit-amip-as-the-new-year-arrives/}{b}
\B{Fkol ftxo\-lä\-ngey na eyk\-tan tu\-tan\-ti a tì\-eyk\-tan\-ìri ke lu pxan kaw\-'it.} A man was chosen as leader who isn't worthy of leadership at all.~\LINK{http://naviteri.org/2016/12/tengkrr-perahem-zisit-amip-as-the-new-year-arrives/}{b}
(\ding{43}~\HL{tute1}{\B{tu\-te}}\textsuperscript{\on{1}})

% tute
\hi
\HT{tute1}{\B{\cp{tu}te}}\textsuperscript{\on{1}} \I{["tu.tE]} \emph{n.} person. 
\B{Ke tsun tu\-te a\-'aw tsa\-tskxe\-ti a\-ku'\-up kxi\-vel\-tek nì\-'aw\-tu.} One person alone can't lift that heavy rock.~\LINK{http://naviteri.org/2012/01/mipa-zisit-ayliu-amip-new-words-for-the-new-year/}{b}
\B{Ke\-tsran tu\-tel 'i\-vem, tsa\-fne\-tsngan lu ftxì\-vä'.} That kind of meat is gross no matter who cooks it.~\LINK{http://naviteri.org/2013/03/whoever-whatever-whenever/}{b}
\B{Sìl\-pey oe, tsa\-me\-su\-tel a mun\-txa so\-li nge\-yä ay\-lì\-'ut a\-lor tsì\-ye\-ve\-'a.} I hope, those two people who got married, will see your beautiful words.~\LINK{http://naviteri.org/2014/02/barter-and-exchange/comment-page-1/\#comment-2619}{b}
\B{Oel tse\-'a 'a\-'aw\-a tu\-tet.} I see several people.~\LINK{http://naviteri.org/2010/07/vocabulary-update/}{b} 
\B{Ngal oe\-yä 'upxa\-ret ay\-su\-te\-ru fpo\-le'.} You sent my message to people.~\LINK{http://wiki.learnnavi.org/index.php?title=Canon\#Extracts_from_various_emails}{wiki}
\B{Txam\-pxì\-ri su\-te\-yä, ke tsun fko ni\-vu\-me fte nì\-lì'\-fya a\-mip pi\-vll\-txe nìl\-tsan fa fwa tìng mik\-yun ho\-ren\-ur nì\-'aw.} For the majority of people, it's not possible to learn to speak a new language well simply by listening to rules.~\LINK{http://naviteri.org/2012/10/audio-and-video-learning-materials-for-navi-101/}{b}
\B{Tu\-te\-ri tì\-txur lu txew\-nga'.} There are limits to a person's strength.~\LINK{http://naviteri.org/2012/03/spring-vocabulary-part-1/}{b}
\B{Ik\-ran\-ìri krr\-a hu tu\-te tsa\-heyl si, ftang li\-vu yrr.} When an \emph{ikran} bonds with a person, it ceases to be wild.~\LINK{http://naviteri.org/2013/01/awvea-posti-zisita-amip-first-post-of-the-new-year/}{b}
\B{Tì\-'e\-fu\-mì oe\-yä su\-te\-kip nì\-wotx, ftxey Na'\-vi ftxey Saw\-tu\-te, lu sìl\-tsan lu kawng.} In my opinion among all people, whether Na'vi or Skypeople, there are good and bad.~\LINK{http://naviteri.org/2010/07/ting-mikyun-fte-tslivam-listening-comprehension-1-3/}{b}
(\ding{43}~\HL{-tu}{\B{-tu}}, \HL{-yu}{\B{-yu}}) \\
\textsc{Derivatives}: \ding{43}
\on{1}.~\HL{tuteo}{\B{tu\-teo}}, someone, somebody.
\on{2}.~\HL{tutan}{\B{tu\-tan}}, male person. 
\on{3}.~\HL{tute2}{\B{tu\-té}}, female person. 
\on{4}.~\HL{tawtute}{\B{taw\-tu\-te}}, human, skyperson.
\on{5}.~\HL{tutsena}{\B{tut\-sena}}, stretcher. 
\on{6}.~\HL{syeptute}{\B{syep\-tu\-te}}, hyneman.
\on{7}.~\HL{tiretu}{\B{ti\-re\-tu}}, shaman.
\on{8}.~\HL{nawmtu}{\B{nawm\-tu}}, noble person.
\on{9}.~\HL{tsantu}{\B{tsan\-tu}}, good person.
\on{10}.~\HL{kawngtu}{\B{kawng\-tu}}, bad person.
\on{11}.~\HL{matu}{\B{ma\-tu}}, ``excuse me, hey''.

% tuté
\hi
\HT{tute2}{\B{tu\cp{te}}}\textsuperscript{\on{2}} \emph{or} \B{tu\cp{té}} \I{[tu."tE]} \emph{n.} female (person), woman.
\B{Lu po lor\-a tu\-té a\-lor.} She's an extremely beautiful woman.~\LINK{http://naviteri.org/2013/02/vospxi-ayol-posti-apup-short-post-for-a-short-month/}{b}
\B{Fkxi\-le pewn\-ta tu\-té\-yä kur.} The bib necklace hangs from the woman's neck.~\LINK{http://naviteri.org/2013/04/wheres-the-bathroom-and-other-useful-things/}{b}
(\ding{43}~\HL{tute1}{\B{tu\-te}}\textsuperscript{\on{1}})

% tuteo
\hi
\HT{tuteo}{\B{\cp{tu}teo}} \I{["tu.tE.o]} \emph{n., pn.} someone, somebody. 
\B{Txo tu\-te\-o ni\-vew hi\-vum, po\-ru pll\-txe san ru\-txe 'i\-vì\-'awn.} If someone wants to leave, tell them to please stay.~\LINK{http://naviteri.org/2013/03/whoever-whatever-whenever/}{b}
\B{Lu su\-te\-o a lu 'e\-wan\-a ay\-nu\-me\-yu, ul\-te kop su\-te\-o a\-la\-he a tok fra\-kll\-pxìl\-tut ki\-fkey\-ä ul\-te ke lu 'e\-wan kaw\-'it.} There are some people who are young students, and also some other people in every territory of the world who are not young at all.~\LINK{http://forum.learnnavi.org/language-updates/ke-kawit/}{ln} 
(\ding{43}~\HL{tute1}{\B{tu\-te}}\textsuperscript{\on{1}}, \HL{-o}{\B{-o}})

% tutepe
\hi
\HT{tutepe}{\B{tu\cp{te}pe}} \emph{or} \HL{peste}{\B{pe\cp{ste}}} \I{[tu."tE.pE]} \emph{inter.} who? what woman?~\LINK{http://naviteri.org/2011/07/number-in-na\%E2\%80\%99vi/}{b}
(\ding{43}~\HL{-pe}{\B{pe}}\texttt{+}, \HL{tute2}{\B{tu\-te}}\textsuperscript{\on{2}}) 

% tutsena
\hi
\HT{tutsena}{\B{\cp{tut}sena}} \I{["tut\textcorner.sE.na]} \emph{or} \I{[tUt\textcorner.]} \emph{n.} stretcher.
(\ding{43}~\HL{tute1}{\B{tu\-te}}\textsuperscript{\on{1}}, \HL{sAhena}{\B{sä\-he\-na}})

% tuvom
\hi
\HT{tuvom}{\B{tu\cp{vom}}} \I{[tu."vom]} \emph{adj.} greatest of all, exceedingly great.
\B{En\-tu lu tu\-vom\-a ta\-ron\-yu. Kaw\-tut na po ke tso\-le\-'a oel mì sì\-rey.} Entu is an incredible hunter. I've never seen anyone like him before.~\LINK{http://naviteri.org/2015/04/some-new-words-for-may-day/}{b}
\B{Sa'\-sem\-ìri lu 'eveng\-ä kxitx yeng\-wal a\-tu\-vom.} For a parent, the greatest sorrow is the death of a child.~\LINK{http://naviteri.org/2015/04/some-new-words-for-may-day/}{b}

% tuvon
\hi
\HT{tuvon}{\B{\cp{tu}von}} \I{["tu.von]} \emph{vin.} lean.
\B{Ton\-ìri a\-la\-he, aw\-ngal yom wu\-tsot teng\-krr he\-reyn nì\-pxim, teng\-krr te\-ru\-von, ke tsran\-ten.} As for all other nights, it doesn't matter whether we eat dinner while sitting upright or leaning.~\LINK{http://whyisthisnight.com/addl-more.html\#na\%27vi}{tn}

\end{multicols}

% ===== Section TX ===== 

 \pdfbookmark[0]{TX}{TX}
\begin{center}
\section*{TX \I{[t']} (\emph{TxeTx})}
\end{center}

\begin{multicols}{2}

% txa'
\hi
\HT{txa'}{\B{txa'}} \I{[t'aP]} \emph{adj.} hard. 
\B{Teng\-krr hu pa\-lu\-kan\-tsyìp uvan se\-ri zo\-lìm oel mì sa'\-leng a 'u\-ot a lu txa' sì ekx\-txu.} While playing with my cat I felt something hard and rough on his skin.~\LINK{http://naviteri.org/2012/11/renu-ayinanfyaya-the-senses-paradigm/}{b}
(\ding{214}~\HL{hewne}{\B{hew\-ne}})

% txakrrfpìl
\hi 
\HT{txakrrfpIl}{\B{txa\cp{krr}fpìl}} \I{[t'a."k\s{r}.fp$\cdot$Il]} \emph{vtr., inf.}\on{33} consider, ponder. 
\B{Oel sä\-mok\-ti nge\-yä txa\-krr\-fpo\-lìl.} I have considered your suggestion.~\LINK{http://naviteri.org/2023/12/mipa-zisit-ayliu-amip-new-year-new-words/}{b} 
(\ding{43}~\HL{txan}{\B{txan}}, \HL{krr}{\B{krr}}, \HL{fpIl}{\B{fpìl}}) 

% txal
\hi
\HT{txal}{\B{txal}} \I{[t'al]} \emph{n.} back (of the body).
\B{Sran, sran, oe\-ri pì\-tìk tsa\-tse\-nget a mì tal! Fkxa\-ke nì\-ftxan kum\-a ter\-kup.} Yeah, yeah, scratch that place on my back! It itches like crazy!~\LINK{http://naviteri.org/2020/07/vospxivol-lefpom-happy-august/}{b}
(\emph{a. proverb}) \B{Txìm a\-'aw ke tsun hi\-veyn mì tal me\-fa'\-li\-yä.} You can't take both positions or sit on the fence; you need to decide. [lit.: One butt can't sit on the backs of two direhorses]~\LINK{http://naviteri.org/2013/02/vospxi-ayol-posti-apup-short-post-for-a-short-month/}{b} 

% txampam
\hi 
\HT{txampam}{\B{\cp{txam}pam}} \I{["t'am.pam]} \emph{n., fauna} soundblast colossus. 
(\ding{43}~\HL{txan}{\B{txan}}, \HL{pa}{\B{pam}}) 

% txampay
\hi
\HT{txampay}{\B{txam\cp{pay}}} \I{[t'am."paj]} \emph{n.} sea, ocean. 
\B{'Rr\-ta\-mì a tam\-pay lu ho\-et.} The oceans on Earth are vast.~\LINK{http://naviteri.org/2012/06/spring-vocabulary-part-3/}{b}
\B{Lu oer tì\-new a tse\-'a txam\-pay\-it.} I have a desire to see the ocean.~\LINK{http://naviteri.org/2014/11/twenty-before-the-holidays/}{b}
\B{Oe 'o\-long\-okx mì sray a txam\-pay\-ìri sim.} I was born in a town near the ocean.~\LINK{http://naviteri.org/2010/09/getting-to-know-you-part-2/}{b}
\B{Ne\-ni lew si fì\-txa\-yor vay txam\-pay nì\-wotx.} Sand covers this expanse all the way to the ocean.~\LINK{http://naviteri.org/2014/07/tsawlultxamaw-a-ayliu-post-meetup-words/}{b}
(\ding{43}~\HL{txan}{\B{txan}}, \HL{pay}{\B{pay}})

% txampxì
\hi
\HT{txampxI}{\B{txam\cp{pxì}}} \I{[t'am."p'I]} \emph{n.} majority, most, large part. 
\B{Hìm\-pxì Saw\-tu\-te\-yä lu tstun\-wi, slä fe\-yä txam\-pxì lä\-ngu kawng\-lan.} A minority of the Sky People are kind, but the majority are malicious.~\LINK{http://naviteri.org/2011/11/\%E2\%80\%99a\%E2\%80\%99awa-ayli\%E2\%80\%99u-amip-ni\%E2\%80\%99aw-\%E2\%80\%94-a-few-new-words-only/}{b}
\B{Txam\-pxì\-ri su\-te\-yä, ke tsun fko ni\-vu\-me fte nì\-lì'\-fya a\-mip pi\-vll\-txe nìl\-tsan fa fwa tìng mik\-yun ho\-ren\-ur nì\-'aw.} For the majority of people, it's not possible to learn to speak a new language well simply by listening to rules.~\LINK{http://naviteri.org/2012/10/audio-and-video-learning-materials-for-navi-101/}{b}
\B{Slä mì tam\-pxì krr\-ä, ke sar fkol kea hem\-lì\-'u\-vit fì\-lì\-'u\-mì.} But most of the time, one doesn't use infixes in this word.~\LINK{http://naviteri.org/2011/04/\%E2\%80\%99a\%E2\%80\%99awa-li\%E2\%80\%99fyavi-amip\%E2\%80\%94a-few-new-expressions/comment-page-1/\#comment-677}{b}
(\ding{214}~\HL{hImpxI}{\B{hìm\-pxì}}, \ding{43}~\HL{txan}{\B{txan}}, \HL{hapxI}{\B{ha\-pxì}})

% txan
\hi
\HT{txan}{\B{txan}} \I{[t'an]} \emph{adj.} (\emph{a. in quantity}) great, much. 
\B{Hu\-fwa tì\-sraw lu txan, tsun ay\-oe tsat li\-vewn.} Although the pain is great, we can endure it.~\LINK{http://naviteri.org/2021/03/mipa-ayliu-si-aylifyavi-new-words-and-expressions/}{b}
\B{Eo ay\-oeng lu txan\-a tì\-kawng.} A great evil is upon us.~\LINK{http://wiki.learnnavi.org/index.php?title=Na\%27vi_from_Avatar_Movie\#Jake_back_at_hometree}{wiki\slash a}
\B{Txan\-a srung so\-li po oer.} He helped me a lot.~\LINK{http://forum.learnnavi.org/language-updates/si-verbs-modification/}{ln}
\B{Txan\-a ira\-yo.} Much thanks.~\LINK{http://naviteri.org/2011/12/ulte-yora\%E2\%80\%99tu-leiu-and-the-winner-is/}{b}
\B{Na\-ri si! Kll\-te\-sìn lu pay a\-txan. Fwi rä\-'ä!} Be careful! There’s a lot of water on the ground. Don't slip!~\LINK{http://naviteri.org/2013/04/wheres-the-bathroom-and-other-useful-things/}{b}
\B{Slär\-ìl nga' fwep\-it a\-txan.} The cave is very dusty.~\LINK{http://naviteri.org/2014/07/tsawlultxamaw-a-ayliu-post-meetup-words/}{b} 
(\ding{214}~\HL{hIm}{\B{hìm}})
(\emph{b. of time}) long, much. \B{Yol\-a krr, txan\-a krr, ke tsran\-ten.} It doesn't matter how long it takes.~\LINK{http://naviteri.org/2013/03/whoever-whatever-whenever/}{b\slash a\textsuperscript{c}}
\B{Krr a fol txepit tok lu txan nìtam srak?} Have they been in the fire long enough?~\LINK{http://forum.learnnavi.org/intermediate/translation-exercise-a-navajo-coyote-tale-in-navi/}{ln}
\B{Ke lu oer set krr a\-txan.} I don't have much time now.~\LINK{http://forum.learnnavi.org/language-updates/verb-for-write/}{ln}
(\ding{214}~\HL{yol}{\B{yol}}) \\
\textsc{Derivatives}: \ding{43} 
\on{1}.~\HL{fItxan}{\B{fì\-txan}}, so, this much. 
\on{2}.~\HL{hImtxan}{\B{hìm\-txan}}, quantity.  
\on{3}.~\HL{nIftxan}{\B{nì\-ftxan}}, this much. 
\on{4}.~\HL{nItxan}{\B{nì\-txan}}, very.   
\on{5}.~\HL{txampay}{\B{txam\-pay}}, ocean, sea. 
\on{6}.~\HL{txampxI}{\B{txam\-pxì}}, majority. 
\on{7}.~\HL{txanatan}{\B{txa\-na\-tan}}, bright, vivid. 
\on{8}.~\HL{txanew}{\B{txa\-new}}, greedy. 
\on{9}.~\HL{txanfwerwI}{\B{txan\-fwer\-wì}}, blizzard. 
\on{10}.~\HL{txankrr}{\B{txan\-krr}}, for a long time. 
\on{11}.~\HL{txanro'a}{\B{txan\-ro\-'a}}, be famous. 
\on{12}.~\HL{txantompa}{\B{txan\-tom\-pa}}, rainstorm. 
\on{13}.~\HL{txantur}{\B{txan\-tur}}, powerful. 
\on{14}.~\HL{txantxew}{\B{txan\-txew}}, maximum. 
\on{15}.~\HL{txantsan}{\B{txan\-tsan}}, excellent. 
\on{16}.~\HL{txantslusam}{\B{txan\-tslu\-sam}}, wise. 
\on{17}.~\HL{txantstew}{\B{txan\-tstew}}, hero, heroine. 
\on{18}.~\HL{txanwawe}{\B{txan\-wa\-we}}, significant. 
\on{19}.~\HL{txasom}{\B{txa\-som}}, very hot. 
\on{20}.~\HL{txawew}{\B{txa\-wew}}, very cold. 
\on{21}.~\HL{txanlal}{\B{txan\-lal}}, ancient, very old. 
\on{22}.~\HL{txanlokxe}{\B{txan\-lo\-kxe}}, clan or tribal domain; country.
\on{23}.~\HL{txanso'hayu}{\B{txan\-so'\-ha\-yu}}, fan, enthusiast. 
\on{24}.~\HL{txankeltrrtrr}{\B{txan\-kel\-trr\-trr}}, extraordinary.
\on{25}.~\HL{txansngum}{\B{txan\-sngum}}, desperation.
\on{26}.~\HL{txantsawl}{\B{txan\-tsawl}}, huge, giant.
\on{27}.~\HL{txavA'}{\B{txa\-vä'}}, disgusting.
\on{28}.~\HL{txasunu}{\B{txa\-su\-nu}}, love greatly, enjoy tremendously.
\on{29}.~\HL{txanfwIngtu}{\B{txanfwìngtu}}, bastard, loser (\emph{insult})
\on{30}.~\HL{txantsawtsray}{\B{txan\-tsaw\-tsray}}, large city, metropolis.
\on{31}.~\HL{txantxewm}{\B{txan\-txewm}}, terrifying.
\on{32}.~\HL{txakrrfpIl}{\B{txa\-krr\-fpìl}}, consider, ponder.
\on{33}.~\HL{txampam}{\B{txam\-pam}}, soundblast colossus.

% txanatan
\hi
\HT{txanatan}{\B{\cp{txa}natan}} \I{["t'a.na.tan]} \emph{adj.} bright, vivid.
\B{Lu po\-ru tì\-ron\-srel a\-txa\-na\-tan.} She has a vivid imagination.~\LINK{http://naviteri.org/2011/02/new-vocabulary-part-2/}{b}
(\ding{43}~\HL{txan}{\B{txan}}, \HL{atan}{\B{atan}})

% txanew
\hi
\HT{txanew}{\B{\cp{txa}new}} \I{["t'a.nEw]} \emph{adj.} greedy.
(\ding{43}~\HL{txan}{\B{txan}}, \HL{new}{\B{new}}) \\
\textsc{Derivatives}: \ding{43}
\on{1}.~\HL{tItxanew}{\B{tì\-txa\-new}}, greed.

% txanfwerwì
\hi
\HT{txanfwerwI}{\B{txan\cp{fwer}wì}} \I{[t'an."fwER.wI]} \emph{n.} blizzard.
(\ding{43}~\HL{txan}{\B{txan}}, \HL{hufwe}{\B{hu\-fwe}}, \HL{herwI}{\B{her\-wì}})

% txanfwìngtu
\hi 
\HT{txanfwIngtu}{\B{txan\cp{fwìng}tu}} \I{[t'an."fwIN.tu]} \emph{n.} bastard, loser (\emph{insult}). 
(\ding{43}~\HL{txan}{\B{txan}}, \HL{fwIng}{\B{fwìng}})

% txankeltrrtrr
\hi 
\HT{txankeltrrtrr}{\B{\cp{txan}kel\cp{trr}trr}} \I{["t'an.kEl."t\s{r}.t\s{r}]} \emph{adj.} extraordinary. 
(\ding{43}~\HL{txan}{\B{txan}}, \HL{keltrrtrr}{\B{kel\-trr\-trr}}) \\
\textsc{Derivatives}: \ding{43}
\on{1}.~\HL{nItxankeltrrtrr}{\B{nì\-txan\-kel\-trr\-trr}}, extraordinarily.

% txankrr
\hi
\HT{txankrr}{\B{txan\cp{krr}}} \I{[t'an."k\s{r}]} \emph{adv.} for a long time.
\B{Txan\-krr ngal ke lawk pot kaw\-'it.} You haven't mentioned a thing about him in a long time.~\LINK{http://naviteri.org/2011/09/\%E2\%80\%9Cby-the-way-what-are-you-reading\%E2\%80\%9D/}{b}
\B{Fì\-sä\-omum txan\-krr lo\-lu oer.} I had this information for a long time.~\LINK{http://naviteri.org/2013/08/fmawn-a-fkol-kilmulat-this-just-in/}{b}
(\ding{43}~\HL{txan}{\B{txan}}, \HL{krr}{\B{krr}})

% txanlal
\hi
\HT{txanlal}{\B{\cp{txan}lal}} \I{["t'an.lal]}  \emph{adj.} ancient, very old (\emph{not for people}).
\B{Po\-leng ay\-oe\-ru ko\-ak\-tel vur\-it a\-txan\-lal.} The old woman told us an ancient story.~\LINK{http://naviteri.org/2015/08/aylifyavi-lereyfya-2-cultural-terms-2-and-more/}{b} 
(\ding{43}~\HL{txan}{\B{txan}}, \HL{lal}{\B{lal}})

% txanlokxe
\hi
\HT{txanlokxe}{\B{txanlo\cp{kxe}}} \I{[t'an.lo."k'E]} \emph{n.}  (\emph{a. on Pandora}) clan or tribal domain. 
(\emph{b. on Earth}) country.
\B{Po yaw\-ne lu snor nì\-'aw; fpom txan\-lo\-kxe\-yä ke tsran\-ten.} He only loves himself; the well\hyp{}being of the country doesn't matter.~\LINK{http://naviteri.org/2016/12/tengkrr-perahem-zisit-amip-as-the-new-year-arrives/}{b}
(\ding{43}~\HL{txan}{\B{txan}}, \HL{olo'}{\B{o\-lo'}}, \HL{atxkxe}{\B{atx\-kxe}})

% txanro'a
\hi
\HT{txanro'a}{\B{txan\cp{ro}'a}} \I{[t'an."R$\cdot$o.P$\cdot$a]} \emph{vin., inf.}\on{23} be famous. 
\B{Fì\-'um\-tsat to\-lìng 'e\-veng\-ur a\-spxin zey\-ko\-yul a txan\-ro\-'a.} This medicine was given to the sick child by a famous healer.~\LINK{http://naviteri.org/2019/06/50a-liu-amip-40-new-words/}{b}
\B{Vay fwa zo\-la\-'u Tsyeyk\-Su\-li, To\-ruk Mak\-to alu pi\-za\-yu Ney\-ti\-ri\-yä txan\-rar\-mo\-'a fra\-to kip ay\-ha\-pxì\-tu Oma\-ti\-ka\-yaä.} Until Jake Sully arrived, Neytiri's ancestor was the most famous Toruk Makto among the Omatikaya.~\LINK{http://naviteri.org/2012/03/spring-vocabulary-part-1/}{b} 
(\ding{43}~\HL{txan}{\B{txan}}, \HL{ro'a}{\B{ro\-'a}}) \\
\textsc{Derivatives}: \ding{43} 
\on{1}.~\HL{txanro'tu}{\B{txan\-ro'\-tu}}, famous person, celebrity.
\on{2}.\HL{tItxanro'a}{\B{tì\-txan\-ro\-'a}}, fame, glory.

% txanro'tu
\hi 
\HT{txanro'atu}{\B{txan\cp{ro'}tu}} \I{[t'an."RoP.tu]} \emph{n.} famous person, celebrity. 
\B{Krr\-a ul\-txa si nga tsa\-txan\-ro'\-tu\-hu, ze\-rok fu\-ta ze\-ne he\-nga si\-vi, ma 'e\-veng.} When you meet that famous person, remember that you have to be formally polite, child.~\LINK{http://naviteri.org/2022/02/lifyengteri-concerning-honorific-language/}{b} 
(\ding{43}~\HL{txanro'a}{\B{txan\-ro\-'a}}, \HL{tute1}{\B{tu\-te}})

% txansngum / si
\hi 
\HT{txansngum}{\B{txan\cp{sngum}}} \I{[t'an."{}sNum]} \emph{or} \I{[."{}sNUm]} \textsc{i}. \emph{n.} desperation, feeling of great worry.  \\
\textsc{ii}. \B{txan\cp{sngum} si} \I{[t'an."{}sNum si]} \emph{or} \I{[."{}sNum]} \emph{vin.} feel desperate.
\B{Ke lu syu\-ve ul\-te tu\-te a\-pxay txan\-sngum si.} There is no food, and many people are desperate.~\LINK{http://naviteri.org/2019/06/50a-liu-amip-40-new-words/}{b}
(\ding{43}~\HL{txan}{\B{txan}}, \HL{sngum}{\B{sngum}})

% txanso'hayu
\hi 
\HT{txanso'hayu}{\B{txan\cp{so'}hayu}} \I{[t'an."{}soP.ha.ju]} \emph{n.} fan, enthusiast; (\emph{coll. shortened to} \ding{43}~\HL{so'yu}{\B{so'\-yu}}). 
\B{Lu pxay\-a txan\-so'\-ha\-yu tsa\-rel\-ä a\-ru\-sikx alu Unil\-tì\-ran\-tokx ki\-fkey\-ka nì\-wotx.} There are many fans of \emph{Avatar} all over the world.~\LINK{http://naviteri.org/2017/12/zisit-amip-lefpom-ma-eylan-happy-new-year-friends/}{b} 
(\ding{43}~\HL{txan}{\B{txan}}, \HL{so'ha}{\B{so'\-ha}})

% txantompa
\hi
\HT{txantompa}{\B{txan\cp{tom}pa}} \I{[t'an."tom.pa]} \emph{n.} rainstorm, heavy rain.
\B{Maw txan\-tom\-pa, pxay\-a rìk\-äo la\-mu tska\-lep pe\-yä, ha tsat ku\-lat ay\-oel.} After the rainstorm, his crossbow was under a lot of leaves, so we uncovered it.~\LINK{http://naviteri.org/2011/10/more-vocabulary-a-bit-of-grammar/}{b}
(\ding{43}~\HL{txan}{\B{txan}}, \HL{tompa}{\B{tom\-pa}})

% txantur
\hi
\HT{txantur}{\B{\cp{txan}tur}} \I{["t'an.tuR]} \emph{or} \I{[.tUR]} \emph{adj.} powerful.
\B{Tsa\-tu\-tan la\-yä\-ngu ye'\-rìn eyk\-tan a txan\-tur fra\-to 'Rr\-ta\-mì.} That man will soon be the most powerful leader on Earth.~\LINK{http://naviteri.org/2016/12/tengkrr-perahem-zisit-amip-as-the-new-year-arrives/}{b}
\B{tu\-té a\-txan\-tur}, a powerful woman.~\LINK{http://naviteri.org/2016/12/tengkrr-perahem-zisit-amip-as-the-new-year-arrives/}{b}
(\ding{43}~\HL{txan}{\B{txan}}, \HL{txur}{\B{txur}})

% txantxew
\hi
\HT{txantxew}{\B{txan\cp{txew}}} \I{[t'an."t'Ew]} \emph{n.} maximum. 
(\ding{214}~\HL{hImtxew}{\B{hìm\-txew}}, \ding{43}~\HL{txan}{\B{txan}}, \HL{txew}{\B{txew}}) \\
\textsc{Derivatives}: \ding{43}
\on{1}.~\HL{txantxewvay}{\B{txan\-txew\-vay}}, maximally.

% txantxewm
\hi 
\HT{txantxewm}{\B{\cp{txan}txewm}} \I{["t'an.t'Ewm]} \emph{adj.} terrifying.~\LINK{http://naviteri.org/2020/11/vospxivopeya-ayliu-amip-novembers-new-words/}{b} 
(\ding{43}~\HL{txan}{\B{txan}}, \HL{txewm}{\B{txewm}})

% txantxewvay
\hi
\HT{txantxewvay}{\B{txan\cp{txew}vay}} \I{[t'an."t'Ew.vaj]} \emph{adv.} (\emph{a.}) maximally. 
(\emph{b. used with adj.s and adv.s to mean}) as … as possible.
\B{Tì\-ran nì\-fnu txan\-txew\-vay fte\-ke ay\-ye\-rik\-ìl aw\-nga\-ti sti\-vawm.} Walk as quietly as possible so the hexapedes won't hear us.~\LINK{http://naviteri.org/2012/10/snewsyea-ftxoza-halowina-a-spooky-halloween/}{b}
\B{Ay\-nga\-ri ftxo\-zä Hä\-lo\-win\-ä li\-vu 'o' sì snew\-sye txan\-txew\-vay.} May your Halloween be as fun and spooky as possible.~\LINK{http://naviteri.org/2012/10/snewsyea-ftxoza-halowina-a-spooky-halloween/}{b}
(\ding{214}~\HL{hImtxewvay}{\B{hìm\-txew\-vay}}, \ding{43}~\HL{txantxew}{\B{txan\-txew}}, \HL{vay}{\B{vay}})

% txantsan
\hi
\HT{txantsan}{\B{\cp{txan}tsan}} \I{["t'an.\t{ts}an]} \emph{adj.} excellent. 
\B{Tsa\-lì\-'u\-ri alu yo\-ra'\-tu, nge\-yä sä\-pll\-txe\-vi lu txan\-tsan.} As for the word \emph{yo\-ra'\-tu} your comment is excellent.~\LINK{http://naviteri.org/2011/12/ulte-yora\%E2\%80\%99tu-leiu-and-the-winner-is/comment-page-1/\#comment-1267}{b}
\B{Nge\-yä txan\-tsan\-a tì\-pawm\-ìri nga\-ru sei\-yi oe ira\-yo.} I thank you for your excellent question.~\LINK{http://forum.learnnavi.org/language-updates/another-email-from-frommer/}{ln}
\B{Rel\-tseo\-tu a\-txan\-tsan lu nga!} You are an excellent artist.~\LINK{http://forum.learnnavi.org/language-updates/art-related-vocab/}{ln}
(\ding{43}~\HL{txan}{\B{txan}}, \HL{sIltsan}{\B{sìl\-tsan}})

% txantsawl
\hi 
\HT{txantsawl}{\B{\cp{txan}tsawl}} \I{["t'an.\t{ts}awl]} \emph{adj.} giant, huge. 
(\ding{43}~\HL{txan}{\B{txan}}, \HL{tsawl}{\B{tsawl}})

% txantsawtsray
\hi 
\HT{txantsawtsray}{\B{txan\cp{tsaw}tsray}} \I{[t'an."\t{ts}aw.\t{ts}Raj]} \emph{n.} large city, metropolis.
\B{Slä set tä\-ngok kop tan\-tsaw\-tsray\-ti ay\-oe\-yä 'uol a\-la\-he: tsak\-tap.} But now, something else is also at our large cities: violence.~\LINK{http://naviteri.org/2020/06/mi-tanlokxe-oeya-srr-afpxamo-terrible-days-in-my-country/}{b}
(\ding{43}~\HL{txan}{\B{txan}}, \HL{tsawtsray}{\B{tsaw\-tsray}} )

% txantslusam 
\hi
\HT{txantslusam}{\B{\cp{txan}tslusam}} \I{["t'an.\t{ts}lu.sam]} \emph{adj.} wise, much\hyp{}knowing.   
\B{Lu nga win sì txur | Lu nga txan\-tslu\-sam.} You are fast and strong | You are wise.~\LINK{http://wiki.learnnavi.org/index.php/Hunting_Song}{wiki}
\B{Lu fra\-eyk\-tan\-ur a\-sìl\-tsan txan\-tslu\-sam\-a ay\-mo\-war\-si\-yu.} Every good leader has wise advisors.~\LINK{http://naviteri.org/2014/05/mipa-ayliu-mipa-aysafpil-new-words-new-ideas/}{b}
(\ding{43}~\HL{txan}{\B{txan}}, \HL{tslam}{\B{tslam}}, \HL{lafyon}{\B{la\-fyon}}, \HL{hafyonga'}{\B{ha\-fyo\-nga'}}) \\
\textsc{Derivatives}: \ding{43}
\on{1}.~\HL{tItxantslusam}{\B{tì\-txan\-tslu\-sam}}, wisdom.

% txantstew
\hi
\HT{txantstew}{\B{\cp{txan}tstew}} \I{["t'an.\t{ts}tEw]} \emph{n.} hero, heroine. 
\B{Txan\-tstew sä\-ro\-'a si, fna\-we'\-tu ke si.} A hero does great deeds, a coward does not.~\LINK{http://naviteri.org/2012/03/spring-vocabulary-part-1/}{b}
(\ding{43}~\HL{txan}{\B{txan}}, \HL{tstew}{\B{tstew}})

% txanwawe
\hi
\HT{txanwawe}{\B{txanwa\cp{we}}} \I{[t'an.wa."wE]} \emph{adj.} personally meaningful, significant. 
\B{Kì\-rey\-sì\-ri lu tsa\-puk\-ä ay\-vur a te\-ri 'Rr\-ta txan\-wa\-we nì\-ngay, slä oe\-ri, hu\-fwa el\-tur tì\-txen si, lu ral\-nga' nì\-'aw.} For Grace, the stories in that book relating to Earth are personally meaningful, but for me, although interesting, they’re simply instructive.~\LINK{http://naviteri.org/2012/07/fivospxiya-aylifyavi-amip-this-months-new-expressions-pt-2/}{b}
\B{Na\-ri si!  Tsa\-tsko lu spu\-win ul\-te ke lu mi txur, slä oe\-ri txan\-wa\-we le\-iu. Lu stxe\-li a sem\-pul\-ta.} Careful! That bow is old and no longer strong, but it means a lot to me. It was a gift from my father.~\LINK{http://naviteri.org/2011/05/some-miscellaneous-vocabulary/}{b}
(\ding{43}~\HL{txan}{\B{txan}}, \HL{wawe}{\B{wa\-we}}) \\
\textsc{Derivatives}: \ding{43}
\on{1}.~\HL{txanwetseng}{\B{txan\-we\-tseng}}, personally significant or beloved place, \emph{heimat}.

% txanwetseng
\hi 
\HT{txanwetseng}{\B{txan\cp{we}tseng}} \I{[t'an."wE.\t{ts}EN]} \emph{n.} personally significant or beloved place, \emph{heimat}. 
(\ding{43}~\HL{txanwawe}{\B{txan\-wa\-we}}, \HL{tseng}{\B{tseng}})

% txap
\hi 
\HT{txap}{\B{txap}} \I{[t'ap\textcorner]} \emph{vtr.} press, press on, apply pressure to. 
\B{Txap skxir\-it fte\-ke rey\-pay wrr\-zi\-va\-'u.} Apply pressure to the wound so that the blood won't flow.~\LINK{http://naviteri.org/2019/06/50a-liu-amip-40-new-words/}{b}  \\
\textsc{Derivatives}: \ding{43}
\on{1}.~\HL{tItxap}{\B{tì\-txap}}, pressure. 
\on{2}.~\HL{pxawtxap}{\B{pxaw\-txap}}, squeeze.

% txasom
\hi
\HT{txasom}{\B{\cp{txa}som}} \I{["t'a.som]} \emph{adj.} very hot. 
\B{Ya\-ri som\-wew\-pe (set)?---Ya lu txa\-som.} What's the temperature?---It's very hot.~\LINK{http://naviteri.org/2012/10/audio-and-video-learning-materials-for-navi-101/}{b}
(\ding{214}~\HL{txawew}{\B{txa\-wew}}, \ding{43}~\HL{txan}{\B{txan}}, \HL{som}{\B{som}})

% txasunu
\hi 
\HT{txasunu}{\B{txa\cp{su}nu}} \I{[t'a."{}su.nu]} \emph{vin.} love greatly, enjoy tremendously. 
\B{Txa\-su\-nu oe\-ru tey\-lu!} I really love \emph{teylu}!~\LINK{http://naviteri.org/2019/06/50a-liu-amip-40-new-words/}{b} 
\B{Txa\-su\-nu nga\-ru oey way\-te\-lem, ke\-fyak?} You love my songcord, don't you?~\LINK{https://www.youtube.com/watch?v=13aPkyXDuEE}{yt\slash fop}
(\ding{43}~\HL{txan}{\B{txan}}, \HL{sunu}{\B{su\-nu}})

% txatx
\hi
\HT{txatx}{\B{txatx}} \I{["t'at']} \emph{n.} bubble.
\B{Yo\-sìn kil\-van\-ä lu tatx.} There are bubbles on the surface of the river.~\LINK{http://naviteri.org/2015/11/vomuna-liu-amip-ten-new-words/}{b}

% txavä'
\hi 
\HT{txavA'}{\B{txa\cp{vä'}}} \I{[t'a."v\ae{}P]} \emph{adj.} disgusting. 
\B{Lu tsa\-kem txa\-vä', ma tsmuk.} That's disgusting, bro.~\LINK{http://naviteri.org/2019/06/50a-liu-amip-40-new-words/}{b} 
(\ding{43}~\HL{txan}{\B{txan}}, \HL{vA'}{\B{vä'}})

% txawew
\hi
\HT{txawew}{\B{\cp{txa}wew}} \I{["t'a.wEw]} \emph{adj.} very cold. 
\B{Ya\-ri som\-wew\-pe (set)?---Ya lu txa\-wew.} What's the temperature?---It's very cold.~\LINK{http://naviteri.org/2012/10/audio-and-video-learning-materials-for-navi-101/}{b}
\B{Mì zì\-sì\-krr a\-txa\-wew slu ay\-ora leyr.} In the very cold season, the lakes freeze.~\LINK{http://naviteri.org/2018/07/ta-sulfatu-a-ayliu-niul-more-words-from-our-experts/}{b}
(\ding{214}~\HL{txasom}{\B{txa\-som}}, \ding{43}~\HL{txan}{\B{txan}}, \HL{wew}{\B{wew}})

% txayo
\hi
\HT{txayo}{\B{\cp{txa}yo}} \I{["t'a.jo]} \emph{n.} field, plain, open terrain. 
\B{Oe lu Va'\-ru a ftu txa\-yo zo\-la\-'u.} I'm Va'ru from the plains.~\LINK{http://naviteri.org/2010/09/getting-to-know-you-part-1/}{b}
\B{Ne\-ni lew si fì\-txa\-yor vay txam\-pay nì\-wotx.} Sand covers this expanse all the way to the ocean.~\LINK{http://naviteri.org/2014/07/tsawlultxamaw-a-ayliu-post-meetup-words/}{b}
\B{Maw sä\-tsway\-on a\-yol ay\-oe kll\-po\-lä mì ta\-yo a lu ro\-fa kil\-van.} After a short flight we landed in a field beside the river.~\LINK{http://naviteri.org/2012/01/mipa-zisit-ayliu-amip-new-words-for-the-new-year/}{b}
(\ding{43}~\HL{txay}{\B{txay}}, \HL{yo}{\B{yo}}) \\
\textsc{Derivatives}: \ding{43}
\on{1}.~\HL{meitayo}{\B{mei\-ta\-yo}}, wetlands.

% txaw
\hi 
\HT{txaw}{\B{txaw}} \I{[t'aw]} \emph{vtr.} punish. 
\B{Sem\-pul\-ìl a\-sìl\-tsan sney eveng\-it ke txaw nì\-ka\-kan.} A good father doesn't punish his children roughly.~\LINK{http://naviteri.org/2021/09/aawa-ayliu-amip-a-few-new-words/}{b} \\
\textsc{Derivatives}: \ding{43} 
\on{1}.~\HL{tItxaw}{\B{tì\-txaw}}, punishment.

% txawnulsrung
\hi 
\HT{txawnulsrung}{\B{txaw\cp{nul}srung}} \I{[t'aw."nul.sRuN]} \emph{or} \I{[."nUl.sRUN]} (\textsc{rn}: \B{txaw\cp{nùl}srùng} \I{[daw."nUl.sRUN]}) \emph{n.} machine. \\
(i.) \B{txaw\-\cp{nul}\-srung a yur}, washing machine.~\LINK{https://naviteri.org/2024/03/fleltrra-ayliu-words-for-april-fools-day/}{b} \\
(ii.) \B{txaw\-\cp{nul}\-srung a tsway\-on}, airplane.~\LINK{https://naviteri.org/2024/03/fleltrra-ayliu-words-for-april-fools-day/}{b} 
(\ding{43}~\HL{txula}{\B{txu\-la}}, \HL{sAsrung}{\B{sä\-srung}})

% txay
\hi
\HT{txay}{\B{txay}} \I{[t'aj]} \emph{vin.} be horizontal, lie flat. \\
\textsc{Derivatives}: \ding{43}
\on{1}.~\HL{nItxay}{\B{nì\-txay}}, horizontally. 
\on{2}.~\HL{klltxay}{\B{kll\-txay}}, lie on the ground. 
\on{3}.~\HL{txayo}{\B{txa\-yo}}, field, plain.

% txärem
\hi
\HT{txArem}{\B{\cp{txä}rem}} \I{[t'\ae.REm]} \emph{n.} bone.
\B{neyn na txä\-rem}, the light color of bone.~\LINK{http://naviteri.org/2010/08/mipa-ayopin-mipa-ayliu-new-colors-new-words/}{b} 
\B{Tì\-kawng a su\-tel ngop var ri\-vey, tsìl\-tsan\-it pxìm kll\-yem fkol fe\-yä tä\-rem\-hu.} The evil that men do lives after them, The good is oft interred with their bones [\emph{Julius Caesar, Act III, Sc. 2}].~\LINK{http://naviteri.org/2017/09/zisikrr-amip-ayliu-amip-new-words-for-the-new-season/}{b} \\
\textsc{Derivatives}: \ding{43}
\on{1}.~\HL{snatxArem}{\B{sna\-txä\-rem}}, skeleton. 
\on{2}.~\HL{txArpawk}{\B{txär\-pawk}}, Pa\-lu\-lu\-kan Bone Horn.

% txärpawk
\hi
\HT{txArpawk}{\B{\cp{txär}pawk}} \I{["t'\ae{}R.pawk\textcorner]} \emph{n.} Pa\-lu\-lu\-kan Bone Horn.
(\ding{43}~\HL{txArem}{\B{txärem}}, \HL{pawk}{\B{pawk}})

% txe'lan
\hi
\HT{txe'lan}{\B{txe'\cp{lan}}} \I{[t'EP."lan]} \emph{n.} heart. 
\B{Oe\-yä tuk\-rul txe'\-lan\-it ti\-va\-kuk | Oe\-ri tì\-ngay\-ìl txe'\-lan\-it ti\-va\-kuk | Oe\-yä txe'\-lan li\-vu ngay.} Let my spear strike the heart | Let the ruth strike my heart | Let my heart be true.~\LINK{http://wiki.learnnavi.org/index.php/Hunting_Song}{wiki}
\B{Ay\-nga\-ri teng\-krr ya wur sle\-ru nì\-'ul, sìl\-pey oe, li\-vu hel\-ku sang ul\-te te'\-lan le\-fpom.} As for you all, while it's getting cooler, I hope your homes may be warm and your hearts happy.~\LINK{http://naviteri.org/2013/09/on-si-salewfya-shapes-and-directions/}{b} 
\B{Ay\-lì\-'u na ay\-skxe mì te'\-lan.} The words (are) like stones in (my) heart.~\LINK{http://wiki.learnnavi.org/index.php?title=Corpus\#Behind_the_Scenes}{ln\slash a\textsuperscript{c}}
\B{Sì 'e\-kong te'\-lan\-ä.} And the beat of the hearts.~\LINK{http://wiki.learnnavi.org/index.php/Weaving_Song}{wiki}
(\emph{a. idiom}) \B{Nga\-ri txe'\-lan ma\-wey li\-vu.\slash Txe'\-lan ma\-wey.} ``Don't worry (about it).''~\LINK{http://forum.learnnavi.org/language-updates/fmawno/}{ln} 
(\emph{b. idiom}) \B{txe'\-lan\-ti wrr\-zä\-rìp} \emph{vtr.} to greatly move emotionally.
\B{Oe\-ri pe\-yä ay\-lì\-'ul txe'\-lan\-ti wrr\-zo\-lä\-rìp.} Her words moved me greatly.~\LINK{http://naviteri.org/2021/03/mipa-ayliu-si-aylifyavi-new-words-and-expressions/}{b} \\
\textsc{Derivatives}: \ding{43}
\on{1}.~\HL{kawnglan}{\B{kawng\-lan}}, malicious, bad hearted. 
\on{2}.~\HL{txe'lankong}{\B{txe'\-lan\-kong}}, heartbeat.
\on{3}.~\HL{tstunlan}{\B{tstun\-lan}}, kind\hyp{}hearted.
\on{4}.~\HL{txeylan}{\B{txey\-lan}}, best friend.

% txe'lankong
\hi 
\HT{txe'lankong}{\B{txe'\cp{lan}kong}} \I{[t'E."lan.koN]} \emph{n.} heartbeat. 
(\ding{43}~\HL{txe'lan}{\B{txe'\-lan}}, \HL{'ekong}{\B{'ekong}})

% txekxumpay
\hi
\HT{txekxumpay}{\B{txe\cp{kxum}pay}} \I{[t'E."k'um.paj]} \emph{or} \I{[."k'Um.]} (\textsc{rn}: \B{txe\cp{kxùm}pay} \I{[dE."gUm.paj]}) \emph{n.} magma, lava.
\B{Txep\-ram pxor a krr, txan\-a txe\-kxum\-pay wrr\-za\-'u.} When a volcano erupts, a lot of lava comes out.~\LINK{http://naviteri.org/2015/03/kap-si-ayunil-saylahe-threats-dreams-and-other-things/}{b} 
(\ding{43}~\HL{txep}{\B{txep}}, \HL{kxumpay}{\B{kxum\-pay}})

% txele
\hi
\HT{txele}{\B{\cp{txe}le}} \I{["t'E.lE]} \emph{n.} matter (subject), topic.  
\B{Tsa\-hìk\-ur txe\-le lu.} This is a matter for the Tsahik.~\LINK{http://wiki.learnnavi.org/index.php?title=Na\%27vi_from_Avatar_Movie\#Neytiri_talking_to_Tsu.27tey}{wiki\slash a}
\B{Lu aw\-ngar ay\-te\-le a\-pxay a te\-ri sa\-'u pi\-vll\-txe.} We have many matters to talk about.~\LINK{http://forum.learnnavi.org/language-updates/txantsana-pxesisung-kefyak/}{ln}
\B{So\-rewn\-ìl kan\-'ìn tì\-'em\-it nì\-txan ul\-te kxawm tsa\-txe\-le me\-nga\-ne za\-'a\-tsu nì\-'eng.} Sorewn is really into cooking and you two might share that in common.\LINK{http://naviteri.org/2010/09/getting-to-know-you-part-1/}{b}
\B{Nge\-yä te\-ri fay\-te\-le a ay\-sä\-nu\-me\-ri ngar ira\-yo sei\-yi ay\-oe nì\-wotx.} We all thank you for your teachings about these topics.~\LINK{http://forum.learnnavi.org/language-updates/combined-answers-1-\%28feb-16\%29/}{ln}
\B{Fì\-txe\-le\-ri tì\-mok ke tam; ze\-ne fko fngi\-vo'.} In this matter, suggesting won't cut it; you need to demand.~\LINK{http://naviteri.org/2012/02/trr-asawnung-lefpom-happy-leap-day/}{b}
\B{Fì\-txe\-le\-ri tì\-nui ke lu tì\-ftxey.} In this matter, failure is not an option.~\LINK{http://naviteri.org/2013/04/wheres-the-bathroom-and-other-useful-things/}{b} \\
\textsc{Derivatives}: \ding{43}
\on{1}.~\HL{mIftxele}{\B{mì\-ftxe\-le}}, by the way. 
\on{2}.~\HL{mawftxele}{\B{maw\-ftxe\-le}}, belatedly, in hindsight.

% txen
\hi
\HT{txen}{\B{txen}} \I{[t'En]} \emph{adj.} awake.
\B{Nga lum\-pe lu txen nì\-mun, ma 'e\-vi?} Why are you awake again, my child?~\LINK{https://www.youtube.com/watch?v=13aPkyXDuEE}{yt\slash fop}
(\ding{214}~\HL{ngeyn}{\B{ngeyn}}) \\
\textsc{Derivatives}: \ding{43}
\on{1}.~\HL{tItxen}{\B{tì\-txen}}, awakeness.

% txep
\hi
\HT{txep}{\B{txep}} \I{[t'Ep\textcorner]} \emph{n.} fire. 
\B{Txep a\-hì\-'i mì tep\-tseng par\-ma\-lon.} A little fire was burning in the fireplace.~\LINK{http://naviteri.org/2018/03/100a-liu-amip-64-new-words-part-1/}{b}
\B{Txep\-ìl tìng le\-fpom\-a sä\-nrr\-ti.} The fire gives a pleasant glow.~\LINK{http://naviteri.org/2012/07/fivospxiya-aylifyavi-amip-this-months-new-expressions-pt-2/}{b}
\B{Txep\-ìri, yrr\-a rìn\-ti rä\-'ä sar, ma skxawng.} Don't use that green wood for a fire, you fool.~\LINK{http://naviteri.org/2013/01/awvea-posti-zisita-amip-first-post-of-the-new-year/}{b} 
\B{Maw\-krr\-a ngal txep\-it to\-long tsun hi\-vum.} After you've put out the fire you can leave.~\LINK{http://naviteri.org/2018/03/100a-liu-amip-64-new-words-part-1/}{b} \\
\textsc{Derivatives}: \ding{43}
\on{1}.~\HL{txepIva}{\B{txe\-pì\-va}}, ash, cinder. 
\on{2}.~\HL{txeptseng}{\B{txep\-tseng}}, fire pit. 
\on{3}.~\HL{txepvi}{\B{txep\-vi}}, spark. 
\on{4}.~\HL{ylltxep}{\B{yll\-txep}}, communal fire. 
\on{5}.~\HL{txepram}{\B{txep\-ram}}, volcano. 
\on{6}.~\HL{txekxumpay}{\B{txe\-kxum\-pay}}, magma, lava.
\on{7}.~\HL{vrrtep}{\B{vrr\-tep}}, demon. 
\on{8}.~\HL{txeptun}{\B{txep\-tun}}, orange.
\on{9}.~\HL{txeptsyIp}{\B{txep\-tsyìp}}, flame.
\on{10}.~\HL{txeptsyal}{\B{txep\-tsyal}}, coronis.

% txepìva
\hi
\HT{txepIva}{\B{\cp{txe}pìva}} \I{["t'E.pI.va]} \emph{n.} ash, cinder.
(\ding{43}~\HL{txep}{\B{txep}}, \HL{Ilva}{\B{ìl\-va}})

% txepram
\hi
\HT{txepram}{\B{txep\cp{ram}}} \I{[t'Ep\textcorner."Ram]} \emph{n.} volcano.
\B{Txep\-ram pxor a krr, txan\-a txe\-kxum\-pay wrr\-za\-'u.} When a volcano erupts, a lot of lava comes out.~\LINK{http://naviteri.org/2015/03/kap-si-ayunil-saylahe-threats-dreams-and-other-things/}{b} 
(\ding{43}~\HL{txep}{\B{txep}}, \HL{ram}{\B{ram}})

% txeptun
\hi 
\HT{txeptun}{\B{\cp{txep}tun}} \I{["t'Ep\textcorner.tun]} \emph{or} \I{[.tUn]} \emph{adj.} orange. 
(\ding{43}~\HL{txep}{\B{txep}}, \HL{tun}{\B{tun}})

% txeptseng
\hi
\HT{txeptseng}{\B{\cp{txep}tseng}} \I{["t'Ep\textcorner.\t{ts}EN]} \emph{n.} place where a fire is burning or has burned.   
\B{Tsa\-txep\-tseng\-mì lä\-ngu ay\-ut\-ral a\-ke\-ru\-sey nì\-'aw.} Sadly, there are only dead trees where the fire has been.~\LINK{http://naviteri.org/2012/07/fivospxiya-aylifyavi-amip-this-months-new-expressions-pt-2/}{b}
(\ding{43}~\HL{txep}{\B{txep}}, \HL{tseng}{\B{tse\-nge}})

% txeptsyal
\hi 
\HT{txeptsyal}{\B{\cp{txep}tsyal}} \I{["t'Ep\textcorner.\t{ts}jal]} \emph{n., fauna} coronis [lit.: `firewing']. 
(\ding{43}~\HL{txep}{\B{txep}}, \HL{tsyal}{\B{tsyal}})

% txeptsyìp 
\hi 
\HT{txeptsyIp}{\B{\cp{txep}tsyìp}} \I{["t'Ep'.\t{ts}jIp\textcorner]} \emph{n.} flame. 
\B{Tong txep\-tsyìp\-it.} Extinguish the flame.~\LINK{http://naviteri.org/2020/07/vospxivol-lefpom-happy-august/}{b}
\B{Lo\-nu\-sye tong txep\-tsyìp\-it.} Blow out the flame.~\LINK{http://naviteri.org/2020/07/vospxivol-lefpom-happy-august/}{b}
(\ding{43}~\HL{txep}{\B{txep}}) 

% txepvi
\hi
\HT{txepvi}{\B{\cp{txep}vi}} \I{["t'Ep\textcorner.vi]} \emph{n.} spark. 
(\ding{43}~\HL{txep}{\B{txep}}) \\
\textsc{Derivatives}: \ding{43} 
\on{1}.~\HL{txepvispxam}{\B{txep\-vi\-spxam}}, sparkle pod. 

% txepvispxam
\hi 
\HT{txepvispxam}{\B{\cp{txep}vispxam}} \I{["t'Ep\textcorner.vi.sp'am]} \emph{n., funga} sparkle pod. 
(\ding{43}~\HL{txepvi}{\B{txep\-vi}}, \HL{spxam}{\B{spxam}}) 

% txew
\hi
\HT{txew}{\B{txew}} \I{[t'Ew]} \emph{n.} edge, brink, limit, border, end.
\B{Sra\-ke pol la\-yok txew\-ti na'\-rìng\-ä?} Will he approach the edge of the forest?~\LINK{http://naviteri.org/2012/03/spring-vocabulary-part-1/}{b}
\B{Me\-tew\-it fì\-swek\-ä ey\-ki\-vil.} Pull the two ends of this bar together.~\LINK{http://naviteri.org/2013/04/wheres-the-bathroom-and-other-useful-things/}{b} 
\B{Ik\-ran ya\-wo\-lo ftu txew 'awkx\-ä.} The banshee took to the air from the edge of a cliff.~\LINK{http://naviteri.org/2012/03/spring-vocabulary-part-1/}{b}  
(\emph{a. idiom}) \B{Ke tsun aw\-nga pi\-vey nul\-krr--txew lok.} We can't wait any longer--time is almost up.~\LINK{http://naviteri.org/2012/03/spring-vocabulary-part-1/}{b} \\
\textsc{Derivatives}: \ding{43}
\on{1}.~\HL{txewluke}{\B{txew\-lu\-ke}}, endless, boundless. 
\on{2}.~\HL{txantxew}{\B{txan\-txew}}, maximum. 
\on{3}.~\HL{hImtxew}{\B{hìm\-txew}}, minimum. 
\on{4}.~\HL{txewnga'}{\B{txew\-nga'}}, finite. 
\on{5}.~\HL{tawtxew}{\B{taw\-txew}}, horizon, skyline.

% txewì
\hi
\B{\cp{Txe}wì} \I{["t'E.wI]} \emph{n. male name.} 
\B{Oe\-ru syaw (fko) Txe\-wì.} My name is Txe\-wì.~\LINK{http://naviteri.org/2010/09/getting-to-know-you-part-3/}{b}
\B{Txe\-wìl 'o\-la\-ku fwep\-it ftum\-fa kel\-ku sne\-yä.} Txewì removed the dust from inside his house.~\LINK{http://naviteri.org/2014/07/tsawlultxamaw-a-ayliu-post-meetup-words/}{b}
\B{Su\-nu Txe\-wì\-ru tok\-tor Kì\-rey\-sì.} Txe\-wì likes doctor Grace.~\LINK{http://naviteri.org/2010/07/ting-mikyun-fte-tslivam-listening-comprehension-1-3/}{b}
\B{Ni\-nat\-ìri tì\-ru\-sol Txe\-wì\-yä lu vä\-pam.} To Ni\-nat, Txe\-wì's singing is noise.~\LINK{http://naviteri.org/2014/11/twenty-before-the-holidays/}{b}

% txewk
\hi
\HT{txewk}{\B{txewk}} \I{[t'Ewk\textcorner]} \emph{n.} \ding{246} club (\emph{weapon}). 

% txewluke
\hi
\HT{txewluke}{\B{\cp{txew}luke}} \I{["t'Ew.lu.kE]} \emph{adj.} endless, boundless, without limit (\emph{for uncountable nouns}).
\B{Spaw Na'\-vil fu\-ta tì\-yawn Ey\-wa\-yä lu txew\-lu\-ke.} The Na'vi believe that Eywa's love is boundless.~\LINK{http://naviteri.org/2014/11/twenty-before-the-holidays/}{b} 
(\ding{43}~\HL{ketsuktiam}{\B{ke\-tsuk\-ti\-am}})

% txewm
\hi
\HT{txewm}{\B{txewm}} \I{[t'Ewm]} \emph{adj.} scary, frightening.   
\B{Slä ma sa'\-nu, ik\-ran txewm lu! Oe txo\-pu si!} But Mommy, the banshee is scary! I'm afraid!~\LINK{http://naviteri.org/2011/11/\%E2\%80\%99a\%E2\%80\%99awa-ayli\%E2\%80\%99u-amip-ni\%E2\%80\%99aw-\%E2\%80\%94-a-few-new-words-only/}{b}
\B{Txon\-am teng\-krr tar\-mì\-ran oe kxam\-lä na'\-rìng, sro\-ler eo ut\-ral a\-tsawl txewm\-a vrr\-tep.} Last night as I was walking through the forest, a frightening demon appeared in front of a big tree.~\LINK{http://naviteri.org/2012/03/spring-vocabulary-part-1/}{b} \\
\textsc{Derivatives}: \ding{43}
\on{1}.~\HL{txantxewm}{\B{txan\-txewm}}, terrifying.

% txewnga'
\hi
\HT{txenga'}{\B{\cp{txew}nga'}} \I{["t'Ew.NaP]} \emph{adj.} having a limit, not without bounds, finite.
\B{Tu\-te\-ri tì\-txur lu txew\-nga'.} There are limits to a person's strength.~\LINK{http://naviteri.org/2012/03/spring-vocabulary-part-1/}{b}
(\ding{43}~\HL{txew}{\B{txew}})

% txey
\hi
\HT{txey}{\B{txey}} \I{[t'Ej]} \emph{vin.} halt (\emph{usually used in military context}). 

% txeylan
\hi 
\HT{txeylan}{\B{\cp{txey}lan}} \I{["t'Ej.lan]} \emph{n.} best friend. 
(\ding{43}~\HL{txe'lan}{\B{txe'\-lan}}, \HL{'eylan}{\B{'ey\-lan}}) 

% txeym
\hi 
\HT{txeym}{\B{txeym}} \I{[t'Ejm]} \emph{adj.} interested (\emph{used with} \ding{43}~\HL{'efu}{\B{'e\-fu}}). 
\B{Tsa\-vur\-ìri oe 'e\-fu txeym.} I'm interested in that story.~\LINK{https://naviteri.org/2024/09/pxevola-liu-amip-two-dozen-new-words/}{b} \\
\textsc{Derivatives}: \ding{43} 
\on{1}.~\HL{tItxeym}{\B{tì\-txeym}}, interest.

% txi
\hi
\HT{txi}{\B{txi}} \I{[t'i]} \emph{n.} hurry, hurriedness, frenzy. \\
\textsc{Derivatives}: \ding{43}
\on{1}.~\HL{letxi}{\B{le\-txi}}, hurried, frenzied. 
\on{2}.~\HL{nItxi}{\B{nì\-txi}}, hurriedly, in a frenzied way. 
\on{3}.~\HL{letxiluke}{\B{le\-txi\-lu\-ke}}, unhurried. 
\on{4}.~\HL{nItxiluke}{\B{nì\-txi\-lu\-ke}}, unhurriedly, leisurely

% txikx
\hi 
\HT{txikx}{\B{txikx}} \I{[t'ik']} \emph{vtr.} chew. 

% Txilte
\hi
\B{\cp{Txil}te} \I{["t'il.tE]} \emph{n. female name}
\B{Txil\-te Ri\-ni\-sì tä\-pa\-re fì\-tsap nì\-soaia, slä tsal\-su\-ngay ke nì\-olo' ta\-krr\-a Ri\-ni mun\-txa slo\-lu.} Txilte and Rini are related by blood, but nevertheless not by clan since the time Rini got married.~\LINK{http://naviteri.org/2012/06/spring-vocabulary-part-3/}{b}

% txim
\hi 
\HT{txim}{\B{txim}} \I{[t'im]} \emph{n.} spike, thorn of a plant. 
(\emph{a. proverbial expression}) \B{na txim a txìm\-mì}, s.t. extremely annoying [lit.: `like a thorn in the butt'].~\LINK{http://naviteri.org/2023/09/aawa-ayliu-si-aylifyavi-amip-a-few-new-words-and-expressions/}{b}  \\
\textsc{Derivatives}: \ding{43} 
\on{1}.~\HL{naritxim}{\B{na\-ri\-txim}}, eyethorn.

% txin
\hi
\HT{txin}{\B{txin}} \I{[t'in]} \emph{adj.} main, primary. 
\B{Na'\-vi\-ri txin\-a sä\-'o tì\-tu\-sa\-ron\-ä lu tsko swi\-zaw.} For the Na'vi the bow and arrow is the main hunting tool.~\LINK{http://naviteri.org/2011/01/new-year-new-vocabulary/}{b}
\B{Hu\-fwa lu oe fì\-trr kar\-yu a\-txin, le\-iu oer kop pxe\-srung\-si\-yu a fì\-lì'\-fya\-ri slo\-lu tsul\-fä\-tu.} Although I'm the main teacher today, I'm happy to say I also have three assistants who have become masters of the language.~\LINK{http://naviteri.org/2012/07/tskxekengtsyip-a-mikyunfpi-a-little-listening-exercise/}{b} \\
\textsc{Derivatives}: \ding{43}
\on{1}.~\HL{txintIn}{\B{txin\-tìn}}, occupation. 
\on{2}.~\HL{txintseng}{\B{txin\-tseng}}, base of operations.
\on{3}.~\HL{txinfpIl}{\B{txin\-fpìl}}, main point.

% txinfpìl
\hi 
\HT{txinfpIl}{\B{\cp{txin}fpìl}} \I{["t'in.fpIl]} \emph{n.} main point (\emph{the primary idea or thesis statement of a presentation or argument}). 
\B{Oel nge\-yä txin\-fpìl\-it mi ke tslam.} I still don't understand your main point.~\LINK{http://naviteri.org/2021/03/mipa-ayliu-si-aylifyavi-new-words-and-expressions/}{b}
(\ding{43}~\HL{txin}{\B{txin}}, \HL{sAfpIl}{\B{sä\-fpìl}}; \HL{ngrrfpIl}{\B{ngrr\-fpìl}}) 

% txintìn
\hi
\HT{txintIn}{\B{\cp{txin}tìn}} \I{["t'in.tIn]} \emph{n.} occupation, primary role in society.   
\B{Nga\-ru lu pe\-fne\-txin\-tìn nì\-trr\-trr?} What is your primary role (in society)?~\LINK{http://naviteri.org/2010/09/getting-to-know-you-part-2/}{b}
\B{Oe\-yä txin\-tìn lu fwa stä'\-nì fay\-oang\-it.} My central societal occupation is to catch fish.~\LINK{http://naviteri.org/2010/09/getting-to-know-you-part-2/}{b}
(\ding{43}~\HL{txin}{\B{txin}}, \HL{tIn}{\B{tìn}})

% txintseng
\hi 
\HT{txintseng}{\B{\cp{txin}tseng}} \I{["t'in.\t{ts}EN]} \emph{n.} base of operations. 
(\ding{43}~\HL{txin}{\B{txin}}, \HL{tseng}{\B{tseng}}) \\
\textsc{Derivatives}: \ding{43}
\on{1}.~\HL{txintseng sawtuteyA}{\B{Txin\-tseng Saw\-tu\-te\-yä}}, Hell's Gate.

% Txintseng Sawtuteyä
\hi 
\HT{txintseng sawtuteyA}{\B{\cp{Txin}tseng \cp{Saw}tuteyä}} \I{["t'in.\t{ts}EN "{}saw.tu.tE.j\ae]} \emph{prop. n.} Hell's Gate (\emph{the main human base on Pandora}). 
(\ding{43}~\HL{txintseng}{\B{txin\-tseng}}, \HL{tawtute}{\B{taw\-tu\-te}})

% txìm
\hi
\HT{txIm}{\B{txìm}} \I{[t'Im]} \emph{n.} butt, ass, rear end (of the body). 
\B{Ean\-a txìm a\-tsawl.} Big, blue butt.~\LINK{http://www.pbs.org/wgbh/nova/blogs/secretlife/entertainment/paul-frommer/}{nova}
(\emph{a. proverb}) \B{Txìm a\-'aw ke tsun hi\-veyn mì tal me\-fa'\-li\-yä.} You can't take both positions or sit on the fence; you need to decide. [lit.: One butt can't sit on the backs of two direhorses]~\LINK{http://naviteri.org/2013/02/vospxi-ayol-posti-apup-short-post-for-a-short-month/}{b} 
(\emph{b. proverbial expression}) \B{na txim a txìm\-mì}, s.t. extremely annoying [lit.: `like a thorn in the butt'].~\LINK{http://naviteri.org/2023/09/aawa-ayliu-si-aylifyavi-amip-a-few-new-words-and-expressions/}{b} \\
\textsc{Derivatives}: \ding{43}
\on{1}.~\HL{txImya}{\B{txìm\-ya}}, fart.

% txìmya / si
\hi 
\HT{txImya}{\B{\cp{txìm}ya}} \I{["t'Im.ja]} \textsc{i}. \emph{n.} fart, gas emitted from the anus.~\LINK{https://forum.learnnavi.org/language-updates/new-words-fart-to-fart/}{ln} \\
\textsc{ii}. \B{txìm\-ya si} \I{["t'Im.ja si]} \emph{vin.} to fart.
(\ding{43}~\HL{txIm}{\B{txìm}}, \HL{ya}{\B{ya}})

% txìng
\hi
\HT{txIng}{\B{txìng}} \I{[t'IN]} \emph{vtr.} leave, abandon. 
\B{Txìng ay\-ioang\-it! Hum ne wrr\-pa! Tul!} Leave the animals! Get outside! Run!~\LINK{http://forum.learnnavi.org/language-updates/txing-and-hum/}{ln\slash a\textsuperscript{c}} 
\B{Swey\-lu txo ngal ke txi\-vìng sä\-fpìl\-it le\-tsim.} You shouldn't abandon your original idea.~\LINK{http://naviteri.org/2011/10/more-vocabulary-a-bit-of-grammar/}{b}

% txll'u
\hi
\HT{txll'u}{\B{\cp{txll}'u}} \I{["t'\s{l}.Pu]} \emph{n.}, \ding{96} hookagourd, \emph{flaska ascendens}. 

% txo
\hi
\HT{txo}{\B{txo}} \I{[t'o]} \emph{conj.} if (\emph{conditional}), \emph{optionally together with} \B{tsa\-krr}, then, \emph{or the imperative}.
(\emph{a. factual conditions})
\B{Txo new nga ri\-vey, oe\-hu!} Come with me, if you want to live!~\LINK{http://wiki.learnnavi.org/index.php?title=Corpus\#UGO_Movie_Blog_.28Twitter_Questions.29}{wiki}
\B{Tì\-tu\-sa\-ron\-ìri txo new fko sli\-vu tsul\-fä\-tu, ze\-ne smar\-to li\-vu wa\-lak.} If you want to become a master hunter, you have to be more active than your prey.~\LINK{http://naviteri.org/2011/10/more-vocabulary-a-bit-of-grammar/}{b}
\B{Txo new nga tsli\-vam, ze\-ne ki\-van\-fpìl.} If you want to understand, you have to concentrate.~\LINK{http://naviteri.org/2011/10/more-vocabulary-a-bit-of-grammar/}{b}
(\emph{b. hypothetical conditions}) 
\B{Pxan li\-vu txo nì\-'aw oe nga\-ri | Tsa\-krr nga Na'\-vi\-ru yom\-tì\-yìng.} Only if I am worthy of you | Will you feed the People.~\LINK{http://wiki.learnnavi.org/index.php/Hunting_Song}{wiki}
\B{Txo kxey\-ey\-ti ay\-ngal tsi\-ve\-'a, ru\-txe oe\-ru pi\-veng.} If you see a mistake, please tell me.~\LINK{http://naviteri.org/2011/05/some-miscellaneous-vocabulary/}{b}
\B{Txo ke tsi\-ye\-vun oe tì\-ftang si\-vi for.}  If I should not succeed in stopping them.~\LINK{http://forum.learnnavi.org/language-updates/stop!/}{ln}
\B{Txo fkol ke fyi\-vel u\-ran\-it pay\-wä, ze\-ne fko sli\-ve\-le.} If one does not seal a boat against water, one must swim.~\LINK{http://naviteri.org/2012/06/spring-vocabulary-part-3/}{b}
(\emph{c. with} \B{sweylu} \emph{and} ‹\B{iv}› \emph{for s.t. that hasn't yet happened}) should.
\B{Fì\-na\-er ftxì\-vä' lu nì\-hawng, ha swey\-lu txo ngal tsat kxi\-vukx nì\-win.} This drink tastes horrible, so you'd better swallow it quickly.~\LINK{http://naviteri.org/2011/11/\%E2\%80\%99a\%E2\%80\%99awa-ayli\%E2\%80\%99u-amip-ni\%E2\%80\%99aw-\%E2\%80\%94-a-few-new-words-only/}{b}
\B{Swey\-lu txo nga ki\-vä.\slash Nga swey\-lu txo ki\-vä.} You should go.~\LINK{http://naviteri.org/2011/04/\%E2\%80\%99a\%E2\%80\%99awa-li\%E2\%80\%99fyavi-amip\%E2\%80\%94a-few-new-expressions/}{b}
\B{Swey\-lu txo nga ke ki\-vä.\slash Nga swey\-lu txo ke ki\-vä.} You shouldn't go.~\LINK{http://naviteri.org/2011/04/\%E2\%80\%99a\%E2\%80\%99awa-li\%E2\%80\%99fyavi-amip\%E2\%80\%94a-few-new-expressions/}{b}
\B{Swey\-lu txo ki\-vä.} I\slash you\slash she\slash one\slash etc. should go.~\LINK{http://naviteri.org/2011/04/\%E2\%80\%99a\%E2\%80\%99awa-li\%E2\%80\%99fyavi-amip\%E2\%80\%94a-few-new-expressions/}{b}
(\emph{d. proverb}) \B{Txo ke nì\-yo' tsa\-krr nì\-yol.} ``If you can't be flawless, at least be brief.''~\LINK{http://naviteri.org/2012/01/mipa-zisit-ayliu-amip-new-words-for-the-new-year/}{b} 
(\ding{43}~\B{zun}) \\
\textsc{Derivatives}: \ding{43}
\on{1}.~\HL{txokefyaw}{\B{txo\-ke\-fyaw}}, if not, or else, otherwise.
\on{2}.~\HL{mungwrrtxo}{\B{mung\-wrr\-txo}}, unless, except if.
\on{3}.~\HL{txotsafya}{\B{txo\-tsa\-fya}}, if that's the case.

% txoa
\hi
\HT{txoa}{\B{\cp{txo}a}} \I{["t'o.a]} \emph{n.} forgiveness. 
\B{Säm\-yam\-ìl po\-ru wa\-yìn\-txu fu\-ta nga\-ta lo\-lu li txo\-a.} A hug will show him that you've already forgiven him.~\LINK{http://naviteri.org/2011/10/more-vocabulary-a-bit-of-grammar/}{b}
\B{Tì\-'al lu zop\-lo a tsa\-ri ke tsun txoa li\-vu nì\-ftue.} Wastlefulness is an offense that cannot easily be forgiven.~\LINK{http://naviteri.org/2014/07/tsawlultxamaw-a-ayliu-post-meetup-words/}{b} 
\B{Tu\-te a ke tsun txoa\-ti ti\-vìng, ke tsun tì\-yawn\-it ti\-vìng nì\-teng.} A person who cannot forgive, cannot love either.~\LINK{http://forum.learnnavi.org/navi-customs-and-culture/navi-proverbs/msg582576/\#msg582576}{ln}
(\emph{a. conv. phrase}) \B{Oe\-ru txoa li\-vu.} ``Forgive me.'' [lit.: may I have forgiveness] 
\B{Oe\-ru txoa li\-vu, ma 'ey\-lan. Rä\-'ä sti\-vi. Lu hì\-'i\-a sä\-'i\-pu nì\-'aw.} I'm sorry, friend. Don't be angry. It was just a small bit of humor.~\LINK{http://naviteri.org/2012/01/mipa-zisit-ayliu-amip-new-words-for-the-new-year/}{b} \\
\textsc{Derivatives}: \ding{43}
\on{1}.~\HL{hItxoa}{\B{hì\-txo\-a}}, excuse me, pardon me. 
\on{2}.~\HL{ngaytxoa}{\B{ngay\-txo\-a}}, forgive me, sorry. 

% txokefyaw
\hi
\HT{txokefyaw}{\B{\cp{txo}kefyaw}} \I{["t'o.kE.fjaw]} \emph{conj.} if not, or else, otherwise. 
\B{Kll\-te lu ekx\-txu. Na\-ri si txo\-ke\-fyaw tì\-ran nì\-nu.} The ground is rough. Be careful or you'll trip.~\LINK{http://naviteri.org/2013/04/wheres-the-bathroom-and-other-useful-things/}{b}
\B{Txo\-ke\-fyaw lì\-ye\-vu kxawm 'a\-'aw\-a tu\-te a pll\-txe san po tì\-kxey so\-li.} Otherwise, there may be several people who say, ``He messed up.''~\LINK{http://naviteri.org/2013/01/awvea-posti-zisita-amip-first-post-of-the-new-year/comment-page-1/\#comment-1951}{b}
(\ding{43}~\HL{txo}{\B{txo}}, \HL{fya'o}{\B{fya'o}}) 

% txolar
\hi
\HT{txolar}{\B{\cp{txo}lar}} \I{["t'o.laR]} \emph{n.}, \textsc{lw} dollar (\emph{currency}).
\B{Po\-el haw\-re'\-tsyìp\-it ko\-la\-nom yoa txo\-lar a\-me\-vol.} She bought a little cap for \$\on{16}.~\LINK{http://naviteri.org/2014/02/barter-and-exchange/}{b}

% txon
\hi
\HT{txon}{\B{txon}} \I{[t'on]} \emph{n.}; \emph{adv.} night.   
\B{Tom\-pa\-yä ka\-to | Tsaw\-ke\-yä ka\-to | Trr\-ä sì txon\-ä.} The rhythm of the rain | The rhythm of the sun | Of day and night.~\LINK{http://wiki.learnnavi.org/index.php/Weaving_Song}{wiki}
\B{Fì\-txon na ton a\-la\-he nì\-wotx pe\-lun ke lu teng?} Why is this night not the same like all other nights?~\LINK{http://whyisthisnight.com/addl-more.html\#na\%27vi}{tn}
\B{Txon le\-fpom (li\-vu ngar).} (May you have a) Good night.~\LINK{http://wiki.learnnavi.org/index.php?title=Canon\#Extracts_from_various_emails}{wiki} \B{Ton\-ìri a\-la\-he, aw\-ngal yom fkxen\-ti le\-rìk nì\-wotx.} As for other nights, we eat leafy vegetable\hyp{}food.~\LINK{http://whyisthisnight.com/addl-more.html\#na\%27vi}{tn}  \\
(i.) \B{txo\cp{nam}} \I{[t'o."nam]} \emph{adv., n.} last night. \B{Txon\-am teng\-krr tar\-mì\-ran oe kxam\-lä na'\-rìng, sro\-ler eo ut\-ral a\-tsawl txewm\-a vrr\-tep.} Last night as I was walking through the forest, a frightening demon appeared in front of a big tree.~\LINK{http://naviteri.org/2012/03/spring-vocabulary-part-1/}{b} \\
(ii.) \B{txo\cp{nay}} \I{[t'o."naj]} \emph{adv., n.} tomorrow night. \\  
\textsc{Derivatives}: \ding{43}
\on{1}.~\HL{trrtxon}{\B{trr\-txon}}, day night cycle.   
\on{2}.~\HL{txon'ong}{\B{txon\-'ong}}, nightfall.   
\on{3}.~\HL{txonkrr}{\B{txon\-krr}}, at night.   
\on{4}.~\HL{kxamtxon}{\B{kxam\-txon}}, midnight. 

% txon'ong
\hi
\HT{txon'ong}{\B{txon\cp{'ong}}} \I{[t'on."PoN]} \emph{n.}; \emph{adv.} nightfall, sunset, dusk. 
\B{Set sre\-fey oel fu\-ta tsam\-po\-ngu tä\-txaw maw txon\-'ong.} I'm currently expecting the war party back after nightfall.~\LINK{http://naviteri.org/2011/02/new-vocabulary-part-2/}{b}
\B{Trr\-'ong\-ta Txon\-'ong\-vay po to\-lì\-ran.} He walked from dawn until dusk.~\LINK{http://naviteri.org/2011/01/new-year-new-vocabulary/}{b} \\
(i.) \B{txon'o\cp{ngam}} \I{[t'on.Po."Nam]} yesterday's nightfall\slash sunset. \\
(ii.) \B{txon'o\cp{ngay}} \I{[t'on.Po."Naj]} tomorrow's nightfall\slash sunset. \B{Ul\-txa si\-vi oeng sìn ram\-tsyìp txon\-'ong\-ay.} Let's meet on the hill tomorrow at nightfall.~\LINK{http://naviteri.org/2012/01/mipa-zisit-ayliu-amip-new-words-for-the-new-year/}{b}
(\ding{43}~\B{txon}, \B{'ong}) \\  
\textsc{Derivatives}: \ding{43}
\on{1}.~\HL{sreton'ong}{\B{sre\-ton\-'ong}}, before nightfall. 
\on{2}.~\HL{txon'ongmaw}{\B{txon\-'ong\-maw}}, after nightfall.

% txon'ongmaw
\hi
\HT{txon'ongmaw}{\B{txon\cp{'ong}maw}} \I{[t'on."PoN.maw]} \emph{n.}; \emph{adv.} after nightfall\slash sunset\slash dusk. 
(\ding{43}~\HL{sreton'ong}{\B{sreton'ong}}, \HL{txon'ong}{\B{txon'ong}}, \HL{maw}{\B{maw}})

% txonkrr
\hi
\HT{txonkrr}{\B{\cp{txon}krr}} \I{["t'on.k\s{r}]} \emph{adv.} at night.
\B{Txon\-krr lu syu\-ra\-tan na'\-rìng\-ä Ey\-we\-veng\-ä lor nì\-txan.} At night, the bioluminescence of the Pandoran forest is very beautiful.~\LINK{http://naviteri.org/2012/03/spring-vocabulary-part-1/}{b}
\B{Na'\-rìng sngä\-'i ni\-vrr txon\-krr syu\-ra\-tan\-fa.} The forest begins to glow at night with bioluminescence.~\LINK{http://naviteri.org/2012/07/fivospxiya-aylifyavi-amip-this-months-new-expressions-pt-2/}{b}
(\ding{43}~\HL{txon}{\B{txon}}, \HL{krr}{\B{krr}})

% txopu
\hi
\HT{txopu}{\B{\cp{txo}pu}} \I{["t'o.pu]} \textsc{i}. \emph{countable n.} fear.
\B{Txo\-pul pe\-yä vll fu\-ta kaw\-krr ke sla\-yu tsam\-si\-yu.} His fear indicates he'll never become a warrior.~\LINK{http://naviteri.org/2012/06/spring-vocabulary-part-3/}{b} 
\B{Ru\-txe, ay\-nga\-ri fì\-kem fe\-yä 'e\-'a\-la to\-pur ka\-ngay sì\-yi nì\-'aw.} Please, this will only confirm their worst fears about you.~\LINK{http://forum.learnnavi.org/language-updates/two-words-in-use-eal-and-kangay-si/msg564284/\#msg564284}{ln} \\
\textsc{ii}.a \B{\cp{txo}pu si} \I{["t'o.pu si]} \emph{vin.} be afraid (\emph{with the topic} of). 
\B{Txo\-pu rä\-'ä si, lu ke\-tu\-wong\-o nì\-'aw.} Don't be afraid, it's just an alien.~\LINK{http://forum.learnnavi.org/language-updates/a-collection/}{ln\slash a}    
\B{Slä ma sa'\-nu, ik\-ran txewm lu! Oe txo\-pu si!} But Mommy, the banshee is scary! I'm afraid!~\LINK{http://naviteri.org/2011/11/\%E2\%80\%99a\%E2\%80\%99awa-ayli\%E2\%80\%99u-amip-ni\%E2\%80\%99aw-\%E2\%80\%94-a-few-new-words-only/}{b}
\B{Ay\-ha\-pxì\-tu po\-ngu\-ä txo\-pu si nì\-nän ta\-krr\-a Va'\-rul pxe\-ku\-tut lä\-txayn.} The members of the group are less afraid since Va'\-ru defeated three of the enemy.~\LINK{http://naviteri.org/2012/02/trr-asawnung-lefpom-happy-leap-day/}{b} 
\B{Ul\-te zu\-saw\-krr\-ìri txo\-pu si tu\-te a\-pxay.} And many people are afraid of the future.~\LINK{http://naviteri.org/2020/03/teri-tifkeytok-lefkrr-about-the-present-situation/}{b} \\
\textsc{ii}.b \B{\cp{txo}pu sley\cp{ku}} \I{["t'o.pu slEj."ku]} \emph{vtr.} frighten, scare, produce fear. \B{'Uol ik\-ran\-it txo\-pu sley\-ko\-la\-tsu, ta\-lun\-a po tsìk ya\-wo.} Something must have frightened the banshee, because it suddenly took to the air.~\LINK{http://naviteri.org/2012/01/mipa-zisit-ayliu-amip-new-words-for-the-new-year/}{b}
\B{Nge\-yä ay\-lì\-'ul a\-kxap\-nga' txo\-pu ke sley\-ku oet kaw\-'it.} Your threatening words don't scare me one bit.~\LINK{http://naviteri.org/2015/03/kap-si-ayunil-saylahe-threats-dreams-and-other-things/}{b} \\
\textsc{Derivatives}: \ding{43} 
\on{1}.~\HL{txopunil}{\B{txo\-pu\-nil}}, nightmare.

% txopunil
\hi
\HT{txopunil}{\B{\cp{txo}punil}} \I{["t'o.pu.nil]} \emph{n.} nightmare.
(\ding{43}~\HL{txopu}{\B{txo\-pu}}, \HL{unil}{\B{u\-nil}})

% txotsafya
\hi 
\HT{txotsafya}{\B{\cp{txo}tsafya}} \I{["t'o.\t{ts}a.fja]} \emph{adv.} if that's the case, if that's so. 
\B{Ke su\-nu ngar tey\-lu srak? Txo\-tsa\-fya, tìng oer pum\-it nge\-yä!} You don’t like teylu? If that’s the case, give me yours!~\LINK{http://naviteri.org/2022/12/trr-anawm-polaheiem-the-great-day-has-arrived/}{b}  
(\ding{43}~\HL{txo}{\B{txo}}, \HL{tsafya}{\B{tsa\-fya}}, \HL{txokefyaw}{\B{txo\-ke\-fyaw}})

% txukx
\hi
\HT{txukx}{\B{txukx}} \I{[t'uk']} \emph{or} \I{[t'Uk']} \emph{adj.} (\emph{a. physical}) deep. 
(\emph{b. metaphorically for thoughts, ideas, analyses, etc.}) deep, thorough. 
\B{Ay\-sä\-fpìl a\-txukx.} Deep thoughts.~\LINK{http://naviteri.org/2010/07/vocabulary-update/}{b} 
\B{Tsun sä\-fpìl a\-txukx nì\-'aw tsa\-tì\-pawm\-it ti\-vìng, ke\-fyak?} Only deep thought can give that question, right?~\LINK{http://naviteri.org/2010/08/mipa-ayopin-mipa-ayliu-new-colors-new-words/comment-page-1/\#comment-225}{b}
(\ding{214}~\HL{reng}{\B{reng}}) \\
\textsc{Derivatives}: \ding{43}
\on{1}.~\HL{nItxukx}{\B{nì\-txukx}}, deeply.
\on{2}.~\HL{txukxefu}{\B{txu\-kxe\-fu}}, care, be concerned about.

% txukxefu
\hi 
\HT{txukxefu}{\B{txu\cp{kxe}fu}} \I{[t'u."k'$\cdot$E.f$\cdot$u]} \emph{vin., inf.}\on{23} care, be concerned about, have deep feelings for (with the topical or \ding{43}~\HL{teri}{\B{teri}}). 
\B{Nga\-ri po ke txu\-kxe\-fu kaw\-'it.} He doesn't care a bit about you.~\LINK{http://naviteri.org/2022/12/trr-anawm-polaheiem-the-great-day-has-arrived/}{b}
\B{Fu\-ria te\-ri lì'\-fya aw\-nge\-yä nga txu\-kxe\-fu nì\-ftxan, sei\-yi i\-ra\-yo.} Thank you for caring so much about our language.~\LINK{http://naviteri.org/2022/12/trr-anawm-polaheiem-the-great-day-has-arrived/}{b}  
(\ding{43}~\HL{txukx}{\B{txukx}}, \HL{'efu}{\B{'e\-fu}}) \\
\textsc{Derivatives}: \ding{43} 
\on{1}.~\HL{tItxukxefu}{\B{tì\-txu\-kxe\-fu}}, care, concern. 

% txula
\hi
\HT{txula}{\B{\cp{txu}la}} \I{["t'u.la]} (\textsc{rn}: \B{\cp{txù}la} \I{["dU.la]}) \emph{vtr.} build, construct.
\B{Fu\-ri\-a txu\-la tsa\-lew\-ti lu srä sìl\-tsan to fngap.} For constructing that cover, woven cloth is better than metal.~\LINK{http://naviteri.org/2013/07/pximaw-tsawlultxa-just-after-the-meet-up/}{b}
\B{Fì\-fne\-ya\-yol tsrul\-it txu\-la mì song\-ropx ut\-ral\-ä fte ay\-li\-nit hi\-vaw\-nu wä sar\-ni\-o\-ang.} This kind of bird builds its nest in a tree cavity so as to protect its young from predators.~\LINK{http://naviteri.org/2013/04/wheres-the-bathroom-and-other-useful-things/}{b}
\B{Txo ngal tsat txi\-vu\-la, (tsa\-krr) fo za\-ya\-'u.} If you build that (then) they will come.~\LINK{http://naviteri.org/2013/04/zun-zel-counterfactual-conditionals/}{b}
\B{Maw fwa fkol Kel\-ut\-ral\-it sko\-la\-'a, lem\-rey\-a ha\-pxì\-tul tsa\-soai\-ä txo\-lu\-la mip\-a kel\-ku\-ti mì ta\-yo.} After Hometree was destroyed, the surviving members of that family built a new home on the plains.~\LINK{http://naviteri.org/2011/05/some-miscellaneous-vocabulary/}{b}
\B{Fì\-stxe\-lit fol txe\-ru\-la fpi olo'\-eyk\-tan.} They're constructing this gift for the chief.~\LINK{http://naviteri.org/2012/01/mipa-zisit-ayliu-amip-new-words-for-the-new-year/}{b}
\B{tuk\-ru a txo\-lu\-la fkol ta rìn}, a spear made of wood.~\LINK{http://naviteri.org/2017/09/zisikrr-amip-ayliu-amip-new-words-for-the-new-season/}{b}
(\emph{a. with} \B{txep}) make a fire(place). \B{Frr\-nen hì\-'i lu nì\-txan a krr, apxa txep\-it txu\-la oel.} When my babies are very little, I build a big fire.~\LINK{http://forum.learnnavi.org/intermediate/translation-exercise-a-navajo-coyote-tale-in-navi/}{ln} \\
\textsc{Derivatives}: \ding{43}
\on{1}.~\HL{tItxula}{\B{tì\-txu\-la}}, construction, constructed thing.
\on{2}.~\HL{txawnulsrung}{\B{txaw\-nul\-srung}}, machine.

% txum
\hi
\HT{txum}{\B{txum}} \I{[t'um]} \emph{or} \I{[t'Um]} (\textsc{rn}: \B{txùm} \I{[dUm]}) \emph{n.} poison. 
\B{Oe\-ri lu swoa fne\-txum.} I'm allergic to intoxicating beverage.~\LINK{http://forum.learnnavi.org/language-updates/new-words-from-the-2012-avatarmeet-navi-class/msg551479/\#msg551479}{b}  \\
\textsc{Derivatives}: \ding{43}
\on{1}.~\HL{txumnga'}{\B{txum\-nga'}}, poisonous. 
\on{2}.~\HL{txumpaywll}{\B{txum\-pay\-wll}}, scorpion thistle.  
\on{3}.~\HL{txumtsA'wll}{\B{txum\-tsä'\-wll}}, baja tickler.
\on{4}.~\HL{wAtum}{\B{wä\-tum}}, antidote.

% txumnga'
\hi
\HT{txumnga'}{\B{\cp{txum}nga'}} \I{["t'um.NaP]} \emph{or} \I{["t'Um.]} (\textsc{rn}: \B{\cp{txù}nga'} \I{["dU.NaP]}) \emph{adj.} poisonous.
(\ding{43}~\HL{txum}{\B{txum}}, \HL{-nga'}{\B{-nga'}})

% txumpaywll
\hi
\HT{txumpaywll}{\B{\cp{txum}paywll}} \I{["t'um.paj.w\s{l}]} \emph{or} \I{["t'Um.]} (\textsc{rn}: \B{\cp{txùm}paywll} \I{["dUm.paj.w\s{l}]}) \emph{n.}, \ding{96} scorpion thistle, ``poison water plant'', \emph{scorpioflora maxima}. 
(\ding{43}~\HL{txum}{\B{txum}}, \HL{pay}{\B{pay}}, \HL{'ewll}{\B{'ewll}})

% txumtsä'wll
\hi
\HT{txumtsA'wll}{\B{\cp{txum}tsä'wll}} \I{["t'um.\t{ts}\ae{}P.w\s{l}]} \emph{or} \I{["t'Um.]} (\textsc{rn}: \B{\cp{txùm}tsä'wll} \I{["dUm.\t{ts}E.w\s{l}]}) \emph{n.}, \ding{96} baja tickler, ``poison squirt plant'', \emph{flaska reclinata}. 
(\ding{43}~\HL{txum}{\B{txum}}, \HL{tsA'}{\B{tsä'}}, \HL{'ewll}{\B{'ewll}})

% txung
\hi 
\HT{txung}{\B{txung}} \I{[t'uN]} \emph{or} \I{[t'UN]} \emph{vtr.} disturb, disrupt, bother, affect negatively. 
\B{Fì\-sä\-spxin\-ìl ke txa\-yung aw\-ngey sì\-rey\-ti tì\-'i'\-a\-vay krr\-ä.} This illness will not disrupt our lifes forever.~\LINK{http://naviteri.org/2020/03/teri-tifkeytok-lefkrr-about-the-present-situation/}{b}
\B{Oey fpom\-it txung rä\-'ä!} Don't bother me!~\LINK{http://naviteri.org/2020/03/teri-tifkeytok-lefkrr-about-the-present-situation/}{b} 

% txur
\hi
\HT{txur}{\B{txur}} \I{[t'uR]} \emph{or} \I{[t'UR]} \emph{adj.} strong.   
\B{Lu nga win sì txur | Lu nga txan\-tslu\-sam.} You are fast and strong | You are wise.~\LINK{http://wiki.learnnavi.org/index.php/Hunting_Song}{wiki}
\B{Nga nì\-aw\-no\-mum to oe\-tsyìp lu txur nì\-txan.} As everyone knows, you're a lot stronger than little old me.~\LINK{http://naviteri.org/2010/07/diminutives-conversational-expressions/}{b}
\B{Krro lu 'Ìng\-lì\-sì txur nì\-hawng mì re\-'o.} Sometimes English is too strong in (my) head.~\LINK{http://naviteri.org/2011/09/\%E2\%80\%9Cby-the-way-what-are-you-reading\%E2\%80\%9D/comment-page-1/\#comment-1092}{b}
\B{Tsa\-sä\-lä\-txayn Na'\-vi\-ru srung so\-li nì\-'aw fte sli\-vu txur nì\-'ul.} That defeat only helped the People become stronger.~\LINK{http://naviteri.org/2012/01/more-additions-to-the-lexicon/}{b}
\B{Nga pam\-rel so\-li san sley\-ki\-vu ut\-ral\-ìl lì'\-fya\-yä le\-Na'\-vi pxay\-a ay\-vur\-it a\-txur sìk.} You wrote, ``May the Na'vi language tree produce many strong stories.''~\LINK{http://forum.learnnavi.org/language-updates/life-death-and-gerunds/}{ln} \\
\textsc{Derivatives}: \ding{43} 
\on{1}.~\HL{tItxur}{\B{tì\-txur}}, strength. 
\on{2}.~\HL{txantur}{\B{txan\-tur}}, powerful. 
\on{3}.~\HL{txurtel}{\B{txur\-tel}}, rope. 
\on{4}.~\HL{txurtu}{\B{txur\-tu}}, strongman\slash strongwoman, brawny person.
\on{5}.~\HL{txurtseng}{\B{txur\-tseng}}, fortification, fortress, bulwark.

% txurtel
\hi
\HT{txurtel}{\B{\cp{txur}tel}} \I{["t'uR.tEl]} \emph{or} \I{["t'UR.]} \emph{adj.} rope.
\B{Na\-ri si fte\-ke ay\-tur\-tel\-ur tì\-zin si\-vi!} Be careful not to tangle the ropes.~\LINK{http://naviteri.org/2021/04/mipa-ayliu-mipa-sioeykting-new-words-new-explanations/}{b} 
\B{Txur\-tel\-mì fo fta so\-li fte\-ke ka tsyokx fwi\-vi.} They tied a knot in the rope so it wouldn't slip through their hands.~\LINK{http://naviteri.org/2013/04/wheres-the-bathroom-and-other-useful-things/}{b}
(\ding{43}~\HL{txur}{\B{txur}}, \HL{telem}{\B{te\-lem}})

% txurtu
\hi 
\HT{txurtu}{\B{\cp{txur}tu}} \I{["t'uR.tu]} \emph{or} \I{["t'UR.]} \emph{n.} strongman\slash strongwoman, brawny person, \emph{a person of strong physique or character}.
(\ding{43}~\HL{txur}{\B{txur}}, \HL{-tu}{\B{-tu}}, \ding{43}~\HL{meyptu}{\B{meyptu}})

% txurtseng
\hi 
\HT{txurtseng}{\B{\cp{txur}tseng}} \I{["t'uR.\t{ts}EN]} \emph{or} \I{["t'UR.]} \emph{n.} fortification, fortress, bulwark. 
\B{Aw\-nga\-ri ke\-tsran tseng\-ne ki\-vä, fì\-soaia lu txur\-tseng.} Wherever we go, this family is our fortress.~\LINK{naviteri.org/2022/05/results-of-the-tttc/}{b} 
(\ding{43}~\HL{txur}{\B{txur}}, \HL{tseng}{\B{tse\-nge}}) 

\end{multicols}

% ===== Section TS =====  

 \pdfbookmark[0]{TS}{TS}
\begin{center}
\section*{TS \I{[\t{ts}]} (\emph{Tsä})}
\end{center}

\begin{multicols}{2}

% tsa-
\hi 
\HT{tsa-}{\B{tsa-}} \I{[\t{ts}a.]} \textsc{i}. \emph{prod. pref. on n.s} that (there) (\emph{distal singular dem. pref.}). \B{vur}, (a) story, \B{tsa\-vur}, that story. (\ding{43}~\B{fì-}) \\
\textsc{ii}. \emph{shortened form of} \B{tsa'u\slash tsaw} \emph{with enclitic adp.s} \B{Po tsa\-ne kar\-mä a tseng\-it ke tsì\-me\-'a oel.} I didn't see where she was going. [lit.: I didn't see the place that she was going to (it)]~\LINK{http://forum.learnnavi.org/language-updates/combined-answers-1-\%28feb-16\%29/}{ln}
\B{Oel nge\-y\"a tì\-hawl\-it na\-txu ulte tsa\-w\"a wa\-syem.} I disapprove of your plan and will oppose (fight against) it.~\LINK{http://naviteri.org/2020/02/some-words-for-leap-year-day/}{b} 

% tsa'u
\hi
\HT{tsa'u}{\B{\cp{tsa}'u}} \I{["\t{ts}a.Pu]} (\emph{pl.} \B{(ay)sa'u}) \textsc{i}. \emph{n.} that thing, it. 
\B{Tsa\-'u el\-tur tì\-txen si.} That's interesting.~\LINK{http://forum.learnnavi.org/language-updates/txelanit-hivawl/}{ln}
\B{Tsa\-'u la\-tsu pe\-u?} What on earth is that?~\LINK{http://naviteri.org/2011/10/more-vocabulary-a-bit-of-grammar/}{b}
\B{Ke lu oer am\-'a, tsa\-'u po\-lä\-hem a krr, su\-nu oe\-ru nì\-txan.} I have no doubt that when it arrives, I'm going to enjoy it very much.~\LINK{http://naviteri.org/2011/05/weather-part-2-and-a-bit-more-2/}{b}
\B{Ngal new a tsa\-'ut rä\-'ä wi\-vo, ma 'evi. Vi\-vin.} Don't reach for what you want, child. Ask for it.~\LINK{http://naviteri.org/2012/01/mipa-zisit-ayliu-amip-new-words-for-the-new-year/}{b}
\B{Txo tsi\-ve\-'a ay\-ngal key\-eyt, ru\-txe oe\-ru pi\-veng fte tsi\-vun oe sa\-'ut ley\-ki\-va\-tem.} If you see mistakes, please tell me so that I can change them.~\LINK{http://naviteri.org/2010/06/first-post/}{b}
\B{Tsa\-'u\-ri po pll\-ngay.} He admits it.~\LINK{http://naviteri.org/2011/07/txantsana-ultxa-mi-siatll-great-meeting-in-seattle/}{b} \\
\textsc{ii}. \B{tsa'u} \emph{or} \B{tsaw} \I{[\t{ts}aw]} \emph{pn.} it, that. 
\B{Lä\-ngu fì\-seyn kel\-ho\-an nì\-ngay. Tsun oe hi\-veyn tsaw\-sìn 'a\-'aw\-a swaw\-tsyìp nì\-'aw.} This chair is really uncomfortable. I can only sit on it a few seconds.~\LINK{http://naviteri.org/2015/08/aylifyavi-lereyfya-2-cultural-terms-2-and-more/}{b}
(\emph{shortened forms:} \B{tsal, tsat(i), tsar(u), tse\-yä, tsa\-ri})
\B{Tsat rä\-'ä yi\-vom!} Don't eat that.~\LINK{http://naviteri.org/2014/10/tson-si-fpomron-obligation-and-mental-health/}{b}
\B{Ke tsun oe fì\-'awkx\-it tsyi\-vìl. Ke lu tsa\-ru kea yì.} I can't scale this cliff. It doesn't have any footholds.~\LINK{http://naviteri.org/2013/01/lifyari-po-peyi-how-good-is-her-navi/}{b}
\B{Fra\-'u\-ä ran ngä\-pop fa fra\-syon tse\-yä.} The \emph{ran} of each thing arises from the totality of its attributes.~\LINK{http://naviteri.org/2012/10/mipa-vospxi-mipa-ayliu-new-words-for-the-new-month/}{b}
\B{Ye'\-rìn 'ì\-yi'\-a sä\-nu\-me a tsa\-ri kll\-fro' oe.} Soon the teaching that I am responsible for will end.~\LINK{http://forum.learnnavi.org/language-updates/trr-rrtaya/}{ln}

% tsay+
\hi 
\HT{tsay}{\B{tsay}}\texttt{+} \I{[\t{ts}aj.]} \emph{prod. pref. on n.s} those (there) (\emph{distal plural dem. pref.}). 
\B{Ke pol\-txe pol tsay\-lì\-'ut.} She didn't say that (i.e., those words).~\LINK{http://naviteri.org/2011/08/reported-speech-reported-questions/}{b}
(\ding{43}~\HL{tsa-}{\B{tsa-}}; \HL{fay}{\B{fay}\texttt{+}}) 

% tsafya
\hi
\HT{tsafya}{\B{\cp{tsa}fya}} \I{["\t{ts}a.fja]} \emph{adv.} that way, like that. 
\B{Sra\-ke tsun oe fay\-upxa\-ret tsli\-vam nì\-ftue? Tse … ze\-ne pi\-vll\-txe san pxìm tsa\-fya lu sìk.} Can I understand these messages easily? Well…  I must say ``Often that way''.~\LINK{http://forum.learnnavi.org/language-updates/verb-for-write/msg175592/\#msg175592}{ln}
\B{Slä tsa\-fya pam\-rel rä\-'ä si!} But don't write like that.~\LINK{http://naviteri.org/2011/09/\%E2\%80\%9Cby-the-way-what-are-you-reading\%E2\%80\%9D/comment-page-1/\#comment-1092}{b}
(\emph{a. idiom}) \B{fì\-fya tsa\-fya}, one way or another. \B{Slä fì\-fya tsa\-fya, nga\-ri ti\-re\-al tsa\-tse\-nget ta\-yok.} But one way or the other your spirit will be there.~\LINK{http://naviteri.org/2013/04/sin-asok-sin-zusawkrra-recent-and-upcoming-activities/comment-page-1/\#comment-2316}{b}

% tsa-hey
\hi
\HT{tsa-hey}{\B{tsa-hey}} \I{["\t{ts}a."hEj]} \emph{intj. expression of warning or frustration}.

% tsaheylu / tsaheyl si
\hi
\HT{tsaheylu}{\B{tsa\cp{hey}lu}} \I{[\t{ts}a."hEj.lu]} \textsc{i}. \emph{n.} bond, neural connection. 
\B{Sä\-tsway\-on a\-'aw\-ve tsa\-heyl\-ur ha\-sey si; ke tsun nga pi\-vey.} First flight seals the bond; you cannot wait.~\LINK{http://naviteri.org/2012/06/spring-vocabulary-part-3/}{b} \\
\textsc{ii}. \B{tsa\cp{heyl} si} \I{[\t{ts}a."hEjl si]} \emph{vin.} bond (\B{hu}, with), establish a neural connection. 
\B{Ik\-ran\-ìri krr\-a hu tu\-te tsa\-heyl si, ftang li\-vu yrr.} When an \emph{ikran} bonds with a person, it ceases to be wild.~\LINK{http://naviteri.org/2013/01/awvea-posti-zisita-amip-first-post-of-the-new-year/}{b}
\B{Tsyìl Ik\-ni\-ma\-yat ul\-te tsa\-heyl si ik\-ran\-hu a fì\-'u lu tì\-fme\-tok a ze\-ne fra\-ta\-ron\-yu a\-'e\-wan em\-zi\-va\-'u.} Scaling Iknimaya and bonding with a banshee is a test that every young hunter must pass.~\LINK{http://naviteri.org/2012/01/more-additions-to-the-lexicon/}{b}

% tsahìk
\hi
\HT{tsahIk}{\B{\cp{tsa}hìk}} \I{["\t{ts}a.hIk]} \emph{n.} Tsa\-hik, high priestess (\emph{who interprets the will of Eywa}). 
\B{Mo\-'at \mbox{alu} Tsa\-hìk lu O\-ma\-ti\-ka\-ya\-ä le\-'awa ha\-pxì\-tu a \mbox{ioi} sä\-pi fa 'a\-re.} Mo'at, the Tsahik, is the only member of the Omatikaya who wears a poncho.~\LINK{http://naviteri.org/2011/08/new-vocabulary-clothing/}{b}
\B{Sko Sa\-hìk ke tsun oe mì\-ftxe\-le tsngi\-vaw\-vìk, sko sa'\-nok tsun.} As Tsahik, I cannot weep over this matter, as mother I can.~\LINK{http://naviteri.org/2012/03/spring-vocabulary-part-2/}{b}
\B{Tsa\-hìk\-ur txe\-le lu.} This is a matter for the Tsahik.~\LINK{http://wiki.learnnavi.org/index.php?title=Na\%27vi_from_Avatar_Movie\#Neytiri_talking_to_Tsu.27tey}{a\slash wiki}

% tsakem / tsakem si
\hi
\HT{tsakem}{\B{\cp{tsa}kem}} \I{["\t{ts}a.kEm]} \emph{or} \B{tsa\cp{kem}} \I{[\t{ts}a."kEm]} \textsc{i}. \emph{dem., pn.} that (action). 
\B{Tsa\-kem kop krr\-nekx.} That also takes time.~\LINK{http://forum.learnnavi.org/language-updates/just-a-few-interesting-little-words/}{ln} 
\B{Txo tsa\-kem li\-ven mì ftaw\-nem\-krr.} If that happens in the past.~\LINK{http://naviteri.org/2013/01/awvea-posti-zisita-amip-first-post-of-the-new-year/comment-page-1/\#comment-1910}{b} \\ 
\textsc{ii}. \B{tsakem si} \I{["\t{ts}a.kEm si]} \emph{vin.} do that (action). 
\B{Tsa\-kem so\-li oe nga\-ru.} I did it to you.~\LINK{http://forum.learnnavi.org/language-updates/i-have-something-to-say/}{ln}
\B{Sìl\-pey oe tsnì tsì\-ye\-vun nì\-'i'\-a tsa\-kem si\-vi fì\-pì\-lok\-fa.} I hope that I can finally do that through this blog.~\LINK{http://naviteri.org/2010/06/first-post/}{b}
\B{Zun ay\-oe li\-vu tsam\-si\-yu, zel tsa\-kem ke sim\-vi.} If we were warriors, we wouldn't have done that.~\LINK{http://naviteri.org/2013/04/zun-zel-counterfactual-conditionals/}{b}
(\emph{a. proverb}) \B{Kem\-ìri a nga\-ru prr\-te' ke lu, tsa\-kem rä\-'ä si\-vi ay\-la\-he\-ru.} As for any action that you don't like, don't act that way to others.~\LINK{http://wiki.learnnavi.org/index.php?title=Corpus\#Good_Morning_America}{wiki}
(\ding{43}~\HL{kem}{\B{kem}}, \HL{tsa-}{\B{tsa-}})

% tsakrr 
\hi
\HT{tsakrr}{\B{tsa\cp{krr}}} \I{[\t{ts}a."k\s{r}]} \textsc{i}. \emph{n.} that time. 
\B{Mun\-trr\-am oe ko\-lä\-teng hu Ra\-lu ul\-te su\-nu oer tsa\-krr.} I spent last weekend with Ralu and had a good time [lit.: that time was pleasant to me].~\LINK{http://naviteri.org/2015/08/aylifyavi-lereyfya-2-cultural-terms-2-and-more/}{b} \\ 
\textsc{ii}. \emph{adv.} at that time, then.
\B{Tì\-'eyng\-it oel to\-lel a krr, tsa\-krr pa\-ye\-'un swey\-a fya\-'ot a za\-mi\-vu\-nge oel ay\-ngar ay\-lì\-'ut ho\-ren\-ti\-sì lì'\-fya\-yä le\-Na'\-vi.} When I receive an answer, I will then decide the best way to bring you the words and rules of Na'vi.~\LINK{http://forum.learnnavi.org/news-announcements/a-response-from-paul-frommer!/}{ln}
\B{Oe\-yä ik\-ran sli\-vu nga, tsa\-krr oeng 'aw\-si\-teng mi\-vak\-to.} Be my ikran and let's ride together.~\LINK{http://wiki.learnnavi.org/index.php?title=Corpus\#Lemondrop.com}{ld}
\B{Tsa\-krr syay ay\-nge\-yä, syay o\-lo'\-ä oe\-yä la\-yu teng.} Then you will suffer the same fate as my clan.~\LINK{http://naviteri.org/2010/07/vocabulary-update/}{b} 
(\emph{a. optionally with} \B{txo}) \B{Pxan li\-vu txo nì\-'aw oe nga\-ri | Tsa\-krr nga Na'\-vi\-ru yom\-tì\-yìng.} Only if I am worthy of you | Will you feed the People.~\LINK{http://wiki.learnnavi.org/index.php/Hunting_Song}{wiki}
(\emph{b. proverb}) \B{Txo ke nì\-yo' tsa\-krr nì\-yol.} `If you can't be flawless, at least be brief.'~\LINK{http://naviteri.org/2012/01/mipa-zisit-ayliu-amip-new-words-for-the-new-year/}{b}
(\ding{43}~\HL{krr}{\B{krr}}, \HL{tsa-}{\B{tsa-}}) \\
\textsc{Derivatives}: \ding{43}
\on{1}.~\HL{tsakrrvay}{\B{tsakrrvay}}, until then, in the meantime.

% tsakrrvay
\hi
\HT{tsakrrvay}{\B{tsa\cp{krr}vay}} \I{[\t{ts}a."k\s{r}.vaj]} (\emph{a. adv.}) until then, in the meantime.
\B{Tsa\-krr\-vay, ay\-nge\-yä tìm\-wey\-pey\-ri ira\-yo sei\-yi oe, ul\-te fì\-trr\-ä ftxo\-zä\-ri, sìl\-pey oe, ay\-nga\-ru prr\-te' li\-vu.} In the meantime, I thank you for your patience, and I hope you enjoy today's celebration.~\LINK{http://forum.learnnavi.org/language-updates/trr-rrtaya/}{ln}
(\emph{b. conv.}) until then, see you then. \B{Ul\-txa si\-vi oeng sìn ram\-tsyìp txon\-'ong\-ay. -- Ye\-io'! Tsa\-krr\-vay ko!} Let's meet on the hill tomorrow at nightfall. -- Perfect! See you then.~\LINK{http://naviteri.org/2012/01/mipa-zisit-ayliu-amip-new-words-for-the-new-year/}{b}
(\ding{43}~\HL{tsakrr}{\B{tsa\-krr}}, \HL{vay}{\B{vay}})

% tsaktap / si
\hi 
\HT{tsaktap}{\B{\cp{tsak}tap}} \I{["\t{ts}ak\textcorner.tap\textcorner]} \textsc{i}. \emph{n.} violence. 
\B{Slä set tä\-ngok kop tan\-tsaw\-tsray\-ti ay\-oe\-yä 'uol a\-la\-he: tsak\-tap.} But now, something else is also at our large cities: violence.~\LINK{http://naviteri.org/2020/06/mi-tanlokxe-oeya-srr-afpxamo-terrible-days-in-my-country/}{b} \\
\textsc{ii}. \B{\cp{tsak}tap si} \I{["\t{ts}ak\textcorner.tap\textcorner{} si]} \emph{vin.} be violent, use violence.
\B{Tsak\-tap r\"a'\"a si kaw\-krr mung\-wrr\-txo ke li\-vu kea fya\-'o a\-la\-he.} Never use violence unless there is no other way. ~\LINK{http://naviteri.org/2020/02/some-words-for-leap-year-day/}{b} \\
\textsc{Derivatives}: \ding{43}
\on{1}.~\HL{letsaktap}{\B{le\-tsak\-tap}}, violent.

% tsal
\hi
\B{tsal} \I{[\t{ts}al]} : \HL{tsa'u}{\B{tsa'u}}

% tsala
\hi
\HT{tsala}{\B{\cp{tsa}la}}, \B{tsa\-\cp{'ul} a}, \emph{or} \B{a tsa\-\cp{'ul}} \I{["\t{ts}a.la]} \emph{conj.} that (\emph{to introduce an agentive subordinate clause}). 
(\ding{43}~\HL{fula}{\B{fu\-la}})

% tsalsungay
\hi
\HT{tsalsungay}{\B{tsalsu\cp{ngay}}} \I{[\t{ts}al.su."Naj]} \emph{adv.} nevertheless, even so.
\B{Slä tsal\-su\-ngay, txo tì\-tslam for, lu txa\-yo na'\-rìng\-to sìl\-tsan.} But even so, if they're smart they'll take open terrain over bush.~\LINK{http://naviteri.org/2011/09/\%E2\%80\%9Cby-the-way-what-are-you-reading\%E2\%80\%9D/}{b}
\B{Ta\-fral kì\-syar a ay\-sä\-num\-vi sì ay\-lì'\-fya\-vi lu fyin. Tsal\-su\-ngay sìl\-pey oe, nu\-mul\-txa le\-sar lì\-ye\-vu.} That's why I'll teach lessons and expressions that are simple. Nevertheless, I hope the class will be useful.~\LINK{http://naviteri.org/2012/07/tskxekengtsyip-a-mikyunfpi-a-little-listening-exercise/}{b} 
\B{Hu\-fwa lu fil\-ur Va'\-ru\-ä fne\-fe'\-ran\-vi, tsal\-su\-ngay fpìl fu\-ta say\-rìp lu nì\-txan.} Although Va'ru's facial stripes are rather uneven, I still think he's very handsome.~\LINK{http://naviteri.org/2012/10/mipa-vospxi-mipa-ayliu-new-words-for-the-new-month/}{b} 

% tsam
\hi
\HT{tsam}{\B{tsam}} \I{[\t{ts}am]} \emph{n.} war. 
\B{Hu Saw\-tu\-te a tsam lu ha\-sey.} The war with the Skypeople is over.~\LINK{http://naviteri.org/2010/07/ting-mikyun-fte-tslivam-listening-comprehension-1-3/}{b} 
\B{Sra\-ke fpìl ay\-ngal fu\-ta aw\-nga tsun hu Saw\-tu\-te a tsam\-it 'i\-vaw\-nìm fa fwa päng\-kxo fo\-hu nì\-'aw, ma smuk?} Do you think we can avoid a war with the Sky People solely by talking with them, my brothers?~\LINK{http://naviteri.org/2012/10/snewsyea-ftxoza-halowina-a-spooky-halloween/comment-page-1/\#comment-1755}{b}
\B{Po\-ri tì\-kak\-rel tì\-kak\-pam\-sì kum tsam\-ä lu.} His blindness and deafness are a result of the war.~\LINK{http://naviteri.org/2014/05/mipa-ayliu-mipa-aysafpil-new-words-new-ideas/}{b} \\
\textsc{Derivatives}: \ding{43}
\on{1}.~\HL{tsamsiyu}{\B{tsam\-si\-yu}}, warrior. 
\on{2}.~\HL{tsampongu}{\B{tsam\-po\-ngu}}, war party. 
\on{3}.~\HL{tsamopin}{\B{tsa\-mo\-pin}}, war paint. 
\on{4}.~\HL{tsamkuk}{\B{tsam\-kuk}}, war drum.
\on{5}.~\HL{tsamsA'o}{\B{tsam\-sä\-'o}}, weapon of war.

% tsame-
\hi 
\HT{tsame}{\B{tsame}\texttt{+}} \I{[\t{ts}a.mE.]} \emph{prod. pref. on n.s} those two (there) (\emph{distal dual dem. pref.}). 
\B{Fpìl oel fu\-ta tsun ay\-nga tsli\-vam teyng\-ta tsa\-me\-stxo za\-'u ftu pe\-sim.} I think you understand where those two names come from.~\LINK{http://naviteri.org/2015/09/tskxekengtsyip-a-mikyunfpi-a-little-listening-exercise-2/}{b}
(\ding{43}~\HL{tsa-}{\B{tsa-}}; \HL{me}{\B{me}\texttt{+}}) 

% tsamkuk
\hi
\HT{tsamkuk}{\B{\cp{tsam}kuk}} \I{["\t{ts}am.kuk]} \emph{or} \I{[.kUk]} (\textsc{rn}: \B{\cp{tsam}kùk} \I{["\t{ts}am.kUk\textcorner]}) \emph{n.} war drum. 
(\ding{43}~\HL{tsam}{\B{tsam}}, \HL{takuk}{\B{ta\-kuk}})

% tsamopin
\hi
\HT{tsamopin}{\B{\cp{tsa}mopin}} \I{["\t{ts}a.mo.pin]} \emph{n.} warpaint.
(\ding{43}~\HL{tsam}{\B{tsam}}, \HL{'opin}{\B{'o\-pin}})

% tsampongu
\hi
\HT{tsampongu}{\B{\cp{tsam}pongu}} \I{["\t{ts}am.po.Nu]} \emph{n.} war party. 
\B{Set sre\-fey oel fu\-ta tsam\-po\-ngu tä\-txaw maw txon\-'ong.} I'm currently expecting the war party back after nightfall.~\LINK{http://naviteri.org/2011/02/new-vocabulary-part-2/}{b} 
\B{Lu tsa\-tsam\-si\-yu le\-'aw\-a ha\-pxì\-tu tsam\-po\-ngu\-ä a mal lu moer.} That warrior is the only member of the war party that we both trust.~\LINK{http://naviteri.org/2012/01/more-additions-to-the-lexicon/}{b}
(\ding{43}~\HL{tsam}{\B{tsam}}, \HL{pongu}{\B{po\-ngu}})

% tsamsä'o
\hi
\HT{tsamsA'o}{\B{\cp{tsam}sä'o}} \I{["\t{ts}am.s\ae.Po]} \emph{n.} weapon of war.
\B{Po\-ri ay\-sam\-sä\-'o a\-tì\-txur\-nga' fra\-to mì hi\-fkey la\-yu mì syokx.} He will have the most powerful weapons of war in the world in his hands.~\LINK{http://naviteri.org/2016/12/tengkrr-perahem-zisit-amip-as-the-new-year-arrives/}{b}
(\ding{43}~\HL{tsam}{\B{tsam}}, \HL{sA'o}{\B{sä'o}})

% tsamsiyu
\hi
\HT{tsamsiyu}{\B{\cp{tsam}siyu}} \I{["\t{ts}am.si.ju]} \emph{n.} warrior. 
\B{Oe lu tsam\-si\-yu.} I am a worrior.~\LINK{http://naviteri.org/2010/09/getting-to-know-you-part-2/}{b} 
\B{Ke lu po tsam\-si\-yu kaw\-'it!} There's not a warrior's bone in his whole body!~\LINK{http://forum.learnnavi.org/language-updates/ke-kawit/msg174572/\#msg174572}{ln}
\B{Nì\-trr\-trr pxì\-mun\-'i sam\-si\-yul ay\-swi\-zaw\-it ku\-tu\-ä a\-law\-nä\-txayn sno\-kip nì\-'eng.} Warriors typically share the arrows of their defeated enemies among themselves.~\LINK{http://naviteri.org/2012/01/more-additions-to-the-lexicon/}{b}
\B{Po\-ri wem\-tswo fra\-tsam\-si\-yur ro\-lo\-'a nì\-txan.} His ability to fight greatly impressed all the warriors.~\LINK{http://naviteri.org/2012/03/spring-vocabulary-part-2/}{b}
\B{Tsran\-ten fra\-to a syon tsam\-si\-yu\-ä lu tì\-tstew.} The most important characteristic of a warrior is bravery.~\LINK{http://naviteri.org/2012/10/mipa-vospxi-mipa-ayliu-new-words-for-the-new-month/}{b}
\B{Tsam\-si\-yu\-ri lu tì\-yo\-ra'\-ä tì\-new le\-kin.} A warrior must have the desire for victory.~\LINK{http://naviteri.org/2014/11/twenty-before-the-holidays/}{b}
\B{Lar\-mu tsa\-tsam\-si\-yu\-hu tì\-vawm a lu stxong ay\-oer.} That warrior carried with him a darkness unknown to us.~\LINK{http://naviteri.org/2010/07/vocabulary-update/}{b}
(\ding{43}~\HL{tsam}{\B{tsam}})

% tsan'ul
\hi
\HT{tsan'ul}{\B{\cp{tsan}'ul}} \I{["\t{ts}an.P$\cdot$ul]} \emph{vin., inf.}\on{22} improve, get better.
\B{Saw\-tu\-te\-ri tì\-fkey\-tok ke tsan\-'o\-lul kaw\-'it.} The situation with the Skypeople hasn't improved one bit.~\LINK{http://naviteri.org/2013/03/tsanerul-feerul-getting-better-getting-worse/}{b}
\B{Lì'\-fya\-ri le\-Na'\-vi nga tsan\-'e\-re\-i\-ul fra\-trr.} I'm delighted that your Na'vi is improving every day.~\LINK{http://naviteri.org/2013/03/tsanerul-feerul-getting-better-getting-worse/}{b}
\B{Ru\-txe fì\-tì\-oeyk\-tìng\-it tsan\-'ey\-ki\-vul. Ke lu law kaw\-'it.} Please improve this explanation. It's not at all clear.~\LINK{http://naviteri.org/2013/03/tsanerul-feerul-getting-better-getting-worse/}{b}
(\emph{a. rhyming expression with stress change}) \B{Tslikx, tì\-ran, tul; | Ftu yì ne yì tsan\-\cp{'ul}.} \emph{when learning s.t. new, you have to proceed from step to step: baby steps first, then bigger ones} [lit.: `Crawl, walk, run; from level to level get better.']~\LINK{http://naviteri.org/2023/09/aawa-ayliu-si-aylifyavi-amip-a-few-new-words-and-expressions/}{b}
(\ding{214}~\B{fe\-'ul}, \ding{43}~\B{sìl\-tsan}, \B{'ul}) \\
\textsc{Derivatives}: \ding{43} 
\on{1}.~\HL{tItsan'ul}{\B{tì\-tsan\-'ul}}, improvement (abstract).
\on{2}.~\HL{sAtsan'ul}{\B{sä\-tsan\-'ul}}, improvement (specific).

% tsankum
\hi
\HT{tsankum}{\B{\cp{tsan}kum}} \I{["\t{ts}an.kum]} \emph{or} \I{[.kUm]}, \emph{coll. pronounced as} \I{["\t{ts}aN.kum]} \emph{n.} advantage, benefit, upside, gain.
\B{Tsa\-kem\-ìri tsan\-kum lu law.} The advantage of that action is clear.~\LINK{http://naviteri.org/2015/08/aylifyavi-lereyfya-2-cultural-terms-2-and-more/}{b} 
(\ding{43}~\HL{kum}{\B{kum}}, \ding{214}~\HL{fekum}{\B{fe\-kum}}) \\
\textsc{Derivatives}: \ding{43}
\on{1}.~\HL{tsankumnga'}{\B{tsan\-kum\-nga'}}, adventageous.

% tsankumnga'
\hi
\HT{tsankumnga'}{\B{\cp{tsan}kumnga'}} \I{["\t{ts}an.kum.NaP]} \emph{or} \I{[.kUm.]}, \emph{coll. pronounced as} \I{["\t{ts}aN.kum.NaP]} \emph{adj.} advantageous.
(\ding{43}~\HL{tsankum}{\B{tsan\-kum}}, \ding{214}~\HL{fekumnga'}{\B{fe\-kum\-nga'}})

% tsantu
\hi 
\HT{tsantu}{\B{\cp{tsan}tu}} \I{["\t{ts}an.tu]} \emph{n.} good person, `good guy'. 
\B{Lal\-a tsa\-rel\-mì a\-ru\-sikx, yem\-stokx tsan\-tul haw\-re'\-ti a\-teyr, kawng\-tul pum\-it a\-la\-yon.} In that old movie, the good guy wears a white hat, the bad guy a black one.~\LINK{http://naviteri.org/2021/12/zolau-niprrte-ma-3746-welcome-2022/}{b} 
(\ding{214}~\HL{kawngtu}{\B{kawng\-tu}}; \ding{43}~\HL{sIltsan}{\B{sìl\-tsan}}, \HL{tute1}{\B{tu\-te}})

% tsantxäl / si
\hi 
\HT{tsantxAl}{\B{tsan\cp{txäl}}} \I{[\t{ts}an."t'\ae{}l]} \textsc{i}. \emph{n.} invitation. 
\B{Nge\-yä fì\-tsan\-txäl\-ìri a\-tì\-tstun\-wi\-nga' ira\-yo, slä ke tsä\-ngun oe zi\-va\-'u.} Thank you for this kind invitation, but unfortunately I cannot come.~\LINK{http://naviteri.org/2021/12/zolau-niprrte-ma-3746-welcome-2022/}{b} \\
\textsc{ii}. \B{tsan\cp{txäl} si} \I{[\t{ts}an."t'\ae{}l si]} \emph{vin.} invite (\emph{with} \ding{43}~\HL{tsnI}{\B{tsnì}}).
\B{Po tsan\-txäl so\-li oer tsnì zi\-va\-'u kel\-ku\-ne.} She invited me to come to her home.~\LINK{http://naviteri.org/2021/12/zolau-niprrte-ma-3746-welcome-2022/}{b}
(\ding{43}~\HL{sIlstan}{\B{sìl\-tsan}}, \HL{AtxAle}{\B{ä\-txä\-le}}) 

% tsang
\hi
\HT{tsang}{\B{tsang}} \I{[\t{ts}aN]} \emph{n.} a piercing. \\
\textsc{Derivatives}: \ding{43}
\on{1}.~\HL{miktsang}{\B{mik\-tsang}}, earring.
\on{2}.~\HL{ontsang}{\B{on\-tsang}}, nose ring.

% tsap'alute / si 
\hi
\HT{tsap'alute}{\B{tsap\cp{'a}lute}} \I{[\t{ts}ap\textcorner."Pa.lu.tE]} \textsc{i}. (\emph{a. n.}) apology (for, \emph{with the topic}).
\B{Fì\-kxey\-ey\-ri tsap\-'a\-lu\-te. Ngi\-an la\-he\-a lì\-'u alu Pol ke lu kxey\-ey.} Apologies for this mistake. However the other word `pol' is not a mistake.~\LINK{http://naviteri.org/2011/02/new-vocabulary-ii\%E2\%80\%94part-1/comment-page-1/\#comment-520}{ln}
(\emph{b. intj.}) ``(my) apologies''. \B{Tsap\-'a\-lu\-te, mì upxa\-re a fpo\-le' oel ay\-nga\-ru trr\-am lu kxey\-ey.} My apologies, in the message I sent you yesterday is a mistake.~\LINK{http://forum.learnnavi.org/language-updates/another-email-from-frommer/msg68133/\#msg68133}{b} \\
\textsc{ii}. \B{tsap\cp{'a}lute si} \I{[\t{ts}ap\textcorner."Pa.lu.tE si]} \emph{vin.} apologize. 
\B{Fì\-skxawng\-ìri tsap\-'a\-lu\-te se\-ngi oe.} I apologize for this moron.~\LINK{http://wiki.learnnavi.org/index.php?title=Corpus\#Times_Online}{wiki}
\B{Po tsap\-'a\-lu\-te so\-li nìk\-mar.} He apologized at the right time.~\LINK{http://naviteri.org/2011/10/more-vocabulary-a-bit-of-grammar/}{b}
\B{Nga tsap\-'alu\-te so\-li srak? -- Soli.} Did you apologize? -- (Yes,) I have.~\LINK{http://forum.learnnavi.org/language-updates/auxilary-verb-si-possessive-dative-krr/}{ln}

% tsapo
\hi
\HT{tsapo}{\B{\cp{tsa}po}} \I{["\t{ts}a.po]} \emph{dem. pn.} that one.
\B{Tu\-te a ke\-ftxo fra\-to lu tsa\-po a tì\-yawn\-ur lie ke so\-li kaw\-krr.} The saddest person of all is the one who has never experienced love.~\LINK{http://naviteri.org/2018/04/100a-liu-amip-64-new-words-part-2/}{b}
\B{Ri\-ni tsa\-po\-hu ho\-lum nì\-mal nì\-wotx.} Rini left with that guy without thinking twice about it.~\LINK{http://naviteri.org/2012/01/more-additions-to-the-lexicon/}{b}
(\emph{a. in contrastive comparisons}) \B{Fì\-kar\-yu alu fì\-po lu tsul\-fä\-tu; tsa\-kar\-yu alu tsa\-po lu \mbox{skxawng}.} \emph{This} teacher is a master; \emph{that} teacher is a fool.~\LINK{http://naviteri.org/2011/12/one-more-for-2011/}{b}
(\ding{43}~\HL{po}{\B{po}}, \HL{tsa-}{\B{tsa-}})

% tsapxe-
\hi 
\HT{tsapxe}{\B{tsapxe}\texttt{+}} \I{[\t{ts}a.p'E.]} \emph{prod. pref. on n.s} those three (there) (\emph{distal trial dem. pref.}). 
\B{Tsa\-pxe\-sam\-si\-yu lu pe\-su\slash tu\-pe?} Who are those three warriors?~\LINK{http://naviteri.org/2011/07/number-in-na\%E2\%80\%99vi/}{b}
(\ding{43}~\HL{tsa-}{\B{tsa-}}; \HL{pxe}{\B{pxe}\texttt{+}}) 

% tsar / tsaru
\hi
\B{tsar} \I{["\t{ts}aR]} \emph{or} \B{\cp{tsa}ru} \I{["\t{ts}a.Ru]}~:~\HL{tsa'u}{\B{tsa'u}}.

% tsara
\hi
\HT{tsara}{\B{\cp{tsa}ra}}, \B{tsa\-\cp{'ur} a}, \emph{or} \B{a tsa\-\cp{{'ur}}} \I{["\t{ts}a.Ra]} \emph{conj.} that (\emph{to introduce a dative subordinate clause}). 
(\ding{43}~\HL{fura}{\B{fu\-ra}}) 

% tsari
\hi
\B{\cp{tsa}ri} \I{["\t{ts}a.Ri]}~:~\HL{tsa'u}{\B{tsa'u}}.

% tsaria
\hi
\HT{tsaria}{\B{\cp{tsa}ria}}, \B{tsa\-\cp{'u}ri a}, \emph{or} \B{a tsa\-\cp{{'u}ri}} \I{["\t{ts}a.Ri.a]} \emph{conj.} that (\emph{to introduce a topic subordinate clause}). 

% tsat / tsati
\hi
\B{tsat} \I{[\t{ts}at\textcorner]} \emph{or} \B{\cp{tsa}ti} \I{["\t{ts}a.ti]} : \HL{tsa'u}{\B{tsa'u}}.

% tsata
\hi
\HT{tsata}{\B{\cp{tsa}ta}}, \B{tsa\-\cp{'ut} a}, \emph{or} \B{a tsa\-\cp{'ut}} \I{["\t{ts}a.ta]} \emph{conj.} that (\emph{to introduce a patientive subordinate clause}). 

% tsatu
\hi
\HT{tsatu}{\B{\cp{tsa}tu}} \I{["\t{ts}a.tu]} \emph{pn.} that person. 
(\ding{43}~\HL{tute1}{\B{tu\-te}}, \HL{tsa-}{\B{tsa-}})

% tsatseng
\hi 
\HT{tsatseng}{\B{\cp{tsa}tseng}} \I{["\t{ts}a.\t{ts}EN]} \emph{or} \B{tsa\cp{tseng}} \I{[\t{ts}a."\t{ts}EN]} \textsc{i}. \emph{n.} there, that place. 
\B{Tsa\-tse\-nge le\-hrr\-ap lu fì\-txan kum\-a tsa\-ne ke kä aw\-nga kaw\-krr.} That place is so dangerous we never go there.~\LINK{http://naviteri.org/2012/06/spring-vocabulary-part-3/}{b}
\B{Slä am\-'a\-lu\-ke nge\-yä ti\-re\-al tsa\-tse\-nget ta\-ye\-i\-ok.} But without a doubt your spirit will be there.~\LINK{http://naviteri.org/2012/07/fivospxiya-aylifyavi-amip-this-months-new-expressions-pt-2/comment-page-1/\#comment-1569}{b}
\B{Fwam\-pop\-ä tem\-rey tsa\-tseng\-mì ngä\-zìk lu nì\-wotx.} It's very hard for a tapirus to survive there.~\LINK{http://naviteri.org/2011/05/some-miscellaneous-vocabulary/}{b} 
\B{Tsun oe nga\-hu tsa\-tse\-nge\-ne ki\-vä, slä nul\-new fu\-ta si\-vop oeng nì\-txi\-lu\-ke.} I can go there with you, but I prefer to travel leisurely.~\LINK{http://naviteri.org/2011/10/more-vocabulary-a-bit-of-grammar/}{b}
\textsc{ii}. \emph{adv.} there. \B{Kil\-van\-ìri tsa\-tseng slo\-snep\-pe?} How wide is the river there?~\LINK{http://naviteri.org/2012/06/spring-vocabulary-part-3/}{b} 
(\ding{43}~\HL{tseng}{\B{tse\-nge}}, \HL{tsa-}{\B{tsa-}}) \\
\textsc{Derivatives}: \ding{43}
\on{1}.~\HL{kAsatseng}{\B{kä\-sa\-tseng}}, out there.

% tsatsenge
\hi
\B{tsa\cp{tse}nge} \I{[\t{ts}a."\t{ts}E.NE]} : \HL{tseng}{\B{tsatseng}, \B{tsenge}} 

% tsaw
\hi
\B{tsaw} \I{[\t{ts}aw]} : \HL{tsa'u}{\B{tsa'u}}

% tsawa
\hi
\HT{tsawa}{\B{\cp{tsa}wa}}, \B{tsa\-\cp{'u} a}, \emph{or} \B{a tsa\-\cp{{'u}}} \I{["\t{ts}a.wa]} \emph{conj.} that (\emph{to introduce a subjective subordinate clause}). 

% tsawke
\hi
\HT{tsawke}{\B{\cp{tsaw}ke}} \I{["\t{ts}aw.kE]} \emph{n., astr.} sun. 
\B{Tom\-pa\-yä ka\-to | Tsaw\-ke\-yä ka\-to | Ka\-tot tä\-ftxu oel | Nì\-ean nì\-rim.} The rhythm of the rain | The rhythm of the sun | I weave the rhythm | In blue and yellow.~\LINK{http://wiki.learnnavi.org/index.php/Weaving_Song}{wiki}
\B{Fä\-za'u tsaw\-ke krr\-pe?} When will the sun come up?~\LINK{http://naviteri.org/2012/01/more-additions-to-the-lexicon/}{b}
\B{Tsaw\-ke fpxe\-rä\-kìm (nem\-fa taw).} The sun is rising.~\LINK{http://naviteri.org/2012/01/more-additions-to-the-lexicon/}{b} \\
\textsc{Derivatives}: \ding{43}
\on{1}.~\HL{tsawkenay}{\B{tsaw\-ke\-nay}}, Alpha Centauri B. 
\on{2}.~\HL{tsawksyul}{\B{tsawk\-syul}}, sun lily.

% Tsawkenay
\hi 
\HT{tsawkenay}{\B{Tsawke\cp{nay}}} \I{[\t{ts}aw.kE."naj]} \emph{n., astr.} Alpha Centauri B. 
(\ding{43}~\HL{tsawke}{\B{tsaw\-ke}})

% tsawksyul
\hi 
\HT{tsawksyul}{\B{\cp{tsawk}syul}} \I{["\t{ts}awk\textcorner.sjul]} \emph{or} \I{[.sjUl]} \emph{n.} \ding{96} sun lily, \emph{stella lilliam}.~\LINK{https://forum.learnnavi.org/language-updates/expansion-to-flora-fauna-and-places-from-the-valley-of-mo'ara/}{ln} 
(\ding{43}~\HL{tsawke}{\B{tsaw\-ke}}, \HL{syulang}{\B{syu\-lang}})

% tsawl (tsawl slu)
\hi
\HT{tsawl}{\B{tsawl}} \I{[\t{ts}awl]} \textsc{i}. \emph{adj.} big (in size or stature); tall.
\B{Kel\-ku nge\-yä lu tsawl; pum oe\-yä lu hì\-'i.} Your house is large; mine is small.~\LINK{http://naviteri.org/2010/06/first-post/}{b}
\B{Frìp, frìp, ma Evi\-tsyìp, To\-ruk\-ä tsawl\-a kxa!} The big mouth of the Great Leonopteryx bites, bites, children.~\LINK{http://forum.learnnavi.org/language-updates/learnings-from-the-siatlla-song-translations/msg473566/\#msg473566}{ln}
\B{Oel hu Txe\-wì trr\-am na'\-rìng\-it tar\-mok, tso\-le\-'a syep\-tu\-tet a\-tsawl fra\-to mì sì\-rey.} Yesterday I was with Txewì in the forest and we saw the biggest Trapper I've ever seen.~\LINK{http://wiki.learnnavi.org/index.php?title=Corpus\#New_York_Times}{wiki}
\B{Txon\-am teng\-krr tar\-mì\-ran oe kxam\-lä na'\-rìng, sro\-ler eo ut\-ral a\-tsawl txewm\-a vrr\-tep.} Last night as I was walking through the forest, a frightening demon appeared in front of a big tree.~\LINK{http://naviteri.org/2012/03/spring-vocabulary-part-1/}{b} 
\B{Po oe\-to lu tsawl.} He's taller than I am.~\LINK{http://naviteri.org/2014/03/value-and-worth/}{b} \\
\textsc{ii}. \B{tsawl slu} \I{[\t{ts}awl slu]} \emph{vin.} grow. 
\B{Fì\-ut\-ral paw kil\-van\-lok nì\-'aw. Tsawl slu nì\-win nì\-txan.} This tree only grows (i.e., germinates, develops) near a river. It grows (i.e., gets big) very quickly.~\LINK{http://naviteri.org/2018/10/aysrr-ayvospxi-ayzisikrr-days-months-seasons/}{b} \\
\textsc{Derivatives}: \ding{43}
\on{1}.~\HL{tsawlhI'}{\B{tsawl\-hì'}}, size. 
\on{2}.~\HL{tsawlultxa}{\B{tsawl\-ul\-txa}}, large gathering, conference. 
\on{3}.~\HL{tsawlapxangrr}{\B{tsawl\-apxa\-ngrr}}, unidelta tree. 
\on{4}.~\HL{tsawsngem}{\B{tsaw\-sngem}}, muscular. 
\on{5}.~\HL{txantsawl}{\B{txan\-tsawl}}, huge, giant.
\on{6}.~\HL{tsawtsray}{\B{tsaw\-tsray}}, small or medium\hyp{}sized city.
\on{7}.~\HL{flrrtsawl}{\B{flrr\-tsawl}}, sailfin goliath. 

% tsawlapx
\hi
\HT{tsawlapx}{\B{\cp{tsawl}apx}} \I{["\t{ts}awl.ap']} : \HL{tsawlapxangrr}{\B{tsawl\-apxa\-ngrr}}

% tsawlapxangrr
\hi
\HT{tsawlapxangrr}{\B{\cp{tsawl}apxangrr}} \I{["\t{ts}awl.a.p'a.N\s{r}]} \emph{n.}, \ding{96} uni\-del\-ta tree, ``tall large root'', \emph{magellum deltoids}. (\emph{short form} \ding{43}~\B{tsawl\-apx}) \LINK{http://forum.learnnavi.org/language-updates/ewll/}{ln}
(\ding{43}~\HL{tsawl}{\B{tsawl}}, \HL{apxa}{\B{a\-pxa}}, \HL{ngrr}{\B{ngrr}})

% tsawlhì'
\hi
\HT{tsawlhI'}{\B{tsawl\cp{hì'}}} \I{[\t{ts}awl."hIP]} \emph{n.} size. 
(\ding{43}~\HL{tsawl}{\B{tsawl}}, \HL{hI'i}{\B{hì\-'i}})

% tsawlultxa 
\hi
\HT{tsawlultxa}{\B{tsawlul\cp{txa}}} \I{[\t{ts}awl.ul."t'a]} \emph{or} \I{[.Ul.]} (\textsc{rn}: \B{tsawlùl\cp{txa}} \I{[\t{ts}awl.Ul."da]}) \emph{n.} large gathering, conference.   
\B{Krr\-ka tsawl\-ul\-txa vo\-spxì\-am.} During the great meeting last month.~\LINK{http://naviteri.org/2014/08/stxeli-alor-a-beautiful-gift/}{b}
\B{Tsawl\-ul\-txa\-maw a Ay\-lì\-'u}, Post\hyp{}meetup Words.~\LINK{http://naviteri.org/2014/07/tsawlultxamaw-a-ayliu-post-meetup-words/}{b}
\B{Tsa\-tsawl\-ul\-txa\-ri alu SETI\-kon a\-2ve, no\-lu\-me oe nì\-txan te\-ri tì\-tsun\-slu tì\-rey\-ä a hi\-fkey\-mì a\-la\-he.} At the \textsc{seti}con \textsc{ii} conference I learned a lot about the possibility of life on other worlds.~\LINK{http://naviteri.org/2012/07/meetings-waterfalls-and-more/}{b}
(\ding{43}~\HL{tsawl}{\B{tsawl}}, \HL{ultxa}{\B{ul\-txa}})

% tsawn
\hi 
\HT{tsawn}{\B{tsawn}} \I{[\t{ts}awn]} \emph{vtr.} gather growing food from the forest; pick; (\emph{in agriculture}) harvest. \\
\textsc{Derivatives}: \ding{43}
\on{1}.~\HL{zIskrrtsawn}{\B{zì\-skrr\-tsawn}}, fall, autumn. 
\on{2}.~\HL{sAtsawn}{\B{sä\-tsawn}}, harvest.

% tsawng
\hi 
\HT{tsawng}{\B{tsawng}} \I{[\t{ts}awN]} \emph{vin.} shatter, break into pieces. 
\B{Ma sem\-pu, oey yom\-yo tso\-lawng!} Daddy, my plate broke!~\LINK{http://naviteri.org/2020/04/aawa-u-amip-a-few-new-things/}{b}
\B{Ma En\-tu, ngal lum\-pe ngey tsmu\-ke\-yä yom\-yot tsey\-ko\-lawng?} En\-tu, why did you break your sister's plate?~\LINK{http://naviteri.org/2020/04/aawa-u-amip-a-few-new-things/}{b}
\B{Kxew! Oey tsko tsìl\-mawng!} Damn! My bow just broke into pieces!~\LINK{https://naviteri.org/2025/11/kaltxi-nimun-hello-again/}{b}

% tsawsngem
\hi
\HT{tsawsngem}{\B{tsaw\cp{sngem}}} \I{[\t{ts}aw."sNEm]} \emph{adj.} muscular.
\B{Ak\-wey ke lu tsaw\-sngem kaw\-'it slä lu say\-rìp nì\-txan.} Akwey isn't at all muscular but he's very handsome.~\LINK{http://naviteri.org/2015/04/some-new-words-for-may-day/}{b}
(\ding{43}~\HL{tsngem}{\B{tsngem}}, \HL{tsawl}{\B{tsawl}})

% tsawtsray
\hi
\HT{tsawtsray}{\B{\cp{tsaw}tsray}} \I{["\t{ts}aw.\t{ts}Raj]} \emph{n.} small or medium\hyp{}sized city.
(\ding{43}~\HL{tsawl}{\B{tsawl}}, \HL{tsray}{\B{tsray}}) \\
\textsc{Derivatives}: \ding{43}
\on{1}.~\HL{txantsawtsray}{\B{txan\-tsaw\-tsray}}, large city, metropolis.

% tsä'
\hi
\HT{tsA'}{\B{tsä'}} \I{[\t{ts}\ae{}P]} \emph{vin.} squirt. \\
\textsc{Derivatives}: \ding{43}
\on{1}.~\HL{txumtsA'wll}{\B{txum\-tsä'\-wll}}, baja tickler.

% tsäfpìlyewn
\hi 
\B{tsä\cp{fpìl}yewn} \I{[\t{ts}\ae."fpIl.jEwn]}~:~\HL{tIsAfpIlyewn}{\B{tì\-sä\-fpìl\-yewn}}

% tse
\hi
\HT{tse}{\B{tse}} \I{[\t{ts}E]} \emph{intj.} well … (\emph{to indicate resumption of conversation, or introduce a remark}). 
\B{Tse … za\-'u fra\-'u ne tu\-te lem\-wey\-pey, ke\-fyak?} Well, everything comes to the patient people, right?~\LINK{http://naviteri.org/2012/03/spring-vocabulary-part-2/}{b}

% tse'a
\hi
\HT{tse'a}{\B{tse\cp{'a}}} \I{[\t{ts}E."Pa]} \emph{vtr.} (\texttt{-}\emph{control}) see (\emph{physical sense}).
\B{Pe\-ut tse\-'a ngal?} What do you see?~\LINK{http://naviteri.org/2012/11/renu-ayinanfyaya-the-senses-paradigm/}{b}
\B{Oel tse\-'a 'a\-'aw\-a tu\-tet.} I see several people.~\LINK{http://naviteri.org/2010/07/vocabulary-update/}{b}
\B{Fpìr\-mìl oel fu\-ta ay\-nga na\-tsew tsi\-ve\-'a fì\-'ut.} I was just thinking that you might want to see this.~\LINK{http://forum.learnnavi.org/language-updates/the-evidential-at-last/msg106055/\#msg106055}{ln} 
\B{Spaw oel fu\-ta Mo\-'at\-ìl tso\-le\-'a Ney\-ti\-rit.} I believe Mo\-'at saw Ney\-ti\-ri.~\LINK{http://naviteri.org/2011/03/word-order-and-case-marking-with-modals/}{b}
\B{Ayfo tsä\-pa\-ye\-'a.} They will see themselves.~\LINK{http://forum.learnnavi.org/language-updates/the-reflexive-\%28from-a-frommer-email!!!!\%29/}{ln}
(\emph{a. idiom}) \B{Tse\-'a tìk yaw\-ne.} `Love at first sight.'~\LINK{http://naviteri.org/2021/12/zolau-niprrte-ma-3746-welcome-2022/}{b}
(\ding{43}~\HL{nIn}{\B{nìn}}, \HL{tIng}{\B{tìng na\-ri}}; \HL{kame}{\B{ka\-me}}) \\
\textsc{Derivatives}: \ding{43}
\on{1}.~\HL{tse'atswo}{\B{tse\-'a\-tswo}}, sight, vision. 
\on{2}.~\HL{srese'a}{\B{sre\-se\-'a}}, predict. 
\on{3}.~\HL{uniltsa}{\B{unil\-tsa}}, dream of\slash about.
\on{4}.~\HL{sAtse'a}{\B{sä\-tse\-'a}}, sight.

% tse'atswo
\hi
\HT{tse'atswo}{\B{tse\cp{'a}tswo}} \I{[\t{ts}E."Pa.\t{ts}wo]} \emph{n.} (sense of) sight, vision.
(\ding{43}~\HL{tse'a}{\B{tse\-'a}}, \HL{tsu'o}{\B{tsu\-'o}})

% tsefta / si
\hi 
\HT{tsefta}{\B{\cp{tse}fta}} \I{["\t{ts}E.fta]} \textsc{i}. \emph{n.} vengeance, revenge. 
\B{Omum oel fu\-ta ngal pot ve'\-kì, slä tse\-fta ke lu tì\-'eyng a\-muiä.} I know you hate him, but vengeance is not a proper response.~\LINK{http://naviteri.org/2022/11/krr-ao-an-exciting-time/}{b} \\
\textsc{ii}. \B{\cp{tse}fta si} \I{["\t{ts}E.fta si]} \emph{vin.} take revenge.
\B{Tu\-tan\-ur a eyk\-ta\-nay\-ti tspo\-lang oe\-yä sem\-pul tse\-fta sa\-yi.} My father will take revenge on the man who killed the deputy.~\LINK{http://naviteri.org/2022/11/krr-ao-an-exciting-time/}{b} \\
\textsc{Derivatives}: \ding{43} \on{1}.~\HL{tseftanga'}{\B{tse\-fta\-nga'}}, vengeful. 

% tseftanga'
\hi 
\HT{tseftanga'}{\B{\cp{tse}ftanga'}} \I{["\t{ts}E.fta.NaP]} \emph{adj.} vengeful. 
\B{tu\-te a\-tse\-fta\-nga'}, vengeful person.~\LINK{http://naviteri.org/2022/11/krr-ao-an-exciting-time/}{b} 
\B{ay\-lì\-'u a\-tse\-fta\-nga'}, vengeful words.~\LINK{http://naviteri.org/2022/11/krr-ao-an-exciting-time/}{b} 
(\ding{43}~\HL{tsefta}{\B{tse\-fta}})

% tseltsul
\hi
\HT{tseltsul}{\B{\cp{tsel}tsul}} \I{["\t{ts}El.\t{ts}ul]} \emph{or} \I{[.\t{ts}Ul]} (\textsc{rn}: \B{\cp{tsel}tsùl} \I{["\t{ts}El.\t{ts}Ul]}) \emph{n.} whitewater rapids (\emph{countable, but only rarely}). 

% Tsenu 
\hi
\B{\cp{Tse}nu} \I{["\t{ts}E.nu]} \emph{n., female name}.
\B{Kal\-txì. Nga lu So\-rewn, ke\-fyak? Oer syaw fko Tse\-nu.} Hi. You're So\-rewn, right? I'm Tse\-nu.~\LINK{http://naviteri.org/2010/09/getting-to-know-you-part-1/}{b}
\B{'En si oe, lor\-a tsa\-fkxi\-le\-ri a\-pxay\-opin so\-lar Tse\-nul srok\-it a\-vo\-zam.} I would guess that Tsenu used a thousand (lit. \on{512}) beads for that beautiful multi\hyp{}colored bib necklace.~\LINK{http://naviteri.org/2013/07/pximaw-tsawlultxa-just-after-the-meet-up/}{b}
\B{Lu Tse\-nur yawn\-yew\-la a lam fwa Va'\-rul pot txì\-yìng.} Tse\-nu is broken hearted that Va'\-ru appears to be about to dump her.~\LINK{http://naviteri.org/2012/10/mipa-vospxi-mipa-ayliu-new-words-for-the-new-month/}{b}
\B{Tse\-nur Lo\-ak ke hä\-nga'.} Loak isn't good for Tsenu.~\LINK{http://naviteri.org/2012/10/mipa-vospxi-mipa-ayliu-new-words-for-the-new-month/}{b}

% tsenga
\hi
\HT{tsenga}{\B{\cp{tse}nga}} \I{["\t{ts}E.Na]} \emph{conj.} where, place where.
\B{Tseng\-a 'aw\-steng\-yä\-pem fì\-me\-kem\-yo lu mek\-tseng.} Where these two walls come together there's a gap.~\LINK{http://naviteri.org/2013/04/wheres-the-bathroom-and-other-useful-things/}{b}
(\ding{43}~\HL{tseng}{\B{tse\-nge}}, \HL{a}{\B{a}})

% tseng / tsenge
\hi
\HT{tseng}{\B{tseng}} \I{[\t{ts}EN]} \emph{or} \B{\cp{tse}nge}\I{["\t{ts}E.NE]} \textsc{i}. \emph{n.} place. 
\B{Lu poe tstun\-wi, ul\-te evi\-ru peng el\-tur tì\-txen si a ay\-vur\-it te\-ri 'Rr\-ta, alu tse\-nge a\-stxong nì\-ngay.} She is kind, and tells the kids interesting stories about Earth, a truly strange place.~\LINK{http://naviteri.org/2010/07/ting-mikyun-fte-tslivam-listening-comprehension-1-3/}{b}  \\
\textsc{ii}. \B{tìng tseng} \I{[tIN \t{ts}EN]} \emph{vin.} back down.  \\
\textsc{iii}. \B{\cp{tseng}pe} \I{["\t{ts}EN.pE]} \emph{or} \B{pe\cp{seng}} \I{[pE."{}sEN]} \emph{inter.} where? what place? 
\B{Nga z(ol)a\-'u ftu tseng\-pe\slash tseng\-pe\-ftu?} Where are you from?~\LINK{http://naviteri.org/2010/09/getting-to-know-you-part-2/}{b} \\
\textsc{Derivatives}: \ding{43}
\on{1}.~\HL{fItseng}{\B{fì\-tseng(e)}}, here. 
\on{2}.~\HL{tsatseng}{\B{tsa\-tseng(e)}}, there. 
\on{3}.~\HL{tsenga}{\B{tse\-nga}}, where, place where. 
\on{4}.~\HL{tsengo}{\B{tse\-ngo}}, somewhere. 
\on{5}.~\HL{fratseng}{\B{fra\-tseng}}, everywhere. 
\on{6}.~\HL{kawtseng}{\B{kaw\-tseng}}, nowhere.  
\on{7}.~\HL{klltseng}{\B{kll\-tseng}}, position. 
\on{8}.~\HL{numtseng}{\B{num\-tseng}}, school. 
\on{9}.~\HL{mektseng}{\B{mek\-tseng}}, gap, breach. 
\on{10}.~\HL{txeptseng}{\B{txep\-tseng}}, fire pit. 
\on{11}.~\HL{kxamtseng}{\B{kxam\-tseng}}, center. 
\on{12}.~\HL{zongtseng}{\B{zong\-tseng}}, safe place, refuge. 
\on{13}.~\HL{sngeltseng}{\B{sngel\-tseng}}, rubbish place, garbage dump. 
\on{14}.~\HL{sngA'itseng}{\B{sngä\-'i\-tseng}}, beginning, starting position. 
\on{15}.~\HL{fngA'tseng}{\B{fngä'\-tseng}}, restroom.
\on{16}.~\HL{txurtseng}{\B{txur\-tseng}}, fortification, fortress, bulwark.
\on{17}.~\HL{wantseng}{\B{wan\-tseng}}, (temporary) hiding place.

% tsengo
\hi
\HT{tsengo}{\B{\cp{tse}ngo}} \I{["\t{ts}E.No]} \emph{n.} somewhere.
\B{Sti\-wi rä\-'ä si, ma 'e\-veng! U\-van si mì se\-ngo a\-la\-he.} Don't be naughty, child! Play somewhere else.~\LINK{http://naviteri.org/2016/06/mrrvola-lifyavi-amip-forty-new-expressions/}{b}
(\ding{43}~\HL{tseng}{\B{tse\-nge}}) 

% tseo
\hi
\HT{tseo}{\B{\cp{tse}o}} \I{["\t{ts}E.o]} \emph{n.} art.
\B{tseo tì\-ru\-sol\-ä}, art of singing.~\LINK{http://forum.learnnavi.org/language-updates/life-death-and-gerunds/}{ln}
\B{Fì\-tseo\-ri ke tsun kaw\-tu pi\-vä\-hem tì\-yo'\-ne; tsran\-ten tì\-pä\-hem\-ä tì\-fmi nì\-'aw.} In this art it's impossible to arrive at perfection; the only thing that matters is the attempt to arrive there.~\LINK{http://naviteri.org/2012/01/mipa-zisit-ayliu-amip-new-words-for-the-new-year/}{b} \\
\textsc{Derivatives}: \ding{43}
\on{1}.~\HL{tseotu}{\B{tseo\-tu}}, artist. 
\on{2}.~\HL{pamtseo}{\B{pam\-tseo}}, music. 
\on{3}.~\HL{reltseo}{\B{rel\-tseo}}, visual art.
\on{4}.~\HL{tseovi}{\B{tseo\-vi}}, work of art.

% tseotu
\hi
\HT{tseotu}{\B{\cp{tse}otu}} \I{["\t{ts}E.o.tu]} \emph{n.} artist (\emph{generic term}).
\B{Fu\-ri\-a lo\-lu oer tsa\-skxom a tseo\-tur mok\-ri\-yä srung si, nit\-ram 'efu oe nì\-txan.} I am very happy that I had the chance to help the voice artist.~\LINK{http://naviteri.org/2015/05/cirque-du-soleils-toruk-the-first-flight-teaser-video-online/comment-page-1/\#comment-19897}{b}
(\ding{43}~\HL{tseo}{\B{tseo}}, \HL{-tu}{\B{-tu}})

% tseovi
\hi 
\HT{tseovi}{\B{\cp{tse}ovi}} \I{["\t{ts}E.o.vi]} \emph{n.} work of art. 
\B{Fì\-fkxi\-le lu tseo\-vi a\-lor.} This necklace is a beautiful work of art.~\LINK{https://naviteri.org/2024/09/pxevola-liu-amip-two-dozen-new-words/}{b} 
(\ding{43}~\HL{tseo}{\B{tseo}}) 

% tseri
\hi
\HT{tseri}{\B{\cp{tse}ri}} \I{["\t{ts}E.Ri]} \emph{vtr.} note, notice.
\B{Pey\-ral\-ä mik\-tsang\-it a\-mip ngal tso\-le\-ri srak?} Did you notice Pey\-ral's new earring?~\LINK{http://naviteri.org/2015/03/kap-si-ayunil-saylahe-threats-dreams-and-other-things/}{b}
\B{Oe\-ru txoa li\-vu. Ke tso\-le\-rä\-ngi oel fu\-ta nga\-ri kxe\-tse eo oe lu.} Forgive me. I didn't notice that your tail was in front of me.~\LINK{http://naviteri.org/2015/03/kap-si-ayunil-saylahe-threats-dreams-and-other-things/}{b}
(\ding{214}~\HL{yAkx}{\B{yäkx}}) \\
\textsc{Derivatives}: \ding{43}
\on{1}.~\HL{tItseri}{\B{tì\-tse\-ri}}, awareness, notice. 
\on{2}.~\HL{tIktseri}{\B{tìk\-tse\-ri}}, unawareness, lack of notice. 
\on{3}.~\HL{sAtseri}{\B{sä\-tse\-ri}}, observation, something noticed.

% tsewtx
\hi
\HT{tsewtx}{\B{tsewtx}} \I{[\t{ts}Ewt']} \emph{adj.} dirty. 
(\ding{214}~\HL{laro}{\B{la\-ro}})

% tsiki
\hi 
\HT{tsiki}{\B{\cp{tsi}ki}} \I{["\t{ts}i.ki]} \emph{n., fauna} reef tick.~\LINK{https://forum.learnnavi.org/language-updates/expansion-to-flora-fauna-and-places-from-the-valley-of-mo'ara/}{ln}

% tsim
\hi
\HT{tsim}{\B{tsim}} \I{[\t{ts}im]} \emph{n.} source.
\B{Fpìl oel fu\-ta tsun ay\-nga tsli\-vam teyng\-ta tsa\-me\-stxo za\-'u ftu pe\-sim.} I think that you can understand which source those two names come from.~\LINK{http://naviteri.org/2015/09/tskxekengtsyip-a-mikyunfpi-a-little-listening-exercise-2/}{b} \\  
\textsc{Derivatives}: \ding{43} 
\on{1}.~\HL{letsim}{\B{le\-tsim}}, original, unique. 
\on{2}.~\HL{nItsim}{\B{nì\-tsim}}, originally. 
\on{3}.~\HL{ingyentsim}{\B{ing\-yen\-tsim}}, mystery, riddle. 
\on{4}.~\HL{sngumtsim}{\B{sngum\-tsim}}, worrisome matter.   
\on{5}.~\HL{yayayrtsim}{\B{ya\-yayr\-tsim}}, s.t. confusing.
\on{6}.~\HL{paytsim}{\B{pay\-tsim}}, water fountain.
\on{7}.~\HL{lawnoltsim}{\B{law\-nol\-tsim}}, source of (great) joy.
\on{8}.~\HL{tItxeymsim}{\B{tìtxeymtsim}}, interest (\emph{countable}).

% tsin 
\hi 
\HT{tsin}{\B{tsin}} \I{[\t{ts}in]} \emph{n.} (\emph{a. body part}) nail. 
(\emph{b. of an animal}) claw. 
\B{Pa\-lu\-kan\-ìl oe\-yä po\-ti fa tsin tso\-lä\-ngupx.} Unfortunately my cat scratched her with its claw.~\LINK{http://naviteri.org/2020/07/vospxivol-lefpom-happy-august/}{b} 

% Tsireya
\hi 
\HT{Tsireya}{\B{Tsi\cp{rey}a}} \I{[\t{ts}i."REj.a]} \emph{n. female name} (\emph{short form:} \ding{43}~\HL{Reya}{\B{Re\-ya}}).

% tsì-
\hi
\HT{tsI-}{\B{tsì-}} \I{[\t{ts}I.]} \emph{num. pref.} four (\emph{base form of} \ding{43}~\HL{tsIng}{\B{tsìng}} \emph{to form higher numbers or fractions}). 
\B{\cp{tsì}vol}, \on{32} (decimal), \on{40} (octal); 
\B{tsì\cp{pxì}}, one\hyp{}fourth, $\frac{1}{4}$; 
\B{tsì\-pxì a\-pxey}, three\hyp{}fourth, $\frac{3}{4}$.

% tsìk
\hi
\HT{tsIk}{\B{tsìk}} \I{[\t{ts}Ik\textcorner]} \emph{adv.} suddenly, without warning. 
\B{'Uol ik\-ran\-it txo\-pu sley\-ko\-la\-tsu, ta\-lun\-a po tsìk ya\-wo.} Something must have frightened the banshee, because it suddenly took to the air.~\LINK{http://naviteri.org/2012/01/mipa-zisit-ayliu-amip-new-words-for-the-new-year/}{b}

% tsìlpey
\hi
\HT{tsIlpey}{\B{tsìl\cp{pey}}} \I{[\t{ts}Il."pEj]} \emph{n.} hope (\emph{abstract idea}).
\B{Tsìl\-pey\-lu\-ke ke tsun kaw\-tu ri\-vey.} No one can live without hope.~\LINK{http://naviteri.org/2013/01/awvea-posti-zisita-amip-first-post-of-the-new-year/}{b}
\B{Tsìl\-pey\-ìl tok txe'\-lan\-it.} Hope lives within the heart.~\LINK{http://naviteri.org/2013/01/awvea-posti-zisita-amip-first-post-of-the-new-year/}{b}
(\ding{43}~\HL{sIlpey}{\B{sìl\-pey}}, \HL{sAsIlpey}{\B{sä\-sìl\-pey}})

% tsìltsan
\hi 
\HT{tsIltsan}{\B{tsìl\cp{tsan}}} \I{[\t{ts}Il."\t{ts}an]} \emph{n} good (\emph{abstract concept}), goodness. 
\B{Tì\-kawng a su\-tel ngop var ri\-vey, tsìl\-tsan\-it pxìm kll\-yem fkol fe\-yä tä\-rem\-hu.} The evil that men do lives after them, The good is oft interred with their bones [\emph{Julius Caesar, Act III, Sc. 2}].~\LINK{http://naviteri.org/2017/09/zisikrr-amip-ayliu-amip-new-words-for-the-new-season/}{b}  
(\ding{43}~\HL{sIltsan}{\B{sìl\-tsan}})

% tsìng
\hi
\HT{tsIng}{\B{tsìng}} \I{[\t{ts}IN]} \emph{num., adj.} four. (\emph{base}: \ding{43}~\HL{tsI-}{\B{tsì-}}, \emph{rest}: \ding{43}~\HL{-sIng}{\B{\mbox{-sìng}}})
\B{Ay\-nga\-ru tsìng\-a yo\-ra'\-tut!} May I present the four winners!~\LINK{http://naviteri.org/2014/12/ayngaru-tsinga-yoratu-may-i-present-the-four-winners/}{b}
\B{Tsyokx\-ìri ke lu zek\-wä a\-tsìng ki a\-mrr.} The hand doesn't have four but five digits.~\LINK{http://naviteri.org/2010/07/ting-mikyun-fte-tslivam-listening-comprehension-1-3/}{b} \\
\textsc{Derivatives}: \ding{43}
\on{1}.~\HL{vospxItsIng}{\B{Vo\-spxì\-tsìng}}, April.

% tsìsyì 
\hi
\HT{tsIsyI}{\B{\cp{tsì}syì}} \I{["\t{ts}I.sjI]} \emph{vin.} whisper.   
\B{Ru\-txe pi\-vll\-txe nì\-wok nì\-'it, oel nga\-ti stum ke stä\-ngawm krr\-a nga fì\-fya tsye\-rì\-syì!} Please speak up a bit, I can barely hear you when you're whispering like this!~\LINK{http://naviteri.org/2011/09/miscellaneous-vocabulary/}{b} \\
\textsc{Derivatives}: \ding{43}
\on{1}.~\HL{sAtsIsyI}{\B{sä\-tsì\-syì}}, whisper.
\on{2}.~\HL{nItsyIsyI}{\B{nì\-tsì\-syì}}, by whispering.

% tsìve
\hi
\HT{tsIve}{\B{\cp{tsì}ve}} \I{["\t{ts}I.vE]} \emph{adj., ord. num.} fourth.

% tsìvol
\hi
\HT{tsIvol}{\B{\cp{tsì}vol}} \I{["\t{ts}I.vol]} \emph{num., adj.} \on{32} (decimal), \on{40} (octal).

% tskaha
\hi
\B{\cp{Tska}ha} \I{["\t{ts}ka.ha]} \emph{n.} \emph{name, family name}.
\B{Ay\-nge\-nga\-ru ohe\-yä tsmuk\-it alu Ne\-wey te Tska\-ha So\-rewn\-'ite.} Allow me to introduce my sister, Ne\-wey te Tska\-ha So\-rewn\-'ite. (highly formal)~\LINK{http://naviteri.org/2010/09/getting-to-know-you-part-1/}{b}

% tskalep
\hi
\HT{tskalep}{\B{\cp{tska}lep}} \I{["\t{ts}ka.lEp\textcorner]} \emph{n.}, \ding{246} crossbow.
\B{Maw txan\-tom\-pa, pxay\-a rìk\-äo la\-mu tska\-lep pe\-yä, ha tsat ku\-lat ay\-oel.} After the rainstorm, his crossbow was under a lot of leaves, so we uncovered it.~\LINK{http://naviteri.org/2011/10/more-vocabulary-a-bit-of-grammar/}{b}
\B{Tska\-lep\-it oel to\-lìng oe\-yä tsmu\-kan\-ur alu Ìstaw.} I gave the crossbow to my brother Istaw.~\LINK{http://naviteri.org/2010/07/vocabulary-update/}{b}

% tskawr
\hi 
\HT{tskawr}{\B{tskawr}} \I{[\t{ts}kawR]} \emph{vin.} limp. 
\B{Oel tse\-ri fu\-ta nga tskawr. Sra\-ke ngal ve\-nut tì\-sraw sey\-ko\-li?} I see you're limping. Did you hurt your foot?~\LINK{http://naviteri.org/2019/06/50a-liu-amip-40-new-words/}{b} 

% tski
\hi 
\HT{tski}{\B{tski}} \I{[\t{ts}ki]} \emph{n.} jaw. 

% tsko (tsko swizaw)
\hi
\HT{tsko}{\B{tsko}} \I{[\t{ts}ko]} \emph{n.}, \ding{246} bow (weapon). 
\B{Na\-ri si!  Tsa\-tsko lu spu\-win ul\-te ke lu mi txur, slä oe\-ri txan\-wa\-we le\-iu. Lu stxe\-li a sem\-pul\-ta.} Careful! That bow is old and no longer strong, but it means a lot to me. It was a gift from my father.~\LINK{http://naviteri.org/2011/05/some-miscellaneous-vocabulary/}{b}
\B{Tskot sngo\-lä\-'i po si\-var 'a\-'aw\-a trr\-kam.} He started to use the bow several days ago.~\LINK{http://naviteri.org/2011/09/miscellaneous-vocabulary/}{b}
\B{Fu\-ri\-a sne\-yä tsko\-ti ngop po kan\-fpìl.} He's concentrating on making his bow.~\LINK{http://naviteri.org/2011/10/more-vocabulary-a-bit-of-grammar/}{b} \\
(i.) \B{tsko swi\cp{zaw}} \I{[\t{ts}ko swi."zaw]} \emph{n.}, \ding{246} bow and arrow; archery. 
\B{Na'\-vi\-ri txin\-a sä\-'o tì\-tu\-sa\-ron\-ä lu tsko swi\-zaw.} For the Na'vi the bow and arrow is the main hunting tool.~\LINK{http://naviteri.org/2011/01/new-year-new-vocabulary/}{b}
\B{Fu\-ri\-a fì\-txan fnä\-ngan Ul\-rey\-ìl tsko swi\-zaw\-it, lu oe\-ru fmokx.} I'm jealous of the fact that Ul\-rey is such an excellent archer.~\LINK{http://naviteri.org/2011/10/more-vocabulary-a-bit-of-grammar/}{b}
\B{Tsko swi\-zaw\-fa a tì\-ta\-ron\-ìri ze\-ne fko ni\-vu\-me.} You must learn to hunt with a bow and arrow.~\LINK{http://forum.learnnavi.org/language-updates/ordinals-nume/}{ln\slash a\textsuperscript{c}}
\B{Oe tskxe\-keng si sä\-su\-lìn\-ur alu tsko swi\-zaw.} I practice my hobby, which is archery.~\LINK{http://naviteri.org/2010/09/getting-to-know-you-part-3/}{b}

% tskxaytsyìp
\hi
\HT{tskxaytsyIp}{\B{\cp{tskxay}tsyìp}} \I{["\t{ts}k'aj.\t{ts}jIp\textcorner]} \emph{n.} hail. 
(\ding{43}~\HL{tskxepay}{\B{tskxe\-pay}}, \HL{-tsyIp}{\B{-tsyìp}})

% tskxe
\hi
\HT{tskxe}{\B{tskxe}} \I{[\t{ts}k'E]} \emph{n.} rock, stone. 
\B{Ay\-skxe a mì sa\-sna\-'o ku'\-up lu nì\-txan.} The rocks in that pile are very heavy.~\LINK{http://naviteri.org/2012/03/spring-vocabulary-part-2/}{b}
\B{Lì\-'u\-ri alu tskxe pam\-rel fya\-pe?--Pam\-rel\-fya lu na Tsä, KxeKx, E.} How do you spell the word \emph{tskxe}?--It's spelled ts, kx, e.~\LINK{http://naviteri.org/2010/08/a-na\%E2\%80\%99vi-alphabet/}{b} 
\B{He\-tu\-wong\-ìl aw\-nge\-yä swo\-tut ska\-'a, fte kll\-ki\-vu\-lat ke\-ru\-sey\-a tskxet.} The aliens destroy our sacred place to dig up dead rock.~\LINK{http://forum.learnnavi.org/language-updates/participles/}{ln}
\B{Fì\-tskxe\-ri fa\-'o lu yey; ke lu ko\-um.} This rock has straight sides; it's not rounded.~\LINK{http://naviteri.org/2013/09/on-si-salewfya-shapes-and-directions/}{b}
\B{ngul na tskxe}, the drab color of stone.~\LINK{http://naviteri.org/2010/08/mipa-ayopin-mipa-ayliu-new-colors-new-words/}{b}
(\emph{a. idiom}) \B{Pay\-ìl a lipx tskxe\-ti ripx.} \emph{persistent effort can accomplish unexpected and amazing things} [lit.: `dripping water pierces a stone']~\LINK{http://naviteri.org/2020/04/aawa-u-amip-a-few-new-things/}{b} \\
\textsc{Derivatives}: \ding{43}
\on{1}.~\HL{letskxe}{\B{le\-tskxe}}, stony, of stone.
\on{2}.~\HL{tskxepay}{\B{tskxe\-pay}}, ice. 
\on{3}.~\HL{tskxevi}{\B{tskxe\-vi}}, pebble.
\on{4}.~\HL{lIngtskxe}{\B{lìng\-tskxe}}, unobtanium.
\on{5}.~\HL{tskxemauti}{\B{tskxe\-ma\-u\-ti}}, nut.

% tskxemauti
\hi 
\HT{tskxemauti}{\B{\cp{tskxe}mauti}} \I{["\t{ts}k'E.ma.u.ti]} \emph{n.} nut. 
(\ding{43}~\HL{tskxe}{\B{tskxe}}, \HL{mauti}{\B{mau\-ti}}) 

% tskxekeng / si 
\hi
\HT{tskxekeng}{\B{\cp{tskxe}keng}} \I{["\t{ts}k'E.kEN]} \textsc{i}. \emph{n.} training, exercise.
\B{Tskxe\-keng\-luke ke lu kea tì\-tsan\-'ul.} Without practice there is no improvement.~\LINK{http://naviteri.org/2013/03/tsanerul-feerul-getting-better-getting-worse/}{b}
\B{Mll\-te oe nì\-wotx: fì\-fne\-tskxe\-keng le\-sar lu nì\-txan.} I agree completely: this kind of exercise is very useful.~\LINK{http://naviteri.org/2014/08/stxeli-alor-a-beautiful-gift/comment-page-1/\#comment-9348}{b} \\
\textsc{ii}. \B{\cp{tskxe}keng si} \I{["\t{ts}k'E.kEN si]} \emph{vin.} exercise, train.
\B{Oe tskxe\-keng si sä\-su\-lìn\-ur alu tsko swi\-zaw.} I practice my hobby, which is archery.~\LINK{http://naviteri.org/2010/09/getting-to-know-you-part-3/}{b}
\B{'Ul tskxe\-keng si, 'ul fnan.} The more you practice, the better you'll get.~\LINK{http://naviteri.org/2012/02/trr-asawnung-lefpom-happy-leap-day/}{b}

% tskxepay
\hi
\HT{tskxepay}{\B{\cp{tskxe}pay}} \I{["\t{ts}k'E.paj]} \emph{n.} ice. 
\B{Sre fwa sìn tskxe\-pay tì\-ran, ze\-ne fko flì\-nutx\-it sti\-ve\-ftxaw.} It's necessary to check the thickness of the ice before walking on it.~\LINK{http://naviteri.org/2011/11/\%E2\%80\%99a\%E2\%80\%99awa-ayli\%E2\%80\%99u-amip-ni\%E2\%80\%99aw-\%E2\%80\%94-a-few-new-words-only/}{b}
(\ding{43}~\HL{tskxe}{\B{tskxe}}, \HL{pay}{\B{pay}}) \\
\textsc{Derivatives}: \ding{43}
\on{1}.~\HL{tskxaytsyIp}{\B{tskxay\-tsyìp}}, hail.

% tskxevi
\hi
\HT{tskxevi}{\B{\cp{tskxe}vi}} \I{["\t{ts}k'E.vi]} \emph{n.} pebble.
\B{'Eveng\-ìl skxe\-vit 'ey\-krr\-ko sko uvan.} Children roll pebbles as a game.~\LINK{http://naviteri.org/2015/04/some-new-words-for-may-day/}{b}
(\ding{43}~\HL{tskxe}{\B{tskxe}})

% tslam
\hi
\HT{tslam}{\B{tslam}} \I{[\t{ts}lam]} \emph{vtr.} understand. 
\B{Maw\-krra Tsyeyk ftxo\-lu\-lì\-'u, tslam fra\-pol fu\-ta slu po O\-lo'\-eyk\-tan a\-mip.} After Jake spoke, everyone understood that he had become the new Clan Leader.~\LINK{http://naviteri.org/2012/06/spring-vocabulary-part-3/}{b}
\B{Sra\-ke tsun oe fay\-upxa\-ret tsli\-vam nì\-ftue?} Can I understand these messages easily?~\LINK{http://forum.learnnavi.org/language-updates/verb-for-write/msg175592/\#msg175592}{ln}
\B{Hì\-txoa, ke tslolam. Ru\-txe li\-veyn.} Sorry, I didn't get that. Could you repeat it, please?~\LINK{http://naviteri.org/2010/09/getting-to-know-you-part-2/}{b}
\B{Ngal ke tslä\-ngam teyng\-ta fì\-tì\-ngä\-zìk\-ìri pe\-fyin\-ep\-'ang.} Unfortunately you don't understand how complex this problem is.~\LINK{http://naviteri.org/2012/07/fivospxiya-aylifyavi-amip-this-months-new-expressions-pt-2/}{b} 
(\emph{a. conv.}) \B{tslo\cp{lam}} ``understood''. 
\B{Oe\-ri lu swoa fne\-txum.--Tslo\-lam.} I am allergic to alcohol.--Understood.~\LINK{http://naviteri.org/2012/10/audio-and-video-learning-materials-for-navi-101/}{b} \\
\textsc{Derivatives}: \ding{43}
\on{1}.~\HL{tItslam}{\B{tì\-tslam}}, understanding, intelligence. 
\on{2}.~\HL{leytslam}{\B{ley\-tslam}}, appreciate.

% tsleng
\hi
\HT{tsleng}{\B{tsleng}} \I{[\t{ts}lEN]} \emph{adj.} false. 
\B{Fì\-po lu vrr\-tep a mì sokx a\-tsleng!} This one is a demon in a false body!~\LINK{http://wiki.learnnavi.org/index.php?title=Na\%27vi_from_Avatar_Movie\#Tsu.27tey_attacks_Jake}{wiki\slash a}
(\ding{214}~\HL{ngay}{\B{ngay}}) \\
\textsc{Derivatives}: \ding{43}
\on{1}.~\HL{tItsleng}{\B{tì\-tsleng}}, falsehood, falsity. 

% tslikx
\hi 
\HT{tslikx}{\B{tslikx}} \I{[\t{ts}lik']} \emph{vin.} crawl. 
(\emph{a. rhyming expression}) \B{Tslikx, tì\-ran, tul; | Ftu yì ne yì tsan\-'ul.} \emph{when learning s.t. new, you have to proceed from step to step: baby steps first, then bigger ones} [lit.: `Crawl, walk, run; from level to level get better.']~\LINK{http://naviteri.org/2023/09/aawa-ayliu-si-aylifyavi-amip-a-few-new-words-and-expressions/}{b} \\
\textsc{Derivatives}: \ding{43} 
\on{1}.~\HL{tslikxyulatopin}{\B{tslikx\-yu la\-to\-pin}}, chamelion crawler.   
\on{2}.~\HL{tslikxyutsawlak}{\B{tslikx\-yu tsaw\-lak}}, scarab crawler.

% tslikxyu latopin
\hi 
\HT{tslikxyulatopin}{\B{\cp{tslikx}yu \cp{la}topin}} \I{["\t{ts}lik'.ju "la.to.pin]} \emph{n., fauna} chamelion crawler. 
(\ding{43}~\HL{tslikx}{\B{tslikx}}, \HL{latem}{\B{la\-tem}}, \HL{'opin}{\B{'opin}})

% tslikxyu tsawlak
\hi 
\HT{tslikxyutsawlak}{\B{\cp{tslikx}yu \cp{tsaw}lak}} \I{["\t{ts}lik'.ju "\t{ts}aw.lak\textcorner]} \emph{n., fauna} scarab crawler. 
(\ding{43}~\HL{tslikx}{\B{tslikx}}, \HL{tsawl}{\B{tsawl}}, \HL{lak}{\B{lak}})

% tsmi
\hi
\HT{tsmi}{\B{tsmi}} \I{[\t{ts}mi]} \emph{n.} nectar. \\
\textsc{Derivatives}: \ding{43}
\on{1}.~\HL{tsmisnrr}{\B{tsmi\-snrr}}, nectar lantern, bladder lantern.

% tsmisnrr
\hi
\HT{tsmisnrr}{\B{\cp{tsmi}snrr}} \I{["\t{ts}mi.sn\s{r}]} \emph{n.} bladder lantern, nectar lantern. 
(\ding{43}~\HL{tsmi}{\B{tsmi}}, \HL{sAnrr}{\B{sä\-nrr}})

% tsmìm
\hi
\HT{tsmIm}{\B{tsmìm}} \I{[\t{ts}mIm]} \emph{n.} track (of animal or human origin).
\B{Äo fì\-ut\-ral lu tsmìm 'ang\-tsìk\-ä.} There's a hammerhead track under this tree.~\LINK{http://naviteri.org/2013/02/vospxi-ayol-posti-apup-short-post-for-a-short-month/}{b}
\B{Tsmìm\-ìri wätx fol ti\-nan\-it nì\-ngay.} They're really bad at reading animal tracks.~\LINK{http://naviteri.org/2011/09/\%E2\%80\%9Cby-the-way-what-are-you-reading\%E2\%80\%9D/}{b}
\B{Ioang\-ìri, tu\-te\-ri, tsun fko si\-var lì\-'ut alu tsmìm.} One can use the word `track' for animals and people.~\LINK{http://naviteri.org/2011/09/\%E2\%80\%9Cby-the-way-what-are-you-reading\%E2\%80\%9D/comment-page-1/\#comment-1118}{b}

% tsmuk
\hi
\HT{tsmuk}{\B{tsmuk}} \I{[\t{ts}muk\textcorner]} \emph{or} \I{[\t{ts}mUk\textcorner]} \emph{n.} sibling. 
\B{Fya\-pe fko syaw ngar, ma tsmuk?--Oe\-ru syaw Tsyey\-son.} What are you called, brother (or sister)?--I'm called Jason.~\LINK{http://naviteri.org/2012/10/audio-and-video-learning-materials-for-navi-101/}{b}
\B{Ay\-nga\-ru sei\-yi ira\-yo, ma smuk!} Thank you all, siblings!~\LINK{http://naviteri.org/2011/07/txantsana-ultxa-mi-siatll-great-meeting-in-seattle/}{b}
\B{Ira\-yo, ma pxe\-smuk, ul\-te sey\-kxel sì nit\-ram.} Thank you, siblings, and congratulation!~\LINK{http://naviteri.org/2010/12/navi-writing-contest-the-first-place-winners/}{b}
\B{Ma fra\-po, ay\-nga\-ru oe\-yä tsmuk\-it alu Ne\-wey.} Everybody, please allow me to introduce (to you) my sister, Newey.~\LINK{http://naviteri.org/2010/09/getting-to-know-you-part-1/}{b}
(\ding{43}~\HL{tsmuktu}{\B{tsmuk\-tu}}, \HL{tsmukan}{\B{tsmu\-kan}}, \HL{tsmuke}{\B{tsmu\-ke}})

% tsmukan
\hi
\HT{tsmukan}{\B{\cp{tsmu}kan}} \I{["\t{ts}mu.kan]} \emph{n.} brother. 
\B{Tsa\-tu\-te a ngal tse\-'a lu tsmu\-kan oe\-yä.} That person you see is my brother.~\LINK{http://naviteri.org/2014/11/twenty-before-the-holidays/comment-page-1/\#comment-15463}{b}
\B{Oe\-yä tsmu\-kan\-ìl mip\-a tsko\-ti kì\-ma\-ne\-i\-om.} My brother just got a new bow, I'm happy to say.~\LINK{http://naviteri.org/2012/01/more-additions-to-the-lexicon/}{b}
\B{Tska\-lep\-it oel to\-lìng oe\-yä tsmu\-kan\-ur alu Ìstaw.} I gave the crossbow to my brother Istaw.~\LINK{http://naviteri.org/2010/07/vocabulary-update/}{b}
(\ding{43}~\HL{tsmuk}{\B{tsmuk}})

% tsmuke
\hi
\HT{tsmuke}{\B{\cp{tsmu}ke}} \I{["\t{ts}mu.kE]} \emph{n.} sister. 
\B{Ke lu po\-ru kea smu\-ke.} He doesn't have any sisters.~\LINK{http://naviteri.org/2010/07/ting-mikyun-fte-tslivam-listening-comprehension-1-3/}{b}
\B{Tsa\-tsmu\-kel alu Ri\-ni mo\-lok fu\-ta oe ngar mu\-wä\-pi\-vìn\-txu.} Sister Rini over there suggested that I introduce myself.~\LINK{http://naviteri.org/2010/09/getting-to-know-you-part-1/}{b}
\B{Sa'\-nok\-ìl oe\-yä tsmu\-ket a\-mip no\-lokx trr\-am.} Mom gave birth to my new sister yesterday.~\LINK{http://naviteri.org/2010/09/getting-to-know-you-part-2/}{b}
(\ding{43}~\HL{tsmuk}{\B{tsmuk}})

% tsmuktu
\hi
\HT{tsmuktu}{\B{\cp{tsmuk}tu}} \I{["\t{ts}muk\textcorner.tu]} \emph{or} \I{["\t{ts}mUk\textcorner.]} \emph{n.} sibling.
\B{Rì\-'ìr rä\-'ä si\-vi tsmuk\-tur!} Don't imitate your sibling!~\LINK{http://naviteri.org/2011/09/miscellaneous-vocabulary/}{b}
(\ding{43}~\HL{tsmuk}{\B{tsmuk}})

% tsnì
\hi
\HT{tsnI}{\B{tsnì}} \I{[\t{ts}nI]} \emph{conj.} that (\emph{to introduce subordinate clauses after the v.s} \ding{43}~\HL{sIlpey}{\B{sìl\-pey}}, \HL{AtxAle}{\B{ätxä\-le si}}, \HL{Azan}{\B{ä\-zan si}}, \HL{rangal}{\B{rang\-al}}, \HL{mowar}{\B{mo\-war si}}, \HL{leymkem}{\B{leym\-kem}}, \HL{leymfe'}{\B{leym\-fe'}}, \HL{fe'pey}{\B{fe'\-pey}}, \HL{la'um}{\B{la\-'um}}, \HL{pe'ngay}{\B{pe'\-ngay}}, \HL{tsantxAl}{\B{tsan\-txäl si}} \emph{and optionally} \HL{srefpIl}{\B{sre\-fpìl}}, \HL{srefey}{\B{sre\-fey}}). 
\B{Ä\-txä\-le si tsnì li\-vu o\-he\-ru U\-nil\-ta\-ron.} I respectfully request the Dream Hunt.~\LINK{http://wiki.learnnavi.org/index.php?title=Corpus\#Behind_the_Scenes}{wiki} 
\B{Sìl\-pey oe tsnì fì\-tì\-oeyk\-tìng law li\-vu nga\-ru set.} I hope that this explanation is clear to you now.~\LINK{http://forum.learnnavi.org/language-updates/a-long-answer-to-a-quick-question/}{ln}
\B{Set sre\-fey oe tsnì tsam\-po\-ngu tä\-txaw maw txon\-'ong.} I'm currently expecting the war party back after nightfall.~\LINK{http://naviteri.org/2011/02/new-vocabulary-part-2/}{b}
\B{Poe mo\-war so\-li poan\-ur tsnì hi\-vum.} She advised him to leave.~\LINK{http://naviteri.org/2014/05/mipa-ayliu-mipa-aysafpil-new-words-new-ideas/}{b}
\B{Sre\-fpìl Oma\-ti\-ka\-ya tsnì Tsyeyk kaw\-krr ke ta\-yä\-txaw maw ka\-vuk sne\-yä.} The Omaticaya assumed that Jake would never return after his treachery.~\LINK{http://naviteri.org/2014/10/tson-si-fpomron-obligation-and-mental-health/}{b}
(\emph{proverb}) \B{Ätxä\-le si pa\-lu\-lu\-kan\-ur tsnì smar\-it li\-vo\-nu.} \emph{a futile gesture, an attempt to achieve something that might be desirable but will clearly not happen} [lit.: Ask a thanator to release its prey.]~\LINK{http://naviteri.org/2010/07/vocabulary-update/}{b}

% tsngal
\hi
\HT{tsgnal}{\B{tsngal}} \I{[\t{ts}Nal]} \emph{n.} cup. 
\B{Nge\-yä tsngal lum\-pe lu mek? Näk nì\-'ul ko!} Why is your cup empty? Drink up!~\LINK{http://naviteri.org/2012/02/trr-asawnung-lefpom-happy-leap-day/}{b}
\B{'Ä'! \mbox{Oel} tsngal\-it tì\-mung\-zup.} Whoops! I just dropped the cup.~\LINK{http://naviteri.org/2011/09/miscellaneous-vocabulary/}{b}  \\
\textsc{Derivatives}: \ding{43}
\on{1}.~\HL{tawtsngal}{\B{taw\-tsngal}}, panopyra.

% tsngan
\hi
\HT{tsngan}{\B{tsngan}} \I{[\t{ts}Nan]} \emph{n.} meat. 
\B{Ke\-tsran tu\-tel 'i\-vem, tsa\-fne\-tsngan lu ftxì\-vä'.} That kind of meat is gross no matter who cooks it.~\LINK{http://naviteri.org/2013/03/whoever-whatever-whenever/}{b}
\B{Fol tsngan\-it pxì\-mo\-lun\-'i nì\-'eng.} They shared the meat.~\LINK{http://naviteri.org/2012/01/more-additions-to-the-lexicon/}{b}
\B{Fol tsngan\-it pxì\-mo\-lun\-'i nì\-yll.} They shared the meat with the entire clan.~\LINK{http://naviteri.org/2012/01/more-additions-to-the-lexicon/}{b}
\B{Ke new oe mi\-vay' tsngan\-ti a 'o\-lem Ri\-nil.} I don't want to taste the meat that Rini cooked.~\LINK{http://naviteri.org/2012/11/renu-ayinanfyaya-the-senses-paradigm/}{b}
\B{Fì\-tsngan\-ur a sno\-läm lä\-ngu fa\-hew a\-kxä\-näng.} This rotten meat has a putrid smell.~\LINK{http://naviteri.org/2015/11/vomuna-liu-amip-ten-new-words/}{b}

% tsngawpay
\hi
\HT{tsngawpay}{\B{\cp{tsngaw}pay}} \I{["\t{ts}Naw.paj]} \emph{n.} tears.
\B{Mì Na'\-rìng lu tsngaw\-pay.} There are tears in the forest.~\LINK{http://naviteri.org/2017/07/way-tiretua-the-shamans-song/}{b}
(\ding{43}~\HL{tsngawvIk}{\B{tsngaw\-vìk}}, \HL{pay}{\B{pay}}) \\
\textsc{Derivatives}: \ding{43}
\on{1}.~\HL{tsngawpayvi}{\B{tsngaw\-pay\-vi}}, teardrop.

% tsngawpayvi
\hi
\HT{tsngawpayvi}{\B{\cp{tsngaw}payvi}} \I{["\t{ts}Naw.paj.vi]} \emph{n.} teardrop.
(\ding{43}~\HL{tsngawpay}{\B{tsngaw\-pay}})

% tsngawvìk
\hi
\HT{tsngawvIk}{\B{\cp{tsngaw}vìk}} \I{["\t{ts}Naw.vIk\textcorner]} \emph{vin.} cry, weep. 
\B{Sko Sa\-hìk ke tsun oe mì\-ftxe\-le tsngi\-vaw\-vìk, sko sa'\-nok tsun.} As Tsahik, I cannot weep over this matter, as mother I can.~\LINK{http://naviteri.org/2012/03/spring-vocabulary-part-2/}{b}
\B{Fpxa\-mo\-a fì\-kem\-ìl a\-fpxa\-mo ay\-oe\-ti tsngey\-ko\-law\-vìk, tsa\-krr ley\-ko\-leym\-kem.} This terrible, terrible action has made us cry, then protest.~\LINK{http://naviteri.org/2020/06/mi-tanlokxe-oeya-srr-afpxamo-terrible-days-in-my-country/}{b}

% tsngem
\hi
\HT{tsngem}{\B{tsngem}} \I{[\t{ts}NEm]} \emph{n.} muscle. 
\B{Lu pa'\-lir sngem a\-txur.} A direhorse has strong muscles.~\LINK{http://naviteri.org/2015/04/some-new-words-for-may-day/}{b} \\
\textsc{Derivatives}: \ding{43}
\on{1}.~\HL{tsawsngem}{\B{tsaw\-sngem}}, muscular.

% tson
\hi
\HT{tson}{\B{tson}} \I{[\t{ts}on]} \emph{n.} obligation, duty, imposed requirement, task.
\B{Lä\-ngu oe\-ru tson a fì\-fmawn\-it pi\-veng ngar.} I'm sorry that I'm obliged to tell you this news.~\LINK{http://naviteri.org/2014/10/tson-si-fpomron-obligation-and-mental-health/}{b}
\B{Za\-'u tsa\-tson ta Ey\-wa.} That obligation comes from Eywa.~\LINK{http://naviteri.org/2014/10/tson-si-fpomron-obligation-and-mental-health/}{b}
\B{Lu Ney\-ti\-rir ta Mo\-'at a tson a kar Tsyeyk\-ur ay\-fya\-'ot Na'\-vi\-yä.} Neytiri is under obligation by\slash from Mo'at to teach Jake the ways of the Na'vi.~\LINK{http://naviteri.org/2014/10/tson-si-fpomron-obligation-and-mental-health/}{b} 
\B{Po tì\-kang\-kem so\-li val nì\-txan fte tsa\-tson\-ur ha\-sey si\-vi.} She worked very hard to complete the task.~\LINK{http://naviteri.org/2022/12/trr-anawm-polaheiem-the-great-day-has-arrived/}{b} \\
\textsc{Derivatives}: \ding{43}
\on{1}.~\HL{nItson}{\B{nì\-tson}}, dutifully, as an obligation. 
\on{2}.~\HL{tsonta}{\B{tson\-ta}}, to (\emph{with} \HL{kxIm}{\B{kxìm}})

% tsonta
\hi
\HT{tsonta}{\B{\cp{tson}ta}} \I{["\t{ts}on.ta]} \emph{conj.} to (with \ding{43}~\HL{kxIm}{\B{kxìm}}), (be obliged) to do a task; (\emph{short form of} \B{tson\-it a})
\B{Ay\-eveng\-ur kxo\-lìm sa'\-nok\-ìl tson\-ta pay\-it za\-mu\-nge.} The children were told by their mother to fetch water.~\LINK{http://naviteri.org/2014/10/tson-si-fpomron-obligation-and-mental-health/}{b}
(\ding{43}~\HL{tson}{\B{tson}})

% tsong 
\hi
\HT{tsong}{\B{tsong}} \I{[\t{ts}oN]} \emph{n.} valley.   
\B{Aw\-nga tsong\-ne ki\-vä fte sti\-var\-sìm tey\-lut.} Let's go to the valley to gather beetle larvae.~\LINK{http://naviteri.org/2012/03/spring-vocabulary-part-1/}{b} \\
\textsc{Derivatives}: \ding{43}
\on{1}.~\HL{tsongropx}{\B{tsong\-ropx}}, hole, cavity. 
\on{2}.~\HL{tsongtsyIp}{\B{tsong\-tsyìp}}, dimple.

% tsongropx
\hi
\HT{tsongropx}{\B{\cp{tsong}ropx}} \I{["\t{ts}oN.Rop']} \emph{n.} hole, cavity, excavation with a bottom (\emph{visible or presumed}).
\B{Fì\-fne\-ya\-yol tsrul\-it txu\-la mì song\-ropx ut\-ral\-ä fte ay\-li\-nit hi\-vaw\-nu wä sar\-ni\-o\-ang. Fì\-tì\-kan\-ìri ropx ke ha'.} This kind of bird builds its nest in a tree cavity so as to protect its young from predators. For this purpose, a hole going right through the tree trunk isn't suitable.~\LINK{http://naviteri.org/2013/04/wheres-the-bathroom-and-other-useful-things/}{b}
(\ding{43}~\HL{tsong}{\B{tsong}}, \HL{ropx}{\B{ropx}})

% tsongtsyìp
\hi
\HT{tsongtsyIp}{\B{\cp{tsong}tsyìp}} \I{["\t{ts}oN.\t{ts}jIp\textcorner]} \emph{n.} (\emph{of the body}) dimple. 
\B{Prr\-nen lrr\-tok si a krr, srer me\-song\-tsyìp a\-ho\-na.} When the baby smiles, two adorable dimples appear.~\LINK{http://naviteri.org/2012/03/spring-vocabulary-part-1/}{b}
(\ding{43}~\HL{tsong}{\B{tsong}}, \HL{-tsyIp}{\B{-tsyìp}})

% tsopì 
\hi
\HT{tsopI}{\B{\cp{tso}pì}} \I{["\t{ts}o.pI]} \emph{n.} lung (\emph{of the body}). 

% tspang
\hi
\HT{tspang}{\B{tspang}} \I{[\t{ts}paN]} \emph{vtr.} kill. 
\B{Sna\-fpìl\-fya\-ri le\-Na'\-vi krra smar\-it fkol tspang, tsran\-ten nì\-txan fwa po ke ngä\-'än nì\-kel\-kin.} It's important in Na'vi philosophy that the prey not suffer unnecessarily when it's killed.~\LINK{http://naviteri.org/2013/02/vospxi-ayol-posti-apup-short-post-for-a-short-month/}{b}
\B{Tsa\-ik\-ran\-ìri ta\-lun\-a new nga\-ti tspi\-vang, law lu fwa mi ke lu zä\-fi.} Since that ikran wants to kill you, it's clear it's still not tame.~\LINK{http://naviteri.org/2013/01/awvea-posti-zisita-amip-first-post-of-the-new-year/}{b}
\B{Fì\-po\-ti oel tspì\-yang, fte tì\-ke\-nong li\-ye\-vu ay\-la\-ru.} I will kill this one as a lesson to the others.~\LINK{http://forum.learnnavi.org/language-updates/capturing-avatar/}{ln\slash a\textsuperscript{c}}
\B{Tsa\-ye\-rik\-tsyìp\-ìl li 'an\-la sa'\-nok\-it a fkol tspì\-mang.} That little hexapede is already yearning for its mother that's just been killed.~\LINK{http://naviteri.org/2013/01/awvea-posti-zisita-amip-first-post-of-the-new-year/}{b}
(\emph{a. idiom}) \B{Ke tsun fko tspi\-vang to\-ruk\-it fa fwa pewn\-ti snew.} \emph{Proverbial expression for a method that will not work.} [lit.: You can't kill a great leonopteryx by constricting its throat]~\LINK{http://naviteri.org/2012/10/snewsyea-ftxoza-halowina-a-spooky-halloween/}{b} \\
\textsc{Derivatives}: \ding{43}
\on{1}.~\HL{tspangyu}{\B{tspang\-yu}}, killer.

% tspanyu
\hi 
\HT{tspangyu}{\B{\cp{tspang}yu}} \I{["\t{ts}paN.ju]} \emph{n.} killer. 
\B{Tsa\-tu\-tan\-ur a fkol tspo\-lang lo\-lu ta'\-leng a\-kll\-vawm; tspang\-yur lu pum a\-teyr.} That person that was killed had brown skin, the killer had white.~\LINK{http://naviteri.org/2020/06/mi-tanlokxe-oeya-srr-afpxamo-terrible-days-in-my-country/}{b}
(\ding{43}~\HL{tspang}{\B{tspang}}) 

% tspìng
\hi
\HT{tspIng}{\B{tspìng}} \I{[\t{ts}pIN]} \emph{n., fauna} austrapede. 

% tspu'
\hi 
\HT{tspu'}{\B{tspu'}} \I{[\t{ts}puP]} \emph{or} \I{[\t{ts}pUP]} (\emph{a. vin.}) extend (\emph{occupying a certain place, stretching to a certain point}). 
\B{Atx\-kxe ay\-oe\-yä tspu' nuä ay\-ram.} Our land extends beyond the mountains.~\LINK{https://naviteri.org/2026/04/tsivola-liu-amip-thirty-two-new-words/}{b}
(\emph{b. vtr.}) hold out (\emph{often to offer}).
\B{Pol tspo\-lu' oer u\-tu\-ma\-ut\-it tsa\-krr oel mu\-nge.} She held a banana fruit out to me and I took it.~\LINK{https://naviteri.org/2026/04/tsivola-liu-amip-thirty-two-new-words/}{b}

% tsranten
\hi
\HT{tsranten}{\B{\cp{tsran}ten}} \I{["\t{ts}Ran.tEn]} \emph{vin.} matter, be of importance\slash import. 
\B{Ton\-ìri a\-la\-he, aw\-ngal yom wu\-tsot teng\-krr he\-reyn nì\-pxim, teng\-krr te\-ru\-von, ke tsran\-ten.} As for all other nights, it doesn't matter whether we eat dinner while sitting upright or leaning.~\LINK{http://whyisthisnight.com/addl-more.html\#na\%27vi}{tn} 
\B{Fì\-tseo\-ri ke tsun kaw\-tu pi\-vä\-hem tì\-yo'\-ne; tsran\-ten tì\-pä\-hem\-ä tì\-fmi nì\-'aw.} In this art it's impossible to arrive at perfection; the only thing that matters is the attempt to arrive there.~\LINK{http://naviteri.org/2012/01/mipa-zisit-ayliu-amip-new-words-for-the-new-year/}{b} 
\B{Nga oer tsran\-ten fra\-to.} You matter the most to me.~\LINK{http://forum.learnnavi.org/language-updates/confirmations-\%28comparisons-gerund\%29/}{ln} \\
\textsc{Derivatives}: \ding{43}
\on{1}.~\HL{letsranten}{\B{le\-tsran\-ten}}, important. 
\on{2}.~\HL{ketsran}{\B{ke\-tsran}}, whatever, whoever, however.

% tsray
\hi
\HT{tsray}{\B{tsray}} \I{[\t{ts}Raj]} \emph{n.} village. 
\B{Mì sray apxa kel\-ku si so\-a\-i\-a\-hu.} (He) lives with his family in a large village.~\LINK{http://naviteri.org/2010/07/ting-mikyun-fte-tslivam-listening-comprehension-1-3/}{b} \\
\textsc{Derivatives}: \ding{43}
\on{1}.~\HL{tsawtsray}{\B{tsaw\-tsray}}, small or medium\hyp{}sized city. 
\on{2}.~\HL{txantsawtsray}{\B{txan\-tsaw\-tsray}}, large city, metropolis.

% tsre'i
\hi
\HT{tsre'i}{\B{\cp{tsre}'i}} \I{["\t{ts}RE.Pi]} \emph{vtr.} throw (\B{ne}, to, towards; \HL{wA}{\B{wä}}, at). 
\B{Oel tsre\-'i tskxet ne nga.} I throw a stone towards you.~\LINK{http://forum.learnnavi.org/language-updates/to-throw-at-x/}{ln}
\B{Oel tsre\-'i tskxet nga\-ru.} I throw a stone to you.~\LINK{http://forum.learnnavi.org/language-updates/to-throw-at-x/}{ln}
\B{Oel tsre\-'i tskxet wä nga.} I throw a stone at you.~\LINK{http://forum.learnnavi.org/language-updates/to-throw-at-x/}{ln} \\
\textsc{Derivatives}: \ding{43}
\on{1}.~\HL{kxapaytsre'}{\B{kxa\-pay\-tsre'}}, spit.

% tsrul
\hi
\HT{tsrul}{\B{tsrul}} \I{[\t{ts}Rul]} \emph{or} \I{[\t{ts}RUl]} \emph{n.} nest;  \emph{protected area serving as the home of Pandoran fauna}.
\B{Fì\-fne\-ya\-yol tsrul\-it txu\-la mì song\-ropx ut\-ral\-ä fte ay\-li\-nit hi\-vaw\-nu wä sar\-ni\-o\-ang.} This kind of bird builds its nest in a tree cavity so as to protect its young from predators.~\LINK{http://naviteri.org/2013/04/wheres-the-bathroom-and-other-useful-things/}{b}
\B{Ro\-lun ay\-oel tsrul\-mì hì\-'i\-a pxe\-lo\-it a\-teyr.} We found three little white eggs in the nest.~\LINK{http://naviteri.org/2013/06/looking-back-looking-forward/}{b} \\
\textsc{Derivatives}: \ding{43}
\on{1}.~\HL{yayotsrul}{\B{ya\-yo\-tsrul}}, bird's nest.

% tstal
\hi
\HT{tstal}{\B{tstal}} \I{[\t{ts}tal]} \emph{n.} knife. 
\B{El\-tu si! Tsa\-tstal a\-fwem lu litx nì\-txan.} Pay attention! That blunt knife is very sharp.~\LINK{http://naviteri.org/2012/03/spring-vocabulary-part-1/}{b}
\B{Oe\-yä tstal\-ìl tok pe\-se\-nget? -- Lu yo\-sìn.} Where's my knife? -- It's on the table.~\LINK{http://naviteri.org/2015/08/aylifyavi-lereyfya-2-cultural-terms-2-and-more/}{b}
\B{Oel tstal\-it fyo\-lep fa aysre'.} I held the knife in my teeth.~\LINK{http://naviteri.org/2012/03/spring-vocabulary-part-1/}{b} \\
\textsc{Derivatives}: \ding{43}
\on{1}.~\HL{tstalsena}{\B{tstal\-se\-na}}, knife sheath.

% tstalsena
\hi 
\HT{tstalsena}{\B{\cp{tstal}sena}} \I{["\t{ts}tal.sE.na]} \emph{n.} knife sheath.
\B{Pol tstal\-it wrr\-zo\-lä\-rìp tstal\-se\-na\-ftu.} He pulled the knife out of its sheath.~\LINK{http://naviteri.org/2021/03/mipa-ayliu-si-aylifyavi-new-words-and-expressions/}{b} 
(\ding{43}~\HL{tstal}{\B{tstal}}, \HL{sAhena}{\B{sä\-he\-na}})

% tstew 
\hi
\HT{tstew}{\B{tstew}} \I{[\t{ts}Ew]} \emph{adj.} (\emph{of a person}) brave. 
\B{Tsam\-si\-yu a\-sìl\-tsan lu tstew slä\-kop swa\-ran.} A good warrior is courageous but also humble.~\LINK{http://naviteri.org/2020/11/vospxivopeya-ayliu-amip-novembers-new-words/}{b}
\B{Kem a\-tì\-tstew\-nga' si tu\-te a\-tstew.} A brave person does brave deeds.~\LINK{http://naviteri.org/2011/09/miscellaneous-vocabulary/}{b} 
\B{'Ur fki\-van tstew!} Do not show fear. [lit.: ``May your appearance come to the senses as brave!'']~\LINK{https://naviteri.org/2026/01/follow-up-to-txep-si-txeva/}{b\slash a\on{3}} \\
\textsc{Derivatives}: \ding{43}
\on{1}.~\HL{tItstew}{\B{tì\-tstew}}, courage, bravery. 
\on{2}.~\HL{txantstew}{\B{txan\-tstew}}, hero. 
\on{3}.~\HL{nItstew}{\B{nì\-tstew}}, bravely, courageously.

% tstir
\hi 
\HT{tstir}{\B{tstir}} \I{[\t{ts}tiR]} \emph{n.} palm of a hand. 

% tstu
\hi
\HT{tstu}{\B{tstu}} \I{[\t{ts}tu]} \textsc{i}. \emph{adj.} (\emph{a.}) closed, shut.    (\emph{b. of the weather}) completely overcast, covered with clouds. \\
\textsc{ii}. \B{tstu si} \I{[\t{ts}tu si]} \emph{vin.} close. 
\B{Oe\-ri me\-na\-ri\-ru tstu si.} I close my eyes.~\LINK{http://forum.learnnavi.org/language-updates/open-close/}{ln} \\
\textsc{iii}. \B{tstu sä\cp{pi}} \I{[\t{ts}tu s\ae."pi]} \emph{vin.} close itself\slash by itself\slash on its own.
\B{Po\-ri me\-na\-ri tstu sä\-po\-li.} His eyes closed.~\LINK{http://forum.learnnavi.org/language-updates/open-close/}{ln}

% tstunkem / si
\hi
\HT{tstunkem}{\B{\cp{tstun}kem}} \I{["\t{ts}tun.kEm]} \emph{or} \I{["\t{ts}tUn.]} \emph{coll. often pronounced} \I{["\t{ts}tuN.]}  \textsc{i}. \emph{n.} favor, act of kindness.  
\B{Tung oer fu\-ta vin tstun\-kem\-it nga\-ta.} Let me ask you a favor.~\LINK{http://naviteri.org/2020/03/mipa-sawasultsyip-a-new-contest/}{b}
(\ding{43}~\HL{tstunwi}{\B{tstu\-nwi}}, \HL{kem}{\B{kem}}) \\
\textsc{ii}. \B{\cp{tstun}kem si} \I{["\t{ts}tun.kEm si]} \emph{or} \I{["\t{ts}tUn.]} \emph{vin.} do a favor.
\B{Tstun\-kem si oer ru\-txe.} Please do me a favor.~\LINK{http://naviteri.org/2020/03/mipa-sawasultsyip-a-new-contest/}{b} \\
\textsc{Derivatives}: \ding{43}
\on{1}.~\HL{tstunkemtsyIp}{\B{tstun\-kem\-tsyìp}}, little favor.

% tstunkemtsyìp
\hi
\HT{tstunkemtsyIp}{\B{\cp{tstun}kemtsyìp}} \I{["\t{ts}tun.kEm.\t{ts}jIp\textcorner]} \emph{or} \I{["\t{ts}tUn.]} \emph{coll. often pronounced} \I{["\t{ts}tuN.]} \emph{n.} a little favor, act of kindness. 
(\ding{43}~\HL{tstunkem}{\B{tstun\-kem}}, \HL{-tsyIp}{\B{-tsyìp}})

% tstunlan
\hi 
\HT{tstunlan}{\B{\cp{tstun}lan}} \I{["\t{ts}tun.lan]} \emph{or} \I{["\t{ts}tUn.]} \emph{adj.} kind\hyp{}hearted. 
(\ding{43}~\HL{tstunwi}{\B{tstun\-wi}}, \HL{txe'lan}{\B{txe'\-lan}}; \ding{214}~\HL{kawnglan}{\B{kawng\-lan}})

% tstunwi
\hi
\HT{tstunwi}{\B{\cp{tstun}wi}} \I{["\t{ts}tun.wi]} \emph{or} \I{["\t{ts}tUn.]} \emph{adj., ofp.} kind, thoughtful, considerate.  
\B{Le\-iu ay\-nga tstun\-wi nì\-txan nì\-wotx.} You are all very kind.~\LINK{http://forum.learnnavi.org/news-announcements/happy-birthday-to-karyu-pawl/msg556899/\#msg556899}{ln} 
\B{Lu nga tstun\-wi\-a tsmuk a\-tstun\-wi.} You are a truly kind sibling.~\LINK{http://naviteri.org/2013/03/tsanerul-feerul-getting-better-getting-worse/comment-page-1/\#comment-2136}{b}
(\emph{a. idiom}) \B{Tstun\-wi.} ``It's kind of you to say that.'' (\emph{in response to a compliment}).~\LINK{http://naviteri.org/2010/07/diminutives-conversational-expressions/}{b} 
(\ding{43}~\HL{tItstunwinga'}{\B{tì\-tstun\-wi\-nga'}}) \\
\textsc{Derivatives}: \ding{43}
\on{1}.~\HL{tItstunwi}{\B{tì\-tstun\-wi}}, kindness. 
\on{2}.~\HL{tstunkem}{\B{tstun\-kem}}, favor, act of kindness.
\on{3}.~\HL{tstunlan}{\B{tstun\-lan}}, kind\hyp{}hearted.
\on{4}.~\HL{ketstun}{\B{ke\-tstun}}, unkind, emotionally cold.

% tstxa'a
\hi
\HT{tstxa'a}{\B{\cp{tstxa}'a}} \I{["\t{ts}t'a.Pa]} \emph{n.}, \ding{96} Canalyd, \emph{canalydium limacineum}.

% tstxo 
\hi
\HT{tstxo}{\B{tstxo}} \I{[\t{ts}t'o]} \emph{n.} name.
\B{Nga\-ri tstxo lu pe\-lì\-'u?} What is your name?~\LINK{http://naviteri.org/2010/09/getting-to-know-you-part-3/}{b}
\B{Tstxo oe\-yä lu X.} My name is X.~\LINK{http://forum.learnnavi.org/language-updates/the-reflexive-\%28from-a-frommer-email!!!!\%29/}{ln}
\B{Tstxo nge\-yä su\-nu oer nì\-ngay.} I truly like your name.~\LINK{http://naviteri.org/2014/08/stxeli-alor-a-beautiful-gift/comment-page-1/\#comment-9344}{b}
\B{Pol tstxo\-ti oe\-yä pol\-txe\-i\-e nì\-tsì\-syì.} He whispered my name, I'm happy to say.~\LINK{http://naviteri.org/2011/09/miscellaneous-vocabulary/}{b} \\
\textsc{Derivatives}: \ding{43}
\on{1}.~\HL{tstxolI'u}{\B{tstxo\-lì\-'u}}, noun. 
\on{2}.~\HL{tstxolI'ukIngvi}{\B{tstxo\-lì\-'u\-kìng\-vi}}, noun phrase.

% tstxolì'u
\hi
\HT{tstxolI'u}{\B{\cp{tstxo}lì'u}} \I{["\t{ts}t'o.lI.Pu]} \emph{n.} noun. 
\B{Sra\-ne, lì\-'ul alu mì fra\-krr ley\-ka\-tem lì\-'ut a\-hay tsa\-fya--fwa li\-vu tstxo\-lì\-'u, li\-vu la\-he\-a fne\-lì\-'u ke tsran\-ten.} Yes, the word \emph{in} always changes the following word like that--whether it be a noun or another kind of word doesn't matter.~\LINK{http://forum.learnnavi.org/language-updates/fmawno/msg293595/\#msg293595}{ln}
(\ding{43}~\HL{tstxo}{\B{tstxo}}, \HL{lI'u}{\B{lì\-'u}})

% tstxolì'ukìngvi
\hi 
\HT{tstxolI'ukIngvi}{\B{tstxo\cp{lì}'ukìngvi}} \I{[\t{ts}t'o."lI.Pu.kIN.vi]} \emph{n.} noun phrase (\emph{a phrase whose pivotal or central word is a noun}). 
\B{Nì\-fkey\-to\-ngay, tsa\-pxe\-lì\-'u alu zì\-sìt a\-mip le\-fpom ke lu tstxo\-lì\-'u\-kìng\-vi.} In fact those three words \emph{zì\-sìt a\-mip le\-fpom} are not a noun phrase.~\LINK{http://naviteri.org/2017/12/zisit-amip-lefpom-ma-eylan-happy-new-year-friends/\#comment-27964}{b} 
(\ding{43}~\HL{tstxo}{\B{tstxo}}, \HL{lI'ukIngvi}{\B{lì\-'u\-kìng\-vi}})

% tsu'o
\hi
\HT{tsu'o}{\B{\cp{tsu}'o}} \I{["\t{ts}u.Po]} \emph{n.} ability.   
\B{Tì\-ru\-sol\-ìri ke lu po\-ru kea tsu\-'o kaw\-'it.} As for singing, he has no ability whatsoever.~\LINK{http://naviteri.org/2012/03/spring-vocabulary-part-2/}{b} 
\B{Ze\-ne fko tsko swi\-zaw\-it si\-var nì\-trr\-trr fte\-ke fì\-tsu\-'o sni\-väm.} One must use a bow and arrow regularly to prevent this ability degrading over time.~\LINK{http://naviteri.org/2015/11/vomuna-liu-amip-ten-new-words/}{b}
(\ding{43}~\HL{-tswo}{\B{-tswo}})

% Tsu'tey
\hi
\B{Tsu'\cp{tey}} \I{[\t{ts}uP."tEj]} \emph{n.} \emph{male name}.
\B{Ma Tsu'\-tey, kem\-pe si nga?} Tsu'\-tey, what are you doing?~\LINK{http://wiki.learnnavi.org/index.php?title=Na\%27vi_from_Avatar_Movie\#Neytiri_talking_to_Tsu.27tey}{wiki\slash a}
\B{Tsu'\-tey\-ìl to\-lìng oer mawl\-it smar\-ä.} Tsu'\-tey gave me a half of the prey.~\LINK{http://naviteri.org/2011/02/new-vocabulary-part-2/}{b}
\B{Fo lu 'ey\-lan Tsu'\-tey\-ä.} They're Tsu'\-tey's friends.~\LINK{http://naviteri.org/2011/07/number-in-na\%E2\%80\%99vi/}{b}
\B{Tsu'\-tey\-ri lu fo\-ru nrr\-a a fì\-re\-won yo\-lo\-ra'.} They're proud of Tsu'\-tey for having won this morning.~\LINK{http://naviteri.org/2012/07/fivospxiya-aylifyavi-amip-this-months-new-expressions-pt-2/}{b}

% tsuk-
\hi
\HT{tsuk}{\B{tsuk-}} \I{[\t{ts}uk\textcorner.]} \emph{or} \I{[\t{ts}Uk\textcorner.]} (\textsc{rn}: \B{tsùk-} \I{[\t{ts}Uk\textcorner.]})\emph{prod. pref. on verbs to form receptive ability, i.e., s.t. is capable of receiving the action of the verb}, ``…-able'' .
\B{Fì\-i\-o\-ang lu tsuk\-yom.} This animal is edible\slash can be eaten.~\LINK{http://naviteri.org/2011/03/\%E2\%80\%9Creceptive-ability\%E2\%80\%9D-and-hesitation/}{b}
\B{Tsuk\-yom\-a ioang lu le\-sar.} An edible animal is useful.~\LINK{http://naviteri.org/2011/03/\%E2\%80\%9Creceptive-ability\%E2\%80\%9D-and-hesitation/}{b}
(\emph{a. more prod. than in English}) \B{Lu na'\-rìng tsuk\-ha\-haw.} One can sleep in the forest.\slash It's possible to sleep in the forest.\slash The forest is ``sleepable''.~\LINK{http://naviteri.org/2011/03/\%E2\%80\%9Creceptive-ability\%E2\%80\%9D-and-hesitation/}{b}
\B{Fì\-tseng lu tsuk\-tsu\-rokx.} One can rest here.\slash It's possible to rest here.\slash This place is ``restable''.~\LINK{http://naviteri.org/2011/03/\%E2\%80\%9Creceptive-ability\%E2\%80\%9D-and-hesitation/}{b}
(\ding{214}~\HL{ketsuk-}{\B{ke\-tsuk-}}) 

% tsukanom
\hi 
\HT{tsukanom}{\B{tsu\cp{ka}nom}} \I{[\t{ts}u."ka.nom]} (\textsc{rn}: \B{tsù\cp{ka}nom} \I{[\t{ts}U."ka.nom]}) \emph{adj.} available, obtainable. 
\B{Tsay\-fa\-suk tsu\-ka\-nom lu krr\-ka fì\-zì\-sì\-krr nì\-'aw.} Those berries are available during this season only.~\LINK{http://naviteri.org/2021/03/mipa-ayliu-si-aylifyavi-new-words-and-expressions/}{b} 
\B{Wä fì\-sä\-spxin a 'um\-tsa le\-iu set tsu\-ka\-nom.} Medicine against this disease is happily now available.~\LINK{http://naviteri.org/2021/03/mipa-ayliu-si-aylifyavi-new-words-and-expressions/}{b} 
(\ding{43}~\HL{tsuk}{\B{tsuk-}}, \HL{kanom}{\B{ka\-nom}}; \ding{214}~\HL{ketsukanom}{\B{ke\-tsu\-ka\-nom}}) \\
\textsc{Derivatives}: \ding{43}
\on{1}.~\HL{tItsukanom}{\B{tì\-tsu\-ka\-nom}}, availability.

% tsukmong
\hi 
\HT{tsukmong}{\B{tsuk\cp{mong}}} \I{[\t{ts}uk\textcorner."moN]} \emph{or} \I{[\t{ts}Uk\textcorner.]} (\textsc{rn}: \B{tsùk\cp{mong}} \I{[\t{ts}Uk\textcorner."moN]}) \emph{adj.} reliable, dependable. 
\B{tu\-te a\-tsuk\-mong}, dependable person.~\LINK{http://naviteri.org/2022/06/solalew-maul-zisita-half-the-year-is-over/}{b} 
\B{ay\-sì\-oeyk\-tìng a\-tsuk\-mong}, reliable explanations.~\LINK{http://naviteri.org/2022/06/solalew-maul-zisita-half-the-year-is-over/}{b} 
(\ding{43}~\HL{mong}{\B{mong}})

% tsuksìm
\hi
\HT{tsuksIm}{\B{\cp{tsuk}sìm}} \I{["\t{ts}uk\textcorner.sIm]} \emph{or} \I{[\t{ts}Uk\textcorner.]} \emph{n.} chin. 
\B{ean na tsuk\-sìm to\-ruk\-ä}, blue like the chin of a great leonopteryx.~\LINK{http://naviteri.org/2010/08/mipa-ayopin-mipa-ayliu-new-colors-new-words/}{b}

% tsukx
\hi 
\HT{tsukx}{\B{tsukx}} \I{[\t{ts}uk']} \emph{or} \I{[\t{ts}Uk']} (\textsc{rn}: \B{tsùkx} \I{[\t{ts}Uk']}) \emph{vtr.} stab. 
\B{Ney\-ti\-ril nan\-tang\-it tso\-lukx fte pe\-yä tì\-sraw\-ti 'ey\-ki\-vi\-'a.} Neytiri stabbed the viperwolf to end its pain.~\LINK{http://naviteri.org/2017/12/zisit-amip-lefpom-ma-eylan-happy-new-year-friends/}{b}

% tsulfä 
\hi
\HT{tsulfA}{\B{tsul\cp{fä}}} \I{[\t{ts}ul."f\ae]} \emph{or} \I{[\t{ts}Ul.]} (\textsc{rn}: \B{tsùl\cp{fä}} \I{[\t{ts}Ul."f\ae]}) \textsc{i}. \emph{n.} mastery.
\B{Pxe\-nge\-yä tsul\-fä lì'\-fya\-yä le\-Na'\-vi oe\-ru te\-ya si.} Your mastery of the Na'vi language fills me with joy.~\LINK{http://naviteri.org/2010/12/navi-writing-contest-the-first-place-winners/}{b}
(\emph{a. idiom}) \B{Nga\-ru tsul\-fä.} ``Thank you.'' [lit.: to you the mastery] (\emph{in response to a compliment, in a situation where the s.o. praising you is better at the activity you're being praised for than you are yourself}).~\LINK{http://naviteri.org/2010/07/diminutives-conversational-expressions/}{b} \\
\textsc{ii}. \B{tsul\cp{fä} si} \I{[\t{ts}ul."f\ae{} si]} (\textsc{rn}: \B{tsùl\cp{fä} si} \I{[\t{ts}Ul."f\ae{} si]}) \emph{vin.} master. \\
\textsc{Derivatives}: \ding{43}
\on{1}.~\HL{tsulfAtu}{\B{tsul\-fä\-tu}}, master of an art.

% tsulfätu
\hi
\HT{tsulfAtu}{\B{tsul\cp{fä}tu}} \I{[\t{ts}ul."f\ae.tu]} \emph{or} \I{[\t{ts}Ul.]} (\textsc{rn}: \B{tsùl\cp{fä}tu} \I{[\t{ts}Ul."f\ae.tu]}) \emph{n.} master of an art, craft, or skill; expert.
\B{Tew\-ti, nga lu tsul\-fä\-tu i\-'en\-ä.} Wow, you are a master on the \emph{i'en}.~\LINK{http://naviteri.org/2011/01/new-year-new-vocabulary/}{b}
\B{Fì\-me\-rel\-it 'o\-ngo\-lop aw\-nge\-yä tsul\-fä\-tul rel\-tseo\-ä alu Älìn.} The two portraits were created by our Master of Visual Arts, Alan.~\LINK{http://naviteri.org/2014/09/stxeli-alor-the-text/}{b} 
\B{Fay\-sul\-fä\-tu\-ä tì\-kang\-kem ohe\-ru meuia lu\-yu nì\-ngay.} The work of these masters truly honors me.~\LINK{http://naviteri.org/2010/06/first-post/}{b}
(\ding{43}~\HL{tsulfA}{\B{tsul\-fä}}, \HL{-tu}{\B{-tu}}) \\
\textsc{Derivatives}: \ding{43}
\on{1}.~\HL{tsulfAtunay}{\B{tsul\-fä\-tu\-nay}}, near\hyp{}master.

% tsulfätunay
\hi
\HT{tsulfAtunay}{\B{tsulfätu\cp{nay}}} \I{[\t{ts}ul.f\ae.tu."naj]} \emph{or} \I{[\t{ts}Ul.]} (\textsc{rn}: \B{tsùlfätun\cp{ay}} \I{[\t{ts}Ul.fE.tu."naj]}) \emph{n.} near\hyp{}master, \emph{s.o. who is one step away from becoming a master}.
(\ding{43}~\HL{tsulfAtu}{\B{tsul\-fä\-tu}})

% tsun 1
\hi
\HT{tsun}{\B{tsun}}\textsuperscript{\on{1}} \I{[\t{ts}un]} \emph{or} \I{[\t{ts}Un]} (\textsc{rn}: \B{tsùn} \I{[\t{ts}Un]}) \emph{vin., modal} be able, can.
\B{Tsun oe nga\-hu nì\-Na'\-vi pi\-väng\-kxo a fì\-'u oe\-ru prr\-te' lu.} It's a pleasure to be able to chat with you in Na'vi.~\LINK{http://wiki.learnnavi.org/index.php?title=Corpus\#Times_Online}{wiki}
\B{Tsat ke tsun oe ron\-srel\-ngi\-vop.} I can't imagine that.~\LINK{http://naviteri.org/2011/02/new-vocabulary-part-2/}{b}
\B{Txo tsi\-ve\-'a ay\-ngal key\-eyt, ru\-txe oe\-ru pi\-veng fte tsi\-vun oe sa\-'ut ley\-ki\-va\-tem.} If you see mistakes, please tell me so that I can change them.~\LINK{http://naviteri.org/2010/06/first-post/}{b}
\B{Lo\-lu tok\-tor Kì\-rey\-sì kar\-yu a\-sìl\-tsan ul\-te po\-e tso\-lun nì\-Na'\-vi pi\-vll\-txe.} Doctor Grace was a good teacher and she could speak Na'vi.~\LINK{http://naviteri.org/2010/07/ting-mikyun-fte-tslivam-listening-comprehension-1-3/}{b} 
(\emph{a. idiom}) \B{Tsun ti\-vam.} ``That's fine.'' [lit.: can suffice]~\LINK{http://wiki.learnnavi.org/index.php?title=Na\%27vi_from_Avatar_Movie\#Grace_and_Norm_in_the_Link_room}{wiki\slash a} 
\B{Nga tsun ni\-väk wew\-a pay\-it. Ul\-te ke ra\-you! -- Tse\-i\-un ti\-vam.} You can drink cold water. And not get drunk! -- That's just fine.~\LINK{http://naviteri.org/2012/10/audio-and-video-learning-materials-for-navi-101/}{b} 
\B{Tsun pe\-hem?} (\emph{short for} \B{Tsun fko pe\-hem si\-vi?}) ``What can one do?'' 
\B{Ri\-ni yaw\-ne lu oer, slä oe yaw\-ne ke lä\-ngu por kaw\-'it.---Tsun pe\-hem?} I'm in love with Rini, but she doesn't love me one bit.---What are you gonna do? That's life.~\LINK{http://naviteri.org/2016/06/mrrvola-lifyavi-amip-forty-new-expressions/}{b}
(\emph{b. proverb}) \B{Ke tsun fko tspi\-vang to\-ruk\-it fa fwa pewn\-ti snew.} \emph{Proverbial expression for a method that will not work.} [lit.: You can't kill a great leonopteryx by constricting its throat]~\LINK{http://naviteri.org/2012/10/snewsyea-ftxoza-halowina-a-spooky-halloween/}{b} 
\B{Txìm a\-'aw ke tsun hi\-veyn mì tal me\-fa'\-li\-yä.} You can't take both positions or sit on the fence; you need to decide. [lit.: One butt can't sit on the backs of two direhorses]~\LINK{http://naviteri.org/2013/02/vospxi-ayol-posti-apup-short-post-for-a-short-month/}{b}  \\
\textsc{Derivatives}: \ding{43}
\on{1}.~\HL{tsunslu}{\B{tsun\-slu}}, be possible. 

% tsun 2
\hi 
\HT{tsun2}{\B{tsun}}\textsuperscript{\on{2}} \I{[\t{ts}un]} \emph{or} \I{[\t{ts}Un]} \emph{n.} heel. 
\B{Oe\-ri teng\-krr te\-rul mì na'\-rìng, tsun\-it a\-ski\-en tì\-sraw sey\-ko\-li.} While I was running in the forest, I hurt my right heel.~\LINK{http://naviteri.org/2022/12/trr-anawm-polaheiem-the-great-day-has-arrived/}{b}  

% tsunslu
\hi
\HT{tsunslu}{\B{\cp{tsun}slu}} \I{["\t{ts}un.slu]} \emph{or} \I{["\t{ts}Un.]} (\textsc{rn}: \B{\cp{tsùn}slu} \I{["\t{ts}Un.slu]}) \emph{vin.} may, be possible.
\B{Kehe, ke tsun\-slu fwa yem me\-syon\-lì\-'ut uo tstxo\-lì\-'u.} No, it's not possible to put two adjectives after a noun.~\LINK{http://naviteri.org/2017/12/zisit-amip-lefpom-ma-eylan-happy-new-year-friends/\#comment-27964}{b} 
(\ding{43}~\HL{tsun}{\B{tsun}}, \HL{slu}{\B{slu}}) \\
\textsc{Derivatives}: \ding{43}
\on{1}.~\HL{tItsunslu}{\B{tì\-tsun\-slu}}, possibility.
\on{2}.~\HL{letsunslu}{\B{le\-tsun\-slu}}, possible.

% tsup
\hi 
\HT{tsup}{\B{tsup}} \I{[\t{ts}up\textcorner]} \emph{or} \I{[\t{ts}Up\textcorner]} \emph{n.} chasm, ravine. 
(\emph{a. proverb}) \B{Rä\-'ä zup nem\-fa tsup.} \emph{don't fall into the trap that s.o. has laid for you} [lit.: ``don't fall into the chasm''].~\LINK{https://naviteri.org/2025/04/mevosinga-liu-amip-twenty-new-words/}{b}

% tsupx
\hi 
\HT{tsupx}{\B{tsupx}} \I{[\t{ts}up']} \emph{or} \I{[\t{ts}Up']} (\textsc{rn}: \B{tsùpx} \I{[\t{ts}Up']}) \emph{vtr.} scratch harmfully (\emph{as with a claw}). 
\B{Pa\-lu\-kan\-ìl oe\-yä po\-ti fa tsin tso\-lä\-ngupx.} Unfortunately my cat scratched her with its claw.~\LINK{http://naviteri.org/2020/07/vospxivol-lefpom-happy-august/}{b} 
(\ding{43}~\HL{pItIk}{\B{pì\-tìk}}) 

% tsurak 
\hi
\HT{tsurak}{\B{\cp{tsu}rak}} \I{["\t{ts}u.Rak\textcorner]} \emph{n., fauna} Skimwing (\emph{temperamental Pandoran sea creature}). 

% tsurokx 
\hi
\HT{tsurokx}{\B{tsu\cp{rokx}}} \I{[\t{ts}u."Rok']} \emph{vin.} rest. 
\B{Sa'\-nur le\-ru haw\-tsyìp. Tsi\-vu\-rokx ko.} Mommy is taking a nap. Let her rest.~\LINK{http://naviteri.org/2011/07/txantsana-ultxa-mi-siatll-great-meeting-in-seattle/}{b}
\B{New oe ri\-vun a\-sim tì\-fnu\-nga'\-a tseng\-it a tsa\-ro tsun syi\-vor tsi\-vu\-rokx fte spä\-pi\-veng.} I want to find a quiet place nearby where I can chill out and rest to get my head back on straight.~\LINK{http://naviteri.org/2013/01/awvea-posti-zisita-amip-first-post-of-the-new-year/}{b}

% tswa' 
\hi
\HT{tswa'}{\B{tswa'}} \I{[\t{ts}waP]} \emph{vtr.} forget.
\B{Tok pe\-se\-nget pam\-rel\-ìl? Tsat ngal tswo\-lä\-nga' nì\-lam. -- Nì\-fkey\-to\-ngay ke tswo\-la' kaw\-'it.} Where is the text? Apparently you've forgotten it. -- As a matter of fact, I haven't.~\LINK{http://naviteri.org/2014/06/teri-tsaliu-alu-nifkeytongay-about-nifkeytongay/}{b}
\B{Fay\-lì\-'u a\-lor oe\-ru te\-ya si nì\-txan, ul\-te ke tswa\-ya' oel sat kaw\-krr.} These beautiful words fill me with joy, and I will never forget them.~\LINK{http://forum.learnnavi.org/language-updates/life-death-and-gerunds/}{ln}
(\emph{a. idiom}) \B{Teng\-krr ftxo\-zä se\-rei\-yi aw\-nga, ke tswi\-va' ay\-lom\-tu\-ti ko!} A toast: To absent friends. [lit.: While we are celebrating, let us not forget those who we wish could also be here (but can't)]~\LINK{http://naviteri.org/2011/07/txantsana-ultxa-mi-siatll-great-meeting-in-seattle/}{b}

% tswayon  
\hi
\HT{tswayon}{\B{\cp{tsway}on}} \I{["\t{ts}waj.on]} (\emph{a. vin.}) fly (\HL{fa}{\B{fa}}, with\slash by means of). 
\B{Oe tsway\-on io Ay\-vit\-ra\-yä Ra\-mu\-nong.} I fly over the Tree of Souls.~\LINK{http://forum.learnnavi.org/language-updates/over-under-and-quickly/}{ln}
\B{Ri\-ti tswo\-lay\-on ftum\-fa slär.} The stingbat flew out of the cave.~\LINK{http://naviteri.org/2013/04/wheres-the-bathroom-and-other-useful-things/}{b} 
\B{Tseyk tswa\-may\-on fa ik\-ran sre\-krr; ta\-fral fmo\-li fì\-kem si\-vi fa to\-ruk nì\-steng.} Jake flew with an ikran before; therefore he tried to do it with a toruk in a similar fashion.~\LINK{http://naviteri.org/2011/09/miscellaneous-vocabulary/}{b}
\HT{tsweykayon}{(\emph{b. \ding{246} as})} \B{swi\-zaw\-ti tswey\-kay\-on}, shoot an arrow.
\B{Swi\-zaw\-ti tswey\-kay\-on ne\-fä ne taw, tse\-nga zup ke lu law.} I shot an arrow into the air, | It fell to earth, I know not where.~\LINK{http://naviteri.org/2017/03/titusemteri-concerning-shooting/}{b} 
(\emph{c. idiom}) \B{Nge\-yä ik\-ran kaw\-krr ke tswi\-vay\-on nì\-'aw\-tu.} May your banshee never fly alone.~\LINK{http://forum.learnnavi.org/language-updates/idiomatic-expressions/}{ln}  \\
(i.) \B{txaw\-nul\-srung a tsway\-on}, airplane.~\LINK{https://naviteri.org/2024/03/fleltrra-ayliu-words-for-april-fools-day/}{b}
\on{1}.~\HL{sAtswayon}{\B{sätswayon}}, flight.

% tswal
\hi 
\HT{tswal}{\B{tswal}} \I{[\t{ts}wal]} \emph{n.} power (\emph{physical, psychological, emotional, or political}). \\
\textsc{Derivatives}: \ding{43}
\on{1}.~\HL{letswal}{\B{le\-tswal}}, powerful (\emph{ofp.}). 
\on{2}.~\HL{tswalnga'}{\B{tswal\-nga'}}, powerful (\emph{nfp.}).

% tswalnga'
\hi 
\HT{tswalnga'}{\B{\cp{tswal}nga'}} \I{["\t{ts}wal.NaP]} \emph{adj., nfp.} powerful. 
(\ding{43}~\HL{tswal}{\B{tswal}}, \HL{letswal}{\B{le\-tswal}})

% tswesya / si
\hi 
\HT{tswesya}{\B{\cp{tswe}sya}} \I{["\t{ts}wE.sja]} \textsc{i}. \emph{n.} current. \\
\textsc{ii}. \B{\cp{tswe}sya si} \I{["\t{ts}wE.sja si]} \emph{vin.} flow.
\B{Na\-ri si, ma 'itan. Kil\-van tswe\-sya si nì\-win nì\-txan.} Be careful, son. The river is flowing very swiftly.~\LINK{http://naviteri.org/2018/04/100a-liu-amip-64-new-words-part-2/}{b}

% tswin
\hi
\HT{tswin}{\B{tswin}} \I{[\t{ts}win]} \emph{n.} queue, braid (\emph{neural queue of the Na'vi to enable them to make} \B{tsa\-hey\-lu}). 
\B{Nga\-ri tswin\-tsyìp se\-vin nì\-txan lu nang!} What a pretty little queue you have!~\LINK{http://naviteri.org/2010/07/diminutives-conversational-expressions/}{b}
(\ding{43}~\HL{kuru}{\B{ku\-ru}})

% tswìk
\hi
\HT{tswIk}{\B{tswìk}} \I{[\t{ts}wIk\textcorner]} \emph{vtr.} suck. \\
(i.) \B{tswìk \cp{kxe}nerit} \I{[\t{ts}wIk\textcorner{} "k'E.nE.Rit\textcorner]} \emph{vin.} smoke (\emph{a cigarette})
\B{Ke tung Na'\-vil fu\-ta tswìk kxe\-ner\-it Mo\-'a\-ra\-ka nì\-wotx.} The Na'vi have banned smoking in Mo'ara.~\LINK{http://naviteri.org/2017/06/ayliu-a-ta-eywaeveng-kifkey-uniltirantokxa-words-from-disneys-pandora-the-world-of-avatar/}{b} \\
\textsc{Derivatives}: \ding{43} 
\on{1}.~\HL{reyptswIk}{\B{reyp\-tswìk}}, wolf tick.

% -tswo
\hi
\HT{-tswo}{\B{-tswo}} \I{[.\t{ts}wo]} \emph{prod. suf. on verbs and non\hyp{}verbal elements of} \B{si}\hyp{}\emph{constructs to indicate the ability to perform the action of the verb.} 
\B{Po\-ri wem\-tswo fra\-tsam\-si\-yur ro\-lo\-'a nì\-txan.} His ability to fight greatly impressed all the warriors.~\LINK{http://naviteri.org/2012/03/spring-vocabulary-part-2/}{b}
\B{Kxa\-ri tstu\-tswo tsran\-ten krra ke lu kea sä\-fpìl le\-sar.} When one has no useful thoughts, the ability to close one's mouth is important.~\LINK{http://naviteri.org/2012/03/spring-vocabulary-part-2/}{b}
(\ding{43}~\HL{tsu'o}{\B{tsu\-'o}})

% tsyafe
\hi
\HT{tsyafe}{\B{\cp{tsya}fe}} \I{["\t{ts}ja.fE]} \emph{adj.} mild, moderate, comfortable.
\B{Ya\-ri som\-wew\-pe (set)? Ya lu tsya\-fe.} What's the temperature? It's nice.~\LINK{http://naviteri.org/2012/10/audio-and-video-learning-materials-for-navi-101/}{b} 

% tsyal
\hi
\HT{tsyal}{\B{tsyal}} \I{[\t{ts}jal]} \emph{n.} wing.
\B{me\-syal\-hu a ik\-ran}, an ikran with two wings.~\LINK{http://wiki.learnnavi.org/index.php?title=Canon\#More_extracts_from_various_emails}{wiki}
\B{Kun\-sìp\-ìri txan\-a tì\-meyp lu tsyal a mìn.} The gunship's main weakness is the rotor system. [lit.: wing that turns]~\LINK{http://forum.learnnavi.org/language-updates/ftarpa-and-skienpa-added-to-dict/msg570362/\#msg570362}{ln\slash a\textsuperscript{c}} \\
\textsc{Derviations}: \ding{43} 
\on{1}.~\HL{lortsyal}{\B{lor\-tsyal}}, shimmyfly.
\on{2}.~\HL{txeptsyal}{\B{txep\-tsyal}}, coronis.
\on{3}.~\HL{slotsyal}{\B{slo\-tsyal}}, stormglider.

% tsyang
\hi 
\HT{tsyang}{\B{tsyang}} \I{[\t{ts}jaN]} \emph{n.} swarm (\emph{of insects and metaphorically of people, annoying and threatening}). 
\B{tsyang ay\-hì\-'ang\-ä}, a swarm of insects.~\LINK{https://naviteri.org/2024/02/mipa-ayliu-si-aylifyavi-niul-more-new-words-and-expressions/}{b} 
\B{tsyang su\-te\-yä}, a swarm of people.~\LINK{https://naviteri.org/2024/02/mipa-ayliu-si-aylifyavi-niul-more-new-words-and-expressions/}{b} 
(\ding{43}~\HL{sna}{\B{sna-}})

% tsyär
\hi
\HT{tsyAr}{\B{tsyär}} \I{[\t{ts}j\ae{}R]} \emph{vtr.} reject.
\B{Ngal pe\-lun fay\-stxe\-nut fra\-krr tsyär?} Why do you always reject these offers?~\LINK{http://naviteri.org/2012/03/spring-vocabulary-part-1/}{b}
\B{Ngal lum\-pe oe\-yä stxe\-nut tsyo\-lä\-ngär?} Why did you reject my offer?~\LINK{http://naviteri.org/2011/09/miscellaneous-vocabulary/}{b} 
(\ding{214}~\HL{mll'an}{\B{mll\-'an}}) \\
\textsc{Derivatives}: \ding{43}
\on{1}.~\HL{tItsyAr}{\B{tì\-tsyär}}, rejection.

% tsyey
\hi
\HT{tsyey}{\B{tsyey}} \I{[\t{ts}jEj]} \emph{n.} snack, light meal.
\B{Ke 'efu oe ohakx nì\-hawng; tam tsyey.} I'm not too hungry; a snack will do.~\LINK{http://naviteri.org/2015/07/aylifyavi-lereyfya-1-cultural-terms-1/}{b} \\
\textsc{Derivatives}: \ding{43}
\on{1}.~\HL{tsyeytsyIp}{\B{tsyey\-tsyìp}}, tiny bite. 
\on{2}.~\HL{niktsyey}{\B{nik\-tsyey}}, food wrap.

% tsyeytsyìp
\hi
\HT{tsyeytsyIp}{\B{\cp{tsyey}tsyìp}} \I{["\t{ts}jEj.\t{ts}jIp\textcorner]} \emph{n.} tiny bite.
(\ding{43}~\HL{tsyey}{\B{tsyey}}, \HL{-tsyIp}{\B{-tsyìp}})

% tsyeym
\hi
\HT{tsyeym}{\B{tsyeym}} \I{[\t{ts}jEjm]} \emph{n.} treasure, \emph{s.t. rare and of great value}.
\B{Tsyan\-ìri sì oe\-ri lo\-le\-iu tsa\-kaym tsyeym a\-ngay.} There was a true treasure that evening for John and I.~\LINK{http://naviteri.org/2014/08/stxeli-alor-a-beautiful-gift/}{b}

% tsyìl
\hi
\HT{tsyIl}{\B{tsyìl}} \I{[\t{ts}jIl]} \emph{vtr.} climb, scale.   
\B{Tsyìl Ik\-ni\-ma\-yat ul\-te tsa\-heyl si ik\-ran\-hu a fì\-'u lu tì\-fme\-tok a ze\-ne fra\-ta\-ron\-yu a\-'e\-wan em\-zi\-va\-'u.} Scaling Iknimaya and bonding with a banshee is a test that every young hunter must pass.~\LINK{http://naviteri.org/2012/01/more-additions-to-the-lexicon/}{b}
\B{Ke tsun oe fì\-'awkx\-it tsyi\-vìl.} I can't scale this cliff.~\LINK{http://naviteri.org/2013/01/lifyari-po-peyi-how-good-is-her-navi/}{b}
\B{Sìn yì sngä\-'i\-yu\-ä, slä tsye\-rìl (hay\-a yì\-ne) nì\-'ul\-'ul.} They're at a beginner's level, but they're getting better and better.~\LINK{http://naviteri.org/2013/01/lifyari-po-peyi-how-good-is-her-navi/}{b}

% tsyìmawnun'i
\hi 
\HT{tsyImawnun'i}{\B{tsyìmawnun\cp{'i}}} \I{[\t{ts}jI.maw.nun."Pi]} \emph{or} \I{[.nUn.]} (\textsc{rn}: \B{tsyìmawnùn\cp{'i}} \I{[\t{tS}I.maw.nUn."Pi]}) \emph{n., shortened form of} \ding{43}~\HL{tIsyImawnun'i}{\B{tì\-syì\-maw\-nun\-'i}}, independence,  freedom from a pre\hyp{}existing relationship. 
\B{Ze\-ne fra\-'e\-veng tsyì\-maw\-nun\-'it a ta sa'\-sem nì\-'i'\-a mi\-vu'\-ni.} Every child must eventually achieve independence from his or her parents.~\LINK{http://naviteri.org/2023/07/trr-tsyimawnuniya-lefpom-happy-independence-day/}{b} 
\B{Trr Tsyì\-maw\-nun\-'i\-yä Le\-fpom!} Happy Independence Day!~\LINK{http://naviteri.org/2023/07/trr-tsyimawnuniya-lefpom-happy-independence-day/}{b} 
(\ding{43}~\HL{syImawnun'i}{\B{syì\-maw\-nun\-'i}}) 

% -tsyìp
\hi
\HT{-tsyIp}{\B{-tsyìp}} \I{[.\t{ts}jIp\textcorner]} \emph{prod. suf. to n.s and pronouns to form the diminutive form}.
(\emph{a. to form new lexical items that originated as small versions of a n.}) e.g. \ding{43} \HL{puktsyIp}{\B{puk\-tsyìp}}, booklet, pamphlet;  
\HL{utraltsyIp}{\B{ut\-ral\-tsyìp}}, bush; 
\HL{sAspxintsyIp}{\B{sä\-spxin\-tsyìp}}, minor ailment.
(\emph{b. to express affection or endearment})
\B{Za\-'u fì\-tseng, ma 'ite\-tsyìp.} Come here, little daughter.~\LINK{http://naviteri.org/2010/07/diminutives-conversational-expressions/}{b}
\B{Kam\-tsyìp\-ìl wu\-tsot ye\-rom.} Little Kamun is having dinner.~\LINK{http://naviteri.org/2010/07/diminutives-conversational-expressions/}{b}
\B{Nga\-tsyìp yaw\-ne lu oer.} I love you, little one.~\LINK{http://naviteri.org/2010/07/diminutives-conversational-expressions/}{b}
\B{Nga\-ri tswin\-tsyìp se\-vin nì\-txan lu nang!} What a pretty little queue you have!~\LINK{http://naviteri.org/2010/07/diminutives-conversational-expressions/}{b}
(\emph{c. to express disparagement or insult})
\B{Nga nì\-aw\-no\-mum to oe\-tsyìp lu txur nì\-txan.} As everyone knows, you're a lot stronger than little old me.~\LINK{http://naviteri.org/2010/07/diminutives-conversational-expressions/}{b}
\B{Nga\-tsyìp\-ìl new pe\-ut ta oe?} What does little you want from me?~\LINK{http://naviteri.org/2010/07/diminutives-conversational-expressions/}{b}
\B{Fì\-ta\-ron\-yu\-tsyìp ke tsun ke\-'ut sti\-vä'\-nì.} This little hunter can't catch anything.~\LINK{http://naviteri.org/2010/07/diminutives-conversational-expressions/}{b}

% tsyo
\hi
\HT{tsyo}{\B{tsyo}} \I{[\t{ts}jo]} \emph{n.} flour. \\
\textsc{Derivatives}: \ding{43}
\on{1}.~\HL{tsyosyu}{\B{tsyo\-syu}}, food made from flour. 
\on{2}.~\HL{tsyorina'wll}{\B{tsyo\-ri\-na'\-wll}}, Cycad.

% tsyokx
\hi
\HT{tsyokx}{\B{tsyokx}} \I{[\t{ts}jok']} \emph{n.} hand. 
\B{Oe\-ri nì\-'i'a tsyokx zo\-slo\-lu.} My hand is finally healed.~\LINK{http://naviteri.org/2010/07/vocabulary-update/}{b}
\B{Oe\-ri skxir a mì syokx tì\-sraw se\-ngi.} The wound on my hand hurts.~\LINK{http://naviteri.org/2011/05/some-miscellaneous-vocabulary/}{b}
\B{Slä lu 'a\-'aw\-a tì\-ke\-teng, nat\-ke\-nong, tsyokx\-ìri ke lu zek\-wä a\-tsìng ki a\-mrr.} But there are several differences, for example, the hand doesn't have four but five digits.~\LINK{http://naviteri.org/2010/07/ting-mikyun-fte-tslivam-listening-comprehension-1-3/}{b}
\B{Txur\-tel\-mì fo fta so\-li fte\-ke ka tsyokx fwi\-vi.} They tied a knot in the rope so it wouldn't slip through their hands.~\LINK{http://naviteri.org/2013/04/wheres-the-bathroom-and-other-useful-things/}{b} \\
\textsc{Derivatives}: \ding{43}
\on{1}.~\HL{hawntsyokx}{\B{hawn\-tsyokx}}, glove. 
\on{2}.~\HL{taksyokx}{\B{tak\-syokx}}, clap hands.
\on{3}.~\HL{tsyokxtil}{\B{tsyokx\-til}}, wrist.
\on{4}.~\HL{rumtsyokx}{\B{rum\-tsyokx}}, fist.
\on{5}.~\HL{rIktsyokx}{\B{rìk\-tsyokx}}, flat hand.

% tsyokxtil
\hi 
\HT{tsyokxtil}{\B{\cp{tsyokx}til}} \I{["\t{ts}jok'.til]} \emph{usually} \I{["\t{ts}jok\textcorner.til]} \emph{n.} wrist. 
(\ding{43}~\HL{tsyokx}{\B{tsyokx}}, \HL{til}{\B{til}})

% tsyorina'wll
\hi
\HT{tsyorina'wll}{\B{tsyori\cp{na'}wll}} \I{[\t{ts}jo.ri."naP.w\s{l}]} \emph{n.}, \ding{96} Cycad, ``flour seed plant'', \emph{pseudocycas altissima}. 
(\ding{43}~\HL{tsyo}{\B{tsyo}}, \HL{rina'}{\B{ri\-na'}}, \HL{'ewll}{\B{'e\-wll}})

% tsyosyu
\hi
\HT{tsyosyu}{\B{\cp{tsyo}syu}} \I{["\t{ts}jo.sju]} \emph{n.} food made from flour.
(\ding{43}~\HL{tsyo}{\B{tsyo}}, \HL{syuve}{\B{syu\-ve}})

% tsyul
\hi 
\HT{tsyul}{\B{tsyul}} \I{[\t{ts}jul]} \emph{or} \I{[\t{ts}jUl]} \emph{vtr.} begin, start. 
\B{Pol tì\-kang\-kem\-it tsyo\-lul.} He began the work.
\B{Tì\-kang\-kem tsyä\-po\-lul.} The work began.~\LINK{http://naviteri.org/2018/08/fmawnti-stolawm-srak-have-you-heard-the-news/}{b}
(\ding{43}~\HL{sngA'i}{\B{sngä'i}}) \\
\textsc{Derivatives}: \ding{43}
\on{1}.~\HL{tItsyul}{\B{tì\-tsyul}}, beginning, start.

% Tsyungwen
\hi
\HT{tsyungwen}{\B{\cp{Tsyung}wen}} \I{["\t{ts}juN.wEn]} \emph{or} \I{[."\t{ts}jUN.]} (\textsc{rn}: \B{\cp{Tsyùng}wen} \I{["\t{tS}UN.wEn]}) \emph{n.} \textsc{lw} Mandarin Chinese.

\end{multicols}


% ===== Section U =====

 \pdfbookmark[0]{U}{U}
\begin{center}
\section*{U \I{[u\slash U]}}
\end{center}

\begin{multicols}{2}

% ue'
\hi 
\HT{ue'}{\B{u\cp{e'}}} \I{[u."{}EP]} (\emph{a. vin.}) vomit.
\B{Oey nan\-tang\-tsyìp o\-lu\-e' ta\-lu\-na yom nì\-hawng.} My dog vomited because it ate too much.~\LINK{http://naviteri.org/2016/06/mrrvola-lifyavi-amip-forty-new-expressions/}{b}
(\emph{b. vtr.}) vomit up. 
\B{Prr\-nen\-ìl wu\-tsot o\-lu\-e'.} The baby vomited up its meal.~\LINK{http://naviteri.org/2016/06/mrrvola-lifyavi-amip-forty-new-expressions/}{b} 

% uk
\hi
\HT{uk}{\B{uk}} \I{[uk\textcorner]} \emph{or} \I{[Uk\textcorner]} (\textsc{rn}: \B{ùk} \I{[Uk\textcorner]}) \emph{n.} shadow.
\B{vawm na uk}, dark\hyp{}shadow color.~\LINK{http://naviteri.org/2010/08/mipa-ayopin-mipa-ayliu-new-colors-new-words/}{b} \\
\textsc{Derivations}: \ding{43}
\on{1}.~\HL{toruk}{\B{to\-ruk}}, leonopteryx, the Last Shadow.
\on{2}.~\HL{ukyom}{\B{uk\-yom}}, eclipse.

% ukyom / si
\hi 
\HT{ukyom}{\B{\cp{uk}yom}} \I{["{}uk\textcorner.jom]} \emph{or} \I{["{}Uk\textcorner.]} (\textsc{rn}: \B{\cp{ùk}yom} \I{["{}Uk\textcorner.jom]}) \textsc{i}. \emph{n.} eclipse. \\ 
\textsc{ii}. \B{\cp{uk}yom si} \I{["{}uk\textcorner.jom si]} \emph{vin.} to eclipse s.t. (\emph{also metaphorically})
\B{Po\-ri tsa\-kxey\-ey a\-'aw fra\-kem\-ur a\-muiä uk\-yom so\-lä\-ngi.} Sadly, that one mistake eclipsed all of his good deeds.~\LINK{http://naviteri.org/2023/12/mipa-zisit-ayliu-amip-new-year-new-words/}{b} 
(\ding{43}~\HL{uk}{\B{uk}}, \HL{yom}{\B{yom}}) 

% ukxo
\hi
\HT{ukxo}{\B{u\cp{kxo}}} \I{[u."k'o]} \emph{adj.} dry. 
\B{Ta\-kuk set tsaw\-ke. Sley\-ku fra\-'ut u\-kxo.} Now the sun bears down. (It) makes everything dry.~\LINK{http://forum.learnnavi.org/language-updates/learnings-from-the-siatlla-song-translations/msg475008/\#msg475008}{ln}
\B{Rìk a\-ukxo mì hu\-fwe\-tsyìp ka\-fi sar\-mi.} The dry leaf was sailing in the breeze.~\LINK{http://naviteri.org/2023/12/mipa-zisit-ayliu-amip-new-year-new-words/}{b}
(\ding{214}~\HL{mei}{\B{mei}})

% ulte
\hi
\HT{ulte}{\B{\cp{ul}te}} \I{["ul.tE]} \emph{or} \I{["Ul.]} \emph{conj.} and (\emph{to connect sentences}) 
\B{Oel ay\-nga\-ti ka\-me\-ie, ma oe\-yä ey\-lan, ul\-te ay\-nga\-ru sei\-yi ira\-yo.} I \textsc{see} you, my friends, and I thank you.~\LINK{http://forum.learnnavi.org/news-announcements/a-response-from-paul-frommer!/}{ln}  
\B{Kì\-ye\-va\-me ul\-te Ey\-wa nga\-hu.} \textsc{See} you again, and may Ey\-wa be with you.~\LINK{http://forum.learnnavi.org/news-announcements/a-response-from-paul-frommer!/}{ln} 
(\ding{43}~\HL{sI}{\B{sì}})

% ultxa / si
\hi
\HT{ultxa}{\B{ul\cp{txa}}} \I{[ul."t'a]} \emph{or} \I{[Ul.]} (\textsc{rn}: \B{ùl\cp{txa}} \I{[Ul."da]}) \textsc{i}. \emph{n.} meeting.
\B{ul\-txa a mì na'\-rìng}, the meeting in the forest.~\LINK{http://naviteri.org/2010/12/catching-up/}{b}
\B{Am\-'a\-lu\-ke fì\-ul\-txa 'o' la\-yu nì\-txan.} Without a doubt this meeting will be very exciting.~\LINK{http://naviteri.org/2012/06/spring-vocabulary-part-3/}{b}
\B{Pol mo\-lä\-kxu ul\-txa\-ti.} He interrupted the meeting.~\LINK{http://naviteri.org/2010/09/getting-to-know-you-part-1/}{b}
\B{Sìl\-pey oe, tsun oeng ul\-txa si\-vi ro hay\-a ul\-txa!} I hope we can meet at the next meeting.~\LINK{http://naviteri.org/2011/07/txantsana-ultxa-mi-siatll-great-meeting-in-seattle/comment-page-1/\#comment-844}{b} \\
\textsc{ii}. \B{ultxa si} \I{[ul."t'a si]} (\textsc{rn}: \B{ùl\cp{txa} si} \I{[Ul."da{} si]}) \emph{vin.} meet (intentionally with s.o., \B{hu}). 
\B{Tsa\-ria nga\-hu ye'\-rìn ul\-txa si nì\-mun, oe sre\-fe\-re\-i\-ey nì\-prr\-te'.} I'm looking forward to getting together with you again soon.~\LINK{http://naviteri.org/2011/02/new-vocabulary-part-2/}{b}
\B{Ul\-txa si\-vi oeng sìn ram\-tsyìp txon\-'ong\-ay.} Let's meet on the hill tomorrow at nightfall.~\LINK{http://naviteri.org/2012/01/mipa-zisit-ayliu-amip-new-words-for-the-new-year/}{b}
\B{Vol\-a ta\-ron\-yu mì na'\-rìng ul\-txa so\-li.} Eight hunters met (together) in the forest.~\LINK{http://naviteri.org/2012/03/spring-vocabulary-part-2/comment-page-1/\#comment-1456}{b} \\
\textsc{Derivations}: \ding{43}
\on{1}.~\HL{ultxatu}{\B{ul\-txa\-tu}}, meeting participant. 
\on{2}.~\HL{ultxarun}{\B{ul\-txa\-run}}, meet (by chance). 
\on{3}.~\HL{tsawlultxa}{\B{tsawl\-ul\-txa}}, great gathering, conference. 
\on{4}.~\HL{numultxa}{\B{nu\-mul\-txa}}, class (for instruction).

% ultxarun
\hi
\HT{ultxarun}{\B{ul\cp{txa}run}} \I{[ul."t'a.R$\cdot$un]} \emph{or} \I{[Ul."t'a.R$\cdot$Un]} (\textsc{rn}: \B{ùl\cp{txa}run} \I{[Ul."da.R$\cdot$un]}) \emph{vtr., inf.}\on{33} meet by chance, encounter. 
\B{Ul\-txa\-ro\-lun pol ye\-rik\-it.} He met a Hexapede.~\LINK{http://forum.learnnavi.org/intermediate/translation-exercise-a-navajo-coyote-tale-in-navi/}{ln} 
(\ding{43}~\HL{ultxa}{\B{ul\-txa}}, \HL{run}{\B{run}})

% ultxatu
\hi
\HT{ultxatu}{\B{ul\cp{txa}tu}} \I{[ul."t'a.tu]} \emph{or} \I{[Ul.]} (\textsc{rn}: \B{ùl\cp{txa}tu} \I{[Ul."da.tu]}) \emph{n.} meeting participant. 
\B{Kaym\-o zo\-la\-'u fray\-ul\-txa\-tu ne kel\-ku moe\-yä fte yi\-vom wu\-tsot, fti\-vi\-a nì\-'it lì'\-fya\-ti le\-Na'\-vi, ul\-te ki\-vä\-teng nì\-'o'.} All the participants came to our house one evening to have dinner, study a little Na'\-vi, hang out together, and have fun.~\LINK{http://naviteri.org/2014/09/stxeli-alor-the-text/}{b}
(\ding{43}~\HL{ultxa}{\B{ul\-txa}}, \HL{-tu}{\B{-tu}})

% um
\hi
\HT{um}{\B{um}} \I{[um]} \emph{or} \I{[Um]} (\textsc{rn}: \B{ùm} \I{[Um]}) \emph{adj.} loose.   
(\ding{214}~\HL{'ekxin}{\B{'ekxin}}) \\
\textsc{Derivations}: \ding{43}
\on{1}.~\HL{'ekxinum}{\B{'e\-kxi\-num}}, tightness, looseness.

% um'a
\hi 
\HT{um'a}{\B{\cp{um}'a}} \I{["{}um.Pa]} \emph{or} \I{["{}Um.Pa]} \emph{adv.} surprisingly (\emph{for unexpectedly negative outcomes}). 
\B{Pol tì\-kang\-kem\-it tsyo\-lul nì\-so'\-ha, slä tsa\-'ur ha\-sey ke so\-li, um\-'a.} He began the work enthusiastically, but surprisingly, he didn't finish it.~\LINK{http://naviteri.org/2022/11/krr-ao-an-exciting-time/}{b}  
(\ding{214}~\HL{ti'a}{\B{ti\-'a}}) 

% unil
\hi
\HT{unil}{\B{\cp{u}nil}} \I{["u.nil]} \textsc{i}. \emph{n.} dream.
\B{Kap sì ay\-u\-nil say\-la\-he.} Threats, dreams, and other things.~\LINK{http://naviteri.org/2015/03/kap-si-ayunil-saylahe-threats-dreams-and-other-things/}{b} \\  
\textsc{ii}. \B{unil si} \I{["u.nil si]} \emph{vin.} dream.
\B{Tì\-txen si, ma 'ite! Unil sar\-mi nga teng\-krr ze\-rawng. Lu fpom srak?} Wake up, daughter! You were dreaming and screaming. Are you okay?~\LINK{http://naviteri.org/2015/03/kap-si-ayunil-saylahe-threats-dreams-and-other-things/}{b} 
(\ding{43}~\HL{uniltsa}{\B{unil\-tsa}}) \\
\textsc{Derivations}: \ding{43}
\on{1}.~\HL{uniltaron}{\B{uniltaron}}, dreamhunt. 
\on{2}.~\HL{uniltIrantokx}{\B{u\-nil\-tì\-ran\-tokx}}, dreamwalker body, avatar; (the movie) \emph{Avatar}. 
\on{3}.~\HL{uniltIranyu}{\B{u\-nil\-tì\-ran\-yu}}, dreamwalker.
\on{4}.~\HL{uniltsa}{\B{u\-nil\-tsa}}, dream of\slash about. 
\on{5}.~\HL{txopunil}{\B{txo\-pu\-nil}}, nightmare.

% uniltaron
\hi
\HT{uniltaron}{\B{unil\cp{ta}ron}} \I{[""{}u.nil."ta.Ron]} \emph{n.} Dreamhunt. 
\B{Ätxä\-le si tsnì li\-vu ohe\-ru Unil\-ta\-ron.} \emph{or}
\B{Vu\-yin ohel Unil\-ta\-ron\-it.} I respectfully request the Dreamhunt.~\LINK{http://naviteri.org/2011/08/reported-speech-reported-questions/}{b}
(\ding{43}~\HL{unil}{\B{unil}}, \HL{taron}{\B{ta\-ron}})

% uniltìrantokx
\hi
\HT{uniltIrantokx}{\B{uniltì\cp{ran}tokx}} \I{[""{}u.nil.tI."Ran.tok']} \emph{n.} (\emph{a.}) dream\-walk\-er body, avatar.  
\B{Me\-unil\-tì\-ran\-tokx Tok\-tor Kì\-rey\-sä sì Tsyeyk\-ä ke lu teng ki steng.} The avatars of Grace and Jake aren't the same, but they are similar.~\LINK{http://naviteri.org/2011/09/miscellaneous-vocabulary/}{b}
\B{Poe\-ri unil\-tì\-ran\-tokx\-it tar\-mok a krr lam stum nì\-ay\-fo.} When she is in her Avatar body, she appears almost like them.~\LINK{http://naviteri.org/2010/07/ting-mikyun-fte-tslivam-listening-comprehension-1-3/}{b}
(\emph{b. the movie} Avatar (\on{2009})) \B{Oe\-ri Unil\-tì\-ran\-tokx\-ìl tì\-rey\-ti ley\-ko\-la\-te\-i\-em nì\-txan.} \emph{Avatar} has greatly changed my life for the better.~\LINK{http://forum.learnnavi.org/language-updates/pseudo-canon-from-bd-live-content/msg350756/\#msg350756}{ln}
\B{Ze\-ne oe 'aw\-si\-teng tì\-kang\-kem si\-vi fo\-hu a Unil\-tì\-ran\-tokx\-it sì ki\-fkey\-ti Ey\-wa\-'e\-veng\-ä za\-mo\-lu\-nge aw\-ngar.} I must work together with those who have brought us ``Avatar'' and the world of Pandora.~\LINK{http://forum.learnnavi.org/news-announcements/a-response-from-paul-frommer!/}{ln} 
\B{Unil\-tì\-ran\-tokx\-ä pam\-tseo\-ngop\-yu a\-yaw\-ne to\-ler\-kä\-ngup.}The dear music creator of \emph{Avatar} has died.~\LINK{http://naviteri.org/2015/06/na-ayskxe-mi-telan-sad-news/}{b}
(\ding{43}~\HL{unil}{\B{unil}}, \HL{tIranyu}{\B{tì\-ran\-yu}}, \HL{tokx}{\B{tokx}}) \\
\textsc{Derivations}: \ding{43}
\on{1}.~\HL{uniltIrantokxolo'}{\B{u\-nil\-tì\-ran\-tokx\-olo'}}, \emph{Avatar} community.

% uniltìrantokxolo'
\hi
\HT{uniltIrantokxolo'}{\B{uniltì\cp{ran}tokxolo'}} \I{[""u.nil.tI."Ran.tok'.o.loP]} \emph{n.} \emph{Avatar} community.
\B{Krr\-ka tsawl\-ul\-txa Unil\-tì\-ran\-tokx\-olo'\-ä vo\-spxì\-am.} During the Avatar Community Meet\hyp{}up last month.~\LINK{http://naviteri.org/2014/08/stxeli-alor-a-beautiful-gift/}{b}
(\ding{43}~\HL{uniltIrantokx}{\B{unil\-tì\-ran\-tokx}}, \HL{olo'}{\B{o\-lo'}})

% uniltìranyu
\hi
\HT{uniltIranyu}{\B{uniltì\cp{ran}yu}} \I{[""{}u.nil.tI."Ran.ju]} \emph{n.} dreamwalker. 
\B{Pol\-txe oe san ze\-ne kea unil\-tì\-ran\-yu ke zi\-va\-'u fì\-tseng.} I said no dreamwalker is to come here.~\LINK{http://wiki.learnnavi.org/index.php?title=Na\%27vi_from_Avatar_Movie\#Scene_in_Hometree}{wiki\slash a}
(\ding{43}~\HL{unil}{\B{unil}}, \HL{tIranyu}{\B{tì\-ran\-yu}})

% uniltsa
\hi
\HT{uniltsa}{\B{\cp{u}niltsa}} \I{["u.nil.\t{ts}$\cdot$a]} \emph{vtr., inf.}\on{33} dream of\slash about, dream (that).
\B{Nì\-trr\-trr oel unil\-tsa sa'\-nu\-ä tey\-lut.} I regularly dream of my mom's \emph{teylu}.~\LINK{http://naviteri.org/2015/03/kap-si-ayunil-saylahe-threats-dreams-and-other-things/}{b}
\B{Unil\-tso\-la oel txon\-am fu\-ta tsway\-on Ney\-ti\-ri\-hu.} Last night I dreamed I was flying with Neytiri.~\LINK{http://naviteri.org/2015/03/kap-si-ayunil-saylahe-threats-dreams-and-other-things/}{b}
(\ding{43}~\HL{unil}{\B{unil}}, \HL{tse'a}{\B{tse\-'a}})

% unyor
\hi
\HT{unyor}{\B{un\cp{yor}}} \I{[un."joR]} \emph{adj.} sweetly aromatic, \emph{a flowery or aromatic woody sort of smell}. 

% uo
\hi
\HT{uo}{\B{\cp{u}o}} \I{["u.o]} \emph{adp.} behind. 
\B{Fnu, fnu, wä\-pan wä ku | Uo tsa\-ìpxa.} Quiet, quiet, hide yourself from harm | Behind that fern.~\LINK{http://forum.learnnavi.org/language-updates/learnings-from-the-siatlla-song-translations/msg473566/\#msg473566}{ln}
\B{Ut\-ral\-ti a tsa\-uo Lo\-ak wä\-par\-man pol vol fa kxe\-tse.} She used her tail to point out the tree Loak was hiding behind.~\LINK{http://naviteri.org/2012/06/spring-vocabulary-part-3/}{b}
(\ding{214}~\HL{eo}{\B{eo}}) \\
\textsc{Derivations}: \ding{43}
\on{1}.~\HL{uolI'uvi}{\B{uo\-lì\-'u\-vi}}, suffix. 
\on{2}.~\HL{ftuopa}{\B{ftuo\-pa}}, from behind.

% uolì'uvi
\hi 
\HT{uolI'uvi}{\B{\cp{u}olì'uvi}} \I{["{}u.o.lI.Pu.vi]} \emph{n.} suffix. 
\B{Lu tsa\-lì\-'ur alu fì\-haw\-re'\-ti me\-lì\-'u\-vi alu 'aw\-a eo\-lì\-'u\-vi sì 'aw\-a uo\-lì\-'u\-vi.} The word \emph{fì\-haw\-re'\-ti} has two affixes--one prefix and one suffix.~\LINK{http://naviteri.org/2017/10/more-language-for-talking-about-language/}{b}
(\ding{43}~\HL{uo}{\B{uo}}, \HL{lI'uvi}{\B{lì\-'u\-vi}})

% -ur
\hi
\HT{-ur}{\B{-ur}} \I{[.uR]} \emph{or} \I{[.UR]} \emph{suf. to mark the dative case of a n. ending in a consonant, the pseudo\hyp{}vowels ll and rr, or the diphthongs ey, ay, and aw}.
(\ding{43}~\HL{-r}{\B{-r}}, \HL{-ru}{\B{-ru}})

% uran
\hi
\HT{uran}{\B{\cp{u}ran}} \I{["u.Ran]} \emph{n.} boat. 
\B{Ay\-fo so\-lop ìlä hil\-van fa u\-ran.} They traveled along (up, down) the river by boat.~\LINK{http://naviteri.org/2011/07/txantsana-ultxa-mi-siatll-great-meeting-in-seattle/}{b}
\B{Txo fkol ke fyi\-vel u\-ran\-it pay\-wä, ze\-ne fko sli\-ve\-le.} If one does not seal a boat against water, one must swim.~\LINK{http://naviteri.org/2012/06/spring-vocabulary-part-3/}{b} \\
\textsc{Derivations}: \ding{43} 
\on{1}.~\HL{kafiuran}{\B{ka\-fi\-uran}}, sailboat.

% ‹us›
\hi
‹\HT{us}{\B{us}}› \I{[u.s]} \emph{inf. in} ‹1› (to form the active participle) 
(\emph{a. always used attributively}) \B{Pa\-lu\-lu\-kan a\-tu\-sa\-ron lu le\-hrr\-ap.} A hunting thanator is dangerous.~\LINK{http://forum.learnnavi.org/language-updates/participles/}{ln}
\B{pay a\-zu\-sup}, falling water.~\LINK{http://naviteri.org/2012/07/meetings-waterfalls-and-more/}{b}
(\emph{b. to productively form the gerund with} \HL{tI-}{\B{tì-}}) 
\B{Tì\-ru\-sol lu oe\-ru mo\-wan.} Singing is enjoyable to me.~\LINK{http://naviteri.org/2012/02/trr-asawnung-lefpom-happy-leap-day/}{b}
\B{Ay\-nge\-yä tì\-ftu\-si\-a 'o' li\-vu nì\-'aw!} May your studying be exciting only.~\LINK{http://naviteri.org/2012/10/sound-files-added-to-previous-post/}{b}
\B{Ran tì\-ru\-sol\-ä pe\-yä lu fyo\-le.} The \emph{ran} of her singing is sublime.~\LINK{http://naviteri.org/2012/10/mipa-vospxi-mipa-ayliu-new-words-for-the-new-month/}{b}

% utral 
\hi
\HT{utral}{\B{\cp{ut}ral}} \I{["ut\textcorner.Ral]} \emph{or} \I{["Ut\textcorner.]} \emph{n.} \ding{96} tree. 
\B{Ut\-ral a lu few pay\-fya a eo kel\-ku oe\-yä tsawl lu nì\-txan.} The tree on the other side of the stream in front of my house is very tall.~\LINK{http://naviteri.org/2011/01/new-year-new-vocabulary/}{b} 
\B{Nga pam\-rel so\-li san sley\-ki\-vu ut\-ral\-ìl lì'\-fya\-yä le\-Na'\-vi pxay\-a ay\-vur\-it a\-txur sìk.} You wrote, ``May the Na'vi language tree produce many strong stories.''~\LINK{http://forum.learnnavi.org/language-updates/life-death-and-gerunds/}{ln}
\B{Ngä\-zìk lu fwa var tsko\-ti fyi\-vep teng\-krr ut\-ral\-it tsye\-rìl.} It's hard to keep holding a bow while climbing a tree.~\LINK{http://naviteri.org/2012/03/spring-vocabulary-part-1/}{b}
\B{Ay\-ut\-ral\-ur lu frr\-ne\-syul}, the trees have buds.~\LINK{http://forum.learnnavi.org/news-announcements/christmas-song-ninavi-on-youtube/}{ln}
\B{Ut\-ral\-ä A\-nawm | Ay\-ri\-na' lu ay\-oeng.} We are all seeds | Of the Great Tree.~\LINK{http://www.pandorapedia.com/navi/music/ritual_music}{pp} 
\B{Tsa\-txon ay\-ut\-ral\-äo krr a pol oe\-ti no\-lìn nì\-wa\-we, olo\-me\-i\-um oel fa key\-rel fu\-ta lu yaw\-ne oe po\-ru.} That night beneath the trees when she looked at me meaningfully, I knew by the expression on her face that she loves me.~\LINK{http://naviteri.org/2011/05/some-miscellaneous-vocabulary/}{b} \\
(i.) \B{\cp{Ut}ral Ay\cp{mok}riyä} \I{["ut\textcorner.Ral aj."mok\textcorner.Ri.j\ae]} \emph{n.} Tree of Voices.
\B{Ke lu kaw\-tu a nul\-ni\-vew oe po\-hu ti\-rea\-pi\-väng\-kxo äo Ut\-ral Ay\-mok\-ri\-yä.} There's nobody I'd rather commune with under the Tree of Voices.~\LINK{http://wiki.learnnavi.org/Corpus\#Lemondrop.com}{ld\slash wiki} \\  
\textsc{Derivations}: \ding{43}
\on{1}.~\HL{utraltsyIp}{\B{ut\-ral\-tsyìp}}, bush. 
\on{2}.~\HL{kelutral}{\B{kel\-ut\-ral}}, Hometree. 
\on{3}.~\HL{tautral}{\B{tau\-tral}}, sky tree.  
\on{4}.~\HL{rumut}{\B{ru\-mut}}, puffball tree. 
\on{5}.~\HL{pxiut}{\B{pxi\-ut}}, razor palm. 
\on{6}.~\HL{vitrautral}{\B{vit\-ra\-ut\-ral}}, Tree of Souls. 
\on{7}.~\HL{utral aymokriyA}{\B{Ut\-ral Ay\-mok\-ri\-yä}}, Tree of Voices.
\on{8}.~\HL{vArumut}{\B{vä\-ru\-mut}}, vein pod. 
\on{9}.~\HL{vozampasukut}{\B{vo\-zam\-pa\-su\-kut}}, grinch tree.
\on{10}.~\HL{fyIpmaut}{\B{fyìp\-ma\-ut}}, squid fruit tree.
\on{11}.~\HL{koaktutral}{\B{ko\-ak\-tut\-ral}}, goblin thistle.
\on{12}.~\HL{lanutral}{\B{lanutral}}, dandetiger.
\on{13}.~\HL{rumaut}{\B{ru\-ma\-ut}}, cannonball fruit tree.

% Utral Aymokriyä
\hi
\HT{utral aymokriyA}{\B{\cp{Ut}ral Ay\cp{mok}riyä}} \emph{or} \B{\cp{Ut}raya \cp{Mok}ri} \I{["ut\textcorner.Ral aj."mok\textcorner.Ri.j\ae]} \emph{or} \I{["Ut\textcorner.]} \emph{n.} \ding{96} Tree of Voices (\emph{important spiritual site}).
\B{Ke lu kaw\-tu a nul\-ni\-vew oe po\-hu ti\-re\-a\-pi\-väng\-kxo äo Ut\-ral Ay\-mok\-ri\-yä.} There's nobody I'd rather commune with under the Tree of Voices.~\LINK{http://wiki.learnnavi.org/index.php?title=Corpus\#Lemondrop.com}{wiki}
(\ding{43}~\HL{utral}{\B{ut\-ral}}, \HL{mokri}{\B{mok\-ri}})

% utraltsyìp 
\hi
\HT{utraltsyIp}{\B{\cp{ut}raltsyìp}} \I{["ut\textcorner.Ral.\t{ts}jIp\textcorner]} \emph{or} \I{["Ut\textcorner.]} \emph{n.} \ding{96} bush. 
(\ding{43}~\HL{utral}{\B{ut\-ral}}, \HL{-tsyIp}{\B{-tsyìp}})

% Utraya Mokri
\hi
\HT{utraya mokri}{\B{\cp{Ut}raya \cp{Mok}ri}} \I{["ut\textcorner.Ra.ja "mok\textcorner.Ri]} : \HL{utral aymokriyA}{\B{Utral Aymokriyä}}

% utu
\hi
\HT{utu}{\B{\cp{u}tu}} \I{["u.tu]} \emph{n.} forest canopy.
(\emph{a. proverb}) \B{Kxìm utu\-ftu fna\-we'\-tu.} ``A coward commands from the canopy,'' i.e. \emph{a real leader will have `boots on the ground' and will help out, whereas a coward will only tell people (from afar) what to do}.~\LINK{http://naviteri.org/2021/08/ulte-ayyoratu-leiu-and-the-winners-are-2/}{b} \\
\textsc{Derivations}: \ding{43}
\on{1}.~\HL{utumauti}{\B{utu\-mau\-ti}}, banana fruit.

% utumauti
\hi
\HT{utumauti}{\B{\cp{u}tumauti}} \I{["u.tu.ma.u.ti]} \emph{n.} banana fruit.
(\ding{43}~\HL{utu}{\B{utu}}, \HL{mauti}{\B{ma\-u\-ti}})

% uturu
\hi 
\HT{uturu}{\B{u\cp{tu}ru}} \I{[u."tu.Ru]} (\textsc{rn}: \B{u\cp{tù}ru} \I{[u."tU.Ru]}) \emph{n.} sanctuary, place of refuge. 
\B{Vu\-yin o\-hel u\-tu\-rut.} I respectfully request sanctuary.~\LINK{http://naviteri.org/2022/12/trr-anawm-polaheiem-the-great-day-has-arrived/}{b} 
\B{Nga ke tsun wä\-pi\-van; nga\-ri ke lu ke\-a u\-tu\-ru kaw\-tseng.} You cannot hide; there is no sanctuary for you anywhere.~\LINK{http://naviteri.org/2022/12/trr-anawm-polaheiem-the-great-day-has-arrived/}{b} 

% uvan / si
\hi
\HT{uvan}{\B{u\cp{van}}} \I{[u."van]} \textsc{i}. \emph{n.} game.
\B{U\-van Vo\-mun\-a Pam\-rel\-vi\-yä.} The Ten Letter Game.~\LINK{http://www.pbs.org/wgbh/nova/blogs/secretlife/blogposts/learn-navi-with-avatars-paul-frommer/}{nova}
\B{Fì\-u\-van\-ìl oe\-ti syey\-kor nì\-txan.} I find this game very relaxing.~\LINK{http://naviteri.org/2013/01/awvea-posti-zisita-amip-first-post-of-the-new-year/}{b}
\B{Yì\-mo\-ra' Tsu'\-tey\-ìl uvan\-it.} Tsu'tey just won the game.~\LINK{http://naviteri.org/2011/05/some-miscellaneous-vocabulary/}{b}
\B{Trr\-am\-ä ay\-u\-van\-ìri mak\-to Ak\-wey nì\-yew\-la ha snaytx.} Akwey rode disappointingly in yesterday's games so he lost.~\LINK{http://naviteri.org/2012/10/mipa-vospxi-mipa-ayliu-new-words-for-the-new-month/}{b}
\B{Fra\-u\-van\-ìri lu yo\-ra'\-tu, lu snay\-tu.} For every game, there's a winner and a loser.~\LINK{http://naviteri.org/2011/12/a-note-on-the-word-\%E2\%80\%9Cyora\%E2\%80\%99tu\%E2\%80\%9D/}{b}
\B{Saw\-no'\-ha ioit ko\-la\-ne\-i\-om oel u\-van\-fa!} I got (won) the prized piece of jewelry in the game!~\LINK{http://naviteri.org/2012/06/spring-vocabulary-part-3/}{b} \\
(i.) \B{u\cp{van} le\cp{tokx}} physical game\slash activity, sport.
\B{Ay\-u\-van le\-tokx 'o' lu nì\-txan.} Sports are great fun.~\LINK{http://naviteri.org/2010/09/getting-to-know-you-part-3/}{b} \\
\textsc{ii}. \B{u\cp{van} si} \I{[u."van si]} \emph{vin.} play (a game).
\B{Nì\-'aw\-ve oeng yo\-lom wu\-tsot; maw\-krr uvan si.} First we had dinner; afterwards we played.~\LINK{http://naviteri.org/2011/08/new-vocabulary-clothing/comment-page-1/\#comment-917}{b}
\B{Teng\-krr hu pa\-lu\-kan\-tsyìp uvan se\-ri zo\-lìm oel mì sa'\-leng a 'u\-ot a lu txa' sì ekx\-txu.} While playing with my cat I felt something hard and rough on his skin.~\LINK{http://naviteri.org/2012/11/renu-ayinanfyaya-the-senses-paradigm/}{b} 
\B{win\-a u\-van si}, play a quick game. \LINK{http://wiki.learnnavi.org/Canon/2010/UltxaAyharyu\%C3\%A4\#Refinements_of_si-construction_verbs}{wiki} \\
\textsc{Derivations}: \ding{43}
\on{1}.~\HL{lI'uvan}{\B{lì\-'u\-van}}, wordplay, pun. 
\on{2}.~\HL{pamuvan}{\B{pam\-u\-van}}, soundplay.

% ‹uy›
\hi
‹\HT{uy}{\B{uy}}› \I{[u.j]} \emph{inf. in} ‹2› (to indicate formal or ceremonial speech).   
\B{Fay\-sul\-fä\-tu\-ä tì\-kang\-kem ohe\-ru meuia lu\-yu nì\-ngay.} The work of these masters truly honors me.~\LINK{http://naviteri.org/2010/06/first-post/}{b}
\B{Ä\-txä\-le su\-yi o\-he pi\-vawm, pe\-o\-lo' lu\-yu pum nge\-nge\-yä?} May I ask what tribe you belong to?~\LINK{http://naviteri.org/2010/09/getting-to-know-you-part-2/}{b} 
(\emph{a. in storytelling}) \emph{opening formula} \B{Ay\-fi\-za\-yu pll\-txu\-ye san trro la\-mu krr a …}, Once upon a time … [lit.: the ancestors say, once upon a time] 
\emph{closing formula} \B{Fì\-fya pll\-txu\-ye ay\-fi\-za\-yu.} Thus say the ancestors.~\LINK{http://forum.learnnavi.org/intermediate/translation-exercise-a-navajo-coyote-tale-in-navi/}{ln}
\\
\end{multicols}

% ===== Section V =====

 \pdfbookmark[0]{V}{V}
\begin{center}
\section*{V \I{[v]} (\emph{Vä})}
\end{center}

\begin{multicols}{2}

% Va'ru
\hi
\B{\cp{Va'}ru} \I{["vaP.Ru]} \emph{n., male name}.
\B{Fì\-tey\-lu\-ri ke nar\-mew Va'\-ru yi\-vom nì\-yll.} Va'ru didn't want to share this \emph{teylu} with the Omatikaya.~\LINK{http://naviteri.org/2012/01/more-additions-to-the-lexicon/}{b}
\B{Ay\-ha\-pxì\-tu po\-ngu\-ä txo\-pu si nì\-nän ta\-krr\-a Va'\-rul pxe\-ku\-tut lä\-txayn.} The members of the group are less afraid since Va'ru defeated three of the enemy.~\LINK{http://naviteri.org/2012/02/trr-asawnung-lefpom-happy-leap-day/}{b}
\B{Hu\-fwa lu fil\-ur Va'\-ru\-ä fne\-fe'\-ran\-vi, tsal\-su\-ngay fpìl fu\-ta say\-rìp lu nì\-txan.} Although Va'ru's facial stripes are rather uneven, I still think he's very handsome.~\LINK{http://naviteri.org/2012/10/mipa-vospxi-mipa-ayliu-new-words-for-the-new-month/}{b}

% vakx
\hi 
\HT{vakx}{\B{vakx}} \I{[vak']} \emph{n., fauna} snake.  

% val
\hi 
\HT{val}{\B{val}} \I{[val]} \emph{adv.}, \textsc{rn} diligently, hard, with effort. 
\B{Mak\-to val!} Ride hard!~\LINK{http://naviteri.org/2022/12/trr-anawm-polaheiem-the-great-day-has-arrived/}{b} 
\B{Po tì\-kang\-kem so\-li val nì\-txan fte tsa\-tson\-ur ha\-sey si\-vi.} She worked very hard to complete the task.~\LINK{http://naviteri.org/2022/12/trr-anawm-polaheiem-the-great-day-has-arrived/}{b}
\B{Tì\-kang\-kem si val!\slash Kang\-kem val!} Work hard!~\LINK{http://naviteri.org/2022/12/trr-anawm-polaheiem-the-great-day-has-arrived/}{b}
(\ding{43}~\HL{kawl}{\B{kawl}}) 

% var 
\hi
\HT{var}{\B{var}} \I{[vaR]} (\emph{a. vin.}) to persist in a state, to continue to perform an action.   
\B{Pol trr\-am ye\-rik\-it tar\-ma\-ron. -- Fì\-trr var.} He was hunting hexapede yesterday. -- He still is today.~\LINK{http://forum.learnnavi.org/language-updates/as-x-as-y-keep-on-keepin-on/}{ln}
\B{Aw\-nga ze\-ne vi\-var ìlä fì\-sa\-lew\-fya. Zen\-ke yak si\-vi.} We must continue in this direction. We must not go astray.~\LINK{http://naviteri.org/2013/09/on-si-salewfya-shapes-and-directions/}{b}
(\emph{b. modal}) continue to, keep on. 
\B{Var ni\-vu\-me ko!} Let's keep learning!~\LINK{http://forum.learnnavi.org/language-updates/as-x-as-y-keep-on-keepin-on/}{ln}
\B{Oe\-ri ay\-som\-pì\-va sìn re\-'o var zi\-vup.} Raindrops keep falling on my head.~\LINK{http://naviteri.org/2011/04/yafkeykiri-plltxe-frapo-everyone-talks-about-the-weather/}{b}
\B{Ke lu kaw\-tu a nga\-to kem so\-li nì\-'ul fte tsi\-vun lì'\-fya le\-Na'\-vi vi\-var 'i\-vong.} There is nobody that does more than you so that the Na'vi language can continue to bloom.~\LINK{http://naviteri.org/2012/07/meetings-waterfalls-and-more/comment-page-1/\#comment-1547}{b}
(\emph{c. idiom}) \B{Pol sä\-fpìl\-it ve\-rar wi\-van.} He's keeping his idea a secret.~\LINK{http://naviteri.org/2011/02/new-vocabulary-ii\%E2\%80\%94part-1/}{b}

% Varang
\hi
\HT{Varang}{\B{\cp{Va}rang}} \I{["va.RaN]} \emph{n., female name}.

% vawm 
\hi
\HT{vawm}{\B{vawm}} \I{[vawm]} \emph{adj.} deep dark colors including browns.
\B{vawm na nik\-re}, the dark color of Na'\-vi hair.~\LINK{http://naviteri.org/2010/08/mipa-ayopin-mipa-ayliu-new-colors-new-words/}{b}
\B{vawm na uk}, dark\hyp{}shadow color.~\LINK{http://naviteri.org/2010/08/mipa-ayopin-mipa-ayliu-new-colors-new-words/}{b} 
\B{Tsaw\-ke slo\-lu vawm ta\-lun fwopx.} The sun became dark due to dust in the air.~\LINK{http://naviteri.org/2014/07/tsawlultxamaw-a-ayliu-post-meetup-words/}{b} \\
\textsc{Derivatives}: \ding{43}
\on{1}.~\HL{tIvawm}{\B{tì\-vawm}}, darkness. 
\on{2}.~\HL{vawmpin}{\B{vawm\-pin}}, the color \B{vawm}. 
\on{3}.~\HL{kllvawm}{\B{kll\-vawm}}, brown.

% vawmpin
\hi
\HT{vawmpin}{\B{\cp{vawm}pin}} \I{["vawm.pin]} \emph{n.} the color \ding{43}~\HL{vawm}{\B{vawm}}.
(\ding{43}~\HL{'opin}{\B{'o\-pin}})

% vawt 
\hi
\HT{vawt}{\B{vawt}} \I{[vawt\textcorner]} \emph{adj.} solid.
\B{Fì\-ut\-ral\-ìri ta\-ngek\-ä zir fkan vawt, hu\-fwa ke rey.} The trunk of the tree feels solid, although it's dead.~\LINK{http://naviteri.org/2013/09/on-si-salewfya-shapes-and-directions/}{b}
(\ding{214}~\HL{momek}{\B{mo\-mek}})

% vay
\hi
\HT{vay}{\B{vay}} \I{[vaj]} \emph{adp.} (\emph{a. local}) up to. 
\B{Ne\-ni lew si fì\-txa\-yor vay txam\-pay nì\-wotx.} Sand covers this expanse all the way to the ocean.~\LINK{http://naviteri.org/2014/07/tsawlultxamaw-a-ayliu-post-meetup-words/}{b}
(\emph{b. temporal}) until.
\B{Trr\-'ong\-ta Txon\-'ong\-vay po to\-lì\-ran.} He walked from dawn until dusk.~\LINK{http://naviteri.org/2011/01/new-year-new-vocabulary/}{b}
\B{Vay fwa zo\-la\-'u Tsyeyk\-Su\-li, To\-ruk Mak\-to alu pi\-za\-yu Ney\-ti\-ri\-yä txan\-rar\-mo\-'a fra\-to kip ay\-ha\-pxì\-tu Oma\-ti\-ka\-yaä.} Until Jake Sully arrived, Neytiri's ancestor was the most famous Toruk Makto among the Omatikaya.~\LINK{http://naviteri.org/2012/03/spring-vocabulary-part-1/}{b} \\
(\emph{c. with numbers}) (\emph{count up}) to. 
\B{Ti\-am (ta pxey) vay vol.} Count (from three) to eight.~\LINK{https://forum.learnnavi.org/language-updates/'count-to'-wa-and-sequential-genitive-(discussion)/}{ln} \\
(i.) \B{vay set}, up to now. \B{Fo\-ru 'upxa\-ret oel fpo\-le', slä vay set ke pa\-mä\-hä\-ngem kea tì\-'eyng.} I have sent them a message, but up to now no answer has arrived.~\LINK{http://forum.learnnavi.org/news-announcements/a-response-from-paul-frommer!/}{ln}
\B{Lì'\-fya\-ri le\-Na'\-vi 'Rr\-ta\-mì, vay set 'al\-mong a fra\-'u ze\-ra\-'u ta ngrr\-po\-ngu.} Everything that has gone on with (blossomed regarding) Na'vi until now on Earth has come from a grassroots movement.~\LINK{http://wiki.learnnavi.org/index.php?title=Corpus\#Good_Morning_America}{wiki} \\
(ii.) \B{tsa\cp{krr}vay}, \emph{conv.} until then, see you then. \B{Ul\-txa si\-vi oeng sìn ram\-tsyìp txon\-'ong\-ay. -- Ye\-io'! Tsa\-krr\-vay ko!} Let's meet on the hill tomorrow at nightfall. -- Perfect! See you then.~\LINK{http://naviteri.org/2012/01/mipa-zisit-ayliu-amip-new-words-for-the-new-year/}{b} \\
(iii.) \B{txan\cp{txew}vay}, maximally, as … as possible. \B{Tì\-ran nì\-fnu txan\-txew\-vay fte\-ke ay\-ye\-rik\-ìl aw\-nga\-ti sti\-vawm.} Walk as quietly as possible so the hexapedes won't hear us.~\LINK{http://naviteri.org/2012/10/snewsyea-ftxoza-halowina-a-spooky-halloween/}{b} \\
(iv.) \B{ha\cp{ya}lovay}, until next time.~\LINK{http://naviteri.org/2010/07/diminutives-conversational-expressions/}{b} 
(\ding{214}~\HL{raw}{\B{raw}}) \\
\textsc{Derivatives}: \ding{43}
\on{1}.~\HL{vaykrr}{\B{vay\-krr}}, until (\emph{conj.}).

% vaykrr
\hi
\HT{vaykrr}{\B{vay\cp{krr}}} \I{[vaj."k\s{r}]} \emph{conj.} until.
(\ding{43}~\HL{vay}{\B{vay}}, \HL{krr}{\B{krr}})

% vä'
\hi
\HT{vA'}{\B{vä'}} \I{[v\ae{}P]} \emph{adj.} unpleasant to the senses. 
\B{Fay\-sä\-re\-ngop\-it a\-vä' oe\-ru rä\-'ä wìn\-txu nì\-mun, ru\-txe. Ke su\-nu oer keng nì\-'it.} Please don't show me these ugly designs again. I don't like them one bit.~\LINK{http://naviteri.org/2013/01/awvea-posti-zisita-amip-first-post-of-the-new-year/}{b}
(\ding{214}~\HL{lor}{\B{lor}}) \\
\textsc{Derivatives}: \ding{43}
\on{1}.~\HL{ftxIvA'}{\B{ftxì\-vä'}}, bad\hyp{}tasting, disgusting. 
\on{2}.~\HL{onvA'}{\B{on\-vä'}}, bad\hyp{}smelling. 
\on{3}.~\HL{mikvA'}{\B{mikvä'}}, bad\hyp{}sounding. 
\on{4}.~\HL{narvA'}{\B{nar\-vä'}}, ugly, unsightly.  
\on{5}.~\HL{vApam}{\B{vä\-pam}}, noise, screech. 
\on{6}.~\HL{vAfewll}{\B{vä\-fe\-wll}}, Centipede. 
\on{7}.~\HL{vArumut}{\B{vä\-ru\-mut}}, vein pod.
\on{8}.~\HL{txavA'}{\B{txa\-vä'}}, disgusting.

% väfewll
\hi
\HT{vAfewll}{\B{\cp{vä}fewll}} \I{["v\ae.fE.w\s{l}]} \emph{n.}, \ding{96} Centipede, ``bad\hyp{}smell plant'', \emph{scolopendra virens}. 
(\ding{43}~\HL{vA'}{\B{vä'}}, \HL{fe'}{\B{fe'}}, \HL{'ewll}{\B{'e\-wll}})

% väng 
\hi
\HT{vAng}{\B{väng}} \I{[v\ae{}N]} \emph{adj.} thirsty.
\B{Nga 'e\-fu väng srak? -- Sran, 'e\-fu nì\-txan.} Are you thirsty? -- Yep, I'm parched.~\LINK{http://naviteri.org/2012/10/audio-and-video-learning-materials-for-navi-101/}{b}
\B{Sre\-krr 'a\-me\-fu väng, set ye\-väng.} Before, I was thirsty; now my thirst has been quenched.~\LINK{http://naviteri.org/2011/04/\%E2\%80\%99a\%E2\%80\%99awa-li\%E2\%80\%99fyavi-amip\%E2\%80\%94a-few-new-expressions/}{b} \\
\textsc{Derivatives}: \ding{43}
\on{1}.~\HL{tIvAng}{\B{tì\-väng}}, thirst. 
\on{2}.~\HL{yevAng}{\B{ye\-väng}}, satisfied, quenched.

% väpam
\hi
\HT{vApam}{\B{\cp{vä}pam}} \I{["v\ae.pam]} \emph{n.} noise (\emph{ugly or unpleasant sound}), screech.
\B{Ni\-nat\-ìri tì\-ru\-sol Txe\-wì\-yä lu vä\-pam.} To Ninat, Txewì's singing is noise.~\LINK{http://naviteri.org/2014/11/twenty-before-the-holidays/}{b}
(\ding{43}~\HL{hawmpam}{\B{hawm\-pam}}, \HL{vA'}{\B{vä'}}, \HL{pam}{\B{pam}})

% värumut
\hi
\HT{vArumut}{\B{\cp{vä}rumut}} \I{["v\ae.Ru.mut\textcorner]} \emph{or} \I{[.mUt\textcorner]} \emph{n.}, \ding{96} vein pod, ``bad ball tree'', \emph{obesus fenestralis}. 
(\ding{43}~\HL{vA'}{\B{vä'}}, \HL{rum}{\B{rum}}, \HL{utral}{\B{ut\-ral}})

% -ve
\hi
\HT{-ve}{\B{-ve}} \I{[.vE]} \emph{suf. to form ordinal numbers}, e.g. \B{'aw\-ve}, first, \B{mu\-ve}, second, etc.
(\ding{43}~\HL{ve'o}{\B{ve\-'o}})

% ve'kì
\hi
\HT{ve'kI}{\B{ve'\cp{kì}}} \I{[vEP."kI]} \emph{vtr.} hate.
\B{Omum oel fu\-ta ngal pot ve'\-kì, slä tse\-fta ke lu tì\-'eyng a\-muiä.} I know you hate him, but vengeance is not a proper response.~\LINK{http://naviteri.org/2022/11/krr-ao-an-exciting-time/}{b}
\B{Txe\-wì\-yä sem\-pul\-ìri sngä\-'i\-krr Saw\-tu\-tet ve\-re'\-kì.} In the beginning, Txe\-wì's father hated the Skypeople.~\LINK{http://naviteri.org/2010/07/ting-mikyun-fte-tslivam-listening-comprehension-1-3/}{b} \\
\textsc{Derivatives}: \ding{43}
\on{1}.~\HL{tIve'kI}{\B{tì\-ve'\-kì}}, hatred.

% ve'o
\hi
\HT{ve'o}{\B{\cp{ve}'o}} \I{["vEP.o]} \emph{n.} order (\emph{as opposed to disorder or chaos}), organization.   
\B{Maw\-krr\-a Saw\-tu\-te\-yä txam\-pxì ho\-lum, tä\-txaw Na'\-vi\-ne nì\-'i'\-a ve\-'o.} After most of the Sky People left, order finally returned to the People.~\LINK{http://naviteri.org/2012/10/mipa-vospxi-mipa-ayliu-new-words-for-the-new-month/}{b} \\
(i.) \B{\cp{ve}'o ayru\cp{sey}ä}, ecosystem.~\LINK{https://naviteri.org/2026/04/tsivola-liu-amip-thirty-two-new-words/}{b}
(\ding{214}~\HL{keve'o}{\B{ke\-ve\-'o}}) \\
\textsc{Derivatives}: \ding{43}
\on{1}.~\HL{vezo}{\B{ve\-zo}}, be in order, be organzied.   
\on{2}.~\HL{vezeyko}{\B{ve\-zey\-ko}}, put in order, organize. 
\on{3}.~\HL{vefya}{\B{ve\-fya}}, system, process, procedure, approach. 
\on{4}.~\HL{velke}{\B{vel\-ke}}, chaotic, messy, disorganized, in shambles. 
\on{5}.~\HL{venga'}{\B{ve\-nga'}}, organized, `on top of things'. 
\on{6}.~\HL{keve'o}{\B{ke\-ve\-'o}}, chaos, disorder.
\on{7}.~\HL{velun}{\B{ve\-lun}}, logic.

% vefya
\hi
\HT{vefya}{\B{\cp{ve}fya}} \I{["vE.fja]} \emph{n.} system, process, procedure, approach.   
\B{Ney\-ti\-ril Tsyey\-kur wa\-mìn\-txu Oma\-ti\-ka\-ya\-ä ve\-fyat tì\-tu\-sa\-ron\-ä.} Ney\-ti\-ri showed Jake the Omatikaya's approach to hunting.~\LINK{http://naviteri.org/2012/10/mipa-vospxi-mipa-ayliu-new-words-for-the-new-month/}{b}
(\ding{43}~\HL{ve'o}{\B{ve\-'o}}, \HL{fya'o}{\B{fya\-'o}})

% vekreng
\hi 
\HT{vekreng}{\B{\cp{vek}reng}} \I{["vEk\textcorner.REN]} \emph{n., fauna} cloaked panther. 

% velek
\hi
\HT{velek}{\B{\cp{ve}lek}} \I{["vE.lEk\textcorner]} \emph{vin.} give up, surrender, concede defeat.   
\B{Tì\-'i'\-a\-ri tsam\-ä ze\-ne Saw\-tu\-te vi\-ve\-lek ta\-lun tì\-txur Ey\-wa\-yä.} At the end of the war, the Sky People had to give up due to the power of Ey\-wa.~\LINK{http://naviteri.org/2012/03/spring-vocabulary-part-1/}{b}
(\ding{43}~\HL{ra'un}{\B{ra\-'un}})

% velke
\hi
\HT{velke}{\B{\cp{vel}ke}} \I{["vEl.kE]} \emph{adj.} chaotic, messy, disorganized, in shambles.   
\B{Eyk Ka\-mun a fra\-lo lä\-ngu tsa\-sä\-ta\-ron vel\-ke nì\-wotx. Ta\-ron\-yut yom smar\-ìl!} Every time Ka\-mun is in charge, the hunt is a mess. Everything goes wrong that can!~\LINK{http://naviteri.org/2012/10/mipa-vospxi-mipa-ayliu-new-words-for-the-new-month/}{b}
(\ding{43}~\HL{ve'o}{\B{ve'o}}, \HL{luke}{\B{lu\-ke}})

% velun
\hi 
\HT{velun}{\B{ve\cp{lun}}} \I{[vE."lun]} \emph{or} \I{[."lUn]} \emph{n.} logic. 
\B{Nge\-yä tì\-hawl\-ìri ke lä\-ngu kea ve\-lun.} I'm sorry to say there's no logic in your plan.~\LINK{https://naviteri.org/2024/09/pxevola-liu-amip-two-dozen-new-words/}{b}
(\ding{43}~\HL{ve'o}{\B{ve\-'o}}, \HL{lun}{\B{lun}})  

% ventil
\hi
\HT{ventil}{\B{\cp{ven}til}} \I{["vEn.til]} \emph{n.} ankle.
(\ding{43}~\HL{venu}{\B{ve\-nu}}, \HL{til}{\B{til}})

% venu
\hi
\HT{venu}{\B{\cp{ve}nu}} \I{["vE.nu]} \emph{n.} foot.
\B{Ri\-ni fmar\-mi hi\-vi\-fwo slä ve\-nu no\-lip äo tskxe.} Ri\-ni tried to escape but her foot got caught under a rock.~\LINK{http://naviteri.org/2015/04/some-new-words-for-may-day/}{b}
\B{Oel tse\-ri fu\-ta nga tskawr. Sra\-ke ngal ve\-nut tì\-sraw sey\-ko\-li?} I see you're limping. Did you hurt your foot?~\LINK{http://naviteri.org/2019/06/50a-liu-amip-40-new-words/}{b}  \\
\textsc{Derivatives}: \ding{43}
\on{1}.~\HL{venzek}{\B{ven\-zek}}, toe. \\
\on{2}.~\HL{hawnven}{\B{hawn\-ven}}, shoe. \\
\on{3}.~\HL{ventil}{\B{ven\-til}}, ankle.

% venzek
\hi
\HT{venzek}{\B{\cp{ven}zek}} \I{["vEn.zEk\textcorner]} \emph{n.} toe. 
(\ding{43}~\HL{venu}{\B{ve\-nu}}, \HL{zekwA}{\B{zek\-wä}})

% venga'
\hi
\HT{venga'}{\B{\cp{ve}nga'}} \I{["vE.NaP]} \emph{adj.} organized, `on top of things'.
\B{Txo ni\-vew fko sä\-ro\-'a si\-vi, ze\-ne nì\-'aw\-ve ve\-nga' li\-vu.} If you want to accomplish great things, you first have to be organized.~\LINK{http://naviteri.org/2012/10/mipa-vospxi-mipa-ayliu-new-words-for-the-new-month/}{b}
\B{Fwa fpìl nì\-fya\-'o a\-ve\-nga' fì\-tì\-zin\-mì lu kel\-tsun.} It's impossible to think clearly (i.e., in an organized manner) in this muddled situation.~\LINK{http://naviteri.org/2021/04/mipa-ayliu-mipa-sioeykting-new-words-new-explanations/\#comment-33483}{b}
(\ding{43}~\HL{ve'o}{\B{ve'o}}, \HL{-nga'}{\B{-nga'}})

% vewng
\hi
\HT{vewng}{\B{vewng}} \I{[vEwN]} \emph{vtr.} look after, take care of.
\B{Oel vewng frr\-nen\-it.} I look after the infants.~\LINK{http://naviteri.org/2010/09/getting-to-know-you-part-2/}{b}
\B{Oel vewng ay\-sngel\-it.} I tend to the refuse.~\LINK{http://naviteri.org/2010/09/getting-to-know-you-part-2/}{b}
\B{Oel vewng fu\-ta ay\-e\-veng ni\-vu\-me te\-ri ay\-ewll na'\-rìng\-ä.} I see to it that the children learn about the forest plants.~\LINK{http://naviteri.org/2010/09/getting-to-know-you-part-2/}{b}
\B{Mong prr\-nen\-ìl fu\-ta sa'\-nul ve\-rewng nì\-tut.} The infant depends on her Mommy to look after her constantly.~\LINK{http://naviteri.org/2012/07/meetings-waterfalls-and-more/}{b}

% vey
\hi
\HT{vey}{\B{vey}} \I{[vEj]} \emph{n.} food of animal origin, meat; flesh.
\B{Ngay\-txoa, ke tsun oe yi\-vom vey\-ti.} Sorry, I can't eat meat.~\LINK{http://naviteri.org/2012/10/audio-and-video-learning-materials-for-navi-101/}{b} \\
\textsc{Derivatives}: \ding{43}
\on{1}.~\HL{yomvey}{\B{yom\-vey}}, dine on flesh, be carnivorous.

% vezeyko
\hi
\HT{vezeyko}{\B{vezey\cp{ko}}} \I{[vE.zEj."k$\cdot$o]} \emph{vtr., inf.}\on{33} put in order, organize.
\B{Nga\-ri sno\-mot krr\-pe ve\-zey\-ko, ma 'itan?} When are you going to organize your room, son?~\LINK{http://naviteri.org/2012/10/mipa-vospxi-mipa-ayliu-new-words-for-the-new-month/}{b} 
(\ding{43}~\HL{vezo}{\B{ve\-zo}})

% vezo
\hi
\HT{vezo}{\B{ve\cp{zo}}} \I{[vE."z$\cdot$o]} \emph{vin., inf.}\on{22} be in order, be organized.
(\ding{43}~\HL{ve'o}{\B{ve\-'o}}, \HL{zo}{\B{zo}})

% -vi
\hi
\HT{-vi}{\B{-vi}} \I{[.vi]} \emph{unprod. suf. used loosely for the spawn of s.t. bigger.}

% vin 
\hi
\HT{vin}{\B{vin}} \I{[vin]} \emph{vtr.} (\emph{a.}) ask for, request.   
\B{Pol vo\-lin mip\-a tska\-lep\-it.} He asked for a new crossbow.~\LINK{http://naviteri.org/2011/08/reported-speech-reported-questions/}{b}
\B{Vu\-yin ohel Unil\-ta\-ron\-it.} I respectfully request the Dreamhunt.~\LINK{http://naviteri.org/2011/08/reported-speech-reported-questions/}{b}
\B{Ngal new a tsa\-'ut rä\-'ä wi\-vo, ma 'evi. Vi\-vin.} Don't reach for what you want, child. Ask for it.~\LINK{http://naviteri.org/2012/01/mipa-zisit-ayliu-amip-new-words-for-the-new-year/}{b}
(\emph{b. with indirect questions}) 
\B{Vo\-lin oel tì\-'eyng\-it a\slash teyng\-ta Ney\-ti\-ri kä pe\-seng\-ne.} I asked where Neytiri was going.~\LINK{http://naviteri.org/2011/08/reported-speech-reported-questions/}{b}
(\ding{43}~\HL{AtxAle}{\B{ä\-txä\-le si}})

% vingkap
\hi
\HT{vingkap}{\B{\cp{ving}kap}} \I{["viN.kap\textcorner]} \emph{vtr.} occur to one, strike one, pop into one's mind. 
\B{Vì\-ming\-kap oe\-ti fu\-la poe ke li ke pol\-txe san oe za\-sya\-'u.} It just occurred to me that she hasn't yet said she's coming.~\LINK{http://naviteri.org/2011/09/\%E2\%80\%9Cby-the-way-what-are-you-reading\%E2\%80\%9D/}{b} \\
\textsc{Derivatives}: \ding{43}
\on{1}.~\HL{nIvingkap}{\B{nì\-ving\-kap}}, by the way, incidentally.

% virä
\hi
\HT{virA}{\B{vi\cp{rä}}} \I{[vi."R\ae{}]} \emph{vin.} spread, proliferate. 
\B{Fì\-tseng pay\-fya vi\-rä ka ngip a\-reng, ha tsun aw\-nga tsat si\-var sko em\-kä\-fya.} Here the stream spreads over a shallow area, so we can use it as a ford.~\LINK{http://naviteri.org/2019/06/50a-liu-amip-40-new-words!}{b} 
\B{Tsran\-ten nì\-txan fwa ftang\-len aw\-ngal fu\-ta fì\-sä\-spxin vi\-vi\-rä.} It's very important that we prevent this disease from spreading.~\LINK{http://naviteri.org/2020/05/aawa-liu-amip-si-vurway-alor-a-few-new-words-and-a-beautiful-poem/}{b} \\
\textsc{Derivatives}: \ding{43}
\on{1}.~\HL{tIvirA}{\B{tì\-vi\-rä}}, spread, proliferation.

% vitra
\hi
\HT{vitra}{\B{vit\cp{ra}}} \I{[vit\textcorner."Ra]} \emph{n.} soul.
(i.) \B{Ayvit\cp{ra}yä Ra\cp{mu}nong}, the Well of Souls (\emph{the closest connection to Eywa on Pandora, a point of extreme spiritual significance to the Na'vi}).
\B{Oe tsway\-on io Ay\-vit\-ra\-yä Ra\-mu\-nong.} I fly over the Tree of Souls.~\LINK{http://forum.learnnavi.org/language-updates/over-under-and-quickly/}{ln} \\
\textsc{Derivatives}: \ding{43}
\on{1}.~\HL{vitronvA'}{\B{vit\-ron\-vä'}}, asshole, dickhead.

% vitrautral
\hi
\HT{vitrautral}{\B{vit\cp{ra}utral}} \I{[vit\textcorner."Ra.ut\textcorner.Ral]} \emph{n.} Tree of Souls.
(\ding{43}~\HL{vitra}{\B{vit\-ra}}, \HL{utral}{\B{ut\-ral}})

% vitronvä'
\hi
\HT{vitronvA'}{\B{vitron\cp{vä'}}} \I{[vit\textcorner.Ron."v\ae{}P]} \emph{n., vulg.} butthole, asshole, dickhead (\emph{highly abusive and vulgar, and is never used in polite society}). \emph{Shortened to} \ding{43}~\HL{vonvA'}{\B{von\-vä'}}.
(\ding{43}~\HL{vitra}{\B{vit\-ra}}, \HL{onvA'}{\B{on\-vä'}})

% vll
\hi
\HT{vll}{\B{vll}} \I{[v\s{l}]} (\textsc{perf}, \B{vol})  (\emph{a. vin.}) indicate, point at.   
\B{Vll eyk\-yu ne\-fä fte po\-ngu fä\-ki\-vä.} The leader signals the party to ascend by pointing upwards.~\LINK{http://naviteri.org/2012/06/spring-vocabulary-part-3/}{b}
(\emph{b. vtr.}) \B{Txo\-pul pe\-yä vll fu\-ta kaw\-krr ke sla\-yu tsam\-si\-yu.} His fear indicates he'll never become a warrior.~\LINK{http://naviteri.org/2012/06/spring-vocabulary-part-3/}{b}
\B{Ut\-ral\-ti a tsa\-uo Lo\-ak wä\-par\-man pol vol fa kxe\-tse.} She used her tail to point out the tree Loak was hiding behind.~\LINK{http://naviteri.org/2012/06/spring-vocabulary-part-3/}{b}
(\emph{c. metaphorically}) \B{Txo\-pul pe\-yä vll fu\-ta kaw\-krr ke sla\-yu tsam\-si\-yu.} His fear indicates he'll never become a warrior.~\LINK{http://naviteri.org/2012/06/spring-vocabulary-part-3/}{b} \\
\textsc{Derivatives}: \ding{43}
\on{1}.~\HL{sAvll}{\B{sä\-vll}}, sign, indication, signal.

% vo-
\hi
\HT{vo-}{\B{vo-}} \I{[vo.]} \emph{num. pref.} eight (\emph{base form of} \ding{43}~\HL{vol}{\B{vol}} \emph{to form higher numbers and fractions}).
\B{vo\-\cp{mun}}, \on{10} (decimal), \on{12} (octal);
\B{me\-vo\cp{mun}}, \on{18} (decimal), \on{22} (octal); 
\B{vo\cp{pxì}}, one\hyp{}eighth, $\frac{1}{8}$.

% vofu
\hi
\HT{vofu}{\B{vo\cp{fu}}} \I{[vo."fu]} \emph{num., adj.} fourteen, \on{14} (decimal), \on{16} (octal).  

% vohin
\hi
\HT{vohin}{\B{vo\cp{hin}}} \I{[vo."hin]} \emph{num., adj.} fifteen, \on{15} (decimal), \on{17} (octal).  

% voìk / si
\hi 
\HT{voIk}{\B{\cp{vo}ìk}} \I{["vo.Ik\textcorner]} \textsc{i}. \emph{n.} behavior, \emph{how one conducts oneself}. 
\B{Nga sìl\-mi a tsa\-kem ke lu vo\-ìk a\-mu\-iä!} What you just did was not proper behavior!~\LINK{http://naviteri.org/2020/11/vospxivopeya-ayliu-amip-novembers-new-words/}{b}  \\
\textsc{ii}. \B{\cp{vo}ìk si} \I{["vo.Ik\textcorner{} si]} \emph{vin.} behave.
\B{Ney\-ti\-ril Tsyey\-kur oeyk\-to\-lìng teyng\-ta fya\-pe ze\-ne vo\-ìk si\-vi tsa\-tì\-fkey\-tok\-mì.} Neytiri explained to Jake how to behave in that situation.~\LINK{http://naviteri.org/2020/11/vospxivopeya-ayliu-amip-novembers-new-words/}{b} 

% vol
\hi
\HT{vol}{\B{vol}}\textsuperscript{\on{1}} \I{[vol]} \emph{num., adj.} eight (decimal), ten (octal). (\emph{base}: \ding{43}~\B{vo-})
\B{Re\-rol teng\-krr ke\-rä | Ìlä fya\-'o a\-vol | Ne kxam\-tseng.} We sing our way, down the eight paths, to the center.~\LINK{http://www.pandorapedia.com/navi/music/ritual_music}{pp}
\B{Pol ke tslam stum ke\-'ut, o\-mum lì\-'ut a\-vol nì\-'aw.} He understands almost nothing and only knows eight words.~\LINK{http://naviteri.org/2013/01/lifyari-po-peyi-how-good-is-her-navi/}{b} \\
\textsc{Derivatives}: \ding{43}
\on{1}.~\HL{vospxIvol}{\B{Vo\-spxì\-vol}}, August.

\B{vol}\textsuperscript{\on{2}} : \HL{vll}{\B{vll}}

% volaw
\hi
\HT{volaw}{\B{vo\cp{law}}} \I{[vo."law]} \emph{num., adj.} nine (decimal), eleven (octal). 

% volve
\hi
\HT{volve}{\B{\cp{vol}ve}} \I{["vol.vE]} \emph{ord. num., adj.} eighth. 

% vomrr
\hi
\HT{vomrr}{\B{vo\cp{mrr}}} \I{[vo."m\s{r}]} \emph{num., adj.} thirteen, \on{13} (decimal), \on{15} (octal).
\B{Ay\-wayl yìm ki\-fkey\-ä | 'Ì\-he\-yut a\-vo\-mrr | Sìn ti\-rea\-fya\-'o a\-vol | Na way\-te\-lem\-ä hìng.} The songs bind | the thirteen spirals of the solid world | To the eight spirit paths | Like the threads of a Songcord.~\LINK{http://www.pandorapedia.com/navi/music/ritual_music}{pp}

% vomun
\hi
\HT{vomun}{\B{vo\cp{mun}}} \I{[vo."mun]} \emph{or} \I{[.mUn]} \emph{num., adj.} ten, \on{10} (decimal), \on{12} (octal).
\B{U\-van Vo\-mun\-a Pam\-rel\-vi\-yä.} The Ten Letter Game.~\LINK{http://www.pbs.org/wgbh/nova/blogs/secretlife/blogposts/learn-navi-with-avatars-paul-frommer/}{nova}
\B{Txe\-wì lu 'e\-veng\-an a so\-la\-lew zì\-sìt a\-vo\-mun.} Txewì is a ten\hyp{}year\hyp{}old boy.~\LINK{http://naviteri.org/2010/07/ting-mikyun-fte-tslivam-listening-comprehension-1-3/}{b} 

% vonvä'
\hi 
\HT{vonvA'}{\B{von\cp{vä'}}} \I{[von."v\ae{}P]}, \emph{coll. pronounced as} \I{[v\~{o}."v\ae{}P]} \emph{n., vulg.} butthole, asshole, dickhead (\emph{highly abusive and vulgar, and is never used in polite society}). \emph{Shortened from} \HL{vitronvA'}{\B{vit\-ron\-vä'}}.
(\ding{43}~\HL{vitra}{\B{vit\-ra}}, \HL{onvA'}{\B{on\-vä'}})

% vopey
\hi
\HT{vopey}{\B{vo\cp{pey}}} \I{[vo."pEj]} \emph{num., adj.} eleven, \on{11} (decimal), \on{13} (octal).

% vopxì
\hi
\HT{vopxI}{\B{vo\cp{pxì}}} \I{[vo."p'I]} \emph{n., num.} one\hyp{}eighth, $\frac{1}{8}$.

% vosìng
\hi
\HT{vosIng}{\B{vo\cp{sìng}}} \I{[vo."sIN]} \emph{num., adj.} twelve, \on{12} (decimal), \on{14} (octal).

% vospxey
\hi 
\HT{vospxey}{\B{Vo\cp{spxey}}} \I{[vo."{}sp'Ej]} \emph{n.} March.
(\ding{43}~\HL{vosxpI}{\B{vospxì}}, \HL{pxey}{\B{pxey}})

% vospxì
\hi
\HT{vospxI}{\B{vo\cp{spxì}}} \I{[vo."{}sp'I]} \emph{n.} month (\emph{from} \B{vo\-sì\-pxì zì\-sìt\-ä}, a twelfth of the year).
\B{Mip\-a Vo\-spxì, Mip\-a Ay\-lì\-'u.} New Words for the New Month.~\LINK{http://naviteri.org/2012/10/mipa-vospxi-mipa-ayliu-new-words-for-the-new-month/}{b}
\B{Ta\-lu\-na vo\-spxì\-o a\-mrr ke zo\-lup tom\-pa, lä\-ngu yì kil\-van\-ä tìm nì\-txan.} Because it hasn't rained for five months, the river level is very low.~\LINK{http://naviteri.org/2013/01/lifyari-po-peyi-how-good-is-her-navi/}{b}
\B{Fì\-vo\-spxì\-yä Ay\-lì'\-fya\-vi A\-mip.} This Month's New Expressions.~\LINK{http://naviteri.org/2012/07/fivospxiya-aylifyavi-amip-this-months-new-expressions-pt-2/}{b} \\
\textsc{Derivatives}: \ding{43}
\on{1}. \B{vospxì\cp{am}} \I{[vo.sp'I."am]} \emph{n., adv.} last month. \B{Krr\-ka tsawl\-ul\-txa vo\-spxì\-am.} During the great meeting last month.~\LINK{http://naviteri.org/2014/08/stxeli-alor-a-beautiful-gift/}{b} \\
\on{2}.	\B{vospxì\cp{ay}} \I{[vo.sp'I."aj]} \emph{n., adv.} next month. \B{Vo\-spxì\-ay\-vay!} Until next month.~\LINK{http://naviteri.org/2013/02/vospxi-ayol-posti-apup-short-post-for-a-short-month/}{b} 

% vospxì'aw
\hi 
\HT{vospxI'aw}{\B{Vospxì\cp{'aw}}} \I{[vo.sp'I."Paw]} \emph{also coll.} \I{[.spI.]} \emph{n.} January.  
(\ding{43}~\HL{vospxI}{\B{vo\-spxì}}, \HL{'aw}{\B{'aw}})

% vospxìkin
\hi 
\HT{vospxIkin}{\B{Vospxì\cp{kin}}} \I{[vo.sp'I."kin]} \emph{also coll.} \I{[.spI.]} \emph{n.} July. 
\B{Sìl\-pey oe, ay\-nga\-ri nì\-wotx Vo\-spxì\-kin sngil\-vä\-'i nìl\-tsan.} I hope that as for your all July has started well.~\LINK{http://naviteri.org/2023/07/trr-tsyimawnuniya-lefpom-happy-independence-day/}{b} 
\B{Vo\-spì\-kin Le\-fpom!} Happy July!~\LINK{https://naviteri.org/2025/06/vospikin-lefpom-happy-july/}{b}
(\ding{43}~\HL{vospxI}{\B{vo\-spxì}}, \HL{kinA}{\B{ki\-nä}})

% vospxìmrr
\hi 
\HT{vospxImrr}{\B{Vospxì\cp{mrr}}} \I{[vo.sp'I."m\s{r}]} \emph{also coll.} \I{[.spI.]} \emph{n.} May. 
(\ding{43}~\HL{vospxI}{\B{vo\-spxì}}, \HL{mrr}{\B{mrr}})

% vospxìmun
\hi 
\HT{vospxImun}{\B{Vospxì\cp{mun}}} \I{[vo.sp'I."mun]} \emph{or} \I{[."mUn]} \emph{also coll.} \I{[.spI.]} \emph{n.} February. 
(\ding{43}~\HL{vospxI}{\B{vo\-spxì}}, \HL{mune}{\B{mu\-ne}})

% vospxìpuk
\hi 
\HT{vospxIpuk}{\B{Vospxì\cp{puk}}} \I{[vo.sp'I."puk\textcorner]} \emph{or} \I{[."pUk\textcorner]} \emph{also coll.} \I{[.spI.]} \emph{n.} June.  
(\ding{43}~\HL{vospxI}{\B{vo\-spxì}}, \HL{pukap}{\B{pu\-kap}})

% vospxìtsìng
\hi 
\HT{vospxItsIng}{\B{Vospxì\cp{tsìng}}} \I{[vo.sp'I."\t{ts}IN]} \emph{also coll.} \I{[.spI.]} \emph{n.} April. 
(\ding{43}~\HL{vospxI}{\B{vo\-spxì}}, \HL{tsIng}{\B{tsìng}})

% vospxìvol
\hi 
\HT{vospxIvol}{\B{Vospxì\cp{vol}}} \I{[vo.sp'I."vol]} \emph{also coll.} \I{[.spI.]} \emph{n.} August.  
(\ding{43}~\HL{vospxI}{\B{vo\-spxì}}, \HL{vol}{\B{vol}})

% vospxìvolaw
\hi 
\HT{vospxIvolaw}{\B{Vospxìvo\cp{law}}} \I{[vo.sp'I.vo."law]} \emph{also coll.} \I{[.spI.]} \emph{n.} September. 
(\ding{43}~\HL{vospxI}{\B{vo\-spxì}}, \HL{volaw}{\B{volaw}})

% vospxìvomun
\hi 
\HT{vospxIvomun}{\B{Vospxìvo\cp{mun}}} \I{[vo.sp'I.vo."mun]} \emph{or} \I{[."mUn]} \emph{also coll.} \I{[.spI.]} \emph{n.} October. 
(\ding{43}~\HL{vospxI}{\B{vo\-spxì}}, \HL{vomun}{\B{vo\-mun}})

% vospxìvopey
\hi 
\HT{vospxIvopey}{\B{Vospxìvo\cp{pey}}} \I{[vo.sp'I.vo."pEj]} \emph{also coll.} \I{[.spI.]} \emph{n.} November. 
\B{Vo\-spxì\-vo\-pey\-ä ay\-lì\-'u a\-mip} November's new words.~\LINK{http://naviteri.org/2020/11/vospxivopeya-ayliu-amip-novembers-new-words/}{b}
(\ding{43}~\HL{vospxI}{\B{vo\-spxì}}, \HL{vopey}{\B{vo\-pey}})

% vospxìvosìng
\hi 
\HT{vospxIvosIng}{\B{Vospxìvo\cp{sìng}}} \I{[vo.sp'I.vo."{}sIN]} \emph{also coll.} \I{[.spI.]} \emph{n.} December. 
(\ding{43}~\HL{vospxI}{\B{vo\-spxì}}, \HL{vosIng}{\B{vo\-sìng}})

% vozam
\hi
\HT{vozam}{\B{vo\cp{zam}}} \I{[vo."zam]} \emph{num., adj.} \on{512} (decimal), \on{1.000} (octal).  
\B{'En si oe, lor\-a tsa\-fkxi\-le\-ri a\-pxay\-opin so\-lar Tse\-nul srok\-it a\-vo\-zam.} I would guess that Tsenu used a thousand (lit. \on{512}) beads for that beautiful multi\hyp{}colored bib necklace.~\LINK{http://naviteri.org/2013/07/pximaw-tsawlultxa-just-after-the-meet-up/}{b} \\
\textsc{Derivatives}: \ding{43}
\on{1}.~\HL{vozampasukut}{\B{vo\-zam\-pa\-su\-kut}}, grinch tree.

% vozampasukut
\hi 
\HT{vozampasukut}{\B{\cp{vo}zampasukut}} \I{["vo.zam.pa.suk\textcorner.ut\textcorner]} \emph{or} \I{[.sUk\textcorner.Ut\textcorner]} (\textsc{rn}: \B{\cp{vo}zampasùkut} \I{["vo.zam.pa.sU.kut\textcorner]}) \emph{n.}, \ding{96} grinch tree; thousand berry tree.
(\ding{43}~\HL{vozam}{\B{vo\-zam}}, \HL{pasuk}{\B{pa\-suk}}, \HL{utral}{\B{ut\-ral}})

% vrrìn
\hi
\HT{vrrIn}{\B{\cp{vrr}ìn}} \I{["v\s{r}.In]} \emph{vin.} be busy (\emph{negative sense}), \emph{be tired out and overwhelmed by an activity}.
\B{Tì\-kang\-kem\-ìri var\-mrr\-ìn oe nì\-wotx.} I was completely swamped (overwhelmed) at work.~\LINK{http://naviteri.org/2010/09/getting-to-know-you-part-3/}{b}
\B{Pe\-yä tì\-kang\-kem\-ìl mi vey\-krr\-ei\-yìn pot nì\-wotx.} His work is still completely overwhelming him (and I'm glad).~\LINK{http://naviteri.org/2010/09/getting-to-know-you-part-3/}{b}
(\ding{43}~\HL{'In}{\B{'ìn}}, \ding{214}~\HL{sulIn}{\B{su\-lìn}})

% vrrtep
\hi
\HT{vrrtep}{\B{\cp{vrr}tep}} \I{["v\s{r}.tEp\textcorner]} \emph{n.} demon.   
\B{Txon\-am teng\-krr tar\-mì\-ran oe kxam\-lä na'\-rìng, sro\-ler eo ut\-ral a\-tsawl txewm\-a vrr\-tep.} Last night as I was walking through the forest, a frightening demon appeared in front of a big tree.~\LINK{http://naviteri.org/2012/03/spring-vocabulary-part-1/}{b}
\B{Saw\-tu\-te lu ay\-vrr\-tep nì\-wotx a sä\-fpìl\-it oel wä\-te.} I dispute the idea that the Skypeople are all demons.~\LINK{http://naviteri.org/2011/01/new-year-new-vocabulary/}{b}
\B{Fay\-vrr\-tep fì\-tse\-nge lu kxa\-nì.} These demons are forbidden here.~\LINK{http://wiki.learnnavi.org/index.php?title=Corpus\#Times_Online}{wiki\slash a}
(\ding{43}~\HL{txep}{\B{txep}})

% vul
\hi
\HT{vul}{\B{vul}} \I{[vul]} \emph{or} \I{[vUl]} \emph{n.} \ding{96} branch (of a tree). 
\B{Txo vul ivil nì\-hawng kxì\-ye\-vakx.} If a branch bends too much, it might break.~\LINK{http://naviteri.org/2013/04/wheres-the-bathroom-and-other-useful-things/}{b}
\B{Krro krro, flì\-a vul a\-ru\-sey to nutx\-a pum a\-ke\-ru\-sey lu txur.} A thin living branch is sometimes stronger than a thick dead one.~\LINK{http://naviteri.org/2011/11/\%E2\%80\%99a\%E2\%80\%99awa-ayli\%E2\%80\%99u-amip-ni\%E2\%80\%99aw-\%E2\%80\%94-a-few-new-words-only/}{b}
\B{'Ali\-'ät vul\-sìn key\-kur za\-'u fì\-tseng.} Hang the choker on the branch and come here.~\LINK{http://naviteri.org/2013/04/wheres-the-bathroom-and-other-useful-things/}{b}
\B{A pe\-yä tì\-txur … Mì pun | Na ay\-vul a\-hu\-saw\-nu.} Whose strength is in our arms | As sheltering branches.~\LINK{http://www.pandorapedia.com/navi/music/ritual_music}{pp} \\
\textsc{Derivatives}: \ding{43}
\on{1}.~\HL{vultsyIp}{\B{vul\-tsyìp}}, stick. 
\on{2}.~\HL{pamrelvul}{\B{pam\-rel\-vul}}, pen, pencil. 
\on{3}.~\HL{relvul}{\B{rel\-vul}}, pen, pencil.

% vultsyìp
\hi 
\HT{vultsyIp}{\B{\cp{vul}tsyìp}} \I{["vul.\t{ts}jIp\textcorner]} \emph{or} \I{["vUl.]} \emph{n.} \ding{96} stick (\emph{of a tree, branch}). 
(\ding{43}~\HL{vul}{\B{vul}}) \\
\textsc{Derivatives}: \ding{43}
\on{1}.~\HL{'opinvultsyIp}{\B{'opin\-vul\-tsyìp}}, crayon.

% vun
\hi 
\HT{vun}{\B{vun}} \I{[vun]} \emph{or} \I{[vUn]} \emph{vtr.} provide. 
\B{Sem\-pul\-ìl a\-sìl\-tsan vun syu\-vet soaia\-ru sne\-yä.} A good father provides food to his family.~\LINK{http://naviteri.org/2020/05/aawa-liu-amip-si-vurway-alor-a-few-new-words-and-a-beautiful-poem/}{b}
\B{Ey\-wa va\-yun.} Ey\-wa will provide.~\LINK{http://naviteri.org/2020/05/aawa-liu-amip-si-vurway-alor-a-few-new-words-and-a-beautiful-poem/}{b}  

% vur
\hi
\HT{vur}{\B{vur}} \I{[vuR]} \emph{or} \I{[vUR]} \emph{n.} story. 
\B{Tsa\-vur oe\-yä el\-tur tì\-txen sol\-i.} That story intrigued me.~\LINK{http://forum.learnnavi.org/language-updates/txelanit-hivawl/}{ln}
\B{Nga pam\-rel so\-li san sley\-ki\-vu ut\-ral\-ìl lì'\-fya\-yä le\-Na'\-vi pxay\-a ay\-vur\-it a\-txur sìk.} You wrote, ``May the Na'vi language tree produce many strong stories.''~\LINK{http://forum.learnnavi.org/language-updates/life-death-and-gerunds/}{ln}
\B{Nga\-ri wa\-we fì\-vur\-ä lu 'u\-pe?} What is the significance of this story to you?~\LINK{http://naviteri.org/2011/05/some-miscellaneous-vocabulary/}{b} \\
\textsc{Derivatives}: \ding{43}
\on{1}.~\HL{vurvi}{\B{vur\-vi}}, summary. 
\on{2}.~\HL{'okvur}{\B{'ok\-vur}}, (recent) history. 
\on{3}.~\HL{hangvur}{\B{hang\-vur}}, joke. 
\on{4}.~\HL{vurway}{\B{vur\-way}}, story poem, narrative poem. 
\on{5}.~\HL{vurtu}{\B{vur\-tu}}, fictional character.

% vurtu
\hi 
\HT{vurtu}{\B{\cp{vur}tu}} \I{["vuR.tu]} \emph{or} \I{["vUR.]} \emph{n.} fictional character. 
\B{Txon Ki\-ho\-te lu vur\-tu nì\-'aw; ke fkey\-to\-lok kaw\-krr.} Don Quixote is a fictional character; he never existed.~\LINK{https://naviteri.org/2024/09/pxevola-liu-amip-two-dozen-new-words/}{b} 
(\ding{43}~\HL{vur}{\B{vur}}) 

% vurway
\hi
\HT{vurway}{\B{\cp{vur}way}} \I{["vuR.waj]} \emph{or} \I{["vUR.]} \emph{n.} story poem, narrative poem. 
\B{Vur\-way a\-lor.} A beautiful narrative poem.~\LINK{http://naviteri.org/2017/10/vurway-alor-a-beautiful-narrative-poem/}{b}
\B{Fì\-vur\-wayt ivi\-nan nì\-'o' nì\-'aw!} Have fun reading this narrative poem!~\LINK{http://naviteri.org/2017/10/vurway-alor-a-beautiful-narrative-poem/}{b}

% vurvi
\hi
\HT{vurvi}{\B{vurvi}} \I{["vuR.vi]} \emph{or} \I{["vUR.]} \emph{n.} summary, synopsis. 
(\ding{43}~\HL{vur}{\B{vur}})  \\

\end{multicols}

% ===== Section W =====

 \pdfbookmark[0]{W}{W}
\begin{center}
\section*{W \I{[w]} (\emph{Wä})}
\end{center}

\begin{multicols}{2}

% walak
\hi 
\HT{walak}{\B{\cp{wa}lak}} \I{["wa.lak\textcorner]} \emph{adj.} energetic, active.   
\B{Tì\-tu\-sa\-ron\-ìri txo new fko sli\-vu tsul\-fä\-tu, ze\-ne smar\-to li\-vu wa\-lak.} If you want to become a master hunter, you have to be more active than your prey.~\LINK{http://naviteri.org/2011/10/more-vocabulary-a-bit-of-grammar/}{b}
\B{Ftu\-e lu fwa ta\-ron ngong\-a io\-ang\-it to fwa ta\-ron pum\-it a lu wa\-lak sì win.} It's easier to hunt lethargic animals than to hunt perky, speedy ones.~\LINK{http://naviteri.org/2011/10/more-vocabulary-a-bit-of-grammar/}{b}
(\ding{214}~\HL{ngong}{\B{ngong}})

% walew
\hi 
\HT{walew}{\B{wa\cp{lew}}} \I{[wa."lEw]} \emph{vin.} get over, accept some fact, reconcile oneself, move on. 
\B{Fu\-ri\-a oe yaw\-ne ke lu Va'\-rur nul\-krr, ke tsä\-ngun oe wi\-va\-lew.} I can't get over the fact that Va'ru no longer loves me.~\LINK{http://naviteri.org/2018/07/ta-sulfatu-a-ayliu-niul-more-words-from-our-experts/}{b}
\B{Tì\-ska\-'a\-ri Kel\-ut\-ral\-ä Na'\-vi wa\-ya\-lew pe\-fya?} With the destruction of Hometree, how will the Na’vi ever move on?~\LINK{http://naviteri.org/2018/07/ta-sulfatu-a-ayliu-niul-more-words-from-our-experts/}{b}

% wan
\hi 
\HT{wan}{\B{wan}} \I{[wan]} \emph{vtr.} hide.   
\B{Pol sä\-'o\-ti wo\-lan äo ay\-rìk.} he hid his tool under the leaves.~\LINK{http://naviteri.org/2011/02/new-vocabulary-ii\%E2\%80\%94part-1/}{b}
\B{Nì\-sngä\-'i fmawn\-it fo nar\-mew wi\-van, slä nì\-'i'\-a fra\-por lo\-lo\-nu.} They originally wanted to hide the news, but in the end they revealed it everyone.~\LINK{http://naviteri.org/2011/10/more-vocabulary-a-bit-of-grammar/}{b}
\B{Fnu, fnu, wä\-pan wä ku | Uo tsa\-ìpxa.} Quiet, quiet, hide yourself from harm | Behind that fern.~\LINK{http://forum.learnnavi.org/language-updates/learnings-from-the-siatlla-song-translations/msg473566/\#msg473566}{ln}
\B{Nga pe\-lun wä\-pe\-ran?} Why are you hiding?~\LINK{http://naviteri.org/2011/02/new-vocabulary-ii\%E2\%80\%94part-1/}{b}
\B{Ya\-mì tsun fko tsi\-ve\-'a lo\-ran\-it re\-nu\-ä kil\-van\-ä slä kll\-te\-sìn wä\-pan.} From the air you can see the grace of the river's form but from the ground it's hidden.~\LINK{http://naviteri.org/2012/10/mipa-vospxi-mipa-ayliu-new-words-for-the-new-month/}{b} 
(\emph{a. idiom}) \B{Pol sä\-fpìl\-it ve\-rar wi\-van.} He's keeping his idea a secret.~\LINK{http://naviteri.org/2011/02/new-vocabulary-ii\%E2\%80\%94part-1/}{b} \\  
\textsc{Derivatives}: \ding{43}
\on{1}.~\HL{tIwan}{\B{tì\-wan}}, obfuscation, cover\hyp{}up.   
\on{2}.~\HL{nIwan}{\B{nì\-wan}}, secretly; in hiding, by hiding.   
\on{3}.~\HL{letwan}{\B{let\-wan}}, dodgy, sneaky.
\on{4}.~\HL{wantseng}{\B{wan\-tseng}}, (temporary) hiding place.

% wantseng
\hi 
\HT{wantseng}{\B{\cp{wan}tseng}} \I{["wan.\t{ts}EN]} \emph{n.} (temporary) hiding place. 
\B{Tsa\-ye\-rik\-ìl mu'\-nit\-kan\-a wan\-tseng\-it ro\-lun; ke tsun fko pot tsi\-ve\-'a kaw\-'it.} That hexapede has found an effective hiding place. He can't be seen at all.~\LINK{https://naviteri.org/2025/04/mevosinga-liu-amip-twenty-new-words/}{b} 
(\ding{43}~\HL{wan}{\B{wan}}, \HL{tseng}{\B{tseng}})

% wapx
\hi 
\HT{wapx}{\B{wapx}} \I{[wap']} \emph{vin.} sink. 
\B{Ke omum teyng\-ta fì\-u\-ran a\-ku'\-up fu fwa\-yum fu wa\-yapx.} I don't know (or: It's not known) whether this heavy boat will float or sink.~\LINK{http://naviteri.org/2021/12/zolau-niprrte-ma-3746-welcome-2022/}{b}

% wawe
\hi 
\HT{wawe}{\B{wa\cp{we}}} \I{[wa."wE]} \emph{n.} meaning, importance, significance.   
\B{Nga\-ri wa\-we fì\-vur\-ä lu 'u\-pe?} What is the significance of this story to you?~\LINK{http://naviteri.org/2011/05/some-miscellaneous-vocabulary/}{b}
\B{Wa\-we\-ti ke tsun fko ral\-pi\-veng.} One can't explain (or: put into words) \emph{wawe}.~\LINK{http://naviteri.org/2011/05/some-miscellaneous-vocabulary/}{b} 
\B{Unil\-tì\-ran\-tokx\-ìri lì'\-fya\-yä le\-Na'\-vi wa\-wet fol tslam.} They understand the importance of the Na'vi language for \emph{Avatar}.~\LINK{http://naviteri.org/2013/08/fmawn-a-fkol-kilmulat-this-just-in/}{b} \\  
\textsc{Derivatives}: \ding{43}
\on{1}.~\HL{txanwawe}{\B{txan\-wa\-we}}, personally meaningful, significant.   
\on{2}.~\HL{nIwawe}{\B{nì\-wa\-we}}, meaningfully, significantly.

% way / si
\hi 
\HT{way}{\B{way}} \I{[waj]} \textsc{i}. \emph{n.} song, poetry.
\B{Ay\-wayl yìm ki\-fkey\-ä | 'Ì\-he\-yut a\-vo\-mrr | Sìn ti\-rea\-fya\-'o a\-vol.} The songs bind | the thirteen spirals of the solid world | To the eight spirit paths.~\LINK{http://www.pandorapedia.com/navi/music/ritual_music}{pp}
\B{Fì\-way\-ri hì\-no\-a re\-nut ngam\-pam\-ä ke tsä\-ngun oe tsl\-ivam.} I'm afraid I can't understand the intricate rhyme scheme of this poem.~\LINK{http://naviteri.org/2012/01/mipa-zisit-ayliu-amip-new-words-for-the-new-year/}{b} 
\B{Sìl\-pey oe, tsi\-vun oe nge\-yä ay\-wayt a\-mip i\-vi\-nan ye'\-rìn!} I hope, I can read your new songs soon.~\LINK{http://naviteri.org/2021/04/mipa-ayliu-mipa-sioeykting-new-words-new-explanations/\#comment-33459}{b}
\B{Fì\-way\-ri hì\-no\-a re\-nut ngam\-pam\-ä ke tsä\-ngun oe tsli\-vam.} I'm afraid I can't understand the intricate rhyme scheme of this poem.~\LINK{http://naviteri.org/2012/01/mipa-zisit-ayliu-amip-new-words-for-the-new-year/}{b} 
(i.) \B{way a pll\-\cp{txe}}, spoken poem.~\LINK{http://naviteri.org/2011/12/ulte-yora\%E2\%80\%99tu-leiu-and-the-winner-is/}{b} 
(ii.) \B{way a rol}, song.~\LINK{http://naviteri.org/2011/12/ulte-yora\%E2\%80\%99tu-leiu-and-the-winner-is/}{b} 
(iii.) \B{\cp{way}ä ay\cp{lì}'u}, words of a poem, lyrics of a song.~\LINK{http://naviteri.org/2011/12/ulte-yora\%E2\%80\%99tu-leiu-and-the-winner-is/}{b} \\
\textsc{ii}. \B{\cp{way} si} \I{["waj si]} \emph{vin.} sing (singing of ancient, traditional song)
(\ding{43}~\HL{rol}{\B{rol}}) \\  
\textsc{Derivatives}: \ding{43}
\on{1}.~\HL{waytelem}{\B{way\-te\-lem}}, songcord. 
\on{2}.~\HL{taronway}{\B{ta\-ron\-way}}, hunt song. 
\on{4}.~\HL{vurway}{\B{vur\-way}}, story poem, narrative poem.
\on{5}.~\HL{haway}{\B{ha\-way}}, lullaby.

% waytelem
\hi 
\HT{waytelem}{\B{\cp{way}telem}} \I{["waj.tE.lEm]} \emph{n.} songcord. 
\B{Ay\-wayl yìm ki\-fkey\-ä | 'Ì\-he\-yut a\-vo\-mrr | Sìn ti\-rea\-fya\-'o a\-vol | Na way\-te\-lem\-ä hìng.} The songs bind | the thirteen spirals of the solid world | To the eight spirit paths | Like the threads of a Songcord.~\LINK{http://www.pandorapedia.com/navi/music/ritual_music}{pp} 
\B{Trro la\-yu ngar lor\-a way\-te\-lem a\-ngim fte vur\-it nge\-yä pi\-veng.} One day you will have a long and beautiful songcord to tell your story.~\LINK{https://www.youtube.com/watch?v=13aPkyXDuEE}{yt\slash fop}
(\ding{43}~\HL{way}{\B{way}}, \HL{telem}{\B{te\-lem}})

% wä 
\hi 
\HT{wA}{\B{wä}} \I{[w\ae{}]} \emph{adp.}\texttt{+} (\emph{a.}) against, opposing. 
(\emph{a.1}) \B{fyel wä}, seal from\slash against. \B{Txo fkol ke fyi\-vel u\-ran\-it pay\-wä, ze\-ne fko sli\-ve\-le.} If one does not seal a boat against water, one must swim.~\LINK{http://naviteri.org/2012/06/spring-vocabulary-part-3/}{b}
(\emph{a.2}) \B{haw\-nu wä}, protect from. \B{Fì\-fne\-ya\-yol tsrul\-it txu\-la mì song\-ropx ut\-ral\-ä fte ay\-li\-nit hi\-vaw\-nu wä sar\-ni\-o\-ang. Fì\-tì\-kan\-ìri ropx ke ha'.} This kind of bird builds its nest in a tree cavity so as to protect its young from predators. For this purpose, a hole going right through the tree trunk isn't suitable.~\LINK{http://naviteri.org/2013/04/wheres-the-bathroom-and-other-useful-things/}{b}
(\emph{a.3}) \B{kan wä}, aim s.t. against s.t.\slash s.o.
\B{Pol tsko\-ti ko\-lan wä ku\-tu.} He raised his bow against the enemy.~\LINK{http://naviteri.org/2015/11/vomuna-liu-amip-ten-new-words/comment-page-1/\#comment-22792}{g\slash b}
(\emph{a.4}) \B{kxap wä}, a threat to s.t.\slash s.o., threaten.
\B{Krr\-a kxap lar\-mu sì\-rey\-wä fe\-yä nì\-wotx.} When their lives were threatened.~\LINK{http://naviteri.org/2015/05/cirque-du-soleils-toruk-the-first-flight-teaser-video-online/}{b}
(\emph{a.5}) \B{tsre\-'i wä}, throw at. \B{Oel tsre\-'i tskxet wä nga.} I throw a stone at you.~\LINK{http://forum.learnnavi.org/language-updates/to-throw-at-x/}{ln}
(\emph{a.6}) \B{w(ä\-p)an wä}, hide (oneself) from. \B{Fnu, fnu, wä\-pan wä ku | Uo tsa\-ìpxa.} Quiet, quiet, hide yourself from harm | Behind that fern.~\LINK{http://forum.learnnavi.org/language-updates/learnings-from-the-siatlla-song-translations/msg473566/\#msg473566}{ln}
(\emph{a.7}) \B{wem wä}, fight against. \B{Peyä tsa\-tì\-pe\-'un a swey\-lu txo wi\-vem ayoeng Oma\-ti\-ka\-ya\-wä lu fe'.} His decision to fight against the Omaticaya was a bad one.~\LINK{http://naviteri.org/2011/07/txantsana-ultxa-mi-siatll-great-meeting-in-seattle/}{b} 
(\emph{b.}) contrary to (what somebody said).
\B{Wä Fey\-ral, mun\-txa ke so\-li Ra\-lu sì Ne\-wey nì\-wan me\-srr\-am.} Contrary to what Peyral reported, Ralu and Newey were not secretly married the day before yesterday.~\LINK{https://forum.learnnavi.org/language-updates/'count-to'-wa-and-sequential-genitive-(discussion)/}{ln}
(\emph{b.1}) \B{tì\-'e\-fu\-wä oe\-yä}, ``contrary to my opinion''.
\B{Tì\-'efu\-wä oe\-yä, fpìl Pey\-ral\-ìl fu\-ta ke ze\-ne ay\-oeng ki\-vä.} I think we have to go, but Peyral doesn't.~\LINK{https://forum.learnnavi.org/language-updates/'count-to'-wa-and-sequential-genitive-(discussion)/}{ln} (\ding{214}~\HL{IlA}{\B{ìlä}}) \\
\textsc{Derivatives}: \ding{43}
\on{1}.~\HL{wAsul}{\B{wä\-sul}}, compete.  
\on{2}.~\HL{wAte}{\B{wä\-te}}, argue, dispute. 
\on{3}.~\HL{wAtu}{\B{wä\-tu}}, opponent. 
\on{4}.~\HL{nIwA}{\B{nì\-wä}}, on the contrary, conversely. 
\on{5}.~\HL{wAralI'u}{\B{wä\-ra\-lì\-'u}}, antonym.
\on{6}.~\HL{wAtum}{\B{wä\-tum}}, antidote.

% wäralì'u
\hi 
\HT{wAralI'u}{\B{\cp{wä}ralì'u}} \I{["w\ae.Ra.lI.Pu]} \emph{n.} antonym. 
(\ding{43}~\HL{wA}{\B{wä}}, \HL{ral}{\B{ral}}, \HL{lI'u}{\B{lì\-'u}}; \ding{214}~\HL{tengralI'u}{\B{teng\-ra\-lì\-'u}}) 

% wäsul
\hi 
\HT{wAsul}{\B{\cp{wä}sul}} \I{["w\ae{}.s$\cdot$ul]} \emph{vin., inf.}\on{22} compete. 
\B{Oe new nga\-hu 'aw\-si\-teng tì\-kang\-kem si\-vi -- ke new fu\-ta wä\-si\-vul oeng.} I want to work together with you -- I don't want us to compete.~\LINK{http://naviteri.org/2012/10/wina-sawasultsyip-ahii-a-quick-little-contest/}{b} 
(\ding{43}~\HL{wA}{\B{wä}}, \HL{tul}{\B{tul}}) \\
\textsc{Derivatives}: \ding{43}
\on{1}.~\HL{tIwAsul}{\B{tì\-wä\-sul}}, competition, competing.
\on{2}.~\HL{sAwAsul}{\B{sä\-wä\-sul}}, a competition.

% wäte
\hi 
\HT{wAte}{\B{wä\cp{te}}} \I{[w\ae{}."t$\cdot$E]} \emph{vtr., inf.}\on{22} argue, dispute.   
\B{Saw\-tu\-te lu ay\-vrr\-tep nì\-wotx a sä\-fpìl\-it oel wä\-te.} I dispute the idea that the Skypeople are all demons.~\LINK{http://naviteri.org/2011/01/new-year-new-vocabulary/}{b}
(\ding{43}~\HL{wA}{\B{wä}}, \HL{plltxe}{\B{pll\-txe}}) \\  
\textsc{Derivatives}: \ding{43}
\on{1}.~\HL{tIwAte}{\B{tì\-wä\-te}}, dispute, argument. 
\on{2}.~\HL{sAwAte}{\B{sä\-wä\-te}}, point of contention, source of argument. 
\on{3}.~\HL{lewAte}{\B{le\-wä\-te}}, disagreeable, argumentative. 
\on{4}.~\HL{nIwAte}{\B{nì\-wä\-te}}, disagreeably, begrudgingly.

% wätu
\hi 
\HT{wAtu}{\B{\cp{wä}tu}} \I{["w\ae{}.tu]} \emph{n.} opponent.
\B{Oe pll\-ngay san mo\-lak\-to oe nì\-fe', ta\-fral sno\-laytx; wä\-tu lu oe\-to txur.} I acknowledge that I rode badly, so I lost; my opponent was stronger than I was.~\LINK{http://naviteri.org/2011/07/txantsana-ultxa-mi-siatll-great-meeting-in-seattle/}{b} 
\B{Krr\-a wä\-tut tse\-'a, pe\-yä sä\-sìl\-pey a yo\-ra' 'o\-lìp.} When he saw his opponent, his hope of winning vanished.~\LINK{http://naviteri.org/2013/01/awvea-posti-zisita-amip-first-post-of-the-new-year/}{b}
(\ding{43}~\HL{wA}{\B{wä}}, \HL{tute1}{\B{tu\-te}}) 

% wätx
\hi 
\HT{wAtx}{\B{wätx}} \I{[w\ae{}t']} \emph{vtr.} be bad at. 
\B{Sra\-ke fnan ngal lì'\-fya\-ti le\-Na'\-vi? -- Ke\-he. Slä hu\-fwa oel (tsat) wätx nì\-'aw, tsal\-su\-ngay yaw\-ne lu oer nì\-wotx!} Are you good at Na'vi? -- No, I totally suck at it, but I still love it to death.~\LINK{http://naviteri.org/2010/09/getting-to-know-you-part-3/}{b}
\B{Tsmìm\-ìri wätx fol ti\-nan\-it nì\-ngay.} They're really bad at reading animal tracks.~\LINK{http://naviteri.org/2011/09/\%E2\%80\%9Cby-the-way-what-are-you-reading\%E2\%80\%9D/}{b}
\B{Oe\-ri lu tì\-ftxu\-lì\-'u ngä\-zìk nì\-txan. Wätx nì\-'aw.} Public speaking is very difficult for me. I'm hopelessly bad at it.~\LINK{http://naviteri.org/2012/06/spring-vocabulary-part-3/}{b}
(\ding{214}~\HL{fnan}{\B{fnan}})

% wätum
\hi 
\HT{wAtum}{\B{\cp{wä}tum}} \I{["w\ae.tum]} \emph{or} \I{[.tUm]} (\textsc{rn}: \B{\cp{wä}tùm} \I{["w\ae.tUm]}) \emph{n.} antidote. 
\B{Pe\-yä ay\-lì\-'u\-ri a\-txum\-nga' wä\-tum\-ìl pe\-se\-nget?} Where is the antidote for his poisonous words?~\LINK{https://naviteri.org/2025/04/mevosinga-liu-amip-twenty-new-words/}{b} 
(\ding{43}~\HL{wA}{\B{wä}}, \HL{txum}{\B{txum}}) 

% we'ay
\hi 
\HT{we'ay}{\B{\cp{we}'ay}} \I{["wE.Paj]} \emph{adj.} sour.

% wem
\hi 
\HT{wem}{\B{wem}} \I{[wEm]} \emph{vin.} fight (with\slash alongside, \HL{hu}{\B{hu}}; against, oppose, \HL{wA}{\B{wä}}).   
\B{Ru\-txe wi\-vem nì\-nän.} Please fight less.~\LINK{http://naviteri.org/2012/02/trr-asawnung-lefpom-happy-leap-day/}{b}
\B{Peyä tsa\-tì\-pe\-'un a swey\-lu txo wi\-vem ayoeng Oma\-ti\-ka\-ya\-wä lu fe'.} His decision to fight against the Omaticaya was a bad one.~\LINK{http://naviteri.org/2011/07/txantsana-ultxa-mi-siatll-great-meeting-in-seattle/}{b}
\B{Ay\-oe ke wa\-syem.} We don't intend to fight.~\LINK{http://forum.learnnavi.org/language-updates/ltasygt-with-ke-and-nga/}{ln\slash g} 
\B{Po\-ri wem\-tswo fra\-tsam\-si\-yur ro\-lo\-'a nì\-txan.} His ability to fight greatly impressed all the warriors.~\LINK{http://naviteri.org/2012/03/spring-vocabulary-part-2/}{b}
\B{Oel nge\-y\"a tì\-hawl\-it na\-txu ulte tsa\-w\"a wa\-syem.} I disapprove of your plan and will oppose (fight against) it.~\LINK{http://naviteri.org/2020/02/some-words-for-leap-year-day/}{b}  \\  
\textsc{Derivatives}: \ding{43}
\on{1}.~\HL{wempongu}{\B{wem\-po\-ngu}}, squad, battle party.

% wempongu
\hi 
\HT{wempongu}{\B{\cp{wem}pongu}} \I{["wEm.po.Nu]} \emph{n.} squad, military clan, battle party. 
(\ding{43}~\HL{wem}{\B{wem}}, \HL{pongu}{\B{po\-ngu}})

% wew
\hi 
\HT{wew}{\B{wew}} \I{[wEw]} \emph{adj.} cold. 
\B{Ya\-ri som\-wew\-pe (set)? Ya lu wew.} What's the temperature? It's cold.~\LINK{http://naviteri.org/2012/10/audio-and-video-learning-materials-for-navi-101/}{b}
\B{Nga tsun ni\-väk wew\-a pay\-it.} You can drink cold water.~\LINK{http://naviteri.org/2012/10/audio-and-video-learning-materials-for-navi-101/}{b} 
(\ding{214}~\HL{}{\B{som}}) \\  
\textsc{Derivatives}: \ding{43}
\on{1}.~\HL{tIwew}{\B{tì\-wew}}, coldness, cold. 
\on{2}.~\HL{txawew}{\B{txa\-wew}}, very cold. 
\on{3}.~\HL{somwew}{\B{som\-wew}}, temperature.
\on{4}.~\HL{zIskrrwew}{\B{zì\-skrr\-wew}}, winter.

% weyn
\hi 
\HT{weyn}{\B{weyn}} \I{[wEjn]} \emph{vtr.} draw, illustrate. 
\B{Wo\-leyn Ìstawl yey\-fyat mì hll\-te fte oeyk\-ti\-vìng fra\-po\-ru tì\-hawl\-te\-ri sne\-yä.} Ìstaw drew a line on the ground to explain his plan to everyone.~\LINK{http://naviteri.org/2013/09/on-si-salewfya-shapes-and-directions/}{b}
\B{Pol rel\-it wo\-leyn fa pin\-vul.} He drew a picture with a crayon.~\LINK{http://naviteri.org/2018/03/100a-liu-amip-64-new-words-part-1/\#comment-28355}{b}

% weopx
\hi 
\HT{weopx}{\B{we\cp{opx}}} \I{[wE."{}op']} \emph{n.} wave (of water). 
\B{Krr\-a hu\-fwe tul nì\-win, tsun fko tsi\-ve\-'a ay\-we\-opx\-it a sìn yo pay\-ä.} When there is strong wind, you can see waves on the water.~\LINK{http://naviteri.org/2017/09/zisikrr-amip-ayliu-amip-new-words-for-the-new-season/}{b}  \\  
\textsc{Derivatives}: \ding{43}
\on{1}.~\HL{leweopx}{\B{le\-we\-opx}}, wave\hyp{}like.

% win / si / säpi
\hi 
\HT{win}{\B{win}} \I{[win]} \textsc{i}. \emph{adj.} fast, quick, rapid.
\B{Lu nga win sì txur | Lu nga txan\-tslu\-sam.} You are fast and strong | You are wise.~\LINK{http://naviteri.org/2013/08/taronway-the-hunt-song/}{b} 
\B{Po lu win fra\-to.} She is the fastest.~\LINK{http://forum.learnnavi.org/language-updates/confirmations-\%28comparisons-gerund\%29/}{ln}
\B{Oel tso\-le\-'a fwam\-pop\-it a\-win fra\-to.} I saw the fastest Tapirus.~\LINK{http://forum.learnnavi.org/language-updates/confirmations-\%28comparisons-gerund\%29/}{ln} 
\B{Win\-a sä\-wä\-sul\-tsyìp a\-hì\-'i.} A quick little contest.~\LINK{http://naviteri.org/2012/10/wina-sawasultsyip-ahii-a-quick-little-contest/}{b} 
\B{win\-a u\-van si}, play a quick game. \LINK{http://wiki.learnnavi.org/Canon/2010/UltxaAyharyu\%C3\%A4\#Refinements_of_si-construction_verbs}{wiki} \\
\textsc{ii}.a \B{win si} \I{[win si]} \emph{vin.} rush s.t., make s.t. fast. \\
\textsc{ii}.b \B{win sä\cp{pi}} \I{[win s\ae{}."pi]} \emph{vin.} hurry oneself. 
\B{Ke ze\-ne win sä\-pi\-vi. 'Ivong nìk\-'ong.} Take your time; don't rush. Slow is fine.~\LINK{http://naviteri.org/2010/09/quick-follow-up/}{b} \\  
\textsc{Derivatives}: \ding{43} \on{1}.~\HL{nIwin}{\B{nì\-win}}, fast, quickly, rapidly. 
\on{2}.~\HL{winzaw}{\B{win\-zaw}}, arrow deer.

% winzaw
\hi 
\HT{winzaw}{\B{\cp{win}zaw}} \I{["win.zaw]} \emph{n., fauna} arrow deer. 
(\ding{43}~\HL{win}{\B{win}}, \HL{swizaw}{\B{swi\-zaw}}) 

% wion
\hi 
\HT{wion}{\B{\cp{wi}on}} \I{["wi.on]} \emph{n.} reef. 
\B{Lì'\-fya Wi\-on\-ä}, Reef Na'vi, \emph{a dialect of the Na'vi language family} [lit.: `language of the reef'].~\LINK{http://naviteri.org/2023/01/reef-navi-part-1-phonetics-and-phonology/}{b}

% wip
\hi 
\HT{wip}{\B{wip}} \I{[wip\textcorner]} \emph{adj.} salty.

% wiya
\hi 
\HT{wiya}{\B{\cp{wi}ya}} \I{["wi.ja]} \emph{intj. expression of warning or frustration}.

% wìntxu
\hi 
\HT{wIntxu}{\B{wìn\cp{txu}}} \I{[wIn."t'u]} \emph{vtr.} show. 
\B{Fay\-sä\-re\-ngop\-it a\-vä' oe\-ru rä\-'ä wìn\-txu nì\-mun, ru\-txe. Ke su\-nu oer keng nì\-'it.} Please don't show me these ugly designs again. I don't like them one bit.~\LINK{http://naviteri.org/2013/01/awvea-posti-zisita-amip-first-post-of-the-new-year/}{b}
\B{Nar\-mew oe fo\-ru na'\-rìng\-ä tì\-lor\-it wi\-vìn\-txu, slä ke tsä\-ngun fo tsli\-vam.} I wanted to show them the beauty of the forest, but sadly, they weren't able to understand.~\LINK{http://naviteri.org/2014/11/twenty-before-the-holidays/}{b}
\B{Oel lum\-pe tsat so\-lung mì upxa\-re? Fte wi\-vìn\-txu fu\-ta le\-sar lu fì\-lì\-'u nì\-txan.} Why have I added it to the message? To show that this word is very useful.~\LINK{http://naviteri.org/2010/07/vocabulary-update/comment-page-1/\#comment-95}{b} 
\B{Tsa\-ta txe'\-lan\-it nge\-yä tok wo\-lìn\-txu ngal ay\-oe\-ru nì\-ngay.} You truly showed us where your heart is.~\LINK{http://masempul.org/2010/07/\%E2\%80\%99ite-polaheiem/}{sempul} 
\B{Ney\-ti\-ril Tsyey\-kur wa\-mìn\-txu Oma\-ti\-ka\-ya\-ä ve\-fyat tì\-tu\-sa\-ron\-ä.} Ney\-ti\-ri showed Jake the Omatikaya's approach to hunting.~\LINK{http://naviteri.org/2012/10/mipa-vospxi-mipa-ayliu-new-words-for-the-new-month/}{b} \\  
\textsc{Derivatives}: \ding{43}
\on{1}.~\HL{sAwIntxu}{\B{sä\-wìn\-txu}}, show, showing.  
\on{2}.~\HL{muwIntxu}{\B{mu\-wìn\-txu}}, introduce, present.  
\on{3}.~\HL{fyawIntxu}{\B{fya\-wìn\-txu}}, guide. 
\on{4}.~\HL{wIngay}{\B{wì\-ngay}}, prove.

% wìngay
\hi 
\HT{wIngay}{\B{wì\cp{ngay}}} \I{[w$\cdot$I."Naj]} \emph{vtr., inf.}\on{11} prove.
\B{Fa fwa tsyìl kxem\-yot a\-kxayl fra\-to, pol ay\-oer wo\-lì\-ngay fu\-ta tsyìl\-tswo tsan\-'o\-lul.} By scaling the highest wall, he proved to us that his climbing ability had improved.~\LINK{http://naviteri.org/2015/04/some-new-words-for-may-day/}{b}
(\ding{43}~\HL{wIntxu}{\B{wìn\-txu}}, \HL{ngay}{\B{ngay}}) \\  
\textsc{Derivatives}: \ding{43}
\on{1}.~\HL{tIwIngay}{\B{tì\-wì\-ngay}}, proving (abstract). 
\on{2}.~\HL{sAwIngay}{\B{sä\-wì\-ngay}}, proof.

% wo
\hi 
\HT{wo}{\B{wo}} \I{[wo]} \emph{vtr.} reach for.   
\B{Ngal new a tsa\-'ut rä\-'ä wi\-vo, ma 'evi. Vi\-vin.} Don't reach for what you want, child. Ask for it.~\LINK{http://naviteri.org/2012/01/mipa-zisit-ayliu-amip-new-words-for-the-new-year/}{b} \\  
\textsc{Derivatives}: \ding{43}
\on{1}.~\HL{kllwo}{\B{kll\-wo}}, alight, land, approach for land. 
\on{2}.~\HL{yawo}{\B{ya\-wo}}, take off, launch. 

% wok
\hi 
\HT{wok}{\B{wok}} \I{[wok\textcorner]} \emph{adj.} loud.
\B{Lu Saw\-tu\-te wok sì pì\-saw nì\-wotx, na prr\-nen.} The skypeople are all loud and clumsy, like a baby.~\LINK{http://naviteri.org/2017/09/zisikrr-amip-ayliu-amip-new-words-for-the-new-season/}{b} 
\B{Zo\-la\-'u ftu kä\-pxì num\-tseng\-vi\-yä sä\-mewn a\-wok.} From the back of the classroom came a loud sigh.~\LINK{https://naviteri.org/2025/04/mevosinga-liu-amip-twenty-new-words/}{b} 
(\ding{214}~\HL{'ango}{\B{'ango}}) \\  
\textsc{Derivatives}: \ding{43}
\on{1}.~\HL{nIwok}{\B{nì\-wok}}, loudly. 
\on{2}.~\HL{wokau}{\B{wok\-au}}, pendulum drum.

% wokau
\hi 
\HT{wokau}{\B{wok\cp{a}u}} \I{[wok."a.u]} \emph{n.} pendulum drum. 
(\ding{43}~\HL{wok}{\B{wok}}, \HL{au}{\B{au}})

% wotx
\hi 
\HT{wotx}{\B{wotx}} \I{[wot']} \emph{n.} totality, whole. 
(\emph{a. idiom}) \B{ka wotx}, generally, for the most part. 
\B{Hu\-fwa ro\-lun oel 'a\-'aw\-a kxey\-ey\-ti, fì\-tì\-kang\-kem\-vi lu txan\-tsan ka wotx.} Although I found a few errors, this piece of work is generally excellent.~\LINK{http://naviteri.org/2012/07/fivospxiya-aylifyavi-amip-this-months-new-expressions-pt-2/}{b} \\  
\textsc{Derivatives}: \ding{43}
\on{1}.~\HL{nIwotx}{\B{nì\-wotx}}, all (of), completely.
\on{2}.~\HL{tengwotx}{\B{teng\-wotx}}, uniform, non\hyp{}diverse.

% wou
\hi 
\HT{wou}{\B{\cp{wo}u}} \I{["wo.u]} \emph{vin., slang} be amazing, be fascinating (\emph{slang}).   
\B{Wou Ney\-ti\-ri!} Neytiri is awesome!~\LINK{http://forum.learnnavi.org/language-updates/txelanit-hivawl/}{ln}
\B{Oe\-ru wo\-eiu lì'\-fya le\-Na'\-vi.} I'm just fascinated by the Na'vi language.~\LINK{http://forum.learnnavi.org/language-updates/txelanit-hivawl/}{ln}
\B{Mì na'\-rìng Tsyeyk\-ìl lo\-re\-yu\-ti 'o\-lam\-pi a tsa\-swaw wou oer nì\-wotx!} That instant when Jake touched the \emph{Helicoradiam Spirale} in the forest I was totally amazed.~\LINK{http://forum.learnnavi.org/language-updates/txelanit-hivawl/}{ln}
\B{Por war\-mou fra\-fne\-i\-o\-ang nì\-tut, slä nì\-pxi pxe\-syì\-rä\-fì.} He just couldn't get over all the animals, but was especially taken with the three giraffe.~\LINK{http://forum.learnnavi.org/language-updates/txelanit-hivawl/}{ln}

% wrrkä
\hi 
\HT{wrrkA}{\B{wrr\cp{kä}}} \I{[w\s{r}."k$\cdot$\ae]} \emph{vin., inf.}\on{22} go out, go outside. 
\B{Ke tsun wrr\-ki\-vä fte tì\-kang\-kem si\-vi.} (We) cannot go out to work.~\LINK{http://naviteri.org/2020/03/teri-tifkeytok-lefkrr-about-the-present-situation/}{b} 
\B{A\-ha'\-ri wrr\-ko\-lä fte li\-mang si\-vu\-nu por.} Aha'ri is out enjoying the moot.~\LINK{https://forum.learnnavi.org/language-updates/limangmoot-from-afop-dialogue/}{ln\slash fop}
(\ding{43}~\HL{wrrpa}{\B{wrr\-pa}}, \HL{kA}{\B{kä}}; \ding{214}~\HL{wrrza'u}{\B{wrr\-za\-'u}})

% wrrpa
\hi 
\HT{wrrpa}{\B{\cp{wrr}pa}} \I{["w\s{r}.pa]} \emph{n., adv.} outside. 
\B{Txìng ay\-ioang\-it! Hum ne wrr\-pa! Tul!} Leave the animals! Get outside! Run!~\LINK{http://forum.learnnavi.org/language-updates/txing-and-hum/}{ln\slash a\textsuperscript{c}} 
\B{Ke za\-syup lì\-'O\-na ne kxu\-tu a mì\-fa fu a wrr\-pa.} The l'Ona will not perish to the enemy within or the enemy without.~\LINK{http://forum.learnnavi.org/language-updates/ltasygt-with-ke-and-nga/}{ln\slash g}
(\ding{214}~\HL{mIfa}{\B{mì\-fa}}, \ding{43}~\HL{pa'o}{\B{pa\-'o}}) \\  
\textsc{Derivatives}: \ding{43}
\on{1}.~\HL{tawrrpa}{\B{ta\-wrr\-pa}}, external.

% wrrza'u
\hi 
\HT{wrrza'u}{\B{wrr\cp{za}'u}} \I{[w\s{r}."z$\cdot$a.P$\cdot$u]} \emph{vin., inf.}\on{23} come out, emerge.
\B{Txep\-ram pxor a krr, txan\-a txe\-kxum\-pay wrr\-za\-'u.} When a volcano erupts, a lot of lava comes out.~\LINK{http://naviteri.org/2015/03/kap-si-ayunil-saylahe-threats-dreams-and-other-things/}{b} 
(\ding{43}~\HL{za'u}{\B{za\-'u}}; \ding{214}~\HL{wrrkA}{\B{wrr\-kä}})

% wrrzärìp
\hi 
\HT{wrrzArIp}{\B{wrr\cp{zä}rìp}} \I{[w\s{r}."z$\cdot$\ae.R$\cdot$Ip\textcorner]} \emph{vtr., inf.}\on{23} (\emph{a. physical}) pull out, extract. 
\B{Pol tstal\-it wrr\-zo\-lä\-rìp tstal\-se\-na\-ftu.} He pulled the knife out of its sheath.~\LINK{http://naviteri.org/2021/03/mipa-ayliu-si-aylifyavi-new-words-and-expressions/}{b} 
(\emph{b. fig.}) (1.) \B{txe'\-lan\-ti wrr\-zä\-rìp} \emph{vtr.} to greatly move emotionally.
\B{Oe\-ri pe\-yä ay\-lì\-'ul txe'\-lan\-ti wrr\-zo\-lä\-rìp.} Her words moved me greatly.~\LINK{http://naviteri.org/2021/03/mipa-ayliu-si-aylifyavi-new-words-and-expressions/}{b}
(2.) \B{tì\-pe'\-ngayt wrr\-zä\-rìp} \emph{vtr.} infer.
\B{Ngey ay\-lì\-'u\-ftu wrr\-zä\-rìp oel tì\-pe'\-ngayt a lu ngar yew\-la.} From your words, I infer that you're disappointed.~\LINK{http://naviteri.org/2021/03/mipa-ayliu-si-aylifyavi-new-words-and-expressions/}{b} 
(3.) \B{tì\-pe'\-ngayt wrr\-zey\-kä\-rìp} \emph{vtr.} imply.
\B{Ngey ay\-lì\-'ul wrr\-zey\-kä\-rìp tì\-pe'\-ngayt a lu ngar yew\-la.} Your words imply that you're disappointed.~\LINK{http://naviteri.org/2021/03/mipa-ayliu-si-aylifyavi-new-words-and-expressions/}{b}
(\ding{43}~\HL{wrrpa}{\B{wrr\-pa}}, \HL{za'ArIp}{\B{za\-'ä\-rìp}})

% wum
\hi 
\HT{wum}{\B{wum}} \I{[wum]} \emph{or} \I{[wUm]} (\textsc{rn}: \B{wùm} \I{[wUm]}) \emph{adv.} approximately, roughly.
\B{Oe\-ri so\-la\-lew wum zì\-sìt °a14 a krr, fo\-lrr\-fen spo\-not alo a\-'aw\-ve.} When I was about 12 years old I visited an island for the first time.~\LINK{http://naviteri.org/2012/10/mipa-vospxi-mipa-ayliu-new-words-for-the-new-month/}{b}

% wur
\hi 
\HT{wur}{\B{wur}} \I{[wuR]} \emph{or} \I{[wUR]} \emph{adj.} cool, a bit chilly. 
\B{Ya\-ri som\-wew\-pe (set)? Ya lu wur.} What's the temperature? It's chilly.~\LINK{http://naviteri.org/2012/10/audio-and-video-learning-materials-for-navi-101/}{b}
\B{Ay\-nga\-ri teng\-krr ya wur sle\-ru nì\-'ul, sìl\-pey oe, li\-vu hel\-ku sang ul\-te te'\-lan le\-fpom.} As for you all, while it's getting cooler, I hope your homes may be warm and your hearts happy.~\LINK{http://naviteri.org/2013/09/on-si-salewfya-shapes-and-directions/}{b}
\B{Pey\-ral\-ìl zet wur\-a wu\-tsot a\-'aw\-nem pxel sngel.} Peyral treats a cool cooked meal like garbage.~\LINK{http://naviteri.org/2012/10/mipa-vospxi-mipa-ayliu-new-words-for-the-new-month/}{b}
(\ding{214}~\HL{sang}{\B{sang}})

% wutso
\hi 
\HT{wutso}{\B{\cp{wu}tso}} \I{["wu.\t{ts}o]} \emph{n.} dinner, served meal.
\B{Nì\-fra\-krr fol 'o\-lem a wu\-tso ftxì\-vä' lu nì\-ngay.} As always, the dinner they cooked tasted really terrible.~\LINK{http://naviteri.org/2011/05/weather-part-2-and-a-bit-more-2/}{b}
\B{Nì\-'aw\-ve oeng yo\-lom wu\-tsot; maw\-krr uvan si.} First we had dinner; afterwards we played.~\LINK{http://naviteri.org/2011/08/new-vocabulary-clothing/comment-page-1/\#comment-917}{b} 
\B{Pey\-ral\-ìl zet wur\-a wu\-tsot a\-'aw\-nem pxel sngel.} Peyral treats a cool cooked meal like garbage.~\LINK{http://naviteri.org/2012/10/mipa-vospxi-mipa-ayliu-new-words-for-the-new-month/}{b}

% wuwuk
\hi 
\HT{wuwuk}{\B{\cp{wu}wuk}} \I{["wu.wuk\textcorner]} \emph{or} \I{[.wUk\textcorner]} (\textsc{rn}: \B{\cp{wù}wùk} \I{["wU.wUk\textcorner]}) \emph{n., fauna} lizard, \emph{any of a variety of lizard\hyp{}like creatures}. 
\\
\end{multicols}




% ===== Section Y =====

 \pdfbookmark[0]{Y}{Y}
\begin{center}
\section*{Y \I{[j]} (\emph{Yä})}
\end{center}

\begin{multicols}{2}

% ya
\hi 
\HT{ya}{\B{ya}} \I{[ja]} \emph{n.} air.
\B{Ya\-ri som\-wew\-pe (set)? -- Ya lu txa\-wew\slash wew\slash wur\slash tsya\-fe\slash sang\slash som\slash txa\-som.} What's the temperature? -- It's very cold\slash cold\slash cool\slash mild\slash warm\slash hot\slash very hot.~\LINK{http://naviteri.org/2011/05/weather-part-2-and-a-bit-more-2/}{b}
\B{Ya\-mì tsun fko tsi\-ve\-'a lo\-ran\-it re\-nu\-ä kil\-van\-ä slä kll\-te\-sìn wä\-pan.} From the air you can see the grace of the river's form but from the ground it's hidden.~\LINK{http://naviteri.org/2012/10/mipa-vospxi-mipa-ayliu-new-words-for-the-new-month/}{b} \\  
\textsc{Derivatives}: \ding{43}
\on{1}.~\HL{yafkeyk}{\B{ya\-fkeyk}}, weather.  
\on{2}.~\HL{yapay}{\B{ya\-pay}}, mist, fog, steam. 
\on{3}.~\HL{yawo}{\B{ya\-wo}}, take off, launch.  
\on{4}.~\HL{yayo}{\B{ya\-yo}}, bird. 
\on{5}.~\HL{yrrap}{\B{yrr\-ap}}, storm. 
\on{6}.~\HL{txImya}{\B{txìm\-ya}}, fart.

% -ya
\hi 
\HT{ya2}{\B{-ya}} \I{[.ja]} \emph{suf. vocative for collective nouns}.
\B{Ma\-wey, Na'\-vi\-ya, ma\-wey!} Be calm, people, be calm.~\LINK{https://wiki.learnnavi.org/Na\%27vi_from_Avatar_Movie\#Neytiri_talking_to_Jake}{a\slash wiki}
(\ding{43}~\HL{ma}{\B{ma}})

% yafkeyk
\hi 
\HT{yafkeyk}{\B{\cp{ya}fkeyk}} \I{["ja.fkEjk\textcorner]} \emph{n.} weather. 
\B{Ya\-fkeyk (za\-'u) fya\-pe?} How's the weather?~\LINK{http://naviteri.org/2011/04/yafkeykiri-plltxe-frapo-everyone-talks-about-the-weather/}{b}
\B{Ya\-fkeyk fya\-pe set? -- Taw lu pi\-ak.} What's the weather like now? -- The sky is clear.~\LINK{http://naviteri.org/2012/10/audio-and-video-learning-materials-for-navi-101/}{b}
\B{Tsun fko ay\-on\-ti fì\-wopx\-ä ni\-vìn fte ya\-fkeyk\-it sre\-si\-ve\-'a.} The shapes of clouds can be used to predict the weather.~\LINK{http://naviteri.org/2013/09/on-si-salewfya-shapes-and-directions/}{b}
\B{Ya\-fkeyk\-ìri päng\-kxo fra\-po.} Everyone talks about the weather.~\LINK{http://naviteri.org/2011/04/yafkeykiri-plltxe-frapo-everyone-talks-about-the-weather/}{b}
(\ding{43}~\HL{ya}{\B{ya}}, \HL{tIfkeytok}{\B{tì\-fkey\-tok}})

% yak / si
\hi 
\HT{yak}{\B{yak}} \I{[jak\textcorner]} \textsc{i}. \emph{n.} fork, branch, point of divergence.
\B{Hay\-a yak\-ro fti\-vang. Sa\-lew rä\-'ä.} Stop at the next fork. Do not proceed further.~\LINK{http://naviteri.org/2013/09/on-si-salewfya-shapes-and-directions/}{b} \\
\textsc{ii}. \B{yak si} \I{[jak\textcorner si]} \emph{vin.} diverge, change direction, go astray.~\LINK{http://naviteri.org/2013/09/on-si-salewfya-shapes-and-directions/}{b}
\B{Ne 'o\-ra\-tsyìp po\-lä\-hem, yak si nì\-ftär.} When you arrive at the pond, turn to the left.~\LINK{http://naviteri.org/2013/09/on-si-salewfya-shapes-and-directions/}{b}
\B{Aw\-nga ze\-ne vi\-var ìlä fì\-sa\-lew\-fya. Zen\-ke yak si\-vi.} We must continue in this direction. We must not go astray.~\LINK{http://naviteri.org/2013/09/on-si-salewfya-shapes-and-directions/}{b}

% yaney
\hi 
\HT{yaney}{\B{ya\cp{ney}}} \I{[ja."nEj]} \emph{n.} canoe. \\  
\textsc{Derivatives}: \ding{43}
\on{1}.~\HL{spulyaney}{\B{spul\-ya\-ney}}, canoe paddle.

% yapay
\hi 
\HT{yapay}{\B{\cp{ya}pay}} \I{["ja.paj]} \emph{n.} mist, fog, steam.
\B{neyn na ya\-pay}, the light, nondescript color of mist or fog.~\LINK{http://naviteri.org/2010/08/mipa-ayopin-mipa-ayliu-new-colors-new-words/}{b}
(\ding{43}~\HL{ya}{\B{ya}}, \HL{pay}{\B{pay}}) \\  
\textsc{Derivatives}: \ding{43}
\on{1}.~\HL{leyapay}{\B{le\-ya\-pay}}, foggy, misty.

% yawne
\hi 
\HT{yawne}{\B{\cp{yaw}ne}} \I{["jaw.nE]} \textsc{i}. \emph{adj.} (\emph{a.}) beloved.  
\B{Ma oe\-yä ey\-lan a\-yaw\-ne a lì'\-fya\-olo'\-mì.} My beloved friends in the language community.~\LINK{http://forum.learnnavi.org/news-announcements/happy-birthday-to-karyu-pawl/msg556899/\#msg556899}{ln} 
(\emph{b. to express}) to love s.o.\slash s.t. \B{Nga yaw\-ne lu oer.} I love you.~\LINK{http://wiki.learnnavi.org/index.php?title=Corpus\#Lemondrop.com}{ld\slash wiki}
\B{Ay\-lì'\-fya yaw\-ne le\-ru oer ta\-krr\-a 'e\-veng la\-mu.} I've loved languages since I was a child.~\LINK{http://naviteri.org/2011/01/new-year-new-vocabulary/}{b} 
\B{Po yaw\-ne lu por.} He\slash she loves him\slash her.~\LINK{http://naviteri.org/2011/12/one-more-for-2011/}{b}
\B{Po yaw\-ne lu snor.} He loves himself.~\LINK{http://naviteri.org/2011/12/one-more-for-2011/}{b}
\B{Me\-fo yaw\-ne lu (snor) fì\-tsap.} They (= those two) love each other.~\LINK{http://naviteri.org/2011/12/one-more-for-2011/}{b} 
(\emph{c. idiom}) \B{Tse\-'a tìk yaw\-ne.} `Love at first sight.'~\LINK{http://naviteri.org/2021/12/zolau-niprrte-ma-3746-welcome-2022/}{b} \\
\textsc{ii}. \B{\cp{yaw}ne slu} \I{["jaw.nE slu]} \emph{vin.} become beloved, fall in love.
\B{Lu poe se\-vin nì\-ftxan kuma yaw\-ne slo\-lu oer.} She was so beautiful that I fell in love with her.~\LINK{http://naviteri.org/2012/06/spring-vocabulary-part-3/}{b}
\B{Ma mun\-txa\-tu, oeng yaw\-ne lu oe\-nga\-ru fì\-tsap, ke\-fyak?} We love each other, don't we, my spouse?~\LINK{http://naviteri.org/2011/12/one-more-for-2011/}{b} \\  
\textsc{Derivatives}: \ding{43}
\on{1}.~\HL{tIyawn}{\B{tì\-yawn}}, love. 
\on{2}.~\HL{yawnetu}{\B{yaw\-ne\-tu\slash yawn\-tu}}, loved one. 
\on{3}.~\HL{yawnyewla}{\B{yawn\-yew\-la}}, broken heart, broken\hyp{}heartedness. 

% yawnetu / yawntu
\hi 
\HT{yawnetu}{\B{\cp{yaw}netu}} \I{["jaw.nE.tu]} \emph{or} \B{\cp{yawn}tu} \I{["jawn.tu]}  \emph{n.} loved one. 
\B{Hu\-fwa me\-fo le\-ru mun\-txa\-tu txan\-krr, mi lu mun\-snar ho\-na tì\-flrr a na pum me\-yaw\-ne\-tu\-ä a\-mip nì\-wotx.} Although the two of them have been mates a long time, they still have all the adorable tenderness of new sweethearts.~\LINK{http://naviteri.org/2013/02/vospxi-ayol-posti-apup-short-post-for-a-short-month/}{b}
(\ding{43}~\HL{yawne}{\B{yaw\-ne}}, \HL{tute}{\B{tu\-te}}) \\  
\textsc{Derivatives}: \ding{43}
\on{1}.~\HL{yawntutsyIp}{\B{yawn\-tu\-tsyìp}}, darling.

% yawntutsyìp
\hi 
\HT{yawntutsyIp}{\B{\cp{yawn}tutsyìp}} \I{["jawn.tu.\t{ts}jIp\textcorner]} \emph{n.} darling, little loved one (\emph{often reflects parental or familial love}). 
\B{Sem\-pu\-ti rä\-'ä srätx, ma yawn\-tu\-tsyìp. Tì\-kang\-kem se\-ri.} Don't bother daddy, little one. He's working.~\LINK{http://naviteri.org/2017/09/zisikrr-amip-ayliu-amip-new-words-for-the-new-season/}{b}  
(\ding{43}~\HL{yawntu}{\B{yawn\-tu}})

% yawnyewla
\hi 
\HT{yawnyewla}{\B{yawn\cp{yew}la}} \I{[jawn."jEw.la]} \emph{n.} broken heart; broken\hyp{}heartedness.   
\B{Lu Tse\-nur yawn\-yew\-la a lam fwa Va'\-rul pot txì\-yìng.} Tse\-nu is broken hearted that Va'\-ru appears to be about to dump her.~\LINK{http://naviteri.org/2012/10/mipa-vospxi-mipa-ayliu-new-words-for-the-new-month/}{b}
(\ding{43}~\HL{yawne}{\B{yaw\-ne}}, \HL{yewla}{\B{yew\-la}})

% yawr
\hi 
\HT{yawr}{\B{yawr}} \I{[jawR]} \emph{n.} feather. \\  
\textsc{Derivatives}: \ding{43}
\on{1}.~\HL{yawrwll}{\B{yawr\-wll}}, feather blade.

% yawrwll
\hi 
\HT{yawrwll}{\B{\cp{yawr}wll}} \I{["jawR.w\s{l}]} \emph{n.} \ding{96} feather blade [lit.: `feather plant']. 
(\ding{43}~\HL{yawr}{\B{yawr}}, \HL{'ewll}{\B{'e\-wll}}) 

% yayl
\hi 
\HT{yayl}{\B{yayl}} \I{[jajl]} \emph{n.} nonsense, gibberish. 
\B{Fay\-lì\-'u a pol\-txe nga lu yayl nì\-'aw.} What you said is nothing but nonsense.~\LINK{https://naviteri.org/2024/06/ayhapxi-tokxa-si-ayliu-alahe-parts-of-the-body-and-more/}{b}
(\emph{a. vulg.}) \B{Yayl!} `Bullshit!'  \\  
\textsc{Derivatives}: \ding{43}
\on{1}.~\HL{leyayl}{\B{le\-yayl}}, nonsensical.

% yawo
\hi 
\HT{yawo}{\B{ya\cp{wo}}} \I{[ja."w$\cdot$o]} \emph{vin., inf.}\on{22} take off, launch.   
\B{Fwa ya\-wo ftu kll\-te to fwa tsway\-on ftu 'awkx lu ngä\-zìk.} Taking off from the ground is harder than flying off a cliff.~\LINK{http://naviteri.org/2012/01/mipa-zisit-ayliu-amip-new-words-for-the-new-year/}{b}
\B{'Uol ik\-ran\-it txo\-pu sley\-ko\-la\-tsu, ta\-lun\-a po tsìk ya\-wo.} Something must have frightened the banshee, because it suddenly took to the air.~\LINK{http://naviteri.org/2012/01/mipa-zisit-ayliu-amip-new-words-for-the-new-year/}{b}
\B{Ik\-ran ya\-wo\-lo ftu txew 'awkx\-ä.} The banshee took to the air from the edge of a cliff.~\LINK{http://naviteri.org/2012/03/spring-vocabulary-part-1/}{b} 
(\ding{214}~\HL{kllwo}{\B{kll\-wo}}, \ding{43}~\HL{ya}{\B{ya}}, \HL{wo}{\B{wo}})

% yayayr 
\hi 
\HT{yayayr}{\B{ya\cp{yayr}}} \I{[ja."jajR]} \emph{n.} (\emph{a.}) confusion.   
\B{Fì\-ya\-yayr\-ìri oe\-ru txoa li\-vu.} Forgive this confusion.~\LINK{http://naviteri.org/2013/01/awvea-posti-zisita-amip-first-post-of-the-new-year/comment-page-1/\#comment-1910}{b}
(\emph{b. with dative and} \B{lu}) to be confused. \B{Saw\-tu\-te\-yä hem\-ìri lu aw\-nga\-ru ya\-yayr.} The Skypeople's actions confuse us.~\LINK{http://naviteri.org/2011/01/new-year-new-vocabulary/}{b} \\  
\textsc{Derivatives}: \ding{43}
\on{1}.~\HL{yayayrnga'}{\B{ya\-yayr\-nga'}}, confusing. 
\on{2}.~\HL{yayayrtsim}{\B{ya\-yayr\-tsim}}, s.t. confusing.

% yayayrnga'
\hi 
\HT{yayayrnga'}{\B{ya\cp{yayr}nga'}} \I{[ja."jajR.NaP]} \emph{or coll.} \I{[ja."jaj.NaP]} (\emph{in rapid speech}) \emph{adj.} confusing.
\B{Tsa\-tì\-oeyk\-tìng a\-ya\-yayr\-nga' srung ke so\-li oer fte tsli\-vam teyng\-ta kem\-pe ze\-ne si\-vi.} That confusing explanation didn't help me understand what I have to do.~\LINK{http://naviteri.org/2013/01/awvea-posti-zisita-amip-first-post-of-the-new-year/}{b} 
(\ding{43}~\HL{yayayr}{\B{ya\-yayr}}, \HL{-nga'}{\B{-nga'}})

% yayayrtsim
\hi 
\HT{yayayrtsim}{\B{ya\cp{yayr}tsim}} \I{[ja."jajR.\t{ts}im]} \emph{n.} s.t. confusing, source of confusion.
(\ding{43}~\HL{yayayr}{\B{ya\-yayr}}, \HL{tsim}{\B{tsim}})

% yaymak
\hi 
\HT{yaymak}{\B{\cp{yay}mak}} \I{["jaj.mak\textcorner]} \emph{adj.} foolish, ignorant.
\B{Ki\-fkey\-ri fì\-tu\-tan\-ìl a\-yay\-mak ke tslam stum ke\-'ut.} This foolish man understands almost nothing of the world.~\LINK{http://naviteri.org/2016/12/tengkrr-perahem-zisit-amip-as-the-new-year-arrives/}{b}

% yayo 
\hi 
\HT{yayo}{\B{\cp{ya}yo}} \I{["ja.jo]} \emph{n., fauna} bird. 
\B{Fì\-fne\-ya\-yol tsrul\-it txu\-la mì song\-ropx ut\-ral\-ä fte ay\-li\-nit hi\-vaw\-nu wä sar\-ni\-o\-ang. Fì\-tì\-kan\-ìri ropx ke ha'.} This kind of bird builds its nest in a tree cavity so as to protect its young from predators. For this purpose, a hole going right through the tree trunk isn't suitable.~\LINK{http://naviteri.org/2013/04/wheres-the-bathroom-and-other-useful-things/}{b} 
(\ding{43}~\HL{ya}{\B{ya}}, \HL{ioang}{\B{ioang}}) \\  
\textsc{Derivatives}: \ding{43}
\on{1}.~\HL{yayotsrul}{\B{ya\-yo\-tsrul}}, bird's nest.

% yayotsrul
\hi 
\HT{yayotsrul}{\B{\cp{ya}yotsrul}} \I{["ja.jo.\t{ts}Rul]} \emph{or} \I{[.\t{ts}RUl]} \emph{n.} bird's nest. 
\B{Tsa\-ya\-yo\-tsrul\-it sweyn, ma 'i\-tan.} Don't disturb that bird's nest, son. ~\LINK{http://naviteri.org/2021/09/aawa-ayliu-amip-a-few-new-words/}{b} 
(\ding{43}~\HL{yayo}{\B{ya\-yo}}, \HL{tsrul}{\B{tsrul}})

% -yä
\hi 
\HT{-yA}{\B{-yä}} \I{[.j\ae{}]} \emph{suf. to mark the genitive case of a n. ending in a vowel other than u and o}. 
(\ding{43}~\HL{-A}{\B{-ä}})

% yäkx
\hi 
\HT{yAkx}{\B{yäkx}} \I{[j\ae{}k']} \emph{vtr.} not notice; ignore, snub, overlook intentionally.
\B{Sra\-ke fo hang\-ham ta\-lu\-na nga snaytx? Fo\-ti yäkx.} They're laughing because you lost? Ignore them.~\LINK{http://naviteri.org/2015/03/kap-si-ayunil-saylahe-threats-dreams-and-other-things/}{b}
\B{Tsam\-si\-yu ze\-ne tsi\-vun yi\-väkx sne\-yä tì\-sraw\-it.} Warriors must be able to ignore their own pain.~\LINK{http://naviteri.org/2015/03/kap-si-ayunil-saylahe-threats-dreams-and-other-things/}{b}
(\ding{214}~\HL{tseri}{\B{tse\-ri}}) \\  
\textsc{Derivatives}: \ding{43}
\on{1}.~\HL{tIyAkx}{\B{tì\-yäkx}}, lack of notice; snubbing. 
\on{2}.~\HL{sAyAkx}{\B{sä\-yäkx}}, snub.

% yän
\hi 
\HT{yAn}{\B{yän}} \I{[j\ae{}n]} \emph{vtr.} fasten, tie down. 

% ye
\hi 
\HT{ye}{\B{ye}} \I{[jE]} \emph{adj.} satisfied, content; satiated, ``full'' (\emph{used with} \HL{'efu}{\B{'efu}}).   
\B{Tsa\-ri\-a fkol po\-le\-'un fu\-ta Lo\-ak slu ta\-ron\-yu, sem\-pul 'efu ye.} Father is content that it's been decided Loak will be a hunter.~\LINK{http://naviteri.org/2011/04/\%E2\%80\%99a\%E2\%80\%99awa-li\%E2\%80\%99fyavi-amip\%E2\%80\%94a-few-new-expressions/}{b}
\B{Nge\-yä tì\-kang\-kem\-ìri 'e\-fe\-iu oe ye nì\-txan. Sey\-so\-nìl\-tsan!} I'm very satisfied with your work. Well done!~\LINK{http://naviteri.org/2011/04/\%E2\%80\%99a\%E2\%80\%99awa-li\%E2\%80\%99fyavi-amip\%E2\%80\%94a-few-new-expressions/}{b}
\B{Sra\-ke ye\-hakx? -- Ye. Tsun ti\-vam.} Did you get enough to eat? -- Yes. That'll be enough.~\LINK{http://naviteri.org/2011/04/\%E2\%80\%99a\%E2\%80\%99awa-li\%E2\%80\%99fyavi-amip\%E2\%80\%94a-few-new-expressions/}{b} \\  
\textsc{Derivatives}: \ding{43}
\on{1}.~\HL{yehakx}{\B{ye\-hakx}}, satisfied from hunger by food. 
\on{2}.~\HL{yevAng}{\B{ye\-väng}}, satisfied from thirst by drink.

% ye'krr
\hi 
\HT{ye'krr}{\B{\cp{ye'}krr}} \I{["jEP.k\s{r}]} \emph{adv.} early.
\B{Kar\-yul fngo\-lo' fu\-ta ay\-nu\-me\-yu pi\-va\-te ye'\-krr.} The teacher required the students to arrive early.~\LINK{http://naviteri.org/2012/01/mipa-zisit-ayliu-amip-new-words-for-the-new-year/}{b} \\  
\textsc{Derivatives}: \ding{43}
\on{1}.~\HL{leye'krr}{\B{le\-ye'\-krr}}, early.

% ye'rìn
\hi 
\HT{ye'rIn}{\B{\cp{ye'}rìn}} \I{["jEP.RIn]} \emph{adv.} soon.   
\B{Sìl\-pey oe, la\-yu oe\-ru ye'\-rìn sìl\-tsan\-a fmawn.} I hope, I will soon have good news.~\LINK{http://forum.learnnavi.org/news-announcements/a-response-from-paul-frommer!/}{ln} 
\B{Tsa\-ria nga\-hu ye'\-rìn ul\-txa si nì\-mun, oe sre\-fe\-re\-i\-ey nì\-prr\-te'.} I'm looking forward to getting together with you again soon.~\LINK{http://naviteri.org/2011/02/new-vocabulary-part-2/}{b}

% yehaw
\hi 
\HT{yehaw}{\B{\cp{ye}haw}} \I{["jE.haw]} \emph{adj.} well\hyp{}rested, having had enough sleep. 
\B{Lam fra\-kem le\-tsun\-slu krr\-a fko 'e\-fu ye\-haw.} Everything seems possible when one is well\hyp{}rested.~\LINK{https://naviteri.org/2024/09/pxevola-liu-amip-two-dozen-new-words/}{b} 
(\ding{43}~\HL{ye}{\B{ye}}, \HL{hahaw}{\B{ha\-haw}})

% yehakx
\hi 
\HT{yehakx}{\B{\cp{ye}hakx}} \I{["jE.hak']} \emph{adj.} satisfied from hunger by food, ``full stomach''.  
\B{Sra\-ke ye\-hakx? -- Stum.\slash Ye. Tsun ti\-vam.\slash Nì\-txan!\slash Nì\-haawwwng.} Did you get enough to eat? -- Almost.\slash Yes. That'll be enough.\slash Very! I'm quite full.\slash Oooh. I ate too much.~\LINK{http://naviteri.org/2011/04/\%E2\%80\%99a\%E2\%80\%99awa-li\%E2\%80\%99fyavi-amip\%E2\%80\%94a-few-new-expressions/}{b}
(\ding{43}~\HL{ye}{\B{ye}}, \HL{ohakx}{\B{o\-hakx}})

% yem
\hi 
\HT{yem}{\B{yem}} \I{[jEm]} \emph{vtr.} put, place. 
\B{Ngey tsko\-ti yem tsa\-yì\-sìn tsa\-krr za\-'u fì\-tseng!} Put your bow on that ledge and come here!~\LINK{http://naviteri.org/2013/01/lifyari-po-peyi-how-good-is-her-navi/}{b} \\  
\textsc{Derivatives}: \ding{43}
\on{1}.~\HL{yemfpay}{\B{yem\-fpay}}, dipping. 
\on{2}.~\HL{yemstokx}{\B{yem\-stokx}}, put on (clothing). 
\on{3}.~\HL{kllyem}{\B{kll\-yem}}, bury. 
\on{4}.~\HL{'awstengyem}{\B{'aw\-steng\-yem}}, join together.

% yemfpay / si
\hi 
\HT{yemfpay}{\B{yem\cp{fpay}}} \I{[jEm."fpaj]} \textsc{i}. \emph{n.} dipping, immersion (into a liquid). \\
\textsc{ii}. \B{yem\cp{fpay} si} \I{[jEm."fpaj si]} \emph{vin.} dip (into a liquid). 
\B{Ton\-ìri a\-la\-he, aw\-nga ke yem\-fpay si keng 'aw\-lo; fì\-txon yem\-fpay si me\-lo.} As for other nights, we don't even dip once; this night (we) dip twice.~\LINK{http://whyisthisnight.com/addl-more.html\#na\%27vi}{tn} 
(\ding{43}~\HL{yem}{\B{yem}}, \HL{nemfa}{\B{nem\-fa}} \HL{pay}{\B{pay}})

% yemstokx
\hi
\HT{yemstokx}{\B{\cp{yem}stokx}} \I{["j$\cdot$Em.stok']} \emph{vtr., inf.}\on{11} put on (functional clothing), don; wear. (\emph{from} \B{yem sìn tokx}, ``put on the body'')
\B{Pen\-it yem\-stokx!} Get dressed!~\LINK{http://naviteri.org/2011/08/new-vocabulary-clothing/}{b}
\B{Su\-nu oer haw\-re' a ngal yo\-lem\-stokx.} I like the hat you're wearing.~\LINK{http://naviteri.org/2011/08/new-vocabulary-clothing/}{b}
\B{Fì\-re\-won ngal lum\-pe kea haw\-re'\-it ke yo\-lem\-stokx?} Why didn't you put on a hat this morning?~\LINK{http://naviteri.org/2011/08/new-vocabulary-clothing/}{b}
\B{Fì\-ra\-spu' 'e\-kxin lu nì\-hawng. Ke tsun oe yi\-vem\-stokx.} These leggings are too tight. I can't wear them.~\LINK{http://naviteri.org/2012/07/meetings-waterfalls-and-more/}{b}
(\ding{43}~\HL{ioi si}{\B{ioi si}}, \ding{214}~\HL{'aku}{\B{'a\-ku}})

% yengwal
\hi 
\HT{yengwal}{\B{yeng\cp{wal}}} \I{[jEN."wal]} \emph{n.} sorrow
\B{Sa'\-sem\-ìri lu 'eveng\-ä kxitx yeng\-wal a\-tu\-vom.} For a parent, the greatest sorrow is the death of a child.~\LINK{http://naviteri.org/2015/04/some-new-words-for-may-day/}{b}

% yerik
\hi 
\HT{yerik}{\B{\cp{ye}rik}} \I{["jE.Rik\textcorner]} \emph{n., fauna} yerik, hexapede, \emph{sexcruscervus caeruleus}.  
\B{Ye\-rik lu swi\-rä a\-nim nì\-txan.} The hexapede is a very timid creature.~\LINK{http://naviteri.org/2011/09/miscellaneous-vocabulary/}{b}
\B{Tì\-ran nì\-fnu txan\-txew\-vay fte\-ke ay\-ye\-rik\-ìl aw\-nga\-ti sti\-vawm.} Walk as quietly as possible so the hexapedes won't hear us.~\LINK{http://naviteri.org/2012/10/snewsyea-ftxoza-halowina-a-spooky-halloween/}{b} 
\B{Pol ye\-rik\-it ta\-ron nìl\-tsan fra\-to.} He hunts yerik the best of all.~\LINK{http://forum.learnnavi.org/language-updates/confirmations-\%28comparisons-gerund\%29/}{ln} 
\B{Zawr ye\-rik\-ä lu 'a\-ngo.} The call of the hexapede is quiet.~\LINK{http://naviteri.org/2014/11/twenty-before-the-holidays/}{b}
\B{Ra\-lu\-ri fkan na ye\-rik.} Ralu smells (looks\slash sounds) like a hexapede.~\LINK{http://naviteri.org/2012/11/renu-ayinanfyaya-the-senses-paradigm/}{b}
(\emph{a. idiom}) \B{ye\-rik (a) mì yrr\-ap} \emph{extreme panic, anxiety or timidity}, [lit.: Hexapede in a storm].~\LINK{http://naviteri.org/2021/08/ulte-ayyoratu-leiu-and-the-winners-are-2/}{b}

% yeväng
\hi 
\HT{yevAng}{\B{\cp{ye}väng}} \I{["jE.v\ae{}N]} \emph{adj.} satisfied from thirst by drink, feeling quenched\slash slaked.
\B{Sre\-krr 'a\-me\-fu väng, set ye\-väng.} Before, I was thirsty; now my thirst has been quenched.~\LINK{http://naviteri.org/2011/04/\%E2\%80\%99a\%E2\%80\%99awa-li\%E2\%80\%99fyavi-amip\%E2\%80\%94a-few-new-expressions/}{b} 
(\ding{43}~\HL{ye}{\B{ye}}, \HL{vAng}{\B{väng}}; \ding{214}~\HL{vAng}{\B{väng}})

% yewla / si
\hi 
\HT{yewla}{\B{\cp{yew}la}} \I{["jEw.la]} \textsc{i}. \emph{n.} disappointment, emotional let\hyp{}down, failed expectation.   
(\emph{a. with the dative and} \HL{lu}{\B{lu}}) to be disappointed. 
\B{Oer lu txan\-a yew\-la a ke tsun nga oe\-hu ki\-vä\-teng me\-srr\-ay.} I'm very disappointed you can't hang out with me the day after tomorrow.
(\emph{b. conv. idiom}) \B{Yew\-la!} Bummer! That's a shame! What a shame!~\LINK{http://naviteri.org/2012/10/mipa-vospxi-mipa-ayliu-new-words-for-the-new-month/}{b} \\
\textsc{ii}. \B{\cp{yew}la si} \I{["jEw.la si]} \emph{vin.} disappoint s.o.
\B{O\-mum oel fu\-ta sä\-nui pe\-yä nga\-ru yew\-la so\-li nì\-txan.} I know that his failure disappointed you greatly.~\LINK{http://naviteri.org/2023/12/mipa-zisit-ayliu-amip-new-year-new-words/}{b}  \\  
\textsc{Derivatives}: \ding{43}
\on{1}.~\HL{leyewla}{\B{le\-yew\-la}}, disappointing. 
\on{2}.~\HL{nIyewla}{\B{nì\-yew\-la}}, in a disappointing fashion. 
\on{3}.~\HL{yawnyewla}{\B{yawn\-yew\-la}}, broken heart.

% yewn
\hi 
\HT{yewn}{\B{yewn}} \I{[jEwn]} \emph{vtr.} express, convey (\emph{a thought or feeling}). 
\B{Oe new oey sì\-'e\-fut yi\-vewn po\-e\-ru, slä ke tsä\-ngun.} I want to express my feelings to her, but, sadly, I can't.~\LINK{http://naviteri.org/2020/11/vospxivopeya-ayliu-amip-novembers-new-words/}{b} \\  
\textsc{Derivatives}: \ding{43}
\on{1}.~\HL{tIyewn}{\B{tì\-yewn}}, expression (\emph{of a thought or feeling}).
\on{2}.~\HL{sAfpIlyewn}{\B{sä\-fpìl\-yewn}}, communicate. 

% yey 
\hi 
\HT{yey}{\B{yey}} \I{[jEj]} \emph{adj.} straight.
\B{Fì\-tskxe\-ri fa\-'o lu yey; ke lu ko\-um.} This rock has straight sides; it's not rounded.~\LINK{http://naviteri.org/2013/09/on-si-salewfya-shapes-and-directions/}{b}
(\ding{214}~\HL{lIktap}{\B{lìk\-tap}}) \\  
\textsc{Derivatives}: \ding{43}
\on{1}.~\HL{nIyey}{\B{nì\-yey}}, directly, straight to the point; just. 
\on{2}.~\HL{yeyfya}{\B{yey\-fya}}, straight line.

% yeyfya
\hi 
\HT{yeyfya}{\B{\cp{yey}fya}} \I{["jEj.fja]} \emph{n.} straight line.
\B{Wo\-leyn Ìstawl yey\-fyat mì hll\-te fte oeyk\-ti\-vìng fra\-po\-ru tì\-hawl\-te\-ri sne\-yä.} Ìstaw drew a line on the ground to explain his plan to everyone.~\LINK{http://naviteri.org/2013/09/on-si-salewfya-shapes-and-directions/}{b}
(\ding{43}~\HL{yey}{\B{yey}}, \HL{fya'o}{\B{fya'o}})  \\  
\textsc{Derivatives}: \ding{43}
\on{1}.~\HL{nIyeyfya}{\B{nì\-yey\-fya}}, straight ahead, in a straight line.

% yì
\hi 
\HT{yI}{\B{yì}} \I{[jI]} \emph{n.} (\emph{a.}) shelf, ledge, level, step, rung, \emph{a small, flat area on which one can stand or place an object, such as a foot}.
\B{Ta\-lun\-a vo\-spxì\-o a\-mrr ke zo\-lup tom\-pa, lä\-ngu yì kil\-van\-ä tìm nì\-txan.} Because it hasn't rained for five months, the river level is very low.~\LINK{http://naviteri.org/2013/01/lifyari-po-peyi-how-good-is-her-navi/}{b}
\B{Ngey tsko\-ti yem tsa\-yì\-sìn tsa\-krr za\-'u fì\-tseng!} Put your bow on that ledge and come here!~\LINK{http://naviteri.org/2013/01/lifyari-po-peyi-how-good-is-her-navi/}{b}
(\emph{b. metaphorically with} \HL{kllkxem}{\B{kll\-kxem}} \emph{to refer to the level of anything scalable: water level, temperature, talent, anger, etc.}) 
\B{Lì'\-fya\-ri po (kll\-kxem) sìn pe\-yì?} [\emph{coll.}: \I{[sIm.pE."jI]}]\slash \B{Lì'\-fyari po pe\-yì?} How good is her Na'vi?~\LINK{http://naviteri.org/2013/01/lifyari-po-peyi-how-good-is-her-navi/}{b}
\B{Sìn yì sngä\-'i\-yu\-ä, slä tsye\-rìl (hay\-a yì\-ne) nì\-'ul\-'ul.} They're at a beginner's level, but they're getting better and better.~\LINK{http://naviteri.org/2013/01/lifyari-po-peyi-how-good-is-her-navi/}{b}
\B{Sìn yì a ke tsun kaw\-tu spi\-vaw, nì\-win fra\-to.} He runs at an incredible level, faster than anyone else.~\LINK{http://naviteri.org/2013/01/lifyari-po-peyi-how-good-is-her-navi/}{b} 
(\emph{c. rhyming expression}) \B{Tslikx, tì\-ran, tul; | Ftu yì ne yì tsan\-'ul.} \emph{when learning s.t. new, you have to proceed from step to step: baby steps first, then bigger ones} [lit.: `Crawl, walk, run; from level to level get better.']~\LINK{http://naviteri.org/2023/09/aawa-ayliu-si-aylifyavi-amip-a-few-new-words-and-expressions/}{b} \\  
\textsc{Derivatives}: \ding{43}
\on{1}.~\HL{snayI}{\B{sna\-yì}}, staircase, series of step\hyp{}like levels. 
\on{2}.~\HL{kxaylyI}{\B{kxayl\-yì}}, high level. 
\on{3}.~\HL{kxamyI}{\B{kxam\-yì}}, intermediate level. 
\on{4}.~\HL{tImyI}{\B{tìm\-yì}}, low level. 
\on{5}.~\HL{fyanyI}{\B{fyan\-yì}}, shelf.
\on{6}.~\HL{yIspul}{\B{yì\-spul}}, mermaid tail.
\on{7}.~\HL{heynyI}{\B{heyn\-yì}}, lap.
\on{8}.~\HL{swayI}{\B{swa\-yì}}, generation within a family.

% yìm
\hi 
\HT{yIm}{\B{yìm}} \I{[jIm]} \emph{vtr.} bind. 
\B{Ay\-wayl yìm ki\-fkey\-ä | 'Ì\-he\-yut a\-vo\-mrr | Sìn ti\-rea\-fya\-'o a\-vol | Na way\-te\-lem\-ä hìng.} The songs bind | the thirteen spirals of the solid world | To the eight spirit paths | Like the threads of a Songcord.~\LINK{http://www.pandorapedia.com/navi/music/ritual_music}{pp}
\B{Me\-fo\-ti yìm.} Tie them up.~\LINK{https://wiki.learnnavi.org/Na\%27vi_from_Avatar_Movie\#Attack_on_hometree}{a} \\  
\textsc{Derivatives}: \ding{43}
\on{1}.~\HL{sAyIm}{\B{sä\-yìm}}, tie, s.t. used for binding.

% yìspul
\hi 
\HT{yIspul}{\B{\cp{yì}spul}} \I{["jI.spul]} \emph{or} \I{[.spUl]} \emph{n.} \ding{96} mermaid tail. 
(\ding{43}~\HL{yI}{\B{yì}}, \HL{spule}{\B{spu\-le}})

% yll
\hi 
\HT{yll}{\B{yll}} \I{[j\s{l}]} \emph{adj.} communal. 
\B{Kxawm set 'u a\-ngä\-zìk fra\-to lu la\-'a a\-yll.} Maybe the most difficult thing now is social distance.~\LINK{http://naviteri.org/2020/03/teri-tifkeytok-lefkrr-about-the-present-situation/}{b} \\
(i.) \B{swoa\-sey a\-yll}, large social kava bowl.~\LINK{http://naviteri.org/2015/07/aylifyavi-lereyfya-1-cultural-terms-1/}{b} \\
(ii.) \B{ko\-ren a\-yll} law, societal rule.
\B{Pol aw\-nga\-ru ho\-ren\-it a\-yll ley\-kek.} He makes us obey the laws.~\LINK{http://naviteri.org/2020/06/mi-tanlokxe-oeya-srr-afpxamo-terrible-days-in-my-country/}{b} \\
\textsc{Derivatives}: \ding{43}
\on{1}.~\HL{nIyll}{\B{nì\-yll}}, communally, in a communal manner. 
\on{2}.~\HL{ylltxep}{\B{yll\-txep}}, communal fire.

% ylltxep
\hi 
\HT{ylltxep}{\B{\cp{yll}txep}} \I{["j\s{l}.t'Ep\textcorner]} \emph{n.} communal fire or fire pit.
\B{Nì\-trr\-trr yom Na'\-vil wu\-tsot 'aw\-si\-teng pxaw yll\-txep.} The Na'vi regularly eat dinner together around a communal fire.~\LINK{http://naviteri.org/2012/07/fivospxiya-aylifyavi-amip-this-months-new-expressions-pt-2/}{b}
(\ding{43}~\HL{yll}{\B{yll}}, \HL{txep}{\B{txep}}, \HL{txeptseng}{\B{txep\-tseng}}) 

% yo
\hi 
\HT{yo}{\B{yo}} \I{[jo]} \emph{n.} (\emph{a.}) surface.
\B{Yo\-sìn kil\-van\-ä lu tatx.} There are bubbles on the surface of the river.~\LINK{http://naviteri.org/2015/11/vomuna-liu-amip-ten-new-words/}{b}
\B{Krr\-a hu\-fwe tul nì\-win, tsun fko tsi\-ve\-'a ay\-we\-opx\-it a sìn yo pay\-ä.} When there is strong wind, you can see waves on the water.~\LINK{http://naviteri.org/2017/09/zisikrr-amip-ayliu-amip-new-words-for-the-new-season/}{b}
(\emph{b. coll. form of} \B{fyan\-yo}) table. 
\B{Yo\-ti fwal ru\-txe. Mei slo\-lu maw tom\-pa.} Please wipe the table. It's gotten wet after the rain.~\LINK{http://naviteri.org/2020/07/vospxivol-lefpom-happy-august/}{b}
\B{Oe\-yä tstal\-ìl tok pe\-se\-nget? -- Lu yo\-sìn.} Where's my knife? -- It's on the table.~\LINK{http://naviteri.org/2015/08/aylifyavi-lereyfya-2-cultural-terms-2-and-more/}{b} \\  
\textsc{Derivatives}: \ding{43}
\on{1}.~\HL{txayo}{\B{txa\-yo}}, field, plain. 
\on{2}.~\HL{kxemyo}{\B{kxem\-yo}}, wall.
\on{3}.~\HL{fyanyo}{\B{fyan\-yo}}, table. 
\on{4}.~\HL{yomyo}{\B{yom\-yo}}, plate (for food).
\on{5}.~\HL{fApyo}{\B{fäp\-yo}}, roof.

% yo'
\hi 
\HT{yo'}{\B{yo'}} \I{[joP]} \emph{vin.} be perfect, flawless.
\B{Tì\-hawl le\-sngä\-'i lu tì\-kang\-kem\-vi skxawng\-ä, slä pum alu fì\-'u yo' nì\-'aw.} The original plan was the work of an idiot, but this one is just perfect.~\LINK{http://naviteri.org/2012/01/mipa-zisit-ayliu-amip-new-words-for-the-new-year/}{b}
\B{Fì\-stxel\-it fol txe\-ru\-la fpi olo'\-eyk\-tan. Ze\-ne yi\-vo' lu\-ke kxey\-ey\-o kaw\-'it.} They're constructing this gift for the chief. It must be perfect without a single flaw.~\LINK{http://naviteri.org/2012/01/mipa-zisit-ayliu-amip-new-words-for-the-new-year/}{b}
\B{Ri\-ni\-ri nik\-re yä\-ngo' nì\-tut.} Rini's hair is always perfect.~\LINK{http://naviteri.org/2012/01/mipa-zisit-ayliu-amip-new-words-for-the-new-year/}{b}
(\emph{b. conv. idiom}) \B{Ye\-io'}, perfect. 
\B{Ul\-txa si\-vi oeng sìn ram\-tsyìp txon\-'ong\-ay. -- Ye\-io'! Tsa\-krr\-vay ko!} Let's meet on the hill tomorrow at nightfall. -- Perfect! See you then.~\LINK{http://naviteri.org/2012/01/mipa-zisit-ayliu-amip-new-words-for-the-new-year/}{b} \\  
\textsc{Derivatives}: \ding{43}
\on{1}.~\HL{tIyo'}{\B{tì\-yo'}}, perfection. 
\on{2}.~\HL{nIyo'}{\B{nì\-yo'}}, pefectly, flawlessly. 
\on{3}.~\HL{yo'ko}{\B{yo'\-ko}}, circle. 

% yo'ko
\hi 
\HT{yo'ko}{\B{\cp{yo'}ko}} \I{["joP.ko]} \emph{n.} circle, \emph{a perfectly circular ring}. 
(\ding{43}~\HL{yo'}{\B{yo'}}, \HL{koum}{\B{ko\-um}}) \\  
\textsc{Derivatives}: \ding{43}
\on{1}.~\HL{yo'kofya}{\B{yo'\-ko\-fya}}, cycle.

% yo'kofya
\hi
\HT{yo'kofya}{\B{\cp{yo'}kofya}} \I{["joP.ko.fja]} \emph{n.} cycle.
\B{Hum zì\-sìt a\-lal, pä\-hem pum a\-mip. Yo'\-ko\-fya a\-tì\-'i\-lu\-ke.} The old year leaves, the new one arrives. An endless cycle.~\LINK{http://naviteri.org/2017/12/zisit-amip-lefpom-ma-eylan-happy-new-year-friends/}{b}
(\ding{43}~\HL{yo'ko}{\B{yo'\-ko}}, \HL{fya'o}{\B{fya\-'o}})

% yoa
\hi 
\HT{yoa}{\B{\cp{yo}a}} \I{["jo.a]} \emph{adp.} in exchange for (\emph{usually used with} \B{tìng}, \B{tel}, \B{ka\-nom}, \B{stxe\-nu\-tìng}, \B{mll\-'an})
\B{Oel to\-lìng nga\-ru tsngan\-it yoa fkxen.} I traded you meat for vegetables.~\LINK{http://naviteri.org/2014/02/barter-and-exchange/}{b}
\B{Fol ko\-la\-nom po\-ta ay\-srok\-it fay\-o\-ang\-yoa.} They bartered fish with him for beads.~\LINK{http://naviteri.org/2014/02/barter-and-exchange/}{b}
\B{Ta\-yel Tse\-nul pxe\-swi\-zaw\-ti yoa mun\-sna\-hawn\-ven.} Tsenu will receive three arrows in exchange for a pair of shoes.~\LINK{http://naviteri.org/2014/02/barter-and-exchange/}{b}
\B{Kä\-sro\-lìn oel nik\-ro\-it Pey\-ral\-ur yoa fwa po rol oer.} I loaned Peyral a hair ornament in exchange for her singing to me.~\LINK{http://naviteri.org/2014/02/barter-and-exchange/}{b}
\B{Fu\-ta nga\-ta tel pxen\-yoa srät, mll\-'e\-i\-an oel.} I'm happy to take cloth from you for finished garments.~\LINK{http://naviteri.org/2014/02/barter-and-exchange/}{b}
\B{Poel haw\-re'\-tsyìp\-it ko\-la\-nom yoa txo\-lar a\-me\-vol.} She bought a little cap for \$\on{16}.~\LINK{http://naviteri.org/2014/02/barter-and-exchange/}{b}
\B{Kì\-ma\-nom oel mip\-a el\-tut le\-fngap yoa ew\-ro.} I just bought a new computer (for an unspecified amount of euros).~\LINK{http://naviteri.org/2014/02/barter-and-exchange/}{b}

% yokx
\hi 
\HT{yokx}{\B{yokx}} \I{[jok']} \emph{n.} shield.

% yol
\hi 
\HT{yol}{\B{yol}} \I{[jol]} \emph{adj.} short (of time).
\B{Yol\-a krr, txan\-a krr, ke tsran\-ten.} It doesn't matter how long it takes.~\LINK{http://forum.learnnavi.org/language-updates/a-collection/}{ln\slash a\textsuperscript{c}}
\B{Sra\-ne, lu yol\-a lì'\-fya\-vi.} Yes, there is a short expression.~\LINK{http://naviteri.org/2012/10/snewsyea-ftxoza-halowina-a-spooky-halloween/comment-page-1/\#comment-1755}{b}
\B{Maw sä\-tsway\-on a\-yol ay\-oe kll\-po\-lä mì ta\-yo a lu ro\-fa kil\-van.} After a short flight we landed in a field beside the river.~\LINK{http://naviteri.org/2012/01/mipa-zisit-ayliu-amip-new-words-for-the-new-year/}{b}
(\ding{214}~\HL{nun}{\B{nun}}, \HL{txan}{\B{txan}}) 

% yom 
\hi 
\HT{yom}{\B{yom}} \I{[jom]} \emph{vtr.} eat.
\B{Yom tey\-lut keng oel!}  Even \emph{I} eat teylu!~\LINK{http://forum.learnnavi.org/language-updates/even-keng/}{ln}
\B{Oe new yi\-vom tey\-lut.} I want to eat teylu.~\LINK{http://wiki.learnnavi.org/index.php?title=Canon\#Extracts_from_various_emails}{wiki} 
\B{Oe ye\-rom set.} I'm eating now.~\LINK{http://naviteri.org/2012/03/spring-vocabulary-part-1/}{b}
\B{Nì\-'aw\-ve oeng yo\-lom wu\-tsot; maw\-krr uvan si.} First we had dinner; afterwards we played.~\LINK{http://naviteri.org/2011/08/new-vocabulary-clothing/comment-page-1/\#comment-917}{b}
\B{Oe nga\-to yom nì\-'ul.} I eat more than you.~\LINK{http://forum.learnnavi.org/language-updates/prefer-x-to-y-24474/}{ln}
(i.) \B{mo a yom}, dining room, dining area.
(\emph{a. proverb}) \B{Ta\-ron\-yut yom smar\-ìl.} \emph{A totally botched situation, chaos} [lit.: the prey eats the hunter].
\B{Eyk Ka\-mun a fra\-lo lä\-ngu tsa\-sä\-ta\-ron vel\-ke nì\-wotx. Ta\-ron\-yut yom smar\-ìl!} Every time Kamun is in charge, the hunt is a mess. Everything goes wrong that can!~\LINK{http://naviteri.org/2012/10/mipa-vospxi-mipa-ayliu-new-words-for-the-new-month/}{b} \\  
\textsc{Derivatives}: \ding{43}
\on{1}.~\HL{yomhI'ang}{\B{yom\-hì\-'ang}}, dakteron.  
\on{2}.~\HL{yomioang}{\B{yom\-io\-ang}}, chalice plant. 
\on{3}.~\HL{yomtIng}{\B{yom\-tìng}}, feed. 
\on{4}.~\HL{yomyo}{\B{yom\-yo}}, plate (for food).
\on{5}.~\HL{yomvey}{\B{yom\-vey}}, dine on flesh, be carnivorous.
\on{6}.~\HL{yomzeswa}{\B{yom\-ze\-swa}}, graze.
\on{7}.~\HL{ukyom}{\B{uk\-yom}}, eclipse.

% yomhì'ang
\hi 
\HT{yomhI'ang}{\B{yom\cp{hì'}ang}} \I{[jom."hIP.aN]} \emph{n.}, \ding{96} dakteron, `insectivore', \emph{pseudopenthes coralis}.
(\ding{43}~\HL{yom}{\B{yom}}, \HL{hI'ang}{\B{hì\-'ang}})

% yomioang
\hi 
\HT{yomioang}{\B{yomi\cp{o}ang}} \I{[jom.i"o.aN]} \emph{n.}, \ding{96} chalice plant, \emph{pseudocenia rosea}.
(\ding{43}~\HL{yom}{\B{yom}}, \HL{ioang}{\B{ioang}})

% yomtìng
\hi 
\HT{yomtIng}{\B{\cp{yom}tìng}} \I{["jom.t$\cdot$IN]} \emph{vin., inf.}\on{22} feed.   
\B{Sa'\-nok prr\-nen\-ur yom\-tìng fa okup sne\-yä.} A mother feeds an infant with her milk.~\LINK{http://naviteri.org/2013/06/looking-back-looking-forward/}{b}
\B{Ta\-ron\-yu a\-swey lu tsa\-tu\-te a ta\-ron fte yom\-ti\-vìng ay\-la\-ru.} The best hunter is the one who hunts to feeds others.~\LINK{http://forum.learnnavi.org/language-updates/clarification-on-the-use-of-awpo/}{ln}
\B{Ha ftxey | 'aw\-pot set ay\-ngal a lu ay\-nga\-kip | 'Aw\-pot a Na'\-vi\-ru yom\-tì\-yìng.} So choose | one among you | Who will feed the People.~\LINK{http://naviteri.org/2013/08/taronway-the-hunt-song/}{b}
(\ding{43}~\HL{yom}{\B{yom}}, \HL{tIng}{\B{tìng}})

% yomvey
\hi 
\HT{yomvey}{\B{yom\cp{vey}}} \I{[j$\cdot$om."vEj]} \emph{vin., inf.}\on{11} dine on flesh, be carnivorous. 
\B{Pa\-lu\-kan\-tsyìp yom\-vey nì\-wotx.} All cats are carnivorous.~\LINK{http://naviteri.org/2019/06/50a-liu-amip-40-new-words/}{b} 
(\ding{43}~\HL{yom}{\B{yom}}, \HL{vey}{\B{vey}})

% yomyo
\hi 
\HT{yomyo}{\B{\cp{yom}yo}} \I{["jom.jo]} \emph{n.} plate (for food).
\B{Ma sem\-pu, oey yom\-yo tso\-lawng!} Daddy, my plate broke!~\LINK{http://naviteri.org/2020/04/aawa-u-amip-a-few-new-things/}{b}
\B{Ma En\-tu, ngal lum\-pe ngey tsmu\-ke\-yä yom\-yot tsey\-ko\-lawng?} En\-tu, why did you break your sister's plate?~\LINK{http://naviteri.org/2020/04/aawa-u-amip-a-few-new-things/}{b}
(\ding{43}~\HL{yom}{\B{yom}}, \HL{yo}{\B{yo}}) \\  
\textsc{Derivatives}: \ding{43}
\on{1}.~\HL{yomyolerIk}{\B{yom\-yo le\-rìk}}, leaf plate.

% yomyo lerìk
\hi 
\HT{yomyolerIk}{\B{\cp{yom}yo le\cp{rìk}}} \I{["jom.jo{} lE."rIk\textcorner]} \emph{n.} leaf plate; (\emph{coll.}: \HL{rIk}{\B{rìk}})
(\ding{43}~\HL{yomyo}{\B{yom\-yo}}, \HL{lerIk}{\B{le\-rìk}})

% yomzeswa
\hi 
\HT{yomzeswa}{\B{yom\cp{ze}swa}} \I{[j$\cdot$om."zE.swa]} \emph{vin., inf.}\on{11} graze. 
\B{Sna\-ye\-rik ye\-rom\-ze\-swa mì ta\-yo.} A herd of hexapedes are grazing in the field.~\LINK{http://naviteri.org/2019/06/50a-liu-amip-40-new-words/}{b} 
(\ding{43}~\HL{yom}{\B{yom}}, \HL{zeswa}{\B{ze\-swa}})

% yora'
\hi 
\HT{yora'}{\B{yo\cp{ra'}}} \I{[jo."RaP]} \emph{vtr.} win.
\B{Krr\-a wä\-tut tse\-'a, pe\-yä sä\-sìl\-pey a yo\-ra' 'o\-lìp.} When he saw his opponent, his hope of winning vanished.~\LINK{http://naviteri.org/2013/01/awvea-posti-zisita-amip-first-post-of-the-new-year/}{b}
\B{Tsu'\-tey\-ri lu fo\-ru nrr\-a a fì\-re\-won yo\-lo\-ra'.} They're proud of Tsu'\-tey for having won this morning.~\LINK{http://naviteri.org/2012/07/fivospxiya-aylifyavi-amip-this-months-new-expressions-pt-2/}{b}  
\B{Am\-'a\-lu\-ke sna\-yaytx Saw\-tu\-te, ya\-yo\-ra' Na'\-vi.} Without a doubt the Sky People will lose and the People will win.~\LINK{http://naviteri.org/2012/06/spring-vocabulary-part-3/}{b}
\B{Yì\-mo\-ra' Tsu'\-tey\-ìl uvan\-it.} Tsu'\-tey just won the game.~\LINK{http://naviteri.org/2011/05/some-miscellaneous-vocabulary/}{b}
(\ding{214}~\HL{snaytx}{\B{snaytx}}) \\  
\textsc{Derivatives}: \ding{43}
\on{1}.~\HL{yora'tu}{\B{yo\-ra'\-tu}}, winner. 
\on{2}.~\HL{tIyora'}{\B{tì\-yo\-ra'}}, victory, a win.

% yora'tu
\hi 
\HT{yora'tu}{\B{yo\cp{ra'}tu}} \I{[jo."RaP.tu]} \emph{n.} winner. 
\B{Fra\-u\-van\-ìri lu yo\-ra'\-tu, lu snay\-tu.} For every game, there's a winner and a loser.~\LINK{http://naviteri.org/2011/12/a-note-on-the-word-\%E2\%80\%9Cyora\%E2\%80\%99tu\%E2\%80\%9D/}{b}
\B{Ay\-nga\-ru tsìng\-a yo\-ra'\-tut!} May I present the four winners!~\LINK{http://naviteri.org/2014/12/ayngaru-tsinga-yoratu-may-i-present-the-four-winners/}{b}
(\ding{214}~\HL{snaytu}{\B{snay\-tu}}, \ding{43}~\HL{yora'}{\B{yo\-ra'}})

% yoten
\hi 
\HT{yoten}{\B{\cp{yo}ten}} \I{["jo.tEn]} \emph{n., fauna} yoten. 

% yrr
\hi 
\HT{yrr}{\B{yrr}} \I{[j\s{r}]} \emph{adj.} wild, natural (\emph{s.t. in its original, unmodified, untampered\hyp{}with natural state}), (\emph{of food}) raw.
\B{Ik\-ran\-ìri krr\-a hu tu\-te tsa\-heyl si, ftang li\-vu yrr.} When an \emph{ikran} bonds with a person, it ceases to be wild.~\LINK{http://naviteri.org/2013/01/awvea-posti-zisita-amip-first-post-of-the-new-year/}{b}
\B{Lu tsa\-fne\-pay\-oang ftxì\-lor fra\-to krr\-a lu yrr.} That kind of fish is most tasty when eaten as \emph{sashimi}.~\LINK{http://naviteri.org/2013/01/awvea-posti-zisita-amip-first-post-of-the-new-year/}{b}
\B{Txep\-ìri, yrr\-a rìn\-ti rä\-'ä sar, ma skxawng.} Don't use that green wood for a fire, you fool.~\LINK{http://naviteri.org/2013/01/awvea-posti-zisita-amip-first-post-of-the-new-year/}{b}
(\ding{214}~\HL{zAfi}{\B{zä\-fi}})

% yrrap
\hi 
\HT{yrrap}{\B{\cp{yrr}ap}} \I{["j\s{r}.ap\textcorner]} \emph{n.} storm.   
\B{Fra\-po ne mì\-fa! Le\-rok yrr\-ap apxa!} Everybody inside! A big storm is approaching!~\LINK{http://naviteri.org/2012/07/meetings-waterfalls-and-more/}{b}
(\emph{a. idiom}) \B{ye\-rik (a) mì yrr\-ap} \emph{extreme panic, anxiety or timidity}, [lit.: Hexapede in a storm].~\LINK{http://naviteri.org/2021/08/ulte-ayyoratu-leiu-and-the-winners-are-2/}{b}
(\ding{43}~\HL{ya}{\B{ya}}, \HL{hrrap}{\B{hrr\-ap}})

% -yu
\hi
\HT{-yu}{\B{-yu}} \I{[.ju]} \emph{suf. to turn verbs into an agent, e.g.:} \B{ta\-ron\-yu}, hunter; \B{ma\-wey\-pey\-yu}, patient person, \emph{etc.}
(\ding{43}~\HL{-tu}{\B{-tu}}) 

% yuey
\hi 
\HT{yuey}{\B{\cp{yu}ey}} \I{["ju.Ej]} \emph{adj.} beautiful (\emph{inner beauty, only for people}).
\B{Lu poe lor, lu yu\-ey nì\-teng.} She's beautiful on the outside \emph{and} the inside.~\LINK{http://naviteri.org/2012/03/spring-vocabulary-part-1/}{b}
(\ding{43}~\HL{lor}{\B{lor}})

% yune
\hi 
\HT{yune}{\B{\cp{yu}ne}} \I{["ju.nE]} \emph{vtr.} listen to (\texttt{+}\emph{control}). 
\B{Nga ze\-ne ay\-lì\-'ut kar\-yu\-ä yi\-vu\-ne, ma 'evi.} You must listen to your teacher, my son.~\LINK{http://naviteri.org/2012/11/renu-ayinanfyaya-the-senses-paradigm/}{b}
\B{Yìl\-mu\-ne\-ie oel ay\-nge\-yä stä'\-nì\-pam\-it sok\-a tsa\-pos\-tì\-yä oe\-yä.} I was happy to listen to your recording of that recent blog post of mine.~\LINK{https://forum.learnnavi.org/projects/an-audio-message-for-pawl/msg670815/\#msg670815}{ln}
(\ding{43}~\HL{mikyun}{\B{tìng mik\-yun}}, \HL{stawm}{\B{stawm}})

% yur
\hi 
\HT{yur}{\B{yur}} \I{[juR]} \emph{or} \I{[jUR]} \emph{vtr.} wash, clean.  
\B{Ze\-ne nga yi\-vur pxìm fì\-skxir\-it.} You must clean the wound frequently.~\LINK{http://naviteri.org/2011/05/some-miscellaneous-vocabulary/}{b} 
\B{Oe yä\-pur.} I wash myself.~\LINK{http://forum.learnnavi.org/language-updates/the-reflexive-\%28from-a-frommer-email!!!!\%29/}{ln}
\B{Oe yä\-pe\-rur.} I'm washing myself.~\LINK{http://forum.learnnavi.org/language-updates/the-reflexive-\%28from-a-frommer-email!!!!\%29/}{ln} \\
\\
(i.) \B{txaw\-nul\-srung a yur}, washing machine.~\LINK{https://naviteri.org/2024/03/fleltrra-ayliu-words-for-april-fools-day/}{b}
\\
\end{multicols}

% ===== Section Z =====  HT/HL - 

 \pdfbookmark[0]{Z}{Z}
\begin{center}
\section*{Z \I{[z]} (\emph{Zä})}
\end{center}

\begin{multicols}{2}

% za'ärìp
\hi
\HT{za'ArIp}{\B{za\cp{'ä}rìp}} \I{[za."P$\cdot$\ae{}.R$\cdot$Ip\textcorner]} \emph{vtr., inf.}\on{23} pull (\emph{coll. also} \ding{43}~\HL{zArIp}{\B{zä\-rìp}}). 
(\ding{214}~\HL{kA'ArIp}{\B{kä\-'ä\-rìp}}, \ding{43}~\HL{'ArIp}{\B{'ärìp}}) \\
\textsc{Derivatives}: \ding{43}
\on{1}.~\HL{wrrzArIp}{\B{wrr\-zä\-rìp}}, pull out, extract.

% za'u 
\hi
\HT{za'u}{\B{\cp{za}'u}} \I{["za.Pu]} \emph{vin.} come (\HL{ne}{\B{ne}}, to\slash towards (\emph{can be optionally omitted if the destination comes after the verb}); \HL{ftu}{\B{ftu}}, from (a place); \HL{ta}{\B{ta}}, from (\emph{fig.})).   
\B{Nga za\-'u fì\-tseng pxìm srak?} Do you come here often?~\LINK{http://www.lemondrop.com/2010/01/26/avatar-pick-up-lines-talk-dirty-in-Navi/}{ld}
\B{Ftxey nga za\-'u fu\-ke?} Are you coming (or not)?~\LINK{http://forum.learnnavi.org/language-updates/if-vs-whether-ftxey-fuke/msg156629/\#msg156629}{ln}
\B{Za\-'u kal\-txì si ko!} Come and say hello!~\LINK{http://naviteri.org/2010/09/getting-to-know-you-part-3/}{b} 
\B{Za\-'u fra\-'u ne tu\-te lem\-wey\-pey, ke\-fyak?} Everything comes to the patient people, right?~\LINK{http://naviteri.org/2012/03/spring-vocabulary-part-2/}{b}
\B{Hawng\-krr rä\-'ä zi\-va\-'u!} Don't come late!~\LINK{http://naviteri.org/2010/07/vocabulary-update/}{b} 
\B{Nga zo\-la\-'u ftu pe\-seng?} Where are you from?~\LINK{http://naviteri.org/2010/09/getting-to-know-you-part-2/}{b}
\B{Pxay\-a a\-to\-ki\-ri\-na' ne fì\-ke\-tu\-wong zì\-ma\-'u.} Many \emph{atokirina} just came to this alien.~\LINK{http://forum.learnnavi.org/language-updates/pseudo-canon-from-bd-live-content/}{a\textsuperscript{c}\slash ln}
(\emph{a. idiom}) \B{zo\-la\-'u nì\-prr\-te'}, ``welcome'' (\B{ne}, to) \LINK{http://naviteri.org/2010/06/first-post/}{b}
\B{Zo\-la\-'u nì\-prr\-te' ay\-nga\-ru nì\-wotx!} Welcome to you all!~\LINK{http://naviteri.org/2010/06/first-post/}{b}
\B{Kal\-txì, ma fra\-po, ul\-te zo\-la\-'u nì\-prr\-te' aw\-nge\-yä nu\-mul\-txa\-ne a te\-ri lì'\-fya le\-Na'\-vi.} Hello, everyone and welcome to our Na'\-vi language class.~\LINK{http://naviteri.org/2012/07/tskxekengtsyip-a-mikyunfpi-a-little-listening-exercise/}{b}
(\emph{b. idiom}) \B{za\-'u nì\-'eng ne}, share an interest in common with s.o. 
\B{Tì\-ru\-sol za\-'u ne fo nì\-'eng.} They share an interest in singing.~\LINK{http://naviteri.org/2010/09/getting-to-know-you-part-1/}{b}
\B{So\-rewn\-ìl kan\-'ìn tì\-'em\-it nì\-txan ul\-te kxawm tsa\-txe\-le me\-nga\-ne za\-'a\-tsu nì\-'eng.} Sorewn is really into cooking and you two might share that in common.~\LINK{http://naviteri.org/2010/09/getting-to-know-you-part-1/}{b}
(\emph{c. of the weather}) \B{Ya\-fkeyk (za\-'u) fya\-pe?} How's the weather?~\LINK{http://naviteri.org/2011/04/yafkeykiri-plltxe-frapo-everyone-talks-about-the-weather/}{b} 
(\emph{d. idiom}) \emph{with} \HL{ne}{\B{ne}}, \emph{an adv. and a sense organ to convey} `s.o. looks\slash sounds\slash tastes\slash smells good or bad'
\B{Nga za\-'u nìl\-tsan ne na\-ri.} You look good.~\LINK{https://naviteri.org/2025/07/ayu-leketengsyon-various-things/}{b}
\B{Fì\-syu\-ve za\-'u nì\-fe' ne ftxì.} This food tastes bad.~\LINK{https://naviteri.org/2025/07/ayu-leketengsyon-various-things/}{b} 
\B{Nga za\-'u nì\-fpxa\-mo ne mik\-yun. Mok\-ri\-ri kem\-pe lo\-len?} You sound terrible. What happened to your voice?~\LINK{https://naviteri.org/2025/07/ayu-leketengsyon-various-things/}{b} \\
\textsc{Derivatives}: \ding{43}
\on{1}.~\HL{emza'u}{\B{emz\-a\-'u}}, pass, overcome.  
\on{2}.~\HL{fAza'u}{\B{fä\-za\-'u}}, come up, ascend. 
\on{3}.~\HL{kllza'u}{\B{kll\-za\-'u}}, come down, descend. 
\on{4}.~\HL{wrrza'u}{\B{wrr\-za\-'u}}, come out, emerge. 
\on{5}.~\HL{zIma'uyu}{\B{zì\-ma\-'u\-yu}}, newcomer. 
\on{6}.~\HL{zawprrte'}{\B{zaw\-prr\-te'}}, be enjoyable. 
\on{7}.~\HL{zusawkrr}{\B{zu\-saw\-krr}}, future.

% zakru
\hi 
\HT{zakru}{\B{\cp{zak}ru}} \I{["zak\textcorner.Ru]} \emph{n., fauna} zakru. 

% zal
\hi 
\HT{zal}{\B{zal}} \I{[zal]} \emph{n.} (natural) clay (\emph{usually for pottery}) or baked clay. 
\B{Fkol fì\-fne\-tsngal\-it txu\-la ta zal.} This kind of cup is made of clay.~\LINK{https://naviteri.org/2024/09/pxevola-liu-amip-two-dozen-new-words/}{b} \\
\textsc{Derivatives}: \ding{43}
\on{1}.~\HL{zAfizal}{\B{zä\-fi\-zal}}, fired clay, ceramic material.

% zam
\hi
\HT{zam}{\B{zam}} \I{[zam]} \emph{num., adj.} \on{64} (decimal), \on{100} (octal). (\emph{base}: \B{za-})
\B{Zam\-a Lì\-'u A\-mip.} \on{64} new words.~\LINK{http://naviteri.org/2018/04/100a-liu-amip-64-new-words-part-2/}{b}

% zam mì zam
\hi 
\HT{zammIzam}{\B{zam mì zam}} \I{[zam mI zam]} \emph{adv.} completely, one hundred percent. 
\B{Nga\-hu mll\-te oe zam mì zam.} I agree with you one hundred percent.~\LINK{http://naviteri.org/2021/12/zolau-niprrte-ma-3746-welcome-2022/}{b} 

% zamunge
\hi
\HT{zamunge}{\B{za\cp{mu}nge}} \I{[za."m$\cdot$u.N$\cdot$E]} (\textsc{rn}: \B{za\cp{mù}nge} \I{[za."m$\cdot$U.N$\cdot$E]}) \emph{vtr., inf.}\on{23} bring.
\B{Ay\-eveng\-ur kxo\-lìm sa'\-nok\-ìl tson\-ta pay\-it za\-mu\-nge.} The children were told by their mother to fetch water.~\LINK{http://naviteri.org/2014/10/tson-si-fpomron-obligation-and-mental-health/}{b}
\B{Pa\-ye\-'un swey\-a fya\-'ot a za\-mi\-vu\-nge oel ay\-ngar ay\-lì\-'ut ho\-ren\-ti\-sì lì'\-fya\-yä le\-Na'\-vi.} I will decide the best way to bring you the words and rules of Na'vi.~\LINK{http://forum.learnnavi.org/news-announcements/a-response-from-paul-frommer!/}{ln}
\B{Ay\-fol za\-mo\-lu\-nge aw\-ngar ay\-ri\-na'\-it fa sä\-he\-na a\-pxa.} They brought us the seeds in a large container.~\LINK{http://naviteri.org/2014/05/mipa-ayliu-mipa-aysafpil-new-words-new-ideas/}{b}
(\ding{43}~\HL{munge}{\B{mu\-nge}})

% zapxì   
\hi
\HT{zapxI}{\B{za\cp{pxì}}} \I{[za."p'I]} \emph{n.} front (part or section). 
(\ding{43}~\HL{hapxI}{\B{ha\-pxì}}) \\
\textsc{Derivatives}: \ding{43}
\on{1}.~\HL{zare'}{\B{zare'}}, forehead, brow.

% zare'
\hi
\HT{zare'}{\B{za\cp{re'}}} \I{[za."REP]} \emph{n.} forehead, brow.  
(\ding{43}~\HL{zapxI}{\B{za\-pxì}}, \HL{re'o}{\B{re\-'o}})

% zasrìn 
\hi
\HT{zasrIn}{\B{za\cp{srìn}}} \I{[za."sR$\cdot$In]} \emph{vtr., inf.}\on{22} borrow. 
\B{New Va'\-ru tskot Ey\-tu\-kan\-ä za\-sri\-vìn.} Va'ru wants to borrow Eytukan's bow.~\LINK{http://naviteri.org/2012/06/spring-vocabulary-part-3/}{b}
\B{Fra\-syu\-ra\-ti fkol za\-sro\-lìn nì\-'aw ul\-te trr\-o ze\-ne tey\-ki\-vä\-txaw.} All energy is only borrowed, and one day it will have to be given back.~\LINK{http://naviteri.org/2012/03/spring-vocabulary-part-1/}{b} 
(\ding{43}~\HL{srIn}{\B{srìn}})

% zazam 
\hi
\HT{zazam}{\B{\cp{za}zam}} \I{["za.zam]} \emph{num., adj.} \on{4096} (decimal), \on{10,000} (octal). 

% zawng
\hi
\HT{zawng}{\B{zawng}} \I{[zawN]} \emph{vin.} scream. 
\B{Tì\-txen si, ma 'ite! Unil sar\-mi nga teng\-krr ze\-rawng. Lu fpom srak?} Wake up, daughter! You were dreaming and screaming. Are you okay?~\LINK{http://naviteri.org/2015/03/kap-si-ayunil-saylahe-threats-dreams-and-other-things/}{b}

% zawprrte'
\hi
\HT{zawprrte'}{\B{zaw\cp{prr}te'}} \I{[z$\cdot$aw."p\s{r}.tEP]} \emph{vin., inf.}\on{11} be enjoyable (\emph{deriving physical or emotional pleasure from s.t.}); \emph{used with} \ding{43}~\HL{ne}{\B{ne}}. 
\B{Tsa\-fne\-pam\-tseo ke zaw\-prr\-te' oe\-ne.} I don't enjoy that kind of music.~\LINK{http://naviteri.org/2015/04/some-new-words-for-may-day/}{b}
\B{Sìl\-pey oe tsnì fpom li\-vu ay\-nga\-ru nì\-wotx ul\-te snge\-rä\-'i a zì\-sì\-krr a\-sang zi\-vaw\-prr\-te' ay\-nga\-ne.} I hope that you are all well and that you are enjoying the beginning warm season.~\LINK{http://naviteri.org/2015/04/some-new-words-for-may-day/}{b}
\B{zaw\-prr\-te' nì\-layl fko\-ne}, be pleasurable to one in an innocent way.~\LINK{http://naviteri.org/2023/09/aawa-ayliu-si-aylifyavi-amip-a-few-new-words-and-expressions/}{b}
\B{zaw\-prr\-te' nì\-to\-kat fko\-ne}, be pleasurable to one in a guilty way.~\LINK{http://naviteri.org/2023/09/aawa-ayliu-si-aylifyavi-amip-a-few-new-words-and-expressions/}{b}
(\ding{43}~\HL{za'u}{\B{za\-'u}}, \HL{nIprrte'}{\B{nì\-prr\-te'}})

% zawr
\hi
\HT{zawr}{\B{zawr}} \I{[zawR]} \emph{n.} animal call, an animal cry. 
\B{Zawr ye\-rik\-ä lu 'a\-ngo.} The call of the hexapede is quiet.~\LINK{http://naviteri.org/2014/11/twenty-before-the-holidays/}{b}
(\emph{a. idiom}) \B{zawr (a) mì na'\-rìng}, ``Old news.'' [lit.: an animal cry in the forest].~\LINK{http://naviteri.org/2021/08/ulte-ayyoratu-leiu-and-the-winners-are-2/}{b}
\B{Sra\-ke ngal sto\-lawm fu\-ta Tse\-nu Ra\-lur mo\-wan lu nì\-txan? -- Zawr mì na'\-rìng, ma tsmuk! Tsat o\-mum oel kin\-trr\-o.} Have you heard that Ralu has the hots for Tsenu? -- Old news, brother! I've known it for a week.~\LINK{http://naviteri.org/2021/08/ulte-ayyoratu-leiu-and-the-winners-are-2/}{b}

% zawsena 
\hi
\HT{zawsena}{\B{\cp{zaw}sena}} \I{["zaw.sE.na]} \emph{n., coll.} quiver (\emph{attached to the ik\-ran's saddle}), \emph{shortened form of} \ding{43}~\HL{swizawsena}{\B{swi\-zaw\-se\-na}}.

% zäfi
\hi
\HT{zAfi}{\B{\cp{zä}fi}} \I{["z\ae.fi]} \emph{adj.} modified, interfered with, no longer in a natural state, tamed, (\emph{of food}) processed, cooked.
\B{Oel yom tsngan\-it a\-zä\-fi nì\-'aw.} I only eat cooked meat.~\LINK{http://naviteri.org/2013/01/awvea-posti-zisita-amip-first-post-of-the-new-year/}{b}
\B{Tsa\-ik\-ran\-ìri ta\-lun\-a new nga\-ti tspi\-vang, law lu fwa mi ke lu zä\-fi.} Since that ikran wants to kill you, it's clear it's still not tame.~\LINK{http://naviteri.org/2013/01/awvea-posti-zisita-amip-first-post-of-the-new-year/}{b}
(\ding{214}~\HL{yrr}{\B{yrr}}) \\
\textsc{Derivatives}: \ding{43} 
\on{1}.~\HL{zAfizal}{\B{zä\-fi\-zal}}, fired clay, ceramic material.

% zäfizal
\hi 
\HT{zAfizal}{\B{\cp{zä}fizal}} \I{["z\ae.fi.zal]} \emph{n.} fired clay, ceramic material. 
\B{Fì\-tsngal\-it fkol txu\-la ta zä\-fi\-zal.} This is a ceramic cup.\slash This cup is made of fired clay\slash ceramic material.~\LINK{https://naviteri.org/2025/06/vospikin-lefpom-happy-july/}{b} 
(\ding{43}~\HL{zAfi}{\B{zä\-fi}}, \HL{zal}{\B{zal}}) 

% zärìp
\hi 
\HT{zArIp}{\B{\cp{zä}rìp}} \I{["z\ae.RIp\textcorner]} \emph{vtr.} (\emph{coll. form of}) pull. 
(\ding{43}~\HL{za'ArIp}{\B{za\-'ä\-rìp}})

% zeklor
\hi 
\HT{zeklor}{\B{zek\cp{lor}}} \I{[zEk\textcorner."loR]} \emph{adj.} pleasant to the touch. 
\B{Su\-nu oer fì\-fne\-srä ta\-lu\-na lu zek\-lor.} I like this kind of cloth because it feels good.~\LINK{https://naviteri.org/2024/06/ayhapxi-tokxa-si-ayliu-alahe-parts-of-the-body-and-more/}{b}
(\ding{43}~\HL{zekwA}{\B{zek\-wä}}, \HL{lor}{\B{lor}})

% zektil
\hi 
\HT{zektil}{\B{\cp{zek}til}} \I{["zEk\textcorner.til]} \emph{n.} knuckle. 
(\ding{43}~\HL{zekwA}{\B{zek\-wä}}, \HL{til}{\B{til}}) 

% zekwä 
\hi
\HT{zekwA}{\B{\cp{zek}wä}} \I{["zEk\textcorner.w\ae{}]} \emph{n.} finger. 
\B{Tsyokx\-ìri ke lu zek\-wä a\-tsìng ki a\-mrr.} The hand doesn't have four but five digits.~\LINK{http://naviteri.org/2010/07/ting-mikyun-fte-tslivam-listening-comprehension-1-3/}{b}
\B{Hol\-pxay ay\-zek\-wä fe\-yä.} The number of their fingers.~\LINK{http://forum.learnnavi.org/language-updates/talun-taweyk-tafral-lun-oeyk-holpxay-hol/}{ln} \\
(i.) \B{tìng zekwä}, \emph{vin.} (\texttt{+}\emph{control}) touch (intentionally). \B{Rä\-'ä tìng zek\-wä oer!} Don't touch me!~\LINK{http://naviteri.org/2015/04/some-new-words-for-may-day/}{b}
\B{Lu fì\-vul su\-syang nì\-txan. Txo fko ti\-vìng zek\-wä kxakx.} This branch is very fragile. If you touch it, it'll break.~\LINK{http://naviteri.org/2015/04/some-new-words-for-may-day/}{b} 
(\ding{43}~\HL{zIm}{\B{zìm}}, \HL{'ampi}{\B{'am\-pi}}) \\
\textsc{Derivatives}: \ding{43}
\on{1}.~\HL{venzek}{\B{ven\-zek}}, toe. 
\on{2}.~\HL{pxayzekwA}{\B{pxay\-zek\-wä}}, spiny whips.
\on{3}.~\HL{zeklor}{\B{zek\-lor}}, pleasant to the touch.
\on{4}.~\HL{zektil}{\B{zek\-til}}, knuckle.

% zel
\hi
\HT{zel}{\B{zel}} \I{[zEl]} \emph{conj.} then (\emph{for counterfactuals})
\B{Zun tom\-pa zì\-ye\-vup trr\-ay, zel fo srì\-ye\-vew.\slash … zel fo srew.} If it rained tomorrow, they'd do a dance.~\LINK{http://naviteri.org/2013/04/zun-zel-counterfactual-conditionals/}{b}
\B{Zun ay\-oe li\-vu tsam\-si\-yu, zel tsa\-kem ke sim\-vi.} If we were warriors, we wouldn't have done that.~\LINK{http://naviteri.org/2013/04/zun-zel-counterfactual-conditionals/}{b}
\B{Zel swey\-lu zun nga kil\-vä.} You should have gone [lit.: `it would have been better if you had gone'].~\LINK{https://naviteri.org/2024/12/odds-and-ends/}{b}
(\emph{can be left out in casual conversation if both clauses have the same tempus}) \B{Swey\-lu zun nga kil\-vä.} You should have gone [lit.: `it would have been better if you had gone'].~\LINK{https://naviteri.org/2024/12/odds-and-ends/}{b}
(\ding{43}~\HL{zun}{\B{zun}})

% zene
\hi
\HT{zene}{\B{\cp{ze}ne}} \I{["zE.nE]} \emph{vin., modal} must. 
\B{Slä nì\-aw\-no\-mum, ze\-ne oe 'aw\-si\-teng tì\-kang\-kem si\-vi fo\-hu.} But as you know, I must work together with them.~\LINK{http://forum.learnnavi.org/news-announcements/a-response-from-paul-frommer!/}{ln} 
\B{Nga ze\-ne ki\-vä.} You must go.~\LINK{http://naviteri.org/2011/04/\%E2\%80\%99a\%E2\%80\%99awa-li\%E2\%80\%99fyavi-amip\%E2\%80\%94a-few-new-expressions/}{b}
\B{Ke ze\-ne nga ki\-vä.} It's not necessary\slash obligatory that you go.~\LINK{http://naviteri.org/2011/04/\%E2\%80\%99a\%E2\%80\%99awa-li\%E2\%80\%99fyavi-amip\%E2\%80\%94a-few-new-expressions/}{b} \\
\textsc{Derivatives}: \ding{43}
\on{1}.~\HL{zenke}{\B{zen\-ke}}, must not. 
\on{2}.~\HL{kezemplltxe}{\B{ke\-zem\-pll\-txe}}, of course.
\on{3}.~\HL{nIzen}{\B{nì\-zen}}, necessarily.

% zenke
\hi
\HT{zenke}{\B{\cp{zen}ke}} \I{["z$\cdot$EN$\cdot$.kE]} \emph{vin., modal} must not. (!\emph{irregular infix distribution with second position infixes}! \B{\cp{ze}\-näng\-ke}, \B{\cp{ze}\-nei\-ke}, \B{\cp{ze}\-nu\-ye\-ke}, \B{\cp{ze}\-na\-tse\-ke} \I{["zE.n\ae{}N.kE, "zE.nE.i.kE, "zE.nu.jE.kE, "zE.na.\t{ts}E.kE]} \LINK{http://forum.learnnavi.org/language-updates/affective-obligation-\%28zenke\%29/}{ln})
\B{Nga zen\-ke ki\-vä.} You must not go.~\LINK{http://naviteri.org/2011/04/\%E2\%80\%99a\%E2\%80\%99awa-li\%E2\%80\%99fyavi-amip\%E2\%80\%94a-few-new-expressions/}{b}
\B{Hem nge\-yä zen\-ke fko\-ru li\-vu, tsay\-hem a nga\-ru prr\-te' ke lu.} Your actions must not be to others (as) those actions that you don't like.~\LINK{http://wiki.learnnavi.org/index.php?title=Corpus\#Good_Morning_America}{wiki}

% zeng
\hi 
\HT{zeng}{\B{zeng}} \I{[zEN]} \emph{n.} crest (\emph{on the head of a bird or animal}).

% zerok
\hi
\HT{zerok}{\B{\cp{ze}rok}} \I{["zE.Rok\textcorner]} \emph{vtr.} remember.
\B{Maw\-krr la\-ye\-iu oer krr nì\-'ul fte ngi\-vop ay\-lì\-'ut sì tsay\-fne\-sä\-num\-vit a tsun fra\-por srung si\-vi fte ni\-vu\-me sì zi\-ve\-rok nì\-swey.} After that I will have more time to create words and the kinds of lessons that can help everyone best learn and remember.~\LINK{http://forum.learnnavi.org/language-updates/trr-rrtaya/}{ln}
\B{Ze\-rok aw\-ngal nì\-wotx krr\-it a pol\-txe Tseyk san Eo ay\-oeng lu txan\-a tì\-kawng sìk.} We all remember the time that Jake said, `A great evil is upon us.'~\LINK{http://naviteri.org/2016/12/tengkrr-perahem-zisit-amip-as-the-new-year-arrives/}{b}
(\ding{43}~\HL{'ok}{\B{'ok}})

% zeswa
\hi
\HT{zeswa}{\B{\cp{ze}swa}} \I{["zE.swa]} \emph{n.} \ding{96} grass.
\B{Pa\-lu\-kan\-it tso\-le\-'a, ye\-rik lopx hi\-fwo kxam\-lä ze\-swa.} Spotting a thanator, the hexapede panicked and escaped through the grass.~\LINK{http://naviteri.org/2012/07/meetings-waterfalls-and-more/}{b} 
\B{Torukìl ke yom zeswat.} A toruk doesn't eat grass.~\LINK{http://naviteri.org/2019/06/50a-liu-amip-40-new-words}{b} \\
\textsc{Derivatives}: \ding{43}
\on{1}.~\HL{lezeswa}{\B{le\-ze\-swa}}, grassy. 
\on{2}.~\HL{zeswavi}{\B{ze\-swa\-vi}}, blad of grass. 
\on{3}.~\HL{yomzeswa}{\B{yom\-ze\-swa}}, graze.

% zeswavi
\hi
\HT{zeswavi}{\B{\cp{ze}swavi}} \I{["zE.swa.vi]} \emph{n.} blade of grass.
(\ding{43}~\HL{zeswa}{\B{ze\-swa}}, \HL{-vi}{\B{-vi}})

% zet 
\hi
\HT{zet}{\B{zet}} \I{[zEt\textcorner]} \emph{vtr.} treat (\emph{emotionally}), display an attitude towards (\emph{with} \ding{43}~\HL{pxel}{\B{pxel}}).
\B{Pey\-ral\-ìl zet wur\-a wu\-tsot a\-'aw\-nem pxel sngel.} Peyral treats a cool cooked meal like garbage.~\LINK{http://naviteri.org/2012/10/mipa-vospxi-mipa-ayliu-new-words-for-the-new-month/}{b}
\B{Pol zo\-let oe\-yä sä\-fpìl\-it pxel pum a tì\-ngä\-zìk\-it ngop.} He treated my idea as though it created a problem.~\LINK{http://naviteri.org/2012/10/mipa-vospxi-mipa-ayliu-new-words-for-the-new-month/}{b} 
\B{Pol ze\-ret oe\-ti pxel tu\-te a ke inan pot.} He's treating me like I don't know him.~\LINK{http://naviteri.org/2012/10/mipa-vospxi-mipa-ayliu-new-words-for-the-new-month/}{b}
\B{Va'\-rul zä\-nget ik\-ran\-it sn\-eyä pxel ha\-pxì\-tu so\-a\-i\-ä.} Va'ru treats his ikran like a member of the family.~\LINK{http://naviteri.org/2012/10/mipa-vospxi-mipa-ayliu-new-words-for-the-new-month/}{b}

% zey
\hi
\HT{zey}{\B{zey}} \I{[zEj]} \emph{adj.} special, distinct.   
\B{Kel\-ut\-ral lu fne\-ut\-ral a\-zey.} Hometree is a special kind of tree.~\LINK{http://naviteri.org/2012/06/spring-vocabulary-part-3/}{b}

% zeyko
\hi
\HT{zeyko}{\B{zey\cp{ko}}} \I{[zEj."k$\cdot$o]} \emph{vtr., inf.}\on{22} heal, fix.
\B{Ey\-wal zey\-ki\-vo ngat nì\-win.} May Eywa heal you quickly.~\LINK{http://naviteri.org/2010/07/vocabulary-update/}{b}
\B{'Ä'! Ira\-yo! Zey\-ko\-leio.} Oops! Thanks! Gladly fixed!~\LINK{http://naviteri.org/2014/07/tsawlultxamaw-a-ayliu-post-meetup-words/comment-page-1/\#comment-7689}{b}
\B{Nì\-ra\-ngal lir\-vu ay\-oe\-ru eyk\-tan a tsi\-vun srung si\-vi, pum a li\-vu por ay\-lì\-'u a\-zu\-sey\-ko.} I wish to have a leader who can help, one who has healing words.~\LINK{http://naviteri.org/2020/06/mi-tanlokxe-oeya-srr-afpxamo-terrible-days-in-my-country/}{b}
(\ding{43}~\HL{zo}{\B{zo}}) \\
\textsc{Derivatives}: \ding{43}
\on{1}.~\HL{vezeyko}{\B{ve\-zey\-ko}}, put in order, organize. 
\on{2}.~\HL{zeykoyu}{\B{zey\-ko\-yu}}, healer.

% zeykoyu
\hi 
\HT{zeykoyu}{\B{zey\cp{ko}yu}} \I{[zEj."ko.ju]} \emph{n.} healer. 
\B{Fì\-'um\-tsat to\-lìng 'e\-veng\-ur a\-spxin zey\-ko\-yul a txan\-ro\-'a.} This medicine was given to the sick child by a famous healer.~\LINK{http://naviteri.org/2019/06/50a-liu-amip-40-new-words/}{b}  
(\ding{43}~\HL{zeyko}{\B{zey\-ko}})

% zeykup
\hi
\HT{zeykup}{\B{zey\cp{kup}}} \I{[zEj."k$\cdot$up\textcorner]} \emph{or} \I{[."kUp\textcorner]} \emph{vtr., inf.}\on{22} (intentionally\slash deliberately) let fall.
\B{Nge\-yä tsko\-ti zey\-kup! Set!} Drop your bow! Now!~\LINK{http://naviteri.org/2013/04/wheres-the-bathroom-and-other-useful-things/}{b}
(\ding{43}~\HL{zup}{\B{zup}}, \HL{tungzup}{\B{tung\-zup}}) 

% zevakx
\hi 
\HT{zevakx}{\B{\cp{ze}vakx}} \I{["zE.vak']} \emph{adj., ofp.} cruel. 
\B{tu\-te a\-ze\-vakx}, cruel person.~\LINK{http://naviteri.org/2020/06/mi-tanlokxe-oeya-srr-afpxamo-terrible-days-in-my-country/}{b}
(\ding{43}~\HL{tIzevakxnga'}{\B{tì\-ze\-vakx\-nga'}}) \\
\textsc{Derivatives}: \ding{43}
\on{1}.~\HL{tIzevakx}{\B{tì\-ze\-vakx}}, cruelty. 
\on{2}.~\HL{nIzevakx}{\B{nì\-ze\-vakx}}, cruelly.

% zin
\hi 
\HT{zin}{\B{zin}} \I{[zin]} \emph{adj.} tangled. 
\B{Lä\-ngu fay\-hìng zin nì\-wotx; ke tsun sat si\-var.} Unfortunately, these threads are all tangled up; they can't be used.~\LINK{http://naviteri.org/2021/04/mipa-ayliu-mipa-sioeykting-new-words-new-explanations/}{b} (\ding{214}~\HL{kezin}{\B{ke\-zin}}) \\
\textsc{Derivatives}: \ding{43}
\on{1}.~\HL{tIzin}{\B{tì\-zin}}, tanglement. 
\on{2}.~\HL{kezin}{\B{ke\-zin}}, untangled.

% zir
\hi
\HT{zir}{\B{zir}} \I{[ziR]} \emph{n.} touch, feel, texture (\emph{sensation}).
\B{Fì\-ut\-ral\-ìri ta\-ngek\-ä zir fkan vawt, hu\-fwa ke rey.} The trunk of the tree feels solid, although it's dead.~\LINK{http://naviteri.org/2013/09/on-si-salewfya-shapes-and-directions/}{b} \\
\textsc{Derivatives}: \ding{43}
\on{1}.~\HL{kakzir}{\B{kak\-zir}}, numb. 

% zize'
\hi
\HT{zize'}{\B{zi\cp{ze'}}} \I{[zi."zEP]} \emph{n., fauna} hellfire wasp.
\B{Oe\-ri heyr tì\-sraw sä\-ngi ta\-lu\-na zi\-ze'\-ìl oet sngo\-lap tsa\-tseng.} My chest hurts because a hellfire wasp stung me there.~\LINK{http://naviteri.org/2015/03/kap-si-ayunil-saylahe-threats-dreams-and-other-things/}{b}
(\emph{a. idiom}) \B{sre fwa sngap zi\-ze'}, ``as quickly as possible'' [lit: before the hellfire wasp stings].~\LINK{http://naviteri.org/2012/07/fivospxiya-aylifyavi-amip-this-months-new-expressions-pt-2/}{b}

% zìm 
\hi
\HT{zIm}{\B{zìm}} \I{[zIm]} \emph{vtr.} feel (\texttt{-}\emph{control}).  
\B{Teng\-krr hu pa\-lu\-kan\-tsyìp uvan se\-ri zo\-lìm oel mì sa'\-leng a 'u\-ot a lu txa' sì ekx\-txu.} While playing with my cat I felt something hard and rough on his skin.~\LINK{http://naviteri.org/2012/11/renu-ayinanfyaya-the-senses-paradigm/}{b}
(\ding{43}~\HL{'ampi}{\B{'am\-pi}}, \HL{zekwA}{\B{tìng zek\-wä}})

% zìma'uyu
\hi
\HT{zIma'uyu}{\B{zì\cp{ma}'uyu}} \I{[zI."ma.Pu.ju]} \emph{n.} newcomer, s.o. who has just arrived on the scene.   
\B{Lam oer fwa tsa\-zì\-ma\-'u\-yul ke fnan tì\-tu\-sa\-ron\-it.} It seems to me that that newcomer isn't any good at hunting.~\LINK{http://naviteri.org/2011/09/\%E2\%80\%9Cby-the-way-what-are-you-reading\%E2\%80\%9D/}{b} 
(\ding{43}~\HL{za'u}{\B{za\-'u}})

% zìmtswo
\hi
\HT{zImtswo}{\B{\cp{zìm}tswo}} \I{["zIm.\t{ts}wo]} \emph{n.} sense of touch.
(\ding{43}~\HL{zIm}{\B{zìm}}, \HL{-tswo}{\B{-tswo}})

% zìsìkrr
\hi
\HT{zIsIkrr}{\B{\cp{zì}sìkrr}} \I{["zI.sI.k\s{r}]} \emph{n.} season
\B{Fra\-zì\-sì\-krr pay kil\-van\-ä nän nì\-'ul\-'ul.} Every season the river dries up more and more.~\LINK{http://naviteri.org/2012/02/trr-asawnung-lefpom-happy-leap-day/}{b}
\B{Tsay\-fa\-suk tsu\-ka\-nom lu krr\-ka fì\-zì\-sì\-krr nì\-'aw.} Those berries are available during this season only.~\LINK{http://naviteri.org/2021/03/mipa-ayliu-si-aylifyavi-new-words-and-expressions/}{b}
(\ding{43}~\HL{zIsIt}{\B{zì\-sìt}}, \HL{krr}{\B{krr}}) \\
\textsc{Derivatives}: \ding{43}
\on{1}.~\HL{zIskrrsom}{\B{zì\-skrr\-som}}, summer. 
\on{2}.~\HL{zIskrrwew}{\B{zì\-skrr\-wew}}, winter.
\on{3}.~\HL{zIskrrmipaw}{\B{zì\-skrr\-mi\-paw}}, spring (season).
\on{4}.~\HL{zIskrrtsawn}{\B{zì\-skrr\-tsawn}}, fall, autumn.

% zìsìt 
\hi
\HT{zIsIt}{\B{\cp{zì}sìt}} \I{["zI.sIt\textcorner]} \emph{n.} year.   
\B{Mip\-a Zì\-sìt Le\-fpom, ma fra\-po!} Happy New Year, everyone!~\LINK{http://naviteri.org/2011/01/new-year-new-vocabulary/}{b}
\B{Sìl\-pey oe tsnì fì\-zì\-sìt\-ìl ay\-nga\-ru za\-mì\-ye\-vu\-nge txan\-a fpom\-it sì fpom\-tokx\-it.} I hope this year will bring you much happiness and health.~\LINK{http://naviteri.org/2013/12/mipa-zisit-lefpom-happy-new-year/}{b}
\B{Sì ay\-zì\-sìt\-ä ka\-to | Sì ekong te'\-lan\-ä.} And the rhyhtm of the years | And the beats of the hearts.~\LINK{http://wiki.learnnavi.org/index.php/Weaving_Song}{wiki}
\B{Zì\-sìt\-o a\-vol ke tsä\-po\-le\-'a fo fì\-tsap.} They haven't seen each other in eight years.~\LINK{http://naviteri.org/2011/12/one-more-for-2011/}{b}
\B{Nge\-yä tì\-eyk\-tan\-ìri, tì\-slan\-ìri sì tsran\-ten fra\-to a tì\-'ey\-lan\-ìri a ka ay\-zì\-sìt nì\-wotx, sei\-yi oe ira\-yo nì\-txan.} Thank you so much for your leadership, your support, and most importantly your friendship throughout the years.~\LINK{http://naviteri.org/2012/03/spring-vocabulary-part-1/}{b}
(\emph{a. used for stating age}) 
\B{Nga\-ri so\-la\-lew pol\-pxay\-a zì\-sìt?\slash Nga\-ri so\-la\-lew zì\-sìt a\-pol\-pxay?} How old are you?~\LINK{http://naviteri.org/2010/09/getting-to-know-you-part-2/}{b}
\B{Oe\-ri so\-la\-lew (so\-lew) zì\-sìt a\-pxe\-vol.} I'm 24 years old.~\LINK{http://naviteri.org/2010/09/getting-to-know-you-part-2/}{b} \\
(i.) \B{zìsì\cp{tam}} \I{[zI.sI"tam]} \emph{adv.} last year. \\
(ii.) \B{zìsì\cp{tay}} \I{[zI.sI."taj]} \emph{adv.} next year. \B{Nì\-sìl\-pey zì\-sìt\-ay!} Hopefully next year.~\LINK{http://naviteri.org/2011/07/txantsana-ultxa-mi-siatll-great-meeting-in-seattle/}{b} \\
\textsc{Derivatives}: \ding{43}
\on{1}.~\HL{zIsIkrr}{\B{zì\-sì\-krr}}, season. 
\on{2}.~\HL{zIsItsaltrr}{\B{zì\-sìt\-sal\-trr}}, (yearly) anniversary.
\on{3}.~\HL{zIsItnun}{\B{zì\-sìt\-nun}}, leap year.

% zìsìtnun
\hi 
\HT{zIsItnun}{\B{zìsìt\cp{nun}}} \I{[zI.sIt\textcorner."nun]} \emph{or} \I{[."nUn]} (\textsc{rn}: \B{zìsìt\cp{nun}} \I{[zI.sIt\textcorner."nun]}) \emph{n.} leap year. 
(\ding{43}~\HL{zIsIt}{\B{zìsìt}}, \HL{nun}{\B{nun}}) \\
\textsc{Derivatives}: \ding{43}
\on{1}.~\HL{zIsItnuntrr}{\B{zì\-sìt\-nun\-trr}}, leap year day.

% zìsìtnuntrr
\hi 
\HT{zIsItnuntrr}{\B{zìsìt\cp{nun}trr}} \I{[zI.sIt\textcorner."nun.t\s{r}]} \emph{or} \I{[."nUn.]} (\textsc{rn}: \B{zìsìt\cp{nun}trr} \I{[zI.sIt\textcorner."nun.t\s{r}]}) \emph{n.} leap year day. 
\B{Zì\-sìt\-nun\-trr Le\-fpom.} Happy Leap Year Day!~\LINK{https://naviteri.org/2024/02/mipa-ayliu-si-aylifyavi-niul-more-new-words-and-expressions/}{b}
(\ding{43}~\HL{zIsItnun}{\B{zì\-sìt\-nun}}, \HL{trr}{\B{trr}})

% zìsìtsaltrr
\hi 
\HT{zIsItsaltrr}{\B{zìsìt\cp{sal}trr}} \I{[zI.sIt\textcorner."{}sal.t\s{r}]} \emph{n.} (yearly) anniversary; (\emph{coll. shortened to} \ding{43}~\HL{zItsaltrr}{\B{zìt\-sal\-trr}}). 
(\ding{43}~\HL{zIsIt}{\B{zì\-sìt}}, \HL{salew}{\B{sa\-lew}}, \HL{trr}{\B{trr}})

% zìskrrmipaw
\hi 
\HT{zIskrrmipaw}{\B{zìskrr\cp{mi}paw}} \I{[zI.sk\s{r}."mi.paw]} \emph{n.} spring (season of new growth). 
(\ding{43}~\HL{zIsIkrr}{\B{zì\-sì\-krr}}, \HL{mip}{\B{mip}}, \HL{paw}{\B{paw}}; \ding{214}~\HL{zIskrrtsawn}{\B{zì\-skrr\-tsawn}})

% zìskrrsom
\hi 
\HT{zIskrrsom}{\B{zìskrr\cp{som}}} \I{[zI.sk\s{r}."{}som]} \emph{n.} summer (\emph{hot season}). 
\B{Teng\-krr le\-rok zì\-skrr\-som\-ìl tì\-'i'\-at, sìl\-pey oe, ay\-nga\-ri te'\-lan li\-vu le\-fpom ul\-te tì\-rey zi\-vaw\-prr\-te'.} While the summer nears its end, I hope your hearts are well and you enjoy life.~\LINK{http://naviteri.org/2019/08/aawa-lifya-si-lifyavi-amip-a-few-new-words-and-expressions/}{b}
(\ding{43}~\HL{zIsIkrr}{\B{zì\-sì\-krr}}, \HL{som}{\B{som}}; \ding{214}~\HL{zIskrrwew}{\B{zì\-skrr\-wew}})

% zìskrrtsawn
\hi 
\HT{zIskrrtsawn}{\B{zìskrr\cp{tsawn}}} \I{[zI.sk\s{r}."\t{ts}awn]} \emph{n.} fall, autumn. 
\B{Sìl\-pey oe, ay\-nga\-ri zì\-skrr\-tsawn sir\-va\-lew pxay\-a lrr\-tok\-hu a ta Ey\-wa.} I hope that the fall season proceeds with many smiles from Eywa for you all.~\LINK{http://naviteri.org/2018/10/aysrr-ayvospxi-ayzisikrr-days-months-seasons/}{b}
(\ding{43}~\HL{zIsIkrr}{\B{zì\-sì\-krr}}, \HL{tsawn}{\B{tsawn}}; \ding{214}~\HL{zIskrrmipaw}{\B{zì\-skrr\-mi\-paw}})

% zìskrrwew
\hi 
\HT{zIskrrwew}{\B{zìskrr\cp{wew}}} \I{[zI.sk\s{r}."wEw]} \emph{n.} winter (\emph{cold season}).  
(\ding{43}~\HL{zIsIkrr}{\B{zì\-sì\-krr}}, \HL{wew}{\B{wew}}; \ding{214}~\HL{zIskrrsom}{\B{zì\-skrr\-som}})

% zìtsaltrr
\hi 
\HT{zItsaltrr}{\B{zìt\cp{sal}trr}} \I{[zIt\textcorner."{}sal.t\s{r}]} \emph{or} \I{[zI."\t{ts}al.t\s{r}]} \emph{n.} (yearly) anniversary; (\emph{coll. form of} \ding{43}~\HL{zIsItsaltrr}{\B{zì\-sìt\-sal\-trr}}). 
\B{Zì\-tsal\-trr\-ìri tì\-mun\-txa\-yä ay\-lrr\-tok!} Happy anniversary (of your marriage)!~\LINK{http://naviteri.org/2017/12/zisit-amip-lefpom-ma-eylan-happy-new-year-friends/}{b} 

% zo
\hi
\HT{zo}{\B{zo}} \I{[zo]} \emph{vin.} be well, be intact, be as it should be, work correctly or as nature intended.
\B{Nga zo srak?} Are you well? (\emph{e.g. have you recovered from your illness\slash fall, etc.?})~\LINK{http://naviteri.org/2010/07/vocabulary-update/}{b}
(\emph{a. conversational expressions}) \B{'Uo ke zo srak?\slash 'Uk zo srak?} `Is something wrong?'~\LINK{http://naviteri.org/2020/07/vospxivol-lefpom-happy-august/}{b}
\B{'U\-pe ke zo?} `What’s the matter?\slash What's wrong?'~\LINK{http://naviteri.org/2020/07/vospxivol-lefpom-happy-august/}{b}  \\
\textsc{Derivatives}: \ding{43}
\on{1}.~\HL{vezo}{\B{ve\-zo}}, be in order, be organized. 
\on{2}.~\HL{zoslu}{\B{zo\-slu}}, heal, become well, get fixed. 
\on{3}.~\HL{frawzo}{\B{fraw\-zo}}, all is well. 
\on{4}.~\HL{zeyko}{\B{zey\-ko}}, heal, fix. 

% zong
\hi
\HT{zong}{\B{zong}} \I{[zoN]} \emph{vtr.} save, defend. 
% \B{To\-ruk nì\-'aw tsun srung si\-vi fte kel\-ut\-ral\-it zi\-vong.} Only the Last Shadow can help to save Hometree. 
\B{Te\-ya si oer fwa ngal paw\-nu\-nga ay\-ioang\-it zong.} It moves me that you save injured animals.~\LINK{http://naviteri.org/2021/09/aawa-ayliu-amip-a-few-new-words/}{b} 
(\emph{a. idiom}) \B{Mak\-to zong.} ``Take care. Travel safely.'' (\emph{coll. of} \B{Mi\-vak\-to nì\-zaw\-nong})~\LINK{http://forum.learnnavi.org/language-updates/makto-zong!/}{ln} 
(\emph{b. short form of} \HL{nIzawnong}{\B{nì\-zaw\-nong}}) \emph{to answer the question} \B{Mak\-to fya\-pe?} ``How are things? How's everything going? How are you doing?'' \B{Zong.} ``Well.''~\LINK{http://naviteri.org/2020/05/tipusawm-tiuseyng-si-okvur-a-eltur-titxen-si-asking-answering-and-an-interesting-story/}{b} 
(\ding{43}~\HL{tarep}{\B{ta\-rep}}, \HL{'avun}{\B{'a\-vun}}) \\
\textsc{Derivatives}: \ding{43}
\on{1}.~\HL{zongtseng}{\B{zong\-tseng}}, safe place, refuge.
\on{2}.~\HL{nIzawnong}{\B{nì\-zaw\-nong}}, safely.

% zongtseng
\hi
\HT{zongtseng}{\B{\cp{zong}tseng}} \I{["zoN.\t{ts}EN]} \emph{n.} safe place, refuge. 
(\ding{43}~\HL{zong}{\B{zong}}, \HL{tseng}{\B{tse\-nge}})

% zoplo / si
\hi
\HT{zoplo}{\B{zop\cp{lo}}} \I{[zop\textcorner."lo]} \textsc{i}. \emph{n.} offense, insult. 
\B{Tì\-'al lu zop\-lo a tsa\-ri ke tsun txoa li\-vu nì\-ftue.} Wastefulness is an offense that cannot easily be forgiven.~\LINK{http://naviteri.org/2014/07/tsawlultxamaw-a-ayliu-post-meetup-words/}{b} \\
\textsc{ii}. \B{zop\cp{lo} si} \I{[zop\textcorner."lo s$\cdot$i]} \emph{vin.} offend, insult (s.o.)
\B{Pol\-txe oe san ze\-ne kea unil\-tì\-ran\-yu ke zi\-va\-'u fì\-tseng fte po kel\-kur zop\-lo si\-vi sìk.} I have said no dreamwalker will come here to offend our home.~\LINK{http://forum.learnnavi.org/language-updates/little-confirmation-on-teach-and-insult/}{ln}

% zoslu
\hi
\HT{zoslu}{\B{\cp{zo}slu}} \I{["zo.sl$\cdot$u]} \emph{vin., inf.}\on{22} heal, become well, get fixed.   
\B{Fu\-r\-i\-a 'll\-ngo nge\-yä zo\-slu, 'e\-fu oe nit\-ram nì\-'aw.} I'm just happy that your hip is getting better.~\LINK{http://naviteri.org/2013/07/pximaw-tsawlultxa-just-after-the-meet-up/comment-page-1/\#comment-2484}{b}
\B{Oe\-ri nì\-'i'a tsyokx zo\-slo\-lu.} My hand is finally healed.~\LINK{http://naviteri.org/2010/07/vocabulary-update/}{b}
\B{Tì\-kxey zo\-slo\-lu.} The incorrectness has been fixed.~\LINK{http://naviteri.org/2013/12/mipa-zisit-lefpom-happy-new-year/comment-page-1/\#comment-2585}{b}
\B{Tsa\-tì\-ke\-nong\-ìri oel kxey\-eyt zo\-sley\-ko\-lu.} As for that example, I got the mistake fixed.~\LINK{http://naviteri.org/2013/01/awvea-posti-zisita-amip-first-post-of-the-new-year/comment-page-1/\#comment-1903}{b}
(\ding{43}~\HL{zo}{\B{zo}}, \HL{slu}{\B{slu}}, \HL{zeyko}{\B{zey\-ko}})

% zum
\hi
\HT{zum}{\B{zum}} \I{[zum]} \emph{or} \I{[zUm]} (\textsc{rn}: \B{zùm} \I{[zUm]}) \emph{n.} object, thing (\emph{physical or tangible}).
\B{Fay\-zum lu pe\-u?} What are these things?~\LINK{http://naviteri.org/2015/11/vomuna-liu-amip-ten-new-words/}{b}
(\ding{43}~\HL{'u}{\B{'u}})

% zun
\hi
\HT{zun}{\B{zun}} \I{[zun]} \emph{or} \I{[zUn]} \emph{conj.} if (\emph{for counterfactuals})
\B{Zun tom\-pa zì\-ye\-vup trr\-ay, zel fo srì\-ye\-vew.\slash … zel fo srew.} If it rained tomorrow, they'd do a dance.~\LINK{http://naviteri.org/2013/04/zun-zel-counterfactual-conditionals/}{b}
\B{Zun ay\-oe li\-vu tsam\-si\-yu, zel tsa\-kem ke sim\-vi.} If we were warriors, we wouldn't have done that.~\LINK{http://naviteri.org/2013/04/zun-zel-counterfactual-conditionals/}{b}
\emph{a. as counterfactual}) `should have'. 
\B{(Zel) swey\-lu zun nga kil\-vä.} You should have gone [lit.: `it would have been better if you had gone']~\LINK{https://naviteri.org/2024/12/odds-and-ends/}{b}
(\ding{43}~\HL{zel}{\B{zel}})

% zung
\hi 
\HT{zung}{\B{zung}} \I{[zuN]} \emph{or} \I{[zUN]} (\textsc{rn}: \B{zùng} \I{[zUN]}) \emph{vin.} crouch. 
\B{Zo\-lung ay\-oe ne\-kll fte\-ke ay\-ioang tsi\-vun ay\-oe\-ti tsi\-ve\-'a.} We crouched down so that the animals wouldn't be able to see us.~\LINK{http://naviteri.org/2019/06/50a-liu-amip-40-new-words/}{b}

% zup 
\hi
\HT{zup}{\B{zup}} \I{[zup\textcorner]} \emph{or} \I{[zUp\textcorner]} \emph{vin.} (\emph{a.}) fall. 
\B{Po pxun\-fa pxek tsa\-krr me\-fo zup ne\-kll!} He shoved the two of them and they fell down.~\LINK{http://naviteri.org/2014/05/mipa-ayliu-mipa-aysafpil-new-words-new-ideas/}{b}
(\emph{b. of the weather, precipitation}) \B{Zup tompa\slash her\-wì.} It's raining\slash snowing.~\LINK{http://naviteri.org/2012/10/audio-and-video-learning-materials-for-navi-101/}{b}
\B{Oe\-ri ay\-som\-pì\-va sìn re\-'o var zi\-vup.} Raindrops keep falling on my head.~\LINK{http://naviteri.org/2011/04/yafkeykiri-plltxe-frapo-everyone-talks-about-the-weather/}{b}
(\emph{c. with} \ding{43}~\HL{ne}{\B{ne}}) perish. 
\B{Ke za\-syup lì\-'O\-na ne kxu\-tu a mì\-fa fu a wrr\-pa.} The l'Ona will not perish to the enemy within or the enemy without.~\LINK{http://forum.learnnavi.org/language-updates/ltasygt-with-ke-and-nga/}{g\slash ln} 
(\emph{d. proverb}) \B{Rä\-'ä zup nem\-fa tsup.} \emph{don't fall into the trap that s.o. has laid for you} [lit.: ``don't fall into the chasm''].~\LINK{https://naviteri.org/2025/04/mevosinga-liu-amip-twenty-new-words/}{b} \\
\textsc{Derivatives}: \ding{43}
\on{1}.~\HL{zeykup}{\B{zey\-kup}}, let fall (intentionally). 
\on{2}.~\HL{tungzup}{\B{tung\-zup}}, drop, let fall (unintentionally).

% zusawkrr
\hi
\HT{zusawkrr}{\B{zu\cp{saw}krr}} \I{[zu."saw.k\s{r}]} \textsc{i}. \emph{n.} future time, future.
\B{Nì\-sung, txo tsa\-kem li\-ven mì zu\-saw\-krr, ze\-ne fko si\-var kem\-lì\-'u\-vit a\-lu ‹ìyev›.} Furthermore, if that happens in the future, one must use the infix ‹\B{ìyev}›.~\LINK{http://naviteri.org/2013/01/awvea-posti-zisita-amip-first-post-of-the-new-year/comment-page-1/\#comment-1910}{b}
\B{Sìn A\-sok, Sìn Zu\-saw\-krr\-ä}, Recent and Upcoming Activities [lit.: recent activities, future's activities].~\LINK{http://naviteri.org/2013/04/sin-asok-sin-zusawkrra-recent-and-upcoming-activities/}{b}  
\B{Olo'\-ìri po\-an zu\-saw\-krr\-ur 'on so\-li.} He shaped the tribe's future.~\LINK{http://naviteri.org/2018/07/ta-sulfatu-a-ayliu-niul-more-words-from-our-experts/}{b} \\
\textsc{ii}. \emph{adv.} in the future.
\B{'Aw\-a tì\-pawm\-ì\-ri 'i\-veyng oe set;  ay\-la\-ri zu\-saw\-krr 'ay\-eyng.} Let me answer one of the questions now; the others (I'll) answer in the future.~\LINK{http://forum.learnnavi.org/language-updates/email-from-frommer-re-numbers-\%28or-the-full-number-chart!\%29/}{ln}
(\ding{43}~\HL{za'u}{\B{za\-'u}}, \HL{krr}{\B{krr}})
\\
\end{multicols}

 
% ===== Changelog =====

 \clearpage
\pdfbookmark[0]{Change Log}{Change Log}
\begin{center}
\section*{Change Log}
\end{center}

\begin{footnotesize}
\begin{itemize}

%  \HL{}{\B{}}, \HL{}{\B{}}, \HL{}{\B{}}, \HL{}{\B{}}, \HL{}{\B{}}, \HL{}{\B{}}

\item[2024-07-11] added an example to \HL{fpIl}{\B{fpìl}} and \HL{sIlpey}{\B{sìl\-pey}} that shows usage without a patientive subordinate clause (\emph{mesari irayo, ma Tekre!}); changed the stress in \HL{leftxozA}{\B{le\-ftxo\-zä}} as per \LINK{https://forum.learnnavi.org/language-updates/stress-of-leftxoza/}{this} confirmation, fixed the stress in \HL{fya'o}{\B{fya\-'o}} (\emph{tsari irayo ma Wllìm}); added new words from \LINK{https://naviteri.org/2024/02/mipa-ayliu-si-aylifyavi-niul-more-new-words-and-expressions/}{More new words and expressions}, \LINK{https://naviteri.org/2024/03/fleltrra-ayliu-words-for-april-fools-day/}{Flel\-trrä ay\-lì\-'u}, altered the entry of \HL{sko}{\B{sko}}, added new words from \LINK{https://naviteri.org/2024/05/krrteri-about-time/}{Krr\-te\-ri}, and \LINK{https://naviteri.org/2024/06/ayhapxi-tokxa-si-ayliu-alahe-parts-of-the-body-and-more/}{Parts of the body …}; added examples to \HL{mawl}{\B{mawl}}, \HL{nga'}{\B{nga'}}, \HL{tI'i'a}{\B{tì'i'a}}, \HL{sre}{\B{sre}}, \HL{srekamtrr}{\B{sre\-kam\-trr}}, \HL{maw}{\B{maw}}, \HL{kay}{\B{kay}}, \HL{fa}{\B{fa}}, \HL{slu}{\B{slu}}, \HL{pil}{\B{pil}}, \HL{takuk}{\B{ta\-kuk}}; Reef Na'vi variants of \B{u}\slash \B{ù} still not yet completed!

\item[2024-01-06] added an example to \HL{alaksi}{\B{a\-lak\-si}} that shows usage with \emph{lu}; changed the entry for \HL{aylahe}{\B{ay\-la\-he}} due to the confirmation in \LINK{http://naviteri.org/2023/05/reef-navi-part-2-morphology-syntax-lexicon-and-more/}{this blog} post, and added the new vocab items from that post; added new items from the \LINK{http://naviteri.org/2023/07/trr-tsyimawnuniya-lefpom-happy-independence-day/}{Happy Independence Day} blog post; started including the Reef Na'vi variants of \B{u}\slash \B{ù} as per \LINK{http://naviteri.org/2023/07/the-complete-u-u-list-for-reef-navi/}{this} confirming post (\textsc{not} yet completed!); added new items from the \LINK{http://naviteri.org/2023/09/aawa-ayliu-si-aylifyavi-amip-a-few-new-words-and-expressions/}{A few new words} blog post; corrected a mistake in the \HL{fu1}{\B{fu}} entry concerning choice statements vs. choice questions (\emph{tsari irayo nìtxan, ma Neytiri}); added vocab from the \LINK{http://naviteri.org/2023/12/mipa-zisit-ayliu-amip-new-year-new-words/}{New Year, New Words} blog post; various edits behind the scenes (\emph{sari irayo nìtxan, ma Wllìm}); corrected the stress on \HL{tsu'tey}{\B{Tsu'tey}} (\emph{tsari irayo, ma Zángtsuva}); added examples to \HL{vospxIkin}{\B{Vospxìkin}}, \HL{mu'ni}{\B{mu\-'ni}}, \HL{pIsaw}{\B{pì\-saw}}, \HL{sAmok}{\B{sä\-mok}}, \HL{tIran}{\B{tìran}}, \HL{tul}{\B{tul}}, \HL{yI}{\B{yì}}, \HL{tsan'ul}{\B{tsan\-'ul}}, \HL{spe'e}{\B{spe\-'e}}, \HL{sAlatem}{\B{sä\-la\-tem}}, \HL{ukxo}{\B{ukxo}}, \HL{hufwetsyIp}{\B{hu\-fwe\-tsyìp}}, \HL{AtxAle}{\B{ätxäle}}, \HL{palukan}{\B{palukan}}, \HL{zawprrte'}{\B{zaw\-prr\-te'}}.

\item[2023-02-06] corrected some real \emph{skxawng} mistakes in the IPA section of the description of Reef Na'vi (\emph{tsari irayo ngaru ma Zángtsuva!}); incorporated new information about the Reef Na'vi dialect from \LINK{http://naviteri.org/2023/01/2653/}{Clarifications and Announcements} and \LINK{https://youtu.be/Rb_ZF9ZwzAM}{OmatiCon \on{2023}}, including information on the known words that are affected by tense and lax u sounds, i.e. \HL{txula}{\B{txu\-la}}; added a long overlooked example of \HL{sempu}{\B{sem\-pu}}, and \HL{nulkrr}{\B{nul\-krr}}. 

\item[2023-01-10] added examples to \HL{meyp}{\B{meyp}}, \HL{'om}{\B{'om}}, \HL{mikyun}{\B{mik\-yun}}, continued formatting adjustment up to section \B{N}; added new vocabulary items from \LINK{http://naviteri.org/2022/12/neytiriya-waytelem-neytiris-song-cord/}{Neytiriyä Waytelem} and incorporated the new information about Reef Na'vi from \LINK{http://naviteri.org/2023/01/reef-navi-part-1-phonetics-and-phonology/}{Reef Na'vi part 1}: updated the IPA section, Abbreviations, and added information on the known words that are affected by tense and lax u sounds, i.e. \HL{tutan}{\B{tutan}}, \HL{run}{\B{run}}, \HL{ftu}{\B{ftu}}, \HL{hum}{\B{hum}}, \HL{txukx}{\B{(nì)txukx}}, \HL{tsun}{\B{tsun}}, \HL{pum}{\B{pum}}, \HL{'u}{\B{'u}}, and \HL{fra'u}{\B{fra'u}}.

\item[2022-12-21] added a clarifying example with \B{pxel} to \HL{nIftxan}{\B{nìftxan}} and \HL{pxel}{\B{pxel}} (\emph{irayo} to Wllìm for making me aware of that); added new vocabulary items from \LINK{naviteri.org/2022/11/krr-ao-an-exciting-time/}{Krr a'o'!} and \LINK{http://naviteri.org/2022/12/trr-anawm-polaheiem-the-great-day-has-arrived/}{Trr anawm poläheiem!}; added examples to \HL{payoang}{\B{pay\-oang}}, \HL{ve'kI}{\B{ve'\-kì}}, \HL{tson}{\B{tson}}, \HL{tIpe'un}{\B{tì\-pe\-'un}}, \HL{kawtseng}{\B{kaw\-tseng}}, \HL{tIsraw}{\B{tì\-sraw sey\-ki}}.

\item[2022-07-03] added new vocabulary items from \LINK{http://naviteri.org/2022/06/solalew-maul-zisita-half-the-year-is-over/}{Solalew Mawl Zìsìtä}; several errors corrected behind the scenes that affected the hyper references (a massive \emph{irayo} to tsmukan Wllìm); added \HL{tulkun}{\B{tulkun}} and \HL{tsurak}{\B{tsurak}};

\item[2022-06-11] \HL{lawnoltsim}{\B{lawnoltsim}} and \HL{txurtseng}{\B{txurtseng}} from the latest \LINK{http://naviteri.org/2022/05/results-of-the-tttc/}{blog post} added; added examples to \HL{kintrr}{\B{kin\-trr}}, \HL{lefkrr}{\B{le\-fkrr}}, \HL{tIlen}{\B{tì\-len}}, and \HL{fpxamo}{\B{fpxa\-mo}}; minor changes pertaining to style and references.

\item[2022-03-04] transitivity of \HL{nong}{\B{nong}} and \HL{kenong}{\B{ke\-nong}} cleared up, added examples to those; added \HL{txopunil}{\B{txo\-pu\-nil}}; corrected an error in the \HL{yak}{\B{yak}} entry (\emph{irayo ngaru, ma wm.annis}); added new blog entry \LINK{http://naviteri.org/2022/02/lifyengteri-concerning-honorific-language/}{Lì'fyengteri}.

\item[2022-01-03] added the vocab, idioms and proverbs of the latest blog posts (\LINK{http://naviteri.org/2021/09/aawa-ayliu-amip-a-few-new-words/}{'A'awa aylì'u}, \LINK{http://naviteri.org/2021/08/ulte-ayyoratu-leiu-and-the-winners-are-2/}{Ulte ayyora'tu leiu …} and \LINK{http://naviteri.org/2021/12/zolau-niprrte-ma-3746-welcome-2022/}{Zola'u nìprrte' …}), finally cleared up unknown transitivity as per this bit of \LINK{https://forum.learnnavi.org/language-updates/transitivity-clarification-of-all-remaining-v-verbs/}{news}; added a long overlooked example of \HL{ekxan}{\B{ekxan}}, added an example to \HL{virA}{\B{virä}}; refined the definition of \HL{lIng}{\B{lìng}}; added examples to \HL{fnawe'tu}{\B{fna\-we'\-tu}}, \HL{soaia}{\B{soaia}}, \HL{fkxile}{\B{fkxi\-le}}, \HL{yayotsrul}{\B{ya\-yo\-tsrul}}, \HL{fpom}{\B{fpom}}, \HL{layon}{\B{la\-yon}}, \HL{atxkxe}{\B{atx\-kxe}}, \HL{spA}{\B{spä}}, \HL{tse'a}{\B{tse\-'a}}, \HL{yawne}{\B{yaw\-ne}}, \HL{sAlang}{\B{sä\-lang}}, \HL{tIflA}{\B{tì\-flä}}, and \HL{krr}{\B{krr}}.

\item[2021-05-09] added new vocabulary from \LINK{http://naviteri.org/2021/04/mipa-ayliu-mipa-sioeykting-new-words-new-explanations/}{Mipa Aylì'u, Mipa sìoeyktìng}; corrected mistakes in the IPA of \HL{fe'ran}{\B{fe'\-ran}}, and \HL{fe'ranvi}{\B{fe'\-ran\-vi}}, removed erroneous ``\emph{ofp.}'' designation from and added an example to \HL{venga'}{\B{ve\-nga'}} as per \LINK{http://naviteri.org/2021/04/mipa-ayliu-mipa-sioeykting-new-words-new-explanations/\#comment-33505}{this} confirmation; added examples to \HL{pImtxan}{\B{pìm\-txan}}, \HL{lu}{\B{lu}}, \HL{pum}{\B{pum}}, \HL{kIng}{\B{kìng}}, \HL{txurtel}{\B{txurtel}}, \HL{ingyentsim}{\B{ingyentsim}}, \HL{'efu}{\B{'efu}}, \HL{zenke}{\B{zenke}}, \HL{fImuntrr}{\B{fìmuntrr}}, \HL{way}{\B{way}}, \HL{snonivi}{\B{snonivi}}, and \HL{keltsun}{\B{keltsun}}.

% \item[2021-04-05] added an example of topical inalienable possession in \HL{tIrey}{\B{tì\-rey}}, an example of \HL{sunu}{\B{su\-nu}} referring to people, and fixed a broken link in \HL{tAftxuyu}{\B{tä\-ftxu\-yu}} (\emph{tsari irayo ma wm.annis}); added bookmarks for all the lexical sections; changed the example in \HL{fpap}{\B{fpap}} as per \LINK{http://naviteri.org/2020/12/mrra-tipangkxotsyip-five-little-discussions/}{this correction} by Karyu Pawl; finally added \HL{fray}{\B{fray}}\texttt{+} (\emph{totally my mistake}); added an example for \HL{irayo}{\B{ira\-yo}}; added \HL{kot}{\B{kot}} from the OmatiCon event on Jan 02; corrected a few errors pertaining the patientive ending after \B{ey} and errors in the IPA of \emph{vit\-ra} words; added \HL{vitronvA'}{\B{vit\-ron\-vä'}} as the fuller form of \HL{vonvA'}{\B{von\-vä'}}; added the transitivity of \HL{txey}{\B{txey}} as per \LINK{https://forum.learnnavi.org/language-updates/transitivity-of-txey/}{this} update; added an example to \HL{fpIl}{\B{fpìl}} that uses \B{san … sìk}; included new vocabulary from \LINK{http://naviteri.org/2021/03/mipa-ayliu-si-aylifyavi-new-words-and-expressions/}{Mipa aylì'u sì aylì'fyavi}; added examples to \HL{tstalsena}{\B{tstal\-se\-na}}, \HL{txe'lan}{\B{txe'\-lan}}, \HL{txan}{\B{txan}}, \HL{pasuk}{\B{pa\-suk}}, \HL{zIsIkrr}{\B{zì\-sì\-krr}}, and \HL{'umtsa}{\B{'um\-tsa}}; continued adjusting formatting to make the file more consistent up to section \B{M}.

% \item[2020-11-25] added new vocabulary from \LINK{http://naviteri.org/2020/08/hiia-vur-a-teri-mefalukantsyip-a-little-story-about-two-cats/}{Hì'ia vur} and \LINK{http://naviteri.org/2020/11/vospxivopeya-ayliu-amip-novembers-new-words/}{Vospxìvopeyä Aylì'u Amip}; added \HL{paytsim}{\B{pay\-tsim}}, added examples to \HL{lekye'ung}{\B{lek\-ye\-'ung}}, \HL{muiA}{\B{muiä}}, \HL{oeyktIng}{\B{oeyktìng}}, \HL{'u}{\B{'u}}, \HL{tIyawn}{\B{tì\-yawn}}, \HL{hasey si}{\B{hasey si}}, \HL{kar}{\B{kar}}, \HL{tstew}{\B{tstew}}, \HL{kxemyo}{\B{kxem\-yo}}, \HL{srAtx}{\B{srätx}}, and \HL{sAsrAtx}{\B{sä\-srätx}} edited the definition of \HL{slAkop}{\B{slä\-kop}}, fixed a mistake in the ‹\B{ary}› entry (\emph{irayo ma Neytiri!}); began adjusting formatting to make it more consistent up to section \B{I}.

% \item[2020-08-06] added new vocabulary from the blog post \LINK{http://naviteri.org/2020/07/vospxivol-lefpom-happy-august/}{Vospxìvol Lefpom}; confirmation of transitivity of \HL{fkxake}{\B{fkxa\-ke}}, added examples to \HL{yune}{\B{yu\-ne}}, \HL{tIleymkem}{\B{tì\-leym\-kem}}, \HL{yo}{\B{yo}}, \HL{txal}{\B{txal}}, \HL{terkup}{\B{ter\-kup}}, \HL{nIftxan}{\B{nì\-ftxan}}, \HL{kuma}{\B{kuma}}, \HL{tsin}{\B{tsin}}, \HL{ne}{\B{ne}}, \HL{tong}{\B{tong}}, \HL{zo}{\B{zo}}, and \HL{fra'u}{\B{fra'u}}.

% \item[2020-06-04] added new vocabulary from the blog post \LINK{http://naviteri.org/2020/06/mi-tanlokxe-oeya-srr-afpxamo-terrible-days-in-my-country/}{Mì Tanlokxe oeyä \ldots} added examples to \HL{koren}{\B{ko\-ren}}, \HL{yll}{\B{yll}}, \HL{kllvawm}{\B{kll\-vawm}}, \HL{kawng}{\B{kawng}}, \HL{tsngawvIk}{\B{tsngaw\-vìk}}, \HL{txantsawtsray}{\B{txan\-tsaw\-tsray}}, \HL{muiA}{\B{muiä}}, \HL{tsaktap}{\B{tsak\-tap}}, \HL{tImwiA}{\B{tìmwiä}}, \HL{ngA'An}{\B{ngä\-'än}}, \HL{zeyko}{\B{zeyko}}, \HL{paw}{\B{paw}}, and \HL{tsawl}{\B{tsawl slu}}.

% \item[2020-05-30] added \HL{'ak}{\B{'ak}}, a note on \B{ko\-mum} from \LINK{https://forum.learnnavi.org/language-updates/na'vi-slang/}{here} in the \HL{omum}{\B{omum}} entry; edited the entries for \HL{pawm}{\B{pawm}}, \HL{'eyng}{\B{'eyng}}, and \HL{ta}{\B{ta}} as per \LINK{http://naviteri.org/2020/05/tipusawm-tiuseyng-si-okvur-a-eltur-titxen-si-asking-answering-and-an-interesting-story/}{this} blog post, added examples to \HL{makto}{\B{mak\-to}}, \HL{fyape}{\B{fya\-pe}}, \HL{zong}{\B{zong}}, \HL{nIzawnong}{\B{nì\-zaw\-nong}}, \HL{spxin}{\B{spxin}} and \HL{latsi}{\B{la\-tsi}}.

% \item[2020-05-25] added new words from \LINK{http://naviteri.org/2020/05/aawa-liu-amip-si-vurway-alor-a-few-new-words-and-a-beautiful-poem/}{'A'awa Lì'u Amip sì Vurway Alor}; added example to \HL{fwIng}{\B{fwìng}} and \HL{Eywa}{\B{Ey\-wa}}.

% \item[2020-04-24] added new words from \LINK{http://naviteri.org/2020/04/aawa-u-amip-a-few-new-things/}{’A’awa ’U Amip}; corrected a typo in the Abbreviation section (\emph{Irayo nìtxan ma Altafen}); edited the \HL{plltxe}{\B{plltxe}} section to include transitive examples; added examples to \ding{43} \HL{hayalo}{\B{ha\-ya\-lo}}, \HL{lahe}{\B{la\-he}}, \HL{yomyo}{\B{yom\-yo}}, \HL{sey}{\B{sey}}, \HL{pay2}{\B{pay}}, \HL{ripx}{\B{ripx}}, \HL{tskxe}{\B{tskxe}}, \HL{kangay}{\B{ka\-ngay si}}, \HL{mIkam}{\B{mì\-kam}}, \HL{pxawpa}{\B{pxaw\-pa}}, \HL{'etnaw}{\B{'et\-naw}}, \HL{ftAr}{\B{ftär}}, and \HL{sAmunge}{\B{sä\-mu\-nge}}; corrected a few errors in the hyperlinks.

% \item[2020-03-27] added new words from \LINK{http://naviteri.org/2020/03/mipa-sawasultsyip-a-new-contest/}{Mip\-a Sä\-wä\-sul\-tsyìp} and \LINK{http://naviteri.org/2020/04/keltrrtrra-tieylan-an-unusual-friendship/}{Kel\-trr\-trr\-a Tì\-'ey\-lan}; added examples to \HL{keltrrtrr}{\B{kel\-trr\-trr}}.

% \item[2020-03-24] added new words from \LINK{http://naviteri.org/2020/03/teri-tifkeytok-lefkrr-about-the-present-situation/}{Teri Tìfkeytok Lefkrr}; added examples to \HL{latem}{\B{la\-tem}}, \HL{txopu}{\B{txopu si}}, \HL{'u}{\B{'u}}, \HL{la'a}{\B{la'a}}, \HL{yll}{\B{yll}}, and \HL{kaoktu}{\B{koaktu}}.

% \item[2020-03-02] added \HL{yokx}{\B{yokx}}, new words from \LINK{http://naviteri.org/2020/02/some-words-for-leap-year-day/}{Some Words for Leap Year Day}; added examples to \HL{sAspxin}{\B{s\"a\-spxin}}, \HL{mungwrr}{\B{mung\-wrr}} and \HL{tok}{\B{tok}}.

% \item[2020-01-03] added new words from \LINK{http://naviteri.org/2019/12/tengkrr-zisit-leratem-as-the-year-changes/}{Tengkrr Leratem Zìsìt …}. Added to the definition of \B{ke\-ftxo}.

% \item[2019-10-28] added stress for the flora and fauna words from \LINK{http://naviteri.org/2019/07/mipa-ayewll-mipa-ayioang-new-plants-new-animals/}{Mipa Ayewll, Mipa Ayioang}; added blog entries \LINK{http://naviteri.org/2019/08/aawa-lifya-si-lifyavi-amip-a-few-new-words-and-expressions/}{'A'awa Lì'u sì Lì'fyavi Amip}, \LINK{http://naviteri.org/2019/09/choice-statements-vs-choice-questions-and-some-insults/}{Choice Statements vs. Choice Questions.}, \LINK{http://naviteri.org/2019/10/vurway-alor-a-beautiful-narrative-poem-2/}{Vurway Alor}; added examples to: \HL{fu1}{\B{fu}}, \HL{keftxo}{\B{keftxo}}, \HL{tok}{\B{tok}}, \HL{slantire}{\B{slan\-ti\-re}}, \HL{zIskrrsom}{\B{zì\-skrr\-som}}, \HL{tIsop}{\B{tì\-sop}}, \HL{nari}{\B{na\-ri si}}, \HL{kxetse}{\B{kxe\-tse}}, and \HL{onvA'}{\B{on\-vä'}}.

% \item[2019-06-10] names and words from the glossary from the \emph{Dark Horse Comic} book series added; new entries from the new blog post \LINK{naviteri.org/2019/06/50a-liu-amip-40-new-words/}{40 new words} added; added examples to \B{venu}, \B{txanro'a}, \B{fwäkì}, \B{sey}, \B{heyr}, \B{zeswa}, \B{tiam}, \B{virä} (transitivity changed), \B{epxang}; added new flora and fauna words from \LINK{https://forum.learnnavi.org/language-updates/expansion-to-flora-fauna-and-places-from-the-valley-of-mo'ara/}{Expansion to Flora, Fauna, and Places} on LearnNavi.org.

\item[2019-01-03] \textsc{not a draft anymore!}%; added (\emph{b.}) definition of \B{inan}; hyperlinks should work throughout the document now (\emph{txo ayngal keyeyoti rivun, rutxe peng oer}); latest entries from the blog post \LINK{http://naviteri.org/2018/08/fmawnti-stolawm-srak-have-you-heard-the-news/}{Fmawn\-ti sto\-lawm srak?}, \LINK{http://naviteri.org/2018/10/aysrr-ayvospxi-ayzisikrr-days-months-seasons/}{Ay\-srr, Ay\-vo\-spxì, Ay\-zì\-sì\-krr} and the \LINK{https://forum.learnnavi.org/language-updates/new-words-leftxoza-(celebratory)-poston-(boston)-kelni-(koln-cologne)/}{Language Update} added.

% \item[2018-08-04] somehow \B{sälang} never made it into the entries -- that's fixed now; hyperlinks included up to section U now; new entries from the blog post \LINK{http://naviteri.org/2018/07/ta-sulfatu-a-ayliu-niul-more-words-from-our-experts/}{Ta sulfätu a aylì'u nì'ul} and \LINK{https://forum.learnnavi.org/language-updates/'count-to'-wa-and-sequential-genitive-(discussion)/}{this language update}; examples added to \B{wä}, \B{tì\-'efu}, \B{vay}, \B{ta}, \B{tiam}, \B{tì\-ska\-'a}, \B{Na\-ra\-nawm}, \B{flì}, \B{letskxe}, \B{txawew}, and \B{zu\-saw\-krr}.

% \item[2018-05-20] vocabulary from the blog post \LINK{http://naviteri.org/2018/04/100a-liu-amip-64-new-words-part-2/}{64 New Words, Pt. 2} added; hyperlinks include section T now, added examples to \B{äie}, \B{swok}, \B{ngawng}, \B{niä}, \B{snu\-mì\-na} and \B{te\-swo\-tìng} (transitivity confirmed) mainly from the deleted scenes; added examples to \B{tu\-tan}, \B{zam}, \B{numtsengvi}, \B{tìeyktan}, \B{tsapo}, and \B{kilvan}; colloquial pronunciation added to nouns beginning with \B{sä-}.

% \item[2018-04-02] vocabulary from the blog post \LINK{http://naviteri.org/2018/02/negative-questions-in-navi/}{Negative Questions in Na'vi} and \LINK{http://naviteri.org/2018/03/100a-liu-amip-64-new-words-part-1/}{64 New Words} added; examples added to \B{tì\-txur}, \B{mun\-txa si}, \B{lä\-txayn}, \B{ftxu\-lì\-'u}, \B{fpom\-tokx}, \B{kxitx}, \B{txep}, \B{rìn}, \B{tìng} and \B{mauti}.

% \item[2018-01-03] error in \B{mllte} example corrected (\emph{irayo wm.annisur}), vocabulary from the blog post \LINK{http://naviteri.org/2017/12/zisit-amip-lefpom-ma-eylan-happy-new-year-friends/}{Zìsìt Amip Lefpom} added; examples added to \B{lal}, \B{tì\-mun\-txa}, \B{lrr\-tok}, \B{'i'a}, \B{kxe\-ner}, \B{eyk\-yu}, \B{tsun\-slu}, \B{syon\-lì'u}, \B{srey}, and \B{pup}.

% \item[2017-12-13] vocabulary from the blog post \LINK{http://naviteri.org/2017/10/vurway-alor-a-beautiful-narrative-poem/}{Vurway Alor} added; added `later' to the definition of \B{maw\-krr}, examples to \B{txu\-la}, and \B{ta} (\emph{irayo tsmukur alu wm.annis}).

% \item[2017-10-05] vocabulary from the blog post \LINK{http://naviteri.org/2017/10/more-language-for-talking-about-language/}{More language …} added.

% \item[2017-10-03] vocabulary from the blog post \LINK{http://naviteri.org/2017/09/zisikrr-amip-ayliu-amip-new-words-for-the-new-season/}{Zìsìkrr amip, Aylì'u amip} added; example added to \HL{wok}{\B{wok}}, \HL{txArem}{\B{txä\-rem}}, \HL{fkxile}{\B{fki\-le}}, \HL{yo}{\B{yo}}, \HL{renu}{\B{renu}}, and \HL{neni}{\B{neni}}.

% \item[2017-07-08] example added to \HL{yIm}{\B{yìm}} (a big irayo to wm.annis for pointing that out); mistake in \HL{'ongop}{\B{'o\-ngop}} corrected; example to \HL{tsun}{\B{tsun}} (\emph{b.}) added; confirmation of \HL{fmetok}{\B{fme\-tok}} as \emph{vtr.} and example added; words from the new blog post \LINK{http://naviteri.org/2017/02/ayioang-amip-si-ayu-alahe-new-animals-and-other-things/}{New Animals and other things} and examples from \LINK{http://naviteri.org/2017/03/titusemteri-concerning-shooting/}{Concerning Shooting} added to \HL{tem}{\B{tem}} and \HL{toltem}{\B{tol\-tem}}; example to \HL{tsweykayon}{\B{tswayon} (\emph{b.})} added; definition of \HL{takuk}{\B{ta\-kuk}} changed; proper nouns \B{Met\-ka\-yi\-na} and \B{To\-no\-wa\-ri} added; vocabulary from the blog posts \LINK{http://naviteri.org/2017/06/ayliu-a-ta-eywaeveng-kifkey-uniltirantokxa-words-from-disneys-pandora-the-world-of-avatar/}{Aylì'u a ta ``Eywa'eveng''} and \LINK{http://naviteri.org/2017/07/way-tiretua-the-shamans-song/}{Way Tiretuä} added; examples to \HL{tsngawpay}{\B{tsngaw\-pay}} and \HL{frrtu}{\B{frr\-tu}} added.

% \item[2017-01-05] still major reworking, hyperlinks including section R added; transitivity of \HL{zong}{\B{zong}} confirmed, definition of the \B{nì}- prefix slightly altered as per \LINK{https://forum.learnnavi.org/language-updates/eapressions-of-'like-we'/}{this confirmation}; error in \HL{ne}{\B{ne}} (\emph{f.}) definition and \B{tsan\-kum\-nga'} corrected; \B{tam ke tam} in the \B{tam} entry added; new vocabulary from the blog post \LINK{http://naviteri.org/2016/12/tengkrr-perahem-zisit-amip-as-the-new-year-arrives/}{As the New Year arrives …}; examples added to: \HL{ftxey}{\B{ftxey}}, \HL{txantur}{\B{txan\-tur}}, \HL{muiA}{\B{muiä}} (\emph{a.}), \HL{yaymak}{\B{yay\-mak}}, \HL{mawl}{\B{mawl}}, \HL{ska'a}{\B{ska\-'a}}, \HL{zerok}{\B{ze\-rok}}, \HL{eyktanay}{\B{eyk\-ta\-nay}}, \HL{reng}{\B{reng}} (\emph{b.}), \HL{slan}{\B{slan}}, and \HL{tIlam}{\B{tì\-lam}}; confirmation of \HL{tung}{\B{tung}}'s transitivity and examples added, \HL{pllhrr}{\B{pll\-hrr}} and \HL{sApllhrr}{\B{sä\-pll\-hrr}} added.

\item[2016-07-04] major reworking, hyperlinks up to section N added. % \HL{'u}{\B{'u}} entry decluttered, example to \HL{'awm}{\B{'awm}} added, stress of \HL{lefpomtokx}{\B{le\-fpom\-tokx}} and \HL{kelfpomtokx}{\B{kel\-fpom\-tokx}} corrected; \HL{tsamkuk}{\B{tsam\-kuk}} added, words from new blog post \LINK{http://naviteri.org/2016/06/mrrvola-lifyavi-amip-forty-new-expressions/}{Forty new expressions} added, examples to \HL{somwew}{\B{som\-wew}}, \HL{pxi}{\B{pxi}} added.

% \item[2016-01-06] examples of \HL{to}{\B{to}} refined.

% \item[2015-12-15] confirmation of \HL{sAnume si}{\B{sä\-nu\-me si}} and \HL{zop\-lo si}{\B{zoplo si}} added.

% \item[2015-12-06] words from new blog post \LINK{http://naviteri.org/2015/11/vomuna-liu-amip-ten-new-words/}{Ten new words} added; examples and definition for \HL{kan}{\B{kan}} modified; examples to \HL{takuk}{\B{ta\-kuk}}, \HL{kxAnAng}{\B{kxä\-näng}}, and \HL{yo}{\B{yo}} (\emph{a.}) added; definition of \HL{fpak}{\B{fpak}} changed; \emph{Avatar Game} reference added as \textsc{\cb{g}}.

% \item[2015-11-07] definition of \HL{'llngo}{\B{'llngo}} changed.

\item[2015-10-29] `change log' added. % definition of \B{ew} in the Diphthongs section modified; definition of \HL{zawprrte'}{\B{zaw\-prr\-te'}} modified; derivation of \HL{zeswa}{\B{ze\-swa}} added.

% \item[2015-10-02] \HL{pxawngip}{\B{pxawngip}} from new blog post \LINK{http://naviteri.org/2015/09/tskxekengtsyip-a-mikyunfpi-a-little-listening-exercise-2/}{Little Listening Exercise} added

% \item[2015-09-04] words from new blog post \LINK{http://naviteri.org/2015/08/aylifyavi-lereyfya-2-cultural-terms-2-and-more/}{Cultural Terms 2} added; entry for \HL{}{\B{hufwetsyìp}} added; example to \HL{frrtu}{\B{frrtu}}, \HL{inanfya}{\B{inan\-fya}}, and \HL{tsa\-krr}{\B{tsakrr}} (\emph{n}.) added.

\item[2015-08-06] sections `Introduction to IPA', `word classes', and `abbreviations' added.

% \item[ 2015-04-23] sections Y \& Z added.

% \item[ 2015-04-04] sections V \& W added.

% \item[ 2015-03-28] section U added.

% \item[ 2015-03-27] section TS added.

% \item[ 2015-01-25] section TX added.

% \item[ 2015-01-22] revised draft and letter T added.

% \item[ 2015-01-06] section S added.

% \item[ 2014-10-20] sections PX \& R added.

% \item[ 2014-09-27] section P added.

\item[ 2014-08-31] first draft (work in progress) presented on \LINK{http://forum.learnnavi.org/projects/an-annotated-dictionary-\%28draft\%29/}{LN}.

\item[ 2014-02-01] project start.

\end{itemize}
\end{footnotesize}
	
\end{document}